\documentclass[tikz,border=5pt]{standalone}
\usepackage[UTF8]{ctex}
\usepackage{amsmath,amssymb}
\usetikzlibrary{arrows.meta,calc}
\tikzset{>=Stealth,line cap=round,line join=round,
force/.style={blue!75!black,thick},
pressure/.style={red!75!black,thick},
aux/.style={gray!70,dashed,thin}}
\newcommand{\cyl}{
\fill[gray!8,opacity=.70](-1.1,0)--(-1.1,3)
arc(180:360:1.1 and .34)--(1.1,0)arc(0:-180:1.1 and .34)--cycle;
\draw[aux](-1.1,3)arc(180:0:1.1 and .34);
\draw[thick](-1.1,3)arc(180:360:1.1 and .34);
\draw[aux](-1.1,0)arc(180:0:1.1 and .34);
\draw[thick](-1.1,0)arc(180:360:1.1 and .34);
\draw[thick](-1.1,0)--(-1.1,3);
\draw[thick](1.1,0)--(1.1,3);
\draw[aux](0,-.45)--(0,3.55);
}
\newcommand{\tp}[3]{
\coordinate(#1)at({(2.4-.8*cos(#2))*cos(#3)},
{.5*(2.4-.8*cos(#2))*sin(#3)+.8660254038*.8*sin(#2)});
}
\newcommand{\meridiancurve}[2]{
\draw[#1,domain=#2,samples=80]
plot({2.4-.8*cos(\x)},{.8660254038*.8*sin(\x)});
}
\begin{document}
% 1. The planes are shown separately, so their intersection does not hide either.
\begin{tikzpicture}[font=\small]
\begin{scope}
\filldraw[fill=blue!12,draw=blue!65!black]
(-.85,-.516)--(.85,-.084)--(.85,3.616)--(-.85,3.184)--cycle;
\cyl
\draw[blue!75!black,thick](-.85,-.216)--(-.85,2.784);
\draw[blue!75!black,dashed](.85,.216)--(.85,3.216);
\draw[thin](0,3)--(1.1,3)node[midway,above=2pt]{$r$};
\draw[<->](-1.55,0)--(-1.55,3)node[midway,left]{$l$};
\node at(0,4.1){纵切面：含中轴线};
\end{scope}
\begin{scope}[xshift=5.15cm]
\filldraw[fill=green!12,draw=green!50!black]
(-1.75,1.02)--(1.05,1.02)--(1.75,1.92)--(-1.05,1.92)--cycle;
\cyl
\draw[green!50!black,thick](-1.1,1.5)arc(180:360:1.1 and .34);
\draw[green!50!black,dashed](-1.1,1.5)arc(180:0:1.1 and .34);
\draw[green!50!black](-1.75,1.02)--(1.05,1.02)--(1.75,1.92);
\node at(0,4.1){横切面：垂直中轴线};
\end{scope}
\end{tikzpicture}
% 2. Actual torus surface plus the meridian section using the same angle origin.
\begin{tikzpicture}[font=\small]
\begin{scope}
% Automatically generated vector mesh; parameters and camera in build/torus-mesh-q03.json.
\path[fill=blue!39!white,draw=blue!39!white,line width=.12pt] (0.000000,1.304421)--(0.000000,1.147677)--(-0.300440,1.140298)--(-0.307686,1.296863)--cycle;
\path[fill=blue!39!white,draw=blue!39!white,line width=.12pt] (0.307686,1.296863)--(0.300440,1.140298)--(0.000000,1.147677)--(0.000000,1.304421)--cycle;
\path[fill=blue!40!white,draw=blue!40!white,line width=.12pt] (0.300440,1.140298)--(0.290688,0.985804)--(0.000000,0.992945)--(0.000000,1.147677)--cycle;
\path[fill=blue!40!white,draw=blue!40!white,line width=.12pt] (0.000000,1.147677)--(0.000000,0.992945)--(-0.290688,0.985804)--(-0.300440,1.140298)--cycle;
\path[fill=blue!37!white,draw=blue!37!white,line width=.12pt] (0.312148,1.449484)--(0.307686,1.296863)--(0.000000,1.304421)--(0.000000,1.457152)--cycle;
\path[fill=blue!37!white,draw=blue!37!white,line width=.12pt] (0.000000,1.457152)--(0.000000,1.304421)--(-0.307686,1.296863)--(-0.312148,1.449484)--cycle;
\path[fill=blue!39!white,draw=blue!39!white,line width=.12pt] (0.612409,1.274262)--(0.597986,1.118229)--(0.300440,1.140298)--(0.307686,1.296863)--cycle;
\path[fill=blue!38!white,draw=blue!38!white,line width=.12pt] (-0.307686,1.296863)--(-0.300440,1.140298)--(-0.597986,1.118229)--(-0.612409,1.274262)--cycle;
\path[fill=blue!41!white,draw=blue!41!white,line width=.12pt] (0.597986,1.118229)--(0.578577,0.964452)--(0.290688,0.985804)--(0.300440,1.140298)--cycle;
\path[fill=blue!40!white,draw=blue!40!white,line width=.12pt] (-0.300440,1.140298)--(-0.290688,0.985804)--(-0.578577,0.964452)--(-0.597986,1.118229)--cycle;
\path[fill=blue!38!white,draw=blue!38!white,line width=.12pt] (0.621290,1.426556)--(0.612409,1.274262)--(0.307686,1.296863)--(0.312148,1.449484)--cycle;
\path[fill=blue!36!white,draw=blue!36!white,line width=.12pt] (-0.312148,1.449484)--(-0.307686,1.296863)--(-0.612409,1.274262)--(-0.621290,1.426556)--cycle;
\path[fill=blue!41!white,draw=blue!41!white,line width=.12pt] (0.290688,0.985804)--(0.278805,0.839321)--(0.000000,0.846169)--(0.000000,0.992945)--cycle;
\path[fill=blue!41!white,draw=blue!41!white,line width=.12pt] (0.000000,0.992945)--(0.000000,0.846169)--(-0.278805,0.839321)--(-0.290688,0.985804)--cycle;
\path[fill=blue!35!white,draw=blue!35!white,line width=.12pt] (0.313655,1.592296)--(0.312148,1.449484)--(0.000000,1.457152)--(0.000000,1.600000)--cycle;
\path[fill=blue!35!white,draw=blue!35!white,line width=.12pt] (0.000000,1.600000)--(0.000000,1.457152)--(-0.312148,1.449484)--(-0.313655,1.592296)--cycle;
\path[fill=blue!41!white,draw=blue!41!white,line width=.12pt] (-0.290688,0.985804)--(-0.278805,0.839321)--(-0.554926,0.818841)--(-0.578577,0.964452)--cycle;
\path[fill=blue!41!white,draw=blue!41!white,line width=.12pt] (0.578577,0.964452)--(0.554926,0.818841)--(0.278805,0.839321)--(0.290688,0.985804)--cycle;
\path[fill=blue!38!white,draw=blue!38!white,line width=.12pt] (-0.612409,1.274262)--(-0.597986,1.118229)--(-0.889774,1.081685)--(-0.911234,1.236837)--cycle;
\path[fill=blue!40!white,draw=blue!40!white,line width=.12pt] (0.911234,1.236837)--(0.889774,1.081685)--(0.597986,1.118229)--(0.612409,1.274262)--cycle;
\path[fill=blue!41!white,draw=blue!41!white,line width=.12pt] (0.889774,1.081685)--(0.860893,0.929094)--(0.578577,0.964452)--(0.597986,1.118229)--cycle;
\path[fill=blue!39!white,draw=blue!39!white,line width=.12pt] (-0.597986,1.118229)--(-0.578577,0.964452)--(-0.860893,0.929094)--(-0.889774,1.081685)--cycle;
\path[fill=blue!36!white,draw=blue!36!white,line width=.12pt] (0.624289,1.569256)--(0.621290,1.426556)--(0.312148,1.449484)--(0.313655,1.592296)--cycle;
\path[fill=blue!34!white,draw=blue!34!white,line width=.12pt] (-0.313655,1.592296)--(-0.312148,1.449484)--(-0.621290,1.426556)--(-0.624289,1.569256)--cycle;
\path[fill=blue!38!white,draw=blue!38!white,line width=.12pt] (0.924449,1.388587)--(0.911234,1.236837)--(0.612409,1.274262)--(0.621290,1.426556)--cycle;
\path[fill=blue!36!white,draw=blue!36!white,line width=.12pt] (-0.621290,1.426556)--(-0.612409,1.274262)--(-0.911234,1.236837)--(-0.924449,1.388587)--cycle;
\path[fill=blue!41!white,draw=blue!41!white,line width=.12pt] (0.278805,0.839321)--(0.265249,0.706475)--(0.000000,0.712991)--(0.000000,0.846169)--cycle;
\path[fill=blue!41!white,draw=blue!41!white,line width=.12pt] (0.000000,0.846169)--(0.000000,0.712991)--(-0.265249,0.706475)--(-0.278805,0.839321)--cycle;
\path[fill=blue!42!white,draw=blue!42!white,line width=.12pt] (0.860893,0.929094)--(0.825702,0.784928)--(0.554926,0.818841)--(0.578577,0.964452)--cycle;
\path[fill=blue!40!white,draw=blue!40!white,line width=.12pt] (-0.578577,0.964452)--(-0.554926,0.818841)--(-0.825702,0.784928)--(-0.860893,0.929094)--cycle;
\path[fill=blue!42!white,draw=blue!42!white,line width=.12pt] (0.554926,0.818841)--(0.527943,0.686992)--(0.265249,0.706475)--(0.278805,0.839321)--cycle;
\path[fill=blue!41!white,draw=blue!41!white,line width=.12pt] (-0.278805,0.839321)--(-0.265249,0.706475)--(-0.527943,0.686992)--(-0.554926,0.818841)--cycle;
\path[fill=blue!33!white,draw=blue!33!white,line width=.12pt] (0.312148,1.719809)--(0.313655,1.592296)--(0.000000,1.600000)--(0.000000,1.727477)--cycle;
\path[fill=blue!32!white,draw=blue!32!white,line width=.12pt] (0.000000,1.727477)--(0.000000,1.600000)--(-0.313655,1.592296)--(-0.312148,1.719809)--cycle;
\path[fill=blue!36!white,draw=blue!36!white,line width=.12pt] (0.928911,1.531105)--(0.924449,1.388587)--(0.621290,1.426556)--(0.624289,1.569256)--cycle;
\path[fill=blue!34!white,draw=blue!34!white,line width=.12pt] (-0.624289,1.569256)--(-0.621290,1.426556)--(-0.924449,1.388587)--(-0.928911,1.531105)--cycle;
\path[fill=blue!37!white,draw=blue!37!white,line width=.12pt] (-0.911234,1.236837)--(-0.889774,1.081685)--(-1.172992,1.031016)--(-1.201283,1.184946)--cycle;
\path[fill=blue!40!white,draw=blue!40!white,line width=.12pt] (1.201283,1.184946)--(1.172992,1.031016)--(0.889774,1.081685)--(0.911234,1.236837)--cycle;
\path[fill=blue!41!white,draw=blue!41!white,line width=.12pt] (1.172992,1.031016)--(1.134919,0.880070)--(0.860893,0.929094)--(0.889774,1.081685)--cycle;
\path[fill=blue!39!white,draw=blue!39!white,line width=.12pt] (-0.889774,1.081685)--(-0.860893,0.929094)--(-1.134919,0.880070)--(-1.172992,1.031016)--cycle;
\path[fill=blue!32!white,draw=blue!32!white,line width=.12pt] (-0.312148,1.719809)--(-0.313655,1.592296)--(-0.624289,1.569256)--(-0.621290,1.696881)--cycle;
\path[fill=blue!33!white,draw=blue!33!white,line width=.12pt] (0.621290,1.696881)--(0.624289,1.569256)--(0.313655,1.592296)--(0.312148,1.719809)--cycle;
\path[fill=blue!38!white,draw=blue!38!white,line width=.12pt] (1.218704,1.335944)--(1.201283,1.184946)--(0.911234,1.236837)--(0.924449,1.388587)--cycle;
\path[fill=blue!35!white,draw=blue!35!white,line width=.12pt] (-0.924449,1.388587)--(-0.911234,1.236837)--(-1.201283,1.184946)--(-1.218704,1.335944)--cycle;
\path[fill=blue!42!white,draw=blue!42!white,line width=.12pt] (0.825702,0.784928)--(0.785553,0.654728)--(0.527943,0.686992)--(0.554926,0.818841)--cycle;
\path[fill=blue!41!white,draw=blue!41!white,line width=.12pt] (-0.554926,0.818841)--(-0.527943,0.686992)--(-0.785553,0.654728)--(-0.825702,0.784928)--cycle;
\path[fill=blue!42!white,draw=blue!42!white,line width=.12pt] (1.134919,0.880070)--(1.088526,0.737908)--(0.825702,0.784928)--(0.860893,0.929094)--cycle;
\path[fill=blue!40!white,draw=blue!40!white,line width=.12pt] (-0.860893,0.929094)--(-0.825702,0.784928)--(-1.088526,0.737908)--(-1.134919,0.880070)--cycle;
\path[fill=blue!41!white,draw=blue!41!white,line width=.12pt] (0.265249,0.706475)--(0.250539,0.592374)--(0.000000,0.598528)--(0.000000,0.712991)--cycle;
\path[fill=blue!41!white,draw=blue!41!white,line width=.12pt] (0.000000,0.712991)--(0.000000,0.598528)--(-0.250539,0.592374)--(-0.265249,0.706475)--cycle;
\path[fill=blue!33!white,draw=blue!33!white,line width=.12pt] (0.924449,1.658912)--(0.928911,1.531105)--(0.624289,1.569256)--(0.621290,1.696881)--cycle;
\path[fill=blue!31!white,draw=blue!31!white,line width=.12pt] (-0.621290,1.696881)--(-0.624289,1.569256)--(-0.928911,1.531105)--(-0.924449,1.658912)--cycle;
\path[fill=blue!33!white,draw=blue!33!white,line width=.12pt] (-0.928911,1.531105)--(-0.924449,1.388587)--(-1.218704,1.335944)--(-1.224587,1.478207)--cycle;
\path[fill=blue!36!white,draw=blue!36!white,line width=.12pt] (1.224587,1.478207)--(1.218704,1.335944)--(0.924449,1.388587)--(0.928911,1.531105)--cycle;
\path[fill=blue!41!white,draw=blue!41!white,line width=.12pt] (-0.265249,0.706475)--(-0.250539,0.592374)--(-0.498665,0.573971)--(-0.527943,0.686992)--cycle;
\path[fill=blue!41!white,draw=blue!41!white,line width=.12pt] (0.527943,0.686992)--(0.498665,0.573971)--(0.250539,0.592374)--(0.265249,0.706475)--cycle;
\path[fill=blue!30!white,draw=blue!30!white,line width=.12pt] (0.000000,1.834683)--(0.000000,1.727477)--(-0.312148,1.719809)--(-0.307686,1.827125)--cycle;
\path[fill=blue!30!white,draw=blue!30!white,line width=.12pt] (0.307686,1.827125)--(0.312148,1.719809)--(0.000000,1.727477)--(0.000000,1.834683)--cycle;
\path[fill=blue!40!white,draw=blue!40!white,line width=.12pt] (1.479763,1.119090)--(1.444914,0.966711)--(1.172992,1.031016)--(1.201283,1.184946)--cycle;
\path[fill=blue!37!white,draw=blue!37!white,line width=.12pt] (-1.201283,1.184946)--(-1.172992,1.031016)--(-1.444914,0.966711)--(-1.479763,1.119090)--cycle;
\path[fill=blue!42!white,draw=blue!42!white,line width=.12pt] (1.088526,0.737908)--(1.035598,0.609994)--(0.785553,0.654728)--(0.825702,0.784928)--cycle;
\path[fill=blue!40!white,draw=blue!40!white,line width=.12pt] (-0.825702,0.784928)--(-0.785553,0.654728)--(-1.035598,0.609994)--(-1.088526,0.737908)--cycle;
\path[fill=blue!41!white,draw=blue!41!white,line width=.12pt] (1.444914,0.966711)--(1.398014,0.817853)--(1.134919,0.880070)--(1.172992,1.031016)--cycle;
\path[fill=blue!38!white,draw=blue!38!white,line width=.12pt] (-1.172992,1.031016)--(-1.134919,0.880070)--(-1.398014,0.817853)--(-1.444914,0.966711)--cycle;
\path[fill=blue!29!white,draw=blue!29!white,line width=.12pt] (-0.307686,1.827125)--(-0.312148,1.719809)--(-0.621290,1.696881)--(-0.612409,1.804524)--cycle;
\path[fill=blue!30!white,draw=blue!30!white,line width=.12pt] (0.612409,1.804524)--(0.621290,1.696881)--(0.312148,1.719809)--(0.307686,1.827125)--cycle;
\path[fill=blue!41!white,draw=blue!41!white,line width=.12pt] (0.785553,0.654728)--(0.741989,0.543496)--(0.498665,0.573971)--(0.527943,0.686992)--cycle;
\path[fill=blue!41!white,draw=blue!41!white,line width=.12pt] (-0.527943,0.686992)--(-0.498665,0.573971)--(-0.741989,0.543496)--(-0.785553,0.654728)--cycle;
\path[fill=blue!38!white,draw=blue!38!white,line width=.12pt] (1.501223,1.269133)--(1.479763,1.119090)--(1.201283,1.184946)--(1.218704,1.335944)--cycle;
\path[fill=blue!35!white,draw=blue!35!white,line width=.12pt] (-1.218704,1.335944)--(-1.201283,1.184946)--(-1.479763,1.119090)--(-1.501223,1.269133)--cycle;
\path[fill=blue!42!white,draw=blue!42!white,line width=.12pt] (1.398014,0.817853)--(1.340867,0.678234)--(1.088526,0.737908)--(1.134919,0.880070)--cycle;
\path[fill=blue!39!white,draw=blue!39!white,line width=.12pt] (-1.134919,0.880070)--(-1.088526,0.737908)--(-1.340867,0.678234)--(-1.398014,0.817853)--cycle;
\path[fill=blue!31!white,draw=blue!31!white,line width=.12pt] (-0.924449,1.658912)--(-0.928911,1.531105)--(-1.224587,1.478207)--(-1.218704,1.606269)--cycle;
\path[fill=blue!34!white,draw=blue!34!white,line width=.12pt] (1.218704,1.606269)--(1.224587,1.478207)--(0.928911,1.531105)--(0.924449,1.658912)--cycle;
\path[fill=blue!31!white,draw=blue!31!white,line width=.12pt] (0.911234,1.767098)--(0.924449,1.658912)--(0.621290,1.696881)--(0.612409,1.804524)--cycle;
\path[fill=blue!29!white,draw=blue!29!white,line width=.12pt] (-0.612409,1.804524)--(-0.621290,1.696881)--(-0.924449,1.658912)--(-0.911234,1.767098)--cycle;
\path[fill=blue!40!white,draw=blue!40!white,line width=.12pt] (0.250539,0.592374)--(0.235241,0.501401)--(0.000000,0.507180)--(0.000000,0.598528)--cycle;
\path[fill=blue!40!white,draw=blue!40!white,line width=.12pt] (0.000000,0.598528)--(0.000000,0.507180)--(-0.235241,0.501401)--(-0.250539,0.592374)--cycle;
\path[fill=blue!36!white,draw=blue!36!white,line width=.12pt] (1.508470,1.411074)--(1.501223,1.269133)--(1.218704,1.335944)--(1.224587,1.478207)--cycle;
\path[fill=blue!32!white,draw=blue!32!white,line width=.12pt] (-1.224587,1.478207)--(-1.218704,1.335944)--(-1.501223,1.269133)--(-1.508470,1.411074)--cycle;
\path[fill=blue!40!white,draw=blue!40!white,line width=.12pt] (-0.785553,0.654728)--(-0.741989,0.543496)--(-0.978167,0.501243)--(-1.035598,0.609994)--cycle;
\path[fill=blue!41!white,draw=blue!41!white,line width=.12pt] (1.035598,0.609994)--(0.978167,0.501243)--(0.741989,0.543496)--(0.785553,0.654728)--cycle;
\path[fill=blue!42!white,draw=blue!42!white,line width=.12pt] (1.340867,0.678234)--(1.275669,0.553222)--(1.035598,0.609994)--(1.088526,0.737908)--cycle;
\path[fill=blue!40!white,draw=blue!40!white,line width=.12pt] (-1.088526,0.737908)--(-1.035598,0.609994)--(-1.275669,0.553222)--(-1.340867,0.678234)--cycle;
\path[fill=blue!40!white,draw=blue!40!white,line width=.12pt] (-0.250539,0.592374)--(-0.235241,0.501401)--(-0.468217,0.484122)--(-0.498665,0.573971)--cycle;
\path[fill=blue!40!white,draw=blue!40!white,line width=.12pt] (0.498665,0.573971)--(0.468217,0.484122)--(0.235241,0.501401)--(0.250539,0.592374)--cycle;
\path[fill=blue!27!white,draw=blue!27!white,line width=.12pt] (0.000000,1.917498)--(0.000000,1.834683)--(-0.307686,1.827125)--(-0.300440,1.910118)--cycle;
\path[fill=blue!27!white,draw=blue!27!white,line width=.12pt] (0.300440,1.910118)--(0.307686,1.827125)--(0.000000,1.834683)--(0.000000,1.917498)--cycle;
\path[fill=blue!41!white,draw=blue!41!white,line width=.12pt] (1.702920,0.889390)--(1.647647,0.743041)--(1.398014,0.817853)--(1.444914,0.966711)--cycle;
\path[fill=blue!38!white,draw=blue!38!white,line width=.12pt] (-1.444914,0.966711)--(-1.398014,0.817853)--(-1.647647,0.743041)--(-1.702920,0.889390)--cycle;
\path[fill=blue!40!white,draw=blue!40!white,line width=.12pt] (1.743993,1.039904)--(1.702920,0.889390)--(1.444914,0.966711)--(1.479763,1.119090)--cycle;
\path[fill=blue!36!white,draw=blue!36!white,line width=.12pt] (-1.479763,1.119090)--(-1.444914,0.966711)--(-1.702920,0.889390)--(-1.743993,1.039904)--cycle;
\path[fill=blue!40!white,draw=blue!40!white,line width=.12pt] (0.741989,0.543496)--(0.696683,0.455508)--(0.468217,0.484122)--(0.498665,0.573971)--cycle;
\path[fill=blue!40!white,draw=blue!40!white,line width=.12pt] (-0.498665,0.573971)--(-0.468217,0.484122)--(-0.696683,0.455508)--(-0.741989,0.543496)--cycle;
\path[fill=blue!28!white,draw=blue!28!white,line width=.12pt] (0.597986,1.888050)--(0.612409,1.804524)--(0.307686,1.827125)--(0.300440,1.910118)--cycle;
\path[fill=blue!26!white,draw=blue!26!white,line width=.12pt] (-0.300440,1.910118)--(-0.307686,1.827125)--(-0.612409,1.804524)--(-0.597986,1.888050)--cycle;
\path[fill=blue!39!white,draw=blue!39!white,line width=.12pt] (-1.398014,0.817853)--(-1.340867,0.678234)--(-1.580295,0.606480)--(-1.647647,0.743041)--cycle;
\path[fill=blue!42!white,draw=blue!42!white,line width=.12pt] (1.647647,0.743041)--(1.580295,0.606480)--(1.340867,0.678234)--(1.398014,0.817853)--cycle;
\path[fill=blue!34!white,draw=blue!34!white,line width=.12pt] (-1.501223,1.269133)--(-1.479763,1.119090)--(-1.743993,1.039904)--(-1.769285,1.188798)--cycle;
\path[fill=blue!38!white,draw=blue!38!white,line width=.12pt] (1.769285,1.188798)--(1.743993,1.039904)--(1.479763,1.119090)--(1.501223,1.269133)--cycle;
\path[fill=blue!31!white,draw=blue!31!white,line width=.12pt] (1.201283,1.715208)--(1.218704,1.606269)--(0.924449,1.658912)--(0.911234,1.767098)--cycle;
\path[fill=blue!28!white,draw=blue!28!white,line width=.12pt] (-0.911234,1.767098)--(-0.924449,1.658912)--(-1.218704,1.606269)--(-1.201283,1.715208)--cycle;
\path[fill=blue!30!white,draw=blue!30!white,line width=.12pt] (-1.218704,1.606269)--(-1.224587,1.478207)--(-1.508470,1.411074)--(-1.501223,1.539458)--cycle;
\path[fill=blue!34!white,draw=blue!34!white,line width=.12pt] (1.501223,1.539458)--(1.508470,1.411074)--(1.224587,1.478207)--(1.218704,1.606269)--cycle;
\path[fill=blue!40!white,draw=blue!40!white,line width=.12pt] (-1.035598,0.609994)--(-0.978167,0.501243)--(-1.204924,0.447619)--(-1.275669,0.553222)--cycle;
\path[fill=blue!41!white,draw=blue!41!white,line width=.12pt] (1.275669,0.553222)--(1.204924,0.447619)--(0.978167,0.501243)--(1.035598,0.609994)--cycle;
\path[fill=blue!26!white,draw=blue!26!white,line width=.12pt] (-0.597986,1.888050)--(-0.612409,1.804524)--(-0.911234,1.767098)--(-0.889774,1.851505)--cycle;
\path[fill=blue!28!white,draw=blue!28!white,line width=.12pt] (0.889774,1.851505)--(0.911234,1.767098)--(0.612409,1.804524)--(0.597986,1.888050)--cycle;
\path[fill=blue!40!white,draw=blue!40!white,line width=.12pt] (0.978167,0.501243)--(0.918440,0.415835)--(0.696683,0.455508)--(0.741989,0.543496)--cycle;
\path[fill=blue!40!white,draw=blue!40!white,line width=.12pt] (-0.741989,0.543496)--(-0.696683,0.455508)--(-0.918440,0.415835)--(-0.978167,0.501243)--cycle;
\path[fill=blue!40!white,draw=blue!40!white,line width=.12pt] (-1.340867,0.678234)--(-1.275669,0.553222)--(-1.503455,0.484957)--(-1.580295,0.606480)--cycle;
\path[fill=blue!42!white,draw=blue!42!white,line width=.12pt] (1.580295,0.606480)--(1.503455,0.484957)--(1.275669,0.553222)--(1.340867,0.678234)--cycle;
\path[fill=blue!32!white,draw=blue!32!white,line width=.12pt] (-1.508470,1.411074)--(-1.501223,1.269133)--(-1.769285,1.188798)--(-1.777825,1.330351)--cycle;
\path[fill=blue!36!white,draw=blue!36!white,line width=.12pt] (1.777825,1.330351)--(1.769285,1.188798)--(1.501223,1.269133)--(1.508470,1.411074)--cycle;
\path[fill=blue!39!white,draw=blue!39!white,line width=.12pt] (0.000000,0.507180)--(0.000000,0.442456)--(-0.219943,0.437053)--(-0.235241,0.501401)--cycle;
\path[fill=blue!39!white,draw=blue!39!white,line width=.12pt] (0.235241,0.501401)--(0.219943,0.437053)--(0.000000,0.442456)--(0.000000,0.507180)--cycle;
\path[fill=blue!39!white,draw=blue!39!white,line width=.12pt] (-0.235241,0.501401)--(-0.219943,0.437053)--(-0.437769,0.420898)--(-0.468217,0.484122)--cycle;
\path[fill=blue!39!white,draw=blue!39!white,line width=.12pt] (0.468217,0.484122)--(0.437769,0.420898)--(0.219943,0.437053)--(0.235241,0.501401)--cycle;
\path[fill=blue!31!white,draw=blue!31!white,line width=.12pt] (1.479763,1.649352)--(1.501223,1.539458)--(1.218704,1.606269)--(1.201283,1.715208)--cycle;
\path[fill=blue!27!white,draw=blue!27!white,line width=.12pt] (-1.201283,1.715208)--(-1.218704,1.606269)--(-1.501223,1.539458)--(-1.479763,1.649352)--cycle;
\path[fill=blue!24!white,draw=blue!24!white,line width=.12pt] (0.290688,1.965600)--(0.300440,1.910118)--(0.000000,1.917498)--(0.000000,1.972741)--cycle;
\path[fill=blue!24!white,draw=blue!24!white,line width=.12pt] (0.000000,1.972741)--(0.000000,1.917498)--(-0.300440,1.910118)--(-0.290688,1.965600)--cycle;
\path[fill=blue!41!white,draw=blue!41!white,line width=.12pt] (1.944527,0.799796)--(1.881411,0.656355)--(1.647647,0.743041)--(1.702920,0.889390)--cycle;
\path[fill=blue!37!white,draw=blue!37!white,line width=.12pt] (-1.702920,0.889390)--(-1.647647,0.743041)--(-1.881411,0.656355)--(-1.944527,0.799796)--cycle;
\path[fill=blue!39!white,draw=blue!39!white,line width=.12pt] (0.696683,0.455508)--(0.651378,0.394145)--(0.437769,0.420898)--(0.468217,0.484122)--cycle;
\path[fill=blue!39!white,draw=blue!39!white,line width=.12pt] (-0.468217,0.484122)--(-0.437769,0.420898)--(-0.651378,0.394145)--(-0.696683,0.455508)--cycle;
\path[fill=blue!40!white,draw=blue!40!white,line width=.12pt] (1.991426,0.948149)--(1.944527,0.799796)--(1.702920,0.889390)--(1.743993,1.039904)--cycle;
\path[fill=blue!35!white,draw=blue!35!white,line width=.12pt] (-1.743993,1.039904)--(-1.702920,0.889390)--(-1.944527,0.799796)--(-1.991426,0.948149)--cycle;
\path[fill=blue!28!white,draw=blue!28!white,line width=.12pt] (1.172992,1.800837)--(1.201283,1.715208)--(0.911234,1.767098)--(0.889774,1.851505)--cycle;
\path[fill=blue!25!white,draw=blue!25!white,line width=.12pt] (-0.889774,1.851505)--(-0.911234,1.767098)--(-1.201283,1.715208)--(-1.172992,1.800837)--cycle;
\path[fill=blue!40!white,draw=blue!40!white,line width=.12pt] (1.204924,0.447619)--(1.131352,0.365485)--(0.918440,0.415835)--(0.978167,0.501243)--cycle;
\path[fill=blue!40!white,draw=blue!40!white,line width=.12pt] (-0.978167,0.501243)--(-0.918440,0.415835)--(-1.131352,0.365485)--(-1.204924,0.447619)--cycle;
\path[fill=blue!40!white,draw=blue!40!white,line width=.12pt] (-1.275669,0.553222)--(-1.204924,0.447619)--(-1.420078,0.383140)--(-1.503455,0.484957)--cycle;
\path[fill=blue!41!white,draw=blue!41!white,line width=.12pt] (1.503455,0.484957)--(1.420078,0.383140)--(1.204924,0.447619)--(1.275669,0.553222)--cycle;
\path[fill=blue!29!white,draw=blue!29!white,line width=.12pt] (-1.501223,1.539458)--(-1.508470,1.411074)--(-1.777825,1.330351)--(-1.769285,1.459123)--cycle;
\path[fill=blue!34!white,draw=blue!34!white,line width=.12pt] (1.769285,1.459123)--(1.777825,1.330351)--(1.508470,1.411074)--(1.501223,1.539458)--cycle;
\path[fill=blue!25!white,draw=blue!25!white,line width=.12pt] (0.578577,1.944248)--(0.597986,1.888050)--(0.300440,1.910118)--(0.290688,1.965600)--cycle;
\path[fill=blue!24!white,draw=blue!24!white,line width=.12pt] (-0.290688,1.965600)--(-0.300440,1.910118)--(-0.597986,1.888050)--(-0.578577,1.944248)--cycle;
\path[fill=blue!42!white,draw=blue!42!white,line width=.12pt] (1.881411,0.656355)--(1.804504,0.523338)--(1.580295,0.606480)--(1.647647,0.743041)--cycle;
\path[fill=blue!38!white,draw=blue!38!white,line width=.12pt] (-1.647647,0.743041)--(-1.580295,0.606480)--(-1.804504,0.523338)--(-1.881411,0.656355)--cycle;
\path[fill=blue!38!white,draw=blue!38!white,line width=.12pt] (2.020307,1.095713)--(1.991426,0.948149)--(1.743993,1.039904)--(1.769285,1.188798)--cycle;
\path[fill=blue!33!white,draw=blue!33!white,line width=.12pt] (-1.769285,1.188798)--(-1.743993,1.039904)--(-1.991426,0.948149)--(-2.020307,1.095713)--cycle;
\path[fill=blue!39!white,draw=blue!39!white,line width=.12pt] (0.918440,0.415835)--(0.858714,0.357051)--(0.651378,0.394145)--(0.696683,0.455508)--cycle;
\path[fill=blue!39!white,draw=blue!39!white,line width=.12pt] (-0.696683,0.455508)--(-0.651378,0.394145)--(-0.858714,0.357051)--(-0.918440,0.415835)--cycle;
\path[fill=blue!39!white,draw=blue!39!white,line width=.12pt] (-1.580295,0.606480)--(-1.503455,0.484957)--(-1.716761,0.405857)--(-1.804504,0.523338)--cycle;
\path[fill=blue!42!white,draw=blue!42!white,line width=.12pt] (1.804504,0.523338)--(1.716761,0.405857)--(1.503455,0.484957)--(1.580295,0.606480)--cycle;
\path[fill=blue!25!white,draw=blue!25!white,line width=.12pt] (0.860893,1.908890)--(0.889774,1.851505)--(0.597986,1.888050)--(0.578577,1.944248)--cycle;
\path[fill=blue!23!white,draw=blue!23!white,line width=.12pt] (-0.578577,1.944248)--(-0.597986,1.888050)--(-0.889774,1.851505)--(-0.860893,1.908890)--cycle;
\path[fill=blue!31!white,draw=blue!31!white,line width=.12pt] (-1.777825,1.330351)--(-1.769285,1.188798)--(-2.020307,1.095713)--(-2.030059,1.236817)--cycle;
\path[fill=blue!36!white,draw=blue!36!white,line width=.12pt] (2.030059,1.236817)--(2.020307,1.095713)--(1.769285,1.188798)--(1.777825,1.330351)--cycle;
\path[fill=blue!37!white,draw=blue!37!white,line width=.12pt] (0.000000,0.442456)--(0.000000,0.406844)--(-0.205234,0.401803)--(-0.219943,0.437053)--cycle;
\path[fill=blue!37!white,draw=blue!37!white,line width=.12pt] (0.219943,0.437053)--(0.205234,0.401803)--(0.000000,0.406844)--(0.000000,0.442456)--cycle;
\path[fill=blue!28!white,draw=blue!28!white,line width=.12pt] (1.444914,1.736532)--(1.479763,1.649352)--(1.201283,1.715208)--(1.172992,1.800837)--cycle;
\path[fill=blue!25!white,draw=blue!25!white,line width=.12pt] (-1.172992,1.800837)--(-1.201283,1.715208)--(-1.479763,1.649352)--(-1.444914,1.736532)--cycle;
\path[fill=blue!37!white,draw=blue!37!white,line width=.12pt] (-0.219943,0.437053)--(-0.205234,0.401803)--(-0.408491,0.386728)--(-0.437769,0.420898)--cycle;
\path[fill=blue!37!white,draw=blue!37!white,line width=.12pt] (0.437769,0.420898)--(0.408491,0.386728)--(0.205234,0.401803)--(0.219943,0.437053)--cycle;
\path[fill=blue!40!white,draw=blue!40!white,line width=.12pt] (-1.204924,0.447619)--(-1.131352,0.365485)--(-1.333369,0.304943)--(-1.420078,0.383140)--cycle;
\path[fill=blue!40!white,draw=blue!40!white,line width=.12pt] (1.420078,0.383140)--(1.333369,0.304943)--(1.131352,0.365485)--(1.204924,0.447619)--cycle;
\path[fill=blue!27!white,draw=blue!27!white,line width=.12pt] (-1.479763,1.649352)--(-1.501223,1.539458)--(-1.769285,1.459123)--(-1.743993,1.570165)--cycle;
\path[fill=blue!31!white,draw=blue!31!white,line width=.12pt] (1.743993,1.570165)--(1.769285,1.459123)--(1.501223,1.539458)--(1.479763,1.649352)--cycle;
\path[fill=blue!39!white,draw=blue!39!white,line width=.12pt] (1.131352,0.365485)--(1.057780,0.309976)--(0.858714,0.357051)--(0.918440,0.415835)--cycle;
\path[fill=blue!39!white,draw=blue!39!white,line width=.12pt] (-0.918440,0.415835)--(-0.858714,0.357051)--(-1.057780,0.309976)--(-1.131352,0.365485)--cycle;
\path[fill=blue!40!white,draw=blue!40!white,line width=.12pt] (-1.503455,0.484957)--(-1.420078,0.383140)--(-1.621555,0.308427)--(-1.716761,0.405857)--cycle;
\path[fill=blue!41!white,draw=blue!41!white,line width=.12pt] (1.716761,0.405857)--(1.621555,0.308427)--(1.420078,0.383140)--(1.503455,0.484957)--cycle;
\path[fill=blue!25!white,draw=blue!25!white,line width=.12pt] (1.134919,1.859866)--(1.172992,1.800837)--(0.889774,1.851505)--(0.860893,1.908890)--cycle;
\path[fill=blue!23!white,draw=blue!23!white,line width=.12pt] (-0.860893,1.908890)--(-0.889774,1.851505)--(-1.172992,1.800837)--(-1.134919,1.859866)--cycle;
\path[fill=blue!22!white,draw=blue!22!white,line width=.12pt] (0.278805,1.991439)--(0.290688,1.965600)--(0.000000,1.972741)--(0.000000,1.998287)--cycle;
\path[fill=blue!21!white,draw=blue!21!white,line width=.12pt] (0.000000,1.998287)--(0.000000,1.972741)--(-0.290688,1.965600)--(-0.278805,1.991439)--cycle;
\path[fill=blue!37!white,draw=blue!37!white,line width=.12pt] (0.651378,0.394145)--(0.607814,0.361764)--(0.408491,0.386728)--(0.437769,0.420898)--cycle;
\path[fill=blue!37!white,draw=blue!37!white,line width=.12pt] (-0.437769,0.420898)--(-0.408491,0.386728)--(-0.607814,0.361764)--(-0.651378,0.394145)--cycle;
\path[fill=blue!41!white,draw=blue!41!white,line width=.12pt] (2.167407,0.698793)--(2.097056,0.558630)--(1.881411,0.656355)--(1.944527,0.799796)--cycle;
\path[fill=blue!37!white,draw=blue!37!white,line width=.12pt] (-1.944527,0.799796)--(-1.881411,0.656355)--(-2.097056,0.558630)--(-2.167407,0.698793)--cycle;
\path[fill=blue!34!white,draw=blue!34!white,line width=.12pt] (2.020307,1.366038)--(2.030059,1.236817)--(1.777825,1.330351)--(1.769285,1.459123)--cycle;
\path[fill=blue!28!white,draw=blue!28!white,line width=.12pt] (-1.769285,1.459123)--(-1.777825,1.330351)--(-2.030059,1.236817)--(-2.020307,1.366038)--cycle;
\path[fill=blue!22!white,draw=blue!22!white,line width=.12pt] (0.554926,1.970959)--(0.578577,1.944248)--(0.290688,1.965600)--(0.278805,1.991439)--cycle;
\path[fill=blue!21!white,draw=blue!21!white,line width=.12pt] (-0.278805,1.991439)--(-0.290688,1.965600)--(-0.578577,1.944248)--(-0.554926,1.970959)--cycle;
\path[fill=blue!40!white,draw=blue!40!white,line width=.12pt] (2.219681,0.844710)--(2.167407,0.698793)--(1.944527,0.799796)--(1.991426,0.948149)--cycle;
\path[fill=blue!35!white,draw=blue!35!white,line width=.12pt] (-1.991426,0.948149)--(-1.944527,0.799796)--(-2.167407,0.698793)--(-2.219681,0.844710)--cycle;
\path[fill=blue!42!white,draw=blue!42!white,line width=.12pt] (2.097056,0.558630)--(2.011334,0.429608)--(1.804504,0.523338)--(1.881411,0.656355)--cycle;
\path[fill=blue!38!white,draw=blue!38!white,line width=.12pt] (-1.881411,0.656355)--(-1.804504,0.523338)--(-2.011334,0.429608)--(-2.097056,0.558630)--cycle;
\path[fill=blue!37!white,draw=blue!37!white,line width=.12pt] (0.858714,0.357051)--(0.801283,0.327152)--(0.607814,0.361764)--(0.651378,0.394145)--cycle;
\path[fill=blue!38!white,draw=blue!38!white,line width=.12pt] (-0.651378,0.394145)--(-0.607814,0.361764)--(-0.801283,0.327152)--(-0.858714,0.357051)--cycle;
\path[fill=blue!33!white,draw=blue!33!white,line width=.12pt] (-2.020307,1.095713)--(-1.991426,0.948149)--(-2.219681,0.844710)--(-2.251872,0.990774)--cycle;
\path[fill=blue!38!white,draw=blue!38!white,line width=.12pt] (2.251872,0.990774)--(2.219681,0.844710)--(1.991426,0.948149)--(2.020307,1.095713)--cycle;
\path[fill=blue!22!white,draw=blue!22!white,line width=.12pt] (0.825702,1.937046)--(0.860893,1.908890)--(0.578577,1.944248)--(0.554926,1.970959)--cycle;
\path[fill=blue!21!white,draw=blue!21!white,line width=.12pt] (-0.554926,1.970959)--(-0.578577,1.944248)--(-0.860893,1.908890)--(-0.825702,1.937046)--cycle;
\path[fill=blue!39!white,draw=blue!39!white,line width=.12pt] (-1.804504,0.523338)--(-1.716761,0.405857)--(-1.913535,0.316685)--(-2.011334,0.429608)--cycle;
\path[fill=blue!42!white,draw=blue!42!white,line width=.12pt] (2.011334,0.429608)--(1.913535,0.316685)--(1.716761,0.405857)--(1.804504,0.523338)--cycle;
\path[fill=blue!39!white,draw=blue!39!white,line width=.12pt] (-1.131352,0.365485)--(-1.057780,0.309976)--(-1.246659,0.253371)--(-1.333369,0.304943)--cycle;
\path[fill=blue!39!white,draw=blue!39!white,line width=.12pt] (1.333369,0.304943)--(1.246659,0.253371)--(1.057780,0.309976)--(1.131352,0.365485)--cycle;
\path[fill=blue!40!white,draw=blue!40!white,line width=.12pt] (-1.420078,0.383140)--(-1.333369,0.304943)--(-1.522544,0.234792)--(-1.621555,0.308427)--cycle;
\path[fill=blue!40!white,draw=blue!40!white,line width=.12pt] (1.621555,0.308427)--(1.522544,0.234792)--(1.333369,0.304943)--(1.420078,0.383140)--cycle;
\path[fill=blue!28!white,draw=blue!28!white,line width=.12pt] (1.702920,1.659211)--(1.743993,1.570165)--(1.479763,1.649352)--(1.444914,1.736532)--cycle;
\path[fill=blue!24!white,draw=blue!24!white,line width=.12pt] (-1.444914,1.736532)--(-1.479763,1.649352)--(-1.743993,1.570165)--(-1.702920,1.659211)--cycle;
\path[fill=blue!25!white,draw=blue!25!white,line width=.12pt] (1.398014,1.797648)--(1.444914,1.736532)--(1.172992,1.800837)--(1.134919,1.859866)--cycle;
\path[fill=blue!22!white,draw=blue!22!white,line width=.12pt] (-1.134919,1.859866)--(-1.172992,1.800837)--(-1.444914,1.736532)--(-1.398014,1.797648)--cycle;
\path[fill=blue!35!white,draw=blue!35!white,line width=.12pt] (0.205234,0.401803)--(0.191677,0.397005)--(0.000000,0.401713)--(0.000000,0.406844)--cycle;
\path[fill=blue!35!white,draw=blue!35!white,line width=.12pt] (0.000000,0.406844)--(0.000000,0.401713)--(-0.191677,0.397005)--(-0.205234,0.401803)--cycle;
\path[fill=blue!30!white,draw=blue!30!white,line width=.12pt] (-2.030059,1.236817)--(-2.020307,1.095713)--(-2.251872,0.990774)--(-2.262742,1.131371)--cycle;
\path[fill=blue!36!white,draw=blue!36!white,line width=.12pt] (2.262742,1.131371)--(2.251872,0.990774)--(2.020307,1.095713)--(2.030059,1.236817)--cycle;
\path[fill=blue!35!white,draw=blue!35!white,line width=.12pt] (0.408491,0.386728)--(0.381508,0.382925)--(0.191677,0.397005)--(0.205234,0.401803)--cycle;
\path[fill=blue!35!white,draw=blue!35!white,line width=.12pt] (-0.205234,0.401803)--(-0.191677,0.397005)--(-0.381508,0.382925)--(-0.408491,0.386728)--cycle;
\path[fill=blue!31!white,draw=blue!31!white,line width=.12pt] (1.991426,1.478411)--(2.020307,1.366038)--(1.769285,1.459123)--(1.743993,1.570165)--cycle;
\path[fill=blue!26!white,draw=blue!26!white,line width=.12pt] (-1.743993,1.570165)--(-1.769285,1.459123)--(-2.020307,1.366038)--(-1.991426,1.478411)--cycle;
\path[fill=blue!37!white,draw=blue!37!white,line width=.12pt] (1.057780,0.309976)--(0.987036,0.283224)--(0.801283,0.327152)--(0.858714,0.357051)--cycle;
\path[fill=blue!38!white,draw=blue!38!white,line width=.12pt] (-0.858714,0.357051)--(-0.801283,0.327152)--(-0.987036,0.283224)--(-1.057780,0.309976)--cycle;
\path[fill=blue!41!white,draw=blue!41!white,line width=.12pt] (1.913535,0.316685)--(1.807416,0.224200)--(1.621555,0.308427)--(1.716761,0.405857)--cycle;
\path[fill=blue!40!white,draw=blue!40!white,line width=.12pt] (-1.716761,0.405857)--(-1.621555,0.308427)--(-1.807416,0.224200)--(-1.913535,0.316685)--cycle;
\path[fill=blue!20!white,draw=blue!20!white,line width=.12pt] (-0.825702,1.937046)--(-0.860893,1.908890)--(-1.134919,1.859866)--(-1.088526,1.890026)--cycle;
\path[fill=blue!22!white,draw=blue!22!white,line width=.12pt] (1.088526,1.890026)--(1.134919,1.859866)--(0.860893,1.908890)--(0.825702,1.937046)--cycle;
\path[fill=blue!35!white,draw=blue!35!white,line width=.12pt] (0.607814,0.361764)--(0.567664,0.359610)--(0.381508,0.382925)--(0.408491,0.386728)--cycle;
\path[fill=blue!36!white,draw=blue!36!white,line width=.12pt] (-0.408491,0.386728)--(-0.381508,0.382925)--(-0.567664,0.359610)--(-0.607814,0.361764)--cycle;
\path[fill=blue!19!white,draw=blue!19!white,line width=.12pt] (0.265249,1.986640)--(0.278805,1.991439)--(0.000000,1.998287)--(0.000000,1.993156)--cycle;
\path[fill=blue!19!white,draw=blue!19!white,line width=.12pt] (0.000000,1.993156)--(0.000000,1.998287)--(-0.278805,1.991439)--(-0.265249,1.986640)--cycle;
\path[fill=blue!19!white,draw=blue!19!white,line width=.12pt] (0.527943,1.967157)--(0.554926,1.970959)--(0.278805,1.991439)--(0.265249,1.986640)--cycle;
\path[fill=blue!19!white,draw=blue!19!white,line width=.12pt] (-0.265249,1.986640)--(-0.278805,1.991439)--(-0.554926,1.970959)--(-0.527943,1.967157)--cycle;
\path[fill=blue!39!white,draw=blue!39!white,line width=.12pt] (-1.333369,0.304943)--(-1.246659,0.253371)--(-1.423533,0.187782)--(-1.522544,0.234792)--cycle;
\path[fill=blue!39!white,draw=blue!39!white,line width=.12pt] (1.522544,0.234792)--(1.423533,0.187782)--(1.246659,0.253371)--(1.333369,0.304943)--cycle;
\path[fill=blue!28!white,draw=blue!28!white,line width=.12pt] (-2.020307,1.366038)--(-2.030059,1.236817)--(-2.262742,1.131371)--(-2.251872,1.261099)--cycle;
\path[fill=blue!34!white,draw=blue!34!white,line width=.12pt] (2.251872,1.261099)--(2.262742,1.131371)--(2.030059,1.236817)--(2.020307,1.366038)--cycle;
\path[fill=blue!41!white,draw=blue!41!white,line width=.12pt] (2.369413,0.587353)--(2.292506,0.450808)--(2.097056,0.558630)--(2.167407,0.698793)--cycle;
\path[fill=blue!36!white,draw=blue!36!white,line width=.12pt] (-2.167407,0.698793)--(-2.097056,0.558630)--(-2.292506,0.450808)--(-2.369413,0.587353)--cycle;
\path[fill=blue!34!white,draw=blue!34!white,line width=.12pt] (0.801283,0.327152)--(0.748354,0.327284)--(0.567664,0.359610)--(0.607814,0.361764)--cycle;
\path[fill=blue!36!white,draw=blue!36!white,line width=.12pt] (-0.607814,0.361764)--(-0.567664,0.359610)--(-0.748354,0.327284)--(-0.801283,0.327152)--cycle;
\path[fill=blue!42!white,draw=blue!42!white,line width=.12pt] (2.292506,0.450808)--(2.198794,0.326193)--(2.011334,0.429608)--(2.097056,0.558630)--cycle;
\path[fill=blue!38!white,draw=blue!38!white,line width=.12pt] (-2.097056,0.558630)--(-2.011334,0.429608)--(-2.198794,0.326193)--(-2.292506,0.450808)--cycle;
\path[fill=blue!37!white,draw=blue!37!white,line width=.12pt] (1.246659,0.253371)--(1.163283,0.230405)--(0.987036,0.283224)--(1.057780,0.309976)--cycle;
\path[fill=blue!38!white,draw=blue!38!white,line width=.12pt] (-1.057780,0.309976)--(-0.987036,0.283224)--(-1.163283,0.230405)--(-1.246659,0.253371)--cycle;
\path[fill=blue!40!white,draw=blue!40!white,line width=.12pt] (2.426560,0.730582)--(2.369413,0.587353)--(2.167407,0.698793)--(2.219681,0.844710)--cycle;
\path[fill=blue!34!white,draw=blue!34!white,line width=.12pt] (-2.219681,0.844710)--(-2.167407,0.698793)--(-2.369413,0.587353)--(-2.426560,0.730582)--cycle;
\path[fill=blue!25!white,draw=blue!25!white,line width=.12pt] (1.647647,1.722837)--(1.702920,1.659211)--(1.444914,1.736532)--(1.398014,1.797648)--cycle;
\path[fill=blue!22!white,draw=blue!22!white,line width=.12pt] (-1.398014,1.797648)--(-1.444914,1.736532)--(-1.702920,1.659211)--(-1.647647,1.722837)--cycle;
\path[fill=blue!40!white,draw=blue!40!white,line width=.12pt] (-1.621555,0.308427)--(-1.522544,0.234792)--(-1.697056,0.155708)--(-1.807416,0.224200)--cycle;
\path[fill=blue!40!white,draw=blue!40!white,line width=.12pt] (1.807416,0.224200)--(1.697056,0.155708)--(1.522544,0.234792)--(1.621555,0.308427)--cycle;
\path[fill=blue!18!white,draw=blue!18!white,line width=.12pt] (-0.527943,1.967157)--(-0.554926,1.970959)--(-0.825702,1.937046)--(-0.785553,1.934893)--cycle;
\path[fill=blue!19!white,draw=blue!19!white,line width=.12pt] (0.785553,1.934893)--(0.825702,1.937046)--(0.554926,1.970959)--(0.527943,1.967157)--cycle;
\path[fill=blue!28!white,draw=blue!28!white,line width=.12pt] (1.944527,1.569617)--(1.991426,1.478411)--(1.743993,1.570165)--(1.702920,1.659211)--cycle;
\path[fill=blue!23!white,draw=blue!23!white,line width=.12pt] (-1.702920,1.659211)--(-1.743993,1.570165)--(-1.991426,1.478411)--(-1.944527,1.569617)--cycle;
\path[fill=blue!22!white,draw=blue!22!white,line width=.12pt] (1.340867,1.830352)--(1.398014,1.797648)--(1.134919,1.859866)--(1.088526,1.890026)--cycle;
\path[fill=blue!20!white,draw=blue!20!white,line width=.12pt] (-1.088526,1.890026)--(-1.134919,1.859866)--(-1.398014,1.797648)--(-1.340867,1.830352)--cycle;
\path[fill=blue!39!white,draw=blue!39!white,line width=.12pt] (-2.011334,0.429608)--(-1.913535,0.316685)--(-2.091880,0.218298)--(-2.198794,0.326193)--cycle;
\path[fill=blue!42!white,draw=blue!42!white,line width=.12pt] (2.198794,0.326193)--(2.091880,0.218298)--(1.913535,0.316685)--(2.011334,0.429608)--cycle;
\path[fill=blue!32!white,draw=blue!32!white,line width=.12pt] (-2.251872,0.990774)--(-2.219681,0.844710)--(-2.426560,0.730582)--(-2.461751,0.874991)--cycle;
\path[fill=blue!38!white,draw=blue!38!white,line width=.12pt] (2.461751,0.874991)--(2.426560,0.730582)--(2.219681,0.844710)--(2.251872,0.990774)--cycle;
\path[fill=blue!32!white,draw=blue!32!white,line width=.12pt] (0.191677,0.397005)--(0.179794,0.422843)--(0.000000,0.427259)--(0.000000,0.401713)--cycle;
\path[fill=blue!33!white,draw=blue!33!white,line width=.12pt] (0.000000,0.401713)--(0.000000,0.427259)--(-0.179794,0.422843)--(-0.191677,0.397005)--cycle;
\path[fill=blue!36!white,draw=blue!36!white,line width=.12pt] (-0.801283,0.327152)--(-0.748354,0.327284)--(-0.921837,0.286259)--(-0.987036,0.283224)--cycle;
\path[fill=blue!34!white,draw=blue!34!white,line width=.12pt] (0.987036,0.283224)--(0.921837,0.286259)--(0.748354,0.327284)--(0.801283,0.327152)--cycle;
\path[fill=blue!33!white,draw=blue!33!white,line width=.12pt] (-0.191677,0.397005)--(-0.179794,0.422843)--(-0.357857,0.409636)--(-0.381508,0.382925)--cycle;
\path[fill=blue!32!white,draw=blue!32!white,line width=.12pt] (0.381508,0.382925)--(0.357857,0.409636)--(0.179794,0.422843)--(0.191677,0.397005)--cycle;
\path[fill=blue!25!white,draw=blue!25!white,line width=.12pt] (-1.991426,1.478411)--(-2.020307,1.366038)--(-2.251872,1.261099)--(-2.219681,1.374972)--cycle;
\path[fill=blue!31!white,draw=blue!31!white,line width=.12pt] (2.219681,1.374972)--(2.251872,1.261099)--(2.020307,1.366038)--(1.991426,1.478411)--cycle;
\path[fill=blue!39!white,draw=blue!39!white,line width=.12pt] (-1.913535,0.316685)--(-1.807416,0.224200)--(-1.975871,0.131270)--(-2.091880,0.218298)--cycle;
\path[fill=blue!41!white,draw=blue!41!white,line width=.12pt] (2.091880,0.218298)--(1.975871,0.131270)--(1.807416,0.224200)--(1.913535,0.316685)--cycle;
\path[fill=blue!20!white,draw=blue!20!white,line width=.12pt] (1.035598,1.890159)--(1.088526,1.890026)--(0.825702,1.937046)--(0.785553,1.934893)--cycle;
\path[fill=blue!18!white,draw=blue!18!white,line width=.12pt] (-0.785553,1.934893)--(-0.825702,1.937046)--(-1.088526,1.890026)--(-1.035598,1.890159)--cycle;
\path[fill=blue!32!white,draw=blue!32!white,line width=.12pt] (0.567664,0.359610)--(0.532473,0.387767)--(0.357857,0.409636)--(0.381508,0.382925)--cycle;
\path[fill=blue!33!white,draw=blue!33!white,line width=.12pt] (-0.381508,0.382925)--(-0.357857,0.409636)--(-0.532473,0.387767)--(-0.567664,0.359610)--cycle;
\path[fill=blue!29!white,draw=blue!29!white,line width=.12pt] (-2.262742,1.131371)--(-2.251872,0.990774)--(-2.461751,0.874991)--(-2.473633,1.015029)--cycle;
\path[fill=blue!36!white,draw=blue!36!white,line width=.12pt] (2.473633,1.015029)--(2.461751,0.874991)--(2.251872,0.990774)--(2.262742,1.131371)--cycle;
\path[fill=blue!17!white,draw=blue!17!white,line width=.12pt] (0.000000,1.957544)--(0.000000,1.993156)--(-0.265249,1.986640)--(-0.250539,1.951390)--cycle;
\path[fill=blue!17!white,draw=blue!17!white,line width=.12pt] (0.250539,1.951390)--(0.265249,1.986640)--(0.000000,1.993156)--(0.000000,1.957544)--cycle;
\path[fill=blue!37!white,draw=blue!37!white,line width=.12pt] (1.423533,0.187782)--(1.328326,0.169203)--(1.163283,0.230405)--(1.246659,0.253371)--cycle;
\path[fill=blue!38!white,draw=blue!38!white,line width=.12pt] (-1.246659,0.253371)--(-1.163283,0.230405)--(-1.328326,0.169203)--(-1.423533,0.187782)--cycle;
\path[fill=blue!39!white,draw=blue!39!white,line width=.12pt] (-1.522544,0.234792)--(-1.423533,0.187782)--(-1.586697,0.113840)--(-1.697056,0.155708)--cycle;
\path[fill=blue!39!white,draw=blue!39!white,line width=.12pt] (1.697056,0.155708)--(1.586697,0.113840)--(1.423533,0.187782)--(1.522544,0.234792)--cycle;
\path[fill=blue!17!white,draw=blue!17!white,line width=.12pt] (0.498665,1.932987)--(0.527943,1.967157)--(0.265249,1.986640)--(0.250539,1.951390)--cycle;
\path[fill=blue!17!white,draw=blue!17!white,line width=.12pt] (-0.250539,1.951390)--(-0.265249,1.986640)--(-0.527943,1.967157)--(-0.498665,1.932987)--cycle;
\path[fill=blue!36!white,draw=blue!36!white,line width=.12pt] (-0.987036,0.283224)--(-0.921837,0.286259)--(-1.086442,0.236929)--(-1.163283,0.230405)--cycle;
\path[fill=blue!34!white,draw=blue!34!white,line width=.12pt] (1.163283,0.230405)--(1.086442,0.236929)--(0.921837,0.286259)--(0.987036,0.283224)--cycle;
\path[fill=blue!34!white,draw=blue!34!white,line width=.12pt] (-0.567664,0.359610)--(-0.532473,0.387767)--(-0.701962,0.357445)--(-0.748354,0.327284)--cycle;
\path[fill=blue!32!white,draw=blue!32!white,line width=.12pt] (0.748354,0.327284)--(0.701962,0.357445)--(0.532473,0.387767)--(0.567664,0.359610)--cycle;
\path[fill=blue!19!white,draw=blue!19!white,line width=.12pt] (-1.340867,1.830352)--(-1.398014,1.797648)--(-1.647647,1.722837)--(-1.580295,1.758598)--cycle;
\path[fill=blue!22!white,draw=blue!22!white,line width=.12pt] (1.580295,1.758598)--(1.647647,1.722837)--(1.398014,1.797648)--(1.340867,1.830352)--cycle;
\path[fill=blue!21!white,draw=blue!21!white,line width=.12pt] (-1.647647,1.722837)--(-1.702920,1.659211)--(-1.944527,1.569617)--(-1.881411,1.636151)--cycle;
\path[fill=blue!25!white,draw=blue!25!white,line width=.12pt] (1.881411,1.636151)--(1.944527,1.569617)--(1.702920,1.659211)--(1.647647,1.722837)--cycle;
\path[fill=blue!40!white,draw=blue!40!white,line width=.12pt] (1.975871,0.131270)--(1.855225,0.068452)--(1.697056,0.155708)--(1.807416,0.224200)--cycle;
\path[fill=blue!40!white,draw=blue!40!white,line width=.12pt] (-1.807416,0.224200)--(-1.697056,0.155708)--(-1.855225,0.068452)--(-1.975871,0.131270)--cycle;
\path[fill=blue!17!white,draw=blue!17!white,line width=.12pt] (0.741989,1.902512)--(0.785553,1.934893)--(0.527943,1.967157)--(0.498665,1.932987)--cycle;
\path[fill=blue!17!white,draw=blue!17!white,line width=.12pt] (-0.498665,1.932987)--(-0.527943,1.967157)--(-0.785553,1.934893)--(-0.741989,1.902512)--cycle;
\path[fill=blue!20!white,draw=blue!20!white,line width=.12pt] (1.275669,1.833387)--(1.340867,1.830352)--(1.088526,1.890026)--(1.035598,1.890159)--cycle;
\path[fill=blue!18!white,draw=blue!18!white,line width=.12pt] (-1.035598,1.890159)--(-1.088526,1.890026)--(-1.340867,1.830352)--(-1.275669,1.833387)--cycle;
\path[fill=blue!34!white,draw=blue!34!white,line width=.12pt] (2.461751,1.145316)--(2.473633,1.015029)--(2.262742,1.131371)--(2.251872,1.261099)--cycle;
\path[fill=blue!27!white,draw=blue!27!white,line width=.12pt] (-2.251872,1.261099)--(-2.262742,1.131371)--(-2.473633,1.015029)--(-2.461751,1.145316)--cycle;
\path[fill=blue!41!white,draw=blue!41!white,line width=.12pt] (2.548600,0.466550)--(2.465877,0.333925)--(2.292506,0.450808)--(2.369413,0.587353)--cycle;
\path[fill=blue!35!white,draw=blue!35!white,line width=.12pt] (-2.369413,0.587353)--(-2.292506,0.450808)--(-2.465877,0.333925)--(-2.548600,0.466550)--cycle;
\path[fill=blue!42!white,draw=blue!42!white,line width=.12pt] (2.465877,0.333925)--(2.365079,0.214089)--(2.198794,0.326193)--(2.292506,0.450808)--cycle;
\path[fill=blue!37!white,draw=blue!37!white,line width=.12pt] (-2.292506,0.450808)--(-2.198794,0.326193)--(-2.365079,0.214089)--(-2.465877,0.333925)--cycle;
\path[fill=blue!28!white,draw=blue!28!white,line width=.12pt] (2.167407,1.468614)--(2.219681,1.374972)--(1.991426,1.478411)--(1.944527,1.569617)--cycle;
\path[fill=blue!23!white,draw=blue!23!white,line width=.12pt] (-1.944527,1.569617)--(-1.991426,1.478411)--(-2.219681,1.374972)--(-2.167407,1.468614)--cycle;
\path[fill=blue!30!white,draw=blue!30!white,line width=.12pt] (0.179794,0.422843)--(0.170043,0.478325)--(0.000000,0.482502)--(0.000000,0.427259)--cycle;
\path[fill=blue!30!white,draw=blue!30!white,line width=.12pt] (0.000000,0.427259)--(0.000000,0.482502)--(-0.170043,0.478325)--(-0.179794,0.422843)--cycle;
\path[fill=blue!32!white,draw=blue!32!white,line width=.12pt] (0.921837,0.286259)--(0.864690,0.318963)--(0.701962,0.357445)--(0.748354,0.327284)--cycle;
\path[fill=blue!34!white,draw=blue!34!white,line width=.12pt] (-0.748354,0.327284)--(-0.701962,0.357445)--(-0.864690,0.318963)--(-0.921837,0.286259)--cycle;
\path[fill=blue!40!white,draw=blue!40!white,line width=.12pt] (2.610069,0.606865)--(2.548600,0.466550)--(2.369413,0.587353)--(2.426560,0.730582)--cycle;
\path[fill=blue!33!white,draw=blue!33!white,line width=.12pt] (-2.426560,0.730582)--(-2.369413,0.587353)--(-2.548600,0.466550)--(-2.610069,0.606865)--cycle;
\path[fill=blue!42!white,draw=blue!42!white,line width=.12pt] (2.365079,0.214089)--(2.250079,0.111645)--(2.091880,0.218298)--(2.198794,0.326193)--cycle;
\path[fill=blue!38!white,draw=blue!38!white,line width=.12pt] (-2.198794,0.326193)--(-2.091880,0.218298)--(-2.250079,0.111645)--(-2.365079,0.214089)--cycle;
\path[fill=blue!29!white,draw=blue!29!white,line width=.12pt] (0.357857,0.409636)--(0.338447,0.465835)--(0.170043,0.478325)--(0.179794,0.422843)--cycle;
\path[fill=blue!30!white,draw=blue!30!white,line width=.12pt] (-0.179794,0.422843)--(-0.170043,0.478325)--(-0.338447,0.465835)--(-0.357857,0.409636)--cycle;
\path[fill=blue!38!white,draw=blue!38!white,line width=.12pt] (-1.423533,0.187782)--(-1.328326,0.169203)--(-1.480578,0.100206)--(-1.586697,0.113840)--cycle;
\path[fill=blue!37!white,draw=blue!37!white,line width=.12pt] (1.586697,0.113840)--(1.480578,0.100206)--(1.328326,0.169203)--(1.423533,0.187782)--cycle;
\path[fill=blue!16!white,draw=blue!16!white,line width=.12pt] (-0.741989,1.902512)--(-0.785553,1.934893)--(-1.035598,1.890159)--(-0.978167,1.860259)--cycle;
\path[fill=blue!17!white,draw=blue!17!white,line width=.12pt] (0.978167,1.860259)--(1.035598,1.890159)--(0.785553,1.934893)--(0.741989,1.902512)--cycle;
\path[fill=blue!37!white,draw=blue!37!white,line width=.12pt] (-1.163283,0.230405)--(-1.086442,0.236929)--(-1.240584,0.179769)--(-1.328326,0.169203)--cycle;
\path[fill=blue!34!white,draw=blue!34!white,line width=.12pt] (1.328326,0.169203)--(1.240584,0.179769)--(1.086442,0.236929)--(1.163283,0.230405)--cycle;
\path[fill=blue!29!white,draw=blue!29!white,line width=.12pt] (0.532473,0.387767)--(0.503593,0.445151)--(0.338447,0.465835)--(0.357857,0.409636)--cycle;
\path[fill=blue!31!white,draw=blue!31!white,line width=.12pt] (-0.357857,0.409636)--(-0.338447,0.465835)--(-0.503593,0.445151)--(-0.532473,0.387767)--cycle;
\path[fill=blue!39!white,draw=blue!39!white,line width=.12pt] (1.855225,0.068452)--(1.734580,0.032258)--(1.586697,0.113840)--(1.697056,0.155708)--cycle;
\path[fill=blue!39!white,draw=blue!39!white,line width=.12pt] (-1.697056,0.155708)--(-1.586697,0.113840)--(-1.734580,0.032258)--(-1.855225,0.068452)--cycle;
\path[fill=blue!38!white,draw=blue!38!white,line width=.12pt] (2.647922,0.749480)--(2.610069,0.606865)--(2.426560,0.730582)--(2.461751,0.874991)--cycle;
\path[fill=blue!31!white,draw=blue!31!white,line width=.12pt] (-2.461751,0.874991)--(-2.426560,0.730582)--(-2.610069,0.606865)--(-2.647922,0.749480)--cycle;
\path[fill=blue!15!white,draw=blue!15!white,line width=.12pt] (0.000000,1.892820)--(0.000000,1.957544)--(-0.250539,1.951390)--(-0.235241,1.887042)--cycle;
\path[fill=blue!15!white,draw=blue!15!white,line width=.12pt] (0.235241,1.887042)--(0.250539,1.951390)--(0.000000,1.957544)--(0.000000,1.892820)--cycle;
\path[fill=blue!41!white,draw=blue!41!white,line width=.12pt] (2.250079,0.111645)--(2.125296,0.030531)--(1.975871,0.131270)--(2.091880,0.218298)--cycle;
\path[fill=blue!39!white,draw=blue!39!white,line width=.12pt] (-2.091880,0.218298)--(-1.975871,0.131270)--(-2.125296,0.030531)--(-2.250079,0.111645)--cycle;
\path[fill=blue!15!white,draw=blue!15!white,line width=.12pt] (-0.235241,1.887042)--(-0.250539,1.951390)--(-0.498665,1.932987)--(-0.468217,1.869763)--cycle;
\path[fill=blue!15!white,draw=blue!15!white,line width=.12pt] (0.468217,1.869763)--(0.498665,1.932987)--(0.250539,1.951390)--(0.235241,1.887042)--cycle;
\path[fill=blue!31!white,draw=blue!31!white,line width=.12pt] (2.426560,1.260844)--(2.461751,1.145316)--(2.251872,1.261099)--(2.219681,1.374972)--cycle;
\path[fill=blue!24!white,draw=blue!24!white,line width=.12pt] (-2.219681,1.374972)--(-2.251872,1.261099)--(-2.461751,1.145316)--(-2.426560,1.260844)--cycle;
\path[fill=blue!35!white,draw=blue!35!white,line width=.12pt] (-0.921837,0.286259)--(-0.864690,0.318963)--(-1.019091,0.272690)--(-1.086442,0.236929)--cycle;
\path[fill=blue!32!white,draw=blue!32!white,line width=.12pt] (1.086442,0.236929)--(1.019091,0.272690)--(0.864690,0.318963)--(0.921837,0.286259)--cycle;
\path[fill=blue!22!white,draw=blue!22!white,line width=.12pt] (1.804504,1.675456)--(1.881411,1.636151)--(1.647647,1.722837)--(1.580295,1.758598)--cycle;
\path[fill=blue!19!white,draw=blue!19!white,line width=.12pt] (-1.580295,1.758598)--(-1.647647,1.722837)--(-1.881411,1.636151)--(-1.804504,1.675456)--cycle;
\path[fill=blue!18!white,draw=blue!18!white,line width=.12pt] (-1.275669,1.833387)--(-1.340867,1.830352)--(-1.580295,1.758598)--(-1.503455,1.765122)--cycle;
\path[fill=blue!20!white,draw=blue!20!white,line width=.12pt] (1.503455,1.765122)--(1.580295,1.758598)--(1.340867,1.830352)--(1.275669,1.833387)--cycle;
\path[fill=blue!29!white,draw=blue!29!white,line width=.12pt] (0.701962,0.357445)--(0.663889,0.416474)--(0.503593,0.445151)--(0.532473,0.387767)--cycle;
\path[fill=blue!31!white,draw=blue!31!white,line width=.12pt] (-0.532473,0.387767)--(-0.503593,0.445151)--(-0.663889,0.416474)--(-0.701962,0.357445)--cycle;
\path[fill=blue!36!white,draw=blue!36!white,line width=.12pt] (2.660703,0.888912)--(2.647922,0.749480)--(2.461751,0.874991)--(2.473633,1.015029)--cycle;
\path[fill=blue!29!white,draw=blue!29!white,line width=.12pt] (-2.473633,1.015029)--(-2.461751,0.874991)--(-2.647922,0.749480)--(-2.660703,0.888912)--cycle;
\path[fill=blue!15!white,draw=blue!15!white,line width=.12pt] (0.696683,1.841149)--(0.741989,1.902512)--(0.498665,1.932987)--(0.468217,1.869763)--cycle;
\path[fill=blue!15!white,draw=blue!15!white,line width=.12pt] (-0.468217,1.869763)--(-0.498665,1.932987)--(-0.741989,1.902512)--(-0.696683,1.841149)--cycle;
\path[fill=blue!25!white,draw=blue!25!white,line width=.12pt] (2.097056,1.538426)--(2.167407,1.468614)--(1.944527,1.569617)--(1.881411,1.636151)--cycle;
\path[fill=blue!21!white,draw=blue!21!white,line width=.12pt] (-1.881411,1.636151)--(-1.944527,1.569617)--(-2.167407,1.468614)--(-2.097056,1.538426)--cycle;
\path[fill=blue!17!white,draw=blue!17!white,line width=.12pt] (1.204924,1.806635)--(1.275669,1.833387)--(1.035598,1.890159)--(0.978167,1.860259)--cycle;
\path[fill=blue!16!white,draw=blue!16!white,line width=.12pt] (-0.978167,1.860259)--(-1.035598,1.890159)--(-1.275669,1.833387)--(-1.204924,1.806635)--cycle;
\path[fill=blue!40!white,draw=blue!40!white,line width=.12pt] (2.125296,0.030531)--(1.995527,-0.026136)--(1.855225,0.068452)--(1.975871,0.131270)--cycle;
\path[fill=blue!40!white,draw=blue!40!white,line width=.12pt] (-1.975871,0.131270)--(-1.855225,0.068452)--(-1.995527,-0.026136)--(-2.125296,0.030531)--cycle;
\path[fill=blue!27!white,draw=blue!27!white,line width=.12pt] (0.000000,0.482502)--(0.000000,0.565317)--(-0.162796,0.561318)--(-0.170043,0.478325)--cycle;
\path[fill=blue!27!white,draw=blue!27!white,line width=.12pt] (0.170043,0.478325)--(0.162796,0.561318)--(0.000000,0.565317)--(0.000000,0.482502)--cycle;
\path[fill=blue!37!white,draw=blue!37!white,line width=.12pt] (-1.328326,0.169203)--(-1.240584,0.179769)--(-1.382778,0.115330)--(-1.480578,0.100206)--cycle;
\path[fill=blue!34!white,draw=blue!34!white,line width=.12pt] (1.480578,0.100206)--(1.382778,0.115330)--(1.240584,0.179769)--(1.328326,0.169203)--cycle;
\path[fill=blue!29!white,draw=blue!29!white,line width=.12pt] (0.864690,0.318963)--(0.817791,0.380079)--(0.663889,0.416474)--(0.701962,0.357445)--cycle;
\path[fill=blue!32!white,draw=blue!32!white,line width=.12pt] (-0.701962,0.357445)--(-0.663889,0.416474)--(-0.817791,0.380079)--(-0.864690,0.318963)--cycle;
\path[fill=blue!39!white,draw=blue!39!white,line width=.12pt] (-1.586697,0.113840)--(-1.480578,0.100206)--(-1.618570,0.024081)--(-1.734580,0.032258)--cycle;
\path[fill=blue!37!white,draw=blue!37!white,line width=.12pt] (1.734580,0.032258)--(1.618570,0.024081)--(1.480578,0.100206)--(1.586697,0.113840)--cycle;
\path[fill=blue!26!white,draw=blue!26!white,line width=.12pt] (0.338447,0.465835)--(0.324025,0.549360)--(0.162796,0.561318)--(0.170043,0.478325)--cycle;
\path[fill=blue!28!white,draw=blue!28!white,line width=.12pt] (-0.170043,0.478325)--(-0.162796,0.561318)--(-0.324025,0.549360)--(-0.338447,0.465835)--cycle;
\path[fill=blue!32!white,draw=blue!32!white,line width=.12pt] (1.240584,0.179769)--(1.163677,0.219074)--(1.019091,0.272690)--(1.086442,0.236929)--cycle;
\path[fill=blue!35!white,draw=blue!35!white,line width=.12pt] (-1.086442,0.236929)--(-1.019091,0.272690)--(-1.163677,0.219074)--(-1.240584,0.179769)--cycle;
\path[fill=blue!15!white,draw=blue!15!white,line width=.12pt] (0.918440,1.801476)--(0.978167,1.860259)--(0.741989,1.902512)--(0.696683,1.841149)--cycle;
\path[fill=blue!15!white,draw=blue!15!white,line width=.12pt] (-0.696683,1.841149)--(-0.741989,1.902512)--(-0.978167,1.860259)--(-0.918440,1.801476)--cycle;
\path[fill=blue!28!white,draw=blue!28!white,line width=.12pt] (2.369413,1.357174)--(2.426560,1.260844)--(2.219681,1.374972)--(2.167407,1.468614)--cycle;
\path[fill=blue!22!white,draw=blue!22!white,line width=.12pt] (-2.167407,1.468614)--(-2.219681,1.374972)--(-2.426560,1.260844)--(-2.369413,1.357174)--cycle;
\path[fill=blue!28!white,draw=blue!28!white,line width=.12pt] (-0.338447,0.465835)--(-0.324025,0.549360)--(-0.482133,0.529559)--(-0.503593,0.445151)--cycle;
\path[fill=blue!26!white,draw=blue!26!white,line width=.12pt] (0.503593,0.445151)--(0.482133,0.529559)--(0.324025,0.549360)--(0.338447,0.465835)--cycle;
\path[fill=blue!14!white,draw=blue!14!white,line width=.12pt] (0.219943,1.796069)--(0.235241,1.887042)--(0.000000,1.892820)--(0.000000,1.801472)--cycle;
\path[fill=blue!14!white,draw=blue!14!white,line width=.12pt] (0.000000,1.801472)--(0.000000,1.892820)--(-0.235241,1.887042)--(-0.219943,1.796069)--cycle;
\path[fill=blue!33!white,draw=blue!33!white,line width=.12pt] (2.647922,1.019805)--(2.660703,0.888912)--(2.473633,1.015029)--(2.461751,1.145316)--cycle;
\path[fill=blue!26!white,draw=blue!26!white,line width=.12pt] (-2.461751,1.145316)--(-2.473633,1.015029)--(-2.660703,0.888912)--(-2.647922,1.019805)--cycle;
\path[fill=blue!37!white,draw=blue!37!white,line width=.12pt] (-2.465877,0.333925)--(-2.365079,0.214089)--(-2.508586,0.094375)--(-2.615501,0.209109)--cycle;
\path[fill=blue!41!white,draw=blue!41!white,line width=.12pt] (2.615501,0.209109)--(2.508586,0.094375)--(2.365079,0.214089)--(2.465877,0.333925)--cycle;
\path[fill=blue!35!white,draw=blue!35!white,line width=.12pt] (-2.548600,0.466550)--(-2.465877,0.333925)--(-2.615501,0.209109)--(-2.703244,0.337547)--cycle;
\path[fill=blue!41!white,draw=blue!41!white,line width=.12pt] (2.703244,0.337547)--(2.615501,0.209109)--(2.465877,0.333925)--(2.548600,0.466550)--cycle;
\path[fill=blue!39!white,draw=blue!39!white,line width=.12pt] (1.995527,-0.026136)--(1.865758,-0.056178)--(1.734580,0.032258)--(1.855225,0.068452)--cycle;
\path[fill=blue!40!white,draw=blue!40!white,line width=.12pt] (-1.855225,0.068452)--(-1.734580,0.032258)--(-1.865758,-0.056178)--(-1.995527,-0.026136)--cycle;
\path[fill=blue!14!white,draw=blue!14!white,line width=.12pt] (0.437769,1.779914)--(0.468217,1.869763)--(0.235241,1.887042)--(0.219943,1.796069)--cycle;
\path[fill=blue!14!white,draw=blue!14!white,line width=.12pt] (-0.219943,1.796069)--(-0.235241,1.887042)--(-0.468217,1.869763)--(-0.437769,1.779914)--cycle;
\path[fill=blue!38!white,draw=blue!38!white,line width=.12pt] (-2.365079,0.214089)--(-2.250079,0.111645)--(-2.386608,-0.002248)--(-2.508586,0.094375)--cycle;
\path[fill=blue!42!white,draw=blue!42!white,line width=.12pt] (2.508586,0.094375)--(2.386608,-0.002248)--(2.250079,0.111645)--(2.365079,0.214089)--cycle;
\path[fill=blue!20!white,draw=blue!20!white,line width=.12pt] (1.716761,1.686022)--(1.804504,1.675456)--(1.580295,1.758598)--(1.503455,1.765122)--cycle;
\path[fill=blue!17!white,draw=blue!17!white,line width=.12pt] (-1.503455,1.765122)--(-1.580295,1.758598)--(-1.804504,1.675456)--(-1.716761,1.686022)--cycle;
\path[fill=blue!29!white,draw=blue!29!white,line width=.12pt] (1.019091,0.272690)--(0.963817,0.336316)--(0.817791,0.380079)--(0.864690,0.318963)--cycle;
\path[fill=blue!32!white,draw=blue!32!white,line width=.12pt] (-0.864690,0.318963)--(-0.817791,0.380079)--(-0.963817,0.336316)--(-1.019091,0.272690)--cycle;
\path[fill=blue!16!white,draw=blue!16!white,line width=.12pt] (-1.204924,1.806635)--(-1.275669,1.833387)--(-1.503455,1.765122)--(-1.420078,1.742156)--cycle;
\path[fill=blue!17!white,draw=blue!17!white,line width=.12pt] (1.420078,1.742156)--(1.503455,1.765122)--(1.275669,1.833387)--(1.204924,1.806635)--cycle;
\path[fill=blue!29!white,draw=blue!29!white,line width=.12pt] (-0.503593,0.445151)--(-0.482133,0.529559)--(-0.635598,0.502103)--(-0.663889,0.416474)--cycle;
\path[fill=blue!26!white,draw=blue!26!white,line width=.12pt] (0.663889,0.416474)--(0.635598,0.502103)--(0.482133,0.529559)--(0.503593,0.445151)--cycle;
\path[fill=blue!33!white,draw=blue!33!white,line width=.12pt] (-2.610069,0.606865)--(-2.548600,0.466550)--(-2.703244,0.337547)--(-2.768442,0.474751)--cycle;
\path[fill=blue!39!white,draw=blue!39!white,line width=.12pt] (2.768442,0.474751)--(2.703244,0.337547)--(2.548600,0.466550)--(2.610069,0.606865)--cycle;
\path[fill=blue!19!white,draw=blue!19!white,line width=.12pt] (-1.804504,1.675456)--(-1.881411,1.636151)--(-2.097056,1.538426)--(-2.011334,1.581726)--cycle;
\path[fill=blue!22!white,draw=blue!22!white,line width=.12pt] (2.011334,1.581726)--(2.097056,1.538426)--(1.881411,1.636151)--(1.804504,1.675456)--cycle;
\path[fill=blue!14!white,draw=blue!14!white,line width=.12pt] (0.651378,1.753160)--(0.696683,1.841149)--(0.468217,1.869763)--(0.437769,1.779914)--cycle;
\path[fill=blue!14!white,draw=blue!14!white,line width=.12pt] (-0.437769,1.779914)--(-0.468217,1.869763)--(-0.696683,1.841149)--(-0.651378,1.753160)--cycle;
\path[fill=blue!39!white,draw=blue!39!white,line width=.12pt] (-2.250079,0.111645)--(-2.125296,0.030531)--(-2.254254,-0.077046)--(-2.386608,-0.002248)--cycle;
\path[fill=blue!41!white,draw=blue!41!white,line width=.12pt] (2.386608,-0.002248)--(2.254254,-0.077046)--(2.125296,0.030531)--(2.250079,0.111645)--cycle;
\path[fill=blue!15!white,draw=blue!15!white,line width=.12pt] (1.131352,1.751126)--(1.204924,1.806635)--(0.978167,1.860259)--(0.918440,1.801476)--cycle;
\path[fill=blue!15!white,draw=blue!15!white,line width=.12pt] (-0.918440,1.801476)--(-0.978167,1.860259)--(-1.204924,1.806635)--(-1.131352,1.751126)--cycle;
\path[fill=blue!24!white,draw=blue!24!white,line width=.12pt] (0.162796,0.561318)--(0.158334,0.668634)--(0.000000,0.672523)--(0.000000,0.565317)--cycle;
\path[fill=blue!24!white,draw=blue!24!white,line width=.12pt] (0.000000,0.565317)--(0.000000,0.672523)--(-0.158334,0.668634)--(-0.162796,0.561318)--cycle;
\path[fill=blue!34!white,draw=blue!34!white,line width=.12pt] (1.618570,0.024081)--(1.511656,0.044233)--(1.382778,0.115330)--(1.480578,0.100206)--cycle;
\path[fill=blue!38!white,draw=blue!38!white,line width=.12pt] (-1.480578,0.100206)--(-1.382778,0.115330)--(-1.511656,0.044233)--(-1.618570,0.024081)--cycle;
\path[fill=blue!35!white,draw=blue!35!white,line width=.12pt] (-1.240584,0.179769)--(-1.163677,0.219074)--(-1.297056,0.158630)--(-1.382778,0.115330)--cycle;
\path[fill=blue!32!white,draw=blue!32!white,line width=.12pt] (1.382778,0.115330)--(1.297056,0.158630)--(1.163677,0.219074)--(1.240584,0.179769)--cycle;
\path[fill=blue!30!white,draw=blue!30!white,line width=.12pt] (-2.647922,0.749480)--(-2.610069,0.606865)--(-2.768442,0.474751)--(-2.808591,0.615449)--cycle;
\path[fill=blue!38!white,draw=blue!38!white,line width=.12pt] (2.808591,0.615449)--(2.768442,0.474751)--(2.610069,0.606865)--(2.647922,0.749480)--cycle;
\path[fill=blue!24!white,draw=blue!24!white,line width=.12pt] (0.324025,0.549360)--(0.315143,0.657004)--(0.158334,0.668634)--(0.162796,0.561318)--cycle;
\path[fill=blue!25!white,draw=blue!25!white,line width=.12pt] (-0.162796,0.561318)--(-0.158334,0.668634)--(-0.315143,0.657004)--(-0.324025,0.549360)--cycle;
\path[fill=blue!26!white,draw=blue!26!white,line width=.12pt] (0.817791,0.380079)--(0.782941,0.467259)--(0.635598,0.502103)--(0.663889,0.416474)--cycle;
\path[fill=blue!29!white,draw=blue!29!white,line width=.12pt] (-0.663889,0.416474)--(-0.635598,0.502103)--(-0.782941,0.467259)--(-0.817791,0.380079)--cycle;
\path[fill=blue!31!white,draw=blue!31!white,line width=.12pt] (2.610069,1.137127)--(2.647922,1.019805)--(2.461751,1.145316)--(2.426560,1.260844)--cycle;
\path[fill=blue!24!white,draw=blue!24!white,line width=.12pt] (-2.426560,1.260844)--(-2.461751,1.145316)--(-2.647922,1.019805)--(-2.610069,1.137127)--cycle;
\path[fill=blue!25!white,draw=blue!25!white,line width=.12pt] (2.292506,1.430603)--(2.369413,1.357174)--(2.167407,1.468614)--(2.097056,1.538426)--cycle;
\path[fill=blue!20!white,draw=blue!20!white,line width=.12pt] (-2.097056,1.538426)--(-2.167407,1.468614)--(-2.369413,1.357174)--(-2.292506,1.430603)--cycle;
\path[fill=blue!37!white,draw=blue!37!white,line width=.12pt] (1.865758,-0.056178)--(1.740975,-0.058441)--(1.618570,0.024081)--(1.734580,0.032258)--cycle;
\path[fill=blue!39!white,draw=blue!39!white,line width=.12pt] (-1.734580,0.032258)--(-1.618570,0.024081)--(-1.740975,-0.058441)--(-1.865758,-0.056178)--cycle;
\path[fill=blue!40!white,draw=blue!40!white,line width=.12pt] (2.254254,-0.077046)--(2.116611,-0.127144)--(1.995527,-0.026136)--(2.125296,0.030531)--cycle;
\path[fill=blue!40!white,draw=blue!40!white,line width=.12pt] (-2.125296,0.030531)--(-1.995527,-0.026136)--(-2.116611,-0.127144)--(-2.254254,-0.077046)--cycle;
\path[fill=blue!14!white,draw=blue!14!white,line width=.12pt] (0.858714,1.716067)--(0.918440,1.801476)--(0.696683,1.841149)--(0.651378,1.753160)--cycle;
\path[fill=blue!14!white,draw=blue!14!white,line width=.12pt] (-0.651378,1.753160)--(-0.696683,1.841149)--(-0.918440,1.801476)--(-0.858714,1.716067)--cycle;
\path[fill=blue!13!white,draw=blue!13!white,line width=.12pt] (0.000000,1.687009)--(0.000000,1.801472)--(-0.219943,1.796069)--(-0.205234,1.681968)--cycle;
\path[fill=blue!13!white,draw=blue!13!white,line width=.12pt] (0.205234,1.681968)--(0.219943,1.796069)--(0.000000,1.801472)--(0.000000,1.687009)--cycle;
\path[fill=blue!33!white,draw=blue!33!white,line width=.12pt] (-1.019091,0.272690)--(-0.963817,0.336316)--(-1.100561,0.285608)--(-1.163677,0.219074)--cycle;
\path[fill=blue!29!white,draw=blue!29!white,line width=.12pt] (1.163677,0.219074)--(1.100561,0.285608)--(0.963817,0.336316)--(1.019091,0.272690)--cycle;
\path[fill=blue!25!white,draw=blue!25!white,line width=.12pt] (-0.324025,0.549360)--(-0.315143,0.657004)--(-0.468918,0.637745)--(-0.482133,0.529559)--cycle;
\path[fill=blue!23!white,draw=blue!23!white,line width=.12pt] (0.482133,0.529559)--(0.468918,0.637745)--(0.315143,0.657004)--(0.324025,0.549360)--cycle;
\path[fill=blue!13!white,draw=blue!13!white,line width=.12pt] (0.408491,1.666893)--(0.437769,1.779914)--(0.219943,1.796069)--(0.205234,1.681968)--cycle;
\path[fill=blue!13!white,draw=blue!13!white,line width=.12pt] (-0.205234,1.681968)--(-0.219943,1.796069)--(-0.437769,1.779914)--(-0.408491,1.666893)--cycle;
\path[fill=blue!28!white,draw=blue!28!white,line width=.12pt] (-2.660703,0.888912)--(-2.647922,0.749480)--(-2.808591,0.615449)--(-2.822148,0.754235)--cycle;
\path[fill=blue!36!white,draw=blue!36!white,line width=.12pt] (2.822148,0.754235)--(2.808591,0.615449)--(2.647922,0.749480)--(2.660703,0.888912)--cycle;
\path[fill=blue!23!white,draw=blue!23!white,line width=.12pt] (0.635598,0.502103)--(0.618176,0.611042)--(0.468918,0.637745)--(0.482133,0.529559)--cycle;
\path[fill=blue!26!white,draw=blue!26!white,line width=.12pt] (-0.482133,0.529559)--(-0.468918,0.637745)--(-0.618176,0.611042)--(-0.635598,0.502103)--cycle;
\path[fill=blue!30!white,draw=blue!30!white,line width=.12pt] (-0.817791,0.380079)--(-0.782941,0.467259)--(-0.922745,0.425362)--(-0.963817,0.336316)--cycle;
\path[fill=blue!26!white,draw=blue!26!white,line width=.12pt] (0.963817,0.336316)--(0.922745,0.425362)--(0.782941,0.467259)--(0.817791,0.380079)--cycle;
\path[fill=blue!16!white,draw=blue!16!white,line width=.12pt] (-1.420078,1.742156)--(-1.503455,1.765122)--(-1.716761,1.686022)--(-1.621555,1.667443)--cycle;
\path[fill=blue!17!white,draw=blue!17!white,line width=.12pt] (1.621555,1.667443)--(1.716761,1.686022)--(1.503455,1.765122)--(1.420078,1.742156)--cycle;
\path[fill=blue!15!white,draw=blue!15!white,line width=.12pt] (-1.131352,1.751126)--(-1.204924,1.806635)--(-1.420078,1.742156)--(-1.333369,1.690584)--cycle;
\path[fill=blue!15!white,draw=blue!15!white,line width=.12pt] (1.333369,1.690584)--(1.420078,1.742156)--(1.204924,1.806635)--(1.131352,1.751126)--cycle;
\path[fill=blue!13!white,draw=blue!13!white,line width=.12pt] (0.607814,1.641929)--(0.651378,1.753160)--(0.437769,1.779914)--(0.408491,1.666893)--cycle;
\path[fill=blue!13!white,draw=blue!13!white,line width=.12pt] (-0.408491,1.666893)--(-0.437769,1.779914)--(-0.651378,1.753160)--(-0.607814,1.641929)--cycle;
\path[fill=blue!40!white,draw=blue!40!white,line width=.12pt] (-1.995527,-0.026136)--(-1.865758,-0.056178)--(-1.978968,-0.150618)--(-2.116611,-0.127144)--cycle;
\path[fill=blue!39!white,draw=blue!39!white,line width=.12pt] (2.116611,-0.127144)--(1.978968,-0.150618)--(1.865758,-0.056178)--(1.995527,-0.026136)--cycle;
\path[fill=blue!22!white,draw=blue!22!white,line width=.12pt] (0.000000,0.672523)--(0.000000,0.800000)--(-0.156827,0.796148)--(-0.158334,0.668634)--cycle;
\path[fill=blue!21!white,draw=blue!21!white,line width=.12pt] (0.158334,0.668634)--(0.156827,0.796148)--(0.000000,0.800000)--(0.000000,0.672523)--cycle;
\path[fill=blue!17!white,draw=blue!17!white,line width=.12pt] (-1.716761,1.686022)--(-1.804504,1.675456)--(-2.011334,1.581726)--(-1.913535,1.596850)--cycle;
\path[fill=blue!20!white,draw=blue!20!white,line width=.12pt] (1.913535,1.596850)--(2.011334,1.581726)--(1.804504,1.675456)--(1.716761,1.686022)--cycle;
\path[fill=blue!14!white,draw=blue!14!white,line width=.12pt] (1.057780,1.668992)--(1.131352,1.751126)--(0.918440,1.801476)--(0.858714,1.716067)--cycle;
\path[fill=blue!14!white,draw=blue!14!white,line width=.12pt] (-0.858714,1.716067)--(-0.918440,1.801476)--(-1.131352,1.751126)--(-1.057780,1.668992)--cycle;
\path[fill=blue!21!white,draw=blue!21!white,line width=.12pt] (0.315143,0.657004)--(0.312145,0.784628)--(0.156827,0.796148)--(0.158334,0.668634)--cycle;
\path[fill=blue!22!white,draw=blue!22!white,line width=.12pt] (-0.158334,0.668634)--(-0.156827,0.796148)--(-0.312145,0.784628)--(-0.315143,0.657004)--cycle;
\path[fill=blue!28!white,draw=blue!28!white,line width=.12pt] (2.548600,1.236371)--(2.610069,1.137127)--(2.426560,1.260844)--(2.369413,1.357174)--cycle;
\path[fill=blue!21!white,draw=blue!21!white,line width=.12pt] (-2.369413,1.357174)--(-2.426560,1.260844)--(-2.610069,1.137127)--(-2.548600,1.236371)--cycle;
\path[fill=blue!32!white,draw=blue!32!white,line width=.12pt] (1.511656,0.044233)--(1.417944,0.091940)--(1.297056,0.158630)--(1.382778,0.115330)--cycle;
\path[fill=blue!36!white,draw=blue!36!white,line width=.12pt] (-1.382778,0.115330)--(-1.297056,0.158630)--(-1.417944,0.091940)--(-1.511656,0.044233)--cycle;
\path[fill=blue!13!white,draw=blue!13!white,line width=.12pt] (0.191677,1.549123)--(0.205234,1.681968)--(0.000000,1.687009)--(0.000000,1.553831)--cycle;
\path[fill=blue!13!white,draw=blue!13!white,line width=.12pt] (0.000000,1.553831)--(0.000000,1.687009)--(-0.205234,1.681968)--(-0.191677,1.549123)--cycle;
\path[fill=blue!23!white,draw=blue!23!white,line width=.12pt] (0.782941,0.467259)--(0.761481,0.577153)--(0.618176,0.611042)--(0.635598,0.502103)--cycle;
\path[fill=blue!27!white,draw=blue!27!white,line width=.12pt] (-0.635598,0.502103)--(-0.618176,0.611042)--(-0.761481,0.577153)--(-0.782941,0.467259)--cycle;
\path[fill=blue!29!white,draw=blue!29!white,line width=.12pt] (1.297056,0.158630)--(1.226706,0.228443)--(1.100561,0.285608)--(1.163677,0.219074)--cycle;
\path[fill=blue!33!white,draw=blue!33!white,line width=.12pt] (-1.163677,0.219074)--(-1.100561,0.285608)--(-1.226706,0.228443)--(-1.297056,0.158630)--cycle;
\path[fill=blue!35!white,draw=blue!35!white,line width=.12pt] (1.740975,-0.058441)--(1.625975,-0.032838)--(1.511656,0.044233)--(1.618570,0.024081)--cycle;
\path[fill=blue!38!white,draw=blue!38!white,line width=.12pt] (-1.618570,0.024081)--(-1.511656,0.044233)--(-1.625975,-0.032838)--(-1.740975,-0.058441)--cycle;
\path[fill=blue!23!white,draw=blue!23!white,line width=.12pt] (-0.315143,0.657004)--(-0.312145,0.784628)--(-0.464455,0.765552)--(-0.468918,0.637745)--cycle;
\path[fill=blue!21!white,draw=blue!21!white,line width=.12pt] (0.468918,0.637745)--(0.464455,0.765552)--(0.312145,0.784628)--(0.315143,0.657004)--cycle;
\path[fill=blue!13!white,draw=blue!13!white,line width=.12pt] (0.801283,1.607317)--(0.858714,1.716067)--(0.651378,1.753160)--(0.607814,1.641929)--cycle;
\path[fill=blue!14!white,draw=blue!14!white,line width=.12pt] (-0.607814,1.641929)--(-0.651378,1.753160)--(-0.858714,1.716067)--(-0.801283,1.607317)--cycle;
\path[fill=blue!36!white,draw=blue!36!white,line width=.12pt] (-2.615501,0.209109)--(-2.508586,0.094375)--(-2.627935,-0.031796)--(-2.739936,0.077561)--cycle;
\path[fill=blue!41!white,draw=blue!41!white,line width=.12pt] (2.739936,0.077561)--(2.627935,-0.031796)--(2.508586,0.094375)--(2.615501,0.209109)--cycle;
\path[fill=blue!12!white,draw=blue!12!white,line width=.12pt] (0.381508,1.535043)--(0.408491,1.666893)--(0.205234,1.681968)--(0.191677,1.549123)--cycle;
\path[fill=blue!13!white,draw=blue!13!white,line width=.12pt] (-0.191677,1.549123)--(-0.205234,1.681968)--(-0.408491,1.666893)--(-0.381508,1.535043)--cycle;
\path[fill=blue!33!white,draw=blue!33!white,line width=.12pt] (2.808591,0.885774)--(2.822148,0.754235)--(2.660703,0.888912)--(2.647922,1.019805)--cycle;
\path[fill=blue!25!white,draw=blue!25!white,line width=.12pt] (-2.647922,1.019805)--(-2.660703,0.888912)--(-2.822148,0.754235)--(-2.808591,0.885774)--cycle;
\path[fill=blue!18!white,draw=blue!18!white,line width=.12pt] (-2.011334,1.581726)--(-2.097056,1.538426)--(-2.292506,1.430603)--(-2.198794,1.478311)--cycle;
\path[fill=blue!22!white,draw=blue!22!white,line width=.12pt] (2.198794,1.478311)--(2.292506,1.430603)--(2.097056,1.538426)--(2.011334,1.581726)--cycle;
\path[fill=blue!41!white,draw=blue!41!white,line width=.12pt] (2.627935,-0.031796)--(2.500154,-0.122284)--(2.386608,-0.002248)--(2.508586,0.094375)--cycle;
\path[fill=blue!38!white,draw=blue!38!white,line width=.12pt] (-2.508586,0.094375)--(-2.386608,-0.002248)--(-2.500154,-0.122284)--(-2.627935,-0.031796)--cycle;
\path[fill=blue!34!white,draw=blue!34!white,line width=.12pt] (-2.703244,0.337547)--(-2.615501,0.209109)--(-2.739936,0.077561)--(-2.831853,0.201586)--cycle;
\path[fill=blue!40!white,draw=blue!40!white,line width=.12pt] (2.831853,0.201586)--(2.739936,0.077561)--(2.615501,0.209109)--(2.703244,0.337547)--cycle;
\path[fill=blue!31!white,draw=blue!31!white,line width=.12pt] (-0.963817,0.336316)--(-0.922745,0.425362)--(-1.053662,0.376814)--(-1.100561,0.285608)--cycle;
\path[fill=blue!26!white,draw=blue!26!white,line width=.12pt] (1.100561,0.285608)--(1.053662,0.376814)--(0.922745,0.425362)--(0.963817,0.336316)--cycle;
\path[fill=blue!39!white,draw=blue!39!white,line width=.12pt] (-2.386608,-0.002248)--(-2.254254,-0.077046)--(-2.361503,-0.190425)--(-2.500154,-0.122284)--cycle;
\path[fill=blue!41!white,draw=blue!41!white,line width=.12pt] (2.500154,-0.122284)--(2.361503,-0.190425)--(2.254254,-0.077046)--(2.386608,-0.002248)--cycle;
\path[fill=blue!19!white,draw=blue!19!white,line width=.12pt] (0.156827,0.796148)--(0.158334,0.938959)--(0.000000,0.942848)--(0.000000,0.800000)--cycle;
\path[fill=blue!19!white,draw=blue!19!white,line width=.12pt] (0.000000,0.800000)--(0.000000,0.942848)--(-0.158334,0.938959)--(-0.156827,0.796148)--cycle;
\path[fill=blue!20!white,draw=blue!20!white,line width=.12pt] (0.618176,0.611042)--(0.612293,0.739104)--(0.464455,0.765552)--(0.468918,0.637745)--cycle;
\path[fill=blue!23!white,draw=blue!23!white,line width=.12pt] (-0.468918,0.637745)--(-0.464455,0.765552)--(-0.612293,0.739104)--(-0.618176,0.611042)--cycle;
\path[fill=blue!13!white,draw=blue!13!white,line width=.12pt] (-0.381508,1.535043)--(-0.408491,1.666893)--(-0.607814,1.641929)--(-0.567664,1.511728)--cycle;
\path[fill=blue!12!white,draw=blue!12!white,line width=.12pt] (0.567664,1.511728)--(0.607814,1.641929)--(0.408491,1.666893)--(0.381508,1.535043)--cycle;
\path[fill=blue!37!white,draw=blue!37!white,line width=.12pt] (1.978968,-0.150618)--(1.846614,-0.146565)--(1.740975,-0.058441)--(1.865758,-0.056178)--cycle;
\path[fill=blue!39!white,draw=blue!39!white,line width=.12pt] (-1.865758,-0.056178)--(-1.740975,-0.058441)--(-1.846614,-0.146565)--(-1.978968,-0.150618)--cycle;
\path[fill=blue!18!white,draw=blue!18!white,line width=.12pt] (0.312145,0.784628)--(0.315143,0.927329)--(0.158334,0.938959)--(0.156827,0.796148)--cycle;
\path[fill=blue!20!white,draw=blue!20!white,line width=.12pt] (-0.156827,0.796148)--(-0.158334,0.938959)--(-0.315143,0.927329)--(-0.312145,0.784628)--cycle;
\path[fill=blue!32!white,draw=blue!32!white,line width=.12pt] (-2.768442,0.474751)--(-2.703244,0.337547)--(-2.831853,0.201586)--(-2.900154,0.335511)--cycle;
\path[fill=blue!39!white,draw=blue!39!white,line width=.12pt] (2.900154,0.335511)--(2.831853,0.201586)--(2.703244,0.337547)--(2.768442,0.474751)--cycle;
\path[fill=blue!27!white,draw=blue!27!white,line width=.12pt] (-0.782941,0.467259)--(-0.761481,0.577153)--(-0.897452,0.536404)--(-0.922745,0.425362)--cycle;
\path[fill=blue!23!white,draw=blue!23!white,line width=.12pt] (0.922745,0.425362)--(0.897452,0.536404)--(0.761481,0.577153)--(0.782941,0.467259)--cycle;
\path[fill=blue!13!white,draw=blue!13!white,line width=.12pt] (0.179794,1.402639)--(0.191677,1.549123)--(0.000000,1.553831)--(0.000000,1.407055)--cycle;
\path[fill=blue!13!white,draw=blue!13!white,line width=.12pt] (0.000000,1.407055)--(0.000000,1.553831)--(-0.191677,1.549123)--(-0.179794,1.402639)--cycle;
\path[fill=blue!14!white,draw=blue!14!white,line width=.12pt] (-1.057780,1.668992)--(-1.131352,1.751126)--(-1.333369,1.690584)--(-1.246659,1.612387)--cycle;
\path[fill=blue!14!white,draw=blue!14!white,line width=.12pt] (1.246659,1.612387)--(1.333369,1.690584)--(1.131352,1.751126)--(1.057780,1.668992)--cycle;
\path[fill=blue!15!white,draw=blue!15!white,line width=.12pt] (1.522544,1.620433)--(1.621555,1.667443)--(1.420078,1.742156)--(1.333369,1.690584)--cycle;
\path[fill=blue!15!white,draw=blue!15!white,line width=.12pt] (-1.333369,1.690584)--(-1.420078,1.742156)--(-1.621555,1.667443)--(-1.522544,1.620433)--cycle;
\path[fill=blue!14!white,draw=blue!14!white,line width=.12pt] (-0.801283,1.607317)--(-0.858714,1.716067)--(-1.057780,1.668992)--(-0.987036,1.563389)--cycle;
\path[fill=blue!13!white,draw=blue!13!white,line width=.12pt] (0.987036,1.563389)--(1.057780,1.668992)--(0.858714,1.716067)--(0.801283,1.607317)--cycle;
\path[fill=blue!13!white,draw=blue!13!white,line width=.12pt] (0.357857,1.389432)--(0.381508,1.535043)--(0.191677,1.549123)--(0.179794,1.402639)--cycle;
\path[fill=blue!13!white,draw=blue!13!white,line width=.12pt] (-0.179794,1.402639)--(-0.191677,1.549123)--(-0.381508,1.535043)--(-0.357857,1.389432)--cycle;
\path[fill=blue!40!white,draw=blue!40!white,line width=.12pt] (2.361503,-0.190425)--(2.217311,-0.233600)--(2.116611,-0.127144)--(2.254254,-0.077046)--cycle;
\path[fill=blue!39!white,draw=blue!39!white,line width=.12pt] (-2.254254,-0.077046)--(-2.116611,-0.127144)--(-2.217311,-0.233600)--(-2.361503,-0.190425)--cycle;
\path[fill=blue!20!white,draw=blue!20!white,line width=.12pt] (-0.312145,0.784628)--(-0.315143,0.927329)--(-0.468918,0.908070)--(-0.464455,0.765552)--cycle;
\path[fill=blue!18!white,draw=blue!18!white,line width=.12pt] (0.464455,0.765552)--(0.468918,0.908070)--(0.315143,0.927329)--(0.312145,0.784628)--cycle;
\path[fill=blue!16!white,draw=blue!16!white,line width=.12pt] (-1.621555,1.667443)--(-1.716761,1.686022)--(-1.913535,1.596850)--(-1.807416,1.583216)--cycle;
\path[fill=blue!17!white,draw=blue!17!white,line width=.12pt] (1.807416,1.583216)--(1.913535,1.596850)--(1.716761,1.686022)--(1.621555,1.667443)--cycle;
\path[fill=blue!25!white,draw=blue!25!white,line width=.12pt] (2.465877,1.313721)--(2.548600,1.236371)--(2.369413,1.357174)--(2.292506,1.430603)--cycle;
\path[fill=blue!19!white,draw=blue!19!white,line width=.12pt] (-2.292506,1.430603)--(-2.369413,1.357174)--(-2.548600,1.236371)--(-2.465877,1.313721)--cycle;
\path[fill=blue!17!white,draw=blue!17!white,line width=.12pt] (0.158334,0.938959)--(0.162796,1.091580)--(0.000000,1.095579)--(0.000000,0.942848)--cycle;
\path[fill=blue!17!white,draw=blue!17!white,line width=.12pt] (0.000000,0.942848)--(0.000000,1.095579)--(-0.162796,1.091580)--(-0.158334,0.938959)--cycle;
\path[fill=blue!12!white,draw=blue!12!white,line width=.12pt] (0.748354,1.479402)--(0.801283,1.607317)--(0.607814,1.641929)--(0.567664,1.511728)--cycle;
\path[fill=blue!14!white,draw=blue!14!white,line width=.12pt] (-0.567664,1.511728)--(-0.607814,1.641929)--(-0.801283,1.607317)--(-0.748354,1.479402)--cycle;
\path[fill=blue!24!white,draw=blue!24!white,line width=.12pt] (-0.618176,0.611042)--(-0.612293,0.739104)--(-0.754235,0.705537)--(-0.761481,0.577153)--cycle;
\path[fill=blue!20!white,draw=blue!20!white,line width=.12pt] (0.761481,0.577153)--(0.754235,0.705537)--(0.612293,0.739104)--(0.618176,0.611042)--cycle;
\path[fill=blue!30!white,draw=blue!30!white,line width=.12pt] (2.768442,1.005012)--(2.808591,0.885774)--(2.647922,1.019805)--(2.610069,1.137127)--cycle;
\path[fill=blue!23!white,draw=blue!23!white,line width=.12pt] (-2.610069,1.137127)--(-2.647922,1.019805)--(-2.808591,0.885774)--(-2.768442,1.005012)--cycle;
\path[fill=blue!14!white,draw=blue!14!white,line width=.12pt] (0.000000,1.252323)--(0.000000,1.407055)--(-0.179794,1.402639)--(-0.170043,1.248146)--cycle;
\path[fill=blue!14!white,draw=blue!14!white,line width=.12pt] (0.170043,1.248146)--(0.179794,1.402639)--(0.000000,1.407055)--(0.000000,1.252323)--cycle;
\path[fill=blue!29!white,draw=blue!29!white,line width=.12pt] (1.417944,0.091940)--(1.341037,0.165370)--(1.226706,0.228443)--(1.297056,0.158630)--cycle;
\path[fill=blue!34!white,draw=blue!34!white,line width=.12pt] (-1.297056,0.158630)--(-1.226706,0.228443)--(-1.341037,0.165370)--(-1.417944,0.091940)--cycle;
\path[fill=blue!16!white,draw=blue!16!white,line width=.12pt] (0.315143,0.927329)--(0.324025,1.079622)--(0.162796,1.091580)--(0.158334,0.938959)--cycle;
\path[fill=blue!18!white,draw=blue!18!white,line width=.12pt] (-0.158334,0.938959)--(-0.162796,1.091580)--(-0.324025,1.079622)--(-0.315143,0.927329)--cycle;
\path[fill=blue!30!white,draw=blue!30!white,line width=.12pt] (-2.808591,0.615449)--(-2.768442,0.474751)--(-2.900154,0.335511)--(-2.942213,0.474190)--cycle;
\path[fill=blue!37!white,draw=blue!37!white,line width=.12pt] (2.942213,0.474190)--(2.900154,0.335511)--(2.768442,0.474751)--(2.808591,0.615449)--cycle;
\path[fill=blue!32!white,draw=blue!32!white,line width=.12pt] (1.625975,-0.032838)--(1.525177,0.019647)--(1.417944,0.091940)--(1.511656,0.044233)--cycle;
\path[fill=blue!36!white,draw=blue!36!white,line width=.12pt] (-1.511656,0.044233)--(-1.417944,0.091940)--(-1.525177,0.019647)--(-1.625975,-0.032838)--cycle;
\path[fill=blue!14!white,draw=blue!14!white,line width=.12pt] (-0.357857,1.389432)--(-0.381508,1.535043)--(-0.567664,1.511728)--(-0.532473,1.367563)--cycle;
\path[fill=blue!12!white,draw=blue!12!white,line width=.12pt] (0.532473,1.367563)--(0.567664,1.511728)--(0.381508,1.535043)--(0.357857,1.389432)--cycle;
\path[fill=blue!15!white,draw=blue!15!white,line width=.12pt] (0.162796,1.091580)--(0.170043,1.248146)--(0.000000,1.252323)--(0.000000,1.095579)--cycle;
\path[fill=blue!15!white,draw=blue!15!white,line width=.12pt] (0.000000,1.095579)--(0.000000,1.252323)--(-0.170043,1.248146)--(-0.162796,1.091580)--cycle;
\path[fill=blue!26!white,draw=blue!26!white,line width=.12pt] (1.226706,0.228443)--(1.174431,0.322085)--(1.053662,0.376814)--(1.100561,0.285608)--cycle;
\path[fill=blue!31!white,draw=blue!31!white,line width=.12pt] (-1.100561,0.285608)--(-1.053662,0.376814)--(-1.174431,0.322085)--(-1.226706,0.228443)--cycle;
\path[fill=blue!13!white,draw=blue!13!white,line width=.12pt] (0.338447,1.235655)--(0.357857,1.389432)--(0.179794,1.402639)--(0.170043,1.248146)--cycle;
\path[fill=blue!14!white,draw=blue!14!white,line width=.12pt] (-0.170043,1.248146)--(-0.179794,1.402639)--(-0.357857,1.389432)--(-0.338447,1.235655)--cycle;
\path[fill=blue!18!white,draw=blue!18!white,line width=.12pt] (0.612293,0.739104)--(0.618176,0.881367)--(0.468918,0.908070)--(0.464455,0.765552)--cycle;
\path[fill=blue!21!white,draw=blue!21!white,line width=.12pt] (-0.464455,0.765552)--(-0.468918,0.908070)--(-0.618176,0.881367)--(-0.612293,0.739104)--cycle;
\path[fill=blue!16!white,draw=blue!16!white,line width=.12pt] (-0.162796,1.091580)--(-0.170043,1.248146)--(-0.338447,1.235655)--(-0.324025,1.079622)--cycle;
\path[fill=blue!15!white,draw=blue!15!white,line width=.12pt] (0.324025,1.079622)--(0.338447,1.235655)--(0.170043,1.248146)--(0.162796,1.091580)--cycle;
\path[fill=blue!16!white,draw=blue!16!white,line width=.12pt] (-1.913535,1.596850)--(-2.011334,1.581726)--(-2.198794,1.478311)--(-2.091880,1.498463)--cycle;
\path[fill=blue!20!white,draw=blue!20!white,line width=.12pt] (2.091880,1.498463)--(2.198794,1.478311)--(2.011334,1.581726)--(1.913535,1.596850)--cycle;
\path[fill=blue!16!white,draw=blue!16!white,line width=.12pt] (0.468918,0.908070)--(0.482133,1.059820)--(0.324025,1.079622)--(0.315143,0.927329)--cycle;
\path[fill=blue!18!white,draw=blue!18!white,line width=.12pt] (-0.315143,0.927329)--(-0.324025,1.079622)--(-0.482133,1.059820)--(-0.468918,0.908070)--cycle;
\path[fill=blue!40!white,draw=blue!40!white,line width=.12pt] (-2.116611,-0.127144)--(-1.978968,-0.150618)--(-2.073119,-0.250151)--(-2.217311,-0.233600)--cycle;
\path[fill=blue!39!white,draw=blue!39!white,line width=.12pt] (2.217311,-0.233600)--(2.073119,-0.250151)--(1.978968,-0.150618)--(2.116611,-0.127144)--cycle;
\path[fill=blue!23!white,draw=blue!23!white,line width=.12pt] (1.053662,0.376814)--(1.024781,0.489187)--(0.897452,0.536404)--(0.922745,0.425362)--cycle;
\path[fill=blue!28!white,draw=blue!28!white,line width=.12pt] (-0.922745,0.425362)--(-0.897452,0.536404)--(-1.024781,0.489187)--(-1.053662,0.376814)--cycle;
\path[fill=blue!38!white,draw=blue!38!white,line width=.12pt] (-1.740975,-0.058441)--(-1.625975,-0.032838)--(-1.724636,-0.115141)--(-1.846614,-0.146565)--cycle;
\path[fill=blue!35!white,draw=blue!35!white,line width=.12pt] (1.846614,-0.146565)--(1.724636,-0.115141)--(1.625975,-0.032838)--(1.740975,-0.058441)--cycle;
\path[fill=blue!15!white,draw=blue!15!white,line width=.12pt] (-0.338447,1.235655)--(-0.357857,1.389432)--(-0.532473,1.367563)--(-0.503593,1.214972)--cycle;
\path[fill=blue!13!white,draw=blue!13!white,line width=.12pt] (0.503593,1.214972)--(0.532473,1.367563)--(0.357857,1.389432)--(0.338447,1.235655)--cycle;
\path[fill=blue!14!white,draw=blue!14!white,line width=.12pt] (0.482133,1.059820)--(0.503593,1.214972)--(0.338447,1.235655)--(0.324025,1.079622)--cycle;
\path[fill=blue!16!white,draw=blue!16!white,line width=.12pt] (-0.324025,1.079622)--(-0.338447,1.235655)--(-0.503593,1.214972)--(-0.482133,1.059820)--cycle;
\path[fill=blue!12!white,draw=blue!12!white,line width=.12pt] (0.701962,1.337241)--(0.748354,1.479402)--(0.567664,1.511728)--(0.532473,1.367563)--cycle;
\path[fill=blue!14!white,draw=blue!14!white,line width=.12pt] (-0.532473,1.367563)--(-0.567664,1.511728)--(-0.748354,1.479402)--(-0.701962,1.337241)--cycle;
\path[fill=blue!12!white,draw=blue!12!white,line width=.12pt] (0.921837,1.438377)--(0.987036,1.563389)--(0.801283,1.607317)--(0.748354,1.479402)--cycle;
\path[fill=blue!14!white,draw=blue!14!white,line width=.12pt] (-0.748354,1.479402)--(-0.801283,1.607317)--(-0.987036,1.563389)--(-0.921837,1.438377)--cycle;
\path[fill=blue!20!white,draw=blue!20!white,line width=.12pt] (0.897452,0.536404)--(0.888912,0.665176)--(0.754235,0.705537)--(0.761481,0.577153)--cycle;
\path[fill=blue!25!white,draw=blue!25!white,line width=.12pt] (-0.761481,0.577153)--(-0.754235,0.705537)--(-0.888912,0.665176)--(-0.897452,0.536404)--cycle;
\path[fill=blue!13!white,draw=blue!13!white,line width=.12pt] (1.163283,1.510570)--(1.246659,1.612387)--(1.057780,1.668992)--(0.987036,1.563389)--cycle;
\path[fill=blue!14!white,draw=blue!14!white,line width=.12pt] (-0.987036,1.563389)--(-1.057780,1.668992)--(-1.246659,1.612387)--(-1.163283,1.510570)--cycle;
\path[fill=blue!27!white,draw=blue!27!white,line width=.12pt] (-2.822148,0.754235)--(-2.808591,0.615449)--(-2.942213,0.474190)--(-2.956415,0.612293)--cycle;
\path[fill=blue!35!white,draw=blue!35!white,line width=.12pt] (2.956415,0.612293)--(2.942213,0.474190)--(2.808591,0.615449)--(2.822148,0.754235)--cycle;
\path[fill=blue!19!white,draw=blue!19!white,line width=.12pt] (-0.468918,0.908070)--(-0.482133,1.059820)--(-0.635598,1.032365)--(-0.618176,0.881367)--cycle;
\path[fill=blue!16!white,draw=blue!16!white,line width=.12pt] (0.618176,0.881367)--(0.635598,1.032365)--(0.482133,1.059820)--(0.468918,0.908070)--cycle;
\path[fill=blue!22!white,draw=blue!22!white,line width=.12pt] (-0.612293,0.739104)--(-0.618176,0.881367)--(-0.761481,0.847478)--(-0.754235,0.705537)--cycle;
\path[fill=blue!18!white,draw=blue!18!white,line width=.12pt] (0.754235,0.705537)--(0.761481,0.847478)--(0.618176,0.881367)--(0.612293,0.739104)--cycle;
\path[fill=blue!14!white,draw=blue!14!white,line width=.12pt] (1.423533,1.546798)--(1.522544,1.620433)--(1.333369,1.690584)--(1.246659,1.612387)--cycle;
\path[fill=blue!14!white,draw=blue!14!white,line width=.12pt] (-1.246659,1.612387)--(-1.333369,1.690584)--(-1.522544,1.620433)--(-1.423533,1.546798)--cycle;
\path[fill=blue!15!white,draw=blue!15!white,line width=.12pt] (-0.503593,1.214972)--(-0.532473,1.367563)--(-0.701962,1.337241)--(-0.663889,1.186295)--cycle;
\path[fill=blue!13!white,draw=blue!13!white,line width=.12pt] (0.663889,1.186295)--(0.701962,1.337241)--(0.532473,1.367563)--(0.503593,1.214972)--cycle;
\path[fill=blue!14!white,draw=blue!14!white,line width=.12pt] (0.635598,1.032365)--(0.663889,1.186295)--(0.503593,1.214972)--(0.482133,1.059820)--cycle;
\path[fill=blue!17!white,draw=blue!17!white,line width=.12pt] (-0.482133,1.059820)--(-0.503593,1.214972)--(-0.663889,1.186295)--(-0.635598,1.032365)--cycle;
\path[fill=blue!15!white,draw=blue!15!white,line width=.12pt] (-1.522544,1.620433)--(-1.621555,1.667443)--(-1.807416,1.583216)--(-1.697056,1.541348)--cycle;
\path[fill=blue!15!white,draw=blue!15!white,line width=.12pt] (1.697056,1.541348)--(1.807416,1.583216)--(1.621555,1.667443)--(1.522544,1.620433)--cycle;
\path[fill=blue!21!white,draw=blue!21!white,line width=.12pt] (-2.548600,1.236371)--(-2.610069,1.137127)--(-2.768442,1.005012)--(-2.703244,1.107367)--cycle;
\path[fill=blue!28!white,draw=blue!28!white,line width=.12pt] (2.703244,1.107367)--(2.768442,1.005012)--(2.610069,1.137127)--(2.548600,1.236371)--cycle;
\path[fill=blue!18!white,draw=blue!18!white,line width=.12pt] (-2.198794,1.478311)--(-2.292506,1.430603)--(-2.465877,1.313721)--(-2.365079,1.366207)--cycle;
\path[fill=blue!22!white,draw=blue!22!white,line width=.12pt] (2.365079,1.366207)--(2.465877,1.313721)--(2.292506,1.430603)--(2.198794,1.478311)--cycle;
\path[fill=blue!39!white,draw=blue!39!white,line width=.12pt] (-1.978968,-0.150618)--(-1.846614,-0.146565)--(-1.934468,-0.239441)--(-2.073119,-0.250151)--cycle;
\path[fill=blue!37!white,draw=blue!37!white,line width=.12pt] (2.073119,-0.250151)--(1.934468,-0.239441)--(1.846614,-0.146565)--(1.978968,-0.150618)--cycle;
\path[fill=blue!12!white,draw=blue!12!white,line width=.12pt] (0.864690,1.298758)--(0.921837,1.438377)--(0.748354,1.479402)--(0.701962,1.337241)--cycle;
\path[fill=blue!15!white,draw=blue!15!white,line width=.12pt] (-0.701962,1.337241)--(-0.748354,1.479402)--(-0.921837,1.438377)--(-0.864690,1.298758)--cycle;
\path[fill=blue!32!white,draw=blue!32!white,line width=.12pt] (-1.226706,0.228443)--(-1.174431,0.322085)--(-1.283890,0.261700)--(-1.341037,0.165370)--cycle;
\path[fill=blue!26!white,draw=blue!26!white,line width=.12pt] (1.341037,0.165370)--(1.283890,0.261700)--(1.174431,0.322085)--(1.226706,0.228443)--cycle;
\path[fill=blue!29!white,draw=blue!29!white,line width=.12pt] (1.525177,0.019647)--(1.442454,0.096998)--(1.341037,0.165370)--(1.417944,0.091940)--cycle;
\path[fill=blue!35!white,draw=blue!35!white,line width=.12pt] (-1.417944,0.091940)--(-1.341037,0.165370)--(-1.442454,0.096998)--(-1.525177,0.019647)--cycle;
\path[fill=blue!23!white,draw=blue!23!white,line width=.12pt] (1.174431,0.322085)--(1.142240,0.435958)--(1.024781,0.489187)--(1.053662,0.376814)--cycle;
\path[fill=blue!29!white,draw=blue!29!white,line width=.12pt] (-1.053662,0.376814)--(-1.024781,0.489187)--(-1.142240,0.435958)--(-1.174431,0.322085)--cycle;
\path[fill=blue!16!white,draw=blue!16!white,line width=.12pt] (0.761481,0.847478)--(0.782941,0.997521)--(0.635598,1.032365)--(0.618176,0.881367)--cycle;
\path[fill=blue!19!white,draw=blue!19!white,line width=.12pt] (-0.618176,0.881367)--(-0.635598,1.032365)--(-0.782941,0.997521)--(-0.761481,0.847478)--cycle;
\path[fill=blue!37!white,draw=blue!37!white,line width=.12pt] (-2.627935,-0.031796)--(-2.500154,-0.122284)--(-2.589621,-0.247306)--(-2.721975,-0.163208)--cycle;
\path[fill=blue!41!white,draw=blue!41!white,line width=.12pt] (2.721975,-0.163208)--(2.589621,-0.247306)--(2.500154,-0.122284)--(2.627935,-0.031796)--cycle;
\path[fill=blue!20!white,draw=blue!20!white,line width=.12pt] (1.024781,0.489187)--(1.015029,0.618408)--(0.888912,0.665176)--(0.897452,0.536404)--cycle;
\path[fill=blue!26!white,draw=blue!26!white,line width=.12pt] (-0.897452,0.536404)--(-0.888912,0.665176)--(-1.015029,0.618408)--(-1.024781,0.489187)--cycle;
\path[fill=blue!41!white,draw=blue!41!white,line width=.12pt] (2.837984,-0.059451)--(2.721975,-0.163208)--(2.627935,-0.031796)--(2.739936,0.077561)--cycle;
\path[fill=blue!36!white,draw=blue!36!white,line width=.12pt] (-2.739936,0.077561)--(-2.627935,-0.031796)--(-2.721975,-0.163208)--(-2.837984,-0.059451)--cycle;
\path[fill=blue!22!white,draw=blue!22!white,line width=.12pt] (-0.754235,0.705537)--(-0.761481,0.847478)--(-0.897452,0.806729)--(-0.888912,0.665176)--cycle;
\path[fill=blue!18!white,draw=blue!18!white,line width=.12pt] (0.888912,0.665176)--(0.897452,0.806729)--(0.761481,0.847478)--(0.754235,0.705537)--cycle;
\path[fill=blue!12!white,draw=blue!12!white,line width=.12pt] (1.086442,1.389047)--(1.163283,1.510570)--(0.987036,1.563389)--(0.921837,1.438377)--cycle;
\path[fill=blue!14!white,draw=blue!14!white,line width=.12pt] (-0.921837,1.438377)--(-0.987036,1.563389)--(-1.163283,1.510570)--(-1.086442,1.389047)--cycle;
\path[fill=blue!32!white,draw=blue!32!white,line width=.12pt] (1.724636,-0.115141)--(1.617721,-0.057553)--(1.525177,0.019647)--(1.625975,-0.032838)--cycle;
\path[fill=blue!37!white,draw=blue!37!white,line width=.12pt] (-1.625975,-0.032838)--(-1.525177,0.019647)--(-1.617721,-0.057553)--(-1.724636,-0.115141)--cycle;
\path[fill=blue!38!white,draw=blue!38!white,line width=.12pt] (-2.500154,-0.122284)--(-2.361503,-0.190425)--(-2.446009,-0.308514)--(-2.589621,-0.247306)--cycle;
\path[fill=blue!41!white,draw=blue!41!white,line width=.12pt] (2.589621,-0.247306)--(2.446009,-0.308514)--(2.361503,-0.190425)--(2.500154,-0.122284)--cycle;
\path[fill=blue!17!white,draw=blue!17!white,line width=.12pt] (1.975871,1.490286)--(2.091880,1.498463)--(1.913535,1.596850)--(1.807416,1.583216)--cycle;
\path[fill=blue!15!white,draw=blue!15!white,line width=.12pt] (-1.807416,1.583216)--(-1.913535,1.596850)--(-2.091880,1.498463)--(-1.975871,1.490286)--cycle;
\path[fill=blue!13!white,draw=blue!13!white,line width=.12pt] (0.817791,1.149900)--(0.864690,1.298758)--(0.701962,1.337241)--(0.663889,1.186295)--cycle;
\path[fill=blue!16!white,draw=blue!16!white,line width=.12pt] (-0.663889,1.186295)--(-0.701962,1.337241)--(-0.864690,1.298758)--(-0.817791,1.149900)--cycle;
\path[fill=blue!25!white,draw=blue!25!white,line width=.12pt] (-2.808591,0.885774)--(-2.822148,0.754235)--(-2.956415,0.612293)--(-2.942213,0.744515)--cycle;
\path[fill=blue!33!white,draw=blue!33!white,line width=.12pt] (2.942213,0.744515)--(2.956415,0.612293)--(2.822148,0.754235)--(2.808591,0.885774)--cycle;
\path[fill=blue!17!white,draw=blue!17!white,line width=.12pt] (-0.635598,1.032365)--(-0.663889,1.186295)--(-0.817791,1.149900)--(-0.782941,0.997521)--cycle;
\path[fill=blue!14!white,draw=blue!14!white,line width=.12pt] (0.782941,0.997521)--(0.817791,1.149900)--(0.663889,1.186295)--(0.635598,1.032365)--cycle;
\path[fill=blue!33!white,draw=blue!33!white,line width=.12pt] (-2.831853,0.201586)--(-2.739936,0.077561)--(-2.837984,-0.059451)--(-2.933190,0.059976)--cycle;
\path[fill=blue!40!white,draw=blue!40!white,line width=.12pt] (2.933190,0.059976)--(2.837984,-0.059451)--(2.739936,0.077561)--(2.831853,0.201586)--cycle;
\path[fill=blue!14!white,draw=blue!14!white,line width=.12pt] (-1.163283,1.510570)--(-1.246659,1.612387)--(-1.423533,1.546798)--(-1.328326,1.449368)--cycle;
\path[fill=blue!13!white,draw=blue!13!white,line width=.12pt] (1.328326,1.449368)--(1.423533,1.546798)--(1.246659,1.612387)--(1.163283,1.510570)--cycle;
\path[fill=blue!39!white,draw=blue!39!white,line width=.12pt] (-2.361503,-0.190425)--(-2.217311,-0.233600)--(-2.296657,-0.344479)--(-2.446009,-0.308514)--cycle;
\path[fill=blue!40!white,draw=blue!40!white,line width=.12pt] (2.446009,-0.308514)--(2.296657,-0.344479)--(2.217311,-0.233600)--(2.361503,-0.190425)--cycle;
\path[fill=blue!39!white,draw=blue!39!white,line width=.12pt] (-1.846614,-0.146565)--(-1.724636,-0.115141)--(-1.806687,-0.201882)--(-1.934468,-0.239441)--cycle;
\path[fill=blue!35!white,draw=blue!35!white,line width=.12pt] (1.934468,-0.239441)--(1.806687,-0.201882)--(1.724636,-0.115141)--(1.846614,-0.146565)--cycle;
\path[fill=blue!20!white,draw=blue!20!white,line width=.12pt] (-0.761481,0.847478)--(-0.782941,0.997521)--(-0.922745,0.955623)--(-0.897452,0.806729)--cycle;
\path[fill=blue!16!white,draw=blue!16!white,line width=.12pt] (0.897452,0.806729)--(0.922745,0.955623)--(0.782941,0.997521)--(0.761481,0.847478)--cycle;
\path[fill=blue!12!white,draw=blue!12!white,line width=.12pt] (1.019091,1.252486)--(1.086442,1.389047)--(0.921837,1.438377)--(0.864690,1.298758)--cycle;
\path[fill=blue!15!white,draw=blue!15!white,line width=.12pt] (-0.864690,1.298758)--(-0.921837,1.438377)--(-1.086442,1.389047)--(-1.019091,1.252486)--cycle;
\path[fill=blue!14!white,draw=blue!14!white,line width=.12pt] (1.586697,1.472856)--(1.697056,1.541348)--(1.522544,1.620433)--(1.423533,1.546798)--cycle;
\path[fill=blue!14!white,draw=blue!14!white,line width=.12pt] (-1.423533,1.546798)--(-1.522544,1.620433)--(-1.697056,1.541348)--(-1.586697,1.472856)--cycle;
\path[fill=blue!31!white,draw=blue!31!white,line width=.12pt] (-2.900154,0.335511)--(-2.831853,0.201586)--(-2.933190,0.059976)--(-3.003935,0.190486)--cycle;
\path[fill=blue!39!white,draw=blue!39!white,line width=.12pt] (3.003935,0.190486)--(2.933190,0.059976)--(2.831853,0.201586)--(2.900154,0.335511)--cycle;
\path[fill=blue!25!white,draw=blue!25!white,line width=.12pt] (2.615501,1.188905)--(2.703244,1.107367)--(2.548600,1.236371)--(2.465877,1.313721)--cycle;
\path[fill=blue!19!white,draw=blue!19!white,line width=.12pt] (-2.465877,1.313721)--(-2.548600,1.236371)--(-2.703244,1.107367)--(-2.615501,1.188905)--cycle;
\path[fill=blue!19!white,draw=blue!19!white,line width=.12pt] (2.250079,1.391810)--(2.365079,1.366207)--(2.198794,1.478311)--(2.091880,1.498463)--cycle;
\path[fill=blue!16!white,draw=blue!16!white,line width=.12pt] (-2.091880,1.498463)--(-2.198794,1.478311)--(-2.365079,1.366207)--(-2.250079,1.391810)--cycle;
\path[fill=blue!23!white,draw=blue!23!white,line width=.12pt] (-0.888912,0.665176)--(-0.897452,0.806729)--(-1.024781,0.759512)--(-1.015029,0.618408)--cycle;
\path[fill=blue!18!white,draw=blue!18!white,line width=.12pt] (1.015029,0.618408)--(1.024781,0.759512)--(0.897452,0.806729)--(0.888912,0.665176)--cycle;
\path[fill=blue!40!white,draw=blue!40!white,line width=.12pt] (-2.217311,-0.233600)--(-2.073119,-0.250151)--(-2.147305,-0.353819)--(-2.296657,-0.344479)--cycle;
\path[fill=blue!39!white,draw=blue!39!white,line width=.12pt] (2.296657,-0.344479)--(2.147305,-0.353819)--(2.073119,-0.250151)--(2.217311,-0.233600)--cycle;
\path[fill=blue!14!white,draw=blue!14!white,line width=.12pt] (0.922745,0.955623)--(0.963817,1.106137)--(0.817791,1.149900)--(0.782941,0.997521)--cycle;
\path[fill=blue!18!white,draw=blue!18!white,line width=.12pt] (-0.782941,0.997521)--(-0.817791,1.149900)--(-0.963817,1.106137)--(-0.922745,0.955623)--cycle;
\path[fill=blue!13!white,draw=blue!13!white,line width=.12pt] (0.963817,1.106137)--(1.019091,1.252486)--(0.864690,1.298758)--(0.817791,1.149900)--cycle;
\path[fill=blue!16!white,draw=blue!16!white,line width=.12pt] (-0.817791,1.149900)--(-0.864690,1.298758)--(-1.019091,1.252486)--(-0.963817,1.106137)--cycle;
\path[fill=blue!26!white,draw=blue!26!white,line width=.12pt] (-1.024781,0.489187)--(-1.015029,0.618408)--(-1.131371,0.565685)--(-1.142240,0.435958)--cycle;
\path[fill=blue!20!white,draw=blue!20!white,line width=.12pt] (1.142240,0.435958)--(1.131371,0.565685)--(1.015029,0.618408)--(1.024781,0.489187)--cycle;
\path[fill=blue!30!white,draw=blue!30!white,line width=.12pt] (-1.174431,0.322085)--(-1.142240,0.435958)--(-1.248699,0.377228)--(-1.283890,0.261700)--cycle;
\path[fill=blue!23!white,draw=blue!23!white,line width=.12pt] (1.283890,0.261700)--(1.248699,0.377228)--(1.142240,0.435958)--(1.174431,0.322085)--cycle;
\path[fill=blue!26!white,draw=blue!26!white,line width=.12pt] (1.442454,0.096998)--(1.380985,0.196241)--(1.283890,0.261700)--(1.341037,0.165370)--cycle;
\path[fill=blue!33!white,draw=blue!33!white,line width=.12pt] (-1.341037,0.165370)--(-1.283890,0.261700)--(-1.380985,0.196241)--(-1.442454,0.096998)--cycle;
\path[fill=blue!15!white,draw=blue!15!white,line width=.12pt] (-1.086442,1.389047)--(-1.163283,1.510570)--(-1.328326,1.449368)--(-1.240584,1.331887)--cycle;
\path[fill=blue!12!white,draw=blue!12!white,line width=.12pt] (1.240584,1.331887)--(1.328326,1.449368)--(1.163283,1.510570)--(1.086442,1.389047)--cycle;
\path[fill=blue!22!white,draw=blue!22!white,line width=.12pt] (-2.768442,1.005012)--(-2.808591,0.885774)--(-2.942213,0.744515)--(-2.900154,0.865772)--cycle;
\path[fill=blue!30!white,draw=blue!30!white,line width=.12pt] (2.900154,0.865772)--(2.942213,0.744515)--(2.808591,0.885774)--(2.768442,1.005012)--cycle;
\path[fill=blue!29!white,draw=blue!29!white,line width=.12pt] (1.617721,-0.057553)--(1.529978,0.023985)--(1.442454,0.096998)--(1.525177,0.019647)--cycle;
\path[fill=blue!35!white,draw=blue!35!white,line width=.12pt] (-1.525177,0.019647)--(-1.442454,0.096998)--(-1.529978,0.023985)--(-1.617721,-0.057553)--cycle;
\path[fill=blue!15!white,draw=blue!15!white,line width=.12pt] (1.855225,1.454092)--(1.975871,1.490286)--(1.807416,1.583216)--(1.697056,1.541348)--cycle;
\path[fill=blue!15!white,draw=blue!15!white,line width=.12pt] (-1.697056,1.541348)--(-1.807416,1.583216)--(-1.975871,1.490286)--(-1.855225,1.454092)--cycle;
\path[fill=blue!29!white,draw=blue!29!white,line width=.12pt] (-2.942213,0.474190)--(-2.900154,0.335511)--(-3.003935,0.190486)--(-3.047499,0.327062)--cycle;
\path[fill=blue!37!white,draw=blue!37!white,line width=.12pt] (3.047499,0.327062)--(3.003935,0.190486)--(2.900154,0.335511)--(2.942213,0.474190)--cycle;
\path[fill=blue!21!white,draw=blue!21!white,line width=.12pt] (-0.897452,0.806729)--(-0.922745,0.955623)--(-1.053662,0.907076)--(-1.024781,0.759512)--cycle;
\path[fill=blue!16!white,draw=blue!16!white,line width=.12pt] (1.024781,0.759512)--(1.053662,0.907076)--(0.922745,0.955623)--(0.897452,0.806729)--cycle;
\path[fill=blue!40!white,draw=blue!40!white,line width=.12pt] (-2.073119,-0.250151)--(-1.934468,-0.239441)--(-2.003693,-0.336176)--(-2.147305,-0.353819)--cycle;
\path[fill=blue!37!white,draw=blue!37!white,line width=.12pt] (2.147305,-0.353819)--(2.003693,-0.336176)--(1.934468,-0.239441)--(2.073119,-0.250151)--cycle;
\path[fill=blue!14!white,draw=blue!14!white,line width=.12pt] (-1.328326,1.449368)--(-1.423533,1.546798)--(-1.586697,1.472856)--(-1.480578,1.380371)--cycle;
\path[fill=blue!13!white,draw=blue!13!white,line width=.12pt] (1.480578,1.380371)--(1.586697,1.472856)--(1.423533,1.546798)--(1.328326,1.449368)--cycle;
\path[fill=blue!12!white,draw=blue!12!white,line width=.12pt] (1.163677,1.198870)--(1.240584,1.331887)--(1.086442,1.389047)--(1.019091,1.252486)--cycle;
\path[fill=blue!16!white,draw=blue!16!white,line width=.12pt] (-1.019091,1.252486)--(-1.086442,1.389047)--(-1.240584,1.331887)--(-1.163677,1.198870)--cycle;
\path[fill=blue!37!white,draw=blue!37!white,line width=.12pt] (-1.724636,-0.115141)--(-1.617721,-0.057553)--(-1.694686,-0.138917)--(-1.806687,-0.201882)--cycle;
\path[fill=blue!32!white,draw=blue!32!white,line width=.12pt] (1.806687,-0.201882)--(1.694686,-0.138917)--(1.617721,-0.057553)--(1.724636,-0.115141)--cycle;
\path[fill=blue!19!white,draw=blue!19!white,line width=.12pt] (-0.922745,0.955623)--(-0.963817,1.106137)--(-1.100561,1.055429)--(-1.053662,0.907076)--cycle;
\path[fill=blue!14!white,draw=blue!14!white,line width=.12pt] (1.053662,0.907076)--(1.100561,1.055429)--(0.963817,1.106137)--(0.922745,0.955623)--cycle;
\path[fill=blue!24!white,draw=blue!24!white,line width=.12pt] (-1.015029,0.618408)--(-1.024781,0.759512)--(-1.142240,0.706283)--(-1.131371,0.565685)--cycle;
\path[fill=blue!18!white,draw=blue!18!white,line width=.12pt] (1.131371,0.565685)--(1.142240,0.706283)--(1.024781,0.759512)--(1.015029,0.618408)--cycle;
\path[fill=blue!17!white,draw=blue!17!white,line width=.12pt] (-0.963817,1.106137)--(-1.019091,1.252486)--(-1.163677,1.198870)--(-1.100561,1.055429)--cycle;
\path[fill=blue!13!white,draw=blue!13!white,line width=.12pt] (1.100561,1.055429)--(1.163677,1.198870)--(1.019091,1.252486)--(0.963817,1.106137)--cycle;
\path[fill=blue!22!white,draw=blue!22!white,line width=.12pt] (2.508586,1.246493)--(2.615501,1.188905)--(2.465877,1.313721)--(2.365079,1.366207)--cycle;
\path[fill=blue!17!white,draw=blue!17!white,line width=.12pt] (-2.365079,1.366207)--(-2.465877,1.313721)--(-2.615501,1.188905)--(-2.508586,1.246493)--cycle;
\path[fill=blue!15!white,draw=blue!15!white,line width=.12pt] (-1.975871,1.490286)--(-2.091880,1.498463)--(-2.250079,1.391810)--(-2.125296,1.389547)--cycle;
\path[fill=blue!17!white,draw=blue!17!white,line width=.12pt] (2.125296,1.389547)--(2.250079,1.391810)--(2.091880,1.498463)--(1.975871,1.490286)--cycle;
\path[fill=blue!27!white,draw=blue!27!white,line width=.12pt] (-1.142240,0.435958)--(-1.131371,0.565685)--(-1.236817,0.507515)--(-1.248699,0.377228)--cycle;
\path[fill=blue!20!white,draw=blue!20!white,line width=.12pt] (1.248699,0.377228)--(1.236817,0.507515)--(1.131371,0.565685)--(1.142240,0.435958)--cycle;
\path[fill=blue!26!white,draw=blue!26!white,line width=.12pt] (-2.956415,0.612293)--(-2.942213,0.474190)--(-3.047499,0.327062)--(-3.062209,0.464455)--cycle;
\path[fill=blue!35!white,draw=blue!35!white,line width=.12pt] (3.062209,0.464455)--(3.047499,0.327062)--(2.942213,0.474190)--(2.956415,0.612293)--cycle;
\path[fill=blue!23!white,draw=blue!23!white,line width=.12pt] (1.380985,0.196241)--(1.343133,0.313564)--(1.248699,0.377228)--(1.283890,0.261700)--cycle;
\path[fill=blue!30!white,draw=blue!30!white,line width=.12pt] (-1.283890,0.261700)--(-1.248699,0.377228)--(-1.343133,0.313564)--(-1.380985,0.196241)--cycle;
\path[fill=blue!20!white,draw=blue!20!white,line width=.12pt] (-2.703244,1.107367)--(-2.768442,1.005012)--(-2.900154,0.865772)--(-2.831853,0.971406)--cycle;
\path[fill=blue!27!white,draw=blue!27!white,line width=.12pt] (2.831853,0.971406)--(2.900154,0.865772)--(2.768442,1.005012)--(2.703244,1.107367)--cycle;
\path[fill=blue!14!white,draw=blue!14!white,line width=.12pt] (1.734580,1.391274)--(1.855225,1.454092)--(1.697056,1.541348)--(1.586697,1.472856)--cycle;
\path[fill=blue!14!white,draw=blue!14!white,line width=.12pt] (-1.586697,1.472856)--(-1.697056,1.541348)--(-1.855225,1.454092)--(-1.734580,1.391274)--cycle;
\path[fill=blue!15!white,draw=blue!15!white,line width=.12pt] (-1.240584,1.331887)--(-1.328326,1.449368)--(-1.480578,1.380371)--(-1.382778,1.267448)--cycle;
\path[fill=blue!12!white,draw=blue!12!white,line width=.12pt] (1.382778,1.267448)--(1.480578,1.380371)--(1.328326,1.449368)--(1.240584,1.331887)--cycle;
\path[fill=blue!33!white,draw=blue!33!white,line width=.12pt] (-1.442454,0.096998)--(-1.380985,0.196241)--(-1.464780,0.126340)--(-1.529978,0.023985)--cycle;
\path[fill=blue!26!white,draw=blue!26!white,line width=.12pt] (1.529978,0.023985)--(1.464780,0.126340)--(1.380985,0.196241)--(1.442454,0.096998)--cycle;
\path[fill=blue!39!white,draw=blue!39!white,line width=.12pt] (-1.934468,-0.239441)--(-1.806687,-0.201882)--(-1.871339,-0.292227)--(-2.003693,-0.336176)--cycle;
\path[fill=blue!35!white,draw=blue!35!white,line width=.12pt] (2.003693,-0.336176)--(1.871339,-0.292227)--(1.806687,-0.201882)--(1.934468,-0.239441)--cycle;
\path[fill=blue!37!white,draw=blue!37!white,line width=.12pt] (-2.721975,-0.163208)--(-2.589621,-0.247306)--(-2.654149,-0.376111)--(-2.789801,-0.298596)--cycle;
\path[fill=blue!41!white,draw=blue!41!white,line width=.12pt] (2.789801,-0.298596)--(2.654149,-0.376111)--(2.589621,-0.247306)--(2.721975,-0.163208)--cycle;
\path[fill=blue!16!white,draw=blue!16!white,line width=.12pt] (1.142240,0.706283)--(1.174431,0.852346)--(1.053662,0.907076)--(1.024781,0.759512)--cycle;
\path[fill=blue!21!white,draw=blue!21!white,line width=.12pt] (-1.024781,0.759512)--(-1.053662,0.907076)--(-1.174431,0.852346)--(-1.142240,0.706283)--cycle;
\path[fill=blue!38!white,draw=blue!38!white,line width=.12pt] (-2.589621,-0.247306)--(-2.446009,-0.308514)--(-2.506958,-0.430175)--(-2.654149,-0.376111)--cycle;
\path[fill=blue!41!white,draw=blue!41!white,line width=.12pt] (2.654149,-0.376111)--(2.506958,-0.430175)--(2.446009,-0.308514)--(2.589621,-0.247306)--cycle;
\path[fill=blue!35!white,draw=blue!35!white,line width=.12pt] (-2.837984,-0.059451)--(-2.721975,-0.163208)--(-2.789801,-0.298596)--(-2.908701,-0.200610)--cycle;
\path[fill=blue!41!white,draw=blue!41!white,line width=.12pt] (2.908701,-0.200610)--(2.789801,-0.298596)--(2.721975,-0.163208)--(2.837984,-0.059451)--cycle;
\path[fill=blue!39!white,draw=blue!39!white,line width=.12pt] (-2.446009,-0.308514)--(-2.296657,-0.344479)--(-2.353885,-0.458712)--(-2.506958,-0.430175)--cycle;
\path[fill=blue!40!white,draw=blue!40!white,line width=.12pt] (2.506958,-0.430175)--(2.353885,-0.458712)--(2.296657,-0.344479)--(2.446009,-0.308514)--cycle;
\path[fill=blue!36!white,draw=blue!36!white,line width=.12pt] (-1.617721,-0.057553)--(-1.529978,0.023985)--(-1.602769,-0.052966)--(-1.694686,-0.138917)--cycle;
\path[fill=blue!30!white,draw=blue!30!white,line width=.12pt] (1.694686,-0.138917)--(1.602769,-0.052966)--(1.529978,0.023985)--(1.617721,-0.057553)--cycle;
\path[fill=blue!16!white,draw=blue!16!white,line width=.12pt] (-1.163677,1.198870)--(-1.240584,1.331887)--(-1.382778,1.267448)--(-1.297056,1.138426)--cycle;
\path[fill=blue!12!white,draw=blue!12!white,line width=.12pt] (1.297056,1.138426)--(1.382778,1.267448)--(1.240584,1.331887)--(1.163677,1.198870)--cycle;
\path[fill=blue!14!white,draw=blue!14!white,line width=.12pt] (1.174431,0.852346)--(1.226706,0.998263)--(1.100561,1.055429)--(1.053662,0.907076)--cycle;
\path[fill=blue!19!white,draw=blue!19!white,line width=.12pt] (-1.053662,0.907076)--(-1.100561,1.055429)--(-1.226706,0.998263)--(-1.174431,0.852346)--cycle;
\path[fill=blue!18!white,draw=blue!18!white,line width=.12pt] (1.236817,0.507515)--(1.248699,0.647553)--(1.142240,0.706283)--(1.131371,0.565685)--cycle;
\path[fill=blue!25!white,draw=blue!25!white,line width=.12pt] (-1.131371,0.565685)--(-1.142240,0.706283)--(-1.248699,0.647553)--(-1.236817,0.507515)--cycle;
\path[fill=blue!24!white,draw=blue!24!white,line width=.12pt] (-2.942213,0.744515)--(-2.956415,0.612293)--(-3.062209,0.464455)--(-3.047499,0.597387)--cycle;
\path[fill=blue!32!white,draw=blue!32!white,line width=.12pt] (3.047499,0.597387)--(3.062209,0.464455)--(2.956415,0.612293)--(2.942213,0.744515)--cycle;
\path[fill=blue!17!white,draw=blue!17!white,line width=.12pt] (-1.100561,1.055429)--(-1.163677,1.198870)--(-1.297056,1.138426)--(-1.226706,0.998263)--cycle;
\path[fill=blue!13!white,draw=blue!13!white,line width=.12pt] (1.226706,0.998263)--(1.297056,1.138426)--(1.163677,1.198870)--(1.100561,1.055429)--cycle;
\path[fill=blue!33!white,draw=blue!33!white,line width=.12pt] (-2.933190,0.059976)--(-2.837984,-0.059451)--(-2.908701,-0.200610)--(-3.006279,-0.085917)--cycle;
\path[fill=blue!40!white,draw=blue!40!white,line width=.12pt] (3.006279,-0.085917)--(2.908701,-0.200610)--(2.837984,-0.059451)--(2.933190,0.059976)--cycle;
\path[fill=blue!40!white,draw=blue!40!white,line width=.12pt] (-2.296657,-0.344479)--(-2.147305,-0.353819)--(-2.200811,-0.460624)--(-2.353885,-0.458712)--cycle;
\path[fill=blue!39!white,draw=blue!39!white,line width=.12pt] (2.353885,-0.458712)--(2.200811,-0.460624)--(2.147305,-0.353819)--(2.296657,-0.344479)--cycle;
\path[fill=blue!19!white,draw=blue!19!white,line width=.12pt] (2.386608,1.277917)--(2.508586,1.246493)--(2.365079,1.366207)--(2.250079,1.391810)--cycle;
\path[fill=blue!16!white,draw=blue!16!white,line width=.12pt] (-2.250079,1.391810)--(-2.365079,1.366207)--(-2.508586,1.246493)--(-2.386608,1.277917)--cycle;
\path[fill=blue!14!white,draw=blue!14!white,line width=.12pt] (-1.855225,1.454092)--(-1.975871,1.490286)--(-2.125296,1.389547)--(-1.995527,1.359505)--cycle;
\path[fill=blue!15!white,draw=blue!15!white,line width=.12pt] (1.995527,1.359505)--(2.125296,1.389547)--(1.975871,1.490286)--(1.855225,1.454092)--cycle;
\path[fill=blue!15!white,draw=blue!15!white,line width=.12pt] (-1.480578,1.380371)--(-1.586697,1.472856)--(-1.734580,1.391274)--(-1.618570,1.304246)--cycle;
\path[fill=blue!13!white,draw=blue!13!white,line width=.12pt] (1.618570,1.304246)--(1.734580,1.391274)--(1.586697,1.472856)--(1.480578,1.380371)--cycle;
\path[fill=blue!28!white,draw=blue!28!white,line width=.12pt] (-1.248699,0.377228)--(-1.236817,0.507515)--(-1.330351,0.444456)--(-1.343133,0.313564)--cycle;
\path[fill=blue!21!white,draw=blue!21!white,line width=.12pt] (1.343133,0.313564)--(1.330351,0.444456)--(1.236817,0.507515)--(1.248699,0.377228)--cycle;
\path[fill=blue!18!white,draw=blue!18!white,line width=.12pt] (-2.615501,1.188905)--(-2.703244,1.107367)--(-2.831853,0.971406)--(-2.739936,1.057357)--cycle;
\path[fill=blue!24!white,draw=blue!24!white,line width=.12pt] (2.739936,1.057357)--(2.831853,0.971406)--(2.703244,1.107367)--(2.615501,1.188905)--cycle;
\path[fill=blue!38!white,draw=blue!38!white,line width=.12pt] (-1.806687,-0.201882)--(-1.694686,-0.138917)--(-1.755330,-0.223661)--(-1.871339,-0.292227)--cycle;
\path[fill=blue!33!white,draw=blue!33!white,line width=.12pt] (1.871339,-0.292227)--(1.755330,-0.223661)--(1.694686,-0.138917)--(1.806687,-0.201882)--cycle;
\path[fill=blue!24!white,draw=blue!24!white,line width=.12pt] (1.464780,0.126340)--(1.424631,0.245578)--(1.343133,0.313564)--(1.380985,0.196241)--cycle;
\path[fill=blue!31!white,draw=blue!31!white,line width=.12pt] (-1.380985,0.196241)--(-1.343133,0.313564)--(-1.424631,0.245578)--(-1.464780,0.126340)--cycle;
\path[fill=blue!31!white,draw=blue!31!white,line width=.12pt] (-3.003935,0.190486)--(-2.933190,0.059976)--(-3.006279,-0.085917)--(-3.078787,0.041074)--cycle;
\path[fill=blue!38!white,draw=blue!38!white,line width=.12pt] (3.078787,0.041074)--(3.006279,-0.085917)--(2.933190,0.059976)--(3.003935,0.190486)--cycle;
\path[fill=blue!40!white,draw=blue!40!white,line width=.12pt] (-2.147305,-0.353819)--(-2.003693,-0.336176)--(-2.053620,-0.435837)--(-2.200811,-0.460624)--cycle;
\path[fill=blue!37!white,draw=blue!37!white,line width=.12pt] (2.200811,-0.460624)--(2.053620,-0.435837)--(2.003693,-0.336176)--(2.147305,-0.353819)--cycle;
\path[fill=blue!16!white,draw=blue!16!white,line width=.12pt] (1.248699,0.647553)--(1.283890,0.791962)--(1.174431,0.852346)--(1.142240,0.706283)--cycle;
\path[fill=blue!22!white,draw=blue!22!white,line width=.12pt] (-1.142240,0.706283)--(-1.174431,0.852346)--(-1.283890,0.791962)--(-1.248699,0.647553)--cycle;
\path[fill=blue!15!white,draw=blue!15!white,line width=.12pt] (-1.382778,1.267448)--(-1.480578,1.380371)--(-1.618570,1.304246)--(-1.511656,1.196351)--cycle;
\path[fill=blue!12!white,draw=blue!12!white,line width=.12pt] (1.511656,1.196351)--(1.618570,1.304246)--(1.480578,1.380371)--(1.382778,1.267448)--cycle;
\path[fill=blue!34!white,draw=blue!34!white,line width=.12pt] (-1.529978,0.023985)--(-1.464780,0.126340)--(-1.534468,0.052668)--(-1.602769,-0.052966)--cycle;
\path[fill=blue!27!white,draw=blue!27!white,line width=.12pt] (1.602769,-0.052966)--(1.534468,0.052668)--(1.464780,0.126340)--(1.529978,0.023985)--cycle;
\path[fill=blue!21!white,draw=blue!21!white,line width=.12pt] (-2.900154,0.865772)--(-2.942213,0.744515)--(-3.047499,0.597387)--(-3.003935,0.720748)--cycle;
\path[fill=blue!30!white,draw=blue!30!white,line width=.12pt] (3.003935,0.720748)--(3.047499,0.597387)--(2.942213,0.744515)--(2.900154,0.865772)--cycle;
\path[fill=blue!14!white,draw=blue!14!white,line width=.12pt] (-1.734580,1.391274)--(-1.855225,1.454092)--(-1.995527,1.359505)--(-1.865758,1.302838)--cycle;
\path[fill=blue!14!white,draw=blue!14!white,line width=.12pt] (1.865758,1.302838)--(1.995527,1.359505)--(1.855225,1.454092)--(1.734580,1.391274)--cycle;
\path[fill=blue!14!white,draw=blue!14!white,line width=.12pt] (1.283890,0.791962)--(1.341037,0.935191)--(1.226706,0.998263)--(1.174431,0.852346)--cycle;
\path[fill=blue!20!white,draw=blue!20!white,line width=.12pt] (-1.174431,0.852346)--(-1.226706,0.998263)--(-1.341037,0.935191)--(-1.283890,0.791962)--cycle;
\path[fill=blue!16!white,draw=blue!16!white,line width=.12pt] (-1.297056,1.138426)--(-1.382778,1.267448)--(-1.511656,1.196351)--(-1.417944,1.071736)--cycle;
\path[fill=blue!12!white,draw=blue!12!white,line width=.12pt] (1.417944,1.071736)--(1.511656,1.196351)--(1.382778,1.267448)--(1.297056,1.138426)--cycle;
\path[fill=blue!25!white,draw=blue!25!white,line width=.12pt] (-1.236817,0.507515)--(-1.248699,0.647553)--(-1.343133,0.583889)--(-1.330351,0.444456)--cycle;
\path[fill=blue!18!white,draw=blue!18!white,line width=.12pt] (1.330351,0.444456)--(1.343133,0.583889)--(1.248699,0.647553)--(1.236817,0.507515)--cycle;
\path[fill=blue!15!white,draw=blue!15!white,line width=.12pt] (-2.125296,1.389547)--(-2.250079,1.391810)--(-2.386608,1.277917)--(-2.254254,1.281970)--cycle;
\path[fill=blue!17!white,draw=blue!17!white,line width=.12pt] (2.254254,1.281970)--(2.386608,1.277917)--(2.250079,1.391810)--(2.125296,1.389547)--cycle;
\path[fill=blue!13!white,draw=blue!13!white,line width=.12pt] (1.341037,0.935191)--(1.417944,1.071736)--(1.297056,1.138426)--(1.226706,0.998263)--cycle;
\path[fill=blue!18!white,draw=blue!18!white,line width=.12pt] (-1.226706,0.998263)--(-1.297056,1.138426)--(-1.417944,1.071736)--(-1.341037,0.935191)--cycle;
\path[fill=blue!28!white,draw=blue!28!white,line width=.12pt] (-3.047499,0.327062)--(-3.003935,0.190486)--(-3.078787,0.041074)--(-3.123436,0.175483)--cycle;
\path[fill=blue!36!white,draw=blue!36!white,line width=.12pt] (3.123436,0.175483)--(3.078787,0.041074)--(3.003935,0.190486)--(3.047499,0.327062)--cycle;
\path[fill=blue!39!white,draw=blue!39!white,line width=.12pt] (-2.003693,-0.336176)--(-1.871339,-0.292227)--(-1.917969,-0.385305)--(-2.053620,-0.435837)--cycle;
\path[fill=blue!35!white,draw=blue!35!white,line width=.12pt] (2.053620,-0.435837)--(1.917969,-0.385305)--(1.871339,-0.292227)--(2.003693,-0.336176)--cycle;
\path[fill=blue!37!white,draw=blue!37!white,line width=.12pt] (-1.694686,-0.138917)--(-1.602769,-0.052966)--(-1.660123,-0.133114)--(-1.755330,-0.223661)--cycle;
\path[fill=blue!30!white,draw=blue!30!white,line width=.12pt] (1.755330,-0.223661)--(1.660123,-0.133114)--(1.602769,-0.052966)--(1.694686,-0.138917)--cycle;
\path[fill=blue!17!white,draw=blue!17!white,line width=.12pt] (-2.508586,1.246493)--(-2.615501,1.188905)--(-2.739936,1.057357)--(-2.627935,1.120322)--cycle;
\path[fill=blue!22!white,draw=blue!22!white,line width=.12pt] (2.627935,1.120322)--(2.739936,1.057357)--(2.615501,1.188905)--(2.508586,1.246493)--cycle;
\path[fill=blue!29!white,draw=blue!29!white,line width=.12pt] (-1.343133,0.313564)--(-1.330351,0.444456)--(-1.411074,0.377117)--(-1.424631,0.245578)--cycle;
\path[fill=blue!21!white,draw=blue!21!white,line width=.12pt] (1.424631,0.245578)--(1.411074,0.377117)--(1.330351,0.444456)--(1.343133,0.313564)--cycle;
\path[fill=blue!13!white,draw=blue!13!white,line width=.12pt] (1.740975,1.221724)--(1.865758,1.302838)--(1.734580,1.391274)--(1.618570,1.304246)--cycle;
\path[fill=blue!15!white,draw=blue!15!white,line width=.12pt] (-1.618570,1.304246)--(-1.734580,1.391274)--(-1.865758,1.302838)--(-1.740975,1.221724)--cycle;
\path[fill=blue!16!white,draw=blue!16!white,line width=.12pt] (1.343133,0.583889)--(1.380985,0.726503)--(1.283890,0.791962)--(1.248699,0.647553)--cycle;
\path[fill=blue!23!white,draw=blue!23!white,line width=.12pt] (-1.248699,0.647553)--(-1.283890,0.791962)--(-1.380985,0.726503)--(-1.343133,0.583889)--cycle;
\path[fill=blue!32!white,draw=blue!32!white,line width=.12pt] (-1.464780,0.126340)--(-1.424631,0.245578)--(-1.492409,0.173925)--(-1.534468,0.052668)--cycle;
\path[fill=blue!24!white,draw=blue!24!white,line width=.12pt] (1.534468,0.052668)--(1.492409,0.173925)--(1.424631,0.245578)--(1.464780,0.126340)--cycle;
\path[fill=blue!19!white,draw=blue!19!white,line width=.12pt] (-2.831853,0.971406)--(-2.900154,0.865772)--(-3.003935,0.720748)--(-2.933190,0.829797)--cycle;
\path[fill=blue!27!white,draw=blue!27!white,line width=.12pt] (2.933190,0.829797)--(3.003935,0.720748)--(2.900154,0.865772)--(2.831853,0.971406)--cycle;
\path[fill=blue!38!white,draw=blue!38!white,line width=.12pt] (-2.654149,-0.376111)--(-2.506958,-0.430175)--(-2.543764,-0.554239)--(-2.693116,-0.507458)--cycle;
\path[fill=blue!41!white,draw=blue!41!white,line width=.12pt] (2.693116,-0.507458)--(2.543764,-0.554239)--(2.506958,-0.430175)--(2.654149,-0.376111)--cycle;
\path[fill=blue!39!white,draw=blue!39!white,line width=.12pt] (-2.506958,-0.430175)--(-2.353885,-0.458712)--(-2.388443,-0.575200)--(-2.543764,-0.554239)--cycle;
\path[fill=blue!40!white,draw=blue!40!white,line width=.12pt] (2.543764,-0.554239)--(2.388443,-0.575200)--(2.353885,-0.458712)--(2.506958,-0.430175)--cycle;
\path[fill=blue!26!white,draw=blue!26!white,line width=.12pt] (-3.062209,0.464455)--(-3.047499,0.327062)--(-3.123436,0.175483)--(-3.138513,0.312145)--cycle;
\path[fill=blue!34!white,draw=blue!34!white,line width=.12pt] (3.138513,0.312145)--(3.123436,0.175483)--(3.047499,0.327062)--(3.062209,0.464455)--cycle;
\path[fill=blue!38!white,draw=blue!38!white,line width=.12pt] (-1.871339,-0.292227)--(-1.755330,-0.223661)--(-1.799069,-0.310969)--(-1.917969,-0.385305)--cycle;
\path[fill=blue!33!white,draw=blue!33!white,line width=.12pt] (1.917969,-0.385305)--(1.799069,-0.310969)--(1.755330,-0.223661)--(1.871339,-0.292227)--cycle;
\path[fill=blue!14!white,draw=blue!14!white,line width=.12pt] (-1.995527,1.359505)--(-2.125296,1.389547)--(-2.254254,1.281970)--(-2.116611,1.258496)--cycle;
\path[fill=blue!15!white,draw=blue!15!white,line width=.12pt] (2.116611,1.258496)--(2.254254,1.281970)--(2.125296,1.389547)--(1.995527,1.359505)--cycle;
\path[fill=blue!36!white,draw=blue!36!white,line width=.12pt] (-2.789801,-0.298596)--(-2.654149,-0.376111)--(-2.693116,-0.507458)--(-2.830759,-0.436656)--cycle;
\path[fill=blue!41!white,draw=blue!41!white,line width=.12pt] (2.830759,-0.436656)--(2.693116,-0.507458)--(2.654149,-0.376111)--(2.789801,-0.298596)--cycle;
\path[fill=blue!12!white,draw=blue!12!white,line width=.12pt] (1.625975,1.119280)--(1.740975,1.221724)--(1.618570,1.304246)--(1.511656,1.196351)--cycle;
\path[fill=blue!16!white,draw=blue!16!white,line width=.12pt] (-1.511656,1.196351)--(-1.618570,1.304246)--(-1.740975,1.221724)--(-1.625975,1.119280)--cycle;
\path[fill=blue!40!white,draw=blue!40!white,line width=.12pt] (-2.353885,-0.458712)--(-2.200811,-0.460624)--(-2.233123,-0.569536)--(-2.388443,-0.575200)--cycle;
\path[fill=blue!39!white,draw=blue!39!white,line width=.12pt] (2.388443,-0.575200)--(2.233123,-0.569536)--(2.200811,-0.460624)--(2.353885,-0.458712)--cycle;
\path[fill=blue!14!white,draw=blue!14!white,line width=.12pt] (1.380985,0.726503)--(1.442454,0.866819)--(1.341037,0.935191)--(1.283890,0.791962)--cycle;
\path[fill=blue!21!white,draw=blue!21!white,line width=.12pt] (-1.283890,0.791962)--(-1.341037,0.935191)--(-1.442454,0.866819)--(-1.380985,0.726503)--cycle;
\path[fill=blue!35!white,draw=blue!35!white,line width=.12pt] (-1.602769,-0.052966)--(-1.534468,0.052668)--(-1.589379,-0.024064)--(-1.660123,-0.133114)--cycle;
\path[fill=blue!27!white,draw=blue!27!white,line width=.12pt] (1.660123,-0.133114)--(1.589379,-0.024064)--(1.534468,0.052668)--(1.602769,-0.052966)--cycle;
\path[fill=blue!35!white,draw=blue!35!white,line width=.12pt] (-2.908701,-0.200610)--(-2.789801,-0.298596)--(-2.830759,-0.436656)--(-2.951405,-0.344554)--cycle;
\path[fill=blue!40!white,draw=blue!40!white,line width=.12pt] (2.951405,-0.344554)--(2.830759,-0.436656)--(2.789801,-0.298596)--(2.908701,-0.200610)--cycle;
\path[fill=blue!26!white,draw=blue!26!white,line width=.12pt] (-1.330351,0.444456)--(-1.343133,0.583889)--(-1.424631,0.515903)--(-1.411074,0.377117)--cycle;
\path[fill=blue!18!white,draw=blue!18!white,line width=.12pt] (1.411074,0.377117)--(1.424631,0.515903)--(1.343133,0.583889)--(1.330351,0.444456)--cycle;
\path[fill=blue!15!white,draw=blue!15!white,line width=.12pt] (-2.386608,1.277917)--(-2.508586,1.246493)--(-2.627935,1.120322)--(-2.500154,1.157881)--cycle;
\path[fill=blue!19!white,draw=blue!19!white,line width=.12pt] (2.500154,1.157881)--(2.627935,1.120322)--(2.508586,1.246493)--(2.386608,1.277917)--cycle;
\path[fill=blue!12!white,draw=blue!12!white,line width=.12pt] (1.525177,0.999443)--(1.625975,1.119280)--(1.511656,1.196351)--(1.417944,1.071736)--cycle;
\path[fill=blue!17!white,draw=blue!17!white,line width=.12pt] (-1.417944,1.071736)--(-1.511656,1.196351)--(-1.625975,1.119280)--(-1.525177,0.999443)--cycle;
\path[fill=blue!19!white,draw=blue!19!white,line width=.12pt] (-1.341037,0.935191)--(-1.417944,1.071736)--(-1.525177,0.999443)--(-1.442454,0.866819)--cycle;
\path[fill=blue!13!white,draw=blue!13!white,line width=.12pt] (1.442454,0.866819)--(1.525177,0.999443)--(1.417944,1.071736)--(1.341037,0.935191)--cycle;
\path[fill=blue!40!white,draw=blue!40!white,line width=.12pt] (-2.200811,-0.460624)--(-2.053620,-0.435837)--(-2.083771,-0.537466)--(-2.233123,-0.569536)--cycle;
\path[fill=blue!37!white,draw=blue!37!white,line width=.12pt] (2.233123,-0.569536)--(2.083771,-0.537466)--(2.053620,-0.435837)--(2.200811,-0.460624)--cycle;
\path[fill=blue!14!white,draw=blue!14!white,line width=.12pt] (-1.865758,1.302838)--(-1.995527,1.359505)--(-2.116611,1.258496)--(-1.978968,1.208398)--cycle;
\path[fill=blue!14!white,draw=blue!14!white,line width=.12pt] (1.978968,1.208398)--(2.116611,1.258496)--(1.995527,1.359505)--(1.865758,1.302838)--cycle;
\path[fill=blue!32!white,draw=blue!32!white,line width=.12pt] (-3.006279,-0.085917)--(-2.908701,-0.200610)--(-2.951405,-0.344554)--(-3.050416,-0.234690)--cycle;
\path[fill=blue!39!white,draw=blue!39!white,line width=.12pt] (3.050416,-0.234690)--(2.951405,-0.344554)--(2.908701,-0.200610)--(3.006279,-0.085917)--cycle;
\path[fill=blue!29!white,draw=blue!29!white,line width=.12pt] (-1.424631,0.245578)--(-1.411074,0.377117)--(-1.478207,0.306147)--(-1.492409,0.173925)--cycle;
\path[fill=blue!21!white,draw=blue!21!white,line width=.12pt] (1.492409,0.173925)--(1.478207,0.306147)--(1.411074,0.377117)--(1.424631,0.245578)--cycle;
\path[fill=blue!23!white,draw=blue!23!white,line width=.12pt] (-3.047499,0.597387)--(-3.062209,0.464455)--(-3.138513,0.312145)--(-3.123436,0.445808)--cycle;
\path[fill=blue!32!white,draw=blue!32!white,line width=.12pt] (3.123436,0.445808)--(3.138513,0.312145)--(3.062209,0.464455)--(3.047499,0.597387)--cycle;
\path[fill=blue!37!white,draw=blue!37!white,line width=.12pt] (-1.755330,-0.223661)--(-1.660123,-0.133114)--(-1.701490,-0.215687)--(-1.799069,-0.310969)--cycle;
\path[fill=blue!30!white,draw=blue!30!white,line width=.12pt] (1.799069,-0.310969)--(1.701490,-0.215687)--(1.660123,-0.133114)--(1.755330,-0.223661)--cycle;
\path[fill=blue!17!white,draw=blue!17!white,line width=.12pt] (-2.739936,1.057357)--(-2.831853,0.971406)--(-2.933190,0.829797)--(-2.837984,0.920344)--cycle;
\path[fill=blue!24!white,draw=blue!24!white,line width=.12pt] (2.837984,0.920344)--(2.933190,0.829797)--(2.831853,0.971406)--(2.739936,1.057357)--cycle;
\path[fill=blue!24!white,draw=blue!24!white,line width=.12pt] (-1.343133,0.583889)--(-1.380985,0.726503)--(-1.464780,0.656601)--(-1.424631,0.515903)--cycle;
\path[fill=blue!16!white,draw=blue!16!white,line width=.12pt] (1.424631,0.515903)--(1.464780,0.656601)--(1.380985,0.726503)--(1.343133,0.583889)--cycle;
\path[fill=blue!40!white,draw=blue!40!white,line width=.12pt] (-2.053620,-0.435837)--(-1.917969,-0.385305)--(-1.946127,-0.480221)--(-2.083771,-0.537466)--cycle;
\path[fill=blue!36!white,draw=blue!36!white,line width=.12pt] (2.083771,-0.537466)--(1.946127,-0.480221)--(1.917969,-0.385305)--(2.053620,-0.435837)--cycle;
\path[fill=blue!24!white,draw=blue!24!white,line width=.12pt] (1.589379,-0.024064)--(1.545814,0.099296)--(1.492409,0.173925)--(1.534468,0.052668)--cycle;
\path[fill=blue!33!white,draw=blue!33!white,line width=.12pt] (-1.534468,0.052668)--(-1.492409,0.173925)--(-1.545814,0.099296)--(-1.589379,-0.024064)--cycle;
\path[fill=blue!15!white,draw=blue!15!white,line width=.12pt] (-2.254254,1.281970)--(-2.386608,1.277917)--(-2.500154,1.157881)--(-2.361503,1.168591)--cycle;
\path[fill=blue!17!white,draw=blue!17!white,line width=.12pt] (2.361503,1.168591)--(2.500154,1.157881)--(2.386608,1.277917)--(2.254254,1.281970)--cycle;
\path[fill=blue!15!white,draw=blue!15!white,line width=.12pt] (-1.740975,1.221724)--(-1.865758,1.302838)--(-1.978968,1.208398)--(-1.846614,1.133600)--cycle;
\path[fill=blue!13!white,draw=blue!13!white,line width=.12pt] (1.846614,1.133600)--(1.978968,1.208398)--(1.865758,1.302838)--(1.740975,1.221724)--cycle;
\path[fill=blue!30!white,draw=blue!30!white,line width=.12pt] (-3.078787,0.041074)--(-3.006279,-0.085917)--(-3.050416,-0.234690)--(-3.123988,-0.111288)--cycle;
\path[fill=blue!38!white,draw=blue!38!white,line width=.12pt] (3.123988,-0.111288)--(3.050416,-0.234690)--(3.006279,-0.085917)--(3.078787,0.041074)--cycle;
\path[fill=blue!15!white,draw=blue!15!white,line width=.12pt] (1.464780,0.656601)--(1.529978,0.793806)--(1.442454,0.866819)--(1.380985,0.726503)--cycle;
\path[fill=blue!21!white,draw=blue!21!white,line width=.12pt] (-1.380985,0.726503)--(-1.442454,0.866819)--(-1.529978,0.793806)--(-1.464780,0.656601)--cycle;
\path[fill=blue!27!white,draw=blue!27!white,line width=.12pt] (-1.411074,0.377117)--(-1.424631,0.515903)--(-1.492409,0.444251)--(-1.478207,0.306147)--cycle;
\path[fill=blue!19!white,draw=blue!19!white,line width=.12pt] (1.478207,0.306147)--(1.492409,0.444251)--(1.424631,0.515903)--(1.411074,0.377117)--cycle;
\path[fill=blue!12!white,draw=blue!12!white,line width=.12pt] (1.724636,1.036978)--(1.846614,1.133600)--(1.740975,1.221724)--(1.625975,1.119280)--cycle;
\path[fill=blue!16!white,draw=blue!16!white,line width=.12pt] (-1.625975,1.119280)--(-1.740975,1.221724)--(-1.846614,1.133600)--(-1.724636,1.036978)--cycle;
\path[fill=blue!21!white,draw=blue!21!white,line width=.12pt] (-3.003935,0.720748)--(-3.047499,0.597387)--(-3.123436,0.445808)--(-3.078787,0.571335)--cycle;
\path[fill=blue!29!white,draw=blue!29!white,line width=.12pt] (3.078787,0.571335)--(3.123436,0.445808)--(3.047499,0.597387)--(3.003935,0.720748)--cycle;
\path[fill=blue!39!white,draw=blue!39!white,line width=.12pt] (-1.917969,-0.385305)--(-1.799069,-0.310969)--(-1.825482,-0.400001)--(-1.946127,-0.480221)--cycle;
\path[fill=blue!33!white,draw=blue!33!white,line width=.12pt] (1.946127,-0.480221)--(1.825482,-0.400001)--(1.799069,-0.310969)--(1.917969,-0.385305)--cycle;
\path[fill=blue!35!white,draw=blue!35!white,line width=.12pt] (-1.660123,-0.133114)--(-1.589379,-0.024064)--(-1.628983,-0.103118)--(-1.701490,-0.215687)--cycle;
\path[fill=blue!28!white,draw=blue!28!white,line width=.12pt] (1.701490,-0.215687)--(1.628983,-0.103118)--(1.589379,-0.024064)--(1.660123,-0.133114)--cycle;
\path[fill=blue!19!white,draw=blue!19!white,line width=.12pt] (-1.442454,0.866819)--(-1.525177,0.999443)--(-1.617721,0.922243)--(-1.529978,0.793806)--cycle;
\path[fill=blue!13!white,draw=blue!13!white,line width=.12pt] (1.529978,0.793806)--(1.617721,0.922243)--(1.525177,0.999443)--(1.442454,0.866819)--cycle;
\path[fill=blue!21!white,draw=blue!21!white,line width=.12pt] (2.721975,0.988910)--(2.837984,0.920344)--(2.739936,1.057357)--(2.627935,1.120322)--cycle;
\path[fill=blue!16!white,draw=blue!16!white,line width=.12pt] (-2.627935,1.120322)--(-2.739936,1.057357)--(-2.837984,0.920344)--(-2.721975,0.988910)--cycle;
\path[fill=blue!17!white,draw=blue!17!white,line width=.12pt] (-1.525177,0.999443)--(-1.625975,1.119280)--(-1.724636,1.036978)--(-1.617721,0.922243)--cycle;
\path[fill=blue!13!white,draw=blue!13!white,line width=.12pt] (1.617721,0.922243)--(1.724636,1.036978)--(1.625975,1.119280)--(1.525177,0.999443)--cycle;
\path[fill=blue!14!white,draw=blue!14!white,line width=.12pt] (-2.116611,1.258496)--(-2.254254,1.281970)--(-2.361503,1.168591)--(-2.217311,1.152040)--cycle;
\path[fill=blue!15!white,draw=blue!15!white,line width=.12pt] (2.217311,1.152040)--(2.361503,1.168591)--(2.254254,1.281970)--(2.116611,1.258496)--cycle;
\path[fill=blue!27!white,draw=blue!27!white,line width=.12pt] (-3.123436,0.175483)--(-3.078787,0.041074)--(-3.123988,-0.111288)--(-3.169293,0.020912)--cycle;
\path[fill=blue!36!white,draw=blue!36!white,line width=.12pt] (3.169293,0.020912)--(3.123988,-0.111288)--(3.078787,0.041074)--(3.123436,0.175483)--cycle;
\path[fill=blue!30!white,draw=blue!30!white,line width=.12pt] (-1.492409,0.173925)--(-1.478207,0.306147)--(-1.531105,0.232228)--(-1.545814,0.099296)--cycle;
\path[fill=blue!22!white,draw=blue!22!white,line width=.12pt] (1.545814,0.099296)--(1.531105,0.232228)--(1.478207,0.306147)--(1.492409,0.173925)--cycle;
\path[fill=blue!39!white,draw=blue!39!white,line width=.12pt] (-2.543764,-0.554239)--(-2.388443,-0.575200)--(-2.400000,-0.692820)--(-2.556072,-0.679508)--cycle;
\path[fill=blue!40!white,draw=blue!40!white,line width=.12pt] (2.556072,-0.679508)--(2.400000,-0.692820)--(2.388443,-0.575200)--(2.543764,-0.554239)--cycle;
\path[fill=blue!40!white,draw=blue!40!white,line width=.12pt] (-2.388443,-0.575200)--(-2.233123,-0.569536)--(-2.243928,-0.679508)--(-2.400000,-0.692820)--cycle;
\path[fill=blue!39!white,draw=blue!39!white,line width=.12pt] (2.400000,-0.692820)--(2.243928,-0.679508)--(2.233123,-0.569536)--(2.388443,-0.575200)--cycle;
\path[fill=blue!38!white,draw=blue!38!white,line width=.12pt] (-2.693116,-0.507458)--(-2.543764,-0.554239)--(-2.556072,-0.679508)--(-2.706147,-0.640083)--cycle;
\path[fill=blue!40!white,draw=blue!40!white,line width=.12pt] (2.706147,-0.640083)--(2.556072,-0.679508)--(2.543764,-0.554239)--(2.693116,-0.507458)--cycle;
\path[fill=blue!24!white,draw=blue!24!white,line width=.12pt] (-1.424631,0.515903)--(-1.464780,0.656601)--(-1.534468,0.582930)--(-1.492409,0.444251)--cycle;
\path[fill=blue!17!white,draw=blue!17!white,line width=.12pt] (1.492409,0.444251)--(1.534468,0.582930)--(1.464780,0.656601)--(1.424631,0.515903)--cycle;
\path[fill=blue!38!white,draw=blue!38!white,line width=.12pt] (-1.799069,-0.310969)--(-1.701490,-0.215687)--(-1.726471,-0.299889)--(-1.825482,-0.400001)--cycle;
\path[fill=blue!31!white,draw=blue!31!white,line width=.12pt] (1.825482,-0.400001)--(1.726471,-0.299889)--(1.701490,-0.215687)--(1.799069,-0.310969)--cycle;
\path[fill=blue!40!white,draw=blue!40!white,line width=.12pt] (-2.233123,-0.569536)--(-2.083771,-0.537466)--(-2.093853,-0.640083)--(-2.243928,-0.679508)--cycle;
\path[fill=blue!38!white,draw=blue!38!white,line width=.12pt] (2.243928,-0.679508)--(2.093853,-0.640083)--(2.083771,-0.537466)--(2.233123,-0.569536)--cycle;
\path[fill=blue!36!white,draw=blue!36!white,line width=.12pt] (-2.830759,-0.436656)--(-2.693116,-0.507458)--(-2.706147,-0.640083)--(-2.844456,-0.576059)--cycle;
\path[fill=blue!40!white,draw=blue!40!white,line width=.12pt] (2.844456,-0.576059)--(2.706147,-0.640083)--(2.693116,-0.507458)--(2.830759,-0.436656)--cycle;
\path[fill=blue!19!white,draw=blue!19!white,line width=.12pt] (-2.933190,0.829797)--(-3.003935,0.720748)--(-3.078787,0.571335)--(-3.006279,0.683903)--cycle;
\path[fill=blue!26!white,draw=blue!26!white,line width=.12pt] (3.006279,0.683903)--(3.078787,0.571335)--(3.003935,0.720748)--(2.933190,0.829797)--cycle;
\path[fill=blue!15!white,draw=blue!15!white,line width=.12pt] (-1.978968,1.208398)--(-2.116611,1.258496)--(-2.217311,1.152040)--(-2.073119,1.108865)--cycle;
\path[fill=blue!14!white,draw=blue!14!white,line width=.12pt] (2.073119,1.108865)--(2.217311,1.152040)--(2.116611,1.258496)--(1.978968,1.208398)--cycle;
\path[fill=blue!33!white,draw=blue!33!white,line width=.12pt] (-1.589379,-0.024064)--(-1.545814,0.099296)--(-1.584333,0.022409)--(-1.628983,-0.103118)--cycle;
\path[fill=blue!25!white,draw=blue!25!white,line width=.12pt] (1.628983,-0.103118)--(1.584333,0.022409)--(1.545814,0.099296)--(1.589379,-0.024064)--cycle;
\path[fill=blue!15!white,draw=blue!15!white,line width=.12pt] (-2.500154,1.157881)--(-2.627935,1.120322)--(-2.721975,0.988910)--(-2.589621,1.032859)--cycle;
\path[fill=blue!19!white,draw=blue!19!white,line width=.12pt] (2.589621,1.032859)--(2.721975,0.988910)--(2.627935,1.120322)--(2.500154,1.157881)--cycle;
\path[fill=blue!25!white,draw=blue!25!white,line width=.12pt] (-3.138513,0.312145)--(-3.123436,0.175483)--(-3.169293,0.020912)--(-3.184591,0.156827)--cycle;
\path[fill=blue!34!white,draw=blue!34!white,line width=.12pt] (3.184591,0.156827)--(3.169293,0.020912)--(3.123436,0.175483)--(3.138513,0.312145)--cycle;
\path[fill=blue!22!white,draw=blue!22!white,line width=.12pt] (-1.464780,0.656601)--(-1.529978,0.793806)--(-1.602769,0.716855)--(-1.534468,0.582930)--cycle;
\path[fill=blue!15!white,draw=blue!15!white,line width=.12pt] (1.534468,0.582930)--(1.602769,0.716855)--(1.529978,0.793806)--(1.464780,0.656601)--cycle;
\path[fill=blue!40!white,draw=blue!40!white,line width=.12pt] (-2.083771,-0.537466)--(-1.946127,-0.480221)--(-1.955544,-0.576059)--(-2.093853,-0.640083)--cycle;
\path[fill=blue!36!white,draw=blue!36!white,line width=.12pt] (2.093853,-0.640083)--(1.955544,-0.576059)--(1.946127,-0.480221)--(2.083771,-0.537466)--cycle;
\path[fill=blue!34!white,draw=blue!34!white,line width=.12pt] (-2.951405,-0.344554)--(-2.830759,-0.436656)--(-2.844456,-0.576059)--(-2.965685,-0.489898)--cycle;
\path[fill=blue!40!white,draw=blue!40!white,line width=.12pt] (2.965685,-0.489898)--(2.844456,-0.576059)--(2.830759,-0.436656)--(2.951405,-0.344554)--cycle;
\path[fill=blue!28!white,draw=blue!28!white,line width=.12pt] (-1.478207,0.306147)--(-1.492409,0.444251)--(-1.545814,0.369621)--(-1.531105,0.232228)--cycle;
\path[fill=blue!19!white,draw=blue!19!white,line width=.12pt] (1.531105,0.232228)--(1.545814,0.369621)--(1.492409,0.444251)--(1.478207,0.306147)--cycle;
\path[fill=blue!15!white,draw=blue!15!white,line width=.12pt] (-1.846614,1.133600)--(-1.978968,1.208398)--(-2.073119,1.108865)--(-1.934468,1.040724)--cycle;
\path[fill=blue!13!white,draw=blue!13!white,line width=.12pt] (1.934468,1.040724)--(2.073119,1.108865)--(1.978968,1.208398)--(1.846614,1.133600)--cycle;
\path[fill=blue!20!white,draw=blue!20!white,line width=.12pt] (-1.529978,0.793806)--(-1.617721,0.922243)--(-1.694686,0.840879)--(-1.602769,0.716855)--cycle;
\path[fill=blue!14!white,draw=blue!14!white,line width=.12pt] (1.602769,0.716855)--(1.694686,0.840879)--(1.617721,0.922243)--(1.529978,0.793806)--cycle;
\path[fill=blue!36!white,draw=blue!36!white,line width=.12pt] (-1.701490,-0.215687)--(-1.628983,-0.103118)--(-1.652899,-0.183733)--(-1.726471,-0.299889)--cycle;
\path[fill=blue!28!white,draw=blue!28!white,line width=.12pt] (1.726471,-0.299889)--(1.652899,-0.183733)--(1.628983,-0.103118)--(1.701490,-0.215687)--cycle;
\path[fill=blue!16!white,draw=blue!16!white,line width=.12pt] (-1.724636,1.036978)--(-1.846614,1.133600)--(-1.934468,1.040724)--(-1.806687,0.950236)--cycle;
\path[fill=blue!13!white,draw=blue!13!white,line width=.12pt] (1.806687,0.950236)--(1.934468,1.040724)--(1.846614,1.133600)--(1.724636,1.036978)--cycle;
\path[fill=blue!18!white,draw=blue!18!white,line width=.12pt] (-1.617721,0.922243)--(-1.724636,1.036978)--(-1.806687,0.950236)--(-1.694686,0.840879)--cycle;
\path[fill=blue!13!white,draw=blue!13!white,line width=.12pt] (1.694686,0.840879)--(1.806687,0.950236)--(1.724636,1.036978)--(1.617721,0.922243)--cycle;
\path[fill=blue!39!white,draw=blue!39!white,line width=.12pt] (-1.946127,-0.480221)--(-1.825482,-0.400001)--(-1.834315,-0.489898)--(-1.955544,-0.576059)--cycle;
\path[fill=blue!34!white,draw=blue!34!white,line width=.12pt] (1.955544,-0.576059)--(1.834315,-0.489898)--(1.825482,-0.400001)--(1.946127,-0.480221)--cycle;
\path[fill=blue!14!white,draw=blue!14!white,line width=.12pt] (-2.361503,1.168591)--(-2.500154,1.157881)--(-2.589621,1.032859)--(-2.446009,1.050502)--cycle;
\path[fill=blue!17!white,draw=blue!17!white,line width=.12pt] (2.446009,1.050502)--(2.589621,1.032859)--(2.500154,1.157881)--(2.361503,1.168591)--cycle;
\path[fill=blue!17!white,draw=blue!17!white,line width=.12pt] (-2.837984,0.920344)--(-2.933190,0.829797)--(-3.006279,0.683903)--(-2.908701,0.779186)--cycle;
\path[fill=blue!24!white,draw=blue!24!white,line width=.12pt] (2.908701,0.779186)--(3.006279,0.683903)--(2.933190,0.829797)--(2.837984,0.920344)--cycle;
\path[fill=blue!32!white,draw=blue!32!white,line width=.12pt] (-3.050416,-0.234690)--(-2.951405,-0.344554)--(-2.965685,-0.489898)--(-3.065176,-0.384910)--cycle;
\path[fill=blue!39!white,draw=blue!39!white,line width=.12pt] (3.065176,-0.384910)--(2.965685,-0.489898)--(2.951405,-0.344554)--(3.050416,-0.234690)--cycle;
\path[fill=blue!31!white,draw=blue!31!white,line width=.12pt] (-1.545814,0.099296)--(-1.531105,0.232228)--(-1.569256,0.156072)--(-1.584333,0.022409)--cycle;
\path[fill=blue!22!white,draw=blue!22!white,line width=.12pt] (1.584333,0.022409)--(1.569256,0.156072)--(1.531105,0.232228)--(1.545814,0.099296)--cycle;
\path[fill=blue!22!white,draw=blue!22!white,line width=.12pt] (-3.123436,0.445808)--(-3.138513,0.312145)--(-3.184591,0.156827)--(-3.169293,0.291237)--cycle;
\path[fill=blue!31!white,draw=blue!31!white,line width=.12pt] (3.169293,0.291237)--(3.184591,0.156827)--(3.138513,0.312145)--(3.123436,0.445808)--cycle;
\path[fill=blue!25!white,draw=blue!25!white,line width=.12pt] (-1.492409,0.444251)--(-1.534468,0.582930)--(-1.589379,0.506197)--(-1.545814,0.369621)--cycle;
\path[fill=blue!17!white,draw=blue!17!white,line width=.12pt] (1.545814,0.369621)--(1.589379,0.506197)--(1.534468,0.582930)--(1.492409,0.444251)--cycle;
\path[fill=blue!38!white,draw=blue!38!white,line width=.12pt] (-1.825482,-0.400001)--(-1.726471,-0.299889)--(-1.734824,-0.384910)--(-1.834315,-0.489898)--cycle;
\path[fill=blue!31!white,draw=blue!31!white,line width=.12pt] (1.834315,-0.489898)--(1.734824,-0.384910)--(1.726471,-0.299889)--(1.825482,-0.400001)--cycle;
\path[fill=blue!14!white,draw=blue!14!white,line width=.12pt] (-2.217311,1.152040)--(-2.361503,1.168591)--(-2.446009,1.050502)--(-2.296657,1.041162)--cycle;
\path[fill=blue!15!white,draw=blue!15!white,line width=.12pt] (2.296657,1.041162)--(2.446009,1.050502)--(2.361503,1.168591)--(2.217311,1.152040)--cycle;
\path[fill=blue!34!white,draw=blue!34!white,line width=.12pt] (-1.628983,-0.103118)--(-1.584333,0.022409)--(-1.607593,-0.055995)--(-1.652899,-0.183733)--cycle;
\path[fill=blue!25!white,draw=blue!25!white,line width=.12pt] (1.652899,-0.183733)--(1.607593,-0.055995)--(1.584333,0.022409)--(1.628983,-0.103118)--cycle;
\path[fill=blue!29!white,draw=blue!29!white,line width=.12pt] (-3.123988,-0.111288)--(-3.050416,-0.234690)--(-3.065176,-0.384910)--(-3.139104,-0.265131)--cycle;
\path[fill=blue!37!white,draw=blue!37!white,line width=.12pt] (3.139104,-0.265131)--(3.065176,-0.384910)--(3.050416,-0.234690)--(3.123988,-0.111288)--cycle;
\path[fill=blue!23!white,draw=blue!23!white,line width=.12pt] (-1.534468,0.582930)--(-1.602769,0.716855)--(-1.660123,0.636707)--(-1.589379,0.506197)--cycle;
\path[fill=blue!15!white,draw=blue!15!white,line width=.12pt] (1.589379,0.506197)--(1.660123,0.636707)--(1.602769,0.716855)--(1.534468,0.582930)--cycle;
\path[fill=blue!16!white,draw=blue!16!white,line width=.12pt] (-2.721975,0.988910)--(-2.837984,0.920344)--(-2.908701,0.779186)--(-2.789801,0.853522)--cycle;
\path[fill=blue!21!white,draw=blue!21!white,line width=.12pt] (2.789801,0.853522)--(2.908701,0.779186)--(2.837984,0.920344)--(2.721975,0.988910)--cycle;
\path[fill=blue!28!white,draw=blue!28!white,line width=.12pt] (-1.531105,0.232228)--(-1.545814,0.369621)--(-1.584333,0.292734)--(-1.569256,0.156072)--cycle;
\path[fill=blue!20!white,draw=blue!20!white,line width=.12pt] (1.569256,0.156072)--(1.584333,0.292734)--(1.545814,0.369621)--(1.531105,0.232228)--cycle;
\path[fill=blue!40!white,draw=blue!40!white,line width=.12pt] (-2.400000,-0.692820)--(-2.243928,-0.679508)--(-2.233123,-0.789480)--(-2.388443,-0.810441)--cycle;
\path[fill=blue!39!white,draw=blue!39!white,line width=.12pt] (2.388443,-0.810441)--(2.233123,-0.789480)--(2.243928,-0.679508)--(2.400000,-0.692820)--cycle;
\path[fill=blue!39!white,draw=blue!39!white,line width=.12pt] (-2.556072,-0.679508)--(-2.400000,-0.692820)--(-2.388443,-0.810441)--(-2.543764,-0.804777)--cycle;
\path[fill=blue!40!white,draw=blue!40!white,line width=.12pt] (2.543764,-0.804777)--(2.388443,-0.810441)--(2.400000,-0.692820)--(2.556072,-0.679508)--cycle;
\path[fill=blue!40!white,draw=blue!40!white,line width=.12pt] (-2.243928,-0.679508)--(-2.093853,-0.640083)--(-2.083771,-0.742699)--(-2.233123,-0.789480)--cycle;
\path[fill=blue!38!white,draw=blue!38!white,line width=.12pt] (2.233123,-0.789480)--(2.083771,-0.742699)--(2.093853,-0.640083)--(2.243928,-0.679508)--cycle;
\path[fill=blue!20!white,draw=blue!20!white,line width=.12pt] (-3.078787,0.571335)--(-3.123436,0.445808)--(-3.169293,0.291237)--(-3.123988,0.418974)--cycle;
\path[fill=blue!29!white,draw=blue!29!white,line width=.12pt] (3.123988,0.418974)--(3.169293,0.291237)--(3.123436,0.445808)--(3.078787,0.571335)--cycle;
\path[fill=blue!15!white,draw=blue!15!white,line width=.12pt] (-2.073119,1.108865)--(-2.217311,1.152040)--(-2.296657,1.041162)--(-2.147305,1.005197)--cycle;
\path[fill=blue!14!white,draw=blue!14!white,line width=.12pt] (2.147305,1.005197)--(2.296657,1.041162)--(2.217311,1.152040)--(2.073119,1.108865)--cycle;
\path[fill=blue!38!white,draw=blue!38!white,line width=.12pt] (-2.706147,-0.640083)--(-2.556072,-0.679508)--(-2.543764,-0.804777)--(-2.693116,-0.772707)--cycle;
\path[fill=blue!40!white,draw=blue!40!white,line width=.12pt] (2.693116,-0.772707)--(2.543764,-0.804777)--(2.556072,-0.679508)--(2.706147,-0.640083)--cycle;
\path[fill=blue!21!white,draw=blue!21!white,line width=.12pt] (-1.602769,0.716855)--(-1.694686,0.840879)--(-1.755330,0.756135)--(-1.660123,0.636707)--cycle;
\path[fill=blue!14!white,draw=blue!14!white,line width=.12pt] (1.660123,0.636707)--(1.755330,0.756135)--(1.694686,0.840879)--(1.602769,0.716855)--cycle;
\path[fill=blue!40!white,draw=blue!40!white,line width=.12pt] (-2.093853,-0.640083)--(-1.955544,-0.576059)--(-1.946127,-0.671897)--(-2.083771,-0.742699)--cycle;
\path[fill=blue!36!white,draw=blue!36!white,line width=.12pt] (2.083771,-0.742699)--(1.946127,-0.671897)--(1.955544,-0.576059)--(2.093853,-0.640083)--cycle;
\path[fill=blue!37!white,draw=blue!37!white,line width=.12pt] (-1.726471,-0.299889)--(-1.652899,-0.183733)--(-1.660896,-0.265131)--(-1.734824,-0.384910)--cycle;
\path[fill=blue!29!white,draw=blue!29!white,line width=.12pt] (1.734824,-0.384910)--(1.660896,-0.265131)--(1.652899,-0.183733)--(1.726471,-0.299889)--cycle;
\path[fill=blue!16!white,draw=blue!16!white,line width=.12pt] (-1.934468,1.040724)--(-2.073119,1.108865)--(-2.147305,1.005197)--(-2.003693,0.943989)--cycle;
\path[fill=blue!13!white,draw=blue!13!white,line width=.12pt] (2.003693,0.943989)--(2.147305,1.005197)--(2.073119,1.108865)--(1.934468,1.040724)--cycle;
\path[fill=blue!27!white,draw=blue!27!white,line width=.12pt] (-3.169293,0.020912)--(-3.123988,-0.111288)--(-3.139104,-0.265131)--(-3.184628,-0.135163)--cycle;
\path[fill=blue!35!white,draw=blue!35!white,line width=.12pt] (3.184628,-0.135163)--(3.139104,-0.265131)--(3.123988,-0.111288)--(3.169293,0.020912)--cycle;
\path[fill=blue!18!white,draw=blue!18!white,line width=.12pt] (-1.694686,0.840879)--(-1.806687,0.950236)--(-1.871339,0.859891)--(-1.755330,0.756135)--cycle;
\path[fill=blue!13!white,draw=blue!13!white,line width=.12pt] (1.755330,0.756135)--(1.871339,0.859891)--(1.806687,0.950236)--(1.694686,0.840879)--cycle;
\path[fill=blue!32!white,draw=blue!32!white,line width=.12pt] (-1.584333,0.022409)--(-1.569256,0.156072)--(-1.592296,0.078414)--(-1.607593,-0.055995)--cycle;
\path[fill=blue!23!white,draw=blue!23!white,line width=.12pt] (1.607593,-0.055995)--(1.592296,0.078414)--(1.569256,0.156072)--(1.584333,0.022409)--cycle;
\path[fill=blue!17!white,draw=blue!17!white,line width=.12pt] (-1.806687,0.950236)--(-1.934468,1.040724)--(-2.003693,0.943989)--(-1.871339,0.859891)--cycle;
\path[fill=blue!13!white,draw=blue!13!white,line width=.12pt] (1.871339,0.859891)--(2.003693,0.943989)--(1.934468,1.040724)--(1.806687,0.950236)--cycle;
\path[fill=blue!36!white,draw=blue!36!white,line width=.12pt] (-2.844456,-0.576059)--(-2.706147,-0.640083)--(-2.693116,-0.772707)--(-2.830759,-0.715462)--cycle;
\path[fill=blue!40!white,draw=blue!40!white,line width=.12pt] (2.830759,-0.715462)--(2.693116,-0.772707)--(2.706147,-0.640083)--(2.844456,-0.576059)--cycle;
\path[fill=blue!15!white,draw=blue!15!white,line width=.12pt] (-2.589621,1.032859)--(-2.721975,0.988910)--(-2.789801,0.853522)--(-2.654149,0.904054)--cycle;
\path[fill=blue!19!white,draw=blue!19!white,line width=.12pt] (2.654149,0.904054)--(2.789801,0.853522)--(2.721975,0.988910)--(2.589621,1.032859)--cycle;
\path[fill=blue!26!white,draw=blue!26!white,line width=.12pt] (-1.545814,0.369621)--(-1.589379,0.506197)--(-1.628983,0.427143)--(-1.584333,0.292734)--cycle;
\path[fill=blue!18!white,draw=blue!18!white,line width=.12pt] (1.584333,0.292734)--(1.628983,0.427143)--(1.589379,0.506197)--(1.545814,0.369621)--cycle;
\path[fill=blue!40!white,draw=blue!40!white,line width=.12pt] (-1.955544,-0.576059)--(-1.834315,-0.489898)--(-1.825482,-0.579795)--(-1.946127,-0.671897)--cycle;
\path[fill=blue!34!white,draw=blue!34!white,line width=.12pt] (1.946127,-0.671897)--(1.825482,-0.579795)--(1.834315,-0.489898)--(1.955544,-0.576059)--cycle;
\path[fill=blue!18!white,draw=blue!18!white,line width=.12pt] (-3.006279,0.683903)--(-3.078787,0.571335)--(-3.123988,0.418974)--(-3.050416,0.535130)--cycle;
\path[fill=blue!26!white,draw=blue!26!white,line width=.12pt] (3.050416,0.535130)--(3.123988,0.418974)--(3.078787,0.571335)--(3.006279,0.683903)--cycle;
\path[fill=blue!35!white,draw=blue!35!white,line width=.12pt] (-1.652899,-0.183733)--(-1.607593,-0.055995)--(-1.615372,-0.135163)--(-1.660896,-0.265131)--cycle;
\path[fill=blue!26!white,draw=blue!26!white,line width=.12pt] (1.660896,-0.265131)--(1.615372,-0.135163)--(1.607593,-0.055995)--(1.652899,-0.183733)--cycle;
\path[fill=blue!34!white,draw=blue!34!white,line width=.12pt] (-2.965685,-0.489898)--(-2.844456,-0.576059)--(-2.830759,-0.715462)--(-2.951405,-0.635242)--cycle;
\path[fill=blue!39!white,draw=blue!39!white,line width=.12pt] (2.951405,-0.635242)--(2.830759,-0.715462)--(2.844456,-0.576059)--(2.965685,-0.489898)--cycle;
\path[fill=blue!39!white,draw=blue!39!white,line width=.12pt] (-1.834315,-0.489898)--(-1.734824,-0.384910)--(-1.726471,-0.469932)--(-1.825482,-0.579795)--cycle;
\path[fill=blue!32!white,draw=blue!32!white,line width=.12pt] (1.825482,-0.579795)--(1.726471,-0.469932)--(1.734824,-0.384910)--(1.834315,-0.489898)--cycle;
\path[fill=blue!24!white,draw=blue!24!white,line width=.12pt] (-3.184591,0.156827)--(-3.169293,0.020912)--(-3.184628,-0.135163)--(-3.200000,0.000000)--cycle;
\path[fill=blue!33!white,draw=blue!33!white,line width=.12pt] (3.200000,0.000000)--(3.184628,-0.135163)--(3.169293,0.020912)--(3.184591,0.156827)--cycle;
\path[fill=blue!14!white,draw=blue!14!white,line width=.12pt] (-2.446009,1.050502)--(-2.589621,1.032859)--(-2.654149,0.904054)--(-2.506958,0.928840)--cycle;
\path[fill=blue!17!white,draw=blue!17!white,line width=.12pt] (2.506958,0.928840)--(2.654149,0.904054)--(2.589621,1.032859)--(2.446009,1.050502)--cycle;
\path[fill=blue!23!white,draw=blue!23!white,line width=.12pt] (-1.589379,0.506197)--(-1.660123,0.636707)--(-1.701490,0.554134)--(-1.628983,0.427143)--cycle;
\path[fill=blue!16!white,draw=blue!16!white,line width=.12pt] (1.628983,0.427143)--(1.701490,0.554134)--(1.660123,0.636707)--(1.589379,0.506197)--cycle;
\path[fill=blue!29!white,draw=blue!29!white,line width=.12pt] (-1.569256,0.156072)--(-1.584333,0.292734)--(-1.607593,0.214330)--(-1.592296,0.078414)--cycle;
\path[fill=blue!20!white,draw=blue!20!white,line width=.12pt] (1.592296,0.078414)--(1.607593,0.214330)--(1.584333,0.292734)--(1.569256,0.156072)--cycle;
\path[fill=blue!16!white,draw=blue!16!white,line width=.12pt] (-2.908701,0.779186)--(-3.006279,0.683903)--(-3.050416,0.535130)--(-2.951405,0.635242)--cycle;
\path[fill=blue!23!white,draw=blue!23!white,line width=.12pt] (2.951405,0.635242)--(3.050416,0.535130)--(3.006279,0.683903)--(2.908701,0.779186)--cycle;
\path[fill=blue!14!white,draw=blue!14!white,line width=.12pt] (-2.296657,1.041162)--(-2.446009,1.050502)--(-2.506958,0.928840)--(-2.353885,0.926929)--cycle;
\path[fill=blue!15!white,draw=blue!15!white,line width=.12pt] (2.353885,0.926929)--(2.506958,0.928840)--(2.446009,1.050502)--(2.296657,1.041162)--cycle;
\path[fill=blue!31!white,draw=blue!31!white,line width=.12pt] (-3.065176,-0.384910)--(-2.965685,-0.489898)--(-2.951405,-0.635242)--(-3.050416,-0.535130)--cycle;
\path[fill=blue!38!white,draw=blue!38!white,line width=.12pt] (3.050416,-0.535130)--(2.951405,-0.635242)--(2.965685,-0.489898)--(3.065176,-0.384910)--cycle;
\path[fill=blue!21!white,draw=blue!21!white,line width=.12pt] (-1.660123,0.636707)--(-1.755330,0.756135)--(-1.799069,0.668826)--(-1.701490,0.554134)--cycle;
\path[fill=blue!14!white,draw=blue!14!white,line width=.12pt] (1.701490,0.554134)--(1.799069,0.668826)--(1.755330,0.756135)--(1.660123,0.636707)--cycle;
\path[fill=blue!37!white,draw=blue!37!white,line width=.12pt] (-1.734824,-0.384910)--(-1.660896,-0.265131)--(-1.652899,-0.346529)--(-1.726471,-0.469932)--cycle;
\path[fill=blue!29!white,draw=blue!29!white,line width=.12pt] (1.726471,-0.469932)--(1.652899,-0.346529)--(1.660896,-0.265131)--(1.734824,-0.384910)--cycle;
\path[fill=blue!32!white,draw=blue!32!white,line width=.12pt] (-1.607593,-0.055995)--(-1.592296,0.078414)--(-1.600000,-0.000000)--(-1.615372,-0.135163)--cycle;
\path[fill=blue!23!white,draw=blue!23!white,line width=.12pt] (1.615372,-0.135163)--(1.600000,-0.000000)--(1.592296,0.078414)--(1.607593,-0.055995)--cycle;
\path[fill=blue!41!white,draw=blue!41!white,line width=.12pt] (-2.233123,-0.789480)--(-2.083771,-0.742699)--(-2.053620,-0.844328)--(-2.200811,-0.898392)--cycle;
\path[fill=blue!38!white,draw=blue!38!white,line width=.12pt] (2.200811,-0.898392)--(2.053620,-0.844328)--(2.083771,-0.742699)--(2.233123,-0.789480)--cycle;
\path[fill=blue!40!white,draw=blue!40!white,line width=.12pt] (-2.388443,-0.810441)--(-2.233123,-0.789480)--(-2.200811,-0.898392)--(-2.353885,-0.926929)--cycle;
\path[fill=blue!39!white,draw=blue!39!white,line width=.12pt] (2.353885,-0.926929)--(2.200811,-0.898392)--(2.233123,-0.789480)--(2.388443,-0.810441)--cycle;
\path[fill=blue!22!white,draw=blue!22!white,line width=.12pt] (-3.169293,0.291237)--(-3.184591,0.156827)--(-3.200000,0.000000)--(-3.184628,0.135163)--cycle;
\path[fill=blue!31!white,draw=blue!31!white,line width=.12pt] (3.184628,0.135163)--(3.200000,0.000000)--(3.184591,0.156827)--(3.169293,0.291237)--cycle;
\path[fill=blue!41!white,draw=blue!41!white,line width=.12pt] (-2.083771,-0.742699)--(-1.946127,-0.671897)--(-1.917969,-0.766813)--(-2.053620,-0.844328)--cycle;
\path[fill=blue!36!white,draw=blue!36!white,line width=.12pt] (2.053620,-0.844328)--(1.917969,-0.766813)--(1.946127,-0.671897)--(2.083771,-0.742699)--cycle;
\path[fill=blue!27!white,draw=blue!27!white,line width=.12pt] (-1.584333,0.292734)--(-1.628983,0.427143)--(-1.652899,0.346529)--(-1.607593,0.214330)--cycle;
\path[fill=blue!18!white,draw=blue!18!white,line width=.12pt] (1.607593,0.214330)--(1.652899,0.346529)--(1.628983,0.427143)--(1.584333,0.292734)--cycle;
\path[fill=blue!15!white,draw=blue!15!white,line width=.12pt] (-2.147305,1.005197)--(-2.296657,1.041162)--(-2.353885,0.926929)--(-2.200811,0.898392)--cycle;
\path[fill=blue!14!white,draw=blue!14!white,line width=.12pt] (2.200811,0.898392)--(2.353885,0.926929)--(2.296657,1.041162)--(2.147305,1.005197)--cycle;
\path[fill=blue!19!white,draw=blue!19!white,line width=.12pt] (-1.755330,0.756135)--(-1.871339,0.859891)--(-1.917969,0.766813)--(-1.799069,0.668826)--cycle;
\path[fill=blue!13!white,draw=blue!13!white,line width=.12pt] (1.799069,0.668826)--(1.917969,0.766813)--(1.871339,0.859891)--(1.755330,0.756135)--cycle;
\path[fill=blue!39!white,draw=blue!39!white,line width=.12pt] (-2.543764,-0.804777)--(-2.388443,-0.810441)--(-2.353885,-0.926929)--(-2.506958,-0.928840)--cycle;
\path[fill=blue!40!white,draw=blue!40!white,line width=.12pt] (2.506958,-0.928840)--(2.353885,-0.926929)--(2.388443,-0.810441)--(2.543764,-0.804777)--cycle;
\path[fill=blue!16!white,draw=blue!16!white,line width=.12pt] (-2.003693,0.943989)--(-2.147305,1.005197)--(-2.200811,0.898392)--(-2.053620,0.844328)--cycle;
\path[fill=blue!13!white,draw=blue!13!white,line width=.12pt] (2.053620,0.844328)--(2.200811,0.898392)--(2.147305,1.005197)--(2.003693,0.943989)--cycle;
\path[fill=blue!17!white,draw=blue!17!white,line width=.12pt] (-1.871339,0.859891)--(-2.003693,0.943989)--(-2.053620,0.844328)--(-1.917969,0.766813)--cycle;
\path[fill=blue!13!white,draw=blue!13!white,line width=.12pt] (1.917969,0.766813)--(2.053620,0.844328)--(2.003693,0.943989)--(1.871339,0.859891)--cycle;
\path[fill=blue!40!white,draw=blue!40!white,line width=.12pt] (-1.946127,-0.671897)--(-1.825482,-0.579795)--(-1.799069,-0.668826)--(-1.917969,-0.766813)--cycle;
\path[fill=blue!35!white,draw=blue!35!white,line width=.12pt] (1.917969,-0.766813)--(1.799069,-0.668826)--(1.825482,-0.579795)--(1.946127,-0.671897)--cycle;
\path[fill=blue!15!white,draw=blue!15!white,line width=.12pt] (-2.789801,0.853522)--(-2.908701,0.779186)--(-2.951405,0.635242)--(-2.830759,0.715462)--cycle;
\path[fill=blue!21!white,draw=blue!21!white,line width=.12pt] (2.830759,0.715462)--(2.951405,0.635242)--(2.908701,0.779186)--(2.789801,0.853522)--cycle;
\path[fill=blue!29!white,draw=blue!29!white,line width=.12pt] (-3.139104,-0.265131)--(-3.065176,-0.384910)--(-3.050416,-0.535130)--(-3.123988,-0.418974)--cycle;
\path[fill=blue!37!white,draw=blue!37!white,line width=.12pt] (3.123988,-0.418974)--(3.050416,-0.535130)--(3.065176,-0.384910)--(3.139104,-0.265131)--cycle;
\path[fill=blue!37!white,draw=blue!37!white,line width=.12pt] (-2.693116,-0.772707)--(-2.543764,-0.804777)--(-2.506958,-0.928840)--(-2.654149,-0.904054)--cycle;
\path[fill=blue!40!white,draw=blue!40!white,line width=.12pt] (2.654149,-0.904054)--(2.506958,-0.928840)--(2.543764,-0.804777)--(2.693116,-0.772707)--cycle;
\path[fill=blue!35!white,draw=blue!35!white,line width=.12pt] (-1.660896,-0.265131)--(-1.615372,-0.135163)--(-1.607593,-0.214330)--(-1.652899,-0.346529)--cycle;
\path[fill=blue!27!white,draw=blue!27!white,line width=.12pt] (1.652899,-0.346529)--(1.607593,-0.214330)--(1.615372,-0.135163)--(1.660896,-0.265131)--cycle;
\path[fill=blue!30!white,draw=blue!30!white,line width=.12pt] (-1.592296,0.078414)--(-1.607593,0.214330)--(-1.615372,0.135163)--(-1.600000,0.000000)--cycle;
\path[fill=blue!21!white,draw=blue!21!white,line width=.12pt] (1.600000,0.000000)--(1.615372,0.135163)--(1.607593,0.214330)--(1.592296,0.078414)--cycle;
\path[fill=blue!24!white,draw=blue!24!white,line width=.12pt] (-1.628983,0.427143)--(-1.701490,0.554134)--(-1.726471,0.469932)--(-1.652899,0.346529)--cycle;
\path[fill=blue!16!white,draw=blue!16!white,line width=.12pt] (1.652899,0.346529)--(1.726471,0.469932)--(1.701490,0.554134)--(1.628983,0.427143)--cycle;
\path[fill=blue!39!white,draw=blue!39!white,line width=.12pt] (-1.825482,-0.579795)--(-1.726471,-0.469932)--(-1.701490,-0.554134)--(-1.799069,-0.668826)--cycle;
\path[fill=blue!32!white,draw=blue!32!white,line width=.12pt] (1.799069,-0.668826)--(1.701490,-0.554134)--(1.726471,-0.469932)--(1.825482,-0.579795)--cycle;
\path[fill=blue!19!white,draw=blue!19!white,line width=.12pt] (-3.123988,0.418974)--(-3.169293,0.291237)--(-3.184628,0.135163)--(-3.139104,0.265131)--cycle;
\path[fill=blue!28!white,draw=blue!28!white,line width=.12pt] (3.139104,0.265131)--(3.184628,0.135163)--(3.169293,0.291237)--(3.123988,0.418974)--cycle;
\path[fill=blue!14!white,draw=blue!14!white,line width=.12pt] (-2.654149,0.904054)--(-2.789801,0.853522)--(-2.830759,0.715462)--(-2.693116,0.772707)--cycle;
\path[fill=blue!18!white,draw=blue!18!white,line width=.12pt] (2.693116,0.772707)--(2.830759,0.715462)--(2.789801,0.853522)--(2.654149,0.904054)--cycle;
\path[fill=blue!36!white,draw=blue!36!white,line width=.12pt] (-2.830759,-0.715462)--(-2.693116,-0.772707)--(-2.654149,-0.904054)--(-2.789801,-0.853522)--cycle;
\path[fill=blue!40!white,draw=blue!40!white,line width=.12pt] (2.789801,-0.853522)--(2.654149,-0.904054)--(2.693116,-0.772707)--(2.830759,-0.715462)--cycle;
\path[fill=blue!26!white,draw=blue!26!white,line width=.12pt] (-3.184628,-0.135163)--(-3.139104,-0.265131)--(-3.123988,-0.418974)--(-3.169293,-0.291237)--cycle;
\path[fill=blue!35!white,draw=blue!35!white,line width=.12pt] (3.169293,-0.291237)--(3.123988,-0.418974)--(3.139104,-0.265131)--(3.184628,-0.135163)--cycle;
\path[fill=blue!22!white,draw=blue!22!white,line width=.12pt] (-1.701490,0.554134)--(-1.799069,0.668826)--(-1.825482,0.579795)--(-1.726471,0.469932)--cycle;
\path[fill=blue!15!white,draw=blue!15!white,line width=.12pt] (1.726471,0.469932)--(1.825482,0.579795)--(1.799069,0.668826)--(1.701490,0.554134)--cycle;
\path[fill=blue!38!white,draw=blue!38!white,line width=.12pt] (-1.726471,-0.469932)--(-1.652899,-0.346529)--(-1.628983,-0.427143)--(-1.701490,-0.554134)--cycle;
\path[fill=blue!30!white,draw=blue!30!white,line width=.12pt] (1.701490,-0.554134)--(1.628983,-0.427143)--(1.652899,-0.346529)--(1.726471,-0.469932)--cycle;
\path[fill=blue!33!white,draw=blue!33!white,line width=.12pt] (-1.615372,-0.135163)--(-1.600000,-0.000000)--(-1.592296,-0.078414)--(-1.607593,-0.214330)--cycle;
\path[fill=blue!24!white,draw=blue!24!white,line width=.12pt] (1.607593,-0.214330)--(1.592296,-0.078414)--(1.600000,-0.000000)--(1.615372,-0.135163)--cycle;
\path[fill=blue!27!white,draw=blue!27!white,line width=.12pt] (-1.607593,0.214330)--(-1.652899,0.346529)--(-1.660896,0.265131)--(-1.615372,0.135163)--cycle;
\path[fill=blue!19!white,draw=blue!19!white,line width=.12pt] (1.615372,0.135163)--(1.660896,0.265131)--(1.652899,0.346529)--(1.607593,0.214330)--cycle;
\path[fill=blue!14!white,draw=blue!14!white,line width=.12pt] (-2.506958,0.928840)--(-2.654149,0.904054)--(-2.693116,0.772707)--(-2.543764,0.804777)--cycle;
\path[fill=blue!17!white,draw=blue!17!white,line width=.12pt] (2.543764,0.804777)--(2.693116,0.772707)--(2.654149,0.904054)--(2.506958,0.928840)--cycle;
\path[fill=blue!17!white,draw=blue!17!white,line width=.12pt] (-3.050416,0.535130)--(-3.123988,0.418974)--(-3.139104,0.265131)--(-3.065176,0.384910)--cycle;
\path[fill=blue!25!white,draw=blue!25!white,line width=.12pt] (3.065176,0.384910)--(3.139104,0.265131)--(3.123988,0.418974)--(3.050416,0.535130)--cycle;
\path[fill=blue!33!white,draw=blue!33!white,line width=.12pt] (-2.951405,-0.635242)--(-2.830759,-0.715462)--(-2.789801,-0.853522)--(-2.908701,-0.779186)--cycle;
\path[fill=blue!39!white,draw=blue!39!white,line width=.12pt] (2.908701,-0.779186)--(2.789801,-0.853522)--(2.830759,-0.715462)--(2.951405,-0.635242)--cycle;
\path[fill=blue!19!white,draw=blue!19!white,line width=.12pt] (-1.799069,0.668826)--(-1.917969,0.766813)--(-1.946127,0.671897)--(-1.825482,0.579795)--cycle;
\path[fill=blue!14!white,draw=blue!14!white,line width=.12pt] (1.825482,0.579795)--(1.946127,0.671897)--(1.917969,0.766813)--(1.799069,0.668826)--cycle;
\path[fill=blue!41!white,draw=blue!41!white,line width=.12pt] (-2.053620,-0.844328)--(-1.917969,-0.766813)--(-1.871339,-0.859891)--(-2.003693,-0.943989)--cycle;
\path[fill=blue!37!white,draw=blue!37!white,line width=.12pt] (2.003693,-0.943989)--(1.871339,-0.859891)--(1.917969,-0.766813)--(2.053620,-0.844328)--cycle;
\path[fill=blue!41!white,draw=blue!41!white,line width=.12pt] (-2.200811,-0.898392)--(-2.053620,-0.844328)--(-2.003693,-0.943989)--(-2.147305,-1.005197)--cycle;
\path[fill=blue!38!white,draw=blue!38!white,line width=.12pt] (2.147305,-1.005197)--(2.003693,-0.943989)--(2.053620,-0.844328)--(2.200811,-0.898392)--cycle;
\path[fill=blue!14!white,draw=blue!14!white,line width=.12pt] (-2.353885,0.926929)--(-2.506958,0.928840)--(-2.543764,0.804777)--(-2.388443,0.810441)--cycle;
\path[fill=blue!15!white,draw=blue!15!white,line width=.12pt] (2.388443,0.810441)--(2.543764,0.804777)--(2.506958,0.928840)--(2.353885,0.926929)--cycle;
\path[fill=blue!41!white,draw=blue!41!white,line width=.12pt] (-1.917969,-0.766813)--(-1.799069,-0.668826)--(-1.755330,-0.756135)--(-1.871339,-0.859891)--cycle;
\path[fill=blue!35!white,draw=blue!35!white,line width=.12pt] (1.871339,-0.859891)--(1.755330,-0.756135)--(1.799069,-0.668826)--(1.917969,-0.766813)--cycle;
\path[fill=blue!40!white,draw=blue!40!white,line width=.12pt] (-2.353885,-0.926929)--(-2.200811,-0.898392)--(-2.147305,-1.005197)--(-2.296657,-1.041162)--cycle;
\path[fill=blue!39!white,draw=blue!39!white,line width=.12pt] (2.296657,-1.041162)--(2.147305,-1.005197)--(2.200811,-0.898392)--(2.353885,-0.926929)--cycle;
\path[fill=blue!36!white,draw=blue!36!white,line width=.12pt] (-1.652899,-0.346529)--(-1.607593,-0.214330)--(-1.584333,-0.292734)--(-1.628983,-0.427143)--cycle;
\path[fill=blue!27!white,draw=blue!27!white,line width=.12pt] (1.628983,-0.427143)--(1.584333,-0.292734)--(1.607593,-0.214330)--(1.652899,-0.346529)--cycle;
\path[fill=blue!18!white,draw=blue!18!white,line width=.12pt] (-1.917969,0.766813)--(-2.053620,0.844328)--(-2.083771,0.742699)--(-1.946127,0.671897)--cycle;
\path[fill=blue!13!white,draw=blue!13!white,line width=.12pt] (1.946127,0.671897)--(2.083771,0.742699)--(2.053620,0.844328)--(1.917969,0.766813)--cycle;
\path[fill=blue!23!white,draw=blue!23!white,line width=.12pt] (-3.200000,0.000000)--(-3.184628,-0.135163)--(-3.169293,-0.291237)--(-3.184591,-0.156827)--cycle;
\path[fill=blue!32!white,draw=blue!32!white,line width=.12pt] (3.184591,-0.156827)--(3.169293,-0.291237)--(3.184628,-0.135163)--(3.200000,-0.000000)--cycle;
\path[fill=blue!15!white,draw=blue!15!white,line width=.12pt] (-2.200811,0.898392)--(-2.353885,0.926929)--(-2.388443,0.810441)--(-2.233123,0.789480)--cycle;
\path[fill=blue!14!white,draw=blue!14!white,line width=.12pt] (2.233123,0.789480)--(2.388443,0.810441)--(2.353885,0.926929)--(2.200811,0.898392)--cycle;
\path[fill=blue!16!white,draw=blue!16!white,line width=.12pt] (-2.053620,0.844328)--(-2.200811,0.898392)--(-2.233123,0.789480)--(-2.083771,0.742699)--cycle;
\path[fill=blue!13!white,draw=blue!13!white,line width=.12pt] (2.083771,0.742699)--(2.233123,0.789480)--(2.200811,0.898392)--(2.053620,0.844328)--cycle;
\path[fill=blue!31!white,draw=blue!31!white,line width=.12pt] (-1.600000,0.000000)--(-1.615372,0.135163)--(-1.607593,0.055995)--(-1.592296,-0.078414)--cycle;
\path[fill=blue!22!white,draw=blue!22!white,line width=.12pt] (1.592296,-0.078414)--(1.607593,0.055995)--(1.615372,0.135163)--(1.600000,-0.000000)--cycle;
\path[fill=blue!25!white,draw=blue!25!white,line width=.12pt] (-1.652899,0.346529)--(-1.726471,0.469932)--(-1.734824,0.384910)--(-1.660896,0.265131)--cycle;
\path[fill=blue!17!white,draw=blue!17!white,line width=.12pt] (1.660896,0.265131)--(1.734824,0.384910)--(1.726471,0.469932)--(1.652899,0.346529)--cycle;
\path[fill=blue!40!white,draw=blue!40!white,line width=.12pt] (-1.799069,-0.668826)--(-1.701490,-0.554134)--(-1.660123,-0.636707)--(-1.755330,-0.756135)--cycle;
\path[fill=blue!33!white,draw=blue!33!white,line width=.12pt] (1.755330,-0.756135)--(1.660123,-0.636707)--(1.701490,-0.554134)--(1.799069,-0.668826)--cycle;
\path[fill=blue!16!white,draw=blue!16!white,line width=.12pt] (-2.951405,0.635242)--(-3.050416,0.535130)--(-3.065176,0.384910)--(-2.965685,0.489898)--cycle;
\path[fill=blue!23!white,draw=blue!23!white,line width=.12pt] (2.965685,0.489898)--(3.065176,0.384910)--(3.050416,0.535130)--(2.951405,0.635242)--cycle;
\path[fill=blue!39!white,draw=blue!39!white,line width=.12pt] (-2.506958,-0.928840)--(-2.353885,-0.926929)--(-2.296657,-1.041162)--(-2.446009,-1.050502)--cycle;
\path[fill=blue!40!white,draw=blue!40!white,line width=.12pt] (2.446009,-1.050502)--(2.296657,-1.041162)--(2.353885,-0.926929)--(2.506958,-0.928840)--cycle;
\path[fill=blue!31!white,draw=blue!31!white,line width=.12pt] (-3.050416,-0.535130)--(-2.951405,-0.635242)--(-2.908701,-0.779186)--(-3.006279,-0.683903)--cycle;
\path[fill=blue!38!white,draw=blue!38!white,line width=.12pt] (3.006279,-0.683903)--(2.908701,-0.779186)--(2.951405,-0.635242)--(3.050416,-0.535130)--cycle;
\path[fill=blue!34!white,draw=blue!34!white,line width=.12pt] (-1.607593,-0.214330)--(-1.592296,-0.078414)--(-1.569256,-0.156072)--(-1.584333,-0.292734)--cycle;
\path[fill=blue!25!white,draw=blue!25!white,line width=.12pt] (1.584333,-0.292734)--(1.569256,-0.156072)--(1.592296,-0.078414)--(1.607593,-0.214330)--cycle;
\path[fill=blue!38!white,draw=blue!38!white,line width=.12pt] (-1.701490,-0.554134)--(-1.628983,-0.427143)--(-1.589379,-0.506197)--(-1.660123,-0.636707)--cycle;
\path[fill=blue!31!white,draw=blue!31!white,line width=.12pt] (1.660123,-0.636707)--(1.589379,-0.506197)--(1.628983,-0.427143)--(1.701490,-0.554134)--cycle;
\path[fill=blue!37!white,draw=blue!37!white,line width=.12pt] (-2.654149,-0.904054)--(-2.506958,-0.928840)--(-2.446009,-1.050502)--(-2.589621,-1.032859)--cycle;
\path[fill=blue!40!white,draw=blue!40!white,line width=.12pt] (2.589621,-1.032859)--(2.446009,-1.050502)--(2.506958,-0.928840)--(2.654149,-0.904054)--cycle;
\path[fill=blue!21!white,draw=blue!21!white,line width=.12pt] (-3.184628,0.135163)--(-3.200000,0.000000)--(-3.184591,-0.156827)--(-3.169293,-0.020912)--cycle;
\path[fill=blue!30!white,draw=blue!30!white,line width=.12pt] (3.169293,-0.020912)--(3.184591,-0.156827)--(3.200000,-0.000000)--(3.184628,0.135163)--cycle;
\path[fill=blue!22!white,draw=blue!22!white,line width=.12pt] (-1.726471,0.469932)--(-1.825482,0.579795)--(-1.834315,0.489898)--(-1.734824,0.384910)--cycle;
\path[fill=blue!15!white,draw=blue!15!white,line width=.12pt] (1.734824,0.384910)--(1.834315,0.489898)--(1.825482,0.579795)--(1.726471,0.469932)--cycle;
\path[fill=blue!15!white,draw=blue!15!white,line width=.12pt] (-2.830759,0.715462)--(-2.951405,0.635242)--(-2.965685,0.489898)--(-2.844456,0.576059)--cycle;
\path[fill=blue!20!white,draw=blue!20!white,line width=.12pt] (2.844456,0.576059)--(2.965685,0.489898)--(2.951405,0.635242)--(2.830759,0.715462)--cycle;
\path[fill=blue!28!white,draw=blue!28!white,line width=.12pt] (-1.615372,0.135163)--(-1.660896,0.265131)--(-1.652899,0.183733)--(-1.607593,0.055995)--cycle;
\path[fill=blue!19!white,draw=blue!19!white,line width=.12pt] (1.607593,0.055995)--(1.652899,0.183733)--(1.660896,0.265131)--(1.615372,0.135163)--cycle;
\path[fill=blue!28!white,draw=blue!28!white,line width=.12pt] (-3.123988,-0.418974)--(-3.050416,-0.535130)--(-3.006279,-0.683903)--(-3.078787,-0.571335)--cycle;
\path[fill=blue!36!white,draw=blue!36!white,line width=.12pt] (3.078787,-0.571335)--(3.006279,-0.683903)--(3.050416,-0.535130)--(3.123988,-0.418974)--cycle;
\path[fill=blue!20!white,draw=blue!20!white,line width=.12pt] (-1.825482,0.579795)--(-1.946127,0.671897)--(-1.955544,0.576059)--(-1.834315,0.489898)--cycle;
\path[fill=blue!14!white,draw=blue!14!white,line width=.12pt] (1.834315,0.489898)--(1.955544,0.576059)--(1.946127,0.671897)--(1.825482,0.579795)--cycle;
\path[fill=blue!36!white,draw=blue!36!white,line width=.12pt] (-1.628983,-0.427143)--(-1.584333,-0.292734)--(-1.545814,-0.369621)--(-1.589379,-0.506197)--cycle;
\path[fill=blue!28!white,draw=blue!28!white,line width=.12pt] (1.589379,-0.506197)--(1.545814,-0.369621)--(1.584333,-0.292734)--(1.628983,-0.427143)--cycle;
\path[fill=blue!35!white,draw=blue!35!white,line width=.12pt] (-2.789801,-0.853522)--(-2.654149,-0.904054)--(-2.589621,-1.032859)--(-2.721975,-0.988910)--cycle;
\path[fill=blue!39!white,draw=blue!39!white,line width=.12pt] (2.721975,-0.988910)--(2.589621,-1.032859)--(2.654149,-0.904054)--(2.789801,-0.853522)--cycle;
\path[fill=blue!14!white,draw=blue!14!white,line width=.12pt] (-2.693116,0.772707)--(-2.830759,0.715462)--(-2.844456,0.576059)--(-2.706147,0.640083)--cycle;
\path[fill=blue!18!white,draw=blue!18!white,line width=.12pt] (2.706147,0.640083)--(2.844456,0.576059)--(2.830759,0.715462)--(2.693116,0.772707)--cycle;
\path[fill=blue!41!white,draw=blue!41!white,line width=.12pt] (-2.003693,-0.943989)--(-1.871339,-0.859891)--(-1.806687,-0.950236)--(-1.934468,-1.040724)--cycle;
\path[fill=blue!37!white,draw=blue!37!white,line width=.12pt] (1.934468,-1.040724)--(1.806687,-0.950236)--(1.871339,-0.859891)--(2.003693,-0.943989)--cycle;
\path[fill=blue!31!white,draw=blue!31!white,line width=.12pt] (-1.592296,-0.078414)--(-1.607593,0.055995)--(-1.584333,-0.022409)--(-1.569256,-0.156072)--cycle;
\path[fill=blue!22!white,draw=blue!22!white,line width=.12pt] (1.569256,-0.156072)--(1.584333,-0.022409)--(1.607593,0.055995)--(1.592296,-0.078414)--cycle;
\path[fill=blue!41!white,draw=blue!41!white,line width=.12pt] (-1.871339,-0.859891)--(-1.755330,-0.756135)--(-1.694686,-0.840879)--(-1.806687,-0.950236)--cycle;
\path[fill=blue!36!white,draw=blue!36!white,line width=.12pt] (1.806687,-0.950236)--(1.694686,-0.840879)--(1.755330,-0.756135)--(1.871339,-0.859891)--cycle;
\path[fill=blue!19!white,draw=blue!19!white,line width=.12pt] (-3.139104,0.265131)--(-3.184628,0.135163)--(-3.169293,-0.020912)--(-3.123988,0.111288)--cycle;
\path[fill=blue!27!white,draw=blue!27!white,line width=.12pt] (3.123988,0.111288)--(3.169293,-0.020912)--(3.184628,0.135163)--(3.139104,0.265131)--cycle;
\path[fill=blue!41!white,draw=blue!41!white,line width=.12pt] (-2.147305,-1.005197)--(-2.003693,-0.943989)--(-1.934468,-1.040724)--(-2.073119,-1.108865)--cycle;
\path[fill=blue!38!white,draw=blue!38!white,line width=.12pt] (2.073119,-1.108865)--(1.934468,-1.040724)--(2.003693,-0.943989)--(2.147305,-1.005197)--cycle;
\path[fill=blue!25!white,draw=blue!25!white,line width=.12pt] (-1.660896,0.265131)--(-1.734824,0.384910)--(-1.726471,0.299889)--(-1.652899,0.183733)--cycle;
\path[fill=blue!17!white,draw=blue!17!white,line width=.12pt] (1.652899,0.183733)--(1.726471,0.299889)--(1.734824,0.384910)--(1.660896,0.265131)--cycle;
\path[fill=blue!18!white,draw=blue!18!white,line width=.12pt] (-1.946127,0.671897)--(-2.083771,0.742699)--(-2.093853,0.640083)--(-1.955544,0.576059)--cycle;
\path[fill=blue!14!white,draw=blue!14!white,line width=.12pt] (1.955544,0.576059)--(2.093853,0.640083)--(2.083771,0.742699)--(1.946127,0.671897)--cycle;
\path[fill=blue!40!white,draw=blue!40!white,line width=.12pt] (-1.755330,-0.756135)--(-1.660123,-0.636707)--(-1.602769,-0.716855)--(-1.694686,-0.840879)--cycle;
\path[fill=blue!33!white,draw=blue!33!white,line width=.12pt] (1.694686,-0.840879)--(1.602769,-0.716855)--(1.660123,-0.636707)--(1.755330,-0.756135)--cycle;
\path[fill=blue!14!white,draw=blue!14!white,line width=.12pt] (-2.543764,0.804777)--(-2.693116,0.772707)--(-2.706147,0.640083)--(-2.556072,0.679508)--cycle;
\path[fill=blue!16!white,draw=blue!16!white,line width=.12pt] (2.556072,0.679508)--(2.706147,0.640083)--(2.693116,0.772707)--(2.543764,0.804777)--cycle;
\path[fill=blue!40!white,draw=blue!40!white,line width=.12pt] (-2.296657,-1.041162)--(-2.147305,-1.005197)--(-2.073119,-1.108865)--(-2.217311,-1.152040)--cycle;
\path[fill=blue!39!white,draw=blue!39!white,line width=.12pt] (2.217311,-1.152040)--(2.073119,-1.108865)--(2.147305,-1.005197)--(2.296657,-1.041162)--cycle;
\path[fill=blue!25!white,draw=blue!25!white,line width=.12pt] (-3.169293,-0.291237)--(-3.123988,-0.418974)--(-3.078787,-0.571335)--(-3.123436,-0.445808)--cycle;
\path[fill=blue!34!white,draw=blue!34!white,line width=.12pt] (3.123436,-0.445808)--(3.078787,-0.571335)--(3.123988,-0.418974)--(3.169293,-0.291237)--cycle;
\path[fill=blue!16!white,draw=blue!16!white,line width=.12pt] (-2.083771,0.742699)--(-2.233123,0.789480)--(-2.243928,0.679508)--(-2.093853,0.640083)--cycle;
\path[fill=blue!14!white,draw=blue!14!white,line width=.12pt] (2.093853,0.640083)--(2.243928,0.679508)--(2.233123,0.789480)--(2.083771,0.742699)--cycle;
\path[fill=blue!14!white,draw=blue!14!white,line width=.12pt] (-2.388443,0.810441)--(-2.543764,0.804777)--(-2.556072,0.679508)--(-2.400000,0.692820)--cycle;
\path[fill=blue!15!white,draw=blue!15!white,line width=.12pt] (2.400000,0.692820)--(2.556072,0.679508)--(2.543764,0.804777)--(2.388443,0.810441)--cycle;
\path[fill=blue!15!white,draw=blue!15!white,line width=.12pt] (-2.233123,0.789480)--(-2.388443,0.810441)--(-2.400000,0.692820)--(-2.243928,0.679508)--cycle;
\path[fill=blue!14!white,draw=blue!14!white,line width=.12pt] (2.243928,0.679508)--(2.400000,0.692820)--(2.388443,0.810441)--(2.233123,0.789480)--cycle;
\path[fill=blue!34!white,draw=blue!34!white,line width=.12pt] (-1.584333,-0.292734)--(-1.569256,-0.156072)--(-1.531105,-0.232228)--(-1.545814,-0.369621)--cycle;
\path[fill=blue!26!white,draw=blue!26!white,line width=.12pt] (1.545814,-0.369621)--(1.531105,-0.232228)--(1.569256,-0.156072)--(1.584333,-0.292734)--cycle;
\path[fill=blue!33!white,draw=blue!33!white,line width=.12pt] (-2.908701,-0.779186)--(-2.789801,-0.853522)--(-2.721975,-0.988910)--(-2.837984,-0.920344)--cycle;
\path[fill=blue!38!white,draw=blue!38!white,line width=.12pt] (2.837984,-0.920344)--(2.721975,-0.988910)--(2.789801,-0.853522)--(2.908701,-0.779186)--cycle;
\path[fill=blue!39!white,draw=blue!39!white,line width=.12pt] (-1.660123,-0.636707)--(-1.589379,-0.506197)--(-1.534468,-0.582930)--(-1.602769,-0.716855)--cycle;
\path[fill=blue!31!white,draw=blue!31!white,line width=.12pt] (1.602769,-0.716855)--(1.534468,-0.582930)--(1.589379,-0.506197)--(1.660123,-0.636707)--cycle;
\path[fill=blue!17!white,draw=blue!17!white,line width=.12pt] (-3.065176,0.384910)--(-3.139104,0.265131)--(-3.123988,0.111288)--(-3.050416,0.234690)--cycle;
\path[fill=blue!25!white,draw=blue!25!white,line width=.12pt] (3.050416,0.234690)--(3.123988,0.111288)--(3.139104,0.265131)--(3.065176,0.384910)--cycle;
\path[fill=blue!29!white,draw=blue!29!white,line width=.12pt] (-1.607593,0.055995)--(-1.652899,0.183733)--(-1.628983,0.103118)--(-1.584333,-0.022409)--cycle;
\path[fill=blue!20!white,draw=blue!20!white,line width=.12pt] (1.584333,-0.022409)--(1.628983,0.103118)--(1.652899,0.183733)--(1.607593,0.055995)--cycle;
\path[fill=blue!39!white,draw=blue!39!white,line width=.12pt] (-2.446009,-1.050502)--(-2.296657,-1.041162)--(-2.217311,-1.152040)--(-2.361503,-1.168591)--cycle;
\path[fill=blue!40!white,draw=blue!40!white,line width=.12pt] (2.361503,-1.168591)--(2.217311,-1.152040)--(2.296657,-1.041162)--(2.446009,-1.050502)--cycle;
\path[fill=blue!23!white,draw=blue!23!white,line width=.12pt] (-1.734824,0.384910)--(-1.834315,0.489898)--(-1.825482,0.400001)--(-1.726471,0.299889)--cycle;
\path[fill=blue!16!white,draw=blue!16!white,line width=.12pt] (1.726471,0.299889)--(1.825482,0.400001)--(1.834315,0.489898)--(1.734824,0.384910)--cycle;
\path[fill=blue!37!white,draw=blue!37!white,line width=.12pt] (-1.589379,-0.506197)--(-1.545814,-0.369621)--(-1.492409,-0.444251)--(-1.534468,-0.582930)--cycle;
\path[fill=blue!29!white,draw=blue!29!white,line width=.12pt] (1.534468,-0.582930)--(1.492409,-0.444251)--(1.545814,-0.369621)--(1.589379,-0.506197)--cycle;
\path[fill=blue!23!white,draw=blue!23!white,line width=.12pt] (-3.184591,-0.156827)--(-3.169293,-0.291237)--(-3.123436,-0.445808)--(-3.138513,-0.312145)--cycle;
\path[fill=blue!32!white,draw=blue!32!white,line width=.12pt] (3.138513,-0.312145)--(3.123436,-0.445808)--(3.169293,-0.291237)--(3.184591,-0.156827)--cycle;
\path[fill=blue!32!white,draw=blue!32!white,line width=.12pt] (-1.569256,-0.156072)--(-1.584333,-0.022409)--(-1.545814,-0.099296)--(-1.531105,-0.232228)--cycle;
\path[fill=blue!23!white,draw=blue!23!white,line width=.12pt] (1.531105,-0.232228)--(1.545814,-0.099296)--(1.584333,-0.022409)--(1.569256,-0.156072)--cycle;
\path[fill=blue!15!white,draw=blue!15!white,line width=.12pt] (-2.965685,0.489898)--(-3.065176,0.384910)--(-3.050416,0.234690)--(-2.951405,0.344554)--cycle;
\path[fill=blue!22!white,draw=blue!22!white,line width=.12pt] (2.951405,0.344554)--(3.050416,0.234690)--(3.065176,0.384910)--(2.965685,0.489898)--cycle;
\path[fill=blue!30!white,draw=blue!30!white,line width=.12pt] (-3.006279,-0.683903)--(-2.908701,-0.779186)--(-2.837984,-0.920344)--(-2.933190,-0.829797)--cycle;
\path[fill=blue!37!white,draw=blue!37!white,line width=.12pt] (2.933190,-0.829797)--(2.837984,-0.920344)--(2.908701,-0.779186)--(3.006279,-0.683903)--cycle;
\path[fill=blue!37!white,draw=blue!37!white,line width=.12pt] (-2.589621,-1.032859)--(-2.446009,-1.050502)--(-2.361503,-1.168591)--(-2.500154,-1.157881)--cycle;
\path[fill=blue!40!white,draw=blue!40!white,line width=.12pt] (2.500154,-1.157881)--(2.361503,-1.168591)--(2.446009,-1.050502)--(2.589621,-1.032859)--cycle;
\path[fill=blue!20!white,draw=blue!20!white,line width=.12pt] (-1.834315,0.489898)--(-1.955544,0.576059)--(-1.946127,0.480221)--(-1.825482,0.400001)--cycle;
\path[fill=blue!15!white,draw=blue!15!white,line width=.12pt] (1.825482,0.400001)--(1.946127,0.480221)--(1.955544,0.576059)--(1.834315,0.489898)--cycle;
\path[fill=blue!41!white,draw=blue!41!white,line width=.12pt] (-1.806687,-0.950236)--(-1.694686,-0.840879)--(-1.617721,-0.922243)--(-1.724636,-1.036978)--cycle;
\path[fill=blue!36!white,draw=blue!36!white,line width=.12pt] (1.724636,-1.036978)--(1.617721,-0.922243)--(1.694686,-0.840879)--(1.806687,-0.950236)--cycle;
\path[fill=blue!41!white,draw=blue!41!white,line width=.12pt] (-1.934468,-1.040724)--(-1.806687,-0.950236)--(-1.724636,-1.036978)--(-1.846614,-1.133600)--cycle;
\path[fill=blue!38!white,draw=blue!38!white,line width=.12pt] (1.846614,-1.133600)--(1.724636,-1.036978)--(1.806687,-0.950236)--(1.934468,-1.040724)--cycle;
\path[fill=blue!26!white,draw=blue!26!white,line width=.12pt] (-1.652899,0.183733)--(-1.726471,0.299889)--(-1.701490,0.215687)--(-1.628983,0.103118)--cycle;
\path[fill=blue!18!white,draw=blue!18!white,line width=.12pt] (1.628983,0.103118)--(1.701490,0.215687)--(1.726471,0.299889)--(1.652899,0.183733)--cycle;
\path[fill=blue!40!white,draw=blue!40!white,line width=.12pt] (-1.694686,-0.840879)--(-1.602769,-0.716855)--(-1.529978,-0.793806)--(-1.617721,-0.922243)--cycle;
\path[fill=blue!34!white,draw=blue!34!white,line width=.12pt] (1.617721,-0.922243)--(1.529978,-0.793806)--(1.602769,-0.716855)--(1.694686,-0.840879)--cycle;
\path[fill=blue!41!white,draw=blue!41!white,line width=.12pt] (-2.073119,-1.108865)--(-1.934468,-1.040724)--(-1.846614,-1.133600)--(-1.978968,-1.208398)--cycle;
\path[fill=blue!39!white,draw=blue!39!white,line width=.12pt] (1.978968,-1.208398)--(1.846614,-1.133600)--(1.934468,-1.040724)--(2.073119,-1.108865)--cycle;
\path[fill=blue!35!white,draw=blue!35!white,line width=.12pt] (-1.545814,-0.369621)--(-1.531105,-0.232228)--(-1.478207,-0.306147)--(-1.492409,-0.444251)--cycle;
\path[fill=blue!26!white,draw=blue!26!white,line width=.12pt] (1.492409,-0.444251)--(1.478207,-0.306147)--(1.531105,-0.232228)--(1.545814,-0.369621)--cycle;
\path[fill=blue!14!white,draw=blue!14!white,line width=.12pt] (-2.844456,0.576059)--(-2.965685,0.489898)--(-2.951405,0.344554)--(-2.830759,0.436656)--cycle;
\path[fill=blue!20!white,draw=blue!20!white,line width=.12pt] (2.830759,0.436656)--(2.951405,0.344554)--(2.965685,0.489898)--(2.844456,0.576059)--cycle;
\path[fill=blue!18!white,draw=blue!18!white,line width=.12pt] (-1.955544,0.576059)--(-2.093853,0.640083)--(-2.083771,0.537466)--(-1.946127,0.480221)--cycle;
\path[fill=blue!14!white,draw=blue!14!white,line width=.12pt] (1.946127,0.480221)--(2.083771,0.537466)--(2.093853,0.640083)--(1.955544,0.576059)--cycle;
\path[fill=blue!39!white,draw=blue!39!white,line width=.12pt] (-1.602769,-0.716855)--(-1.534468,-0.582930)--(-1.464780,-0.656601)--(-1.529978,-0.793806)--cycle;
\path[fill=blue!32!white,draw=blue!32!white,line width=.12pt] (1.529978,-0.793806)--(1.464780,-0.656601)--(1.534468,-0.582930)--(1.602769,-0.716855)--cycle;
\path[fill=blue!20!white,draw=blue!20!white,line width=.12pt] (-3.169293,-0.020912)--(-3.184591,-0.156827)--(-3.138513,-0.312145)--(-3.123436,-0.175483)--cycle;
\path[fill=blue!29!white,draw=blue!29!white,line width=.12pt] (3.123436,-0.175483)--(3.138513,-0.312145)--(3.184591,-0.156827)--(3.169293,-0.020912)--cycle;
\path[fill=blue!35!white,draw=blue!35!white,line width=.12pt] (-2.721975,-0.988910)--(-2.589621,-1.032859)--(-2.500154,-1.157881)--(-2.627935,-1.120322)--cycle;
\path[fill=blue!39!white,draw=blue!39!white,line width=.12pt] (2.627935,-1.120322)--(2.500154,-1.157881)--(2.589621,-1.032859)--(2.721975,-0.988910)--cycle;
\path[fill=blue!29!white,draw=blue!29!white,line width=.12pt] (-1.584333,-0.022409)--(-1.628983,0.103118)--(-1.589379,0.024064)--(-1.545814,-0.099296)--cycle;
\path[fill=blue!21!white,draw=blue!21!white,line width=.12pt] (1.545814,-0.099296)--(1.589379,0.024064)--(1.628983,0.103118)--(1.584333,-0.022409)--cycle;
\path[fill=blue!40!white,draw=blue!40!white,line width=.12pt] (-2.217311,-1.152040)--(-2.073119,-1.108865)--(-1.978968,-1.208398)--(-2.116611,-1.258496)--cycle;
\path[fill=blue!39!white,draw=blue!39!white,line width=.12pt] (2.116611,-1.258496)--(1.978968,-1.208398)--(2.073119,-1.108865)--(2.217311,-1.152040)--cycle;
\path[fill=blue!28!white,draw=blue!28!white,line width=.12pt] (-3.078787,-0.571335)--(-3.006279,-0.683903)--(-2.933190,-0.829797)--(-3.003935,-0.720748)--cycle;
\path[fill=blue!35!white,draw=blue!35!white,line width=.12pt] (3.003935,-0.720748)--(2.933190,-0.829797)--(3.006279,-0.683903)--(3.078787,-0.571335)--cycle;
\path[fill=blue!14!white,draw=blue!14!white,line width=.12pt] (-2.706147,0.640083)--(-2.844456,0.576059)--(-2.830759,0.436656)--(-2.693116,0.507458)--cycle;
\path[fill=blue!18!white,draw=blue!18!white,line width=.12pt] (2.693116,0.507458)--(2.830759,0.436656)--(2.844456,0.576059)--(2.706147,0.640083)--cycle;
\path[fill=blue!16!white,draw=blue!16!white,line width=.12pt] (-2.093853,0.640083)--(-2.243928,0.679508)--(-2.233123,0.569536)--(-2.083771,0.537466)--cycle;
\path[fill=blue!14!white,draw=blue!14!white,line width=.12pt] (2.083771,0.537466)--(2.233123,0.569536)--(2.243928,0.679508)--(2.093853,0.640083)--cycle;
\path[fill=blue!23!white,draw=blue!23!white,line width=.12pt] (-1.726471,0.299889)--(-1.825482,0.400001)--(-1.799069,0.310969)--(-1.701490,0.215687)--cycle;
\path[fill=blue!16!white,draw=blue!16!white,line width=.12pt] (1.701490,0.215687)--(1.799069,0.310969)--(1.825482,0.400001)--(1.726471,0.299889)--cycle;
\path[fill=blue!37!white,draw=blue!37!white,line width=.12pt] (-1.534468,-0.582930)--(-1.492409,-0.444251)--(-1.424631,-0.515903)--(-1.464780,-0.656601)--cycle;
\path[fill=blue!30!white,draw=blue!30!white,line width=.12pt] (1.464780,-0.656601)--(1.424631,-0.515903)--(1.492409,-0.444251)--(1.534468,-0.582930)--cycle;
\path[fill=blue!14!white,draw=blue!14!white,line width=.12pt] (-2.556072,0.679508)--(-2.706147,0.640083)--(-2.693116,0.507458)--(-2.543764,0.554239)--cycle;
\path[fill=blue!16!white,draw=blue!16!white,line width=.12pt] (2.543764,0.554239)--(2.693116,0.507458)--(2.706147,0.640083)--(2.556072,0.679508)--cycle;
\path[fill=blue!15!white,draw=blue!15!white,line width=.12pt] (-2.243928,0.679508)--(-2.400000,0.692820)--(-2.388443,0.575200)--(-2.233123,0.569536)--cycle;
\path[fill=blue!14!white,draw=blue!14!white,line width=.12pt] (2.233123,0.569536)--(2.388443,0.575200)--(2.400000,0.692820)--(2.243928,0.679508)--cycle;
\path[fill=blue!14!white,draw=blue!14!white,line width=.12pt] (-2.400000,0.692820)--(-2.556072,0.679508)--(-2.543764,0.554239)--(-2.388443,0.575200)--cycle;
\path[fill=blue!15!white,draw=blue!15!white,line width=.12pt] (2.388443,0.575200)--(2.543764,0.554239)--(2.556072,0.679508)--(2.400000,0.692820)--cycle;
\path[fill=blue!32!white,draw=blue!32!white,line width=.12pt] (-1.531105,-0.232228)--(-1.545814,-0.099296)--(-1.492409,-0.173925)--(-1.478207,-0.306147)--cycle;
\path[fill=blue!24!white,draw=blue!24!white,line width=.12pt] (1.478207,-0.306147)--(1.492409,-0.173925)--(1.545814,-0.099296)--(1.531105,-0.232228)--cycle;
\path[fill=blue!18!white,draw=blue!18!white,line width=.12pt] (-3.123988,0.111288)--(-3.169293,-0.020912)--(-3.123436,-0.175483)--(-3.078787,-0.041074)--cycle;
\path[fill=blue!27!white,draw=blue!27!white,line width=.12pt] (3.078787,-0.041074)--(3.123436,-0.175483)--(3.169293,-0.020912)--(3.123988,0.111288)--cycle;
\path[fill=blue!39!white,draw=blue!39!white,line width=.12pt] (-2.361503,-1.168591)--(-2.217311,-1.152040)--(-2.116611,-1.258496)--(-2.254254,-1.281970)--cycle;
\path[fill=blue!40!white,draw=blue!40!white,line width=.12pt] (2.254254,-1.281970)--(2.116611,-1.258496)--(2.217311,-1.152040)--(2.361503,-1.168591)--cycle;
\path[fill=blue!41!white,draw=blue!41!white,line width=.12pt] (-1.724636,-1.036978)--(-1.617721,-0.922243)--(-1.525177,-0.999443)--(-1.625975,-1.119280)--cycle;
\path[fill=blue!37!white,draw=blue!37!white,line width=.12pt] (1.625975,-1.119280)--(1.525177,-0.999443)--(1.617721,-0.922243)--(1.724636,-1.036978)--cycle;
\path[fill=blue!33!white,draw=blue!33!white,line width=.12pt] (-2.837984,-0.920344)--(-2.721975,-0.988910)--(-2.627935,-1.120322)--(-2.739936,-1.057357)--cycle;
\path[fill=blue!38!white,draw=blue!38!white,line width=.12pt] (2.739936,-1.057357)--(2.627935,-1.120322)--(2.721975,-0.988910)--(2.837984,-0.920344)--cycle;
\path[fill=blue!41!white,draw=blue!41!white,line width=.12pt] (-1.617721,-0.922243)--(-1.529978,-0.793806)--(-1.442454,-0.866819)--(-1.525177,-0.999443)--cycle;
\path[fill=blue!35!white,draw=blue!35!white,line width=.12pt] (1.525177,-0.999443)--(1.442454,-0.866819)--(1.529978,-0.793806)--(1.617721,-0.922243)--cycle;
\path[fill=blue!26!white,draw=blue!26!white,line width=.12pt] (-1.628983,0.103118)--(-1.701490,0.215687)--(-1.660123,0.133114)--(-1.589379,0.024064)--cycle;
\path[fill=blue!19!white,draw=blue!19!white,line width=.12pt] (1.589379,0.024064)--(1.660123,0.133114)--(1.701490,0.215687)--(1.628983,0.103118)--cycle;
\path[fill=blue!21!white,draw=blue!21!white,line width=.12pt] (-1.825482,0.400001)--(-1.946127,0.480221)--(-1.917969,0.385305)--(-1.799069,0.310969)--cycle;
\path[fill=blue!15!white,draw=blue!15!white,line width=.12pt] (1.799069,0.310969)--(1.917969,0.385305)--(1.946127,0.480221)--(1.825482,0.400001)--cycle;
\path[fill=blue!25!white,draw=blue!25!white,line width=.12pt] (-3.123436,-0.445808)--(-3.078787,-0.571335)--(-3.003935,-0.720748)--(-3.047499,-0.597387)--cycle;
\path[fill=blue!33!white,draw=blue!33!white,line width=.12pt] (3.047499,-0.597387)--(3.003935,-0.720748)--(3.078787,-0.571335)--(3.123436,-0.445808)--cycle;
\path[fill=blue!42!white,draw=blue!42!white,line width=.12pt] (-1.846614,-1.133600)--(-1.724636,-1.036978)--(-1.625975,-1.119280)--(-1.740975,-1.221724)--cycle;
\path[fill=blue!38!white,draw=blue!38!white,line width=.12pt] (1.740975,-1.221724)--(1.625975,-1.119280)--(1.724636,-1.036978)--(1.846614,-1.133600)--cycle;
\path[fill=blue!35!white,draw=blue!35!white,line width=.12pt] (-1.492409,-0.444251)--(-1.478207,-0.306147)--(-1.411074,-0.377117)--(-1.424631,-0.515903)--cycle;
\path[fill=blue!27!white,draw=blue!27!white,line width=.12pt] (1.424631,-0.515903)--(1.411074,-0.377117)--(1.478207,-0.306147)--(1.492409,-0.444251)--cycle;
\path[fill=blue!39!white,draw=blue!39!white,line width=.12pt] (-1.529978,-0.793806)--(-1.464780,-0.656601)--(-1.380985,-0.726503)--(-1.442454,-0.866819)--cycle;
\path[fill=blue!33!white,draw=blue!33!white,line width=.12pt] (1.442454,-0.866819)--(1.380985,-0.726503)--(1.464780,-0.656601)--(1.529978,-0.793806)--cycle;
\path[fill=blue!16!white,draw=blue!16!white,line width=.12pt] (-3.050416,0.234690)--(-3.123988,0.111288)--(-3.078787,-0.041074)--(-3.006279,0.085917)--cycle;
\path[fill=blue!24!white,draw=blue!24!white,line width=.12pt] (3.006279,0.085917)--(3.078787,-0.041074)--(3.123988,0.111288)--(3.050416,0.234690)--cycle;
\path[fill=blue!41!white,draw=blue!41!white,line width=.12pt] (-1.978968,-1.208398)--(-1.846614,-1.133600)--(-1.740975,-1.221724)--(-1.865758,-1.302838)--cycle;
\path[fill=blue!39!white,draw=blue!39!white,line width=.12pt] (1.865758,-1.302838)--(1.740975,-1.221724)--(1.846614,-1.133600)--(1.978968,-1.208398)--cycle;
\path[fill=blue!37!white,draw=blue!37!white,line width=.12pt] (-2.500154,-1.157881)--(-2.361503,-1.168591)--(-2.254254,-1.281970)--(-2.386608,-1.277917)--cycle;
\path[fill=blue!39!white,draw=blue!39!white,line width=.12pt] (2.386608,-1.277917)--(2.254254,-1.281970)--(2.361503,-1.168591)--(2.500154,-1.157881)--cycle;
\path[fill=blue!30!white,draw=blue!30!white,line width=.12pt] (-1.545814,-0.099296)--(-1.589379,0.024064)--(-1.534468,-0.052668)--(-1.492409,-0.173925)--cycle;
\path[fill=blue!21!white,draw=blue!21!white,line width=.12pt] (1.492409,-0.173925)--(1.534468,-0.052668)--(1.589379,0.024064)--(1.545814,-0.099296)--cycle;
\path[fill=blue!18!white,draw=blue!18!white,line width=.12pt] (-1.946127,0.480221)--(-2.083771,0.537466)--(-2.053620,0.435837)--(-1.917969,0.385305)--cycle;
\path[fill=blue!14!white,draw=blue!14!white,line width=.12pt] (1.917969,0.385305)--(2.053620,0.435837)--(2.083771,0.537466)--(1.946127,0.480221)--cycle;
\path[fill=blue!38!white,draw=blue!38!white,line width=.12pt] (-1.464780,-0.656601)--(-1.424631,-0.515903)--(-1.343133,-0.583889)--(-1.380985,-0.726503)--cycle;
\path[fill=blue!30!white,draw=blue!30!white,line width=.12pt] (1.380985,-0.726503)--(1.343133,-0.583889)--(1.424631,-0.515903)--(1.464780,-0.656601)--cycle;
\path[fill=blue!30!white,draw=blue!30!white,line width=.12pt] (-2.933190,-0.829797)--(-2.837984,-0.920344)--(-2.739936,-1.057357)--(-2.831853,-0.971406)--cycle;
\path[fill=blue!37!white,draw=blue!37!white,line width=.12pt] (2.831853,-0.971406)--(2.739936,-1.057357)--(2.837984,-0.920344)--(2.933190,-0.829797)--cycle;
\path[fill=blue!24!white,draw=blue!24!white,line width=.12pt] (-1.701490,0.215687)--(-1.799069,0.310969)--(-1.755330,0.223661)--(-1.660123,0.133114)--cycle;
\path[fill=blue!17!white,draw=blue!17!white,line width=.12pt] (1.660123,0.133114)--(1.755330,0.223661)--(1.799069,0.310969)--(1.701490,0.215687)--cycle;
\path[fill=blue!22!white,draw=blue!22!white,line width=.12pt] (-3.138513,-0.312145)--(-3.123436,-0.445808)--(-3.047499,-0.597387)--(-3.062209,-0.464455)--cycle;
\path[fill=blue!31!white,draw=blue!31!white,line width=.12pt] (3.062209,-0.464455)--(3.047499,-0.597387)--(3.123436,-0.445808)--(3.138513,-0.312145)--cycle;
\path[fill=blue!33!white,draw=blue!33!white,line width=.12pt] (-1.478207,-0.306147)--(-1.492409,-0.173925)--(-1.424631,-0.245578)--(-1.411074,-0.377117)--cycle;
\path[fill=blue!25!white,draw=blue!25!white,line width=.12pt] (1.411074,-0.377117)--(1.424631,-0.245578)--(1.492409,-0.173925)--(1.478207,-0.306147)--cycle;
\path[fill=blue!15!white,draw=blue!15!white,line width=.12pt] (-2.951405,0.344554)--(-3.050416,0.234690)--(-3.006279,0.085917)--(-2.908701,0.200610)--cycle;
\path[fill=blue!22!white,draw=blue!22!white,line width=.12pt] (2.908701,0.200610)--(3.006279,0.085917)--(3.050416,0.234690)--(2.951405,0.344554)--cycle;
\path[fill=blue!40!white,draw=blue!40!white,line width=.12pt] (-2.116611,-1.258496)--(-1.978968,-1.208398)--(-1.865758,-1.302838)--(-1.995527,-1.359505)--cycle;
\path[fill=blue!40!white,draw=blue!40!white,line width=.12pt] (1.995527,-1.359505)--(1.865758,-1.302838)--(1.978968,-1.208398)--(2.116611,-1.258496)--cycle;
\path[fill=blue!17!white,draw=blue!17!white,line width=.12pt] (-2.083771,0.537466)--(-2.233123,0.569536)--(-2.200811,0.460624)--(-2.053620,0.435837)--cycle;
\path[fill=blue!14!white,draw=blue!14!white,line width=.12pt] (2.053620,0.435837)--(2.200811,0.460624)--(2.233123,0.569536)--(2.083771,0.537466)--cycle;
\path[fill=blue!41!white,draw=blue!41!white,line width=.12pt] (-1.525177,-0.999443)--(-1.442454,-0.866819)--(-1.341037,-0.935191)--(-1.417944,-1.071736)--cycle;
\path[fill=blue!35!white,draw=blue!35!white,line width=.12pt] (1.417944,-1.071736)--(1.341037,-0.935191)--(1.442454,-0.866819)--(1.525177,-0.999443)--cycle;
\path[fill=blue!42!white,draw=blue!42!white,line width=.12pt] (-1.625975,-1.119280)--(-1.525177,-0.999443)--(-1.417944,-1.071736)--(-1.511656,-1.196351)--cycle;
\path[fill=blue!37!white,draw=blue!37!white,line width=.12pt] (1.511656,-1.196351)--(1.417944,-1.071736)--(1.525177,-0.999443)--(1.625975,-1.119280)--cycle;
\path[fill=blue!35!white,draw=blue!35!white,line width=.12pt] (-2.627935,-1.120322)--(-2.500154,-1.157881)--(-2.386608,-1.277917)--(-2.508586,-1.246493)--cycle;
\path[fill=blue!39!white,draw=blue!39!white,line width=.12pt] (2.508586,-1.246493)--(2.386608,-1.277917)--(2.500154,-1.157881)--(2.627935,-1.120322)--cycle;
\path[fill=blue!36!white,draw=blue!36!white,line width=.12pt] (-1.424631,-0.515903)--(-1.411074,-0.377117)--(-1.330351,-0.444456)--(-1.343133,-0.583889)--cycle;
\path[fill=blue!28!white,draw=blue!28!white,line width=.12pt] (1.343133,-0.583889)--(1.330351,-0.444456)--(1.411074,-0.377117)--(1.424631,-0.515903)--cycle;
\path[fill=blue!14!white,draw=blue!14!white,line width=.12pt] (-2.830759,0.436656)--(-2.951405,0.344554)--(-2.908701,0.200610)--(-2.789801,0.298596)--cycle;
\path[fill=blue!19!white,draw=blue!19!white,line width=.12pt] (2.789801,0.298596)--(2.908701,0.200610)--(2.951405,0.344554)--(2.830759,0.436656)--cycle;
\path[fill=blue!27!white,draw=blue!27!white,line width=.12pt] (-1.589379,0.024064)--(-1.660123,0.133114)--(-1.602769,0.052966)--(-1.534468,-0.052668)--cycle;
\path[fill=blue!19!white,draw=blue!19!white,line width=.12pt] (1.534468,-0.052668)--(1.602769,0.052966)--(1.660123,0.133114)--(1.589379,0.024064)--cycle;
\path[fill=blue!40!white,draw=blue!40!white,line width=.12pt] (-1.442454,-0.866819)--(-1.380985,-0.726503)--(-1.283890,-0.791962)--(-1.341037,-0.935191)--cycle;
\path[fill=blue!33!white,draw=blue!33!white,line width=.12pt] (1.341037,-0.935191)--(1.283890,-0.791962)--(1.380985,-0.726503)--(1.442454,-0.866819)--cycle;
\path[fill=blue!15!white,draw=blue!15!white,line width=.12pt] (-2.233123,0.569536)--(-2.388443,0.575200)--(-2.353885,0.458712)--(-2.200811,0.460624)--cycle;
\path[fill=blue!14!white,draw=blue!14!white,line width=.12pt] (2.200811,0.460624)--(2.353885,0.458712)--(2.388443,0.575200)--(2.233123,0.569536)--cycle;
\path[fill=blue!42!white,draw=blue!42!white,line width=.12pt] (-1.740975,-1.221724)--(-1.625975,-1.119280)--(-1.511656,-1.196351)--(-1.618570,-1.304246)--cycle;
\path[fill=blue!38!white,draw=blue!38!white,line width=.12pt] (1.618570,-1.304246)--(1.511656,-1.196351)--(1.625975,-1.119280)--(1.740975,-1.221724)--cycle;
\path[fill=blue!13!white,draw=blue!13!white,line width=.12pt] (-2.693116,0.507458)--(-2.830759,0.436656)--(-2.789801,0.298596)--(-2.654149,0.376111)--cycle;
\path[fill=blue!18!white,draw=blue!18!white,line width=.12pt] (2.654149,0.376111)--(2.789801,0.298596)--(2.830759,0.436656)--(2.693116,0.507458)--cycle;
\path[fill=blue!39!white,draw=blue!39!white,line width=.12pt] (-2.254254,-1.281970)--(-2.116611,-1.258496)--(-1.995527,-1.359505)--(-2.125296,-1.389547)--cycle;
\path[fill=blue!40!white,draw=blue!40!white,line width=.12pt] (2.125296,-1.389547)--(1.995527,-1.359505)--(2.116611,-1.258496)--(2.254254,-1.281970)--cycle;
\path[fill=blue!21!white,draw=blue!21!white,line width=.12pt] (-1.799069,0.310969)--(-1.917969,0.385305)--(-1.871339,0.292227)--(-1.755330,0.223661)--cycle;
\path[fill=blue!16!white,draw=blue!16!white,line width=.12pt] (1.755330,0.223661)--(1.871339,0.292227)--(1.917969,0.385305)--(1.799069,0.310969)--cycle;
\path[fill=blue!20!white,draw=blue!20!white,line width=.12pt] (-3.123436,-0.175483)--(-3.138513,-0.312145)--(-3.062209,-0.464455)--(-3.047499,-0.327062)--cycle;
\path[fill=blue!28!white,draw=blue!28!white,line width=.12pt] (3.047499,-0.327062)--(3.062209,-0.464455)--(3.138513,-0.312145)--(3.123436,-0.175483)--cycle;
\path[fill=blue!14!white,draw=blue!14!white,line width=.12pt] (-2.388443,0.575200)--(-2.543764,0.554239)--(-2.506958,0.430175)--(-2.353885,0.458712)--cycle;
\path[fill=blue!15!white,draw=blue!15!white,line width=.12pt] (2.353885,0.458712)--(2.506958,0.430175)--(2.543764,0.554239)--(2.388443,0.575200)--cycle;
\path[fill=blue!13!white,draw=blue!13!white,line width=.12pt] (-2.543764,0.554239)--(-2.693116,0.507458)--(-2.654149,0.376111)--(-2.506958,0.430175)--cycle;
\path[fill=blue!16!white,draw=blue!16!white,line width=.12pt] (2.506958,0.430175)--(2.654149,0.376111)--(2.693116,0.507458)--(2.543764,0.554239)--cycle;
\path[fill=blue!27!white,draw=blue!27!white,line width=.12pt] (-3.003935,-0.720748)--(-2.933190,-0.829797)--(-2.831853,-0.971406)--(-2.900154,-0.865772)--cycle;
\path[fill=blue!35!white,draw=blue!35!white,line width=.12pt] (2.900154,-0.865772)--(2.831853,-0.971406)--(2.933190,-0.829797)--(3.003935,-0.720748)--cycle;
\path[fill=blue!30!white,draw=blue!30!white,line width=.12pt] (-1.492409,-0.173925)--(-1.534468,-0.052668)--(-1.464780,-0.126340)--(-1.424631,-0.245578)--cycle;
\path[fill=blue!22!white,draw=blue!22!white,line width=.12pt] (1.424631,-0.245578)--(1.464780,-0.126340)--(1.534468,-0.052668)--(1.492409,-0.173925)--cycle;
\path[fill=blue!38!white,draw=blue!38!white,line width=.12pt] (-1.380985,-0.726503)--(-1.343133,-0.583889)--(-1.248699,-0.647553)--(-1.283890,-0.791962)--cycle;
\path[fill=blue!31!white,draw=blue!31!white,line width=.12pt] (1.283890,-0.791962)--(1.248699,-0.647553)--(1.343133,-0.583889)--(1.380985,-0.726503)--cycle;
\path[fill=blue!41!white,draw=blue!41!white,line width=.12pt] (-1.865758,-1.302838)--(-1.740975,-1.221724)--(-1.618570,-1.304246)--(-1.734580,-1.391274)--cycle;
\path[fill=blue!39!white,draw=blue!39!white,line width=.12pt] (1.734580,-1.391274)--(1.618570,-1.304246)--(1.740975,-1.221724)--(1.865758,-1.302838)--cycle;
\path[fill=blue!33!white,draw=blue!33!white,line width=.12pt] (-1.411074,-0.377117)--(-1.424631,-0.245578)--(-1.343133,-0.313564)--(-1.330351,-0.444456)--cycle;
\path[fill=blue!25!white,draw=blue!25!white,line width=.12pt] (1.330351,-0.444456)--(1.343133,-0.313564)--(1.424631,-0.245578)--(1.411074,-0.377117)--cycle;
\path[fill=blue!32!white,draw=blue!32!white,line width=.12pt] (-2.739936,-1.057357)--(-2.627935,-1.120322)--(-2.508586,-1.246493)--(-2.615501,-1.188905)--cycle;
\path[fill=blue!37!white,draw=blue!37!white,line width=.12pt] (2.615501,-1.188905)--(2.508586,-1.246493)--(2.627935,-1.120322)--(2.739936,-1.057357)--cycle;
\path[fill=blue!24!white,draw=blue!24!white,line width=.12pt] (-1.660123,0.133114)--(-1.755330,0.223661)--(-1.694686,0.138917)--(-1.602769,0.052966)--cycle;
\path[fill=blue!17!white,draw=blue!17!white,line width=.12pt] (1.602769,0.052966)--(1.694686,0.138917)--(1.755330,0.223661)--(1.660123,0.133114)--cycle;
\path[fill=blue!19!white,draw=blue!19!white,line width=.12pt] (-1.917969,0.385305)--(-2.053620,0.435837)--(-2.003693,0.336176)--(-1.871339,0.292227)--cycle;
\path[fill=blue!15!white,draw=blue!15!white,line width=.12pt] (1.871339,0.292227)--(2.003693,0.336176)--(2.053620,0.435837)--(1.917969,0.385305)--cycle;
\path[fill=blue!18!white,draw=blue!18!white,line width=.12pt] (-3.078787,-0.041074)--(-3.123436,-0.175483)--(-3.047499,-0.327062)--(-3.003935,-0.190486)--cycle;
\path[fill=blue!26!white,draw=blue!26!white,line width=.12pt] (3.003935,-0.190486)--(3.047499,-0.327062)--(3.123436,-0.175483)--(3.078787,-0.041074)--cycle;
\path[fill=blue!41!white,draw=blue!41!white,line width=.12pt] (-1.417944,-1.071736)--(-1.341037,-0.935191)--(-1.226706,-0.998263)--(-1.297056,-1.138426)--cycle;
\path[fill=blue!36!white,draw=blue!36!white,line width=.12pt] (1.297056,-1.138426)--(1.226706,-0.998263)--(1.341037,-0.935191)--(1.417944,-1.071736)--cycle;
\path[fill=blue!37!white,draw=blue!37!white,line width=.12pt] (-2.386608,-1.277917)--(-2.254254,-1.281970)--(-2.125296,-1.389547)--(-2.250079,-1.391810)--cycle;
\path[fill=blue!39!white,draw=blue!39!white,line width=.12pt] (2.250079,-1.391810)--(2.125296,-1.389547)--(2.254254,-1.281970)--(2.386608,-1.277917)--cycle;
\path[fill=blue!36!white,draw=blue!36!white,line width=.12pt] (-1.343133,-0.583889)--(-1.330351,-0.444456)--(-1.236817,-0.507515)--(-1.248699,-0.647553)--cycle;
\path[fill=blue!29!white,draw=blue!29!white,line width=.12pt] (1.248699,-0.647553)--(1.236817,-0.507515)--(1.330351,-0.444456)--(1.343133,-0.583889)--cycle;
\path[fill=blue!42!white,draw=blue!42!white,line width=.12pt] (-1.511656,-1.196351)--(-1.417944,-1.071736)--(-1.297056,-1.138426)--(-1.382778,-1.267448)--cycle;
\path[fill=blue!38!white,draw=blue!38!white,line width=.12pt] (1.382778,-1.267448)--(1.297056,-1.138426)--(1.417944,-1.071736)--(1.511656,-1.196351)--cycle;
\path[fill=blue!40!white,draw=blue!40!white,line width=.12pt] (-1.341037,-0.935191)--(-1.283890,-0.791962)--(-1.174431,-0.852346)--(-1.226706,-0.998263)--cycle;
\path[fill=blue!34!white,draw=blue!34!white,line width=.12pt] (1.226706,-0.998263)--(1.174431,-0.852346)--(1.283890,-0.791962)--(1.341037,-0.935191)--cycle;
\path[fill=blue!40!white,draw=blue!40!white,line width=.12pt] (-1.995527,-1.359505)--(-1.865758,-1.302838)--(-1.734580,-1.391274)--(-1.855225,-1.454092)--cycle;
\path[fill=blue!40!white,draw=blue!40!white,line width=.12pt] (1.855225,-1.454092)--(1.734580,-1.391274)--(1.865758,-1.302838)--(1.995527,-1.359505)--cycle;
\path[fill=blue!24!white,draw=blue!24!white,line width=.12pt] (-3.047499,-0.597387)--(-3.003935,-0.720748)--(-2.900154,-0.865772)--(-2.942213,-0.744515)--cycle;
\path[fill=blue!33!white,draw=blue!33!white,line width=.12pt] (2.942213,-0.744515)--(2.900154,-0.865772)--(3.003935,-0.720748)--(3.047499,-0.597387)--cycle;
\path[fill=blue!27!white,draw=blue!27!white,line width=.12pt] (-1.534468,-0.052668)--(-1.602769,0.052966)--(-1.529978,-0.023985)--(-1.464780,-0.126340)--cycle;
\path[fill=blue!20!white,draw=blue!20!white,line width=.12pt] (1.464780,-0.126340)--(1.529978,-0.023985)--(1.602769,0.052966)--(1.534468,-0.052668)--cycle;
\path[fill=blue!42!white,draw=blue!42!white,line width=.12pt] (-1.618570,-1.304246)--(-1.511656,-1.196351)--(-1.382778,-1.267448)--(-1.480578,-1.380371)--cycle;
\path[fill=blue!39!white,draw=blue!39!white,line width=.12pt] (1.480578,-1.380371)--(1.382778,-1.267448)--(1.511656,-1.196351)--(1.618570,-1.304246)--cycle;
\path[fill=blue!38!white,draw=blue!38!white,line width=.12pt] (-1.283890,-0.791962)--(-1.248699,-0.647553)--(-1.142240,-0.706283)--(-1.174431,-0.852346)--cycle;
\path[fill=blue!32!white,draw=blue!32!white,line width=.12pt] (1.174431,-0.852346)--(1.142240,-0.706283)--(1.248699,-0.647553)--(1.283890,-0.791962)--cycle;
\path[fill=blue!17!white,draw=blue!17!white,line width=.12pt] (-2.053620,0.435837)--(-2.200811,0.460624)--(-2.147305,0.353819)--(-2.003693,0.336176)--cycle;
\path[fill=blue!14!white,draw=blue!14!white,line width=.12pt] (2.003693,0.336176)--(2.147305,0.353819)--(2.200811,0.460624)--(2.053620,0.435837)--cycle;
\path[fill=blue!16!white,draw=blue!16!white,line width=.12pt] (-3.006279,0.085917)--(-3.078787,-0.041074)--(-3.003935,-0.190486)--(-2.933190,-0.059976)--cycle;
\path[fill=blue!23!white,draw=blue!23!white,line width=.12pt] (2.933190,-0.059976)--(3.003935,-0.190486)--(3.078787,-0.041074)--(3.006279,0.085917)--cycle;
\path[fill=blue!30!white,draw=blue!30!white,line width=.12pt] (-1.424631,-0.245578)--(-1.464780,-0.126340)--(-1.380985,-0.196241)--(-1.343133,-0.313564)--cycle;
\path[fill=blue!23!white,draw=blue!23!white,line width=.12pt] (1.343133,-0.313564)--(1.380985,-0.196241)--(1.464780,-0.126340)--(1.424631,-0.245578)--cycle;
\path[fill=blue!21!white,draw=blue!21!white,line width=.12pt] (-1.755330,0.223661)--(-1.871339,0.292227)--(-1.806687,0.201882)--(-1.694686,0.138917)--cycle;
\path[fill=blue!16!white,draw=blue!16!white,line width=.12pt] (1.694686,0.138917)--(1.806687,0.201882)--(1.871339,0.292227)--(1.755330,0.223661)--cycle;
\path[fill=blue!30!white,draw=blue!30!white,line width=.12pt] (-2.831853,-0.971406)--(-2.739936,-1.057357)--(-2.615501,-1.188905)--(-2.703244,-1.107367)--cycle;
\path[fill=blue!36!white,draw=blue!36!white,line width=.12pt] (2.703244,-1.107367)--(2.615501,-1.188905)--(2.739936,-1.057357)--(2.831853,-0.971406)--cycle;
\path[fill=blue!33!white,draw=blue!33!white,line width=.12pt] (-1.330351,-0.444456)--(-1.343133,-0.313564)--(-1.248699,-0.377228)--(-1.236817,-0.507515)--cycle;
\path[fill=blue!26!white,draw=blue!26!white,line width=.12pt] (1.236817,-0.507515)--(1.248699,-0.377228)--(1.343133,-0.313564)--(1.330351,-0.444456)--cycle;
\path[fill=blue!41!white,draw=blue!41!white,line width=.12pt] (-1.734580,-1.391274)--(-1.618570,-1.304246)--(-1.480578,-1.380371)--(-1.586697,-1.472856)--cycle;
\path[fill=blue!39!white,draw=blue!39!white,line width=.12pt] (1.586697,-1.472856)--(1.480578,-1.380371)--(1.618570,-1.304246)--(1.734580,-1.391274)--cycle;
\path[fill=blue!39!white,draw=blue!39!white,line width=.12pt] (-2.125296,-1.389547)--(-1.995527,-1.359505)--(-1.855225,-1.454092)--(-1.975871,-1.490286)--cycle;
\path[fill=blue!40!white,draw=blue!40!white,line width=.12pt] (1.975871,-1.490286)--(1.855225,-1.454092)--(1.995527,-1.359505)--(2.125296,-1.389547)--cycle;
\path[fill=blue!35!white,draw=blue!35!white,line width=.12pt] (-2.508586,-1.246493)--(-2.386608,-1.277917)--(-2.250079,-1.391810)--(-2.365079,-1.366207)--cycle;
\path[fill=blue!38!white,draw=blue!38!white,line width=.12pt] (2.365079,-1.366207)--(2.250079,-1.391810)--(2.386608,-1.277917)--(2.508586,-1.246493)--cycle;
\path[fill=blue!15!white,draw=blue!15!white,line width=.12pt] (-2.200811,0.460624)--(-2.353885,0.458712)--(-2.296657,0.344479)--(-2.147305,0.353819)--cycle;
\path[fill=blue!14!white,draw=blue!14!white,line width=.12pt] (2.147305,0.353819)--(2.296657,0.344479)--(2.353885,0.458712)--(2.200811,0.460624)--cycle;
\path[fill=blue!14!white,draw=blue!14!white,line width=.12pt] (-2.908701,0.200610)--(-3.006279,0.085917)--(-2.933190,-0.059976)--(-2.837984,0.059451)--cycle;
\path[fill=blue!21!white,draw=blue!21!white,line width=.12pt] (2.837984,0.059451)--(2.933190,-0.059976)--(3.006279,0.085917)--(2.908701,0.200610)--cycle;
\path[fill=blue!41!white,draw=blue!41!white,line width=.12pt] (-1.297056,-1.138426)--(-1.226706,-0.998263)--(-1.100561,-1.055429)--(-1.163677,-1.198870)--cycle;
\path[fill=blue!37!white,draw=blue!37!white,line width=.12pt] (1.163677,-1.198870)--(1.100561,-1.055429)--(1.226706,-0.998263)--(1.297056,-1.138426)--cycle;
\path[fill=blue!22!white,draw=blue!22!white,line width=.12pt] (-3.062209,-0.464455)--(-3.047499,-0.597387)--(-2.942213,-0.744515)--(-2.956415,-0.612293)--cycle;
\path[fill=blue!30!white,draw=blue!30!white,line width=.12pt] (2.956415,-0.612293)--(2.942213,-0.744515)--(3.047499,-0.597387)--(3.062209,-0.464455)--cycle;
\path[fill=blue!36!white,draw=blue!36!white,line width=.12pt] (-1.248699,-0.647553)--(-1.236817,-0.507515)--(-1.131371,-0.565685)--(-1.142240,-0.706283)--cycle;
\path[fill=blue!29!white,draw=blue!29!white,line width=.12pt] (1.142240,-0.706283)--(1.131371,-0.565685)--(1.236817,-0.507515)--(1.248699,-0.647553)--cycle;
\path[fill=blue!40!white,draw=blue!40!white,line width=.12pt] (-1.226706,-0.998263)--(-1.174431,-0.852346)--(-1.053662,-0.907076)--(-1.100561,-1.055429)--cycle;
\path[fill=blue!35!white,draw=blue!35!white,line width=.12pt] (1.100561,-1.055429)--(1.053662,-0.907076)--(1.174431,-0.852346)--(1.226706,-0.998263)--cycle;
\path[fill=blue!42!white,draw=blue!42!white,line width=.12pt] (-1.382778,-1.267448)--(-1.297056,-1.138426)--(-1.163677,-1.198870)--(-1.240584,-1.331887)--cycle;
\path[fill=blue!38!white,draw=blue!38!white,line width=.12pt] (1.240584,-1.331887)--(1.163677,-1.198870)--(1.297056,-1.138426)--(1.382778,-1.267448)--cycle;
\path[fill=blue!24!white,draw=blue!24!white,line width=.12pt] (-1.602769,0.052966)--(-1.694686,0.138917)--(-1.617721,0.057553)--(-1.529978,-0.023985)--cycle;
\path[fill=blue!18!white,draw=blue!18!white,line width=.12pt] (1.529978,-0.023985)--(1.617721,0.057553)--(1.694686,0.138917)--(1.602769,0.052966)--cycle;
\path[fill=blue!14!white,draw=blue!14!white,line width=.12pt] (-2.353885,0.458712)--(-2.506958,0.430175)--(-2.446009,0.308514)--(-2.296657,0.344479)--cycle;
\path[fill=blue!15!white,draw=blue!15!white,line width=.12pt] (2.296657,0.344479)--(2.446009,0.308514)--(2.506958,0.430175)--(2.353885,0.458712)--cycle;
\path[fill=blue!13!white,draw=blue!13!white,line width=.12pt] (-2.789801,0.298596)--(-2.908701,0.200610)--(-2.837984,0.059451)--(-2.721975,0.163208)--cycle;
\path[fill=blue!19!white,draw=blue!19!white,line width=.12pt] (2.721975,0.163208)--(2.837984,0.059451)--(2.908701,0.200610)--(2.789801,0.298596)--cycle;
\path[fill=blue!13!white,draw=blue!13!white,line width=.12pt] (-2.506958,0.430175)--(-2.654149,0.376111)--(-2.589621,0.247306)--(-2.446009,0.308514)--cycle;
\path[fill=blue!16!white,draw=blue!16!white,line width=.12pt] (2.446009,0.308514)--(2.589621,0.247306)--(2.654149,0.376111)--(2.506958,0.430175)--cycle;
\path[fill=blue!38!white,draw=blue!38!white,line width=.12pt] (-1.174431,-0.852346)--(-1.142240,-0.706283)--(-1.024781,-0.759512)--(-1.053662,-0.907076)--cycle;
\path[fill=blue!33!white,draw=blue!33!white,line width=.12pt] (1.053662,-0.907076)--(1.024781,-0.759512)--(1.142240,-0.706283)--(1.174431,-0.852346)--cycle;
\path[fill=blue!13!white,draw=blue!13!white,line width=.12pt] (-2.654149,0.376111)--(-2.789801,0.298596)--(-2.721975,0.163208)--(-2.589621,0.247306)--cycle;
\path[fill=blue!17!white,draw=blue!17!white,line width=.12pt] (2.589621,0.247306)--(2.721975,0.163208)--(2.789801,0.298596)--(2.654149,0.376111)--cycle;
\path[fill=blue!19!white,draw=blue!19!white,line width=.12pt] (-1.871339,0.292227)--(-2.003693,0.336176)--(-1.934468,0.239441)--(-1.806687,0.201882)--cycle;
\path[fill=blue!15!white,draw=blue!15!white,line width=.12pt] (1.806687,0.201882)--(1.934468,0.239441)--(2.003693,0.336176)--(1.871339,0.292227)--cycle;
\path[fill=blue!28!white,draw=blue!28!white,line width=.12pt] (-1.464780,-0.126340)--(-1.529978,-0.023985)--(-1.442454,-0.096998)--(-1.380985,-0.196241)--cycle;
\path[fill=blue!21!white,draw=blue!21!white,line width=.12pt] (1.380985,-0.196241)--(1.442454,-0.096998)--(1.529978,-0.023985)--(1.464780,-0.126340)--cycle;
\path[fill=blue!42!white,draw=blue!42!white,line width=.12pt] (-1.480578,-1.380371)--(-1.382778,-1.267448)--(-1.240584,-1.331887)--(-1.328326,-1.449368)--cycle;
\path[fill=blue!39!white,draw=blue!39!white,line width=.12pt] (1.328326,-1.449368)--(1.240584,-1.331887)--(1.382778,-1.267448)--(1.480578,-1.380371)--cycle;
\path[fill=blue!40!white,draw=blue!40!white,line width=.12pt] (-1.855225,-1.454092)--(-1.734580,-1.391274)--(-1.586697,-1.472856)--(-1.697056,-1.541348)--cycle;
\path[fill=blue!40!white,draw=blue!40!white,line width=.12pt] (1.697056,-1.541348)--(1.586697,-1.472856)--(1.734580,-1.391274)--(1.855225,-1.454092)--cycle;
\path[fill=blue!27!white,draw=blue!27!white,line width=.12pt] (-2.900154,-0.865772)--(-2.831853,-0.971406)--(-2.703244,-1.107367)--(-2.768442,-1.005012)--cycle;
\path[fill=blue!34!white,draw=blue!34!white,line width=.12pt] (2.768442,-1.005012)--(2.703244,-1.107367)--(2.831853,-0.971406)--(2.900154,-0.865772)--cycle;
\path[fill=blue!31!white,draw=blue!31!white,line width=.12pt] (-1.343133,-0.313564)--(-1.380985,-0.196241)--(-1.283890,-0.261700)--(-1.248699,-0.377228)--cycle;
\path[fill=blue!24!white,draw=blue!24!white,line width=.12pt] (1.248699,-0.377228)--(1.283890,-0.261700)--(1.380985,-0.196241)--(1.343133,-0.313564)--cycle;
\path[fill=blue!19!white,draw=blue!19!white,line width=.12pt] (-3.047499,-0.327062)--(-3.062209,-0.464455)--(-2.956415,-0.612293)--(-2.942213,-0.474190)--cycle;
\path[fill=blue!28!white,draw=blue!28!white,line width=.12pt] (2.942213,-0.474190)--(2.956415,-0.612293)--(3.062209,-0.464455)--(3.047499,-0.327062)--cycle;
\path[fill=blue!34!white,draw=blue!34!white,line width=.12pt] (-1.236817,-0.507515)--(-1.248699,-0.377228)--(-1.142240,-0.435958)--(-1.131371,-0.565685)--cycle;
\path[fill=blue!27!white,draw=blue!27!white,line width=.12pt] (1.131371,-0.565685)--(1.142240,-0.435958)--(1.248699,-0.377228)--(1.236817,-0.507515)--cycle;
\path[fill=blue!37!white,draw=blue!37!white,line width=.12pt] (-2.250079,-1.391810)--(-2.125296,-1.389547)--(-1.975871,-1.490286)--(-2.091880,-1.498463)--cycle;
\path[fill=blue!39!white,draw=blue!39!white,line width=.12pt] (2.091880,-1.498463)--(1.975871,-1.490286)--(2.125296,-1.389547)--(2.250079,-1.391810)--cycle;
\path[fill=blue!32!white,draw=blue!32!white,line width=.12pt] (-2.615501,-1.188905)--(-2.508586,-1.246493)--(-2.365079,-1.366207)--(-2.465877,-1.313721)--cycle;
\path[fill=blue!37!white,draw=blue!37!white,line width=.12pt] (2.465877,-1.313721)--(2.365079,-1.366207)--(2.508586,-1.246493)--(2.615501,-1.188905)--cycle;
\path[fill=blue!41!white,draw=blue!41!white,line width=.12pt] (-1.163677,-1.198870)--(-1.100561,-1.055429)--(-0.963817,-1.106137)--(-1.019091,-1.252486)--cycle;
\path[fill=blue!37!white,draw=blue!37!white,line width=.12pt] (1.019091,-1.252486)--(0.963817,-1.106137)--(1.100561,-1.055429)--(1.163677,-1.198870)--cycle;
\path[fill=blue!36!white,draw=blue!36!white,line width=.12pt] (-1.142240,-0.706283)--(-1.131371,-0.565685)--(-1.015029,-0.618408)--(-1.024781,-0.759512)--cycle;
\path[fill=blue!30!white,draw=blue!30!white,line width=.12pt] (1.024781,-0.759512)--(1.015029,-0.618408)--(1.131371,-0.565685)--(1.142240,-0.706283)--cycle;
\path[fill=blue!40!white,draw=blue!40!white,line width=.12pt] (-1.100561,-1.055429)--(-1.053662,-0.907076)--(-0.922745,-0.955623)--(-0.963817,-1.106137)--cycle;
\path[fill=blue!35!white,draw=blue!35!white,line width=.12pt] (0.963817,-1.106137)--(0.922745,-0.955623)--(1.053662,-0.907076)--(1.100561,-1.055429)--cycle;
\path[fill=blue!22!white,draw=blue!22!white,line width=.12pt] (-1.694686,0.138917)--(-1.806687,0.201882)--(-1.724636,0.115141)--(-1.617721,0.057553)--cycle;
\path[fill=blue!17!white,draw=blue!17!white,line width=.12pt] (1.617721,0.057553)--(1.724636,0.115141)--(1.806687,0.201882)--(1.694686,0.138917)--cycle;
\path[fill=blue!42!white,draw=blue!42!white,line width=.12pt] (-1.240584,-1.331887)--(-1.163677,-1.198870)--(-1.019091,-1.252486)--(-1.086442,-1.389047)--cycle;
\path[fill=blue!38!white,draw=blue!38!white,line width=.12pt] (1.086442,-1.389047)--(1.019091,-1.252486)--(1.163677,-1.198870)--(1.240584,-1.331887)--cycle;
\path[fill=blue!41!white,draw=blue!41!white,line width=.12pt] (-1.586697,-1.472856)--(-1.480578,-1.380371)--(-1.328326,-1.449368)--(-1.423533,-1.546798)--cycle;
\path[fill=blue!40!white,draw=blue!40!white,line width=.12pt] (1.423533,-1.546798)--(1.328326,-1.449368)--(1.480578,-1.380371)--(1.586697,-1.472856)--cycle;
\path[fill=blue!17!white,draw=blue!17!white,line width=.12pt] (-2.003693,0.336176)--(-2.147305,0.353819)--(-2.073119,0.250151)--(-1.934468,0.239441)--cycle;
\path[fill=blue!14!white,draw=blue!14!white,line width=.12pt] (1.934468,0.239441)--(2.073119,0.250151)--(2.147305,0.353819)--(2.003693,0.336176)--cycle;
\path[fill=blue!38!white,draw=blue!38!white,line width=.12pt] (-1.053662,-0.907076)--(-1.024781,-0.759512)--(-0.897452,-0.806729)--(-0.922745,-0.955623)--cycle;
\path[fill=blue!33!white,draw=blue!33!white,line width=.12pt] (0.922745,-0.955623)--(0.897452,-0.806729)--(1.024781,-0.759512)--(1.053662,-0.907076)--cycle;
\path[fill=blue!17!white,draw=blue!17!white,line width=.12pt] (-3.003935,-0.190486)--(-3.047499,-0.327062)--(-2.942213,-0.474190)--(-2.900154,-0.335511)--cycle;
\path[fill=blue!25!white,draw=blue!25!white,line width=.12pt] (2.900154,-0.335511)--(2.942213,-0.474190)--(3.047499,-0.327062)--(3.003935,-0.190486)--cycle;
\path[fill=blue!39!white,draw=blue!39!white,line width=.12pt] (-1.975871,-1.490286)--(-1.855225,-1.454092)--(-1.697056,-1.541348)--(-1.807416,-1.583216)--cycle;
\path[fill=blue!39!white,draw=blue!39!white,line width=.12pt] (1.807416,-1.583216)--(1.697056,-1.541348)--(1.855225,-1.454092)--(1.975871,-1.490286)--cycle;
\path[fill=blue!25!white,draw=blue!25!white,line width=.12pt] (-1.529978,-0.023985)--(-1.617721,0.057553)--(-1.525177,-0.019647)--(-1.442454,-0.096998)--cycle;
\path[fill=blue!19!white,draw=blue!19!white,line width=.12pt] (1.442454,-0.096998)--(1.525177,-0.019647)--(1.617721,0.057553)--(1.529978,-0.023985)--cycle;
\path[fill=blue!24!white,draw=blue!24!white,line width=.12pt] (-2.942213,-0.744515)--(-2.900154,-0.865772)--(-2.768442,-1.005012)--(-2.808591,-0.885774)--cycle;
\path[fill=blue!32!white,draw=blue!32!white,line width=.12pt] (2.808591,-0.885774)--(2.768442,-1.005012)--(2.900154,-0.865772)--(2.942213,-0.744515)--cycle;
\path[fill=blue!42!white,draw=blue!42!white,line width=.12pt] (-1.328326,-1.449368)--(-1.240584,-1.331887)--(-1.086442,-1.389047)--(-1.163283,-1.510570)--cycle;
\path[fill=blue!39!white,draw=blue!39!white,line width=.12pt] (1.163283,-1.510570)--(1.086442,-1.389047)--(1.240584,-1.331887)--(1.328326,-1.449368)--cycle;
\path[fill=blue!28!white,draw=blue!28!white,line width=.12pt] (-1.380985,-0.196241)--(-1.442454,-0.096998)--(-1.341037,-0.165370)--(-1.283890,-0.261700)--cycle;
\path[fill=blue!21!white,draw=blue!21!white,line width=.12pt] (1.283890,-0.261700)--(1.341037,-0.165370)--(1.442454,-0.096998)--(1.380985,-0.196241)--cycle;
\path[fill=blue!31!white,draw=blue!31!white,line width=.12pt] (-1.248699,-0.377228)--(-1.283890,-0.261700)--(-1.174431,-0.322085)--(-1.142240,-0.435958)--cycle;
\path[fill=blue!24!white,draw=blue!24!white,line width=.12pt] (1.142240,-0.435958)--(1.174431,-0.322085)--(1.283890,-0.261700)--(1.248699,-0.377228)--cycle;
\path[fill=blue!34!white,draw=blue!34!white,line width=.12pt] (-1.131371,-0.565685)--(-1.142240,-0.435958)--(-1.024781,-0.489187)--(-1.015029,-0.618408)--cycle;
\path[fill=blue!28!white,draw=blue!28!white,line width=.12pt] (1.015029,-0.618408)--(1.024781,-0.489187)--(1.142240,-0.435958)--(1.131371,-0.565685)--cycle;
\path[fill=blue!41!white,draw=blue!41!white,line width=.12pt] (-1.019091,-1.252486)--(-0.963817,-1.106137)--(-0.817791,-1.149900)--(-0.864690,-1.298758)--cycle;
\path[fill=blue!38!white,draw=blue!38!white,line width=.12pt] (0.864690,-1.298758)--(0.817791,-1.149900)--(0.963817,-1.106137)--(1.019091,-1.252486)--cycle;
\path[fill=blue!40!white,draw=blue!40!white,line width=.12pt] (-0.963817,-1.106137)--(-0.922745,-0.955623)--(-0.782941,-0.997521)--(-0.817791,-1.149900)--cycle;
\path[fill=blue!36!white,draw=blue!36!white,line width=.12pt] (0.817791,-1.149900)--(0.782941,-0.997521)--(0.922745,-0.955623)--(0.963817,-1.106137)--cycle;
\path[fill=blue!15!white,draw=blue!15!white,line width=.12pt] (-2.147305,0.353819)--(-2.296657,0.344479)--(-2.217311,0.233600)--(-2.073119,0.250151)--cycle;
\path[fill=blue!14!white,draw=blue!14!white,line width=.12pt] (2.073119,0.250151)--(2.217311,0.233600)--(2.296657,0.344479)--(2.147305,0.353819)--cycle;
\path[fill=blue!36!white,draw=blue!36!white,line width=.12pt] (-1.024781,-0.759512)--(-1.015029,-0.618408)--(-0.888912,-0.665176)--(-0.897452,-0.806729)--cycle;
\path[fill=blue!31!white,draw=blue!31!white,line width=.12pt] (0.897452,-0.806729)--(0.888912,-0.665176)--(1.015029,-0.618408)--(1.024781,-0.759512)--cycle;
\path[fill=blue!35!white,draw=blue!35!white,line width=.12pt] (-2.365079,-1.366207)--(-2.250079,-1.391810)--(-2.091880,-1.498463)--(-2.198794,-1.478311)--cycle;
\path[fill=blue!38!white,draw=blue!38!white,line width=.12pt] (2.198794,-1.478311)--(2.091880,-1.498463)--(2.250079,-1.391810)--(2.365079,-1.366207)--cycle;
\path[fill=blue!29!white,draw=blue!29!white,line width=.12pt] (-2.703244,-1.107367)--(-2.615501,-1.188905)--(-2.465877,-1.313721)--(-2.548600,-1.236371)--cycle;
\path[fill=blue!35!white,draw=blue!35!white,line width=.12pt] (2.548600,-1.236371)--(2.465877,-1.313721)--(2.615501,-1.188905)--(2.703244,-1.107367)--cycle;
\path[fill=blue!15!white,draw=blue!15!white,line width=.12pt] (-2.933190,-0.059976)--(-3.003935,-0.190486)--(-2.900154,-0.335511)--(-2.831853,-0.201586)--cycle;
\path[fill=blue!23!white,draw=blue!23!white,line width=.12pt] (2.831853,-0.201586)--(2.900154,-0.335511)--(3.003935,-0.190486)--(2.933190,-0.059976)--cycle;
\path[fill=blue!40!white,draw=blue!40!white,line width=.12pt] (-1.697056,-1.541348)--(-1.586697,-1.472856)--(-1.423533,-1.546798)--(-1.522544,-1.620433)--cycle;
\path[fill=blue!40!white,draw=blue!40!white,line width=.12pt] (1.522544,-1.620433)--(1.423533,-1.546798)--(1.586697,-1.472856)--(1.697056,-1.541348)--cycle;
\path[fill=blue!42!white,draw=blue!42!white,line width=.12pt] (-1.086442,-1.389047)--(-1.019091,-1.252486)--(-0.864690,-1.298758)--(-0.921837,-1.438377)--cycle;
\path[fill=blue!39!white,draw=blue!39!white,line width=.12pt] (0.921837,-1.438377)--(0.864690,-1.298758)--(1.019091,-1.252486)--(1.086442,-1.389047)--cycle;
\path[fill=blue!38!white,draw=blue!38!white,line width=.12pt] (-0.922745,-0.955623)--(-0.897452,-0.806729)--(-0.761481,-0.847478)--(-0.782941,-0.997521)--cycle;
\path[fill=blue!34!white,draw=blue!34!white,line width=.12pt] (0.782941,-0.997521)--(0.761481,-0.847478)--(0.897452,-0.806729)--(0.922745,-0.955623)--cycle;
\path[fill=blue!19!white,draw=blue!19!white,line width=.12pt] (-1.806687,0.201882)--(-1.934468,0.239441)--(-1.846614,0.146565)--(-1.724636,0.115141)--cycle;
\path[fill=blue!15!white,draw=blue!15!white,line width=.12pt] (1.724636,0.115141)--(1.846614,0.146565)--(1.934468,0.239441)--(1.806687,0.201882)--cycle;
\path[fill=blue!14!white,draw=blue!14!white,line width=.12pt] (-2.296657,0.344479)--(-2.446009,0.308514)--(-2.361503,0.190425)--(-2.217311,0.233600)--cycle;
\path[fill=blue!15!white,draw=blue!15!white,line width=.12pt] (2.217311,0.233600)--(2.361503,0.190425)--(2.446009,0.308514)--(2.296657,0.344479)--cycle;
\path[fill=blue!41!white,draw=blue!41!white,line width=.12pt] (-1.423533,-1.546798)--(-1.328326,-1.449368)--(-1.163283,-1.510570)--(-1.246659,-1.612387)--cycle;
\path[fill=blue!40!white,draw=blue!40!white,line width=.12pt] (1.246659,-1.612387)--(1.163283,-1.510570)--(1.328326,-1.449368)--(1.423533,-1.546798)--cycle;
\path[fill=blue!14!white,draw=blue!14!white,line width=.12pt] (-2.837984,0.059451)--(-2.933190,-0.059976)--(-2.831853,-0.201586)--(-2.739936,-0.077561)--cycle;
\path[fill=blue!21!white,draw=blue!21!white,line width=.12pt] (2.739936,-0.077561)--(2.831853,-0.201586)--(2.933190,-0.059976)--(2.837984,0.059451)--cycle;
\path[fill=blue!40!white,draw=blue!40!white,line width=.12pt] (-0.817791,-1.149900)--(-0.782941,-0.997521)--(-0.635598,-1.032365)--(-0.663889,-1.186295)--cycle;
\path[fill=blue!37!white,draw=blue!37!white,line width=.12pt] (0.663889,-1.186295)--(0.635598,-1.032365)--(0.782941,-0.997521)--(0.817791,-1.149900)--cycle;
\path[fill=blue!21!white,draw=blue!21!white,line width=.12pt] (-2.956415,-0.612293)--(-2.942213,-0.744515)--(-2.808591,-0.885774)--(-2.822148,-0.754235)--cycle;
\path[fill=blue!29!white,draw=blue!29!white,line width=.12pt] (2.822148,-0.754235)--(2.808591,-0.885774)--(2.942213,-0.744515)--(2.956415,-0.612293)--cycle;
\path[fill=blue!41!white,draw=blue!41!white,line width=.12pt] (-0.864690,-1.298758)--(-0.817791,-1.149900)--(-0.663889,-1.186295)--(-0.701962,-1.337241)--cycle;
\path[fill=blue!38!white,draw=blue!38!white,line width=.12pt] (0.701962,-1.337241)--(0.663889,-1.186295)--(0.817791,-1.149900)--(0.864690,-1.298758)--cycle;
\path[fill=blue!37!white,draw=blue!37!white,line width=.12pt] (-2.091880,-1.498463)--(-1.975871,-1.490286)--(-1.807416,-1.583216)--(-1.913535,-1.596850)--cycle;
\path[fill=blue!39!white,draw=blue!39!white,line width=.12pt] (1.913535,-1.596850)--(1.807416,-1.583216)--(1.975871,-1.490286)--(2.091880,-1.498463)--cycle;
\path[fill=blue!13!white,draw=blue!13!white,line width=.12pt] (-2.446009,0.308514)--(-2.589621,0.247306)--(-2.500154,0.122284)--(-2.361503,0.190425)--cycle;
\path[fill=blue!16!white,draw=blue!16!white,line width=.12pt] (2.361503,0.190425)--(2.500154,0.122284)--(2.589621,0.247306)--(2.446009,0.308514)--cycle;
\path[fill=blue!22!white,draw=blue!22!white,line width=.12pt] (-1.617721,0.057553)--(-1.724636,0.115141)--(-1.625975,0.032838)--(-1.525177,-0.019647)--cycle;
\path[fill=blue!17!white,draw=blue!17!white,line width=.12pt] (1.525177,-0.019647)--(1.625975,0.032838)--(1.724636,0.115141)--(1.617721,0.057553)--cycle;
\path[fill=blue!42!white,draw=blue!42!white,line width=.12pt] (-1.163283,-1.510570)--(-1.086442,-1.389047)--(-0.921837,-1.438377)--(-0.987036,-1.563389)--cycle;
\path[fill=blue!40!white,draw=blue!40!white,line width=.12pt] (0.987036,-1.563389)--(0.921837,-1.438377)--(1.086442,-1.389047)--(1.163283,-1.510570)--cycle;
\path[fill=blue!36!white,draw=blue!36!white,line width=.12pt] (-0.897452,-0.806729)--(-0.888912,-0.665176)--(-0.754235,-0.705537)--(-0.761481,-0.847478)--cycle;
\path[fill=blue!32!white,draw=blue!32!white,line width=.12pt] (0.761481,-0.847478)--(0.754235,-0.705537)--(0.888912,-0.665176)--(0.897452,-0.806729)--cycle;
\path[fill=blue!13!white,draw=blue!13!white,line width=.12pt] (-2.721975,0.163208)--(-2.837984,0.059451)--(-2.739936,-0.077561)--(-2.627935,0.031796)--cycle;
\path[fill=blue!18!white,draw=blue!18!white,line width=.12pt] (2.627935,0.031796)--(2.739936,-0.077561)--(2.837984,0.059451)--(2.721975,0.163208)--cycle;
\path[fill=blue!34!white,draw=blue!34!white,line width=.12pt] (-1.015029,-0.618408)--(-1.024781,-0.489187)--(-0.897452,-0.536404)--(-0.888912,-0.665176)--cycle;
\path[fill=blue!28!white,draw=blue!28!white,line width=.12pt] (0.888912,-0.665176)--(0.897452,-0.536404)--(1.024781,-0.489187)--(1.015029,-0.618408)--cycle;
\path[fill=blue!13!white,draw=blue!13!white,line width=.12pt] (-2.589621,0.247306)--(-2.721975,0.163208)--(-2.627935,0.031796)--(-2.500154,0.122284)--cycle;
\path[fill=blue!17!white,draw=blue!17!white,line width=.12pt] (2.500154,0.122284)--(2.627935,0.031796)--(2.721975,0.163208)--(2.589621,0.247306)--cycle;
\path[fill=blue!38!white,draw=blue!38!white,line width=.12pt] (-0.782941,-0.997521)--(-0.761481,-0.847478)--(-0.618176,-0.881367)--(-0.635598,-1.032365)--cycle;
\path[fill=blue!35!white,draw=blue!35!white,line width=.12pt] (0.635598,-1.032365)--(0.618176,-0.881367)--(0.761481,-0.847478)--(0.782941,-0.997521)--cycle;
\path[fill=blue!31!white,draw=blue!31!white,line width=.12pt] (-1.142240,-0.435958)--(-1.174431,-0.322085)--(-1.053662,-0.376814)--(-1.024781,-0.489187)--cycle;
\path[fill=blue!25!white,draw=blue!25!white,line width=.12pt] (1.024781,-0.489187)--(1.053662,-0.376814)--(1.174431,-0.322085)--(1.142240,-0.435958)--cycle;
\path[fill=blue!25!white,draw=blue!25!white,line width=.12pt] (-1.442454,-0.096998)--(-1.525177,-0.019647)--(-1.417944,-0.091940)--(-1.341037,-0.165370)--cycle;
\path[fill=blue!19!white,draw=blue!19!white,line width=.12pt] (1.341037,-0.165370)--(1.417944,-0.091940)--(1.525177,-0.019647)--(1.442454,-0.096998)--cycle;
\path[fill=blue!28!white,draw=blue!28!white,line width=.12pt] (-1.283890,-0.261700)--(-1.341037,-0.165370)--(-1.226706,-0.228443)--(-1.174431,-0.322085)--cycle;
\path[fill=blue!22!white,draw=blue!22!white,line width=.12pt] (1.174431,-0.322085)--(1.226706,-0.228443)--(1.341037,-0.165370)--(1.283890,-0.261700)--cycle;
\path[fill=blue!42!white,draw=blue!42!white,line width=.12pt] (-0.921837,-1.438377)--(-0.864690,-1.298758)--(-0.701962,-1.337241)--(-0.748354,-1.479402)--cycle;
\path[fill=blue!39!white,draw=blue!39!white,line width=.12pt] (0.748354,-1.479402)--(0.701962,-1.337241)--(0.864690,-1.298758)--(0.921837,-1.438377)--cycle;
\path[fill=blue!17!white,draw=blue!17!white,line width=.12pt] (-1.934468,0.239441)--(-2.073119,0.250151)--(-1.978968,0.150618)--(-1.846614,0.146565)--cycle;
\path[fill=blue!15!white,draw=blue!15!white,line width=.12pt] (1.846614,0.146565)--(1.978968,0.150618)--(2.073119,0.250151)--(1.934468,0.239441)--cycle;
\path[fill=blue!32!white,draw=blue!32!white,line width=.12pt] (-2.465877,-1.313721)--(-2.365079,-1.366207)--(-2.198794,-1.478311)--(-2.292506,-1.430603)--cycle;
\path[fill=blue!36!white,draw=blue!36!white,line width=.12pt] (2.292506,-1.430603)--(2.198794,-1.478311)--(2.365079,-1.366207)--(2.465877,-1.313721)--cycle;
\path[fill=blue!26!white,draw=blue!26!white,line width=.12pt] (-2.768442,-1.005012)--(-2.703244,-1.107367)--(-2.548600,-1.236371)--(-2.610069,-1.137127)--cycle;
\path[fill=blue!33!white,draw=blue!33!white,line width=.12pt] (2.610069,-1.137127)--(2.548600,-1.236371)--(2.703244,-1.107367)--(2.768442,-1.005012)--cycle;
\path[fill=blue!39!white,draw=blue!39!white,line width=.12pt] (-1.807416,-1.583216)--(-1.697056,-1.541348)--(-1.522544,-1.620433)--(-1.621555,-1.667443)--cycle;
\path[fill=blue!39!white,draw=blue!39!white,line width=.12pt] (1.621555,-1.667443)--(1.522544,-1.620433)--(1.697056,-1.541348)--(1.807416,-1.583216)--cycle;
\path[fill=blue!40!white,draw=blue!40!white,line width=.12pt] (-0.663889,-1.186295)--(-0.635598,-1.032365)--(-0.482133,-1.059820)--(-0.503593,-1.214972)--cycle;
\path[fill=blue!37!white,draw=blue!37!white,line width=.12pt] (0.503593,-1.214972)--(0.482133,-1.059820)--(0.635598,-1.032365)--(0.663889,-1.186295)--cycle;
\path[fill=blue!41!white,draw=blue!41!white,line width=.12pt] (-0.701962,-1.337241)--(-0.663889,-1.186295)--(-0.503593,-1.214972)--(-0.532473,-1.367563)--cycle;
\path[fill=blue!39!white,draw=blue!39!white,line width=.12pt] (0.532473,-1.367563)--(0.503593,-1.214972)--(0.663889,-1.186295)--(0.701962,-1.337241)--cycle;
\path[fill=blue!40!white,draw=blue!40!white,line width=.12pt] (-1.522544,-1.620433)--(-1.423533,-1.546798)--(-1.246659,-1.612387)--(-1.333369,-1.690584)--cycle;
\path[fill=blue!40!white,draw=blue!40!white,line width=.12pt] (1.333369,-1.690584)--(1.246659,-1.612387)--(1.423533,-1.546798)--(1.522544,-1.620433)--cycle;
\path[fill=blue!36!white,draw=blue!36!white,line width=.12pt] (-0.761481,-0.847478)--(-0.754235,-0.705537)--(-0.612293,-0.739104)--(-0.618176,-0.881367)--cycle;
\path[fill=blue!32!white,draw=blue!32!white,line width=.12pt] (0.618176,-0.881367)--(0.612293,-0.739104)--(0.754235,-0.705537)--(0.761481,-0.847478)--cycle;
\path[fill=blue!38!white,draw=blue!38!white,line width=.12pt] (-0.635598,-1.032365)--(-0.618176,-0.881367)--(-0.468918,-0.908070)--(-0.482133,-1.059820)--cycle;
\path[fill=blue!35!white,draw=blue!35!white,line width=.12pt] (0.482133,-1.059820)--(0.468918,-0.908070)--(0.618176,-0.881367)--(0.635598,-1.032365)--cycle;
\path[fill=blue!19!white,draw=blue!19!white,line width=.12pt] (-2.942213,-0.474190)--(-2.956415,-0.612293)--(-2.822148,-0.754235)--(-2.808591,-0.615449)--cycle;
\path[fill=blue!27!white,draw=blue!27!white,line width=.12pt] (2.808591,-0.615449)--(2.822148,-0.754235)--(2.956415,-0.612293)--(2.942213,-0.474190)--cycle;
\path[fill=blue!41!white,draw=blue!41!white,line width=.12pt] (-1.246659,-1.612387)--(-1.163283,-1.510570)--(-0.987036,-1.563389)--(-1.057780,-1.668992)--cycle;
\path[fill=blue!40!white,draw=blue!40!white,line width=.12pt] (1.057780,-1.668992)--(0.987036,-1.563389)--(1.163283,-1.510570)--(1.246659,-1.612387)--cycle;
\path[fill=blue!34!white,draw=blue!34!white,line width=.12pt] (-0.888912,-0.665176)--(-0.897452,-0.536404)--(-0.761481,-0.577153)--(-0.754235,-0.705537)--cycle;
\path[fill=blue!29!white,draw=blue!29!white,line width=.12pt] (0.754235,-0.705537)--(0.761481,-0.577153)--(0.897452,-0.536404)--(0.888912,-0.665176)--cycle;
\path[fill=blue!42!white,draw=blue!42!white,line width=.12pt] (-0.987036,-1.563389)--(-0.921837,-1.438377)--(-0.748354,-1.479402)--(-0.801283,-1.607317)--cycle;
\path[fill=blue!40!white,draw=blue!40!white,line width=.12pt] (0.801283,-1.607317)--(0.748354,-1.479402)--(0.921837,-1.438377)--(0.987036,-1.563389)--cycle;
\path[fill=blue!42!white,draw=blue!42!white,line width=.12pt] (-0.748354,-1.479402)--(-0.701962,-1.337241)--(-0.532473,-1.367563)--(-0.567664,-1.511728)--cycle;
\path[fill=blue!40!white,draw=blue!40!white,line width=.12pt] (0.567664,-1.511728)--(0.532473,-1.367563)--(0.701962,-1.337241)--(0.748354,-1.479402)--cycle;
\path[fill=blue!40!white,draw=blue!40!white,line width=.12pt] (-0.503593,-1.214972)--(-0.482133,-1.059820)--(-0.324025,-1.079622)--(-0.338447,-1.235655)--cycle;
\path[fill=blue!38!white,draw=blue!38!white,line width=.12pt] (0.338447,-1.235655)--(0.324025,-1.079622)--(0.482133,-1.059820)--(0.503593,-1.214972)--cycle;
\path[fill=blue!41!white,draw=blue!41!white,line width=.12pt] (-0.532473,-1.367563)--(-0.503593,-1.214972)--(-0.338447,-1.235655)--(-0.357857,-1.389432)--cycle;
\path[fill=blue!39!white,draw=blue!39!white,line width=.12pt] (0.357857,-1.389432)--(0.338447,-1.235655)--(0.503593,-1.214972)--(0.532473,-1.367563)--cycle;
\path[fill=blue!19!white,draw=blue!19!white,line width=.12pt] (-1.724636,0.115141)--(-1.846614,0.146565)--(-1.740975,0.058441)--(-1.625975,0.032838)--cycle;
\path[fill=blue!16!white,draw=blue!16!white,line width=.12pt] (1.625975,0.032838)--(1.740975,0.058441)--(1.846614,0.146565)--(1.724636,0.115141)--cycle;
\path[fill=blue!31!white,draw=blue!31!white,line width=.12pt] (-1.024781,-0.489187)--(-1.053662,-0.376814)--(-0.922745,-0.425362)--(-0.897452,-0.536404)--cycle;
\path[fill=blue!26!white,draw=blue!26!white,line width=.12pt] (0.897452,-0.536404)--(0.922745,-0.425362)--(1.053662,-0.376814)--(1.024781,-0.489187)--cycle;
\path[fill=blue!15!white,draw=blue!15!white,line width=.12pt] (-2.073119,0.250151)--(-2.217311,0.233600)--(-2.116611,0.127144)--(-1.978968,0.150618)--cycle;
\path[fill=blue!14!white,draw=blue!14!white,line width=.12pt] (1.978968,0.150618)--(2.116611,0.127144)--(2.217311,0.233600)--(2.073119,0.250151)--cycle;
\path[fill=blue!38!white,draw=blue!38!white,line width=.12pt] (-0.482133,-1.059820)--(-0.468918,-0.908070)--(-0.315143,-0.927329)--(-0.324025,-1.079622)--cycle;
\path[fill=blue!36!white,draw=blue!36!white,line width=.12pt] (0.324025,-1.079622)--(0.315143,-0.927329)--(0.468918,-0.908070)--(0.482133,-1.059820)--cycle;
\path[fill=blue!34!white,draw=blue!34!white,line width=.12pt] (-2.198794,-1.478311)--(-2.091880,-1.498463)--(-1.913535,-1.596850)--(-2.011334,-1.581726)--cycle;
\path[fill=blue!38!white,draw=blue!38!white,line width=.12pt] (2.011334,-1.581726)--(1.913535,-1.596850)--(2.091880,-1.498463)--(2.198794,-1.478311)--cycle;
\path[fill=blue!39!white,draw=blue!39!white,line width=.12pt] (-0.338447,-1.235655)--(-0.324025,-1.079622)--(-0.162796,-1.091580)--(-0.170043,-1.248146)--cycle;
\path[fill=blue!38!white,draw=blue!38!white,line width=.12pt] (0.170043,-1.248146)--(0.162796,-1.091580)--(0.324025,-1.079622)--(0.338447,-1.235655)--cycle;
\path[fill=blue!36!white,draw=blue!36!white,line width=.12pt] (-0.618176,-0.881367)--(-0.612293,-0.739104)--(-0.464455,-0.765552)--(-0.468918,-0.908070)--cycle;
\path[fill=blue!33!white,draw=blue!33!white,line width=.12pt] (0.468918,-0.908070)--(0.464455,-0.765552)--(0.612293,-0.739104)--(0.618176,-0.881367)--cycle;
\path[fill=blue!41!white,draw=blue!41!white,line width=.12pt] (-0.357857,-1.389432)--(-0.338447,-1.235655)--(-0.170043,-1.248146)--(-0.179794,-1.402639)--cycle;
\path[fill=blue!40!white,draw=blue!40!white,line width=.12pt] (0.179794,-1.402639)--(0.170043,-1.248146)--(0.338447,-1.235655)--(0.357857,-1.389432)--cycle;
\path[fill=blue!28!white,draw=blue!28!white,line width=.12pt] (-1.174431,-0.322085)--(-1.226706,-0.228443)--(-1.100561,-0.285608)--(-1.053662,-0.376814)--cycle;
\path[fill=blue!23!white,draw=blue!23!white,line width=.12pt] (1.053662,-0.376814)--(1.100561,-0.285608)--(1.226706,-0.228443)--(1.174431,-0.322085)--cycle;
\path[fill=blue!39!white,draw=blue!39!white,line width=.12pt] (-0.170043,-1.248146)--(-0.162796,-1.091580)--(-0.000000,-1.095579)--(-0.000000,-1.252323)--cycle;
\path[fill=blue!39!white,draw=blue!39!white,line width=.12pt] (-0.000000,-1.252323)--(-0.000000,-1.095579)--(0.162796,-1.091580)--(0.170043,-1.248146)--cycle;
\path[fill=blue!42!white,draw=blue!42!white,line width=.12pt] (-0.567664,-1.511728)--(-0.532473,-1.367563)--(-0.357857,-1.389432)--(-0.381508,-1.535043)--cycle;
\path[fill=blue!40!white,draw=blue!40!white,line width=.12pt] (0.381508,-1.535043)--(0.357857,-1.389432)--(0.532473,-1.367563)--(0.567664,-1.511728)--cycle;
\path[fill=blue!22!white,draw=blue!22!white,line width=.12pt] (-1.525177,-0.019647)--(-1.625975,0.032838)--(-1.511656,-0.044233)--(-1.417944,-0.091940)--cycle;
\path[fill=blue!18!white,draw=blue!18!white,line width=.12pt] (1.417944,-0.091940)--(1.511656,-0.044233)--(1.625975,0.032838)--(1.525177,-0.019647)--cycle;
\path[fill=blue!17!white,draw=blue!17!white,line width=.12pt] (-2.900154,-0.335511)--(-2.942213,-0.474190)--(-2.808591,-0.615449)--(-2.768442,-0.474751)--cycle;
\path[fill=blue!24!white,draw=blue!24!white,line width=.12pt] (2.768442,-0.474751)--(2.808591,-0.615449)--(2.942213,-0.474190)--(2.900154,-0.335511)--cycle;
\path[fill=blue!38!white,draw=blue!38!white,line width=.12pt] (-0.324025,-1.079622)--(-0.315143,-0.927329)--(-0.158334,-0.938959)--(-0.162796,-1.091580)--cycle;
\path[fill=blue!36!white,draw=blue!36!white,line width=.12pt] (0.162796,-1.091580)--(0.158334,-0.938959)--(0.315143,-0.927329)--(0.324025,-1.079622)--cycle;
\path[fill=blue!25!white,draw=blue!25!white,line width=.12pt] (-1.341037,-0.165370)--(-1.417944,-0.091940)--(-1.297056,-0.158630)--(-1.226706,-0.228443)--cycle;
\path[fill=blue!20!white,draw=blue!20!white,line width=.12pt] (1.226706,-0.228443)--(1.297056,-0.158630)--(1.417944,-0.091940)--(1.341037,-0.165370)--cycle;
\path[fill=blue!40!white,draw=blue!40!white,line width=.12pt] (-0.179794,-1.402639)--(-0.170043,-1.248146)--(-0.000000,-1.252323)--(-0.000000,-1.407055)--cycle;
\path[fill=blue!40!white,draw=blue!40!white,line width=.12pt] (-0.000000,-1.407055)--(-0.000000,-1.252323)--(0.170043,-1.248146)--(0.179794,-1.402639)--cycle;
\path[fill=blue!24!white,draw=blue!24!white,line width=.12pt] (-2.808591,-0.885774)--(-2.768442,-1.005012)--(-2.610069,-1.137127)--(-2.647922,-1.019805)--cycle;
\path[fill=blue!31!white,draw=blue!31!white,line width=.12pt] (2.647922,-1.019805)--(2.610069,-1.137127)--(2.768442,-1.005012)--(2.808591,-0.885774)--cycle;
\path[fill=blue!34!white,draw=blue!34!white,line width=.12pt] (-0.754235,-0.705537)--(-0.761481,-0.577153)--(-0.618176,-0.611042)--(-0.612293,-0.739104)--cycle;
\path[fill=blue!30!white,draw=blue!30!white,line width=.12pt] (0.612293,-0.739104)--(0.618176,-0.611042)--(0.761481,-0.577153)--(0.754235,-0.705537)--cycle;
\path[fill=blue!42!white,draw=blue!42!white,line width=.12pt] (-0.801283,-1.607317)--(-0.748354,-1.479402)--(-0.567664,-1.511728)--(-0.607814,-1.641929)--cycle;
\path[fill=blue!40!white,draw=blue!40!white,line width=.12pt] (0.607814,-1.641929)--(0.567664,-1.511728)--(0.748354,-1.479402)--(0.801283,-1.607317)--cycle;
\path[fill=blue!37!white,draw=blue!37!white,line width=.12pt] (-0.162796,-1.091580)--(-0.158334,-0.938959)--(-0.000000,-0.942848)--(-0.000000,-1.095579)--cycle;
\path[fill=blue!37!white,draw=blue!37!white,line width=.12pt] (-0.000000,-1.095579)--(-0.000000,-0.942848)--(0.158334,-0.938959)--(0.162796,-1.091580)--cycle;
\path[fill=blue!29!white,draw=blue!29!white,line width=.12pt] (-2.548600,-1.236371)--(-2.465877,-1.313721)--(-2.292506,-1.430603)--(-2.369413,-1.357174)--cycle;
\path[fill=blue!35!white,draw=blue!35!white,line width=.12pt] (2.369413,-1.357174)--(2.292506,-1.430603)--(2.465877,-1.313721)--(2.548600,-1.236371)--cycle;
\path[fill=blue!37!white,draw=blue!37!white,line width=.12pt] (-1.913535,-1.596850)--(-1.807416,-1.583216)--(-1.621555,-1.667443)--(-1.716761,-1.686022)--cycle;
\path[fill=blue!38!white,draw=blue!38!white,line width=.12pt] (1.716761,-1.686022)--(1.621555,-1.667443)--(1.807416,-1.583216)--(1.913535,-1.596850)--cycle;
\path[fill=blue!36!white,draw=blue!36!white,line width=.12pt] (-0.468918,-0.908070)--(-0.464455,-0.765552)--(-0.312145,-0.784628)--(-0.315143,-0.927329)--cycle;
\path[fill=blue!34!white,draw=blue!34!white,line width=.12pt] (0.315143,-0.927329)--(0.312145,-0.784628)--(0.464455,-0.765552)--(0.468918,-0.908070)--cycle;
\path[fill=blue!14!white,draw=blue!14!white,line width=.12pt] (-2.217311,0.233600)--(-2.361503,0.190425)--(-2.254254,0.077046)--(-2.116611,0.127144)--cycle;
\path[fill=blue!15!white,draw=blue!15!white,line width=.12pt] (2.116611,0.127144)--(2.254254,0.077046)--(2.361503,0.190425)--(2.217311,0.233600)--cycle;
\path[fill=blue!41!white,draw=blue!41!white,line width=.12pt] (-0.381508,-1.535043)--(-0.357857,-1.389432)--(-0.179794,-1.402639)--(-0.191677,-1.549123)--cycle;
\path[fill=blue!41!white,draw=blue!41!white,line width=.12pt] (0.191677,-1.549123)--(0.179794,-1.402639)--(0.357857,-1.389432)--(0.381508,-1.535043)--cycle;
\path[fill=blue!41!white,draw=blue!41!white,line width=.12pt] (-1.057780,-1.668992)--(-0.987036,-1.563389)--(-0.801283,-1.607317)--(-0.858714,-1.716067)--cycle;
\path[fill=blue!40!white,draw=blue!40!white,line width=.12pt] (0.858714,-1.716067)--(0.801283,-1.607317)--(0.987036,-1.563389)--(1.057780,-1.668992)--cycle;
\path[fill=blue!39!white,draw=blue!39!white,line width=.12pt] (-1.621555,-1.667443)--(-1.522544,-1.620433)--(-1.333369,-1.690584)--(-1.420078,-1.742156)--cycle;
\path[fill=blue!39!white,draw=blue!39!white,line width=.12pt] (1.420078,-1.742156)--(1.333369,-1.690584)--(1.522544,-1.620433)--(1.621555,-1.667443)--cycle;
\path[fill=blue!40!white,draw=blue!40!white,line width=.12pt] (-1.333369,-1.690584)--(-1.246659,-1.612387)--(-1.057780,-1.668992)--(-1.131352,-1.751126)--cycle;
\path[fill=blue!40!white,draw=blue!40!white,line width=.12pt] (1.131352,-1.751126)--(1.057780,-1.668992)--(1.246659,-1.612387)--(1.333369,-1.690584)--cycle;
\path[fill=blue!41!white,draw=blue!41!white,line width=.12pt] (-0.000000,-1.553831)--(-0.000000,-1.407055)--(0.179794,-1.402639)--(0.191677,-1.549123)--cycle;
\path[fill=blue!41!white,draw=blue!41!white,line width=.12pt] (-0.191677,-1.549123)--(-0.179794,-1.402639)--(-0.000000,-1.407055)--(-0.000000,-1.553831)--cycle;
\path[fill=blue!31!white,draw=blue!31!white,line width=.12pt] (-0.897452,-0.536404)--(-0.922745,-0.425362)--(-0.782941,-0.467259)--(-0.761481,-0.577153)--cycle;
\path[fill=blue!27!white,draw=blue!27!white,line width=.12pt] (0.761481,-0.577153)--(0.782941,-0.467259)--(0.922745,-0.425362)--(0.897452,-0.536404)--cycle;
\path[fill=blue!15!white,draw=blue!15!white,line width=.12pt] (-2.831853,-0.201586)--(-2.900154,-0.335511)--(-2.768442,-0.474751)--(-2.703244,-0.337547)--cycle;
\path[fill=blue!22!white,draw=blue!22!white,line width=.12pt] (2.703244,-0.337547)--(2.768442,-0.474751)--(2.900154,-0.335511)--(2.831853,-0.201586)--cycle;
\path[fill=blue!36!white,draw=blue!36!white,line width=.12pt] (-0.315143,-0.927329)--(-0.312145,-0.784628)--(-0.156827,-0.796148)--(-0.158334,-0.938959)--cycle;
\path[fill=blue!34!white,draw=blue!34!white,line width=.12pt] (0.158334,-0.938959)--(0.156827,-0.796148)--(0.312145,-0.784628)--(0.315143,-0.927329)--cycle;
\path[fill=blue!17!white,draw=blue!17!white,line width=.12pt] (-1.846614,0.146565)--(-1.978968,0.150618)--(-1.865758,0.056178)--(-1.740975,0.058441)--cycle;
\path[fill=blue!15!white,draw=blue!15!white,line width=.12pt] (1.740975,0.058441)--(1.865758,0.056178)--(1.978968,0.150618)--(1.846614,0.146565)--cycle;
\path[fill=blue!42!white,draw=blue!42!white,line width=.12pt] (-0.607814,-1.641929)--(-0.567664,-1.511728)--(-0.381508,-1.535043)--(-0.408491,-1.666893)--cycle;
\path[fill=blue!41!white,draw=blue!41!white,line width=.12pt] (0.408491,-1.666893)--(0.381508,-1.535043)--(0.567664,-1.511728)--(0.607814,-1.641929)--cycle;
\path[fill=blue!34!white,draw=blue!34!white,line width=.12pt] (-0.612293,-0.739104)--(-0.618176,-0.611042)--(-0.468918,-0.637745)--(-0.464455,-0.765552)--cycle;
\path[fill=blue!31!white,draw=blue!31!white,line width=.12pt] (0.464455,-0.765552)--(0.468918,-0.637745)--(0.618176,-0.611042)--(0.612293,-0.739104)--cycle;
\path[fill=blue!35!white,draw=blue!35!white,line width=.12pt] (-0.158334,-0.938959)--(-0.156827,-0.796148)--(-0.000000,-0.800000)--(-0.000000,-0.942848)--cycle;
\path[fill=blue!35!white,draw=blue!35!white,line width=.12pt] (-0.000000,-0.942848)--(-0.000000,-0.800000)--(0.156827,-0.796148)--(0.158334,-0.938959)--cycle;
\path[fill=blue!13!white,draw=blue!13!white,line width=.12pt] (-2.361503,0.190425)--(-2.500154,0.122284)--(-2.386608,0.002248)--(-2.254254,0.077046)--cycle;
\path[fill=blue!15!white,draw=blue!15!white,line width=.12pt] (2.254254,0.077046)--(2.386608,0.002248)--(2.500154,0.122284)--(2.361503,0.190425)--cycle;
\path[fill=blue!28!white,draw=blue!28!white,line width=.12pt] (-1.053662,-0.376814)--(-1.100561,-0.285608)--(-0.963817,-0.336316)--(-0.922745,-0.425362)--cycle;
\path[fill=blue!23!white,draw=blue!23!white,line width=.12pt] (0.922745,-0.425362)--(0.963817,-0.336316)--(1.100561,-0.285608)--(1.053662,-0.376814)--cycle;
\path[fill=blue!14!white,draw=blue!14!white,line width=.12pt] (-2.739936,-0.077561)--(-2.831853,-0.201586)--(-2.703244,-0.337547)--(-2.615501,-0.209109)--cycle;
\path[fill=blue!20!white,draw=blue!20!white,line width=.12pt] (2.615501,-0.209109)--(2.703244,-0.337547)--(2.831853,-0.201586)--(2.739936,-0.077561)--cycle;
\path[fill=blue!13!white,draw=blue!13!white,line width=.12pt] (-2.500154,0.122284)--(-2.627935,0.031796)--(-2.508586,-0.094375)--(-2.386608,0.002248)--cycle;
\path[fill=blue!16!white,draw=blue!16!white,line width=.12pt] (2.386608,0.002248)--(2.508586,-0.094375)--(2.627935,0.031796)--(2.500154,0.122284)--cycle;
\path[fill=blue!32!white,draw=blue!32!white,line width=.12pt] (-2.292506,-1.430603)--(-2.198794,-1.478311)--(-2.011334,-1.581726)--(-2.097056,-1.538426)--cycle;
\path[fill=blue!36!white,draw=blue!36!white,line width=.12pt] (2.097056,-1.538426)--(2.011334,-1.581726)--(2.198794,-1.478311)--(2.292506,-1.430603)--cycle;
\path[fill=blue!21!white,draw=blue!21!white,line width=.12pt] (-2.822148,-0.754235)--(-2.808591,-0.885774)--(-2.647922,-1.019805)--(-2.660703,-0.888912)--cycle;
\path[fill=blue!29!white,draw=blue!29!white,line width=.12pt] (2.660703,-0.888912)--(2.647922,-1.019805)--(2.808591,-0.885774)--(2.822148,-0.754235)--cycle;
\path[fill=blue!42!white,draw=blue!42!white,line width=.12pt] (-0.408491,-1.666893)--(-0.381508,-1.535043)--(-0.191677,-1.549123)--(-0.205234,-1.681968)--cycle;
\path[fill=blue!41!white,draw=blue!41!white,line width=.12pt] (0.205234,-1.681968)--(0.191677,-1.549123)--(0.381508,-1.535043)--(0.408491,-1.666893)--cycle;
\path[fill=blue!13!white,draw=blue!13!white,line width=.12pt] (-2.627935,0.031796)--(-2.739936,-0.077561)--(-2.615501,-0.209109)--(-2.508586,-0.094375)--cycle;
\path[fill=blue!18!white,draw=blue!18!white,line width=.12pt] (2.508586,-0.094375)--(2.615501,-0.209109)--(2.739936,-0.077561)--(2.627935,0.031796)--cycle;
\path[fill=blue!41!white,draw=blue!41!white,line width=.12pt] (-0.858714,-1.716067)--(-0.801283,-1.607317)--(-0.607814,-1.641929)--(-0.651378,-1.753160)--cycle;
\path[fill=blue!40!white,draw=blue!40!white,line width=.12pt] (0.651378,-1.753160)--(0.607814,-1.641929)--(0.801283,-1.607317)--(0.858714,-1.716067)--cycle;
\path[fill=blue!33!white,draw=blue!33!white,line width=.12pt] (-0.464455,-0.765552)--(-0.468918,-0.637745)--(-0.315143,-0.657004)--(-0.312145,-0.784628)--cycle;
\path[fill=blue!31!white,draw=blue!31!white,line width=.12pt] (0.312145,-0.784628)--(0.315143,-0.657004)--(0.468918,-0.637745)--(0.464455,-0.765552)--cycle;
\path[fill=blue!19!white,draw=blue!19!white,line width=.12pt] (-1.625975,0.032838)--(-1.740975,0.058441)--(-1.618570,-0.024081)--(-1.511656,-0.044233)--cycle;
\path[fill=blue!16!white,draw=blue!16!white,line width=.12pt] (1.511656,-0.044233)--(1.618570,-0.024081)--(1.740975,0.058441)--(1.625975,0.032838)--cycle;
\path[fill=blue!25!white,draw=blue!25!white,line width=.12pt] (-1.226706,-0.228443)--(-1.297056,-0.158630)--(-1.163677,-0.219074)--(-1.100561,-0.285608)--cycle;
\path[fill=blue!21!white,draw=blue!21!white,line width=.12pt] (1.100561,-0.285608)--(1.163677,-0.219074)--(1.297056,-0.158630)--(1.226706,-0.228443)--cycle;
\path[fill=blue!31!white,draw=blue!31!white,line width=.12pt] (-0.761481,-0.577153)--(-0.782941,-0.467259)--(-0.635598,-0.502103)--(-0.618176,-0.611042)--cycle;
\path[fill=blue!27!white,draw=blue!27!white,line width=.12pt] (0.618176,-0.611042)--(0.635598,-0.502103)--(0.782941,-0.467259)--(0.761481,-0.577153)--cycle;
\path[fill=blue!41!white,draw=blue!41!white,line width=.12pt] (-0.205234,-1.681968)--(-0.191677,-1.549123)--(-0.000000,-1.553831)--(-0.000000,-1.687009)--cycle;
\path[fill=blue!41!white,draw=blue!41!white,line width=.12pt] (-0.000000,-1.687009)--(-0.000000,-1.553831)--(0.191677,-1.549123)--(0.205234,-1.681968)--cycle;
\path[fill=blue!22!white,draw=blue!22!white,line width=.12pt] (-1.417944,-0.091940)--(-1.511656,-0.044233)--(-1.382778,-0.115330)--(-1.297056,-0.158630)--cycle;
\path[fill=blue!18!white,draw=blue!18!white,line width=.12pt] (1.297056,-0.158630)--(1.382778,-0.115330)--(1.511656,-0.044233)--(1.417944,-0.091940)--cycle;
\path[fill=blue!26!white,draw=blue!26!white,line width=.12pt] (-2.610069,-1.137127)--(-2.548600,-1.236371)--(-2.369413,-1.357174)--(-2.426560,-1.260844)--cycle;
\path[fill=blue!33!white,draw=blue!33!white,line width=.12pt] (2.426560,-1.260844)--(2.369413,-1.357174)--(2.548600,-1.236371)--(2.610069,-1.137127)--cycle;
\path[fill=blue!33!white,draw=blue!33!white,line width=.12pt] (-0.312145,-0.784628)--(-0.315143,-0.657004)--(-0.158334,-0.668634)--(-0.156827,-0.796148)--cycle;
\path[fill=blue!32!white,draw=blue!32!white,line width=.12pt] (0.156827,-0.796148)--(0.158334,-0.668634)--(0.315143,-0.657004)--(0.312145,-0.784628)--cycle;
\path[fill=blue!40!white,draw=blue!40!white,line width=.12pt] (-1.131352,-1.751126)--(-1.057780,-1.668992)--(-0.858714,-1.716067)--(-0.918440,-1.801476)--cycle;
\path[fill=blue!40!white,draw=blue!40!white,line width=.12pt] (0.918440,-1.801476)--(0.858714,-1.716067)--(1.057780,-1.668992)--(1.131352,-1.751126)--cycle;
\path[fill=blue!34!white,draw=blue!34!white,line width=.12pt] (-2.011334,-1.581726)--(-1.913535,-1.596850)--(-1.716761,-1.686022)--(-1.804504,-1.675456)--cycle;
\path[fill=blue!37!white,draw=blue!37!white,line width=.12pt] (1.804504,-1.675456)--(1.716761,-1.686022)--(1.913535,-1.596850)--(2.011334,-1.581726)--cycle;
\path[fill=blue!33!white,draw=blue!33!white,line width=.12pt] (-0.156827,-0.796148)--(-0.158334,-0.668634)--(-0.000000,-0.672523)--(-0.000000,-0.800000)--cycle;
\path[fill=blue!32!white,draw=blue!32!white,line width=.12pt] (-0.000000,-0.800000)--(-0.000000,-0.672523)--(0.158334,-0.668634)--(0.156827,-0.796148)--cycle;
\path[fill=blue!15!white,draw=blue!15!white,line width=.12pt] (-1.978968,0.150618)--(-2.116611,0.127144)--(-1.995527,0.026136)--(-1.865758,0.056178)--cycle;
\path[fill=blue!14!white,draw=blue!14!white,line width=.12pt] (1.865758,0.056178)--(1.995527,0.026136)--(2.116611,0.127144)--(1.978968,0.150618)--cycle;
\path[fill=blue!41!white,draw=blue!41!white,line width=.12pt] (-0.651378,-1.753160)--(-0.607814,-1.641929)--(-0.408491,-1.666893)--(-0.437769,-1.779914)--cycle;
\path[fill=blue!41!white,draw=blue!41!white,line width=.12pt] (0.437769,-1.779914)--(0.408491,-1.666893)--(0.607814,-1.641929)--(0.651378,-1.753160)--cycle;
\path[fill=blue!39!white,draw=blue!39!white,line width=.12pt] (-1.420078,-1.742156)--(-1.333369,-1.690584)--(-1.131352,-1.751126)--(-1.204924,-1.806635)--cycle;
\path[fill=blue!39!white,draw=blue!39!white,line width=.12pt] (1.204924,-1.806635)--(1.131352,-1.751126)--(1.333369,-1.690584)--(1.420078,-1.742156)--cycle;
\path[fill=blue!37!white,draw=blue!37!white,line width=.12pt] (-1.716761,-1.686022)--(-1.621555,-1.667443)--(-1.420078,-1.742156)--(-1.503455,-1.765122)--cycle;
\path[fill=blue!38!white,draw=blue!38!white,line width=.12pt] (1.503455,-1.765122)--(1.420078,-1.742156)--(1.621555,-1.667443)--(1.716761,-1.686022)--cycle;
\path[fill=blue!28!white,draw=blue!28!white,line width=.12pt] (-0.922745,-0.425362)--(-0.963817,-0.336316)--(-0.817791,-0.380079)--(-0.782941,-0.467259)--cycle;
\path[fill=blue!24!white,draw=blue!24!white,line width=.12pt] (0.782941,-0.467259)--(0.817791,-0.380079)--(0.963817,-0.336316)--(0.922745,-0.425362)--cycle;
\path[fill=blue!31!white,draw=blue!31!white,line width=.12pt] (-0.618176,-0.611042)--(-0.635598,-0.502103)--(-0.482133,-0.529559)--(-0.468918,-0.637745)--cycle;
\path[fill=blue!28!white,draw=blue!28!white,line width=.12pt] (0.468918,-0.637745)--(0.482133,-0.529559)--(0.635598,-0.502103)--(0.618176,-0.611042)--cycle;
\path[fill=blue!18!white,draw=blue!18!white,line width=.12pt] (-2.808591,-0.615449)--(-2.822148,-0.754235)--(-2.660703,-0.888912)--(-2.647922,-0.749480)--cycle;
\path[fill=blue!26!white,draw=blue!26!white,line width=.12pt] (2.647922,-0.749480)--(2.660703,-0.888912)--(2.822148,-0.754235)--(2.808591,-0.615449)--cycle;
\path[fill=blue!41!white,draw=blue!41!white,line width=.12pt] (-0.437769,-1.779914)--(-0.408491,-1.666893)--(-0.205234,-1.681968)--(-0.219943,-1.796069)--cycle;
\path[fill=blue!41!white,draw=blue!41!white,line width=.12pt] (0.219943,-1.796069)--(0.205234,-1.681968)--(0.408491,-1.666893)--(0.437769,-1.779914)--cycle;
\path[fill=blue!31!white,draw=blue!31!white,line width=.12pt] (-0.468918,-0.637745)--(-0.482133,-0.529559)--(-0.324025,-0.549360)--(-0.315143,-0.657004)--cycle;
\path[fill=blue!29!white,draw=blue!29!white,line width=.12pt] (0.315143,-0.657004)--(0.324025,-0.549360)--(0.482133,-0.529559)--(0.468918,-0.637745)--cycle;
\path[fill=blue!25!white,draw=blue!25!white,line width=.12pt] (-1.100561,-0.285608)--(-1.163677,-0.219074)--(-1.019091,-0.272690)--(-0.963817,-0.336316)--cycle;
\path[fill=blue!21!white,draw=blue!21!white,line width=.12pt] (0.963817,-0.336316)--(1.019091,-0.272690)--(1.163677,-0.219074)--(1.100561,-0.285608)--cycle;
\path[fill=blue!41!white,draw=blue!41!white,line width=.12pt] (-0.219943,-1.796069)--(-0.205234,-1.681968)--(-0.000000,-1.687009)--(-0.000000,-1.801472)--cycle;
\path[fill=blue!41!white,draw=blue!41!white,line width=.12pt] (-0.000000,-1.801472)--(-0.000000,-1.687009)--(0.205234,-1.681968)--(0.219943,-1.796069)--cycle;
\path[fill=blue!40!white,draw=blue!40!white,line width=.12pt] (-0.918440,-1.801476)--(-0.858714,-1.716067)--(-0.651378,-1.753160)--(-0.696683,-1.841149)--cycle;
\path[fill=blue!40!white,draw=blue!40!white,line width=.12pt] (0.696683,-1.841149)--(0.651378,-1.753160)--(0.858714,-1.716067)--(0.918440,-1.801476)--cycle;
\path[fill=blue!14!white,draw=blue!14!white,line width=.12pt] (-2.116611,0.127144)--(-2.254254,0.077046)--(-2.125296,-0.030531)--(-1.995527,0.026136)--cycle;
\path[fill=blue!14!white,draw=blue!14!white,line width=.12pt] (1.995527,0.026136)--(2.125296,-0.030531)--(2.254254,0.077046)--(2.116611,0.127144)--cycle;
\path[fill=blue!17!white,draw=blue!17!white,line width=.12pt] (-1.740975,0.058441)--(-1.865758,0.056178)--(-1.734580,-0.032258)--(-1.618570,-0.024081)--cycle;
\path[fill=blue!15!white,draw=blue!15!white,line width=.12pt] (1.618570,-0.024081)--(1.734580,-0.032258)--(1.865758,0.056178)--(1.740975,0.058441)--cycle;
\path[fill=blue!29!white,draw=blue!29!white,line width=.12pt] (-2.369413,-1.357174)--(-2.292506,-1.430603)--(-2.097056,-1.538426)--(-2.167407,-1.468614)--cycle;
\path[fill=blue!34!white,draw=blue!34!white,line width=.12pt] (2.167407,-1.468614)--(2.097056,-1.538426)--(2.292506,-1.430603)--(2.369413,-1.357174)--cycle;
\path[fill=blue!23!white,draw=blue!23!white,line width=.12pt] (-2.647922,-1.019805)--(-2.610069,-1.137127)--(-2.426560,-1.260844)--(-2.461751,-1.145316)--cycle;
\path[fill=blue!30!white,draw=blue!30!white,line width=.12pt] (2.461751,-1.145316)--(2.426560,-1.260844)--(2.610069,-1.137127)--(2.647922,-1.019805)--cycle;
\path[fill=blue!28!white,draw=blue!28!white,line width=.12pt] (-0.782941,-0.467259)--(-0.817791,-0.380079)--(-0.663889,-0.416474)--(-0.635598,-0.502103)--cycle;
\path[fill=blue!25!white,draw=blue!25!white,line width=.12pt] (0.635598,-0.502103)--(0.663889,-0.416474)--(0.817791,-0.380079)--(0.782941,-0.467259)--cycle;
\path[fill=blue!30!white,draw=blue!30!white,line width=.12pt] (-0.315143,-0.657004)--(-0.324025,-0.549360)--(-0.162796,-0.561318)--(-0.158334,-0.668634)--cycle;
\path[fill=blue!29!white,draw=blue!29!white,line width=.12pt] (0.158334,-0.668634)--(0.162796,-0.561318)--(0.324025,-0.549360)--(0.315143,-0.657004)--cycle;
\path[fill=blue!16!white,draw=blue!16!white,line width=.12pt] (-2.768442,-0.474751)--(-2.808591,-0.615449)--(-2.647922,-0.749480)--(-2.610069,-0.606865)--cycle;
\path[fill=blue!24!white,draw=blue!24!white,line width=.12pt] (2.610069,-0.606865)--(2.647922,-0.749480)--(2.808591,-0.615449)--(2.768442,-0.474751)--cycle;
\path[fill=blue!22!white,draw=blue!22!white,line width=.12pt] (-1.297056,-0.158630)--(-1.382778,-0.115330)--(-1.240584,-0.179769)--(-1.163677,-0.219074)--cycle;
\path[fill=blue!19!white,draw=blue!19!white,line width=.12pt] (1.163677,-0.219074)--(1.240584,-0.179769)--(1.382778,-0.115330)--(1.297056,-0.158630)--cycle;
\path[fill=blue!20!white,draw=blue!20!white,line width=.12pt] (-1.511656,-0.044233)--(-1.618570,-0.024081)--(-1.480578,-0.100206)--(-1.382778,-0.115330)--cycle;
\path[fill=blue!16!white,draw=blue!16!white,line width=.12pt] (1.382778,-0.115330)--(1.480578,-0.100206)--(1.618570,-0.024081)--(1.511656,-0.044233)--cycle;
\path[fill=blue!30!white,draw=blue!30!white,line width=.12pt] (-0.158334,-0.668634)--(-0.162796,-0.561318)--(-0.000000,-0.565317)--(-0.000000,-0.672523)--cycle;
\path[fill=blue!30!white,draw=blue!30!white,line width=.12pt] (-0.000000,-0.672523)--(-0.000000,-0.565317)--(0.162796,-0.561318)--(0.158334,-0.668634)--cycle;
\path[fill=blue!39!white,draw=blue!39!white,line width=.12pt] (-1.204924,-1.806635)--(-1.131352,-1.751126)--(-0.918440,-1.801476)--(-0.978167,-1.860259)--cycle;
\path[fill=blue!39!white,draw=blue!39!white,line width=.12pt] (0.978167,-1.860259)--(0.918440,-1.801476)--(1.131352,-1.751126)--(1.204924,-1.806635)--cycle;
\path[fill=blue!13!white,draw=blue!13!white,line width=.12pt] (-2.254254,0.077046)--(-2.386608,0.002248)--(-2.250079,-0.111645)--(-2.125296,-0.030531)--cycle;
\path[fill=blue!15!white,draw=blue!15!white,line width=.12pt] (2.125296,-0.030531)--(2.250079,-0.111645)--(2.386608,0.002248)--(2.254254,0.077046)--cycle;
\path[fill=blue!40!white,draw=blue!40!white,line width=.12pt] (-0.696683,-1.841149)--(-0.651378,-1.753160)--(-0.437769,-1.779914)--(-0.468217,-1.869763)--cycle;
\path[fill=blue!40!white,draw=blue!40!white,line width=.12pt] (0.468217,-1.869763)--(0.437769,-1.779914)--(0.651378,-1.753160)--(0.696683,-1.841149)--cycle;
\path[fill=blue!32!white,draw=blue!32!white,line width=.12pt] (-2.097056,-1.538426)--(-2.011334,-1.581726)--(-1.804504,-1.675456)--(-1.881411,-1.636151)--cycle;
\path[fill=blue!35!white,draw=blue!35!white,line width=.12pt] (1.881411,-1.636151)--(1.804504,-1.675456)--(2.011334,-1.581726)--(2.097056,-1.538426)--cycle;
\path[fill=blue!15!white,draw=blue!15!white,line width=.12pt] (-2.703244,-0.337547)--(-2.768442,-0.474751)--(-2.610069,-0.606865)--(-2.548600,-0.466550)--cycle;
\path[fill=blue!21!white,draw=blue!21!white,line width=.12pt] (2.548600,-0.466550)--(2.610069,-0.606865)--(2.768442,-0.474751)--(2.703244,-0.337547)--cycle;
\path[fill=blue!28!white,draw=blue!28!white,line width=.12pt] (-0.635598,-0.502103)--(-0.663889,-0.416474)--(-0.503593,-0.445151)--(-0.482133,-0.529559)--cycle;
\path[fill=blue!25!white,draw=blue!25!white,line width=.12pt] (0.482133,-0.529559)--(0.503593,-0.445151)--(0.663889,-0.416474)--(0.635598,-0.502103)--cycle;
\path[fill=blue!37!white,draw=blue!37!white,line width=.12pt] (-1.503455,-1.765122)--(-1.420078,-1.742156)--(-1.204924,-1.806635)--(-1.275669,-1.833387)--cycle;
\path[fill=blue!38!white,draw=blue!38!white,line width=.12pt] (1.275669,-1.833387)--(1.204924,-1.806635)--(1.420078,-1.742156)--(1.503455,-1.765122)--cycle;
\path[fill=blue!25!white,draw=blue!25!white,line width=.12pt] (-0.963817,-0.336316)--(-1.019091,-0.272690)--(-0.864690,-0.318963)--(-0.817791,-0.380079)--cycle;
\path[fill=blue!22!white,draw=blue!22!white,line width=.12pt] (0.817791,-0.380079)--(0.864690,-0.318963)--(1.019091,-0.272690)--(0.963817,-0.336316)--cycle;
\path[fill=blue!34!white,draw=blue!34!white,line width=.12pt] (-1.804504,-1.675456)--(-1.716761,-1.686022)--(-1.503455,-1.765122)--(-1.580295,-1.758598)--cycle;
\path[fill=blue!37!white,draw=blue!37!white,line width=.12pt] (1.580295,-1.758598)--(1.503455,-1.765122)--(1.716761,-1.686022)--(1.804504,-1.675456)--cycle;
\path[fill=blue!12!white,draw=blue!12!white,line width=.12pt] (-2.386608,0.002248)--(-2.508586,-0.094375)--(-2.365079,-0.214089)--(-2.250079,-0.111645)--cycle;
\path[fill=blue!16!white,draw=blue!16!white,line width=.12pt] (2.250079,-0.111645)--(2.365079,-0.214089)--(2.508586,-0.094375)--(2.386608,0.002248)--cycle;
\path[fill=blue!40!white,draw=blue!40!white,line width=.12pt] (-0.468217,-1.869763)--(-0.437769,-1.779914)--(-0.219943,-1.796069)--(-0.235241,-1.887042)--cycle;
\path[fill=blue!40!white,draw=blue!40!white,line width=.12pt] (0.235241,-1.887042)--(0.219943,-1.796069)--(0.437769,-1.779914)--(0.468217,-1.869763)--cycle;
\path[fill=blue!15!white,draw=blue!15!white,line width=.12pt] (-1.865758,0.056178)--(-1.995527,0.026136)--(-1.855225,-0.068452)--(-1.734580,-0.032258)--cycle;
\path[fill=blue!14!white,draw=blue!14!white,line width=.12pt] (1.734580,-0.032258)--(1.855225,-0.068452)--(1.995527,0.026136)--(1.865758,0.056178)--cycle;
\path[fill=blue!13!white,draw=blue!13!white,line width=.12pt] (-2.615501,-0.209109)--(-2.703244,-0.337547)--(-2.548600,-0.466550)--(-2.465877,-0.333925)--cycle;
\path[fill=blue!19!white,draw=blue!19!white,line width=.12pt] (2.465877,-0.333925)--(2.548600,-0.466550)--(2.703244,-0.337547)--(2.615501,-0.209109)--cycle;
\path[fill=blue!13!white,draw=blue!13!white,line width=.12pt] (-2.508586,-0.094375)--(-2.615501,-0.209109)--(-2.465877,-0.333925)--(-2.365079,-0.214089)--cycle;
\path[fill=blue!17!white,draw=blue!17!white,line width=.12pt] (2.365079,-0.214089)--(2.465877,-0.333925)--(2.615501,-0.209109)--(2.508586,-0.094375)--cycle;
\path[fill=blue!21!white,draw=blue!21!white,line width=.12pt] (-2.660703,-0.888912)--(-2.647922,-1.019805)--(-2.461751,-1.145316)--(-2.473633,-1.015029)--cycle;
\path[fill=blue!28!white,draw=blue!28!white,line width=.12pt] (2.473633,-1.015029)--(2.461751,-1.145316)--(2.647922,-1.019805)--(2.660703,-0.888912)--cycle;
\path[fill=blue!40!white,draw=blue!40!white,line width=.12pt] (-0.235241,-1.887042)--(-0.219943,-1.796069)--(-0.000000,-1.801472)--(-0.000000,-1.892820)--cycle;
\path[fill=blue!40!white,draw=blue!40!white,line width=.12pt] (-0.000000,-1.892820)--(-0.000000,-1.801472)--(0.219943,-1.796069)--(0.235241,-1.887042)--cycle;
\path[fill=blue!28!white,draw=blue!28!white,line width=.12pt] (-0.482133,-0.529559)--(-0.503593,-0.445151)--(-0.338447,-0.465835)--(-0.324025,-0.549360)--cycle;
\path[fill=blue!26!white,draw=blue!26!white,line width=.12pt] (0.324025,-0.549360)--(0.338447,-0.465835)--(0.503593,-0.445151)--(0.482133,-0.529559)--cycle;
\path[fill=blue!26!white,draw=blue!26!white,line width=.12pt] (-2.426560,-1.260844)--(-2.369413,-1.357174)--(-2.167407,-1.468614)--(-2.219681,-1.374972)--cycle;
\path[fill=blue!32!white,draw=blue!32!white,line width=.12pt] (2.219681,-1.374972)--(2.167407,-1.468614)--(2.369413,-1.357174)--(2.426560,-1.260844)--cycle;
\path[fill=blue!39!white,draw=blue!39!white,line width=.12pt] (-0.978167,-1.860259)--(-0.918440,-1.801476)--(-0.696683,-1.841149)--(-0.741989,-1.902512)--cycle;
\path[fill=blue!39!white,draw=blue!39!white,line width=.12pt] (0.741989,-1.902512)--(0.696683,-1.841149)--(0.918440,-1.801476)--(0.978167,-1.860259)--cycle;
\path[fill=blue!22!white,draw=blue!22!white,line width=.12pt] (-1.163677,-0.219074)--(-1.240584,-0.179769)--(-1.086442,-0.236929)--(-1.019091,-0.272690)--cycle;
\path[fill=blue!19!white,draw=blue!19!white,line width=.12pt] (1.019091,-0.272690)--(1.086442,-0.236929)--(1.240584,-0.179769)--(1.163677,-0.219074)--cycle;
\path[fill=blue!28!white,draw=blue!28!white,line width=.12pt] (-0.324025,-0.549360)--(-0.338447,-0.465835)--(-0.170043,-0.478325)--(-0.162796,-0.561318)--cycle;
\path[fill=blue!26!white,draw=blue!26!white,line width=.12pt] (0.162796,-0.561318)--(0.170043,-0.478325)--(0.338447,-0.465835)--(0.324025,-0.549360)--cycle;
\path[fill=blue!17!white,draw=blue!17!white,line width=.12pt] (-1.618570,-0.024081)--(-1.734580,-0.032258)--(-1.586697,-0.113840)--(-1.480578,-0.100206)--cycle;
\path[fill=blue!15!white,draw=blue!15!white,line width=.12pt] (1.480578,-0.100206)--(1.586697,-0.113840)--(1.734580,-0.032258)--(1.618570,-0.024081)--cycle;
\path[fill=blue!25!white,draw=blue!25!white,line width=.12pt] (-0.817791,-0.380079)--(-0.864690,-0.318963)--(-0.701962,-0.357445)--(-0.663889,-0.416474)--cycle;
\path[fill=blue!22!white,draw=blue!22!white,line width=.12pt] (0.663889,-0.416474)--(0.701962,-0.357445)--(0.864690,-0.318963)--(0.817791,-0.380079)--cycle;
\path[fill=blue!20!white,draw=blue!20!white,line width=.12pt] (-1.382778,-0.115330)--(-1.480578,-0.100206)--(-1.328326,-0.169203)--(-1.240584,-0.179769)--cycle;
\path[fill=blue!17!white,draw=blue!17!white,line width=.12pt] (1.240584,-0.179769)--(1.328326,-0.169203)--(1.480578,-0.100206)--(1.382778,-0.115330)--cycle;
\path[fill=blue!27!white,draw=blue!27!white,line width=.12pt] (-0.162796,-0.561318)--(-0.170043,-0.478325)--(-0.000000,-0.482502)--(-0.000000,-0.565317)--cycle;
\path[fill=blue!27!white,draw=blue!27!white,line width=.12pt] (-0.000000,-0.565317)--(-0.000000,-0.482502)--(0.170043,-0.478325)--(0.162796,-0.561318)--cycle;
\path[fill=blue!14!white,draw=blue!14!white,line width=.12pt] (-1.995527,0.026136)--(-2.125296,-0.030531)--(-1.975871,-0.131270)--(-1.855225,-0.068452)--cycle;
\path[fill=blue!14!white,draw=blue!14!white,line width=.12pt] (1.855225,-0.068452)--(1.975871,-0.131270)--(2.125296,-0.030531)--(1.995527,0.026136)--cycle;
\path[fill=blue!37!white,draw=blue!37!white,line width=.12pt] (-1.275669,-1.833387)--(-1.204924,-1.806635)--(-0.978167,-1.860259)--(-1.035598,-1.890159)--cycle;
\path[fill=blue!38!white,draw=blue!38!white,line width=.12pt] (1.035598,-1.890159)--(0.978167,-1.860259)--(1.204924,-1.806635)--(1.275669,-1.833387)--cycle;
\path[fill=blue!29!white,draw=blue!29!white,line width=.12pt] (-2.167407,-1.468614)--(-2.097056,-1.538426)--(-1.881411,-1.636151)--(-1.944527,-1.569617)--cycle;
\path[fill=blue!33!white,draw=blue!33!white,line width=.12pt] (1.944527,-1.569617)--(1.881411,-1.636151)--(2.097056,-1.538426)--(2.167407,-1.468614)--cycle;
\path[fill=blue!39!white,draw=blue!39!white,line width=.12pt] (-0.741989,-1.902512)--(-0.696683,-1.841149)--(-0.468217,-1.869763)--(-0.498665,-1.932987)--cycle;
\path[fill=blue!39!white,draw=blue!39!white,line width=.12pt] (0.498665,-1.932987)--(0.468217,-1.869763)--(0.696683,-1.841149)--(0.741989,-1.902512)--cycle;
\path[fill=blue!18!white,draw=blue!18!white,line width=.12pt] (-2.647922,-0.749480)--(-2.660703,-0.888912)--(-2.473633,-1.015029)--(-2.461751,-0.874991)--cycle;
\path[fill=blue!25!white,draw=blue!25!white,line width=.12pt] (2.461751,-0.874991)--(2.473633,-1.015029)--(2.660703,-0.888912)--(2.647922,-0.749480)--cycle;
\path[fill=blue!25!white,draw=blue!25!white,line width=.12pt] (-0.663889,-0.416474)--(-0.701962,-0.357445)--(-0.532473,-0.387767)--(-0.503593,-0.445151)--cycle;
\path[fill=blue!23!white,draw=blue!23!white,line width=.12pt] (0.503593,-0.445151)--(0.532473,-0.387767)--(0.701962,-0.357445)--(0.663889,-0.416474)--cycle;
\path[fill=blue!34!white,draw=blue!34!white,line width=.12pt] (-1.580295,-1.758598)--(-1.503455,-1.765122)--(-1.275669,-1.833387)--(-1.340867,-1.830352)--cycle;
\path[fill=blue!36!white,draw=blue!36!white,line width=.12pt] (1.340867,-1.830352)--(1.275669,-1.833387)--(1.503455,-1.765122)--(1.580295,-1.758598)--cycle;
\path[fill=blue!32!white,draw=blue!32!white,line width=.12pt] (-1.881411,-1.636151)--(-1.804504,-1.675456)--(-1.580295,-1.758598)--(-1.647647,-1.722837)--cycle;
\path[fill=blue!35!white,draw=blue!35!white,line width=.12pt] (1.647647,-1.722837)--(1.580295,-1.758598)--(1.804504,-1.675456)--(1.881411,-1.636151)--cycle;
\path[fill=blue!22!white,draw=blue!22!white,line width=.12pt] (-1.019091,-0.272690)--(-1.086442,-0.236929)--(-0.921837,-0.286259)--(-0.864690,-0.318963)--cycle;
\path[fill=blue!19!white,draw=blue!19!white,line width=.12pt] (0.864690,-0.318963)--(0.921837,-0.286259)--(1.086442,-0.236929)--(1.019091,-0.272690)--cycle;
\path[fill=blue!23!white,draw=blue!23!white,line width=.12pt] (-2.461751,-1.145316)--(-2.426560,-1.260844)--(-2.219681,-1.374972)--(-2.251872,-1.261099)--cycle;
\path[fill=blue!30!white,draw=blue!30!white,line width=.12pt] (2.251872,-1.261099)--(2.219681,-1.374972)--(2.426560,-1.260844)--(2.461751,-1.145316)--cycle;
\path[fill=blue!39!white,draw=blue!39!white,line width=.12pt] (-0.498665,-1.932987)--(-0.468217,-1.869763)--(-0.235241,-1.887042)--(-0.250539,-1.951390)--cycle;
\path[fill=blue!39!white,draw=blue!39!white,line width=.12pt] (0.250539,-1.951390)--(0.235241,-1.887042)--(0.468217,-1.869763)--(0.498665,-1.932987)--cycle;
\path[fill=blue!13!white,draw=blue!13!white,line width=.12pt] (-2.125296,-0.030531)--(-2.250079,-0.111645)--(-2.091880,-0.218298)--(-1.975871,-0.131270)--cycle;
\path[fill=blue!15!white,draw=blue!15!white,line width=.12pt] (1.975871,-0.131270)--(2.091880,-0.218298)--(2.250079,-0.111645)--(2.125296,-0.030531)--cycle;
\path[fill=blue!39!white,draw=blue!39!white,line width=.12pt] (-0.250539,-1.951390)--(-0.235241,-1.887042)--(-0.000000,-1.892820)--(-0.000000,-1.957544)--cycle;
\path[fill=blue!39!white,draw=blue!39!white,line width=.12pt] (-0.000000,-1.957544)--(-0.000000,-1.892820)--(0.235241,-1.887042)--(0.250539,-1.951390)--cycle;
\path[fill=blue!16!white,draw=blue!16!white,line width=.12pt] (-2.610069,-0.606865)--(-2.647922,-0.749480)--(-2.461751,-0.874991)--(-2.426560,-0.730582)--cycle;
\path[fill=blue!23!white,draw=blue!23!white,line width=.12pt] (2.426560,-0.730582)--(2.461751,-0.874991)--(2.647922,-0.749480)--(2.610069,-0.606865)--cycle;
\path[fill=blue!15!white,draw=blue!15!white,line width=.12pt] (-1.734580,-0.032258)--(-1.855225,-0.068452)--(-1.697056,-0.155708)--(-1.586697,-0.113840)--cycle;
\path[fill=blue!15!white,draw=blue!15!white,line width=.12pt] (1.586697,-0.113840)--(1.697056,-0.155708)--(1.855225,-0.068452)--(1.734580,-0.032258)--cycle;
\path[fill=blue!25!white,draw=blue!25!white,line width=.12pt] (-0.503593,-0.445151)--(-0.532473,-0.387767)--(-0.357857,-0.409636)--(-0.338447,-0.465835)--cycle;
\path[fill=blue!23!white,draw=blue!23!white,line width=.12pt] (0.338447,-0.465835)--(0.357857,-0.409636)--(0.532473,-0.387767)--(0.503593,-0.445151)--cycle;
\path[fill=blue!20!white,draw=blue!20!white,line width=.12pt] (-1.240584,-0.179769)--(-1.328326,-0.169203)--(-1.163283,-0.230405)--(-1.086442,-0.236929)--cycle;
\path[fill=blue!17!white,draw=blue!17!white,line width=.12pt] (1.086442,-0.236929)--(1.163283,-0.230405)--(1.328326,-0.169203)--(1.240584,-0.179769)--cycle;
\path[fill=blue!37!white,draw=blue!37!white,line width=.12pt] (-1.035598,-1.890159)--(-0.978167,-1.860259)--(-0.741989,-1.902512)--(-0.785553,-1.934893)--cycle;
\path[fill=blue!38!white,draw=blue!38!white,line width=.12pt] (0.785553,-1.934893)--(0.741989,-1.902512)--(0.978167,-1.860259)--(1.035598,-1.890159)--cycle;
\path[fill=blue!17!white,draw=blue!17!white,line width=.12pt] (-1.480578,-0.100206)--(-1.586697,-0.113840)--(-1.423533,-0.187782)--(-1.328326,-0.169203)--cycle;
\path[fill=blue!16!white,draw=blue!16!white,line width=.12pt] (1.328326,-0.169203)--(1.423533,-0.187782)--(1.586697,-0.113840)--(1.480578,-0.100206)--cycle;
\path[fill=blue!25!white,draw=blue!25!white,line width=.12pt] (-0.338447,-0.465835)--(-0.357857,-0.409636)--(-0.179794,-0.422843)--(-0.170043,-0.478325)--cycle;
\path[fill=blue!24!white,draw=blue!24!white,line width=.12pt] (0.170043,-0.478325)--(0.179794,-0.422843)--(0.357857,-0.409636)--(0.338447,-0.465835)--cycle;
\path[fill=blue!12!white,draw=blue!12!white,line width=.12pt] (-2.250079,-0.111645)--(-2.365079,-0.214089)--(-2.198794,-0.326193)--(-2.091880,-0.218298)--cycle;
\path[fill=blue!16!white,draw=blue!16!white,line width=.12pt] (2.091880,-0.218298)--(2.198794,-0.326193)--(2.365079,-0.214089)--(2.250079,-0.111645)--cycle;
\path[fill=blue!14!white,draw=blue!14!white,line width=.12pt] (-2.548600,-0.466550)--(-2.610069,-0.606865)--(-2.426560,-0.730582)--(-2.369413,-0.587353)--cycle;
\path[fill=blue!21!white,draw=blue!21!white,line width=.12pt] (2.369413,-0.587353)--(2.426560,-0.730582)--(2.610069,-0.606865)--(2.548600,-0.466550)--cycle;
\path[fill=blue!22!white,draw=blue!22!white,line width=.12pt] (-0.864690,-0.318963)--(-0.921837,-0.286259)--(-0.748354,-0.327284)--(-0.701962,-0.357445)--cycle;
\path[fill=blue!20!white,draw=blue!20!white,line width=.12pt] (0.701962,-0.357445)--(0.748354,-0.327284)--(0.921837,-0.286259)--(0.864690,-0.318963)--cycle;
\path[fill=blue!24!white,draw=blue!24!white,line width=.12pt] (-0.170043,-0.478325)--(-0.179794,-0.422843)--(-0.000000,-0.427259)--(-0.000000,-0.482502)--cycle;
\path[fill=blue!24!white,draw=blue!24!white,line width=.12pt] (-0.000000,-0.482502)--(-0.000000,-0.427259)--(0.179794,-0.422843)--(0.170043,-0.478325)--cycle;
\path[fill=blue!26!white,draw=blue!26!white,line width=.12pt] (-2.219681,-1.374972)--(-2.167407,-1.468614)--(-1.944527,-1.569617)--(-1.991426,-1.478411)--cycle;
\path[fill=blue!31!white,draw=blue!31!white,line width=.12pt] (1.991426,-1.478411)--(1.944527,-1.569617)--(2.167407,-1.468614)--(2.219681,-1.374972)--cycle;
\path[fill=blue!12!white,draw=blue!12!white,line width=.12pt] (-2.365079,-0.214089)--(-2.465877,-0.333925)--(-2.292506,-0.450808)--(-2.198794,-0.326193)--cycle;
\path[fill=blue!17!white,draw=blue!17!white,line width=.12pt] (2.198794,-0.326193)--(2.292506,-0.450808)--(2.465877,-0.333925)--(2.365079,-0.214089)--cycle;
\path[fill=blue!13!white,draw=blue!13!white,line width=.12pt] (-2.465877,-0.333925)--(-2.548600,-0.466550)--(-2.369413,-0.587353)--(-2.292506,-0.450808)--cycle;
\path[fill=blue!19!white,draw=blue!19!white,line width=.12pt] (2.292506,-0.450808)--(2.369413,-0.587353)--(2.548600,-0.466550)--(2.465877,-0.333925)--cycle;
\path[fill=blue!20!white,draw=blue!20!white,line width=.12pt] (-2.473633,-1.015029)--(-2.461751,-1.145316)--(-2.251872,-1.261099)--(-2.262742,-1.131371)--cycle;
\path[fill=blue!27!white,draw=blue!27!white,line width=.12pt] (2.262742,-1.131371)--(2.251872,-1.261099)--(2.461751,-1.145316)--(2.473633,-1.015029)--cycle;
\path[fill=blue!34!white,draw=blue!34!white,line width=.12pt] (-1.340867,-1.830352)--(-1.275669,-1.833387)--(-1.035598,-1.890159)--(-1.088526,-1.890026)--cycle;
\path[fill=blue!36!white,draw=blue!36!white,line width=.12pt] (1.088526,-1.890026)--(1.035598,-1.890159)--(1.275669,-1.833387)--(1.340867,-1.830352)--cycle;
\path[fill=blue!37!white,draw=blue!37!white,line width=.12pt] (-0.785553,-1.934893)--(-0.741989,-1.902512)--(-0.498665,-1.932987)--(-0.527943,-1.967157)--cycle;
\path[fill=blue!37!white,draw=blue!37!white,line width=.12pt] (0.527943,-1.967157)--(0.498665,-1.932987)--(0.741989,-1.902512)--(0.785553,-1.934893)--cycle;
\path[fill=blue!14!white,draw=blue!14!white,line width=.12pt] (-1.855225,-0.068452)--(-1.975871,-0.131270)--(-1.807416,-0.224200)--(-1.697056,-0.155708)--cycle;
\path[fill=blue!14!white,draw=blue!14!white,line width=.12pt] (1.697056,-0.155708)--(1.807416,-0.224200)--(1.975871,-0.131270)--(1.855225,-0.068452)--cycle;
\path[fill=blue!29!white,draw=blue!29!white,line width=.12pt] (-1.944527,-1.569617)--(-1.881411,-1.636151)--(-1.647647,-1.722837)--(-1.702920,-1.659211)--cycle;
\path[fill=blue!33!white,draw=blue!33!white,line width=.12pt] (1.702920,-1.659211)--(1.647647,-1.722837)--(1.881411,-1.636151)--(1.944527,-1.569617)--cycle;
\path[fill=blue!32!white,draw=blue!32!white,line width=.12pt] (-1.647647,-1.722837)--(-1.580295,-1.758598)--(-1.340867,-1.830352)--(-1.398014,-1.797648)--cycle;
\path[fill=blue!35!white,draw=blue!35!white,line width=.12pt] (1.398014,-1.797648)--(1.340867,-1.830352)--(1.580295,-1.758598)--(1.647647,-1.722837)--cycle;
\path[fill=blue!22!white,draw=blue!22!white,line width=.12pt] (-0.701962,-0.357445)--(-0.748354,-0.327284)--(-0.567664,-0.359610)--(-0.532473,-0.387767)--cycle;
\path[fill=blue!20!white,draw=blue!20!white,line width=.12pt] (0.532473,-0.387767)--(0.567664,-0.359610)--(0.748354,-0.327284)--(0.701962,-0.357445)--cycle;
\path[fill=blue!20!white,draw=blue!20!white,line width=.12pt] (-1.086442,-0.236929)--(-1.163283,-0.230405)--(-0.987036,-0.283224)--(-0.921837,-0.286259)--cycle;
\path[fill=blue!18!white,draw=blue!18!white,line width=.12pt] (0.921837,-0.286259)--(0.987036,-0.283224)--(1.163283,-0.230405)--(1.086442,-0.236929)--cycle;
\path[fill=blue!37!white,draw=blue!37!white,line width=.12pt] (-0.527943,-1.967157)--(-0.498665,-1.932987)--(-0.250539,-1.951390)--(-0.265249,-1.986640)--cycle;
\path[fill=blue!37!white,draw=blue!37!white,line width=.12pt] (0.265249,-1.986640)--(0.250539,-1.951390)--(0.498665,-1.932987)--(0.527943,-1.967157)--cycle;
\path[fill=blue!15!white,draw=blue!15!white,line width=.12pt] (-1.586697,-0.113840)--(-1.697056,-0.155708)--(-1.522544,-0.234792)--(-1.423533,-0.187782)--cycle;
\path[fill=blue!15!white,draw=blue!15!white,line width=.12pt] (1.423533,-0.187782)--(1.522544,-0.234792)--(1.697056,-0.155708)--(1.586697,-0.113840)--cycle;
\path[fill=blue!17!white,draw=blue!17!white,line width=.12pt] (-1.328326,-0.169203)--(-1.423533,-0.187782)--(-1.246659,-0.253371)--(-1.163283,-0.230405)--cycle;
\path[fill=blue!16!white,draw=blue!16!white,line width=.12pt] (1.163283,-0.230405)--(1.246659,-0.253371)--(1.423533,-0.187782)--(1.328326,-0.169203)--cycle;
\path[fill=blue!37!white,draw=blue!37!white,line width=.12pt] (-0.265249,-1.986640)--(-0.250539,-1.951390)--(-0.000000,-1.957544)--(-0.000000,-1.993156)--cycle;
\path[fill=blue!37!white,draw=blue!37!white,line width=.12pt] (-0.000000,-1.993156)--(-0.000000,-1.957544)--(0.250539,-1.951390)--(0.265249,-1.986640)--cycle;
\path[fill=blue!18!white,draw=blue!18!white,line width=.12pt] (-2.461751,-0.874991)--(-2.473633,-1.015029)--(-2.262742,-1.131371)--(-2.251872,-0.990774)--cycle;
\path[fill=blue!25!white,draw=blue!25!white,line width=.12pt] (2.251872,-0.990774)--(2.262742,-1.131371)--(2.473633,-1.015029)--(2.461751,-0.874991)--cycle;
\path[fill=blue!22!white,draw=blue!22!white,line width=.12pt] (-0.532473,-0.387767)--(-0.567664,-0.359610)--(-0.381508,-0.382925)--(-0.357857,-0.409636)--cycle;
\path[fill=blue!21!white,draw=blue!21!white,line width=.12pt] (0.357857,-0.409636)--(0.381508,-0.382925)--(0.567664,-0.359610)--(0.532473,-0.387767)--cycle;
\path[fill=blue!34!white,draw=blue!34!white,line width=.12pt] (-1.088526,-1.890026)--(-1.035598,-1.890159)--(-0.785553,-1.934893)--(-0.825702,-1.937046)--cycle;
\path[fill=blue!36!white,draw=blue!36!white,line width=.12pt] (0.825702,-1.937046)--(0.785553,-1.934893)--(1.035598,-1.890159)--(1.088526,-1.890026)--cycle;
\path[fill=blue!13!white,draw=blue!13!white,line width=.12pt] (-1.975871,-0.131270)--(-2.091880,-0.218298)--(-1.913535,-0.316685)--(-1.807416,-0.224200)--cycle;
\path[fill=blue!15!white,draw=blue!15!white,line width=.12pt] (1.807416,-0.224200)--(1.913535,-0.316685)--(2.091880,-0.218298)--(1.975871,-0.131270)--cycle;
\path[fill=blue!23!white,draw=blue!23!white,line width=.12pt] (-2.251872,-1.261099)--(-2.219681,-1.374972)--(-1.991426,-1.478411)--(-2.020307,-1.366038)--cycle;
\path[fill=blue!29!white,draw=blue!29!white,line width=.12pt] (2.020307,-1.366038)--(1.991426,-1.478411)--(2.219681,-1.374972)--(2.251872,-1.261099)--cycle;
\path[fill=blue!22!white,draw=blue!22!white,line width=.12pt] (-0.357857,-0.409636)--(-0.381508,-0.382925)--(-0.191677,-0.397005)--(-0.179794,-0.422843)--cycle;
\path[fill=blue!21!white,draw=blue!21!white,line width=.12pt] (0.179794,-0.422843)--(0.191677,-0.397005)--(0.381508,-0.382925)--(0.357857,-0.409636)--cycle;
\path[fill=blue!20!white,draw=blue!20!white,line width=.12pt] (-0.921837,-0.286259)--(-0.987036,-0.283224)--(-0.801283,-0.327152)--(-0.748354,-0.327284)--cycle;
\path[fill=blue!18!white,draw=blue!18!white,line width=.12pt] (0.748354,-0.327284)--(0.801283,-0.327152)--(0.987036,-0.283224)--(0.921837,-0.286259)--cycle;
\path[fill=blue!22!white,draw=blue!22!white,line width=.12pt] (-0.179794,-0.422843)--(-0.191677,-0.397005)--(-0.000000,-0.401713)--(-0.000000,-0.427259)--cycle;
\path[fill=blue!21!white,draw=blue!21!white,line width=.12pt] (-0.000000,-0.427259)--(-0.000000,-0.401713)--(0.191677,-0.397005)--(0.179794,-0.422843)--cycle;
\path[fill=blue!16!white,draw=blue!16!white,line width=.12pt] (-2.426560,-0.730582)--(-2.461751,-0.874991)--(-2.251872,-0.990774)--(-2.219681,-0.844710)--cycle;
\path[fill=blue!22!white,draw=blue!22!white,line width=.12pt] (2.219681,-0.844710)--(2.251872,-0.990774)--(2.461751,-0.874991)--(2.426560,-0.730582)--cycle;
\path[fill=blue!12!white,draw=blue!12!white,line width=.12pt] (-2.091880,-0.218298)--(-2.198794,-0.326193)--(-2.011334,-0.429608)--(-1.913535,-0.316685)--cycle;
\path[fill=blue!15!white,draw=blue!15!white,line width=.12pt] (1.913535,-0.316685)--(2.011334,-0.429608)--(2.198794,-0.326193)--(2.091880,-0.218298)--cycle;
\path[fill=blue!32!white,draw=blue!32!white,line width=.12pt] (-1.398014,-1.797648)--(-1.340867,-1.830352)--(-1.088526,-1.890026)--(-1.134919,-1.859866)--cycle;
\path[fill=blue!34!white,draw=blue!34!white,line width=.12pt] (1.134919,-1.859866)--(1.088526,-1.890026)--(1.340867,-1.830352)--(1.398014,-1.797648)--cycle;
\path[fill=blue!26!white,draw=blue!26!white,line width=.12pt] (-1.991426,-1.478411)--(-1.944527,-1.569617)--(-1.702920,-1.659211)--(-1.743993,-1.570165)--cycle;
\path[fill=blue!31!white,draw=blue!31!white,line width=.12pt] (1.743993,-1.570165)--(1.702920,-1.659211)--(1.944527,-1.569617)--(1.991426,-1.478411)--cycle;
\path[fill=blue!35!white,draw=blue!35!white,line width=.12pt] (-0.825702,-1.937046)--(-0.785553,-1.934893)--(-0.527943,-1.967157)--(-0.554926,-1.970959)--cycle;
\path[fill=blue!36!white,draw=blue!36!white,line width=.12pt] (0.554926,-1.970959)--(0.527943,-1.967157)--(0.785553,-1.934893)--(0.825702,-1.937046)--cycle;
\path[fill=blue!14!white,draw=blue!14!white,line width=.12pt] (-1.697056,-0.155708)--(-1.807416,-0.224200)--(-1.621555,-0.308427)--(-1.522544,-0.234792)--cycle;
\path[fill=blue!14!white,draw=blue!14!white,line width=.12pt] (1.522544,-0.234792)--(1.621555,-0.308427)--(1.807416,-0.224200)--(1.697056,-0.155708)--cycle;
\path[fill=blue!29!white,draw=blue!29!white,line width=.12pt] (-1.702920,-1.659211)--(-1.647647,-1.722837)--(-1.398014,-1.797648)--(-1.444914,-1.736532)--cycle;
\path[fill=blue!32!white,draw=blue!32!white,line width=.12pt] (1.444914,-1.736532)--(1.398014,-1.797648)--(1.647647,-1.722837)--(1.702920,-1.659211)--cycle;
\path[fill=blue!14!white,draw=blue!14!white,line width=.12pt] (-2.369413,-0.587353)--(-2.426560,-0.730582)--(-2.219681,-0.844710)--(-2.167407,-0.698793)--cycle;
\path[fill=blue!20!white,draw=blue!20!white,line width=.12pt] (2.167407,-0.698793)--(2.219681,-0.844710)--(2.426560,-0.730582)--(2.369413,-0.587353)--cycle;
\path[fill=blue!17!white,draw=blue!17!white,line width=.12pt] (-1.163283,-0.230405)--(-1.246659,-0.253371)--(-1.057780,-0.309976)--(-0.987036,-0.283224)--cycle;
\path[fill=blue!16!white,draw=blue!16!white,line width=.12pt] (0.987036,-0.283224)--(1.057780,-0.309976)--(1.246659,-0.253371)--(1.163283,-0.230405)--cycle;
\path[fill=blue!12!white,draw=blue!12!white,line width=.12pt] (-2.198794,-0.326193)--(-2.292506,-0.450808)--(-2.097056,-0.558630)--(-2.011334,-0.429608)--cycle;
\path[fill=blue!16!white,draw=blue!16!white,line width=.12pt] (2.011334,-0.429608)--(2.097056,-0.558630)--(2.292506,-0.450808)--(2.198794,-0.326193)--cycle;
\path[fill=blue!18!white,draw=blue!18!white,line width=.12pt] (0.567664,-0.359610)--(0.607814,-0.361764)--(0.801283,-0.327152)--(0.748354,-0.327284)--cycle;
\path[fill=blue!20!white,draw=blue!20!white,line width=.12pt] (-0.748354,-0.327284)--(-0.801283,-0.327152)--(-0.607814,-0.361764)--(-0.567664,-0.359610)--cycle;
\path[fill=blue!13!white,draw=blue!13!white,line width=.12pt] (-2.292506,-0.450808)--(-2.369413,-0.587353)--(-2.167407,-0.698793)--(-2.097056,-0.558630)--cycle;
\path[fill=blue!18!white,draw=blue!18!white,line width=.12pt] (2.097056,-0.558630)--(2.167407,-0.698793)--(2.369413,-0.587353)--(2.292506,-0.450808)--cycle;
\path[fill=blue!20!white,draw=blue!20!white,line width=.12pt] (-2.262742,-1.131371)--(-2.251872,-1.261099)--(-2.020307,-1.366038)--(-2.030059,-1.236817)--cycle;
\path[fill=blue!26!white,draw=blue!26!white,line width=.12pt] (2.030059,-1.236817)--(2.020307,-1.366038)--(2.251872,-1.261099)--(2.262742,-1.131371)--cycle;
\path[fill=blue!15!white,draw=blue!15!white,line width=.12pt] (-1.423533,-0.187782)--(-1.522544,-0.234792)--(-1.333369,-0.304943)--(-1.246659,-0.253371)--cycle;
\path[fill=blue!15!white,draw=blue!15!white,line width=.12pt] (1.246659,-0.253371)--(1.333369,-0.304943)--(1.522544,-0.234792)--(1.423533,-0.187782)--cycle;
\path[fill=blue!35!white,draw=blue!35!white,line width=.12pt] (-0.554926,-1.970959)--(-0.527943,-1.967157)--(-0.265249,-1.986640)--(-0.278805,-1.991439)--cycle;
\path[fill=blue!35!white,draw=blue!35!white,line width=.12pt] (0.278805,-1.991439)--(0.265249,-1.986640)--(0.527943,-1.967157)--(0.554926,-1.970959)--cycle;
\path[fill=blue!35!white,draw=blue!35!white,line width=.12pt] (-0.278805,-1.991439)--(-0.265249,-1.986640)--(-0.000000,-1.993156)--(-0.000000,-1.998287)--cycle;
\path[fill=blue!35!white,draw=blue!35!white,line width=.12pt] (-0.000000,-1.998287)--(-0.000000,-1.993156)--(0.265249,-1.986640)--(0.278805,-1.991439)--cycle;
\path[fill=blue!19!white,draw=blue!19!white,line width=.12pt] (-0.567664,-0.359610)--(-0.607814,-0.361764)--(-0.408491,-0.386728)--(-0.381508,-0.382925)--cycle;
\path[fill=blue!18!white,draw=blue!18!white,line width=.12pt] (0.381508,-0.382925)--(0.408491,-0.386728)--(0.607814,-0.361764)--(0.567664,-0.359610)--cycle;
\path[fill=blue!34!white,draw=blue!34!white,line width=.12pt] (0.860893,-1.908890)--(0.825702,-1.937046)--(1.088526,-1.890026)--(1.134919,-1.859866)--cycle;
\path[fill=blue!32!white,draw=blue!32!white,line width=.12pt] (-1.134919,-1.859866)--(-1.088526,-1.890026)--(-0.825702,-1.937046)--(-0.860893,-1.908890)--cycle;
\path[fill=blue!13!white,draw=blue!13!white,line width=.12pt] (-1.807416,-0.224200)--(-1.913535,-0.316685)--(-1.716761,-0.405857)--(-1.621555,-0.308427)--cycle;
\path[fill=blue!14!white,draw=blue!14!white,line width=.12pt] (1.621555,-0.308427)--(1.716761,-0.405857)--(1.913535,-0.316685)--(1.807416,-0.224200)--cycle;
\path[fill=blue!17!white,draw=blue!17!white,line width=.12pt] (-0.987036,-0.283224)--(-1.057780,-0.309976)--(-0.858714,-0.357051)--(-0.801283,-0.327152)--cycle;
\path[fill=blue!16!white,draw=blue!16!white,line width=.12pt] (0.801283,-0.327152)--(0.858714,-0.357051)--(1.057780,-0.309976)--(0.987036,-0.283224)--cycle;
\path[fill=blue!23!white,draw=blue!23!white,line width=.12pt] (-2.020307,-1.366038)--(-1.991426,-1.478411)--(-1.743993,-1.570165)--(-1.769285,-1.459123)--cycle;
\path[fill=blue!28!white,draw=blue!28!white,line width=.12pt] (1.769285,-1.459123)--(1.743993,-1.570165)--(1.991426,-1.478411)--(2.020307,-1.366038)--cycle;
\path[fill=blue!19!white,draw=blue!19!white,line width=.12pt] (-0.381508,-0.382925)--(-0.408491,-0.386728)--(-0.205234,-0.401803)--(-0.191677,-0.397005)--cycle;
\path[fill=blue!19!white,draw=blue!19!white,line width=.12pt] (0.191677,-0.397005)--(0.205234,-0.401803)--(0.408491,-0.386728)--(0.381508,-0.382925)--cycle;
\path[fill=blue!18!white,draw=blue!18!white,line width=.12pt] (-2.251872,-0.990774)--(-2.262742,-1.131371)--(-2.030059,-1.236817)--(-2.020307,-1.095713)--cycle;
\path[fill=blue!24!white,draw=blue!24!white,line width=.12pt] (2.020307,-1.095713)--(2.030059,-1.236817)--(2.262742,-1.131371)--(2.251872,-0.990774)--cycle;
\path[fill=blue!19!white,draw=blue!19!white,line width=.12pt] (-0.191677,-0.397005)--(-0.205234,-0.401803)--(-0.000000,-0.406844)--(-0.000000,-0.401713)--cycle;
\path[fill=blue!19!white,draw=blue!19!white,line width=.12pt] (-0.000000,-0.401713)--(-0.000000,-0.406844)--(0.205234,-0.401803)--(0.191677,-0.397005)--cycle;
\path[fill=blue!29!white,draw=blue!29!white,line width=.12pt] (-1.444914,-1.736532)--(-1.398014,-1.797648)--(-1.134919,-1.859866)--(-1.172992,-1.800837)--cycle;
\path[fill=blue!32!white,draw=blue!32!white,line width=.12pt] (1.172992,-1.800837)--(1.134919,-1.859866)--(1.398014,-1.797648)--(1.444914,-1.736532)--cycle;
\path[fill=blue!26!white,draw=blue!26!white,line width=.12pt] (-1.743993,-1.570165)--(-1.702920,-1.659211)--(-1.444914,-1.736532)--(-1.479763,-1.649352)--cycle;
\path[fill=blue!30!white,draw=blue!30!white,line width=.12pt] (1.479763,-1.649352)--(1.444914,-1.736532)--(1.702920,-1.659211)--(1.743993,-1.570165)--cycle;
\path[fill=blue!14!white,draw=blue!14!white,line width=.12pt] (-1.522544,-0.234792)--(-1.621555,-0.308427)--(-1.420078,-0.383140)--(-1.333369,-0.304943)--cycle;
\path[fill=blue!14!white,draw=blue!14!white,line width=.12pt] (1.333369,-0.304943)--(1.420078,-0.383140)--(1.621555,-0.308427)--(1.522544,-0.234792)--cycle;
\path[fill=blue!15!white,draw=blue!15!white,line width=.12pt] (-1.246659,-0.253371)--(-1.333369,-0.304943)--(-1.131352,-0.365485)--(-1.057780,-0.309976)--cycle;
\path[fill=blue!15!white,draw=blue!15!white,line width=.12pt] (1.057780,-0.309976)--(1.131352,-0.365485)--(1.333369,-0.304943)--(1.246659,-0.253371)--cycle;
\path[fill=blue!12!white,draw=blue!12!white,line width=.12pt] (-1.913535,-0.316685)--(-2.011334,-0.429608)--(-1.804504,-0.523338)--(-1.716761,-0.405857)--cycle;
\path[fill=blue!15!white,draw=blue!15!white,line width=.12pt] (1.716761,-0.405857)--(1.804504,-0.523338)--(2.011334,-0.429608)--(1.913535,-0.316685)--cycle;
\path[fill=blue!32!white,draw=blue!32!white,line width=.12pt] (-0.860893,-1.908890)--(-0.825702,-1.937046)--(-0.554926,-1.970959)--(-0.578577,-1.944248)--cycle;
\path[fill=blue!33!white,draw=blue!33!white,line width=.12pt] (0.578577,-1.944248)--(0.554926,-1.970959)--(0.825702,-1.937046)--(0.860893,-1.908890)--cycle;
\path[fill=blue!16!white,draw=blue!16!white,line width=.12pt] (-2.219681,-0.844710)--(-2.251872,-0.990774)--(-2.020307,-1.095713)--(-1.991426,-0.948149)--cycle;
\path[fill=blue!21!white,draw=blue!21!white,line width=.12pt] (1.991426,-0.948149)--(2.020307,-1.095713)--(2.251872,-0.990774)--(2.219681,-0.844710)--cycle;
\path[fill=blue!17!white,draw=blue!17!white,line width=.12pt] (-0.801283,-0.327152)--(-0.858714,-0.357051)--(-0.651378,-0.394145)--(-0.607814,-0.361764)--cycle;
\path[fill=blue!16!white,draw=blue!16!white,line width=.12pt] (0.607814,-0.361764)--(0.651378,-0.394145)--(0.858714,-0.357051)--(0.801283,-0.327152)--cycle;
\path[fill=blue!12!white,draw=blue!12!white,line width=.12pt] (-2.011334,-0.429608)--(-2.097056,-0.558630)--(-1.881411,-0.656355)--(-1.804504,-0.523338)--cycle;
\path[fill=blue!16!white,draw=blue!16!white,line width=.12pt] (1.804504,-0.523338)--(1.881411,-0.656355)--(2.097056,-0.558630)--(2.011334,-0.429608)--cycle;
\path[fill=blue!14!white,draw=blue!14!white,line width=.12pt] (-2.167407,-0.698793)--(-2.219681,-0.844710)--(-1.991426,-0.948149)--(-1.944527,-0.799796)--cycle;
\path[fill=blue!19!white,draw=blue!19!white,line width=.12pt] (1.944527,-0.799796)--(1.991426,-0.948149)--(2.219681,-0.844710)--(2.167407,-0.698793)--cycle;
\path[fill=blue!32!white,draw=blue!32!white,line width=.12pt] (-0.578577,-1.944248)--(-0.554926,-1.970959)--(-0.278805,-1.991439)--(-0.290688,-1.965600)--cycle;
\path[fill=blue!33!white,draw=blue!33!white,line width=.12pt] (0.290688,-1.965600)--(0.278805,-1.991439)--(0.554926,-1.970959)--(0.578577,-1.944248)--cycle;
\path[fill=blue!20!white,draw=blue!20!white,line width=.12pt] (-2.030059,-1.236817)--(-2.020307,-1.366038)--(-1.769285,-1.459123)--(-1.777825,-1.330351)--cycle;
\path[fill=blue!26!white,draw=blue!26!white,line width=.12pt] (1.777825,-1.330351)--(1.769285,-1.459123)--(2.020307,-1.366038)--(2.030059,-1.236817)--cycle;
\path[fill=blue!13!white,draw=blue!13!white,line width=.12pt] (-2.097056,-0.558630)--(-2.167407,-0.698793)--(-1.944527,-0.799796)--(-1.881411,-0.656355)--cycle;
\path[fill=blue!17!white,draw=blue!17!white,line width=.12pt] (1.881411,-0.656355)--(1.944527,-0.799796)--(2.167407,-0.698793)--(2.097056,-0.558630)--cycle;
\path[fill=blue!17!white,draw=blue!17!white,line width=.12pt] (-0.607814,-0.361764)--(-0.651378,-0.394145)--(-0.437769,-0.420898)--(-0.408491,-0.386728)--cycle;
\path[fill=blue!17!white,draw=blue!17!white,line width=.12pt] (0.408491,-0.386728)--(0.437769,-0.420898)--(0.651378,-0.394145)--(0.607814,-0.361764)--cycle;
\path[fill=blue!32!white,draw=blue!32!white,line width=.12pt] (-0.290688,-1.965600)--(-0.278805,-1.991439)--(-0.000000,-1.998287)--(-0.000000,-1.972741)--cycle;
\path[fill=blue!33!white,draw=blue!33!white,line width=.12pt] (-0.000000,-1.972741)--(-0.000000,-1.998287)--(0.278805,-1.991439)--(0.290688,-1.965600)--cycle;
\path[fill=blue!29!white,draw=blue!29!white,line width=.12pt] (-1.172992,-1.800837)--(-1.134919,-1.859866)--(-0.860893,-1.908890)--(-0.889774,-1.851505)--cycle;
\path[fill=blue!31!white,draw=blue!31!white,line width=.12pt] (0.889774,-1.851505)--(0.860893,-1.908890)--(1.134919,-1.859866)--(1.172992,-1.800837)--cycle;
\path[fill=blue!13!white,draw=blue!13!white,line width=.12pt] (-1.621555,-0.308427)--(-1.716761,-0.405857)--(-1.503455,-0.484957)--(-1.420078,-0.383140)--cycle;
\path[fill=blue!14!white,draw=blue!14!white,line width=.12pt] (1.420078,-0.383140)--(1.503455,-0.484957)--(1.716761,-0.405857)--(1.621555,-0.308427)--cycle;
\path[fill=blue!15!white,draw=blue!15!white,line width=.12pt] (-1.057780,-0.309976)--(-1.131352,-0.365485)--(-0.918440,-0.415835)--(-0.858714,-0.357051)--cycle;
\path[fill=blue!15!white,draw=blue!15!white,line width=.12pt] (0.858714,-0.357051)--(0.918440,-0.415835)--(1.131352,-0.365485)--(1.057780,-0.309976)--cycle;
\path[fill=blue!23!white,draw=blue!23!white,line width=.12pt] (-1.769285,-1.459123)--(-1.743993,-1.570165)--(-1.479763,-1.649352)--(-1.501223,-1.539458)--cycle;
\path[fill=blue!27!white,draw=blue!27!white,line width=.12pt] (1.501223,-1.539458)--(1.479763,-1.649352)--(1.743993,-1.570165)--(1.769285,-1.459123)--cycle;
\path[fill=blue!14!white,draw=blue!14!white,line width=.12pt] (-1.333369,-0.304943)--(-1.420078,-0.383140)--(-1.204924,-0.447619)--(-1.131352,-0.365485)--cycle;
\path[fill=blue!14!white,draw=blue!14!white,line width=.12pt] (1.131352,-0.365485)--(1.204924,-0.447619)--(1.420078,-0.383140)--(1.333369,-0.304943)--cycle;
\path[fill=blue!17!white,draw=blue!17!white,line width=.12pt] (0.205234,-0.401803)--(0.219943,-0.437053)--(0.437769,-0.420898)--(0.408491,-0.386728)--cycle;
\path[fill=blue!17!white,draw=blue!17!white,line width=.12pt] (-0.408491,-0.386728)--(-0.437769,-0.420898)--(-0.219943,-0.437053)--(-0.205234,-0.401803)--cycle;
\path[fill=blue!26!white,draw=blue!26!white,line width=.12pt] (-1.479763,-1.649352)--(-1.444914,-1.736532)--(-1.172992,-1.800837)--(-1.201283,-1.715208)--cycle;
\path[fill=blue!29!white,draw=blue!29!white,line width=.12pt] (1.201283,-1.715208)--(1.172992,-1.800837)--(1.444914,-1.736532)--(1.479763,-1.649352)--cycle;
\path[fill=blue!17!white,draw=blue!17!white,line width=.12pt] (-0.205234,-0.401803)--(-0.219943,-0.437053)--(-0.000000,-0.442456)--(-0.000000,-0.406844)--cycle;
\path[fill=blue!17!white,draw=blue!17!white,line width=.12pt] (-0.000000,-0.406844)--(-0.000000,-0.442456)--(0.219943,-0.437053)--(0.205234,-0.401803)--cycle;
\path[fill=blue!18!white,draw=blue!18!white,line width=.12pt] (-2.020307,-1.095713)--(-2.030059,-1.236817)--(-1.777825,-1.330351)--(-1.769285,-1.188798)--cycle;
\path[fill=blue!23!white,draw=blue!23!white,line width=.12pt] (1.769285,-1.188798)--(1.777825,-1.330351)--(2.030059,-1.236817)--(2.020307,-1.095713)--cycle;
\path[fill=blue!29!white,draw=blue!29!white,line width=.12pt] (-0.889774,-1.851505)--(-0.860893,-1.908890)--(-0.578577,-1.944248)--(-0.597986,-1.888050)--cycle;
\path[fill=blue!31!white,draw=blue!31!white,line width=.12pt] (0.597986,-1.888050)--(0.578577,-1.944248)--(0.860893,-1.908890)--(0.889774,-1.851505)--cycle;
\path[fill=blue!12!white,draw=blue!12!white,line width=.12pt] (-1.716761,-0.405857)--(-1.804504,-0.523338)--(-1.580295,-0.606480)--(-1.503455,-0.484957)--cycle;
\path[fill=blue!15!white,draw=blue!15!white,line width=.12pt] (1.503455,-0.484957)--(1.580295,-0.606480)--(1.804504,-0.523338)--(1.716761,-0.405857)--cycle;
\path[fill=blue!15!white,draw=blue!15!white,line width=.12pt] (-0.858714,-0.357051)--(-0.918440,-0.415835)--(-0.696683,-0.455508)--(-0.651378,-0.394145)--cycle;
\path[fill=blue!15!white,draw=blue!15!white,line width=.12pt] (0.651378,-0.394145)--(0.696683,-0.455508)--(0.918440,-0.415835)--(0.858714,-0.357051)--cycle;
\path[fill=blue!16!white,draw=blue!16!white,line width=.12pt] (-1.991426,-0.948149)--(-2.020307,-1.095713)--(-1.769285,-1.188798)--(-1.743993,-1.039904)--cycle;
\path[fill=blue!21!white,draw=blue!21!white,line width=.12pt] (1.743993,-1.039904)--(1.769285,-1.188798)--(2.020307,-1.095713)--(1.991426,-0.948149)--cycle;
\path[fill=blue!12!white,draw=blue!12!white,line width=.12pt] (-1.804504,-0.523338)--(-1.881411,-0.656355)--(-1.647647,-0.743041)--(-1.580295,-0.606480)--cycle;
\path[fill=blue!16!white,draw=blue!16!white,line width=.12pt] (1.580295,-0.606480)--(1.647647,-0.743041)--(1.881411,-0.656355)--(1.804504,-0.523338)--cycle;
\path[fill=blue!29!white,draw=blue!29!white,line width=.12pt] (-0.597986,-1.888050)--(-0.578577,-1.944248)--(-0.290688,-1.965600)--(-0.300440,-1.910118)--cycle;
\path[fill=blue!30!white,draw=blue!30!white,line width=.12pt] (0.300440,-1.910118)--(0.290688,-1.965600)--(0.578577,-1.944248)--(0.597986,-1.888050)--cycle;
\path[fill=blue!20!white,draw=blue!20!white,line width=.12pt] (-1.777825,-1.330351)--(-1.769285,-1.459123)--(-1.501223,-1.539458)--(-1.508470,-1.411074)--cycle;
\path[fill=blue!25!white,draw=blue!25!white,line width=.12pt] (1.508470,-1.411074)--(1.501223,-1.539458)--(1.769285,-1.459123)--(1.777825,-1.330351)--cycle;
\path[fill=blue!13!white,draw=blue!13!white,line width=.12pt] (-1.420078,-0.383140)--(-1.503455,-0.484957)--(-1.275669,-0.553222)--(-1.204924,-0.447619)--cycle;
\path[fill=blue!14!white,draw=blue!14!white,line width=.12pt] (1.204924,-0.447619)--(1.275669,-0.553222)--(1.503455,-0.484957)--(1.420078,-0.383140)--cycle;
\path[fill=blue!14!white,draw=blue!14!white,line width=.12pt] (-1.131352,-0.365485)--(-1.204924,-0.447619)--(-0.978167,-0.501243)--(-0.918440,-0.415835)--cycle;
\path[fill=blue!14!white,draw=blue!14!white,line width=.12pt] (0.918440,-0.415835)--(0.978167,-0.501243)--(1.204924,-0.447619)--(1.131352,-0.365485)--cycle;
\path[fill=blue!26!white,draw=blue!26!white,line width=.12pt] (-1.201283,-1.715208)--(-1.172992,-1.800837)--(-0.889774,-1.851505)--(-0.911234,-1.767098)--cycle;
\path[fill=blue!29!white,draw=blue!29!white,line width=.12pt] (0.911234,-1.767098)--(0.889774,-1.851505)--(1.172992,-1.800837)--(1.201283,-1.715208)--cycle;
\path[fill=blue!14!white,draw=blue!14!white,line width=.12pt] (-1.944527,-0.799796)--(-1.991426,-0.948149)--(-1.743993,-1.039904)--(-1.702920,-0.889390)--cycle;
\path[fill=blue!19!white,draw=blue!19!white,line width=.12pt] (1.702920,-0.889390)--(1.743993,-1.039904)--(1.991426,-0.948149)--(1.944527,-0.799796)--cycle;
\path[fill=blue!15!white,draw=blue!15!white,line width=.12pt] (-0.651378,-0.394145)--(-0.696683,-0.455508)--(-0.468217,-0.484122)--(-0.437769,-0.420898)--cycle;
\path[fill=blue!15!white,draw=blue!15!white,line width=.12pt] (0.437769,-0.420898)--(0.468217,-0.484122)--(0.696683,-0.455508)--(0.651378,-0.394145)--cycle;
\path[fill=blue!13!white,draw=blue!13!white,line width=.12pt] (-1.881411,-0.656355)--(-1.944527,-0.799796)--(-1.702920,-0.889390)--(-1.647647,-0.743041)--cycle;
\path[fill=blue!17!white,draw=blue!17!white,line width=.12pt] (1.647647,-0.743041)--(1.702920,-0.889390)--(1.944527,-0.799796)--(1.881411,-0.656355)--cycle;
\path[fill=blue!30!white,draw=blue!30!white,line width=.12pt] (-0.300440,-1.910118)--(-0.290688,-1.965600)--(-0.000000,-1.972741)--(-0.000000,-1.917498)--cycle;
\path[fill=blue!30!white,draw=blue!30!white,line width=.12pt] (-0.000000,-1.917498)--(-0.000000,-1.972741)--(0.290688,-1.965600)--(0.300440,-1.910118)--cycle;
\path[fill=blue!23!white,draw=blue!23!white,line width=.12pt] (-1.501223,-1.539458)--(-1.479763,-1.649352)--(-1.201283,-1.715208)--(-1.218704,-1.606269)--cycle;
\path[fill=blue!27!white,draw=blue!27!white,line width=.12pt] (1.218704,-1.606269)--(1.201283,-1.715208)--(1.479763,-1.649352)--(1.501223,-1.539458)--cycle;
\path[fill=blue!15!white,draw=blue!15!white,line width=.12pt] (-0.437769,-0.420898)--(-0.468217,-0.484122)--(-0.235241,-0.501401)--(-0.219943,-0.437053)--cycle;
\path[fill=blue!15!white,draw=blue!15!white,line width=.12pt] (0.219943,-0.437053)--(0.235241,-0.501401)--(0.468217,-0.484122)--(0.437769,-0.420898)--cycle;
\path[fill=blue!15!white,draw=blue!15!white,line width=.12pt] (-0.219943,-0.437053)--(-0.235241,-0.501401)--(-0.000000,-0.507180)--(-0.000000,-0.442456)--cycle;
\path[fill=blue!15!white,draw=blue!15!white,line width=.12pt] (-0.000000,-0.442456)--(-0.000000,-0.507180)--(0.235241,-0.501401)--(0.219943,-0.437053)--cycle;
\path[fill=blue!18!white,draw=blue!18!white,line width=.12pt] (-1.769285,-1.188798)--(-1.777825,-1.330351)--(-1.508470,-1.411074)--(-1.501223,-1.269133)--cycle;
\path[fill=blue!22!white,draw=blue!22!white,line width=.12pt] (1.501223,-1.269133)--(1.508470,-1.411074)--(1.777825,-1.330351)--(1.769285,-1.188798)--cycle;
\path[fill=blue!12!white,draw=blue!12!white,line width=.12pt] (-1.503455,-0.484957)--(-1.580295,-0.606480)--(-1.340867,-0.678234)--(-1.275669,-0.553222)--cycle;
\path[fill=blue!14!white,draw=blue!14!white,line width=.12pt] (1.275669,-0.553222)--(1.340867,-0.678234)--(1.580295,-0.606480)--(1.503455,-0.484957)--cycle;
\path[fill=blue!14!white,draw=blue!14!white,line width=.12pt] (-0.918440,-0.415835)--(-0.978167,-0.501243)--(-0.741989,-0.543496)--(-0.696683,-0.455508)--cycle;
\path[fill=blue!14!white,draw=blue!14!white,line width=.12pt] (0.696683,-0.455508)--(0.741989,-0.543496)--(0.978167,-0.501243)--(0.918440,-0.415835)--cycle;
\path[fill=blue!26!white,draw=blue!26!white,line width=.12pt] (-0.911234,-1.767098)--(-0.889774,-1.851505)--(-0.597986,-1.888050)--(-0.612409,-1.804524)--cycle;
\path[fill=blue!28!white,draw=blue!28!white,line width=.12pt] (0.612409,-1.804524)--(0.597986,-1.888050)--(0.889774,-1.851505)--(0.911234,-1.767098)--cycle;
\path[fill=blue!13!white,draw=blue!13!white,line width=.12pt] (-1.204924,-0.447619)--(-1.275669,-0.553222)--(-1.035598,-0.609994)--(-0.978167,-0.501243)--cycle;
\path[fill=blue!14!white,draw=blue!14!white,line width=.12pt] (0.978167,-0.501243)--(1.035598,-0.609994)--(1.275669,-0.553222)--(1.204924,-0.447619)--cycle;
\path[fill=blue!20!white,draw=blue!20!white,line width=.12pt] (-1.508470,-1.411074)--(-1.501223,-1.539458)--(-1.218704,-1.606269)--(-1.224587,-1.478207)--cycle;
\path[fill=blue!24!white,draw=blue!24!white,line width=.12pt] (1.224587,-1.478207)--(1.218704,-1.606269)--(1.501223,-1.539458)--(1.508470,-1.411074)--cycle;
\path[fill=blue!23!white,draw=blue!23!white,line width=.12pt] (-1.218704,-1.606269)--(-1.201283,-1.715208)--(-0.911234,-1.767098)--(-0.924449,-1.658912)--cycle;
\path[fill=blue!26!white,draw=blue!26!white,line width=.12pt] (0.924449,-1.658912)--(0.911234,-1.767098)--(1.201283,-1.715208)--(1.218704,-1.606269)--cycle;
\path[fill=blue!16!white,draw=blue!16!white,line width=.12pt] (-1.743993,-1.039904)--(-1.769285,-1.188798)--(-1.501223,-1.269133)--(-1.479763,-1.119090)--cycle;
\path[fill=blue!20!white,draw=blue!20!white,line width=.12pt] (1.479763,-1.119090)--(1.501223,-1.269133)--(1.769285,-1.188798)--(1.743993,-1.039904)--cycle;
\path[fill=blue!12!white,draw=blue!12!white,line width=.12pt] (-1.580295,-0.606480)--(-1.647647,-0.743041)--(-1.398014,-0.817853)--(-1.340867,-0.678234)--cycle;
\path[fill=blue!15!white,draw=blue!15!white,line width=.12pt] (1.340867,-0.678234)--(1.398014,-0.817853)--(1.647647,-0.743041)--(1.580295,-0.606480)--cycle;
\path[fill=blue!26!white,draw=blue!26!white,line width=.12pt] (-0.612409,-1.804524)--(-0.597986,-1.888050)--(-0.300440,-1.910118)--(-0.307686,-1.827125)--cycle;
\path[fill=blue!28!white,draw=blue!28!white,line width=.12pt] (0.307686,-1.827125)--(0.300440,-1.910118)--(0.597986,-1.888050)--(0.612409,-1.804524)--cycle;
\path[fill=blue!14!white,draw=blue!14!white,line width=.12pt] (-0.696683,-0.455508)--(-0.741989,-0.543496)--(-0.498665,-0.573971)--(-0.468217,-0.484122)--cycle;
\path[fill=blue!14!white,draw=blue!14!white,line width=.12pt] (0.468217,-0.484122)--(0.498665,-0.573971)--(0.741989,-0.543496)--(0.696683,-0.455508)--cycle;
\path[fill=blue!14!white,draw=blue!14!white,line width=.12pt] (-1.702920,-0.889390)--(-1.743993,-1.039904)--(-1.479763,-1.119090)--(-1.444914,-0.966711)--cycle;
\path[fill=blue!18!white,draw=blue!18!white,line width=.12pt] (1.444914,-0.966711)--(1.479763,-1.119090)--(1.743993,-1.039904)--(1.702920,-0.889390)--cycle;
\path[fill=blue!13!white,draw=blue!13!white,line width=.12pt] (-1.647647,-0.743041)--(-1.702920,-0.889390)--(-1.444914,-0.966711)--(-1.398014,-0.817853)--cycle;
\path[fill=blue!16!white,draw=blue!16!white,line width=.12pt] (1.398014,-0.817853)--(1.444914,-0.966711)--(1.702920,-0.889390)--(1.647647,-0.743041)--cycle;
\path[fill=blue!27!white,draw=blue!27!white,line width=.12pt] (-0.307686,-1.827125)--(-0.300440,-1.910118)--(-0.000000,-1.917498)--(-0.000000,-1.834683)--cycle;
\path[fill=blue!27!white,draw=blue!27!white,line width=.12pt] (-0.000000,-1.834683)--(-0.000000,-1.917498)--(0.300440,-1.910118)--(0.307686,-1.827125)--cycle;
\path[fill=blue!14!white,draw=blue!14!white,line width=.12pt] (-0.468217,-0.484122)--(-0.498665,-0.573971)--(-0.250539,-0.592374)--(-0.235241,-0.501401)--cycle;
\path[fill=blue!14!white,draw=blue!14!white,line width=.12pt] (0.235241,-0.501401)--(0.250539,-0.592374)--(0.498665,-0.573971)--(0.468217,-0.484122)--cycle;
\path[fill=blue!12!white,draw=blue!12!white,line width=.12pt] (-1.275669,-0.553222)--(-1.340867,-0.678234)--(-1.088526,-0.737908)--(-1.035598,-0.609994)--cycle;
\path[fill=blue!14!white,draw=blue!14!white,line width=.12pt] (1.035598,-0.609994)--(1.088526,-0.737908)--(1.340867,-0.678234)--(1.275669,-0.553222)--cycle;
\path[fill=blue!13!white,draw=blue!13!white,line width=.12pt] (-0.978167,-0.501243)--(-1.035598,-0.609994)--(-0.785553,-0.654728)--(-0.741989,-0.543496)--cycle;
\path[fill=blue!14!white,draw=blue!14!white,line width=.12pt] (0.741989,-0.543496)--(0.785553,-0.654728)--(1.035598,-0.609994)--(0.978167,-0.501243)--cycle;
\path[fill=blue!18!white,draw=blue!18!white,line width=.12pt] (-1.501223,-1.269133)--(-1.508470,-1.411074)--(-1.224587,-1.478207)--(-1.218704,-1.335944)--cycle;
\path[fill=blue!22!white,draw=blue!22!white,line width=.12pt] (1.218704,-1.335944)--(1.224587,-1.478207)--(1.508470,-1.411074)--(1.501223,-1.269133)--cycle;
\path[fill=blue!14!white,draw=blue!14!white,line width=.12pt] (-0.235241,-0.501401)--(-0.250539,-0.592374)--(-0.000000,-0.598528)--(-0.000000,-0.507180)--cycle;
\path[fill=blue!14!white,draw=blue!14!white,line width=.12pt] (-0.000000,-0.507180)--(-0.000000,-0.598528)--(0.250539,-0.592374)--(0.235241,-0.501401)--cycle;
\path[fill=blue!23!white,draw=blue!23!white,line width=.12pt] (-0.924449,-1.658912)--(-0.911234,-1.767098)--(-0.612409,-1.804524)--(-0.621290,-1.696881)--cycle;
\path[fill=blue!25!white,draw=blue!25!white,line width=.12pt] (0.621290,-1.696881)--(0.612409,-1.804524)--(0.911234,-1.767098)--(0.924449,-1.658912)--cycle;
\path[fill=blue!20!white,draw=blue!20!white,line width=.12pt] (-1.224587,-1.478207)--(-1.218704,-1.606269)--(-0.924449,-1.658912)--(-0.928911,-1.531105)--cycle;
\path[fill=blue!23!white,draw=blue!23!white,line width=.12pt] (0.928911,-1.531105)--(0.924449,-1.658912)--(1.218704,-1.606269)--(1.224587,-1.478207)--cycle;
\path[fill=blue!12!white,draw=blue!12!white,line width=.12pt] (-1.340867,-0.678234)--(-1.398014,-0.817853)--(-1.134919,-0.880070)--(-1.088526,-0.737908)--cycle;
\path[fill=blue!15!white,draw=blue!15!white,line width=.12pt] (1.088526,-0.737908)--(1.134919,-0.880070)--(1.398014,-0.817853)--(1.340867,-0.678234)--cycle;
\path[fill=blue!16!white,draw=blue!16!white,line width=.12pt] (-1.479763,-1.119090)--(-1.501223,-1.269133)--(-1.218704,-1.335944)--(-1.201283,-1.184946)--cycle;
\path[fill=blue!19!white,draw=blue!19!white,line width=.12pt] (1.201283,-1.184946)--(1.218704,-1.335944)--(1.501223,-1.269133)--(1.479763,-1.119090)--cycle;
\path[fill=blue!13!white,draw=blue!13!white,line width=.12pt] (-0.741989,-0.543496)--(-0.785553,-0.654728)--(-0.527943,-0.686992)--(-0.498665,-0.573971)--cycle;
\path[fill=blue!13!white,draw=blue!13!white,line width=.12pt] (0.498665,-0.573971)--(0.527943,-0.686992)--(0.785553,-0.654728)--(0.741989,-0.543496)--cycle;
\path[fill=blue!24!white,draw=blue!24!white,line width=.12pt] (-0.621290,-1.696881)--(-0.612409,-1.804524)--(-0.307686,-1.827125)--(-0.312148,-1.719809)--cycle;
\path[fill=blue!25!white,draw=blue!25!white,line width=.12pt] (0.312148,-1.719809)--(0.307686,-1.827125)--(0.612409,-1.804524)--(0.621290,-1.696881)--cycle;
\path[fill=blue!13!white,draw=blue!13!white,line width=.12pt] (-1.398014,-0.817853)--(-1.444914,-0.966711)--(-1.172992,-1.031016)--(-1.134919,-0.880070)--cycle;
\path[fill=blue!16!white,draw=blue!16!white,line width=.12pt] (1.134919,-0.880070)--(1.172992,-1.031016)--(1.444914,-0.966711)--(1.398014,-0.817853)--cycle;
\path[fill=blue!12!white,draw=blue!12!white,line width=.12pt] (-1.035598,-0.609994)--(-1.088526,-0.737908)--(-0.825702,-0.784928)--(-0.785553,-0.654728)--cycle;
\path[fill=blue!14!white,draw=blue!14!white,line width=.12pt] (0.785553,-0.654728)--(0.825702,-0.784928)--(1.088526,-0.737908)--(1.035598,-0.609994)--cycle;
\path[fill=blue!14!white,draw=blue!14!white,line width=.12pt] (-1.444914,-0.966711)--(-1.479763,-1.119090)--(-1.201283,-1.184946)--(-1.172992,-1.031016)--cycle;
\path[fill=blue!17!white,draw=blue!17!white,line width=.12pt] (1.172992,-1.031016)--(1.201283,-1.184946)--(1.479763,-1.119090)--(1.444914,-0.966711)--cycle;
\path[fill=blue!24!white,draw=blue!24!white,line width=.12pt] (-0.312148,-1.719809)--(-0.307686,-1.827125)--(-0.000000,-1.834683)--(-0.000000,-1.727477)--cycle;
\path[fill=blue!24!white,draw=blue!24!white,line width=.12pt] (-0.000000,-1.727477)--(-0.000000,-1.834683)--(0.307686,-1.827125)--(0.312148,-1.719809)--cycle;
\path[fill=blue!13!white,draw=blue!13!white,line width=.12pt] (-0.498665,-0.573971)--(-0.527943,-0.686992)--(-0.265249,-0.706475)--(-0.250539,-0.592374)--cycle;
\path[fill=blue!13!white,draw=blue!13!white,line width=.12pt] (0.250539,-0.592374)--(0.265249,-0.706475)--(0.527943,-0.686992)--(0.498665,-0.573971)--cycle;
\path[fill=blue!18!white,draw=blue!18!white,line width=.12pt] (-1.218704,-1.335944)--(-1.224587,-1.478207)--(-0.928911,-1.531105)--(-0.924449,-1.388587)--cycle;
\path[fill=blue!21!white,draw=blue!21!white,line width=.12pt] (0.924449,-1.388587)--(0.928911,-1.531105)--(1.224587,-1.478207)--(1.218704,-1.335944)--cycle;
\path[fill=blue!21!white,draw=blue!21!white,line width=.12pt] (-0.928911,-1.531105)--(-0.924449,-1.658912)--(-0.621290,-1.696881)--(-0.624289,-1.569256)--cycle;
\path[fill=blue!23!white,draw=blue!23!white,line width=.12pt] (0.624289,-1.569256)--(0.621290,-1.696881)--(0.924449,-1.658912)--(0.928911,-1.531105)--cycle;
\path[fill=blue!13!white,draw=blue!13!white,line width=.12pt] (-0.250539,-0.592374)--(-0.265249,-0.706475)--(-0.000000,-0.712991)--(-0.000000,-0.598528)--cycle;
\path[fill=blue!13!white,draw=blue!13!white,line width=.12pt] (-0.000000,-0.598528)--(-0.000000,-0.712991)--(0.265249,-0.706475)--(0.250539,-0.592374)--cycle;
\path[fill=blue!14!white,draw=blue!14!white,line width=.12pt] (0.825702,-0.784928)--(0.860893,-0.929094)--(1.134919,-0.880070)--(1.088526,-0.737908)--cycle;
\path[fill=blue!12!white,draw=blue!12!white,line width=.12pt] (-1.088526,-0.737908)--(-1.134919,-0.880070)--(-0.860893,-0.929094)--(-0.825702,-0.784928)--cycle;
\path[fill=blue!12!white,draw=blue!12!white,line width=.12pt] (-0.785553,-0.654728)--(-0.825702,-0.784928)--(-0.554926,-0.818841)--(-0.527943,-0.686992)--cycle;
\path[fill=blue!13!white,draw=blue!13!white,line width=.12pt] (0.527943,-0.686992)--(0.554926,-0.818841)--(0.825702,-0.784928)--(0.785553,-0.654728)--cycle;
\path[fill=blue!16!white,draw=blue!16!white,line width=.12pt] (-1.201283,-1.184946)--(-1.218704,-1.335944)--(-0.924449,-1.388587)--(-0.911234,-1.236837)--cycle;
\path[fill=blue!19!white,draw=blue!19!white,line width=.12pt] (0.911234,-1.236837)--(0.924449,-1.388587)--(1.218704,-1.335944)--(1.201283,-1.184946)--cycle;
\path[fill=blue!21!white,draw=blue!21!white,line width=.12pt] (-0.624289,-1.569256)--(-0.621290,-1.696881)--(-0.312148,-1.719809)--(-0.313655,-1.592296)--cycle;
\path[fill=blue!22!white,draw=blue!22!white,line width=.12pt] (0.313655,-1.592296)--(0.312148,-1.719809)--(0.621290,-1.696881)--(0.624289,-1.569256)--cycle;
\path[fill=blue!13!white,draw=blue!13!white,line width=.12pt] (-1.134919,-0.880070)--(-1.172992,-1.031016)--(-0.889774,-1.081685)--(-0.860893,-0.929094)--cycle;
\path[fill=blue!15!white,draw=blue!15!white,line width=.12pt] (0.860893,-0.929094)--(0.889774,-1.081685)--(1.172992,-1.031016)--(1.134919,-0.880070)--cycle;
\path[fill=blue!14!white,draw=blue!14!white,line width=.12pt] (-1.172992,-1.031016)--(-1.201283,-1.184946)--(-0.911234,-1.236837)--(-0.889774,-1.081685)--cycle;
\path[fill=blue!17!white,draw=blue!17!white,line width=.12pt] (0.889774,-1.081685)--(0.911234,-1.236837)--(1.201283,-1.184946)--(1.172992,-1.031016)--cycle;
\path[fill=blue!18!white,draw=blue!18!white,line width=.12pt] (-0.924449,-1.388587)--(-0.928911,-1.531105)--(-0.624289,-1.569256)--(-0.621290,-1.426556)--cycle;
\path[fill=blue!20!white,draw=blue!20!white,line width=.12pt] (0.621290,-1.426556)--(0.624289,-1.569256)--(0.928911,-1.531105)--(0.924449,-1.388587)--cycle;
\path[fill=blue!21!white,draw=blue!21!white,line width=.12pt] (-0.313655,-1.592296)--(-0.312148,-1.719809)--(-0.000000,-1.727477)--(-0.000000,-1.600000)--cycle;
\path[fill=blue!22!white,draw=blue!22!white,line width=.12pt] (-0.000000,-1.600000)--(-0.000000,-1.727477)--(0.312148,-1.719809)--(0.313655,-1.592296)--cycle;
\path[fill=blue!12!white,draw=blue!12!white,line width=.12pt] (-0.527943,-0.686992)--(-0.554926,-0.818841)--(-0.278805,-0.839321)--(-0.265249,-0.706475)--cycle;
\path[fill=blue!13!white,draw=blue!13!white,line width=.12pt] (0.265249,-0.706475)--(0.278805,-0.839321)--(0.554926,-0.818841)--(0.527943,-0.686992)--cycle;
\path[fill=blue!12!white,draw=blue!12!white,line width=.12pt] (-0.825702,-0.784928)--(-0.860893,-0.929094)--(-0.578577,-0.964452)--(-0.554926,-0.818841)--cycle;
\path[fill=blue!14!white,draw=blue!14!white,line width=.12pt] (0.554926,-0.818841)--(0.578577,-0.964452)--(0.860893,-0.929094)--(0.825702,-0.784928)--cycle;
\path[fill=blue!13!white,draw=blue!13!white,line width=.12pt] (-0.000000,-0.712991)--(-0.000000,-0.846169)--(0.278805,-0.839321)--(0.265249,-0.706475)--cycle;
\path[fill=blue!13!white,draw=blue!13!white,line width=.12pt] (-0.265249,-0.706475)--(-0.278805,-0.839321)--(-0.000000,-0.846169)--(-0.000000,-0.712991)--cycle;
\path[fill=blue!16!white,draw=blue!16!white,line width=.12pt] (-0.911234,-1.236837)--(-0.924449,-1.388587)--(-0.621290,-1.426556)--(-0.612409,-1.274262)--cycle;
\path[fill=blue!18!white,draw=blue!18!white,line width=.12pt] (0.612409,-1.274262)--(0.621290,-1.426556)--(0.924449,-1.388587)--(0.911234,-1.236837)--cycle;
\path[fill=blue!18!white,draw=blue!18!white,line width=.12pt] (-0.621290,-1.426556)--(-0.624289,-1.569256)--(-0.313655,-1.592296)--(-0.312148,-1.449484)--cycle;
\path[fill=blue!20!white,draw=blue!20!white,line width=.12pt] (0.312148,-1.449484)--(0.313655,-1.592296)--(0.624289,-1.569256)--(0.621290,-1.426556)--cycle;
\path[fill=blue!13!white,draw=blue!13!white,line width=.12pt] (-0.860893,-0.929094)--(-0.889774,-1.081685)--(-0.597986,-1.118229)--(-0.578577,-0.964452)--cycle;
\path[fill=blue!15!white,draw=blue!15!white,line width=.12pt] (0.578577,-0.964452)--(0.597986,-1.118229)--(0.889774,-1.081685)--(0.860893,-0.929094)--cycle;
\path[fill=blue!14!white,draw=blue!14!white,line width=.12pt] (-0.889774,-1.081685)--(-0.911234,-1.236837)--(-0.612409,-1.274262)--(-0.597986,-1.118229)--cycle;
\path[fill=blue!16!white,draw=blue!16!white,line width=.12pt] (0.597986,-1.118229)--(0.612409,-1.274262)--(0.911234,-1.236837)--(0.889774,-1.081685)--cycle;
\path[fill=blue!13!white,draw=blue!13!white,line width=.12pt] (-0.554926,-0.818841)--(-0.578577,-0.964452)--(-0.290688,-0.985804)--(-0.278805,-0.839321)--cycle;
\path[fill=blue!13!white,draw=blue!13!white,line width=.12pt] (0.278805,-0.839321)--(0.290688,-0.985804)--(0.578577,-0.964452)--(0.554926,-0.818841)--cycle;
\path[fill=blue!19!white,draw=blue!19!white,line width=.12pt] (-0.312148,-1.449484)--(-0.313655,-1.592296)--(-0.000000,-1.600000)--(-0.000000,-1.457152)--cycle;
\path[fill=blue!19!white,draw=blue!19!white,line width=.12pt] (-0.000000,-1.457152)--(-0.000000,-1.600000)--(0.313655,-1.592296)--(0.312148,-1.449484)--cycle;
\path[fill=blue!13!white,draw=blue!13!white,line width=.12pt] (-0.278805,-0.839321)--(-0.290688,-0.985804)--(-0.000000,-0.992945)--(-0.000000,-0.846169)--cycle;
\path[fill=blue!13!white,draw=blue!13!white,line width=.12pt] (-0.000000,-0.846169)--(-0.000000,-0.992945)--(0.290688,-0.985804)--(0.278805,-0.839321)--cycle;
\path[fill=blue!16!white,draw=blue!16!white,line width=.12pt] (-0.612409,-1.274262)--(-0.621290,-1.426556)--(-0.312148,-1.449484)--(-0.307686,-1.296863)--cycle;
\path[fill=blue!18!white,draw=blue!18!white,line width=.12pt] (0.307686,-1.296863)--(0.312148,-1.449484)--(0.621290,-1.426556)--(0.612409,-1.274262)--cycle;
\path[fill=blue!13!white,draw=blue!13!white,line width=.12pt] (-0.578577,-0.964452)--(-0.597986,-1.118229)--(-0.300440,-1.140298)--(-0.290688,-0.985804)--cycle;
\path[fill=blue!14!white,draw=blue!14!white,line width=.12pt] (0.290688,-0.985804)--(0.300440,-1.140298)--(0.597986,-1.118229)--(0.578577,-0.964452)--cycle;
\path[fill=blue!15!white,draw=blue!15!white,line width=.12pt] (-0.597986,-1.118229)--(-0.612409,-1.274262)--(-0.307686,-1.296863)--(-0.300440,-1.140298)--cycle;
\path[fill=blue!16!white,draw=blue!16!white,line width=.12pt] (0.300440,-1.140298)--(0.307686,-1.296863)--(0.612409,-1.274262)--(0.597986,-1.118229)--cycle;
\path[fill=blue!17!white,draw=blue!17!white,line width=.12pt] (-0.000000,-1.304421)--(-0.000000,-1.457152)--(0.312148,-1.449484)--(0.307686,-1.296863)--cycle;
\path[fill=blue!17!white,draw=blue!17!white,line width=.12pt] (-0.307686,-1.296863)--(-0.312148,-1.449484)--(-0.000000,-1.457152)--(-0.000000,-1.304421)--cycle;
\path[fill=blue!14!white,draw=blue!14!white,line width=.12pt] (-0.290688,-0.985804)--(-0.300440,-1.140298)--(-0.000000,-1.147677)--(-0.000000,-0.992945)--cycle;
\path[fill=blue!14!white,draw=blue!14!white,line width=.12pt] (-0.000000,-0.992945)--(-0.000000,-1.147677)--(0.300440,-1.140298)--(0.290688,-0.985804)--cycle;
\path[fill=blue!15!white,draw=blue!15!white,line width=.12pt] (-0.300440,-1.140298)--(-0.307686,-1.296863)--(-0.000000,-1.304421)--(-0.000000,-1.147677)--cycle;
\path[fill=blue!15!white,draw=blue!15!white,line width=.12pt] (-0.000000,-1.147677)--(-0.000000,-1.304421)--(0.307686,-1.296863)--(0.300440,-1.140298)--cycle;

\draw[aux](0,-2.35)--(0,2.75);
\draw[->](0,2.75)--(0,3.2)node[above]{$z$};
\draw[aux](0,0)--(2.4,0);
\fill(0,0)circle(1pt)node[below left]{$O$};
\fill(2.4,0)circle(1pt)node[below]{$C$};
\node[fill=white,inner sep=1pt]at(1.08,-.20){$R$};
\meridiancurve{red!75!black,thick}{0:180}
\meridiancurve{red!75!black,dashed}{180:360}
\tp{A}{40}{0}
\draw[red!75!black,thin](2.4,0)--(A);
\fill(A)circle(1.5pt)node[above left]{$A$};
\node[anchor=west]at(3.35,.85){子午横截面};
\draw[thin](3.31,.69)--(2.96,.48);
\end{scope}
\begin{scope}[xshift=7.2cm,yshift=.55cm]
\draw[->](-1.72,0)--(1.90,0)node[right]{$\mathbf e_\rho$};
\draw[->](0,-1.40)--(0,1.67)node[above]{$z$};
\draw[thick](0,0)circle(1.12);
\coordinate(C)at(0,0);
\coordinate(Ac)at({-1.12*cos(40)},{1.12*sin(40)});
\draw[red!75!black](C)--(Ac)node[midway,above right]{$r$};
\draw[aux](C)--(-1.12,0);
\draw(-.42,0)arc(180:140:.42);
\node at(-.63,.21){$\theta$};
\fill(C)circle(1.2pt)node[below right]{$C$};
\fill(Ac)circle(1.5pt)node[above left]{$A$};
\node at(0,-1.90){$s=R-r\cos\theta,\quad z=r\sin\theta$};
\node at(0,2.65){子午截面；内侧在左};
\end{scope}
\end{tikzpicture}
% 3. Cylinder force balances: cut-line length, not area stress.
\begin{tikzpicture}[font=\small]
\begin{scope}
\fill[blue!5](0,-1.35)arc(-90:90:1.35)--cycle;
\draw[thick](0,-1.35)arc(-90:90:1.35);
\draw[aux](0,-1.35)--(0,1.35);
\draw[thin](0,0)--(1.35,0)node[midway,below]{$r$};
\foreach \ang in {-60,0,60}{
\draw[pressure,->]({1.35*cos(\ang)},{1.35*sin(\ang)})
--({2.0*cos(\ang)},{2.0*sin(\ang)});
}
\node[red!75!black]at(1.25,2.05){$p$};
\draw[force,->](0,1.35)--(-1.2,1.35)node[left]{$\sigma_1h$};
\draw[force,->](0,-1.35)--(-1.2,-1.35)node[left]{$\sigma_1h$};
\node at(.15,2.65){半筒横截面；轴向长 $h$};
\node at(.15,-2.3){$F_{p,x}=2prh$};
\end{scope}
\begin{scope}[xshift=6.1cm,yshift=-1.0cm]
\fill[gray!8](-1.1,0)--(-1.1,2.2)arc(180:360:1.1 and .33)
--(1.1,0)arc(0:-180:1.1 and .33)--cycle;
\draw[thick](-1.1,0)--(-1.1,2.2);
\draw[thick](1.1,0)--(1.1,2.2);
\filldraw[fill=gray!12,thick](0,2.2)ellipse(1.1 and .33);
\draw[blue!75!black,thick](-1.1,0)arc(180:360:1.1 and .33);
\draw[blue!75!black,dashed](-1.1,0)arc(180:0:1.1 and .33);
\foreach \xx in {-.7,0,.7}{
\draw[pressure,->](\xx,2.3)--(\xx,3.05);
\draw[force,->](\xx,-.25)--(\xx,-.98);
}
\node[red!75!black]at(0,3.40){$p\pi r^2$};
\node[blue!75!black]at(0,-1.40){$2\pi r\,\sigma_2$};
\draw[->](1.75,.6)--(1.75,1.45)node[above]{$z$};
\node at(0,3.85){闭端筒段：轴向平衡};
\end{scope}
\end{tikzpicture}
% 4. Highlighted inner upper band is a subset of the real surface, not a disc.
\begin{tikzpicture}[font=\small]
\begin{scope}[yshift=.8cm]
% Automatically generated vector mesh; parameters and camera in build/torus-mesh-q03.json.
\path[fill=blue!39!white,draw=blue!39!white,line width=.12pt] (0.000000,1.304421)--(0.000000,1.147677)--(-0.300440,1.140298)--(-0.307686,1.296863)--cycle;
\path[fill=blue!39!white,draw=blue!39!white,line width=.12pt] (0.307686,1.296863)--(0.300440,1.140298)--(0.000000,1.147677)--(0.000000,1.304421)--cycle;
\path[fill=blue!40!white,draw=blue!40!white,line width=.12pt] (0.300440,1.140298)--(0.290688,0.985804)--(0.000000,0.992945)--(0.000000,1.147677)--cycle;
\path[fill=blue!40!white,draw=blue!40!white,line width=.12pt] (0.000000,1.147677)--(0.000000,0.992945)--(-0.290688,0.985804)--(-0.300440,1.140298)--cycle;
\path[fill=blue!37!white,draw=blue!37!white,line width=.12pt] (0.312148,1.449484)--(0.307686,1.296863)--(0.000000,1.304421)--(0.000000,1.457152)--cycle;
\path[fill=blue!37!white,draw=blue!37!white,line width=.12pt] (0.000000,1.457152)--(0.000000,1.304421)--(-0.307686,1.296863)--(-0.312148,1.449484)--cycle;
\path[fill=blue!39!white,draw=blue!39!white,line width=.12pt] (0.612409,1.274262)--(0.597986,1.118229)--(0.300440,1.140298)--(0.307686,1.296863)--cycle;
\path[fill=blue!38!white,draw=blue!38!white,line width=.12pt] (-0.307686,1.296863)--(-0.300440,1.140298)--(-0.597986,1.118229)--(-0.612409,1.274262)--cycle;
\path[fill=blue!41!white,draw=blue!41!white,line width=.12pt] (0.597986,1.118229)--(0.578577,0.964452)--(0.290688,0.985804)--(0.300440,1.140298)--cycle;
\path[fill=blue!40!white,draw=blue!40!white,line width=.12pt] (-0.300440,1.140298)--(-0.290688,0.985804)--(-0.578577,0.964452)--(-0.597986,1.118229)--cycle;
\path[fill=blue!38!white,draw=blue!38!white,line width=.12pt] (0.621290,1.426556)--(0.612409,1.274262)--(0.307686,1.296863)--(0.312148,1.449484)--cycle;
\path[fill=blue!36!white,draw=blue!36!white,line width=.12pt] (-0.312148,1.449484)--(-0.307686,1.296863)--(-0.612409,1.274262)--(-0.621290,1.426556)--cycle;
\path[fill=blue!41!white,draw=blue!41!white,line width=.12pt] (0.290688,0.985804)--(0.278805,0.839321)--(0.000000,0.846169)--(0.000000,0.992945)--cycle;
\path[fill=blue!41!white,draw=blue!41!white,line width=.12pt] (0.000000,0.992945)--(0.000000,0.846169)--(-0.278805,0.839321)--(-0.290688,0.985804)--cycle;
\path[fill=blue!35!white,draw=blue!35!white,line width=.12pt] (0.313655,1.592296)--(0.312148,1.449484)--(0.000000,1.457152)--(0.000000,1.600000)--cycle;
\path[fill=blue!35!white,draw=blue!35!white,line width=.12pt] (0.000000,1.600000)--(0.000000,1.457152)--(-0.312148,1.449484)--(-0.313655,1.592296)--cycle;
\path[fill=blue!41!white,draw=blue!41!white,line width=.12pt] (-0.290688,0.985804)--(-0.278805,0.839321)--(-0.554926,0.818841)--(-0.578577,0.964452)--cycle;
\path[fill=blue!41!white,draw=blue!41!white,line width=.12pt] (0.578577,0.964452)--(0.554926,0.818841)--(0.278805,0.839321)--(0.290688,0.985804)--cycle;
\path[fill=blue!38!white,draw=blue!38!white,line width=.12pt] (-0.612409,1.274262)--(-0.597986,1.118229)--(-0.889774,1.081685)--(-0.911234,1.236837)--cycle;
\path[fill=blue!40!white,draw=blue!40!white,line width=.12pt] (0.911234,1.236837)--(0.889774,1.081685)--(0.597986,1.118229)--(0.612409,1.274262)--cycle;
\path[fill=blue!41!white,draw=blue!41!white,line width=.12pt] (0.889774,1.081685)--(0.860893,0.929094)--(0.578577,0.964452)--(0.597986,1.118229)--cycle;
\path[fill=blue!39!white,draw=blue!39!white,line width=.12pt] (-0.597986,1.118229)--(-0.578577,0.964452)--(-0.860893,0.929094)--(-0.889774,1.081685)--cycle;
\path[fill=blue!36!white,draw=blue!36!white,line width=.12pt] (0.624289,1.569256)--(0.621290,1.426556)--(0.312148,1.449484)--(0.313655,1.592296)--cycle;
\path[fill=blue!34!white,draw=blue!34!white,line width=.12pt] (-0.313655,1.592296)--(-0.312148,1.449484)--(-0.621290,1.426556)--(-0.624289,1.569256)--cycle;
\path[fill=blue!38!white,draw=blue!38!white,line width=.12pt] (0.924449,1.388587)--(0.911234,1.236837)--(0.612409,1.274262)--(0.621290,1.426556)--cycle;
\path[fill=blue!36!white,draw=blue!36!white,line width=.12pt] (-0.621290,1.426556)--(-0.612409,1.274262)--(-0.911234,1.236837)--(-0.924449,1.388587)--cycle;
\path[fill=blue!41!white,draw=blue!41!white,line width=.12pt] (0.278805,0.839321)--(0.265249,0.706475)--(0.000000,0.712991)--(0.000000,0.846169)--cycle;
\path[fill=blue!41!white,draw=blue!41!white,line width=.12pt] (0.000000,0.846169)--(0.000000,0.712991)--(-0.265249,0.706475)--(-0.278805,0.839321)--cycle;
\path[fill=blue!42!white,draw=blue!42!white,line width=.12pt] (0.860893,0.929094)--(0.825702,0.784928)--(0.554926,0.818841)--(0.578577,0.964452)--cycle;
\path[fill=blue!40!white,draw=blue!40!white,line width=.12pt] (-0.578577,0.964452)--(-0.554926,0.818841)--(-0.825702,0.784928)--(-0.860893,0.929094)--cycle;
\path[fill=blue!42!white,draw=blue!42!white,line width=.12pt] (0.554926,0.818841)--(0.527943,0.686992)--(0.265249,0.706475)--(0.278805,0.839321)--cycle;
\path[fill=blue!41!white,draw=blue!41!white,line width=.12pt] (-0.278805,0.839321)--(-0.265249,0.706475)--(-0.527943,0.686992)--(-0.554926,0.818841)--cycle;
\path[fill=blue!33!white,draw=blue!33!white,line width=.12pt] (0.312148,1.719809)--(0.313655,1.592296)--(0.000000,1.600000)--(0.000000,1.727477)--cycle;
\path[fill=blue!32!white,draw=blue!32!white,line width=.12pt] (0.000000,1.727477)--(0.000000,1.600000)--(-0.313655,1.592296)--(-0.312148,1.719809)--cycle;
\path[fill=blue!36!white,draw=blue!36!white,line width=.12pt] (0.928911,1.531105)--(0.924449,1.388587)--(0.621290,1.426556)--(0.624289,1.569256)--cycle;
\path[fill=blue!34!white,draw=blue!34!white,line width=.12pt] (-0.624289,1.569256)--(-0.621290,1.426556)--(-0.924449,1.388587)--(-0.928911,1.531105)--cycle;
\path[fill=blue!37!white,draw=blue!37!white,line width=.12pt] (-0.911234,1.236837)--(-0.889774,1.081685)--(-1.172992,1.031016)--(-1.201283,1.184946)--cycle;
\path[fill=blue!40!white,draw=blue!40!white,line width=.12pt] (1.201283,1.184946)--(1.172992,1.031016)--(0.889774,1.081685)--(0.911234,1.236837)--cycle;
\path[fill=blue!41!white,draw=blue!41!white,line width=.12pt] (1.172992,1.031016)--(1.134919,0.880070)--(0.860893,0.929094)--(0.889774,1.081685)--cycle;
\path[fill=blue!39!white,draw=blue!39!white,line width=.12pt] (-0.889774,1.081685)--(-0.860893,0.929094)--(-1.134919,0.880070)--(-1.172992,1.031016)--cycle;
\path[fill=blue!32!white,draw=blue!32!white,line width=.12pt] (-0.312148,1.719809)--(-0.313655,1.592296)--(-0.624289,1.569256)--(-0.621290,1.696881)--cycle;
\path[fill=blue!33!white,draw=blue!33!white,line width=.12pt] (0.621290,1.696881)--(0.624289,1.569256)--(0.313655,1.592296)--(0.312148,1.719809)--cycle;
\path[fill=blue!38!white,draw=blue!38!white,line width=.12pt] (1.218704,1.335944)--(1.201283,1.184946)--(0.911234,1.236837)--(0.924449,1.388587)--cycle;
\path[fill=blue!35!white,draw=blue!35!white,line width=.12pt] (-0.924449,1.388587)--(-0.911234,1.236837)--(-1.201283,1.184946)--(-1.218704,1.335944)--cycle;
\path[fill=blue!42!white,draw=blue!42!white,line width=.12pt] (0.825702,0.784928)--(0.785553,0.654728)--(0.527943,0.686992)--(0.554926,0.818841)--cycle;
\path[fill=blue!41!white,draw=blue!41!white,line width=.12pt] (-0.554926,0.818841)--(-0.527943,0.686992)--(-0.785553,0.654728)--(-0.825702,0.784928)--cycle;
\path[fill=blue!42!white,draw=blue!42!white,line width=.12pt] (1.134919,0.880070)--(1.088526,0.737908)--(0.825702,0.784928)--(0.860893,0.929094)--cycle;
\path[fill=blue!40!white,draw=blue!40!white,line width=.12pt] (-0.860893,0.929094)--(-0.825702,0.784928)--(-1.088526,0.737908)--(-1.134919,0.880070)--cycle;
\path[fill=blue!41!white,draw=blue!41!white,line width=.12pt] (0.265249,0.706475)--(0.250539,0.592374)--(0.000000,0.598528)--(0.000000,0.712991)--cycle;
\path[fill=blue!41!white,draw=blue!41!white,line width=.12pt] (0.000000,0.712991)--(0.000000,0.598528)--(-0.250539,0.592374)--(-0.265249,0.706475)--cycle;
\path[fill=blue!33!white,draw=blue!33!white,line width=.12pt] (0.924449,1.658912)--(0.928911,1.531105)--(0.624289,1.569256)--(0.621290,1.696881)--cycle;
\path[fill=blue!31!white,draw=blue!31!white,line width=.12pt] (-0.621290,1.696881)--(-0.624289,1.569256)--(-0.928911,1.531105)--(-0.924449,1.658912)--cycle;
\path[fill=blue!33!white,draw=blue!33!white,line width=.12pt] (-0.928911,1.531105)--(-0.924449,1.388587)--(-1.218704,1.335944)--(-1.224587,1.478207)--cycle;
\path[fill=blue!36!white,draw=blue!36!white,line width=.12pt] (1.224587,1.478207)--(1.218704,1.335944)--(0.924449,1.388587)--(0.928911,1.531105)--cycle;
\path[fill=blue!41!white,draw=blue!41!white,line width=.12pt] (-0.265249,0.706475)--(-0.250539,0.592374)--(-0.498665,0.573971)--(-0.527943,0.686992)--cycle;
\path[fill=blue!41!white,draw=blue!41!white,line width=.12pt] (0.527943,0.686992)--(0.498665,0.573971)--(0.250539,0.592374)--(0.265249,0.706475)--cycle;
\path[fill=blue!30!white,draw=blue!30!white,line width=.12pt] (0.000000,1.834683)--(0.000000,1.727477)--(-0.312148,1.719809)--(-0.307686,1.827125)--cycle;
\path[fill=blue!30!white,draw=blue!30!white,line width=.12pt] (0.307686,1.827125)--(0.312148,1.719809)--(0.000000,1.727477)--(0.000000,1.834683)--cycle;
\path[fill=blue!40!white,draw=blue!40!white,line width=.12pt] (1.479763,1.119090)--(1.444914,0.966711)--(1.172992,1.031016)--(1.201283,1.184946)--cycle;
\path[fill=blue!37!white,draw=blue!37!white,line width=.12pt] (-1.201283,1.184946)--(-1.172992,1.031016)--(-1.444914,0.966711)--(-1.479763,1.119090)--cycle;
\path[fill=blue!42!white,draw=blue!42!white,line width=.12pt] (1.088526,0.737908)--(1.035598,0.609994)--(0.785553,0.654728)--(0.825702,0.784928)--cycle;
\path[fill=blue!40!white,draw=blue!40!white,line width=.12pt] (-0.825702,0.784928)--(-0.785553,0.654728)--(-1.035598,0.609994)--(-1.088526,0.737908)--cycle;
\path[fill=blue!41!white,draw=blue!41!white,line width=.12pt] (1.444914,0.966711)--(1.398014,0.817853)--(1.134919,0.880070)--(1.172992,1.031016)--cycle;
\path[fill=blue!38!white,draw=blue!38!white,line width=.12pt] (-1.172992,1.031016)--(-1.134919,0.880070)--(-1.398014,0.817853)--(-1.444914,0.966711)--cycle;
\path[fill=blue!29!white,draw=blue!29!white,line width=.12pt] (-0.307686,1.827125)--(-0.312148,1.719809)--(-0.621290,1.696881)--(-0.612409,1.804524)--cycle;
\path[fill=blue!30!white,draw=blue!30!white,line width=.12pt] (0.612409,1.804524)--(0.621290,1.696881)--(0.312148,1.719809)--(0.307686,1.827125)--cycle;
\path[fill=blue!41!white,draw=blue!41!white,line width=.12pt] (0.785553,0.654728)--(0.741989,0.543496)--(0.498665,0.573971)--(0.527943,0.686992)--cycle;
\path[fill=blue!41!white,draw=blue!41!white,line width=.12pt] (-0.527943,0.686992)--(-0.498665,0.573971)--(-0.741989,0.543496)--(-0.785553,0.654728)--cycle;
\path[fill=blue!38!white,draw=blue!38!white,line width=.12pt] (1.501223,1.269133)--(1.479763,1.119090)--(1.201283,1.184946)--(1.218704,1.335944)--cycle;
\path[fill=blue!35!white,draw=blue!35!white,line width=.12pt] (-1.218704,1.335944)--(-1.201283,1.184946)--(-1.479763,1.119090)--(-1.501223,1.269133)--cycle;
\path[fill=blue!42!white,draw=blue!42!white,line width=.12pt] (1.398014,0.817853)--(1.340867,0.678234)--(1.088526,0.737908)--(1.134919,0.880070)--cycle;
\path[fill=blue!39!white,draw=blue!39!white,line width=.12pt] (-1.134919,0.880070)--(-1.088526,0.737908)--(-1.340867,0.678234)--(-1.398014,0.817853)--cycle;
\path[fill=blue!31!white,draw=blue!31!white,line width=.12pt] (-0.924449,1.658912)--(-0.928911,1.531105)--(-1.224587,1.478207)--(-1.218704,1.606269)--cycle;
\path[fill=blue!34!white,draw=blue!34!white,line width=.12pt] (1.218704,1.606269)--(1.224587,1.478207)--(0.928911,1.531105)--(0.924449,1.658912)--cycle;
\path[fill=blue!31!white,draw=blue!31!white,line width=.12pt] (0.911234,1.767098)--(0.924449,1.658912)--(0.621290,1.696881)--(0.612409,1.804524)--cycle;
\path[fill=blue!29!white,draw=blue!29!white,line width=.12pt] (-0.612409,1.804524)--(-0.621290,1.696881)--(-0.924449,1.658912)--(-0.911234,1.767098)--cycle;
\path[fill=blue!40!white,draw=blue!40!white,line width=.12pt] (0.250539,0.592374)--(0.235241,0.501401)--(0.000000,0.507180)--(0.000000,0.598528)--cycle;
\path[fill=blue!40!white,draw=blue!40!white,line width=.12pt] (0.000000,0.598528)--(0.000000,0.507180)--(-0.235241,0.501401)--(-0.250539,0.592374)--cycle;
\path[fill=blue!36!white,draw=blue!36!white,line width=.12pt] (1.508470,1.411074)--(1.501223,1.269133)--(1.218704,1.335944)--(1.224587,1.478207)--cycle;
\path[fill=blue!32!white,draw=blue!32!white,line width=.12pt] (-1.224587,1.478207)--(-1.218704,1.335944)--(-1.501223,1.269133)--(-1.508470,1.411074)--cycle;
\path[fill=blue!40!white,draw=blue!40!white,line width=.12pt] (-0.785553,0.654728)--(-0.741989,0.543496)--(-0.978167,0.501243)--(-1.035598,0.609994)--cycle;
\path[fill=blue!41!white,draw=blue!41!white,line width=.12pt] (1.035598,0.609994)--(0.978167,0.501243)--(0.741989,0.543496)--(0.785553,0.654728)--cycle;
\path[fill=blue!42!white,draw=blue!42!white,line width=.12pt] (1.340867,0.678234)--(1.275669,0.553222)--(1.035598,0.609994)--(1.088526,0.737908)--cycle;
\path[fill=blue!40!white,draw=blue!40!white,line width=.12pt] (-1.088526,0.737908)--(-1.035598,0.609994)--(-1.275669,0.553222)--(-1.340867,0.678234)--cycle;
\path[fill=blue!40!white,draw=blue!40!white,line width=.12pt] (-0.250539,0.592374)--(-0.235241,0.501401)--(-0.468217,0.484122)--(-0.498665,0.573971)--cycle;
\path[fill=blue!40!white,draw=blue!40!white,line width=.12pt] (0.498665,0.573971)--(0.468217,0.484122)--(0.235241,0.501401)--(0.250539,0.592374)--cycle;
\path[fill=blue!27!white,draw=blue!27!white,line width=.12pt] (0.000000,1.917498)--(0.000000,1.834683)--(-0.307686,1.827125)--(-0.300440,1.910118)--cycle;
\path[fill=blue!27!white,draw=blue!27!white,line width=.12pt] (0.300440,1.910118)--(0.307686,1.827125)--(0.000000,1.834683)--(0.000000,1.917498)--cycle;
\path[fill=blue!41!white,draw=blue!41!white,line width=.12pt] (1.702920,0.889390)--(1.647647,0.743041)--(1.398014,0.817853)--(1.444914,0.966711)--cycle;
\path[fill=blue!38!white,draw=blue!38!white,line width=.12pt] (-1.444914,0.966711)--(-1.398014,0.817853)--(-1.647647,0.743041)--(-1.702920,0.889390)--cycle;
\path[fill=blue!40!white,draw=blue!40!white,line width=.12pt] (1.743993,1.039904)--(1.702920,0.889390)--(1.444914,0.966711)--(1.479763,1.119090)--cycle;
\path[fill=blue!36!white,draw=blue!36!white,line width=.12pt] (-1.479763,1.119090)--(-1.444914,0.966711)--(-1.702920,0.889390)--(-1.743993,1.039904)--cycle;
\path[fill=blue!40!white,draw=blue!40!white,line width=.12pt] (0.741989,0.543496)--(0.696683,0.455508)--(0.468217,0.484122)--(0.498665,0.573971)--cycle;
\path[fill=blue!40!white,draw=blue!40!white,line width=.12pt] (-0.498665,0.573971)--(-0.468217,0.484122)--(-0.696683,0.455508)--(-0.741989,0.543496)--cycle;
\path[fill=blue!28!white,draw=blue!28!white,line width=.12pt] (0.597986,1.888050)--(0.612409,1.804524)--(0.307686,1.827125)--(0.300440,1.910118)--cycle;
\path[fill=blue!26!white,draw=blue!26!white,line width=.12pt] (-0.300440,1.910118)--(-0.307686,1.827125)--(-0.612409,1.804524)--(-0.597986,1.888050)--cycle;
\path[fill=blue!39!white,draw=blue!39!white,line width=.12pt] (-1.398014,0.817853)--(-1.340867,0.678234)--(-1.580295,0.606480)--(-1.647647,0.743041)--cycle;
\path[fill=blue!42!white,draw=blue!42!white,line width=.12pt] (1.647647,0.743041)--(1.580295,0.606480)--(1.340867,0.678234)--(1.398014,0.817853)--cycle;
\path[fill=blue!34!white,draw=blue!34!white,line width=.12pt] (-1.501223,1.269133)--(-1.479763,1.119090)--(-1.743993,1.039904)--(-1.769285,1.188798)--cycle;
\path[fill=blue!38!white,draw=blue!38!white,line width=.12pt] (1.769285,1.188798)--(1.743993,1.039904)--(1.479763,1.119090)--(1.501223,1.269133)--cycle;
\path[fill=blue!31!white,draw=blue!31!white,line width=.12pt] (1.201283,1.715208)--(1.218704,1.606269)--(0.924449,1.658912)--(0.911234,1.767098)--cycle;
\path[fill=blue!28!white,draw=blue!28!white,line width=.12pt] (-0.911234,1.767098)--(-0.924449,1.658912)--(-1.218704,1.606269)--(-1.201283,1.715208)--cycle;
\path[fill=blue!30!white,draw=blue!30!white,line width=.12pt] (-1.218704,1.606269)--(-1.224587,1.478207)--(-1.508470,1.411074)--(-1.501223,1.539458)--cycle;
\path[fill=blue!34!white,draw=blue!34!white,line width=.12pt] (1.501223,1.539458)--(1.508470,1.411074)--(1.224587,1.478207)--(1.218704,1.606269)--cycle;
\path[fill=blue!40!white,draw=blue!40!white,line width=.12pt] (-1.035598,0.609994)--(-0.978167,0.501243)--(-1.204924,0.447619)--(-1.275669,0.553222)--cycle;
\path[fill=blue!41!white,draw=blue!41!white,line width=.12pt] (1.275669,0.553222)--(1.204924,0.447619)--(0.978167,0.501243)--(1.035598,0.609994)--cycle;
\path[fill=blue!26!white,draw=blue!26!white,line width=.12pt] (-0.597986,1.888050)--(-0.612409,1.804524)--(-0.911234,1.767098)--(-0.889774,1.851505)--cycle;
\path[fill=blue!28!white,draw=blue!28!white,line width=.12pt] (0.889774,1.851505)--(0.911234,1.767098)--(0.612409,1.804524)--(0.597986,1.888050)--cycle;
\path[fill=blue!40!white,draw=blue!40!white,line width=.12pt] (0.978167,0.501243)--(0.918440,0.415835)--(0.696683,0.455508)--(0.741989,0.543496)--cycle;
\path[fill=blue!40!white,draw=blue!40!white,line width=.12pt] (-0.741989,0.543496)--(-0.696683,0.455508)--(-0.918440,0.415835)--(-0.978167,0.501243)--cycle;
\path[fill=blue!40!white,draw=blue!40!white,line width=.12pt] (-1.340867,0.678234)--(-1.275669,0.553222)--(-1.503455,0.484957)--(-1.580295,0.606480)--cycle;
\path[fill=blue!42!white,draw=blue!42!white,line width=.12pt] (1.580295,0.606480)--(1.503455,0.484957)--(1.275669,0.553222)--(1.340867,0.678234)--cycle;
\path[fill=blue!32!white,draw=blue!32!white,line width=.12pt] (-1.508470,1.411074)--(-1.501223,1.269133)--(-1.769285,1.188798)--(-1.777825,1.330351)--cycle;
\path[fill=blue!36!white,draw=blue!36!white,line width=.12pt] (1.777825,1.330351)--(1.769285,1.188798)--(1.501223,1.269133)--(1.508470,1.411074)--cycle;
\path[fill=blue!39!white,draw=blue!39!white,line width=.12pt] (0.000000,0.507180)--(0.000000,0.442456)--(-0.219943,0.437053)--(-0.235241,0.501401)--cycle;
\path[fill=blue!39!white,draw=blue!39!white,line width=.12pt] (0.235241,0.501401)--(0.219943,0.437053)--(0.000000,0.442456)--(0.000000,0.507180)--cycle;
\path[fill=blue!39!white,draw=blue!39!white,line width=.12pt] (-0.235241,0.501401)--(-0.219943,0.437053)--(-0.437769,0.420898)--(-0.468217,0.484122)--cycle;
\path[fill=blue!39!white,draw=blue!39!white,line width=.12pt] (0.468217,0.484122)--(0.437769,0.420898)--(0.219943,0.437053)--(0.235241,0.501401)--cycle;
\path[fill=blue!31!white,draw=blue!31!white,line width=.12pt] (1.479763,1.649352)--(1.501223,1.539458)--(1.218704,1.606269)--(1.201283,1.715208)--cycle;
\path[fill=blue!27!white,draw=blue!27!white,line width=.12pt] (-1.201283,1.715208)--(-1.218704,1.606269)--(-1.501223,1.539458)--(-1.479763,1.649352)--cycle;
\path[fill=blue!24!white,draw=blue!24!white,line width=.12pt] (0.290688,1.965600)--(0.300440,1.910118)--(0.000000,1.917498)--(0.000000,1.972741)--cycle;
\path[fill=blue!24!white,draw=blue!24!white,line width=.12pt] (0.000000,1.972741)--(0.000000,1.917498)--(-0.300440,1.910118)--(-0.290688,1.965600)--cycle;
\path[fill=blue!41!white,draw=blue!41!white,line width=.12pt] (1.944527,0.799796)--(1.881411,0.656355)--(1.647647,0.743041)--(1.702920,0.889390)--cycle;
\path[fill=blue!37!white,draw=blue!37!white,line width=.12pt] (-1.702920,0.889390)--(-1.647647,0.743041)--(-1.881411,0.656355)--(-1.944527,0.799796)--cycle;
\path[fill=blue!39!white,draw=blue!39!white,line width=.12pt] (0.696683,0.455508)--(0.651378,0.394145)--(0.437769,0.420898)--(0.468217,0.484122)--cycle;
\path[fill=blue!39!white,draw=blue!39!white,line width=.12pt] (-0.468217,0.484122)--(-0.437769,0.420898)--(-0.651378,0.394145)--(-0.696683,0.455508)--cycle;
\path[fill=blue!40!white,draw=blue!40!white,line width=.12pt] (1.991426,0.948149)--(1.944527,0.799796)--(1.702920,0.889390)--(1.743993,1.039904)--cycle;
\path[fill=blue!35!white,draw=blue!35!white,line width=.12pt] (-1.743993,1.039904)--(-1.702920,0.889390)--(-1.944527,0.799796)--(-1.991426,0.948149)--cycle;
\path[fill=blue!28!white,draw=blue!28!white,line width=.12pt] (1.172992,1.800837)--(1.201283,1.715208)--(0.911234,1.767098)--(0.889774,1.851505)--cycle;
\path[fill=blue!25!white,draw=blue!25!white,line width=.12pt] (-0.889774,1.851505)--(-0.911234,1.767098)--(-1.201283,1.715208)--(-1.172992,1.800837)--cycle;
\path[fill=blue!40!white,draw=blue!40!white,line width=.12pt] (1.204924,0.447619)--(1.131352,0.365485)--(0.918440,0.415835)--(0.978167,0.501243)--cycle;
\path[fill=blue!40!white,draw=blue!40!white,line width=.12pt] (-0.978167,0.501243)--(-0.918440,0.415835)--(-1.131352,0.365485)--(-1.204924,0.447619)--cycle;
\path[fill=blue!40!white,draw=blue!40!white,line width=.12pt] (-1.275669,0.553222)--(-1.204924,0.447619)--(-1.420078,0.383140)--(-1.503455,0.484957)--cycle;
\path[fill=blue!41!white,draw=blue!41!white,line width=.12pt] (1.503455,0.484957)--(1.420078,0.383140)--(1.204924,0.447619)--(1.275669,0.553222)--cycle;
\path[fill=blue!29!white,draw=blue!29!white,line width=.12pt] (-1.501223,1.539458)--(-1.508470,1.411074)--(-1.777825,1.330351)--(-1.769285,1.459123)--cycle;
\path[fill=blue!34!white,draw=blue!34!white,line width=.12pt] (1.769285,1.459123)--(1.777825,1.330351)--(1.508470,1.411074)--(1.501223,1.539458)--cycle;
\path[fill=blue!25!white,draw=blue!25!white,line width=.12pt] (0.578577,1.944248)--(0.597986,1.888050)--(0.300440,1.910118)--(0.290688,1.965600)--cycle;
\path[fill=blue!24!white,draw=blue!24!white,line width=.12pt] (-0.290688,1.965600)--(-0.300440,1.910118)--(-0.597986,1.888050)--(-0.578577,1.944248)--cycle;
\path[fill=blue!42!white,draw=blue!42!white,line width=.12pt] (1.881411,0.656355)--(1.804504,0.523338)--(1.580295,0.606480)--(1.647647,0.743041)--cycle;
\path[fill=blue!38!white,draw=blue!38!white,line width=.12pt] (-1.647647,0.743041)--(-1.580295,0.606480)--(-1.804504,0.523338)--(-1.881411,0.656355)--cycle;
\path[fill=blue!38!white,draw=blue!38!white,line width=.12pt] (2.020307,1.095713)--(1.991426,0.948149)--(1.743993,1.039904)--(1.769285,1.188798)--cycle;
\path[fill=blue!33!white,draw=blue!33!white,line width=.12pt] (-1.769285,1.188798)--(-1.743993,1.039904)--(-1.991426,0.948149)--(-2.020307,1.095713)--cycle;
\path[fill=blue!39!white,draw=blue!39!white,line width=.12pt] (0.918440,0.415835)--(0.858714,0.357051)--(0.651378,0.394145)--(0.696683,0.455508)--cycle;
\path[fill=blue!39!white,draw=blue!39!white,line width=.12pt] (-0.696683,0.455508)--(-0.651378,0.394145)--(-0.858714,0.357051)--(-0.918440,0.415835)--cycle;
\path[fill=blue!39!white,draw=blue!39!white,line width=.12pt] (-1.580295,0.606480)--(-1.503455,0.484957)--(-1.716761,0.405857)--(-1.804504,0.523338)--cycle;
\path[fill=blue!42!white,draw=blue!42!white,line width=.12pt] (1.804504,0.523338)--(1.716761,0.405857)--(1.503455,0.484957)--(1.580295,0.606480)--cycle;
\path[fill=blue!25!white,draw=blue!25!white,line width=.12pt] (0.860893,1.908890)--(0.889774,1.851505)--(0.597986,1.888050)--(0.578577,1.944248)--cycle;
\path[fill=blue!23!white,draw=blue!23!white,line width=.12pt] (-0.578577,1.944248)--(-0.597986,1.888050)--(-0.889774,1.851505)--(-0.860893,1.908890)--cycle;
\path[fill=blue!31!white,draw=blue!31!white,line width=.12pt] (-1.777825,1.330351)--(-1.769285,1.188798)--(-2.020307,1.095713)--(-2.030059,1.236817)--cycle;
\path[fill=blue!36!white,draw=blue!36!white,line width=.12pt] (2.030059,1.236817)--(2.020307,1.095713)--(1.769285,1.188798)--(1.777825,1.330351)--cycle;
\path[fill=blue!37!white,draw=blue!37!white,line width=.12pt] (0.000000,0.442456)--(0.000000,0.406844)--(-0.205234,0.401803)--(-0.219943,0.437053)--cycle;
\path[fill=blue!37!white,draw=blue!37!white,line width=.12pt] (0.219943,0.437053)--(0.205234,0.401803)--(0.000000,0.406844)--(0.000000,0.442456)--cycle;
\path[fill=blue!28!white,draw=blue!28!white,line width=.12pt] (1.444914,1.736532)--(1.479763,1.649352)--(1.201283,1.715208)--(1.172992,1.800837)--cycle;
\path[fill=blue!25!white,draw=blue!25!white,line width=.12pt] (-1.172992,1.800837)--(-1.201283,1.715208)--(-1.479763,1.649352)--(-1.444914,1.736532)--cycle;
\path[fill=blue!37!white,draw=blue!37!white,line width=.12pt] (-0.219943,0.437053)--(-0.205234,0.401803)--(-0.408491,0.386728)--(-0.437769,0.420898)--cycle;
\path[fill=blue!37!white,draw=blue!37!white,line width=.12pt] (0.437769,0.420898)--(0.408491,0.386728)--(0.205234,0.401803)--(0.219943,0.437053)--cycle;
\path[fill=blue!40!white,draw=blue!40!white,line width=.12pt] (-1.204924,0.447619)--(-1.131352,0.365485)--(-1.333369,0.304943)--(-1.420078,0.383140)--cycle;
\path[fill=blue!40!white,draw=blue!40!white,line width=.12pt] (1.420078,0.383140)--(1.333369,0.304943)--(1.131352,0.365485)--(1.204924,0.447619)--cycle;
\path[fill=blue!27!white,draw=blue!27!white,line width=.12pt] (-1.479763,1.649352)--(-1.501223,1.539458)--(-1.769285,1.459123)--(-1.743993,1.570165)--cycle;
\path[fill=blue!31!white,draw=blue!31!white,line width=.12pt] (1.743993,1.570165)--(1.769285,1.459123)--(1.501223,1.539458)--(1.479763,1.649352)--cycle;
\path[fill=blue!39!white,draw=blue!39!white,line width=.12pt] (1.131352,0.365485)--(1.057780,0.309976)--(0.858714,0.357051)--(0.918440,0.415835)--cycle;
\path[fill=blue!39!white,draw=blue!39!white,line width=.12pt] (-0.918440,0.415835)--(-0.858714,0.357051)--(-1.057780,0.309976)--(-1.131352,0.365485)--cycle;
\path[fill=blue!40!white,draw=blue!40!white,line width=.12pt] (-1.503455,0.484957)--(-1.420078,0.383140)--(-1.621555,0.308427)--(-1.716761,0.405857)--cycle;
\path[fill=blue!41!white,draw=blue!41!white,line width=.12pt] (1.716761,0.405857)--(1.621555,0.308427)--(1.420078,0.383140)--(1.503455,0.484957)--cycle;
\path[fill=blue!25!white,draw=blue!25!white,line width=.12pt] (1.134919,1.859866)--(1.172992,1.800837)--(0.889774,1.851505)--(0.860893,1.908890)--cycle;
\path[fill=blue!23!white,draw=blue!23!white,line width=.12pt] (-0.860893,1.908890)--(-0.889774,1.851505)--(-1.172992,1.800837)--(-1.134919,1.859866)--cycle;
\path[fill=blue!22!white,draw=blue!22!white,line width=.12pt] (0.278805,1.991439)--(0.290688,1.965600)--(0.000000,1.972741)--(0.000000,1.998287)--cycle;
\path[fill=blue!21!white,draw=blue!21!white,line width=.12pt] (0.000000,1.998287)--(0.000000,1.972741)--(-0.290688,1.965600)--(-0.278805,1.991439)--cycle;
\path[fill=blue!37!white,draw=blue!37!white,line width=.12pt] (0.651378,0.394145)--(0.607814,0.361764)--(0.408491,0.386728)--(0.437769,0.420898)--cycle;
\path[fill=blue!37!white,draw=blue!37!white,line width=.12pt] (-0.437769,0.420898)--(-0.408491,0.386728)--(-0.607814,0.361764)--(-0.651378,0.394145)--cycle;
\path[fill=blue!41!white,draw=blue!41!white,line width=.12pt] (2.167407,0.698793)--(2.097056,0.558630)--(1.881411,0.656355)--(1.944527,0.799796)--cycle;
\path[fill=blue!37!white,draw=blue!37!white,line width=.12pt] (-1.944527,0.799796)--(-1.881411,0.656355)--(-2.097056,0.558630)--(-2.167407,0.698793)--cycle;
\path[fill=blue!34!white,draw=blue!34!white,line width=.12pt] (2.020307,1.366038)--(2.030059,1.236817)--(1.777825,1.330351)--(1.769285,1.459123)--cycle;
\path[fill=blue!28!white,draw=blue!28!white,line width=.12pt] (-1.769285,1.459123)--(-1.777825,1.330351)--(-2.030059,1.236817)--(-2.020307,1.366038)--cycle;
\path[fill=blue!22!white,draw=blue!22!white,line width=.12pt] (0.554926,1.970959)--(0.578577,1.944248)--(0.290688,1.965600)--(0.278805,1.991439)--cycle;
\path[fill=blue!21!white,draw=blue!21!white,line width=.12pt] (-0.278805,1.991439)--(-0.290688,1.965600)--(-0.578577,1.944248)--(-0.554926,1.970959)--cycle;
\path[fill=blue!40!white,draw=blue!40!white,line width=.12pt] (2.219681,0.844710)--(2.167407,0.698793)--(1.944527,0.799796)--(1.991426,0.948149)--cycle;
\path[fill=blue!35!white,draw=blue!35!white,line width=.12pt] (-1.991426,0.948149)--(-1.944527,0.799796)--(-2.167407,0.698793)--(-2.219681,0.844710)--cycle;
\path[fill=blue!42!white,draw=blue!42!white,line width=.12pt] (2.097056,0.558630)--(2.011334,0.429608)--(1.804504,0.523338)--(1.881411,0.656355)--cycle;
\path[fill=blue!38!white,draw=blue!38!white,line width=.12pt] (-1.881411,0.656355)--(-1.804504,0.523338)--(-2.011334,0.429608)--(-2.097056,0.558630)--cycle;
\path[fill=blue!37!white,draw=blue!37!white,line width=.12pt] (0.858714,0.357051)--(0.801283,0.327152)--(0.607814,0.361764)--(0.651378,0.394145)--cycle;
\path[fill=blue!38!white,draw=blue!38!white,line width=.12pt] (-0.651378,0.394145)--(-0.607814,0.361764)--(-0.801283,0.327152)--(-0.858714,0.357051)--cycle;
\path[fill=blue!33!white,draw=blue!33!white,line width=.12pt] (-2.020307,1.095713)--(-1.991426,0.948149)--(-2.219681,0.844710)--(-2.251872,0.990774)--cycle;
\path[fill=blue!38!white,draw=blue!38!white,line width=.12pt] (2.251872,0.990774)--(2.219681,0.844710)--(1.991426,0.948149)--(2.020307,1.095713)--cycle;
\path[fill=blue!22!white,draw=blue!22!white,line width=.12pt] (0.825702,1.937046)--(0.860893,1.908890)--(0.578577,1.944248)--(0.554926,1.970959)--cycle;
\path[fill=blue!21!white,draw=blue!21!white,line width=.12pt] (-0.554926,1.970959)--(-0.578577,1.944248)--(-0.860893,1.908890)--(-0.825702,1.937046)--cycle;
\path[fill=blue!39!white,draw=blue!39!white,line width=.12pt] (-1.804504,0.523338)--(-1.716761,0.405857)--(-1.913535,0.316685)--(-2.011334,0.429608)--cycle;
\path[fill=blue!42!white,draw=blue!42!white,line width=.12pt] (2.011334,0.429608)--(1.913535,0.316685)--(1.716761,0.405857)--(1.804504,0.523338)--cycle;
\path[fill=blue!39!white,draw=blue!39!white,line width=.12pt] (-1.131352,0.365485)--(-1.057780,0.309976)--(-1.246659,0.253371)--(-1.333369,0.304943)--cycle;
\path[fill=blue!39!white,draw=blue!39!white,line width=.12pt] (1.333369,0.304943)--(1.246659,0.253371)--(1.057780,0.309976)--(1.131352,0.365485)--cycle;
\path[fill=blue!40!white,draw=blue!40!white,line width=.12pt] (-1.420078,0.383140)--(-1.333369,0.304943)--(-1.522544,0.234792)--(-1.621555,0.308427)--cycle;
\path[fill=blue!40!white,draw=blue!40!white,line width=.12pt] (1.621555,0.308427)--(1.522544,0.234792)--(1.333369,0.304943)--(1.420078,0.383140)--cycle;
\path[fill=blue!28!white,draw=blue!28!white,line width=.12pt] (1.702920,1.659211)--(1.743993,1.570165)--(1.479763,1.649352)--(1.444914,1.736532)--cycle;
\path[fill=blue!24!white,draw=blue!24!white,line width=.12pt] (-1.444914,1.736532)--(-1.479763,1.649352)--(-1.743993,1.570165)--(-1.702920,1.659211)--cycle;
\path[fill=blue!25!white,draw=blue!25!white,line width=.12pt] (1.398014,1.797648)--(1.444914,1.736532)--(1.172992,1.800837)--(1.134919,1.859866)--cycle;
\path[fill=blue!22!white,draw=blue!22!white,line width=.12pt] (-1.134919,1.859866)--(-1.172992,1.800837)--(-1.444914,1.736532)--(-1.398014,1.797648)--cycle;
\path[fill=blue!35!white,draw=blue!35!white,line width=.12pt] (0.205234,0.401803)--(0.191677,0.397005)--(0.000000,0.401713)--(0.000000,0.406844)--cycle;
\path[fill=blue!35!white,draw=blue!35!white,line width=.12pt] (0.000000,0.406844)--(0.000000,0.401713)--(-0.191677,0.397005)--(-0.205234,0.401803)--cycle;
\path[fill=blue!30!white,draw=blue!30!white,line width=.12pt] (-2.030059,1.236817)--(-2.020307,1.095713)--(-2.251872,0.990774)--(-2.262742,1.131371)--cycle;
\path[fill=blue!36!white,draw=blue!36!white,line width=.12pt] (2.262742,1.131371)--(2.251872,0.990774)--(2.020307,1.095713)--(2.030059,1.236817)--cycle;
\path[fill=blue!35!white,draw=blue!35!white,line width=.12pt] (0.408491,0.386728)--(0.381508,0.382925)--(0.191677,0.397005)--(0.205234,0.401803)--cycle;
\path[fill=blue!35!white,draw=blue!35!white,line width=.12pt] (-0.205234,0.401803)--(-0.191677,0.397005)--(-0.381508,0.382925)--(-0.408491,0.386728)--cycle;
\path[fill=blue!31!white,draw=blue!31!white,line width=.12pt] (1.991426,1.478411)--(2.020307,1.366038)--(1.769285,1.459123)--(1.743993,1.570165)--cycle;
\path[fill=blue!26!white,draw=blue!26!white,line width=.12pt] (-1.743993,1.570165)--(-1.769285,1.459123)--(-2.020307,1.366038)--(-1.991426,1.478411)--cycle;
\path[fill=blue!37!white,draw=blue!37!white,line width=.12pt] (1.057780,0.309976)--(0.987036,0.283224)--(0.801283,0.327152)--(0.858714,0.357051)--cycle;
\path[fill=blue!38!white,draw=blue!38!white,line width=.12pt] (-0.858714,0.357051)--(-0.801283,0.327152)--(-0.987036,0.283224)--(-1.057780,0.309976)--cycle;
\path[fill=blue!41!white,draw=blue!41!white,line width=.12pt] (1.913535,0.316685)--(1.807416,0.224200)--(1.621555,0.308427)--(1.716761,0.405857)--cycle;
\path[fill=blue!40!white,draw=blue!40!white,line width=.12pt] (-1.716761,0.405857)--(-1.621555,0.308427)--(-1.807416,0.224200)--(-1.913535,0.316685)--cycle;
\path[fill=blue!20!white,draw=blue!20!white,line width=.12pt] (-0.825702,1.937046)--(-0.860893,1.908890)--(-1.134919,1.859866)--(-1.088526,1.890026)--cycle;
\path[fill=blue!22!white,draw=blue!22!white,line width=.12pt] (1.088526,1.890026)--(1.134919,1.859866)--(0.860893,1.908890)--(0.825702,1.937046)--cycle;
\path[fill=blue!35!white,draw=blue!35!white,line width=.12pt] (0.607814,0.361764)--(0.567664,0.359610)--(0.381508,0.382925)--(0.408491,0.386728)--cycle;
\path[fill=blue!36!white,draw=blue!36!white,line width=.12pt] (-0.408491,0.386728)--(-0.381508,0.382925)--(-0.567664,0.359610)--(-0.607814,0.361764)--cycle;
\path[fill=blue!19!white,draw=blue!19!white,line width=.12pt] (0.265249,1.986640)--(0.278805,1.991439)--(0.000000,1.998287)--(0.000000,1.993156)--cycle;
\path[fill=blue!19!white,draw=blue!19!white,line width=.12pt] (0.000000,1.993156)--(0.000000,1.998287)--(-0.278805,1.991439)--(-0.265249,1.986640)--cycle;
\path[fill=blue!19!white,draw=blue!19!white,line width=.12pt] (0.527943,1.967157)--(0.554926,1.970959)--(0.278805,1.991439)--(0.265249,1.986640)--cycle;
\path[fill=blue!19!white,draw=blue!19!white,line width=.12pt] (-0.265249,1.986640)--(-0.278805,1.991439)--(-0.554926,1.970959)--(-0.527943,1.967157)--cycle;
\path[fill=blue!39!white,draw=blue!39!white,line width=.12pt] (-1.333369,0.304943)--(-1.246659,0.253371)--(-1.423533,0.187782)--(-1.522544,0.234792)--cycle;
\path[fill=blue!39!white,draw=blue!39!white,line width=.12pt] (1.522544,0.234792)--(1.423533,0.187782)--(1.246659,0.253371)--(1.333369,0.304943)--cycle;
\path[fill=blue!28!white,draw=blue!28!white,line width=.12pt] (-2.020307,1.366038)--(-2.030059,1.236817)--(-2.262742,1.131371)--(-2.251872,1.261099)--cycle;
\path[fill=blue!34!white,draw=blue!34!white,line width=.12pt] (2.251872,1.261099)--(2.262742,1.131371)--(2.030059,1.236817)--(2.020307,1.366038)--cycle;
\path[fill=blue!41!white,draw=blue!41!white,line width=.12pt] (2.369413,0.587353)--(2.292506,0.450808)--(2.097056,0.558630)--(2.167407,0.698793)--cycle;
\path[fill=blue!36!white,draw=blue!36!white,line width=.12pt] (-2.167407,0.698793)--(-2.097056,0.558630)--(-2.292506,0.450808)--(-2.369413,0.587353)--cycle;
\path[fill=blue!34!white,draw=blue!34!white,line width=.12pt] (0.801283,0.327152)--(0.748354,0.327284)--(0.567664,0.359610)--(0.607814,0.361764)--cycle;
\path[fill=blue!36!white,draw=blue!36!white,line width=.12pt] (-0.607814,0.361764)--(-0.567664,0.359610)--(-0.748354,0.327284)--(-0.801283,0.327152)--cycle;
\path[fill=blue!42!white,draw=blue!42!white,line width=.12pt] (2.292506,0.450808)--(2.198794,0.326193)--(2.011334,0.429608)--(2.097056,0.558630)--cycle;
\path[fill=blue!38!white,draw=blue!38!white,line width=.12pt] (-2.097056,0.558630)--(-2.011334,0.429608)--(-2.198794,0.326193)--(-2.292506,0.450808)--cycle;
\path[fill=blue!37!white,draw=blue!37!white,line width=.12pt] (1.246659,0.253371)--(1.163283,0.230405)--(0.987036,0.283224)--(1.057780,0.309976)--cycle;
\path[fill=blue!38!white,draw=blue!38!white,line width=.12pt] (-1.057780,0.309976)--(-0.987036,0.283224)--(-1.163283,0.230405)--(-1.246659,0.253371)--cycle;
\path[fill=blue!40!white,draw=blue!40!white,line width=.12pt] (2.426560,0.730582)--(2.369413,0.587353)--(2.167407,0.698793)--(2.219681,0.844710)--cycle;
\path[fill=blue!34!white,draw=blue!34!white,line width=.12pt] (-2.219681,0.844710)--(-2.167407,0.698793)--(-2.369413,0.587353)--(-2.426560,0.730582)--cycle;
\path[fill=blue!25!white,draw=blue!25!white,line width=.12pt] (1.647647,1.722837)--(1.702920,1.659211)--(1.444914,1.736532)--(1.398014,1.797648)--cycle;
\path[fill=blue!22!white,draw=blue!22!white,line width=.12pt] (-1.398014,1.797648)--(-1.444914,1.736532)--(-1.702920,1.659211)--(-1.647647,1.722837)--cycle;
\path[fill=blue!40!white,draw=blue!40!white,line width=.12pt] (-1.621555,0.308427)--(-1.522544,0.234792)--(-1.697056,0.155708)--(-1.807416,0.224200)--cycle;
\path[fill=blue!40!white,draw=blue!40!white,line width=.12pt] (1.807416,0.224200)--(1.697056,0.155708)--(1.522544,0.234792)--(1.621555,0.308427)--cycle;
\path[fill=blue!18!white,draw=blue!18!white,line width=.12pt] (-0.527943,1.967157)--(-0.554926,1.970959)--(-0.825702,1.937046)--(-0.785553,1.934893)--cycle;
\path[fill=blue!19!white,draw=blue!19!white,line width=.12pt] (0.785553,1.934893)--(0.825702,1.937046)--(0.554926,1.970959)--(0.527943,1.967157)--cycle;
\path[fill=blue!28!white,draw=blue!28!white,line width=.12pt] (1.944527,1.569617)--(1.991426,1.478411)--(1.743993,1.570165)--(1.702920,1.659211)--cycle;
\path[fill=blue!23!white,draw=blue!23!white,line width=.12pt] (-1.702920,1.659211)--(-1.743993,1.570165)--(-1.991426,1.478411)--(-1.944527,1.569617)--cycle;
\path[fill=blue!22!white,draw=blue!22!white,line width=.12pt] (1.340867,1.830352)--(1.398014,1.797648)--(1.134919,1.859866)--(1.088526,1.890026)--cycle;
\path[fill=blue!20!white,draw=blue!20!white,line width=.12pt] (-1.088526,1.890026)--(-1.134919,1.859866)--(-1.398014,1.797648)--(-1.340867,1.830352)--cycle;
\path[fill=blue!39!white,draw=blue!39!white,line width=.12pt] (-2.011334,0.429608)--(-1.913535,0.316685)--(-2.091880,0.218298)--(-2.198794,0.326193)--cycle;
\path[fill=blue!42!white,draw=blue!42!white,line width=.12pt] (2.198794,0.326193)--(2.091880,0.218298)--(1.913535,0.316685)--(2.011334,0.429608)--cycle;
\path[fill=blue!32!white,draw=blue!32!white,line width=.12pt] (-2.251872,0.990774)--(-2.219681,0.844710)--(-2.426560,0.730582)--(-2.461751,0.874991)--cycle;
\path[fill=blue!38!white,draw=blue!38!white,line width=.12pt] (2.461751,0.874991)--(2.426560,0.730582)--(2.219681,0.844710)--(2.251872,0.990774)--cycle;
\path[fill=blue!32!white,draw=blue!32!white,line width=.12pt] (0.191677,0.397005)--(0.179794,0.422843)--(0.000000,0.427259)--(0.000000,0.401713)--cycle;
\path[fill=blue!33!white,draw=blue!33!white,line width=.12pt] (0.000000,0.401713)--(0.000000,0.427259)--(-0.179794,0.422843)--(-0.191677,0.397005)--cycle;
\path[fill=blue!36!white,draw=blue!36!white,line width=.12pt] (-0.801283,0.327152)--(-0.748354,0.327284)--(-0.921837,0.286259)--(-0.987036,0.283224)--cycle;
\path[fill=blue!34!white,draw=blue!34!white,line width=.12pt] (0.987036,0.283224)--(0.921837,0.286259)--(0.748354,0.327284)--(0.801283,0.327152)--cycle;
\path[fill=blue!33!white,draw=blue!33!white,line width=.12pt] (-0.191677,0.397005)--(-0.179794,0.422843)--(-0.357857,0.409636)--(-0.381508,0.382925)--cycle;
\path[fill=blue!32!white,draw=blue!32!white,line width=.12pt] (0.381508,0.382925)--(0.357857,0.409636)--(0.179794,0.422843)--(0.191677,0.397005)--cycle;
\path[fill=blue!25!white,draw=blue!25!white,line width=.12pt] (-1.991426,1.478411)--(-2.020307,1.366038)--(-2.251872,1.261099)--(-2.219681,1.374972)--cycle;
\path[fill=blue!31!white,draw=blue!31!white,line width=.12pt] (2.219681,1.374972)--(2.251872,1.261099)--(2.020307,1.366038)--(1.991426,1.478411)--cycle;
\path[fill=blue!39!white,draw=blue!39!white,line width=.12pt] (-1.913535,0.316685)--(-1.807416,0.224200)--(-1.975871,0.131270)--(-2.091880,0.218298)--cycle;
\path[fill=blue!41!white,draw=blue!41!white,line width=.12pt] (2.091880,0.218298)--(1.975871,0.131270)--(1.807416,0.224200)--(1.913535,0.316685)--cycle;
\path[fill=blue!20!white,draw=blue!20!white,line width=.12pt] (1.035598,1.890159)--(1.088526,1.890026)--(0.825702,1.937046)--(0.785553,1.934893)--cycle;
\path[fill=blue!18!white,draw=blue!18!white,line width=.12pt] (-0.785553,1.934893)--(-0.825702,1.937046)--(-1.088526,1.890026)--(-1.035598,1.890159)--cycle;
\path[fill=blue!32!white,draw=blue!32!white,line width=.12pt] (0.567664,0.359610)--(0.532473,0.387767)--(0.357857,0.409636)--(0.381508,0.382925)--cycle;
\path[fill=blue!33!white,draw=blue!33!white,line width=.12pt] (-0.381508,0.382925)--(-0.357857,0.409636)--(-0.532473,0.387767)--(-0.567664,0.359610)--cycle;
\path[fill=blue!29!white,draw=blue!29!white,line width=.12pt] (-2.262742,1.131371)--(-2.251872,0.990774)--(-2.461751,0.874991)--(-2.473633,1.015029)--cycle;
\path[fill=blue!36!white,draw=blue!36!white,line width=.12pt] (2.473633,1.015029)--(2.461751,0.874991)--(2.251872,0.990774)--(2.262742,1.131371)--cycle;
\path[fill=blue!17!white,draw=blue!17!white,line width=.12pt] (0.000000,1.957544)--(0.000000,1.993156)--(-0.265249,1.986640)--(-0.250539,1.951390)--cycle;
\path[fill=blue!17!white,draw=blue!17!white,line width=.12pt] (0.250539,1.951390)--(0.265249,1.986640)--(0.000000,1.993156)--(0.000000,1.957544)--cycle;
\path[fill=blue!37!white,draw=blue!37!white,line width=.12pt] (1.423533,0.187782)--(1.328326,0.169203)--(1.163283,0.230405)--(1.246659,0.253371)--cycle;
\path[fill=blue!38!white,draw=blue!38!white,line width=.12pt] (-1.246659,0.253371)--(-1.163283,0.230405)--(-1.328326,0.169203)--(-1.423533,0.187782)--cycle;
\path[fill=blue!39!white,draw=blue!39!white,line width=.12pt] (-1.522544,0.234792)--(-1.423533,0.187782)--(-1.586697,0.113840)--(-1.697056,0.155708)--cycle;
\path[fill=blue!39!white,draw=blue!39!white,line width=.12pt] (1.697056,0.155708)--(1.586697,0.113840)--(1.423533,0.187782)--(1.522544,0.234792)--cycle;
\path[fill=blue!17!white,draw=blue!17!white,line width=.12pt] (0.498665,1.932987)--(0.527943,1.967157)--(0.265249,1.986640)--(0.250539,1.951390)--cycle;
\path[fill=blue!17!white,draw=blue!17!white,line width=.12pt] (-0.250539,1.951390)--(-0.265249,1.986640)--(-0.527943,1.967157)--(-0.498665,1.932987)--cycle;
\path[fill=blue!36!white,draw=blue!36!white,line width=.12pt] (-0.987036,0.283224)--(-0.921837,0.286259)--(-1.086442,0.236929)--(-1.163283,0.230405)--cycle;
\path[fill=blue!34!white,draw=blue!34!white,line width=.12pt] (1.163283,0.230405)--(1.086442,0.236929)--(0.921837,0.286259)--(0.987036,0.283224)--cycle;
\path[fill=blue!34!white,draw=blue!34!white,line width=.12pt] (-0.567664,0.359610)--(-0.532473,0.387767)--(-0.701962,0.357445)--(-0.748354,0.327284)--cycle;
\path[fill=blue!32!white,draw=blue!32!white,line width=.12pt] (0.748354,0.327284)--(0.701962,0.357445)--(0.532473,0.387767)--(0.567664,0.359610)--cycle;
\path[fill=blue!19!white,draw=blue!19!white,line width=.12pt] (-1.340867,1.830352)--(-1.398014,1.797648)--(-1.647647,1.722837)--(-1.580295,1.758598)--cycle;
\path[fill=blue!22!white,draw=blue!22!white,line width=.12pt] (1.580295,1.758598)--(1.647647,1.722837)--(1.398014,1.797648)--(1.340867,1.830352)--cycle;
\path[fill=blue!21!white,draw=blue!21!white,line width=.12pt] (-1.647647,1.722837)--(-1.702920,1.659211)--(-1.944527,1.569617)--(-1.881411,1.636151)--cycle;
\path[fill=blue!25!white,draw=blue!25!white,line width=.12pt] (1.881411,1.636151)--(1.944527,1.569617)--(1.702920,1.659211)--(1.647647,1.722837)--cycle;
\path[fill=blue!40!white,draw=blue!40!white,line width=.12pt] (1.975871,0.131270)--(1.855225,0.068452)--(1.697056,0.155708)--(1.807416,0.224200)--cycle;
\path[fill=blue!40!white,draw=blue!40!white,line width=.12pt] (-1.807416,0.224200)--(-1.697056,0.155708)--(-1.855225,0.068452)--(-1.975871,0.131270)--cycle;
\path[fill=blue!17!white,draw=blue!17!white,line width=.12pt] (0.741989,1.902512)--(0.785553,1.934893)--(0.527943,1.967157)--(0.498665,1.932987)--cycle;
\path[fill=blue!17!white,draw=blue!17!white,line width=.12pt] (-0.498665,1.932987)--(-0.527943,1.967157)--(-0.785553,1.934893)--(-0.741989,1.902512)--cycle;
\path[fill=blue!20!white,draw=blue!20!white,line width=.12pt] (1.275669,1.833387)--(1.340867,1.830352)--(1.088526,1.890026)--(1.035598,1.890159)--cycle;
\path[fill=blue!18!white,draw=blue!18!white,line width=.12pt] (-1.035598,1.890159)--(-1.088526,1.890026)--(-1.340867,1.830352)--(-1.275669,1.833387)--cycle;
\path[fill=blue!34!white,draw=blue!34!white,line width=.12pt] (2.461751,1.145316)--(2.473633,1.015029)--(2.262742,1.131371)--(2.251872,1.261099)--cycle;
\path[fill=blue!27!white,draw=blue!27!white,line width=.12pt] (-2.251872,1.261099)--(-2.262742,1.131371)--(-2.473633,1.015029)--(-2.461751,1.145316)--cycle;
\path[fill=blue!41!white,draw=blue!41!white,line width=.12pt] (2.548600,0.466550)--(2.465877,0.333925)--(2.292506,0.450808)--(2.369413,0.587353)--cycle;
\path[fill=blue!35!white,draw=blue!35!white,line width=.12pt] (-2.369413,0.587353)--(-2.292506,0.450808)--(-2.465877,0.333925)--(-2.548600,0.466550)--cycle;
\path[fill=blue!42!white,draw=blue!42!white,line width=.12pt] (2.465877,0.333925)--(2.365079,0.214089)--(2.198794,0.326193)--(2.292506,0.450808)--cycle;
\path[fill=blue!37!white,draw=blue!37!white,line width=.12pt] (-2.292506,0.450808)--(-2.198794,0.326193)--(-2.365079,0.214089)--(-2.465877,0.333925)--cycle;
\path[fill=blue!28!white,draw=blue!28!white,line width=.12pt] (2.167407,1.468614)--(2.219681,1.374972)--(1.991426,1.478411)--(1.944527,1.569617)--cycle;
\path[fill=blue!23!white,draw=blue!23!white,line width=.12pt] (-1.944527,1.569617)--(-1.991426,1.478411)--(-2.219681,1.374972)--(-2.167407,1.468614)--cycle;
\path[fill=blue!30!white,draw=blue!30!white,line width=.12pt] (0.179794,0.422843)--(0.170043,0.478325)--(0.000000,0.482502)--(0.000000,0.427259)--cycle;
\path[fill=blue!30!white,draw=blue!30!white,line width=.12pt] (0.000000,0.427259)--(0.000000,0.482502)--(-0.170043,0.478325)--(-0.179794,0.422843)--cycle;
\path[fill=blue!32!white,draw=blue!32!white,line width=.12pt] (0.921837,0.286259)--(0.864690,0.318963)--(0.701962,0.357445)--(0.748354,0.327284)--cycle;
\path[fill=blue!34!white,draw=blue!34!white,line width=.12pt] (-0.748354,0.327284)--(-0.701962,0.357445)--(-0.864690,0.318963)--(-0.921837,0.286259)--cycle;
\path[fill=blue!40!white,draw=blue!40!white,line width=.12pt] (2.610069,0.606865)--(2.548600,0.466550)--(2.369413,0.587353)--(2.426560,0.730582)--cycle;
\path[fill=blue!33!white,draw=blue!33!white,line width=.12pt] (-2.426560,0.730582)--(-2.369413,0.587353)--(-2.548600,0.466550)--(-2.610069,0.606865)--cycle;
\path[fill=blue!42!white,draw=blue!42!white,line width=.12pt] (2.365079,0.214089)--(2.250079,0.111645)--(2.091880,0.218298)--(2.198794,0.326193)--cycle;
\path[fill=blue!38!white,draw=blue!38!white,line width=.12pt] (-2.198794,0.326193)--(-2.091880,0.218298)--(-2.250079,0.111645)--(-2.365079,0.214089)--cycle;
\path[fill=blue!29!white,draw=blue!29!white,line width=.12pt] (0.357857,0.409636)--(0.338447,0.465835)--(0.170043,0.478325)--(0.179794,0.422843)--cycle;
\path[fill=blue!30!white,draw=blue!30!white,line width=.12pt] (-0.179794,0.422843)--(-0.170043,0.478325)--(-0.338447,0.465835)--(-0.357857,0.409636)--cycle;
\path[fill=blue!38!white,draw=blue!38!white,line width=.12pt] (-1.423533,0.187782)--(-1.328326,0.169203)--(-1.480578,0.100206)--(-1.586697,0.113840)--cycle;
\path[fill=blue!37!white,draw=blue!37!white,line width=.12pt] (1.586697,0.113840)--(1.480578,0.100206)--(1.328326,0.169203)--(1.423533,0.187782)--cycle;
\path[fill=blue!16!white,draw=blue!16!white,line width=.12pt] (-0.741989,1.902512)--(-0.785553,1.934893)--(-1.035598,1.890159)--(-0.978167,1.860259)--cycle;
\path[fill=blue!17!white,draw=blue!17!white,line width=.12pt] (0.978167,1.860259)--(1.035598,1.890159)--(0.785553,1.934893)--(0.741989,1.902512)--cycle;
\path[fill=blue!37!white,draw=blue!37!white,line width=.12pt] (-1.163283,0.230405)--(-1.086442,0.236929)--(-1.240584,0.179769)--(-1.328326,0.169203)--cycle;
\path[fill=blue!34!white,draw=blue!34!white,line width=.12pt] (1.328326,0.169203)--(1.240584,0.179769)--(1.086442,0.236929)--(1.163283,0.230405)--cycle;
\path[fill=blue!29!white,draw=blue!29!white,line width=.12pt] (0.532473,0.387767)--(0.503593,0.445151)--(0.338447,0.465835)--(0.357857,0.409636)--cycle;
\path[fill=blue!31!white,draw=blue!31!white,line width=.12pt] (-0.357857,0.409636)--(-0.338447,0.465835)--(-0.503593,0.445151)--(-0.532473,0.387767)--cycle;
\path[fill=blue!39!white,draw=blue!39!white,line width=.12pt] (1.855225,0.068452)--(1.734580,0.032258)--(1.586697,0.113840)--(1.697056,0.155708)--cycle;
\path[fill=blue!39!white,draw=blue!39!white,line width=.12pt] (-1.697056,0.155708)--(-1.586697,0.113840)--(-1.734580,0.032258)--(-1.855225,0.068452)--cycle;
\path[fill=blue!38!white,draw=blue!38!white,line width=.12pt] (2.647922,0.749480)--(2.610069,0.606865)--(2.426560,0.730582)--(2.461751,0.874991)--cycle;
\path[fill=blue!31!white,draw=blue!31!white,line width=.12pt] (-2.461751,0.874991)--(-2.426560,0.730582)--(-2.610069,0.606865)--(-2.647922,0.749480)--cycle;
\path[fill=blue!15!white,draw=blue!15!white,line width=.12pt] (0.000000,1.892820)--(0.000000,1.957544)--(-0.250539,1.951390)--(-0.235241,1.887042)--cycle;
\path[fill=blue!15!white,draw=blue!15!white,line width=.12pt] (0.235241,1.887042)--(0.250539,1.951390)--(0.000000,1.957544)--(0.000000,1.892820)--cycle;
\path[fill=blue!41!white,draw=blue!41!white,line width=.12pt] (2.250079,0.111645)--(2.125296,0.030531)--(1.975871,0.131270)--(2.091880,0.218298)--cycle;
\path[fill=blue!39!white,draw=blue!39!white,line width=.12pt] (-2.091880,0.218298)--(-1.975871,0.131270)--(-2.125296,0.030531)--(-2.250079,0.111645)--cycle;
\path[fill=blue!15!white,draw=blue!15!white,line width=.12pt] (-0.235241,1.887042)--(-0.250539,1.951390)--(-0.498665,1.932987)--(-0.468217,1.869763)--cycle;
\path[fill=blue!15!white,draw=blue!15!white,line width=.12pt] (0.468217,1.869763)--(0.498665,1.932987)--(0.250539,1.951390)--(0.235241,1.887042)--cycle;
\path[fill=blue!31!white,draw=blue!31!white,line width=.12pt] (2.426560,1.260844)--(2.461751,1.145316)--(2.251872,1.261099)--(2.219681,1.374972)--cycle;
\path[fill=blue!24!white,draw=blue!24!white,line width=.12pt] (-2.219681,1.374972)--(-2.251872,1.261099)--(-2.461751,1.145316)--(-2.426560,1.260844)--cycle;
\path[fill=blue!35!white,draw=blue!35!white,line width=.12pt] (-0.921837,0.286259)--(-0.864690,0.318963)--(-1.019091,0.272690)--(-1.086442,0.236929)--cycle;
\path[fill=blue!32!white,draw=blue!32!white,line width=.12pt] (1.086442,0.236929)--(1.019091,0.272690)--(0.864690,0.318963)--(0.921837,0.286259)--cycle;
\path[fill=blue!22!white,draw=blue!22!white,line width=.12pt] (1.804504,1.675456)--(1.881411,1.636151)--(1.647647,1.722837)--(1.580295,1.758598)--cycle;
\path[fill=blue!19!white,draw=blue!19!white,line width=.12pt] (-1.580295,1.758598)--(-1.647647,1.722837)--(-1.881411,1.636151)--(-1.804504,1.675456)--cycle;
\path[fill=blue!18!white,draw=blue!18!white,line width=.12pt] (-1.275669,1.833387)--(-1.340867,1.830352)--(-1.580295,1.758598)--(-1.503455,1.765122)--cycle;
\path[fill=blue!20!white,draw=blue!20!white,line width=.12pt] (1.503455,1.765122)--(1.580295,1.758598)--(1.340867,1.830352)--(1.275669,1.833387)--cycle;
\path[fill=blue!29!white,draw=blue!29!white,line width=.12pt] (0.701962,0.357445)--(0.663889,0.416474)--(0.503593,0.445151)--(0.532473,0.387767)--cycle;
\path[fill=blue!31!white,draw=blue!31!white,line width=.12pt] (-0.532473,0.387767)--(-0.503593,0.445151)--(-0.663889,0.416474)--(-0.701962,0.357445)--cycle;
\path[fill=blue!36!white,draw=blue!36!white,line width=.12pt] (2.660703,0.888912)--(2.647922,0.749480)--(2.461751,0.874991)--(2.473633,1.015029)--cycle;
\path[fill=blue!29!white,draw=blue!29!white,line width=.12pt] (-2.473633,1.015029)--(-2.461751,0.874991)--(-2.647922,0.749480)--(-2.660703,0.888912)--cycle;
\path[fill=blue!15!white,draw=blue!15!white,line width=.12pt] (0.696683,1.841149)--(0.741989,1.902512)--(0.498665,1.932987)--(0.468217,1.869763)--cycle;
\path[fill=blue!15!white,draw=blue!15!white,line width=.12pt] (-0.468217,1.869763)--(-0.498665,1.932987)--(-0.741989,1.902512)--(-0.696683,1.841149)--cycle;
\path[fill=blue!25!white,draw=blue!25!white,line width=.12pt] (2.097056,1.538426)--(2.167407,1.468614)--(1.944527,1.569617)--(1.881411,1.636151)--cycle;
\path[fill=blue!21!white,draw=blue!21!white,line width=.12pt] (-1.881411,1.636151)--(-1.944527,1.569617)--(-2.167407,1.468614)--(-2.097056,1.538426)--cycle;
\path[fill=blue!17!white,draw=blue!17!white,line width=.12pt] (1.204924,1.806635)--(1.275669,1.833387)--(1.035598,1.890159)--(0.978167,1.860259)--cycle;
\path[fill=blue!16!white,draw=blue!16!white,line width=.12pt] (-0.978167,1.860259)--(-1.035598,1.890159)--(-1.275669,1.833387)--(-1.204924,1.806635)--cycle;
\path[fill=blue!40!white,draw=blue!40!white,line width=.12pt] (2.125296,0.030531)--(1.995527,-0.026136)--(1.855225,0.068452)--(1.975871,0.131270)--cycle;
\path[fill=blue!40!white,draw=blue!40!white,line width=.12pt] (-1.975871,0.131270)--(-1.855225,0.068452)--(-1.995527,-0.026136)--(-2.125296,0.030531)--cycle;
\path[fill=blue!27!white,draw=blue!27!white,line width=.12pt] (0.000000,0.482502)--(0.000000,0.565317)--(-0.162796,0.561318)--(-0.170043,0.478325)--cycle;
\path[fill=blue!27!white,draw=blue!27!white,line width=.12pt] (0.170043,0.478325)--(0.162796,0.561318)--(0.000000,0.565317)--(0.000000,0.482502)--cycle;
\path[fill=blue!37!white,draw=blue!37!white,line width=.12pt] (-1.328326,0.169203)--(-1.240584,0.179769)--(-1.382778,0.115330)--(-1.480578,0.100206)--cycle;
\path[fill=blue!34!white,draw=blue!34!white,line width=.12pt] (1.480578,0.100206)--(1.382778,0.115330)--(1.240584,0.179769)--(1.328326,0.169203)--cycle;
\path[fill=blue!29!white,draw=blue!29!white,line width=.12pt] (0.864690,0.318963)--(0.817791,0.380079)--(0.663889,0.416474)--(0.701962,0.357445)--cycle;
\path[fill=blue!32!white,draw=blue!32!white,line width=.12pt] (-0.701962,0.357445)--(-0.663889,0.416474)--(-0.817791,0.380079)--(-0.864690,0.318963)--cycle;
\path[fill=blue!39!white,draw=blue!39!white,line width=.12pt] (-1.586697,0.113840)--(-1.480578,0.100206)--(-1.618570,0.024081)--(-1.734580,0.032258)--cycle;
\path[fill=blue!37!white,draw=blue!37!white,line width=.12pt] (1.734580,0.032258)--(1.618570,0.024081)--(1.480578,0.100206)--(1.586697,0.113840)--cycle;
\path[fill=blue!26!white,draw=blue!26!white,line width=.12pt] (0.338447,0.465835)--(0.324025,0.549360)--(0.162796,0.561318)--(0.170043,0.478325)--cycle;
\path[fill=blue!28!white,draw=blue!28!white,line width=.12pt] (-0.170043,0.478325)--(-0.162796,0.561318)--(-0.324025,0.549360)--(-0.338447,0.465835)--cycle;
\path[fill=blue!32!white,draw=blue!32!white,line width=.12pt] (1.240584,0.179769)--(1.163677,0.219074)--(1.019091,0.272690)--(1.086442,0.236929)--cycle;
\path[fill=blue!35!white,draw=blue!35!white,line width=.12pt] (-1.086442,0.236929)--(-1.019091,0.272690)--(-1.163677,0.219074)--(-1.240584,0.179769)--cycle;
\path[fill=blue!15!white,draw=blue!15!white,line width=.12pt] (0.918440,1.801476)--(0.978167,1.860259)--(0.741989,1.902512)--(0.696683,1.841149)--cycle;
\path[fill=blue!15!white,draw=blue!15!white,line width=.12pt] (-0.696683,1.841149)--(-0.741989,1.902512)--(-0.978167,1.860259)--(-0.918440,1.801476)--cycle;
\path[fill=blue!28!white,draw=blue!28!white,line width=.12pt] (2.369413,1.357174)--(2.426560,1.260844)--(2.219681,1.374972)--(2.167407,1.468614)--cycle;
\path[fill=blue!22!white,draw=blue!22!white,line width=.12pt] (-2.167407,1.468614)--(-2.219681,1.374972)--(-2.426560,1.260844)--(-2.369413,1.357174)--cycle;
\path[fill=blue!28!white,draw=blue!28!white,line width=.12pt] (-0.338447,0.465835)--(-0.324025,0.549360)--(-0.482133,0.529559)--(-0.503593,0.445151)--cycle;
\path[fill=blue!26!white,draw=blue!26!white,line width=.12pt] (0.503593,0.445151)--(0.482133,0.529559)--(0.324025,0.549360)--(0.338447,0.465835)--cycle;
\path[fill=green!65!black!29!white,draw=green!65!black!29!white,line width=.12pt] (0.219943,1.796069)--(0.235241,1.887042)--(0.000000,1.892820)--(0.000000,1.801472)--cycle;
\path[fill=green!65!black!29!white,draw=green!65!black!29!white,line width=.12pt] (0.000000,1.801472)--(0.000000,1.892820)--(-0.235241,1.887042)--(-0.219943,1.796069)--cycle;
\path[fill=blue!33!white,draw=blue!33!white,line width=.12pt] (2.647922,1.019805)--(2.660703,0.888912)--(2.473633,1.015029)--(2.461751,1.145316)--cycle;
\path[fill=blue!26!white,draw=blue!26!white,line width=.12pt] (-2.461751,1.145316)--(-2.473633,1.015029)--(-2.660703,0.888912)--(-2.647922,1.019805)--cycle;
\path[fill=blue!37!white,draw=blue!37!white,line width=.12pt] (-2.465877,0.333925)--(-2.365079,0.214089)--(-2.508586,0.094375)--(-2.615501,0.209109)--cycle;
\path[fill=blue!41!white,draw=blue!41!white,line width=.12pt] (2.615501,0.209109)--(2.508586,0.094375)--(2.365079,0.214089)--(2.465877,0.333925)--cycle;
\path[fill=blue!35!white,draw=blue!35!white,line width=.12pt] (-2.548600,0.466550)--(-2.465877,0.333925)--(-2.615501,0.209109)--(-2.703244,0.337547)--cycle;
\path[fill=blue!41!white,draw=blue!41!white,line width=.12pt] (2.703244,0.337547)--(2.615501,0.209109)--(2.465877,0.333925)--(2.548600,0.466550)--cycle;
\path[fill=blue!39!white,draw=blue!39!white,line width=.12pt] (1.995527,-0.026136)--(1.865758,-0.056178)--(1.734580,0.032258)--(1.855225,0.068452)--cycle;
\path[fill=blue!40!white,draw=blue!40!white,line width=.12pt] (-1.855225,0.068452)--(-1.734580,0.032258)--(-1.865758,-0.056178)--(-1.995527,-0.026136)--cycle;
\path[fill=green!65!black!29!white,draw=green!65!black!29!white,line width=.12pt] (0.437769,1.779914)--(0.468217,1.869763)--(0.235241,1.887042)--(0.219943,1.796069)--cycle;
\path[fill=green!65!black!29!white,draw=green!65!black!29!white,line width=.12pt] (-0.219943,1.796069)--(-0.235241,1.887042)--(-0.468217,1.869763)--(-0.437769,1.779914)--cycle;
\path[fill=blue!38!white,draw=blue!38!white,line width=.12pt] (-2.365079,0.214089)--(-2.250079,0.111645)--(-2.386608,-0.002248)--(-2.508586,0.094375)--cycle;
\path[fill=blue!42!white,draw=blue!42!white,line width=.12pt] (2.508586,0.094375)--(2.386608,-0.002248)--(2.250079,0.111645)--(2.365079,0.214089)--cycle;
\path[fill=blue!20!white,draw=blue!20!white,line width=.12pt] (1.716761,1.686022)--(1.804504,1.675456)--(1.580295,1.758598)--(1.503455,1.765122)--cycle;
\path[fill=blue!17!white,draw=blue!17!white,line width=.12pt] (-1.503455,1.765122)--(-1.580295,1.758598)--(-1.804504,1.675456)--(-1.716761,1.686022)--cycle;
\path[fill=blue!29!white,draw=blue!29!white,line width=.12pt] (1.019091,0.272690)--(0.963817,0.336316)--(0.817791,0.380079)--(0.864690,0.318963)--cycle;
\path[fill=blue!32!white,draw=blue!32!white,line width=.12pt] (-0.864690,0.318963)--(-0.817791,0.380079)--(-0.963817,0.336316)--(-1.019091,0.272690)--cycle;
\path[fill=blue!16!white,draw=blue!16!white,line width=.12pt] (-1.204924,1.806635)--(-1.275669,1.833387)--(-1.503455,1.765122)--(-1.420078,1.742156)--cycle;
\path[fill=blue!17!white,draw=blue!17!white,line width=.12pt] (1.420078,1.742156)--(1.503455,1.765122)--(1.275669,1.833387)--(1.204924,1.806635)--cycle;
\path[fill=blue!29!white,draw=blue!29!white,line width=.12pt] (-0.503593,0.445151)--(-0.482133,0.529559)--(-0.635598,0.502103)--(-0.663889,0.416474)--cycle;
\path[fill=blue!26!white,draw=blue!26!white,line width=.12pt] (0.663889,0.416474)--(0.635598,0.502103)--(0.482133,0.529559)--(0.503593,0.445151)--cycle;
\path[fill=blue!33!white,draw=blue!33!white,line width=.12pt] (-2.610069,0.606865)--(-2.548600,0.466550)--(-2.703244,0.337547)--(-2.768442,0.474751)--cycle;
\path[fill=blue!39!white,draw=blue!39!white,line width=.12pt] (2.768442,0.474751)--(2.703244,0.337547)--(2.548600,0.466550)--(2.610069,0.606865)--cycle;
\path[fill=blue!19!white,draw=blue!19!white,line width=.12pt] (-1.804504,1.675456)--(-1.881411,1.636151)--(-2.097056,1.538426)--(-2.011334,1.581726)--cycle;
\path[fill=blue!22!white,draw=blue!22!white,line width=.12pt] (2.011334,1.581726)--(2.097056,1.538426)--(1.881411,1.636151)--(1.804504,1.675456)--cycle;
\path[fill=green!65!black!29!white,draw=green!65!black!29!white,line width=.12pt] (0.651378,1.753160)--(0.696683,1.841149)--(0.468217,1.869763)--(0.437769,1.779914)--cycle;
\path[fill=green!65!black!29!white,draw=green!65!black!29!white,line width=.12pt] (-0.437769,1.779914)--(-0.468217,1.869763)--(-0.696683,1.841149)--(-0.651378,1.753160)--cycle;
\path[fill=blue!39!white,draw=blue!39!white,line width=.12pt] (-2.250079,0.111645)--(-2.125296,0.030531)--(-2.254254,-0.077046)--(-2.386608,-0.002248)--cycle;
\path[fill=blue!41!white,draw=blue!41!white,line width=.12pt] (2.386608,-0.002248)--(2.254254,-0.077046)--(2.125296,0.030531)--(2.250079,0.111645)--cycle;
\path[fill=blue!15!white,draw=blue!15!white,line width=.12pt] (1.131352,1.751126)--(1.204924,1.806635)--(0.978167,1.860259)--(0.918440,1.801476)--cycle;
\path[fill=blue!15!white,draw=blue!15!white,line width=.12pt] (-0.918440,1.801476)--(-0.978167,1.860259)--(-1.204924,1.806635)--(-1.131352,1.751126)--cycle;
\path[fill=blue!24!white,draw=blue!24!white,line width=.12pt] (0.162796,0.561318)--(0.158334,0.668634)--(0.000000,0.672523)--(0.000000,0.565317)--cycle;
\path[fill=blue!24!white,draw=blue!24!white,line width=.12pt] (0.000000,0.565317)--(0.000000,0.672523)--(-0.158334,0.668634)--(-0.162796,0.561318)--cycle;
\path[fill=blue!34!white,draw=blue!34!white,line width=.12pt] (1.618570,0.024081)--(1.511656,0.044233)--(1.382778,0.115330)--(1.480578,0.100206)--cycle;
\path[fill=blue!38!white,draw=blue!38!white,line width=.12pt] (-1.480578,0.100206)--(-1.382778,0.115330)--(-1.511656,0.044233)--(-1.618570,0.024081)--cycle;
\path[fill=blue!35!white,draw=blue!35!white,line width=.12pt] (-1.240584,0.179769)--(-1.163677,0.219074)--(-1.297056,0.158630)--(-1.382778,0.115330)--cycle;
\path[fill=blue!32!white,draw=blue!32!white,line width=.12pt] (1.382778,0.115330)--(1.297056,0.158630)--(1.163677,0.219074)--(1.240584,0.179769)--cycle;
\path[fill=blue!30!white,draw=blue!30!white,line width=.12pt] (-2.647922,0.749480)--(-2.610069,0.606865)--(-2.768442,0.474751)--(-2.808591,0.615449)--cycle;
\path[fill=blue!38!white,draw=blue!38!white,line width=.12pt] (2.808591,0.615449)--(2.768442,0.474751)--(2.610069,0.606865)--(2.647922,0.749480)--cycle;
\path[fill=blue!24!white,draw=blue!24!white,line width=.12pt] (0.324025,0.549360)--(0.315143,0.657004)--(0.158334,0.668634)--(0.162796,0.561318)--cycle;
\path[fill=blue!25!white,draw=blue!25!white,line width=.12pt] (-0.162796,0.561318)--(-0.158334,0.668634)--(-0.315143,0.657004)--(-0.324025,0.549360)--cycle;
\path[fill=blue!26!white,draw=blue!26!white,line width=.12pt] (0.817791,0.380079)--(0.782941,0.467259)--(0.635598,0.502103)--(0.663889,0.416474)--cycle;
\path[fill=blue!29!white,draw=blue!29!white,line width=.12pt] (-0.663889,0.416474)--(-0.635598,0.502103)--(-0.782941,0.467259)--(-0.817791,0.380079)--cycle;
\path[fill=blue!31!white,draw=blue!31!white,line width=.12pt] (2.610069,1.137127)--(2.647922,1.019805)--(2.461751,1.145316)--(2.426560,1.260844)--cycle;
\path[fill=blue!24!white,draw=blue!24!white,line width=.12pt] (-2.426560,1.260844)--(-2.461751,1.145316)--(-2.647922,1.019805)--(-2.610069,1.137127)--cycle;
\path[fill=blue!25!white,draw=blue!25!white,line width=.12pt] (2.292506,1.430603)--(2.369413,1.357174)--(2.167407,1.468614)--(2.097056,1.538426)--cycle;
\path[fill=blue!20!white,draw=blue!20!white,line width=.12pt] (-2.097056,1.538426)--(-2.167407,1.468614)--(-2.369413,1.357174)--(-2.292506,1.430603)--cycle;
\path[fill=blue!37!white,draw=blue!37!white,line width=.12pt] (1.865758,-0.056178)--(1.740975,-0.058441)--(1.618570,0.024081)--(1.734580,0.032258)--cycle;
\path[fill=blue!39!white,draw=blue!39!white,line width=.12pt] (-1.734580,0.032258)--(-1.618570,0.024081)--(-1.740975,-0.058441)--(-1.865758,-0.056178)--cycle;
\path[fill=blue!40!white,draw=blue!40!white,line width=.12pt] (2.254254,-0.077046)--(2.116611,-0.127144)--(1.995527,-0.026136)--(2.125296,0.030531)--cycle;
\path[fill=blue!40!white,draw=blue!40!white,line width=.12pt] (-2.125296,0.030531)--(-1.995527,-0.026136)--(-2.116611,-0.127144)--(-2.254254,-0.077046)--cycle;
\path[fill=green!65!black!29!white,draw=green!65!black!29!white,line width=.12pt] (0.858714,1.716067)--(0.918440,1.801476)--(0.696683,1.841149)--(0.651378,1.753160)--cycle;
\path[fill=green!65!black!29!white,draw=green!65!black!29!white,line width=.12pt] (-0.651378,1.753160)--(-0.696683,1.841149)--(-0.918440,1.801476)--(-0.858714,1.716067)--cycle;
\path[fill=green!65!black!28!white,draw=green!65!black!28!white,line width=.12pt] (0.000000,1.687009)--(0.000000,1.801472)--(-0.219943,1.796069)--(-0.205234,1.681968)--cycle;
\path[fill=green!65!black!28!white,draw=green!65!black!28!white,line width=.12pt] (0.205234,1.681968)--(0.219943,1.796069)--(0.000000,1.801472)--(0.000000,1.687009)--cycle;
\path[fill=blue!33!white,draw=blue!33!white,line width=.12pt] (-1.019091,0.272690)--(-0.963817,0.336316)--(-1.100561,0.285608)--(-1.163677,0.219074)--cycle;
\path[fill=blue!29!white,draw=blue!29!white,line width=.12pt] (1.163677,0.219074)--(1.100561,0.285608)--(0.963817,0.336316)--(1.019091,0.272690)--cycle;
\path[fill=blue!25!white,draw=blue!25!white,line width=.12pt] (-0.324025,0.549360)--(-0.315143,0.657004)--(-0.468918,0.637745)--(-0.482133,0.529559)--cycle;
\path[fill=blue!23!white,draw=blue!23!white,line width=.12pt] (0.482133,0.529559)--(0.468918,0.637745)--(0.315143,0.657004)--(0.324025,0.549360)--cycle;
\path[fill=green!65!black!28!white,draw=green!65!black!28!white,line width=.12pt] (0.408491,1.666893)--(0.437769,1.779914)--(0.219943,1.796069)--(0.205234,1.681968)--cycle;
\path[fill=green!65!black!28!white,draw=green!65!black!28!white,line width=.12pt] (-0.205234,1.681968)--(-0.219943,1.796069)--(-0.437769,1.779914)--(-0.408491,1.666893)--cycle;
\path[fill=blue!28!white,draw=blue!28!white,line width=.12pt] (-2.660703,0.888912)--(-2.647922,0.749480)--(-2.808591,0.615449)--(-2.822148,0.754235)--cycle;
\path[fill=blue!36!white,draw=blue!36!white,line width=.12pt] (2.822148,0.754235)--(2.808591,0.615449)--(2.647922,0.749480)--(2.660703,0.888912)--cycle;
\path[fill=blue!23!white,draw=blue!23!white,line width=.12pt] (0.635598,0.502103)--(0.618176,0.611042)--(0.468918,0.637745)--(0.482133,0.529559)--cycle;
\path[fill=blue!26!white,draw=blue!26!white,line width=.12pt] (-0.482133,0.529559)--(-0.468918,0.637745)--(-0.618176,0.611042)--(-0.635598,0.502103)--cycle;
\path[fill=blue!30!white,draw=blue!30!white,line width=.12pt] (-0.817791,0.380079)--(-0.782941,0.467259)--(-0.922745,0.425362)--(-0.963817,0.336316)--cycle;
\path[fill=blue!26!white,draw=blue!26!white,line width=.12pt] (0.963817,0.336316)--(0.922745,0.425362)--(0.782941,0.467259)--(0.817791,0.380079)--cycle;
\path[fill=blue!16!white,draw=blue!16!white,line width=.12pt] (-1.420078,1.742156)--(-1.503455,1.765122)--(-1.716761,1.686022)--(-1.621555,1.667443)--cycle;
\path[fill=blue!17!white,draw=blue!17!white,line width=.12pt] (1.621555,1.667443)--(1.716761,1.686022)--(1.503455,1.765122)--(1.420078,1.742156)--cycle;
\path[fill=blue!15!white,draw=blue!15!white,line width=.12pt] (-1.131352,1.751126)--(-1.204924,1.806635)--(-1.420078,1.742156)--(-1.333369,1.690584)--cycle;
\path[fill=blue!15!white,draw=blue!15!white,line width=.12pt] (1.333369,1.690584)--(1.420078,1.742156)--(1.204924,1.806635)--(1.131352,1.751126)--cycle;
\path[fill=green!65!black!28!white,draw=green!65!black!28!white,line width=.12pt] (0.607814,1.641929)--(0.651378,1.753160)--(0.437769,1.779914)--(0.408491,1.666893)--cycle;
\path[fill=green!65!black!28!white,draw=green!65!black!28!white,line width=.12pt] (-0.408491,1.666893)--(-0.437769,1.779914)--(-0.651378,1.753160)--(-0.607814,1.641929)--cycle;
\path[fill=blue!40!white,draw=blue!40!white,line width=.12pt] (-1.995527,-0.026136)--(-1.865758,-0.056178)--(-1.978968,-0.150618)--(-2.116611,-0.127144)--cycle;
\path[fill=blue!39!white,draw=blue!39!white,line width=.12pt] (2.116611,-0.127144)--(1.978968,-0.150618)--(1.865758,-0.056178)--(1.995527,-0.026136)--cycle;
\path[fill=blue!22!white,draw=blue!22!white,line width=.12pt] (0.000000,0.672523)--(0.000000,0.800000)--(-0.156827,0.796148)--(-0.158334,0.668634)--cycle;
\path[fill=blue!21!white,draw=blue!21!white,line width=.12pt] (0.158334,0.668634)--(0.156827,0.796148)--(0.000000,0.800000)--(0.000000,0.672523)--cycle;
\path[fill=blue!17!white,draw=blue!17!white,line width=.12pt] (-1.716761,1.686022)--(-1.804504,1.675456)--(-2.011334,1.581726)--(-1.913535,1.596850)--cycle;
\path[fill=blue!20!white,draw=blue!20!white,line width=.12pt] (1.913535,1.596850)--(2.011334,1.581726)--(1.804504,1.675456)--(1.716761,1.686022)--cycle;
\path[fill=green!65!black!29!white,draw=green!65!black!29!white,line width=.12pt] (1.057780,1.668992)--(1.131352,1.751126)--(0.918440,1.801476)--(0.858714,1.716067)--cycle;
\path[fill=green!65!black!29!white,draw=green!65!black!29!white,line width=.12pt] (-0.858714,1.716067)--(-0.918440,1.801476)--(-1.131352,1.751126)--(-1.057780,1.668992)--cycle;
\path[fill=blue!21!white,draw=blue!21!white,line width=.12pt] (0.315143,0.657004)--(0.312145,0.784628)--(0.156827,0.796148)--(0.158334,0.668634)--cycle;
\path[fill=blue!22!white,draw=blue!22!white,line width=.12pt] (-0.158334,0.668634)--(-0.156827,0.796148)--(-0.312145,0.784628)--(-0.315143,0.657004)--cycle;
\path[fill=blue!28!white,draw=blue!28!white,line width=.12pt] (2.548600,1.236371)--(2.610069,1.137127)--(2.426560,1.260844)--(2.369413,1.357174)--cycle;
\path[fill=blue!21!white,draw=blue!21!white,line width=.12pt] (-2.369413,1.357174)--(-2.426560,1.260844)--(-2.610069,1.137127)--(-2.548600,1.236371)--cycle;
\path[fill=blue!32!white,draw=blue!32!white,line width=.12pt] (1.511656,0.044233)--(1.417944,0.091940)--(1.297056,0.158630)--(1.382778,0.115330)--cycle;
\path[fill=blue!36!white,draw=blue!36!white,line width=.12pt] (-1.382778,0.115330)--(-1.297056,0.158630)--(-1.417944,0.091940)--(-1.511656,0.044233)--cycle;
\path[fill=green!65!black!28!white,draw=green!65!black!28!white,line width=.12pt] (0.191677,1.549123)--(0.205234,1.681968)--(0.000000,1.687009)--(0.000000,1.553831)--cycle;
\path[fill=green!65!black!28!white,draw=green!65!black!28!white,line width=.12pt] (0.000000,1.553831)--(0.000000,1.687009)--(-0.205234,1.681968)--(-0.191677,1.549123)--cycle;
\path[fill=blue!23!white,draw=blue!23!white,line width=.12pt] (0.782941,0.467259)--(0.761481,0.577153)--(0.618176,0.611042)--(0.635598,0.502103)--cycle;
\path[fill=blue!27!white,draw=blue!27!white,line width=.12pt] (-0.635598,0.502103)--(-0.618176,0.611042)--(-0.761481,0.577153)--(-0.782941,0.467259)--cycle;
\path[fill=blue!29!white,draw=blue!29!white,line width=.12pt] (1.297056,0.158630)--(1.226706,0.228443)--(1.100561,0.285608)--(1.163677,0.219074)--cycle;
\path[fill=blue!33!white,draw=blue!33!white,line width=.12pt] (-1.163677,0.219074)--(-1.100561,0.285608)--(-1.226706,0.228443)--(-1.297056,0.158630)--cycle;
\path[fill=blue!35!white,draw=blue!35!white,line width=.12pt] (1.740975,-0.058441)--(1.625975,-0.032838)--(1.511656,0.044233)--(1.618570,0.024081)--cycle;
\path[fill=blue!38!white,draw=blue!38!white,line width=.12pt] (-1.618570,0.024081)--(-1.511656,0.044233)--(-1.625975,-0.032838)--(-1.740975,-0.058441)--cycle;
\path[fill=blue!23!white,draw=blue!23!white,line width=.12pt] (-0.315143,0.657004)--(-0.312145,0.784628)--(-0.464455,0.765552)--(-0.468918,0.637745)--cycle;
\path[fill=blue!21!white,draw=blue!21!white,line width=.12pt] (0.468918,0.637745)--(0.464455,0.765552)--(0.312145,0.784628)--(0.315143,0.657004)--cycle;
\path[fill=green!65!black!28!white,draw=green!65!black!28!white,line width=.12pt] (0.801283,1.607317)--(0.858714,1.716067)--(0.651378,1.753160)--(0.607814,1.641929)--cycle;
\path[fill=green!65!black!29!white,draw=green!65!black!29!white,line width=.12pt] (-0.607814,1.641929)--(-0.651378,1.753160)--(-0.858714,1.716067)--(-0.801283,1.607317)--cycle;
\path[fill=blue!36!white,draw=blue!36!white,line width=.12pt] (-2.615501,0.209109)--(-2.508586,0.094375)--(-2.627935,-0.031796)--(-2.739936,0.077561)--cycle;
\path[fill=blue!41!white,draw=blue!41!white,line width=.12pt] (2.739936,0.077561)--(2.627935,-0.031796)--(2.508586,0.094375)--(2.615501,0.209109)--cycle;
\path[fill=green!65!black!27!white,draw=green!65!black!27!white,line width=.12pt] (0.381508,1.535043)--(0.408491,1.666893)--(0.205234,1.681968)--(0.191677,1.549123)--cycle;
\path[fill=green!65!black!28!white,draw=green!65!black!28!white,line width=.12pt] (-0.191677,1.549123)--(-0.205234,1.681968)--(-0.408491,1.666893)--(-0.381508,1.535043)--cycle;
\path[fill=blue!33!white,draw=blue!33!white,line width=.12pt] (2.808591,0.885774)--(2.822148,0.754235)--(2.660703,0.888912)--(2.647922,1.019805)--cycle;
\path[fill=blue!25!white,draw=blue!25!white,line width=.12pt] (-2.647922,1.019805)--(-2.660703,0.888912)--(-2.822148,0.754235)--(-2.808591,0.885774)--cycle;
\path[fill=blue!18!white,draw=blue!18!white,line width=.12pt] (-2.011334,1.581726)--(-2.097056,1.538426)--(-2.292506,1.430603)--(-2.198794,1.478311)--cycle;
\path[fill=blue!22!white,draw=blue!22!white,line width=.12pt] (2.198794,1.478311)--(2.292506,1.430603)--(2.097056,1.538426)--(2.011334,1.581726)--cycle;
\path[fill=blue!41!white,draw=blue!41!white,line width=.12pt] (2.627935,-0.031796)--(2.500154,-0.122284)--(2.386608,-0.002248)--(2.508586,0.094375)--cycle;
\path[fill=blue!38!white,draw=blue!38!white,line width=.12pt] (-2.508586,0.094375)--(-2.386608,-0.002248)--(-2.500154,-0.122284)--(-2.627935,-0.031796)--cycle;
\path[fill=blue!34!white,draw=blue!34!white,line width=.12pt] (-2.703244,0.337547)--(-2.615501,0.209109)--(-2.739936,0.077561)--(-2.831853,0.201586)--cycle;
\path[fill=blue!40!white,draw=blue!40!white,line width=.12pt] (2.831853,0.201586)--(2.739936,0.077561)--(2.615501,0.209109)--(2.703244,0.337547)--cycle;
\path[fill=blue!31!white,draw=blue!31!white,line width=.12pt] (-0.963817,0.336316)--(-0.922745,0.425362)--(-1.053662,0.376814)--(-1.100561,0.285608)--cycle;
\path[fill=blue!26!white,draw=blue!26!white,line width=.12pt] (1.100561,0.285608)--(1.053662,0.376814)--(0.922745,0.425362)--(0.963817,0.336316)--cycle;
\path[fill=blue!39!white,draw=blue!39!white,line width=.12pt] (-2.386608,-0.002248)--(-2.254254,-0.077046)--(-2.361503,-0.190425)--(-2.500154,-0.122284)--cycle;
\path[fill=blue!41!white,draw=blue!41!white,line width=.12pt] (2.500154,-0.122284)--(2.361503,-0.190425)--(2.254254,-0.077046)--(2.386608,-0.002248)--cycle;
\path[fill=blue!19!white,draw=blue!19!white,line width=.12pt] (0.156827,0.796148)--(0.158334,0.938959)--(0.000000,0.942848)--(0.000000,0.800000)--cycle;
\path[fill=blue!19!white,draw=blue!19!white,line width=.12pt] (0.000000,0.800000)--(0.000000,0.942848)--(-0.158334,0.938959)--(-0.156827,0.796148)--cycle;
\path[fill=blue!20!white,draw=blue!20!white,line width=.12pt] (0.618176,0.611042)--(0.612293,0.739104)--(0.464455,0.765552)--(0.468918,0.637745)--cycle;
\path[fill=blue!23!white,draw=blue!23!white,line width=.12pt] (-0.468918,0.637745)--(-0.464455,0.765552)--(-0.612293,0.739104)--(-0.618176,0.611042)--cycle;
\path[fill=green!65!black!28!white,draw=green!65!black!28!white,line width=.12pt] (-0.381508,1.535043)--(-0.408491,1.666893)--(-0.607814,1.641929)--(-0.567664,1.511728)--cycle;
\path[fill=green!65!black!27!white,draw=green!65!black!27!white,line width=.12pt] (0.567664,1.511728)--(0.607814,1.641929)--(0.408491,1.666893)--(0.381508,1.535043)--cycle;
\path[fill=blue!37!white,draw=blue!37!white,line width=.12pt] (1.978968,-0.150618)--(1.846614,-0.146565)--(1.740975,-0.058441)--(1.865758,-0.056178)--cycle;
\path[fill=blue!39!white,draw=blue!39!white,line width=.12pt] (-1.865758,-0.056178)--(-1.740975,-0.058441)--(-1.846614,-0.146565)--(-1.978968,-0.150618)--cycle;
\path[fill=blue!18!white,draw=blue!18!white,line width=.12pt] (0.312145,0.784628)--(0.315143,0.927329)--(0.158334,0.938959)--(0.156827,0.796148)--cycle;
\path[fill=blue!20!white,draw=blue!20!white,line width=.12pt] (-0.156827,0.796148)--(-0.158334,0.938959)--(-0.315143,0.927329)--(-0.312145,0.784628)--cycle;
\path[fill=blue!32!white,draw=blue!32!white,line width=.12pt] (-2.768442,0.474751)--(-2.703244,0.337547)--(-2.831853,0.201586)--(-2.900154,0.335511)--cycle;
\path[fill=blue!39!white,draw=blue!39!white,line width=.12pt] (2.900154,0.335511)--(2.831853,0.201586)--(2.703244,0.337547)--(2.768442,0.474751)--cycle;
\path[fill=blue!27!white,draw=blue!27!white,line width=.12pt] (-0.782941,0.467259)--(-0.761481,0.577153)--(-0.897452,0.536404)--(-0.922745,0.425362)--cycle;
\path[fill=blue!23!white,draw=blue!23!white,line width=.12pt] (0.922745,0.425362)--(0.897452,0.536404)--(0.761481,0.577153)--(0.782941,0.467259)--cycle;
\path[fill=green!65!black!28!white,draw=green!65!black!28!white,line width=.12pt] (0.179794,1.402639)--(0.191677,1.549123)--(0.000000,1.553831)--(0.000000,1.407055)--cycle;
\path[fill=green!65!black!28!white,draw=green!65!black!28!white,line width=.12pt] (0.000000,1.407055)--(0.000000,1.553831)--(-0.191677,1.549123)--(-0.179794,1.402639)--cycle;
\path[fill=green!65!black!29!white,draw=green!65!black!29!white,line width=.12pt] (-1.057780,1.668992)--(-1.131352,1.751126)--(-1.333369,1.690584)--(-1.246659,1.612387)--cycle;
\path[fill=green!65!black!29!white,draw=green!65!black!29!white,line width=.12pt] (1.246659,1.612387)--(1.333369,1.690584)--(1.131352,1.751126)--(1.057780,1.668992)--cycle;
\path[fill=blue!15!white,draw=blue!15!white,line width=.12pt] (1.522544,1.620433)--(1.621555,1.667443)--(1.420078,1.742156)--(1.333369,1.690584)--cycle;
\path[fill=blue!15!white,draw=blue!15!white,line width=.12pt] (-1.333369,1.690584)--(-1.420078,1.742156)--(-1.621555,1.667443)--(-1.522544,1.620433)--cycle;
\path[fill=green!65!black!29!white,draw=green!65!black!29!white,line width=.12pt] (-0.801283,1.607317)--(-0.858714,1.716067)--(-1.057780,1.668992)--(-0.987036,1.563389)--cycle;
\path[fill=green!65!black!28!white,draw=green!65!black!28!white,line width=.12pt] (0.987036,1.563389)--(1.057780,1.668992)--(0.858714,1.716067)--(0.801283,1.607317)--cycle;
\path[fill=green!65!black!28!white,draw=green!65!black!28!white,line width=.12pt] (0.357857,1.389432)--(0.381508,1.535043)--(0.191677,1.549123)--(0.179794,1.402639)--cycle;
\path[fill=green!65!black!28!white,draw=green!65!black!28!white,line width=.12pt] (-0.179794,1.402639)--(-0.191677,1.549123)--(-0.381508,1.535043)--(-0.357857,1.389432)--cycle;
\path[fill=blue!40!white,draw=blue!40!white,line width=.12pt] (2.361503,-0.190425)--(2.217311,-0.233600)--(2.116611,-0.127144)--(2.254254,-0.077046)--cycle;
\path[fill=blue!39!white,draw=blue!39!white,line width=.12pt] (-2.254254,-0.077046)--(-2.116611,-0.127144)--(-2.217311,-0.233600)--(-2.361503,-0.190425)--cycle;
\path[fill=blue!20!white,draw=blue!20!white,line width=.12pt] (-0.312145,0.784628)--(-0.315143,0.927329)--(-0.468918,0.908070)--(-0.464455,0.765552)--cycle;
\path[fill=blue!18!white,draw=blue!18!white,line width=.12pt] (0.464455,0.765552)--(0.468918,0.908070)--(0.315143,0.927329)--(0.312145,0.784628)--cycle;
\path[fill=blue!16!white,draw=blue!16!white,line width=.12pt] (-1.621555,1.667443)--(-1.716761,1.686022)--(-1.913535,1.596850)--(-1.807416,1.583216)--cycle;
\path[fill=blue!17!white,draw=blue!17!white,line width=.12pt] (1.807416,1.583216)--(1.913535,1.596850)--(1.716761,1.686022)--(1.621555,1.667443)--cycle;
\path[fill=blue!25!white,draw=blue!25!white,line width=.12pt] (2.465877,1.313721)--(2.548600,1.236371)--(2.369413,1.357174)--(2.292506,1.430603)--cycle;
\path[fill=blue!19!white,draw=blue!19!white,line width=.12pt] (-2.292506,1.430603)--(-2.369413,1.357174)--(-2.548600,1.236371)--(-2.465877,1.313721)--cycle;
\path[fill=blue!17!white,draw=blue!17!white,line width=.12pt] (0.158334,0.938959)--(0.162796,1.091580)--(0.000000,1.095579)--(0.000000,0.942848)--cycle;
\path[fill=blue!17!white,draw=blue!17!white,line width=.12pt] (0.000000,0.942848)--(0.000000,1.095579)--(-0.162796,1.091580)--(-0.158334,0.938959)--cycle;
\path[fill=green!65!black!27!white,draw=green!65!black!27!white,line width=.12pt] (0.748354,1.479402)--(0.801283,1.607317)--(0.607814,1.641929)--(0.567664,1.511728)--cycle;
\path[fill=green!65!black!29!white,draw=green!65!black!29!white,line width=.12pt] (-0.567664,1.511728)--(-0.607814,1.641929)--(-0.801283,1.607317)--(-0.748354,1.479402)--cycle;
\path[fill=blue!24!white,draw=blue!24!white,line width=.12pt] (-0.618176,0.611042)--(-0.612293,0.739104)--(-0.754235,0.705537)--(-0.761481,0.577153)--cycle;
\path[fill=blue!20!white,draw=blue!20!white,line width=.12pt] (0.761481,0.577153)--(0.754235,0.705537)--(0.612293,0.739104)--(0.618176,0.611042)--cycle;
\path[fill=blue!30!white,draw=blue!30!white,line width=.12pt] (2.768442,1.005012)--(2.808591,0.885774)--(2.647922,1.019805)--(2.610069,1.137127)--cycle;
\path[fill=blue!23!white,draw=blue!23!white,line width=.12pt] (-2.610069,1.137127)--(-2.647922,1.019805)--(-2.808591,0.885774)--(-2.768442,1.005012)--cycle;
\path[fill=blue!14!white,draw=blue!14!white,line width=.12pt] (0.000000,1.252323)--(0.000000,1.407055)--(-0.179794,1.402639)--(-0.170043,1.248146)--cycle;
\path[fill=blue!14!white,draw=blue!14!white,line width=.12pt] (0.170043,1.248146)--(0.179794,1.402639)--(0.000000,1.407055)--(0.000000,1.252323)--cycle;
\path[fill=blue!29!white,draw=blue!29!white,line width=.12pt] (1.417944,0.091940)--(1.341037,0.165370)--(1.226706,0.228443)--(1.297056,0.158630)--cycle;
\path[fill=blue!34!white,draw=blue!34!white,line width=.12pt] (-1.297056,0.158630)--(-1.226706,0.228443)--(-1.341037,0.165370)--(-1.417944,0.091940)--cycle;
\path[fill=blue!16!white,draw=blue!16!white,line width=.12pt] (0.315143,0.927329)--(0.324025,1.079622)--(0.162796,1.091580)--(0.158334,0.938959)--cycle;
\path[fill=blue!18!white,draw=blue!18!white,line width=.12pt] (-0.158334,0.938959)--(-0.162796,1.091580)--(-0.324025,1.079622)--(-0.315143,0.927329)--cycle;
\path[fill=blue!30!white,draw=blue!30!white,line width=.12pt] (-2.808591,0.615449)--(-2.768442,0.474751)--(-2.900154,0.335511)--(-2.942213,0.474190)--cycle;
\path[fill=blue!37!white,draw=blue!37!white,line width=.12pt] (2.942213,0.474190)--(2.900154,0.335511)--(2.768442,0.474751)--(2.808591,0.615449)--cycle;
\path[fill=blue!32!white,draw=blue!32!white,line width=.12pt] (1.625975,-0.032838)--(1.525177,0.019647)--(1.417944,0.091940)--(1.511656,0.044233)--cycle;
\path[fill=blue!36!white,draw=blue!36!white,line width=.12pt] (-1.511656,0.044233)--(-1.417944,0.091940)--(-1.525177,0.019647)--(-1.625975,-0.032838)--cycle;
\path[fill=green!65!black!29!white,draw=green!65!black!29!white,line width=.12pt] (-0.357857,1.389432)--(-0.381508,1.535043)--(-0.567664,1.511728)--(-0.532473,1.367563)--cycle;
\path[fill=green!65!black!27!white,draw=green!65!black!27!white,line width=.12pt] (0.532473,1.367563)--(0.567664,1.511728)--(0.381508,1.535043)--(0.357857,1.389432)--cycle;
\path[fill=blue!15!white,draw=blue!15!white,line width=.12pt] (0.162796,1.091580)--(0.170043,1.248146)--(0.000000,1.252323)--(0.000000,1.095579)--cycle;
\path[fill=blue!15!white,draw=blue!15!white,line width=.12pt] (0.000000,1.095579)--(0.000000,1.252323)--(-0.170043,1.248146)--(-0.162796,1.091580)--cycle;
\path[fill=blue!26!white,draw=blue!26!white,line width=.12pt] (1.226706,0.228443)--(1.174431,0.322085)--(1.053662,0.376814)--(1.100561,0.285608)--cycle;
\path[fill=blue!31!white,draw=blue!31!white,line width=.12pt] (-1.100561,0.285608)--(-1.053662,0.376814)--(-1.174431,0.322085)--(-1.226706,0.228443)--cycle;
\path[fill=blue!13!white,draw=blue!13!white,line width=.12pt] (0.338447,1.235655)--(0.357857,1.389432)--(0.179794,1.402639)--(0.170043,1.248146)--cycle;
\path[fill=blue!14!white,draw=blue!14!white,line width=.12pt] (-0.170043,1.248146)--(-0.179794,1.402639)--(-0.357857,1.389432)--(-0.338447,1.235655)--cycle;
\path[fill=blue!18!white,draw=blue!18!white,line width=.12pt] (0.612293,0.739104)--(0.618176,0.881367)--(0.468918,0.908070)--(0.464455,0.765552)--cycle;
\path[fill=blue!21!white,draw=blue!21!white,line width=.12pt] (-0.464455,0.765552)--(-0.468918,0.908070)--(-0.618176,0.881367)--(-0.612293,0.739104)--cycle;
\path[fill=blue!16!white,draw=blue!16!white,line width=.12pt] (-0.162796,1.091580)--(-0.170043,1.248146)--(-0.338447,1.235655)--(-0.324025,1.079622)--cycle;
\path[fill=blue!15!white,draw=blue!15!white,line width=.12pt] (0.324025,1.079622)--(0.338447,1.235655)--(0.170043,1.248146)--(0.162796,1.091580)--cycle;
\path[fill=blue!16!white,draw=blue!16!white,line width=.12pt] (-1.913535,1.596850)--(-2.011334,1.581726)--(-2.198794,1.478311)--(-2.091880,1.498463)--cycle;
\path[fill=blue!20!white,draw=blue!20!white,line width=.12pt] (2.091880,1.498463)--(2.198794,1.478311)--(2.011334,1.581726)--(1.913535,1.596850)--cycle;
\path[fill=blue!16!white,draw=blue!16!white,line width=.12pt] (0.468918,0.908070)--(0.482133,1.059820)--(0.324025,1.079622)--(0.315143,0.927329)--cycle;
\path[fill=blue!18!white,draw=blue!18!white,line width=.12pt] (-0.315143,0.927329)--(-0.324025,1.079622)--(-0.482133,1.059820)--(-0.468918,0.908070)--cycle;
\path[fill=blue!40!white,draw=blue!40!white,line width=.12pt] (-2.116611,-0.127144)--(-1.978968,-0.150618)--(-2.073119,-0.250151)--(-2.217311,-0.233600)--cycle;
\path[fill=blue!39!white,draw=blue!39!white,line width=.12pt] (2.217311,-0.233600)--(2.073119,-0.250151)--(1.978968,-0.150618)--(2.116611,-0.127144)--cycle;
\path[fill=blue!23!white,draw=blue!23!white,line width=.12pt] (1.053662,0.376814)--(1.024781,0.489187)--(0.897452,0.536404)--(0.922745,0.425362)--cycle;
\path[fill=blue!28!white,draw=blue!28!white,line width=.12pt] (-0.922745,0.425362)--(-0.897452,0.536404)--(-1.024781,0.489187)--(-1.053662,0.376814)--cycle;
\path[fill=blue!38!white,draw=blue!38!white,line width=.12pt] (-1.740975,-0.058441)--(-1.625975,-0.032838)--(-1.724636,-0.115141)--(-1.846614,-0.146565)--cycle;
\path[fill=blue!35!white,draw=blue!35!white,line width=.12pt] (1.846614,-0.146565)--(1.724636,-0.115141)--(1.625975,-0.032838)--(1.740975,-0.058441)--cycle;
\path[fill=blue!15!white,draw=blue!15!white,line width=.12pt] (-0.338447,1.235655)--(-0.357857,1.389432)--(-0.532473,1.367563)--(-0.503593,1.214972)--cycle;
\path[fill=blue!13!white,draw=blue!13!white,line width=.12pt] (0.503593,1.214972)--(0.532473,1.367563)--(0.357857,1.389432)--(0.338447,1.235655)--cycle;
\path[fill=blue!14!white,draw=blue!14!white,line width=.12pt] (0.482133,1.059820)--(0.503593,1.214972)--(0.338447,1.235655)--(0.324025,1.079622)--cycle;
\path[fill=blue!16!white,draw=blue!16!white,line width=.12pt] (-0.324025,1.079622)--(-0.338447,1.235655)--(-0.503593,1.214972)--(-0.482133,1.059820)--cycle;
\path[fill=green!65!black!27!white,draw=green!65!black!27!white,line width=.12pt] (0.701962,1.337241)--(0.748354,1.479402)--(0.567664,1.511728)--(0.532473,1.367563)--cycle;
\path[fill=green!65!black!29!white,draw=green!65!black!29!white,line width=.12pt] (-0.532473,1.367563)--(-0.567664,1.511728)--(-0.748354,1.479402)--(-0.701962,1.337241)--cycle;
\path[fill=green!65!black!27!white,draw=green!65!black!27!white,line width=.12pt] (0.921837,1.438377)--(0.987036,1.563389)--(0.801283,1.607317)--(0.748354,1.479402)--cycle;
\path[fill=green!65!black!29!white,draw=green!65!black!29!white,line width=.12pt] (-0.748354,1.479402)--(-0.801283,1.607317)--(-0.987036,1.563389)--(-0.921837,1.438377)--cycle;
\path[fill=blue!20!white,draw=blue!20!white,line width=.12pt] (0.897452,0.536404)--(0.888912,0.665176)--(0.754235,0.705537)--(0.761481,0.577153)--cycle;
\path[fill=blue!25!white,draw=blue!25!white,line width=.12pt] (-0.761481,0.577153)--(-0.754235,0.705537)--(-0.888912,0.665176)--(-0.897452,0.536404)--cycle;
\path[fill=green!65!black!28!white,draw=green!65!black!28!white,line width=.12pt] (1.163283,1.510570)--(1.246659,1.612387)--(1.057780,1.668992)--(0.987036,1.563389)--cycle;
\path[fill=green!65!black!29!white,draw=green!65!black!29!white,line width=.12pt] (-0.987036,1.563389)--(-1.057780,1.668992)--(-1.246659,1.612387)--(-1.163283,1.510570)--cycle;
\path[fill=blue!27!white,draw=blue!27!white,line width=.12pt] (-2.822148,0.754235)--(-2.808591,0.615449)--(-2.942213,0.474190)--(-2.956415,0.612293)--cycle;
\path[fill=blue!35!white,draw=blue!35!white,line width=.12pt] (2.956415,0.612293)--(2.942213,0.474190)--(2.808591,0.615449)--(2.822148,0.754235)--cycle;
\path[fill=blue!19!white,draw=blue!19!white,line width=.12pt] (-0.468918,0.908070)--(-0.482133,1.059820)--(-0.635598,1.032365)--(-0.618176,0.881367)--cycle;
\path[fill=blue!16!white,draw=blue!16!white,line width=.12pt] (0.618176,0.881367)--(0.635598,1.032365)--(0.482133,1.059820)--(0.468918,0.908070)--cycle;
\path[fill=blue!22!white,draw=blue!22!white,line width=.12pt] (-0.612293,0.739104)--(-0.618176,0.881367)--(-0.761481,0.847478)--(-0.754235,0.705537)--cycle;
\path[fill=blue!18!white,draw=blue!18!white,line width=.12pt] (0.754235,0.705537)--(0.761481,0.847478)--(0.618176,0.881367)--(0.612293,0.739104)--cycle;
\path[fill=green!65!black!29!white,draw=green!65!black!29!white,line width=.12pt] (1.423533,1.546798)--(1.522544,1.620433)--(1.333369,1.690584)--(1.246659,1.612387)--cycle;
\path[fill=green!65!black!29!white,draw=green!65!black!29!white,line width=.12pt] (-1.246659,1.612387)--(-1.333369,1.690584)--(-1.522544,1.620433)--(-1.423533,1.546798)--cycle;
\path[fill=blue!15!white,draw=blue!15!white,line width=.12pt] (-0.503593,1.214972)--(-0.532473,1.367563)--(-0.701962,1.337241)--(-0.663889,1.186295)--cycle;
\path[fill=blue!13!white,draw=blue!13!white,line width=.12pt] (0.663889,1.186295)--(0.701962,1.337241)--(0.532473,1.367563)--(0.503593,1.214972)--cycle;
\path[fill=blue!14!white,draw=blue!14!white,line width=.12pt] (0.635598,1.032365)--(0.663889,1.186295)--(0.503593,1.214972)--(0.482133,1.059820)--cycle;
\path[fill=blue!17!white,draw=blue!17!white,line width=.12pt] (-0.482133,1.059820)--(-0.503593,1.214972)--(-0.663889,1.186295)--(-0.635598,1.032365)--cycle;
\path[fill=blue!15!white,draw=blue!15!white,line width=.12pt] (-1.522544,1.620433)--(-1.621555,1.667443)--(-1.807416,1.583216)--(-1.697056,1.541348)--cycle;
\path[fill=blue!15!white,draw=blue!15!white,line width=.12pt] (1.697056,1.541348)--(1.807416,1.583216)--(1.621555,1.667443)--(1.522544,1.620433)--cycle;
\path[fill=blue!21!white,draw=blue!21!white,line width=.12pt] (-2.548600,1.236371)--(-2.610069,1.137127)--(-2.768442,1.005012)--(-2.703244,1.107367)--cycle;
\path[fill=blue!28!white,draw=blue!28!white,line width=.12pt] (2.703244,1.107367)--(2.768442,1.005012)--(2.610069,1.137127)--(2.548600,1.236371)--cycle;
\path[fill=blue!18!white,draw=blue!18!white,line width=.12pt] (-2.198794,1.478311)--(-2.292506,1.430603)--(-2.465877,1.313721)--(-2.365079,1.366207)--cycle;
\path[fill=blue!22!white,draw=blue!22!white,line width=.12pt] (2.365079,1.366207)--(2.465877,1.313721)--(2.292506,1.430603)--(2.198794,1.478311)--cycle;
\path[fill=blue!39!white,draw=blue!39!white,line width=.12pt] (-1.978968,-0.150618)--(-1.846614,-0.146565)--(-1.934468,-0.239441)--(-2.073119,-0.250151)--cycle;
\path[fill=blue!37!white,draw=blue!37!white,line width=.12pt] (2.073119,-0.250151)--(1.934468,-0.239441)--(1.846614,-0.146565)--(1.978968,-0.150618)--cycle;
\path[fill=green!65!black!27!white,draw=green!65!black!27!white,line width=.12pt] (0.864690,1.298758)--(0.921837,1.438377)--(0.748354,1.479402)--(0.701962,1.337241)--cycle;
\path[fill=green!65!black!30!white,draw=green!65!black!30!white,line width=.12pt] (-0.701962,1.337241)--(-0.748354,1.479402)--(-0.921837,1.438377)--(-0.864690,1.298758)--cycle;
\path[fill=blue!32!white,draw=blue!32!white,line width=.12pt] (-1.226706,0.228443)--(-1.174431,0.322085)--(-1.283890,0.261700)--(-1.341037,0.165370)--cycle;
\path[fill=blue!26!white,draw=blue!26!white,line width=.12pt] (1.341037,0.165370)--(1.283890,0.261700)--(1.174431,0.322085)--(1.226706,0.228443)--cycle;
\path[fill=blue!29!white,draw=blue!29!white,line width=.12pt] (1.525177,0.019647)--(1.442454,0.096998)--(1.341037,0.165370)--(1.417944,0.091940)--cycle;
\path[fill=blue!35!white,draw=blue!35!white,line width=.12pt] (-1.417944,0.091940)--(-1.341037,0.165370)--(-1.442454,0.096998)--(-1.525177,0.019647)--cycle;
\path[fill=blue!23!white,draw=blue!23!white,line width=.12pt] (1.174431,0.322085)--(1.142240,0.435958)--(1.024781,0.489187)--(1.053662,0.376814)--cycle;
\path[fill=blue!29!white,draw=blue!29!white,line width=.12pt] (-1.053662,0.376814)--(-1.024781,0.489187)--(-1.142240,0.435958)--(-1.174431,0.322085)--cycle;
\path[fill=blue!16!white,draw=blue!16!white,line width=.12pt] (0.761481,0.847478)--(0.782941,0.997521)--(0.635598,1.032365)--(0.618176,0.881367)--cycle;
\path[fill=blue!19!white,draw=blue!19!white,line width=.12pt] (-0.618176,0.881367)--(-0.635598,1.032365)--(-0.782941,0.997521)--(-0.761481,0.847478)--cycle;
\path[fill=blue!37!white,draw=blue!37!white,line width=.12pt] (-2.627935,-0.031796)--(-2.500154,-0.122284)--(-2.589621,-0.247306)--(-2.721975,-0.163208)--cycle;
\path[fill=blue!41!white,draw=blue!41!white,line width=.12pt] (2.721975,-0.163208)--(2.589621,-0.247306)--(2.500154,-0.122284)--(2.627935,-0.031796)--cycle;
\path[fill=blue!20!white,draw=blue!20!white,line width=.12pt] (1.024781,0.489187)--(1.015029,0.618408)--(0.888912,0.665176)--(0.897452,0.536404)--cycle;
\path[fill=blue!26!white,draw=blue!26!white,line width=.12pt] (-0.897452,0.536404)--(-0.888912,0.665176)--(-1.015029,0.618408)--(-1.024781,0.489187)--cycle;
\path[fill=blue!41!white,draw=blue!41!white,line width=.12pt] (2.837984,-0.059451)--(2.721975,-0.163208)--(2.627935,-0.031796)--(2.739936,0.077561)--cycle;
\path[fill=blue!36!white,draw=blue!36!white,line width=.12pt] (-2.739936,0.077561)--(-2.627935,-0.031796)--(-2.721975,-0.163208)--(-2.837984,-0.059451)--cycle;
\path[fill=blue!22!white,draw=blue!22!white,line width=.12pt] (-0.754235,0.705537)--(-0.761481,0.847478)--(-0.897452,0.806729)--(-0.888912,0.665176)--cycle;
\path[fill=blue!18!white,draw=blue!18!white,line width=.12pt] (0.888912,0.665176)--(0.897452,0.806729)--(0.761481,0.847478)--(0.754235,0.705537)--cycle;
\path[fill=green!65!black!27!white,draw=green!65!black!27!white,line width=.12pt] (1.086442,1.389047)--(1.163283,1.510570)--(0.987036,1.563389)--(0.921837,1.438377)--cycle;
\path[fill=green!65!black!29!white,draw=green!65!black!29!white,line width=.12pt] (-0.921837,1.438377)--(-0.987036,1.563389)--(-1.163283,1.510570)--(-1.086442,1.389047)--cycle;
\path[fill=blue!32!white,draw=blue!32!white,line width=.12pt] (1.724636,-0.115141)--(1.617721,-0.057553)--(1.525177,0.019647)--(1.625975,-0.032838)--cycle;
\path[fill=blue!37!white,draw=blue!37!white,line width=.12pt] (-1.625975,-0.032838)--(-1.525177,0.019647)--(-1.617721,-0.057553)--(-1.724636,-0.115141)--cycle;
\path[fill=blue!38!white,draw=blue!38!white,line width=.12pt] (-2.500154,-0.122284)--(-2.361503,-0.190425)--(-2.446009,-0.308514)--(-2.589621,-0.247306)--cycle;
\path[fill=blue!41!white,draw=blue!41!white,line width=.12pt] (2.589621,-0.247306)--(2.446009,-0.308514)--(2.361503,-0.190425)--(2.500154,-0.122284)--cycle;
\path[fill=blue!17!white,draw=blue!17!white,line width=.12pt] (1.975871,1.490286)--(2.091880,1.498463)--(1.913535,1.596850)--(1.807416,1.583216)--cycle;
\path[fill=blue!15!white,draw=blue!15!white,line width=.12pt] (-1.807416,1.583216)--(-1.913535,1.596850)--(-2.091880,1.498463)--(-1.975871,1.490286)--cycle;
\path[fill=blue!13!white,draw=blue!13!white,line width=.12pt] (0.817791,1.149900)--(0.864690,1.298758)--(0.701962,1.337241)--(0.663889,1.186295)--cycle;
\path[fill=blue!16!white,draw=blue!16!white,line width=.12pt] (-0.663889,1.186295)--(-0.701962,1.337241)--(-0.864690,1.298758)--(-0.817791,1.149900)--cycle;
\path[fill=blue!25!white,draw=blue!25!white,line width=.12pt] (-2.808591,0.885774)--(-2.822148,0.754235)--(-2.956415,0.612293)--(-2.942213,0.744515)--cycle;
\path[fill=blue!33!white,draw=blue!33!white,line width=.12pt] (2.942213,0.744515)--(2.956415,0.612293)--(2.822148,0.754235)--(2.808591,0.885774)--cycle;
\path[fill=blue!17!white,draw=blue!17!white,line width=.12pt] (-0.635598,1.032365)--(-0.663889,1.186295)--(-0.817791,1.149900)--(-0.782941,0.997521)--cycle;
\path[fill=blue!14!white,draw=blue!14!white,line width=.12pt] (0.782941,0.997521)--(0.817791,1.149900)--(0.663889,1.186295)--(0.635598,1.032365)--cycle;
\path[fill=blue!33!white,draw=blue!33!white,line width=.12pt] (-2.831853,0.201586)--(-2.739936,0.077561)--(-2.837984,-0.059451)--(-2.933190,0.059976)--cycle;
\path[fill=blue!40!white,draw=blue!40!white,line width=.12pt] (2.933190,0.059976)--(2.837984,-0.059451)--(2.739936,0.077561)--(2.831853,0.201586)--cycle;
\path[fill=green!65!black!29!white,draw=green!65!black!29!white,line width=.12pt] (-1.163283,1.510570)--(-1.246659,1.612387)--(-1.423533,1.546798)--(-1.328326,1.449368)--cycle;
\path[fill=green!65!black!28!white,draw=green!65!black!28!white,line width=.12pt] (1.328326,1.449368)--(1.423533,1.546798)--(1.246659,1.612387)--(1.163283,1.510570)--cycle;
\path[fill=blue!39!white,draw=blue!39!white,line width=.12pt] (-2.361503,-0.190425)--(-2.217311,-0.233600)--(-2.296657,-0.344479)--(-2.446009,-0.308514)--cycle;
\path[fill=blue!40!white,draw=blue!40!white,line width=.12pt] (2.446009,-0.308514)--(2.296657,-0.344479)--(2.217311,-0.233600)--(2.361503,-0.190425)--cycle;
\path[fill=blue!39!white,draw=blue!39!white,line width=.12pt] (-1.846614,-0.146565)--(-1.724636,-0.115141)--(-1.806687,-0.201882)--(-1.934468,-0.239441)--cycle;
\path[fill=blue!35!white,draw=blue!35!white,line width=.12pt] (1.934468,-0.239441)--(1.806687,-0.201882)--(1.724636,-0.115141)--(1.846614,-0.146565)--cycle;
\path[fill=blue!20!white,draw=blue!20!white,line width=.12pt] (-0.761481,0.847478)--(-0.782941,0.997521)--(-0.922745,0.955623)--(-0.897452,0.806729)--cycle;
\path[fill=blue!16!white,draw=blue!16!white,line width=.12pt] (0.897452,0.806729)--(0.922745,0.955623)--(0.782941,0.997521)--(0.761481,0.847478)--cycle;
\path[fill=green!65!black!27!white,draw=green!65!black!27!white,line width=.12pt] (1.019091,1.252486)--(1.086442,1.389047)--(0.921837,1.438377)--(0.864690,1.298758)--cycle;
\path[fill=green!65!black!30!white,draw=green!65!black!30!white,line width=.12pt] (-0.864690,1.298758)--(-0.921837,1.438377)--(-1.086442,1.389047)--(-1.019091,1.252486)--cycle;
\path[fill=green!65!black!29!white,draw=green!65!black!29!white,line width=.12pt] (1.586697,1.472856)--(1.697056,1.541348)--(1.522544,1.620433)--(1.423533,1.546798)--cycle;
\path[fill=green!65!black!29!white,draw=green!65!black!29!white,line width=.12pt] (-1.423533,1.546798)--(-1.522544,1.620433)--(-1.697056,1.541348)--(-1.586697,1.472856)--cycle;
\path[fill=blue!31!white,draw=blue!31!white,line width=.12pt] (-2.900154,0.335511)--(-2.831853,0.201586)--(-2.933190,0.059976)--(-3.003935,0.190486)--cycle;
\path[fill=blue!39!white,draw=blue!39!white,line width=.12pt] (3.003935,0.190486)--(2.933190,0.059976)--(2.831853,0.201586)--(2.900154,0.335511)--cycle;
\path[fill=blue!25!white,draw=blue!25!white,line width=.12pt] (2.615501,1.188905)--(2.703244,1.107367)--(2.548600,1.236371)--(2.465877,1.313721)--cycle;
\path[fill=blue!19!white,draw=blue!19!white,line width=.12pt] (-2.465877,1.313721)--(-2.548600,1.236371)--(-2.703244,1.107367)--(-2.615501,1.188905)--cycle;
\path[fill=blue!19!white,draw=blue!19!white,line width=.12pt] (2.250079,1.391810)--(2.365079,1.366207)--(2.198794,1.478311)--(2.091880,1.498463)--cycle;
\path[fill=blue!16!white,draw=blue!16!white,line width=.12pt] (-2.091880,1.498463)--(-2.198794,1.478311)--(-2.365079,1.366207)--(-2.250079,1.391810)--cycle;
\path[fill=blue!23!white,draw=blue!23!white,line width=.12pt] (-0.888912,0.665176)--(-0.897452,0.806729)--(-1.024781,0.759512)--(-1.015029,0.618408)--cycle;
\path[fill=blue!18!white,draw=blue!18!white,line width=.12pt] (1.015029,0.618408)--(1.024781,0.759512)--(0.897452,0.806729)--(0.888912,0.665176)--cycle;
\path[fill=blue!40!white,draw=blue!40!white,line width=.12pt] (-2.217311,-0.233600)--(-2.073119,-0.250151)--(-2.147305,-0.353819)--(-2.296657,-0.344479)--cycle;
\path[fill=blue!39!white,draw=blue!39!white,line width=.12pt] (2.296657,-0.344479)--(2.147305,-0.353819)--(2.073119,-0.250151)--(2.217311,-0.233600)--cycle;
\path[fill=blue!14!white,draw=blue!14!white,line width=.12pt] (0.922745,0.955623)--(0.963817,1.106137)--(0.817791,1.149900)--(0.782941,0.997521)--cycle;
\path[fill=blue!18!white,draw=blue!18!white,line width=.12pt] (-0.782941,0.997521)--(-0.817791,1.149900)--(-0.963817,1.106137)--(-0.922745,0.955623)--cycle;
\path[fill=blue!13!white,draw=blue!13!white,line width=.12pt] (0.963817,1.106137)--(1.019091,1.252486)--(0.864690,1.298758)--(0.817791,1.149900)--cycle;
\path[fill=blue!16!white,draw=blue!16!white,line width=.12pt] (-0.817791,1.149900)--(-0.864690,1.298758)--(-1.019091,1.252486)--(-0.963817,1.106137)--cycle;
\path[fill=blue!26!white,draw=blue!26!white,line width=.12pt] (-1.024781,0.489187)--(-1.015029,0.618408)--(-1.131371,0.565685)--(-1.142240,0.435958)--cycle;
\path[fill=blue!20!white,draw=blue!20!white,line width=.12pt] (1.142240,0.435958)--(1.131371,0.565685)--(1.015029,0.618408)--(1.024781,0.489187)--cycle;
\path[fill=blue!30!white,draw=blue!30!white,line width=.12pt] (-1.174431,0.322085)--(-1.142240,0.435958)--(-1.248699,0.377228)--(-1.283890,0.261700)--cycle;
\path[fill=blue!23!white,draw=blue!23!white,line width=.12pt] (1.283890,0.261700)--(1.248699,0.377228)--(1.142240,0.435958)--(1.174431,0.322085)--cycle;
\path[fill=blue!26!white,draw=blue!26!white,line width=.12pt] (1.442454,0.096998)--(1.380985,0.196241)--(1.283890,0.261700)--(1.341037,0.165370)--cycle;
\path[fill=blue!33!white,draw=blue!33!white,line width=.12pt] (-1.341037,0.165370)--(-1.283890,0.261700)--(-1.380985,0.196241)--(-1.442454,0.096998)--cycle;
\path[fill=green!65!black!30!white,draw=green!65!black!30!white,line width=.12pt] (-1.086442,1.389047)--(-1.163283,1.510570)--(-1.328326,1.449368)--(-1.240584,1.331887)--cycle;
\path[fill=green!65!black!27!white,draw=green!65!black!27!white,line width=.12pt] (1.240584,1.331887)--(1.328326,1.449368)--(1.163283,1.510570)--(1.086442,1.389047)--cycle;
\path[fill=blue!22!white,draw=blue!22!white,line width=.12pt] (-2.768442,1.005012)--(-2.808591,0.885774)--(-2.942213,0.744515)--(-2.900154,0.865772)--cycle;
\path[fill=blue!30!white,draw=blue!30!white,line width=.12pt] (2.900154,0.865772)--(2.942213,0.744515)--(2.808591,0.885774)--(2.768442,1.005012)--cycle;
\path[fill=blue!29!white,draw=blue!29!white,line width=.12pt] (1.617721,-0.057553)--(1.529978,0.023985)--(1.442454,0.096998)--(1.525177,0.019647)--cycle;
\path[fill=blue!35!white,draw=blue!35!white,line width=.12pt] (-1.525177,0.019647)--(-1.442454,0.096998)--(-1.529978,0.023985)--(-1.617721,-0.057553)--cycle;
\path[fill=blue!15!white,draw=blue!15!white,line width=.12pt] (1.855225,1.454092)--(1.975871,1.490286)--(1.807416,1.583216)--(1.697056,1.541348)--cycle;
\path[fill=blue!15!white,draw=blue!15!white,line width=.12pt] (-1.697056,1.541348)--(-1.807416,1.583216)--(-1.975871,1.490286)--(-1.855225,1.454092)--cycle;
\path[fill=blue!29!white,draw=blue!29!white,line width=.12pt] (-2.942213,0.474190)--(-2.900154,0.335511)--(-3.003935,0.190486)--(-3.047499,0.327062)--cycle;
\path[fill=blue!37!white,draw=blue!37!white,line width=.12pt] (3.047499,0.327062)--(3.003935,0.190486)--(2.900154,0.335511)--(2.942213,0.474190)--cycle;
\path[fill=blue!21!white,draw=blue!21!white,line width=.12pt] (-0.897452,0.806729)--(-0.922745,0.955623)--(-1.053662,0.907076)--(-1.024781,0.759512)--cycle;
\path[fill=blue!16!white,draw=blue!16!white,line width=.12pt] (1.024781,0.759512)--(1.053662,0.907076)--(0.922745,0.955623)--(0.897452,0.806729)--cycle;
\path[fill=blue!40!white,draw=blue!40!white,line width=.12pt] (-2.073119,-0.250151)--(-1.934468,-0.239441)--(-2.003693,-0.336176)--(-2.147305,-0.353819)--cycle;
\path[fill=blue!37!white,draw=blue!37!white,line width=.12pt] (2.147305,-0.353819)--(2.003693,-0.336176)--(1.934468,-0.239441)--(2.073119,-0.250151)--cycle;
\path[fill=green!65!black!29!white,draw=green!65!black!29!white,line width=.12pt] (-1.328326,1.449368)--(-1.423533,1.546798)--(-1.586697,1.472856)--(-1.480578,1.380371)--cycle;
\path[fill=green!65!black!28!white,draw=green!65!black!28!white,line width=.12pt] (1.480578,1.380371)--(1.586697,1.472856)--(1.423533,1.546798)--(1.328326,1.449368)--cycle;
\path[fill=green!65!black!27!white,draw=green!65!black!27!white,line width=.12pt] (1.163677,1.198870)--(1.240584,1.331887)--(1.086442,1.389047)--(1.019091,1.252486)--cycle;
\path[fill=green!65!black!31!white,draw=green!65!black!31!white,line width=.12pt] (-1.019091,1.252486)--(-1.086442,1.389047)--(-1.240584,1.331887)--(-1.163677,1.198870)--cycle;
\path[fill=blue!37!white,draw=blue!37!white,line width=.12pt] (-1.724636,-0.115141)--(-1.617721,-0.057553)--(-1.694686,-0.138917)--(-1.806687,-0.201882)--cycle;
\path[fill=blue!32!white,draw=blue!32!white,line width=.12pt] (1.806687,-0.201882)--(1.694686,-0.138917)--(1.617721,-0.057553)--(1.724636,-0.115141)--cycle;
\path[fill=blue!19!white,draw=blue!19!white,line width=.12pt] (-0.922745,0.955623)--(-0.963817,1.106137)--(-1.100561,1.055429)--(-1.053662,0.907076)--cycle;
\path[fill=blue!14!white,draw=blue!14!white,line width=.12pt] (1.053662,0.907076)--(1.100561,1.055429)--(0.963817,1.106137)--(0.922745,0.955623)--cycle;
\path[fill=blue!24!white,draw=blue!24!white,line width=.12pt] (-1.015029,0.618408)--(-1.024781,0.759512)--(-1.142240,0.706283)--(-1.131371,0.565685)--cycle;
\path[fill=blue!18!white,draw=blue!18!white,line width=.12pt] (1.131371,0.565685)--(1.142240,0.706283)--(1.024781,0.759512)--(1.015029,0.618408)--cycle;
\path[fill=blue!17!white,draw=blue!17!white,line width=.12pt] (-0.963817,1.106137)--(-1.019091,1.252486)--(-1.163677,1.198870)--(-1.100561,1.055429)--cycle;
\path[fill=blue!13!white,draw=blue!13!white,line width=.12pt] (1.100561,1.055429)--(1.163677,1.198870)--(1.019091,1.252486)--(0.963817,1.106137)--cycle;
\path[fill=blue!22!white,draw=blue!22!white,line width=.12pt] (2.508586,1.246493)--(2.615501,1.188905)--(2.465877,1.313721)--(2.365079,1.366207)--cycle;
\path[fill=blue!17!white,draw=blue!17!white,line width=.12pt] (-2.365079,1.366207)--(-2.465877,1.313721)--(-2.615501,1.188905)--(-2.508586,1.246493)--cycle;
\path[fill=blue!15!white,draw=blue!15!white,line width=.12pt] (-1.975871,1.490286)--(-2.091880,1.498463)--(-2.250079,1.391810)--(-2.125296,1.389547)--cycle;
\path[fill=blue!17!white,draw=blue!17!white,line width=.12pt] (2.125296,1.389547)--(2.250079,1.391810)--(2.091880,1.498463)--(1.975871,1.490286)--cycle;
\path[fill=blue!27!white,draw=blue!27!white,line width=.12pt] (-1.142240,0.435958)--(-1.131371,0.565685)--(-1.236817,0.507515)--(-1.248699,0.377228)--cycle;
\path[fill=blue!20!white,draw=blue!20!white,line width=.12pt] (1.248699,0.377228)--(1.236817,0.507515)--(1.131371,0.565685)--(1.142240,0.435958)--cycle;
\path[fill=blue!26!white,draw=blue!26!white,line width=.12pt] (-2.956415,0.612293)--(-2.942213,0.474190)--(-3.047499,0.327062)--(-3.062209,0.464455)--cycle;
\path[fill=blue!35!white,draw=blue!35!white,line width=.12pt] (3.062209,0.464455)--(3.047499,0.327062)--(2.942213,0.474190)--(2.956415,0.612293)--cycle;
\path[fill=blue!23!white,draw=blue!23!white,line width=.12pt] (1.380985,0.196241)--(1.343133,0.313564)--(1.248699,0.377228)--(1.283890,0.261700)--cycle;
\path[fill=blue!30!white,draw=blue!30!white,line width=.12pt] (-1.283890,0.261700)--(-1.248699,0.377228)--(-1.343133,0.313564)--(-1.380985,0.196241)--cycle;
\path[fill=blue!20!white,draw=blue!20!white,line width=.12pt] (-2.703244,1.107367)--(-2.768442,1.005012)--(-2.900154,0.865772)--(-2.831853,0.971406)--cycle;
\path[fill=blue!27!white,draw=blue!27!white,line width=.12pt] (2.831853,0.971406)--(2.900154,0.865772)--(2.768442,1.005012)--(2.703244,1.107367)--cycle;
\path[fill=green!65!black!29!white,draw=green!65!black!29!white,line width=.12pt] (1.734580,1.391274)--(1.855225,1.454092)--(1.697056,1.541348)--(1.586697,1.472856)--cycle;
\path[fill=green!65!black!29!white,draw=green!65!black!29!white,line width=.12pt] (-1.586697,1.472856)--(-1.697056,1.541348)--(-1.855225,1.454092)--(-1.734580,1.391274)--cycle;
\path[fill=green!65!black!30!white,draw=green!65!black!30!white,line width=.12pt] (-1.240584,1.331887)--(-1.328326,1.449368)--(-1.480578,1.380371)--(-1.382778,1.267448)--cycle;
\path[fill=green!65!black!27!white,draw=green!65!black!27!white,line width=.12pt] (1.382778,1.267448)--(1.480578,1.380371)--(1.328326,1.449368)--(1.240584,1.331887)--cycle;
\path[fill=blue!33!white,draw=blue!33!white,line width=.12pt] (-1.442454,0.096998)--(-1.380985,0.196241)--(-1.464780,0.126340)--(-1.529978,0.023985)--cycle;
\path[fill=blue!26!white,draw=blue!26!white,line width=.12pt] (1.529978,0.023985)--(1.464780,0.126340)--(1.380985,0.196241)--(1.442454,0.096998)--cycle;
\path[fill=blue!39!white,draw=blue!39!white,line width=.12pt] (-1.934468,-0.239441)--(-1.806687,-0.201882)--(-1.871339,-0.292227)--(-2.003693,-0.336176)--cycle;
\path[fill=blue!35!white,draw=blue!35!white,line width=.12pt] (2.003693,-0.336176)--(1.871339,-0.292227)--(1.806687,-0.201882)--(1.934468,-0.239441)--cycle;
\path[fill=blue!37!white,draw=blue!37!white,line width=.12pt] (-2.721975,-0.163208)--(-2.589621,-0.247306)--(-2.654149,-0.376111)--(-2.789801,-0.298596)--cycle;
\path[fill=blue!41!white,draw=blue!41!white,line width=.12pt] (2.789801,-0.298596)--(2.654149,-0.376111)--(2.589621,-0.247306)--(2.721975,-0.163208)--cycle;
\path[fill=blue!16!white,draw=blue!16!white,line width=.12pt] (1.142240,0.706283)--(1.174431,0.852346)--(1.053662,0.907076)--(1.024781,0.759512)--cycle;
\path[fill=blue!21!white,draw=blue!21!white,line width=.12pt] (-1.024781,0.759512)--(-1.053662,0.907076)--(-1.174431,0.852346)--(-1.142240,0.706283)--cycle;
\path[fill=blue!38!white,draw=blue!38!white,line width=.12pt] (-2.589621,-0.247306)--(-2.446009,-0.308514)--(-2.506958,-0.430175)--(-2.654149,-0.376111)--cycle;
\path[fill=blue!41!white,draw=blue!41!white,line width=.12pt] (2.654149,-0.376111)--(2.506958,-0.430175)--(2.446009,-0.308514)--(2.589621,-0.247306)--cycle;
\path[fill=blue!35!white,draw=blue!35!white,line width=.12pt] (-2.837984,-0.059451)--(-2.721975,-0.163208)--(-2.789801,-0.298596)--(-2.908701,-0.200610)--cycle;
\path[fill=blue!41!white,draw=blue!41!white,line width=.12pt] (2.908701,-0.200610)--(2.789801,-0.298596)--(2.721975,-0.163208)--(2.837984,-0.059451)--cycle;
\path[fill=blue!39!white,draw=blue!39!white,line width=.12pt] (-2.446009,-0.308514)--(-2.296657,-0.344479)--(-2.353885,-0.458712)--(-2.506958,-0.430175)--cycle;
\path[fill=blue!40!white,draw=blue!40!white,line width=.12pt] (2.506958,-0.430175)--(2.353885,-0.458712)--(2.296657,-0.344479)--(2.446009,-0.308514)--cycle;
\path[fill=blue!36!white,draw=blue!36!white,line width=.12pt] (-1.617721,-0.057553)--(-1.529978,0.023985)--(-1.602769,-0.052966)--(-1.694686,-0.138917)--cycle;
\path[fill=blue!30!white,draw=blue!30!white,line width=.12pt] (1.694686,-0.138917)--(1.602769,-0.052966)--(1.529978,0.023985)--(1.617721,-0.057553)--cycle;
\path[fill=green!65!black!31!white,draw=green!65!black!31!white,line width=.12pt] (-1.163677,1.198870)--(-1.240584,1.331887)--(-1.382778,1.267448)--(-1.297056,1.138426)--cycle;
\path[fill=green!65!black!27!white,draw=green!65!black!27!white,line width=.12pt] (1.297056,1.138426)--(1.382778,1.267448)--(1.240584,1.331887)--(1.163677,1.198870)--cycle;
\path[fill=blue!14!white,draw=blue!14!white,line width=.12pt] (1.174431,0.852346)--(1.226706,0.998263)--(1.100561,1.055429)--(1.053662,0.907076)--cycle;
\path[fill=blue!19!white,draw=blue!19!white,line width=.12pt] (-1.053662,0.907076)--(-1.100561,1.055429)--(-1.226706,0.998263)--(-1.174431,0.852346)--cycle;
\path[fill=blue!18!white,draw=blue!18!white,line width=.12pt] (1.236817,0.507515)--(1.248699,0.647553)--(1.142240,0.706283)--(1.131371,0.565685)--cycle;
\path[fill=blue!25!white,draw=blue!25!white,line width=.12pt] (-1.131371,0.565685)--(-1.142240,0.706283)--(-1.248699,0.647553)--(-1.236817,0.507515)--cycle;
\path[fill=blue!24!white,draw=blue!24!white,line width=.12pt] (-2.942213,0.744515)--(-2.956415,0.612293)--(-3.062209,0.464455)--(-3.047499,0.597387)--cycle;
\path[fill=blue!32!white,draw=blue!32!white,line width=.12pt] (3.047499,0.597387)--(3.062209,0.464455)--(2.956415,0.612293)--(2.942213,0.744515)--cycle;
\path[fill=blue!17!white,draw=blue!17!white,line width=.12pt] (-1.100561,1.055429)--(-1.163677,1.198870)--(-1.297056,1.138426)--(-1.226706,0.998263)--cycle;
\path[fill=blue!13!white,draw=blue!13!white,line width=.12pt] (1.226706,0.998263)--(1.297056,1.138426)--(1.163677,1.198870)--(1.100561,1.055429)--cycle;
\path[fill=blue!33!white,draw=blue!33!white,line width=.12pt] (-2.933190,0.059976)--(-2.837984,-0.059451)--(-2.908701,-0.200610)--(-3.006279,-0.085917)--cycle;
\path[fill=blue!40!white,draw=blue!40!white,line width=.12pt] (3.006279,-0.085917)--(2.908701,-0.200610)--(2.837984,-0.059451)--(2.933190,0.059976)--cycle;
\path[fill=blue!40!white,draw=blue!40!white,line width=.12pt] (-2.296657,-0.344479)--(-2.147305,-0.353819)--(-2.200811,-0.460624)--(-2.353885,-0.458712)--cycle;
\path[fill=blue!39!white,draw=blue!39!white,line width=.12pt] (2.353885,-0.458712)--(2.200811,-0.460624)--(2.147305,-0.353819)--(2.296657,-0.344479)--cycle;
\path[fill=blue!19!white,draw=blue!19!white,line width=.12pt] (2.386608,1.277917)--(2.508586,1.246493)--(2.365079,1.366207)--(2.250079,1.391810)--cycle;
\path[fill=blue!16!white,draw=blue!16!white,line width=.12pt] (-2.250079,1.391810)--(-2.365079,1.366207)--(-2.508586,1.246493)--(-2.386608,1.277917)--cycle;
\path[fill=blue!14!white,draw=blue!14!white,line width=.12pt] (-1.855225,1.454092)--(-1.975871,1.490286)--(-2.125296,1.389547)--(-1.995527,1.359505)--cycle;
\path[fill=blue!15!white,draw=blue!15!white,line width=.12pt] (1.995527,1.359505)--(2.125296,1.389547)--(1.975871,1.490286)--(1.855225,1.454092)--cycle;
\path[fill=green!65!black!30!white,draw=green!65!black!30!white,line width=.12pt] (-1.480578,1.380371)--(-1.586697,1.472856)--(-1.734580,1.391274)--(-1.618570,1.304246)--cycle;
\path[fill=green!65!black!28!white,draw=green!65!black!28!white,line width=.12pt] (1.618570,1.304246)--(1.734580,1.391274)--(1.586697,1.472856)--(1.480578,1.380371)--cycle;
\path[fill=blue!28!white,draw=blue!28!white,line width=.12pt] (-1.248699,0.377228)--(-1.236817,0.507515)--(-1.330351,0.444456)--(-1.343133,0.313564)--cycle;
\path[fill=blue!21!white,draw=blue!21!white,line width=.12pt] (1.343133,0.313564)--(1.330351,0.444456)--(1.236817,0.507515)--(1.248699,0.377228)--cycle;
\path[fill=blue!18!white,draw=blue!18!white,line width=.12pt] (-2.615501,1.188905)--(-2.703244,1.107367)--(-2.831853,0.971406)--(-2.739936,1.057357)--cycle;
\path[fill=blue!24!white,draw=blue!24!white,line width=.12pt] (2.739936,1.057357)--(2.831853,0.971406)--(2.703244,1.107367)--(2.615501,1.188905)--cycle;
\path[fill=blue!38!white,draw=blue!38!white,line width=.12pt] (-1.806687,-0.201882)--(-1.694686,-0.138917)--(-1.755330,-0.223661)--(-1.871339,-0.292227)--cycle;
\path[fill=blue!33!white,draw=blue!33!white,line width=.12pt] (1.871339,-0.292227)--(1.755330,-0.223661)--(1.694686,-0.138917)--(1.806687,-0.201882)--cycle;
\path[fill=blue!24!white,draw=blue!24!white,line width=.12pt] (1.464780,0.126340)--(1.424631,0.245578)--(1.343133,0.313564)--(1.380985,0.196241)--cycle;
\path[fill=blue!31!white,draw=blue!31!white,line width=.12pt] (-1.380985,0.196241)--(-1.343133,0.313564)--(-1.424631,0.245578)--(-1.464780,0.126340)--cycle;
\path[fill=blue!31!white,draw=blue!31!white,line width=.12pt] (-3.003935,0.190486)--(-2.933190,0.059976)--(-3.006279,-0.085917)--(-3.078787,0.041074)--cycle;
\path[fill=blue!38!white,draw=blue!38!white,line width=.12pt] (3.078787,0.041074)--(3.006279,-0.085917)--(2.933190,0.059976)--(3.003935,0.190486)--cycle;
\path[fill=blue!40!white,draw=blue!40!white,line width=.12pt] (-2.147305,-0.353819)--(-2.003693,-0.336176)--(-2.053620,-0.435837)--(-2.200811,-0.460624)--cycle;
\path[fill=blue!37!white,draw=blue!37!white,line width=.12pt] (2.200811,-0.460624)--(2.053620,-0.435837)--(2.003693,-0.336176)--(2.147305,-0.353819)--cycle;
\path[fill=blue!16!white,draw=blue!16!white,line width=.12pt] (1.248699,0.647553)--(1.283890,0.791962)--(1.174431,0.852346)--(1.142240,0.706283)--cycle;
\path[fill=blue!22!white,draw=blue!22!white,line width=.12pt] (-1.142240,0.706283)--(-1.174431,0.852346)--(-1.283890,0.791962)--(-1.248699,0.647553)--cycle;
\path[fill=green!65!black!30!white,draw=green!65!black!30!white,line width=.12pt] (-1.382778,1.267448)--(-1.480578,1.380371)--(-1.618570,1.304246)--(-1.511656,1.196351)--cycle;
\path[fill=green!65!black!27!white,draw=green!65!black!27!white,line width=.12pt] (1.511656,1.196351)--(1.618570,1.304246)--(1.480578,1.380371)--(1.382778,1.267448)--cycle;
\path[fill=blue!34!white,draw=blue!34!white,line width=.12pt] (-1.529978,0.023985)--(-1.464780,0.126340)--(-1.534468,0.052668)--(-1.602769,-0.052966)--cycle;
\path[fill=blue!27!white,draw=blue!27!white,line width=.12pt] (1.602769,-0.052966)--(1.534468,0.052668)--(1.464780,0.126340)--(1.529978,0.023985)--cycle;
\path[fill=blue!21!white,draw=blue!21!white,line width=.12pt] (-2.900154,0.865772)--(-2.942213,0.744515)--(-3.047499,0.597387)--(-3.003935,0.720748)--cycle;
\path[fill=blue!30!white,draw=blue!30!white,line width=.12pt] (3.003935,0.720748)--(3.047499,0.597387)--(2.942213,0.744515)--(2.900154,0.865772)--cycle;
\path[fill=green!65!black!29!white,draw=green!65!black!29!white,line width=.12pt] (-1.734580,1.391274)--(-1.855225,1.454092)--(-1.995527,1.359505)--(-1.865758,1.302838)--cycle;
\path[fill=green!65!black!29!white,draw=green!65!black!29!white,line width=.12pt] (1.865758,1.302838)--(1.995527,1.359505)--(1.855225,1.454092)--(1.734580,1.391274)--cycle;
\path[fill=blue!14!white,draw=blue!14!white,line width=.12pt] (1.283890,0.791962)--(1.341037,0.935191)--(1.226706,0.998263)--(1.174431,0.852346)--cycle;
\path[fill=blue!20!white,draw=blue!20!white,line width=.12pt] (-1.174431,0.852346)--(-1.226706,0.998263)--(-1.341037,0.935191)--(-1.283890,0.791962)--cycle;
\path[fill=green!65!black!31!white,draw=green!65!black!31!white,line width=.12pt] (-1.297056,1.138426)--(-1.382778,1.267448)--(-1.511656,1.196351)--(-1.417944,1.071736)--cycle;
\path[fill=green!65!black!27!white,draw=green!65!black!27!white,line width=.12pt] (1.417944,1.071736)--(1.511656,1.196351)--(1.382778,1.267448)--(1.297056,1.138426)--cycle;
\path[fill=blue!25!white,draw=blue!25!white,line width=.12pt] (-1.236817,0.507515)--(-1.248699,0.647553)--(-1.343133,0.583889)--(-1.330351,0.444456)--cycle;
\path[fill=blue!18!white,draw=blue!18!white,line width=.12pt] (1.330351,0.444456)--(1.343133,0.583889)--(1.248699,0.647553)--(1.236817,0.507515)--cycle;
\path[fill=blue!15!white,draw=blue!15!white,line width=.12pt] (-2.125296,1.389547)--(-2.250079,1.391810)--(-2.386608,1.277917)--(-2.254254,1.281970)--cycle;
\path[fill=blue!17!white,draw=blue!17!white,line width=.12pt] (2.254254,1.281970)--(2.386608,1.277917)--(2.250079,1.391810)--(2.125296,1.389547)--cycle;
\path[fill=blue!13!white,draw=blue!13!white,line width=.12pt] (1.341037,0.935191)--(1.417944,1.071736)--(1.297056,1.138426)--(1.226706,0.998263)--cycle;
\path[fill=blue!18!white,draw=blue!18!white,line width=.12pt] (-1.226706,0.998263)--(-1.297056,1.138426)--(-1.417944,1.071736)--(-1.341037,0.935191)--cycle;
\path[fill=blue!28!white,draw=blue!28!white,line width=.12pt] (-3.047499,0.327062)--(-3.003935,0.190486)--(-3.078787,0.041074)--(-3.123436,0.175483)--cycle;
\path[fill=blue!36!white,draw=blue!36!white,line width=.12pt] (3.123436,0.175483)--(3.078787,0.041074)--(3.003935,0.190486)--(3.047499,0.327062)--cycle;
\path[fill=blue!39!white,draw=blue!39!white,line width=.12pt] (-2.003693,-0.336176)--(-1.871339,-0.292227)--(-1.917969,-0.385305)--(-2.053620,-0.435837)--cycle;
\path[fill=blue!35!white,draw=blue!35!white,line width=.12pt] (2.053620,-0.435837)--(1.917969,-0.385305)--(1.871339,-0.292227)--(2.003693,-0.336176)--cycle;
\path[fill=blue!37!white,draw=blue!37!white,line width=.12pt] (-1.694686,-0.138917)--(-1.602769,-0.052966)--(-1.660123,-0.133114)--(-1.755330,-0.223661)--cycle;
\path[fill=blue!30!white,draw=blue!30!white,line width=.12pt] (1.755330,-0.223661)--(1.660123,-0.133114)--(1.602769,-0.052966)--(1.694686,-0.138917)--cycle;
\path[fill=blue!17!white,draw=blue!17!white,line width=.12pt] (-2.508586,1.246493)--(-2.615501,1.188905)--(-2.739936,1.057357)--(-2.627935,1.120322)--cycle;
\path[fill=blue!22!white,draw=blue!22!white,line width=.12pt] (2.627935,1.120322)--(2.739936,1.057357)--(2.615501,1.188905)--(2.508586,1.246493)--cycle;
\path[fill=blue!29!white,draw=blue!29!white,line width=.12pt] (-1.343133,0.313564)--(-1.330351,0.444456)--(-1.411074,0.377117)--(-1.424631,0.245578)--cycle;
\path[fill=blue!21!white,draw=blue!21!white,line width=.12pt] (1.424631,0.245578)--(1.411074,0.377117)--(1.330351,0.444456)--(1.343133,0.313564)--cycle;
\path[fill=green!65!black!28!white,draw=green!65!black!28!white,line width=.12pt] (1.740975,1.221724)--(1.865758,1.302838)--(1.734580,1.391274)--(1.618570,1.304246)--cycle;
\path[fill=green!65!black!30!white,draw=green!65!black!30!white,line width=.12pt] (-1.618570,1.304246)--(-1.734580,1.391274)--(-1.865758,1.302838)--(-1.740975,1.221724)--cycle;
\path[fill=blue!16!white,draw=blue!16!white,line width=.12pt] (1.343133,0.583889)--(1.380985,0.726503)--(1.283890,0.791962)--(1.248699,0.647553)--cycle;
\path[fill=blue!23!white,draw=blue!23!white,line width=.12pt] (-1.248699,0.647553)--(-1.283890,0.791962)--(-1.380985,0.726503)--(-1.343133,0.583889)--cycle;
\path[fill=blue!32!white,draw=blue!32!white,line width=.12pt] (-1.464780,0.126340)--(-1.424631,0.245578)--(-1.492409,0.173925)--(-1.534468,0.052668)--cycle;
\path[fill=blue!24!white,draw=blue!24!white,line width=.12pt] (1.534468,0.052668)--(1.492409,0.173925)--(1.424631,0.245578)--(1.464780,0.126340)--cycle;
\path[fill=blue!19!white,draw=blue!19!white,line width=.12pt] (-2.831853,0.971406)--(-2.900154,0.865772)--(-3.003935,0.720748)--(-2.933190,0.829797)--cycle;
\path[fill=blue!27!white,draw=blue!27!white,line width=.12pt] (2.933190,0.829797)--(3.003935,0.720748)--(2.900154,0.865772)--(2.831853,0.971406)--cycle;
\path[fill=blue!38!white,draw=blue!38!white,line width=.12pt] (-2.654149,-0.376111)--(-2.506958,-0.430175)--(-2.543764,-0.554239)--(-2.693116,-0.507458)--cycle;
\path[fill=blue!41!white,draw=blue!41!white,line width=.12pt] (2.693116,-0.507458)--(2.543764,-0.554239)--(2.506958,-0.430175)--(2.654149,-0.376111)--cycle;
\path[fill=blue!39!white,draw=blue!39!white,line width=.12pt] (-2.506958,-0.430175)--(-2.353885,-0.458712)--(-2.388443,-0.575200)--(-2.543764,-0.554239)--cycle;
\path[fill=blue!40!white,draw=blue!40!white,line width=.12pt] (2.543764,-0.554239)--(2.388443,-0.575200)--(2.353885,-0.458712)--(2.506958,-0.430175)--cycle;
\path[fill=blue!26!white,draw=blue!26!white,line width=.12pt] (-3.062209,0.464455)--(-3.047499,0.327062)--(-3.123436,0.175483)--(-3.138513,0.312145)--cycle;
\path[fill=blue!34!white,draw=blue!34!white,line width=.12pt] (3.138513,0.312145)--(3.123436,0.175483)--(3.047499,0.327062)--(3.062209,0.464455)--cycle;
\path[fill=blue!38!white,draw=blue!38!white,line width=.12pt] (-1.871339,-0.292227)--(-1.755330,-0.223661)--(-1.799069,-0.310969)--(-1.917969,-0.385305)--cycle;
\path[fill=blue!33!white,draw=blue!33!white,line width=.12pt] (1.917969,-0.385305)--(1.799069,-0.310969)--(1.755330,-0.223661)--(1.871339,-0.292227)--cycle;
\path[fill=blue!14!white,draw=blue!14!white,line width=.12pt] (-1.995527,1.359505)--(-2.125296,1.389547)--(-2.254254,1.281970)--(-2.116611,1.258496)--cycle;
\path[fill=blue!15!white,draw=blue!15!white,line width=.12pt] (2.116611,1.258496)--(2.254254,1.281970)--(2.125296,1.389547)--(1.995527,1.359505)--cycle;
\path[fill=blue!36!white,draw=blue!36!white,line width=.12pt] (-2.789801,-0.298596)--(-2.654149,-0.376111)--(-2.693116,-0.507458)--(-2.830759,-0.436656)--cycle;
\path[fill=blue!41!white,draw=blue!41!white,line width=.12pt] (2.830759,-0.436656)--(2.693116,-0.507458)--(2.654149,-0.376111)--(2.789801,-0.298596)--cycle;
\path[fill=green!65!black!27!white,draw=green!65!black!27!white,line width=.12pt] (1.625975,1.119280)--(1.740975,1.221724)--(1.618570,1.304246)--(1.511656,1.196351)--cycle;
\path[fill=green!65!black!31!white,draw=green!65!black!31!white,line width=.12pt] (-1.511656,1.196351)--(-1.618570,1.304246)--(-1.740975,1.221724)--(-1.625975,1.119280)--cycle;
\path[fill=blue!40!white,draw=blue!40!white,line width=.12pt] (-2.353885,-0.458712)--(-2.200811,-0.460624)--(-2.233123,-0.569536)--(-2.388443,-0.575200)--cycle;
\path[fill=blue!39!white,draw=blue!39!white,line width=.12pt] (2.388443,-0.575200)--(2.233123,-0.569536)--(2.200811,-0.460624)--(2.353885,-0.458712)--cycle;
\path[fill=blue!14!white,draw=blue!14!white,line width=.12pt] (1.380985,0.726503)--(1.442454,0.866819)--(1.341037,0.935191)--(1.283890,0.791962)--cycle;
\path[fill=blue!21!white,draw=blue!21!white,line width=.12pt] (-1.283890,0.791962)--(-1.341037,0.935191)--(-1.442454,0.866819)--(-1.380985,0.726503)--cycle;
\path[fill=blue!35!white,draw=blue!35!white,line width=.12pt] (-1.602769,-0.052966)--(-1.534468,0.052668)--(-1.589379,-0.024064)--(-1.660123,-0.133114)--cycle;
\path[fill=blue!27!white,draw=blue!27!white,line width=.12pt] (1.660123,-0.133114)--(1.589379,-0.024064)--(1.534468,0.052668)--(1.602769,-0.052966)--cycle;
\path[fill=blue!35!white,draw=blue!35!white,line width=.12pt] (-2.908701,-0.200610)--(-2.789801,-0.298596)--(-2.830759,-0.436656)--(-2.951405,-0.344554)--cycle;
\path[fill=blue!40!white,draw=blue!40!white,line width=.12pt] (2.951405,-0.344554)--(2.830759,-0.436656)--(2.789801,-0.298596)--(2.908701,-0.200610)--cycle;
\path[fill=blue!26!white,draw=blue!26!white,line width=.12pt] (-1.330351,0.444456)--(-1.343133,0.583889)--(-1.424631,0.515903)--(-1.411074,0.377117)--cycle;
\path[fill=blue!18!white,draw=blue!18!white,line width=.12pt] (1.411074,0.377117)--(1.424631,0.515903)--(1.343133,0.583889)--(1.330351,0.444456)--cycle;
\path[fill=blue!15!white,draw=blue!15!white,line width=.12pt] (-2.386608,1.277917)--(-2.508586,1.246493)--(-2.627935,1.120322)--(-2.500154,1.157881)--cycle;
\path[fill=blue!19!white,draw=blue!19!white,line width=.12pt] (2.500154,1.157881)--(2.627935,1.120322)--(2.508586,1.246493)--(2.386608,1.277917)--cycle;
\path[fill=green!65!black!27!white,draw=green!65!black!27!white,line width=.12pt] (1.525177,0.999443)--(1.625975,1.119280)--(1.511656,1.196351)--(1.417944,1.071736)--cycle;
\path[fill=green!65!black!32!white,draw=green!65!black!32!white,line width=.12pt] (-1.417944,1.071736)--(-1.511656,1.196351)--(-1.625975,1.119280)--(-1.525177,0.999443)--cycle;
\path[fill=blue!19!white,draw=blue!19!white,line width=.12pt] (-1.341037,0.935191)--(-1.417944,1.071736)--(-1.525177,0.999443)--(-1.442454,0.866819)--cycle;
\path[fill=blue!13!white,draw=blue!13!white,line width=.12pt] (1.442454,0.866819)--(1.525177,0.999443)--(1.417944,1.071736)--(1.341037,0.935191)--cycle;
\path[fill=blue!40!white,draw=blue!40!white,line width=.12pt] (-2.200811,-0.460624)--(-2.053620,-0.435837)--(-2.083771,-0.537466)--(-2.233123,-0.569536)--cycle;
\path[fill=blue!37!white,draw=blue!37!white,line width=.12pt] (2.233123,-0.569536)--(2.083771,-0.537466)--(2.053620,-0.435837)--(2.200811,-0.460624)--cycle;
\path[fill=green!65!black!29!white,draw=green!65!black!29!white,line width=.12pt] (-1.865758,1.302838)--(-1.995527,1.359505)--(-2.116611,1.258496)--(-1.978968,1.208398)--cycle;
\path[fill=green!65!black!29!white,draw=green!65!black!29!white,line width=.12pt] (1.978968,1.208398)--(2.116611,1.258496)--(1.995527,1.359505)--(1.865758,1.302838)--cycle;
\path[fill=blue!32!white,draw=blue!32!white,line width=.12pt] (-3.006279,-0.085917)--(-2.908701,-0.200610)--(-2.951405,-0.344554)--(-3.050416,-0.234690)--cycle;
\path[fill=blue!39!white,draw=blue!39!white,line width=.12pt] (3.050416,-0.234690)--(2.951405,-0.344554)--(2.908701,-0.200610)--(3.006279,-0.085917)--cycle;
\path[fill=blue!29!white,draw=blue!29!white,line width=.12pt] (-1.424631,0.245578)--(-1.411074,0.377117)--(-1.478207,0.306147)--(-1.492409,0.173925)--cycle;
\path[fill=blue!21!white,draw=blue!21!white,line width=.12pt] (1.492409,0.173925)--(1.478207,0.306147)--(1.411074,0.377117)--(1.424631,0.245578)--cycle;
\path[fill=blue!23!white,draw=blue!23!white,line width=.12pt] (-3.047499,0.597387)--(-3.062209,0.464455)--(-3.138513,0.312145)--(-3.123436,0.445808)--cycle;
\path[fill=blue!32!white,draw=blue!32!white,line width=.12pt] (3.123436,0.445808)--(3.138513,0.312145)--(3.062209,0.464455)--(3.047499,0.597387)--cycle;
\path[fill=blue!37!white,draw=blue!37!white,line width=.12pt] (-1.755330,-0.223661)--(-1.660123,-0.133114)--(-1.701490,-0.215687)--(-1.799069,-0.310969)--cycle;
\path[fill=blue!30!white,draw=blue!30!white,line width=.12pt] (1.799069,-0.310969)--(1.701490,-0.215687)--(1.660123,-0.133114)--(1.755330,-0.223661)--cycle;
\path[fill=blue!17!white,draw=blue!17!white,line width=.12pt] (-2.739936,1.057357)--(-2.831853,0.971406)--(-2.933190,0.829797)--(-2.837984,0.920344)--cycle;
\path[fill=blue!24!white,draw=blue!24!white,line width=.12pt] (2.837984,0.920344)--(2.933190,0.829797)--(2.831853,0.971406)--(2.739936,1.057357)--cycle;
\path[fill=blue!24!white,draw=blue!24!white,line width=.12pt] (-1.343133,0.583889)--(-1.380985,0.726503)--(-1.464780,0.656601)--(-1.424631,0.515903)--cycle;
\path[fill=blue!16!white,draw=blue!16!white,line width=.12pt] (1.424631,0.515903)--(1.464780,0.656601)--(1.380985,0.726503)--(1.343133,0.583889)--cycle;
\path[fill=blue!40!white,draw=blue!40!white,line width=.12pt] (-2.053620,-0.435837)--(-1.917969,-0.385305)--(-1.946127,-0.480221)--(-2.083771,-0.537466)--cycle;
\path[fill=blue!36!white,draw=blue!36!white,line width=.12pt] (2.083771,-0.537466)--(1.946127,-0.480221)--(1.917969,-0.385305)--(2.053620,-0.435837)--cycle;
\path[fill=blue!24!white,draw=blue!24!white,line width=.12pt] (1.589379,-0.024064)--(1.545814,0.099296)--(1.492409,0.173925)--(1.534468,0.052668)--cycle;
\path[fill=blue!33!white,draw=blue!33!white,line width=.12pt] (-1.534468,0.052668)--(-1.492409,0.173925)--(-1.545814,0.099296)--(-1.589379,-0.024064)--cycle;
\path[fill=blue!15!white,draw=blue!15!white,line width=.12pt] (-2.254254,1.281970)--(-2.386608,1.277917)--(-2.500154,1.157881)--(-2.361503,1.168591)--cycle;
\path[fill=blue!17!white,draw=blue!17!white,line width=.12pt] (2.361503,1.168591)--(2.500154,1.157881)--(2.386608,1.277917)--(2.254254,1.281970)--cycle;
\path[fill=green!65!black!30!white,draw=green!65!black!30!white,line width=.12pt] (-1.740975,1.221724)--(-1.865758,1.302838)--(-1.978968,1.208398)--(-1.846614,1.133600)--cycle;
\path[fill=green!65!black!28!white,draw=green!65!black!28!white,line width=.12pt] (1.846614,1.133600)--(1.978968,1.208398)--(1.865758,1.302838)--(1.740975,1.221724)--cycle;
\path[fill=blue!30!white,draw=blue!30!white,line width=.12pt] (-3.078787,0.041074)--(-3.006279,-0.085917)--(-3.050416,-0.234690)--(-3.123988,-0.111288)--cycle;
\path[fill=blue!38!white,draw=blue!38!white,line width=.12pt] (3.123988,-0.111288)--(3.050416,-0.234690)--(3.006279,-0.085917)--(3.078787,0.041074)--cycle;
\path[fill=blue!15!white,draw=blue!15!white,line width=.12pt] (1.464780,0.656601)--(1.529978,0.793806)--(1.442454,0.866819)--(1.380985,0.726503)--cycle;
\path[fill=blue!21!white,draw=blue!21!white,line width=.12pt] (-1.380985,0.726503)--(-1.442454,0.866819)--(-1.529978,0.793806)--(-1.464780,0.656601)--cycle;
\path[fill=blue!27!white,draw=blue!27!white,line width=.12pt] (-1.411074,0.377117)--(-1.424631,0.515903)--(-1.492409,0.444251)--(-1.478207,0.306147)--cycle;
\path[fill=blue!19!white,draw=blue!19!white,line width=.12pt] (1.478207,0.306147)--(1.492409,0.444251)--(1.424631,0.515903)--(1.411074,0.377117)--cycle;
\path[fill=green!65!black!27!white,draw=green!65!black!27!white,line width=.12pt] (1.724636,1.036978)--(1.846614,1.133600)--(1.740975,1.221724)--(1.625975,1.119280)--cycle;
\path[fill=green!65!black!31!white,draw=green!65!black!31!white,line width=.12pt] (-1.625975,1.119280)--(-1.740975,1.221724)--(-1.846614,1.133600)--(-1.724636,1.036978)--cycle;
\path[fill=blue!21!white,draw=blue!21!white,line width=.12pt] (-3.003935,0.720748)--(-3.047499,0.597387)--(-3.123436,0.445808)--(-3.078787,0.571335)--cycle;
\path[fill=blue!29!white,draw=blue!29!white,line width=.12pt] (3.078787,0.571335)--(3.123436,0.445808)--(3.047499,0.597387)--(3.003935,0.720748)--cycle;
\path[fill=blue!39!white,draw=blue!39!white,line width=.12pt] (-1.917969,-0.385305)--(-1.799069,-0.310969)--(-1.825482,-0.400001)--(-1.946127,-0.480221)--cycle;
\path[fill=blue!33!white,draw=blue!33!white,line width=.12pt] (1.946127,-0.480221)--(1.825482,-0.400001)--(1.799069,-0.310969)--(1.917969,-0.385305)--cycle;
\path[fill=blue!35!white,draw=blue!35!white,line width=.12pt] (-1.660123,-0.133114)--(-1.589379,-0.024064)--(-1.628983,-0.103118)--(-1.701490,-0.215687)--cycle;
\path[fill=blue!28!white,draw=blue!28!white,line width=.12pt] (1.701490,-0.215687)--(1.628983,-0.103118)--(1.589379,-0.024064)--(1.660123,-0.133114)--cycle;
\path[fill=blue!19!white,draw=blue!19!white,line width=.12pt] (-1.442454,0.866819)--(-1.525177,0.999443)--(-1.617721,0.922243)--(-1.529978,0.793806)--cycle;
\path[fill=blue!13!white,draw=blue!13!white,line width=.12pt] (1.529978,0.793806)--(1.617721,0.922243)--(1.525177,0.999443)--(1.442454,0.866819)--cycle;
\path[fill=blue!21!white,draw=blue!21!white,line width=.12pt] (2.721975,0.988910)--(2.837984,0.920344)--(2.739936,1.057357)--(2.627935,1.120322)--cycle;
\path[fill=blue!16!white,draw=blue!16!white,line width=.12pt] (-2.627935,1.120322)--(-2.739936,1.057357)--(-2.837984,0.920344)--(-2.721975,0.988910)--cycle;
\path[fill=green!65!black!32!white,draw=green!65!black!32!white,line width=.12pt] (-1.525177,0.999443)--(-1.625975,1.119280)--(-1.724636,1.036978)--(-1.617721,0.922243)--cycle;
\path[fill=green!65!black!28!white,draw=green!65!black!28!white,line width=.12pt] (1.617721,0.922243)--(1.724636,1.036978)--(1.625975,1.119280)--(1.525177,0.999443)--cycle;
\path[fill=blue!14!white,draw=blue!14!white,line width=.12pt] (-2.116611,1.258496)--(-2.254254,1.281970)--(-2.361503,1.168591)--(-2.217311,1.152040)--cycle;
\path[fill=blue!15!white,draw=blue!15!white,line width=.12pt] (2.217311,1.152040)--(2.361503,1.168591)--(2.254254,1.281970)--(2.116611,1.258496)--cycle;
\path[fill=blue!27!white,draw=blue!27!white,line width=.12pt] (-3.123436,0.175483)--(-3.078787,0.041074)--(-3.123988,-0.111288)--(-3.169293,0.020912)--cycle;
\path[fill=blue!36!white,draw=blue!36!white,line width=.12pt] (3.169293,0.020912)--(3.123988,-0.111288)--(3.078787,0.041074)--(3.123436,0.175483)--cycle;
\path[fill=blue!30!white,draw=blue!30!white,line width=.12pt] (-1.492409,0.173925)--(-1.478207,0.306147)--(-1.531105,0.232228)--(-1.545814,0.099296)--cycle;
\path[fill=blue!22!white,draw=blue!22!white,line width=.12pt] (1.545814,0.099296)--(1.531105,0.232228)--(1.478207,0.306147)--(1.492409,0.173925)--cycle;
\path[fill=blue!39!white,draw=blue!39!white,line width=.12pt] (-2.543764,-0.554239)--(-2.388443,-0.575200)--(-2.400000,-0.692820)--(-2.556072,-0.679508)--cycle;
\path[fill=blue!40!white,draw=blue!40!white,line width=.12pt] (2.556072,-0.679508)--(2.400000,-0.692820)--(2.388443,-0.575200)--(2.543764,-0.554239)--cycle;
\path[fill=blue!40!white,draw=blue!40!white,line width=.12pt] (-2.388443,-0.575200)--(-2.233123,-0.569536)--(-2.243928,-0.679508)--(-2.400000,-0.692820)--cycle;
\path[fill=blue!39!white,draw=blue!39!white,line width=.12pt] (2.400000,-0.692820)--(2.243928,-0.679508)--(2.233123,-0.569536)--(2.388443,-0.575200)--cycle;
\path[fill=blue!38!white,draw=blue!38!white,line width=.12pt] (-2.693116,-0.507458)--(-2.543764,-0.554239)--(-2.556072,-0.679508)--(-2.706147,-0.640083)--cycle;
\path[fill=blue!40!white,draw=blue!40!white,line width=.12pt] (2.706147,-0.640083)--(2.556072,-0.679508)--(2.543764,-0.554239)--(2.693116,-0.507458)--cycle;
\path[fill=blue!24!white,draw=blue!24!white,line width=.12pt] (-1.424631,0.515903)--(-1.464780,0.656601)--(-1.534468,0.582930)--(-1.492409,0.444251)--cycle;
\path[fill=blue!17!white,draw=blue!17!white,line width=.12pt] (1.492409,0.444251)--(1.534468,0.582930)--(1.464780,0.656601)--(1.424631,0.515903)--cycle;
\path[fill=blue!38!white,draw=blue!38!white,line width=.12pt] (-1.799069,-0.310969)--(-1.701490,-0.215687)--(-1.726471,-0.299889)--(-1.825482,-0.400001)--cycle;
\path[fill=blue!31!white,draw=blue!31!white,line width=.12pt] (1.825482,-0.400001)--(1.726471,-0.299889)--(1.701490,-0.215687)--(1.799069,-0.310969)--cycle;
\path[fill=blue!40!white,draw=blue!40!white,line width=.12pt] (-2.233123,-0.569536)--(-2.083771,-0.537466)--(-2.093853,-0.640083)--(-2.243928,-0.679508)--cycle;
\path[fill=blue!38!white,draw=blue!38!white,line width=.12pt] (2.243928,-0.679508)--(2.093853,-0.640083)--(2.083771,-0.537466)--(2.233123,-0.569536)--cycle;
\path[fill=blue!36!white,draw=blue!36!white,line width=.12pt] (-2.830759,-0.436656)--(-2.693116,-0.507458)--(-2.706147,-0.640083)--(-2.844456,-0.576059)--cycle;
\path[fill=blue!40!white,draw=blue!40!white,line width=.12pt] (2.844456,-0.576059)--(2.706147,-0.640083)--(2.693116,-0.507458)--(2.830759,-0.436656)--cycle;
\path[fill=blue!19!white,draw=blue!19!white,line width=.12pt] (-2.933190,0.829797)--(-3.003935,0.720748)--(-3.078787,0.571335)--(-3.006279,0.683903)--cycle;
\path[fill=blue!26!white,draw=blue!26!white,line width=.12pt] (3.006279,0.683903)--(3.078787,0.571335)--(3.003935,0.720748)--(2.933190,0.829797)--cycle;
\path[fill=green!65!black!30!white,draw=green!65!black!30!white,line width=.12pt] (-1.978968,1.208398)--(-2.116611,1.258496)--(-2.217311,1.152040)--(-2.073119,1.108865)--cycle;
\path[fill=green!65!black!29!white,draw=green!65!black!29!white,line width=.12pt] (2.073119,1.108865)--(2.217311,1.152040)--(2.116611,1.258496)--(1.978968,1.208398)--cycle;
\path[fill=blue!33!white,draw=blue!33!white,line width=.12pt] (-1.589379,-0.024064)--(-1.545814,0.099296)--(-1.584333,0.022409)--(-1.628983,-0.103118)--cycle;
\path[fill=blue!25!white,draw=blue!25!white,line width=.12pt] (1.628983,-0.103118)--(1.584333,0.022409)--(1.545814,0.099296)--(1.589379,-0.024064)--cycle;
\path[fill=blue!15!white,draw=blue!15!white,line width=.12pt] (-2.500154,1.157881)--(-2.627935,1.120322)--(-2.721975,0.988910)--(-2.589621,1.032859)--cycle;
\path[fill=blue!19!white,draw=blue!19!white,line width=.12pt] (2.589621,1.032859)--(2.721975,0.988910)--(2.627935,1.120322)--(2.500154,1.157881)--cycle;
\path[fill=blue!25!white,draw=blue!25!white,line width=.12pt] (-3.138513,0.312145)--(-3.123436,0.175483)--(-3.169293,0.020912)--(-3.184591,0.156827)--cycle;
\path[fill=blue!34!white,draw=blue!34!white,line width=.12pt] (3.184591,0.156827)--(3.169293,0.020912)--(3.123436,0.175483)--(3.138513,0.312145)--cycle;
\path[fill=blue!22!white,draw=blue!22!white,line width=.12pt] (-1.464780,0.656601)--(-1.529978,0.793806)--(-1.602769,0.716855)--(-1.534468,0.582930)--cycle;
\path[fill=blue!15!white,draw=blue!15!white,line width=.12pt] (1.534468,0.582930)--(1.602769,0.716855)--(1.529978,0.793806)--(1.464780,0.656601)--cycle;
\path[fill=blue!40!white,draw=blue!40!white,line width=.12pt] (-2.083771,-0.537466)--(-1.946127,-0.480221)--(-1.955544,-0.576059)--(-2.093853,-0.640083)--cycle;
\path[fill=blue!36!white,draw=blue!36!white,line width=.12pt] (2.093853,-0.640083)--(1.955544,-0.576059)--(1.946127,-0.480221)--(2.083771,-0.537466)--cycle;
\path[fill=blue!34!white,draw=blue!34!white,line width=.12pt] (-2.951405,-0.344554)--(-2.830759,-0.436656)--(-2.844456,-0.576059)--(-2.965685,-0.489898)--cycle;
\path[fill=blue!40!white,draw=blue!40!white,line width=.12pt] (2.965685,-0.489898)--(2.844456,-0.576059)--(2.830759,-0.436656)--(2.951405,-0.344554)--cycle;
\path[fill=blue!28!white,draw=blue!28!white,line width=.12pt] (-1.478207,0.306147)--(-1.492409,0.444251)--(-1.545814,0.369621)--(-1.531105,0.232228)--cycle;
\path[fill=blue!19!white,draw=blue!19!white,line width=.12pt] (1.531105,0.232228)--(1.545814,0.369621)--(1.492409,0.444251)--(1.478207,0.306147)--cycle;
\path[fill=green!65!black!30!white,draw=green!65!black!30!white,line width=.12pt] (-1.846614,1.133600)--(-1.978968,1.208398)--(-2.073119,1.108865)--(-1.934468,1.040724)--cycle;
\path[fill=green!65!black!28!white,draw=green!65!black!28!white,line width=.12pt] (1.934468,1.040724)--(2.073119,1.108865)--(1.978968,1.208398)--(1.846614,1.133600)--cycle;
\path[fill=blue!20!white,draw=blue!20!white,line width=.12pt] (-1.529978,0.793806)--(-1.617721,0.922243)--(-1.694686,0.840879)--(-1.602769,0.716855)--cycle;
\path[fill=blue!14!white,draw=blue!14!white,line width=.12pt] (1.602769,0.716855)--(1.694686,0.840879)--(1.617721,0.922243)--(1.529978,0.793806)--cycle;
\path[fill=blue!36!white,draw=blue!36!white,line width=.12pt] (-1.701490,-0.215687)--(-1.628983,-0.103118)--(-1.652899,-0.183733)--(-1.726471,-0.299889)--cycle;
\path[fill=blue!28!white,draw=blue!28!white,line width=.12pt] (1.726471,-0.299889)--(1.652899,-0.183733)--(1.628983,-0.103118)--(1.701490,-0.215687)--cycle;
\path[fill=green!65!black!31!white,draw=green!65!black!31!white,line width=.12pt] (-1.724636,1.036978)--(-1.846614,1.133600)--(-1.934468,1.040724)--(-1.806687,0.950236)--cycle;
\path[fill=green!65!black!28!white,draw=green!65!black!28!white,line width=.12pt] (1.806687,0.950236)--(1.934468,1.040724)--(1.846614,1.133600)--(1.724636,1.036978)--cycle;
\path[fill=green!65!black!33!white,draw=green!65!black!33!white,line width=.12pt] (-1.617721,0.922243)--(-1.724636,1.036978)--(-1.806687,0.950236)--(-1.694686,0.840879)--cycle;
\path[fill=green!65!black!28!white,draw=green!65!black!28!white,line width=.12pt] (1.694686,0.840879)--(1.806687,0.950236)--(1.724636,1.036978)--(1.617721,0.922243)--cycle;
\path[fill=blue!39!white,draw=blue!39!white,line width=.12pt] (-1.946127,-0.480221)--(-1.825482,-0.400001)--(-1.834315,-0.489898)--(-1.955544,-0.576059)--cycle;
\path[fill=blue!34!white,draw=blue!34!white,line width=.12pt] (1.955544,-0.576059)--(1.834315,-0.489898)--(1.825482,-0.400001)--(1.946127,-0.480221)--cycle;
\path[fill=blue!14!white,draw=blue!14!white,line width=.12pt] (-2.361503,1.168591)--(-2.500154,1.157881)--(-2.589621,1.032859)--(-2.446009,1.050502)--cycle;
\path[fill=blue!17!white,draw=blue!17!white,line width=.12pt] (2.446009,1.050502)--(2.589621,1.032859)--(2.500154,1.157881)--(2.361503,1.168591)--cycle;
\path[fill=blue!17!white,draw=blue!17!white,line width=.12pt] (-2.837984,0.920344)--(-2.933190,0.829797)--(-3.006279,0.683903)--(-2.908701,0.779186)--cycle;
\path[fill=blue!24!white,draw=blue!24!white,line width=.12pt] (2.908701,0.779186)--(3.006279,0.683903)--(2.933190,0.829797)--(2.837984,0.920344)--cycle;
\path[fill=blue!32!white,draw=blue!32!white,line width=.12pt] (-3.050416,-0.234690)--(-2.951405,-0.344554)--(-2.965685,-0.489898)--(-3.065176,-0.384910)--cycle;
\path[fill=blue!39!white,draw=blue!39!white,line width=.12pt] (3.065176,-0.384910)--(2.965685,-0.489898)--(2.951405,-0.344554)--(3.050416,-0.234690)--cycle;
\path[fill=blue!31!white,draw=blue!31!white,line width=.12pt] (-1.545814,0.099296)--(-1.531105,0.232228)--(-1.569256,0.156072)--(-1.584333,0.022409)--cycle;
\path[fill=blue!22!white,draw=blue!22!white,line width=.12pt] (1.584333,0.022409)--(1.569256,0.156072)--(1.531105,0.232228)--(1.545814,0.099296)--cycle;
\path[fill=blue!22!white,draw=blue!22!white,line width=.12pt] (-3.123436,0.445808)--(-3.138513,0.312145)--(-3.184591,0.156827)--(-3.169293,0.291237)--cycle;
\path[fill=blue!31!white,draw=blue!31!white,line width=.12pt] (3.169293,0.291237)--(3.184591,0.156827)--(3.138513,0.312145)--(3.123436,0.445808)--cycle;
\path[fill=blue!25!white,draw=blue!25!white,line width=.12pt] (-1.492409,0.444251)--(-1.534468,0.582930)--(-1.589379,0.506197)--(-1.545814,0.369621)--cycle;
\path[fill=blue!17!white,draw=blue!17!white,line width=.12pt] (1.545814,0.369621)--(1.589379,0.506197)--(1.534468,0.582930)--(1.492409,0.444251)--cycle;
\path[fill=blue!38!white,draw=blue!38!white,line width=.12pt] (-1.825482,-0.400001)--(-1.726471,-0.299889)--(-1.734824,-0.384910)--(-1.834315,-0.489898)--cycle;
\path[fill=blue!31!white,draw=blue!31!white,line width=.12pt] (1.834315,-0.489898)--(1.734824,-0.384910)--(1.726471,-0.299889)--(1.825482,-0.400001)--cycle;
\path[fill=blue!14!white,draw=blue!14!white,line width=.12pt] (-2.217311,1.152040)--(-2.361503,1.168591)--(-2.446009,1.050502)--(-2.296657,1.041162)--cycle;
\path[fill=blue!15!white,draw=blue!15!white,line width=.12pt] (2.296657,1.041162)--(2.446009,1.050502)--(2.361503,1.168591)--(2.217311,1.152040)--cycle;
\path[fill=blue!34!white,draw=blue!34!white,line width=.12pt] (-1.628983,-0.103118)--(-1.584333,0.022409)--(-1.607593,-0.055995)--(-1.652899,-0.183733)--cycle;
\path[fill=blue!25!white,draw=blue!25!white,line width=.12pt] (1.652899,-0.183733)--(1.607593,-0.055995)--(1.584333,0.022409)--(1.628983,-0.103118)--cycle;
\path[fill=blue!29!white,draw=blue!29!white,line width=.12pt] (-3.123988,-0.111288)--(-3.050416,-0.234690)--(-3.065176,-0.384910)--(-3.139104,-0.265131)--cycle;
\path[fill=blue!37!white,draw=blue!37!white,line width=.12pt] (3.139104,-0.265131)--(3.065176,-0.384910)--(3.050416,-0.234690)--(3.123988,-0.111288)--cycle;
\path[fill=blue!23!white,draw=blue!23!white,line width=.12pt] (-1.534468,0.582930)--(-1.602769,0.716855)--(-1.660123,0.636707)--(-1.589379,0.506197)--cycle;
\path[fill=blue!15!white,draw=blue!15!white,line width=.12pt] (1.589379,0.506197)--(1.660123,0.636707)--(1.602769,0.716855)--(1.534468,0.582930)--cycle;
\path[fill=blue!16!white,draw=blue!16!white,line width=.12pt] (-2.721975,0.988910)--(-2.837984,0.920344)--(-2.908701,0.779186)--(-2.789801,0.853522)--cycle;
\path[fill=blue!21!white,draw=blue!21!white,line width=.12pt] (2.789801,0.853522)--(2.908701,0.779186)--(2.837984,0.920344)--(2.721975,0.988910)--cycle;
\path[fill=blue!28!white,draw=blue!28!white,line width=.12pt] (-1.531105,0.232228)--(-1.545814,0.369621)--(-1.584333,0.292734)--(-1.569256,0.156072)--cycle;
\path[fill=blue!20!white,draw=blue!20!white,line width=.12pt] (1.569256,0.156072)--(1.584333,0.292734)--(1.545814,0.369621)--(1.531105,0.232228)--cycle;
\path[fill=blue!40!white,draw=blue!40!white,line width=.12pt] (-2.400000,-0.692820)--(-2.243928,-0.679508)--(-2.233123,-0.789480)--(-2.388443,-0.810441)--cycle;
\path[fill=blue!39!white,draw=blue!39!white,line width=.12pt] (2.388443,-0.810441)--(2.233123,-0.789480)--(2.243928,-0.679508)--(2.400000,-0.692820)--cycle;
\path[fill=blue!39!white,draw=blue!39!white,line width=.12pt] (-2.556072,-0.679508)--(-2.400000,-0.692820)--(-2.388443,-0.810441)--(-2.543764,-0.804777)--cycle;
\path[fill=blue!40!white,draw=blue!40!white,line width=.12pt] (2.543764,-0.804777)--(2.388443,-0.810441)--(2.400000,-0.692820)--(2.556072,-0.679508)--cycle;
\path[fill=blue!40!white,draw=blue!40!white,line width=.12pt] (-2.243928,-0.679508)--(-2.093853,-0.640083)--(-2.083771,-0.742699)--(-2.233123,-0.789480)--cycle;
\path[fill=blue!38!white,draw=blue!38!white,line width=.12pt] (2.233123,-0.789480)--(2.083771,-0.742699)--(2.093853,-0.640083)--(2.243928,-0.679508)--cycle;
\path[fill=blue!20!white,draw=blue!20!white,line width=.12pt] (-3.078787,0.571335)--(-3.123436,0.445808)--(-3.169293,0.291237)--(-3.123988,0.418974)--cycle;
\path[fill=blue!29!white,draw=blue!29!white,line width=.12pt] (3.123988,0.418974)--(3.169293,0.291237)--(3.123436,0.445808)--(3.078787,0.571335)--cycle;
\path[fill=green!65!black!30!white,draw=green!65!black!30!white,line width=.12pt] (-2.073119,1.108865)--(-2.217311,1.152040)--(-2.296657,1.041162)--(-2.147305,1.005197)--cycle;
\path[fill=green!65!black!29!white,draw=green!65!black!29!white,line width=.12pt] (2.147305,1.005197)--(2.296657,1.041162)--(2.217311,1.152040)--(2.073119,1.108865)--cycle;
\path[fill=blue!38!white,draw=blue!38!white,line width=.12pt] (-2.706147,-0.640083)--(-2.556072,-0.679508)--(-2.543764,-0.804777)--(-2.693116,-0.772707)--cycle;
\path[fill=blue!40!white,draw=blue!40!white,line width=.12pt] (2.693116,-0.772707)--(2.543764,-0.804777)--(2.556072,-0.679508)--(2.706147,-0.640083)--cycle;
\path[fill=blue!21!white,draw=blue!21!white,line width=.12pt] (-1.602769,0.716855)--(-1.694686,0.840879)--(-1.755330,0.756135)--(-1.660123,0.636707)--cycle;
\path[fill=blue!14!white,draw=blue!14!white,line width=.12pt] (1.660123,0.636707)--(1.755330,0.756135)--(1.694686,0.840879)--(1.602769,0.716855)--cycle;
\path[fill=blue!40!white,draw=blue!40!white,line width=.12pt] (-2.093853,-0.640083)--(-1.955544,-0.576059)--(-1.946127,-0.671897)--(-2.083771,-0.742699)--cycle;
\path[fill=blue!36!white,draw=blue!36!white,line width=.12pt] (2.083771,-0.742699)--(1.946127,-0.671897)--(1.955544,-0.576059)--(2.093853,-0.640083)--cycle;
\path[fill=blue!37!white,draw=blue!37!white,line width=.12pt] (-1.726471,-0.299889)--(-1.652899,-0.183733)--(-1.660896,-0.265131)--(-1.734824,-0.384910)--cycle;
\path[fill=blue!29!white,draw=blue!29!white,line width=.12pt] (1.734824,-0.384910)--(1.660896,-0.265131)--(1.652899,-0.183733)--(1.726471,-0.299889)--cycle;
\path[fill=green!65!black!31!white,draw=green!65!black!31!white,line width=.12pt] (-1.934468,1.040724)--(-2.073119,1.108865)--(-2.147305,1.005197)--(-2.003693,0.943989)--cycle;
\path[fill=green!65!black!28!white,draw=green!65!black!28!white,line width=.12pt] (2.003693,0.943989)--(2.147305,1.005197)--(2.073119,1.108865)--(1.934468,1.040724)--cycle;
\path[fill=blue!27!white,draw=blue!27!white,line width=.12pt] (-3.169293,0.020912)--(-3.123988,-0.111288)--(-3.139104,-0.265131)--(-3.184628,-0.135163)--cycle;
\path[fill=blue!35!white,draw=blue!35!white,line width=.12pt] (3.184628,-0.135163)--(3.139104,-0.265131)--(3.123988,-0.111288)--(3.169293,0.020912)--cycle;
\path[fill=green!65!black!33!white,draw=green!65!black!33!white,line width=.12pt] (-1.694686,0.840879)--(-1.806687,0.950236)--(-1.871339,0.859891)--(-1.755330,0.756135)--cycle;
\path[fill=green!65!black!28!white,draw=green!65!black!28!white,line width=.12pt] (1.755330,0.756135)--(1.871339,0.859891)--(1.806687,0.950236)--(1.694686,0.840879)--cycle;
\path[fill=blue!32!white,draw=blue!32!white,line width=.12pt] (-1.584333,0.022409)--(-1.569256,0.156072)--(-1.592296,0.078414)--(-1.607593,-0.055995)--cycle;
\path[fill=blue!23!white,draw=blue!23!white,line width=.12pt] (1.607593,-0.055995)--(1.592296,0.078414)--(1.569256,0.156072)--(1.584333,0.022409)--cycle;
\path[fill=green!65!black!32!white,draw=green!65!black!32!white,line width=.12pt] (-1.806687,0.950236)--(-1.934468,1.040724)--(-2.003693,0.943989)--(-1.871339,0.859891)--cycle;
\path[fill=green!65!black!28!white,draw=green!65!black!28!white,line width=.12pt] (1.871339,0.859891)--(2.003693,0.943989)--(1.934468,1.040724)--(1.806687,0.950236)--cycle;
\path[fill=blue!36!white,draw=blue!36!white,line width=.12pt] (-2.844456,-0.576059)--(-2.706147,-0.640083)--(-2.693116,-0.772707)--(-2.830759,-0.715462)--cycle;
\path[fill=blue!40!white,draw=blue!40!white,line width=.12pt] (2.830759,-0.715462)--(2.693116,-0.772707)--(2.706147,-0.640083)--(2.844456,-0.576059)--cycle;
\path[fill=blue!15!white,draw=blue!15!white,line width=.12pt] (-2.589621,1.032859)--(-2.721975,0.988910)--(-2.789801,0.853522)--(-2.654149,0.904054)--cycle;
\path[fill=blue!19!white,draw=blue!19!white,line width=.12pt] (2.654149,0.904054)--(2.789801,0.853522)--(2.721975,0.988910)--(2.589621,1.032859)--cycle;
\path[fill=blue!26!white,draw=blue!26!white,line width=.12pt] (-1.545814,0.369621)--(-1.589379,0.506197)--(-1.628983,0.427143)--(-1.584333,0.292734)--cycle;
\path[fill=blue!18!white,draw=blue!18!white,line width=.12pt] (1.584333,0.292734)--(1.628983,0.427143)--(1.589379,0.506197)--(1.545814,0.369621)--cycle;
\path[fill=blue!40!white,draw=blue!40!white,line width=.12pt] (-1.955544,-0.576059)--(-1.834315,-0.489898)--(-1.825482,-0.579795)--(-1.946127,-0.671897)--cycle;
\path[fill=blue!34!white,draw=blue!34!white,line width=.12pt] (1.946127,-0.671897)--(1.825482,-0.579795)--(1.834315,-0.489898)--(1.955544,-0.576059)--cycle;
\path[fill=blue!18!white,draw=blue!18!white,line width=.12pt] (-3.006279,0.683903)--(-3.078787,0.571335)--(-3.123988,0.418974)--(-3.050416,0.535130)--cycle;
\path[fill=blue!26!white,draw=blue!26!white,line width=.12pt] (3.050416,0.535130)--(3.123988,0.418974)--(3.078787,0.571335)--(3.006279,0.683903)--cycle;
\path[fill=blue!35!white,draw=blue!35!white,line width=.12pt] (-1.652899,-0.183733)--(-1.607593,-0.055995)--(-1.615372,-0.135163)--(-1.660896,-0.265131)--cycle;
\path[fill=blue!26!white,draw=blue!26!white,line width=.12pt] (1.660896,-0.265131)--(1.615372,-0.135163)--(1.607593,-0.055995)--(1.652899,-0.183733)--cycle;
\path[fill=blue!34!white,draw=blue!34!white,line width=.12pt] (-2.965685,-0.489898)--(-2.844456,-0.576059)--(-2.830759,-0.715462)--(-2.951405,-0.635242)--cycle;
\path[fill=blue!39!white,draw=blue!39!white,line width=.12pt] (2.951405,-0.635242)--(2.830759,-0.715462)--(2.844456,-0.576059)--(2.965685,-0.489898)--cycle;
\path[fill=blue!39!white,draw=blue!39!white,line width=.12pt] (-1.834315,-0.489898)--(-1.734824,-0.384910)--(-1.726471,-0.469932)--(-1.825482,-0.579795)--cycle;
\path[fill=blue!32!white,draw=blue!32!white,line width=.12pt] (1.825482,-0.579795)--(1.726471,-0.469932)--(1.734824,-0.384910)--(1.834315,-0.489898)--cycle;
\path[fill=blue!24!white,draw=blue!24!white,line width=.12pt] (-3.184591,0.156827)--(-3.169293,0.020912)--(-3.184628,-0.135163)--(-3.200000,0.000000)--cycle;
\path[fill=blue!33!white,draw=blue!33!white,line width=.12pt] (3.200000,0.000000)--(3.184628,-0.135163)--(3.169293,0.020912)--(3.184591,0.156827)--cycle;
\path[fill=blue!14!white,draw=blue!14!white,line width=.12pt] (-2.446009,1.050502)--(-2.589621,1.032859)--(-2.654149,0.904054)--(-2.506958,0.928840)--cycle;
\path[fill=blue!17!white,draw=blue!17!white,line width=.12pt] (2.506958,0.928840)--(2.654149,0.904054)--(2.589621,1.032859)--(2.446009,1.050502)--cycle;
\path[fill=blue!23!white,draw=blue!23!white,line width=.12pt] (-1.589379,0.506197)--(-1.660123,0.636707)--(-1.701490,0.554134)--(-1.628983,0.427143)--cycle;
\path[fill=blue!16!white,draw=blue!16!white,line width=.12pt] (1.628983,0.427143)--(1.701490,0.554134)--(1.660123,0.636707)--(1.589379,0.506197)--cycle;
\path[fill=blue!29!white,draw=blue!29!white,line width=.12pt] (-1.569256,0.156072)--(-1.584333,0.292734)--(-1.607593,0.214330)--(-1.592296,0.078414)--cycle;
\path[fill=blue!20!white,draw=blue!20!white,line width=.12pt] (1.592296,0.078414)--(1.607593,0.214330)--(1.584333,0.292734)--(1.569256,0.156072)--cycle;
\path[fill=blue!16!white,draw=blue!16!white,line width=.12pt] (-2.908701,0.779186)--(-3.006279,0.683903)--(-3.050416,0.535130)--(-2.951405,0.635242)--cycle;
\path[fill=blue!23!white,draw=blue!23!white,line width=.12pt] (2.951405,0.635242)--(3.050416,0.535130)--(3.006279,0.683903)--(2.908701,0.779186)--cycle;
\path[fill=blue!14!white,draw=blue!14!white,line width=.12pt] (-2.296657,1.041162)--(-2.446009,1.050502)--(-2.506958,0.928840)--(-2.353885,0.926929)--cycle;
\path[fill=blue!15!white,draw=blue!15!white,line width=.12pt] (2.353885,0.926929)--(2.506958,0.928840)--(2.446009,1.050502)--(2.296657,1.041162)--cycle;
\path[fill=blue!31!white,draw=blue!31!white,line width=.12pt] (-3.065176,-0.384910)--(-2.965685,-0.489898)--(-2.951405,-0.635242)--(-3.050416,-0.535130)--cycle;
\path[fill=blue!38!white,draw=blue!38!white,line width=.12pt] (3.050416,-0.535130)--(2.951405,-0.635242)--(2.965685,-0.489898)--(3.065176,-0.384910)--cycle;
\path[fill=blue!21!white,draw=blue!21!white,line width=.12pt] (-1.660123,0.636707)--(-1.755330,0.756135)--(-1.799069,0.668826)--(-1.701490,0.554134)--cycle;
\path[fill=blue!14!white,draw=blue!14!white,line width=.12pt] (1.701490,0.554134)--(1.799069,0.668826)--(1.755330,0.756135)--(1.660123,0.636707)--cycle;
\path[fill=blue!37!white,draw=blue!37!white,line width=.12pt] (-1.734824,-0.384910)--(-1.660896,-0.265131)--(-1.652899,-0.346529)--(-1.726471,-0.469932)--cycle;
\path[fill=blue!29!white,draw=blue!29!white,line width=.12pt] (1.726471,-0.469932)--(1.652899,-0.346529)--(1.660896,-0.265131)--(1.734824,-0.384910)--cycle;
\path[fill=blue!32!white,draw=blue!32!white,line width=.12pt] (-1.607593,-0.055995)--(-1.592296,0.078414)--(-1.600000,-0.000000)--(-1.615372,-0.135163)--cycle;
\path[fill=blue!23!white,draw=blue!23!white,line width=.12pt] (1.615372,-0.135163)--(1.600000,-0.000000)--(1.592296,0.078414)--(1.607593,-0.055995)--cycle;
\path[fill=blue!41!white,draw=blue!41!white,line width=.12pt] (-2.233123,-0.789480)--(-2.083771,-0.742699)--(-2.053620,-0.844328)--(-2.200811,-0.898392)--cycle;
\path[fill=blue!38!white,draw=blue!38!white,line width=.12pt] (2.200811,-0.898392)--(2.053620,-0.844328)--(2.083771,-0.742699)--(2.233123,-0.789480)--cycle;
\path[fill=blue!40!white,draw=blue!40!white,line width=.12pt] (-2.388443,-0.810441)--(-2.233123,-0.789480)--(-2.200811,-0.898392)--(-2.353885,-0.926929)--cycle;
\path[fill=blue!39!white,draw=blue!39!white,line width=.12pt] (2.353885,-0.926929)--(2.200811,-0.898392)--(2.233123,-0.789480)--(2.388443,-0.810441)--cycle;
\path[fill=blue!22!white,draw=blue!22!white,line width=.12pt] (-3.169293,0.291237)--(-3.184591,0.156827)--(-3.200000,0.000000)--(-3.184628,0.135163)--cycle;
\path[fill=blue!31!white,draw=blue!31!white,line width=.12pt] (3.184628,0.135163)--(3.200000,0.000000)--(3.184591,0.156827)--(3.169293,0.291237)--cycle;
\path[fill=blue!41!white,draw=blue!41!white,line width=.12pt] (-2.083771,-0.742699)--(-1.946127,-0.671897)--(-1.917969,-0.766813)--(-2.053620,-0.844328)--cycle;
\path[fill=blue!36!white,draw=blue!36!white,line width=.12pt] (2.053620,-0.844328)--(1.917969,-0.766813)--(1.946127,-0.671897)--(2.083771,-0.742699)--cycle;
\path[fill=blue!27!white,draw=blue!27!white,line width=.12pt] (-1.584333,0.292734)--(-1.628983,0.427143)--(-1.652899,0.346529)--(-1.607593,0.214330)--cycle;
\path[fill=blue!18!white,draw=blue!18!white,line width=.12pt] (1.607593,0.214330)--(1.652899,0.346529)--(1.628983,0.427143)--(1.584333,0.292734)--cycle;
\path[fill=green!65!black!30!white,draw=green!65!black!30!white,line width=.12pt] (-2.147305,1.005197)--(-2.296657,1.041162)--(-2.353885,0.926929)--(-2.200811,0.898392)--cycle;
\path[fill=green!65!black!29!white,draw=green!65!black!29!white,line width=.12pt] (2.200811,0.898392)--(2.353885,0.926929)--(2.296657,1.041162)--(2.147305,1.005197)--cycle;
\path[fill=green!65!black!34!white,draw=green!65!black!34!white,line width=.12pt] (-1.755330,0.756135)--(-1.871339,0.859891)--(-1.917969,0.766813)--(-1.799069,0.668826)--cycle;
\path[fill=green!65!black!28!white,draw=green!65!black!28!white,line width=.12pt] (1.799069,0.668826)--(1.917969,0.766813)--(1.871339,0.859891)--(1.755330,0.756135)--cycle;
\path[fill=blue!39!white,draw=blue!39!white,line width=.12pt] (-2.543764,-0.804777)--(-2.388443,-0.810441)--(-2.353885,-0.926929)--(-2.506958,-0.928840)--cycle;
\path[fill=blue!40!white,draw=blue!40!white,line width=.12pt] (2.506958,-0.928840)--(2.353885,-0.926929)--(2.388443,-0.810441)--(2.543764,-0.804777)--cycle;
\path[fill=green!65!black!31!white,draw=green!65!black!31!white,line width=.12pt] (-2.003693,0.943989)--(-2.147305,1.005197)--(-2.200811,0.898392)--(-2.053620,0.844328)--cycle;
\path[fill=green!65!black!28!white,draw=green!65!black!28!white,line width=.12pt] (2.053620,0.844328)--(2.200811,0.898392)--(2.147305,1.005197)--(2.003693,0.943989)--cycle;
\path[fill=green!65!black!32!white,draw=green!65!black!32!white,line width=.12pt] (-1.871339,0.859891)--(-2.003693,0.943989)--(-2.053620,0.844328)--(-1.917969,0.766813)--cycle;
\path[fill=green!65!black!28!white,draw=green!65!black!28!white,line width=.12pt] (1.917969,0.766813)--(2.053620,0.844328)--(2.003693,0.943989)--(1.871339,0.859891)--cycle;
\path[fill=blue!40!white,draw=blue!40!white,line width=.12pt] (-1.946127,-0.671897)--(-1.825482,-0.579795)--(-1.799069,-0.668826)--(-1.917969,-0.766813)--cycle;
\path[fill=blue!35!white,draw=blue!35!white,line width=.12pt] (1.917969,-0.766813)--(1.799069,-0.668826)--(1.825482,-0.579795)--(1.946127,-0.671897)--cycle;
\path[fill=blue!15!white,draw=blue!15!white,line width=.12pt] (-2.789801,0.853522)--(-2.908701,0.779186)--(-2.951405,0.635242)--(-2.830759,0.715462)--cycle;
\path[fill=blue!21!white,draw=blue!21!white,line width=.12pt] (2.830759,0.715462)--(2.951405,0.635242)--(2.908701,0.779186)--(2.789801,0.853522)--cycle;
\path[fill=blue!29!white,draw=blue!29!white,line width=.12pt] (-3.139104,-0.265131)--(-3.065176,-0.384910)--(-3.050416,-0.535130)--(-3.123988,-0.418974)--cycle;
\path[fill=blue!37!white,draw=blue!37!white,line width=.12pt] (3.123988,-0.418974)--(3.050416,-0.535130)--(3.065176,-0.384910)--(3.139104,-0.265131)--cycle;
\path[fill=blue!37!white,draw=blue!37!white,line width=.12pt] (-2.693116,-0.772707)--(-2.543764,-0.804777)--(-2.506958,-0.928840)--(-2.654149,-0.904054)--cycle;
\path[fill=blue!40!white,draw=blue!40!white,line width=.12pt] (2.654149,-0.904054)--(2.506958,-0.928840)--(2.543764,-0.804777)--(2.693116,-0.772707)--cycle;
\path[fill=blue!35!white,draw=blue!35!white,line width=.12pt] (-1.660896,-0.265131)--(-1.615372,-0.135163)--(-1.607593,-0.214330)--(-1.652899,-0.346529)--cycle;
\path[fill=blue!27!white,draw=blue!27!white,line width=.12pt] (1.652899,-0.346529)--(1.607593,-0.214330)--(1.615372,-0.135163)--(1.660896,-0.265131)--cycle;
\path[fill=blue!30!white,draw=blue!30!white,line width=.12pt] (-1.592296,0.078414)--(-1.607593,0.214330)--(-1.615372,0.135163)--(-1.600000,0.000000)--cycle;
\path[fill=blue!21!white,draw=blue!21!white,line width=.12pt] (1.600000,0.000000)--(1.615372,0.135163)--(1.607593,0.214330)--(1.592296,0.078414)--cycle;
\path[fill=blue!24!white,draw=blue!24!white,line width=.12pt] (-1.628983,0.427143)--(-1.701490,0.554134)--(-1.726471,0.469932)--(-1.652899,0.346529)--cycle;
\path[fill=blue!16!white,draw=blue!16!white,line width=.12pt] (1.652899,0.346529)--(1.726471,0.469932)--(1.701490,0.554134)--(1.628983,0.427143)--cycle;
\path[fill=blue!39!white,draw=blue!39!white,line width=.12pt] (-1.825482,-0.579795)--(-1.726471,-0.469932)--(-1.701490,-0.554134)--(-1.799069,-0.668826)--cycle;
\path[fill=blue!32!white,draw=blue!32!white,line width=.12pt] (1.799069,-0.668826)--(1.701490,-0.554134)--(1.726471,-0.469932)--(1.825482,-0.579795)--cycle;
\path[fill=blue!19!white,draw=blue!19!white,line width=.12pt] (-3.123988,0.418974)--(-3.169293,0.291237)--(-3.184628,0.135163)--(-3.139104,0.265131)--cycle;
\path[fill=blue!28!white,draw=blue!28!white,line width=.12pt] (3.139104,0.265131)--(3.184628,0.135163)--(3.169293,0.291237)--(3.123988,0.418974)--cycle;
\path[fill=blue!14!white,draw=blue!14!white,line width=.12pt] (-2.654149,0.904054)--(-2.789801,0.853522)--(-2.830759,0.715462)--(-2.693116,0.772707)--cycle;
\path[fill=blue!18!white,draw=blue!18!white,line width=.12pt] (2.693116,0.772707)--(2.830759,0.715462)--(2.789801,0.853522)--(2.654149,0.904054)--cycle;
\path[fill=blue!36!white,draw=blue!36!white,line width=.12pt] (-2.830759,-0.715462)--(-2.693116,-0.772707)--(-2.654149,-0.904054)--(-2.789801,-0.853522)--cycle;
\path[fill=blue!40!white,draw=blue!40!white,line width=.12pt] (2.789801,-0.853522)--(2.654149,-0.904054)--(2.693116,-0.772707)--(2.830759,-0.715462)--cycle;
\path[fill=blue!26!white,draw=blue!26!white,line width=.12pt] (-3.184628,-0.135163)--(-3.139104,-0.265131)--(-3.123988,-0.418974)--(-3.169293,-0.291237)--cycle;
\path[fill=blue!35!white,draw=blue!35!white,line width=.12pt] (3.169293,-0.291237)--(3.123988,-0.418974)--(3.139104,-0.265131)--(3.184628,-0.135163)--cycle;
\path[fill=blue!22!white,draw=blue!22!white,line width=.12pt] (-1.701490,0.554134)--(-1.799069,0.668826)--(-1.825482,0.579795)--(-1.726471,0.469932)--cycle;
\path[fill=blue!15!white,draw=blue!15!white,line width=.12pt] (1.726471,0.469932)--(1.825482,0.579795)--(1.799069,0.668826)--(1.701490,0.554134)--cycle;
\path[fill=blue!38!white,draw=blue!38!white,line width=.12pt] (-1.726471,-0.469932)--(-1.652899,-0.346529)--(-1.628983,-0.427143)--(-1.701490,-0.554134)--cycle;
\path[fill=blue!30!white,draw=blue!30!white,line width=.12pt] (1.701490,-0.554134)--(1.628983,-0.427143)--(1.652899,-0.346529)--(1.726471,-0.469932)--cycle;
\path[fill=blue!33!white,draw=blue!33!white,line width=.12pt] (-1.615372,-0.135163)--(-1.600000,-0.000000)--(-1.592296,-0.078414)--(-1.607593,-0.214330)--cycle;
\path[fill=blue!24!white,draw=blue!24!white,line width=.12pt] (1.607593,-0.214330)--(1.592296,-0.078414)--(1.600000,-0.000000)--(1.615372,-0.135163)--cycle;
\path[fill=blue!27!white,draw=blue!27!white,line width=.12pt] (-1.607593,0.214330)--(-1.652899,0.346529)--(-1.660896,0.265131)--(-1.615372,0.135163)--cycle;
\path[fill=blue!19!white,draw=blue!19!white,line width=.12pt] (1.615372,0.135163)--(1.660896,0.265131)--(1.652899,0.346529)--(1.607593,0.214330)--cycle;
\path[fill=blue!14!white,draw=blue!14!white,line width=.12pt] (-2.506958,0.928840)--(-2.654149,0.904054)--(-2.693116,0.772707)--(-2.543764,0.804777)--cycle;
\path[fill=blue!17!white,draw=blue!17!white,line width=.12pt] (2.543764,0.804777)--(2.693116,0.772707)--(2.654149,0.904054)--(2.506958,0.928840)--cycle;
\path[fill=blue!17!white,draw=blue!17!white,line width=.12pt] (-3.050416,0.535130)--(-3.123988,0.418974)--(-3.139104,0.265131)--(-3.065176,0.384910)--cycle;
\path[fill=blue!25!white,draw=blue!25!white,line width=.12pt] (3.065176,0.384910)--(3.139104,0.265131)--(3.123988,0.418974)--(3.050416,0.535130)--cycle;
\path[fill=blue!33!white,draw=blue!33!white,line width=.12pt] (-2.951405,-0.635242)--(-2.830759,-0.715462)--(-2.789801,-0.853522)--(-2.908701,-0.779186)--cycle;
\path[fill=blue!39!white,draw=blue!39!white,line width=.12pt] (2.908701,-0.779186)--(2.789801,-0.853522)--(2.830759,-0.715462)--(2.951405,-0.635242)--cycle;
\path[fill=green!65!black!34!white,draw=green!65!black!34!white,line width=.12pt] (-1.799069,0.668826)--(-1.917969,0.766813)--(-1.946127,0.671897)--(-1.825482,0.579795)--cycle;
\path[fill=green!65!black!29!white,draw=green!65!black!29!white,line width=.12pt] (1.825482,0.579795)--(1.946127,0.671897)--(1.917969,0.766813)--(1.799069,0.668826)--cycle;
\path[fill=blue!41!white,draw=blue!41!white,line width=.12pt] (-2.053620,-0.844328)--(-1.917969,-0.766813)--(-1.871339,-0.859891)--(-2.003693,-0.943989)--cycle;
\path[fill=blue!37!white,draw=blue!37!white,line width=.12pt] (2.003693,-0.943989)--(1.871339,-0.859891)--(1.917969,-0.766813)--(2.053620,-0.844328)--cycle;
\path[fill=blue!41!white,draw=blue!41!white,line width=.12pt] (-2.200811,-0.898392)--(-2.053620,-0.844328)--(-2.003693,-0.943989)--(-2.147305,-1.005197)--cycle;
\path[fill=blue!38!white,draw=blue!38!white,line width=.12pt] (2.147305,-1.005197)--(2.003693,-0.943989)--(2.053620,-0.844328)--(2.200811,-0.898392)--cycle;
\path[fill=blue!14!white,draw=blue!14!white,line width=.12pt] (-2.353885,0.926929)--(-2.506958,0.928840)--(-2.543764,0.804777)--(-2.388443,0.810441)--cycle;
\path[fill=blue!15!white,draw=blue!15!white,line width=.12pt] (2.388443,0.810441)--(2.543764,0.804777)--(2.506958,0.928840)--(2.353885,0.926929)--cycle;
\path[fill=blue!41!white,draw=blue!41!white,line width=.12pt] (-1.917969,-0.766813)--(-1.799069,-0.668826)--(-1.755330,-0.756135)--(-1.871339,-0.859891)--cycle;
\path[fill=blue!35!white,draw=blue!35!white,line width=.12pt] (1.871339,-0.859891)--(1.755330,-0.756135)--(1.799069,-0.668826)--(1.917969,-0.766813)--cycle;
\path[fill=blue!40!white,draw=blue!40!white,line width=.12pt] (-2.353885,-0.926929)--(-2.200811,-0.898392)--(-2.147305,-1.005197)--(-2.296657,-1.041162)--cycle;
\path[fill=blue!39!white,draw=blue!39!white,line width=.12pt] (2.296657,-1.041162)--(2.147305,-1.005197)--(2.200811,-0.898392)--(2.353885,-0.926929)--cycle;
\path[fill=blue!36!white,draw=blue!36!white,line width=.12pt] (-1.652899,-0.346529)--(-1.607593,-0.214330)--(-1.584333,-0.292734)--(-1.628983,-0.427143)--cycle;
\path[fill=blue!27!white,draw=blue!27!white,line width=.12pt] (1.628983,-0.427143)--(1.584333,-0.292734)--(1.607593,-0.214330)--(1.652899,-0.346529)--cycle;
\path[fill=green!65!black!33!white,draw=green!65!black!33!white,line width=.12pt] (-1.917969,0.766813)--(-2.053620,0.844328)--(-2.083771,0.742699)--(-1.946127,0.671897)--cycle;
\path[fill=green!65!black!28!white,draw=green!65!black!28!white,line width=.12pt] (1.946127,0.671897)--(2.083771,0.742699)--(2.053620,0.844328)--(1.917969,0.766813)--cycle;
\path[fill=blue!23!white,draw=blue!23!white,line width=.12pt] (-3.200000,0.000000)--(-3.184628,-0.135163)--(-3.169293,-0.291237)--(-3.184591,-0.156827)--cycle;
\path[fill=blue!32!white,draw=blue!32!white,line width=.12pt] (3.184591,-0.156827)--(3.169293,-0.291237)--(3.184628,-0.135163)--(3.200000,-0.000000)--cycle;
\path[fill=green!65!black!30!white,draw=green!65!black!30!white,line width=.12pt] (-2.200811,0.898392)--(-2.353885,0.926929)--(-2.388443,0.810441)--(-2.233123,0.789480)--cycle;
\path[fill=green!65!black!29!white,draw=green!65!black!29!white,line width=.12pt] (2.233123,0.789480)--(2.388443,0.810441)--(2.353885,0.926929)--(2.200811,0.898392)--cycle;
\path[fill=green!65!black!31!white,draw=green!65!black!31!white,line width=.12pt] (-2.053620,0.844328)--(-2.200811,0.898392)--(-2.233123,0.789480)--(-2.083771,0.742699)--cycle;
\path[fill=green!65!black!28!white,draw=green!65!black!28!white,line width=.12pt] (2.083771,0.742699)--(2.233123,0.789480)--(2.200811,0.898392)--(2.053620,0.844328)--cycle;
\path[fill=blue!31!white,draw=blue!31!white,line width=.12pt] (-1.600000,0.000000)--(-1.615372,0.135163)--(-1.607593,0.055995)--(-1.592296,-0.078414)--cycle;
\path[fill=blue!22!white,draw=blue!22!white,line width=.12pt] (1.592296,-0.078414)--(1.607593,0.055995)--(1.615372,0.135163)--(1.600000,-0.000000)--cycle;
\path[fill=blue!25!white,draw=blue!25!white,line width=.12pt] (-1.652899,0.346529)--(-1.726471,0.469932)--(-1.734824,0.384910)--(-1.660896,0.265131)--cycle;
\path[fill=blue!17!white,draw=blue!17!white,line width=.12pt] (1.660896,0.265131)--(1.734824,0.384910)--(1.726471,0.469932)--(1.652899,0.346529)--cycle;
\path[fill=blue!40!white,draw=blue!40!white,line width=.12pt] (-1.799069,-0.668826)--(-1.701490,-0.554134)--(-1.660123,-0.636707)--(-1.755330,-0.756135)--cycle;
\path[fill=blue!33!white,draw=blue!33!white,line width=.12pt] (1.755330,-0.756135)--(1.660123,-0.636707)--(1.701490,-0.554134)--(1.799069,-0.668826)--cycle;
\path[fill=blue!16!white,draw=blue!16!white,line width=.12pt] (-2.951405,0.635242)--(-3.050416,0.535130)--(-3.065176,0.384910)--(-2.965685,0.489898)--cycle;
\path[fill=blue!23!white,draw=blue!23!white,line width=.12pt] (2.965685,0.489898)--(3.065176,0.384910)--(3.050416,0.535130)--(2.951405,0.635242)--cycle;
\path[fill=blue!39!white,draw=blue!39!white,line width=.12pt] (-2.506958,-0.928840)--(-2.353885,-0.926929)--(-2.296657,-1.041162)--(-2.446009,-1.050502)--cycle;
\path[fill=blue!40!white,draw=blue!40!white,line width=.12pt] (2.446009,-1.050502)--(2.296657,-1.041162)--(2.353885,-0.926929)--(2.506958,-0.928840)--cycle;
\path[fill=blue!31!white,draw=blue!31!white,line width=.12pt] (-3.050416,-0.535130)--(-2.951405,-0.635242)--(-2.908701,-0.779186)--(-3.006279,-0.683903)--cycle;
\path[fill=blue!38!white,draw=blue!38!white,line width=.12pt] (3.006279,-0.683903)--(2.908701,-0.779186)--(2.951405,-0.635242)--(3.050416,-0.535130)--cycle;
\path[fill=blue!34!white,draw=blue!34!white,line width=.12pt] (-1.607593,-0.214330)--(-1.592296,-0.078414)--(-1.569256,-0.156072)--(-1.584333,-0.292734)--cycle;
\path[fill=blue!25!white,draw=blue!25!white,line width=.12pt] (1.584333,-0.292734)--(1.569256,-0.156072)--(1.592296,-0.078414)--(1.607593,-0.214330)--cycle;
\path[fill=blue!38!white,draw=blue!38!white,line width=.12pt] (-1.701490,-0.554134)--(-1.628983,-0.427143)--(-1.589379,-0.506197)--(-1.660123,-0.636707)--cycle;
\path[fill=blue!31!white,draw=blue!31!white,line width=.12pt] (1.660123,-0.636707)--(1.589379,-0.506197)--(1.628983,-0.427143)--(1.701490,-0.554134)--cycle;
\path[fill=blue!37!white,draw=blue!37!white,line width=.12pt] (-2.654149,-0.904054)--(-2.506958,-0.928840)--(-2.446009,-1.050502)--(-2.589621,-1.032859)--cycle;
\path[fill=blue!40!white,draw=blue!40!white,line width=.12pt] (2.589621,-1.032859)--(2.446009,-1.050502)--(2.506958,-0.928840)--(2.654149,-0.904054)--cycle;
\path[fill=blue!21!white,draw=blue!21!white,line width=.12pt] (-3.184628,0.135163)--(-3.200000,0.000000)--(-3.184591,-0.156827)--(-3.169293,-0.020912)--cycle;
\path[fill=blue!30!white,draw=blue!30!white,line width=.12pt] (3.169293,-0.020912)--(3.184591,-0.156827)--(3.200000,-0.000000)--(3.184628,0.135163)--cycle;
\path[fill=blue!22!white,draw=blue!22!white,line width=.12pt] (-1.726471,0.469932)--(-1.825482,0.579795)--(-1.834315,0.489898)--(-1.734824,0.384910)--cycle;
\path[fill=blue!15!white,draw=blue!15!white,line width=.12pt] (1.734824,0.384910)--(1.834315,0.489898)--(1.825482,0.579795)--(1.726471,0.469932)--cycle;
\path[fill=blue!15!white,draw=blue!15!white,line width=.12pt] (-2.830759,0.715462)--(-2.951405,0.635242)--(-2.965685,0.489898)--(-2.844456,0.576059)--cycle;
\path[fill=blue!20!white,draw=blue!20!white,line width=.12pt] (2.844456,0.576059)--(2.965685,0.489898)--(2.951405,0.635242)--(2.830759,0.715462)--cycle;
\path[fill=blue!28!white,draw=blue!28!white,line width=.12pt] (-1.615372,0.135163)--(-1.660896,0.265131)--(-1.652899,0.183733)--(-1.607593,0.055995)--cycle;
\path[fill=blue!19!white,draw=blue!19!white,line width=.12pt] (1.607593,0.055995)--(1.652899,0.183733)--(1.660896,0.265131)--(1.615372,0.135163)--cycle;
\path[fill=blue!28!white,draw=blue!28!white,line width=.12pt] (-3.123988,-0.418974)--(-3.050416,-0.535130)--(-3.006279,-0.683903)--(-3.078787,-0.571335)--cycle;
\path[fill=blue!36!white,draw=blue!36!white,line width=.12pt] (3.078787,-0.571335)--(3.006279,-0.683903)--(3.050416,-0.535130)--(3.123988,-0.418974)--cycle;
\path[fill=green!65!black!35!white,draw=green!65!black!35!white,line width=.12pt] (-1.825482,0.579795)--(-1.946127,0.671897)--(-1.955544,0.576059)--(-1.834315,0.489898)--cycle;
\path[fill=green!65!black!29!white,draw=green!65!black!29!white,line width=.12pt] (1.834315,0.489898)--(1.955544,0.576059)--(1.946127,0.671897)--(1.825482,0.579795)--cycle;
\path[fill=blue!36!white,draw=blue!36!white,line width=.12pt] (-1.628983,-0.427143)--(-1.584333,-0.292734)--(-1.545814,-0.369621)--(-1.589379,-0.506197)--cycle;
\path[fill=blue!28!white,draw=blue!28!white,line width=.12pt] (1.589379,-0.506197)--(1.545814,-0.369621)--(1.584333,-0.292734)--(1.628983,-0.427143)--cycle;
\path[fill=blue!35!white,draw=blue!35!white,line width=.12pt] (-2.789801,-0.853522)--(-2.654149,-0.904054)--(-2.589621,-1.032859)--(-2.721975,-0.988910)--cycle;
\path[fill=blue!39!white,draw=blue!39!white,line width=.12pt] (2.721975,-0.988910)--(2.589621,-1.032859)--(2.654149,-0.904054)--(2.789801,-0.853522)--cycle;
\path[fill=blue!14!white,draw=blue!14!white,line width=.12pt] (-2.693116,0.772707)--(-2.830759,0.715462)--(-2.844456,0.576059)--(-2.706147,0.640083)--cycle;
\path[fill=blue!18!white,draw=blue!18!white,line width=.12pt] (2.706147,0.640083)--(2.844456,0.576059)--(2.830759,0.715462)--(2.693116,0.772707)--cycle;
\path[fill=blue!41!white,draw=blue!41!white,line width=.12pt] (-2.003693,-0.943989)--(-1.871339,-0.859891)--(-1.806687,-0.950236)--(-1.934468,-1.040724)--cycle;
\path[fill=blue!37!white,draw=blue!37!white,line width=.12pt] (1.934468,-1.040724)--(1.806687,-0.950236)--(1.871339,-0.859891)--(2.003693,-0.943989)--cycle;
\path[fill=blue!31!white,draw=blue!31!white,line width=.12pt] (-1.592296,-0.078414)--(-1.607593,0.055995)--(-1.584333,-0.022409)--(-1.569256,-0.156072)--cycle;
\path[fill=blue!22!white,draw=blue!22!white,line width=.12pt] (1.569256,-0.156072)--(1.584333,-0.022409)--(1.607593,0.055995)--(1.592296,-0.078414)--cycle;
\path[fill=blue!41!white,draw=blue!41!white,line width=.12pt] (-1.871339,-0.859891)--(-1.755330,-0.756135)--(-1.694686,-0.840879)--(-1.806687,-0.950236)--cycle;
\path[fill=blue!36!white,draw=blue!36!white,line width=.12pt] (1.806687,-0.950236)--(1.694686,-0.840879)--(1.755330,-0.756135)--(1.871339,-0.859891)--cycle;
\path[fill=blue!19!white,draw=blue!19!white,line width=.12pt] (-3.139104,0.265131)--(-3.184628,0.135163)--(-3.169293,-0.020912)--(-3.123988,0.111288)--cycle;
\path[fill=blue!27!white,draw=blue!27!white,line width=.12pt] (3.123988,0.111288)--(3.169293,-0.020912)--(3.184628,0.135163)--(3.139104,0.265131)--cycle;
\path[fill=blue!41!white,draw=blue!41!white,line width=.12pt] (-2.147305,-1.005197)--(-2.003693,-0.943989)--(-1.934468,-1.040724)--(-2.073119,-1.108865)--cycle;
\path[fill=blue!38!white,draw=blue!38!white,line width=.12pt] (2.073119,-1.108865)--(1.934468,-1.040724)--(2.003693,-0.943989)--(2.147305,-1.005197)--cycle;
\path[fill=blue!25!white,draw=blue!25!white,line width=.12pt] (-1.660896,0.265131)--(-1.734824,0.384910)--(-1.726471,0.299889)--(-1.652899,0.183733)--cycle;
\path[fill=blue!17!white,draw=blue!17!white,line width=.12pt] (1.652899,0.183733)--(1.726471,0.299889)--(1.734824,0.384910)--(1.660896,0.265131)--cycle;
\path[fill=green!65!black!33!white,draw=green!65!black!33!white,line width=.12pt] (-1.946127,0.671897)--(-2.083771,0.742699)--(-2.093853,0.640083)--(-1.955544,0.576059)--cycle;
\path[fill=green!65!black!29!white,draw=green!65!black!29!white,line width=.12pt] (1.955544,0.576059)--(2.093853,0.640083)--(2.083771,0.742699)--(1.946127,0.671897)--cycle;
\path[fill=blue!40!white,draw=blue!40!white,line width=.12pt] (-1.755330,-0.756135)--(-1.660123,-0.636707)--(-1.602769,-0.716855)--(-1.694686,-0.840879)--cycle;
\path[fill=blue!33!white,draw=blue!33!white,line width=.12pt] (1.694686,-0.840879)--(1.602769,-0.716855)--(1.660123,-0.636707)--(1.755330,-0.756135)--cycle;
\path[fill=blue!14!white,draw=blue!14!white,line width=.12pt] (-2.543764,0.804777)--(-2.693116,0.772707)--(-2.706147,0.640083)--(-2.556072,0.679508)--cycle;
\path[fill=blue!16!white,draw=blue!16!white,line width=.12pt] (2.556072,0.679508)--(2.706147,0.640083)--(2.693116,0.772707)--(2.543764,0.804777)--cycle;
\path[fill=blue!40!white,draw=blue!40!white,line width=.12pt] (-2.296657,-1.041162)--(-2.147305,-1.005197)--(-2.073119,-1.108865)--(-2.217311,-1.152040)--cycle;
\path[fill=blue!39!white,draw=blue!39!white,line width=.12pt] (2.217311,-1.152040)--(2.073119,-1.108865)--(2.147305,-1.005197)--(2.296657,-1.041162)--cycle;
\path[fill=blue!25!white,draw=blue!25!white,line width=.12pt] (-3.169293,-0.291237)--(-3.123988,-0.418974)--(-3.078787,-0.571335)--(-3.123436,-0.445808)--cycle;
\path[fill=blue!34!white,draw=blue!34!white,line width=.12pt] (3.123436,-0.445808)--(3.078787,-0.571335)--(3.123988,-0.418974)--(3.169293,-0.291237)--cycle;
\path[fill=green!65!black!31!white,draw=green!65!black!31!white,line width=.12pt] (-2.083771,0.742699)--(-2.233123,0.789480)--(-2.243928,0.679508)--(-2.093853,0.640083)--cycle;
\path[fill=green!65!black!29!white,draw=green!65!black!29!white,line width=.12pt] (2.093853,0.640083)--(2.243928,0.679508)--(2.233123,0.789480)--(2.083771,0.742699)--cycle;
\path[fill=blue!14!white,draw=blue!14!white,line width=.12pt] (-2.388443,0.810441)--(-2.543764,0.804777)--(-2.556072,0.679508)--(-2.400000,0.692820)--cycle;
\path[fill=blue!15!white,draw=blue!15!white,line width=.12pt] (2.400000,0.692820)--(2.556072,0.679508)--(2.543764,0.804777)--(2.388443,0.810441)--cycle;
\path[fill=green!65!black!30!white,draw=green!65!black!30!white,line width=.12pt] (-2.233123,0.789480)--(-2.388443,0.810441)--(-2.400000,0.692820)--(-2.243928,0.679508)--cycle;
\path[fill=green!65!black!29!white,draw=green!65!black!29!white,line width=.12pt] (2.243928,0.679508)--(2.400000,0.692820)--(2.388443,0.810441)--(2.233123,0.789480)--cycle;
\path[fill=blue!34!white,draw=blue!34!white,line width=.12pt] (-1.584333,-0.292734)--(-1.569256,-0.156072)--(-1.531105,-0.232228)--(-1.545814,-0.369621)--cycle;
\path[fill=blue!26!white,draw=blue!26!white,line width=.12pt] (1.545814,-0.369621)--(1.531105,-0.232228)--(1.569256,-0.156072)--(1.584333,-0.292734)--cycle;
\path[fill=blue!33!white,draw=blue!33!white,line width=.12pt] (-2.908701,-0.779186)--(-2.789801,-0.853522)--(-2.721975,-0.988910)--(-2.837984,-0.920344)--cycle;
\path[fill=blue!38!white,draw=blue!38!white,line width=.12pt] (2.837984,-0.920344)--(2.721975,-0.988910)--(2.789801,-0.853522)--(2.908701,-0.779186)--cycle;
\path[fill=blue!39!white,draw=blue!39!white,line width=.12pt] (-1.660123,-0.636707)--(-1.589379,-0.506197)--(-1.534468,-0.582930)--(-1.602769,-0.716855)--cycle;
\path[fill=blue!31!white,draw=blue!31!white,line width=.12pt] (1.602769,-0.716855)--(1.534468,-0.582930)--(1.589379,-0.506197)--(1.660123,-0.636707)--cycle;
\path[fill=blue!17!white,draw=blue!17!white,line width=.12pt] (-3.065176,0.384910)--(-3.139104,0.265131)--(-3.123988,0.111288)--(-3.050416,0.234690)--cycle;
\path[fill=blue!25!white,draw=blue!25!white,line width=.12pt] (3.050416,0.234690)--(3.123988,0.111288)--(3.139104,0.265131)--(3.065176,0.384910)--cycle;
\path[fill=blue!29!white,draw=blue!29!white,line width=.12pt] (-1.607593,0.055995)--(-1.652899,0.183733)--(-1.628983,0.103118)--(-1.584333,-0.022409)--cycle;
\path[fill=blue!20!white,draw=blue!20!white,line width=.12pt] (1.584333,-0.022409)--(1.628983,0.103118)--(1.652899,0.183733)--(1.607593,0.055995)--cycle;
\path[fill=blue!39!white,draw=blue!39!white,line width=.12pt] (-2.446009,-1.050502)--(-2.296657,-1.041162)--(-2.217311,-1.152040)--(-2.361503,-1.168591)--cycle;
\path[fill=blue!40!white,draw=blue!40!white,line width=.12pt] (2.361503,-1.168591)--(2.217311,-1.152040)--(2.296657,-1.041162)--(2.446009,-1.050502)--cycle;
\path[fill=blue!23!white,draw=blue!23!white,line width=.12pt] (-1.734824,0.384910)--(-1.834315,0.489898)--(-1.825482,0.400001)--(-1.726471,0.299889)--cycle;
\path[fill=blue!16!white,draw=blue!16!white,line width=.12pt] (1.726471,0.299889)--(1.825482,0.400001)--(1.834315,0.489898)--(1.734824,0.384910)--cycle;
\path[fill=blue!37!white,draw=blue!37!white,line width=.12pt] (-1.589379,-0.506197)--(-1.545814,-0.369621)--(-1.492409,-0.444251)--(-1.534468,-0.582930)--cycle;
\path[fill=blue!29!white,draw=blue!29!white,line width=.12pt] (1.534468,-0.582930)--(1.492409,-0.444251)--(1.545814,-0.369621)--(1.589379,-0.506197)--cycle;
\path[fill=blue!23!white,draw=blue!23!white,line width=.12pt] (-3.184591,-0.156827)--(-3.169293,-0.291237)--(-3.123436,-0.445808)--(-3.138513,-0.312145)--cycle;
\path[fill=blue!32!white,draw=blue!32!white,line width=.12pt] (3.138513,-0.312145)--(3.123436,-0.445808)--(3.169293,-0.291237)--(3.184591,-0.156827)--cycle;
\path[fill=blue!32!white,draw=blue!32!white,line width=.12pt] (-1.569256,-0.156072)--(-1.584333,-0.022409)--(-1.545814,-0.099296)--(-1.531105,-0.232228)--cycle;
\path[fill=blue!23!white,draw=blue!23!white,line width=.12pt] (1.531105,-0.232228)--(1.545814,-0.099296)--(1.584333,-0.022409)--(1.569256,-0.156072)--cycle;
\path[fill=blue!15!white,draw=blue!15!white,line width=.12pt] (-2.965685,0.489898)--(-3.065176,0.384910)--(-3.050416,0.234690)--(-2.951405,0.344554)--cycle;
\path[fill=blue!22!white,draw=blue!22!white,line width=.12pt] (2.951405,0.344554)--(3.050416,0.234690)--(3.065176,0.384910)--(2.965685,0.489898)--cycle;
\path[fill=blue!30!white,draw=blue!30!white,line width=.12pt] (-3.006279,-0.683903)--(-2.908701,-0.779186)--(-2.837984,-0.920344)--(-2.933190,-0.829797)--cycle;
\path[fill=blue!37!white,draw=blue!37!white,line width=.12pt] (2.933190,-0.829797)--(2.837984,-0.920344)--(2.908701,-0.779186)--(3.006279,-0.683903)--cycle;
\path[fill=blue!37!white,draw=blue!37!white,line width=.12pt] (-2.589621,-1.032859)--(-2.446009,-1.050502)--(-2.361503,-1.168591)--(-2.500154,-1.157881)--cycle;
\path[fill=blue!40!white,draw=blue!40!white,line width=.12pt] (2.500154,-1.157881)--(2.361503,-1.168591)--(2.446009,-1.050502)--(2.589621,-1.032859)--cycle;
\path[fill=green!65!black!35!white,draw=green!65!black!35!white,line width=.12pt] (-1.834315,0.489898)--(-1.955544,0.576059)--(-1.946127,0.480221)--(-1.825482,0.400001)--cycle;
\path[fill=green!65!black!30!white,draw=green!65!black!30!white,line width=.12pt] (1.825482,0.400001)--(1.946127,0.480221)--(1.955544,0.576059)--(1.834315,0.489898)--cycle;
\path[fill=blue!41!white,draw=blue!41!white,line width=.12pt] (-1.806687,-0.950236)--(-1.694686,-0.840879)--(-1.617721,-0.922243)--(-1.724636,-1.036978)--cycle;
\path[fill=blue!36!white,draw=blue!36!white,line width=.12pt] (1.724636,-1.036978)--(1.617721,-0.922243)--(1.694686,-0.840879)--(1.806687,-0.950236)--cycle;
\path[fill=blue!41!white,draw=blue!41!white,line width=.12pt] (-1.934468,-1.040724)--(-1.806687,-0.950236)--(-1.724636,-1.036978)--(-1.846614,-1.133600)--cycle;
\path[fill=blue!38!white,draw=blue!38!white,line width=.12pt] (1.846614,-1.133600)--(1.724636,-1.036978)--(1.806687,-0.950236)--(1.934468,-1.040724)--cycle;
\path[fill=blue!26!white,draw=blue!26!white,line width=.12pt] (-1.652899,0.183733)--(-1.726471,0.299889)--(-1.701490,0.215687)--(-1.628983,0.103118)--cycle;
\path[fill=blue!18!white,draw=blue!18!white,line width=.12pt] (1.628983,0.103118)--(1.701490,0.215687)--(1.726471,0.299889)--(1.652899,0.183733)--cycle;
\path[fill=blue!40!white,draw=blue!40!white,line width=.12pt] (-1.694686,-0.840879)--(-1.602769,-0.716855)--(-1.529978,-0.793806)--(-1.617721,-0.922243)--cycle;
\path[fill=blue!34!white,draw=blue!34!white,line width=.12pt] (1.617721,-0.922243)--(1.529978,-0.793806)--(1.602769,-0.716855)--(1.694686,-0.840879)--cycle;
\path[fill=blue!41!white,draw=blue!41!white,line width=.12pt] (-2.073119,-1.108865)--(-1.934468,-1.040724)--(-1.846614,-1.133600)--(-1.978968,-1.208398)--cycle;
\path[fill=blue!39!white,draw=blue!39!white,line width=.12pt] (1.978968,-1.208398)--(1.846614,-1.133600)--(1.934468,-1.040724)--(2.073119,-1.108865)--cycle;
\path[fill=blue!35!white,draw=blue!35!white,line width=.12pt] (-1.545814,-0.369621)--(-1.531105,-0.232228)--(-1.478207,-0.306147)--(-1.492409,-0.444251)--cycle;
\path[fill=blue!26!white,draw=blue!26!white,line width=.12pt] (1.492409,-0.444251)--(1.478207,-0.306147)--(1.531105,-0.232228)--(1.545814,-0.369621)--cycle;
\path[fill=blue!14!white,draw=blue!14!white,line width=.12pt] (-2.844456,0.576059)--(-2.965685,0.489898)--(-2.951405,0.344554)--(-2.830759,0.436656)--cycle;
\path[fill=blue!20!white,draw=blue!20!white,line width=.12pt] (2.830759,0.436656)--(2.951405,0.344554)--(2.965685,0.489898)--(2.844456,0.576059)--cycle;
\path[fill=green!65!black!33!white,draw=green!65!black!33!white,line width=.12pt] (-1.955544,0.576059)--(-2.093853,0.640083)--(-2.083771,0.537466)--(-1.946127,0.480221)--cycle;
\path[fill=green!65!black!29!white,draw=green!65!black!29!white,line width=.12pt] (1.946127,0.480221)--(2.083771,0.537466)--(2.093853,0.640083)--(1.955544,0.576059)--cycle;
\path[fill=blue!39!white,draw=blue!39!white,line width=.12pt] (-1.602769,-0.716855)--(-1.534468,-0.582930)--(-1.464780,-0.656601)--(-1.529978,-0.793806)--cycle;
\path[fill=blue!32!white,draw=blue!32!white,line width=.12pt] (1.529978,-0.793806)--(1.464780,-0.656601)--(1.534468,-0.582930)--(1.602769,-0.716855)--cycle;
\path[fill=blue!20!white,draw=blue!20!white,line width=.12pt] (-3.169293,-0.020912)--(-3.184591,-0.156827)--(-3.138513,-0.312145)--(-3.123436,-0.175483)--cycle;
\path[fill=blue!29!white,draw=blue!29!white,line width=.12pt] (3.123436,-0.175483)--(3.138513,-0.312145)--(3.184591,-0.156827)--(3.169293,-0.020912)--cycle;
\path[fill=blue!35!white,draw=blue!35!white,line width=.12pt] (-2.721975,-0.988910)--(-2.589621,-1.032859)--(-2.500154,-1.157881)--(-2.627935,-1.120322)--cycle;
\path[fill=blue!39!white,draw=blue!39!white,line width=.12pt] (2.627935,-1.120322)--(2.500154,-1.157881)--(2.589621,-1.032859)--(2.721975,-0.988910)--cycle;
\path[fill=blue!29!white,draw=blue!29!white,line width=.12pt] (-1.584333,-0.022409)--(-1.628983,0.103118)--(-1.589379,0.024064)--(-1.545814,-0.099296)--cycle;
\path[fill=blue!21!white,draw=blue!21!white,line width=.12pt] (1.545814,-0.099296)--(1.589379,0.024064)--(1.628983,0.103118)--(1.584333,-0.022409)--cycle;
\path[fill=blue!40!white,draw=blue!40!white,line width=.12pt] (-2.217311,-1.152040)--(-2.073119,-1.108865)--(-1.978968,-1.208398)--(-2.116611,-1.258496)--cycle;
\path[fill=blue!39!white,draw=blue!39!white,line width=.12pt] (2.116611,-1.258496)--(1.978968,-1.208398)--(2.073119,-1.108865)--(2.217311,-1.152040)--cycle;
\path[fill=blue!28!white,draw=blue!28!white,line width=.12pt] (-3.078787,-0.571335)--(-3.006279,-0.683903)--(-2.933190,-0.829797)--(-3.003935,-0.720748)--cycle;
\path[fill=blue!35!white,draw=blue!35!white,line width=.12pt] (3.003935,-0.720748)--(2.933190,-0.829797)--(3.006279,-0.683903)--(3.078787,-0.571335)--cycle;
\path[fill=blue!14!white,draw=blue!14!white,line width=.12pt] (-2.706147,0.640083)--(-2.844456,0.576059)--(-2.830759,0.436656)--(-2.693116,0.507458)--cycle;
\path[fill=blue!18!white,draw=blue!18!white,line width=.12pt] (2.693116,0.507458)--(2.830759,0.436656)--(2.844456,0.576059)--(2.706147,0.640083)--cycle;
\path[fill=green!65!black!31!white,draw=green!65!black!31!white,line width=.12pt] (-2.093853,0.640083)--(-2.243928,0.679508)--(-2.233123,0.569536)--(-2.083771,0.537466)--cycle;
\path[fill=green!65!black!29!white,draw=green!65!black!29!white,line width=.12pt] (2.083771,0.537466)--(2.233123,0.569536)--(2.243928,0.679508)--(2.093853,0.640083)--cycle;
\path[fill=blue!23!white,draw=blue!23!white,line width=.12pt] (-1.726471,0.299889)--(-1.825482,0.400001)--(-1.799069,0.310969)--(-1.701490,0.215687)--cycle;
\path[fill=blue!16!white,draw=blue!16!white,line width=.12pt] (1.701490,0.215687)--(1.799069,0.310969)--(1.825482,0.400001)--(1.726471,0.299889)--cycle;
\path[fill=blue!37!white,draw=blue!37!white,line width=.12pt] (-1.534468,-0.582930)--(-1.492409,-0.444251)--(-1.424631,-0.515903)--(-1.464780,-0.656601)--cycle;
\path[fill=blue!30!white,draw=blue!30!white,line width=.12pt] (1.464780,-0.656601)--(1.424631,-0.515903)--(1.492409,-0.444251)--(1.534468,-0.582930)--cycle;
\path[fill=blue!14!white,draw=blue!14!white,line width=.12pt] (-2.556072,0.679508)--(-2.706147,0.640083)--(-2.693116,0.507458)--(-2.543764,0.554239)--cycle;
\path[fill=blue!16!white,draw=blue!16!white,line width=.12pt] (2.543764,0.554239)--(2.693116,0.507458)--(2.706147,0.640083)--(2.556072,0.679508)--cycle;
\path[fill=green!65!black!30!white,draw=green!65!black!30!white,line width=.12pt] (-2.243928,0.679508)--(-2.400000,0.692820)--(-2.388443,0.575200)--(-2.233123,0.569536)--cycle;
\path[fill=green!65!black!29!white,draw=green!65!black!29!white,line width=.12pt] (2.233123,0.569536)--(2.388443,0.575200)--(2.400000,0.692820)--(2.243928,0.679508)--cycle;
\path[fill=blue!14!white,draw=blue!14!white,line width=.12pt] (-2.400000,0.692820)--(-2.556072,0.679508)--(-2.543764,0.554239)--(-2.388443,0.575200)--cycle;
\path[fill=blue!15!white,draw=blue!15!white,line width=.12pt] (2.388443,0.575200)--(2.543764,0.554239)--(2.556072,0.679508)--(2.400000,0.692820)--cycle;
\path[fill=blue!32!white,draw=blue!32!white,line width=.12pt] (-1.531105,-0.232228)--(-1.545814,-0.099296)--(-1.492409,-0.173925)--(-1.478207,-0.306147)--cycle;
\path[fill=blue!24!white,draw=blue!24!white,line width=.12pt] (1.478207,-0.306147)--(1.492409,-0.173925)--(1.545814,-0.099296)--(1.531105,-0.232228)--cycle;
\path[fill=blue!18!white,draw=blue!18!white,line width=.12pt] (-3.123988,0.111288)--(-3.169293,-0.020912)--(-3.123436,-0.175483)--(-3.078787,-0.041074)--cycle;
\path[fill=blue!27!white,draw=blue!27!white,line width=.12pt] (3.078787,-0.041074)--(3.123436,-0.175483)--(3.169293,-0.020912)--(3.123988,0.111288)--cycle;
\path[fill=blue!39!white,draw=blue!39!white,line width=.12pt] (-2.361503,-1.168591)--(-2.217311,-1.152040)--(-2.116611,-1.258496)--(-2.254254,-1.281970)--cycle;
\path[fill=blue!40!white,draw=blue!40!white,line width=.12pt] (2.254254,-1.281970)--(2.116611,-1.258496)--(2.217311,-1.152040)--(2.361503,-1.168591)--cycle;
\path[fill=blue!41!white,draw=blue!41!white,line width=.12pt] (-1.724636,-1.036978)--(-1.617721,-0.922243)--(-1.525177,-0.999443)--(-1.625975,-1.119280)--cycle;
\path[fill=blue!37!white,draw=blue!37!white,line width=.12pt] (1.625975,-1.119280)--(1.525177,-0.999443)--(1.617721,-0.922243)--(1.724636,-1.036978)--cycle;
\path[fill=blue!33!white,draw=blue!33!white,line width=.12pt] (-2.837984,-0.920344)--(-2.721975,-0.988910)--(-2.627935,-1.120322)--(-2.739936,-1.057357)--cycle;
\path[fill=blue!38!white,draw=blue!38!white,line width=.12pt] (2.739936,-1.057357)--(2.627935,-1.120322)--(2.721975,-0.988910)--(2.837984,-0.920344)--cycle;
\path[fill=blue!41!white,draw=blue!41!white,line width=.12pt] (-1.617721,-0.922243)--(-1.529978,-0.793806)--(-1.442454,-0.866819)--(-1.525177,-0.999443)--cycle;
\path[fill=blue!35!white,draw=blue!35!white,line width=.12pt] (1.525177,-0.999443)--(1.442454,-0.866819)--(1.529978,-0.793806)--(1.617721,-0.922243)--cycle;
\path[fill=blue!26!white,draw=blue!26!white,line width=.12pt] (-1.628983,0.103118)--(-1.701490,0.215687)--(-1.660123,0.133114)--(-1.589379,0.024064)--cycle;
\path[fill=blue!19!white,draw=blue!19!white,line width=.12pt] (1.589379,0.024064)--(1.660123,0.133114)--(1.701490,0.215687)--(1.628983,0.103118)--cycle;
\path[fill=green!65!black!36!white,draw=green!65!black!36!white,line width=.12pt] (-1.825482,0.400001)--(-1.946127,0.480221)--(-1.917969,0.385305)--(-1.799069,0.310969)--cycle;
\path[fill=green!65!black!30!white,draw=green!65!black!30!white,line width=.12pt] (1.799069,0.310969)--(1.917969,0.385305)--(1.946127,0.480221)--(1.825482,0.400001)--cycle;
\path[fill=blue!25!white,draw=blue!25!white,line width=.12pt] (-3.123436,-0.445808)--(-3.078787,-0.571335)--(-3.003935,-0.720748)--(-3.047499,-0.597387)--cycle;
\path[fill=blue!33!white,draw=blue!33!white,line width=.12pt] (3.047499,-0.597387)--(3.003935,-0.720748)--(3.078787,-0.571335)--(3.123436,-0.445808)--cycle;
\path[fill=blue!42!white,draw=blue!42!white,line width=.12pt] (-1.846614,-1.133600)--(-1.724636,-1.036978)--(-1.625975,-1.119280)--(-1.740975,-1.221724)--cycle;
\path[fill=blue!38!white,draw=blue!38!white,line width=.12pt] (1.740975,-1.221724)--(1.625975,-1.119280)--(1.724636,-1.036978)--(1.846614,-1.133600)--cycle;
\path[fill=blue!35!white,draw=blue!35!white,line width=.12pt] (-1.492409,-0.444251)--(-1.478207,-0.306147)--(-1.411074,-0.377117)--(-1.424631,-0.515903)--cycle;
\path[fill=blue!27!white,draw=blue!27!white,line width=.12pt] (1.424631,-0.515903)--(1.411074,-0.377117)--(1.478207,-0.306147)--(1.492409,-0.444251)--cycle;
\path[fill=blue!39!white,draw=blue!39!white,line width=.12pt] (-1.529978,-0.793806)--(-1.464780,-0.656601)--(-1.380985,-0.726503)--(-1.442454,-0.866819)--cycle;
\path[fill=blue!33!white,draw=blue!33!white,line width=.12pt] (1.442454,-0.866819)--(1.380985,-0.726503)--(1.464780,-0.656601)--(1.529978,-0.793806)--cycle;
\path[fill=blue!16!white,draw=blue!16!white,line width=.12pt] (-3.050416,0.234690)--(-3.123988,0.111288)--(-3.078787,-0.041074)--(-3.006279,0.085917)--cycle;
\path[fill=blue!24!white,draw=blue!24!white,line width=.12pt] (3.006279,0.085917)--(3.078787,-0.041074)--(3.123988,0.111288)--(3.050416,0.234690)--cycle;
\path[fill=blue!41!white,draw=blue!41!white,line width=.12pt] (-1.978968,-1.208398)--(-1.846614,-1.133600)--(-1.740975,-1.221724)--(-1.865758,-1.302838)--cycle;
\path[fill=blue!39!white,draw=blue!39!white,line width=.12pt] (1.865758,-1.302838)--(1.740975,-1.221724)--(1.846614,-1.133600)--(1.978968,-1.208398)--cycle;
\path[fill=blue!37!white,draw=blue!37!white,line width=.12pt] (-2.500154,-1.157881)--(-2.361503,-1.168591)--(-2.254254,-1.281970)--(-2.386608,-1.277917)--cycle;
\path[fill=blue!39!white,draw=blue!39!white,line width=.12pt] (2.386608,-1.277917)--(2.254254,-1.281970)--(2.361503,-1.168591)--(2.500154,-1.157881)--cycle;
\path[fill=blue!30!white,draw=blue!30!white,line width=.12pt] (-1.545814,-0.099296)--(-1.589379,0.024064)--(-1.534468,-0.052668)--(-1.492409,-0.173925)--cycle;
\path[fill=blue!21!white,draw=blue!21!white,line width=.12pt] (1.492409,-0.173925)--(1.534468,-0.052668)--(1.589379,0.024064)--(1.545814,-0.099296)--cycle;
\path[fill=green!65!black!33!white,draw=green!65!black!33!white,line width=.12pt] (-1.946127,0.480221)--(-2.083771,0.537466)--(-2.053620,0.435837)--(-1.917969,0.385305)--cycle;
\path[fill=green!65!black!29!white,draw=green!65!black!29!white,line width=.12pt] (1.917969,0.385305)--(2.053620,0.435837)--(2.083771,0.537466)--(1.946127,0.480221)--cycle;
\path[fill=blue!38!white,draw=blue!38!white,line width=.12pt] (-1.464780,-0.656601)--(-1.424631,-0.515903)--(-1.343133,-0.583889)--(-1.380985,-0.726503)--cycle;
\path[fill=blue!30!white,draw=blue!30!white,line width=.12pt] (1.380985,-0.726503)--(1.343133,-0.583889)--(1.424631,-0.515903)--(1.464780,-0.656601)--cycle;
\path[fill=blue!30!white,draw=blue!30!white,line width=.12pt] (-2.933190,-0.829797)--(-2.837984,-0.920344)--(-2.739936,-1.057357)--(-2.831853,-0.971406)--cycle;
\path[fill=blue!37!white,draw=blue!37!white,line width=.12pt] (2.831853,-0.971406)--(2.739936,-1.057357)--(2.837984,-0.920344)--(2.933190,-0.829797)--cycle;
\path[fill=blue!24!white,draw=blue!24!white,line width=.12pt] (-1.701490,0.215687)--(-1.799069,0.310969)--(-1.755330,0.223661)--(-1.660123,0.133114)--cycle;
\path[fill=blue!17!white,draw=blue!17!white,line width=.12pt] (1.660123,0.133114)--(1.755330,0.223661)--(1.799069,0.310969)--(1.701490,0.215687)--cycle;
\path[fill=blue!22!white,draw=blue!22!white,line width=.12pt] (-3.138513,-0.312145)--(-3.123436,-0.445808)--(-3.047499,-0.597387)--(-3.062209,-0.464455)--cycle;
\path[fill=blue!31!white,draw=blue!31!white,line width=.12pt] (3.062209,-0.464455)--(3.047499,-0.597387)--(3.123436,-0.445808)--(3.138513,-0.312145)--cycle;
\path[fill=blue!33!white,draw=blue!33!white,line width=.12pt] (-1.478207,-0.306147)--(-1.492409,-0.173925)--(-1.424631,-0.245578)--(-1.411074,-0.377117)--cycle;
\path[fill=blue!25!white,draw=blue!25!white,line width=.12pt] (1.411074,-0.377117)--(1.424631,-0.245578)--(1.492409,-0.173925)--(1.478207,-0.306147)--cycle;
\path[fill=blue!15!white,draw=blue!15!white,line width=.12pt] (-2.951405,0.344554)--(-3.050416,0.234690)--(-3.006279,0.085917)--(-2.908701,0.200610)--cycle;
\path[fill=blue!22!white,draw=blue!22!white,line width=.12pt] (2.908701,0.200610)--(3.006279,0.085917)--(3.050416,0.234690)--(2.951405,0.344554)--cycle;
\path[fill=blue!40!white,draw=blue!40!white,line width=.12pt] (-2.116611,-1.258496)--(-1.978968,-1.208398)--(-1.865758,-1.302838)--(-1.995527,-1.359505)--cycle;
\path[fill=blue!40!white,draw=blue!40!white,line width=.12pt] (1.995527,-1.359505)--(1.865758,-1.302838)--(1.978968,-1.208398)--(2.116611,-1.258496)--cycle;
\path[fill=green!65!black!32!white,draw=green!65!black!32!white,line width=.12pt] (-2.083771,0.537466)--(-2.233123,0.569536)--(-2.200811,0.460624)--(-2.053620,0.435837)--cycle;
\path[fill=green!65!black!29!white,draw=green!65!black!29!white,line width=.12pt] (2.053620,0.435837)--(2.200811,0.460624)--(2.233123,0.569536)--(2.083771,0.537466)--cycle;
\path[fill=blue!41!white,draw=blue!41!white,line width=.12pt] (-1.525177,-0.999443)--(-1.442454,-0.866819)--(-1.341037,-0.935191)--(-1.417944,-1.071736)--cycle;
\path[fill=blue!35!white,draw=blue!35!white,line width=.12pt] (1.417944,-1.071736)--(1.341037,-0.935191)--(1.442454,-0.866819)--(1.525177,-0.999443)--cycle;
\path[fill=blue!42!white,draw=blue!42!white,line width=.12pt] (-1.625975,-1.119280)--(-1.525177,-0.999443)--(-1.417944,-1.071736)--(-1.511656,-1.196351)--cycle;
\path[fill=blue!37!white,draw=blue!37!white,line width=.12pt] (1.511656,-1.196351)--(1.417944,-1.071736)--(1.525177,-0.999443)--(1.625975,-1.119280)--cycle;
\path[fill=blue!35!white,draw=blue!35!white,line width=.12pt] (-2.627935,-1.120322)--(-2.500154,-1.157881)--(-2.386608,-1.277917)--(-2.508586,-1.246493)--cycle;
\path[fill=blue!39!white,draw=blue!39!white,line width=.12pt] (2.508586,-1.246493)--(2.386608,-1.277917)--(2.500154,-1.157881)--(2.627935,-1.120322)--cycle;
\path[fill=blue!36!white,draw=blue!36!white,line width=.12pt] (-1.424631,-0.515903)--(-1.411074,-0.377117)--(-1.330351,-0.444456)--(-1.343133,-0.583889)--cycle;
\path[fill=blue!28!white,draw=blue!28!white,line width=.12pt] (1.343133,-0.583889)--(1.330351,-0.444456)--(1.411074,-0.377117)--(1.424631,-0.515903)--cycle;
\path[fill=blue!14!white,draw=blue!14!white,line width=.12pt] (-2.830759,0.436656)--(-2.951405,0.344554)--(-2.908701,0.200610)--(-2.789801,0.298596)--cycle;
\path[fill=blue!19!white,draw=blue!19!white,line width=.12pt] (2.789801,0.298596)--(2.908701,0.200610)--(2.951405,0.344554)--(2.830759,0.436656)--cycle;
\path[fill=blue!27!white,draw=blue!27!white,line width=.12pt] (-1.589379,0.024064)--(-1.660123,0.133114)--(-1.602769,0.052966)--(-1.534468,-0.052668)--cycle;
\path[fill=blue!19!white,draw=blue!19!white,line width=.12pt] (1.534468,-0.052668)--(1.602769,0.052966)--(1.660123,0.133114)--(1.589379,0.024064)--cycle;
\path[fill=blue!40!white,draw=blue!40!white,line width=.12pt] (-1.442454,-0.866819)--(-1.380985,-0.726503)--(-1.283890,-0.791962)--(-1.341037,-0.935191)--cycle;
\path[fill=blue!33!white,draw=blue!33!white,line width=.12pt] (1.341037,-0.935191)--(1.283890,-0.791962)--(1.380985,-0.726503)--(1.442454,-0.866819)--cycle;
\path[fill=green!65!black!30!white,draw=green!65!black!30!white,line width=.12pt] (-2.233123,0.569536)--(-2.388443,0.575200)--(-2.353885,0.458712)--(-2.200811,0.460624)--cycle;
\path[fill=green!65!black!29!white,draw=green!65!black!29!white,line width=.12pt] (2.200811,0.460624)--(2.353885,0.458712)--(2.388443,0.575200)--(2.233123,0.569536)--cycle;
\path[fill=blue!42!white,draw=blue!42!white,line width=.12pt] (-1.740975,-1.221724)--(-1.625975,-1.119280)--(-1.511656,-1.196351)--(-1.618570,-1.304246)--cycle;
\path[fill=blue!38!white,draw=blue!38!white,line width=.12pt] (1.618570,-1.304246)--(1.511656,-1.196351)--(1.625975,-1.119280)--(1.740975,-1.221724)--cycle;
\path[fill=blue!13!white,draw=blue!13!white,line width=.12pt] (-2.693116,0.507458)--(-2.830759,0.436656)--(-2.789801,0.298596)--(-2.654149,0.376111)--cycle;
\path[fill=blue!18!white,draw=blue!18!white,line width=.12pt] (2.654149,0.376111)--(2.789801,0.298596)--(2.830759,0.436656)--(2.693116,0.507458)--cycle;
\path[fill=blue!39!white,draw=blue!39!white,line width=.12pt] (-2.254254,-1.281970)--(-2.116611,-1.258496)--(-1.995527,-1.359505)--(-2.125296,-1.389547)--cycle;
\path[fill=blue!40!white,draw=blue!40!white,line width=.12pt] (2.125296,-1.389547)--(1.995527,-1.359505)--(2.116611,-1.258496)--(2.254254,-1.281970)--cycle;
\path[fill=green!65!black!36!white,draw=green!65!black!36!white,line width=.12pt] (-1.799069,0.310969)--(-1.917969,0.385305)--(-1.871339,0.292227)--(-1.755330,0.223661)--cycle;
\path[fill=green!65!black!31!white,draw=green!65!black!31!white,line width=.12pt] (1.755330,0.223661)--(1.871339,0.292227)--(1.917969,0.385305)--(1.799069,0.310969)--cycle;
\path[fill=blue!20!white,draw=blue!20!white,line width=.12pt] (-3.123436,-0.175483)--(-3.138513,-0.312145)--(-3.062209,-0.464455)--(-3.047499,-0.327062)--cycle;
\path[fill=blue!28!white,draw=blue!28!white,line width=.12pt] (3.047499,-0.327062)--(3.062209,-0.464455)--(3.138513,-0.312145)--(3.123436,-0.175483)--cycle;
\path[fill=blue!14!white,draw=blue!14!white,line width=.12pt] (-2.388443,0.575200)--(-2.543764,0.554239)--(-2.506958,0.430175)--(-2.353885,0.458712)--cycle;
\path[fill=blue!15!white,draw=blue!15!white,line width=.12pt] (2.353885,0.458712)--(2.506958,0.430175)--(2.543764,0.554239)--(2.388443,0.575200)--cycle;
\path[fill=blue!13!white,draw=blue!13!white,line width=.12pt] (-2.543764,0.554239)--(-2.693116,0.507458)--(-2.654149,0.376111)--(-2.506958,0.430175)--cycle;
\path[fill=blue!16!white,draw=blue!16!white,line width=.12pt] (2.506958,0.430175)--(2.654149,0.376111)--(2.693116,0.507458)--(2.543764,0.554239)--cycle;
\path[fill=blue!27!white,draw=blue!27!white,line width=.12pt] (-3.003935,-0.720748)--(-2.933190,-0.829797)--(-2.831853,-0.971406)--(-2.900154,-0.865772)--cycle;
\path[fill=blue!35!white,draw=blue!35!white,line width=.12pt] (2.900154,-0.865772)--(2.831853,-0.971406)--(2.933190,-0.829797)--(3.003935,-0.720748)--cycle;
\path[fill=blue!30!white,draw=blue!30!white,line width=.12pt] (-1.492409,-0.173925)--(-1.534468,-0.052668)--(-1.464780,-0.126340)--(-1.424631,-0.245578)--cycle;
\path[fill=blue!22!white,draw=blue!22!white,line width=.12pt] (1.424631,-0.245578)--(1.464780,-0.126340)--(1.534468,-0.052668)--(1.492409,-0.173925)--cycle;
\path[fill=blue!38!white,draw=blue!38!white,line width=.12pt] (-1.380985,-0.726503)--(-1.343133,-0.583889)--(-1.248699,-0.647553)--(-1.283890,-0.791962)--cycle;
\path[fill=blue!31!white,draw=blue!31!white,line width=.12pt] (1.283890,-0.791962)--(1.248699,-0.647553)--(1.343133,-0.583889)--(1.380985,-0.726503)--cycle;
\path[fill=blue!41!white,draw=blue!41!white,line width=.12pt] (-1.865758,-1.302838)--(-1.740975,-1.221724)--(-1.618570,-1.304246)--(-1.734580,-1.391274)--cycle;
\path[fill=blue!39!white,draw=blue!39!white,line width=.12pt] (1.734580,-1.391274)--(1.618570,-1.304246)--(1.740975,-1.221724)--(1.865758,-1.302838)--cycle;
\path[fill=blue!33!white,draw=blue!33!white,line width=.12pt] (-1.411074,-0.377117)--(-1.424631,-0.245578)--(-1.343133,-0.313564)--(-1.330351,-0.444456)--cycle;
\path[fill=blue!25!white,draw=blue!25!white,line width=.12pt] (1.330351,-0.444456)--(1.343133,-0.313564)--(1.424631,-0.245578)--(1.411074,-0.377117)--cycle;
\path[fill=blue!32!white,draw=blue!32!white,line width=.12pt] (-2.739936,-1.057357)--(-2.627935,-1.120322)--(-2.508586,-1.246493)--(-2.615501,-1.188905)--cycle;
\path[fill=blue!37!white,draw=blue!37!white,line width=.12pt] (2.615501,-1.188905)--(2.508586,-1.246493)--(2.627935,-1.120322)--(2.739936,-1.057357)--cycle;
\path[fill=blue!24!white,draw=blue!24!white,line width=.12pt] (-1.660123,0.133114)--(-1.755330,0.223661)--(-1.694686,0.138917)--(-1.602769,0.052966)--cycle;
\path[fill=blue!17!white,draw=blue!17!white,line width=.12pt] (1.602769,0.052966)--(1.694686,0.138917)--(1.755330,0.223661)--(1.660123,0.133114)--cycle;
\path[fill=green!65!black!34!white,draw=green!65!black!34!white,line width=.12pt] (-1.917969,0.385305)--(-2.053620,0.435837)--(-2.003693,0.336176)--(-1.871339,0.292227)--cycle;
\path[fill=green!65!black!30!white,draw=green!65!black!30!white,line width=.12pt] (1.871339,0.292227)--(2.003693,0.336176)--(2.053620,0.435837)--(1.917969,0.385305)--cycle;
\path[fill=blue!18!white,draw=blue!18!white,line width=.12pt] (-3.078787,-0.041074)--(-3.123436,-0.175483)--(-3.047499,-0.327062)--(-3.003935,-0.190486)--cycle;
\path[fill=blue!26!white,draw=blue!26!white,line width=.12pt] (3.003935,-0.190486)--(3.047499,-0.327062)--(3.123436,-0.175483)--(3.078787,-0.041074)--cycle;
\path[fill=blue!41!white,draw=blue!41!white,line width=.12pt] (-1.417944,-1.071736)--(-1.341037,-0.935191)--(-1.226706,-0.998263)--(-1.297056,-1.138426)--cycle;
\path[fill=blue!36!white,draw=blue!36!white,line width=.12pt] (1.297056,-1.138426)--(1.226706,-0.998263)--(1.341037,-0.935191)--(1.417944,-1.071736)--cycle;
\path[fill=blue!37!white,draw=blue!37!white,line width=.12pt] (-2.386608,-1.277917)--(-2.254254,-1.281970)--(-2.125296,-1.389547)--(-2.250079,-1.391810)--cycle;
\path[fill=blue!39!white,draw=blue!39!white,line width=.12pt] (2.250079,-1.391810)--(2.125296,-1.389547)--(2.254254,-1.281970)--(2.386608,-1.277917)--cycle;
\path[fill=blue!36!white,draw=blue!36!white,line width=.12pt] (-1.343133,-0.583889)--(-1.330351,-0.444456)--(-1.236817,-0.507515)--(-1.248699,-0.647553)--cycle;
\path[fill=blue!29!white,draw=blue!29!white,line width=.12pt] (1.248699,-0.647553)--(1.236817,-0.507515)--(1.330351,-0.444456)--(1.343133,-0.583889)--cycle;
\path[fill=blue!42!white,draw=blue!42!white,line width=.12pt] (-1.511656,-1.196351)--(-1.417944,-1.071736)--(-1.297056,-1.138426)--(-1.382778,-1.267448)--cycle;
\path[fill=blue!38!white,draw=blue!38!white,line width=.12pt] (1.382778,-1.267448)--(1.297056,-1.138426)--(1.417944,-1.071736)--(1.511656,-1.196351)--cycle;
\path[fill=blue!40!white,draw=blue!40!white,line width=.12pt] (-1.341037,-0.935191)--(-1.283890,-0.791962)--(-1.174431,-0.852346)--(-1.226706,-0.998263)--cycle;
\path[fill=blue!34!white,draw=blue!34!white,line width=.12pt] (1.226706,-0.998263)--(1.174431,-0.852346)--(1.283890,-0.791962)--(1.341037,-0.935191)--cycle;
\path[fill=blue!40!white,draw=blue!40!white,line width=.12pt] (-1.995527,-1.359505)--(-1.865758,-1.302838)--(-1.734580,-1.391274)--(-1.855225,-1.454092)--cycle;
\path[fill=blue!40!white,draw=blue!40!white,line width=.12pt] (1.855225,-1.454092)--(1.734580,-1.391274)--(1.865758,-1.302838)--(1.995527,-1.359505)--cycle;
\path[fill=blue!24!white,draw=blue!24!white,line width=.12pt] (-3.047499,-0.597387)--(-3.003935,-0.720748)--(-2.900154,-0.865772)--(-2.942213,-0.744515)--cycle;
\path[fill=blue!33!white,draw=blue!33!white,line width=.12pt] (2.942213,-0.744515)--(2.900154,-0.865772)--(3.003935,-0.720748)--(3.047499,-0.597387)--cycle;
\path[fill=blue!27!white,draw=blue!27!white,line width=.12pt] (-1.534468,-0.052668)--(-1.602769,0.052966)--(-1.529978,-0.023985)--(-1.464780,-0.126340)--cycle;
\path[fill=blue!20!white,draw=blue!20!white,line width=.12pt] (1.464780,-0.126340)--(1.529978,-0.023985)--(1.602769,0.052966)--(1.534468,-0.052668)--cycle;
\path[fill=blue!42!white,draw=blue!42!white,line width=.12pt] (-1.618570,-1.304246)--(-1.511656,-1.196351)--(-1.382778,-1.267448)--(-1.480578,-1.380371)--cycle;
\path[fill=blue!39!white,draw=blue!39!white,line width=.12pt] (1.480578,-1.380371)--(1.382778,-1.267448)--(1.511656,-1.196351)--(1.618570,-1.304246)--cycle;
\path[fill=blue!38!white,draw=blue!38!white,line width=.12pt] (-1.283890,-0.791962)--(-1.248699,-0.647553)--(-1.142240,-0.706283)--(-1.174431,-0.852346)--cycle;
\path[fill=blue!32!white,draw=blue!32!white,line width=.12pt] (1.174431,-0.852346)--(1.142240,-0.706283)--(1.248699,-0.647553)--(1.283890,-0.791962)--cycle;
\path[fill=green!65!black!32!white,draw=green!65!black!32!white,line width=.12pt] (-2.053620,0.435837)--(-2.200811,0.460624)--(-2.147305,0.353819)--(-2.003693,0.336176)--cycle;
\path[fill=green!65!black!29!white,draw=green!65!black!29!white,line width=.12pt] (2.003693,0.336176)--(2.147305,0.353819)--(2.200811,0.460624)--(2.053620,0.435837)--cycle;
\path[fill=blue!16!white,draw=blue!16!white,line width=.12pt] (-3.006279,0.085917)--(-3.078787,-0.041074)--(-3.003935,-0.190486)--(-2.933190,-0.059976)--cycle;
\path[fill=blue!23!white,draw=blue!23!white,line width=.12pt] (2.933190,-0.059976)--(3.003935,-0.190486)--(3.078787,-0.041074)--(3.006279,0.085917)--cycle;
\path[fill=blue!30!white,draw=blue!30!white,line width=.12pt] (-1.424631,-0.245578)--(-1.464780,-0.126340)--(-1.380985,-0.196241)--(-1.343133,-0.313564)--cycle;
\path[fill=blue!23!white,draw=blue!23!white,line width=.12pt] (1.343133,-0.313564)--(1.380985,-0.196241)--(1.464780,-0.126340)--(1.424631,-0.245578)--cycle;
\path[fill=green!65!black!36!white,draw=green!65!black!36!white,line width=.12pt] (-1.755330,0.223661)--(-1.871339,0.292227)--(-1.806687,0.201882)--(-1.694686,0.138917)--cycle;
\path[fill=green!65!black!31!white,draw=green!65!black!31!white,line width=.12pt] (1.694686,0.138917)--(1.806687,0.201882)--(1.871339,0.292227)--(1.755330,0.223661)--cycle;
\path[fill=blue!30!white,draw=blue!30!white,line width=.12pt] (-2.831853,-0.971406)--(-2.739936,-1.057357)--(-2.615501,-1.188905)--(-2.703244,-1.107367)--cycle;
\path[fill=blue!36!white,draw=blue!36!white,line width=.12pt] (2.703244,-1.107367)--(2.615501,-1.188905)--(2.739936,-1.057357)--(2.831853,-0.971406)--cycle;
\path[fill=blue!33!white,draw=blue!33!white,line width=.12pt] (-1.330351,-0.444456)--(-1.343133,-0.313564)--(-1.248699,-0.377228)--(-1.236817,-0.507515)--cycle;
\path[fill=blue!26!white,draw=blue!26!white,line width=.12pt] (1.236817,-0.507515)--(1.248699,-0.377228)--(1.343133,-0.313564)--(1.330351,-0.444456)--cycle;
\path[fill=blue!41!white,draw=blue!41!white,line width=.12pt] (-1.734580,-1.391274)--(-1.618570,-1.304246)--(-1.480578,-1.380371)--(-1.586697,-1.472856)--cycle;
\path[fill=blue!39!white,draw=blue!39!white,line width=.12pt] (1.586697,-1.472856)--(1.480578,-1.380371)--(1.618570,-1.304246)--(1.734580,-1.391274)--cycle;
\path[fill=blue!39!white,draw=blue!39!white,line width=.12pt] (-2.125296,-1.389547)--(-1.995527,-1.359505)--(-1.855225,-1.454092)--(-1.975871,-1.490286)--cycle;
\path[fill=blue!40!white,draw=blue!40!white,line width=.12pt] (1.975871,-1.490286)--(1.855225,-1.454092)--(1.995527,-1.359505)--(2.125296,-1.389547)--cycle;
\path[fill=blue!35!white,draw=blue!35!white,line width=.12pt] (-2.508586,-1.246493)--(-2.386608,-1.277917)--(-2.250079,-1.391810)--(-2.365079,-1.366207)--cycle;
\path[fill=blue!38!white,draw=blue!38!white,line width=.12pt] (2.365079,-1.366207)--(2.250079,-1.391810)--(2.386608,-1.277917)--(2.508586,-1.246493)--cycle;
\path[fill=green!65!black!30!white,draw=green!65!black!30!white,line width=.12pt] (-2.200811,0.460624)--(-2.353885,0.458712)--(-2.296657,0.344479)--(-2.147305,0.353819)--cycle;
\path[fill=green!65!black!29!white,draw=green!65!black!29!white,line width=.12pt] (2.147305,0.353819)--(2.296657,0.344479)--(2.353885,0.458712)--(2.200811,0.460624)--cycle;
\path[fill=blue!14!white,draw=blue!14!white,line width=.12pt] (-2.908701,0.200610)--(-3.006279,0.085917)--(-2.933190,-0.059976)--(-2.837984,0.059451)--cycle;
\path[fill=blue!21!white,draw=blue!21!white,line width=.12pt] (2.837984,0.059451)--(2.933190,-0.059976)--(3.006279,0.085917)--(2.908701,0.200610)--cycle;
\path[fill=blue!41!white,draw=blue!41!white,line width=.12pt] (-1.297056,-1.138426)--(-1.226706,-0.998263)--(-1.100561,-1.055429)--(-1.163677,-1.198870)--cycle;
\path[fill=blue!37!white,draw=blue!37!white,line width=.12pt] (1.163677,-1.198870)--(1.100561,-1.055429)--(1.226706,-0.998263)--(1.297056,-1.138426)--cycle;
\path[fill=blue!22!white,draw=blue!22!white,line width=.12pt] (-3.062209,-0.464455)--(-3.047499,-0.597387)--(-2.942213,-0.744515)--(-2.956415,-0.612293)--cycle;
\path[fill=blue!30!white,draw=blue!30!white,line width=.12pt] (2.956415,-0.612293)--(2.942213,-0.744515)--(3.047499,-0.597387)--(3.062209,-0.464455)--cycle;
\path[fill=blue!36!white,draw=blue!36!white,line width=.12pt] (-1.248699,-0.647553)--(-1.236817,-0.507515)--(-1.131371,-0.565685)--(-1.142240,-0.706283)--cycle;
\path[fill=blue!29!white,draw=blue!29!white,line width=.12pt] (1.142240,-0.706283)--(1.131371,-0.565685)--(1.236817,-0.507515)--(1.248699,-0.647553)--cycle;
\path[fill=blue!40!white,draw=blue!40!white,line width=.12pt] (-1.226706,-0.998263)--(-1.174431,-0.852346)--(-1.053662,-0.907076)--(-1.100561,-1.055429)--cycle;
\path[fill=blue!35!white,draw=blue!35!white,line width=.12pt] (1.100561,-1.055429)--(1.053662,-0.907076)--(1.174431,-0.852346)--(1.226706,-0.998263)--cycle;
\path[fill=blue!42!white,draw=blue!42!white,line width=.12pt] (-1.382778,-1.267448)--(-1.297056,-1.138426)--(-1.163677,-1.198870)--(-1.240584,-1.331887)--cycle;
\path[fill=blue!38!white,draw=blue!38!white,line width=.12pt] (1.240584,-1.331887)--(1.163677,-1.198870)--(1.297056,-1.138426)--(1.382778,-1.267448)--cycle;
\path[fill=blue!24!white,draw=blue!24!white,line width=.12pt] (-1.602769,0.052966)--(-1.694686,0.138917)--(-1.617721,0.057553)--(-1.529978,-0.023985)--cycle;
\path[fill=blue!18!white,draw=blue!18!white,line width=.12pt] (1.529978,-0.023985)--(1.617721,0.057553)--(1.694686,0.138917)--(1.602769,0.052966)--cycle;
\path[fill=blue!14!white,draw=blue!14!white,line width=.12pt] (-2.353885,0.458712)--(-2.506958,0.430175)--(-2.446009,0.308514)--(-2.296657,0.344479)--cycle;
\path[fill=blue!15!white,draw=blue!15!white,line width=.12pt] (2.296657,0.344479)--(2.446009,0.308514)--(2.506958,0.430175)--(2.353885,0.458712)--cycle;
\path[fill=blue!13!white,draw=blue!13!white,line width=.12pt] (-2.789801,0.298596)--(-2.908701,0.200610)--(-2.837984,0.059451)--(-2.721975,0.163208)--cycle;
\path[fill=blue!19!white,draw=blue!19!white,line width=.12pt] (2.721975,0.163208)--(2.837984,0.059451)--(2.908701,0.200610)--(2.789801,0.298596)--cycle;
\path[fill=blue!13!white,draw=blue!13!white,line width=.12pt] (-2.506958,0.430175)--(-2.654149,0.376111)--(-2.589621,0.247306)--(-2.446009,0.308514)--cycle;
\path[fill=blue!16!white,draw=blue!16!white,line width=.12pt] (2.446009,0.308514)--(2.589621,0.247306)--(2.654149,0.376111)--(2.506958,0.430175)--cycle;
\path[fill=blue!38!white,draw=blue!38!white,line width=.12pt] (-1.174431,-0.852346)--(-1.142240,-0.706283)--(-1.024781,-0.759512)--(-1.053662,-0.907076)--cycle;
\path[fill=blue!33!white,draw=blue!33!white,line width=.12pt] (1.053662,-0.907076)--(1.024781,-0.759512)--(1.142240,-0.706283)--(1.174431,-0.852346)--cycle;
\path[fill=blue!13!white,draw=blue!13!white,line width=.12pt] (-2.654149,0.376111)--(-2.789801,0.298596)--(-2.721975,0.163208)--(-2.589621,0.247306)--cycle;
\path[fill=blue!17!white,draw=blue!17!white,line width=.12pt] (2.589621,0.247306)--(2.721975,0.163208)--(2.789801,0.298596)--(2.654149,0.376111)--cycle;
\path[fill=green!65!black!34!white,draw=green!65!black!34!white,line width=.12pt] (-1.871339,0.292227)--(-2.003693,0.336176)--(-1.934468,0.239441)--(-1.806687,0.201882)--cycle;
\path[fill=green!65!black!30!white,draw=green!65!black!30!white,line width=.12pt] (1.806687,0.201882)--(1.934468,0.239441)--(2.003693,0.336176)--(1.871339,0.292227)--cycle;
\path[fill=blue!28!white,draw=blue!28!white,line width=.12pt] (-1.464780,-0.126340)--(-1.529978,-0.023985)--(-1.442454,-0.096998)--(-1.380985,-0.196241)--cycle;
\path[fill=blue!21!white,draw=blue!21!white,line width=.12pt] (1.380985,-0.196241)--(1.442454,-0.096998)--(1.529978,-0.023985)--(1.464780,-0.126340)--cycle;
\path[fill=blue!42!white,draw=blue!42!white,line width=.12pt] (-1.480578,-1.380371)--(-1.382778,-1.267448)--(-1.240584,-1.331887)--(-1.328326,-1.449368)--cycle;
\path[fill=blue!39!white,draw=blue!39!white,line width=.12pt] (1.328326,-1.449368)--(1.240584,-1.331887)--(1.382778,-1.267448)--(1.480578,-1.380371)--cycle;
\path[fill=blue!40!white,draw=blue!40!white,line width=.12pt] (-1.855225,-1.454092)--(-1.734580,-1.391274)--(-1.586697,-1.472856)--(-1.697056,-1.541348)--cycle;
\path[fill=blue!40!white,draw=blue!40!white,line width=.12pt] (1.697056,-1.541348)--(1.586697,-1.472856)--(1.734580,-1.391274)--(1.855225,-1.454092)--cycle;
\path[fill=blue!27!white,draw=blue!27!white,line width=.12pt] (-2.900154,-0.865772)--(-2.831853,-0.971406)--(-2.703244,-1.107367)--(-2.768442,-1.005012)--cycle;
\path[fill=blue!34!white,draw=blue!34!white,line width=.12pt] (2.768442,-1.005012)--(2.703244,-1.107367)--(2.831853,-0.971406)--(2.900154,-0.865772)--cycle;
\path[fill=blue!31!white,draw=blue!31!white,line width=.12pt] (-1.343133,-0.313564)--(-1.380985,-0.196241)--(-1.283890,-0.261700)--(-1.248699,-0.377228)--cycle;
\path[fill=blue!24!white,draw=blue!24!white,line width=.12pt] (1.248699,-0.377228)--(1.283890,-0.261700)--(1.380985,-0.196241)--(1.343133,-0.313564)--cycle;
\path[fill=blue!19!white,draw=blue!19!white,line width=.12pt] (-3.047499,-0.327062)--(-3.062209,-0.464455)--(-2.956415,-0.612293)--(-2.942213,-0.474190)--cycle;
\path[fill=blue!28!white,draw=blue!28!white,line width=.12pt] (2.942213,-0.474190)--(2.956415,-0.612293)--(3.062209,-0.464455)--(3.047499,-0.327062)--cycle;
\path[fill=blue!34!white,draw=blue!34!white,line width=.12pt] (-1.236817,-0.507515)--(-1.248699,-0.377228)--(-1.142240,-0.435958)--(-1.131371,-0.565685)--cycle;
\path[fill=blue!27!white,draw=blue!27!white,line width=.12pt] (1.131371,-0.565685)--(1.142240,-0.435958)--(1.248699,-0.377228)--(1.236817,-0.507515)--cycle;
\path[fill=blue!37!white,draw=blue!37!white,line width=.12pt] (-2.250079,-1.391810)--(-2.125296,-1.389547)--(-1.975871,-1.490286)--(-2.091880,-1.498463)--cycle;
\path[fill=blue!39!white,draw=blue!39!white,line width=.12pt] (2.091880,-1.498463)--(1.975871,-1.490286)--(2.125296,-1.389547)--(2.250079,-1.391810)--cycle;
\path[fill=blue!32!white,draw=blue!32!white,line width=.12pt] (-2.615501,-1.188905)--(-2.508586,-1.246493)--(-2.365079,-1.366207)--(-2.465877,-1.313721)--cycle;
\path[fill=blue!37!white,draw=blue!37!white,line width=.12pt] (2.465877,-1.313721)--(2.365079,-1.366207)--(2.508586,-1.246493)--(2.615501,-1.188905)--cycle;
\path[fill=blue!41!white,draw=blue!41!white,line width=.12pt] (-1.163677,-1.198870)--(-1.100561,-1.055429)--(-0.963817,-1.106137)--(-1.019091,-1.252486)--cycle;
\path[fill=blue!37!white,draw=blue!37!white,line width=.12pt] (1.019091,-1.252486)--(0.963817,-1.106137)--(1.100561,-1.055429)--(1.163677,-1.198870)--cycle;
\path[fill=blue!36!white,draw=blue!36!white,line width=.12pt] (-1.142240,-0.706283)--(-1.131371,-0.565685)--(-1.015029,-0.618408)--(-1.024781,-0.759512)--cycle;
\path[fill=blue!30!white,draw=blue!30!white,line width=.12pt] (1.024781,-0.759512)--(1.015029,-0.618408)--(1.131371,-0.565685)--(1.142240,-0.706283)--cycle;
\path[fill=blue!40!white,draw=blue!40!white,line width=.12pt] (-1.100561,-1.055429)--(-1.053662,-0.907076)--(-0.922745,-0.955623)--(-0.963817,-1.106137)--cycle;
\path[fill=blue!35!white,draw=blue!35!white,line width=.12pt] (0.963817,-1.106137)--(0.922745,-0.955623)--(1.053662,-0.907076)--(1.100561,-1.055429)--cycle;
\path[fill=green!65!black!37!white,draw=green!65!black!37!white,line width=.12pt] (-1.694686,0.138917)--(-1.806687,0.201882)--(-1.724636,0.115141)--(-1.617721,0.057553)--cycle;
\path[fill=green!65!black!32!white,draw=green!65!black!32!white,line width=.12pt] (1.617721,0.057553)--(1.724636,0.115141)--(1.806687,0.201882)--(1.694686,0.138917)--cycle;
\path[fill=blue!42!white,draw=blue!42!white,line width=.12pt] (-1.240584,-1.331887)--(-1.163677,-1.198870)--(-1.019091,-1.252486)--(-1.086442,-1.389047)--cycle;
\path[fill=blue!38!white,draw=blue!38!white,line width=.12pt] (1.086442,-1.389047)--(1.019091,-1.252486)--(1.163677,-1.198870)--(1.240584,-1.331887)--cycle;
\path[fill=blue!41!white,draw=blue!41!white,line width=.12pt] (-1.586697,-1.472856)--(-1.480578,-1.380371)--(-1.328326,-1.449368)--(-1.423533,-1.546798)--cycle;
\path[fill=blue!40!white,draw=blue!40!white,line width=.12pt] (1.423533,-1.546798)--(1.328326,-1.449368)--(1.480578,-1.380371)--(1.586697,-1.472856)--cycle;
\path[fill=green!65!black!32!white,draw=green!65!black!32!white,line width=.12pt] (-2.003693,0.336176)--(-2.147305,0.353819)--(-2.073119,0.250151)--(-1.934468,0.239441)--cycle;
\path[fill=green!65!black!29!white,draw=green!65!black!29!white,line width=.12pt] (1.934468,0.239441)--(2.073119,0.250151)--(2.147305,0.353819)--(2.003693,0.336176)--cycle;
\path[fill=blue!38!white,draw=blue!38!white,line width=.12pt] (-1.053662,-0.907076)--(-1.024781,-0.759512)--(-0.897452,-0.806729)--(-0.922745,-0.955623)--cycle;
\path[fill=blue!33!white,draw=blue!33!white,line width=.12pt] (0.922745,-0.955623)--(0.897452,-0.806729)--(1.024781,-0.759512)--(1.053662,-0.907076)--cycle;
\path[fill=blue!17!white,draw=blue!17!white,line width=.12pt] (-3.003935,-0.190486)--(-3.047499,-0.327062)--(-2.942213,-0.474190)--(-2.900154,-0.335511)--cycle;
\path[fill=blue!25!white,draw=blue!25!white,line width=.12pt] (2.900154,-0.335511)--(2.942213,-0.474190)--(3.047499,-0.327062)--(3.003935,-0.190486)--cycle;
\path[fill=blue!39!white,draw=blue!39!white,line width=.12pt] (-1.975871,-1.490286)--(-1.855225,-1.454092)--(-1.697056,-1.541348)--(-1.807416,-1.583216)--cycle;
\path[fill=blue!39!white,draw=blue!39!white,line width=.12pt] (1.807416,-1.583216)--(1.697056,-1.541348)--(1.855225,-1.454092)--(1.975871,-1.490286)--cycle;
\path[fill=blue!25!white,draw=blue!25!white,line width=.12pt] (-1.529978,-0.023985)--(-1.617721,0.057553)--(-1.525177,-0.019647)--(-1.442454,-0.096998)--cycle;
\path[fill=blue!19!white,draw=blue!19!white,line width=.12pt] (1.442454,-0.096998)--(1.525177,-0.019647)--(1.617721,0.057553)--(1.529978,-0.023985)--cycle;
\path[fill=blue!24!white,draw=blue!24!white,line width=.12pt] (-2.942213,-0.744515)--(-2.900154,-0.865772)--(-2.768442,-1.005012)--(-2.808591,-0.885774)--cycle;
\path[fill=blue!32!white,draw=blue!32!white,line width=.12pt] (2.808591,-0.885774)--(2.768442,-1.005012)--(2.900154,-0.865772)--(2.942213,-0.744515)--cycle;
\path[fill=blue!42!white,draw=blue!42!white,line width=.12pt] (-1.328326,-1.449368)--(-1.240584,-1.331887)--(-1.086442,-1.389047)--(-1.163283,-1.510570)--cycle;
\path[fill=blue!39!white,draw=blue!39!white,line width=.12pt] (1.163283,-1.510570)--(1.086442,-1.389047)--(1.240584,-1.331887)--(1.328326,-1.449368)--cycle;
\path[fill=blue!28!white,draw=blue!28!white,line width=.12pt] (-1.380985,-0.196241)--(-1.442454,-0.096998)--(-1.341037,-0.165370)--(-1.283890,-0.261700)--cycle;
\path[fill=blue!21!white,draw=blue!21!white,line width=.12pt] (1.283890,-0.261700)--(1.341037,-0.165370)--(1.442454,-0.096998)--(1.380985,-0.196241)--cycle;
\path[fill=blue!31!white,draw=blue!31!white,line width=.12pt] (-1.248699,-0.377228)--(-1.283890,-0.261700)--(-1.174431,-0.322085)--(-1.142240,-0.435958)--cycle;
\path[fill=blue!24!white,draw=blue!24!white,line width=.12pt] (1.142240,-0.435958)--(1.174431,-0.322085)--(1.283890,-0.261700)--(1.248699,-0.377228)--cycle;
\path[fill=blue!34!white,draw=blue!34!white,line width=.12pt] (-1.131371,-0.565685)--(-1.142240,-0.435958)--(-1.024781,-0.489187)--(-1.015029,-0.618408)--cycle;
\path[fill=blue!28!white,draw=blue!28!white,line width=.12pt] (1.015029,-0.618408)--(1.024781,-0.489187)--(1.142240,-0.435958)--(1.131371,-0.565685)--cycle;
\path[fill=blue!41!white,draw=blue!41!white,line width=.12pt] (-1.019091,-1.252486)--(-0.963817,-1.106137)--(-0.817791,-1.149900)--(-0.864690,-1.298758)--cycle;
\path[fill=blue!38!white,draw=blue!38!white,line width=.12pt] (0.864690,-1.298758)--(0.817791,-1.149900)--(0.963817,-1.106137)--(1.019091,-1.252486)--cycle;
\path[fill=blue!40!white,draw=blue!40!white,line width=.12pt] (-0.963817,-1.106137)--(-0.922745,-0.955623)--(-0.782941,-0.997521)--(-0.817791,-1.149900)--cycle;
\path[fill=blue!36!white,draw=blue!36!white,line width=.12pt] (0.817791,-1.149900)--(0.782941,-0.997521)--(0.922745,-0.955623)--(0.963817,-1.106137)--cycle;
\path[fill=green!65!black!30!white,draw=green!65!black!30!white,line width=.12pt] (-2.147305,0.353819)--(-2.296657,0.344479)--(-2.217311,0.233600)--(-2.073119,0.250151)--cycle;
\path[fill=green!65!black!29!white,draw=green!65!black!29!white,line width=.12pt] (2.073119,0.250151)--(2.217311,0.233600)--(2.296657,0.344479)--(2.147305,0.353819)--cycle;
\path[fill=blue!36!white,draw=blue!36!white,line width=.12pt] (-1.024781,-0.759512)--(-1.015029,-0.618408)--(-0.888912,-0.665176)--(-0.897452,-0.806729)--cycle;
\path[fill=blue!31!white,draw=blue!31!white,line width=.12pt] (0.897452,-0.806729)--(0.888912,-0.665176)--(1.015029,-0.618408)--(1.024781,-0.759512)--cycle;
\path[fill=blue!35!white,draw=blue!35!white,line width=.12pt] (-2.365079,-1.366207)--(-2.250079,-1.391810)--(-2.091880,-1.498463)--(-2.198794,-1.478311)--cycle;
\path[fill=blue!38!white,draw=blue!38!white,line width=.12pt] (2.198794,-1.478311)--(2.091880,-1.498463)--(2.250079,-1.391810)--(2.365079,-1.366207)--cycle;
\path[fill=blue!29!white,draw=blue!29!white,line width=.12pt] (-2.703244,-1.107367)--(-2.615501,-1.188905)--(-2.465877,-1.313721)--(-2.548600,-1.236371)--cycle;
\path[fill=blue!35!white,draw=blue!35!white,line width=.12pt] (2.548600,-1.236371)--(2.465877,-1.313721)--(2.615501,-1.188905)--(2.703244,-1.107367)--cycle;
\path[fill=blue!15!white,draw=blue!15!white,line width=.12pt] (-2.933190,-0.059976)--(-3.003935,-0.190486)--(-2.900154,-0.335511)--(-2.831853,-0.201586)--cycle;
\path[fill=blue!23!white,draw=blue!23!white,line width=.12pt] (2.831853,-0.201586)--(2.900154,-0.335511)--(3.003935,-0.190486)--(2.933190,-0.059976)--cycle;
\path[fill=blue!40!white,draw=blue!40!white,line width=.12pt] (-1.697056,-1.541348)--(-1.586697,-1.472856)--(-1.423533,-1.546798)--(-1.522544,-1.620433)--cycle;
\path[fill=blue!40!white,draw=blue!40!white,line width=.12pt] (1.522544,-1.620433)--(1.423533,-1.546798)--(1.586697,-1.472856)--(1.697056,-1.541348)--cycle;
\path[fill=blue!42!white,draw=blue!42!white,line width=.12pt] (-1.086442,-1.389047)--(-1.019091,-1.252486)--(-0.864690,-1.298758)--(-0.921837,-1.438377)--cycle;
\path[fill=blue!39!white,draw=blue!39!white,line width=.12pt] (0.921837,-1.438377)--(0.864690,-1.298758)--(1.019091,-1.252486)--(1.086442,-1.389047)--cycle;
\path[fill=blue!38!white,draw=blue!38!white,line width=.12pt] (-0.922745,-0.955623)--(-0.897452,-0.806729)--(-0.761481,-0.847478)--(-0.782941,-0.997521)--cycle;
\path[fill=blue!34!white,draw=blue!34!white,line width=.12pt] (0.782941,-0.997521)--(0.761481,-0.847478)--(0.897452,-0.806729)--(0.922745,-0.955623)--cycle;
\path[fill=green!65!black!34!white,draw=green!65!black!34!white,line width=.12pt] (-1.806687,0.201882)--(-1.934468,0.239441)--(-1.846614,0.146565)--(-1.724636,0.115141)--cycle;
\path[fill=green!65!black!30!white,draw=green!65!black!30!white,line width=.12pt] (1.724636,0.115141)--(1.846614,0.146565)--(1.934468,0.239441)--(1.806687,0.201882)--cycle;
\path[fill=blue!14!white,draw=blue!14!white,line width=.12pt] (-2.296657,0.344479)--(-2.446009,0.308514)--(-2.361503,0.190425)--(-2.217311,0.233600)--cycle;
\path[fill=blue!15!white,draw=blue!15!white,line width=.12pt] (2.217311,0.233600)--(2.361503,0.190425)--(2.446009,0.308514)--(2.296657,0.344479)--cycle;
\path[fill=blue!41!white,draw=blue!41!white,line width=.12pt] (-1.423533,-1.546798)--(-1.328326,-1.449368)--(-1.163283,-1.510570)--(-1.246659,-1.612387)--cycle;
\path[fill=blue!40!white,draw=blue!40!white,line width=.12pt] (1.246659,-1.612387)--(1.163283,-1.510570)--(1.328326,-1.449368)--(1.423533,-1.546798)--cycle;
\path[fill=blue!14!white,draw=blue!14!white,line width=.12pt] (-2.837984,0.059451)--(-2.933190,-0.059976)--(-2.831853,-0.201586)--(-2.739936,-0.077561)--cycle;
\path[fill=blue!21!white,draw=blue!21!white,line width=.12pt] (2.739936,-0.077561)--(2.831853,-0.201586)--(2.933190,-0.059976)--(2.837984,0.059451)--cycle;
\path[fill=blue!40!white,draw=blue!40!white,line width=.12pt] (-0.817791,-1.149900)--(-0.782941,-0.997521)--(-0.635598,-1.032365)--(-0.663889,-1.186295)--cycle;
\path[fill=blue!37!white,draw=blue!37!white,line width=.12pt] (0.663889,-1.186295)--(0.635598,-1.032365)--(0.782941,-0.997521)--(0.817791,-1.149900)--cycle;
\path[fill=blue!21!white,draw=blue!21!white,line width=.12pt] (-2.956415,-0.612293)--(-2.942213,-0.744515)--(-2.808591,-0.885774)--(-2.822148,-0.754235)--cycle;
\path[fill=blue!29!white,draw=blue!29!white,line width=.12pt] (2.822148,-0.754235)--(2.808591,-0.885774)--(2.942213,-0.744515)--(2.956415,-0.612293)--cycle;
\path[fill=blue!41!white,draw=blue!41!white,line width=.12pt] (-0.864690,-1.298758)--(-0.817791,-1.149900)--(-0.663889,-1.186295)--(-0.701962,-1.337241)--cycle;
\path[fill=blue!38!white,draw=blue!38!white,line width=.12pt] (0.701962,-1.337241)--(0.663889,-1.186295)--(0.817791,-1.149900)--(0.864690,-1.298758)--cycle;
\path[fill=blue!37!white,draw=blue!37!white,line width=.12pt] (-2.091880,-1.498463)--(-1.975871,-1.490286)--(-1.807416,-1.583216)--(-1.913535,-1.596850)--cycle;
\path[fill=blue!39!white,draw=blue!39!white,line width=.12pt] (1.913535,-1.596850)--(1.807416,-1.583216)--(1.975871,-1.490286)--(2.091880,-1.498463)--cycle;
\path[fill=blue!13!white,draw=blue!13!white,line width=.12pt] (-2.446009,0.308514)--(-2.589621,0.247306)--(-2.500154,0.122284)--(-2.361503,0.190425)--cycle;
\path[fill=blue!16!white,draw=blue!16!white,line width=.12pt] (2.361503,0.190425)--(2.500154,0.122284)--(2.589621,0.247306)--(2.446009,0.308514)--cycle;
\path[fill=green!65!black!37!white,draw=green!65!black!37!white,line width=.12pt] (-1.617721,0.057553)--(-1.724636,0.115141)--(-1.625975,0.032838)--(-1.525177,-0.019647)--cycle;
\path[fill=green!65!black!32!white,draw=green!65!black!32!white,line width=.12pt] (1.525177,-0.019647)--(1.625975,0.032838)--(1.724636,0.115141)--(1.617721,0.057553)--cycle;
\path[fill=blue!42!white,draw=blue!42!white,line width=.12pt] (-1.163283,-1.510570)--(-1.086442,-1.389047)--(-0.921837,-1.438377)--(-0.987036,-1.563389)--cycle;
\path[fill=blue!40!white,draw=blue!40!white,line width=.12pt] (0.987036,-1.563389)--(0.921837,-1.438377)--(1.086442,-1.389047)--(1.163283,-1.510570)--cycle;
\path[fill=blue!36!white,draw=blue!36!white,line width=.12pt] (-0.897452,-0.806729)--(-0.888912,-0.665176)--(-0.754235,-0.705537)--(-0.761481,-0.847478)--cycle;
\path[fill=blue!32!white,draw=blue!32!white,line width=.12pt] (0.761481,-0.847478)--(0.754235,-0.705537)--(0.888912,-0.665176)--(0.897452,-0.806729)--cycle;
\path[fill=blue!13!white,draw=blue!13!white,line width=.12pt] (-2.721975,0.163208)--(-2.837984,0.059451)--(-2.739936,-0.077561)--(-2.627935,0.031796)--cycle;
\path[fill=blue!18!white,draw=blue!18!white,line width=.12pt] (2.627935,0.031796)--(2.739936,-0.077561)--(2.837984,0.059451)--(2.721975,0.163208)--cycle;
\path[fill=blue!34!white,draw=blue!34!white,line width=.12pt] (-1.015029,-0.618408)--(-1.024781,-0.489187)--(-0.897452,-0.536404)--(-0.888912,-0.665176)--cycle;
\path[fill=blue!28!white,draw=blue!28!white,line width=.12pt] (0.888912,-0.665176)--(0.897452,-0.536404)--(1.024781,-0.489187)--(1.015029,-0.618408)--cycle;
\path[fill=blue!13!white,draw=blue!13!white,line width=.12pt] (-2.589621,0.247306)--(-2.721975,0.163208)--(-2.627935,0.031796)--(-2.500154,0.122284)--cycle;
\path[fill=blue!17!white,draw=blue!17!white,line width=.12pt] (2.500154,0.122284)--(2.627935,0.031796)--(2.721975,0.163208)--(2.589621,0.247306)--cycle;
\path[fill=blue!38!white,draw=blue!38!white,line width=.12pt] (-0.782941,-0.997521)--(-0.761481,-0.847478)--(-0.618176,-0.881367)--(-0.635598,-1.032365)--cycle;
\path[fill=blue!35!white,draw=blue!35!white,line width=.12pt] (0.635598,-1.032365)--(0.618176,-0.881367)--(0.761481,-0.847478)--(0.782941,-0.997521)--cycle;
\path[fill=blue!31!white,draw=blue!31!white,line width=.12pt] (-1.142240,-0.435958)--(-1.174431,-0.322085)--(-1.053662,-0.376814)--(-1.024781,-0.489187)--cycle;
\path[fill=blue!25!white,draw=blue!25!white,line width=.12pt] (1.024781,-0.489187)--(1.053662,-0.376814)--(1.174431,-0.322085)--(1.142240,-0.435958)--cycle;
\path[fill=blue!25!white,draw=blue!25!white,line width=.12pt] (-1.442454,-0.096998)--(-1.525177,-0.019647)--(-1.417944,-0.091940)--(-1.341037,-0.165370)--cycle;
\path[fill=blue!19!white,draw=blue!19!white,line width=.12pt] (1.341037,-0.165370)--(1.417944,-0.091940)--(1.525177,-0.019647)--(1.442454,-0.096998)--cycle;
\path[fill=blue!28!white,draw=blue!28!white,line width=.12pt] (-1.283890,-0.261700)--(-1.341037,-0.165370)--(-1.226706,-0.228443)--(-1.174431,-0.322085)--cycle;
\path[fill=blue!22!white,draw=blue!22!white,line width=.12pt] (1.174431,-0.322085)--(1.226706,-0.228443)--(1.341037,-0.165370)--(1.283890,-0.261700)--cycle;
\path[fill=blue!42!white,draw=blue!42!white,line width=.12pt] (-0.921837,-1.438377)--(-0.864690,-1.298758)--(-0.701962,-1.337241)--(-0.748354,-1.479402)--cycle;
\path[fill=blue!39!white,draw=blue!39!white,line width=.12pt] (0.748354,-1.479402)--(0.701962,-1.337241)--(0.864690,-1.298758)--(0.921837,-1.438377)--cycle;
\path[fill=green!65!black!32!white,draw=green!65!black!32!white,line width=.12pt] (-1.934468,0.239441)--(-2.073119,0.250151)--(-1.978968,0.150618)--(-1.846614,0.146565)--cycle;
\path[fill=green!65!black!30!white,draw=green!65!black!30!white,line width=.12pt] (1.846614,0.146565)--(1.978968,0.150618)--(2.073119,0.250151)--(1.934468,0.239441)--cycle;
\path[fill=blue!32!white,draw=blue!32!white,line width=.12pt] (-2.465877,-1.313721)--(-2.365079,-1.366207)--(-2.198794,-1.478311)--(-2.292506,-1.430603)--cycle;
\path[fill=blue!36!white,draw=blue!36!white,line width=.12pt] (2.292506,-1.430603)--(2.198794,-1.478311)--(2.365079,-1.366207)--(2.465877,-1.313721)--cycle;
\path[fill=blue!26!white,draw=blue!26!white,line width=.12pt] (-2.768442,-1.005012)--(-2.703244,-1.107367)--(-2.548600,-1.236371)--(-2.610069,-1.137127)--cycle;
\path[fill=blue!33!white,draw=blue!33!white,line width=.12pt] (2.610069,-1.137127)--(2.548600,-1.236371)--(2.703244,-1.107367)--(2.768442,-1.005012)--cycle;
\path[fill=blue!39!white,draw=blue!39!white,line width=.12pt] (-1.807416,-1.583216)--(-1.697056,-1.541348)--(-1.522544,-1.620433)--(-1.621555,-1.667443)--cycle;
\path[fill=blue!39!white,draw=blue!39!white,line width=.12pt] (1.621555,-1.667443)--(1.522544,-1.620433)--(1.697056,-1.541348)--(1.807416,-1.583216)--cycle;
\path[fill=blue!40!white,draw=blue!40!white,line width=.12pt] (-0.663889,-1.186295)--(-0.635598,-1.032365)--(-0.482133,-1.059820)--(-0.503593,-1.214972)--cycle;
\path[fill=blue!37!white,draw=blue!37!white,line width=.12pt] (0.503593,-1.214972)--(0.482133,-1.059820)--(0.635598,-1.032365)--(0.663889,-1.186295)--cycle;
\path[fill=blue!41!white,draw=blue!41!white,line width=.12pt] (-0.701962,-1.337241)--(-0.663889,-1.186295)--(-0.503593,-1.214972)--(-0.532473,-1.367563)--cycle;
\path[fill=blue!39!white,draw=blue!39!white,line width=.12pt] (0.532473,-1.367563)--(0.503593,-1.214972)--(0.663889,-1.186295)--(0.701962,-1.337241)--cycle;
\path[fill=blue!40!white,draw=blue!40!white,line width=.12pt] (-1.522544,-1.620433)--(-1.423533,-1.546798)--(-1.246659,-1.612387)--(-1.333369,-1.690584)--cycle;
\path[fill=blue!40!white,draw=blue!40!white,line width=.12pt] (1.333369,-1.690584)--(1.246659,-1.612387)--(1.423533,-1.546798)--(1.522544,-1.620433)--cycle;
\path[fill=blue!36!white,draw=blue!36!white,line width=.12pt] (-0.761481,-0.847478)--(-0.754235,-0.705537)--(-0.612293,-0.739104)--(-0.618176,-0.881367)--cycle;
\path[fill=blue!32!white,draw=blue!32!white,line width=.12pt] (0.618176,-0.881367)--(0.612293,-0.739104)--(0.754235,-0.705537)--(0.761481,-0.847478)--cycle;
\path[fill=blue!38!white,draw=blue!38!white,line width=.12pt] (-0.635598,-1.032365)--(-0.618176,-0.881367)--(-0.468918,-0.908070)--(-0.482133,-1.059820)--cycle;
\path[fill=blue!35!white,draw=blue!35!white,line width=.12pt] (0.482133,-1.059820)--(0.468918,-0.908070)--(0.618176,-0.881367)--(0.635598,-1.032365)--cycle;
\path[fill=blue!19!white,draw=blue!19!white,line width=.12pt] (-2.942213,-0.474190)--(-2.956415,-0.612293)--(-2.822148,-0.754235)--(-2.808591,-0.615449)--cycle;
\path[fill=blue!27!white,draw=blue!27!white,line width=.12pt] (2.808591,-0.615449)--(2.822148,-0.754235)--(2.956415,-0.612293)--(2.942213,-0.474190)--cycle;
\path[fill=blue!41!white,draw=blue!41!white,line width=.12pt] (-1.246659,-1.612387)--(-1.163283,-1.510570)--(-0.987036,-1.563389)--(-1.057780,-1.668992)--cycle;
\path[fill=blue!40!white,draw=blue!40!white,line width=.12pt] (1.057780,-1.668992)--(0.987036,-1.563389)--(1.163283,-1.510570)--(1.246659,-1.612387)--cycle;
\path[fill=blue!34!white,draw=blue!34!white,line width=.12pt] (-0.888912,-0.665176)--(-0.897452,-0.536404)--(-0.761481,-0.577153)--(-0.754235,-0.705537)--cycle;
\path[fill=blue!29!white,draw=blue!29!white,line width=.12pt] (0.754235,-0.705537)--(0.761481,-0.577153)--(0.897452,-0.536404)--(0.888912,-0.665176)--cycle;
\path[fill=blue!42!white,draw=blue!42!white,line width=.12pt] (-0.987036,-1.563389)--(-0.921837,-1.438377)--(-0.748354,-1.479402)--(-0.801283,-1.607317)--cycle;
\path[fill=blue!40!white,draw=blue!40!white,line width=.12pt] (0.801283,-1.607317)--(0.748354,-1.479402)--(0.921837,-1.438377)--(0.987036,-1.563389)--cycle;
\path[fill=blue!42!white,draw=blue!42!white,line width=.12pt] (-0.748354,-1.479402)--(-0.701962,-1.337241)--(-0.532473,-1.367563)--(-0.567664,-1.511728)--cycle;
\path[fill=blue!40!white,draw=blue!40!white,line width=.12pt] (0.567664,-1.511728)--(0.532473,-1.367563)--(0.701962,-1.337241)--(0.748354,-1.479402)--cycle;
\path[fill=blue!40!white,draw=blue!40!white,line width=.12pt] (-0.503593,-1.214972)--(-0.482133,-1.059820)--(-0.324025,-1.079622)--(-0.338447,-1.235655)--cycle;
\path[fill=blue!38!white,draw=blue!38!white,line width=.12pt] (0.338447,-1.235655)--(0.324025,-1.079622)--(0.482133,-1.059820)--(0.503593,-1.214972)--cycle;
\path[fill=blue!41!white,draw=blue!41!white,line width=.12pt] (-0.532473,-1.367563)--(-0.503593,-1.214972)--(-0.338447,-1.235655)--(-0.357857,-1.389432)--cycle;
\path[fill=blue!39!white,draw=blue!39!white,line width=.12pt] (0.357857,-1.389432)--(0.338447,-1.235655)--(0.503593,-1.214972)--(0.532473,-1.367563)--cycle;
\path[fill=green!65!black!34!white,draw=green!65!black!34!white,line width=.12pt] (-1.724636,0.115141)--(-1.846614,0.146565)--(-1.740975,0.058441)--(-1.625975,0.032838)--cycle;
\path[fill=green!65!black!31!white,draw=green!65!black!31!white,line width=.12pt] (1.625975,0.032838)--(1.740975,0.058441)--(1.846614,0.146565)--(1.724636,0.115141)--cycle;
\path[fill=blue!31!white,draw=blue!31!white,line width=.12pt] (-1.024781,-0.489187)--(-1.053662,-0.376814)--(-0.922745,-0.425362)--(-0.897452,-0.536404)--cycle;
\path[fill=blue!26!white,draw=blue!26!white,line width=.12pt] (0.897452,-0.536404)--(0.922745,-0.425362)--(1.053662,-0.376814)--(1.024781,-0.489187)--cycle;
\path[fill=green!65!black!30!white,draw=green!65!black!30!white,line width=.12pt] (-2.073119,0.250151)--(-2.217311,0.233600)--(-2.116611,0.127144)--(-1.978968,0.150618)--cycle;
\path[fill=green!65!black!29!white,draw=green!65!black!29!white,line width=.12pt] (1.978968,0.150618)--(2.116611,0.127144)--(2.217311,0.233600)--(2.073119,0.250151)--cycle;
\path[fill=blue!38!white,draw=blue!38!white,line width=.12pt] (-0.482133,-1.059820)--(-0.468918,-0.908070)--(-0.315143,-0.927329)--(-0.324025,-1.079622)--cycle;
\path[fill=blue!36!white,draw=blue!36!white,line width=.12pt] (0.324025,-1.079622)--(0.315143,-0.927329)--(0.468918,-0.908070)--(0.482133,-1.059820)--cycle;
\path[fill=blue!34!white,draw=blue!34!white,line width=.12pt] (-2.198794,-1.478311)--(-2.091880,-1.498463)--(-1.913535,-1.596850)--(-2.011334,-1.581726)--cycle;
\path[fill=blue!38!white,draw=blue!38!white,line width=.12pt] (2.011334,-1.581726)--(1.913535,-1.596850)--(2.091880,-1.498463)--(2.198794,-1.478311)--cycle;
\path[fill=blue!39!white,draw=blue!39!white,line width=.12pt] (-0.338447,-1.235655)--(-0.324025,-1.079622)--(-0.162796,-1.091580)--(-0.170043,-1.248146)--cycle;
\path[fill=blue!38!white,draw=blue!38!white,line width=.12pt] (0.170043,-1.248146)--(0.162796,-1.091580)--(0.324025,-1.079622)--(0.338447,-1.235655)--cycle;
\path[fill=blue!36!white,draw=blue!36!white,line width=.12pt] (-0.618176,-0.881367)--(-0.612293,-0.739104)--(-0.464455,-0.765552)--(-0.468918,-0.908070)--cycle;
\path[fill=blue!33!white,draw=blue!33!white,line width=.12pt] (0.468918,-0.908070)--(0.464455,-0.765552)--(0.612293,-0.739104)--(0.618176,-0.881367)--cycle;
\path[fill=blue!41!white,draw=blue!41!white,line width=.12pt] (-0.357857,-1.389432)--(-0.338447,-1.235655)--(-0.170043,-1.248146)--(-0.179794,-1.402639)--cycle;
\path[fill=blue!40!white,draw=blue!40!white,line width=.12pt] (0.179794,-1.402639)--(0.170043,-1.248146)--(0.338447,-1.235655)--(0.357857,-1.389432)--cycle;
\path[fill=blue!28!white,draw=blue!28!white,line width=.12pt] (-1.174431,-0.322085)--(-1.226706,-0.228443)--(-1.100561,-0.285608)--(-1.053662,-0.376814)--cycle;
\path[fill=blue!23!white,draw=blue!23!white,line width=.12pt] (1.053662,-0.376814)--(1.100561,-0.285608)--(1.226706,-0.228443)--(1.174431,-0.322085)--cycle;
\path[fill=blue!39!white,draw=blue!39!white,line width=.12pt] (-0.170043,-1.248146)--(-0.162796,-1.091580)--(-0.000000,-1.095579)--(-0.000000,-1.252323)--cycle;
\path[fill=blue!39!white,draw=blue!39!white,line width=.12pt] (-0.000000,-1.252323)--(-0.000000,-1.095579)--(0.162796,-1.091580)--(0.170043,-1.248146)--cycle;
\path[fill=blue!42!white,draw=blue!42!white,line width=.12pt] (-0.567664,-1.511728)--(-0.532473,-1.367563)--(-0.357857,-1.389432)--(-0.381508,-1.535043)--cycle;
\path[fill=blue!40!white,draw=blue!40!white,line width=.12pt] (0.381508,-1.535043)--(0.357857,-1.389432)--(0.532473,-1.367563)--(0.567664,-1.511728)--cycle;
\path[fill=green!65!black!37!white,draw=green!65!black!37!white,line width=.12pt] (-1.525177,-0.019647)--(-1.625975,0.032838)--(-1.511656,-0.044233)--(-1.417944,-0.091940)--cycle;
\path[fill=green!65!black!33!white,draw=green!65!black!33!white,line width=.12pt] (1.417944,-0.091940)--(1.511656,-0.044233)--(1.625975,0.032838)--(1.525177,-0.019647)--cycle;
\path[fill=blue!17!white,draw=blue!17!white,line width=.12pt] (-2.900154,-0.335511)--(-2.942213,-0.474190)--(-2.808591,-0.615449)--(-2.768442,-0.474751)--cycle;
\path[fill=blue!24!white,draw=blue!24!white,line width=.12pt] (2.768442,-0.474751)--(2.808591,-0.615449)--(2.942213,-0.474190)--(2.900154,-0.335511)--cycle;
\path[fill=blue!38!white,draw=blue!38!white,line width=.12pt] (-0.324025,-1.079622)--(-0.315143,-0.927329)--(-0.158334,-0.938959)--(-0.162796,-1.091580)--cycle;
\path[fill=blue!36!white,draw=blue!36!white,line width=.12pt] (0.162796,-1.091580)--(0.158334,-0.938959)--(0.315143,-0.927329)--(0.324025,-1.079622)--cycle;
\path[fill=blue!25!white,draw=blue!25!white,line width=.12pt] (-1.341037,-0.165370)--(-1.417944,-0.091940)--(-1.297056,-0.158630)--(-1.226706,-0.228443)--cycle;
\path[fill=blue!20!white,draw=blue!20!white,line width=.12pt] (1.226706,-0.228443)--(1.297056,-0.158630)--(1.417944,-0.091940)--(1.341037,-0.165370)--cycle;
\path[fill=blue!40!white,draw=blue!40!white,line width=.12pt] (-0.179794,-1.402639)--(-0.170043,-1.248146)--(-0.000000,-1.252323)--(-0.000000,-1.407055)--cycle;
\path[fill=blue!40!white,draw=blue!40!white,line width=.12pt] (-0.000000,-1.407055)--(-0.000000,-1.252323)--(0.170043,-1.248146)--(0.179794,-1.402639)--cycle;
\path[fill=blue!24!white,draw=blue!24!white,line width=.12pt] (-2.808591,-0.885774)--(-2.768442,-1.005012)--(-2.610069,-1.137127)--(-2.647922,-1.019805)--cycle;
\path[fill=blue!31!white,draw=blue!31!white,line width=.12pt] (2.647922,-1.019805)--(2.610069,-1.137127)--(2.768442,-1.005012)--(2.808591,-0.885774)--cycle;
\path[fill=blue!34!white,draw=blue!34!white,line width=.12pt] (-0.754235,-0.705537)--(-0.761481,-0.577153)--(-0.618176,-0.611042)--(-0.612293,-0.739104)--cycle;
\path[fill=blue!30!white,draw=blue!30!white,line width=.12pt] (0.612293,-0.739104)--(0.618176,-0.611042)--(0.761481,-0.577153)--(0.754235,-0.705537)--cycle;
\path[fill=blue!42!white,draw=blue!42!white,line width=.12pt] (-0.801283,-1.607317)--(-0.748354,-1.479402)--(-0.567664,-1.511728)--(-0.607814,-1.641929)--cycle;
\path[fill=blue!40!white,draw=blue!40!white,line width=.12pt] (0.607814,-1.641929)--(0.567664,-1.511728)--(0.748354,-1.479402)--(0.801283,-1.607317)--cycle;
\path[fill=blue!37!white,draw=blue!37!white,line width=.12pt] (-0.162796,-1.091580)--(-0.158334,-0.938959)--(-0.000000,-0.942848)--(-0.000000,-1.095579)--cycle;
\path[fill=blue!37!white,draw=blue!37!white,line width=.12pt] (-0.000000,-1.095579)--(-0.000000,-0.942848)--(0.158334,-0.938959)--(0.162796,-1.091580)--cycle;
\path[fill=blue!29!white,draw=blue!29!white,line width=.12pt] (-2.548600,-1.236371)--(-2.465877,-1.313721)--(-2.292506,-1.430603)--(-2.369413,-1.357174)--cycle;
\path[fill=blue!35!white,draw=blue!35!white,line width=.12pt] (2.369413,-1.357174)--(2.292506,-1.430603)--(2.465877,-1.313721)--(2.548600,-1.236371)--cycle;
\path[fill=blue!37!white,draw=blue!37!white,line width=.12pt] (-1.913535,-1.596850)--(-1.807416,-1.583216)--(-1.621555,-1.667443)--(-1.716761,-1.686022)--cycle;
\path[fill=blue!38!white,draw=blue!38!white,line width=.12pt] (1.716761,-1.686022)--(1.621555,-1.667443)--(1.807416,-1.583216)--(1.913535,-1.596850)--cycle;
\path[fill=blue!36!white,draw=blue!36!white,line width=.12pt] (-0.468918,-0.908070)--(-0.464455,-0.765552)--(-0.312145,-0.784628)--(-0.315143,-0.927329)--cycle;
\path[fill=blue!34!white,draw=blue!34!white,line width=.12pt] (0.315143,-0.927329)--(0.312145,-0.784628)--(0.464455,-0.765552)--(0.468918,-0.908070)--cycle;
\path[fill=blue!14!white,draw=blue!14!white,line width=.12pt] (-2.217311,0.233600)--(-2.361503,0.190425)--(-2.254254,0.077046)--(-2.116611,0.127144)--cycle;
\path[fill=blue!15!white,draw=blue!15!white,line width=.12pt] (2.116611,0.127144)--(2.254254,0.077046)--(2.361503,0.190425)--(2.217311,0.233600)--cycle;
\path[fill=blue!41!white,draw=blue!41!white,line width=.12pt] (-0.381508,-1.535043)--(-0.357857,-1.389432)--(-0.179794,-1.402639)--(-0.191677,-1.549123)--cycle;
\path[fill=blue!41!white,draw=blue!41!white,line width=.12pt] (0.191677,-1.549123)--(0.179794,-1.402639)--(0.357857,-1.389432)--(0.381508,-1.535043)--cycle;
\path[fill=blue!41!white,draw=blue!41!white,line width=.12pt] (-1.057780,-1.668992)--(-0.987036,-1.563389)--(-0.801283,-1.607317)--(-0.858714,-1.716067)--cycle;
\path[fill=blue!40!white,draw=blue!40!white,line width=.12pt] (0.858714,-1.716067)--(0.801283,-1.607317)--(0.987036,-1.563389)--(1.057780,-1.668992)--cycle;
\path[fill=blue!39!white,draw=blue!39!white,line width=.12pt] (-1.621555,-1.667443)--(-1.522544,-1.620433)--(-1.333369,-1.690584)--(-1.420078,-1.742156)--cycle;
\path[fill=blue!39!white,draw=blue!39!white,line width=.12pt] (1.420078,-1.742156)--(1.333369,-1.690584)--(1.522544,-1.620433)--(1.621555,-1.667443)--cycle;
\path[fill=blue!40!white,draw=blue!40!white,line width=.12pt] (-1.333369,-1.690584)--(-1.246659,-1.612387)--(-1.057780,-1.668992)--(-1.131352,-1.751126)--cycle;
\path[fill=blue!40!white,draw=blue!40!white,line width=.12pt] (1.131352,-1.751126)--(1.057780,-1.668992)--(1.246659,-1.612387)--(1.333369,-1.690584)--cycle;
\path[fill=blue!41!white,draw=blue!41!white,line width=.12pt] (-0.000000,-1.553831)--(-0.000000,-1.407055)--(0.179794,-1.402639)--(0.191677,-1.549123)--cycle;
\path[fill=blue!41!white,draw=blue!41!white,line width=.12pt] (-0.191677,-1.549123)--(-0.179794,-1.402639)--(-0.000000,-1.407055)--(-0.000000,-1.553831)--cycle;
\path[fill=blue!31!white,draw=blue!31!white,line width=.12pt] (-0.897452,-0.536404)--(-0.922745,-0.425362)--(-0.782941,-0.467259)--(-0.761481,-0.577153)--cycle;
\path[fill=blue!27!white,draw=blue!27!white,line width=.12pt] (0.761481,-0.577153)--(0.782941,-0.467259)--(0.922745,-0.425362)--(0.897452,-0.536404)--cycle;
\path[fill=blue!15!white,draw=blue!15!white,line width=.12pt] (-2.831853,-0.201586)--(-2.900154,-0.335511)--(-2.768442,-0.474751)--(-2.703244,-0.337547)--cycle;
\path[fill=blue!22!white,draw=blue!22!white,line width=.12pt] (2.703244,-0.337547)--(2.768442,-0.474751)--(2.900154,-0.335511)--(2.831853,-0.201586)--cycle;
\path[fill=blue!36!white,draw=blue!36!white,line width=.12pt] (-0.315143,-0.927329)--(-0.312145,-0.784628)--(-0.156827,-0.796148)--(-0.158334,-0.938959)--cycle;
\path[fill=blue!34!white,draw=blue!34!white,line width=.12pt] (0.158334,-0.938959)--(0.156827,-0.796148)--(0.312145,-0.784628)--(0.315143,-0.927329)--cycle;
\path[fill=green!65!black!32!white,draw=green!65!black!32!white,line width=.12pt] (-1.846614,0.146565)--(-1.978968,0.150618)--(-1.865758,0.056178)--(-1.740975,0.058441)--cycle;
\path[fill=green!65!black!30!white,draw=green!65!black!30!white,line width=.12pt] (1.740975,0.058441)--(1.865758,0.056178)--(1.978968,0.150618)--(1.846614,0.146565)--cycle;
\path[fill=blue!42!white,draw=blue!42!white,line width=.12pt] (-0.607814,-1.641929)--(-0.567664,-1.511728)--(-0.381508,-1.535043)--(-0.408491,-1.666893)--cycle;
\path[fill=blue!41!white,draw=blue!41!white,line width=.12pt] (0.408491,-1.666893)--(0.381508,-1.535043)--(0.567664,-1.511728)--(0.607814,-1.641929)--cycle;
\path[fill=blue!34!white,draw=blue!34!white,line width=.12pt] (-0.612293,-0.739104)--(-0.618176,-0.611042)--(-0.468918,-0.637745)--(-0.464455,-0.765552)--cycle;
\path[fill=blue!31!white,draw=blue!31!white,line width=.12pt] (0.464455,-0.765552)--(0.468918,-0.637745)--(0.618176,-0.611042)--(0.612293,-0.739104)--cycle;
\path[fill=blue!35!white,draw=blue!35!white,line width=.12pt] (-0.158334,-0.938959)--(-0.156827,-0.796148)--(-0.000000,-0.800000)--(-0.000000,-0.942848)--cycle;
\path[fill=blue!35!white,draw=blue!35!white,line width=.12pt] (-0.000000,-0.942848)--(-0.000000,-0.800000)--(0.156827,-0.796148)--(0.158334,-0.938959)--cycle;
\path[fill=blue!13!white,draw=blue!13!white,line width=.12pt] (-2.361503,0.190425)--(-2.500154,0.122284)--(-2.386608,0.002248)--(-2.254254,0.077046)--cycle;
\path[fill=blue!15!white,draw=blue!15!white,line width=.12pt] (2.254254,0.077046)--(2.386608,0.002248)--(2.500154,0.122284)--(2.361503,0.190425)--cycle;
\path[fill=blue!28!white,draw=blue!28!white,line width=.12pt] (-1.053662,-0.376814)--(-1.100561,-0.285608)--(-0.963817,-0.336316)--(-0.922745,-0.425362)--cycle;
\path[fill=blue!23!white,draw=blue!23!white,line width=.12pt] (0.922745,-0.425362)--(0.963817,-0.336316)--(1.100561,-0.285608)--(1.053662,-0.376814)--cycle;
\path[fill=blue!14!white,draw=blue!14!white,line width=.12pt] (-2.739936,-0.077561)--(-2.831853,-0.201586)--(-2.703244,-0.337547)--(-2.615501,-0.209109)--cycle;
\path[fill=blue!20!white,draw=blue!20!white,line width=.12pt] (2.615501,-0.209109)--(2.703244,-0.337547)--(2.831853,-0.201586)--(2.739936,-0.077561)--cycle;
\path[fill=blue!13!white,draw=blue!13!white,line width=.12pt] (-2.500154,0.122284)--(-2.627935,0.031796)--(-2.508586,-0.094375)--(-2.386608,0.002248)--cycle;
\path[fill=blue!16!white,draw=blue!16!white,line width=.12pt] (2.386608,0.002248)--(2.508586,-0.094375)--(2.627935,0.031796)--(2.500154,0.122284)--cycle;
\path[fill=blue!32!white,draw=blue!32!white,line width=.12pt] (-2.292506,-1.430603)--(-2.198794,-1.478311)--(-2.011334,-1.581726)--(-2.097056,-1.538426)--cycle;
\path[fill=blue!36!white,draw=blue!36!white,line width=.12pt] (2.097056,-1.538426)--(2.011334,-1.581726)--(2.198794,-1.478311)--(2.292506,-1.430603)--cycle;
\path[fill=blue!21!white,draw=blue!21!white,line width=.12pt] (-2.822148,-0.754235)--(-2.808591,-0.885774)--(-2.647922,-1.019805)--(-2.660703,-0.888912)--cycle;
\path[fill=blue!29!white,draw=blue!29!white,line width=.12pt] (2.660703,-0.888912)--(2.647922,-1.019805)--(2.808591,-0.885774)--(2.822148,-0.754235)--cycle;
\path[fill=blue!42!white,draw=blue!42!white,line width=.12pt] (-0.408491,-1.666893)--(-0.381508,-1.535043)--(-0.191677,-1.549123)--(-0.205234,-1.681968)--cycle;
\path[fill=blue!41!white,draw=blue!41!white,line width=.12pt] (0.205234,-1.681968)--(0.191677,-1.549123)--(0.381508,-1.535043)--(0.408491,-1.666893)--cycle;
\path[fill=blue!13!white,draw=blue!13!white,line width=.12pt] (-2.627935,0.031796)--(-2.739936,-0.077561)--(-2.615501,-0.209109)--(-2.508586,-0.094375)--cycle;
\path[fill=blue!18!white,draw=blue!18!white,line width=.12pt] (2.508586,-0.094375)--(2.615501,-0.209109)--(2.739936,-0.077561)--(2.627935,0.031796)--cycle;
\path[fill=blue!41!white,draw=blue!41!white,line width=.12pt] (-0.858714,-1.716067)--(-0.801283,-1.607317)--(-0.607814,-1.641929)--(-0.651378,-1.753160)--cycle;
\path[fill=blue!40!white,draw=blue!40!white,line width=.12pt] (0.651378,-1.753160)--(0.607814,-1.641929)--(0.801283,-1.607317)--(0.858714,-1.716067)--cycle;
\path[fill=blue!33!white,draw=blue!33!white,line width=.12pt] (-0.464455,-0.765552)--(-0.468918,-0.637745)--(-0.315143,-0.657004)--(-0.312145,-0.784628)--cycle;
\path[fill=blue!31!white,draw=blue!31!white,line width=.12pt] (0.312145,-0.784628)--(0.315143,-0.657004)--(0.468918,-0.637745)--(0.464455,-0.765552)--cycle;
\path[fill=green!65!black!34!white,draw=green!65!black!34!white,line width=.12pt] (-1.625975,0.032838)--(-1.740975,0.058441)--(-1.618570,-0.024081)--(-1.511656,-0.044233)--cycle;
\path[fill=green!65!black!31!white,draw=green!65!black!31!white,line width=.12pt] (1.511656,-0.044233)--(1.618570,-0.024081)--(1.740975,0.058441)--(1.625975,0.032838)--cycle;
\path[fill=blue!25!white,draw=blue!25!white,line width=.12pt] (-1.226706,-0.228443)--(-1.297056,-0.158630)--(-1.163677,-0.219074)--(-1.100561,-0.285608)--cycle;
\path[fill=blue!21!white,draw=blue!21!white,line width=.12pt] (1.100561,-0.285608)--(1.163677,-0.219074)--(1.297056,-0.158630)--(1.226706,-0.228443)--cycle;
\path[fill=blue!31!white,draw=blue!31!white,line width=.12pt] (-0.761481,-0.577153)--(-0.782941,-0.467259)--(-0.635598,-0.502103)--(-0.618176,-0.611042)--cycle;
\path[fill=blue!27!white,draw=blue!27!white,line width=.12pt] (0.618176,-0.611042)--(0.635598,-0.502103)--(0.782941,-0.467259)--(0.761481,-0.577153)--cycle;
\path[fill=blue!41!white,draw=blue!41!white,line width=.12pt] (-0.205234,-1.681968)--(-0.191677,-1.549123)--(-0.000000,-1.553831)--(-0.000000,-1.687009)--cycle;
\path[fill=blue!41!white,draw=blue!41!white,line width=.12pt] (-0.000000,-1.687009)--(-0.000000,-1.553831)--(0.191677,-1.549123)--(0.205234,-1.681968)--cycle;
\path[fill=green!65!black!37!white,draw=green!65!black!37!white,line width=.12pt] (-1.417944,-0.091940)--(-1.511656,-0.044233)--(-1.382778,-0.115330)--(-1.297056,-0.158630)--cycle;
\path[fill=green!65!black!33!white,draw=green!65!black!33!white,line width=.12pt] (1.297056,-0.158630)--(1.382778,-0.115330)--(1.511656,-0.044233)--(1.417944,-0.091940)--cycle;
\path[fill=blue!26!white,draw=blue!26!white,line width=.12pt] (-2.610069,-1.137127)--(-2.548600,-1.236371)--(-2.369413,-1.357174)--(-2.426560,-1.260844)--cycle;
\path[fill=blue!33!white,draw=blue!33!white,line width=.12pt] (2.426560,-1.260844)--(2.369413,-1.357174)--(2.548600,-1.236371)--(2.610069,-1.137127)--cycle;
\path[fill=blue!33!white,draw=blue!33!white,line width=.12pt] (-0.312145,-0.784628)--(-0.315143,-0.657004)--(-0.158334,-0.668634)--(-0.156827,-0.796148)--cycle;
\path[fill=blue!32!white,draw=blue!32!white,line width=.12pt] (0.156827,-0.796148)--(0.158334,-0.668634)--(0.315143,-0.657004)--(0.312145,-0.784628)--cycle;
\path[fill=blue!40!white,draw=blue!40!white,line width=.12pt] (-1.131352,-1.751126)--(-1.057780,-1.668992)--(-0.858714,-1.716067)--(-0.918440,-1.801476)--cycle;
\path[fill=blue!40!white,draw=blue!40!white,line width=.12pt] (0.918440,-1.801476)--(0.858714,-1.716067)--(1.057780,-1.668992)--(1.131352,-1.751126)--cycle;
\path[fill=blue!34!white,draw=blue!34!white,line width=.12pt] (-2.011334,-1.581726)--(-1.913535,-1.596850)--(-1.716761,-1.686022)--(-1.804504,-1.675456)--cycle;
\path[fill=blue!37!white,draw=blue!37!white,line width=.12pt] (1.804504,-1.675456)--(1.716761,-1.686022)--(1.913535,-1.596850)--(2.011334,-1.581726)--cycle;
\path[fill=blue!33!white,draw=blue!33!white,line width=.12pt] (-0.156827,-0.796148)--(-0.158334,-0.668634)--(-0.000000,-0.672523)--(-0.000000,-0.800000)--cycle;
\path[fill=blue!32!white,draw=blue!32!white,line width=.12pt] (-0.000000,-0.800000)--(-0.000000,-0.672523)--(0.158334,-0.668634)--(0.156827,-0.796148)--cycle;
\path[fill=green!65!black!30!white,draw=green!65!black!30!white,line width=.12pt] (-1.978968,0.150618)--(-2.116611,0.127144)--(-1.995527,0.026136)--(-1.865758,0.056178)--cycle;
\path[fill=green!65!black!29!white,draw=green!65!black!29!white,line width=.12pt] (1.865758,0.056178)--(1.995527,0.026136)--(2.116611,0.127144)--(1.978968,0.150618)--cycle;
\path[fill=blue!41!white,draw=blue!41!white,line width=.12pt] (-0.651378,-1.753160)--(-0.607814,-1.641929)--(-0.408491,-1.666893)--(-0.437769,-1.779914)--cycle;
\path[fill=blue!41!white,draw=blue!41!white,line width=.12pt] (0.437769,-1.779914)--(0.408491,-1.666893)--(0.607814,-1.641929)--(0.651378,-1.753160)--cycle;
\path[fill=blue!39!white,draw=blue!39!white,line width=.12pt] (-1.420078,-1.742156)--(-1.333369,-1.690584)--(-1.131352,-1.751126)--(-1.204924,-1.806635)--cycle;
\path[fill=blue!39!white,draw=blue!39!white,line width=.12pt] (1.204924,-1.806635)--(1.131352,-1.751126)--(1.333369,-1.690584)--(1.420078,-1.742156)--cycle;
\path[fill=blue!37!white,draw=blue!37!white,line width=.12pt] (-1.716761,-1.686022)--(-1.621555,-1.667443)--(-1.420078,-1.742156)--(-1.503455,-1.765122)--cycle;
\path[fill=blue!38!white,draw=blue!38!white,line width=.12pt] (1.503455,-1.765122)--(1.420078,-1.742156)--(1.621555,-1.667443)--(1.716761,-1.686022)--cycle;
\path[fill=blue!28!white,draw=blue!28!white,line width=.12pt] (-0.922745,-0.425362)--(-0.963817,-0.336316)--(-0.817791,-0.380079)--(-0.782941,-0.467259)--cycle;
\path[fill=blue!24!white,draw=blue!24!white,line width=.12pt] (0.782941,-0.467259)--(0.817791,-0.380079)--(0.963817,-0.336316)--(0.922745,-0.425362)--cycle;
\path[fill=blue!31!white,draw=blue!31!white,line width=.12pt] (-0.618176,-0.611042)--(-0.635598,-0.502103)--(-0.482133,-0.529559)--(-0.468918,-0.637745)--cycle;
\path[fill=blue!28!white,draw=blue!28!white,line width=.12pt] (0.468918,-0.637745)--(0.482133,-0.529559)--(0.635598,-0.502103)--(0.618176,-0.611042)--cycle;
\path[fill=blue!18!white,draw=blue!18!white,line width=.12pt] (-2.808591,-0.615449)--(-2.822148,-0.754235)--(-2.660703,-0.888912)--(-2.647922,-0.749480)--cycle;
\path[fill=blue!26!white,draw=blue!26!white,line width=.12pt] (2.647922,-0.749480)--(2.660703,-0.888912)--(2.822148,-0.754235)--(2.808591,-0.615449)--cycle;
\path[fill=blue!41!white,draw=blue!41!white,line width=.12pt] (-0.437769,-1.779914)--(-0.408491,-1.666893)--(-0.205234,-1.681968)--(-0.219943,-1.796069)--cycle;
\path[fill=blue!41!white,draw=blue!41!white,line width=.12pt] (0.219943,-1.796069)--(0.205234,-1.681968)--(0.408491,-1.666893)--(0.437769,-1.779914)--cycle;
\path[fill=blue!31!white,draw=blue!31!white,line width=.12pt] (-0.468918,-0.637745)--(-0.482133,-0.529559)--(-0.324025,-0.549360)--(-0.315143,-0.657004)--cycle;
\path[fill=blue!29!white,draw=blue!29!white,line width=.12pt] (0.315143,-0.657004)--(0.324025,-0.549360)--(0.482133,-0.529559)--(0.468918,-0.637745)--cycle;
\path[fill=blue!25!white,draw=blue!25!white,line width=.12pt] (-1.100561,-0.285608)--(-1.163677,-0.219074)--(-1.019091,-0.272690)--(-0.963817,-0.336316)--cycle;
\path[fill=blue!21!white,draw=blue!21!white,line width=.12pt] (0.963817,-0.336316)--(1.019091,-0.272690)--(1.163677,-0.219074)--(1.100561,-0.285608)--cycle;
\path[fill=blue!41!white,draw=blue!41!white,line width=.12pt] (-0.219943,-1.796069)--(-0.205234,-1.681968)--(-0.000000,-1.687009)--(-0.000000,-1.801472)--cycle;
\path[fill=blue!41!white,draw=blue!41!white,line width=.12pt] (-0.000000,-1.801472)--(-0.000000,-1.687009)--(0.205234,-1.681968)--(0.219943,-1.796069)--cycle;
\path[fill=blue!40!white,draw=blue!40!white,line width=.12pt] (-0.918440,-1.801476)--(-0.858714,-1.716067)--(-0.651378,-1.753160)--(-0.696683,-1.841149)--cycle;
\path[fill=blue!40!white,draw=blue!40!white,line width=.12pt] (0.696683,-1.841149)--(0.651378,-1.753160)--(0.858714,-1.716067)--(0.918440,-1.801476)--cycle;
\path[fill=blue!14!white,draw=blue!14!white,line width=.12pt] (-2.116611,0.127144)--(-2.254254,0.077046)--(-2.125296,-0.030531)--(-1.995527,0.026136)--cycle;
\path[fill=blue!14!white,draw=blue!14!white,line width=.12pt] (1.995527,0.026136)--(2.125296,-0.030531)--(2.254254,0.077046)--(2.116611,0.127144)--cycle;
\path[fill=green!65!black!32!white,draw=green!65!black!32!white,line width=.12pt] (-1.740975,0.058441)--(-1.865758,0.056178)--(-1.734580,-0.032258)--(-1.618570,-0.024081)--cycle;
\path[fill=green!65!black!30!white,draw=green!65!black!30!white,line width=.12pt] (1.618570,-0.024081)--(1.734580,-0.032258)--(1.865758,0.056178)--(1.740975,0.058441)--cycle;
\path[fill=blue!29!white,draw=blue!29!white,line width=.12pt] (-2.369413,-1.357174)--(-2.292506,-1.430603)--(-2.097056,-1.538426)--(-2.167407,-1.468614)--cycle;
\path[fill=blue!34!white,draw=blue!34!white,line width=.12pt] (2.167407,-1.468614)--(2.097056,-1.538426)--(2.292506,-1.430603)--(2.369413,-1.357174)--cycle;
\path[fill=blue!23!white,draw=blue!23!white,line width=.12pt] (-2.647922,-1.019805)--(-2.610069,-1.137127)--(-2.426560,-1.260844)--(-2.461751,-1.145316)--cycle;
\path[fill=blue!30!white,draw=blue!30!white,line width=.12pt] (2.461751,-1.145316)--(2.426560,-1.260844)--(2.610069,-1.137127)--(2.647922,-1.019805)--cycle;
\path[fill=blue!28!white,draw=blue!28!white,line width=.12pt] (-0.782941,-0.467259)--(-0.817791,-0.380079)--(-0.663889,-0.416474)--(-0.635598,-0.502103)--cycle;
\path[fill=blue!25!white,draw=blue!25!white,line width=.12pt] (0.635598,-0.502103)--(0.663889,-0.416474)--(0.817791,-0.380079)--(0.782941,-0.467259)--cycle;
\path[fill=blue!30!white,draw=blue!30!white,line width=.12pt] (-0.315143,-0.657004)--(-0.324025,-0.549360)--(-0.162796,-0.561318)--(-0.158334,-0.668634)--cycle;
\path[fill=blue!29!white,draw=blue!29!white,line width=.12pt] (0.158334,-0.668634)--(0.162796,-0.561318)--(0.324025,-0.549360)--(0.315143,-0.657004)--cycle;
\path[fill=blue!16!white,draw=blue!16!white,line width=.12pt] (-2.768442,-0.474751)--(-2.808591,-0.615449)--(-2.647922,-0.749480)--(-2.610069,-0.606865)--cycle;
\path[fill=blue!24!white,draw=blue!24!white,line width=.12pt] (2.610069,-0.606865)--(2.647922,-0.749480)--(2.808591,-0.615449)--(2.768442,-0.474751)--cycle;
\path[fill=green!65!black!37!white,draw=green!65!black!37!white,line width=.12pt] (-1.297056,-0.158630)--(-1.382778,-0.115330)--(-1.240584,-0.179769)--(-1.163677,-0.219074)--cycle;
\path[fill=green!65!black!34!white,draw=green!65!black!34!white,line width=.12pt] (1.163677,-0.219074)--(1.240584,-0.179769)--(1.382778,-0.115330)--(1.297056,-0.158630)--cycle;
\path[fill=green!65!black!35!white,draw=green!65!black!35!white,line width=.12pt] (-1.511656,-0.044233)--(-1.618570,-0.024081)--(-1.480578,-0.100206)--(-1.382778,-0.115330)--cycle;
\path[fill=green!65!black!31!white,draw=green!65!black!31!white,line width=.12pt] (1.382778,-0.115330)--(1.480578,-0.100206)--(1.618570,-0.024081)--(1.511656,-0.044233)--cycle;
\path[fill=blue!30!white,draw=blue!30!white,line width=.12pt] (-0.158334,-0.668634)--(-0.162796,-0.561318)--(-0.000000,-0.565317)--(-0.000000,-0.672523)--cycle;
\path[fill=blue!30!white,draw=blue!30!white,line width=.12pt] (-0.000000,-0.672523)--(-0.000000,-0.565317)--(0.162796,-0.561318)--(0.158334,-0.668634)--cycle;
\path[fill=blue!39!white,draw=blue!39!white,line width=.12pt] (-1.204924,-1.806635)--(-1.131352,-1.751126)--(-0.918440,-1.801476)--(-0.978167,-1.860259)--cycle;
\path[fill=blue!39!white,draw=blue!39!white,line width=.12pt] (0.978167,-1.860259)--(0.918440,-1.801476)--(1.131352,-1.751126)--(1.204924,-1.806635)--cycle;
\path[fill=blue!13!white,draw=blue!13!white,line width=.12pt] (-2.254254,0.077046)--(-2.386608,0.002248)--(-2.250079,-0.111645)--(-2.125296,-0.030531)--cycle;
\path[fill=blue!15!white,draw=blue!15!white,line width=.12pt] (2.125296,-0.030531)--(2.250079,-0.111645)--(2.386608,0.002248)--(2.254254,0.077046)--cycle;
\path[fill=blue!40!white,draw=blue!40!white,line width=.12pt] (-0.696683,-1.841149)--(-0.651378,-1.753160)--(-0.437769,-1.779914)--(-0.468217,-1.869763)--cycle;
\path[fill=blue!40!white,draw=blue!40!white,line width=.12pt] (0.468217,-1.869763)--(0.437769,-1.779914)--(0.651378,-1.753160)--(0.696683,-1.841149)--cycle;
\path[fill=blue!32!white,draw=blue!32!white,line width=.12pt] (-2.097056,-1.538426)--(-2.011334,-1.581726)--(-1.804504,-1.675456)--(-1.881411,-1.636151)--cycle;
\path[fill=blue!35!white,draw=blue!35!white,line width=.12pt] (1.881411,-1.636151)--(1.804504,-1.675456)--(2.011334,-1.581726)--(2.097056,-1.538426)--cycle;
\path[fill=blue!15!white,draw=blue!15!white,line width=.12pt] (-2.703244,-0.337547)--(-2.768442,-0.474751)--(-2.610069,-0.606865)--(-2.548600,-0.466550)--cycle;
\path[fill=blue!21!white,draw=blue!21!white,line width=.12pt] (2.548600,-0.466550)--(2.610069,-0.606865)--(2.768442,-0.474751)--(2.703244,-0.337547)--cycle;
\path[fill=blue!28!white,draw=blue!28!white,line width=.12pt] (-0.635598,-0.502103)--(-0.663889,-0.416474)--(-0.503593,-0.445151)--(-0.482133,-0.529559)--cycle;
\path[fill=blue!25!white,draw=blue!25!white,line width=.12pt] (0.482133,-0.529559)--(0.503593,-0.445151)--(0.663889,-0.416474)--(0.635598,-0.502103)--cycle;
\path[fill=blue!37!white,draw=blue!37!white,line width=.12pt] (-1.503455,-1.765122)--(-1.420078,-1.742156)--(-1.204924,-1.806635)--(-1.275669,-1.833387)--cycle;
\path[fill=blue!38!white,draw=blue!38!white,line width=.12pt] (1.275669,-1.833387)--(1.204924,-1.806635)--(1.420078,-1.742156)--(1.503455,-1.765122)--cycle;
\path[fill=blue!25!white,draw=blue!25!white,line width=.12pt] (-0.963817,-0.336316)--(-1.019091,-0.272690)--(-0.864690,-0.318963)--(-0.817791,-0.380079)--cycle;
\path[fill=blue!22!white,draw=blue!22!white,line width=.12pt] (0.817791,-0.380079)--(0.864690,-0.318963)--(1.019091,-0.272690)--(0.963817,-0.336316)--cycle;
\path[fill=blue!34!white,draw=blue!34!white,line width=.12pt] (-1.804504,-1.675456)--(-1.716761,-1.686022)--(-1.503455,-1.765122)--(-1.580295,-1.758598)--cycle;
\path[fill=blue!37!white,draw=blue!37!white,line width=.12pt] (1.580295,-1.758598)--(1.503455,-1.765122)--(1.716761,-1.686022)--(1.804504,-1.675456)--cycle;
\path[fill=blue!12!white,draw=blue!12!white,line width=.12pt] (-2.386608,0.002248)--(-2.508586,-0.094375)--(-2.365079,-0.214089)--(-2.250079,-0.111645)--cycle;
\path[fill=blue!16!white,draw=blue!16!white,line width=.12pt] (2.250079,-0.111645)--(2.365079,-0.214089)--(2.508586,-0.094375)--(2.386608,0.002248)--cycle;
\path[fill=blue!40!white,draw=blue!40!white,line width=.12pt] (-0.468217,-1.869763)--(-0.437769,-1.779914)--(-0.219943,-1.796069)--(-0.235241,-1.887042)--cycle;
\path[fill=blue!40!white,draw=blue!40!white,line width=.12pt] (0.235241,-1.887042)--(0.219943,-1.796069)--(0.437769,-1.779914)--(0.468217,-1.869763)--cycle;
\path[fill=green!65!black!30!white,draw=green!65!black!30!white,line width=.12pt] (-1.865758,0.056178)--(-1.995527,0.026136)--(-1.855225,-0.068452)--(-1.734580,-0.032258)--cycle;
\path[fill=green!65!black!29!white,draw=green!65!black!29!white,line width=.12pt] (1.734580,-0.032258)--(1.855225,-0.068452)--(1.995527,0.026136)--(1.865758,0.056178)--cycle;
\path[fill=blue!13!white,draw=blue!13!white,line width=.12pt] (-2.615501,-0.209109)--(-2.703244,-0.337547)--(-2.548600,-0.466550)--(-2.465877,-0.333925)--cycle;
\path[fill=blue!19!white,draw=blue!19!white,line width=.12pt] (2.465877,-0.333925)--(2.548600,-0.466550)--(2.703244,-0.337547)--(2.615501,-0.209109)--cycle;
\path[fill=blue!13!white,draw=blue!13!white,line width=.12pt] (-2.508586,-0.094375)--(-2.615501,-0.209109)--(-2.465877,-0.333925)--(-2.365079,-0.214089)--cycle;
\path[fill=blue!17!white,draw=blue!17!white,line width=.12pt] (2.365079,-0.214089)--(2.465877,-0.333925)--(2.615501,-0.209109)--(2.508586,-0.094375)--cycle;
\path[fill=blue!21!white,draw=blue!21!white,line width=.12pt] (-2.660703,-0.888912)--(-2.647922,-1.019805)--(-2.461751,-1.145316)--(-2.473633,-1.015029)--cycle;
\path[fill=blue!28!white,draw=blue!28!white,line width=.12pt] (2.473633,-1.015029)--(2.461751,-1.145316)--(2.647922,-1.019805)--(2.660703,-0.888912)--cycle;
\path[fill=blue!40!white,draw=blue!40!white,line width=.12pt] (-0.235241,-1.887042)--(-0.219943,-1.796069)--(-0.000000,-1.801472)--(-0.000000,-1.892820)--cycle;
\path[fill=blue!40!white,draw=blue!40!white,line width=.12pt] (-0.000000,-1.892820)--(-0.000000,-1.801472)--(0.219943,-1.796069)--(0.235241,-1.887042)--cycle;
\path[fill=blue!28!white,draw=blue!28!white,line width=.12pt] (-0.482133,-0.529559)--(-0.503593,-0.445151)--(-0.338447,-0.465835)--(-0.324025,-0.549360)--cycle;
\path[fill=blue!26!white,draw=blue!26!white,line width=.12pt] (0.324025,-0.549360)--(0.338447,-0.465835)--(0.503593,-0.445151)--(0.482133,-0.529559)--cycle;
\path[fill=blue!26!white,draw=blue!26!white,line width=.12pt] (-2.426560,-1.260844)--(-2.369413,-1.357174)--(-2.167407,-1.468614)--(-2.219681,-1.374972)--cycle;
\path[fill=blue!32!white,draw=blue!32!white,line width=.12pt] (2.219681,-1.374972)--(2.167407,-1.468614)--(2.369413,-1.357174)--(2.426560,-1.260844)--cycle;
\path[fill=blue!39!white,draw=blue!39!white,line width=.12pt] (-0.978167,-1.860259)--(-0.918440,-1.801476)--(-0.696683,-1.841149)--(-0.741989,-1.902512)--cycle;
\path[fill=blue!39!white,draw=blue!39!white,line width=.12pt] (0.741989,-1.902512)--(0.696683,-1.841149)--(0.918440,-1.801476)--(0.978167,-1.860259)--cycle;
\path[fill=green!65!black!37!white,draw=green!65!black!37!white,line width=.12pt] (-1.163677,-0.219074)--(-1.240584,-0.179769)--(-1.086442,-0.236929)--(-1.019091,-0.272690)--cycle;
\path[fill=green!65!black!34!white,draw=green!65!black!34!white,line width=.12pt] (1.019091,-0.272690)--(1.086442,-0.236929)--(1.240584,-0.179769)--(1.163677,-0.219074)--cycle;
\path[fill=blue!28!white,draw=blue!28!white,line width=.12pt] (-0.324025,-0.549360)--(-0.338447,-0.465835)--(-0.170043,-0.478325)--(-0.162796,-0.561318)--cycle;
\path[fill=blue!26!white,draw=blue!26!white,line width=.12pt] (0.162796,-0.561318)--(0.170043,-0.478325)--(0.338447,-0.465835)--(0.324025,-0.549360)--cycle;
\path[fill=green!65!black!32!white,draw=green!65!black!32!white,line width=.12pt] (-1.618570,-0.024081)--(-1.734580,-0.032258)--(-1.586697,-0.113840)--(-1.480578,-0.100206)--cycle;
\path[fill=green!65!black!30!white,draw=green!65!black!30!white,line width=.12pt] (1.480578,-0.100206)--(1.586697,-0.113840)--(1.734580,-0.032258)--(1.618570,-0.024081)--cycle;
\path[fill=blue!25!white,draw=blue!25!white,line width=.12pt] (-0.817791,-0.380079)--(-0.864690,-0.318963)--(-0.701962,-0.357445)--(-0.663889,-0.416474)--cycle;
\path[fill=blue!22!white,draw=blue!22!white,line width=.12pt] (0.663889,-0.416474)--(0.701962,-0.357445)--(0.864690,-0.318963)--(0.817791,-0.380079)--cycle;
\path[fill=green!65!black!35!white,draw=green!65!black!35!white,line width=.12pt] (-1.382778,-0.115330)--(-1.480578,-0.100206)--(-1.328326,-0.169203)--(-1.240584,-0.179769)--cycle;
\path[fill=green!65!black!32!white,draw=green!65!black!32!white,line width=.12pt] (1.240584,-0.179769)--(1.328326,-0.169203)--(1.480578,-0.100206)--(1.382778,-0.115330)--cycle;
\path[fill=blue!27!white,draw=blue!27!white,line width=.12pt] (-0.162796,-0.561318)--(-0.170043,-0.478325)--(-0.000000,-0.482502)--(-0.000000,-0.565317)--cycle;
\path[fill=blue!27!white,draw=blue!27!white,line width=.12pt] (-0.000000,-0.565317)--(-0.000000,-0.482502)--(0.170043,-0.478325)--(0.162796,-0.561318)--cycle;
\path[fill=blue!14!white,draw=blue!14!white,line width=.12pt] (-1.995527,0.026136)--(-2.125296,-0.030531)--(-1.975871,-0.131270)--(-1.855225,-0.068452)--cycle;
\path[fill=blue!14!white,draw=blue!14!white,line width=.12pt] (1.855225,-0.068452)--(1.975871,-0.131270)--(2.125296,-0.030531)--(1.995527,0.026136)--cycle;
\path[fill=blue!37!white,draw=blue!37!white,line width=.12pt] (-1.275669,-1.833387)--(-1.204924,-1.806635)--(-0.978167,-1.860259)--(-1.035598,-1.890159)--cycle;
\path[fill=blue!38!white,draw=blue!38!white,line width=.12pt] (1.035598,-1.890159)--(0.978167,-1.860259)--(1.204924,-1.806635)--(1.275669,-1.833387)--cycle;
\path[fill=blue!29!white,draw=blue!29!white,line width=.12pt] (-2.167407,-1.468614)--(-2.097056,-1.538426)--(-1.881411,-1.636151)--(-1.944527,-1.569617)--cycle;
\path[fill=blue!33!white,draw=blue!33!white,line width=.12pt] (1.944527,-1.569617)--(1.881411,-1.636151)--(2.097056,-1.538426)--(2.167407,-1.468614)--cycle;
\path[fill=blue!39!white,draw=blue!39!white,line width=.12pt] (-0.741989,-1.902512)--(-0.696683,-1.841149)--(-0.468217,-1.869763)--(-0.498665,-1.932987)--cycle;
\path[fill=blue!39!white,draw=blue!39!white,line width=.12pt] (0.498665,-1.932987)--(0.468217,-1.869763)--(0.696683,-1.841149)--(0.741989,-1.902512)--cycle;
\path[fill=blue!18!white,draw=blue!18!white,line width=.12pt] (-2.647922,-0.749480)--(-2.660703,-0.888912)--(-2.473633,-1.015029)--(-2.461751,-0.874991)--cycle;
\path[fill=blue!25!white,draw=blue!25!white,line width=.12pt] (2.461751,-0.874991)--(2.473633,-1.015029)--(2.660703,-0.888912)--(2.647922,-0.749480)--cycle;
\path[fill=blue!25!white,draw=blue!25!white,line width=.12pt] (-0.663889,-0.416474)--(-0.701962,-0.357445)--(-0.532473,-0.387767)--(-0.503593,-0.445151)--cycle;
\path[fill=blue!23!white,draw=blue!23!white,line width=.12pt] (0.503593,-0.445151)--(0.532473,-0.387767)--(0.701962,-0.357445)--(0.663889,-0.416474)--cycle;
\path[fill=blue!34!white,draw=blue!34!white,line width=.12pt] (-1.580295,-1.758598)--(-1.503455,-1.765122)--(-1.275669,-1.833387)--(-1.340867,-1.830352)--cycle;
\path[fill=blue!36!white,draw=blue!36!white,line width=.12pt] (1.340867,-1.830352)--(1.275669,-1.833387)--(1.503455,-1.765122)--(1.580295,-1.758598)--cycle;
\path[fill=blue!32!white,draw=blue!32!white,line width=.12pt] (-1.881411,-1.636151)--(-1.804504,-1.675456)--(-1.580295,-1.758598)--(-1.647647,-1.722837)--cycle;
\path[fill=blue!35!white,draw=blue!35!white,line width=.12pt] (1.647647,-1.722837)--(1.580295,-1.758598)--(1.804504,-1.675456)--(1.881411,-1.636151)--cycle;
\path[fill=green!65!black!37!white,draw=green!65!black!37!white,line width=.12pt] (-1.019091,-0.272690)--(-1.086442,-0.236929)--(-0.921837,-0.286259)--(-0.864690,-0.318963)--cycle;
\path[fill=green!65!black!34!white,draw=green!65!black!34!white,line width=.12pt] (0.864690,-0.318963)--(0.921837,-0.286259)--(1.086442,-0.236929)--(1.019091,-0.272690)--cycle;
\path[fill=blue!23!white,draw=blue!23!white,line width=.12pt] (-2.461751,-1.145316)--(-2.426560,-1.260844)--(-2.219681,-1.374972)--(-2.251872,-1.261099)--cycle;
\path[fill=blue!30!white,draw=blue!30!white,line width=.12pt] (2.251872,-1.261099)--(2.219681,-1.374972)--(2.426560,-1.260844)--(2.461751,-1.145316)--cycle;
\path[fill=blue!39!white,draw=blue!39!white,line width=.12pt] (-0.498665,-1.932987)--(-0.468217,-1.869763)--(-0.235241,-1.887042)--(-0.250539,-1.951390)--cycle;
\path[fill=blue!39!white,draw=blue!39!white,line width=.12pt] (0.250539,-1.951390)--(0.235241,-1.887042)--(0.468217,-1.869763)--(0.498665,-1.932987)--cycle;
\path[fill=blue!13!white,draw=blue!13!white,line width=.12pt] (-2.125296,-0.030531)--(-2.250079,-0.111645)--(-2.091880,-0.218298)--(-1.975871,-0.131270)--cycle;
\path[fill=blue!15!white,draw=blue!15!white,line width=.12pt] (1.975871,-0.131270)--(2.091880,-0.218298)--(2.250079,-0.111645)--(2.125296,-0.030531)--cycle;
\path[fill=blue!39!white,draw=blue!39!white,line width=.12pt] (-0.250539,-1.951390)--(-0.235241,-1.887042)--(-0.000000,-1.892820)--(-0.000000,-1.957544)--cycle;
\path[fill=blue!39!white,draw=blue!39!white,line width=.12pt] (-0.000000,-1.957544)--(-0.000000,-1.892820)--(0.235241,-1.887042)--(0.250539,-1.951390)--cycle;
\path[fill=blue!16!white,draw=blue!16!white,line width=.12pt] (-2.610069,-0.606865)--(-2.647922,-0.749480)--(-2.461751,-0.874991)--(-2.426560,-0.730582)--cycle;
\path[fill=blue!23!white,draw=blue!23!white,line width=.12pt] (2.426560,-0.730582)--(2.461751,-0.874991)--(2.647922,-0.749480)--(2.610069,-0.606865)--cycle;
\path[fill=green!65!black!30!white,draw=green!65!black!30!white,line width=.12pt] (-1.734580,-0.032258)--(-1.855225,-0.068452)--(-1.697056,-0.155708)--(-1.586697,-0.113840)--cycle;
\path[fill=green!65!black!30!white,draw=green!65!black!30!white,line width=.12pt] (1.586697,-0.113840)--(1.697056,-0.155708)--(1.855225,-0.068452)--(1.734580,-0.032258)--cycle;
\path[fill=blue!25!white,draw=blue!25!white,line width=.12pt] (-0.503593,-0.445151)--(-0.532473,-0.387767)--(-0.357857,-0.409636)--(-0.338447,-0.465835)--cycle;
\path[fill=blue!23!white,draw=blue!23!white,line width=.12pt] (0.338447,-0.465835)--(0.357857,-0.409636)--(0.532473,-0.387767)--(0.503593,-0.445151)--cycle;
\path[fill=green!65!black!35!white,draw=green!65!black!35!white,line width=.12pt] (-1.240584,-0.179769)--(-1.328326,-0.169203)--(-1.163283,-0.230405)--(-1.086442,-0.236929)--cycle;
\path[fill=green!65!black!32!white,draw=green!65!black!32!white,line width=.12pt] (1.086442,-0.236929)--(1.163283,-0.230405)--(1.328326,-0.169203)--(1.240584,-0.179769)--cycle;
\path[fill=blue!37!white,draw=blue!37!white,line width=.12pt] (-1.035598,-1.890159)--(-0.978167,-1.860259)--(-0.741989,-1.902512)--(-0.785553,-1.934893)--cycle;
\path[fill=blue!38!white,draw=blue!38!white,line width=.12pt] (0.785553,-1.934893)--(0.741989,-1.902512)--(0.978167,-1.860259)--(1.035598,-1.890159)--cycle;
\path[fill=green!65!black!32!white,draw=green!65!black!32!white,line width=.12pt] (-1.480578,-0.100206)--(-1.586697,-0.113840)--(-1.423533,-0.187782)--(-1.328326,-0.169203)--cycle;
\path[fill=green!65!black!31!white,draw=green!65!black!31!white,line width=.12pt] (1.328326,-0.169203)--(1.423533,-0.187782)--(1.586697,-0.113840)--(1.480578,-0.100206)--cycle;
\path[fill=blue!25!white,draw=blue!25!white,line width=.12pt] (-0.338447,-0.465835)--(-0.357857,-0.409636)--(-0.179794,-0.422843)--(-0.170043,-0.478325)--cycle;
\path[fill=blue!24!white,draw=blue!24!white,line width=.12pt] (0.170043,-0.478325)--(0.179794,-0.422843)--(0.357857,-0.409636)--(0.338447,-0.465835)--cycle;
\path[fill=blue!12!white,draw=blue!12!white,line width=.12pt] (-2.250079,-0.111645)--(-2.365079,-0.214089)--(-2.198794,-0.326193)--(-2.091880,-0.218298)--cycle;
\path[fill=blue!16!white,draw=blue!16!white,line width=.12pt] (2.091880,-0.218298)--(2.198794,-0.326193)--(2.365079,-0.214089)--(2.250079,-0.111645)--cycle;
\path[fill=blue!14!white,draw=blue!14!white,line width=.12pt] (-2.548600,-0.466550)--(-2.610069,-0.606865)--(-2.426560,-0.730582)--(-2.369413,-0.587353)--cycle;
\path[fill=blue!21!white,draw=blue!21!white,line width=.12pt] (2.369413,-0.587353)--(2.426560,-0.730582)--(2.610069,-0.606865)--(2.548600,-0.466550)--cycle;
\path[fill=green!65!black!37!white,draw=green!65!black!37!white,line width=.12pt] (-0.864690,-0.318963)--(-0.921837,-0.286259)--(-0.748354,-0.327284)--(-0.701962,-0.357445)--cycle;
\path[fill=green!65!black!35!white,draw=green!65!black!35!white,line width=.12pt] (0.701962,-0.357445)--(0.748354,-0.327284)--(0.921837,-0.286259)--(0.864690,-0.318963)--cycle;
\path[fill=blue!24!white,draw=blue!24!white,line width=.12pt] (-0.170043,-0.478325)--(-0.179794,-0.422843)--(-0.000000,-0.427259)--(-0.000000,-0.482502)--cycle;
\path[fill=blue!24!white,draw=blue!24!white,line width=.12pt] (-0.000000,-0.482502)--(-0.000000,-0.427259)--(0.179794,-0.422843)--(0.170043,-0.478325)--cycle;
\path[fill=blue!26!white,draw=blue!26!white,line width=.12pt] (-2.219681,-1.374972)--(-2.167407,-1.468614)--(-1.944527,-1.569617)--(-1.991426,-1.478411)--cycle;
\path[fill=blue!31!white,draw=blue!31!white,line width=.12pt] (1.991426,-1.478411)--(1.944527,-1.569617)--(2.167407,-1.468614)--(2.219681,-1.374972)--cycle;
\path[fill=blue!12!white,draw=blue!12!white,line width=.12pt] (-2.365079,-0.214089)--(-2.465877,-0.333925)--(-2.292506,-0.450808)--(-2.198794,-0.326193)--cycle;
\path[fill=blue!17!white,draw=blue!17!white,line width=.12pt] (2.198794,-0.326193)--(2.292506,-0.450808)--(2.465877,-0.333925)--(2.365079,-0.214089)--cycle;
\path[fill=blue!13!white,draw=blue!13!white,line width=.12pt] (-2.465877,-0.333925)--(-2.548600,-0.466550)--(-2.369413,-0.587353)--(-2.292506,-0.450808)--cycle;
\path[fill=blue!19!white,draw=blue!19!white,line width=.12pt] (2.292506,-0.450808)--(2.369413,-0.587353)--(2.548600,-0.466550)--(2.465877,-0.333925)--cycle;
\path[fill=blue!20!white,draw=blue!20!white,line width=.12pt] (-2.473633,-1.015029)--(-2.461751,-1.145316)--(-2.251872,-1.261099)--(-2.262742,-1.131371)--cycle;
\path[fill=blue!27!white,draw=blue!27!white,line width=.12pt] (2.262742,-1.131371)--(2.251872,-1.261099)--(2.461751,-1.145316)--(2.473633,-1.015029)--cycle;
\path[fill=blue!34!white,draw=blue!34!white,line width=.12pt] (-1.340867,-1.830352)--(-1.275669,-1.833387)--(-1.035598,-1.890159)--(-1.088526,-1.890026)--cycle;
\path[fill=blue!36!white,draw=blue!36!white,line width=.12pt] (1.088526,-1.890026)--(1.035598,-1.890159)--(1.275669,-1.833387)--(1.340867,-1.830352)--cycle;
\path[fill=blue!37!white,draw=blue!37!white,line width=.12pt] (-0.785553,-1.934893)--(-0.741989,-1.902512)--(-0.498665,-1.932987)--(-0.527943,-1.967157)--cycle;
\path[fill=blue!37!white,draw=blue!37!white,line width=.12pt] (0.527943,-1.967157)--(0.498665,-1.932987)--(0.741989,-1.902512)--(0.785553,-1.934893)--cycle;
\path[fill=blue!14!white,draw=blue!14!white,line width=.12pt] (-1.855225,-0.068452)--(-1.975871,-0.131270)--(-1.807416,-0.224200)--(-1.697056,-0.155708)--cycle;
\path[fill=blue!14!white,draw=blue!14!white,line width=.12pt] (1.697056,-0.155708)--(1.807416,-0.224200)--(1.975871,-0.131270)--(1.855225,-0.068452)--cycle;
\path[fill=blue!29!white,draw=blue!29!white,line width=.12pt] (-1.944527,-1.569617)--(-1.881411,-1.636151)--(-1.647647,-1.722837)--(-1.702920,-1.659211)--cycle;
\path[fill=blue!33!white,draw=blue!33!white,line width=.12pt] (1.702920,-1.659211)--(1.647647,-1.722837)--(1.881411,-1.636151)--(1.944527,-1.569617)--cycle;
\path[fill=blue!32!white,draw=blue!32!white,line width=.12pt] (-1.647647,-1.722837)--(-1.580295,-1.758598)--(-1.340867,-1.830352)--(-1.398014,-1.797648)--cycle;
\path[fill=blue!35!white,draw=blue!35!white,line width=.12pt] (1.398014,-1.797648)--(1.340867,-1.830352)--(1.580295,-1.758598)--(1.647647,-1.722837)--cycle;
\path[fill=green!65!black!37!white,draw=green!65!black!37!white,line width=.12pt] (-0.701962,-0.357445)--(-0.748354,-0.327284)--(-0.567664,-0.359610)--(-0.532473,-0.387767)--cycle;
\path[fill=green!65!black!35!white,draw=green!65!black!35!white,line width=.12pt] (0.532473,-0.387767)--(0.567664,-0.359610)--(0.748354,-0.327284)--(0.701962,-0.357445)--cycle;
\path[fill=green!65!black!35!white,draw=green!65!black!35!white,line width=.12pt] (-1.086442,-0.236929)--(-1.163283,-0.230405)--(-0.987036,-0.283224)--(-0.921837,-0.286259)--cycle;
\path[fill=green!65!black!33!white,draw=green!65!black!33!white,line width=.12pt] (0.921837,-0.286259)--(0.987036,-0.283224)--(1.163283,-0.230405)--(1.086442,-0.236929)--cycle;
\path[fill=blue!37!white,draw=blue!37!white,line width=.12pt] (-0.527943,-1.967157)--(-0.498665,-1.932987)--(-0.250539,-1.951390)--(-0.265249,-1.986640)--cycle;
\path[fill=blue!37!white,draw=blue!37!white,line width=.12pt] (0.265249,-1.986640)--(0.250539,-1.951390)--(0.498665,-1.932987)--(0.527943,-1.967157)--cycle;
\path[fill=green!65!black!30!white,draw=green!65!black!30!white,line width=.12pt] (-1.586697,-0.113840)--(-1.697056,-0.155708)--(-1.522544,-0.234792)--(-1.423533,-0.187782)--cycle;
\path[fill=green!65!black!30!white,draw=green!65!black!30!white,line width=.12pt] (1.423533,-0.187782)--(1.522544,-0.234792)--(1.697056,-0.155708)--(1.586697,-0.113840)--cycle;
\path[fill=green!65!black!32!white,draw=green!65!black!32!white,line width=.12pt] (-1.328326,-0.169203)--(-1.423533,-0.187782)--(-1.246659,-0.253371)--(-1.163283,-0.230405)--cycle;
\path[fill=green!65!black!31!white,draw=green!65!black!31!white,line width=.12pt] (1.163283,-0.230405)--(1.246659,-0.253371)--(1.423533,-0.187782)--(1.328326,-0.169203)--cycle;
\path[fill=blue!37!white,draw=blue!37!white,line width=.12pt] (-0.265249,-1.986640)--(-0.250539,-1.951390)--(-0.000000,-1.957544)--(-0.000000,-1.993156)--cycle;
\path[fill=blue!37!white,draw=blue!37!white,line width=.12pt] (-0.000000,-1.993156)--(-0.000000,-1.957544)--(0.250539,-1.951390)--(0.265249,-1.986640)--cycle;
\path[fill=blue!18!white,draw=blue!18!white,line width=.12pt] (-2.461751,-0.874991)--(-2.473633,-1.015029)--(-2.262742,-1.131371)--(-2.251872,-0.990774)--cycle;
\path[fill=blue!25!white,draw=blue!25!white,line width=.12pt] (2.251872,-0.990774)--(2.262742,-1.131371)--(2.473633,-1.015029)--(2.461751,-0.874991)--cycle;
\path[fill=green!65!black!37!white,draw=green!65!black!37!white,line width=.12pt] (-0.532473,-0.387767)--(-0.567664,-0.359610)--(-0.381508,-0.382925)--(-0.357857,-0.409636)--cycle;
\path[fill=green!65!black!36!white,draw=green!65!black!36!white,line width=.12pt] (0.357857,-0.409636)--(0.381508,-0.382925)--(0.567664,-0.359610)--(0.532473,-0.387767)--cycle;
\path[fill=blue!34!white,draw=blue!34!white,line width=.12pt] (-1.088526,-1.890026)--(-1.035598,-1.890159)--(-0.785553,-1.934893)--(-0.825702,-1.937046)--cycle;
\path[fill=blue!36!white,draw=blue!36!white,line width=.12pt] (0.825702,-1.937046)--(0.785553,-1.934893)--(1.035598,-1.890159)--(1.088526,-1.890026)--cycle;
\path[fill=blue!13!white,draw=blue!13!white,line width=.12pt] (-1.975871,-0.131270)--(-2.091880,-0.218298)--(-1.913535,-0.316685)--(-1.807416,-0.224200)--cycle;
\path[fill=blue!15!white,draw=blue!15!white,line width=.12pt] (1.807416,-0.224200)--(1.913535,-0.316685)--(2.091880,-0.218298)--(1.975871,-0.131270)--cycle;
\path[fill=blue!23!white,draw=blue!23!white,line width=.12pt] (-2.251872,-1.261099)--(-2.219681,-1.374972)--(-1.991426,-1.478411)--(-2.020307,-1.366038)--cycle;
\path[fill=blue!29!white,draw=blue!29!white,line width=.12pt] (2.020307,-1.366038)--(1.991426,-1.478411)--(2.219681,-1.374972)--(2.251872,-1.261099)--cycle;
\path[fill=green!65!black!37!white,draw=green!65!black!37!white,line width=.12pt] (-0.357857,-0.409636)--(-0.381508,-0.382925)--(-0.191677,-0.397005)--(-0.179794,-0.422843)--cycle;
\path[fill=green!65!black!36!white,draw=green!65!black!36!white,line width=.12pt] (0.179794,-0.422843)--(0.191677,-0.397005)--(0.381508,-0.382925)--(0.357857,-0.409636)--cycle;
\path[fill=green!65!black!35!white,draw=green!65!black!35!white,line width=.12pt] (-0.921837,-0.286259)--(-0.987036,-0.283224)--(-0.801283,-0.327152)--(-0.748354,-0.327284)--cycle;
\path[fill=green!65!black!33!white,draw=green!65!black!33!white,line width=.12pt] (0.748354,-0.327284)--(0.801283,-0.327152)--(0.987036,-0.283224)--(0.921837,-0.286259)--cycle;
\path[fill=green!65!black!37!white,draw=green!65!black!37!white,line width=.12pt] (-0.179794,-0.422843)--(-0.191677,-0.397005)--(-0.000000,-0.401713)--(-0.000000,-0.427259)--cycle;
\path[fill=green!65!black!36!white,draw=green!65!black!36!white,line width=.12pt] (-0.000000,-0.427259)--(-0.000000,-0.401713)--(0.191677,-0.397005)--(0.179794,-0.422843)--cycle;
\path[fill=blue!16!white,draw=blue!16!white,line width=.12pt] (-2.426560,-0.730582)--(-2.461751,-0.874991)--(-2.251872,-0.990774)--(-2.219681,-0.844710)--cycle;
\path[fill=blue!22!white,draw=blue!22!white,line width=.12pt] (2.219681,-0.844710)--(2.251872,-0.990774)--(2.461751,-0.874991)--(2.426560,-0.730582)--cycle;
\path[fill=blue!12!white,draw=blue!12!white,line width=.12pt] (-2.091880,-0.218298)--(-2.198794,-0.326193)--(-2.011334,-0.429608)--(-1.913535,-0.316685)--cycle;
\path[fill=blue!15!white,draw=blue!15!white,line width=.12pt] (1.913535,-0.316685)--(2.011334,-0.429608)--(2.198794,-0.326193)--(2.091880,-0.218298)--cycle;
\path[fill=blue!32!white,draw=blue!32!white,line width=.12pt] (-1.398014,-1.797648)--(-1.340867,-1.830352)--(-1.088526,-1.890026)--(-1.134919,-1.859866)--cycle;
\path[fill=blue!34!white,draw=blue!34!white,line width=.12pt] (1.134919,-1.859866)--(1.088526,-1.890026)--(1.340867,-1.830352)--(1.398014,-1.797648)--cycle;
\path[fill=blue!26!white,draw=blue!26!white,line width=.12pt] (-1.991426,-1.478411)--(-1.944527,-1.569617)--(-1.702920,-1.659211)--(-1.743993,-1.570165)--cycle;
\path[fill=blue!31!white,draw=blue!31!white,line width=.12pt] (1.743993,-1.570165)--(1.702920,-1.659211)--(1.944527,-1.569617)--(1.991426,-1.478411)--cycle;
\path[fill=blue!35!white,draw=blue!35!white,line width=.12pt] (-0.825702,-1.937046)--(-0.785553,-1.934893)--(-0.527943,-1.967157)--(-0.554926,-1.970959)--cycle;
\path[fill=blue!36!white,draw=blue!36!white,line width=.12pt] (0.554926,-1.970959)--(0.527943,-1.967157)--(0.785553,-1.934893)--(0.825702,-1.937046)--cycle;
\path[fill=blue!14!white,draw=blue!14!white,line width=.12pt] (-1.697056,-0.155708)--(-1.807416,-0.224200)--(-1.621555,-0.308427)--(-1.522544,-0.234792)--cycle;
\path[fill=blue!14!white,draw=blue!14!white,line width=.12pt] (1.522544,-0.234792)--(1.621555,-0.308427)--(1.807416,-0.224200)--(1.697056,-0.155708)--cycle;
\path[fill=blue!29!white,draw=blue!29!white,line width=.12pt] (-1.702920,-1.659211)--(-1.647647,-1.722837)--(-1.398014,-1.797648)--(-1.444914,-1.736532)--cycle;
\path[fill=blue!32!white,draw=blue!32!white,line width=.12pt] (1.444914,-1.736532)--(1.398014,-1.797648)--(1.647647,-1.722837)--(1.702920,-1.659211)--cycle;
\path[fill=blue!14!white,draw=blue!14!white,line width=.12pt] (-2.369413,-0.587353)--(-2.426560,-0.730582)--(-2.219681,-0.844710)--(-2.167407,-0.698793)--cycle;
\path[fill=blue!20!white,draw=blue!20!white,line width=.12pt] (2.167407,-0.698793)--(2.219681,-0.844710)--(2.426560,-0.730582)--(2.369413,-0.587353)--cycle;
\path[fill=green!65!black!32!white,draw=green!65!black!32!white,line width=.12pt] (-1.163283,-0.230405)--(-1.246659,-0.253371)--(-1.057780,-0.309976)--(-0.987036,-0.283224)--cycle;
\path[fill=green!65!black!31!white,draw=green!65!black!31!white,line width=.12pt] (0.987036,-0.283224)--(1.057780,-0.309976)--(1.246659,-0.253371)--(1.163283,-0.230405)--cycle;
\path[fill=blue!12!white,draw=blue!12!white,line width=.12pt] (-2.198794,-0.326193)--(-2.292506,-0.450808)--(-2.097056,-0.558630)--(-2.011334,-0.429608)--cycle;
\path[fill=blue!16!white,draw=blue!16!white,line width=.12pt] (2.011334,-0.429608)--(2.097056,-0.558630)--(2.292506,-0.450808)--(2.198794,-0.326193)--cycle;
\path[fill=green!65!black!33!white,draw=green!65!black!33!white,line width=.12pt] (0.567664,-0.359610)--(0.607814,-0.361764)--(0.801283,-0.327152)--(0.748354,-0.327284)--cycle;
\path[fill=green!65!black!35!white,draw=green!65!black!35!white,line width=.12pt] (-0.748354,-0.327284)--(-0.801283,-0.327152)--(-0.607814,-0.361764)--(-0.567664,-0.359610)--cycle;
\path[fill=blue!13!white,draw=blue!13!white,line width=.12pt] (-2.292506,-0.450808)--(-2.369413,-0.587353)--(-2.167407,-0.698793)--(-2.097056,-0.558630)--cycle;
\path[fill=blue!18!white,draw=blue!18!white,line width=.12pt] (2.097056,-0.558630)--(2.167407,-0.698793)--(2.369413,-0.587353)--(2.292506,-0.450808)--cycle;
\path[fill=blue!20!white,draw=blue!20!white,line width=.12pt] (-2.262742,-1.131371)--(-2.251872,-1.261099)--(-2.020307,-1.366038)--(-2.030059,-1.236817)--cycle;
\path[fill=blue!26!white,draw=blue!26!white,line width=.12pt] (2.030059,-1.236817)--(2.020307,-1.366038)--(2.251872,-1.261099)--(2.262742,-1.131371)--cycle;
\path[fill=green!65!black!30!white,draw=green!65!black!30!white,line width=.12pt] (-1.423533,-0.187782)--(-1.522544,-0.234792)--(-1.333369,-0.304943)--(-1.246659,-0.253371)--cycle;
\path[fill=green!65!black!30!white,draw=green!65!black!30!white,line width=.12pt] (1.246659,-0.253371)--(1.333369,-0.304943)--(1.522544,-0.234792)--(1.423533,-0.187782)--cycle;
\path[fill=blue!35!white,draw=blue!35!white,line width=.12pt] (-0.554926,-1.970959)--(-0.527943,-1.967157)--(-0.265249,-1.986640)--(-0.278805,-1.991439)--cycle;
\path[fill=blue!35!white,draw=blue!35!white,line width=.12pt] (0.278805,-1.991439)--(0.265249,-1.986640)--(0.527943,-1.967157)--(0.554926,-1.970959)--cycle;
\path[fill=blue!35!white,draw=blue!35!white,line width=.12pt] (-0.278805,-1.991439)--(-0.265249,-1.986640)--(-0.000000,-1.993156)--(-0.000000,-1.998287)--cycle;
\path[fill=blue!35!white,draw=blue!35!white,line width=.12pt] (-0.000000,-1.998287)--(-0.000000,-1.993156)--(0.265249,-1.986640)--(0.278805,-1.991439)--cycle;
\path[fill=green!65!black!34!white,draw=green!65!black!34!white,line width=.12pt] (-0.567664,-0.359610)--(-0.607814,-0.361764)--(-0.408491,-0.386728)--(-0.381508,-0.382925)--cycle;
\path[fill=green!65!black!33!white,draw=green!65!black!33!white,line width=.12pt] (0.381508,-0.382925)--(0.408491,-0.386728)--(0.607814,-0.361764)--(0.567664,-0.359610)--cycle;
\path[fill=blue!34!white,draw=blue!34!white,line width=.12pt] (0.860893,-1.908890)--(0.825702,-1.937046)--(1.088526,-1.890026)--(1.134919,-1.859866)--cycle;
\path[fill=blue!32!white,draw=blue!32!white,line width=.12pt] (-1.134919,-1.859866)--(-1.088526,-1.890026)--(-0.825702,-1.937046)--(-0.860893,-1.908890)--cycle;
\path[fill=blue!13!white,draw=blue!13!white,line width=.12pt] (-1.807416,-0.224200)--(-1.913535,-0.316685)--(-1.716761,-0.405857)--(-1.621555,-0.308427)--cycle;
\path[fill=blue!14!white,draw=blue!14!white,line width=.12pt] (1.621555,-0.308427)--(1.716761,-0.405857)--(1.913535,-0.316685)--(1.807416,-0.224200)--cycle;
\path[fill=green!65!black!32!white,draw=green!65!black!32!white,line width=.12pt] (-0.987036,-0.283224)--(-1.057780,-0.309976)--(-0.858714,-0.357051)--(-0.801283,-0.327152)--cycle;
\path[fill=green!65!black!31!white,draw=green!65!black!31!white,line width=.12pt] (0.801283,-0.327152)--(0.858714,-0.357051)--(1.057780,-0.309976)--(0.987036,-0.283224)--cycle;
\path[fill=blue!23!white,draw=blue!23!white,line width=.12pt] (-2.020307,-1.366038)--(-1.991426,-1.478411)--(-1.743993,-1.570165)--(-1.769285,-1.459123)--cycle;
\path[fill=blue!28!white,draw=blue!28!white,line width=.12pt] (1.769285,-1.459123)--(1.743993,-1.570165)--(1.991426,-1.478411)--(2.020307,-1.366038)--cycle;
\path[fill=green!65!black!34!white,draw=green!65!black!34!white,line width=.12pt] (-0.381508,-0.382925)--(-0.408491,-0.386728)--(-0.205234,-0.401803)--(-0.191677,-0.397005)--cycle;
\path[fill=green!65!black!34!white,draw=green!65!black!34!white,line width=.12pt] (0.191677,-0.397005)--(0.205234,-0.401803)--(0.408491,-0.386728)--(0.381508,-0.382925)--cycle;
\path[fill=blue!18!white,draw=blue!18!white,line width=.12pt] (-2.251872,-0.990774)--(-2.262742,-1.131371)--(-2.030059,-1.236817)--(-2.020307,-1.095713)--cycle;
\path[fill=blue!24!white,draw=blue!24!white,line width=.12pt] (2.020307,-1.095713)--(2.030059,-1.236817)--(2.262742,-1.131371)--(2.251872,-0.990774)--cycle;
\path[fill=green!65!black!34!white,draw=green!65!black!34!white,line width=.12pt] (-0.191677,-0.397005)--(-0.205234,-0.401803)--(-0.000000,-0.406844)--(-0.000000,-0.401713)--cycle;
\path[fill=green!65!black!34!white,draw=green!65!black!34!white,line width=.12pt] (-0.000000,-0.401713)--(-0.000000,-0.406844)--(0.205234,-0.401803)--(0.191677,-0.397005)--cycle;
\path[fill=blue!29!white,draw=blue!29!white,line width=.12pt] (-1.444914,-1.736532)--(-1.398014,-1.797648)--(-1.134919,-1.859866)--(-1.172992,-1.800837)--cycle;
\path[fill=blue!32!white,draw=blue!32!white,line width=.12pt] (1.172992,-1.800837)--(1.134919,-1.859866)--(1.398014,-1.797648)--(1.444914,-1.736532)--cycle;
\path[fill=blue!26!white,draw=blue!26!white,line width=.12pt] (-1.743993,-1.570165)--(-1.702920,-1.659211)--(-1.444914,-1.736532)--(-1.479763,-1.649352)--cycle;
\path[fill=blue!30!white,draw=blue!30!white,line width=.12pt] (1.479763,-1.649352)--(1.444914,-1.736532)--(1.702920,-1.659211)--(1.743993,-1.570165)--cycle;
\path[fill=blue!14!white,draw=blue!14!white,line width=.12pt] (-1.522544,-0.234792)--(-1.621555,-0.308427)--(-1.420078,-0.383140)--(-1.333369,-0.304943)--cycle;
\path[fill=blue!14!white,draw=blue!14!white,line width=.12pt] (1.333369,-0.304943)--(1.420078,-0.383140)--(1.621555,-0.308427)--(1.522544,-0.234792)--cycle;
\path[fill=green!65!black!30!white,draw=green!65!black!30!white,line width=.12pt] (-1.246659,-0.253371)--(-1.333369,-0.304943)--(-1.131352,-0.365485)--(-1.057780,-0.309976)--cycle;
\path[fill=green!65!black!30!white,draw=green!65!black!30!white,line width=.12pt] (1.057780,-0.309976)--(1.131352,-0.365485)--(1.333369,-0.304943)--(1.246659,-0.253371)--cycle;
\path[fill=blue!12!white,draw=blue!12!white,line width=.12pt] (-1.913535,-0.316685)--(-2.011334,-0.429608)--(-1.804504,-0.523338)--(-1.716761,-0.405857)--cycle;
\path[fill=blue!15!white,draw=blue!15!white,line width=.12pt] (1.716761,-0.405857)--(1.804504,-0.523338)--(2.011334,-0.429608)--(1.913535,-0.316685)--cycle;
\path[fill=blue!32!white,draw=blue!32!white,line width=.12pt] (-0.860893,-1.908890)--(-0.825702,-1.937046)--(-0.554926,-1.970959)--(-0.578577,-1.944248)--cycle;
\path[fill=blue!33!white,draw=blue!33!white,line width=.12pt] (0.578577,-1.944248)--(0.554926,-1.970959)--(0.825702,-1.937046)--(0.860893,-1.908890)--cycle;
\path[fill=blue!16!white,draw=blue!16!white,line width=.12pt] (-2.219681,-0.844710)--(-2.251872,-0.990774)--(-2.020307,-1.095713)--(-1.991426,-0.948149)--cycle;
\path[fill=blue!21!white,draw=blue!21!white,line width=.12pt] (1.991426,-0.948149)--(2.020307,-1.095713)--(2.251872,-0.990774)--(2.219681,-0.844710)--cycle;
\path[fill=green!65!black!32!white,draw=green!65!black!32!white,line width=.12pt] (-0.801283,-0.327152)--(-0.858714,-0.357051)--(-0.651378,-0.394145)--(-0.607814,-0.361764)--cycle;
\path[fill=green!65!black!31!white,draw=green!65!black!31!white,line width=.12pt] (0.607814,-0.361764)--(0.651378,-0.394145)--(0.858714,-0.357051)--(0.801283,-0.327152)--cycle;
\path[fill=blue!12!white,draw=blue!12!white,line width=.12pt] (-2.011334,-0.429608)--(-2.097056,-0.558630)--(-1.881411,-0.656355)--(-1.804504,-0.523338)--cycle;
\path[fill=blue!16!white,draw=blue!16!white,line width=.12pt] (1.804504,-0.523338)--(1.881411,-0.656355)--(2.097056,-0.558630)--(2.011334,-0.429608)--cycle;
\path[fill=blue!14!white,draw=blue!14!white,line width=.12pt] (-2.167407,-0.698793)--(-2.219681,-0.844710)--(-1.991426,-0.948149)--(-1.944527,-0.799796)--cycle;
\path[fill=blue!19!white,draw=blue!19!white,line width=.12pt] (1.944527,-0.799796)--(1.991426,-0.948149)--(2.219681,-0.844710)--(2.167407,-0.698793)--cycle;
\path[fill=blue!32!white,draw=blue!32!white,line width=.12pt] (-0.578577,-1.944248)--(-0.554926,-1.970959)--(-0.278805,-1.991439)--(-0.290688,-1.965600)--cycle;
\path[fill=blue!33!white,draw=blue!33!white,line width=.12pt] (0.290688,-1.965600)--(0.278805,-1.991439)--(0.554926,-1.970959)--(0.578577,-1.944248)--cycle;
\path[fill=blue!20!white,draw=blue!20!white,line width=.12pt] (-2.030059,-1.236817)--(-2.020307,-1.366038)--(-1.769285,-1.459123)--(-1.777825,-1.330351)--cycle;
\path[fill=blue!26!white,draw=blue!26!white,line width=.12pt] (1.777825,-1.330351)--(1.769285,-1.459123)--(2.020307,-1.366038)--(2.030059,-1.236817)--cycle;
\path[fill=blue!13!white,draw=blue!13!white,line width=.12pt] (-2.097056,-0.558630)--(-2.167407,-0.698793)--(-1.944527,-0.799796)--(-1.881411,-0.656355)--cycle;
\path[fill=blue!17!white,draw=blue!17!white,line width=.12pt] (1.881411,-0.656355)--(1.944527,-0.799796)--(2.167407,-0.698793)--(2.097056,-0.558630)--cycle;
\path[fill=green!65!black!32!white,draw=green!65!black!32!white,line width=.12pt] (-0.607814,-0.361764)--(-0.651378,-0.394145)--(-0.437769,-0.420898)--(-0.408491,-0.386728)--cycle;
\path[fill=green!65!black!32!white,draw=green!65!black!32!white,line width=.12pt] (0.408491,-0.386728)--(0.437769,-0.420898)--(0.651378,-0.394145)--(0.607814,-0.361764)--cycle;
\path[fill=blue!32!white,draw=blue!32!white,line width=.12pt] (-0.290688,-1.965600)--(-0.278805,-1.991439)--(-0.000000,-1.998287)--(-0.000000,-1.972741)--cycle;
\path[fill=blue!33!white,draw=blue!33!white,line width=.12pt] (-0.000000,-1.972741)--(-0.000000,-1.998287)--(0.278805,-1.991439)--(0.290688,-1.965600)--cycle;
\path[fill=blue!29!white,draw=blue!29!white,line width=.12pt] (-1.172992,-1.800837)--(-1.134919,-1.859866)--(-0.860893,-1.908890)--(-0.889774,-1.851505)--cycle;
\path[fill=blue!31!white,draw=blue!31!white,line width=.12pt] (0.889774,-1.851505)--(0.860893,-1.908890)--(1.134919,-1.859866)--(1.172992,-1.800837)--cycle;
\path[fill=blue!13!white,draw=blue!13!white,line width=.12pt] (-1.621555,-0.308427)--(-1.716761,-0.405857)--(-1.503455,-0.484957)--(-1.420078,-0.383140)--cycle;
\path[fill=blue!14!white,draw=blue!14!white,line width=.12pt] (1.420078,-0.383140)--(1.503455,-0.484957)--(1.716761,-0.405857)--(1.621555,-0.308427)--cycle;
\path[fill=green!65!black!30!white,draw=green!65!black!30!white,line width=.12pt] (-1.057780,-0.309976)--(-1.131352,-0.365485)--(-0.918440,-0.415835)--(-0.858714,-0.357051)--cycle;
\path[fill=green!65!black!30!white,draw=green!65!black!30!white,line width=.12pt] (0.858714,-0.357051)--(0.918440,-0.415835)--(1.131352,-0.365485)--(1.057780,-0.309976)--cycle;
\path[fill=blue!23!white,draw=blue!23!white,line width=.12pt] (-1.769285,-1.459123)--(-1.743993,-1.570165)--(-1.479763,-1.649352)--(-1.501223,-1.539458)--cycle;
\path[fill=blue!27!white,draw=blue!27!white,line width=.12pt] (1.501223,-1.539458)--(1.479763,-1.649352)--(1.743993,-1.570165)--(1.769285,-1.459123)--cycle;
\path[fill=blue!14!white,draw=blue!14!white,line width=.12pt] (-1.333369,-0.304943)--(-1.420078,-0.383140)--(-1.204924,-0.447619)--(-1.131352,-0.365485)--cycle;
\path[fill=blue!14!white,draw=blue!14!white,line width=.12pt] (1.131352,-0.365485)--(1.204924,-0.447619)--(1.420078,-0.383140)--(1.333369,-0.304943)--cycle;
\path[fill=green!65!black!32!white,draw=green!65!black!32!white,line width=.12pt] (0.205234,-0.401803)--(0.219943,-0.437053)--(0.437769,-0.420898)--(0.408491,-0.386728)--cycle;
\path[fill=green!65!black!32!white,draw=green!65!black!32!white,line width=.12pt] (-0.408491,-0.386728)--(-0.437769,-0.420898)--(-0.219943,-0.437053)--(-0.205234,-0.401803)--cycle;
\path[fill=blue!26!white,draw=blue!26!white,line width=.12pt] (-1.479763,-1.649352)--(-1.444914,-1.736532)--(-1.172992,-1.800837)--(-1.201283,-1.715208)--cycle;
\path[fill=blue!29!white,draw=blue!29!white,line width=.12pt] (1.201283,-1.715208)--(1.172992,-1.800837)--(1.444914,-1.736532)--(1.479763,-1.649352)--cycle;
\path[fill=green!65!black!32!white,draw=green!65!black!32!white,line width=.12pt] (-0.205234,-0.401803)--(-0.219943,-0.437053)--(-0.000000,-0.442456)--(-0.000000,-0.406844)--cycle;
\path[fill=green!65!black!32!white,draw=green!65!black!32!white,line width=.12pt] (-0.000000,-0.406844)--(-0.000000,-0.442456)--(0.219943,-0.437053)--(0.205234,-0.401803)--cycle;
\path[fill=blue!18!white,draw=blue!18!white,line width=.12pt] (-2.020307,-1.095713)--(-2.030059,-1.236817)--(-1.777825,-1.330351)--(-1.769285,-1.188798)--cycle;
\path[fill=blue!23!white,draw=blue!23!white,line width=.12pt] (1.769285,-1.188798)--(1.777825,-1.330351)--(2.030059,-1.236817)--(2.020307,-1.095713)--cycle;
\path[fill=blue!29!white,draw=blue!29!white,line width=.12pt] (-0.889774,-1.851505)--(-0.860893,-1.908890)--(-0.578577,-1.944248)--(-0.597986,-1.888050)--cycle;
\path[fill=blue!31!white,draw=blue!31!white,line width=.12pt] (0.597986,-1.888050)--(0.578577,-1.944248)--(0.860893,-1.908890)--(0.889774,-1.851505)--cycle;
\path[fill=blue!12!white,draw=blue!12!white,line width=.12pt] (-1.716761,-0.405857)--(-1.804504,-0.523338)--(-1.580295,-0.606480)--(-1.503455,-0.484957)--cycle;
\path[fill=blue!15!white,draw=blue!15!white,line width=.12pt] (1.503455,-0.484957)--(1.580295,-0.606480)--(1.804504,-0.523338)--(1.716761,-0.405857)--cycle;
\path[fill=green!65!black!30!white,draw=green!65!black!30!white,line width=.12pt] (-0.858714,-0.357051)--(-0.918440,-0.415835)--(-0.696683,-0.455508)--(-0.651378,-0.394145)--cycle;
\path[fill=green!65!black!30!white,draw=green!65!black!30!white,line width=.12pt] (0.651378,-0.394145)--(0.696683,-0.455508)--(0.918440,-0.415835)--(0.858714,-0.357051)--cycle;
\path[fill=blue!16!white,draw=blue!16!white,line width=.12pt] (-1.991426,-0.948149)--(-2.020307,-1.095713)--(-1.769285,-1.188798)--(-1.743993,-1.039904)--cycle;
\path[fill=blue!21!white,draw=blue!21!white,line width=.12pt] (1.743993,-1.039904)--(1.769285,-1.188798)--(2.020307,-1.095713)--(1.991426,-0.948149)--cycle;
\path[fill=blue!12!white,draw=blue!12!white,line width=.12pt] (-1.804504,-0.523338)--(-1.881411,-0.656355)--(-1.647647,-0.743041)--(-1.580295,-0.606480)--cycle;
\path[fill=blue!16!white,draw=blue!16!white,line width=.12pt] (1.580295,-0.606480)--(1.647647,-0.743041)--(1.881411,-0.656355)--(1.804504,-0.523338)--cycle;
\path[fill=blue!29!white,draw=blue!29!white,line width=.12pt] (-0.597986,-1.888050)--(-0.578577,-1.944248)--(-0.290688,-1.965600)--(-0.300440,-1.910118)--cycle;
\path[fill=blue!30!white,draw=blue!30!white,line width=.12pt] (0.300440,-1.910118)--(0.290688,-1.965600)--(0.578577,-1.944248)--(0.597986,-1.888050)--cycle;
\path[fill=blue!20!white,draw=blue!20!white,line width=.12pt] (-1.777825,-1.330351)--(-1.769285,-1.459123)--(-1.501223,-1.539458)--(-1.508470,-1.411074)--cycle;
\path[fill=blue!25!white,draw=blue!25!white,line width=.12pt] (1.508470,-1.411074)--(1.501223,-1.539458)--(1.769285,-1.459123)--(1.777825,-1.330351)--cycle;
\path[fill=blue!13!white,draw=blue!13!white,line width=.12pt] (-1.420078,-0.383140)--(-1.503455,-0.484957)--(-1.275669,-0.553222)--(-1.204924,-0.447619)--cycle;
\path[fill=blue!14!white,draw=blue!14!white,line width=.12pt] (1.204924,-0.447619)--(1.275669,-0.553222)--(1.503455,-0.484957)--(1.420078,-0.383140)--cycle;
\path[fill=blue!14!white,draw=blue!14!white,line width=.12pt] (-1.131352,-0.365485)--(-1.204924,-0.447619)--(-0.978167,-0.501243)--(-0.918440,-0.415835)--cycle;
\path[fill=blue!14!white,draw=blue!14!white,line width=.12pt] (0.918440,-0.415835)--(0.978167,-0.501243)--(1.204924,-0.447619)--(1.131352,-0.365485)--cycle;
\path[fill=blue!26!white,draw=blue!26!white,line width=.12pt] (-1.201283,-1.715208)--(-1.172992,-1.800837)--(-0.889774,-1.851505)--(-0.911234,-1.767098)--cycle;
\path[fill=blue!29!white,draw=blue!29!white,line width=.12pt] (0.911234,-1.767098)--(0.889774,-1.851505)--(1.172992,-1.800837)--(1.201283,-1.715208)--cycle;
\path[fill=blue!14!white,draw=blue!14!white,line width=.12pt] (-1.944527,-0.799796)--(-1.991426,-0.948149)--(-1.743993,-1.039904)--(-1.702920,-0.889390)--cycle;
\path[fill=blue!19!white,draw=blue!19!white,line width=.12pt] (1.702920,-0.889390)--(1.743993,-1.039904)--(1.991426,-0.948149)--(1.944527,-0.799796)--cycle;
\path[fill=green!65!black!30!white,draw=green!65!black!30!white,line width=.12pt] (-0.651378,-0.394145)--(-0.696683,-0.455508)--(-0.468217,-0.484122)--(-0.437769,-0.420898)--cycle;
\path[fill=green!65!black!30!white,draw=green!65!black!30!white,line width=.12pt] (0.437769,-0.420898)--(0.468217,-0.484122)--(0.696683,-0.455508)--(0.651378,-0.394145)--cycle;
\path[fill=blue!13!white,draw=blue!13!white,line width=.12pt] (-1.881411,-0.656355)--(-1.944527,-0.799796)--(-1.702920,-0.889390)--(-1.647647,-0.743041)--cycle;
\path[fill=blue!17!white,draw=blue!17!white,line width=.12pt] (1.647647,-0.743041)--(1.702920,-0.889390)--(1.944527,-0.799796)--(1.881411,-0.656355)--cycle;
\path[fill=blue!30!white,draw=blue!30!white,line width=.12pt] (-0.300440,-1.910118)--(-0.290688,-1.965600)--(-0.000000,-1.972741)--(-0.000000,-1.917498)--cycle;
\path[fill=blue!30!white,draw=blue!30!white,line width=.12pt] (-0.000000,-1.917498)--(-0.000000,-1.972741)--(0.290688,-1.965600)--(0.300440,-1.910118)--cycle;
\path[fill=blue!23!white,draw=blue!23!white,line width=.12pt] (-1.501223,-1.539458)--(-1.479763,-1.649352)--(-1.201283,-1.715208)--(-1.218704,-1.606269)--cycle;
\path[fill=blue!27!white,draw=blue!27!white,line width=.12pt] (1.218704,-1.606269)--(1.201283,-1.715208)--(1.479763,-1.649352)--(1.501223,-1.539458)--cycle;
\path[fill=green!65!black!30!white,draw=green!65!black!30!white,line width=.12pt] (-0.437769,-0.420898)--(-0.468217,-0.484122)--(-0.235241,-0.501401)--(-0.219943,-0.437053)--cycle;
\path[fill=green!65!black!30!white,draw=green!65!black!30!white,line width=.12pt] (0.219943,-0.437053)--(0.235241,-0.501401)--(0.468217,-0.484122)--(0.437769,-0.420898)--cycle;
\path[fill=green!65!black!30!white,draw=green!65!black!30!white,line width=.12pt] (-0.219943,-0.437053)--(-0.235241,-0.501401)--(-0.000000,-0.507180)--(-0.000000,-0.442456)--cycle;
\path[fill=green!65!black!30!white,draw=green!65!black!30!white,line width=.12pt] (-0.000000,-0.442456)--(-0.000000,-0.507180)--(0.235241,-0.501401)--(0.219943,-0.437053)--cycle;
\path[fill=blue!18!white,draw=blue!18!white,line width=.12pt] (-1.769285,-1.188798)--(-1.777825,-1.330351)--(-1.508470,-1.411074)--(-1.501223,-1.269133)--cycle;
\path[fill=blue!22!white,draw=blue!22!white,line width=.12pt] (1.501223,-1.269133)--(1.508470,-1.411074)--(1.777825,-1.330351)--(1.769285,-1.188798)--cycle;
\path[fill=blue!12!white,draw=blue!12!white,line width=.12pt] (-1.503455,-0.484957)--(-1.580295,-0.606480)--(-1.340867,-0.678234)--(-1.275669,-0.553222)--cycle;
\path[fill=blue!14!white,draw=blue!14!white,line width=.12pt] (1.275669,-0.553222)--(1.340867,-0.678234)--(1.580295,-0.606480)--(1.503455,-0.484957)--cycle;
\path[fill=blue!14!white,draw=blue!14!white,line width=.12pt] (-0.918440,-0.415835)--(-0.978167,-0.501243)--(-0.741989,-0.543496)--(-0.696683,-0.455508)--cycle;
\path[fill=blue!14!white,draw=blue!14!white,line width=.12pt] (0.696683,-0.455508)--(0.741989,-0.543496)--(0.978167,-0.501243)--(0.918440,-0.415835)--cycle;
\path[fill=blue!26!white,draw=blue!26!white,line width=.12pt] (-0.911234,-1.767098)--(-0.889774,-1.851505)--(-0.597986,-1.888050)--(-0.612409,-1.804524)--cycle;
\path[fill=blue!28!white,draw=blue!28!white,line width=.12pt] (0.612409,-1.804524)--(0.597986,-1.888050)--(0.889774,-1.851505)--(0.911234,-1.767098)--cycle;
\path[fill=blue!13!white,draw=blue!13!white,line width=.12pt] (-1.204924,-0.447619)--(-1.275669,-0.553222)--(-1.035598,-0.609994)--(-0.978167,-0.501243)--cycle;
\path[fill=blue!14!white,draw=blue!14!white,line width=.12pt] (0.978167,-0.501243)--(1.035598,-0.609994)--(1.275669,-0.553222)--(1.204924,-0.447619)--cycle;
\path[fill=blue!20!white,draw=blue!20!white,line width=.12pt] (-1.508470,-1.411074)--(-1.501223,-1.539458)--(-1.218704,-1.606269)--(-1.224587,-1.478207)--cycle;
\path[fill=blue!24!white,draw=blue!24!white,line width=.12pt] (1.224587,-1.478207)--(1.218704,-1.606269)--(1.501223,-1.539458)--(1.508470,-1.411074)--cycle;
\path[fill=blue!23!white,draw=blue!23!white,line width=.12pt] (-1.218704,-1.606269)--(-1.201283,-1.715208)--(-0.911234,-1.767098)--(-0.924449,-1.658912)--cycle;
\path[fill=blue!26!white,draw=blue!26!white,line width=.12pt] (0.924449,-1.658912)--(0.911234,-1.767098)--(1.201283,-1.715208)--(1.218704,-1.606269)--cycle;
\path[fill=blue!16!white,draw=blue!16!white,line width=.12pt] (-1.743993,-1.039904)--(-1.769285,-1.188798)--(-1.501223,-1.269133)--(-1.479763,-1.119090)--cycle;
\path[fill=blue!20!white,draw=blue!20!white,line width=.12pt] (1.479763,-1.119090)--(1.501223,-1.269133)--(1.769285,-1.188798)--(1.743993,-1.039904)--cycle;
\path[fill=blue!12!white,draw=blue!12!white,line width=.12pt] (-1.580295,-0.606480)--(-1.647647,-0.743041)--(-1.398014,-0.817853)--(-1.340867,-0.678234)--cycle;
\path[fill=blue!15!white,draw=blue!15!white,line width=.12pt] (1.340867,-0.678234)--(1.398014,-0.817853)--(1.647647,-0.743041)--(1.580295,-0.606480)--cycle;
\path[fill=blue!26!white,draw=blue!26!white,line width=.12pt] (-0.612409,-1.804524)--(-0.597986,-1.888050)--(-0.300440,-1.910118)--(-0.307686,-1.827125)--cycle;
\path[fill=blue!28!white,draw=blue!28!white,line width=.12pt] (0.307686,-1.827125)--(0.300440,-1.910118)--(0.597986,-1.888050)--(0.612409,-1.804524)--cycle;
\path[fill=blue!14!white,draw=blue!14!white,line width=.12pt] (-0.696683,-0.455508)--(-0.741989,-0.543496)--(-0.498665,-0.573971)--(-0.468217,-0.484122)--cycle;
\path[fill=blue!14!white,draw=blue!14!white,line width=.12pt] (0.468217,-0.484122)--(0.498665,-0.573971)--(0.741989,-0.543496)--(0.696683,-0.455508)--cycle;
\path[fill=blue!14!white,draw=blue!14!white,line width=.12pt] (-1.702920,-0.889390)--(-1.743993,-1.039904)--(-1.479763,-1.119090)--(-1.444914,-0.966711)--cycle;
\path[fill=blue!18!white,draw=blue!18!white,line width=.12pt] (1.444914,-0.966711)--(1.479763,-1.119090)--(1.743993,-1.039904)--(1.702920,-0.889390)--cycle;
\path[fill=blue!13!white,draw=blue!13!white,line width=.12pt] (-1.647647,-0.743041)--(-1.702920,-0.889390)--(-1.444914,-0.966711)--(-1.398014,-0.817853)--cycle;
\path[fill=blue!16!white,draw=blue!16!white,line width=.12pt] (1.398014,-0.817853)--(1.444914,-0.966711)--(1.702920,-0.889390)--(1.647647,-0.743041)--cycle;
\path[fill=blue!27!white,draw=blue!27!white,line width=.12pt] (-0.307686,-1.827125)--(-0.300440,-1.910118)--(-0.000000,-1.917498)--(-0.000000,-1.834683)--cycle;
\path[fill=blue!27!white,draw=blue!27!white,line width=.12pt] (-0.000000,-1.834683)--(-0.000000,-1.917498)--(0.300440,-1.910118)--(0.307686,-1.827125)--cycle;
\path[fill=blue!14!white,draw=blue!14!white,line width=.12pt] (-0.468217,-0.484122)--(-0.498665,-0.573971)--(-0.250539,-0.592374)--(-0.235241,-0.501401)--cycle;
\path[fill=blue!14!white,draw=blue!14!white,line width=.12pt] (0.235241,-0.501401)--(0.250539,-0.592374)--(0.498665,-0.573971)--(0.468217,-0.484122)--cycle;
\path[fill=blue!12!white,draw=blue!12!white,line width=.12pt] (-1.275669,-0.553222)--(-1.340867,-0.678234)--(-1.088526,-0.737908)--(-1.035598,-0.609994)--cycle;
\path[fill=blue!14!white,draw=blue!14!white,line width=.12pt] (1.035598,-0.609994)--(1.088526,-0.737908)--(1.340867,-0.678234)--(1.275669,-0.553222)--cycle;
\path[fill=blue!13!white,draw=blue!13!white,line width=.12pt] (-0.978167,-0.501243)--(-1.035598,-0.609994)--(-0.785553,-0.654728)--(-0.741989,-0.543496)--cycle;
\path[fill=blue!14!white,draw=blue!14!white,line width=.12pt] (0.741989,-0.543496)--(0.785553,-0.654728)--(1.035598,-0.609994)--(0.978167,-0.501243)--cycle;
\path[fill=blue!18!white,draw=blue!18!white,line width=.12pt] (-1.501223,-1.269133)--(-1.508470,-1.411074)--(-1.224587,-1.478207)--(-1.218704,-1.335944)--cycle;
\path[fill=blue!22!white,draw=blue!22!white,line width=.12pt] (1.218704,-1.335944)--(1.224587,-1.478207)--(1.508470,-1.411074)--(1.501223,-1.269133)--cycle;
\path[fill=blue!14!white,draw=blue!14!white,line width=.12pt] (-0.235241,-0.501401)--(-0.250539,-0.592374)--(-0.000000,-0.598528)--(-0.000000,-0.507180)--cycle;
\path[fill=blue!14!white,draw=blue!14!white,line width=.12pt] (-0.000000,-0.507180)--(-0.000000,-0.598528)--(0.250539,-0.592374)--(0.235241,-0.501401)--cycle;
\path[fill=blue!23!white,draw=blue!23!white,line width=.12pt] (-0.924449,-1.658912)--(-0.911234,-1.767098)--(-0.612409,-1.804524)--(-0.621290,-1.696881)--cycle;
\path[fill=blue!25!white,draw=blue!25!white,line width=.12pt] (0.621290,-1.696881)--(0.612409,-1.804524)--(0.911234,-1.767098)--(0.924449,-1.658912)--cycle;
\path[fill=blue!20!white,draw=blue!20!white,line width=.12pt] (-1.224587,-1.478207)--(-1.218704,-1.606269)--(-0.924449,-1.658912)--(-0.928911,-1.531105)--cycle;
\path[fill=blue!23!white,draw=blue!23!white,line width=.12pt] (0.928911,-1.531105)--(0.924449,-1.658912)--(1.218704,-1.606269)--(1.224587,-1.478207)--cycle;
\path[fill=blue!12!white,draw=blue!12!white,line width=.12pt] (-1.340867,-0.678234)--(-1.398014,-0.817853)--(-1.134919,-0.880070)--(-1.088526,-0.737908)--cycle;
\path[fill=blue!15!white,draw=blue!15!white,line width=.12pt] (1.088526,-0.737908)--(1.134919,-0.880070)--(1.398014,-0.817853)--(1.340867,-0.678234)--cycle;
\path[fill=blue!16!white,draw=blue!16!white,line width=.12pt] (-1.479763,-1.119090)--(-1.501223,-1.269133)--(-1.218704,-1.335944)--(-1.201283,-1.184946)--cycle;
\path[fill=blue!19!white,draw=blue!19!white,line width=.12pt] (1.201283,-1.184946)--(1.218704,-1.335944)--(1.501223,-1.269133)--(1.479763,-1.119090)--cycle;
\path[fill=blue!13!white,draw=blue!13!white,line width=.12pt] (-0.741989,-0.543496)--(-0.785553,-0.654728)--(-0.527943,-0.686992)--(-0.498665,-0.573971)--cycle;
\path[fill=blue!13!white,draw=blue!13!white,line width=.12pt] (0.498665,-0.573971)--(0.527943,-0.686992)--(0.785553,-0.654728)--(0.741989,-0.543496)--cycle;
\path[fill=blue!24!white,draw=blue!24!white,line width=.12pt] (-0.621290,-1.696881)--(-0.612409,-1.804524)--(-0.307686,-1.827125)--(-0.312148,-1.719809)--cycle;
\path[fill=blue!25!white,draw=blue!25!white,line width=.12pt] (0.312148,-1.719809)--(0.307686,-1.827125)--(0.612409,-1.804524)--(0.621290,-1.696881)--cycle;
\path[fill=blue!13!white,draw=blue!13!white,line width=.12pt] (-1.398014,-0.817853)--(-1.444914,-0.966711)--(-1.172992,-1.031016)--(-1.134919,-0.880070)--cycle;
\path[fill=blue!16!white,draw=blue!16!white,line width=.12pt] (1.134919,-0.880070)--(1.172992,-1.031016)--(1.444914,-0.966711)--(1.398014,-0.817853)--cycle;
\path[fill=blue!12!white,draw=blue!12!white,line width=.12pt] (-1.035598,-0.609994)--(-1.088526,-0.737908)--(-0.825702,-0.784928)--(-0.785553,-0.654728)--cycle;
\path[fill=blue!14!white,draw=blue!14!white,line width=.12pt] (0.785553,-0.654728)--(0.825702,-0.784928)--(1.088526,-0.737908)--(1.035598,-0.609994)--cycle;
\path[fill=blue!14!white,draw=blue!14!white,line width=.12pt] (-1.444914,-0.966711)--(-1.479763,-1.119090)--(-1.201283,-1.184946)--(-1.172992,-1.031016)--cycle;
\path[fill=blue!17!white,draw=blue!17!white,line width=.12pt] (1.172992,-1.031016)--(1.201283,-1.184946)--(1.479763,-1.119090)--(1.444914,-0.966711)--cycle;
\path[fill=blue!24!white,draw=blue!24!white,line width=.12pt] (-0.312148,-1.719809)--(-0.307686,-1.827125)--(-0.000000,-1.834683)--(-0.000000,-1.727477)--cycle;
\path[fill=blue!24!white,draw=blue!24!white,line width=.12pt] (-0.000000,-1.727477)--(-0.000000,-1.834683)--(0.307686,-1.827125)--(0.312148,-1.719809)--cycle;
\path[fill=blue!13!white,draw=blue!13!white,line width=.12pt] (-0.498665,-0.573971)--(-0.527943,-0.686992)--(-0.265249,-0.706475)--(-0.250539,-0.592374)--cycle;
\path[fill=blue!13!white,draw=blue!13!white,line width=.12pt] (0.250539,-0.592374)--(0.265249,-0.706475)--(0.527943,-0.686992)--(0.498665,-0.573971)--cycle;
\path[fill=blue!18!white,draw=blue!18!white,line width=.12pt] (-1.218704,-1.335944)--(-1.224587,-1.478207)--(-0.928911,-1.531105)--(-0.924449,-1.388587)--cycle;
\path[fill=blue!21!white,draw=blue!21!white,line width=.12pt] (0.924449,-1.388587)--(0.928911,-1.531105)--(1.224587,-1.478207)--(1.218704,-1.335944)--cycle;
\path[fill=blue!21!white,draw=blue!21!white,line width=.12pt] (-0.928911,-1.531105)--(-0.924449,-1.658912)--(-0.621290,-1.696881)--(-0.624289,-1.569256)--cycle;
\path[fill=blue!23!white,draw=blue!23!white,line width=.12pt] (0.624289,-1.569256)--(0.621290,-1.696881)--(0.924449,-1.658912)--(0.928911,-1.531105)--cycle;
\path[fill=blue!13!white,draw=blue!13!white,line width=.12pt] (-0.250539,-0.592374)--(-0.265249,-0.706475)--(-0.000000,-0.712991)--(-0.000000,-0.598528)--cycle;
\path[fill=blue!13!white,draw=blue!13!white,line width=.12pt] (-0.000000,-0.598528)--(-0.000000,-0.712991)--(0.265249,-0.706475)--(0.250539,-0.592374)--cycle;
\path[fill=blue!14!white,draw=blue!14!white,line width=.12pt] (0.825702,-0.784928)--(0.860893,-0.929094)--(1.134919,-0.880070)--(1.088526,-0.737908)--cycle;
\path[fill=blue!12!white,draw=blue!12!white,line width=.12pt] (-1.088526,-0.737908)--(-1.134919,-0.880070)--(-0.860893,-0.929094)--(-0.825702,-0.784928)--cycle;
\path[fill=blue!12!white,draw=blue!12!white,line width=.12pt] (-0.785553,-0.654728)--(-0.825702,-0.784928)--(-0.554926,-0.818841)--(-0.527943,-0.686992)--cycle;
\path[fill=blue!13!white,draw=blue!13!white,line width=.12pt] (0.527943,-0.686992)--(0.554926,-0.818841)--(0.825702,-0.784928)--(0.785553,-0.654728)--cycle;
\path[fill=blue!16!white,draw=blue!16!white,line width=.12pt] (-1.201283,-1.184946)--(-1.218704,-1.335944)--(-0.924449,-1.388587)--(-0.911234,-1.236837)--cycle;
\path[fill=blue!19!white,draw=blue!19!white,line width=.12pt] (0.911234,-1.236837)--(0.924449,-1.388587)--(1.218704,-1.335944)--(1.201283,-1.184946)--cycle;
\path[fill=blue!21!white,draw=blue!21!white,line width=.12pt] (-0.624289,-1.569256)--(-0.621290,-1.696881)--(-0.312148,-1.719809)--(-0.313655,-1.592296)--cycle;
\path[fill=blue!22!white,draw=blue!22!white,line width=.12pt] (0.313655,-1.592296)--(0.312148,-1.719809)--(0.621290,-1.696881)--(0.624289,-1.569256)--cycle;
\path[fill=blue!13!white,draw=blue!13!white,line width=.12pt] (-1.134919,-0.880070)--(-1.172992,-1.031016)--(-0.889774,-1.081685)--(-0.860893,-0.929094)--cycle;
\path[fill=blue!15!white,draw=blue!15!white,line width=.12pt] (0.860893,-0.929094)--(0.889774,-1.081685)--(1.172992,-1.031016)--(1.134919,-0.880070)--cycle;
\path[fill=blue!14!white,draw=blue!14!white,line width=.12pt] (-1.172992,-1.031016)--(-1.201283,-1.184946)--(-0.911234,-1.236837)--(-0.889774,-1.081685)--cycle;
\path[fill=blue!17!white,draw=blue!17!white,line width=.12pt] (0.889774,-1.081685)--(0.911234,-1.236837)--(1.201283,-1.184946)--(1.172992,-1.031016)--cycle;
\path[fill=blue!18!white,draw=blue!18!white,line width=.12pt] (-0.924449,-1.388587)--(-0.928911,-1.531105)--(-0.624289,-1.569256)--(-0.621290,-1.426556)--cycle;
\path[fill=blue!20!white,draw=blue!20!white,line width=.12pt] (0.621290,-1.426556)--(0.624289,-1.569256)--(0.928911,-1.531105)--(0.924449,-1.388587)--cycle;
\path[fill=blue!21!white,draw=blue!21!white,line width=.12pt] (-0.313655,-1.592296)--(-0.312148,-1.719809)--(-0.000000,-1.727477)--(-0.000000,-1.600000)--cycle;
\path[fill=blue!22!white,draw=blue!22!white,line width=.12pt] (-0.000000,-1.600000)--(-0.000000,-1.727477)--(0.312148,-1.719809)--(0.313655,-1.592296)--cycle;
\path[fill=blue!12!white,draw=blue!12!white,line width=.12pt] (-0.527943,-0.686992)--(-0.554926,-0.818841)--(-0.278805,-0.839321)--(-0.265249,-0.706475)--cycle;
\path[fill=blue!13!white,draw=blue!13!white,line width=.12pt] (0.265249,-0.706475)--(0.278805,-0.839321)--(0.554926,-0.818841)--(0.527943,-0.686992)--cycle;
\path[fill=blue!12!white,draw=blue!12!white,line width=.12pt] (-0.825702,-0.784928)--(-0.860893,-0.929094)--(-0.578577,-0.964452)--(-0.554926,-0.818841)--cycle;
\path[fill=blue!14!white,draw=blue!14!white,line width=.12pt] (0.554926,-0.818841)--(0.578577,-0.964452)--(0.860893,-0.929094)--(0.825702,-0.784928)--cycle;
\path[fill=blue!13!white,draw=blue!13!white,line width=.12pt] (-0.000000,-0.712991)--(-0.000000,-0.846169)--(0.278805,-0.839321)--(0.265249,-0.706475)--cycle;
\path[fill=blue!13!white,draw=blue!13!white,line width=.12pt] (-0.265249,-0.706475)--(-0.278805,-0.839321)--(-0.000000,-0.846169)--(-0.000000,-0.712991)--cycle;
\path[fill=blue!16!white,draw=blue!16!white,line width=.12pt] (-0.911234,-1.236837)--(-0.924449,-1.388587)--(-0.621290,-1.426556)--(-0.612409,-1.274262)--cycle;
\path[fill=blue!18!white,draw=blue!18!white,line width=.12pt] (0.612409,-1.274262)--(0.621290,-1.426556)--(0.924449,-1.388587)--(0.911234,-1.236837)--cycle;
\path[fill=blue!18!white,draw=blue!18!white,line width=.12pt] (-0.621290,-1.426556)--(-0.624289,-1.569256)--(-0.313655,-1.592296)--(-0.312148,-1.449484)--cycle;
\path[fill=blue!20!white,draw=blue!20!white,line width=.12pt] (0.312148,-1.449484)--(0.313655,-1.592296)--(0.624289,-1.569256)--(0.621290,-1.426556)--cycle;
\path[fill=blue!13!white,draw=blue!13!white,line width=.12pt] (-0.860893,-0.929094)--(-0.889774,-1.081685)--(-0.597986,-1.118229)--(-0.578577,-0.964452)--cycle;
\path[fill=blue!15!white,draw=blue!15!white,line width=.12pt] (0.578577,-0.964452)--(0.597986,-1.118229)--(0.889774,-1.081685)--(0.860893,-0.929094)--cycle;
\path[fill=blue!14!white,draw=blue!14!white,line width=.12pt] (-0.889774,-1.081685)--(-0.911234,-1.236837)--(-0.612409,-1.274262)--(-0.597986,-1.118229)--cycle;
\path[fill=blue!16!white,draw=blue!16!white,line width=.12pt] (0.597986,-1.118229)--(0.612409,-1.274262)--(0.911234,-1.236837)--(0.889774,-1.081685)--cycle;
\path[fill=blue!13!white,draw=blue!13!white,line width=.12pt] (-0.554926,-0.818841)--(-0.578577,-0.964452)--(-0.290688,-0.985804)--(-0.278805,-0.839321)--cycle;
\path[fill=blue!13!white,draw=blue!13!white,line width=.12pt] (0.278805,-0.839321)--(0.290688,-0.985804)--(0.578577,-0.964452)--(0.554926,-0.818841)--cycle;
\path[fill=blue!19!white,draw=blue!19!white,line width=.12pt] (-0.312148,-1.449484)--(-0.313655,-1.592296)--(-0.000000,-1.600000)--(-0.000000,-1.457152)--cycle;
\path[fill=blue!19!white,draw=blue!19!white,line width=.12pt] (-0.000000,-1.457152)--(-0.000000,-1.600000)--(0.313655,-1.592296)--(0.312148,-1.449484)--cycle;
\path[fill=blue!13!white,draw=blue!13!white,line width=.12pt] (-0.278805,-0.839321)--(-0.290688,-0.985804)--(-0.000000,-0.992945)--(-0.000000,-0.846169)--cycle;
\path[fill=blue!13!white,draw=blue!13!white,line width=.12pt] (-0.000000,-0.846169)--(-0.000000,-0.992945)--(0.290688,-0.985804)--(0.278805,-0.839321)--cycle;
\path[fill=blue!16!white,draw=blue!16!white,line width=.12pt] (-0.612409,-1.274262)--(-0.621290,-1.426556)--(-0.312148,-1.449484)--(-0.307686,-1.296863)--cycle;
\path[fill=blue!18!white,draw=blue!18!white,line width=.12pt] (0.307686,-1.296863)--(0.312148,-1.449484)--(0.621290,-1.426556)--(0.612409,-1.274262)--cycle;
\path[fill=blue!13!white,draw=blue!13!white,line width=.12pt] (-0.578577,-0.964452)--(-0.597986,-1.118229)--(-0.300440,-1.140298)--(-0.290688,-0.985804)--cycle;
\path[fill=blue!14!white,draw=blue!14!white,line width=.12pt] (0.290688,-0.985804)--(0.300440,-1.140298)--(0.597986,-1.118229)--(0.578577,-0.964452)--cycle;
\path[fill=blue!15!white,draw=blue!15!white,line width=.12pt] (-0.597986,-1.118229)--(-0.612409,-1.274262)--(-0.307686,-1.296863)--(-0.300440,-1.140298)--cycle;
\path[fill=blue!16!white,draw=blue!16!white,line width=.12pt] (0.300440,-1.140298)--(0.307686,-1.296863)--(0.612409,-1.274262)--(0.597986,-1.118229)--cycle;
\path[fill=blue!17!white,draw=blue!17!white,line width=.12pt] (-0.000000,-1.304421)--(-0.000000,-1.457152)--(0.312148,-1.449484)--(0.307686,-1.296863)--cycle;
\path[fill=blue!17!white,draw=blue!17!white,line width=.12pt] (-0.307686,-1.296863)--(-0.312148,-1.449484)--(-0.000000,-1.457152)--(-0.000000,-1.304421)--cycle;
\path[fill=blue!14!white,draw=blue!14!white,line width=.12pt] (-0.290688,-0.985804)--(-0.300440,-1.140298)--(-0.000000,-1.147677)--(-0.000000,-0.992945)--cycle;
\path[fill=blue!14!white,draw=blue!14!white,line width=.12pt] (-0.000000,-0.992945)--(-0.000000,-1.147677)--(0.300440,-1.140298)--(0.290688,-0.985804)--cycle;
\path[fill=blue!15!white,draw=blue!15!white,line width=.12pt] (-0.300440,-1.140298)--(-0.307686,-1.296863)--(-0.000000,-1.304421)--(-0.000000,-1.147677)--cycle;
\path[fill=blue!15!white,draw=blue!15!white,line width=.12pt] (-0.000000,-1.147677)--(-0.000000,-1.304421)--(0.307686,-1.296863)--(0.300440,-1.140298)--cycle;

\draw[green!45!black,thick,domain=0:360,samples=150]
plot({2.4*cos(\x)},{1.2*sin(\x)+.8660254038*.8});
\pgfmathsetmacro{\ss}{2.4-.8*cos(45)}
\draw[green!45!black,thick,domain=-35.264:215.264,samples=110]
plot({\ss*cos(\x)},{.5*\ss*sin(\x)+.8660254038*.8*sin(45)});
\draw[green!45!black,dashed,domain=215.264:324.736,samples=55]
plot({\ss*cos(\x)},{.5*\ss*sin(\x)+.8660254038*.8*sin(45)});
\tp{Aband}{45}{130}
\draw[force,->](Aband)--++({.70*sin(45)*(-cos(130))},
{-.70*(.5*sin(45)*sin(130)+.8660254038*cos(45))});
\node[blue!75!black]at(-1.2,.31){$\sigma_2(\theta)$};
\node at(0,-2.50){内侧上膜带；两边界为平行圈};
\end{scope}
\begin{scope}[xshift=7cm,yshift=1.55cm]
\draw[gray!60](0,0)circle(1.14);
\draw[->](-1.6,0)--(1.75,0)node[right]{$\rho$};
\draw[->](0,-1.38)--(0,1.58)node[above]{$z$};
\coordinate(Lcap)at({-1.14*cos(45)},{1.14*sin(45)});
\fill[green!12](0,0)--(Lcap)arc(135:90:1.14)--cycle;
\draw[green!45!black,very thick](Lcap)arc(135:90:1.14);
\draw[thin](0,0)--(Lcap)node[midway,above right=3pt]{$r$};
\draw(-.38,0)arc(180:135:.38);
\node at(-.55,.22){$\theta$};
\draw[force,->](Lcap)--++(-.52,-.52)node[below left]{$\sigma_2$};
\draw[force,->](0,1.14)--(.95,1.14)node[right]{$\sigma_2$};
\draw[pressure,->]({-1.14*cos(63)},{1.14*sin(63)})--++({-.6*cos(63)},{.6*sin(63)});
\node[red!75!black]at(-.95,1.56){$p\mathbf n$};
\node at(0,-1.8){顶部子午力的 $z$ 分量为零};
\end{scope}
\begin{scope}[xshift=7cm,yshift=-2.15cm]
\filldraw[fill=green!12,draw=green!45!black](0,0)circle(1.14);
\filldraw[fill=white,draw=green!45!black](0,0)circle(.8713);
\draw[thin](0,0)--(1.14,0)node[above,pos=.97]{$R$};
\draw[thin](0,0)--({.8713*cos(225)},{.8713*sin(225)})node[midway,below right]{$s$};
\node at(0,-1.55){投影面积 $\pi(R^2-s^2)$};
\end{scope}
\end{tikzpicture}
% 5. Surface principal directions and a separate meridian cut.
\begin{tikzpicture}[font=\small]
\begin{scope}
% Automatically generated vector mesh; parameters and camera in build/torus-mesh-q03.json.
\path[fill=blue!39!white,draw=blue!39!white,line width=.12pt] (0.000000,1.304421)--(0.000000,1.147677)--(-0.300440,1.140298)--(-0.307686,1.296863)--cycle;
\path[fill=blue!39!white,draw=blue!39!white,line width=.12pt] (0.307686,1.296863)--(0.300440,1.140298)--(0.000000,1.147677)--(0.000000,1.304421)--cycle;
\path[fill=blue!40!white,draw=blue!40!white,line width=.12pt] (0.300440,1.140298)--(0.290688,0.985804)--(0.000000,0.992945)--(0.000000,1.147677)--cycle;
\path[fill=blue!40!white,draw=blue!40!white,line width=.12pt] (0.000000,1.147677)--(0.000000,0.992945)--(-0.290688,0.985804)--(-0.300440,1.140298)--cycle;
\path[fill=blue!37!white,draw=blue!37!white,line width=.12pt] (0.312148,1.449484)--(0.307686,1.296863)--(0.000000,1.304421)--(0.000000,1.457152)--cycle;
\path[fill=blue!37!white,draw=blue!37!white,line width=.12pt] (0.000000,1.457152)--(0.000000,1.304421)--(-0.307686,1.296863)--(-0.312148,1.449484)--cycle;
\path[fill=blue!39!white,draw=blue!39!white,line width=.12pt] (0.612409,1.274262)--(0.597986,1.118229)--(0.300440,1.140298)--(0.307686,1.296863)--cycle;
\path[fill=blue!38!white,draw=blue!38!white,line width=.12pt] (-0.307686,1.296863)--(-0.300440,1.140298)--(-0.597986,1.118229)--(-0.612409,1.274262)--cycle;
\path[fill=blue!41!white,draw=blue!41!white,line width=.12pt] (0.597986,1.118229)--(0.578577,0.964452)--(0.290688,0.985804)--(0.300440,1.140298)--cycle;
\path[fill=blue!40!white,draw=blue!40!white,line width=.12pt] (-0.300440,1.140298)--(-0.290688,0.985804)--(-0.578577,0.964452)--(-0.597986,1.118229)--cycle;
\path[fill=blue!38!white,draw=blue!38!white,line width=.12pt] (0.621290,1.426556)--(0.612409,1.274262)--(0.307686,1.296863)--(0.312148,1.449484)--cycle;
\path[fill=blue!36!white,draw=blue!36!white,line width=.12pt] (-0.312148,1.449484)--(-0.307686,1.296863)--(-0.612409,1.274262)--(-0.621290,1.426556)--cycle;
\path[fill=blue!41!white,draw=blue!41!white,line width=.12pt] (0.290688,0.985804)--(0.278805,0.839321)--(0.000000,0.846169)--(0.000000,0.992945)--cycle;
\path[fill=blue!41!white,draw=blue!41!white,line width=.12pt] (0.000000,0.992945)--(0.000000,0.846169)--(-0.278805,0.839321)--(-0.290688,0.985804)--cycle;
\path[fill=blue!35!white,draw=blue!35!white,line width=.12pt] (0.313655,1.592296)--(0.312148,1.449484)--(0.000000,1.457152)--(0.000000,1.600000)--cycle;
\path[fill=blue!35!white,draw=blue!35!white,line width=.12pt] (0.000000,1.600000)--(0.000000,1.457152)--(-0.312148,1.449484)--(-0.313655,1.592296)--cycle;
\path[fill=blue!41!white,draw=blue!41!white,line width=.12pt] (-0.290688,0.985804)--(-0.278805,0.839321)--(-0.554926,0.818841)--(-0.578577,0.964452)--cycle;
\path[fill=blue!41!white,draw=blue!41!white,line width=.12pt] (0.578577,0.964452)--(0.554926,0.818841)--(0.278805,0.839321)--(0.290688,0.985804)--cycle;
\path[fill=blue!38!white,draw=blue!38!white,line width=.12pt] (-0.612409,1.274262)--(-0.597986,1.118229)--(-0.889774,1.081685)--(-0.911234,1.236837)--cycle;
\path[fill=blue!40!white,draw=blue!40!white,line width=.12pt] (0.911234,1.236837)--(0.889774,1.081685)--(0.597986,1.118229)--(0.612409,1.274262)--cycle;
\path[fill=blue!41!white,draw=blue!41!white,line width=.12pt] (0.889774,1.081685)--(0.860893,0.929094)--(0.578577,0.964452)--(0.597986,1.118229)--cycle;
\path[fill=blue!39!white,draw=blue!39!white,line width=.12pt] (-0.597986,1.118229)--(-0.578577,0.964452)--(-0.860893,0.929094)--(-0.889774,1.081685)--cycle;
\path[fill=blue!36!white,draw=blue!36!white,line width=.12pt] (0.624289,1.569256)--(0.621290,1.426556)--(0.312148,1.449484)--(0.313655,1.592296)--cycle;
\path[fill=blue!34!white,draw=blue!34!white,line width=.12pt] (-0.313655,1.592296)--(-0.312148,1.449484)--(-0.621290,1.426556)--(-0.624289,1.569256)--cycle;
\path[fill=blue!38!white,draw=blue!38!white,line width=.12pt] (0.924449,1.388587)--(0.911234,1.236837)--(0.612409,1.274262)--(0.621290,1.426556)--cycle;
\path[fill=blue!36!white,draw=blue!36!white,line width=.12pt] (-0.621290,1.426556)--(-0.612409,1.274262)--(-0.911234,1.236837)--(-0.924449,1.388587)--cycle;
\path[fill=blue!41!white,draw=blue!41!white,line width=.12pt] (0.278805,0.839321)--(0.265249,0.706475)--(0.000000,0.712991)--(0.000000,0.846169)--cycle;
\path[fill=blue!41!white,draw=blue!41!white,line width=.12pt] (0.000000,0.846169)--(0.000000,0.712991)--(-0.265249,0.706475)--(-0.278805,0.839321)--cycle;
\path[fill=blue!42!white,draw=blue!42!white,line width=.12pt] (0.860893,0.929094)--(0.825702,0.784928)--(0.554926,0.818841)--(0.578577,0.964452)--cycle;
\path[fill=blue!40!white,draw=blue!40!white,line width=.12pt] (-0.578577,0.964452)--(-0.554926,0.818841)--(-0.825702,0.784928)--(-0.860893,0.929094)--cycle;
\path[fill=blue!42!white,draw=blue!42!white,line width=.12pt] (0.554926,0.818841)--(0.527943,0.686992)--(0.265249,0.706475)--(0.278805,0.839321)--cycle;
\path[fill=blue!41!white,draw=blue!41!white,line width=.12pt] (-0.278805,0.839321)--(-0.265249,0.706475)--(-0.527943,0.686992)--(-0.554926,0.818841)--cycle;
\path[fill=blue!33!white,draw=blue!33!white,line width=.12pt] (0.312148,1.719809)--(0.313655,1.592296)--(0.000000,1.600000)--(0.000000,1.727477)--cycle;
\path[fill=blue!32!white,draw=blue!32!white,line width=.12pt] (0.000000,1.727477)--(0.000000,1.600000)--(-0.313655,1.592296)--(-0.312148,1.719809)--cycle;
\path[fill=blue!36!white,draw=blue!36!white,line width=.12pt] (0.928911,1.531105)--(0.924449,1.388587)--(0.621290,1.426556)--(0.624289,1.569256)--cycle;
\path[fill=blue!34!white,draw=blue!34!white,line width=.12pt] (-0.624289,1.569256)--(-0.621290,1.426556)--(-0.924449,1.388587)--(-0.928911,1.531105)--cycle;
\path[fill=blue!37!white,draw=blue!37!white,line width=.12pt] (-0.911234,1.236837)--(-0.889774,1.081685)--(-1.172992,1.031016)--(-1.201283,1.184946)--cycle;
\path[fill=blue!40!white,draw=blue!40!white,line width=.12pt] (1.201283,1.184946)--(1.172992,1.031016)--(0.889774,1.081685)--(0.911234,1.236837)--cycle;
\path[fill=blue!41!white,draw=blue!41!white,line width=.12pt] (1.172992,1.031016)--(1.134919,0.880070)--(0.860893,0.929094)--(0.889774,1.081685)--cycle;
\path[fill=blue!39!white,draw=blue!39!white,line width=.12pt] (-0.889774,1.081685)--(-0.860893,0.929094)--(-1.134919,0.880070)--(-1.172992,1.031016)--cycle;
\path[fill=blue!32!white,draw=blue!32!white,line width=.12pt] (-0.312148,1.719809)--(-0.313655,1.592296)--(-0.624289,1.569256)--(-0.621290,1.696881)--cycle;
\path[fill=blue!33!white,draw=blue!33!white,line width=.12pt] (0.621290,1.696881)--(0.624289,1.569256)--(0.313655,1.592296)--(0.312148,1.719809)--cycle;
\path[fill=blue!38!white,draw=blue!38!white,line width=.12pt] (1.218704,1.335944)--(1.201283,1.184946)--(0.911234,1.236837)--(0.924449,1.388587)--cycle;
\path[fill=blue!35!white,draw=blue!35!white,line width=.12pt] (-0.924449,1.388587)--(-0.911234,1.236837)--(-1.201283,1.184946)--(-1.218704,1.335944)--cycle;
\path[fill=blue!42!white,draw=blue!42!white,line width=.12pt] (0.825702,0.784928)--(0.785553,0.654728)--(0.527943,0.686992)--(0.554926,0.818841)--cycle;
\path[fill=blue!41!white,draw=blue!41!white,line width=.12pt] (-0.554926,0.818841)--(-0.527943,0.686992)--(-0.785553,0.654728)--(-0.825702,0.784928)--cycle;
\path[fill=blue!42!white,draw=blue!42!white,line width=.12pt] (1.134919,0.880070)--(1.088526,0.737908)--(0.825702,0.784928)--(0.860893,0.929094)--cycle;
\path[fill=blue!40!white,draw=blue!40!white,line width=.12pt] (-0.860893,0.929094)--(-0.825702,0.784928)--(-1.088526,0.737908)--(-1.134919,0.880070)--cycle;
\path[fill=blue!41!white,draw=blue!41!white,line width=.12pt] (0.265249,0.706475)--(0.250539,0.592374)--(0.000000,0.598528)--(0.000000,0.712991)--cycle;
\path[fill=blue!41!white,draw=blue!41!white,line width=.12pt] (0.000000,0.712991)--(0.000000,0.598528)--(-0.250539,0.592374)--(-0.265249,0.706475)--cycle;
\path[fill=blue!33!white,draw=blue!33!white,line width=.12pt] (0.924449,1.658912)--(0.928911,1.531105)--(0.624289,1.569256)--(0.621290,1.696881)--cycle;
\path[fill=blue!31!white,draw=blue!31!white,line width=.12pt] (-0.621290,1.696881)--(-0.624289,1.569256)--(-0.928911,1.531105)--(-0.924449,1.658912)--cycle;
\path[fill=blue!33!white,draw=blue!33!white,line width=.12pt] (-0.928911,1.531105)--(-0.924449,1.388587)--(-1.218704,1.335944)--(-1.224587,1.478207)--cycle;
\path[fill=blue!36!white,draw=blue!36!white,line width=.12pt] (1.224587,1.478207)--(1.218704,1.335944)--(0.924449,1.388587)--(0.928911,1.531105)--cycle;
\path[fill=blue!41!white,draw=blue!41!white,line width=.12pt] (-0.265249,0.706475)--(-0.250539,0.592374)--(-0.498665,0.573971)--(-0.527943,0.686992)--cycle;
\path[fill=blue!41!white,draw=blue!41!white,line width=.12pt] (0.527943,0.686992)--(0.498665,0.573971)--(0.250539,0.592374)--(0.265249,0.706475)--cycle;
\path[fill=blue!30!white,draw=blue!30!white,line width=.12pt] (0.000000,1.834683)--(0.000000,1.727477)--(-0.312148,1.719809)--(-0.307686,1.827125)--cycle;
\path[fill=blue!30!white,draw=blue!30!white,line width=.12pt] (0.307686,1.827125)--(0.312148,1.719809)--(0.000000,1.727477)--(0.000000,1.834683)--cycle;
\path[fill=blue!40!white,draw=blue!40!white,line width=.12pt] (1.479763,1.119090)--(1.444914,0.966711)--(1.172992,1.031016)--(1.201283,1.184946)--cycle;
\path[fill=blue!37!white,draw=blue!37!white,line width=.12pt] (-1.201283,1.184946)--(-1.172992,1.031016)--(-1.444914,0.966711)--(-1.479763,1.119090)--cycle;
\path[fill=blue!42!white,draw=blue!42!white,line width=.12pt] (1.088526,0.737908)--(1.035598,0.609994)--(0.785553,0.654728)--(0.825702,0.784928)--cycle;
\path[fill=blue!40!white,draw=blue!40!white,line width=.12pt] (-0.825702,0.784928)--(-0.785553,0.654728)--(-1.035598,0.609994)--(-1.088526,0.737908)--cycle;
\path[fill=blue!41!white,draw=blue!41!white,line width=.12pt] (1.444914,0.966711)--(1.398014,0.817853)--(1.134919,0.880070)--(1.172992,1.031016)--cycle;
\path[fill=blue!38!white,draw=blue!38!white,line width=.12pt] (-1.172992,1.031016)--(-1.134919,0.880070)--(-1.398014,0.817853)--(-1.444914,0.966711)--cycle;
\path[fill=blue!29!white,draw=blue!29!white,line width=.12pt] (-0.307686,1.827125)--(-0.312148,1.719809)--(-0.621290,1.696881)--(-0.612409,1.804524)--cycle;
\path[fill=blue!30!white,draw=blue!30!white,line width=.12pt] (0.612409,1.804524)--(0.621290,1.696881)--(0.312148,1.719809)--(0.307686,1.827125)--cycle;
\path[fill=blue!41!white,draw=blue!41!white,line width=.12pt] (0.785553,0.654728)--(0.741989,0.543496)--(0.498665,0.573971)--(0.527943,0.686992)--cycle;
\path[fill=blue!41!white,draw=blue!41!white,line width=.12pt] (-0.527943,0.686992)--(-0.498665,0.573971)--(-0.741989,0.543496)--(-0.785553,0.654728)--cycle;
\path[fill=blue!38!white,draw=blue!38!white,line width=.12pt] (1.501223,1.269133)--(1.479763,1.119090)--(1.201283,1.184946)--(1.218704,1.335944)--cycle;
\path[fill=blue!35!white,draw=blue!35!white,line width=.12pt] (-1.218704,1.335944)--(-1.201283,1.184946)--(-1.479763,1.119090)--(-1.501223,1.269133)--cycle;
\path[fill=blue!42!white,draw=blue!42!white,line width=.12pt] (1.398014,0.817853)--(1.340867,0.678234)--(1.088526,0.737908)--(1.134919,0.880070)--cycle;
\path[fill=blue!39!white,draw=blue!39!white,line width=.12pt] (-1.134919,0.880070)--(-1.088526,0.737908)--(-1.340867,0.678234)--(-1.398014,0.817853)--cycle;
\path[fill=blue!31!white,draw=blue!31!white,line width=.12pt] (-0.924449,1.658912)--(-0.928911,1.531105)--(-1.224587,1.478207)--(-1.218704,1.606269)--cycle;
\path[fill=blue!34!white,draw=blue!34!white,line width=.12pt] (1.218704,1.606269)--(1.224587,1.478207)--(0.928911,1.531105)--(0.924449,1.658912)--cycle;
\path[fill=blue!31!white,draw=blue!31!white,line width=.12pt] (0.911234,1.767098)--(0.924449,1.658912)--(0.621290,1.696881)--(0.612409,1.804524)--cycle;
\path[fill=blue!29!white,draw=blue!29!white,line width=.12pt] (-0.612409,1.804524)--(-0.621290,1.696881)--(-0.924449,1.658912)--(-0.911234,1.767098)--cycle;
\path[fill=blue!40!white,draw=blue!40!white,line width=.12pt] (0.250539,0.592374)--(0.235241,0.501401)--(0.000000,0.507180)--(0.000000,0.598528)--cycle;
\path[fill=blue!40!white,draw=blue!40!white,line width=.12pt] (0.000000,0.598528)--(0.000000,0.507180)--(-0.235241,0.501401)--(-0.250539,0.592374)--cycle;
\path[fill=blue!36!white,draw=blue!36!white,line width=.12pt] (1.508470,1.411074)--(1.501223,1.269133)--(1.218704,1.335944)--(1.224587,1.478207)--cycle;
\path[fill=blue!32!white,draw=blue!32!white,line width=.12pt] (-1.224587,1.478207)--(-1.218704,1.335944)--(-1.501223,1.269133)--(-1.508470,1.411074)--cycle;
\path[fill=blue!40!white,draw=blue!40!white,line width=.12pt] (-0.785553,0.654728)--(-0.741989,0.543496)--(-0.978167,0.501243)--(-1.035598,0.609994)--cycle;
\path[fill=blue!41!white,draw=blue!41!white,line width=.12pt] (1.035598,0.609994)--(0.978167,0.501243)--(0.741989,0.543496)--(0.785553,0.654728)--cycle;
\path[fill=blue!42!white,draw=blue!42!white,line width=.12pt] (1.340867,0.678234)--(1.275669,0.553222)--(1.035598,0.609994)--(1.088526,0.737908)--cycle;
\path[fill=blue!40!white,draw=blue!40!white,line width=.12pt] (-1.088526,0.737908)--(-1.035598,0.609994)--(-1.275669,0.553222)--(-1.340867,0.678234)--cycle;
\path[fill=blue!40!white,draw=blue!40!white,line width=.12pt] (-0.250539,0.592374)--(-0.235241,0.501401)--(-0.468217,0.484122)--(-0.498665,0.573971)--cycle;
\path[fill=blue!40!white,draw=blue!40!white,line width=.12pt] (0.498665,0.573971)--(0.468217,0.484122)--(0.235241,0.501401)--(0.250539,0.592374)--cycle;
\path[fill=blue!27!white,draw=blue!27!white,line width=.12pt] (0.000000,1.917498)--(0.000000,1.834683)--(-0.307686,1.827125)--(-0.300440,1.910118)--cycle;
\path[fill=blue!27!white,draw=blue!27!white,line width=.12pt] (0.300440,1.910118)--(0.307686,1.827125)--(0.000000,1.834683)--(0.000000,1.917498)--cycle;
\path[fill=blue!41!white,draw=blue!41!white,line width=.12pt] (1.702920,0.889390)--(1.647647,0.743041)--(1.398014,0.817853)--(1.444914,0.966711)--cycle;
\path[fill=blue!38!white,draw=blue!38!white,line width=.12pt] (-1.444914,0.966711)--(-1.398014,0.817853)--(-1.647647,0.743041)--(-1.702920,0.889390)--cycle;
\path[fill=blue!40!white,draw=blue!40!white,line width=.12pt] (1.743993,1.039904)--(1.702920,0.889390)--(1.444914,0.966711)--(1.479763,1.119090)--cycle;
\path[fill=blue!36!white,draw=blue!36!white,line width=.12pt] (-1.479763,1.119090)--(-1.444914,0.966711)--(-1.702920,0.889390)--(-1.743993,1.039904)--cycle;
\path[fill=blue!40!white,draw=blue!40!white,line width=.12pt] (0.741989,0.543496)--(0.696683,0.455508)--(0.468217,0.484122)--(0.498665,0.573971)--cycle;
\path[fill=blue!40!white,draw=blue!40!white,line width=.12pt] (-0.498665,0.573971)--(-0.468217,0.484122)--(-0.696683,0.455508)--(-0.741989,0.543496)--cycle;
\path[fill=blue!28!white,draw=blue!28!white,line width=.12pt] (0.597986,1.888050)--(0.612409,1.804524)--(0.307686,1.827125)--(0.300440,1.910118)--cycle;
\path[fill=blue!26!white,draw=blue!26!white,line width=.12pt] (-0.300440,1.910118)--(-0.307686,1.827125)--(-0.612409,1.804524)--(-0.597986,1.888050)--cycle;
\path[fill=blue!39!white,draw=blue!39!white,line width=.12pt] (-1.398014,0.817853)--(-1.340867,0.678234)--(-1.580295,0.606480)--(-1.647647,0.743041)--cycle;
\path[fill=blue!42!white,draw=blue!42!white,line width=.12pt] (1.647647,0.743041)--(1.580295,0.606480)--(1.340867,0.678234)--(1.398014,0.817853)--cycle;
\path[fill=blue!34!white,draw=blue!34!white,line width=.12pt] (-1.501223,1.269133)--(-1.479763,1.119090)--(-1.743993,1.039904)--(-1.769285,1.188798)--cycle;
\path[fill=blue!38!white,draw=blue!38!white,line width=.12pt] (1.769285,1.188798)--(1.743993,1.039904)--(1.479763,1.119090)--(1.501223,1.269133)--cycle;
\path[fill=blue!31!white,draw=blue!31!white,line width=.12pt] (1.201283,1.715208)--(1.218704,1.606269)--(0.924449,1.658912)--(0.911234,1.767098)--cycle;
\path[fill=blue!28!white,draw=blue!28!white,line width=.12pt] (-0.911234,1.767098)--(-0.924449,1.658912)--(-1.218704,1.606269)--(-1.201283,1.715208)--cycle;
\path[fill=blue!30!white,draw=blue!30!white,line width=.12pt] (-1.218704,1.606269)--(-1.224587,1.478207)--(-1.508470,1.411074)--(-1.501223,1.539458)--cycle;
\path[fill=blue!34!white,draw=blue!34!white,line width=.12pt] (1.501223,1.539458)--(1.508470,1.411074)--(1.224587,1.478207)--(1.218704,1.606269)--cycle;
\path[fill=blue!40!white,draw=blue!40!white,line width=.12pt] (-1.035598,0.609994)--(-0.978167,0.501243)--(-1.204924,0.447619)--(-1.275669,0.553222)--cycle;
\path[fill=blue!41!white,draw=blue!41!white,line width=.12pt] (1.275669,0.553222)--(1.204924,0.447619)--(0.978167,0.501243)--(1.035598,0.609994)--cycle;
\path[fill=blue!26!white,draw=blue!26!white,line width=.12pt] (-0.597986,1.888050)--(-0.612409,1.804524)--(-0.911234,1.767098)--(-0.889774,1.851505)--cycle;
\path[fill=blue!28!white,draw=blue!28!white,line width=.12pt] (0.889774,1.851505)--(0.911234,1.767098)--(0.612409,1.804524)--(0.597986,1.888050)--cycle;
\path[fill=blue!40!white,draw=blue!40!white,line width=.12pt] (0.978167,0.501243)--(0.918440,0.415835)--(0.696683,0.455508)--(0.741989,0.543496)--cycle;
\path[fill=blue!40!white,draw=blue!40!white,line width=.12pt] (-0.741989,0.543496)--(-0.696683,0.455508)--(-0.918440,0.415835)--(-0.978167,0.501243)--cycle;
\path[fill=blue!40!white,draw=blue!40!white,line width=.12pt] (-1.340867,0.678234)--(-1.275669,0.553222)--(-1.503455,0.484957)--(-1.580295,0.606480)--cycle;
\path[fill=blue!42!white,draw=blue!42!white,line width=.12pt] (1.580295,0.606480)--(1.503455,0.484957)--(1.275669,0.553222)--(1.340867,0.678234)--cycle;
\path[fill=blue!32!white,draw=blue!32!white,line width=.12pt] (-1.508470,1.411074)--(-1.501223,1.269133)--(-1.769285,1.188798)--(-1.777825,1.330351)--cycle;
\path[fill=blue!36!white,draw=blue!36!white,line width=.12pt] (1.777825,1.330351)--(1.769285,1.188798)--(1.501223,1.269133)--(1.508470,1.411074)--cycle;
\path[fill=blue!39!white,draw=blue!39!white,line width=.12pt] (0.000000,0.507180)--(0.000000,0.442456)--(-0.219943,0.437053)--(-0.235241,0.501401)--cycle;
\path[fill=blue!39!white,draw=blue!39!white,line width=.12pt] (0.235241,0.501401)--(0.219943,0.437053)--(0.000000,0.442456)--(0.000000,0.507180)--cycle;
\path[fill=blue!39!white,draw=blue!39!white,line width=.12pt] (-0.235241,0.501401)--(-0.219943,0.437053)--(-0.437769,0.420898)--(-0.468217,0.484122)--cycle;
\path[fill=blue!39!white,draw=blue!39!white,line width=.12pt] (0.468217,0.484122)--(0.437769,0.420898)--(0.219943,0.437053)--(0.235241,0.501401)--cycle;
\path[fill=blue!31!white,draw=blue!31!white,line width=.12pt] (1.479763,1.649352)--(1.501223,1.539458)--(1.218704,1.606269)--(1.201283,1.715208)--cycle;
\path[fill=blue!27!white,draw=blue!27!white,line width=.12pt] (-1.201283,1.715208)--(-1.218704,1.606269)--(-1.501223,1.539458)--(-1.479763,1.649352)--cycle;
\path[fill=blue!24!white,draw=blue!24!white,line width=.12pt] (0.290688,1.965600)--(0.300440,1.910118)--(0.000000,1.917498)--(0.000000,1.972741)--cycle;
\path[fill=blue!24!white,draw=blue!24!white,line width=.12pt] (0.000000,1.972741)--(0.000000,1.917498)--(-0.300440,1.910118)--(-0.290688,1.965600)--cycle;
\path[fill=blue!41!white,draw=blue!41!white,line width=.12pt] (1.944527,0.799796)--(1.881411,0.656355)--(1.647647,0.743041)--(1.702920,0.889390)--cycle;
\path[fill=blue!37!white,draw=blue!37!white,line width=.12pt] (-1.702920,0.889390)--(-1.647647,0.743041)--(-1.881411,0.656355)--(-1.944527,0.799796)--cycle;
\path[fill=blue!39!white,draw=blue!39!white,line width=.12pt] (0.696683,0.455508)--(0.651378,0.394145)--(0.437769,0.420898)--(0.468217,0.484122)--cycle;
\path[fill=blue!39!white,draw=blue!39!white,line width=.12pt] (-0.468217,0.484122)--(-0.437769,0.420898)--(-0.651378,0.394145)--(-0.696683,0.455508)--cycle;
\path[fill=blue!40!white,draw=blue!40!white,line width=.12pt] (1.991426,0.948149)--(1.944527,0.799796)--(1.702920,0.889390)--(1.743993,1.039904)--cycle;
\path[fill=blue!35!white,draw=blue!35!white,line width=.12pt] (-1.743993,1.039904)--(-1.702920,0.889390)--(-1.944527,0.799796)--(-1.991426,0.948149)--cycle;
\path[fill=blue!28!white,draw=blue!28!white,line width=.12pt] (1.172992,1.800837)--(1.201283,1.715208)--(0.911234,1.767098)--(0.889774,1.851505)--cycle;
\path[fill=blue!25!white,draw=blue!25!white,line width=.12pt] (-0.889774,1.851505)--(-0.911234,1.767098)--(-1.201283,1.715208)--(-1.172992,1.800837)--cycle;
\path[fill=blue!40!white,draw=blue!40!white,line width=.12pt] (1.204924,0.447619)--(1.131352,0.365485)--(0.918440,0.415835)--(0.978167,0.501243)--cycle;
\path[fill=blue!40!white,draw=blue!40!white,line width=.12pt] (-0.978167,0.501243)--(-0.918440,0.415835)--(-1.131352,0.365485)--(-1.204924,0.447619)--cycle;
\path[fill=blue!40!white,draw=blue!40!white,line width=.12pt] (-1.275669,0.553222)--(-1.204924,0.447619)--(-1.420078,0.383140)--(-1.503455,0.484957)--cycle;
\path[fill=blue!41!white,draw=blue!41!white,line width=.12pt] (1.503455,0.484957)--(1.420078,0.383140)--(1.204924,0.447619)--(1.275669,0.553222)--cycle;
\path[fill=blue!29!white,draw=blue!29!white,line width=.12pt] (-1.501223,1.539458)--(-1.508470,1.411074)--(-1.777825,1.330351)--(-1.769285,1.459123)--cycle;
\path[fill=blue!34!white,draw=blue!34!white,line width=.12pt] (1.769285,1.459123)--(1.777825,1.330351)--(1.508470,1.411074)--(1.501223,1.539458)--cycle;
\path[fill=blue!25!white,draw=blue!25!white,line width=.12pt] (0.578577,1.944248)--(0.597986,1.888050)--(0.300440,1.910118)--(0.290688,1.965600)--cycle;
\path[fill=blue!24!white,draw=blue!24!white,line width=.12pt] (-0.290688,1.965600)--(-0.300440,1.910118)--(-0.597986,1.888050)--(-0.578577,1.944248)--cycle;
\path[fill=blue!42!white,draw=blue!42!white,line width=.12pt] (1.881411,0.656355)--(1.804504,0.523338)--(1.580295,0.606480)--(1.647647,0.743041)--cycle;
\path[fill=blue!38!white,draw=blue!38!white,line width=.12pt] (-1.647647,0.743041)--(-1.580295,0.606480)--(-1.804504,0.523338)--(-1.881411,0.656355)--cycle;
\path[fill=blue!38!white,draw=blue!38!white,line width=.12pt] (2.020307,1.095713)--(1.991426,0.948149)--(1.743993,1.039904)--(1.769285,1.188798)--cycle;
\path[fill=blue!33!white,draw=blue!33!white,line width=.12pt] (-1.769285,1.188798)--(-1.743993,1.039904)--(-1.991426,0.948149)--(-2.020307,1.095713)--cycle;
\path[fill=blue!39!white,draw=blue!39!white,line width=.12pt] (0.918440,0.415835)--(0.858714,0.357051)--(0.651378,0.394145)--(0.696683,0.455508)--cycle;
\path[fill=blue!39!white,draw=blue!39!white,line width=.12pt] (-0.696683,0.455508)--(-0.651378,0.394145)--(-0.858714,0.357051)--(-0.918440,0.415835)--cycle;
\path[fill=blue!39!white,draw=blue!39!white,line width=.12pt] (-1.580295,0.606480)--(-1.503455,0.484957)--(-1.716761,0.405857)--(-1.804504,0.523338)--cycle;
\path[fill=blue!42!white,draw=blue!42!white,line width=.12pt] (1.804504,0.523338)--(1.716761,0.405857)--(1.503455,0.484957)--(1.580295,0.606480)--cycle;
\path[fill=blue!25!white,draw=blue!25!white,line width=.12pt] (0.860893,1.908890)--(0.889774,1.851505)--(0.597986,1.888050)--(0.578577,1.944248)--cycle;
\path[fill=blue!23!white,draw=blue!23!white,line width=.12pt] (-0.578577,1.944248)--(-0.597986,1.888050)--(-0.889774,1.851505)--(-0.860893,1.908890)--cycle;
\path[fill=blue!31!white,draw=blue!31!white,line width=.12pt] (-1.777825,1.330351)--(-1.769285,1.188798)--(-2.020307,1.095713)--(-2.030059,1.236817)--cycle;
\path[fill=blue!36!white,draw=blue!36!white,line width=.12pt] (2.030059,1.236817)--(2.020307,1.095713)--(1.769285,1.188798)--(1.777825,1.330351)--cycle;
\path[fill=blue!37!white,draw=blue!37!white,line width=.12pt] (0.000000,0.442456)--(0.000000,0.406844)--(-0.205234,0.401803)--(-0.219943,0.437053)--cycle;
\path[fill=blue!37!white,draw=blue!37!white,line width=.12pt] (0.219943,0.437053)--(0.205234,0.401803)--(0.000000,0.406844)--(0.000000,0.442456)--cycle;
\path[fill=blue!28!white,draw=blue!28!white,line width=.12pt] (1.444914,1.736532)--(1.479763,1.649352)--(1.201283,1.715208)--(1.172992,1.800837)--cycle;
\path[fill=blue!25!white,draw=blue!25!white,line width=.12pt] (-1.172992,1.800837)--(-1.201283,1.715208)--(-1.479763,1.649352)--(-1.444914,1.736532)--cycle;
\path[fill=blue!37!white,draw=blue!37!white,line width=.12pt] (-0.219943,0.437053)--(-0.205234,0.401803)--(-0.408491,0.386728)--(-0.437769,0.420898)--cycle;
\path[fill=blue!37!white,draw=blue!37!white,line width=.12pt] (0.437769,0.420898)--(0.408491,0.386728)--(0.205234,0.401803)--(0.219943,0.437053)--cycle;
\path[fill=blue!40!white,draw=blue!40!white,line width=.12pt] (-1.204924,0.447619)--(-1.131352,0.365485)--(-1.333369,0.304943)--(-1.420078,0.383140)--cycle;
\path[fill=blue!40!white,draw=blue!40!white,line width=.12pt] (1.420078,0.383140)--(1.333369,0.304943)--(1.131352,0.365485)--(1.204924,0.447619)--cycle;
\path[fill=blue!27!white,draw=blue!27!white,line width=.12pt] (-1.479763,1.649352)--(-1.501223,1.539458)--(-1.769285,1.459123)--(-1.743993,1.570165)--cycle;
\path[fill=blue!31!white,draw=blue!31!white,line width=.12pt] (1.743993,1.570165)--(1.769285,1.459123)--(1.501223,1.539458)--(1.479763,1.649352)--cycle;
\path[fill=blue!39!white,draw=blue!39!white,line width=.12pt] (1.131352,0.365485)--(1.057780,0.309976)--(0.858714,0.357051)--(0.918440,0.415835)--cycle;
\path[fill=blue!39!white,draw=blue!39!white,line width=.12pt] (-0.918440,0.415835)--(-0.858714,0.357051)--(-1.057780,0.309976)--(-1.131352,0.365485)--cycle;
\path[fill=blue!40!white,draw=blue!40!white,line width=.12pt] (-1.503455,0.484957)--(-1.420078,0.383140)--(-1.621555,0.308427)--(-1.716761,0.405857)--cycle;
\path[fill=blue!41!white,draw=blue!41!white,line width=.12pt] (1.716761,0.405857)--(1.621555,0.308427)--(1.420078,0.383140)--(1.503455,0.484957)--cycle;
\path[fill=blue!25!white,draw=blue!25!white,line width=.12pt] (1.134919,1.859866)--(1.172992,1.800837)--(0.889774,1.851505)--(0.860893,1.908890)--cycle;
\path[fill=blue!23!white,draw=blue!23!white,line width=.12pt] (-0.860893,1.908890)--(-0.889774,1.851505)--(-1.172992,1.800837)--(-1.134919,1.859866)--cycle;
\path[fill=blue!22!white,draw=blue!22!white,line width=.12pt] (0.278805,1.991439)--(0.290688,1.965600)--(0.000000,1.972741)--(0.000000,1.998287)--cycle;
\path[fill=blue!21!white,draw=blue!21!white,line width=.12pt] (0.000000,1.998287)--(0.000000,1.972741)--(-0.290688,1.965600)--(-0.278805,1.991439)--cycle;
\path[fill=blue!37!white,draw=blue!37!white,line width=.12pt] (0.651378,0.394145)--(0.607814,0.361764)--(0.408491,0.386728)--(0.437769,0.420898)--cycle;
\path[fill=blue!37!white,draw=blue!37!white,line width=.12pt] (-0.437769,0.420898)--(-0.408491,0.386728)--(-0.607814,0.361764)--(-0.651378,0.394145)--cycle;
\path[fill=blue!41!white,draw=blue!41!white,line width=.12pt] (2.167407,0.698793)--(2.097056,0.558630)--(1.881411,0.656355)--(1.944527,0.799796)--cycle;
\path[fill=blue!37!white,draw=blue!37!white,line width=.12pt] (-1.944527,0.799796)--(-1.881411,0.656355)--(-2.097056,0.558630)--(-2.167407,0.698793)--cycle;
\path[fill=blue!34!white,draw=blue!34!white,line width=.12pt] (2.020307,1.366038)--(2.030059,1.236817)--(1.777825,1.330351)--(1.769285,1.459123)--cycle;
\path[fill=blue!28!white,draw=blue!28!white,line width=.12pt] (-1.769285,1.459123)--(-1.777825,1.330351)--(-2.030059,1.236817)--(-2.020307,1.366038)--cycle;
\path[fill=blue!22!white,draw=blue!22!white,line width=.12pt] (0.554926,1.970959)--(0.578577,1.944248)--(0.290688,1.965600)--(0.278805,1.991439)--cycle;
\path[fill=blue!21!white,draw=blue!21!white,line width=.12pt] (-0.278805,1.991439)--(-0.290688,1.965600)--(-0.578577,1.944248)--(-0.554926,1.970959)--cycle;
\path[fill=blue!40!white,draw=blue!40!white,line width=.12pt] (2.219681,0.844710)--(2.167407,0.698793)--(1.944527,0.799796)--(1.991426,0.948149)--cycle;
\path[fill=blue!35!white,draw=blue!35!white,line width=.12pt] (-1.991426,0.948149)--(-1.944527,0.799796)--(-2.167407,0.698793)--(-2.219681,0.844710)--cycle;
\path[fill=blue!42!white,draw=blue!42!white,line width=.12pt] (2.097056,0.558630)--(2.011334,0.429608)--(1.804504,0.523338)--(1.881411,0.656355)--cycle;
\path[fill=blue!38!white,draw=blue!38!white,line width=.12pt] (-1.881411,0.656355)--(-1.804504,0.523338)--(-2.011334,0.429608)--(-2.097056,0.558630)--cycle;
\path[fill=blue!37!white,draw=blue!37!white,line width=.12pt] (0.858714,0.357051)--(0.801283,0.327152)--(0.607814,0.361764)--(0.651378,0.394145)--cycle;
\path[fill=blue!38!white,draw=blue!38!white,line width=.12pt] (-0.651378,0.394145)--(-0.607814,0.361764)--(-0.801283,0.327152)--(-0.858714,0.357051)--cycle;
\path[fill=blue!33!white,draw=blue!33!white,line width=.12pt] (-2.020307,1.095713)--(-1.991426,0.948149)--(-2.219681,0.844710)--(-2.251872,0.990774)--cycle;
\path[fill=blue!38!white,draw=blue!38!white,line width=.12pt] (2.251872,0.990774)--(2.219681,0.844710)--(1.991426,0.948149)--(2.020307,1.095713)--cycle;
\path[fill=blue!22!white,draw=blue!22!white,line width=.12pt] (0.825702,1.937046)--(0.860893,1.908890)--(0.578577,1.944248)--(0.554926,1.970959)--cycle;
\path[fill=blue!21!white,draw=blue!21!white,line width=.12pt] (-0.554926,1.970959)--(-0.578577,1.944248)--(-0.860893,1.908890)--(-0.825702,1.937046)--cycle;
\path[fill=blue!39!white,draw=blue!39!white,line width=.12pt] (-1.804504,0.523338)--(-1.716761,0.405857)--(-1.913535,0.316685)--(-2.011334,0.429608)--cycle;
\path[fill=blue!42!white,draw=blue!42!white,line width=.12pt] (2.011334,0.429608)--(1.913535,0.316685)--(1.716761,0.405857)--(1.804504,0.523338)--cycle;
\path[fill=blue!39!white,draw=blue!39!white,line width=.12pt] (-1.131352,0.365485)--(-1.057780,0.309976)--(-1.246659,0.253371)--(-1.333369,0.304943)--cycle;
\path[fill=blue!39!white,draw=blue!39!white,line width=.12pt] (1.333369,0.304943)--(1.246659,0.253371)--(1.057780,0.309976)--(1.131352,0.365485)--cycle;
\path[fill=blue!40!white,draw=blue!40!white,line width=.12pt] (-1.420078,0.383140)--(-1.333369,0.304943)--(-1.522544,0.234792)--(-1.621555,0.308427)--cycle;
\path[fill=blue!40!white,draw=blue!40!white,line width=.12pt] (1.621555,0.308427)--(1.522544,0.234792)--(1.333369,0.304943)--(1.420078,0.383140)--cycle;
\path[fill=blue!28!white,draw=blue!28!white,line width=.12pt] (1.702920,1.659211)--(1.743993,1.570165)--(1.479763,1.649352)--(1.444914,1.736532)--cycle;
\path[fill=blue!24!white,draw=blue!24!white,line width=.12pt] (-1.444914,1.736532)--(-1.479763,1.649352)--(-1.743993,1.570165)--(-1.702920,1.659211)--cycle;
\path[fill=blue!25!white,draw=blue!25!white,line width=.12pt] (1.398014,1.797648)--(1.444914,1.736532)--(1.172992,1.800837)--(1.134919,1.859866)--cycle;
\path[fill=blue!22!white,draw=blue!22!white,line width=.12pt] (-1.134919,1.859866)--(-1.172992,1.800837)--(-1.444914,1.736532)--(-1.398014,1.797648)--cycle;
\path[fill=blue!35!white,draw=blue!35!white,line width=.12pt] (0.205234,0.401803)--(0.191677,0.397005)--(0.000000,0.401713)--(0.000000,0.406844)--cycle;
\path[fill=blue!35!white,draw=blue!35!white,line width=.12pt] (0.000000,0.406844)--(0.000000,0.401713)--(-0.191677,0.397005)--(-0.205234,0.401803)--cycle;
\path[fill=blue!30!white,draw=blue!30!white,line width=.12pt] (-2.030059,1.236817)--(-2.020307,1.095713)--(-2.251872,0.990774)--(-2.262742,1.131371)--cycle;
\path[fill=blue!36!white,draw=blue!36!white,line width=.12pt] (2.262742,1.131371)--(2.251872,0.990774)--(2.020307,1.095713)--(2.030059,1.236817)--cycle;
\path[fill=blue!35!white,draw=blue!35!white,line width=.12pt] (0.408491,0.386728)--(0.381508,0.382925)--(0.191677,0.397005)--(0.205234,0.401803)--cycle;
\path[fill=blue!35!white,draw=blue!35!white,line width=.12pt] (-0.205234,0.401803)--(-0.191677,0.397005)--(-0.381508,0.382925)--(-0.408491,0.386728)--cycle;
\path[fill=blue!31!white,draw=blue!31!white,line width=.12pt] (1.991426,1.478411)--(2.020307,1.366038)--(1.769285,1.459123)--(1.743993,1.570165)--cycle;
\path[fill=blue!26!white,draw=blue!26!white,line width=.12pt] (-1.743993,1.570165)--(-1.769285,1.459123)--(-2.020307,1.366038)--(-1.991426,1.478411)--cycle;
\path[fill=blue!37!white,draw=blue!37!white,line width=.12pt] (1.057780,0.309976)--(0.987036,0.283224)--(0.801283,0.327152)--(0.858714,0.357051)--cycle;
\path[fill=blue!38!white,draw=blue!38!white,line width=.12pt] (-0.858714,0.357051)--(-0.801283,0.327152)--(-0.987036,0.283224)--(-1.057780,0.309976)--cycle;
\path[fill=blue!41!white,draw=blue!41!white,line width=.12pt] (1.913535,0.316685)--(1.807416,0.224200)--(1.621555,0.308427)--(1.716761,0.405857)--cycle;
\path[fill=blue!40!white,draw=blue!40!white,line width=.12pt] (-1.716761,0.405857)--(-1.621555,0.308427)--(-1.807416,0.224200)--(-1.913535,0.316685)--cycle;
\path[fill=blue!20!white,draw=blue!20!white,line width=.12pt] (-0.825702,1.937046)--(-0.860893,1.908890)--(-1.134919,1.859866)--(-1.088526,1.890026)--cycle;
\path[fill=blue!22!white,draw=blue!22!white,line width=.12pt] (1.088526,1.890026)--(1.134919,1.859866)--(0.860893,1.908890)--(0.825702,1.937046)--cycle;
\path[fill=blue!35!white,draw=blue!35!white,line width=.12pt] (0.607814,0.361764)--(0.567664,0.359610)--(0.381508,0.382925)--(0.408491,0.386728)--cycle;
\path[fill=blue!36!white,draw=blue!36!white,line width=.12pt] (-0.408491,0.386728)--(-0.381508,0.382925)--(-0.567664,0.359610)--(-0.607814,0.361764)--cycle;
\path[fill=blue!19!white,draw=blue!19!white,line width=.12pt] (0.265249,1.986640)--(0.278805,1.991439)--(0.000000,1.998287)--(0.000000,1.993156)--cycle;
\path[fill=blue!19!white,draw=blue!19!white,line width=.12pt] (0.000000,1.993156)--(0.000000,1.998287)--(-0.278805,1.991439)--(-0.265249,1.986640)--cycle;
\path[fill=blue!19!white,draw=blue!19!white,line width=.12pt] (0.527943,1.967157)--(0.554926,1.970959)--(0.278805,1.991439)--(0.265249,1.986640)--cycle;
\path[fill=blue!19!white,draw=blue!19!white,line width=.12pt] (-0.265249,1.986640)--(-0.278805,1.991439)--(-0.554926,1.970959)--(-0.527943,1.967157)--cycle;
\path[fill=blue!39!white,draw=blue!39!white,line width=.12pt] (-1.333369,0.304943)--(-1.246659,0.253371)--(-1.423533,0.187782)--(-1.522544,0.234792)--cycle;
\path[fill=blue!39!white,draw=blue!39!white,line width=.12pt] (1.522544,0.234792)--(1.423533,0.187782)--(1.246659,0.253371)--(1.333369,0.304943)--cycle;
\path[fill=blue!28!white,draw=blue!28!white,line width=.12pt] (-2.020307,1.366038)--(-2.030059,1.236817)--(-2.262742,1.131371)--(-2.251872,1.261099)--cycle;
\path[fill=blue!34!white,draw=blue!34!white,line width=.12pt] (2.251872,1.261099)--(2.262742,1.131371)--(2.030059,1.236817)--(2.020307,1.366038)--cycle;
\path[fill=blue!41!white,draw=blue!41!white,line width=.12pt] (2.369413,0.587353)--(2.292506,0.450808)--(2.097056,0.558630)--(2.167407,0.698793)--cycle;
\path[fill=blue!36!white,draw=blue!36!white,line width=.12pt] (-2.167407,0.698793)--(-2.097056,0.558630)--(-2.292506,0.450808)--(-2.369413,0.587353)--cycle;
\path[fill=blue!34!white,draw=blue!34!white,line width=.12pt] (0.801283,0.327152)--(0.748354,0.327284)--(0.567664,0.359610)--(0.607814,0.361764)--cycle;
\path[fill=blue!36!white,draw=blue!36!white,line width=.12pt] (-0.607814,0.361764)--(-0.567664,0.359610)--(-0.748354,0.327284)--(-0.801283,0.327152)--cycle;
\path[fill=blue!42!white,draw=blue!42!white,line width=.12pt] (2.292506,0.450808)--(2.198794,0.326193)--(2.011334,0.429608)--(2.097056,0.558630)--cycle;
\path[fill=blue!38!white,draw=blue!38!white,line width=.12pt] (-2.097056,0.558630)--(-2.011334,0.429608)--(-2.198794,0.326193)--(-2.292506,0.450808)--cycle;
\path[fill=blue!37!white,draw=blue!37!white,line width=.12pt] (1.246659,0.253371)--(1.163283,0.230405)--(0.987036,0.283224)--(1.057780,0.309976)--cycle;
\path[fill=blue!38!white,draw=blue!38!white,line width=.12pt] (-1.057780,0.309976)--(-0.987036,0.283224)--(-1.163283,0.230405)--(-1.246659,0.253371)--cycle;
\path[fill=blue!40!white,draw=blue!40!white,line width=.12pt] (2.426560,0.730582)--(2.369413,0.587353)--(2.167407,0.698793)--(2.219681,0.844710)--cycle;
\path[fill=blue!34!white,draw=blue!34!white,line width=.12pt] (-2.219681,0.844710)--(-2.167407,0.698793)--(-2.369413,0.587353)--(-2.426560,0.730582)--cycle;
\path[fill=blue!25!white,draw=blue!25!white,line width=.12pt] (1.647647,1.722837)--(1.702920,1.659211)--(1.444914,1.736532)--(1.398014,1.797648)--cycle;
\path[fill=blue!22!white,draw=blue!22!white,line width=.12pt] (-1.398014,1.797648)--(-1.444914,1.736532)--(-1.702920,1.659211)--(-1.647647,1.722837)--cycle;
\path[fill=blue!40!white,draw=blue!40!white,line width=.12pt] (-1.621555,0.308427)--(-1.522544,0.234792)--(-1.697056,0.155708)--(-1.807416,0.224200)--cycle;
\path[fill=blue!40!white,draw=blue!40!white,line width=.12pt] (1.807416,0.224200)--(1.697056,0.155708)--(1.522544,0.234792)--(1.621555,0.308427)--cycle;
\path[fill=blue!18!white,draw=blue!18!white,line width=.12pt] (-0.527943,1.967157)--(-0.554926,1.970959)--(-0.825702,1.937046)--(-0.785553,1.934893)--cycle;
\path[fill=blue!19!white,draw=blue!19!white,line width=.12pt] (0.785553,1.934893)--(0.825702,1.937046)--(0.554926,1.970959)--(0.527943,1.967157)--cycle;
\path[fill=blue!28!white,draw=blue!28!white,line width=.12pt] (1.944527,1.569617)--(1.991426,1.478411)--(1.743993,1.570165)--(1.702920,1.659211)--cycle;
\path[fill=blue!23!white,draw=blue!23!white,line width=.12pt] (-1.702920,1.659211)--(-1.743993,1.570165)--(-1.991426,1.478411)--(-1.944527,1.569617)--cycle;
\path[fill=blue!22!white,draw=blue!22!white,line width=.12pt] (1.340867,1.830352)--(1.398014,1.797648)--(1.134919,1.859866)--(1.088526,1.890026)--cycle;
\path[fill=blue!20!white,draw=blue!20!white,line width=.12pt] (-1.088526,1.890026)--(-1.134919,1.859866)--(-1.398014,1.797648)--(-1.340867,1.830352)--cycle;
\path[fill=blue!39!white,draw=blue!39!white,line width=.12pt] (-2.011334,0.429608)--(-1.913535,0.316685)--(-2.091880,0.218298)--(-2.198794,0.326193)--cycle;
\path[fill=blue!42!white,draw=blue!42!white,line width=.12pt] (2.198794,0.326193)--(2.091880,0.218298)--(1.913535,0.316685)--(2.011334,0.429608)--cycle;
\path[fill=blue!32!white,draw=blue!32!white,line width=.12pt] (-2.251872,0.990774)--(-2.219681,0.844710)--(-2.426560,0.730582)--(-2.461751,0.874991)--cycle;
\path[fill=blue!38!white,draw=blue!38!white,line width=.12pt] (2.461751,0.874991)--(2.426560,0.730582)--(2.219681,0.844710)--(2.251872,0.990774)--cycle;
\path[fill=blue!32!white,draw=blue!32!white,line width=.12pt] (0.191677,0.397005)--(0.179794,0.422843)--(0.000000,0.427259)--(0.000000,0.401713)--cycle;
\path[fill=blue!33!white,draw=blue!33!white,line width=.12pt] (0.000000,0.401713)--(0.000000,0.427259)--(-0.179794,0.422843)--(-0.191677,0.397005)--cycle;
\path[fill=blue!36!white,draw=blue!36!white,line width=.12pt] (-0.801283,0.327152)--(-0.748354,0.327284)--(-0.921837,0.286259)--(-0.987036,0.283224)--cycle;
\path[fill=blue!34!white,draw=blue!34!white,line width=.12pt] (0.987036,0.283224)--(0.921837,0.286259)--(0.748354,0.327284)--(0.801283,0.327152)--cycle;
\path[fill=blue!33!white,draw=blue!33!white,line width=.12pt] (-0.191677,0.397005)--(-0.179794,0.422843)--(-0.357857,0.409636)--(-0.381508,0.382925)--cycle;
\path[fill=blue!32!white,draw=blue!32!white,line width=.12pt] (0.381508,0.382925)--(0.357857,0.409636)--(0.179794,0.422843)--(0.191677,0.397005)--cycle;
\path[fill=blue!25!white,draw=blue!25!white,line width=.12pt] (-1.991426,1.478411)--(-2.020307,1.366038)--(-2.251872,1.261099)--(-2.219681,1.374972)--cycle;
\path[fill=blue!31!white,draw=blue!31!white,line width=.12pt] (2.219681,1.374972)--(2.251872,1.261099)--(2.020307,1.366038)--(1.991426,1.478411)--cycle;
\path[fill=blue!39!white,draw=blue!39!white,line width=.12pt] (-1.913535,0.316685)--(-1.807416,0.224200)--(-1.975871,0.131270)--(-2.091880,0.218298)--cycle;
\path[fill=blue!41!white,draw=blue!41!white,line width=.12pt] (2.091880,0.218298)--(1.975871,0.131270)--(1.807416,0.224200)--(1.913535,0.316685)--cycle;
\path[fill=blue!20!white,draw=blue!20!white,line width=.12pt] (1.035598,1.890159)--(1.088526,1.890026)--(0.825702,1.937046)--(0.785553,1.934893)--cycle;
\path[fill=blue!18!white,draw=blue!18!white,line width=.12pt] (-0.785553,1.934893)--(-0.825702,1.937046)--(-1.088526,1.890026)--(-1.035598,1.890159)--cycle;
\path[fill=blue!32!white,draw=blue!32!white,line width=.12pt] (0.567664,0.359610)--(0.532473,0.387767)--(0.357857,0.409636)--(0.381508,0.382925)--cycle;
\path[fill=blue!33!white,draw=blue!33!white,line width=.12pt] (-0.381508,0.382925)--(-0.357857,0.409636)--(-0.532473,0.387767)--(-0.567664,0.359610)--cycle;
\path[fill=blue!29!white,draw=blue!29!white,line width=.12pt] (-2.262742,1.131371)--(-2.251872,0.990774)--(-2.461751,0.874991)--(-2.473633,1.015029)--cycle;
\path[fill=blue!36!white,draw=blue!36!white,line width=.12pt] (2.473633,1.015029)--(2.461751,0.874991)--(2.251872,0.990774)--(2.262742,1.131371)--cycle;
\path[fill=blue!17!white,draw=blue!17!white,line width=.12pt] (0.000000,1.957544)--(0.000000,1.993156)--(-0.265249,1.986640)--(-0.250539,1.951390)--cycle;
\path[fill=blue!17!white,draw=blue!17!white,line width=.12pt] (0.250539,1.951390)--(0.265249,1.986640)--(0.000000,1.993156)--(0.000000,1.957544)--cycle;
\path[fill=blue!37!white,draw=blue!37!white,line width=.12pt] (1.423533,0.187782)--(1.328326,0.169203)--(1.163283,0.230405)--(1.246659,0.253371)--cycle;
\path[fill=blue!38!white,draw=blue!38!white,line width=.12pt] (-1.246659,0.253371)--(-1.163283,0.230405)--(-1.328326,0.169203)--(-1.423533,0.187782)--cycle;
\path[fill=blue!39!white,draw=blue!39!white,line width=.12pt] (-1.522544,0.234792)--(-1.423533,0.187782)--(-1.586697,0.113840)--(-1.697056,0.155708)--cycle;
\path[fill=blue!39!white,draw=blue!39!white,line width=.12pt] (1.697056,0.155708)--(1.586697,0.113840)--(1.423533,0.187782)--(1.522544,0.234792)--cycle;
\path[fill=blue!17!white,draw=blue!17!white,line width=.12pt] (0.498665,1.932987)--(0.527943,1.967157)--(0.265249,1.986640)--(0.250539,1.951390)--cycle;
\path[fill=blue!17!white,draw=blue!17!white,line width=.12pt] (-0.250539,1.951390)--(-0.265249,1.986640)--(-0.527943,1.967157)--(-0.498665,1.932987)--cycle;
\path[fill=blue!36!white,draw=blue!36!white,line width=.12pt] (-0.987036,0.283224)--(-0.921837,0.286259)--(-1.086442,0.236929)--(-1.163283,0.230405)--cycle;
\path[fill=blue!34!white,draw=blue!34!white,line width=.12pt] (1.163283,0.230405)--(1.086442,0.236929)--(0.921837,0.286259)--(0.987036,0.283224)--cycle;
\path[fill=blue!34!white,draw=blue!34!white,line width=.12pt] (-0.567664,0.359610)--(-0.532473,0.387767)--(-0.701962,0.357445)--(-0.748354,0.327284)--cycle;
\path[fill=blue!32!white,draw=blue!32!white,line width=.12pt] (0.748354,0.327284)--(0.701962,0.357445)--(0.532473,0.387767)--(0.567664,0.359610)--cycle;
\path[fill=blue!19!white,draw=blue!19!white,line width=.12pt] (-1.340867,1.830352)--(-1.398014,1.797648)--(-1.647647,1.722837)--(-1.580295,1.758598)--cycle;
\path[fill=blue!22!white,draw=blue!22!white,line width=.12pt] (1.580295,1.758598)--(1.647647,1.722837)--(1.398014,1.797648)--(1.340867,1.830352)--cycle;
\path[fill=blue!21!white,draw=blue!21!white,line width=.12pt] (-1.647647,1.722837)--(-1.702920,1.659211)--(-1.944527,1.569617)--(-1.881411,1.636151)--cycle;
\path[fill=blue!25!white,draw=blue!25!white,line width=.12pt] (1.881411,1.636151)--(1.944527,1.569617)--(1.702920,1.659211)--(1.647647,1.722837)--cycle;
\path[fill=blue!40!white,draw=blue!40!white,line width=.12pt] (1.975871,0.131270)--(1.855225,0.068452)--(1.697056,0.155708)--(1.807416,0.224200)--cycle;
\path[fill=blue!40!white,draw=blue!40!white,line width=.12pt] (-1.807416,0.224200)--(-1.697056,0.155708)--(-1.855225,0.068452)--(-1.975871,0.131270)--cycle;
\path[fill=blue!17!white,draw=blue!17!white,line width=.12pt] (0.741989,1.902512)--(0.785553,1.934893)--(0.527943,1.967157)--(0.498665,1.932987)--cycle;
\path[fill=blue!17!white,draw=blue!17!white,line width=.12pt] (-0.498665,1.932987)--(-0.527943,1.967157)--(-0.785553,1.934893)--(-0.741989,1.902512)--cycle;
\path[fill=blue!20!white,draw=blue!20!white,line width=.12pt] (1.275669,1.833387)--(1.340867,1.830352)--(1.088526,1.890026)--(1.035598,1.890159)--cycle;
\path[fill=blue!18!white,draw=blue!18!white,line width=.12pt] (-1.035598,1.890159)--(-1.088526,1.890026)--(-1.340867,1.830352)--(-1.275669,1.833387)--cycle;
\path[fill=blue!34!white,draw=blue!34!white,line width=.12pt] (2.461751,1.145316)--(2.473633,1.015029)--(2.262742,1.131371)--(2.251872,1.261099)--cycle;
\path[fill=blue!27!white,draw=blue!27!white,line width=.12pt] (-2.251872,1.261099)--(-2.262742,1.131371)--(-2.473633,1.015029)--(-2.461751,1.145316)--cycle;
\path[fill=blue!41!white,draw=blue!41!white,line width=.12pt] (2.548600,0.466550)--(2.465877,0.333925)--(2.292506,0.450808)--(2.369413,0.587353)--cycle;
\path[fill=blue!35!white,draw=blue!35!white,line width=.12pt] (-2.369413,0.587353)--(-2.292506,0.450808)--(-2.465877,0.333925)--(-2.548600,0.466550)--cycle;
\path[fill=blue!42!white,draw=blue!42!white,line width=.12pt] (2.465877,0.333925)--(2.365079,0.214089)--(2.198794,0.326193)--(2.292506,0.450808)--cycle;
\path[fill=blue!37!white,draw=blue!37!white,line width=.12pt] (-2.292506,0.450808)--(-2.198794,0.326193)--(-2.365079,0.214089)--(-2.465877,0.333925)--cycle;
\path[fill=blue!28!white,draw=blue!28!white,line width=.12pt] (2.167407,1.468614)--(2.219681,1.374972)--(1.991426,1.478411)--(1.944527,1.569617)--cycle;
\path[fill=blue!23!white,draw=blue!23!white,line width=.12pt] (-1.944527,1.569617)--(-1.991426,1.478411)--(-2.219681,1.374972)--(-2.167407,1.468614)--cycle;
\path[fill=blue!30!white,draw=blue!30!white,line width=.12pt] (0.179794,0.422843)--(0.170043,0.478325)--(0.000000,0.482502)--(0.000000,0.427259)--cycle;
\path[fill=blue!30!white,draw=blue!30!white,line width=.12pt] (0.000000,0.427259)--(0.000000,0.482502)--(-0.170043,0.478325)--(-0.179794,0.422843)--cycle;
\path[fill=blue!32!white,draw=blue!32!white,line width=.12pt] (0.921837,0.286259)--(0.864690,0.318963)--(0.701962,0.357445)--(0.748354,0.327284)--cycle;
\path[fill=blue!34!white,draw=blue!34!white,line width=.12pt] (-0.748354,0.327284)--(-0.701962,0.357445)--(-0.864690,0.318963)--(-0.921837,0.286259)--cycle;
\path[fill=blue!40!white,draw=blue!40!white,line width=.12pt] (2.610069,0.606865)--(2.548600,0.466550)--(2.369413,0.587353)--(2.426560,0.730582)--cycle;
\path[fill=blue!33!white,draw=blue!33!white,line width=.12pt] (-2.426560,0.730582)--(-2.369413,0.587353)--(-2.548600,0.466550)--(-2.610069,0.606865)--cycle;
\path[fill=blue!42!white,draw=blue!42!white,line width=.12pt] (2.365079,0.214089)--(2.250079,0.111645)--(2.091880,0.218298)--(2.198794,0.326193)--cycle;
\path[fill=blue!38!white,draw=blue!38!white,line width=.12pt] (-2.198794,0.326193)--(-2.091880,0.218298)--(-2.250079,0.111645)--(-2.365079,0.214089)--cycle;
\path[fill=blue!29!white,draw=blue!29!white,line width=.12pt] (0.357857,0.409636)--(0.338447,0.465835)--(0.170043,0.478325)--(0.179794,0.422843)--cycle;
\path[fill=blue!30!white,draw=blue!30!white,line width=.12pt] (-0.179794,0.422843)--(-0.170043,0.478325)--(-0.338447,0.465835)--(-0.357857,0.409636)--cycle;
\path[fill=blue!38!white,draw=blue!38!white,line width=.12pt] (-1.423533,0.187782)--(-1.328326,0.169203)--(-1.480578,0.100206)--(-1.586697,0.113840)--cycle;
\path[fill=blue!37!white,draw=blue!37!white,line width=.12pt] (1.586697,0.113840)--(1.480578,0.100206)--(1.328326,0.169203)--(1.423533,0.187782)--cycle;
\path[fill=blue!16!white,draw=blue!16!white,line width=.12pt] (-0.741989,1.902512)--(-0.785553,1.934893)--(-1.035598,1.890159)--(-0.978167,1.860259)--cycle;
\path[fill=blue!17!white,draw=blue!17!white,line width=.12pt] (0.978167,1.860259)--(1.035598,1.890159)--(0.785553,1.934893)--(0.741989,1.902512)--cycle;
\path[fill=blue!37!white,draw=blue!37!white,line width=.12pt] (-1.163283,0.230405)--(-1.086442,0.236929)--(-1.240584,0.179769)--(-1.328326,0.169203)--cycle;
\path[fill=blue!34!white,draw=blue!34!white,line width=.12pt] (1.328326,0.169203)--(1.240584,0.179769)--(1.086442,0.236929)--(1.163283,0.230405)--cycle;
\path[fill=blue!29!white,draw=blue!29!white,line width=.12pt] (0.532473,0.387767)--(0.503593,0.445151)--(0.338447,0.465835)--(0.357857,0.409636)--cycle;
\path[fill=blue!31!white,draw=blue!31!white,line width=.12pt] (-0.357857,0.409636)--(-0.338447,0.465835)--(-0.503593,0.445151)--(-0.532473,0.387767)--cycle;
\path[fill=blue!39!white,draw=blue!39!white,line width=.12pt] (1.855225,0.068452)--(1.734580,0.032258)--(1.586697,0.113840)--(1.697056,0.155708)--cycle;
\path[fill=blue!39!white,draw=blue!39!white,line width=.12pt] (-1.697056,0.155708)--(-1.586697,0.113840)--(-1.734580,0.032258)--(-1.855225,0.068452)--cycle;
\path[fill=blue!38!white,draw=blue!38!white,line width=.12pt] (2.647922,0.749480)--(2.610069,0.606865)--(2.426560,0.730582)--(2.461751,0.874991)--cycle;
\path[fill=blue!31!white,draw=blue!31!white,line width=.12pt] (-2.461751,0.874991)--(-2.426560,0.730582)--(-2.610069,0.606865)--(-2.647922,0.749480)--cycle;
\path[fill=blue!15!white,draw=blue!15!white,line width=.12pt] (0.000000,1.892820)--(0.000000,1.957544)--(-0.250539,1.951390)--(-0.235241,1.887042)--cycle;
\path[fill=blue!15!white,draw=blue!15!white,line width=.12pt] (0.235241,1.887042)--(0.250539,1.951390)--(0.000000,1.957544)--(0.000000,1.892820)--cycle;
\path[fill=blue!41!white,draw=blue!41!white,line width=.12pt] (2.250079,0.111645)--(2.125296,0.030531)--(1.975871,0.131270)--(2.091880,0.218298)--cycle;
\path[fill=blue!39!white,draw=blue!39!white,line width=.12pt] (-2.091880,0.218298)--(-1.975871,0.131270)--(-2.125296,0.030531)--(-2.250079,0.111645)--cycle;
\path[fill=blue!15!white,draw=blue!15!white,line width=.12pt] (-0.235241,1.887042)--(-0.250539,1.951390)--(-0.498665,1.932987)--(-0.468217,1.869763)--cycle;
\path[fill=blue!15!white,draw=blue!15!white,line width=.12pt] (0.468217,1.869763)--(0.498665,1.932987)--(0.250539,1.951390)--(0.235241,1.887042)--cycle;
\path[fill=blue!31!white,draw=blue!31!white,line width=.12pt] (2.426560,1.260844)--(2.461751,1.145316)--(2.251872,1.261099)--(2.219681,1.374972)--cycle;
\path[fill=blue!24!white,draw=blue!24!white,line width=.12pt] (-2.219681,1.374972)--(-2.251872,1.261099)--(-2.461751,1.145316)--(-2.426560,1.260844)--cycle;
\path[fill=blue!35!white,draw=blue!35!white,line width=.12pt] (-0.921837,0.286259)--(-0.864690,0.318963)--(-1.019091,0.272690)--(-1.086442,0.236929)--cycle;
\path[fill=blue!32!white,draw=blue!32!white,line width=.12pt] (1.086442,0.236929)--(1.019091,0.272690)--(0.864690,0.318963)--(0.921837,0.286259)--cycle;
\path[fill=blue!22!white,draw=blue!22!white,line width=.12pt] (1.804504,1.675456)--(1.881411,1.636151)--(1.647647,1.722837)--(1.580295,1.758598)--cycle;
\path[fill=blue!19!white,draw=blue!19!white,line width=.12pt] (-1.580295,1.758598)--(-1.647647,1.722837)--(-1.881411,1.636151)--(-1.804504,1.675456)--cycle;
\path[fill=blue!18!white,draw=blue!18!white,line width=.12pt] (-1.275669,1.833387)--(-1.340867,1.830352)--(-1.580295,1.758598)--(-1.503455,1.765122)--cycle;
\path[fill=blue!20!white,draw=blue!20!white,line width=.12pt] (1.503455,1.765122)--(1.580295,1.758598)--(1.340867,1.830352)--(1.275669,1.833387)--cycle;
\path[fill=blue!29!white,draw=blue!29!white,line width=.12pt] (0.701962,0.357445)--(0.663889,0.416474)--(0.503593,0.445151)--(0.532473,0.387767)--cycle;
\path[fill=blue!31!white,draw=blue!31!white,line width=.12pt] (-0.532473,0.387767)--(-0.503593,0.445151)--(-0.663889,0.416474)--(-0.701962,0.357445)--cycle;
\path[fill=blue!36!white,draw=blue!36!white,line width=.12pt] (2.660703,0.888912)--(2.647922,0.749480)--(2.461751,0.874991)--(2.473633,1.015029)--cycle;
\path[fill=blue!29!white,draw=blue!29!white,line width=.12pt] (-2.473633,1.015029)--(-2.461751,0.874991)--(-2.647922,0.749480)--(-2.660703,0.888912)--cycle;
\path[fill=blue!15!white,draw=blue!15!white,line width=.12pt] (0.696683,1.841149)--(0.741989,1.902512)--(0.498665,1.932987)--(0.468217,1.869763)--cycle;
\path[fill=blue!15!white,draw=blue!15!white,line width=.12pt] (-0.468217,1.869763)--(-0.498665,1.932987)--(-0.741989,1.902512)--(-0.696683,1.841149)--cycle;
\path[fill=blue!25!white,draw=blue!25!white,line width=.12pt] (2.097056,1.538426)--(2.167407,1.468614)--(1.944527,1.569617)--(1.881411,1.636151)--cycle;
\path[fill=blue!21!white,draw=blue!21!white,line width=.12pt] (-1.881411,1.636151)--(-1.944527,1.569617)--(-2.167407,1.468614)--(-2.097056,1.538426)--cycle;
\path[fill=blue!17!white,draw=blue!17!white,line width=.12pt] (1.204924,1.806635)--(1.275669,1.833387)--(1.035598,1.890159)--(0.978167,1.860259)--cycle;
\path[fill=blue!16!white,draw=blue!16!white,line width=.12pt] (-0.978167,1.860259)--(-1.035598,1.890159)--(-1.275669,1.833387)--(-1.204924,1.806635)--cycle;
\path[fill=blue!40!white,draw=blue!40!white,line width=.12pt] (2.125296,0.030531)--(1.995527,-0.026136)--(1.855225,0.068452)--(1.975871,0.131270)--cycle;
\path[fill=blue!40!white,draw=blue!40!white,line width=.12pt] (-1.975871,0.131270)--(-1.855225,0.068452)--(-1.995527,-0.026136)--(-2.125296,0.030531)--cycle;
\path[fill=blue!27!white,draw=blue!27!white,line width=.12pt] (0.000000,0.482502)--(0.000000,0.565317)--(-0.162796,0.561318)--(-0.170043,0.478325)--cycle;
\path[fill=blue!27!white,draw=blue!27!white,line width=.12pt] (0.170043,0.478325)--(0.162796,0.561318)--(0.000000,0.565317)--(0.000000,0.482502)--cycle;
\path[fill=blue!37!white,draw=blue!37!white,line width=.12pt] (-1.328326,0.169203)--(-1.240584,0.179769)--(-1.382778,0.115330)--(-1.480578,0.100206)--cycle;
\path[fill=blue!34!white,draw=blue!34!white,line width=.12pt] (1.480578,0.100206)--(1.382778,0.115330)--(1.240584,0.179769)--(1.328326,0.169203)--cycle;
\path[fill=blue!29!white,draw=blue!29!white,line width=.12pt] (0.864690,0.318963)--(0.817791,0.380079)--(0.663889,0.416474)--(0.701962,0.357445)--cycle;
\path[fill=blue!32!white,draw=blue!32!white,line width=.12pt] (-0.701962,0.357445)--(-0.663889,0.416474)--(-0.817791,0.380079)--(-0.864690,0.318963)--cycle;
\path[fill=blue!39!white,draw=blue!39!white,line width=.12pt] (-1.586697,0.113840)--(-1.480578,0.100206)--(-1.618570,0.024081)--(-1.734580,0.032258)--cycle;
\path[fill=blue!37!white,draw=blue!37!white,line width=.12pt] (1.734580,0.032258)--(1.618570,0.024081)--(1.480578,0.100206)--(1.586697,0.113840)--cycle;
\path[fill=blue!26!white,draw=blue!26!white,line width=.12pt] (0.338447,0.465835)--(0.324025,0.549360)--(0.162796,0.561318)--(0.170043,0.478325)--cycle;
\path[fill=blue!28!white,draw=blue!28!white,line width=.12pt] (-0.170043,0.478325)--(-0.162796,0.561318)--(-0.324025,0.549360)--(-0.338447,0.465835)--cycle;
\path[fill=blue!32!white,draw=blue!32!white,line width=.12pt] (1.240584,0.179769)--(1.163677,0.219074)--(1.019091,0.272690)--(1.086442,0.236929)--cycle;
\path[fill=blue!35!white,draw=blue!35!white,line width=.12pt] (-1.086442,0.236929)--(-1.019091,0.272690)--(-1.163677,0.219074)--(-1.240584,0.179769)--cycle;
\path[fill=blue!15!white,draw=blue!15!white,line width=.12pt] (0.918440,1.801476)--(0.978167,1.860259)--(0.741989,1.902512)--(0.696683,1.841149)--cycle;
\path[fill=blue!15!white,draw=blue!15!white,line width=.12pt] (-0.696683,1.841149)--(-0.741989,1.902512)--(-0.978167,1.860259)--(-0.918440,1.801476)--cycle;
\path[fill=blue!28!white,draw=blue!28!white,line width=.12pt] (2.369413,1.357174)--(2.426560,1.260844)--(2.219681,1.374972)--(2.167407,1.468614)--cycle;
\path[fill=blue!22!white,draw=blue!22!white,line width=.12pt] (-2.167407,1.468614)--(-2.219681,1.374972)--(-2.426560,1.260844)--(-2.369413,1.357174)--cycle;
\path[fill=blue!28!white,draw=blue!28!white,line width=.12pt] (-0.338447,0.465835)--(-0.324025,0.549360)--(-0.482133,0.529559)--(-0.503593,0.445151)--cycle;
\path[fill=blue!26!white,draw=blue!26!white,line width=.12pt] (0.503593,0.445151)--(0.482133,0.529559)--(0.324025,0.549360)--(0.338447,0.465835)--cycle;
\path[fill=blue!14!white,draw=blue!14!white,line width=.12pt] (0.219943,1.796069)--(0.235241,1.887042)--(0.000000,1.892820)--(0.000000,1.801472)--cycle;
\path[fill=blue!14!white,draw=blue!14!white,line width=.12pt] (0.000000,1.801472)--(0.000000,1.892820)--(-0.235241,1.887042)--(-0.219943,1.796069)--cycle;
\path[fill=blue!33!white,draw=blue!33!white,line width=.12pt] (2.647922,1.019805)--(2.660703,0.888912)--(2.473633,1.015029)--(2.461751,1.145316)--cycle;
\path[fill=blue!26!white,draw=blue!26!white,line width=.12pt] (-2.461751,1.145316)--(-2.473633,1.015029)--(-2.660703,0.888912)--(-2.647922,1.019805)--cycle;
\path[fill=blue!37!white,draw=blue!37!white,line width=.12pt] (-2.465877,0.333925)--(-2.365079,0.214089)--(-2.508586,0.094375)--(-2.615501,0.209109)--cycle;
\path[fill=blue!41!white,draw=blue!41!white,line width=.12pt] (2.615501,0.209109)--(2.508586,0.094375)--(2.365079,0.214089)--(2.465877,0.333925)--cycle;
\path[fill=blue!35!white,draw=blue!35!white,line width=.12pt] (-2.548600,0.466550)--(-2.465877,0.333925)--(-2.615501,0.209109)--(-2.703244,0.337547)--cycle;
\path[fill=blue!41!white,draw=blue!41!white,line width=.12pt] (2.703244,0.337547)--(2.615501,0.209109)--(2.465877,0.333925)--(2.548600,0.466550)--cycle;
\path[fill=blue!39!white,draw=blue!39!white,line width=.12pt] (1.995527,-0.026136)--(1.865758,-0.056178)--(1.734580,0.032258)--(1.855225,0.068452)--cycle;
\path[fill=blue!40!white,draw=blue!40!white,line width=.12pt] (-1.855225,0.068452)--(-1.734580,0.032258)--(-1.865758,-0.056178)--(-1.995527,-0.026136)--cycle;
\path[fill=blue!14!white,draw=blue!14!white,line width=.12pt] (0.437769,1.779914)--(0.468217,1.869763)--(0.235241,1.887042)--(0.219943,1.796069)--cycle;
\path[fill=blue!14!white,draw=blue!14!white,line width=.12pt] (-0.219943,1.796069)--(-0.235241,1.887042)--(-0.468217,1.869763)--(-0.437769,1.779914)--cycle;
\path[fill=blue!38!white,draw=blue!38!white,line width=.12pt] (-2.365079,0.214089)--(-2.250079,0.111645)--(-2.386608,-0.002248)--(-2.508586,0.094375)--cycle;
\path[fill=blue!42!white,draw=blue!42!white,line width=.12pt] (2.508586,0.094375)--(2.386608,-0.002248)--(2.250079,0.111645)--(2.365079,0.214089)--cycle;
\path[fill=blue!20!white,draw=blue!20!white,line width=.12pt] (1.716761,1.686022)--(1.804504,1.675456)--(1.580295,1.758598)--(1.503455,1.765122)--cycle;
\path[fill=blue!17!white,draw=blue!17!white,line width=.12pt] (-1.503455,1.765122)--(-1.580295,1.758598)--(-1.804504,1.675456)--(-1.716761,1.686022)--cycle;
\path[fill=blue!29!white,draw=blue!29!white,line width=.12pt] (1.019091,0.272690)--(0.963817,0.336316)--(0.817791,0.380079)--(0.864690,0.318963)--cycle;
\path[fill=blue!32!white,draw=blue!32!white,line width=.12pt] (-0.864690,0.318963)--(-0.817791,0.380079)--(-0.963817,0.336316)--(-1.019091,0.272690)--cycle;
\path[fill=blue!16!white,draw=blue!16!white,line width=.12pt] (-1.204924,1.806635)--(-1.275669,1.833387)--(-1.503455,1.765122)--(-1.420078,1.742156)--cycle;
\path[fill=blue!17!white,draw=blue!17!white,line width=.12pt] (1.420078,1.742156)--(1.503455,1.765122)--(1.275669,1.833387)--(1.204924,1.806635)--cycle;
\path[fill=blue!29!white,draw=blue!29!white,line width=.12pt] (-0.503593,0.445151)--(-0.482133,0.529559)--(-0.635598,0.502103)--(-0.663889,0.416474)--cycle;
\path[fill=blue!26!white,draw=blue!26!white,line width=.12pt] (0.663889,0.416474)--(0.635598,0.502103)--(0.482133,0.529559)--(0.503593,0.445151)--cycle;
\path[fill=blue!33!white,draw=blue!33!white,line width=.12pt] (-2.610069,0.606865)--(-2.548600,0.466550)--(-2.703244,0.337547)--(-2.768442,0.474751)--cycle;
\path[fill=blue!39!white,draw=blue!39!white,line width=.12pt] (2.768442,0.474751)--(2.703244,0.337547)--(2.548600,0.466550)--(2.610069,0.606865)--cycle;
\path[fill=blue!19!white,draw=blue!19!white,line width=.12pt] (-1.804504,1.675456)--(-1.881411,1.636151)--(-2.097056,1.538426)--(-2.011334,1.581726)--cycle;
\path[fill=blue!22!white,draw=blue!22!white,line width=.12pt] (2.011334,1.581726)--(2.097056,1.538426)--(1.881411,1.636151)--(1.804504,1.675456)--cycle;
\path[fill=blue!14!white,draw=blue!14!white,line width=.12pt] (0.651378,1.753160)--(0.696683,1.841149)--(0.468217,1.869763)--(0.437769,1.779914)--cycle;
\path[fill=blue!14!white,draw=blue!14!white,line width=.12pt] (-0.437769,1.779914)--(-0.468217,1.869763)--(-0.696683,1.841149)--(-0.651378,1.753160)--cycle;
\path[fill=blue!39!white,draw=blue!39!white,line width=.12pt] (-2.250079,0.111645)--(-2.125296,0.030531)--(-2.254254,-0.077046)--(-2.386608,-0.002248)--cycle;
\path[fill=blue!41!white,draw=blue!41!white,line width=.12pt] (2.386608,-0.002248)--(2.254254,-0.077046)--(2.125296,0.030531)--(2.250079,0.111645)--cycle;
\path[fill=blue!15!white,draw=blue!15!white,line width=.12pt] (1.131352,1.751126)--(1.204924,1.806635)--(0.978167,1.860259)--(0.918440,1.801476)--cycle;
\path[fill=blue!15!white,draw=blue!15!white,line width=.12pt] (-0.918440,1.801476)--(-0.978167,1.860259)--(-1.204924,1.806635)--(-1.131352,1.751126)--cycle;
\path[fill=blue!24!white,draw=blue!24!white,line width=.12pt] (0.162796,0.561318)--(0.158334,0.668634)--(0.000000,0.672523)--(0.000000,0.565317)--cycle;
\path[fill=blue!24!white,draw=blue!24!white,line width=.12pt] (0.000000,0.565317)--(0.000000,0.672523)--(-0.158334,0.668634)--(-0.162796,0.561318)--cycle;
\path[fill=blue!34!white,draw=blue!34!white,line width=.12pt] (1.618570,0.024081)--(1.511656,0.044233)--(1.382778,0.115330)--(1.480578,0.100206)--cycle;
\path[fill=blue!38!white,draw=blue!38!white,line width=.12pt] (-1.480578,0.100206)--(-1.382778,0.115330)--(-1.511656,0.044233)--(-1.618570,0.024081)--cycle;
\path[fill=blue!35!white,draw=blue!35!white,line width=.12pt] (-1.240584,0.179769)--(-1.163677,0.219074)--(-1.297056,0.158630)--(-1.382778,0.115330)--cycle;
\path[fill=blue!32!white,draw=blue!32!white,line width=.12pt] (1.382778,0.115330)--(1.297056,0.158630)--(1.163677,0.219074)--(1.240584,0.179769)--cycle;
\path[fill=blue!30!white,draw=blue!30!white,line width=.12pt] (-2.647922,0.749480)--(-2.610069,0.606865)--(-2.768442,0.474751)--(-2.808591,0.615449)--cycle;
\path[fill=blue!38!white,draw=blue!38!white,line width=.12pt] (2.808591,0.615449)--(2.768442,0.474751)--(2.610069,0.606865)--(2.647922,0.749480)--cycle;
\path[fill=blue!24!white,draw=blue!24!white,line width=.12pt] (0.324025,0.549360)--(0.315143,0.657004)--(0.158334,0.668634)--(0.162796,0.561318)--cycle;
\path[fill=blue!25!white,draw=blue!25!white,line width=.12pt] (-0.162796,0.561318)--(-0.158334,0.668634)--(-0.315143,0.657004)--(-0.324025,0.549360)--cycle;
\path[fill=blue!26!white,draw=blue!26!white,line width=.12pt] (0.817791,0.380079)--(0.782941,0.467259)--(0.635598,0.502103)--(0.663889,0.416474)--cycle;
\path[fill=blue!29!white,draw=blue!29!white,line width=.12pt] (-0.663889,0.416474)--(-0.635598,0.502103)--(-0.782941,0.467259)--(-0.817791,0.380079)--cycle;
\path[fill=blue!31!white,draw=blue!31!white,line width=.12pt] (2.610069,1.137127)--(2.647922,1.019805)--(2.461751,1.145316)--(2.426560,1.260844)--cycle;
\path[fill=blue!24!white,draw=blue!24!white,line width=.12pt] (-2.426560,1.260844)--(-2.461751,1.145316)--(-2.647922,1.019805)--(-2.610069,1.137127)--cycle;
\path[fill=blue!25!white,draw=blue!25!white,line width=.12pt] (2.292506,1.430603)--(2.369413,1.357174)--(2.167407,1.468614)--(2.097056,1.538426)--cycle;
\path[fill=blue!20!white,draw=blue!20!white,line width=.12pt] (-2.097056,1.538426)--(-2.167407,1.468614)--(-2.369413,1.357174)--(-2.292506,1.430603)--cycle;
\path[fill=blue!37!white,draw=blue!37!white,line width=.12pt] (1.865758,-0.056178)--(1.740975,-0.058441)--(1.618570,0.024081)--(1.734580,0.032258)--cycle;
\path[fill=blue!39!white,draw=blue!39!white,line width=.12pt] (-1.734580,0.032258)--(-1.618570,0.024081)--(-1.740975,-0.058441)--(-1.865758,-0.056178)--cycle;
\path[fill=blue!40!white,draw=blue!40!white,line width=.12pt] (2.254254,-0.077046)--(2.116611,-0.127144)--(1.995527,-0.026136)--(2.125296,0.030531)--cycle;
\path[fill=blue!40!white,draw=blue!40!white,line width=.12pt] (-2.125296,0.030531)--(-1.995527,-0.026136)--(-2.116611,-0.127144)--(-2.254254,-0.077046)--cycle;
\path[fill=blue!14!white,draw=blue!14!white,line width=.12pt] (0.858714,1.716067)--(0.918440,1.801476)--(0.696683,1.841149)--(0.651378,1.753160)--cycle;
\path[fill=blue!14!white,draw=blue!14!white,line width=.12pt] (-0.651378,1.753160)--(-0.696683,1.841149)--(-0.918440,1.801476)--(-0.858714,1.716067)--cycle;
\path[fill=blue!13!white,draw=blue!13!white,line width=.12pt] (0.000000,1.687009)--(0.000000,1.801472)--(-0.219943,1.796069)--(-0.205234,1.681968)--cycle;
\path[fill=blue!13!white,draw=blue!13!white,line width=.12pt] (0.205234,1.681968)--(0.219943,1.796069)--(0.000000,1.801472)--(0.000000,1.687009)--cycle;
\path[fill=blue!33!white,draw=blue!33!white,line width=.12pt] (-1.019091,0.272690)--(-0.963817,0.336316)--(-1.100561,0.285608)--(-1.163677,0.219074)--cycle;
\path[fill=blue!29!white,draw=blue!29!white,line width=.12pt] (1.163677,0.219074)--(1.100561,0.285608)--(0.963817,0.336316)--(1.019091,0.272690)--cycle;
\path[fill=blue!25!white,draw=blue!25!white,line width=.12pt] (-0.324025,0.549360)--(-0.315143,0.657004)--(-0.468918,0.637745)--(-0.482133,0.529559)--cycle;
\path[fill=blue!23!white,draw=blue!23!white,line width=.12pt] (0.482133,0.529559)--(0.468918,0.637745)--(0.315143,0.657004)--(0.324025,0.549360)--cycle;
\path[fill=blue!13!white,draw=blue!13!white,line width=.12pt] (0.408491,1.666893)--(0.437769,1.779914)--(0.219943,1.796069)--(0.205234,1.681968)--cycle;
\path[fill=blue!13!white,draw=blue!13!white,line width=.12pt] (-0.205234,1.681968)--(-0.219943,1.796069)--(-0.437769,1.779914)--(-0.408491,1.666893)--cycle;
\path[fill=blue!28!white,draw=blue!28!white,line width=.12pt] (-2.660703,0.888912)--(-2.647922,0.749480)--(-2.808591,0.615449)--(-2.822148,0.754235)--cycle;
\path[fill=blue!36!white,draw=blue!36!white,line width=.12pt] (2.822148,0.754235)--(2.808591,0.615449)--(2.647922,0.749480)--(2.660703,0.888912)--cycle;
\path[fill=blue!23!white,draw=blue!23!white,line width=.12pt] (0.635598,0.502103)--(0.618176,0.611042)--(0.468918,0.637745)--(0.482133,0.529559)--cycle;
\path[fill=blue!26!white,draw=blue!26!white,line width=.12pt] (-0.482133,0.529559)--(-0.468918,0.637745)--(-0.618176,0.611042)--(-0.635598,0.502103)--cycle;
\path[fill=blue!30!white,draw=blue!30!white,line width=.12pt] (-0.817791,0.380079)--(-0.782941,0.467259)--(-0.922745,0.425362)--(-0.963817,0.336316)--cycle;
\path[fill=blue!26!white,draw=blue!26!white,line width=.12pt] (0.963817,0.336316)--(0.922745,0.425362)--(0.782941,0.467259)--(0.817791,0.380079)--cycle;
\path[fill=blue!16!white,draw=blue!16!white,line width=.12pt] (-1.420078,1.742156)--(-1.503455,1.765122)--(-1.716761,1.686022)--(-1.621555,1.667443)--cycle;
\path[fill=blue!17!white,draw=blue!17!white,line width=.12pt] (1.621555,1.667443)--(1.716761,1.686022)--(1.503455,1.765122)--(1.420078,1.742156)--cycle;
\path[fill=blue!15!white,draw=blue!15!white,line width=.12pt] (-1.131352,1.751126)--(-1.204924,1.806635)--(-1.420078,1.742156)--(-1.333369,1.690584)--cycle;
\path[fill=blue!15!white,draw=blue!15!white,line width=.12pt] (1.333369,1.690584)--(1.420078,1.742156)--(1.204924,1.806635)--(1.131352,1.751126)--cycle;
\path[fill=blue!13!white,draw=blue!13!white,line width=.12pt] (0.607814,1.641929)--(0.651378,1.753160)--(0.437769,1.779914)--(0.408491,1.666893)--cycle;
\path[fill=blue!13!white,draw=blue!13!white,line width=.12pt] (-0.408491,1.666893)--(-0.437769,1.779914)--(-0.651378,1.753160)--(-0.607814,1.641929)--cycle;
\path[fill=blue!40!white,draw=blue!40!white,line width=.12pt] (-1.995527,-0.026136)--(-1.865758,-0.056178)--(-1.978968,-0.150618)--(-2.116611,-0.127144)--cycle;
\path[fill=blue!39!white,draw=blue!39!white,line width=.12pt] (2.116611,-0.127144)--(1.978968,-0.150618)--(1.865758,-0.056178)--(1.995527,-0.026136)--cycle;
\path[fill=blue!22!white,draw=blue!22!white,line width=.12pt] (0.000000,0.672523)--(0.000000,0.800000)--(-0.156827,0.796148)--(-0.158334,0.668634)--cycle;
\path[fill=blue!21!white,draw=blue!21!white,line width=.12pt] (0.158334,0.668634)--(0.156827,0.796148)--(0.000000,0.800000)--(0.000000,0.672523)--cycle;
\path[fill=blue!17!white,draw=blue!17!white,line width=.12pt] (-1.716761,1.686022)--(-1.804504,1.675456)--(-2.011334,1.581726)--(-1.913535,1.596850)--cycle;
\path[fill=blue!20!white,draw=blue!20!white,line width=.12pt] (1.913535,1.596850)--(2.011334,1.581726)--(1.804504,1.675456)--(1.716761,1.686022)--cycle;
\path[fill=blue!14!white,draw=blue!14!white,line width=.12pt] (1.057780,1.668992)--(1.131352,1.751126)--(0.918440,1.801476)--(0.858714,1.716067)--cycle;
\path[fill=blue!14!white,draw=blue!14!white,line width=.12pt] (-0.858714,1.716067)--(-0.918440,1.801476)--(-1.131352,1.751126)--(-1.057780,1.668992)--cycle;
\path[fill=blue!21!white,draw=blue!21!white,line width=.12pt] (0.315143,0.657004)--(0.312145,0.784628)--(0.156827,0.796148)--(0.158334,0.668634)--cycle;
\path[fill=blue!22!white,draw=blue!22!white,line width=.12pt] (-0.158334,0.668634)--(-0.156827,0.796148)--(-0.312145,0.784628)--(-0.315143,0.657004)--cycle;
\path[fill=blue!28!white,draw=blue!28!white,line width=.12pt] (2.548600,1.236371)--(2.610069,1.137127)--(2.426560,1.260844)--(2.369413,1.357174)--cycle;
\path[fill=blue!21!white,draw=blue!21!white,line width=.12pt] (-2.369413,1.357174)--(-2.426560,1.260844)--(-2.610069,1.137127)--(-2.548600,1.236371)--cycle;
\path[fill=blue!32!white,draw=blue!32!white,line width=.12pt] (1.511656,0.044233)--(1.417944,0.091940)--(1.297056,0.158630)--(1.382778,0.115330)--cycle;
\path[fill=blue!36!white,draw=blue!36!white,line width=.12pt] (-1.382778,0.115330)--(-1.297056,0.158630)--(-1.417944,0.091940)--(-1.511656,0.044233)--cycle;
\path[fill=blue!13!white,draw=blue!13!white,line width=.12pt] (0.191677,1.549123)--(0.205234,1.681968)--(0.000000,1.687009)--(0.000000,1.553831)--cycle;
\path[fill=blue!13!white,draw=blue!13!white,line width=.12pt] (0.000000,1.553831)--(0.000000,1.687009)--(-0.205234,1.681968)--(-0.191677,1.549123)--cycle;
\path[fill=blue!23!white,draw=blue!23!white,line width=.12pt] (0.782941,0.467259)--(0.761481,0.577153)--(0.618176,0.611042)--(0.635598,0.502103)--cycle;
\path[fill=blue!27!white,draw=blue!27!white,line width=.12pt] (-0.635598,0.502103)--(-0.618176,0.611042)--(-0.761481,0.577153)--(-0.782941,0.467259)--cycle;
\path[fill=blue!29!white,draw=blue!29!white,line width=.12pt] (1.297056,0.158630)--(1.226706,0.228443)--(1.100561,0.285608)--(1.163677,0.219074)--cycle;
\path[fill=blue!33!white,draw=blue!33!white,line width=.12pt] (-1.163677,0.219074)--(-1.100561,0.285608)--(-1.226706,0.228443)--(-1.297056,0.158630)--cycle;
\path[fill=blue!35!white,draw=blue!35!white,line width=.12pt] (1.740975,-0.058441)--(1.625975,-0.032838)--(1.511656,0.044233)--(1.618570,0.024081)--cycle;
\path[fill=blue!38!white,draw=blue!38!white,line width=.12pt] (-1.618570,0.024081)--(-1.511656,0.044233)--(-1.625975,-0.032838)--(-1.740975,-0.058441)--cycle;
\path[fill=blue!23!white,draw=blue!23!white,line width=.12pt] (-0.315143,0.657004)--(-0.312145,0.784628)--(-0.464455,0.765552)--(-0.468918,0.637745)--cycle;
\path[fill=blue!21!white,draw=blue!21!white,line width=.12pt] (0.468918,0.637745)--(0.464455,0.765552)--(0.312145,0.784628)--(0.315143,0.657004)--cycle;
\path[fill=blue!13!white,draw=blue!13!white,line width=.12pt] (0.801283,1.607317)--(0.858714,1.716067)--(0.651378,1.753160)--(0.607814,1.641929)--cycle;
\path[fill=blue!14!white,draw=blue!14!white,line width=.12pt] (-0.607814,1.641929)--(-0.651378,1.753160)--(-0.858714,1.716067)--(-0.801283,1.607317)--cycle;
\path[fill=blue!36!white,draw=blue!36!white,line width=.12pt] (-2.615501,0.209109)--(-2.508586,0.094375)--(-2.627935,-0.031796)--(-2.739936,0.077561)--cycle;
\path[fill=blue!41!white,draw=blue!41!white,line width=.12pt] (2.739936,0.077561)--(2.627935,-0.031796)--(2.508586,0.094375)--(2.615501,0.209109)--cycle;
\path[fill=blue!12!white,draw=blue!12!white,line width=.12pt] (0.381508,1.535043)--(0.408491,1.666893)--(0.205234,1.681968)--(0.191677,1.549123)--cycle;
\path[fill=blue!13!white,draw=blue!13!white,line width=.12pt] (-0.191677,1.549123)--(-0.205234,1.681968)--(-0.408491,1.666893)--(-0.381508,1.535043)--cycle;
\path[fill=blue!33!white,draw=blue!33!white,line width=.12pt] (2.808591,0.885774)--(2.822148,0.754235)--(2.660703,0.888912)--(2.647922,1.019805)--cycle;
\path[fill=blue!25!white,draw=blue!25!white,line width=.12pt] (-2.647922,1.019805)--(-2.660703,0.888912)--(-2.822148,0.754235)--(-2.808591,0.885774)--cycle;
\path[fill=blue!18!white,draw=blue!18!white,line width=.12pt] (-2.011334,1.581726)--(-2.097056,1.538426)--(-2.292506,1.430603)--(-2.198794,1.478311)--cycle;
\path[fill=blue!22!white,draw=blue!22!white,line width=.12pt] (2.198794,1.478311)--(2.292506,1.430603)--(2.097056,1.538426)--(2.011334,1.581726)--cycle;
\path[fill=blue!41!white,draw=blue!41!white,line width=.12pt] (2.627935,-0.031796)--(2.500154,-0.122284)--(2.386608,-0.002248)--(2.508586,0.094375)--cycle;
\path[fill=blue!38!white,draw=blue!38!white,line width=.12pt] (-2.508586,0.094375)--(-2.386608,-0.002248)--(-2.500154,-0.122284)--(-2.627935,-0.031796)--cycle;
\path[fill=blue!34!white,draw=blue!34!white,line width=.12pt] (-2.703244,0.337547)--(-2.615501,0.209109)--(-2.739936,0.077561)--(-2.831853,0.201586)--cycle;
\path[fill=blue!40!white,draw=blue!40!white,line width=.12pt] (2.831853,0.201586)--(2.739936,0.077561)--(2.615501,0.209109)--(2.703244,0.337547)--cycle;
\path[fill=blue!31!white,draw=blue!31!white,line width=.12pt] (-0.963817,0.336316)--(-0.922745,0.425362)--(-1.053662,0.376814)--(-1.100561,0.285608)--cycle;
\path[fill=blue!26!white,draw=blue!26!white,line width=.12pt] (1.100561,0.285608)--(1.053662,0.376814)--(0.922745,0.425362)--(0.963817,0.336316)--cycle;
\path[fill=blue!39!white,draw=blue!39!white,line width=.12pt] (-2.386608,-0.002248)--(-2.254254,-0.077046)--(-2.361503,-0.190425)--(-2.500154,-0.122284)--cycle;
\path[fill=blue!41!white,draw=blue!41!white,line width=.12pt] (2.500154,-0.122284)--(2.361503,-0.190425)--(2.254254,-0.077046)--(2.386608,-0.002248)--cycle;
\path[fill=blue!19!white,draw=blue!19!white,line width=.12pt] (0.156827,0.796148)--(0.158334,0.938959)--(0.000000,0.942848)--(0.000000,0.800000)--cycle;
\path[fill=blue!19!white,draw=blue!19!white,line width=.12pt] (0.000000,0.800000)--(0.000000,0.942848)--(-0.158334,0.938959)--(-0.156827,0.796148)--cycle;
\path[fill=blue!20!white,draw=blue!20!white,line width=.12pt] (0.618176,0.611042)--(0.612293,0.739104)--(0.464455,0.765552)--(0.468918,0.637745)--cycle;
\path[fill=blue!23!white,draw=blue!23!white,line width=.12pt] (-0.468918,0.637745)--(-0.464455,0.765552)--(-0.612293,0.739104)--(-0.618176,0.611042)--cycle;
\path[fill=blue!13!white,draw=blue!13!white,line width=.12pt] (-0.381508,1.535043)--(-0.408491,1.666893)--(-0.607814,1.641929)--(-0.567664,1.511728)--cycle;
\path[fill=blue!12!white,draw=blue!12!white,line width=.12pt] (0.567664,1.511728)--(0.607814,1.641929)--(0.408491,1.666893)--(0.381508,1.535043)--cycle;
\path[fill=blue!37!white,draw=blue!37!white,line width=.12pt] (1.978968,-0.150618)--(1.846614,-0.146565)--(1.740975,-0.058441)--(1.865758,-0.056178)--cycle;
\path[fill=blue!39!white,draw=blue!39!white,line width=.12pt] (-1.865758,-0.056178)--(-1.740975,-0.058441)--(-1.846614,-0.146565)--(-1.978968,-0.150618)--cycle;
\path[fill=blue!18!white,draw=blue!18!white,line width=.12pt] (0.312145,0.784628)--(0.315143,0.927329)--(0.158334,0.938959)--(0.156827,0.796148)--cycle;
\path[fill=blue!20!white,draw=blue!20!white,line width=.12pt] (-0.156827,0.796148)--(-0.158334,0.938959)--(-0.315143,0.927329)--(-0.312145,0.784628)--cycle;
\path[fill=blue!32!white,draw=blue!32!white,line width=.12pt] (-2.768442,0.474751)--(-2.703244,0.337547)--(-2.831853,0.201586)--(-2.900154,0.335511)--cycle;
\path[fill=blue!39!white,draw=blue!39!white,line width=.12pt] (2.900154,0.335511)--(2.831853,0.201586)--(2.703244,0.337547)--(2.768442,0.474751)--cycle;
\path[fill=blue!27!white,draw=blue!27!white,line width=.12pt] (-0.782941,0.467259)--(-0.761481,0.577153)--(-0.897452,0.536404)--(-0.922745,0.425362)--cycle;
\path[fill=blue!23!white,draw=blue!23!white,line width=.12pt] (0.922745,0.425362)--(0.897452,0.536404)--(0.761481,0.577153)--(0.782941,0.467259)--cycle;
\path[fill=blue!13!white,draw=blue!13!white,line width=.12pt] (0.179794,1.402639)--(0.191677,1.549123)--(0.000000,1.553831)--(0.000000,1.407055)--cycle;
\path[fill=blue!13!white,draw=blue!13!white,line width=.12pt] (0.000000,1.407055)--(0.000000,1.553831)--(-0.191677,1.549123)--(-0.179794,1.402639)--cycle;
\path[fill=blue!14!white,draw=blue!14!white,line width=.12pt] (-1.057780,1.668992)--(-1.131352,1.751126)--(-1.333369,1.690584)--(-1.246659,1.612387)--cycle;
\path[fill=blue!14!white,draw=blue!14!white,line width=.12pt] (1.246659,1.612387)--(1.333369,1.690584)--(1.131352,1.751126)--(1.057780,1.668992)--cycle;
\path[fill=blue!15!white,draw=blue!15!white,line width=.12pt] (1.522544,1.620433)--(1.621555,1.667443)--(1.420078,1.742156)--(1.333369,1.690584)--cycle;
\path[fill=blue!15!white,draw=blue!15!white,line width=.12pt] (-1.333369,1.690584)--(-1.420078,1.742156)--(-1.621555,1.667443)--(-1.522544,1.620433)--cycle;
\path[fill=blue!14!white,draw=blue!14!white,line width=.12pt] (-0.801283,1.607317)--(-0.858714,1.716067)--(-1.057780,1.668992)--(-0.987036,1.563389)--cycle;
\path[fill=blue!13!white,draw=blue!13!white,line width=.12pt] (0.987036,1.563389)--(1.057780,1.668992)--(0.858714,1.716067)--(0.801283,1.607317)--cycle;
\path[fill=blue!13!white,draw=blue!13!white,line width=.12pt] (0.357857,1.389432)--(0.381508,1.535043)--(0.191677,1.549123)--(0.179794,1.402639)--cycle;
\path[fill=blue!13!white,draw=blue!13!white,line width=.12pt] (-0.179794,1.402639)--(-0.191677,1.549123)--(-0.381508,1.535043)--(-0.357857,1.389432)--cycle;
\path[fill=blue!40!white,draw=blue!40!white,line width=.12pt] (2.361503,-0.190425)--(2.217311,-0.233600)--(2.116611,-0.127144)--(2.254254,-0.077046)--cycle;
\path[fill=blue!39!white,draw=blue!39!white,line width=.12pt] (-2.254254,-0.077046)--(-2.116611,-0.127144)--(-2.217311,-0.233600)--(-2.361503,-0.190425)--cycle;
\path[fill=blue!20!white,draw=blue!20!white,line width=.12pt] (-0.312145,0.784628)--(-0.315143,0.927329)--(-0.468918,0.908070)--(-0.464455,0.765552)--cycle;
\path[fill=blue!18!white,draw=blue!18!white,line width=.12pt] (0.464455,0.765552)--(0.468918,0.908070)--(0.315143,0.927329)--(0.312145,0.784628)--cycle;
\path[fill=blue!16!white,draw=blue!16!white,line width=.12pt] (-1.621555,1.667443)--(-1.716761,1.686022)--(-1.913535,1.596850)--(-1.807416,1.583216)--cycle;
\path[fill=blue!17!white,draw=blue!17!white,line width=.12pt] (1.807416,1.583216)--(1.913535,1.596850)--(1.716761,1.686022)--(1.621555,1.667443)--cycle;
\path[fill=blue!25!white,draw=blue!25!white,line width=.12pt] (2.465877,1.313721)--(2.548600,1.236371)--(2.369413,1.357174)--(2.292506,1.430603)--cycle;
\path[fill=blue!19!white,draw=blue!19!white,line width=.12pt] (-2.292506,1.430603)--(-2.369413,1.357174)--(-2.548600,1.236371)--(-2.465877,1.313721)--cycle;
\path[fill=blue!17!white,draw=blue!17!white,line width=.12pt] (0.158334,0.938959)--(0.162796,1.091580)--(0.000000,1.095579)--(0.000000,0.942848)--cycle;
\path[fill=blue!17!white,draw=blue!17!white,line width=.12pt] (0.000000,0.942848)--(0.000000,1.095579)--(-0.162796,1.091580)--(-0.158334,0.938959)--cycle;
\path[fill=blue!12!white,draw=blue!12!white,line width=.12pt] (0.748354,1.479402)--(0.801283,1.607317)--(0.607814,1.641929)--(0.567664,1.511728)--cycle;
\path[fill=blue!14!white,draw=blue!14!white,line width=.12pt] (-0.567664,1.511728)--(-0.607814,1.641929)--(-0.801283,1.607317)--(-0.748354,1.479402)--cycle;
\path[fill=blue!24!white,draw=blue!24!white,line width=.12pt] (-0.618176,0.611042)--(-0.612293,0.739104)--(-0.754235,0.705537)--(-0.761481,0.577153)--cycle;
\path[fill=blue!20!white,draw=blue!20!white,line width=.12pt] (0.761481,0.577153)--(0.754235,0.705537)--(0.612293,0.739104)--(0.618176,0.611042)--cycle;
\path[fill=blue!30!white,draw=blue!30!white,line width=.12pt] (2.768442,1.005012)--(2.808591,0.885774)--(2.647922,1.019805)--(2.610069,1.137127)--cycle;
\path[fill=blue!23!white,draw=blue!23!white,line width=.12pt] (-2.610069,1.137127)--(-2.647922,1.019805)--(-2.808591,0.885774)--(-2.768442,1.005012)--cycle;
\path[fill=blue!14!white,draw=blue!14!white,line width=.12pt] (0.000000,1.252323)--(0.000000,1.407055)--(-0.179794,1.402639)--(-0.170043,1.248146)--cycle;
\path[fill=blue!14!white,draw=blue!14!white,line width=.12pt] (0.170043,1.248146)--(0.179794,1.402639)--(0.000000,1.407055)--(0.000000,1.252323)--cycle;
\path[fill=blue!29!white,draw=blue!29!white,line width=.12pt] (1.417944,0.091940)--(1.341037,0.165370)--(1.226706,0.228443)--(1.297056,0.158630)--cycle;
\path[fill=blue!34!white,draw=blue!34!white,line width=.12pt] (-1.297056,0.158630)--(-1.226706,0.228443)--(-1.341037,0.165370)--(-1.417944,0.091940)--cycle;
\path[fill=blue!16!white,draw=blue!16!white,line width=.12pt] (0.315143,0.927329)--(0.324025,1.079622)--(0.162796,1.091580)--(0.158334,0.938959)--cycle;
\path[fill=blue!18!white,draw=blue!18!white,line width=.12pt] (-0.158334,0.938959)--(-0.162796,1.091580)--(-0.324025,1.079622)--(-0.315143,0.927329)--cycle;
\path[fill=blue!30!white,draw=blue!30!white,line width=.12pt] (-2.808591,0.615449)--(-2.768442,0.474751)--(-2.900154,0.335511)--(-2.942213,0.474190)--cycle;
\path[fill=blue!37!white,draw=blue!37!white,line width=.12pt] (2.942213,0.474190)--(2.900154,0.335511)--(2.768442,0.474751)--(2.808591,0.615449)--cycle;
\path[fill=blue!32!white,draw=blue!32!white,line width=.12pt] (1.625975,-0.032838)--(1.525177,0.019647)--(1.417944,0.091940)--(1.511656,0.044233)--cycle;
\path[fill=blue!36!white,draw=blue!36!white,line width=.12pt] (-1.511656,0.044233)--(-1.417944,0.091940)--(-1.525177,0.019647)--(-1.625975,-0.032838)--cycle;
\path[fill=blue!14!white,draw=blue!14!white,line width=.12pt] (-0.357857,1.389432)--(-0.381508,1.535043)--(-0.567664,1.511728)--(-0.532473,1.367563)--cycle;
\path[fill=blue!12!white,draw=blue!12!white,line width=.12pt] (0.532473,1.367563)--(0.567664,1.511728)--(0.381508,1.535043)--(0.357857,1.389432)--cycle;
\path[fill=blue!15!white,draw=blue!15!white,line width=.12pt] (0.162796,1.091580)--(0.170043,1.248146)--(0.000000,1.252323)--(0.000000,1.095579)--cycle;
\path[fill=blue!15!white,draw=blue!15!white,line width=.12pt] (0.000000,1.095579)--(0.000000,1.252323)--(-0.170043,1.248146)--(-0.162796,1.091580)--cycle;
\path[fill=blue!26!white,draw=blue!26!white,line width=.12pt] (1.226706,0.228443)--(1.174431,0.322085)--(1.053662,0.376814)--(1.100561,0.285608)--cycle;
\path[fill=blue!31!white,draw=blue!31!white,line width=.12pt] (-1.100561,0.285608)--(-1.053662,0.376814)--(-1.174431,0.322085)--(-1.226706,0.228443)--cycle;
\path[fill=blue!13!white,draw=blue!13!white,line width=.12pt] (0.338447,1.235655)--(0.357857,1.389432)--(0.179794,1.402639)--(0.170043,1.248146)--cycle;
\path[fill=blue!14!white,draw=blue!14!white,line width=.12pt] (-0.170043,1.248146)--(-0.179794,1.402639)--(-0.357857,1.389432)--(-0.338447,1.235655)--cycle;
\path[fill=blue!18!white,draw=blue!18!white,line width=.12pt] (0.612293,0.739104)--(0.618176,0.881367)--(0.468918,0.908070)--(0.464455,0.765552)--cycle;
\path[fill=blue!21!white,draw=blue!21!white,line width=.12pt] (-0.464455,0.765552)--(-0.468918,0.908070)--(-0.618176,0.881367)--(-0.612293,0.739104)--cycle;
\path[fill=blue!16!white,draw=blue!16!white,line width=.12pt] (-0.162796,1.091580)--(-0.170043,1.248146)--(-0.338447,1.235655)--(-0.324025,1.079622)--cycle;
\path[fill=blue!15!white,draw=blue!15!white,line width=.12pt] (0.324025,1.079622)--(0.338447,1.235655)--(0.170043,1.248146)--(0.162796,1.091580)--cycle;
\path[fill=blue!16!white,draw=blue!16!white,line width=.12pt] (-1.913535,1.596850)--(-2.011334,1.581726)--(-2.198794,1.478311)--(-2.091880,1.498463)--cycle;
\path[fill=blue!20!white,draw=blue!20!white,line width=.12pt] (2.091880,1.498463)--(2.198794,1.478311)--(2.011334,1.581726)--(1.913535,1.596850)--cycle;
\path[fill=blue!16!white,draw=blue!16!white,line width=.12pt] (0.468918,0.908070)--(0.482133,1.059820)--(0.324025,1.079622)--(0.315143,0.927329)--cycle;
\path[fill=blue!18!white,draw=blue!18!white,line width=.12pt] (-0.315143,0.927329)--(-0.324025,1.079622)--(-0.482133,1.059820)--(-0.468918,0.908070)--cycle;
\path[fill=blue!40!white,draw=blue!40!white,line width=.12pt] (-2.116611,-0.127144)--(-1.978968,-0.150618)--(-2.073119,-0.250151)--(-2.217311,-0.233600)--cycle;
\path[fill=blue!39!white,draw=blue!39!white,line width=.12pt] (2.217311,-0.233600)--(2.073119,-0.250151)--(1.978968,-0.150618)--(2.116611,-0.127144)--cycle;
\path[fill=blue!23!white,draw=blue!23!white,line width=.12pt] (1.053662,0.376814)--(1.024781,0.489187)--(0.897452,0.536404)--(0.922745,0.425362)--cycle;
\path[fill=blue!28!white,draw=blue!28!white,line width=.12pt] (-0.922745,0.425362)--(-0.897452,0.536404)--(-1.024781,0.489187)--(-1.053662,0.376814)--cycle;
\path[fill=blue!38!white,draw=blue!38!white,line width=.12pt] (-1.740975,-0.058441)--(-1.625975,-0.032838)--(-1.724636,-0.115141)--(-1.846614,-0.146565)--cycle;
\path[fill=blue!35!white,draw=blue!35!white,line width=.12pt] (1.846614,-0.146565)--(1.724636,-0.115141)--(1.625975,-0.032838)--(1.740975,-0.058441)--cycle;
\path[fill=blue!15!white,draw=blue!15!white,line width=.12pt] (-0.338447,1.235655)--(-0.357857,1.389432)--(-0.532473,1.367563)--(-0.503593,1.214972)--cycle;
\path[fill=blue!13!white,draw=blue!13!white,line width=.12pt] (0.503593,1.214972)--(0.532473,1.367563)--(0.357857,1.389432)--(0.338447,1.235655)--cycle;
\path[fill=blue!14!white,draw=blue!14!white,line width=.12pt] (0.482133,1.059820)--(0.503593,1.214972)--(0.338447,1.235655)--(0.324025,1.079622)--cycle;
\path[fill=blue!16!white,draw=blue!16!white,line width=.12pt] (-0.324025,1.079622)--(-0.338447,1.235655)--(-0.503593,1.214972)--(-0.482133,1.059820)--cycle;
\path[fill=blue!12!white,draw=blue!12!white,line width=.12pt] (0.701962,1.337241)--(0.748354,1.479402)--(0.567664,1.511728)--(0.532473,1.367563)--cycle;
\path[fill=blue!14!white,draw=blue!14!white,line width=.12pt] (-0.532473,1.367563)--(-0.567664,1.511728)--(-0.748354,1.479402)--(-0.701962,1.337241)--cycle;
\path[fill=blue!12!white,draw=blue!12!white,line width=.12pt] (0.921837,1.438377)--(0.987036,1.563389)--(0.801283,1.607317)--(0.748354,1.479402)--cycle;
\path[fill=blue!14!white,draw=blue!14!white,line width=.12pt] (-0.748354,1.479402)--(-0.801283,1.607317)--(-0.987036,1.563389)--(-0.921837,1.438377)--cycle;
\path[fill=blue!20!white,draw=blue!20!white,line width=.12pt] (0.897452,0.536404)--(0.888912,0.665176)--(0.754235,0.705537)--(0.761481,0.577153)--cycle;
\path[fill=blue!25!white,draw=blue!25!white,line width=.12pt] (-0.761481,0.577153)--(-0.754235,0.705537)--(-0.888912,0.665176)--(-0.897452,0.536404)--cycle;
\path[fill=blue!13!white,draw=blue!13!white,line width=.12pt] (1.163283,1.510570)--(1.246659,1.612387)--(1.057780,1.668992)--(0.987036,1.563389)--cycle;
\path[fill=blue!14!white,draw=blue!14!white,line width=.12pt] (-0.987036,1.563389)--(-1.057780,1.668992)--(-1.246659,1.612387)--(-1.163283,1.510570)--cycle;
\path[fill=blue!27!white,draw=blue!27!white,line width=.12pt] (-2.822148,0.754235)--(-2.808591,0.615449)--(-2.942213,0.474190)--(-2.956415,0.612293)--cycle;
\path[fill=blue!35!white,draw=blue!35!white,line width=.12pt] (2.956415,0.612293)--(2.942213,0.474190)--(2.808591,0.615449)--(2.822148,0.754235)--cycle;
\path[fill=blue!19!white,draw=blue!19!white,line width=.12pt] (-0.468918,0.908070)--(-0.482133,1.059820)--(-0.635598,1.032365)--(-0.618176,0.881367)--cycle;
\path[fill=blue!16!white,draw=blue!16!white,line width=.12pt] (0.618176,0.881367)--(0.635598,1.032365)--(0.482133,1.059820)--(0.468918,0.908070)--cycle;
\path[fill=blue!22!white,draw=blue!22!white,line width=.12pt] (-0.612293,0.739104)--(-0.618176,0.881367)--(-0.761481,0.847478)--(-0.754235,0.705537)--cycle;
\path[fill=blue!18!white,draw=blue!18!white,line width=.12pt] (0.754235,0.705537)--(0.761481,0.847478)--(0.618176,0.881367)--(0.612293,0.739104)--cycle;
\path[fill=blue!14!white,draw=blue!14!white,line width=.12pt] (1.423533,1.546798)--(1.522544,1.620433)--(1.333369,1.690584)--(1.246659,1.612387)--cycle;
\path[fill=blue!14!white,draw=blue!14!white,line width=.12pt] (-1.246659,1.612387)--(-1.333369,1.690584)--(-1.522544,1.620433)--(-1.423533,1.546798)--cycle;
\path[fill=blue!15!white,draw=blue!15!white,line width=.12pt] (-0.503593,1.214972)--(-0.532473,1.367563)--(-0.701962,1.337241)--(-0.663889,1.186295)--cycle;
\path[fill=blue!13!white,draw=blue!13!white,line width=.12pt] (0.663889,1.186295)--(0.701962,1.337241)--(0.532473,1.367563)--(0.503593,1.214972)--cycle;
\path[fill=blue!14!white,draw=blue!14!white,line width=.12pt] (0.635598,1.032365)--(0.663889,1.186295)--(0.503593,1.214972)--(0.482133,1.059820)--cycle;
\path[fill=blue!17!white,draw=blue!17!white,line width=.12pt] (-0.482133,1.059820)--(-0.503593,1.214972)--(-0.663889,1.186295)--(-0.635598,1.032365)--cycle;
\path[fill=blue!15!white,draw=blue!15!white,line width=.12pt] (-1.522544,1.620433)--(-1.621555,1.667443)--(-1.807416,1.583216)--(-1.697056,1.541348)--cycle;
\path[fill=blue!15!white,draw=blue!15!white,line width=.12pt] (1.697056,1.541348)--(1.807416,1.583216)--(1.621555,1.667443)--(1.522544,1.620433)--cycle;
\path[fill=blue!21!white,draw=blue!21!white,line width=.12pt] (-2.548600,1.236371)--(-2.610069,1.137127)--(-2.768442,1.005012)--(-2.703244,1.107367)--cycle;
\path[fill=blue!28!white,draw=blue!28!white,line width=.12pt] (2.703244,1.107367)--(2.768442,1.005012)--(2.610069,1.137127)--(2.548600,1.236371)--cycle;
\path[fill=blue!18!white,draw=blue!18!white,line width=.12pt] (-2.198794,1.478311)--(-2.292506,1.430603)--(-2.465877,1.313721)--(-2.365079,1.366207)--cycle;
\path[fill=blue!22!white,draw=blue!22!white,line width=.12pt] (2.365079,1.366207)--(2.465877,1.313721)--(2.292506,1.430603)--(2.198794,1.478311)--cycle;
\path[fill=blue!39!white,draw=blue!39!white,line width=.12pt] (-1.978968,-0.150618)--(-1.846614,-0.146565)--(-1.934468,-0.239441)--(-2.073119,-0.250151)--cycle;
\path[fill=blue!37!white,draw=blue!37!white,line width=.12pt] (2.073119,-0.250151)--(1.934468,-0.239441)--(1.846614,-0.146565)--(1.978968,-0.150618)--cycle;
\path[fill=blue!12!white,draw=blue!12!white,line width=.12pt] (0.864690,1.298758)--(0.921837,1.438377)--(0.748354,1.479402)--(0.701962,1.337241)--cycle;
\path[fill=blue!15!white,draw=blue!15!white,line width=.12pt] (-0.701962,1.337241)--(-0.748354,1.479402)--(-0.921837,1.438377)--(-0.864690,1.298758)--cycle;
\path[fill=blue!32!white,draw=blue!32!white,line width=.12pt] (-1.226706,0.228443)--(-1.174431,0.322085)--(-1.283890,0.261700)--(-1.341037,0.165370)--cycle;
\path[fill=blue!26!white,draw=blue!26!white,line width=.12pt] (1.341037,0.165370)--(1.283890,0.261700)--(1.174431,0.322085)--(1.226706,0.228443)--cycle;
\path[fill=blue!29!white,draw=blue!29!white,line width=.12pt] (1.525177,0.019647)--(1.442454,0.096998)--(1.341037,0.165370)--(1.417944,0.091940)--cycle;
\path[fill=blue!35!white,draw=blue!35!white,line width=.12pt] (-1.417944,0.091940)--(-1.341037,0.165370)--(-1.442454,0.096998)--(-1.525177,0.019647)--cycle;
\path[fill=blue!23!white,draw=blue!23!white,line width=.12pt] (1.174431,0.322085)--(1.142240,0.435958)--(1.024781,0.489187)--(1.053662,0.376814)--cycle;
\path[fill=blue!29!white,draw=blue!29!white,line width=.12pt] (-1.053662,0.376814)--(-1.024781,0.489187)--(-1.142240,0.435958)--(-1.174431,0.322085)--cycle;
\path[fill=blue!16!white,draw=blue!16!white,line width=.12pt] (0.761481,0.847478)--(0.782941,0.997521)--(0.635598,1.032365)--(0.618176,0.881367)--cycle;
\path[fill=blue!19!white,draw=blue!19!white,line width=.12pt] (-0.618176,0.881367)--(-0.635598,1.032365)--(-0.782941,0.997521)--(-0.761481,0.847478)--cycle;
\path[fill=blue!37!white,draw=blue!37!white,line width=.12pt] (-2.627935,-0.031796)--(-2.500154,-0.122284)--(-2.589621,-0.247306)--(-2.721975,-0.163208)--cycle;
\path[fill=blue!41!white,draw=blue!41!white,line width=.12pt] (2.721975,-0.163208)--(2.589621,-0.247306)--(2.500154,-0.122284)--(2.627935,-0.031796)--cycle;
\path[fill=blue!20!white,draw=blue!20!white,line width=.12pt] (1.024781,0.489187)--(1.015029,0.618408)--(0.888912,0.665176)--(0.897452,0.536404)--cycle;
\path[fill=blue!26!white,draw=blue!26!white,line width=.12pt] (-0.897452,0.536404)--(-0.888912,0.665176)--(-1.015029,0.618408)--(-1.024781,0.489187)--cycle;
\path[fill=blue!41!white,draw=blue!41!white,line width=.12pt] (2.837984,-0.059451)--(2.721975,-0.163208)--(2.627935,-0.031796)--(2.739936,0.077561)--cycle;
\path[fill=blue!36!white,draw=blue!36!white,line width=.12pt] (-2.739936,0.077561)--(-2.627935,-0.031796)--(-2.721975,-0.163208)--(-2.837984,-0.059451)--cycle;
\path[fill=blue!22!white,draw=blue!22!white,line width=.12pt] (-0.754235,0.705537)--(-0.761481,0.847478)--(-0.897452,0.806729)--(-0.888912,0.665176)--cycle;
\path[fill=blue!18!white,draw=blue!18!white,line width=.12pt] (0.888912,0.665176)--(0.897452,0.806729)--(0.761481,0.847478)--(0.754235,0.705537)--cycle;
\path[fill=blue!12!white,draw=blue!12!white,line width=.12pt] (1.086442,1.389047)--(1.163283,1.510570)--(0.987036,1.563389)--(0.921837,1.438377)--cycle;
\path[fill=blue!14!white,draw=blue!14!white,line width=.12pt] (-0.921837,1.438377)--(-0.987036,1.563389)--(-1.163283,1.510570)--(-1.086442,1.389047)--cycle;
\path[fill=blue!32!white,draw=blue!32!white,line width=.12pt] (1.724636,-0.115141)--(1.617721,-0.057553)--(1.525177,0.019647)--(1.625975,-0.032838)--cycle;
\path[fill=blue!37!white,draw=blue!37!white,line width=.12pt] (-1.625975,-0.032838)--(-1.525177,0.019647)--(-1.617721,-0.057553)--(-1.724636,-0.115141)--cycle;
\path[fill=blue!38!white,draw=blue!38!white,line width=.12pt] (-2.500154,-0.122284)--(-2.361503,-0.190425)--(-2.446009,-0.308514)--(-2.589621,-0.247306)--cycle;
\path[fill=blue!41!white,draw=blue!41!white,line width=.12pt] (2.589621,-0.247306)--(2.446009,-0.308514)--(2.361503,-0.190425)--(2.500154,-0.122284)--cycle;
\path[fill=blue!17!white,draw=blue!17!white,line width=.12pt] (1.975871,1.490286)--(2.091880,1.498463)--(1.913535,1.596850)--(1.807416,1.583216)--cycle;
\path[fill=blue!15!white,draw=blue!15!white,line width=.12pt] (-1.807416,1.583216)--(-1.913535,1.596850)--(-2.091880,1.498463)--(-1.975871,1.490286)--cycle;
\path[fill=blue!13!white,draw=blue!13!white,line width=.12pt] (0.817791,1.149900)--(0.864690,1.298758)--(0.701962,1.337241)--(0.663889,1.186295)--cycle;
\path[fill=blue!16!white,draw=blue!16!white,line width=.12pt] (-0.663889,1.186295)--(-0.701962,1.337241)--(-0.864690,1.298758)--(-0.817791,1.149900)--cycle;
\path[fill=blue!25!white,draw=blue!25!white,line width=.12pt] (-2.808591,0.885774)--(-2.822148,0.754235)--(-2.956415,0.612293)--(-2.942213,0.744515)--cycle;
\path[fill=blue!33!white,draw=blue!33!white,line width=.12pt] (2.942213,0.744515)--(2.956415,0.612293)--(2.822148,0.754235)--(2.808591,0.885774)--cycle;
\path[fill=blue!17!white,draw=blue!17!white,line width=.12pt] (-0.635598,1.032365)--(-0.663889,1.186295)--(-0.817791,1.149900)--(-0.782941,0.997521)--cycle;
\path[fill=blue!14!white,draw=blue!14!white,line width=.12pt] (0.782941,0.997521)--(0.817791,1.149900)--(0.663889,1.186295)--(0.635598,1.032365)--cycle;
\path[fill=blue!33!white,draw=blue!33!white,line width=.12pt] (-2.831853,0.201586)--(-2.739936,0.077561)--(-2.837984,-0.059451)--(-2.933190,0.059976)--cycle;
\path[fill=blue!40!white,draw=blue!40!white,line width=.12pt] (2.933190,0.059976)--(2.837984,-0.059451)--(2.739936,0.077561)--(2.831853,0.201586)--cycle;
\path[fill=blue!14!white,draw=blue!14!white,line width=.12pt] (-1.163283,1.510570)--(-1.246659,1.612387)--(-1.423533,1.546798)--(-1.328326,1.449368)--cycle;
\path[fill=blue!13!white,draw=blue!13!white,line width=.12pt] (1.328326,1.449368)--(1.423533,1.546798)--(1.246659,1.612387)--(1.163283,1.510570)--cycle;
\path[fill=blue!39!white,draw=blue!39!white,line width=.12pt] (-2.361503,-0.190425)--(-2.217311,-0.233600)--(-2.296657,-0.344479)--(-2.446009,-0.308514)--cycle;
\path[fill=blue!40!white,draw=blue!40!white,line width=.12pt] (2.446009,-0.308514)--(2.296657,-0.344479)--(2.217311,-0.233600)--(2.361503,-0.190425)--cycle;
\path[fill=blue!39!white,draw=blue!39!white,line width=.12pt] (-1.846614,-0.146565)--(-1.724636,-0.115141)--(-1.806687,-0.201882)--(-1.934468,-0.239441)--cycle;
\path[fill=blue!35!white,draw=blue!35!white,line width=.12pt] (1.934468,-0.239441)--(1.806687,-0.201882)--(1.724636,-0.115141)--(1.846614,-0.146565)--cycle;
\path[fill=blue!20!white,draw=blue!20!white,line width=.12pt] (-0.761481,0.847478)--(-0.782941,0.997521)--(-0.922745,0.955623)--(-0.897452,0.806729)--cycle;
\path[fill=blue!16!white,draw=blue!16!white,line width=.12pt] (0.897452,0.806729)--(0.922745,0.955623)--(0.782941,0.997521)--(0.761481,0.847478)--cycle;
\path[fill=blue!12!white,draw=blue!12!white,line width=.12pt] (1.019091,1.252486)--(1.086442,1.389047)--(0.921837,1.438377)--(0.864690,1.298758)--cycle;
\path[fill=blue!15!white,draw=blue!15!white,line width=.12pt] (-0.864690,1.298758)--(-0.921837,1.438377)--(-1.086442,1.389047)--(-1.019091,1.252486)--cycle;
\path[fill=blue!14!white,draw=blue!14!white,line width=.12pt] (1.586697,1.472856)--(1.697056,1.541348)--(1.522544,1.620433)--(1.423533,1.546798)--cycle;
\path[fill=blue!14!white,draw=blue!14!white,line width=.12pt] (-1.423533,1.546798)--(-1.522544,1.620433)--(-1.697056,1.541348)--(-1.586697,1.472856)--cycle;
\path[fill=blue!31!white,draw=blue!31!white,line width=.12pt] (-2.900154,0.335511)--(-2.831853,0.201586)--(-2.933190,0.059976)--(-3.003935,0.190486)--cycle;
\path[fill=blue!39!white,draw=blue!39!white,line width=.12pt] (3.003935,0.190486)--(2.933190,0.059976)--(2.831853,0.201586)--(2.900154,0.335511)--cycle;
\path[fill=blue!25!white,draw=blue!25!white,line width=.12pt] (2.615501,1.188905)--(2.703244,1.107367)--(2.548600,1.236371)--(2.465877,1.313721)--cycle;
\path[fill=blue!19!white,draw=blue!19!white,line width=.12pt] (-2.465877,1.313721)--(-2.548600,1.236371)--(-2.703244,1.107367)--(-2.615501,1.188905)--cycle;
\path[fill=blue!19!white,draw=blue!19!white,line width=.12pt] (2.250079,1.391810)--(2.365079,1.366207)--(2.198794,1.478311)--(2.091880,1.498463)--cycle;
\path[fill=blue!16!white,draw=blue!16!white,line width=.12pt] (-2.091880,1.498463)--(-2.198794,1.478311)--(-2.365079,1.366207)--(-2.250079,1.391810)--cycle;
\path[fill=blue!23!white,draw=blue!23!white,line width=.12pt] (-0.888912,0.665176)--(-0.897452,0.806729)--(-1.024781,0.759512)--(-1.015029,0.618408)--cycle;
\path[fill=blue!18!white,draw=blue!18!white,line width=.12pt] (1.015029,0.618408)--(1.024781,0.759512)--(0.897452,0.806729)--(0.888912,0.665176)--cycle;
\path[fill=blue!40!white,draw=blue!40!white,line width=.12pt] (-2.217311,-0.233600)--(-2.073119,-0.250151)--(-2.147305,-0.353819)--(-2.296657,-0.344479)--cycle;
\path[fill=blue!39!white,draw=blue!39!white,line width=.12pt] (2.296657,-0.344479)--(2.147305,-0.353819)--(2.073119,-0.250151)--(2.217311,-0.233600)--cycle;
\path[fill=blue!14!white,draw=blue!14!white,line width=.12pt] (0.922745,0.955623)--(0.963817,1.106137)--(0.817791,1.149900)--(0.782941,0.997521)--cycle;
\path[fill=blue!18!white,draw=blue!18!white,line width=.12pt] (-0.782941,0.997521)--(-0.817791,1.149900)--(-0.963817,1.106137)--(-0.922745,0.955623)--cycle;
\path[fill=blue!13!white,draw=blue!13!white,line width=.12pt] (0.963817,1.106137)--(1.019091,1.252486)--(0.864690,1.298758)--(0.817791,1.149900)--cycle;
\path[fill=blue!16!white,draw=blue!16!white,line width=.12pt] (-0.817791,1.149900)--(-0.864690,1.298758)--(-1.019091,1.252486)--(-0.963817,1.106137)--cycle;
\path[fill=blue!26!white,draw=blue!26!white,line width=.12pt] (-1.024781,0.489187)--(-1.015029,0.618408)--(-1.131371,0.565685)--(-1.142240,0.435958)--cycle;
\path[fill=blue!20!white,draw=blue!20!white,line width=.12pt] (1.142240,0.435958)--(1.131371,0.565685)--(1.015029,0.618408)--(1.024781,0.489187)--cycle;
\path[fill=blue!30!white,draw=blue!30!white,line width=.12pt] (-1.174431,0.322085)--(-1.142240,0.435958)--(-1.248699,0.377228)--(-1.283890,0.261700)--cycle;
\path[fill=blue!23!white,draw=blue!23!white,line width=.12pt] (1.283890,0.261700)--(1.248699,0.377228)--(1.142240,0.435958)--(1.174431,0.322085)--cycle;
\path[fill=blue!26!white,draw=blue!26!white,line width=.12pt] (1.442454,0.096998)--(1.380985,0.196241)--(1.283890,0.261700)--(1.341037,0.165370)--cycle;
\path[fill=blue!33!white,draw=blue!33!white,line width=.12pt] (-1.341037,0.165370)--(-1.283890,0.261700)--(-1.380985,0.196241)--(-1.442454,0.096998)--cycle;
\path[fill=blue!15!white,draw=blue!15!white,line width=.12pt] (-1.086442,1.389047)--(-1.163283,1.510570)--(-1.328326,1.449368)--(-1.240584,1.331887)--cycle;
\path[fill=blue!12!white,draw=blue!12!white,line width=.12pt] (1.240584,1.331887)--(1.328326,1.449368)--(1.163283,1.510570)--(1.086442,1.389047)--cycle;
\path[fill=blue!22!white,draw=blue!22!white,line width=.12pt] (-2.768442,1.005012)--(-2.808591,0.885774)--(-2.942213,0.744515)--(-2.900154,0.865772)--cycle;
\path[fill=blue!30!white,draw=blue!30!white,line width=.12pt] (2.900154,0.865772)--(2.942213,0.744515)--(2.808591,0.885774)--(2.768442,1.005012)--cycle;
\path[fill=blue!29!white,draw=blue!29!white,line width=.12pt] (1.617721,-0.057553)--(1.529978,0.023985)--(1.442454,0.096998)--(1.525177,0.019647)--cycle;
\path[fill=blue!35!white,draw=blue!35!white,line width=.12pt] (-1.525177,0.019647)--(-1.442454,0.096998)--(-1.529978,0.023985)--(-1.617721,-0.057553)--cycle;
\path[fill=blue!15!white,draw=blue!15!white,line width=.12pt] (1.855225,1.454092)--(1.975871,1.490286)--(1.807416,1.583216)--(1.697056,1.541348)--cycle;
\path[fill=blue!15!white,draw=blue!15!white,line width=.12pt] (-1.697056,1.541348)--(-1.807416,1.583216)--(-1.975871,1.490286)--(-1.855225,1.454092)--cycle;
\path[fill=blue!29!white,draw=blue!29!white,line width=.12pt] (-2.942213,0.474190)--(-2.900154,0.335511)--(-3.003935,0.190486)--(-3.047499,0.327062)--cycle;
\path[fill=blue!37!white,draw=blue!37!white,line width=.12pt] (3.047499,0.327062)--(3.003935,0.190486)--(2.900154,0.335511)--(2.942213,0.474190)--cycle;
\path[fill=blue!21!white,draw=blue!21!white,line width=.12pt] (-0.897452,0.806729)--(-0.922745,0.955623)--(-1.053662,0.907076)--(-1.024781,0.759512)--cycle;
\path[fill=blue!16!white,draw=blue!16!white,line width=.12pt] (1.024781,0.759512)--(1.053662,0.907076)--(0.922745,0.955623)--(0.897452,0.806729)--cycle;
\path[fill=blue!40!white,draw=blue!40!white,line width=.12pt] (-2.073119,-0.250151)--(-1.934468,-0.239441)--(-2.003693,-0.336176)--(-2.147305,-0.353819)--cycle;
\path[fill=blue!37!white,draw=blue!37!white,line width=.12pt] (2.147305,-0.353819)--(2.003693,-0.336176)--(1.934468,-0.239441)--(2.073119,-0.250151)--cycle;
\path[fill=blue!14!white,draw=blue!14!white,line width=.12pt] (-1.328326,1.449368)--(-1.423533,1.546798)--(-1.586697,1.472856)--(-1.480578,1.380371)--cycle;
\path[fill=blue!13!white,draw=blue!13!white,line width=.12pt] (1.480578,1.380371)--(1.586697,1.472856)--(1.423533,1.546798)--(1.328326,1.449368)--cycle;
\path[fill=blue!12!white,draw=blue!12!white,line width=.12pt] (1.163677,1.198870)--(1.240584,1.331887)--(1.086442,1.389047)--(1.019091,1.252486)--cycle;
\path[fill=blue!16!white,draw=blue!16!white,line width=.12pt] (-1.019091,1.252486)--(-1.086442,1.389047)--(-1.240584,1.331887)--(-1.163677,1.198870)--cycle;
\path[fill=blue!37!white,draw=blue!37!white,line width=.12pt] (-1.724636,-0.115141)--(-1.617721,-0.057553)--(-1.694686,-0.138917)--(-1.806687,-0.201882)--cycle;
\path[fill=blue!32!white,draw=blue!32!white,line width=.12pt] (1.806687,-0.201882)--(1.694686,-0.138917)--(1.617721,-0.057553)--(1.724636,-0.115141)--cycle;
\path[fill=blue!19!white,draw=blue!19!white,line width=.12pt] (-0.922745,0.955623)--(-0.963817,1.106137)--(-1.100561,1.055429)--(-1.053662,0.907076)--cycle;
\path[fill=blue!14!white,draw=blue!14!white,line width=.12pt] (1.053662,0.907076)--(1.100561,1.055429)--(0.963817,1.106137)--(0.922745,0.955623)--cycle;
\path[fill=blue!24!white,draw=blue!24!white,line width=.12pt] (-1.015029,0.618408)--(-1.024781,0.759512)--(-1.142240,0.706283)--(-1.131371,0.565685)--cycle;
\path[fill=blue!18!white,draw=blue!18!white,line width=.12pt] (1.131371,0.565685)--(1.142240,0.706283)--(1.024781,0.759512)--(1.015029,0.618408)--cycle;
\path[fill=blue!17!white,draw=blue!17!white,line width=.12pt] (-0.963817,1.106137)--(-1.019091,1.252486)--(-1.163677,1.198870)--(-1.100561,1.055429)--cycle;
\path[fill=blue!13!white,draw=blue!13!white,line width=.12pt] (1.100561,1.055429)--(1.163677,1.198870)--(1.019091,1.252486)--(0.963817,1.106137)--cycle;
\path[fill=blue!22!white,draw=blue!22!white,line width=.12pt] (2.508586,1.246493)--(2.615501,1.188905)--(2.465877,1.313721)--(2.365079,1.366207)--cycle;
\path[fill=blue!17!white,draw=blue!17!white,line width=.12pt] (-2.365079,1.366207)--(-2.465877,1.313721)--(-2.615501,1.188905)--(-2.508586,1.246493)--cycle;
\path[fill=blue!15!white,draw=blue!15!white,line width=.12pt] (-1.975871,1.490286)--(-2.091880,1.498463)--(-2.250079,1.391810)--(-2.125296,1.389547)--cycle;
\path[fill=blue!17!white,draw=blue!17!white,line width=.12pt] (2.125296,1.389547)--(2.250079,1.391810)--(2.091880,1.498463)--(1.975871,1.490286)--cycle;
\path[fill=blue!27!white,draw=blue!27!white,line width=.12pt] (-1.142240,0.435958)--(-1.131371,0.565685)--(-1.236817,0.507515)--(-1.248699,0.377228)--cycle;
\path[fill=blue!20!white,draw=blue!20!white,line width=.12pt] (1.248699,0.377228)--(1.236817,0.507515)--(1.131371,0.565685)--(1.142240,0.435958)--cycle;
\path[fill=blue!26!white,draw=blue!26!white,line width=.12pt] (-2.956415,0.612293)--(-2.942213,0.474190)--(-3.047499,0.327062)--(-3.062209,0.464455)--cycle;
\path[fill=blue!35!white,draw=blue!35!white,line width=.12pt] (3.062209,0.464455)--(3.047499,0.327062)--(2.942213,0.474190)--(2.956415,0.612293)--cycle;
\path[fill=blue!23!white,draw=blue!23!white,line width=.12pt] (1.380985,0.196241)--(1.343133,0.313564)--(1.248699,0.377228)--(1.283890,0.261700)--cycle;
\path[fill=blue!30!white,draw=blue!30!white,line width=.12pt] (-1.283890,0.261700)--(-1.248699,0.377228)--(-1.343133,0.313564)--(-1.380985,0.196241)--cycle;
\path[fill=blue!20!white,draw=blue!20!white,line width=.12pt] (-2.703244,1.107367)--(-2.768442,1.005012)--(-2.900154,0.865772)--(-2.831853,0.971406)--cycle;
\path[fill=blue!27!white,draw=blue!27!white,line width=.12pt] (2.831853,0.971406)--(2.900154,0.865772)--(2.768442,1.005012)--(2.703244,1.107367)--cycle;
\path[fill=blue!14!white,draw=blue!14!white,line width=.12pt] (1.734580,1.391274)--(1.855225,1.454092)--(1.697056,1.541348)--(1.586697,1.472856)--cycle;
\path[fill=blue!14!white,draw=blue!14!white,line width=.12pt] (-1.586697,1.472856)--(-1.697056,1.541348)--(-1.855225,1.454092)--(-1.734580,1.391274)--cycle;
\path[fill=blue!15!white,draw=blue!15!white,line width=.12pt] (-1.240584,1.331887)--(-1.328326,1.449368)--(-1.480578,1.380371)--(-1.382778,1.267448)--cycle;
\path[fill=blue!12!white,draw=blue!12!white,line width=.12pt] (1.382778,1.267448)--(1.480578,1.380371)--(1.328326,1.449368)--(1.240584,1.331887)--cycle;
\path[fill=blue!33!white,draw=blue!33!white,line width=.12pt] (-1.442454,0.096998)--(-1.380985,0.196241)--(-1.464780,0.126340)--(-1.529978,0.023985)--cycle;
\path[fill=blue!26!white,draw=blue!26!white,line width=.12pt] (1.529978,0.023985)--(1.464780,0.126340)--(1.380985,0.196241)--(1.442454,0.096998)--cycle;
\path[fill=blue!39!white,draw=blue!39!white,line width=.12pt] (-1.934468,-0.239441)--(-1.806687,-0.201882)--(-1.871339,-0.292227)--(-2.003693,-0.336176)--cycle;
\path[fill=blue!35!white,draw=blue!35!white,line width=.12pt] (2.003693,-0.336176)--(1.871339,-0.292227)--(1.806687,-0.201882)--(1.934468,-0.239441)--cycle;
\path[fill=blue!37!white,draw=blue!37!white,line width=.12pt] (-2.721975,-0.163208)--(-2.589621,-0.247306)--(-2.654149,-0.376111)--(-2.789801,-0.298596)--cycle;
\path[fill=blue!41!white,draw=blue!41!white,line width=.12pt] (2.789801,-0.298596)--(2.654149,-0.376111)--(2.589621,-0.247306)--(2.721975,-0.163208)--cycle;
\path[fill=blue!16!white,draw=blue!16!white,line width=.12pt] (1.142240,0.706283)--(1.174431,0.852346)--(1.053662,0.907076)--(1.024781,0.759512)--cycle;
\path[fill=blue!21!white,draw=blue!21!white,line width=.12pt] (-1.024781,0.759512)--(-1.053662,0.907076)--(-1.174431,0.852346)--(-1.142240,0.706283)--cycle;
\path[fill=blue!38!white,draw=blue!38!white,line width=.12pt] (-2.589621,-0.247306)--(-2.446009,-0.308514)--(-2.506958,-0.430175)--(-2.654149,-0.376111)--cycle;
\path[fill=blue!41!white,draw=blue!41!white,line width=.12pt] (2.654149,-0.376111)--(2.506958,-0.430175)--(2.446009,-0.308514)--(2.589621,-0.247306)--cycle;
\path[fill=blue!35!white,draw=blue!35!white,line width=.12pt] (-2.837984,-0.059451)--(-2.721975,-0.163208)--(-2.789801,-0.298596)--(-2.908701,-0.200610)--cycle;
\path[fill=blue!41!white,draw=blue!41!white,line width=.12pt] (2.908701,-0.200610)--(2.789801,-0.298596)--(2.721975,-0.163208)--(2.837984,-0.059451)--cycle;
\path[fill=blue!39!white,draw=blue!39!white,line width=.12pt] (-2.446009,-0.308514)--(-2.296657,-0.344479)--(-2.353885,-0.458712)--(-2.506958,-0.430175)--cycle;
\path[fill=blue!40!white,draw=blue!40!white,line width=.12pt] (2.506958,-0.430175)--(2.353885,-0.458712)--(2.296657,-0.344479)--(2.446009,-0.308514)--cycle;
\path[fill=blue!36!white,draw=blue!36!white,line width=.12pt] (-1.617721,-0.057553)--(-1.529978,0.023985)--(-1.602769,-0.052966)--(-1.694686,-0.138917)--cycle;
\path[fill=blue!30!white,draw=blue!30!white,line width=.12pt] (1.694686,-0.138917)--(1.602769,-0.052966)--(1.529978,0.023985)--(1.617721,-0.057553)--cycle;
\path[fill=blue!16!white,draw=blue!16!white,line width=.12pt] (-1.163677,1.198870)--(-1.240584,1.331887)--(-1.382778,1.267448)--(-1.297056,1.138426)--cycle;
\path[fill=blue!12!white,draw=blue!12!white,line width=.12pt] (1.297056,1.138426)--(1.382778,1.267448)--(1.240584,1.331887)--(1.163677,1.198870)--cycle;
\path[fill=blue!14!white,draw=blue!14!white,line width=.12pt] (1.174431,0.852346)--(1.226706,0.998263)--(1.100561,1.055429)--(1.053662,0.907076)--cycle;
\path[fill=blue!19!white,draw=blue!19!white,line width=.12pt] (-1.053662,0.907076)--(-1.100561,1.055429)--(-1.226706,0.998263)--(-1.174431,0.852346)--cycle;
\path[fill=blue!18!white,draw=blue!18!white,line width=.12pt] (1.236817,0.507515)--(1.248699,0.647553)--(1.142240,0.706283)--(1.131371,0.565685)--cycle;
\path[fill=blue!25!white,draw=blue!25!white,line width=.12pt] (-1.131371,0.565685)--(-1.142240,0.706283)--(-1.248699,0.647553)--(-1.236817,0.507515)--cycle;
\path[fill=blue!24!white,draw=blue!24!white,line width=.12pt] (-2.942213,0.744515)--(-2.956415,0.612293)--(-3.062209,0.464455)--(-3.047499,0.597387)--cycle;
\path[fill=blue!32!white,draw=blue!32!white,line width=.12pt] (3.047499,0.597387)--(3.062209,0.464455)--(2.956415,0.612293)--(2.942213,0.744515)--cycle;
\path[fill=blue!17!white,draw=blue!17!white,line width=.12pt] (-1.100561,1.055429)--(-1.163677,1.198870)--(-1.297056,1.138426)--(-1.226706,0.998263)--cycle;
\path[fill=blue!13!white,draw=blue!13!white,line width=.12pt] (1.226706,0.998263)--(1.297056,1.138426)--(1.163677,1.198870)--(1.100561,1.055429)--cycle;
\path[fill=blue!33!white,draw=blue!33!white,line width=.12pt] (-2.933190,0.059976)--(-2.837984,-0.059451)--(-2.908701,-0.200610)--(-3.006279,-0.085917)--cycle;
\path[fill=blue!40!white,draw=blue!40!white,line width=.12pt] (3.006279,-0.085917)--(2.908701,-0.200610)--(2.837984,-0.059451)--(2.933190,0.059976)--cycle;
\path[fill=blue!40!white,draw=blue!40!white,line width=.12pt] (-2.296657,-0.344479)--(-2.147305,-0.353819)--(-2.200811,-0.460624)--(-2.353885,-0.458712)--cycle;
\path[fill=blue!39!white,draw=blue!39!white,line width=.12pt] (2.353885,-0.458712)--(2.200811,-0.460624)--(2.147305,-0.353819)--(2.296657,-0.344479)--cycle;
\path[fill=blue!19!white,draw=blue!19!white,line width=.12pt] (2.386608,1.277917)--(2.508586,1.246493)--(2.365079,1.366207)--(2.250079,1.391810)--cycle;
\path[fill=blue!16!white,draw=blue!16!white,line width=.12pt] (-2.250079,1.391810)--(-2.365079,1.366207)--(-2.508586,1.246493)--(-2.386608,1.277917)--cycle;
\path[fill=blue!14!white,draw=blue!14!white,line width=.12pt] (-1.855225,1.454092)--(-1.975871,1.490286)--(-2.125296,1.389547)--(-1.995527,1.359505)--cycle;
\path[fill=blue!15!white,draw=blue!15!white,line width=.12pt] (1.995527,1.359505)--(2.125296,1.389547)--(1.975871,1.490286)--(1.855225,1.454092)--cycle;
\path[fill=blue!15!white,draw=blue!15!white,line width=.12pt] (-1.480578,1.380371)--(-1.586697,1.472856)--(-1.734580,1.391274)--(-1.618570,1.304246)--cycle;
\path[fill=blue!13!white,draw=blue!13!white,line width=.12pt] (1.618570,1.304246)--(1.734580,1.391274)--(1.586697,1.472856)--(1.480578,1.380371)--cycle;
\path[fill=blue!28!white,draw=blue!28!white,line width=.12pt] (-1.248699,0.377228)--(-1.236817,0.507515)--(-1.330351,0.444456)--(-1.343133,0.313564)--cycle;
\path[fill=blue!21!white,draw=blue!21!white,line width=.12pt] (1.343133,0.313564)--(1.330351,0.444456)--(1.236817,0.507515)--(1.248699,0.377228)--cycle;
\path[fill=blue!18!white,draw=blue!18!white,line width=.12pt] (-2.615501,1.188905)--(-2.703244,1.107367)--(-2.831853,0.971406)--(-2.739936,1.057357)--cycle;
\path[fill=blue!24!white,draw=blue!24!white,line width=.12pt] (2.739936,1.057357)--(2.831853,0.971406)--(2.703244,1.107367)--(2.615501,1.188905)--cycle;
\path[fill=blue!38!white,draw=blue!38!white,line width=.12pt] (-1.806687,-0.201882)--(-1.694686,-0.138917)--(-1.755330,-0.223661)--(-1.871339,-0.292227)--cycle;
\path[fill=blue!33!white,draw=blue!33!white,line width=.12pt] (1.871339,-0.292227)--(1.755330,-0.223661)--(1.694686,-0.138917)--(1.806687,-0.201882)--cycle;
\path[fill=blue!24!white,draw=blue!24!white,line width=.12pt] (1.464780,0.126340)--(1.424631,0.245578)--(1.343133,0.313564)--(1.380985,0.196241)--cycle;
\path[fill=blue!31!white,draw=blue!31!white,line width=.12pt] (-1.380985,0.196241)--(-1.343133,0.313564)--(-1.424631,0.245578)--(-1.464780,0.126340)--cycle;
\path[fill=blue!31!white,draw=blue!31!white,line width=.12pt] (-3.003935,0.190486)--(-2.933190,0.059976)--(-3.006279,-0.085917)--(-3.078787,0.041074)--cycle;
\path[fill=blue!38!white,draw=blue!38!white,line width=.12pt] (3.078787,0.041074)--(3.006279,-0.085917)--(2.933190,0.059976)--(3.003935,0.190486)--cycle;
\path[fill=blue!40!white,draw=blue!40!white,line width=.12pt] (-2.147305,-0.353819)--(-2.003693,-0.336176)--(-2.053620,-0.435837)--(-2.200811,-0.460624)--cycle;
\path[fill=blue!37!white,draw=blue!37!white,line width=.12pt] (2.200811,-0.460624)--(2.053620,-0.435837)--(2.003693,-0.336176)--(2.147305,-0.353819)--cycle;
\path[fill=blue!16!white,draw=blue!16!white,line width=.12pt] (1.248699,0.647553)--(1.283890,0.791962)--(1.174431,0.852346)--(1.142240,0.706283)--cycle;
\path[fill=blue!22!white,draw=blue!22!white,line width=.12pt] (-1.142240,0.706283)--(-1.174431,0.852346)--(-1.283890,0.791962)--(-1.248699,0.647553)--cycle;
\path[fill=blue!15!white,draw=blue!15!white,line width=.12pt] (-1.382778,1.267448)--(-1.480578,1.380371)--(-1.618570,1.304246)--(-1.511656,1.196351)--cycle;
\path[fill=blue!12!white,draw=blue!12!white,line width=.12pt] (1.511656,1.196351)--(1.618570,1.304246)--(1.480578,1.380371)--(1.382778,1.267448)--cycle;
\path[fill=blue!34!white,draw=blue!34!white,line width=.12pt] (-1.529978,0.023985)--(-1.464780,0.126340)--(-1.534468,0.052668)--(-1.602769,-0.052966)--cycle;
\path[fill=blue!27!white,draw=blue!27!white,line width=.12pt] (1.602769,-0.052966)--(1.534468,0.052668)--(1.464780,0.126340)--(1.529978,0.023985)--cycle;
\path[fill=blue!21!white,draw=blue!21!white,line width=.12pt] (-2.900154,0.865772)--(-2.942213,0.744515)--(-3.047499,0.597387)--(-3.003935,0.720748)--cycle;
\path[fill=blue!30!white,draw=blue!30!white,line width=.12pt] (3.003935,0.720748)--(3.047499,0.597387)--(2.942213,0.744515)--(2.900154,0.865772)--cycle;
\path[fill=blue!14!white,draw=blue!14!white,line width=.12pt] (-1.734580,1.391274)--(-1.855225,1.454092)--(-1.995527,1.359505)--(-1.865758,1.302838)--cycle;
\path[fill=blue!14!white,draw=blue!14!white,line width=.12pt] (1.865758,1.302838)--(1.995527,1.359505)--(1.855225,1.454092)--(1.734580,1.391274)--cycle;
\path[fill=blue!14!white,draw=blue!14!white,line width=.12pt] (1.283890,0.791962)--(1.341037,0.935191)--(1.226706,0.998263)--(1.174431,0.852346)--cycle;
\path[fill=blue!20!white,draw=blue!20!white,line width=.12pt] (-1.174431,0.852346)--(-1.226706,0.998263)--(-1.341037,0.935191)--(-1.283890,0.791962)--cycle;
\path[fill=blue!16!white,draw=blue!16!white,line width=.12pt] (-1.297056,1.138426)--(-1.382778,1.267448)--(-1.511656,1.196351)--(-1.417944,1.071736)--cycle;
\path[fill=blue!12!white,draw=blue!12!white,line width=.12pt] (1.417944,1.071736)--(1.511656,1.196351)--(1.382778,1.267448)--(1.297056,1.138426)--cycle;
\path[fill=blue!25!white,draw=blue!25!white,line width=.12pt] (-1.236817,0.507515)--(-1.248699,0.647553)--(-1.343133,0.583889)--(-1.330351,0.444456)--cycle;
\path[fill=blue!18!white,draw=blue!18!white,line width=.12pt] (1.330351,0.444456)--(1.343133,0.583889)--(1.248699,0.647553)--(1.236817,0.507515)--cycle;
\path[fill=blue!15!white,draw=blue!15!white,line width=.12pt] (-2.125296,1.389547)--(-2.250079,1.391810)--(-2.386608,1.277917)--(-2.254254,1.281970)--cycle;
\path[fill=blue!17!white,draw=blue!17!white,line width=.12pt] (2.254254,1.281970)--(2.386608,1.277917)--(2.250079,1.391810)--(2.125296,1.389547)--cycle;
\path[fill=blue!13!white,draw=blue!13!white,line width=.12pt] (1.341037,0.935191)--(1.417944,1.071736)--(1.297056,1.138426)--(1.226706,0.998263)--cycle;
\path[fill=blue!18!white,draw=blue!18!white,line width=.12pt] (-1.226706,0.998263)--(-1.297056,1.138426)--(-1.417944,1.071736)--(-1.341037,0.935191)--cycle;
\path[fill=blue!28!white,draw=blue!28!white,line width=.12pt] (-3.047499,0.327062)--(-3.003935,0.190486)--(-3.078787,0.041074)--(-3.123436,0.175483)--cycle;
\path[fill=blue!36!white,draw=blue!36!white,line width=.12pt] (3.123436,0.175483)--(3.078787,0.041074)--(3.003935,0.190486)--(3.047499,0.327062)--cycle;
\path[fill=blue!39!white,draw=blue!39!white,line width=.12pt] (-2.003693,-0.336176)--(-1.871339,-0.292227)--(-1.917969,-0.385305)--(-2.053620,-0.435837)--cycle;
\path[fill=blue!35!white,draw=blue!35!white,line width=.12pt] (2.053620,-0.435837)--(1.917969,-0.385305)--(1.871339,-0.292227)--(2.003693,-0.336176)--cycle;
\path[fill=blue!37!white,draw=blue!37!white,line width=.12pt] (-1.694686,-0.138917)--(-1.602769,-0.052966)--(-1.660123,-0.133114)--(-1.755330,-0.223661)--cycle;
\path[fill=blue!30!white,draw=blue!30!white,line width=.12pt] (1.755330,-0.223661)--(1.660123,-0.133114)--(1.602769,-0.052966)--(1.694686,-0.138917)--cycle;
\path[fill=blue!17!white,draw=blue!17!white,line width=.12pt] (-2.508586,1.246493)--(-2.615501,1.188905)--(-2.739936,1.057357)--(-2.627935,1.120322)--cycle;
\path[fill=blue!22!white,draw=blue!22!white,line width=.12pt] (2.627935,1.120322)--(2.739936,1.057357)--(2.615501,1.188905)--(2.508586,1.246493)--cycle;
\path[fill=blue!29!white,draw=blue!29!white,line width=.12pt] (-1.343133,0.313564)--(-1.330351,0.444456)--(-1.411074,0.377117)--(-1.424631,0.245578)--cycle;
\path[fill=blue!21!white,draw=blue!21!white,line width=.12pt] (1.424631,0.245578)--(1.411074,0.377117)--(1.330351,0.444456)--(1.343133,0.313564)--cycle;
\path[fill=blue!13!white,draw=blue!13!white,line width=.12pt] (1.740975,1.221724)--(1.865758,1.302838)--(1.734580,1.391274)--(1.618570,1.304246)--cycle;
\path[fill=blue!15!white,draw=blue!15!white,line width=.12pt] (-1.618570,1.304246)--(-1.734580,1.391274)--(-1.865758,1.302838)--(-1.740975,1.221724)--cycle;
\path[fill=blue!16!white,draw=blue!16!white,line width=.12pt] (1.343133,0.583889)--(1.380985,0.726503)--(1.283890,0.791962)--(1.248699,0.647553)--cycle;
\path[fill=blue!23!white,draw=blue!23!white,line width=.12pt] (-1.248699,0.647553)--(-1.283890,0.791962)--(-1.380985,0.726503)--(-1.343133,0.583889)--cycle;
\path[fill=blue!32!white,draw=blue!32!white,line width=.12pt] (-1.464780,0.126340)--(-1.424631,0.245578)--(-1.492409,0.173925)--(-1.534468,0.052668)--cycle;
\path[fill=blue!24!white,draw=blue!24!white,line width=.12pt] (1.534468,0.052668)--(1.492409,0.173925)--(1.424631,0.245578)--(1.464780,0.126340)--cycle;
\path[fill=blue!19!white,draw=blue!19!white,line width=.12pt] (-2.831853,0.971406)--(-2.900154,0.865772)--(-3.003935,0.720748)--(-2.933190,0.829797)--cycle;
\path[fill=blue!27!white,draw=blue!27!white,line width=.12pt] (2.933190,0.829797)--(3.003935,0.720748)--(2.900154,0.865772)--(2.831853,0.971406)--cycle;
\path[fill=blue!38!white,draw=blue!38!white,line width=.12pt] (-2.654149,-0.376111)--(-2.506958,-0.430175)--(-2.543764,-0.554239)--(-2.693116,-0.507458)--cycle;
\path[fill=blue!41!white,draw=blue!41!white,line width=.12pt] (2.693116,-0.507458)--(2.543764,-0.554239)--(2.506958,-0.430175)--(2.654149,-0.376111)--cycle;
\path[fill=blue!39!white,draw=blue!39!white,line width=.12pt] (-2.506958,-0.430175)--(-2.353885,-0.458712)--(-2.388443,-0.575200)--(-2.543764,-0.554239)--cycle;
\path[fill=blue!40!white,draw=blue!40!white,line width=.12pt] (2.543764,-0.554239)--(2.388443,-0.575200)--(2.353885,-0.458712)--(2.506958,-0.430175)--cycle;
\path[fill=blue!26!white,draw=blue!26!white,line width=.12pt] (-3.062209,0.464455)--(-3.047499,0.327062)--(-3.123436,0.175483)--(-3.138513,0.312145)--cycle;
\path[fill=blue!34!white,draw=blue!34!white,line width=.12pt] (3.138513,0.312145)--(3.123436,0.175483)--(3.047499,0.327062)--(3.062209,0.464455)--cycle;
\path[fill=blue!38!white,draw=blue!38!white,line width=.12pt] (-1.871339,-0.292227)--(-1.755330,-0.223661)--(-1.799069,-0.310969)--(-1.917969,-0.385305)--cycle;
\path[fill=blue!33!white,draw=blue!33!white,line width=.12pt] (1.917969,-0.385305)--(1.799069,-0.310969)--(1.755330,-0.223661)--(1.871339,-0.292227)--cycle;
\path[fill=blue!14!white,draw=blue!14!white,line width=.12pt] (-1.995527,1.359505)--(-2.125296,1.389547)--(-2.254254,1.281970)--(-2.116611,1.258496)--cycle;
\path[fill=blue!15!white,draw=blue!15!white,line width=.12pt] (2.116611,1.258496)--(2.254254,1.281970)--(2.125296,1.389547)--(1.995527,1.359505)--cycle;
\path[fill=blue!36!white,draw=blue!36!white,line width=.12pt] (-2.789801,-0.298596)--(-2.654149,-0.376111)--(-2.693116,-0.507458)--(-2.830759,-0.436656)--cycle;
\path[fill=blue!41!white,draw=blue!41!white,line width=.12pt] (2.830759,-0.436656)--(2.693116,-0.507458)--(2.654149,-0.376111)--(2.789801,-0.298596)--cycle;
\path[fill=blue!12!white,draw=blue!12!white,line width=.12pt] (1.625975,1.119280)--(1.740975,1.221724)--(1.618570,1.304246)--(1.511656,1.196351)--cycle;
\path[fill=blue!16!white,draw=blue!16!white,line width=.12pt] (-1.511656,1.196351)--(-1.618570,1.304246)--(-1.740975,1.221724)--(-1.625975,1.119280)--cycle;
\path[fill=blue!40!white,draw=blue!40!white,line width=.12pt] (-2.353885,-0.458712)--(-2.200811,-0.460624)--(-2.233123,-0.569536)--(-2.388443,-0.575200)--cycle;
\path[fill=blue!39!white,draw=blue!39!white,line width=.12pt] (2.388443,-0.575200)--(2.233123,-0.569536)--(2.200811,-0.460624)--(2.353885,-0.458712)--cycle;
\path[fill=blue!14!white,draw=blue!14!white,line width=.12pt] (1.380985,0.726503)--(1.442454,0.866819)--(1.341037,0.935191)--(1.283890,0.791962)--cycle;
\path[fill=blue!21!white,draw=blue!21!white,line width=.12pt] (-1.283890,0.791962)--(-1.341037,0.935191)--(-1.442454,0.866819)--(-1.380985,0.726503)--cycle;
\path[fill=blue!35!white,draw=blue!35!white,line width=.12pt] (-1.602769,-0.052966)--(-1.534468,0.052668)--(-1.589379,-0.024064)--(-1.660123,-0.133114)--cycle;
\path[fill=blue!27!white,draw=blue!27!white,line width=.12pt] (1.660123,-0.133114)--(1.589379,-0.024064)--(1.534468,0.052668)--(1.602769,-0.052966)--cycle;
\path[fill=blue!35!white,draw=blue!35!white,line width=.12pt] (-2.908701,-0.200610)--(-2.789801,-0.298596)--(-2.830759,-0.436656)--(-2.951405,-0.344554)--cycle;
\path[fill=blue!40!white,draw=blue!40!white,line width=.12pt] (2.951405,-0.344554)--(2.830759,-0.436656)--(2.789801,-0.298596)--(2.908701,-0.200610)--cycle;
\path[fill=blue!26!white,draw=blue!26!white,line width=.12pt] (-1.330351,0.444456)--(-1.343133,0.583889)--(-1.424631,0.515903)--(-1.411074,0.377117)--cycle;
\path[fill=blue!18!white,draw=blue!18!white,line width=.12pt] (1.411074,0.377117)--(1.424631,0.515903)--(1.343133,0.583889)--(1.330351,0.444456)--cycle;
\path[fill=blue!15!white,draw=blue!15!white,line width=.12pt] (-2.386608,1.277917)--(-2.508586,1.246493)--(-2.627935,1.120322)--(-2.500154,1.157881)--cycle;
\path[fill=blue!19!white,draw=blue!19!white,line width=.12pt] (2.500154,1.157881)--(2.627935,1.120322)--(2.508586,1.246493)--(2.386608,1.277917)--cycle;
\path[fill=blue!12!white,draw=blue!12!white,line width=.12pt] (1.525177,0.999443)--(1.625975,1.119280)--(1.511656,1.196351)--(1.417944,1.071736)--cycle;
\path[fill=blue!17!white,draw=blue!17!white,line width=.12pt] (-1.417944,1.071736)--(-1.511656,1.196351)--(-1.625975,1.119280)--(-1.525177,0.999443)--cycle;
\path[fill=blue!19!white,draw=blue!19!white,line width=.12pt] (-1.341037,0.935191)--(-1.417944,1.071736)--(-1.525177,0.999443)--(-1.442454,0.866819)--cycle;
\path[fill=blue!13!white,draw=blue!13!white,line width=.12pt] (1.442454,0.866819)--(1.525177,0.999443)--(1.417944,1.071736)--(1.341037,0.935191)--cycle;
\path[fill=blue!40!white,draw=blue!40!white,line width=.12pt] (-2.200811,-0.460624)--(-2.053620,-0.435837)--(-2.083771,-0.537466)--(-2.233123,-0.569536)--cycle;
\path[fill=blue!37!white,draw=blue!37!white,line width=.12pt] (2.233123,-0.569536)--(2.083771,-0.537466)--(2.053620,-0.435837)--(2.200811,-0.460624)--cycle;
\path[fill=blue!14!white,draw=blue!14!white,line width=.12pt] (-1.865758,1.302838)--(-1.995527,1.359505)--(-2.116611,1.258496)--(-1.978968,1.208398)--cycle;
\path[fill=blue!14!white,draw=blue!14!white,line width=.12pt] (1.978968,1.208398)--(2.116611,1.258496)--(1.995527,1.359505)--(1.865758,1.302838)--cycle;
\path[fill=blue!32!white,draw=blue!32!white,line width=.12pt] (-3.006279,-0.085917)--(-2.908701,-0.200610)--(-2.951405,-0.344554)--(-3.050416,-0.234690)--cycle;
\path[fill=blue!39!white,draw=blue!39!white,line width=.12pt] (3.050416,-0.234690)--(2.951405,-0.344554)--(2.908701,-0.200610)--(3.006279,-0.085917)--cycle;
\path[fill=blue!29!white,draw=blue!29!white,line width=.12pt] (-1.424631,0.245578)--(-1.411074,0.377117)--(-1.478207,0.306147)--(-1.492409,0.173925)--cycle;
\path[fill=blue!21!white,draw=blue!21!white,line width=.12pt] (1.492409,0.173925)--(1.478207,0.306147)--(1.411074,0.377117)--(1.424631,0.245578)--cycle;
\path[fill=blue!23!white,draw=blue!23!white,line width=.12pt] (-3.047499,0.597387)--(-3.062209,0.464455)--(-3.138513,0.312145)--(-3.123436,0.445808)--cycle;
\path[fill=blue!32!white,draw=blue!32!white,line width=.12pt] (3.123436,0.445808)--(3.138513,0.312145)--(3.062209,0.464455)--(3.047499,0.597387)--cycle;
\path[fill=blue!37!white,draw=blue!37!white,line width=.12pt] (-1.755330,-0.223661)--(-1.660123,-0.133114)--(-1.701490,-0.215687)--(-1.799069,-0.310969)--cycle;
\path[fill=blue!30!white,draw=blue!30!white,line width=.12pt] (1.799069,-0.310969)--(1.701490,-0.215687)--(1.660123,-0.133114)--(1.755330,-0.223661)--cycle;
\path[fill=blue!17!white,draw=blue!17!white,line width=.12pt] (-2.739936,1.057357)--(-2.831853,0.971406)--(-2.933190,0.829797)--(-2.837984,0.920344)--cycle;
\path[fill=blue!24!white,draw=blue!24!white,line width=.12pt] (2.837984,0.920344)--(2.933190,0.829797)--(2.831853,0.971406)--(2.739936,1.057357)--cycle;
\path[fill=blue!24!white,draw=blue!24!white,line width=.12pt] (-1.343133,0.583889)--(-1.380985,0.726503)--(-1.464780,0.656601)--(-1.424631,0.515903)--cycle;
\path[fill=blue!16!white,draw=blue!16!white,line width=.12pt] (1.424631,0.515903)--(1.464780,0.656601)--(1.380985,0.726503)--(1.343133,0.583889)--cycle;
\path[fill=blue!40!white,draw=blue!40!white,line width=.12pt] (-2.053620,-0.435837)--(-1.917969,-0.385305)--(-1.946127,-0.480221)--(-2.083771,-0.537466)--cycle;
\path[fill=blue!36!white,draw=blue!36!white,line width=.12pt] (2.083771,-0.537466)--(1.946127,-0.480221)--(1.917969,-0.385305)--(2.053620,-0.435837)--cycle;
\path[fill=blue!24!white,draw=blue!24!white,line width=.12pt] (1.589379,-0.024064)--(1.545814,0.099296)--(1.492409,0.173925)--(1.534468,0.052668)--cycle;
\path[fill=blue!33!white,draw=blue!33!white,line width=.12pt] (-1.534468,0.052668)--(-1.492409,0.173925)--(-1.545814,0.099296)--(-1.589379,-0.024064)--cycle;
\path[fill=blue!15!white,draw=blue!15!white,line width=.12pt] (-2.254254,1.281970)--(-2.386608,1.277917)--(-2.500154,1.157881)--(-2.361503,1.168591)--cycle;
\path[fill=blue!17!white,draw=blue!17!white,line width=.12pt] (2.361503,1.168591)--(2.500154,1.157881)--(2.386608,1.277917)--(2.254254,1.281970)--cycle;
\path[fill=blue!15!white,draw=blue!15!white,line width=.12pt] (-1.740975,1.221724)--(-1.865758,1.302838)--(-1.978968,1.208398)--(-1.846614,1.133600)--cycle;
\path[fill=blue!13!white,draw=blue!13!white,line width=.12pt] (1.846614,1.133600)--(1.978968,1.208398)--(1.865758,1.302838)--(1.740975,1.221724)--cycle;
\path[fill=blue!30!white,draw=blue!30!white,line width=.12pt] (-3.078787,0.041074)--(-3.006279,-0.085917)--(-3.050416,-0.234690)--(-3.123988,-0.111288)--cycle;
\path[fill=blue!38!white,draw=blue!38!white,line width=.12pt] (3.123988,-0.111288)--(3.050416,-0.234690)--(3.006279,-0.085917)--(3.078787,0.041074)--cycle;
\path[fill=blue!15!white,draw=blue!15!white,line width=.12pt] (1.464780,0.656601)--(1.529978,0.793806)--(1.442454,0.866819)--(1.380985,0.726503)--cycle;
\path[fill=blue!21!white,draw=blue!21!white,line width=.12pt] (-1.380985,0.726503)--(-1.442454,0.866819)--(-1.529978,0.793806)--(-1.464780,0.656601)--cycle;
\path[fill=blue!27!white,draw=blue!27!white,line width=.12pt] (-1.411074,0.377117)--(-1.424631,0.515903)--(-1.492409,0.444251)--(-1.478207,0.306147)--cycle;
\path[fill=blue!19!white,draw=blue!19!white,line width=.12pt] (1.478207,0.306147)--(1.492409,0.444251)--(1.424631,0.515903)--(1.411074,0.377117)--cycle;
\path[fill=blue!12!white,draw=blue!12!white,line width=.12pt] (1.724636,1.036978)--(1.846614,1.133600)--(1.740975,1.221724)--(1.625975,1.119280)--cycle;
\path[fill=blue!16!white,draw=blue!16!white,line width=.12pt] (-1.625975,1.119280)--(-1.740975,1.221724)--(-1.846614,1.133600)--(-1.724636,1.036978)--cycle;
\path[fill=blue!21!white,draw=blue!21!white,line width=.12pt] (-3.003935,0.720748)--(-3.047499,0.597387)--(-3.123436,0.445808)--(-3.078787,0.571335)--cycle;
\path[fill=blue!29!white,draw=blue!29!white,line width=.12pt] (3.078787,0.571335)--(3.123436,0.445808)--(3.047499,0.597387)--(3.003935,0.720748)--cycle;
\path[fill=blue!39!white,draw=blue!39!white,line width=.12pt] (-1.917969,-0.385305)--(-1.799069,-0.310969)--(-1.825482,-0.400001)--(-1.946127,-0.480221)--cycle;
\path[fill=blue!33!white,draw=blue!33!white,line width=.12pt] (1.946127,-0.480221)--(1.825482,-0.400001)--(1.799069,-0.310969)--(1.917969,-0.385305)--cycle;
\path[fill=blue!35!white,draw=blue!35!white,line width=.12pt] (-1.660123,-0.133114)--(-1.589379,-0.024064)--(-1.628983,-0.103118)--(-1.701490,-0.215687)--cycle;
\path[fill=blue!28!white,draw=blue!28!white,line width=.12pt] (1.701490,-0.215687)--(1.628983,-0.103118)--(1.589379,-0.024064)--(1.660123,-0.133114)--cycle;
\path[fill=blue!19!white,draw=blue!19!white,line width=.12pt] (-1.442454,0.866819)--(-1.525177,0.999443)--(-1.617721,0.922243)--(-1.529978,0.793806)--cycle;
\path[fill=blue!13!white,draw=blue!13!white,line width=.12pt] (1.529978,0.793806)--(1.617721,0.922243)--(1.525177,0.999443)--(1.442454,0.866819)--cycle;
\path[fill=blue!21!white,draw=blue!21!white,line width=.12pt] (2.721975,0.988910)--(2.837984,0.920344)--(2.739936,1.057357)--(2.627935,1.120322)--cycle;
\path[fill=blue!16!white,draw=blue!16!white,line width=.12pt] (-2.627935,1.120322)--(-2.739936,1.057357)--(-2.837984,0.920344)--(-2.721975,0.988910)--cycle;
\path[fill=blue!17!white,draw=blue!17!white,line width=.12pt] (-1.525177,0.999443)--(-1.625975,1.119280)--(-1.724636,1.036978)--(-1.617721,0.922243)--cycle;
\path[fill=blue!13!white,draw=blue!13!white,line width=.12pt] (1.617721,0.922243)--(1.724636,1.036978)--(1.625975,1.119280)--(1.525177,0.999443)--cycle;
\path[fill=blue!14!white,draw=blue!14!white,line width=.12pt] (-2.116611,1.258496)--(-2.254254,1.281970)--(-2.361503,1.168591)--(-2.217311,1.152040)--cycle;
\path[fill=blue!15!white,draw=blue!15!white,line width=.12pt] (2.217311,1.152040)--(2.361503,1.168591)--(2.254254,1.281970)--(2.116611,1.258496)--cycle;
\path[fill=blue!27!white,draw=blue!27!white,line width=.12pt] (-3.123436,0.175483)--(-3.078787,0.041074)--(-3.123988,-0.111288)--(-3.169293,0.020912)--cycle;
\path[fill=blue!36!white,draw=blue!36!white,line width=.12pt] (3.169293,0.020912)--(3.123988,-0.111288)--(3.078787,0.041074)--(3.123436,0.175483)--cycle;
\path[fill=blue!30!white,draw=blue!30!white,line width=.12pt] (-1.492409,0.173925)--(-1.478207,0.306147)--(-1.531105,0.232228)--(-1.545814,0.099296)--cycle;
\path[fill=blue!22!white,draw=blue!22!white,line width=.12pt] (1.545814,0.099296)--(1.531105,0.232228)--(1.478207,0.306147)--(1.492409,0.173925)--cycle;
\path[fill=blue!39!white,draw=blue!39!white,line width=.12pt] (-2.543764,-0.554239)--(-2.388443,-0.575200)--(-2.400000,-0.692820)--(-2.556072,-0.679508)--cycle;
\path[fill=blue!40!white,draw=blue!40!white,line width=.12pt] (2.556072,-0.679508)--(2.400000,-0.692820)--(2.388443,-0.575200)--(2.543764,-0.554239)--cycle;
\path[fill=blue!40!white,draw=blue!40!white,line width=.12pt] (-2.388443,-0.575200)--(-2.233123,-0.569536)--(-2.243928,-0.679508)--(-2.400000,-0.692820)--cycle;
\path[fill=blue!39!white,draw=blue!39!white,line width=.12pt] (2.400000,-0.692820)--(2.243928,-0.679508)--(2.233123,-0.569536)--(2.388443,-0.575200)--cycle;
\path[fill=blue!38!white,draw=blue!38!white,line width=.12pt] (-2.693116,-0.507458)--(-2.543764,-0.554239)--(-2.556072,-0.679508)--(-2.706147,-0.640083)--cycle;
\path[fill=blue!40!white,draw=blue!40!white,line width=.12pt] (2.706147,-0.640083)--(2.556072,-0.679508)--(2.543764,-0.554239)--(2.693116,-0.507458)--cycle;
\path[fill=blue!24!white,draw=blue!24!white,line width=.12pt] (-1.424631,0.515903)--(-1.464780,0.656601)--(-1.534468,0.582930)--(-1.492409,0.444251)--cycle;
\path[fill=blue!17!white,draw=blue!17!white,line width=.12pt] (1.492409,0.444251)--(1.534468,0.582930)--(1.464780,0.656601)--(1.424631,0.515903)--cycle;
\path[fill=blue!38!white,draw=blue!38!white,line width=.12pt] (-1.799069,-0.310969)--(-1.701490,-0.215687)--(-1.726471,-0.299889)--(-1.825482,-0.400001)--cycle;
\path[fill=blue!31!white,draw=blue!31!white,line width=.12pt] (1.825482,-0.400001)--(1.726471,-0.299889)--(1.701490,-0.215687)--(1.799069,-0.310969)--cycle;
\path[fill=blue!40!white,draw=blue!40!white,line width=.12pt] (-2.233123,-0.569536)--(-2.083771,-0.537466)--(-2.093853,-0.640083)--(-2.243928,-0.679508)--cycle;
\path[fill=blue!38!white,draw=blue!38!white,line width=.12pt] (2.243928,-0.679508)--(2.093853,-0.640083)--(2.083771,-0.537466)--(2.233123,-0.569536)--cycle;
\path[fill=blue!36!white,draw=blue!36!white,line width=.12pt] (-2.830759,-0.436656)--(-2.693116,-0.507458)--(-2.706147,-0.640083)--(-2.844456,-0.576059)--cycle;
\path[fill=blue!40!white,draw=blue!40!white,line width=.12pt] (2.844456,-0.576059)--(2.706147,-0.640083)--(2.693116,-0.507458)--(2.830759,-0.436656)--cycle;
\path[fill=blue!19!white,draw=blue!19!white,line width=.12pt] (-2.933190,0.829797)--(-3.003935,0.720748)--(-3.078787,0.571335)--(-3.006279,0.683903)--cycle;
\path[fill=blue!26!white,draw=blue!26!white,line width=.12pt] (3.006279,0.683903)--(3.078787,0.571335)--(3.003935,0.720748)--(2.933190,0.829797)--cycle;
\path[fill=blue!15!white,draw=blue!15!white,line width=.12pt] (-1.978968,1.208398)--(-2.116611,1.258496)--(-2.217311,1.152040)--(-2.073119,1.108865)--cycle;
\path[fill=blue!14!white,draw=blue!14!white,line width=.12pt] (2.073119,1.108865)--(2.217311,1.152040)--(2.116611,1.258496)--(1.978968,1.208398)--cycle;
\path[fill=blue!33!white,draw=blue!33!white,line width=.12pt] (-1.589379,-0.024064)--(-1.545814,0.099296)--(-1.584333,0.022409)--(-1.628983,-0.103118)--cycle;
\path[fill=blue!25!white,draw=blue!25!white,line width=.12pt] (1.628983,-0.103118)--(1.584333,0.022409)--(1.545814,0.099296)--(1.589379,-0.024064)--cycle;
\path[fill=blue!15!white,draw=blue!15!white,line width=.12pt] (-2.500154,1.157881)--(-2.627935,1.120322)--(-2.721975,0.988910)--(-2.589621,1.032859)--cycle;
\path[fill=blue!19!white,draw=blue!19!white,line width=.12pt] (2.589621,1.032859)--(2.721975,0.988910)--(2.627935,1.120322)--(2.500154,1.157881)--cycle;
\path[fill=blue!25!white,draw=blue!25!white,line width=.12pt] (-3.138513,0.312145)--(-3.123436,0.175483)--(-3.169293,0.020912)--(-3.184591,0.156827)--cycle;
\path[fill=blue!34!white,draw=blue!34!white,line width=.12pt] (3.184591,0.156827)--(3.169293,0.020912)--(3.123436,0.175483)--(3.138513,0.312145)--cycle;
\path[fill=blue!22!white,draw=blue!22!white,line width=.12pt] (-1.464780,0.656601)--(-1.529978,0.793806)--(-1.602769,0.716855)--(-1.534468,0.582930)--cycle;
\path[fill=blue!15!white,draw=blue!15!white,line width=.12pt] (1.534468,0.582930)--(1.602769,0.716855)--(1.529978,0.793806)--(1.464780,0.656601)--cycle;
\path[fill=blue!40!white,draw=blue!40!white,line width=.12pt] (-2.083771,-0.537466)--(-1.946127,-0.480221)--(-1.955544,-0.576059)--(-2.093853,-0.640083)--cycle;
\path[fill=blue!36!white,draw=blue!36!white,line width=.12pt] (2.093853,-0.640083)--(1.955544,-0.576059)--(1.946127,-0.480221)--(2.083771,-0.537466)--cycle;
\path[fill=blue!34!white,draw=blue!34!white,line width=.12pt] (-2.951405,-0.344554)--(-2.830759,-0.436656)--(-2.844456,-0.576059)--(-2.965685,-0.489898)--cycle;
\path[fill=blue!40!white,draw=blue!40!white,line width=.12pt] (2.965685,-0.489898)--(2.844456,-0.576059)--(2.830759,-0.436656)--(2.951405,-0.344554)--cycle;
\path[fill=blue!28!white,draw=blue!28!white,line width=.12pt] (-1.478207,0.306147)--(-1.492409,0.444251)--(-1.545814,0.369621)--(-1.531105,0.232228)--cycle;
\path[fill=blue!19!white,draw=blue!19!white,line width=.12pt] (1.531105,0.232228)--(1.545814,0.369621)--(1.492409,0.444251)--(1.478207,0.306147)--cycle;
\path[fill=blue!15!white,draw=blue!15!white,line width=.12pt] (-1.846614,1.133600)--(-1.978968,1.208398)--(-2.073119,1.108865)--(-1.934468,1.040724)--cycle;
\path[fill=blue!13!white,draw=blue!13!white,line width=.12pt] (1.934468,1.040724)--(2.073119,1.108865)--(1.978968,1.208398)--(1.846614,1.133600)--cycle;
\path[fill=blue!20!white,draw=blue!20!white,line width=.12pt] (-1.529978,0.793806)--(-1.617721,0.922243)--(-1.694686,0.840879)--(-1.602769,0.716855)--cycle;
\path[fill=blue!14!white,draw=blue!14!white,line width=.12pt] (1.602769,0.716855)--(1.694686,0.840879)--(1.617721,0.922243)--(1.529978,0.793806)--cycle;
\path[fill=blue!36!white,draw=blue!36!white,line width=.12pt] (-1.701490,-0.215687)--(-1.628983,-0.103118)--(-1.652899,-0.183733)--(-1.726471,-0.299889)--cycle;
\path[fill=blue!28!white,draw=blue!28!white,line width=.12pt] (1.726471,-0.299889)--(1.652899,-0.183733)--(1.628983,-0.103118)--(1.701490,-0.215687)--cycle;
\path[fill=blue!16!white,draw=blue!16!white,line width=.12pt] (-1.724636,1.036978)--(-1.846614,1.133600)--(-1.934468,1.040724)--(-1.806687,0.950236)--cycle;
\path[fill=blue!13!white,draw=blue!13!white,line width=.12pt] (1.806687,0.950236)--(1.934468,1.040724)--(1.846614,1.133600)--(1.724636,1.036978)--cycle;
\path[fill=blue!18!white,draw=blue!18!white,line width=.12pt] (-1.617721,0.922243)--(-1.724636,1.036978)--(-1.806687,0.950236)--(-1.694686,0.840879)--cycle;
\path[fill=blue!13!white,draw=blue!13!white,line width=.12pt] (1.694686,0.840879)--(1.806687,0.950236)--(1.724636,1.036978)--(1.617721,0.922243)--cycle;
\path[fill=blue!39!white,draw=blue!39!white,line width=.12pt] (-1.946127,-0.480221)--(-1.825482,-0.400001)--(-1.834315,-0.489898)--(-1.955544,-0.576059)--cycle;
\path[fill=blue!34!white,draw=blue!34!white,line width=.12pt] (1.955544,-0.576059)--(1.834315,-0.489898)--(1.825482,-0.400001)--(1.946127,-0.480221)--cycle;
\path[fill=blue!14!white,draw=blue!14!white,line width=.12pt] (-2.361503,1.168591)--(-2.500154,1.157881)--(-2.589621,1.032859)--(-2.446009,1.050502)--cycle;
\path[fill=blue!17!white,draw=blue!17!white,line width=.12pt] (2.446009,1.050502)--(2.589621,1.032859)--(2.500154,1.157881)--(2.361503,1.168591)--cycle;
\path[fill=blue!17!white,draw=blue!17!white,line width=.12pt] (-2.837984,0.920344)--(-2.933190,0.829797)--(-3.006279,0.683903)--(-2.908701,0.779186)--cycle;
\path[fill=blue!24!white,draw=blue!24!white,line width=.12pt] (2.908701,0.779186)--(3.006279,0.683903)--(2.933190,0.829797)--(2.837984,0.920344)--cycle;
\path[fill=blue!32!white,draw=blue!32!white,line width=.12pt] (-3.050416,-0.234690)--(-2.951405,-0.344554)--(-2.965685,-0.489898)--(-3.065176,-0.384910)--cycle;
\path[fill=blue!39!white,draw=blue!39!white,line width=.12pt] (3.065176,-0.384910)--(2.965685,-0.489898)--(2.951405,-0.344554)--(3.050416,-0.234690)--cycle;
\path[fill=blue!31!white,draw=blue!31!white,line width=.12pt] (-1.545814,0.099296)--(-1.531105,0.232228)--(-1.569256,0.156072)--(-1.584333,0.022409)--cycle;
\path[fill=blue!22!white,draw=blue!22!white,line width=.12pt] (1.584333,0.022409)--(1.569256,0.156072)--(1.531105,0.232228)--(1.545814,0.099296)--cycle;
\path[fill=blue!22!white,draw=blue!22!white,line width=.12pt] (-3.123436,0.445808)--(-3.138513,0.312145)--(-3.184591,0.156827)--(-3.169293,0.291237)--cycle;
\path[fill=blue!31!white,draw=blue!31!white,line width=.12pt] (3.169293,0.291237)--(3.184591,0.156827)--(3.138513,0.312145)--(3.123436,0.445808)--cycle;
\path[fill=blue!25!white,draw=blue!25!white,line width=.12pt] (-1.492409,0.444251)--(-1.534468,0.582930)--(-1.589379,0.506197)--(-1.545814,0.369621)--cycle;
\path[fill=blue!17!white,draw=blue!17!white,line width=.12pt] (1.545814,0.369621)--(1.589379,0.506197)--(1.534468,0.582930)--(1.492409,0.444251)--cycle;
\path[fill=blue!38!white,draw=blue!38!white,line width=.12pt] (-1.825482,-0.400001)--(-1.726471,-0.299889)--(-1.734824,-0.384910)--(-1.834315,-0.489898)--cycle;
\path[fill=blue!31!white,draw=blue!31!white,line width=.12pt] (1.834315,-0.489898)--(1.734824,-0.384910)--(1.726471,-0.299889)--(1.825482,-0.400001)--cycle;
\path[fill=blue!14!white,draw=blue!14!white,line width=.12pt] (-2.217311,1.152040)--(-2.361503,1.168591)--(-2.446009,1.050502)--(-2.296657,1.041162)--cycle;
\path[fill=blue!15!white,draw=blue!15!white,line width=.12pt] (2.296657,1.041162)--(2.446009,1.050502)--(2.361503,1.168591)--(2.217311,1.152040)--cycle;
\path[fill=blue!34!white,draw=blue!34!white,line width=.12pt] (-1.628983,-0.103118)--(-1.584333,0.022409)--(-1.607593,-0.055995)--(-1.652899,-0.183733)--cycle;
\path[fill=blue!25!white,draw=blue!25!white,line width=.12pt] (1.652899,-0.183733)--(1.607593,-0.055995)--(1.584333,0.022409)--(1.628983,-0.103118)--cycle;
\path[fill=blue!29!white,draw=blue!29!white,line width=.12pt] (-3.123988,-0.111288)--(-3.050416,-0.234690)--(-3.065176,-0.384910)--(-3.139104,-0.265131)--cycle;
\path[fill=blue!37!white,draw=blue!37!white,line width=.12pt] (3.139104,-0.265131)--(3.065176,-0.384910)--(3.050416,-0.234690)--(3.123988,-0.111288)--cycle;
\path[fill=blue!23!white,draw=blue!23!white,line width=.12pt] (-1.534468,0.582930)--(-1.602769,0.716855)--(-1.660123,0.636707)--(-1.589379,0.506197)--cycle;
\path[fill=blue!15!white,draw=blue!15!white,line width=.12pt] (1.589379,0.506197)--(1.660123,0.636707)--(1.602769,0.716855)--(1.534468,0.582930)--cycle;
\path[fill=blue!16!white,draw=blue!16!white,line width=.12pt] (-2.721975,0.988910)--(-2.837984,0.920344)--(-2.908701,0.779186)--(-2.789801,0.853522)--cycle;
\path[fill=blue!21!white,draw=blue!21!white,line width=.12pt] (2.789801,0.853522)--(2.908701,0.779186)--(2.837984,0.920344)--(2.721975,0.988910)--cycle;
\path[fill=blue!28!white,draw=blue!28!white,line width=.12pt] (-1.531105,0.232228)--(-1.545814,0.369621)--(-1.584333,0.292734)--(-1.569256,0.156072)--cycle;
\path[fill=blue!20!white,draw=blue!20!white,line width=.12pt] (1.569256,0.156072)--(1.584333,0.292734)--(1.545814,0.369621)--(1.531105,0.232228)--cycle;
\path[fill=blue!40!white,draw=blue!40!white,line width=.12pt] (-2.400000,-0.692820)--(-2.243928,-0.679508)--(-2.233123,-0.789480)--(-2.388443,-0.810441)--cycle;
\path[fill=blue!39!white,draw=blue!39!white,line width=.12pt] (2.388443,-0.810441)--(2.233123,-0.789480)--(2.243928,-0.679508)--(2.400000,-0.692820)--cycle;
\path[fill=blue!39!white,draw=blue!39!white,line width=.12pt] (-2.556072,-0.679508)--(-2.400000,-0.692820)--(-2.388443,-0.810441)--(-2.543764,-0.804777)--cycle;
\path[fill=blue!40!white,draw=blue!40!white,line width=.12pt] (2.543764,-0.804777)--(2.388443,-0.810441)--(2.400000,-0.692820)--(2.556072,-0.679508)--cycle;
\path[fill=blue!40!white,draw=blue!40!white,line width=.12pt] (-2.243928,-0.679508)--(-2.093853,-0.640083)--(-2.083771,-0.742699)--(-2.233123,-0.789480)--cycle;
\path[fill=blue!38!white,draw=blue!38!white,line width=.12pt] (2.233123,-0.789480)--(2.083771,-0.742699)--(2.093853,-0.640083)--(2.243928,-0.679508)--cycle;
\path[fill=blue!20!white,draw=blue!20!white,line width=.12pt] (-3.078787,0.571335)--(-3.123436,0.445808)--(-3.169293,0.291237)--(-3.123988,0.418974)--cycle;
\path[fill=blue!29!white,draw=blue!29!white,line width=.12pt] (3.123988,0.418974)--(3.169293,0.291237)--(3.123436,0.445808)--(3.078787,0.571335)--cycle;
\path[fill=blue!15!white,draw=blue!15!white,line width=.12pt] (-2.073119,1.108865)--(-2.217311,1.152040)--(-2.296657,1.041162)--(-2.147305,1.005197)--cycle;
\path[fill=blue!14!white,draw=blue!14!white,line width=.12pt] (2.147305,1.005197)--(2.296657,1.041162)--(2.217311,1.152040)--(2.073119,1.108865)--cycle;
\path[fill=blue!38!white,draw=blue!38!white,line width=.12pt] (-2.706147,-0.640083)--(-2.556072,-0.679508)--(-2.543764,-0.804777)--(-2.693116,-0.772707)--cycle;
\path[fill=blue!40!white,draw=blue!40!white,line width=.12pt] (2.693116,-0.772707)--(2.543764,-0.804777)--(2.556072,-0.679508)--(2.706147,-0.640083)--cycle;
\path[fill=blue!21!white,draw=blue!21!white,line width=.12pt] (-1.602769,0.716855)--(-1.694686,0.840879)--(-1.755330,0.756135)--(-1.660123,0.636707)--cycle;
\path[fill=blue!14!white,draw=blue!14!white,line width=.12pt] (1.660123,0.636707)--(1.755330,0.756135)--(1.694686,0.840879)--(1.602769,0.716855)--cycle;
\path[fill=blue!40!white,draw=blue!40!white,line width=.12pt] (-2.093853,-0.640083)--(-1.955544,-0.576059)--(-1.946127,-0.671897)--(-2.083771,-0.742699)--cycle;
\path[fill=blue!36!white,draw=blue!36!white,line width=.12pt] (2.083771,-0.742699)--(1.946127,-0.671897)--(1.955544,-0.576059)--(2.093853,-0.640083)--cycle;
\path[fill=blue!37!white,draw=blue!37!white,line width=.12pt] (-1.726471,-0.299889)--(-1.652899,-0.183733)--(-1.660896,-0.265131)--(-1.734824,-0.384910)--cycle;
\path[fill=blue!29!white,draw=blue!29!white,line width=.12pt] (1.734824,-0.384910)--(1.660896,-0.265131)--(1.652899,-0.183733)--(1.726471,-0.299889)--cycle;
\path[fill=blue!16!white,draw=blue!16!white,line width=.12pt] (-1.934468,1.040724)--(-2.073119,1.108865)--(-2.147305,1.005197)--(-2.003693,0.943989)--cycle;
\path[fill=blue!13!white,draw=blue!13!white,line width=.12pt] (2.003693,0.943989)--(2.147305,1.005197)--(2.073119,1.108865)--(1.934468,1.040724)--cycle;
\path[fill=blue!27!white,draw=blue!27!white,line width=.12pt] (-3.169293,0.020912)--(-3.123988,-0.111288)--(-3.139104,-0.265131)--(-3.184628,-0.135163)--cycle;
\path[fill=blue!35!white,draw=blue!35!white,line width=.12pt] (3.184628,-0.135163)--(3.139104,-0.265131)--(3.123988,-0.111288)--(3.169293,0.020912)--cycle;
\path[fill=blue!18!white,draw=blue!18!white,line width=.12pt] (-1.694686,0.840879)--(-1.806687,0.950236)--(-1.871339,0.859891)--(-1.755330,0.756135)--cycle;
\path[fill=blue!13!white,draw=blue!13!white,line width=.12pt] (1.755330,0.756135)--(1.871339,0.859891)--(1.806687,0.950236)--(1.694686,0.840879)--cycle;
\path[fill=blue!32!white,draw=blue!32!white,line width=.12pt] (-1.584333,0.022409)--(-1.569256,0.156072)--(-1.592296,0.078414)--(-1.607593,-0.055995)--cycle;
\path[fill=blue!23!white,draw=blue!23!white,line width=.12pt] (1.607593,-0.055995)--(1.592296,0.078414)--(1.569256,0.156072)--(1.584333,0.022409)--cycle;
\path[fill=blue!17!white,draw=blue!17!white,line width=.12pt] (-1.806687,0.950236)--(-1.934468,1.040724)--(-2.003693,0.943989)--(-1.871339,0.859891)--cycle;
\path[fill=blue!13!white,draw=blue!13!white,line width=.12pt] (1.871339,0.859891)--(2.003693,0.943989)--(1.934468,1.040724)--(1.806687,0.950236)--cycle;
\path[fill=blue!36!white,draw=blue!36!white,line width=.12pt] (-2.844456,-0.576059)--(-2.706147,-0.640083)--(-2.693116,-0.772707)--(-2.830759,-0.715462)--cycle;
\path[fill=blue!40!white,draw=blue!40!white,line width=.12pt] (2.830759,-0.715462)--(2.693116,-0.772707)--(2.706147,-0.640083)--(2.844456,-0.576059)--cycle;
\path[fill=blue!15!white,draw=blue!15!white,line width=.12pt] (-2.589621,1.032859)--(-2.721975,0.988910)--(-2.789801,0.853522)--(-2.654149,0.904054)--cycle;
\path[fill=blue!19!white,draw=blue!19!white,line width=.12pt] (2.654149,0.904054)--(2.789801,0.853522)--(2.721975,0.988910)--(2.589621,1.032859)--cycle;
\path[fill=blue!26!white,draw=blue!26!white,line width=.12pt] (-1.545814,0.369621)--(-1.589379,0.506197)--(-1.628983,0.427143)--(-1.584333,0.292734)--cycle;
\path[fill=blue!18!white,draw=blue!18!white,line width=.12pt] (1.584333,0.292734)--(1.628983,0.427143)--(1.589379,0.506197)--(1.545814,0.369621)--cycle;
\path[fill=blue!40!white,draw=blue!40!white,line width=.12pt] (-1.955544,-0.576059)--(-1.834315,-0.489898)--(-1.825482,-0.579795)--(-1.946127,-0.671897)--cycle;
\path[fill=blue!34!white,draw=blue!34!white,line width=.12pt] (1.946127,-0.671897)--(1.825482,-0.579795)--(1.834315,-0.489898)--(1.955544,-0.576059)--cycle;
\path[fill=blue!18!white,draw=blue!18!white,line width=.12pt] (-3.006279,0.683903)--(-3.078787,0.571335)--(-3.123988,0.418974)--(-3.050416,0.535130)--cycle;
\path[fill=blue!26!white,draw=blue!26!white,line width=.12pt] (3.050416,0.535130)--(3.123988,0.418974)--(3.078787,0.571335)--(3.006279,0.683903)--cycle;
\path[fill=blue!35!white,draw=blue!35!white,line width=.12pt] (-1.652899,-0.183733)--(-1.607593,-0.055995)--(-1.615372,-0.135163)--(-1.660896,-0.265131)--cycle;
\path[fill=blue!26!white,draw=blue!26!white,line width=.12pt] (1.660896,-0.265131)--(1.615372,-0.135163)--(1.607593,-0.055995)--(1.652899,-0.183733)--cycle;
\path[fill=blue!34!white,draw=blue!34!white,line width=.12pt] (-2.965685,-0.489898)--(-2.844456,-0.576059)--(-2.830759,-0.715462)--(-2.951405,-0.635242)--cycle;
\path[fill=blue!39!white,draw=blue!39!white,line width=.12pt] (2.951405,-0.635242)--(2.830759,-0.715462)--(2.844456,-0.576059)--(2.965685,-0.489898)--cycle;
\path[fill=blue!39!white,draw=blue!39!white,line width=.12pt] (-1.834315,-0.489898)--(-1.734824,-0.384910)--(-1.726471,-0.469932)--(-1.825482,-0.579795)--cycle;
\path[fill=blue!32!white,draw=blue!32!white,line width=.12pt] (1.825482,-0.579795)--(1.726471,-0.469932)--(1.734824,-0.384910)--(1.834315,-0.489898)--cycle;
\path[fill=blue!24!white,draw=blue!24!white,line width=.12pt] (-3.184591,0.156827)--(-3.169293,0.020912)--(-3.184628,-0.135163)--(-3.200000,0.000000)--cycle;
\path[fill=blue!33!white,draw=blue!33!white,line width=.12pt] (3.200000,0.000000)--(3.184628,-0.135163)--(3.169293,0.020912)--(3.184591,0.156827)--cycle;
\path[fill=blue!14!white,draw=blue!14!white,line width=.12pt] (-2.446009,1.050502)--(-2.589621,1.032859)--(-2.654149,0.904054)--(-2.506958,0.928840)--cycle;
\path[fill=blue!17!white,draw=blue!17!white,line width=.12pt] (2.506958,0.928840)--(2.654149,0.904054)--(2.589621,1.032859)--(2.446009,1.050502)--cycle;
\path[fill=blue!23!white,draw=blue!23!white,line width=.12pt] (-1.589379,0.506197)--(-1.660123,0.636707)--(-1.701490,0.554134)--(-1.628983,0.427143)--cycle;
\path[fill=blue!16!white,draw=blue!16!white,line width=.12pt] (1.628983,0.427143)--(1.701490,0.554134)--(1.660123,0.636707)--(1.589379,0.506197)--cycle;
\path[fill=blue!29!white,draw=blue!29!white,line width=.12pt] (-1.569256,0.156072)--(-1.584333,0.292734)--(-1.607593,0.214330)--(-1.592296,0.078414)--cycle;
\path[fill=blue!20!white,draw=blue!20!white,line width=.12pt] (1.592296,0.078414)--(1.607593,0.214330)--(1.584333,0.292734)--(1.569256,0.156072)--cycle;
\path[fill=blue!16!white,draw=blue!16!white,line width=.12pt] (-2.908701,0.779186)--(-3.006279,0.683903)--(-3.050416,0.535130)--(-2.951405,0.635242)--cycle;
\path[fill=blue!23!white,draw=blue!23!white,line width=.12pt] (2.951405,0.635242)--(3.050416,0.535130)--(3.006279,0.683903)--(2.908701,0.779186)--cycle;
\path[fill=blue!14!white,draw=blue!14!white,line width=.12pt] (-2.296657,1.041162)--(-2.446009,1.050502)--(-2.506958,0.928840)--(-2.353885,0.926929)--cycle;
\path[fill=blue!15!white,draw=blue!15!white,line width=.12pt] (2.353885,0.926929)--(2.506958,0.928840)--(2.446009,1.050502)--(2.296657,1.041162)--cycle;
\path[fill=blue!31!white,draw=blue!31!white,line width=.12pt] (-3.065176,-0.384910)--(-2.965685,-0.489898)--(-2.951405,-0.635242)--(-3.050416,-0.535130)--cycle;
\path[fill=blue!38!white,draw=blue!38!white,line width=.12pt] (3.050416,-0.535130)--(2.951405,-0.635242)--(2.965685,-0.489898)--(3.065176,-0.384910)--cycle;
\path[fill=blue!21!white,draw=blue!21!white,line width=.12pt] (-1.660123,0.636707)--(-1.755330,0.756135)--(-1.799069,0.668826)--(-1.701490,0.554134)--cycle;
\path[fill=blue!14!white,draw=blue!14!white,line width=.12pt] (1.701490,0.554134)--(1.799069,0.668826)--(1.755330,0.756135)--(1.660123,0.636707)--cycle;
\path[fill=blue!37!white,draw=blue!37!white,line width=.12pt] (-1.734824,-0.384910)--(-1.660896,-0.265131)--(-1.652899,-0.346529)--(-1.726471,-0.469932)--cycle;
\path[fill=blue!29!white,draw=blue!29!white,line width=.12pt] (1.726471,-0.469932)--(1.652899,-0.346529)--(1.660896,-0.265131)--(1.734824,-0.384910)--cycle;
\path[fill=blue!32!white,draw=blue!32!white,line width=.12pt] (-1.607593,-0.055995)--(-1.592296,0.078414)--(-1.600000,-0.000000)--(-1.615372,-0.135163)--cycle;
\path[fill=blue!23!white,draw=blue!23!white,line width=.12pt] (1.615372,-0.135163)--(1.600000,-0.000000)--(1.592296,0.078414)--(1.607593,-0.055995)--cycle;
\path[fill=blue!41!white,draw=blue!41!white,line width=.12pt] (-2.233123,-0.789480)--(-2.083771,-0.742699)--(-2.053620,-0.844328)--(-2.200811,-0.898392)--cycle;
\path[fill=blue!38!white,draw=blue!38!white,line width=.12pt] (2.200811,-0.898392)--(2.053620,-0.844328)--(2.083771,-0.742699)--(2.233123,-0.789480)--cycle;
\path[fill=blue!40!white,draw=blue!40!white,line width=.12pt] (-2.388443,-0.810441)--(-2.233123,-0.789480)--(-2.200811,-0.898392)--(-2.353885,-0.926929)--cycle;
\path[fill=blue!39!white,draw=blue!39!white,line width=.12pt] (2.353885,-0.926929)--(2.200811,-0.898392)--(2.233123,-0.789480)--(2.388443,-0.810441)--cycle;
\path[fill=blue!22!white,draw=blue!22!white,line width=.12pt] (-3.169293,0.291237)--(-3.184591,0.156827)--(-3.200000,0.000000)--(-3.184628,0.135163)--cycle;
\path[fill=blue!31!white,draw=blue!31!white,line width=.12pt] (3.184628,0.135163)--(3.200000,0.000000)--(3.184591,0.156827)--(3.169293,0.291237)--cycle;
\path[fill=blue!41!white,draw=blue!41!white,line width=.12pt] (-2.083771,-0.742699)--(-1.946127,-0.671897)--(-1.917969,-0.766813)--(-2.053620,-0.844328)--cycle;
\path[fill=blue!36!white,draw=blue!36!white,line width=.12pt] (2.053620,-0.844328)--(1.917969,-0.766813)--(1.946127,-0.671897)--(2.083771,-0.742699)--cycle;
\path[fill=blue!27!white,draw=blue!27!white,line width=.12pt] (-1.584333,0.292734)--(-1.628983,0.427143)--(-1.652899,0.346529)--(-1.607593,0.214330)--cycle;
\path[fill=blue!18!white,draw=blue!18!white,line width=.12pt] (1.607593,0.214330)--(1.652899,0.346529)--(1.628983,0.427143)--(1.584333,0.292734)--cycle;
\path[fill=blue!15!white,draw=blue!15!white,line width=.12pt] (-2.147305,1.005197)--(-2.296657,1.041162)--(-2.353885,0.926929)--(-2.200811,0.898392)--cycle;
\path[fill=blue!14!white,draw=blue!14!white,line width=.12pt] (2.200811,0.898392)--(2.353885,0.926929)--(2.296657,1.041162)--(2.147305,1.005197)--cycle;
\path[fill=blue!19!white,draw=blue!19!white,line width=.12pt] (-1.755330,0.756135)--(-1.871339,0.859891)--(-1.917969,0.766813)--(-1.799069,0.668826)--cycle;
\path[fill=blue!13!white,draw=blue!13!white,line width=.12pt] (1.799069,0.668826)--(1.917969,0.766813)--(1.871339,0.859891)--(1.755330,0.756135)--cycle;
\path[fill=blue!39!white,draw=blue!39!white,line width=.12pt] (-2.543764,-0.804777)--(-2.388443,-0.810441)--(-2.353885,-0.926929)--(-2.506958,-0.928840)--cycle;
\path[fill=blue!40!white,draw=blue!40!white,line width=.12pt] (2.506958,-0.928840)--(2.353885,-0.926929)--(2.388443,-0.810441)--(2.543764,-0.804777)--cycle;
\path[fill=blue!16!white,draw=blue!16!white,line width=.12pt] (-2.003693,0.943989)--(-2.147305,1.005197)--(-2.200811,0.898392)--(-2.053620,0.844328)--cycle;
\path[fill=blue!13!white,draw=blue!13!white,line width=.12pt] (2.053620,0.844328)--(2.200811,0.898392)--(2.147305,1.005197)--(2.003693,0.943989)--cycle;
\path[fill=blue!17!white,draw=blue!17!white,line width=.12pt] (-1.871339,0.859891)--(-2.003693,0.943989)--(-2.053620,0.844328)--(-1.917969,0.766813)--cycle;
\path[fill=blue!13!white,draw=blue!13!white,line width=.12pt] (1.917969,0.766813)--(2.053620,0.844328)--(2.003693,0.943989)--(1.871339,0.859891)--cycle;
\path[fill=blue!40!white,draw=blue!40!white,line width=.12pt] (-1.946127,-0.671897)--(-1.825482,-0.579795)--(-1.799069,-0.668826)--(-1.917969,-0.766813)--cycle;
\path[fill=blue!35!white,draw=blue!35!white,line width=.12pt] (1.917969,-0.766813)--(1.799069,-0.668826)--(1.825482,-0.579795)--(1.946127,-0.671897)--cycle;
\path[fill=blue!15!white,draw=blue!15!white,line width=.12pt] (-2.789801,0.853522)--(-2.908701,0.779186)--(-2.951405,0.635242)--(-2.830759,0.715462)--cycle;
\path[fill=blue!21!white,draw=blue!21!white,line width=.12pt] (2.830759,0.715462)--(2.951405,0.635242)--(2.908701,0.779186)--(2.789801,0.853522)--cycle;
\path[fill=blue!29!white,draw=blue!29!white,line width=.12pt] (-3.139104,-0.265131)--(-3.065176,-0.384910)--(-3.050416,-0.535130)--(-3.123988,-0.418974)--cycle;
\path[fill=blue!37!white,draw=blue!37!white,line width=.12pt] (3.123988,-0.418974)--(3.050416,-0.535130)--(3.065176,-0.384910)--(3.139104,-0.265131)--cycle;
\path[fill=blue!37!white,draw=blue!37!white,line width=.12pt] (-2.693116,-0.772707)--(-2.543764,-0.804777)--(-2.506958,-0.928840)--(-2.654149,-0.904054)--cycle;
\path[fill=blue!40!white,draw=blue!40!white,line width=.12pt] (2.654149,-0.904054)--(2.506958,-0.928840)--(2.543764,-0.804777)--(2.693116,-0.772707)--cycle;
\path[fill=blue!35!white,draw=blue!35!white,line width=.12pt] (-1.660896,-0.265131)--(-1.615372,-0.135163)--(-1.607593,-0.214330)--(-1.652899,-0.346529)--cycle;
\path[fill=blue!27!white,draw=blue!27!white,line width=.12pt] (1.652899,-0.346529)--(1.607593,-0.214330)--(1.615372,-0.135163)--(1.660896,-0.265131)--cycle;
\path[fill=blue!30!white,draw=blue!30!white,line width=.12pt] (-1.592296,0.078414)--(-1.607593,0.214330)--(-1.615372,0.135163)--(-1.600000,0.000000)--cycle;
\path[fill=blue!21!white,draw=blue!21!white,line width=.12pt] (1.600000,0.000000)--(1.615372,0.135163)--(1.607593,0.214330)--(1.592296,0.078414)--cycle;
\path[fill=blue!24!white,draw=blue!24!white,line width=.12pt] (-1.628983,0.427143)--(-1.701490,0.554134)--(-1.726471,0.469932)--(-1.652899,0.346529)--cycle;
\path[fill=blue!16!white,draw=blue!16!white,line width=.12pt] (1.652899,0.346529)--(1.726471,0.469932)--(1.701490,0.554134)--(1.628983,0.427143)--cycle;
\path[fill=blue!39!white,draw=blue!39!white,line width=.12pt] (-1.825482,-0.579795)--(-1.726471,-0.469932)--(-1.701490,-0.554134)--(-1.799069,-0.668826)--cycle;
\path[fill=blue!32!white,draw=blue!32!white,line width=.12pt] (1.799069,-0.668826)--(1.701490,-0.554134)--(1.726471,-0.469932)--(1.825482,-0.579795)--cycle;
\path[fill=blue!19!white,draw=blue!19!white,line width=.12pt] (-3.123988,0.418974)--(-3.169293,0.291237)--(-3.184628,0.135163)--(-3.139104,0.265131)--cycle;
\path[fill=blue!28!white,draw=blue!28!white,line width=.12pt] (3.139104,0.265131)--(3.184628,0.135163)--(3.169293,0.291237)--(3.123988,0.418974)--cycle;
\path[fill=blue!14!white,draw=blue!14!white,line width=.12pt] (-2.654149,0.904054)--(-2.789801,0.853522)--(-2.830759,0.715462)--(-2.693116,0.772707)--cycle;
\path[fill=blue!18!white,draw=blue!18!white,line width=.12pt] (2.693116,0.772707)--(2.830759,0.715462)--(2.789801,0.853522)--(2.654149,0.904054)--cycle;
\path[fill=blue!36!white,draw=blue!36!white,line width=.12pt] (-2.830759,-0.715462)--(-2.693116,-0.772707)--(-2.654149,-0.904054)--(-2.789801,-0.853522)--cycle;
\path[fill=blue!40!white,draw=blue!40!white,line width=.12pt] (2.789801,-0.853522)--(2.654149,-0.904054)--(2.693116,-0.772707)--(2.830759,-0.715462)--cycle;
\path[fill=blue!26!white,draw=blue!26!white,line width=.12pt] (-3.184628,-0.135163)--(-3.139104,-0.265131)--(-3.123988,-0.418974)--(-3.169293,-0.291237)--cycle;
\path[fill=blue!35!white,draw=blue!35!white,line width=.12pt] (3.169293,-0.291237)--(3.123988,-0.418974)--(3.139104,-0.265131)--(3.184628,-0.135163)--cycle;
\path[fill=blue!22!white,draw=blue!22!white,line width=.12pt] (-1.701490,0.554134)--(-1.799069,0.668826)--(-1.825482,0.579795)--(-1.726471,0.469932)--cycle;
\path[fill=blue!15!white,draw=blue!15!white,line width=.12pt] (1.726471,0.469932)--(1.825482,0.579795)--(1.799069,0.668826)--(1.701490,0.554134)--cycle;
\path[fill=blue!38!white,draw=blue!38!white,line width=.12pt] (-1.726471,-0.469932)--(-1.652899,-0.346529)--(-1.628983,-0.427143)--(-1.701490,-0.554134)--cycle;
\path[fill=blue!30!white,draw=blue!30!white,line width=.12pt] (1.701490,-0.554134)--(1.628983,-0.427143)--(1.652899,-0.346529)--(1.726471,-0.469932)--cycle;
\path[fill=blue!33!white,draw=blue!33!white,line width=.12pt] (-1.615372,-0.135163)--(-1.600000,-0.000000)--(-1.592296,-0.078414)--(-1.607593,-0.214330)--cycle;
\path[fill=blue!24!white,draw=blue!24!white,line width=.12pt] (1.607593,-0.214330)--(1.592296,-0.078414)--(1.600000,-0.000000)--(1.615372,-0.135163)--cycle;
\path[fill=blue!27!white,draw=blue!27!white,line width=.12pt] (-1.607593,0.214330)--(-1.652899,0.346529)--(-1.660896,0.265131)--(-1.615372,0.135163)--cycle;
\path[fill=blue!19!white,draw=blue!19!white,line width=.12pt] (1.615372,0.135163)--(1.660896,0.265131)--(1.652899,0.346529)--(1.607593,0.214330)--cycle;
\path[fill=blue!14!white,draw=blue!14!white,line width=.12pt] (-2.506958,0.928840)--(-2.654149,0.904054)--(-2.693116,0.772707)--(-2.543764,0.804777)--cycle;
\path[fill=blue!17!white,draw=blue!17!white,line width=.12pt] (2.543764,0.804777)--(2.693116,0.772707)--(2.654149,0.904054)--(2.506958,0.928840)--cycle;
\path[fill=blue!17!white,draw=blue!17!white,line width=.12pt] (-3.050416,0.535130)--(-3.123988,0.418974)--(-3.139104,0.265131)--(-3.065176,0.384910)--cycle;
\path[fill=blue!25!white,draw=blue!25!white,line width=.12pt] (3.065176,0.384910)--(3.139104,0.265131)--(3.123988,0.418974)--(3.050416,0.535130)--cycle;
\path[fill=blue!33!white,draw=blue!33!white,line width=.12pt] (-2.951405,-0.635242)--(-2.830759,-0.715462)--(-2.789801,-0.853522)--(-2.908701,-0.779186)--cycle;
\path[fill=blue!39!white,draw=blue!39!white,line width=.12pt] (2.908701,-0.779186)--(2.789801,-0.853522)--(2.830759,-0.715462)--(2.951405,-0.635242)--cycle;
\path[fill=blue!19!white,draw=blue!19!white,line width=.12pt] (-1.799069,0.668826)--(-1.917969,0.766813)--(-1.946127,0.671897)--(-1.825482,0.579795)--cycle;
\path[fill=blue!14!white,draw=blue!14!white,line width=.12pt] (1.825482,0.579795)--(1.946127,0.671897)--(1.917969,0.766813)--(1.799069,0.668826)--cycle;
\path[fill=blue!41!white,draw=blue!41!white,line width=.12pt] (-2.053620,-0.844328)--(-1.917969,-0.766813)--(-1.871339,-0.859891)--(-2.003693,-0.943989)--cycle;
\path[fill=blue!37!white,draw=blue!37!white,line width=.12pt] (2.003693,-0.943989)--(1.871339,-0.859891)--(1.917969,-0.766813)--(2.053620,-0.844328)--cycle;
\path[fill=blue!41!white,draw=blue!41!white,line width=.12pt] (-2.200811,-0.898392)--(-2.053620,-0.844328)--(-2.003693,-0.943989)--(-2.147305,-1.005197)--cycle;
\path[fill=blue!38!white,draw=blue!38!white,line width=.12pt] (2.147305,-1.005197)--(2.003693,-0.943989)--(2.053620,-0.844328)--(2.200811,-0.898392)--cycle;
\path[fill=blue!14!white,draw=blue!14!white,line width=.12pt] (-2.353885,0.926929)--(-2.506958,0.928840)--(-2.543764,0.804777)--(-2.388443,0.810441)--cycle;
\path[fill=blue!15!white,draw=blue!15!white,line width=.12pt] (2.388443,0.810441)--(2.543764,0.804777)--(2.506958,0.928840)--(2.353885,0.926929)--cycle;
\path[fill=blue!41!white,draw=blue!41!white,line width=.12pt] (-1.917969,-0.766813)--(-1.799069,-0.668826)--(-1.755330,-0.756135)--(-1.871339,-0.859891)--cycle;
\path[fill=blue!35!white,draw=blue!35!white,line width=.12pt] (1.871339,-0.859891)--(1.755330,-0.756135)--(1.799069,-0.668826)--(1.917969,-0.766813)--cycle;
\path[fill=blue!40!white,draw=blue!40!white,line width=.12pt] (-2.353885,-0.926929)--(-2.200811,-0.898392)--(-2.147305,-1.005197)--(-2.296657,-1.041162)--cycle;
\path[fill=blue!39!white,draw=blue!39!white,line width=.12pt] (2.296657,-1.041162)--(2.147305,-1.005197)--(2.200811,-0.898392)--(2.353885,-0.926929)--cycle;
\path[fill=blue!36!white,draw=blue!36!white,line width=.12pt] (-1.652899,-0.346529)--(-1.607593,-0.214330)--(-1.584333,-0.292734)--(-1.628983,-0.427143)--cycle;
\path[fill=blue!27!white,draw=blue!27!white,line width=.12pt] (1.628983,-0.427143)--(1.584333,-0.292734)--(1.607593,-0.214330)--(1.652899,-0.346529)--cycle;
\path[fill=blue!18!white,draw=blue!18!white,line width=.12pt] (-1.917969,0.766813)--(-2.053620,0.844328)--(-2.083771,0.742699)--(-1.946127,0.671897)--cycle;
\path[fill=blue!13!white,draw=blue!13!white,line width=.12pt] (1.946127,0.671897)--(2.083771,0.742699)--(2.053620,0.844328)--(1.917969,0.766813)--cycle;
\path[fill=blue!23!white,draw=blue!23!white,line width=.12pt] (-3.200000,0.000000)--(-3.184628,-0.135163)--(-3.169293,-0.291237)--(-3.184591,-0.156827)--cycle;
\path[fill=blue!32!white,draw=blue!32!white,line width=.12pt] (3.184591,-0.156827)--(3.169293,-0.291237)--(3.184628,-0.135163)--(3.200000,-0.000000)--cycle;
\path[fill=blue!15!white,draw=blue!15!white,line width=.12pt] (-2.200811,0.898392)--(-2.353885,0.926929)--(-2.388443,0.810441)--(-2.233123,0.789480)--cycle;
\path[fill=blue!14!white,draw=blue!14!white,line width=.12pt] (2.233123,0.789480)--(2.388443,0.810441)--(2.353885,0.926929)--(2.200811,0.898392)--cycle;
\path[fill=blue!16!white,draw=blue!16!white,line width=.12pt] (-2.053620,0.844328)--(-2.200811,0.898392)--(-2.233123,0.789480)--(-2.083771,0.742699)--cycle;
\path[fill=blue!13!white,draw=blue!13!white,line width=.12pt] (2.083771,0.742699)--(2.233123,0.789480)--(2.200811,0.898392)--(2.053620,0.844328)--cycle;
\path[fill=blue!31!white,draw=blue!31!white,line width=.12pt] (-1.600000,0.000000)--(-1.615372,0.135163)--(-1.607593,0.055995)--(-1.592296,-0.078414)--cycle;
\path[fill=blue!22!white,draw=blue!22!white,line width=.12pt] (1.592296,-0.078414)--(1.607593,0.055995)--(1.615372,0.135163)--(1.600000,-0.000000)--cycle;
\path[fill=blue!25!white,draw=blue!25!white,line width=.12pt] (-1.652899,0.346529)--(-1.726471,0.469932)--(-1.734824,0.384910)--(-1.660896,0.265131)--cycle;
\path[fill=blue!17!white,draw=blue!17!white,line width=.12pt] (1.660896,0.265131)--(1.734824,0.384910)--(1.726471,0.469932)--(1.652899,0.346529)--cycle;
\path[fill=blue!40!white,draw=blue!40!white,line width=.12pt] (-1.799069,-0.668826)--(-1.701490,-0.554134)--(-1.660123,-0.636707)--(-1.755330,-0.756135)--cycle;
\path[fill=blue!33!white,draw=blue!33!white,line width=.12pt] (1.755330,-0.756135)--(1.660123,-0.636707)--(1.701490,-0.554134)--(1.799069,-0.668826)--cycle;
\path[fill=blue!16!white,draw=blue!16!white,line width=.12pt] (-2.951405,0.635242)--(-3.050416,0.535130)--(-3.065176,0.384910)--(-2.965685,0.489898)--cycle;
\path[fill=blue!23!white,draw=blue!23!white,line width=.12pt] (2.965685,0.489898)--(3.065176,0.384910)--(3.050416,0.535130)--(2.951405,0.635242)--cycle;
\path[fill=blue!39!white,draw=blue!39!white,line width=.12pt] (-2.506958,-0.928840)--(-2.353885,-0.926929)--(-2.296657,-1.041162)--(-2.446009,-1.050502)--cycle;
\path[fill=blue!40!white,draw=blue!40!white,line width=.12pt] (2.446009,-1.050502)--(2.296657,-1.041162)--(2.353885,-0.926929)--(2.506958,-0.928840)--cycle;
\path[fill=blue!31!white,draw=blue!31!white,line width=.12pt] (-3.050416,-0.535130)--(-2.951405,-0.635242)--(-2.908701,-0.779186)--(-3.006279,-0.683903)--cycle;
\path[fill=blue!38!white,draw=blue!38!white,line width=.12pt] (3.006279,-0.683903)--(2.908701,-0.779186)--(2.951405,-0.635242)--(3.050416,-0.535130)--cycle;
\path[fill=blue!34!white,draw=blue!34!white,line width=.12pt] (-1.607593,-0.214330)--(-1.592296,-0.078414)--(-1.569256,-0.156072)--(-1.584333,-0.292734)--cycle;
\path[fill=blue!25!white,draw=blue!25!white,line width=.12pt] (1.584333,-0.292734)--(1.569256,-0.156072)--(1.592296,-0.078414)--(1.607593,-0.214330)--cycle;
\path[fill=blue!38!white,draw=blue!38!white,line width=.12pt] (-1.701490,-0.554134)--(-1.628983,-0.427143)--(-1.589379,-0.506197)--(-1.660123,-0.636707)--cycle;
\path[fill=blue!31!white,draw=blue!31!white,line width=.12pt] (1.660123,-0.636707)--(1.589379,-0.506197)--(1.628983,-0.427143)--(1.701490,-0.554134)--cycle;
\path[fill=blue!37!white,draw=blue!37!white,line width=.12pt] (-2.654149,-0.904054)--(-2.506958,-0.928840)--(-2.446009,-1.050502)--(-2.589621,-1.032859)--cycle;
\path[fill=blue!40!white,draw=blue!40!white,line width=.12pt] (2.589621,-1.032859)--(2.446009,-1.050502)--(2.506958,-0.928840)--(2.654149,-0.904054)--cycle;
\path[fill=blue!21!white,draw=blue!21!white,line width=.12pt] (-3.184628,0.135163)--(-3.200000,0.000000)--(-3.184591,-0.156827)--(-3.169293,-0.020912)--cycle;
\path[fill=blue!30!white,draw=blue!30!white,line width=.12pt] (3.169293,-0.020912)--(3.184591,-0.156827)--(3.200000,-0.000000)--(3.184628,0.135163)--cycle;
\path[fill=blue!22!white,draw=blue!22!white,line width=.12pt] (-1.726471,0.469932)--(-1.825482,0.579795)--(-1.834315,0.489898)--(-1.734824,0.384910)--cycle;
\path[fill=blue!15!white,draw=blue!15!white,line width=.12pt] (1.734824,0.384910)--(1.834315,0.489898)--(1.825482,0.579795)--(1.726471,0.469932)--cycle;
\path[fill=blue!15!white,draw=blue!15!white,line width=.12pt] (-2.830759,0.715462)--(-2.951405,0.635242)--(-2.965685,0.489898)--(-2.844456,0.576059)--cycle;
\path[fill=blue!20!white,draw=blue!20!white,line width=.12pt] (2.844456,0.576059)--(2.965685,0.489898)--(2.951405,0.635242)--(2.830759,0.715462)--cycle;
\path[fill=blue!28!white,draw=blue!28!white,line width=.12pt] (-1.615372,0.135163)--(-1.660896,0.265131)--(-1.652899,0.183733)--(-1.607593,0.055995)--cycle;
\path[fill=blue!19!white,draw=blue!19!white,line width=.12pt] (1.607593,0.055995)--(1.652899,0.183733)--(1.660896,0.265131)--(1.615372,0.135163)--cycle;
\path[fill=blue!28!white,draw=blue!28!white,line width=.12pt] (-3.123988,-0.418974)--(-3.050416,-0.535130)--(-3.006279,-0.683903)--(-3.078787,-0.571335)--cycle;
\path[fill=blue!36!white,draw=blue!36!white,line width=.12pt] (3.078787,-0.571335)--(3.006279,-0.683903)--(3.050416,-0.535130)--(3.123988,-0.418974)--cycle;
\path[fill=blue!20!white,draw=blue!20!white,line width=.12pt] (-1.825482,0.579795)--(-1.946127,0.671897)--(-1.955544,0.576059)--(-1.834315,0.489898)--cycle;
\path[fill=blue!14!white,draw=blue!14!white,line width=.12pt] (1.834315,0.489898)--(1.955544,0.576059)--(1.946127,0.671897)--(1.825482,0.579795)--cycle;
\path[fill=blue!36!white,draw=blue!36!white,line width=.12pt] (-1.628983,-0.427143)--(-1.584333,-0.292734)--(-1.545814,-0.369621)--(-1.589379,-0.506197)--cycle;
\path[fill=blue!28!white,draw=blue!28!white,line width=.12pt] (1.589379,-0.506197)--(1.545814,-0.369621)--(1.584333,-0.292734)--(1.628983,-0.427143)--cycle;
\path[fill=blue!35!white,draw=blue!35!white,line width=.12pt] (-2.789801,-0.853522)--(-2.654149,-0.904054)--(-2.589621,-1.032859)--(-2.721975,-0.988910)--cycle;
\path[fill=blue!39!white,draw=blue!39!white,line width=.12pt] (2.721975,-0.988910)--(2.589621,-1.032859)--(2.654149,-0.904054)--(2.789801,-0.853522)--cycle;
\path[fill=blue!14!white,draw=blue!14!white,line width=.12pt] (-2.693116,0.772707)--(-2.830759,0.715462)--(-2.844456,0.576059)--(-2.706147,0.640083)--cycle;
\path[fill=blue!18!white,draw=blue!18!white,line width=.12pt] (2.706147,0.640083)--(2.844456,0.576059)--(2.830759,0.715462)--(2.693116,0.772707)--cycle;
\path[fill=blue!41!white,draw=blue!41!white,line width=.12pt] (-2.003693,-0.943989)--(-1.871339,-0.859891)--(-1.806687,-0.950236)--(-1.934468,-1.040724)--cycle;
\path[fill=blue!37!white,draw=blue!37!white,line width=.12pt] (1.934468,-1.040724)--(1.806687,-0.950236)--(1.871339,-0.859891)--(2.003693,-0.943989)--cycle;
\path[fill=blue!31!white,draw=blue!31!white,line width=.12pt] (-1.592296,-0.078414)--(-1.607593,0.055995)--(-1.584333,-0.022409)--(-1.569256,-0.156072)--cycle;
\path[fill=blue!22!white,draw=blue!22!white,line width=.12pt] (1.569256,-0.156072)--(1.584333,-0.022409)--(1.607593,0.055995)--(1.592296,-0.078414)--cycle;
\path[fill=blue!41!white,draw=blue!41!white,line width=.12pt] (-1.871339,-0.859891)--(-1.755330,-0.756135)--(-1.694686,-0.840879)--(-1.806687,-0.950236)--cycle;
\path[fill=blue!36!white,draw=blue!36!white,line width=.12pt] (1.806687,-0.950236)--(1.694686,-0.840879)--(1.755330,-0.756135)--(1.871339,-0.859891)--cycle;
\path[fill=blue!19!white,draw=blue!19!white,line width=.12pt] (-3.139104,0.265131)--(-3.184628,0.135163)--(-3.169293,-0.020912)--(-3.123988,0.111288)--cycle;
\path[fill=blue!27!white,draw=blue!27!white,line width=.12pt] (3.123988,0.111288)--(3.169293,-0.020912)--(3.184628,0.135163)--(3.139104,0.265131)--cycle;
\path[fill=blue!41!white,draw=blue!41!white,line width=.12pt] (-2.147305,-1.005197)--(-2.003693,-0.943989)--(-1.934468,-1.040724)--(-2.073119,-1.108865)--cycle;
\path[fill=blue!38!white,draw=blue!38!white,line width=.12pt] (2.073119,-1.108865)--(1.934468,-1.040724)--(2.003693,-0.943989)--(2.147305,-1.005197)--cycle;
\path[fill=blue!25!white,draw=blue!25!white,line width=.12pt] (-1.660896,0.265131)--(-1.734824,0.384910)--(-1.726471,0.299889)--(-1.652899,0.183733)--cycle;
\path[fill=blue!17!white,draw=blue!17!white,line width=.12pt] (1.652899,0.183733)--(1.726471,0.299889)--(1.734824,0.384910)--(1.660896,0.265131)--cycle;
\path[fill=blue!18!white,draw=blue!18!white,line width=.12pt] (-1.946127,0.671897)--(-2.083771,0.742699)--(-2.093853,0.640083)--(-1.955544,0.576059)--cycle;
\path[fill=blue!14!white,draw=blue!14!white,line width=.12pt] (1.955544,0.576059)--(2.093853,0.640083)--(2.083771,0.742699)--(1.946127,0.671897)--cycle;
\path[fill=blue!40!white,draw=blue!40!white,line width=.12pt] (-1.755330,-0.756135)--(-1.660123,-0.636707)--(-1.602769,-0.716855)--(-1.694686,-0.840879)--cycle;
\path[fill=blue!33!white,draw=blue!33!white,line width=.12pt] (1.694686,-0.840879)--(1.602769,-0.716855)--(1.660123,-0.636707)--(1.755330,-0.756135)--cycle;
\path[fill=blue!14!white,draw=blue!14!white,line width=.12pt] (-2.543764,0.804777)--(-2.693116,0.772707)--(-2.706147,0.640083)--(-2.556072,0.679508)--cycle;
\path[fill=blue!16!white,draw=blue!16!white,line width=.12pt] (2.556072,0.679508)--(2.706147,0.640083)--(2.693116,0.772707)--(2.543764,0.804777)--cycle;
\path[fill=blue!40!white,draw=blue!40!white,line width=.12pt] (-2.296657,-1.041162)--(-2.147305,-1.005197)--(-2.073119,-1.108865)--(-2.217311,-1.152040)--cycle;
\path[fill=blue!39!white,draw=blue!39!white,line width=.12pt] (2.217311,-1.152040)--(2.073119,-1.108865)--(2.147305,-1.005197)--(2.296657,-1.041162)--cycle;
\path[fill=blue!25!white,draw=blue!25!white,line width=.12pt] (-3.169293,-0.291237)--(-3.123988,-0.418974)--(-3.078787,-0.571335)--(-3.123436,-0.445808)--cycle;
\path[fill=blue!34!white,draw=blue!34!white,line width=.12pt] (3.123436,-0.445808)--(3.078787,-0.571335)--(3.123988,-0.418974)--(3.169293,-0.291237)--cycle;
\path[fill=blue!16!white,draw=blue!16!white,line width=.12pt] (-2.083771,0.742699)--(-2.233123,0.789480)--(-2.243928,0.679508)--(-2.093853,0.640083)--cycle;
\path[fill=blue!14!white,draw=blue!14!white,line width=.12pt] (2.093853,0.640083)--(2.243928,0.679508)--(2.233123,0.789480)--(2.083771,0.742699)--cycle;
\path[fill=blue!14!white,draw=blue!14!white,line width=.12pt] (-2.388443,0.810441)--(-2.543764,0.804777)--(-2.556072,0.679508)--(-2.400000,0.692820)--cycle;
\path[fill=blue!15!white,draw=blue!15!white,line width=.12pt] (2.400000,0.692820)--(2.556072,0.679508)--(2.543764,0.804777)--(2.388443,0.810441)--cycle;
\path[fill=blue!15!white,draw=blue!15!white,line width=.12pt] (-2.233123,0.789480)--(-2.388443,0.810441)--(-2.400000,0.692820)--(-2.243928,0.679508)--cycle;
\path[fill=blue!14!white,draw=blue!14!white,line width=.12pt] (2.243928,0.679508)--(2.400000,0.692820)--(2.388443,0.810441)--(2.233123,0.789480)--cycle;
\path[fill=blue!34!white,draw=blue!34!white,line width=.12pt] (-1.584333,-0.292734)--(-1.569256,-0.156072)--(-1.531105,-0.232228)--(-1.545814,-0.369621)--cycle;
\path[fill=blue!26!white,draw=blue!26!white,line width=.12pt] (1.545814,-0.369621)--(1.531105,-0.232228)--(1.569256,-0.156072)--(1.584333,-0.292734)--cycle;
\path[fill=blue!33!white,draw=blue!33!white,line width=.12pt] (-2.908701,-0.779186)--(-2.789801,-0.853522)--(-2.721975,-0.988910)--(-2.837984,-0.920344)--cycle;
\path[fill=blue!38!white,draw=blue!38!white,line width=.12pt] (2.837984,-0.920344)--(2.721975,-0.988910)--(2.789801,-0.853522)--(2.908701,-0.779186)--cycle;
\path[fill=blue!39!white,draw=blue!39!white,line width=.12pt] (-1.660123,-0.636707)--(-1.589379,-0.506197)--(-1.534468,-0.582930)--(-1.602769,-0.716855)--cycle;
\path[fill=blue!31!white,draw=blue!31!white,line width=.12pt] (1.602769,-0.716855)--(1.534468,-0.582930)--(1.589379,-0.506197)--(1.660123,-0.636707)--cycle;
\path[fill=blue!17!white,draw=blue!17!white,line width=.12pt] (-3.065176,0.384910)--(-3.139104,0.265131)--(-3.123988,0.111288)--(-3.050416,0.234690)--cycle;
\path[fill=blue!25!white,draw=blue!25!white,line width=.12pt] (3.050416,0.234690)--(3.123988,0.111288)--(3.139104,0.265131)--(3.065176,0.384910)--cycle;
\path[fill=blue!29!white,draw=blue!29!white,line width=.12pt] (-1.607593,0.055995)--(-1.652899,0.183733)--(-1.628983,0.103118)--(-1.584333,-0.022409)--cycle;
\path[fill=blue!20!white,draw=blue!20!white,line width=.12pt] (1.584333,-0.022409)--(1.628983,0.103118)--(1.652899,0.183733)--(1.607593,0.055995)--cycle;
\path[fill=blue!39!white,draw=blue!39!white,line width=.12pt] (-2.446009,-1.050502)--(-2.296657,-1.041162)--(-2.217311,-1.152040)--(-2.361503,-1.168591)--cycle;
\path[fill=blue!40!white,draw=blue!40!white,line width=.12pt] (2.361503,-1.168591)--(2.217311,-1.152040)--(2.296657,-1.041162)--(2.446009,-1.050502)--cycle;
\path[fill=blue!23!white,draw=blue!23!white,line width=.12pt] (-1.734824,0.384910)--(-1.834315,0.489898)--(-1.825482,0.400001)--(-1.726471,0.299889)--cycle;
\path[fill=blue!16!white,draw=blue!16!white,line width=.12pt] (1.726471,0.299889)--(1.825482,0.400001)--(1.834315,0.489898)--(1.734824,0.384910)--cycle;
\path[fill=blue!37!white,draw=blue!37!white,line width=.12pt] (-1.589379,-0.506197)--(-1.545814,-0.369621)--(-1.492409,-0.444251)--(-1.534468,-0.582930)--cycle;
\path[fill=blue!29!white,draw=blue!29!white,line width=.12pt] (1.534468,-0.582930)--(1.492409,-0.444251)--(1.545814,-0.369621)--(1.589379,-0.506197)--cycle;
\path[fill=blue!23!white,draw=blue!23!white,line width=.12pt] (-3.184591,-0.156827)--(-3.169293,-0.291237)--(-3.123436,-0.445808)--(-3.138513,-0.312145)--cycle;
\path[fill=blue!32!white,draw=blue!32!white,line width=.12pt] (3.138513,-0.312145)--(3.123436,-0.445808)--(3.169293,-0.291237)--(3.184591,-0.156827)--cycle;
\path[fill=blue!32!white,draw=blue!32!white,line width=.12pt] (-1.569256,-0.156072)--(-1.584333,-0.022409)--(-1.545814,-0.099296)--(-1.531105,-0.232228)--cycle;
\path[fill=blue!23!white,draw=blue!23!white,line width=.12pt] (1.531105,-0.232228)--(1.545814,-0.099296)--(1.584333,-0.022409)--(1.569256,-0.156072)--cycle;
\path[fill=blue!15!white,draw=blue!15!white,line width=.12pt] (-2.965685,0.489898)--(-3.065176,0.384910)--(-3.050416,0.234690)--(-2.951405,0.344554)--cycle;
\path[fill=blue!22!white,draw=blue!22!white,line width=.12pt] (2.951405,0.344554)--(3.050416,0.234690)--(3.065176,0.384910)--(2.965685,0.489898)--cycle;
\path[fill=blue!30!white,draw=blue!30!white,line width=.12pt] (-3.006279,-0.683903)--(-2.908701,-0.779186)--(-2.837984,-0.920344)--(-2.933190,-0.829797)--cycle;
\path[fill=blue!37!white,draw=blue!37!white,line width=.12pt] (2.933190,-0.829797)--(2.837984,-0.920344)--(2.908701,-0.779186)--(3.006279,-0.683903)--cycle;
\path[fill=blue!37!white,draw=blue!37!white,line width=.12pt] (-2.589621,-1.032859)--(-2.446009,-1.050502)--(-2.361503,-1.168591)--(-2.500154,-1.157881)--cycle;
\path[fill=blue!40!white,draw=blue!40!white,line width=.12pt] (2.500154,-1.157881)--(2.361503,-1.168591)--(2.446009,-1.050502)--(2.589621,-1.032859)--cycle;
\path[fill=blue!20!white,draw=blue!20!white,line width=.12pt] (-1.834315,0.489898)--(-1.955544,0.576059)--(-1.946127,0.480221)--(-1.825482,0.400001)--cycle;
\path[fill=blue!15!white,draw=blue!15!white,line width=.12pt] (1.825482,0.400001)--(1.946127,0.480221)--(1.955544,0.576059)--(1.834315,0.489898)--cycle;
\path[fill=blue!41!white,draw=blue!41!white,line width=.12pt] (-1.806687,-0.950236)--(-1.694686,-0.840879)--(-1.617721,-0.922243)--(-1.724636,-1.036978)--cycle;
\path[fill=blue!36!white,draw=blue!36!white,line width=.12pt] (1.724636,-1.036978)--(1.617721,-0.922243)--(1.694686,-0.840879)--(1.806687,-0.950236)--cycle;
\path[fill=blue!41!white,draw=blue!41!white,line width=.12pt] (-1.934468,-1.040724)--(-1.806687,-0.950236)--(-1.724636,-1.036978)--(-1.846614,-1.133600)--cycle;
\path[fill=blue!38!white,draw=blue!38!white,line width=.12pt] (1.846614,-1.133600)--(1.724636,-1.036978)--(1.806687,-0.950236)--(1.934468,-1.040724)--cycle;
\path[fill=blue!26!white,draw=blue!26!white,line width=.12pt] (-1.652899,0.183733)--(-1.726471,0.299889)--(-1.701490,0.215687)--(-1.628983,0.103118)--cycle;
\path[fill=blue!18!white,draw=blue!18!white,line width=.12pt] (1.628983,0.103118)--(1.701490,0.215687)--(1.726471,0.299889)--(1.652899,0.183733)--cycle;
\path[fill=blue!40!white,draw=blue!40!white,line width=.12pt] (-1.694686,-0.840879)--(-1.602769,-0.716855)--(-1.529978,-0.793806)--(-1.617721,-0.922243)--cycle;
\path[fill=blue!34!white,draw=blue!34!white,line width=.12pt] (1.617721,-0.922243)--(1.529978,-0.793806)--(1.602769,-0.716855)--(1.694686,-0.840879)--cycle;
\path[fill=blue!41!white,draw=blue!41!white,line width=.12pt] (-2.073119,-1.108865)--(-1.934468,-1.040724)--(-1.846614,-1.133600)--(-1.978968,-1.208398)--cycle;
\path[fill=blue!39!white,draw=blue!39!white,line width=.12pt] (1.978968,-1.208398)--(1.846614,-1.133600)--(1.934468,-1.040724)--(2.073119,-1.108865)--cycle;
\path[fill=blue!35!white,draw=blue!35!white,line width=.12pt] (-1.545814,-0.369621)--(-1.531105,-0.232228)--(-1.478207,-0.306147)--(-1.492409,-0.444251)--cycle;
\path[fill=blue!26!white,draw=blue!26!white,line width=.12pt] (1.492409,-0.444251)--(1.478207,-0.306147)--(1.531105,-0.232228)--(1.545814,-0.369621)--cycle;
\path[fill=blue!14!white,draw=blue!14!white,line width=.12pt] (-2.844456,0.576059)--(-2.965685,0.489898)--(-2.951405,0.344554)--(-2.830759,0.436656)--cycle;
\path[fill=blue!20!white,draw=blue!20!white,line width=.12pt] (2.830759,0.436656)--(2.951405,0.344554)--(2.965685,0.489898)--(2.844456,0.576059)--cycle;
\path[fill=blue!18!white,draw=blue!18!white,line width=.12pt] (-1.955544,0.576059)--(-2.093853,0.640083)--(-2.083771,0.537466)--(-1.946127,0.480221)--cycle;
\path[fill=blue!14!white,draw=blue!14!white,line width=.12pt] (1.946127,0.480221)--(2.083771,0.537466)--(2.093853,0.640083)--(1.955544,0.576059)--cycle;
\path[fill=blue!39!white,draw=blue!39!white,line width=.12pt] (-1.602769,-0.716855)--(-1.534468,-0.582930)--(-1.464780,-0.656601)--(-1.529978,-0.793806)--cycle;
\path[fill=blue!32!white,draw=blue!32!white,line width=.12pt] (1.529978,-0.793806)--(1.464780,-0.656601)--(1.534468,-0.582930)--(1.602769,-0.716855)--cycle;
\path[fill=blue!20!white,draw=blue!20!white,line width=.12pt] (-3.169293,-0.020912)--(-3.184591,-0.156827)--(-3.138513,-0.312145)--(-3.123436,-0.175483)--cycle;
\path[fill=blue!29!white,draw=blue!29!white,line width=.12pt] (3.123436,-0.175483)--(3.138513,-0.312145)--(3.184591,-0.156827)--(3.169293,-0.020912)--cycle;
\path[fill=blue!35!white,draw=blue!35!white,line width=.12pt] (-2.721975,-0.988910)--(-2.589621,-1.032859)--(-2.500154,-1.157881)--(-2.627935,-1.120322)--cycle;
\path[fill=blue!39!white,draw=blue!39!white,line width=.12pt] (2.627935,-1.120322)--(2.500154,-1.157881)--(2.589621,-1.032859)--(2.721975,-0.988910)--cycle;
\path[fill=blue!29!white,draw=blue!29!white,line width=.12pt] (-1.584333,-0.022409)--(-1.628983,0.103118)--(-1.589379,0.024064)--(-1.545814,-0.099296)--cycle;
\path[fill=blue!21!white,draw=blue!21!white,line width=.12pt] (1.545814,-0.099296)--(1.589379,0.024064)--(1.628983,0.103118)--(1.584333,-0.022409)--cycle;
\path[fill=blue!40!white,draw=blue!40!white,line width=.12pt] (-2.217311,-1.152040)--(-2.073119,-1.108865)--(-1.978968,-1.208398)--(-2.116611,-1.258496)--cycle;
\path[fill=blue!39!white,draw=blue!39!white,line width=.12pt] (2.116611,-1.258496)--(1.978968,-1.208398)--(2.073119,-1.108865)--(2.217311,-1.152040)--cycle;
\path[fill=blue!28!white,draw=blue!28!white,line width=.12pt] (-3.078787,-0.571335)--(-3.006279,-0.683903)--(-2.933190,-0.829797)--(-3.003935,-0.720748)--cycle;
\path[fill=blue!35!white,draw=blue!35!white,line width=.12pt] (3.003935,-0.720748)--(2.933190,-0.829797)--(3.006279,-0.683903)--(3.078787,-0.571335)--cycle;
\path[fill=blue!14!white,draw=blue!14!white,line width=.12pt] (-2.706147,0.640083)--(-2.844456,0.576059)--(-2.830759,0.436656)--(-2.693116,0.507458)--cycle;
\path[fill=blue!18!white,draw=blue!18!white,line width=.12pt] (2.693116,0.507458)--(2.830759,0.436656)--(2.844456,0.576059)--(2.706147,0.640083)--cycle;
\path[fill=blue!16!white,draw=blue!16!white,line width=.12pt] (-2.093853,0.640083)--(-2.243928,0.679508)--(-2.233123,0.569536)--(-2.083771,0.537466)--cycle;
\path[fill=blue!14!white,draw=blue!14!white,line width=.12pt] (2.083771,0.537466)--(2.233123,0.569536)--(2.243928,0.679508)--(2.093853,0.640083)--cycle;
\path[fill=blue!23!white,draw=blue!23!white,line width=.12pt] (-1.726471,0.299889)--(-1.825482,0.400001)--(-1.799069,0.310969)--(-1.701490,0.215687)--cycle;
\path[fill=blue!16!white,draw=blue!16!white,line width=.12pt] (1.701490,0.215687)--(1.799069,0.310969)--(1.825482,0.400001)--(1.726471,0.299889)--cycle;
\path[fill=blue!37!white,draw=blue!37!white,line width=.12pt] (-1.534468,-0.582930)--(-1.492409,-0.444251)--(-1.424631,-0.515903)--(-1.464780,-0.656601)--cycle;
\path[fill=blue!30!white,draw=blue!30!white,line width=.12pt] (1.464780,-0.656601)--(1.424631,-0.515903)--(1.492409,-0.444251)--(1.534468,-0.582930)--cycle;
\path[fill=blue!14!white,draw=blue!14!white,line width=.12pt] (-2.556072,0.679508)--(-2.706147,0.640083)--(-2.693116,0.507458)--(-2.543764,0.554239)--cycle;
\path[fill=blue!16!white,draw=blue!16!white,line width=.12pt] (2.543764,0.554239)--(2.693116,0.507458)--(2.706147,0.640083)--(2.556072,0.679508)--cycle;
\path[fill=blue!15!white,draw=blue!15!white,line width=.12pt] (-2.243928,0.679508)--(-2.400000,0.692820)--(-2.388443,0.575200)--(-2.233123,0.569536)--cycle;
\path[fill=blue!14!white,draw=blue!14!white,line width=.12pt] (2.233123,0.569536)--(2.388443,0.575200)--(2.400000,0.692820)--(2.243928,0.679508)--cycle;
\path[fill=blue!14!white,draw=blue!14!white,line width=.12pt] (-2.400000,0.692820)--(-2.556072,0.679508)--(-2.543764,0.554239)--(-2.388443,0.575200)--cycle;
\path[fill=blue!15!white,draw=blue!15!white,line width=.12pt] (2.388443,0.575200)--(2.543764,0.554239)--(2.556072,0.679508)--(2.400000,0.692820)--cycle;
\path[fill=blue!32!white,draw=blue!32!white,line width=.12pt] (-1.531105,-0.232228)--(-1.545814,-0.099296)--(-1.492409,-0.173925)--(-1.478207,-0.306147)--cycle;
\path[fill=blue!24!white,draw=blue!24!white,line width=.12pt] (1.478207,-0.306147)--(1.492409,-0.173925)--(1.545814,-0.099296)--(1.531105,-0.232228)--cycle;
\path[fill=blue!18!white,draw=blue!18!white,line width=.12pt] (-3.123988,0.111288)--(-3.169293,-0.020912)--(-3.123436,-0.175483)--(-3.078787,-0.041074)--cycle;
\path[fill=blue!27!white,draw=blue!27!white,line width=.12pt] (3.078787,-0.041074)--(3.123436,-0.175483)--(3.169293,-0.020912)--(3.123988,0.111288)--cycle;
\path[fill=blue!39!white,draw=blue!39!white,line width=.12pt] (-2.361503,-1.168591)--(-2.217311,-1.152040)--(-2.116611,-1.258496)--(-2.254254,-1.281970)--cycle;
\path[fill=blue!40!white,draw=blue!40!white,line width=.12pt] (2.254254,-1.281970)--(2.116611,-1.258496)--(2.217311,-1.152040)--(2.361503,-1.168591)--cycle;
\path[fill=blue!41!white,draw=blue!41!white,line width=.12pt] (-1.724636,-1.036978)--(-1.617721,-0.922243)--(-1.525177,-0.999443)--(-1.625975,-1.119280)--cycle;
\path[fill=blue!37!white,draw=blue!37!white,line width=.12pt] (1.625975,-1.119280)--(1.525177,-0.999443)--(1.617721,-0.922243)--(1.724636,-1.036978)--cycle;
\path[fill=blue!33!white,draw=blue!33!white,line width=.12pt] (-2.837984,-0.920344)--(-2.721975,-0.988910)--(-2.627935,-1.120322)--(-2.739936,-1.057357)--cycle;
\path[fill=blue!38!white,draw=blue!38!white,line width=.12pt] (2.739936,-1.057357)--(2.627935,-1.120322)--(2.721975,-0.988910)--(2.837984,-0.920344)--cycle;
\path[fill=blue!41!white,draw=blue!41!white,line width=.12pt] (-1.617721,-0.922243)--(-1.529978,-0.793806)--(-1.442454,-0.866819)--(-1.525177,-0.999443)--cycle;
\path[fill=blue!35!white,draw=blue!35!white,line width=.12pt] (1.525177,-0.999443)--(1.442454,-0.866819)--(1.529978,-0.793806)--(1.617721,-0.922243)--cycle;
\path[fill=blue!26!white,draw=blue!26!white,line width=.12pt] (-1.628983,0.103118)--(-1.701490,0.215687)--(-1.660123,0.133114)--(-1.589379,0.024064)--cycle;
\path[fill=blue!19!white,draw=blue!19!white,line width=.12pt] (1.589379,0.024064)--(1.660123,0.133114)--(1.701490,0.215687)--(1.628983,0.103118)--cycle;
\path[fill=blue!21!white,draw=blue!21!white,line width=.12pt] (-1.825482,0.400001)--(-1.946127,0.480221)--(-1.917969,0.385305)--(-1.799069,0.310969)--cycle;
\path[fill=blue!15!white,draw=blue!15!white,line width=.12pt] (1.799069,0.310969)--(1.917969,0.385305)--(1.946127,0.480221)--(1.825482,0.400001)--cycle;
\path[fill=blue!25!white,draw=blue!25!white,line width=.12pt] (-3.123436,-0.445808)--(-3.078787,-0.571335)--(-3.003935,-0.720748)--(-3.047499,-0.597387)--cycle;
\path[fill=blue!33!white,draw=blue!33!white,line width=.12pt] (3.047499,-0.597387)--(3.003935,-0.720748)--(3.078787,-0.571335)--(3.123436,-0.445808)--cycle;
\path[fill=blue!42!white,draw=blue!42!white,line width=.12pt] (-1.846614,-1.133600)--(-1.724636,-1.036978)--(-1.625975,-1.119280)--(-1.740975,-1.221724)--cycle;
\path[fill=blue!38!white,draw=blue!38!white,line width=.12pt] (1.740975,-1.221724)--(1.625975,-1.119280)--(1.724636,-1.036978)--(1.846614,-1.133600)--cycle;
\path[fill=blue!35!white,draw=blue!35!white,line width=.12pt] (-1.492409,-0.444251)--(-1.478207,-0.306147)--(-1.411074,-0.377117)--(-1.424631,-0.515903)--cycle;
\path[fill=blue!27!white,draw=blue!27!white,line width=.12pt] (1.424631,-0.515903)--(1.411074,-0.377117)--(1.478207,-0.306147)--(1.492409,-0.444251)--cycle;
\path[fill=blue!39!white,draw=blue!39!white,line width=.12pt] (-1.529978,-0.793806)--(-1.464780,-0.656601)--(-1.380985,-0.726503)--(-1.442454,-0.866819)--cycle;
\path[fill=blue!33!white,draw=blue!33!white,line width=.12pt] (1.442454,-0.866819)--(1.380985,-0.726503)--(1.464780,-0.656601)--(1.529978,-0.793806)--cycle;
\path[fill=blue!16!white,draw=blue!16!white,line width=.12pt] (-3.050416,0.234690)--(-3.123988,0.111288)--(-3.078787,-0.041074)--(-3.006279,0.085917)--cycle;
\path[fill=blue!24!white,draw=blue!24!white,line width=.12pt] (3.006279,0.085917)--(3.078787,-0.041074)--(3.123988,0.111288)--(3.050416,0.234690)--cycle;
\path[fill=blue!41!white,draw=blue!41!white,line width=.12pt] (-1.978968,-1.208398)--(-1.846614,-1.133600)--(-1.740975,-1.221724)--(-1.865758,-1.302838)--cycle;
\path[fill=blue!39!white,draw=blue!39!white,line width=.12pt] (1.865758,-1.302838)--(1.740975,-1.221724)--(1.846614,-1.133600)--(1.978968,-1.208398)--cycle;
\path[fill=blue!37!white,draw=blue!37!white,line width=.12pt] (-2.500154,-1.157881)--(-2.361503,-1.168591)--(-2.254254,-1.281970)--(-2.386608,-1.277917)--cycle;
\path[fill=blue!39!white,draw=blue!39!white,line width=.12pt] (2.386608,-1.277917)--(2.254254,-1.281970)--(2.361503,-1.168591)--(2.500154,-1.157881)--cycle;
\path[fill=blue!30!white,draw=blue!30!white,line width=.12pt] (-1.545814,-0.099296)--(-1.589379,0.024064)--(-1.534468,-0.052668)--(-1.492409,-0.173925)--cycle;
\path[fill=blue!21!white,draw=blue!21!white,line width=.12pt] (1.492409,-0.173925)--(1.534468,-0.052668)--(1.589379,0.024064)--(1.545814,-0.099296)--cycle;
\path[fill=blue!18!white,draw=blue!18!white,line width=.12pt] (-1.946127,0.480221)--(-2.083771,0.537466)--(-2.053620,0.435837)--(-1.917969,0.385305)--cycle;
\path[fill=blue!14!white,draw=blue!14!white,line width=.12pt] (1.917969,0.385305)--(2.053620,0.435837)--(2.083771,0.537466)--(1.946127,0.480221)--cycle;
\path[fill=blue!38!white,draw=blue!38!white,line width=.12pt] (-1.464780,-0.656601)--(-1.424631,-0.515903)--(-1.343133,-0.583889)--(-1.380985,-0.726503)--cycle;
\path[fill=blue!30!white,draw=blue!30!white,line width=.12pt] (1.380985,-0.726503)--(1.343133,-0.583889)--(1.424631,-0.515903)--(1.464780,-0.656601)--cycle;
\path[fill=blue!30!white,draw=blue!30!white,line width=.12pt] (-2.933190,-0.829797)--(-2.837984,-0.920344)--(-2.739936,-1.057357)--(-2.831853,-0.971406)--cycle;
\path[fill=blue!37!white,draw=blue!37!white,line width=.12pt] (2.831853,-0.971406)--(2.739936,-1.057357)--(2.837984,-0.920344)--(2.933190,-0.829797)--cycle;
\path[fill=blue!24!white,draw=blue!24!white,line width=.12pt] (-1.701490,0.215687)--(-1.799069,0.310969)--(-1.755330,0.223661)--(-1.660123,0.133114)--cycle;
\path[fill=blue!17!white,draw=blue!17!white,line width=.12pt] (1.660123,0.133114)--(1.755330,0.223661)--(1.799069,0.310969)--(1.701490,0.215687)--cycle;
\path[fill=blue!22!white,draw=blue!22!white,line width=.12pt] (-3.138513,-0.312145)--(-3.123436,-0.445808)--(-3.047499,-0.597387)--(-3.062209,-0.464455)--cycle;
\path[fill=blue!31!white,draw=blue!31!white,line width=.12pt] (3.062209,-0.464455)--(3.047499,-0.597387)--(3.123436,-0.445808)--(3.138513,-0.312145)--cycle;
\path[fill=blue!33!white,draw=blue!33!white,line width=.12pt] (-1.478207,-0.306147)--(-1.492409,-0.173925)--(-1.424631,-0.245578)--(-1.411074,-0.377117)--cycle;
\path[fill=blue!25!white,draw=blue!25!white,line width=.12pt] (1.411074,-0.377117)--(1.424631,-0.245578)--(1.492409,-0.173925)--(1.478207,-0.306147)--cycle;
\path[fill=blue!15!white,draw=blue!15!white,line width=.12pt] (-2.951405,0.344554)--(-3.050416,0.234690)--(-3.006279,0.085917)--(-2.908701,0.200610)--cycle;
\path[fill=blue!22!white,draw=blue!22!white,line width=.12pt] (2.908701,0.200610)--(3.006279,0.085917)--(3.050416,0.234690)--(2.951405,0.344554)--cycle;
\path[fill=blue!40!white,draw=blue!40!white,line width=.12pt] (-2.116611,-1.258496)--(-1.978968,-1.208398)--(-1.865758,-1.302838)--(-1.995527,-1.359505)--cycle;
\path[fill=blue!40!white,draw=blue!40!white,line width=.12pt] (1.995527,-1.359505)--(1.865758,-1.302838)--(1.978968,-1.208398)--(2.116611,-1.258496)--cycle;
\path[fill=blue!17!white,draw=blue!17!white,line width=.12pt] (-2.083771,0.537466)--(-2.233123,0.569536)--(-2.200811,0.460624)--(-2.053620,0.435837)--cycle;
\path[fill=blue!14!white,draw=blue!14!white,line width=.12pt] (2.053620,0.435837)--(2.200811,0.460624)--(2.233123,0.569536)--(2.083771,0.537466)--cycle;
\path[fill=blue!41!white,draw=blue!41!white,line width=.12pt] (-1.525177,-0.999443)--(-1.442454,-0.866819)--(-1.341037,-0.935191)--(-1.417944,-1.071736)--cycle;
\path[fill=blue!35!white,draw=blue!35!white,line width=.12pt] (1.417944,-1.071736)--(1.341037,-0.935191)--(1.442454,-0.866819)--(1.525177,-0.999443)--cycle;
\path[fill=blue!42!white,draw=blue!42!white,line width=.12pt] (-1.625975,-1.119280)--(-1.525177,-0.999443)--(-1.417944,-1.071736)--(-1.511656,-1.196351)--cycle;
\path[fill=blue!37!white,draw=blue!37!white,line width=.12pt] (1.511656,-1.196351)--(1.417944,-1.071736)--(1.525177,-0.999443)--(1.625975,-1.119280)--cycle;
\path[fill=blue!35!white,draw=blue!35!white,line width=.12pt] (-2.627935,-1.120322)--(-2.500154,-1.157881)--(-2.386608,-1.277917)--(-2.508586,-1.246493)--cycle;
\path[fill=blue!39!white,draw=blue!39!white,line width=.12pt] (2.508586,-1.246493)--(2.386608,-1.277917)--(2.500154,-1.157881)--(2.627935,-1.120322)--cycle;
\path[fill=blue!36!white,draw=blue!36!white,line width=.12pt] (-1.424631,-0.515903)--(-1.411074,-0.377117)--(-1.330351,-0.444456)--(-1.343133,-0.583889)--cycle;
\path[fill=blue!28!white,draw=blue!28!white,line width=.12pt] (1.343133,-0.583889)--(1.330351,-0.444456)--(1.411074,-0.377117)--(1.424631,-0.515903)--cycle;
\path[fill=blue!14!white,draw=blue!14!white,line width=.12pt] (-2.830759,0.436656)--(-2.951405,0.344554)--(-2.908701,0.200610)--(-2.789801,0.298596)--cycle;
\path[fill=blue!19!white,draw=blue!19!white,line width=.12pt] (2.789801,0.298596)--(2.908701,0.200610)--(2.951405,0.344554)--(2.830759,0.436656)--cycle;
\path[fill=blue!27!white,draw=blue!27!white,line width=.12pt] (-1.589379,0.024064)--(-1.660123,0.133114)--(-1.602769,0.052966)--(-1.534468,-0.052668)--cycle;
\path[fill=blue!19!white,draw=blue!19!white,line width=.12pt] (1.534468,-0.052668)--(1.602769,0.052966)--(1.660123,0.133114)--(1.589379,0.024064)--cycle;
\path[fill=blue!40!white,draw=blue!40!white,line width=.12pt] (-1.442454,-0.866819)--(-1.380985,-0.726503)--(-1.283890,-0.791962)--(-1.341037,-0.935191)--cycle;
\path[fill=blue!33!white,draw=blue!33!white,line width=.12pt] (1.341037,-0.935191)--(1.283890,-0.791962)--(1.380985,-0.726503)--(1.442454,-0.866819)--cycle;
\path[fill=blue!15!white,draw=blue!15!white,line width=.12pt] (-2.233123,0.569536)--(-2.388443,0.575200)--(-2.353885,0.458712)--(-2.200811,0.460624)--cycle;
\path[fill=blue!14!white,draw=blue!14!white,line width=.12pt] (2.200811,0.460624)--(2.353885,0.458712)--(2.388443,0.575200)--(2.233123,0.569536)--cycle;
\path[fill=blue!42!white,draw=blue!42!white,line width=.12pt] (-1.740975,-1.221724)--(-1.625975,-1.119280)--(-1.511656,-1.196351)--(-1.618570,-1.304246)--cycle;
\path[fill=blue!38!white,draw=blue!38!white,line width=.12pt] (1.618570,-1.304246)--(1.511656,-1.196351)--(1.625975,-1.119280)--(1.740975,-1.221724)--cycle;
\path[fill=blue!13!white,draw=blue!13!white,line width=.12pt] (-2.693116,0.507458)--(-2.830759,0.436656)--(-2.789801,0.298596)--(-2.654149,0.376111)--cycle;
\path[fill=blue!18!white,draw=blue!18!white,line width=.12pt] (2.654149,0.376111)--(2.789801,0.298596)--(2.830759,0.436656)--(2.693116,0.507458)--cycle;
\path[fill=blue!39!white,draw=blue!39!white,line width=.12pt] (-2.254254,-1.281970)--(-2.116611,-1.258496)--(-1.995527,-1.359505)--(-2.125296,-1.389547)--cycle;
\path[fill=blue!40!white,draw=blue!40!white,line width=.12pt] (2.125296,-1.389547)--(1.995527,-1.359505)--(2.116611,-1.258496)--(2.254254,-1.281970)--cycle;
\path[fill=blue!21!white,draw=blue!21!white,line width=.12pt] (-1.799069,0.310969)--(-1.917969,0.385305)--(-1.871339,0.292227)--(-1.755330,0.223661)--cycle;
\path[fill=blue!16!white,draw=blue!16!white,line width=.12pt] (1.755330,0.223661)--(1.871339,0.292227)--(1.917969,0.385305)--(1.799069,0.310969)--cycle;
\path[fill=blue!20!white,draw=blue!20!white,line width=.12pt] (-3.123436,-0.175483)--(-3.138513,-0.312145)--(-3.062209,-0.464455)--(-3.047499,-0.327062)--cycle;
\path[fill=blue!28!white,draw=blue!28!white,line width=.12pt] (3.047499,-0.327062)--(3.062209,-0.464455)--(3.138513,-0.312145)--(3.123436,-0.175483)--cycle;
\path[fill=blue!14!white,draw=blue!14!white,line width=.12pt] (-2.388443,0.575200)--(-2.543764,0.554239)--(-2.506958,0.430175)--(-2.353885,0.458712)--cycle;
\path[fill=blue!15!white,draw=blue!15!white,line width=.12pt] (2.353885,0.458712)--(2.506958,0.430175)--(2.543764,0.554239)--(2.388443,0.575200)--cycle;
\path[fill=blue!13!white,draw=blue!13!white,line width=.12pt] (-2.543764,0.554239)--(-2.693116,0.507458)--(-2.654149,0.376111)--(-2.506958,0.430175)--cycle;
\path[fill=blue!16!white,draw=blue!16!white,line width=.12pt] (2.506958,0.430175)--(2.654149,0.376111)--(2.693116,0.507458)--(2.543764,0.554239)--cycle;
\path[fill=blue!27!white,draw=blue!27!white,line width=.12pt] (-3.003935,-0.720748)--(-2.933190,-0.829797)--(-2.831853,-0.971406)--(-2.900154,-0.865772)--cycle;
\path[fill=blue!35!white,draw=blue!35!white,line width=.12pt] (2.900154,-0.865772)--(2.831853,-0.971406)--(2.933190,-0.829797)--(3.003935,-0.720748)--cycle;
\path[fill=blue!30!white,draw=blue!30!white,line width=.12pt] (-1.492409,-0.173925)--(-1.534468,-0.052668)--(-1.464780,-0.126340)--(-1.424631,-0.245578)--cycle;
\path[fill=blue!22!white,draw=blue!22!white,line width=.12pt] (1.424631,-0.245578)--(1.464780,-0.126340)--(1.534468,-0.052668)--(1.492409,-0.173925)--cycle;
\path[fill=blue!38!white,draw=blue!38!white,line width=.12pt] (-1.380985,-0.726503)--(-1.343133,-0.583889)--(-1.248699,-0.647553)--(-1.283890,-0.791962)--cycle;
\path[fill=blue!31!white,draw=blue!31!white,line width=.12pt] (1.283890,-0.791962)--(1.248699,-0.647553)--(1.343133,-0.583889)--(1.380985,-0.726503)--cycle;
\path[fill=blue!41!white,draw=blue!41!white,line width=.12pt] (-1.865758,-1.302838)--(-1.740975,-1.221724)--(-1.618570,-1.304246)--(-1.734580,-1.391274)--cycle;
\path[fill=blue!39!white,draw=blue!39!white,line width=.12pt] (1.734580,-1.391274)--(1.618570,-1.304246)--(1.740975,-1.221724)--(1.865758,-1.302838)--cycle;
\path[fill=blue!33!white,draw=blue!33!white,line width=.12pt] (-1.411074,-0.377117)--(-1.424631,-0.245578)--(-1.343133,-0.313564)--(-1.330351,-0.444456)--cycle;
\path[fill=blue!25!white,draw=blue!25!white,line width=.12pt] (1.330351,-0.444456)--(1.343133,-0.313564)--(1.424631,-0.245578)--(1.411074,-0.377117)--cycle;
\path[fill=blue!32!white,draw=blue!32!white,line width=.12pt] (-2.739936,-1.057357)--(-2.627935,-1.120322)--(-2.508586,-1.246493)--(-2.615501,-1.188905)--cycle;
\path[fill=blue!37!white,draw=blue!37!white,line width=.12pt] (2.615501,-1.188905)--(2.508586,-1.246493)--(2.627935,-1.120322)--(2.739936,-1.057357)--cycle;
\path[fill=blue!24!white,draw=blue!24!white,line width=.12pt] (-1.660123,0.133114)--(-1.755330,0.223661)--(-1.694686,0.138917)--(-1.602769,0.052966)--cycle;
\path[fill=blue!17!white,draw=blue!17!white,line width=.12pt] (1.602769,0.052966)--(1.694686,0.138917)--(1.755330,0.223661)--(1.660123,0.133114)--cycle;
\path[fill=blue!19!white,draw=blue!19!white,line width=.12pt] (-1.917969,0.385305)--(-2.053620,0.435837)--(-2.003693,0.336176)--(-1.871339,0.292227)--cycle;
\path[fill=blue!15!white,draw=blue!15!white,line width=.12pt] (1.871339,0.292227)--(2.003693,0.336176)--(2.053620,0.435837)--(1.917969,0.385305)--cycle;
\path[fill=blue!18!white,draw=blue!18!white,line width=.12pt] (-3.078787,-0.041074)--(-3.123436,-0.175483)--(-3.047499,-0.327062)--(-3.003935,-0.190486)--cycle;
\path[fill=blue!26!white,draw=blue!26!white,line width=.12pt] (3.003935,-0.190486)--(3.047499,-0.327062)--(3.123436,-0.175483)--(3.078787,-0.041074)--cycle;
\path[fill=blue!41!white,draw=blue!41!white,line width=.12pt] (-1.417944,-1.071736)--(-1.341037,-0.935191)--(-1.226706,-0.998263)--(-1.297056,-1.138426)--cycle;
\path[fill=blue!36!white,draw=blue!36!white,line width=.12pt] (1.297056,-1.138426)--(1.226706,-0.998263)--(1.341037,-0.935191)--(1.417944,-1.071736)--cycle;
\path[fill=blue!37!white,draw=blue!37!white,line width=.12pt] (-2.386608,-1.277917)--(-2.254254,-1.281970)--(-2.125296,-1.389547)--(-2.250079,-1.391810)--cycle;
\path[fill=blue!39!white,draw=blue!39!white,line width=.12pt] (2.250079,-1.391810)--(2.125296,-1.389547)--(2.254254,-1.281970)--(2.386608,-1.277917)--cycle;
\path[fill=blue!36!white,draw=blue!36!white,line width=.12pt] (-1.343133,-0.583889)--(-1.330351,-0.444456)--(-1.236817,-0.507515)--(-1.248699,-0.647553)--cycle;
\path[fill=blue!29!white,draw=blue!29!white,line width=.12pt] (1.248699,-0.647553)--(1.236817,-0.507515)--(1.330351,-0.444456)--(1.343133,-0.583889)--cycle;
\path[fill=blue!42!white,draw=blue!42!white,line width=.12pt] (-1.511656,-1.196351)--(-1.417944,-1.071736)--(-1.297056,-1.138426)--(-1.382778,-1.267448)--cycle;
\path[fill=blue!38!white,draw=blue!38!white,line width=.12pt] (1.382778,-1.267448)--(1.297056,-1.138426)--(1.417944,-1.071736)--(1.511656,-1.196351)--cycle;
\path[fill=blue!40!white,draw=blue!40!white,line width=.12pt] (-1.341037,-0.935191)--(-1.283890,-0.791962)--(-1.174431,-0.852346)--(-1.226706,-0.998263)--cycle;
\path[fill=blue!34!white,draw=blue!34!white,line width=.12pt] (1.226706,-0.998263)--(1.174431,-0.852346)--(1.283890,-0.791962)--(1.341037,-0.935191)--cycle;
\path[fill=blue!40!white,draw=blue!40!white,line width=.12pt] (-1.995527,-1.359505)--(-1.865758,-1.302838)--(-1.734580,-1.391274)--(-1.855225,-1.454092)--cycle;
\path[fill=blue!40!white,draw=blue!40!white,line width=.12pt] (1.855225,-1.454092)--(1.734580,-1.391274)--(1.865758,-1.302838)--(1.995527,-1.359505)--cycle;
\path[fill=blue!24!white,draw=blue!24!white,line width=.12pt] (-3.047499,-0.597387)--(-3.003935,-0.720748)--(-2.900154,-0.865772)--(-2.942213,-0.744515)--cycle;
\path[fill=blue!33!white,draw=blue!33!white,line width=.12pt] (2.942213,-0.744515)--(2.900154,-0.865772)--(3.003935,-0.720748)--(3.047499,-0.597387)--cycle;
\path[fill=blue!27!white,draw=blue!27!white,line width=.12pt] (-1.534468,-0.052668)--(-1.602769,0.052966)--(-1.529978,-0.023985)--(-1.464780,-0.126340)--cycle;
\path[fill=blue!20!white,draw=blue!20!white,line width=.12pt] (1.464780,-0.126340)--(1.529978,-0.023985)--(1.602769,0.052966)--(1.534468,-0.052668)--cycle;
\path[fill=blue!42!white,draw=blue!42!white,line width=.12pt] (-1.618570,-1.304246)--(-1.511656,-1.196351)--(-1.382778,-1.267448)--(-1.480578,-1.380371)--cycle;
\path[fill=blue!39!white,draw=blue!39!white,line width=.12pt] (1.480578,-1.380371)--(1.382778,-1.267448)--(1.511656,-1.196351)--(1.618570,-1.304246)--cycle;
\path[fill=blue!38!white,draw=blue!38!white,line width=.12pt] (-1.283890,-0.791962)--(-1.248699,-0.647553)--(-1.142240,-0.706283)--(-1.174431,-0.852346)--cycle;
\path[fill=blue!32!white,draw=blue!32!white,line width=.12pt] (1.174431,-0.852346)--(1.142240,-0.706283)--(1.248699,-0.647553)--(1.283890,-0.791962)--cycle;
\path[fill=blue!17!white,draw=blue!17!white,line width=.12pt] (-2.053620,0.435837)--(-2.200811,0.460624)--(-2.147305,0.353819)--(-2.003693,0.336176)--cycle;
\path[fill=blue!14!white,draw=blue!14!white,line width=.12pt] (2.003693,0.336176)--(2.147305,0.353819)--(2.200811,0.460624)--(2.053620,0.435837)--cycle;
\path[fill=blue!16!white,draw=blue!16!white,line width=.12pt] (-3.006279,0.085917)--(-3.078787,-0.041074)--(-3.003935,-0.190486)--(-2.933190,-0.059976)--cycle;
\path[fill=blue!23!white,draw=blue!23!white,line width=.12pt] (2.933190,-0.059976)--(3.003935,-0.190486)--(3.078787,-0.041074)--(3.006279,0.085917)--cycle;
\path[fill=blue!30!white,draw=blue!30!white,line width=.12pt] (-1.424631,-0.245578)--(-1.464780,-0.126340)--(-1.380985,-0.196241)--(-1.343133,-0.313564)--cycle;
\path[fill=blue!23!white,draw=blue!23!white,line width=.12pt] (1.343133,-0.313564)--(1.380985,-0.196241)--(1.464780,-0.126340)--(1.424631,-0.245578)--cycle;
\path[fill=blue!21!white,draw=blue!21!white,line width=.12pt] (-1.755330,0.223661)--(-1.871339,0.292227)--(-1.806687,0.201882)--(-1.694686,0.138917)--cycle;
\path[fill=blue!16!white,draw=blue!16!white,line width=.12pt] (1.694686,0.138917)--(1.806687,0.201882)--(1.871339,0.292227)--(1.755330,0.223661)--cycle;
\path[fill=blue!30!white,draw=blue!30!white,line width=.12pt] (-2.831853,-0.971406)--(-2.739936,-1.057357)--(-2.615501,-1.188905)--(-2.703244,-1.107367)--cycle;
\path[fill=blue!36!white,draw=blue!36!white,line width=.12pt] (2.703244,-1.107367)--(2.615501,-1.188905)--(2.739936,-1.057357)--(2.831853,-0.971406)--cycle;
\path[fill=blue!33!white,draw=blue!33!white,line width=.12pt] (-1.330351,-0.444456)--(-1.343133,-0.313564)--(-1.248699,-0.377228)--(-1.236817,-0.507515)--cycle;
\path[fill=blue!26!white,draw=blue!26!white,line width=.12pt] (1.236817,-0.507515)--(1.248699,-0.377228)--(1.343133,-0.313564)--(1.330351,-0.444456)--cycle;
\path[fill=blue!41!white,draw=blue!41!white,line width=.12pt] (-1.734580,-1.391274)--(-1.618570,-1.304246)--(-1.480578,-1.380371)--(-1.586697,-1.472856)--cycle;
\path[fill=blue!39!white,draw=blue!39!white,line width=.12pt] (1.586697,-1.472856)--(1.480578,-1.380371)--(1.618570,-1.304246)--(1.734580,-1.391274)--cycle;
\path[fill=blue!39!white,draw=blue!39!white,line width=.12pt] (-2.125296,-1.389547)--(-1.995527,-1.359505)--(-1.855225,-1.454092)--(-1.975871,-1.490286)--cycle;
\path[fill=blue!40!white,draw=blue!40!white,line width=.12pt] (1.975871,-1.490286)--(1.855225,-1.454092)--(1.995527,-1.359505)--(2.125296,-1.389547)--cycle;
\path[fill=blue!35!white,draw=blue!35!white,line width=.12pt] (-2.508586,-1.246493)--(-2.386608,-1.277917)--(-2.250079,-1.391810)--(-2.365079,-1.366207)--cycle;
\path[fill=blue!38!white,draw=blue!38!white,line width=.12pt] (2.365079,-1.366207)--(2.250079,-1.391810)--(2.386608,-1.277917)--(2.508586,-1.246493)--cycle;
\path[fill=blue!15!white,draw=blue!15!white,line width=.12pt] (-2.200811,0.460624)--(-2.353885,0.458712)--(-2.296657,0.344479)--(-2.147305,0.353819)--cycle;
\path[fill=blue!14!white,draw=blue!14!white,line width=.12pt] (2.147305,0.353819)--(2.296657,0.344479)--(2.353885,0.458712)--(2.200811,0.460624)--cycle;
\path[fill=blue!14!white,draw=blue!14!white,line width=.12pt] (-2.908701,0.200610)--(-3.006279,0.085917)--(-2.933190,-0.059976)--(-2.837984,0.059451)--cycle;
\path[fill=blue!21!white,draw=blue!21!white,line width=.12pt] (2.837984,0.059451)--(2.933190,-0.059976)--(3.006279,0.085917)--(2.908701,0.200610)--cycle;
\path[fill=blue!41!white,draw=blue!41!white,line width=.12pt] (-1.297056,-1.138426)--(-1.226706,-0.998263)--(-1.100561,-1.055429)--(-1.163677,-1.198870)--cycle;
\path[fill=blue!37!white,draw=blue!37!white,line width=.12pt] (1.163677,-1.198870)--(1.100561,-1.055429)--(1.226706,-0.998263)--(1.297056,-1.138426)--cycle;
\path[fill=blue!22!white,draw=blue!22!white,line width=.12pt] (-3.062209,-0.464455)--(-3.047499,-0.597387)--(-2.942213,-0.744515)--(-2.956415,-0.612293)--cycle;
\path[fill=blue!30!white,draw=blue!30!white,line width=.12pt] (2.956415,-0.612293)--(2.942213,-0.744515)--(3.047499,-0.597387)--(3.062209,-0.464455)--cycle;
\path[fill=blue!36!white,draw=blue!36!white,line width=.12pt] (-1.248699,-0.647553)--(-1.236817,-0.507515)--(-1.131371,-0.565685)--(-1.142240,-0.706283)--cycle;
\path[fill=blue!29!white,draw=blue!29!white,line width=.12pt] (1.142240,-0.706283)--(1.131371,-0.565685)--(1.236817,-0.507515)--(1.248699,-0.647553)--cycle;
\path[fill=blue!40!white,draw=blue!40!white,line width=.12pt] (-1.226706,-0.998263)--(-1.174431,-0.852346)--(-1.053662,-0.907076)--(-1.100561,-1.055429)--cycle;
\path[fill=blue!35!white,draw=blue!35!white,line width=.12pt] (1.100561,-1.055429)--(1.053662,-0.907076)--(1.174431,-0.852346)--(1.226706,-0.998263)--cycle;
\path[fill=blue!42!white,draw=blue!42!white,line width=.12pt] (-1.382778,-1.267448)--(-1.297056,-1.138426)--(-1.163677,-1.198870)--(-1.240584,-1.331887)--cycle;
\path[fill=blue!38!white,draw=blue!38!white,line width=.12pt] (1.240584,-1.331887)--(1.163677,-1.198870)--(1.297056,-1.138426)--(1.382778,-1.267448)--cycle;
\path[fill=blue!24!white,draw=blue!24!white,line width=.12pt] (-1.602769,0.052966)--(-1.694686,0.138917)--(-1.617721,0.057553)--(-1.529978,-0.023985)--cycle;
\path[fill=blue!18!white,draw=blue!18!white,line width=.12pt] (1.529978,-0.023985)--(1.617721,0.057553)--(1.694686,0.138917)--(1.602769,0.052966)--cycle;
\path[fill=blue!14!white,draw=blue!14!white,line width=.12pt] (-2.353885,0.458712)--(-2.506958,0.430175)--(-2.446009,0.308514)--(-2.296657,0.344479)--cycle;
\path[fill=blue!15!white,draw=blue!15!white,line width=.12pt] (2.296657,0.344479)--(2.446009,0.308514)--(2.506958,0.430175)--(2.353885,0.458712)--cycle;
\path[fill=blue!13!white,draw=blue!13!white,line width=.12pt] (-2.789801,0.298596)--(-2.908701,0.200610)--(-2.837984,0.059451)--(-2.721975,0.163208)--cycle;
\path[fill=blue!19!white,draw=blue!19!white,line width=.12pt] (2.721975,0.163208)--(2.837984,0.059451)--(2.908701,0.200610)--(2.789801,0.298596)--cycle;
\path[fill=blue!13!white,draw=blue!13!white,line width=.12pt] (-2.506958,0.430175)--(-2.654149,0.376111)--(-2.589621,0.247306)--(-2.446009,0.308514)--cycle;
\path[fill=blue!16!white,draw=blue!16!white,line width=.12pt] (2.446009,0.308514)--(2.589621,0.247306)--(2.654149,0.376111)--(2.506958,0.430175)--cycle;
\path[fill=blue!38!white,draw=blue!38!white,line width=.12pt] (-1.174431,-0.852346)--(-1.142240,-0.706283)--(-1.024781,-0.759512)--(-1.053662,-0.907076)--cycle;
\path[fill=blue!33!white,draw=blue!33!white,line width=.12pt] (1.053662,-0.907076)--(1.024781,-0.759512)--(1.142240,-0.706283)--(1.174431,-0.852346)--cycle;
\path[fill=blue!13!white,draw=blue!13!white,line width=.12pt] (-2.654149,0.376111)--(-2.789801,0.298596)--(-2.721975,0.163208)--(-2.589621,0.247306)--cycle;
\path[fill=blue!17!white,draw=blue!17!white,line width=.12pt] (2.589621,0.247306)--(2.721975,0.163208)--(2.789801,0.298596)--(2.654149,0.376111)--cycle;
\path[fill=blue!19!white,draw=blue!19!white,line width=.12pt] (-1.871339,0.292227)--(-2.003693,0.336176)--(-1.934468,0.239441)--(-1.806687,0.201882)--cycle;
\path[fill=blue!15!white,draw=blue!15!white,line width=.12pt] (1.806687,0.201882)--(1.934468,0.239441)--(2.003693,0.336176)--(1.871339,0.292227)--cycle;
\path[fill=blue!28!white,draw=blue!28!white,line width=.12pt] (-1.464780,-0.126340)--(-1.529978,-0.023985)--(-1.442454,-0.096998)--(-1.380985,-0.196241)--cycle;
\path[fill=blue!21!white,draw=blue!21!white,line width=.12pt] (1.380985,-0.196241)--(1.442454,-0.096998)--(1.529978,-0.023985)--(1.464780,-0.126340)--cycle;
\path[fill=blue!42!white,draw=blue!42!white,line width=.12pt] (-1.480578,-1.380371)--(-1.382778,-1.267448)--(-1.240584,-1.331887)--(-1.328326,-1.449368)--cycle;
\path[fill=blue!39!white,draw=blue!39!white,line width=.12pt] (1.328326,-1.449368)--(1.240584,-1.331887)--(1.382778,-1.267448)--(1.480578,-1.380371)--cycle;
\path[fill=blue!40!white,draw=blue!40!white,line width=.12pt] (-1.855225,-1.454092)--(-1.734580,-1.391274)--(-1.586697,-1.472856)--(-1.697056,-1.541348)--cycle;
\path[fill=blue!40!white,draw=blue!40!white,line width=.12pt] (1.697056,-1.541348)--(1.586697,-1.472856)--(1.734580,-1.391274)--(1.855225,-1.454092)--cycle;
\path[fill=blue!27!white,draw=blue!27!white,line width=.12pt] (-2.900154,-0.865772)--(-2.831853,-0.971406)--(-2.703244,-1.107367)--(-2.768442,-1.005012)--cycle;
\path[fill=blue!34!white,draw=blue!34!white,line width=.12pt] (2.768442,-1.005012)--(2.703244,-1.107367)--(2.831853,-0.971406)--(2.900154,-0.865772)--cycle;
\path[fill=blue!31!white,draw=blue!31!white,line width=.12pt] (-1.343133,-0.313564)--(-1.380985,-0.196241)--(-1.283890,-0.261700)--(-1.248699,-0.377228)--cycle;
\path[fill=blue!24!white,draw=blue!24!white,line width=.12pt] (1.248699,-0.377228)--(1.283890,-0.261700)--(1.380985,-0.196241)--(1.343133,-0.313564)--cycle;
\path[fill=blue!19!white,draw=blue!19!white,line width=.12pt] (-3.047499,-0.327062)--(-3.062209,-0.464455)--(-2.956415,-0.612293)--(-2.942213,-0.474190)--cycle;
\path[fill=blue!28!white,draw=blue!28!white,line width=.12pt] (2.942213,-0.474190)--(2.956415,-0.612293)--(3.062209,-0.464455)--(3.047499,-0.327062)--cycle;
\path[fill=blue!34!white,draw=blue!34!white,line width=.12pt] (-1.236817,-0.507515)--(-1.248699,-0.377228)--(-1.142240,-0.435958)--(-1.131371,-0.565685)--cycle;
\path[fill=blue!27!white,draw=blue!27!white,line width=.12pt] (1.131371,-0.565685)--(1.142240,-0.435958)--(1.248699,-0.377228)--(1.236817,-0.507515)--cycle;
\path[fill=blue!37!white,draw=blue!37!white,line width=.12pt] (-2.250079,-1.391810)--(-2.125296,-1.389547)--(-1.975871,-1.490286)--(-2.091880,-1.498463)--cycle;
\path[fill=blue!39!white,draw=blue!39!white,line width=.12pt] (2.091880,-1.498463)--(1.975871,-1.490286)--(2.125296,-1.389547)--(2.250079,-1.391810)--cycle;
\path[fill=blue!32!white,draw=blue!32!white,line width=.12pt] (-2.615501,-1.188905)--(-2.508586,-1.246493)--(-2.365079,-1.366207)--(-2.465877,-1.313721)--cycle;
\path[fill=blue!37!white,draw=blue!37!white,line width=.12pt] (2.465877,-1.313721)--(2.365079,-1.366207)--(2.508586,-1.246493)--(2.615501,-1.188905)--cycle;
\path[fill=blue!41!white,draw=blue!41!white,line width=.12pt] (-1.163677,-1.198870)--(-1.100561,-1.055429)--(-0.963817,-1.106137)--(-1.019091,-1.252486)--cycle;
\path[fill=blue!37!white,draw=blue!37!white,line width=.12pt] (1.019091,-1.252486)--(0.963817,-1.106137)--(1.100561,-1.055429)--(1.163677,-1.198870)--cycle;
\path[fill=blue!36!white,draw=blue!36!white,line width=.12pt] (-1.142240,-0.706283)--(-1.131371,-0.565685)--(-1.015029,-0.618408)--(-1.024781,-0.759512)--cycle;
\path[fill=blue!30!white,draw=blue!30!white,line width=.12pt] (1.024781,-0.759512)--(1.015029,-0.618408)--(1.131371,-0.565685)--(1.142240,-0.706283)--cycle;
\path[fill=blue!40!white,draw=blue!40!white,line width=.12pt] (-1.100561,-1.055429)--(-1.053662,-0.907076)--(-0.922745,-0.955623)--(-0.963817,-1.106137)--cycle;
\path[fill=blue!35!white,draw=blue!35!white,line width=.12pt] (0.963817,-1.106137)--(0.922745,-0.955623)--(1.053662,-0.907076)--(1.100561,-1.055429)--cycle;
\path[fill=blue!22!white,draw=blue!22!white,line width=.12pt] (-1.694686,0.138917)--(-1.806687,0.201882)--(-1.724636,0.115141)--(-1.617721,0.057553)--cycle;
\path[fill=blue!17!white,draw=blue!17!white,line width=.12pt] (1.617721,0.057553)--(1.724636,0.115141)--(1.806687,0.201882)--(1.694686,0.138917)--cycle;
\path[fill=blue!42!white,draw=blue!42!white,line width=.12pt] (-1.240584,-1.331887)--(-1.163677,-1.198870)--(-1.019091,-1.252486)--(-1.086442,-1.389047)--cycle;
\path[fill=blue!38!white,draw=blue!38!white,line width=.12pt] (1.086442,-1.389047)--(1.019091,-1.252486)--(1.163677,-1.198870)--(1.240584,-1.331887)--cycle;
\path[fill=blue!41!white,draw=blue!41!white,line width=.12pt] (-1.586697,-1.472856)--(-1.480578,-1.380371)--(-1.328326,-1.449368)--(-1.423533,-1.546798)--cycle;
\path[fill=blue!40!white,draw=blue!40!white,line width=.12pt] (1.423533,-1.546798)--(1.328326,-1.449368)--(1.480578,-1.380371)--(1.586697,-1.472856)--cycle;
\path[fill=blue!17!white,draw=blue!17!white,line width=.12pt] (-2.003693,0.336176)--(-2.147305,0.353819)--(-2.073119,0.250151)--(-1.934468,0.239441)--cycle;
\path[fill=blue!14!white,draw=blue!14!white,line width=.12pt] (1.934468,0.239441)--(2.073119,0.250151)--(2.147305,0.353819)--(2.003693,0.336176)--cycle;
\path[fill=blue!38!white,draw=blue!38!white,line width=.12pt] (-1.053662,-0.907076)--(-1.024781,-0.759512)--(-0.897452,-0.806729)--(-0.922745,-0.955623)--cycle;
\path[fill=blue!33!white,draw=blue!33!white,line width=.12pt] (0.922745,-0.955623)--(0.897452,-0.806729)--(1.024781,-0.759512)--(1.053662,-0.907076)--cycle;
\path[fill=blue!17!white,draw=blue!17!white,line width=.12pt] (-3.003935,-0.190486)--(-3.047499,-0.327062)--(-2.942213,-0.474190)--(-2.900154,-0.335511)--cycle;
\path[fill=blue!25!white,draw=blue!25!white,line width=.12pt] (2.900154,-0.335511)--(2.942213,-0.474190)--(3.047499,-0.327062)--(3.003935,-0.190486)--cycle;
\path[fill=blue!39!white,draw=blue!39!white,line width=.12pt] (-1.975871,-1.490286)--(-1.855225,-1.454092)--(-1.697056,-1.541348)--(-1.807416,-1.583216)--cycle;
\path[fill=blue!39!white,draw=blue!39!white,line width=.12pt] (1.807416,-1.583216)--(1.697056,-1.541348)--(1.855225,-1.454092)--(1.975871,-1.490286)--cycle;
\path[fill=blue!25!white,draw=blue!25!white,line width=.12pt] (-1.529978,-0.023985)--(-1.617721,0.057553)--(-1.525177,-0.019647)--(-1.442454,-0.096998)--cycle;
\path[fill=blue!19!white,draw=blue!19!white,line width=.12pt] (1.442454,-0.096998)--(1.525177,-0.019647)--(1.617721,0.057553)--(1.529978,-0.023985)--cycle;
\path[fill=blue!24!white,draw=blue!24!white,line width=.12pt] (-2.942213,-0.744515)--(-2.900154,-0.865772)--(-2.768442,-1.005012)--(-2.808591,-0.885774)--cycle;
\path[fill=blue!32!white,draw=blue!32!white,line width=.12pt] (2.808591,-0.885774)--(2.768442,-1.005012)--(2.900154,-0.865772)--(2.942213,-0.744515)--cycle;
\path[fill=blue!42!white,draw=blue!42!white,line width=.12pt] (-1.328326,-1.449368)--(-1.240584,-1.331887)--(-1.086442,-1.389047)--(-1.163283,-1.510570)--cycle;
\path[fill=blue!39!white,draw=blue!39!white,line width=.12pt] (1.163283,-1.510570)--(1.086442,-1.389047)--(1.240584,-1.331887)--(1.328326,-1.449368)--cycle;
\path[fill=blue!28!white,draw=blue!28!white,line width=.12pt] (-1.380985,-0.196241)--(-1.442454,-0.096998)--(-1.341037,-0.165370)--(-1.283890,-0.261700)--cycle;
\path[fill=blue!21!white,draw=blue!21!white,line width=.12pt] (1.283890,-0.261700)--(1.341037,-0.165370)--(1.442454,-0.096998)--(1.380985,-0.196241)--cycle;
\path[fill=blue!31!white,draw=blue!31!white,line width=.12pt] (-1.248699,-0.377228)--(-1.283890,-0.261700)--(-1.174431,-0.322085)--(-1.142240,-0.435958)--cycle;
\path[fill=blue!24!white,draw=blue!24!white,line width=.12pt] (1.142240,-0.435958)--(1.174431,-0.322085)--(1.283890,-0.261700)--(1.248699,-0.377228)--cycle;
\path[fill=blue!34!white,draw=blue!34!white,line width=.12pt] (-1.131371,-0.565685)--(-1.142240,-0.435958)--(-1.024781,-0.489187)--(-1.015029,-0.618408)--cycle;
\path[fill=blue!28!white,draw=blue!28!white,line width=.12pt] (1.015029,-0.618408)--(1.024781,-0.489187)--(1.142240,-0.435958)--(1.131371,-0.565685)--cycle;
\path[fill=blue!41!white,draw=blue!41!white,line width=.12pt] (-1.019091,-1.252486)--(-0.963817,-1.106137)--(-0.817791,-1.149900)--(-0.864690,-1.298758)--cycle;
\path[fill=blue!38!white,draw=blue!38!white,line width=.12pt] (0.864690,-1.298758)--(0.817791,-1.149900)--(0.963817,-1.106137)--(1.019091,-1.252486)--cycle;
\path[fill=blue!40!white,draw=blue!40!white,line width=.12pt] (-0.963817,-1.106137)--(-0.922745,-0.955623)--(-0.782941,-0.997521)--(-0.817791,-1.149900)--cycle;
\path[fill=blue!36!white,draw=blue!36!white,line width=.12pt] (0.817791,-1.149900)--(0.782941,-0.997521)--(0.922745,-0.955623)--(0.963817,-1.106137)--cycle;
\path[fill=blue!15!white,draw=blue!15!white,line width=.12pt] (-2.147305,0.353819)--(-2.296657,0.344479)--(-2.217311,0.233600)--(-2.073119,0.250151)--cycle;
\path[fill=blue!14!white,draw=blue!14!white,line width=.12pt] (2.073119,0.250151)--(2.217311,0.233600)--(2.296657,0.344479)--(2.147305,0.353819)--cycle;
\path[fill=blue!36!white,draw=blue!36!white,line width=.12pt] (-1.024781,-0.759512)--(-1.015029,-0.618408)--(-0.888912,-0.665176)--(-0.897452,-0.806729)--cycle;
\path[fill=blue!31!white,draw=blue!31!white,line width=.12pt] (0.897452,-0.806729)--(0.888912,-0.665176)--(1.015029,-0.618408)--(1.024781,-0.759512)--cycle;
\path[fill=blue!35!white,draw=blue!35!white,line width=.12pt] (-2.365079,-1.366207)--(-2.250079,-1.391810)--(-2.091880,-1.498463)--(-2.198794,-1.478311)--cycle;
\path[fill=blue!38!white,draw=blue!38!white,line width=.12pt] (2.198794,-1.478311)--(2.091880,-1.498463)--(2.250079,-1.391810)--(2.365079,-1.366207)--cycle;
\path[fill=blue!29!white,draw=blue!29!white,line width=.12pt] (-2.703244,-1.107367)--(-2.615501,-1.188905)--(-2.465877,-1.313721)--(-2.548600,-1.236371)--cycle;
\path[fill=blue!35!white,draw=blue!35!white,line width=.12pt] (2.548600,-1.236371)--(2.465877,-1.313721)--(2.615501,-1.188905)--(2.703244,-1.107367)--cycle;
\path[fill=blue!15!white,draw=blue!15!white,line width=.12pt] (-2.933190,-0.059976)--(-3.003935,-0.190486)--(-2.900154,-0.335511)--(-2.831853,-0.201586)--cycle;
\path[fill=blue!23!white,draw=blue!23!white,line width=.12pt] (2.831853,-0.201586)--(2.900154,-0.335511)--(3.003935,-0.190486)--(2.933190,-0.059976)--cycle;
\path[fill=blue!40!white,draw=blue!40!white,line width=.12pt] (-1.697056,-1.541348)--(-1.586697,-1.472856)--(-1.423533,-1.546798)--(-1.522544,-1.620433)--cycle;
\path[fill=blue!40!white,draw=blue!40!white,line width=.12pt] (1.522544,-1.620433)--(1.423533,-1.546798)--(1.586697,-1.472856)--(1.697056,-1.541348)--cycle;
\path[fill=blue!42!white,draw=blue!42!white,line width=.12pt] (-1.086442,-1.389047)--(-1.019091,-1.252486)--(-0.864690,-1.298758)--(-0.921837,-1.438377)--cycle;
\path[fill=blue!39!white,draw=blue!39!white,line width=.12pt] (0.921837,-1.438377)--(0.864690,-1.298758)--(1.019091,-1.252486)--(1.086442,-1.389047)--cycle;
\path[fill=blue!38!white,draw=blue!38!white,line width=.12pt] (-0.922745,-0.955623)--(-0.897452,-0.806729)--(-0.761481,-0.847478)--(-0.782941,-0.997521)--cycle;
\path[fill=blue!34!white,draw=blue!34!white,line width=.12pt] (0.782941,-0.997521)--(0.761481,-0.847478)--(0.897452,-0.806729)--(0.922745,-0.955623)--cycle;
\path[fill=blue!19!white,draw=blue!19!white,line width=.12pt] (-1.806687,0.201882)--(-1.934468,0.239441)--(-1.846614,0.146565)--(-1.724636,0.115141)--cycle;
\path[fill=blue!15!white,draw=blue!15!white,line width=.12pt] (1.724636,0.115141)--(1.846614,0.146565)--(1.934468,0.239441)--(1.806687,0.201882)--cycle;
\path[fill=blue!14!white,draw=blue!14!white,line width=.12pt] (-2.296657,0.344479)--(-2.446009,0.308514)--(-2.361503,0.190425)--(-2.217311,0.233600)--cycle;
\path[fill=blue!15!white,draw=blue!15!white,line width=.12pt] (2.217311,0.233600)--(2.361503,0.190425)--(2.446009,0.308514)--(2.296657,0.344479)--cycle;
\path[fill=blue!41!white,draw=blue!41!white,line width=.12pt] (-1.423533,-1.546798)--(-1.328326,-1.449368)--(-1.163283,-1.510570)--(-1.246659,-1.612387)--cycle;
\path[fill=blue!40!white,draw=blue!40!white,line width=.12pt] (1.246659,-1.612387)--(1.163283,-1.510570)--(1.328326,-1.449368)--(1.423533,-1.546798)--cycle;
\path[fill=blue!14!white,draw=blue!14!white,line width=.12pt] (-2.837984,0.059451)--(-2.933190,-0.059976)--(-2.831853,-0.201586)--(-2.739936,-0.077561)--cycle;
\path[fill=blue!21!white,draw=blue!21!white,line width=.12pt] (2.739936,-0.077561)--(2.831853,-0.201586)--(2.933190,-0.059976)--(2.837984,0.059451)--cycle;
\path[fill=blue!40!white,draw=blue!40!white,line width=.12pt] (-0.817791,-1.149900)--(-0.782941,-0.997521)--(-0.635598,-1.032365)--(-0.663889,-1.186295)--cycle;
\path[fill=blue!37!white,draw=blue!37!white,line width=.12pt] (0.663889,-1.186295)--(0.635598,-1.032365)--(0.782941,-0.997521)--(0.817791,-1.149900)--cycle;
\path[fill=blue!21!white,draw=blue!21!white,line width=.12pt] (-2.956415,-0.612293)--(-2.942213,-0.744515)--(-2.808591,-0.885774)--(-2.822148,-0.754235)--cycle;
\path[fill=blue!29!white,draw=blue!29!white,line width=.12pt] (2.822148,-0.754235)--(2.808591,-0.885774)--(2.942213,-0.744515)--(2.956415,-0.612293)--cycle;
\path[fill=blue!41!white,draw=blue!41!white,line width=.12pt] (-0.864690,-1.298758)--(-0.817791,-1.149900)--(-0.663889,-1.186295)--(-0.701962,-1.337241)--cycle;
\path[fill=blue!38!white,draw=blue!38!white,line width=.12pt] (0.701962,-1.337241)--(0.663889,-1.186295)--(0.817791,-1.149900)--(0.864690,-1.298758)--cycle;
\path[fill=blue!37!white,draw=blue!37!white,line width=.12pt] (-2.091880,-1.498463)--(-1.975871,-1.490286)--(-1.807416,-1.583216)--(-1.913535,-1.596850)--cycle;
\path[fill=blue!39!white,draw=blue!39!white,line width=.12pt] (1.913535,-1.596850)--(1.807416,-1.583216)--(1.975871,-1.490286)--(2.091880,-1.498463)--cycle;
\path[fill=blue!13!white,draw=blue!13!white,line width=.12pt] (-2.446009,0.308514)--(-2.589621,0.247306)--(-2.500154,0.122284)--(-2.361503,0.190425)--cycle;
\path[fill=blue!16!white,draw=blue!16!white,line width=.12pt] (2.361503,0.190425)--(2.500154,0.122284)--(2.589621,0.247306)--(2.446009,0.308514)--cycle;
\path[fill=blue!22!white,draw=blue!22!white,line width=.12pt] (-1.617721,0.057553)--(-1.724636,0.115141)--(-1.625975,0.032838)--(-1.525177,-0.019647)--cycle;
\path[fill=blue!17!white,draw=blue!17!white,line width=.12pt] (1.525177,-0.019647)--(1.625975,0.032838)--(1.724636,0.115141)--(1.617721,0.057553)--cycle;
\path[fill=blue!42!white,draw=blue!42!white,line width=.12pt] (-1.163283,-1.510570)--(-1.086442,-1.389047)--(-0.921837,-1.438377)--(-0.987036,-1.563389)--cycle;
\path[fill=blue!40!white,draw=blue!40!white,line width=.12pt] (0.987036,-1.563389)--(0.921837,-1.438377)--(1.086442,-1.389047)--(1.163283,-1.510570)--cycle;
\path[fill=blue!36!white,draw=blue!36!white,line width=.12pt] (-0.897452,-0.806729)--(-0.888912,-0.665176)--(-0.754235,-0.705537)--(-0.761481,-0.847478)--cycle;
\path[fill=blue!32!white,draw=blue!32!white,line width=.12pt] (0.761481,-0.847478)--(0.754235,-0.705537)--(0.888912,-0.665176)--(0.897452,-0.806729)--cycle;
\path[fill=blue!13!white,draw=blue!13!white,line width=.12pt] (-2.721975,0.163208)--(-2.837984,0.059451)--(-2.739936,-0.077561)--(-2.627935,0.031796)--cycle;
\path[fill=blue!18!white,draw=blue!18!white,line width=.12pt] (2.627935,0.031796)--(2.739936,-0.077561)--(2.837984,0.059451)--(2.721975,0.163208)--cycle;
\path[fill=blue!34!white,draw=blue!34!white,line width=.12pt] (-1.015029,-0.618408)--(-1.024781,-0.489187)--(-0.897452,-0.536404)--(-0.888912,-0.665176)--cycle;
\path[fill=blue!28!white,draw=blue!28!white,line width=.12pt] (0.888912,-0.665176)--(0.897452,-0.536404)--(1.024781,-0.489187)--(1.015029,-0.618408)--cycle;
\path[fill=blue!13!white,draw=blue!13!white,line width=.12pt] (-2.589621,0.247306)--(-2.721975,0.163208)--(-2.627935,0.031796)--(-2.500154,0.122284)--cycle;
\path[fill=blue!17!white,draw=blue!17!white,line width=.12pt] (2.500154,0.122284)--(2.627935,0.031796)--(2.721975,0.163208)--(2.589621,0.247306)--cycle;
\path[fill=blue!38!white,draw=blue!38!white,line width=.12pt] (-0.782941,-0.997521)--(-0.761481,-0.847478)--(-0.618176,-0.881367)--(-0.635598,-1.032365)--cycle;
\path[fill=blue!35!white,draw=blue!35!white,line width=.12pt] (0.635598,-1.032365)--(0.618176,-0.881367)--(0.761481,-0.847478)--(0.782941,-0.997521)--cycle;
\path[fill=blue!31!white,draw=blue!31!white,line width=.12pt] (-1.142240,-0.435958)--(-1.174431,-0.322085)--(-1.053662,-0.376814)--(-1.024781,-0.489187)--cycle;
\path[fill=blue!25!white,draw=blue!25!white,line width=.12pt] (1.024781,-0.489187)--(1.053662,-0.376814)--(1.174431,-0.322085)--(1.142240,-0.435958)--cycle;
\path[fill=blue!25!white,draw=blue!25!white,line width=.12pt] (-1.442454,-0.096998)--(-1.525177,-0.019647)--(-1.417944,-0.091940)--(-1.341037,-0.165370)--cycle;
\path[fill=blue!19!white,draw=blue!19!white,line width=.12pt] (1.341037,-0.165370)--(1.417944,-0.091940)--(1.525177,-0.019647)--(1.442454,-0.096998)--cycle;
\path[fill=blue!28!white,draw=blue!28!white,line width=.12pt] (-1.283890,-0.261700)--(-1.341037,-0.165370)--(-1.226706,-0.228443)--(-1.174431,-0.322085)--cycle;
\path[fill=blue!22!white,draw=blue!22!white,line width=.12pt] (1.174431,-0.322085)--(1.226706,-0.228443)--(1.341037,-0.165370)--(1.283890,-0.261700)--cycle;
\path[fill=blue!42!white,draw=blue!42!white,line width=.12pt] (-0.921837,-1.438377)--(-0.864690,-1.298758)--(-0.701962,-1.337241)--(-0.748354,-1.479402)--cycle;
\path[fill=blue!39!white,draw=blue!39!white,line width=.12pt] (0.748354,-1.479402)--(0.701962,-1.337241)--(0.864690,-1.298758)--(0.921837,-1.438377)--cycle;
\path[fill=blue!17!white,draw=blue!17!white,line width=.12pt] (-1.934468,0.239441)--(-2.073119,0.250151)--(-1.978968,0.150618)--(-1.846614,0.146565)--cycle;
\path[fill=blue!15!white,draw=blue!15!white,line width=.12pt] (1.846614,0.146565)--(1.978968,0.150618)--(2.073119,0.250151)--(1.934468,0.239441)--cycle;
\path[fill=blue!32!white,draw=blue!32!white,line width=.12pt] (-2.465877,-1.313721)--(-2.365079,-1.366207)--(-2.198794,-1.478311)--(-2.292506,-1.430603)--cycle;
\path[fill=blue!36!white,draw=blue!36!white,line width=.12pt] (2.292506,-1.430603)--(2.198794,-1.478311)--(2.365079,-1.366207)--(2.465877,-1.313721)--cycle;
\path[fill=blue!26!white,draw=blue!26!white,line width=.12pt] (-2.768442,-1.005012)--(-2.703244,-1.107367)--(-2.548600,-1.236371)--(-2.610069,-1.137127)--cycle;
\path[fill=blue!33!white,draw=blue!33!white,line width=.12pt] (2.610069,-1.137127)--(2.548600,-1.236371)--(2.703244,-1.107367)--(2.768442,-1.005012)--cycle;
\path[fill=blue!39!white,draw=blue!39!white,line width=.12pt] (-1.807416,-1.583216)--(-1.697056,-1.541348)--(-1.522544,-1.620433)--(-1.621555,-1.667443)--cycle;
\path[fill=blue!39!white,draw=blue!39!white,line width=.12pt] (1.621555,-1.667443)--(1.522544,-1.620433)--(1.697056,-1.541348)--(1.807416,-1.583216)--cycle;
\path[fill=blue!40!white,draw=blue!40!white,line width=.12pt] (-0.663889,-1.186295)--(-0.635598,-1.032365)--(-0.482133,-1.059820)--(-0.503593,-1.214972)--cycle;
\path[fill=blue!37!white,draw=blue!37!white,line width=.12pt] (0.503593,-1.214972)--(0.482133,-1.059820)--(0.635598,-1.032365)--(0.663889,-1.186295)--cycle;
\path[fill=blue!41!white,draw=blue!41!white,line width=.12pt] (-0.701962,-1.337241)--(-0.663889,-1.186295)--(-0.503593,-1.214972)--(-0.532473,-1.367563)--cycle;
\path[fill=blue!39!white,draw=blue!39!white,line width=.12pt] (0.532473,-1.367563)--(0.503593,-1.214972)--(0.663889,-1.186295)--(0.701962,-1.337241)--cycle;
\path[fill=blue!40!white,draw=blue!40!white,line width=.12pt] (-1.522544,-1.620433)--(-1.423533,-1.546798)--(-1.246659,-1.612387)--(-1.333369,-1.690584)--cycle;
\path[fill=blue!40!white,draw=blue!40!white,line width=.12pt] (1.333369,-1.690584)--(1.246659,-1.612387)--(1.423533,-1.546798)--(1.522544,-1.620433)--cycle;
\path[fill=blue!36!white,draw=blue!36!white,line width=.12pt] (-0.761481,-0.847478)--(-0.754235,-0.705537)--(-0.612293,-0.739104)--(-0.618176,-0.881367)--cycle;
\path[fill=blue!32!white,draw=blue!32!white,line width=.12pt] (0.618176,-0.881367)--(0.612293,-0.739104)--(0.754235,-0.705537)--(0.761481,-0.847478)--cycle;
\path[fill=blue!38!white,draw=blue!38!white,line width=.12pt] (-0.635598,-1.032365)--(-0.618176,-0.881367)--(-0.468918,-0.908070)--(-0.482133,-1.059820)--cycle;
\path[fill=blue!35!white,draw=blue!35!white,line width=.12pt] (0.482133,-1.059820)--(0.468918,-0.908070)--(0.618176,-0.881367)--(0.635598,-1.032365)--cycle;
\path[fill=blue!19!white,draw=blue!19!white,line width=.12pt] (-2.942213,-0.474190)--(-2.956415,-0.612293)--(-2.822148,-0.754235)--(-2.808591,-0.615449)--cycle;
\path[fill=blue!27!white,draw=blue!27!white,line width=.12pt] (2.808591,-0.615449)--(2.822148,-0.754235)--(2.956415,-0.612293)--(2.942213,-0.474190)--cycle;
\path[fill=blue!41!white,draw=blue!41!white,line width=.12pt] (-1.246659,-1.612387)--(-1.163283,-1.510570)--(-0.987036,-1.563389)--(-1.057780,-1.668992)--cycle;
\path[fill=blue!40!white,draw=blue!40!white,line width=.12pt] (1.057780,-1.668992)--(0.987036,-1.563389)--(1.163283,-1.510570)--(1.246659,-1.612387)--cycle;
\path[fill=blue!34!white,draw=blue!34!white,line width=.12pt] (-0.888912,-0.665176)--(-0.897452,-0.536404)--(-0.761481,-0.577153)--(-0.754235,-0.705537)--cycle;
\path[fill=blue!29!white,draw=blue!29!white,line width=.12pt] (0.754235,-0.705537)--(0.761481,-0.577153)--(0.897452,-0.536404)--(0.888912,-0.665176)--cycle;
\path[fill=blue!42!white,draw=blue!42!white,line width=.12pt] (-0.987036,-1.563389)--(-0.921837,-1.438377)--(-0.748354,-1.479402)--(-0.801283,-1.607317)--cycle;
\path[fill=blue!40!white,draw=blue!40!white,line width=.12pt] (0.801283,-1.607317)--(0.748354,-1.479402)--(0.921837,-1.438377)--(0.987036,-1.563389)--cycle;
\path[fill=blue!42!white,draw=blue!42!white,line width=.12pt] (-0.748354,-1.479402)--(-0.701962,-1.337241)--(-0.532473,-1.367563)--(-0.567664,-1.511728)--cycle;
\path[fill=blue!40!white,draw=blue!40!white,line width=.12pt] (0.567664,-1.511728)--(0.532473,-1.367563)--(0.701962,-1.337241)--(0.748354,-1.479402)--cycle;
\path[fill=blue!40!white,draw=blue!40!white,line width=.12pt] (-0.503593,-1.214972)--(-0.482133,-1.059820)--(-0.324025,-1.079622)--(-0.338447,-1.235655)--cycle;
\path[fill=blue!38!white,draw=blue!38!white,line width=.12pt] (0.338447,-1.235655)--(0.324025,-1.079622)--(0.482133,-1.059820)--(0.503593,-1.214972)--cycle;
\path[fill=blue!41!white,draw=blue!41!white,line width=.12pt] (-0.532473,-1.367563)--(-0.503593,-1.214972)--(-0.338447,-1.235655)--(-0.357857,-1.389432)--cycle;
\path[fill=blue!39!white,draw=blue!39!white,line width=.12pt] (0.357857,-1.389432)--(0.338447,-1.235655)--(0.503593,-1.214972)--(0.532473,-1.367563)--cycle;
\path[fill=blue!19!white,draw=blue!19!white,line width=.12pt] (-1.724636,0.115141)--(-1.846614,0.146565)--(-1.740975,0.058441)--(-1.625975,0.032838)--cycle;
\path[fill=blue!16!white,draw=blue!16!white,line width=.12pt] (1.625975,0.032838)--(1.740975,0.058441)--(1.846614,0.146565)--(1.724636,0.115141)--cycle;
\path[fill=blue!31!white,draw=blue!31!white,line width=.12pt] (-1.024781,-0.489187)--(-1.053662,-0.376814)--(-0.922745,-0.425362)--(-0.897452,-0.536404)--cycle;
\path[fill=blue!26!white,draw=blue!26!white,line width=.12pt] (0.897452,-0.536404)--(0.922745,-0.425362)--(1.053662,-0.376814)--(1.024781,-0.489187)--cycle;
\path[fill=blue!15!white,draw=blue!15!white,line width=.12pt] (-2.073119,0.250151)--(-2.217311,0.233600)--(-2.116611,0.127144)--(-1.978968,0.150618)--cycle;
\path[fill=blue!14!white,draw=blue!14!white,line width=.12pt] (1.978968,0.150618)--(2.116611,0.127144)--(2.217311,0.233600)--(2.073119,0.250151)--cycle;
\path[fill=blue!38!white,draw=blue!38!white,line width=.12pt] (-0.482133,-1.059820)--(-0.468918,-0.908070)--(-0.315143,-0.927329)--(-0.324025,-1.079622)--cycle;
\path[fill=blue!36!white,draw=blue!36!white,line width=.12pt] (0.324025,-1.079622)--(0.315143,-0.927329)--(0.468918,-0.908070)--(0.482133,-1.059820)--cycle;
\path[fill=blue!34!white,draw=blue!34!white,line width=.12pt] (-2.198794,-1.478311)--(-2.091880,-1.498463)--(-1.913535,-1.596850)--(-2.011334,-1.581726)--cycle;
\path[fill=blue!38!white,draw=blue!38!white,line width=.12pt] (2.011334,-1.581726)--(1.913535,-1.596850)--(2.091880,-1.498463)--(2.198794,-1.478311)--cycle;
\path[fill=blue!39!white,draw=blue!39!white,line width=.12pt] (-0.338447,-1.235655)--(-0.324025,-1.079622)--(-0.162796,-1.091580)--(-0.170043,-1.248146)--cycle;
\path[fill=blue!38!white,draw=blue!38!white,line width=.12pt] (0.170043,-1.248146)--(0.162796,-1.091580)--(0.324025,-1.079622)--(0.338447,-1.235655)--cycle;
\path[fill=blue!36!white,draw=blue!36!white,line width=.12pt] (-0.618176,-0.881367)--(-0.612293,-0.739104)--(-0.464455,-0.765552)--(-0.468918,-0.908070)--cycle;
\path[fill=blue!33!white,draw=blue!33!white,line width=.12pt] (0.468918,-0.908070)--(0.464455,-0.765552)--(0.612293,-0.739104)--(0.618176,-0.881367)--cycle;
\path[fill=blue!41!white,draw=blue!41!white,line width=.12pt] (-0.357857,-1.389432)--(-0.338447,-1.235655)--(-0.170043,-1.248146)--(-0.179794,-1.402639)--cycle;
\path[fill=blue!40!white,draw=blue!40!white,line width=.12pt] (0.179794,-1.402639)--(0.170043,-1.248146)--(0.338447,-1.235655)--(0.357857,-1.389432)--cycle;
\path[fill=blue!28!white,draw=blue!28!white,line width=.12pt] (-1.174431,-0.322085)--(-1.226706,-0.228443)--(-1.100561,-0.285608)--(-1.053662,-0.376814)--cycle;
\path[fill=blue!23!white,draw=blue!23!white,line width=.12pt] (1.053662,-0.376814)--(1.100561,-0.285608)--(1.226706,-0.228443)--(1.174431,-0.322085)--cycle;
\path[fill=blue!39!white,draw=blue!39!white,line width=.12pt] (-0.170043,-1.248146)--(-0.162796,-1.091580)--(-0.000000,-1.095579)--(-0.000000,-1.252323)--cycle;
\path[fill=blue!39!white,draw=blue!39!white,line width=.12pt] (-0.000000,-1.252323)--(-0.000000,-1.095579)--(0.162796,-1.091580)--(0.170043,-1.248146)--cycle;
\path[fill=blue!42!white,draw=blue!42!white,line width=.12pt] (-0.567664,-1.511728)--(-0.532473,-1.367563)--(-0.357857,-1.389432)--(-0.381508,-1.535043)--cycle;
\path[fill=blue!40!white,draw=blue!40!white,line width=.12pt] (0.381508,-1.535043)--(0.357857,-1.389432)--(0.532473,-1.367563)--(0.567664,-1.511728)--cycle;
\path[fill=blue!22!white,draw=blue!22!white,line width=.12pt] (-1.525177,-0.019647)--(-1.625975,0.032838)--(-1.511656,-0.044233)--(-1.417944,-0.091940)--cycle;
\path[fill=blue!18!white,draw=blue!18!white,line width=.12pt] (1.417944,-0.091940)--(1.511656,-0.044233)--(1.625975,0.032838)--(1.525177,-0.019647)--cycle;
\path[fill=blue!17!white,draw=blue!17!white,line width=.12pt] (-2.900154,-0.335511)--(-2.942213,-0.474190)--(-2.808591,-0.615449)--(-2.768442,-0.474751)--cycle;
\path[fill=blue!24!white,draw=blue!24!white,line width=.12pt] (2.768442,-0.474751)--(2.808591,-0.615449)--(2.942213,-0.474190)--(2.900154,-0.335511)--cycle;
\path[fill=blue!38!white,draw=blue!38!white,line width=.12pt] (-0.324025,-1.079622)--(-0.315143,-0.927329)--(-0.158334,-0.938959)--(-0.162796,-1.091580)--cycle;
\path[fill=blue!36!white,draw=blue!36!white,line width=.12pt] (0.162796,-1.091580)--(0.158334,-0.938959)--(0.315143,-0.927329)--(0.324025,-1.079622)--cycle;
\path[fill=blue!25!white,draw=blue!25!white,line width=.12pt] (-1.341037,-0.165370)--(-1.417944,-0.091940)--(-1.297056,-0.158630)--(-1.226706,-0.228443)--cycle;
\path[fill=blue!20!white,draw=blue!20!white,line width=.12pt] (1.226706,-0.228443)--(1.297056,-0.158630)--(1.417944,-0.091940)--(1.341037,-0.165370)--cycle;
\path[fill=blue!40!white,draw=blue!40!white,line width=.12pt] (-0.179794,-1.402639)--(-0.170043,-1.248146)--(-0.000000,-1.252323)--(-0.000000,-1.407055)--cycle;
\path[fill=blue!40!white,draw=blue!40!white,line width=.12pt] (-0.000000,-1.407055)--(-0.000000,-1.252323)--(0.170043,-1.248146)--(0.179794,-1.402639)--cycle;
\path[fill=blue!24!white,draw=blue!24!white,line width=.12pt] (-2.808591,-0.885774)--(-2.768442,-1.005012)--(-2.610069,-1.137127)--(-2.647922,-1.019805)--cycle;
\path[fill=blue!31!white,draw=blue!31!white,line width=.12pt] (2.647922,-1.019805)--(2.610069,-1.137127)--(2.768442,-1.005012)--(2.808591,-0.885774)--cycle;
\path[fill=blue!34!white,draw=blue!34!white,line width=.12pt] (-0.754235,-0.705537)--(-0.761481,-0.577153)--(-0.618176,-0.611042)--(-0.612293,-0.739104)--cycle;
\path[fill=blue!30!white,draw=blue!30!white,line width=.12pt] (0.612293,-0.739104)--(0.618176,-0.611042)--(0.761481,-0.577153)--(0.754235,-0.705537)--cycle;
\path[fill=blue!42!white,draw=blue!42!white,line width=.12pt] (-0.801283,-1.607317)--(-0.748354,-1.479402)--(-0.567664,-1.511728)--(-0.607814,-1.641929)--cycle;
\path[fill=blue!40!white,draw=blue!40!white,line width=.12pt] (0.607814,-1.641929)--(0.567664,-1.511728)--(0.748354,-1.479402)--(0.801283,-1.607317)--cycle;
\path[fill=blue!37!white,draw=blue!37!white,line width=.12pt] (-0.162796,-1.091580)--(-0.158334,-0.938959)--(-0.000000,-0.942848)--(-0.000000,-1.095579)--cycle;
\path[fill=blue!37!white,draw=blue!37!white,line width=.12pt] (-0.000000,-1.095579)--(-0.000000,-0.942848)--(0.158334,-0.938959)--(0.162796,-1.091580)--cycle;
\path[fill=blue!29!white,draw=blue!29!white,line width=.12pt] (-2.548600,-1.236371)--(-2.465877,-1.313721)--(-2.292506,-1.430603)--(-2.369413,-1.357174)--cycle;
\path[fill=blue!35!white,draw=blue!35!white,line width=.12pt] (2.369413,-1.357174)--(2.292506,-1.430603)--(2.465877,-1.313721)--(2.548600,-1.236371)--cycle;
\path[fill=blue!37!white,draw=blue!37!white,line width=.12pt] (-1.913535,-1.596850)--(-1.807416,-1.583216)--(-1.621555,-1.667443)--(-1.716761,-1.686022)--cycle;
\path[fill=blue!38!white,draw=blue!38!white,line width=.12pt] (1.716761,-1.686022)--(1.621555,-1.667443)--(1.807416,-1.583216)--(1.913535,-1.596850)--cycle;
\path[fill=blue!36!white,draw=blue!36!white,line width=.12pt] (-0.468918,-0.908070)--(-0.464455,-0.765552)--(-0.312145,-0.784628)--(-0.315143,-0.927329)--cycle;
\path[fill=blue!34!white,draw=blue!34!white,line width=.12pt] (0.315143,-0.927329)--(0.312145,-0.784628)--(0.464455,-0.765552)--(0.468918,-0.908070)--cycle;
\path[fill=blue!14!white,draw=blue!14!white,line width=.12pt] (-2.217311,0.233600)--(-2.361503,0.190425)--(-2.254254,0.077046)--(-2.116611,0.127144)--cycle;
\path[fill=blue!15!white,draw=blue!15!white,line width=.12pt] (2.116611,0.127144)--(2.254254,0.077046)--(2.361503,0.190425)--(2.217311,0.233600)--cycle;
\path[fill=blue!41!white,draw=blue!41!white,line width=.12pt] (-0.381508,-1.535043)--(-0.357857,-1.389432)--(-0.179794,-1.402639)--(-0.191677,-1.549123)--cycle;
\path[fill=blue!41!white,draw=blue!41!white,line width=.12pt] (0.191677,-1.549123)--(0.179794,-1.402639)--(0.357857,-1.389432)--(0.381508,-1.535043)--cycle;
\path[fill=blue!41!white,draw=blue!41!white,line width=.12pt] (-1.057780,-1.668992)--(-0.987036,-1.563389)--(-0.801283,-1.607317)--(-0.858714,-1.716067)--cycle;
\path[fill=blue!40!white,draw=blue!40!white,line width=.12pt] (0.858714,-1.716067)--(0.801283,-1.607317)--(0.987036,-1.563389)--(1.057780,-1.668992)--cycle;
\path[fill=blue!39!white,draw=blue!39!white,line width=.12pt] (-1.621555,-1.667443)--(-1.522544,-1.620433)--(-1.333369,-1.690584)--(-1.420078,-1.742156)--cycle;
\path[fill=blue!39!white,draw=blue!39!white,line width=.12pt] (1.420078,-1.742156)--(1.333369,-1.690584)--(1.522544,-1.620433)--(1.621555,-1.667443)--cycle;
\path[fill=blue!40!white,draw=blue!40!white,line width=.12pt] (-1.333369,-1.690584)--(-1.246659,-1.612387)--(-1.057780,-1.668992)--(-1.131352,-1.751126)--cycle;
\path[fill=blue!40!white,draw=blue!40!white,line width=.12pt] (1.131352,-1.751126)--(1.057780,-1.668992)--(1.246659,-1.612387)--(1.333369,-1.690584)--cycle;
\path[fill=blue!41!white,draw=blue!41!white,line width=.12pt] (-0.000000,-1.553831)--(-0.000000,-1.407055)--(0.179794,-1.402639)--(0.191677,-1.549123)--cycle;
\path[fill=blue!41!white,draw=blue!41!white,line width=.12pt] (-0.191677,-1.549123)--(-0.179794,-1.402639)--(-0.000000,-1.407055)--(-0.000000,-1.553831)--cycle;
\path[fill=blue!31!white,draw=blue!31!white,line width=.12pt] (-0.897452,-0.536404)--(-0.922745,-0.425362)--(-0.782941,-0.467259)--(-0.761481,-0.577153)--cycle;
\path[fill=blue!27!white,draw=blue!27!white,line width=.12pt] (0.761481,-0.577153)--(0.782941,-0.467259)--(0.922745,-0.425362)--(0.897452,-0.536404)--cycle;
\path[fill=blue!15!white,draw=blue!15!white,line width=.12pt] (-2.831853,-0.201586)--(-2.900154,-0.335511)--(-2.768442,-0.474751)--(-2.703244,-0.337547)--cycle;
\path[fill=blue!22!white,draw=blue!22!white,line width=.12pt] (2.703244,-0.337547)--(2.768442,-0.474751)--(2.900154,-0.335511)--(2.831853,-0.201586)--cycle;
\path[fill=blue!36!white,draw=blue!36!white,line width=.12pt] (-0.315143,-0.927329)--(-0.312145,-0.784628)--(-0.156827,-0.796148)--(-0.158334,-0.938959)--cycle;
\path[fill=blue!34!white,draw=blue!34!white,line width=.12pt] (0.158334,-0.938959)--(0.156827,-0.796148)--(0.312145,-0.784628)--(0.315143,-0.927329)--cycle;
\path[fill=blue!17!white,draw=blue!17!white,line width=.12pt] (-1.846614,0.146565)--(-1.978968,0.150618)--(-1.865758,0.056178)--(-1.740975,0.058441)--cycle;
\path[fill=blue!15!white,draw=blue!15!white,line width=.12pt] (1.740975,0.058441)--(1.865758,0.056178)--(1.978968,0.150618)--(1.846614,0.146565)--cycle;
\path[fill=blue!42!white,draw=blue!42!white,line width=.12pt] (-0.607814,-1.641929)--(-0.567664,-1.511728)--(-0.381508,-1.535043)--(-0.408491,-1.666893)--cycle;
\path[fill=blue!41!white,draw=blue!41!white,line width=.12pt] (0.408491,-1.666893)--(0.381508,-1.535043)--(0.567664,-1.511728)--(0.607814,-1.641929)--cycle;
\path[fill=blue!34!white,draw=blue!34!white,line width=.12pt] (-0.612293,-0.739104)--(-0.618176,-0.611042)--(-0.468918,-0.637745)--(-0.464455,-0.765552)--cycle;
\path[fill=blue!31!white,draw=blue!31!white,line width=.12pt] (0.464455,-0.765552)--(0.468918,-0.637745)--(0.618176,-0.611042)--(0.612293,-0.739104)--cycle;
\path[fill=blue!35!white,draw=blue!35!white,line width=.12pt] (-0.158334,-0.938959)--(-0.156827,-0.796148)--(-0.000000,-0.800000)--(-0.000000,-0.942848)--cycle;
\path[fill=blue!35!white,draw=blue!35!white,line width=.12pt] (-0.000000,-0.942848)--(-0.000000,-0.800000)--(0.156827,-0.796148)--(0.158334,-0.938959)--cycle;
\path[fill=blue!13!white,draw=blue!13!white,line width=.12pt] (-2.361503,0.190425)--(-2.500154,0.122284)--(-2.386608,0.002248)--(-2.254254,0.077046)--cycle;
\path[fill=blue!15!white,draw=blue!15!white,line width=.12pt] (2.254254,0.077046)--(2.386608,0.002248)--(2.500154,0.122284)--(2.361503,0.190425)--cycle;
\path[fill=blue!28!white,draw=blue!28!white,line width=.12pt] (-1.053662,-0.376814)--(-1.100561,-0.285608)--(-0.963817,-0.336316)--(-0.922745,-0.425362)--cycle;
\path[fill=blue!23!white,draw=blue!23!white,line width=.12pt] (0.922745,-0.425362)--(0.963817,-0.336316)--(1.100561,-0.285608)--(1.053662,-0.376814)--cycle;
\path[fill=blue!14!white,draw=blue!14!white,line width=.12pt] (-2.739936,-0.077561)--(-2.831853,-0.201586)--(-2.703244,-0.337547)--(-2.615501,-0.209109)--cycle;
\path[fill=blue!20!white,draw=blue!20!white,line width=.12pt] (2.615501,-0.209109)--(2.703244,-0.337547)--(2.831853,-0.201586)--(2.739936,-0.077561)--cycle;
\path[fill=blue!13!white,draw=blue!13!white,line width=.12pt] (-2.500154,0.122284)--(-2.627935,0.031796)--(-2.508586,-0.094375)--(-2.386608,0.002248)--cycle;
\path[fill=blue!16!white,draw=blue!16!white,line width=.12pt] (2.386608,0.002248)--(2.508586,-0.094375)--(2.627935,0.031796)--(2.500154,0.122284)--cycle;
\path[fill=blue!32!white,draw=blue!32!white,line width=.12pt] (-2.292506,-1.430603)--(-2.198794,-1.478311)--(-2.011334,-1.581726)--(-2.097056,-1.538426)--cycle;
\path[fill=blue!36!white,draw=blue!36!white,line width=.12pt] (2.097056,-1.538426)--(2.011334,-1.581726)--(2.198794,-1.478311)--(2.292506,-1.430603)--cycle;
\path[fill=blue!21!white,draw=blue!21!white,line width=.12pt] (-2.822148,-0.754235)--(-2.808591,-0.885774)--(-2.647922,-1.019805)--(-2.660703,-0.888912)--cycle;
\path[fill=blue!29!white,draw=blue!29!white,line width=.12pt] (2.660703,-0.888912)--(2.647922,-1.019805)--(2.808591,-0.885774)--(2.822148,-0.754235)--cycle;
\path[fill=blue!42!white,draw=blue!42!white,line width=.12pt] (-0.408491,-1.666893)--(-0.381508,-1.535043)--(-0.191677,-1.549123)--(-0.205234,-1.681968)--cycle;
\path[fill=blue!41!white,draw=blue!41!white,line width=.12pt] (0.205234,-1.681968)--(0.191677,-1.549123)--(0.381508,-1.535043)--(0.408491,-1.666893)--cycle;
\path[fill=blue!13!white,draw=blue!13!white,line width=.12pt] (-2.627935,0.031796)--(-2.739936,-0.077561)--(-2.615501,-0.209109)--(-2.508586,-0.094375)--cycle;
\path[fill=blue!18!white,draw=blue!18!white,line width=.12pt] (2.508586,-0.094375)--(2.615501,-0.209109)--(2.739936,-0.077561)--(2.627935,0.031796)--cycle;
\path[fill=blue!41!white,draw=blue!41!white,line width=.12pt] (-0.858714,-1.716067)--(-0.801283,-1.607317)--(-0.607814,-1.641929)--(-0.651378,-1.753160)--cycle;
\path[fill=blue!40!white,draw=blue!40!white,line width=.12pt] (0.651378,-1.753160)--(0.607814,-1.641929)--(0.801283,-1.607317)--(0.858714,-1.716067)--cycle;
\path[fill=blue!33!white,draw=blue!33!white,line width=.12pt] (-0.464455,-0.765552)--(-0.468918,-0.637745)--(-0.315143,-0.657004)--(-0.312145,-0.784628)--cycle;
\path[fill=blue!31!white,draw=blue!31!white,line width=.12pt] (0.312145,-0.784628)--(0.315143,-0.657004)--(0.468918,-0.637745)--(0.464455,-0.765552)--cycle;
\path[fill=blue!19!white,draw=blue!19!white,line width=.12pt] (-1.625975,0.032838)--(-1.740975,0.058441)--(-1.618570,-0.024081)--(-1.511656,-0.044233)--cycle;
\path[fill=blue!16!white,draw=blue!16!white,line width=.12pt] (1.511656,-0.044233)--(1.618570,-0.024081)--(1.740975,0.058441)--(1.625975,0.032838)--cycle;
\path[fill=blue!25!white,draw=blue!25!white,line width=.12pt] (-1.226706,-0.228443)--(-1.297056,-0.158630)--(-1.163677,-0.219074)--(-1.100561,-0.285608)--cycle;
\path[fill=blue!21!white,draw=blue!21!white,line width=.12pt] (1.100561,-0.285608)--(1.163677,-0.219074)--(1.297056,-0.158630)--(1.226706,-0.228443)--cycle;
\path[fill=blue!31!white,draw=blue!31!white,line width=.12pt] (-0.761481,-0.577153)--(-0.782941,-0.467259)--(-0.635598,-0.502103)--(-0.618176,-0.611042)--cycle;
\path[fill=blue!27!white,draw=blue!27!white,line width=.12pt] (0.618176,-0.611042)--(0.635598,-0.502103)--(0.782941,-0.467259)--(0.761481,-0.577153)--cycle;
\path[fill=blue!41!white,draw=blue!41!white,line width=.12pt] (-0.205234,-1.681968)--(-0.191677,-1.549123)--(-0.000000,-1.553831)--(-0.000000,-1.687009)--cycle;
\path[fill=blue!41!white,draw=blue!41!white,line width=.12pt] (-0.000000,-1.687009)--(-0.000000,-1.553831)--(0.191677,-1.549123)--(0.205234,-1.681968)--cycle;
\path[fill=blue!22!white,draw=blue!22!white,line width=.12pt] (-1.417944,-0.091940)--(-1.511656,-0.044233)--(-1.382778,-0.115330)--(-1.297056,-0.158630)--cycle;
\path[fill=blue!18!white,draw=blue!18!white,line width=.12pt] (1.297056,-0.158630)--(1.382778,-0.115330)--(1.511656,-0.044233)--(1.417944,-0.091940)--cycle;
\path[fill=blue!26!white,draw=blue!26!white,line width=.12pt] (-2.610069,-1.137127)--(-2.548600,-1.236371)--(-2.369413,-1.357174)--(-2.426560,-1.260844)--cycle;
\path[fill=blue!33!white,draw=blue!33!white,line width=.12pt] (2.426560,-1.260844)--(2.369413,-1.357174)--(2.548600,-1.236371)--(2.610069,-1.137127)--cycle;
\path[fill=blue!33!white,draw=blue!33!white,line width=.12pt] (-0.312145,-0.784628)--(-0.315143,-0.657004)--(-0.158334,-0.668634)--(-0.156827,-0.796148)--cycle;
\path[fill=blue!32!white,draw=blue!32!white,line width=.12pt] (0.156827,-0.796148)--(0.158334,-0.668634)--(0.315143,-0.657004)--(0.312145,-0.784628)--cycle;
\path[fill=blue!40!white,draw=blue!40!white,line width=.12pt] (-1.131352,-1.751126)--(-1.057780,-1.668992)--(-0.858714,-1.716067)--(-0.918440,-1.801476)--cycle;
\path[fill=blue!40!white,draw=blue!40!white,line width=.12pt] (0.918440,-1.801476)--(0.858714,-1.716067)--(1.057780,-1.668992)--(1.131352,-1.751126)--cycle;
\path[fill=blue!34!white,draw=blue!34!white,line width=.12pt] (-2.011334,-1.581726)--(-1.913535,-1.596850)--(-1.716761,-1.686022)--(-1.804504,-1.675456)--cycle;
\path[fill=blue!37!white,draw=blue!37!white,line width=.12pt] (1.804504,-1.675456)--(1.716761,-1.686022)--(1.913535,-1.596850)--(2.011334,-1.581726)--cycle;
\path[fill=blue!33!white,draw=blue!33!white,line width=.12pt] (-0.156827,-0.796148)--(-0.158334,-0.668634)--(-0.000000,-0.672523)--(-0.000000,-0.800000)--cycle;
\path[fill=blue!32!white,draw=blue!32!white,line width=.12pt] (-0.000000,-0.800000)--(-0.000000,-0.672523)--(0.158334,-0.668634)--(0.156827,-0.796148)--cycle;
\path[fill=blue!15!white,draw=blue!15!white,line width=.12pt] (-1.978968,0.150618)--(-2.116611,0.127144)--(-1.995527,0.026136)--(-1.865758,0.056178)--cycle;
\path[fill=blue!14!white,draw=blue!14!white,line width=.12pt] (1.865758,0.056178)--(1.995527,0.026136)--(2.116611,0.127144)--(1.978968,0.150618)--cycle;
\path[fill=blue!41!white,draw=blue!41!white,line width=.12pt] (-0.651378,-1.753160)--(-0.607814,-1.641929)--(-0.408491,-1.666893)--(-0.437769,-1.779914)--cycle;
\path[fill=blue!41!white,draw=blue!41!white,line width=.12pt] (0.437769,-1.779914)--(0.408491,-1.666893)--(0.607814,-1.641929)--(0.651378,-1.753160)--cycle;
\path[fill=blue!39!white,draw=blue!39!white,line width=.12pt] (-1.420078,-1.742156)--(-1.333369,-1.690584)--(-1.131352,-1.751126)--(-1.204924,-1.806635)--cycle;
\path[fill=blue!39!white,draw=blue!39!white,line width=.12pt] (1.204924,-1.806635)--(1.131352,-1.751126)--(1.333369,-1.690584)--(1.420078,-1.742156)--cycle;
\path[fill=blue!37!white,draw=blue!37!white,line width=.12pt] (-1.716761,-1.686022)--(-1.621555,-1.667443)--(-1.420078,-1.742156)--(-1.503455,-1.765122)--cycle;
\path[fill=blue!38!white,draw=blue!38!white,line width=.12pt] (1.503455,-1.765122)--(1.420078,-1.742156)--(1.621555,-1.667443)--(1.716761,-1.686022)--cycle;
\path[fill=blue!28!white,draw=blue!28!white,line width=.12pt] (-0.922745,-0.425362)--(-0.963817,-0.336316)--(-0.817791,-0.380079)--(-0.782941,-0.467259)--cycle;
\path[fill=blue!24!white,draw=blue!24!white,line width=.12pt] (0.782941,-0.467259)--(0.817791,-0.380079)--(0.963817,-0.336316)--(0.922745,-0.425362)--cycle;
\path[fill=blue!31!white,draw=blue!31!white,line width=.12pt] (-0.618176,-0.611042)--(-0.635598,-0.502103)--(-0.482133,-0.529559)--(-0.468918,-0.637745)--cycle;
\path[fill=blue!28!white,draw=blue!28!white,line width=.12pt] (0.468918,-0.637745)--(0.482133,-0.529559)--(0.635598,-0.502103)--(0.618176,-0.611042)--cycle;
\path[fill=blue!18!white,draw=blue!18!white,line width=.12pt] (-2.808591,-0.615449)--(-2.822148,-0.754235)--(-2.660703,-0.888912)--(-2.647922,-0.749480)--cycle;
\path[fill=blue!26!white,draw=blue!26!white,line width=.12pt] (2.647922,-0.749480)--(2.660703,-0.888912)--(2.822148,-0.754235)--(2.808591,-0.615449)--cycle;
\path[fill=blue!41!white,draw=blue!41!white,line width=.12pt] (-0.437769,-1.779914)--(-0.408491,-1.666893)--(-0.205234,-1.681968)--(-0.219943,-1.796069)--cycle;
\path[fill=blue!41!white,draw=blue!41!white,line width=.12pt] (0.219943,-1.796069)--(0.205234,-1.681968)--(0.408491,-1.666893)--(0.437769,-1.779914)--cycle;
\path[fill=blue!31!white,draw=blue!31!white,line width=.12pt] (-0.468918,-0.637745)--(-0.482133,-0.529559)--(-0.324025,-0.549360)--(-0.315143,-0.657004)--cycle;
\path[fill=blue!29!white,draw=blue!29!white,line width=.12pt] (0.315143,-0.657004)--(0.324025,-0.549360)--(0.482133,-0.529559)--(0.468918,-0.637745)--cycle;
\path[fill=blue!25!white,draw=blue!25!white,line width=.12pt] (-1.100561,-0.285608)--(-1.163677,-0.219074)--(-1.019091,-0.272690)--(-0.963817,-0.336316)--cycle;
\path[fill=blue!21!white,draw=blue!21!white,line width=.12pt] (0.963817,-0.336316)--(1.019091,-0.272690)--(1.163677,-0.219074)--(1.100561,-0.285608)--cycle;
\path[fill=blue!41!white,draw=blue!41!white,line width=.12pt] (-0.219943,-1.796069)--(-0.205234,-1.681968)--(-0.000000,-1.687009)--(-0.000000,-1.801472)--cycle;
\path[fill=blue!41!white,draw=blue!41!white,line width=.12pt] (-0.000000,-1.801472)--(-0.000000,-1.687009)--(0.205234,-1.681968)--(0.219943,-1.796069)--cycle;
\path[fill=blue!40!white,draw=blue!40!white,line width=.12pt] (-0.918440,-1.801476)--(-0.858714,-1.716067)--(-0.651378,-1.753160)--(-0.696683,-1.841149)--cycle;
\path[fill=blue!40!white,draw=blue!40!white,line width=.12pt] (0.696683,-1.841149)--(0.651378,-1.753160)--(0.858714,-1.716067)--(0.918440,-1.801476)--cycle;
\path[fill=blue!14!white,draw=blue!14!white,line width=.12pt] (-2.116611,0.127144)--(-2.254254,0.077046)--(-2.125296,-0.030531)--(-1.995527,0.026136)--cycle;
\path[fill=blue!14!white,draw=blue!14!white,line width=.12pt] (1.995527,0.026136)--(2.125296,-0.030531)--(2.254254,0.077046)--(2.116611,0.127144)--cycle;
\path[fill=blue!17!white,draw=blue!17!white,line width=.12pt] (-1.740975,0.058441)--(-1.865758,0.056178)--(-1.734580,-0.032258)--(-1.618570,-0.024081)--cycle;
\path[fill=blue!15!white,draw=blue!15!white,line width=.12pt] (1.618570,-0.024081)--(1.734580,-0.032258)--(1.865758,0.056178)--(1.740975,0.058441)--cycle;
\path[fill=blue!29!white,draw=blue!29!white,line width=.12pt] (-2.369413,-1.357174)--(-2.292506,-1.430603)--(-2.097056,-1.538426)--(-2.167407,-1.468614)--cycle;
\path[fill=blue!34!white,draw=blue!34!white,line width=.12pt] (2.167407,-1.468614)--(2.097056,-1.538426)--(2.292506,-1.430603)--(2.369413,-1.357174)--cycle;
\path[fill=blue!23!white,draw=blue!23!white,line width=.12pt] (-2.647922,-1.019805)--(-2.610069,-1.137127)--(-2.426560,-1.260844)--(-2.461751,-1.145316)--cycle;
\path[fill=blue!30!white,draw=blue!30!white,line width=.12pt] (2.461751,-1.145316)--(2.426560,-1.260844)--(2.610069,-1.137127)--(2.647922,-1.019805)--cycle;
\path[fill=blue!28!white,draw=blue!28!white,line width=.12pt] (-0.782941,-0.467259)--(-0.817791,-0.380079)--(-0.663889,-0.416474)--(-0.635598,-0.502103)--cycle;
\path[fill=blue!25!white,draw=blue!25!white,line width=.12pt] (0.635598,-0.502103)--(0.663889,-0.416474)--(0.817791,-0.380079)--(0.782941,-0.467259)--cycle;
\path[fill=blue!30!white,draw=blue!30!white,line width=.12pt] (-0.315143,-0.657004)--(-0.324025,-0.549360)--(-0.162796,-0.561318)--(-0.158334,-0.668634)--cycle;
\path[fill=blue!29!white,draw=blue!29!white,line width=.12pt] (0.158334,-0.668634)--(0.162796,-0.561318)--(0.324025,-0.549360)--(0.315143,-0.657004)--cycle;
\path[fill=blue!16!white,draw=blue!16!white,line width=.12pt] (-2.768442,-0.474751)--(-2.808591,-0.615449)--(-2.647922,-0.749480)--(-2.610069,-0.606865)--cycle;
\path[fill=blue!24!white,draw=blue!24!white,line width=.12pt] (2.610069,-0.606865)--(2.647922,-0.749480)--(2.808591,-0.615449)--(2.768442,-0.474751)--cycle;
\path[fill=blue!22!white,draw=blue!22!white,line width=.12pt] (-1.297056,-0.158630)--(-1.382778,-0.115330)--(-1.240584,-0.179769)--(-1.163677,-0.219074)--cycle;
\path[fill=blue!19!white,draw=blue!19!white,line width=.12pt] (1.163677,-0.219074)--(1.240584,-0.179769)--(1.382778,-0.115330)--(1.297056,-0.158630)--cycle;
\path[fill=blue!20!white,draw=blue!20!white,line width=.12pt] (-1.511656,-0.044233)--(-1.618570,-0.024081)--(-1.480578,-0.100206)--(-1.382778,-0.115330)--cycle;
\path[fill=blue!16!white,draw=blue!16!white,line width=.12pt] (1.382778,-0.115330)--(1.480578,-0.100206)--(1.618570,-0.024081)--(1.511656,-0.044233)--cycle;
\path[fill=blue!30!white,draw=blue!30!white,line width=.12pt] (-0.158334,-0.668634)--(-0.162796,-0.561318)--(-0.000000,-0.565317)--(-0.000000,-0.672523)--cycle;
\path[fill=blue!30!white,draw=blue!30!white,line width=.12pt] (-0.000000,-0.672523)--(-0.000000,-0.565317)--(0.162796,-0.561318)--(0.158334,-0.668634)--cycle;
\path[fill=blue!39!white,draw=blue!39!white,line width=.12pt] (-1.204924,-1.806635)--(-1.131352,-1.751126)--(-0.918440,-1.801476)--(-0.978167,-1.860259)--cycle;
\path[fill=blue!39!white,draw=blue!39!white,line width=.12pt] (0.978167,-1.860259)--(0.918440,-1.801476)--(1.131352,-1.751126)--(1.204924,-1.806635)--cycle;
\path[fill=blue!13!white,draw=blue!13!white,line width=.12pt] (-2.254254,0.077046)--(-2.386608,0.002248)--(-2.250079,-0.111645)--(-2.125296,-0.030531)--cycle;
\path[fill=blue!15!white,draw=blue!15!white,line width=.12pt] (2.125296,-0.030531)--(2.250079,-0.111645)--(2.386608,0.002248)--(2.254254,0.077046)--cycle;
\path[fill=blue!40!white,draw=blue!40!white,line width=.12pt] (-0.696683,-1.841149)--(-0.651378,-1.753160)--(-0.437769,-1.779914)--(-0.468217,-1.869763)--cycle;
\path[fill=blue!40!white,draw=blue!40!white,line width=.12pt] (0.468217,-1.869763)--(0.437769,-1.779914)--(0.651378,-1.753160)--(0.696683,-1.841149)--cycle;
\path[fill=blue!32!white,draw=blue!32!white,line width=.12pt] (-2.097056,-1.538426)--(-2.011334,-1.581726)--(-1.804504,-1.675456)--(-1.881411,-1.636151)--cycle;
\path[fill=blue!35!white,draw=blue!35!white,line width=.12pt] (1.881411,-1.636151)--(1.804504,-1.675456)--(2.011334,-1.581726)--(2.097056,-1.538426)--cycle;
\path[fill=blue!15!white,draw=blue!15!white,line width=.12pt] (-2.703244,-0.337547)--(-2.768442,-0.474751)--(-2.610069,-0.606865)--(-2.548600,-0.466550)--cycle;
\path[fill=blue!21!white,draw=blue!21!white,line width=.12pt] (2.548600,-0.466550)--(2.610069,-0.606865)--(2.768442,-0.474751)--(2.703244,-0.337547)--cycle;
\path[fill=blue!28!white,draw=blue!28!white,line width=.12pt] (-0.635598,-0.502103)--(-0.663889,-0.416474)--(-0.503593,-0.445151)--(-0.482133,-0.529559)--cycle;
\path[fill=blue!25!white,draw=blue!25!white,line width=.12pt] (0.482133,-0.529559)--(0.503593,-0.445151)--(0.663889,-0.416474)--(0.635598,-0.502103)--cycle;
\path[fill=blue!37!white,draw=blue!37!white,line width=.12pt] (-1.503455,-1.765122)--(-1.420078,-1.742156)--(-1.204924,-1.806635)--(-1.275669,-1.833387)--cycle;
\path[fill=blue!38!white,draw=blue!38!white,line width=.12pt] (1.275669,-1.833387)--(1.204924,-1.806635)--(1.420078,-1.742156)--(1.503455,-1.765122)--cycle;
\path[fill=blue!25!white,draw=blue!25!white,line width=.12pt] (-0.963817,-0.336316)--(-1.019091,-0.272690)--(-0.864690,-0.318963)--(-0.817791,-0.380079)--cycle;
\path[fill=blue!22!white,draw=blue!22!white,line width=.12pt] (0.817791,-0.380079)--(0.864690,-0.318963)--(1.019091,-0.272690)--(0.963817,-0.336316)--cycle;
\path[fill=blue!34!white,draw=blue!34!white,line width=.12pt] (-1.804504,-1.675456)--(-1.716761,-1.686022)--(-1.503455,-1.765122)--(-1.580295,-1.758598)--cycle;
\path[fill=blue!37!white,draw=blue!37!white,line width=.12pt] (1.580295,-1.758598)--(1.503455,-1.765122)--(1.716761,-1.686022)--(1.804504,-1.675456)--cycle;
\path[fill=blue!12!white,draw=blue!12!white,line width=.12pt] (-2.386608,0.002248)--(-2.508586,-0.094375)--(-2.365079,-0.214089)--(-2.250079,-0.111645)--cycle;
\path[fill=blue!16!white,draw=blue!16!white,line width=.12pt] (2.250079,-0.111645)--(2.365079,-0.214089)--(2.508586,-0.094375)--(2.386608,0.002248)--cycle;
\path[fill=blue!40!white,draw=blue!40!white,line width=.12pt] (-0.468217,-1.869763)--(-0.437769,-1.779914)--(-0.219943,-1.796069)--(-0.235241,-1.887042)--cycle;
\path[fill=blue!40!white,draw=blue!40!white,line width=.12pt] (0.235241,-1.887042)--(0.219943,-1.796069)--(0.437769,-1.779914)--(0.468217,-1.869763)--cycle;
\path[fill=blue!15!white,draw=blue!15!white,line width=.12pt] (-1.865758,0.056178)--(-1.995527,0.026136)--(-1.855225,-0.068452)--(-1.734580,-0.032258)--cycle;
\path[fill=blue!14!white,draw=blue!14!white,line width=.12pt] (1.734580,-0.032258)--(1.855225,-0.068452)--(1.995527,0.026136)--(1.865758,0.056178)--cycle;
\path[fill=blue!13!white,draw=blue!13!white,line width=.12pt] (-2.615501,-0.209109)--(-2.703244,-0.337547)--(-2.548600,-0.466550)--(-2.465877,-0.333925)--cycle;
\path[fill=blue!19!white,draw=blue!19!white,line width=.12pt] (2.465877,-0.333925)--(2.548600,-0.466550)--(2.703244,-0.337547)--(2.615501,-0.209109)--cycle;
\path[fill=blue!13!white,draw=blue!13!white,line width=.12pt] (-2.508586,-0.094375)--(-2.615501,-0.209109)--(-2.465877,-0.333925)--(-2.365079,-0.214089)--cycle;
\path[fill=blue!17!white,draw=blue!17!white,line width=.12pt] (2.365079,-0.214089)--(2.465877,-0.333925)--(2.615501,-0.209109)--(2.508586,-0.094375)--cycle;
\path[fill=blue!21!white,draw=blue!21!white,line width=.12pt] (-2.660703,-0.888912)--(-2.647922,-1.019805)--(-2.461751,-1.145316)--(-2.473633,-1.015029)--cycle;
\path[fill=blue!28!white,draw=blue!28!white,line width=.12pt] (2.473633,-1.015029)--(2.461751,-1.145316)--(2.647922,-1.019805)--(2.660703,-0.888912)--cycle;
\path[fill=blue!40!white,draw=blue!40!white,line width=.12pt] (-0.235241,-1.887042)--(-0.219943,-1.796069)--(-0.000000,-1.801472)--(-0.000000,-1.892820)--cycle;
\path[fill=blue!40!white,draw=blue!40!white,line width=.12pt] (-0.000000,-1.892820)--(-0.000000,-1.801472)--(0.219943,-1.796069)--(0.235241,-1.887042)--cycle;
\path[fill=blue!28!white,draw=blue!28!white,line width=.12pt] (-0.482133,-0.529559)--(-0.503593,-0.445151)--(-0.338447,-0.465835)--(-0.324025,-0.549360)--cycle;
\path[fill=blue!26!white,draw=blue!26!white,line width=.12pt] (0.324025,-0.549360)--(0.338447,-0.465835)--(0.503593,-0.445151)--(0.482133,-0.529559)--cycle;
\path[fill=blue!26!white,draw=blue!26!white,line width=.12pt] (-2.426560,-1.260844)--(-2.369413,-1.357174)--(-2.167407,-1.468614)--(-2.219681,-1.374972)--cycle;
\path[fill=blue!32!white,draw=blue!32!white,line width=.12pt] (2.219681,-1.374972)--(2.167407,-1.468614)--(2.369413,-1.357174)--(2.426560,-1.260844)--cycle;
\path[fill=blue!39!white,draw=blue!39!white,line width=.12pt] (-0.978167,-1.860259)--(-0.918440,-1.801476)--(-0.696683,-1.841149)--(-0.741989,-1.902512)--cycle;
\path[fill=blue!39!white,draw=blue!39!white,line width=.12pt] (0.741989,-1.902512)--(0.696683,-1.841149)--(0.918440,-1.801476)--(0.978167,-1.860259)--cycle;
\path[fill=blue!22!white,draw=blue!22!white,line width=.12pt] (-1.163677,-0.219074)--(-1.240584,-0.179769)--(-1.086442,-0.236929)--(-1.019091,-0.272690)--cycle;
\path[fill=blue!19!white,draw=blue!19!white,line width=.12pt] (1.019091,-0.272690)--(1.086442,-0.236929)--(1.240584,-0.179769)--(1.163677,-0.219074)--cycle;
\path[fill=blue!28!white,draw=blue!28!white,line width=.12pt] (-0.324025,-0.549360)--(-0.338447,-0.465835)--(-0.170043,-0.478325)--(-0.162796,-0.561318)--cycle;
\path[fill=blue!26!white,draw=blue!26!white,line width=.12pt] (0.162796,-0.561318)--(0.170043,-0.478325)--(0.338447,-0.465835)--(0.324025,-0.549360)--cycle;
\path[fill=blue!17!white,draw=blue!17!white,line width=.12pt] (-1.618570,-0.024081)--(-1.734580,-0.032258)--(-1.586697,-0.113840)--(-1.480578,-0.100206)--cycle;
\path[fill=blue!15!white,draw=blue!15!white,line width=.12pt] (1.480578,-0.100206)--(1.586697,-0.113840)--(1.734580,-0.032258)--(1.618570,-0.024081)--cycle;
\path[fill=blue!25!white,draw=blue!25!white,line width=.12pt] (-0.817791,-0.380079)--(-0.864690,-0.318963)--(-0.701962,-0.357445)--(-0.663889,-0.416474)--cycle;
\path[fill=blue!22!white,draw=blue!22!white,line width=.12pt] (0.663889,-0.416474)--(0.701962,-0.357445)--(0.864690,-0.318963)--(0.817791,-0.380079)--cycle;
\path[fill=blue!20!white,draw=blue!20!white,line width=.12pt] (-1.382778,-0.115330)--(-1.480578,-0.100206)--(-1.328326,-0.169203)--(-1.240584,-0.179769)--cycle;
\path[fill=blue!17!white,draw=blue!17!white,line width=.12pt] (1.240584,-0.179769)--(1.328326,-0.169203)--(1.480578,-0.100206)--(1.382778,-0.115330)--cycle;
\path[fill=blue!27!white,draw=blue!27!white,line width=.12pt] (-0.162796,-0.561318)--(-0.170043,-0.478325)--(-0.000000,-0.482502)--(-0.000000,-0.565317)--cycle;
\path[fill=blue!27!white,draw=blue!27!white,line width=.12pt] (-0.000000,-0.565317)--(-0.000000,-0.482502)--(0.170043,-0.478325)--(0.162796,-0.561318)--cycle;
\path[fill=blue!14!white,draw=blue!14!white,line width=.12pt] (-1.995527,0.026136)--(-2.125296,-0.030531)--(-1.975871,-0.131270)--(-1.855225,-0.068452)--cycle;
\path[fill=blue!14!white,draw=blue!14!white,line width=.12pt] (1.855225,-0.068452)--(1.975871,-0.131270)--(2.125296,-0.030531)--(1.995527,0.026136)--cycle;
\path[fill=blue!37!white,draw=blue!37!white,line width=.12pt] (-1.275669,-1.833387)--(-1.204924,-1.806635)--(-0.978167,-1.860259)--(-1.035598,-1.890159)--cycle;
\path[fill=blue!38!white,draw=blue!38!white,line width=.12pt] (1.035598,-1.890159)--(0.978167,-1.860259)--(1.204924,-1.806635)--(1.275669,-1.833387)--cycle;
\path[fill=blue!29!white,draw=blue!29!white,line width=.12pt] (-2.167407,-1.468614)--(-2.097056,-1.538426)--(-1.881411,-1.636151)--(-1.944527,-1.569617)--cycle;
\path[fill=blue!33!white,draw=blue!33!white,line width=.12pt] (1.944527,-1.569617)--(1.881411,-1.636151)--(2.097056,-1.538426)--(2.167407,-1.468614)--cycle;
\path[fill=blue!39!white,draw=blue!39!white,line width=.12pt] (-0.741989,-1.902512)--(-0.696683,-1.841149)--(-0.468217,-1.869763)--(-0.498665,-1.932987)--cycle;
\path[fill=blue!39!white,draw=blue!39!white,line width=.12pt] (0.498665,-1.932987)--(0.468217,-1.869763)--(0.696683,-1.841149)--(0.741989,-1.902512)--cycle;
\path[fill=blue!18!white,draw=blue!18!white,line width=.12pt] (-2.647922,-0.749480)--(-2.660703,-0.888912)--(-2.473633,-1.015029)--(-2.461751,-0.874991)--cycle;
\path[fill=blue!25!white,draw=blue!25!white,line width=.12pt] (2.461751,-0.874991)--(2.473633,-1.015029)--(2.660703,-0.888912)--(2.647922,-0.749480)--cycle;
\path[fill=blue!25!white,draw=blue!25!white,line width=.12pt] (-0.663889,-0.416474)--(-0.701962,-0.357445)--(-0.532473,-0.387767)--(-0.503593,-0.445151)--cycle;
\path[fill=blue!23!white,draw=blue!23!white,line width=.12pt] (0.503593,-0.445151)--(0.532473,-0.387767)--(0.701962,-0.357445)--(0.663889,-0.416474)--cycle;
\path[fill=blue!34!white,draw=blue!34!white,line width=.12pt] (-1.580295,-1.758598)--(-1.503455,-1.765122)--(-1.275669,-1.833387)--(-1.340867,-1.830352)--cycle;
\path[fill=blue!36!white,draw=blue!36!white,line width=.12pt] (1.340867,-1.830352)--(1.275669,-1.833387)--(1.503455,-1.765122)--(1.580295,-1.758598)--cycle;
\path[fill=blue!32!white,draw=blue!32!white,line width=.12pt] (-1.881411,-1.636151)--(-1.804504,-1.675456)--(-1.580295,-1.758598)--(-1.647647,-1.722837)--cycle;
\path[fill=blue!35!white,draw=blue!35!white,line width=.12pt] (1.647647,-1.722837)--(1.580295,-1.758598)--(1.804504,-1.675456)--(1.881411,-1.636151)--cycle;
\path[fill=blue!22!white,draw=blue!22!white,line width=.12pt] (-1.019091,-0.272690)--(-1.086442,-0.236929)--(-0.921837,-0.286259)--(-0.864690,-0.318963)--cycle;
\path[fill=blue!19!white,draw=blue!19!white,line width=.12pt] (0.864690,-0.318963)--(0.921837,-0.286259)--(1.086442,-0.236929)--(1.019091,-0.272690)--cycle;
\path[fill=blue!23!white,draw=blue!23!white,line width=.12pt] (-2.461751,-1.145316)--(-2.426560,-1.260844)--(-2.219681,-1.374972)--(-2.251872,-1.261099)--cycle;
\path[fill=blue!30!white,draw=blue!30!white,line width=.12pt] (2.251872,-1.261099)--(2.219681,-1.374972)--(2.426560,-1.260844)--(2.461751,-1.145316)--cycle;
\path[fill=blue!39!white,draw=blue!39!white,line width=.12pt] (-0.498665,-1.932987)--(-0.468217,-1.869763)--(-0.235241,-1.887042)--(-0.250539,-1.951390)--cycle;
\path[fill=blue!39!white,draw=blue!39!white,line width=.12pt] (0.250539,-1.951390)--(0.235241,-1.887042)--(0.468217,-1.869763)--(0.498665,-1.932987)--cycle;
\path[fill=blue!13!white,draw=blue!13!white,line width=.12pt] (-2.125296,-0.030531)--(-2.250079,-0.111645)--(-2.091880,-0.218298)--(-1.975871,-0.131270)--cycle;
\path[fill=blue!15!white,draw=blue!15!white,line width=.12pt] (1.975871,-0.131270)--(2.091880,-0.218298)--(2.250079,-0.111645)--(2.125296,-0.030531)--cycle;
\path[fill=blue!39!white,draw=blue!39!white,line width=.12pt] (-0.250539,-1.951390)--(-0.235241,-1.887042)--(-0.000000,-1.892820)--(-0.000000,-1.957544)--cycle;
\path[fill=blue!39!white,draw=blue!39!white,line width=.12pt] (-0.000000,-1.957544)--(-0.000000,-1.892820)--(0.235241,-1.887042)--(0.250539,-1.951390)--cycle;
\path[fill=blue!16!white,draw=blue!16!white,line width=.12pt] (-2.610069,-0.606865)--(-2.647922,-0.749480)--(-2.461751,-0.874991)--(-2.426560,-0.730582)--cycle;
\path[fill=blue!23!white,draw=blue!23!white,line width=.12pt] (2.426560,-0.730582)--(2.461751,-0.874991)--(2.647922,-0.749480)--(2.610069,-0.606865)--cycle;
\path[fill=blue!15!white,draw=blue!15!white,line width=.12pt] (-1.734580,-0.032258)--(-1.855225,-0.068452)--(-1.697056,-0.155708)--(-1.586697,-0.113840)--cycle;
\path[fill=blue!15!white,draw=blue!15!white,line width=.12pt] (1.586697,-0.113840)--(1.697056,-0.155708)--(1.855225,-0.068452)--(1.734580,-0.032258)--cycle;
\path[fill=blue!25!white,draw=blue!25!white,line width=.12pt] (-0.503593,-0.445151)--(-0.532473,-0.387767)--(-0.357857,-0.409636)--(-0.338447,-0.465835)--cycle;
\path[fill=blue!23!white,draw=blue!23!white,line width=.12pt] (0.338447,-0.465835)--(0.357857,-0.409636)--(0.532473,-0.387767)--(0.503593,-0.445151)--cycle;
\path[fill=blue!20!white,draw=blue!20!white,line width=.12pt] (-1.240584,-0.179769)--(-1.328326,-0.169203)--(-1.163283,-0.230405)--(-1.086442,-0.236929)--cycle;
\path[fill=blue!17!white,draw=blue!17!white,line width=.12pt] (1.086442,-0.236929)--(1.163283,-0.230405)--(1.328326,-0.169203)--(1.240584,-0.179769)--cycle;
\path[fill=blue!37!white,draw=blue!37!white,line width=.12pt] (-1.035598,-1.890159)--(-0.978167,-1.860259)--(-0.741989,-1.902512)--(-0.785553,-1.934893)--cycle;
\path[fill=blue!38!white,draw=blue!38!white,line width=.12pt] (0.785553,-1.934893)--(0.741989,-1.902512)--(0.978167,-1.860259)--(1.035598,-1.890159)--cycle;
\path[fill=blue!17!white,draw=blue!17!white,line width=.12pt] (-1.480578,-0.100206)--(-1.586697,-0.113840)--(-1.423533,-0.187782)--(-1.328326,-0.169203)--cycle;
\path[fill=blue!16!white,draw=blue!16!white,line width=.12pt] (1.328326,-0.169203)--(1.423533,-0.187782)--(1.586697,-0.113840)--(1.480578,-0.100206)--cycle;
\path[fill=blue!25!white,draw=blue!25!white,line width=.12pt] (-0.338447,-0.465835)--(-0.357857,-0.409636)--(-0.179794,-0.422843)--(-0.170043,-0.478325)--cycle;
\path[fill=blue!24!white,draw=blue!24!white,line width=.12pt] (0.170043,-0.478325)--(0.179794,-0.422843)--(0.357857,-0.409636)--(0.338447,-0.465835)--cycle;
\path[fill=blue!12!white,draw=blue!12!white,line width=.12pt] (-2.250079,-0.111645)--(-2.365079,-0.214089)--(-2.198794,-0.326193)--(-2.091880,-0.218298)--cycle;
\path[fill=blue!16!white,draw=blue!16!white,line width=.12pt] (2.091880,-0.218298)--(2.198794,-0.326193)--(2.365079,-0.214089)--(2.250079,-0.111645)--cycle;
\path[fill=blue!14!white,draw=blue!14!white,line width=.12pt] (-2.548600,-0.466550)--(-2.610069,-0.606865)--(-2.426560,-0.730582)--(-2.369413,-0.587353)--cycle;
\path[fill=blue!21!white,draw=blue!21!white,line width=.12pt] (2.369413,-0.587353)--(2.426560,-0.730582)--(2.610069,-0.606865)--(2.548600,-0.466550)--cycle;
\path[fill=blue!22!white,draw=blue!22!white,line width=.12pt] (-0.864690,-0.318963)--(-0.921837,-0.286259)--(-0.748354,-0.327284)--(-0.701962,-0.357445)--cycle;
\path[fill=blue!20!white,draw=blue!20!white,line width=.12pt] (0.701962,-0.357445)--(0.748354,-0.327284)--(0.921837,-0.286259)--(0.864690,-0.318963)--cycle;
\path[fill=blue!24!white,draw=blue!24!white,line width=.12pt] (-0.170043,-0.478325)--(-0.179794,-0.422843)--(-0.000000,-0.427259)--(-0.000000,-0.482502)--cycle;
\path[fill=blue!24!white,draw=blue!24!white,line width=.12pt] (-0.000000,-0.482502)--(-0.000000,-0.427259)--(0.179794,-0.422843)--(0.170043,-0.478325)--cycle;
\path[fill=blue!26!white,draw=blue!26!white,line width=.12pt] (-2.219681,-1.374972)--(-2.167407,-1.468614)--(-1.944527,-1.569617)--(-1.991426,-1.478411)--cycle;
\path[fill=blue!31!white,draw=blue!31!white,line width=.12pt] (1.991426,-1.478411)--(1.944527,-1.569617)--(2.167407,-1.468614)--(2.219681,-1.374972)--cycle;
\path[fill=blue!12!white,draw=blue!12!white,line width=.12pt] (-2.365079,-0.214089)--(-2.465877,-0.333925)--(-2.292506,-0.450808)--(-2.198794,-0.326193)--cycle;
\path[fill=blue!17!white,draw=blue!17!white,line width=.12pt] (2.198794,-0.326193)--(2.292506,-0.450808)--(2.465877,-0.333925)--(2.365079,-0.214089)--cycle;
\path[fill=blue!13!white,draw=blue!13!white,line width=.12pt] (-2.465877,-0.333925)--(-2.548600,-0.466550)--(-2.369413,-0.587353)--(-2.292506,-0.450808)--cycle;
\path[fill=blue!19!white,draw=blue!19!white,line width=.12pt] (2.292506,-0.450808)--(2.369413,-0.587353)--(2.548600,-0.466550)--(2.465877,-0.333925)--cycle;
\path[fill=blue!20!white,draw=blue!20!white,line width=.12pt] (-2.473633,-1.015029)--(-2.461751,-1.145316)--(-2.251872,-1.261099)--(-2.262742,-1.131371)--cycle;
\path[fill=blue!27!white,draw=blue!27!white,line width=.12pt] (2.262742,-1.131371)--(2.251872,-1.261099)--(2.461751,-1.145316)--(2.473633,-1.015029)--cycle;
\path[fill=blue!34!white,draw=blue!34!white,line width=.12pt] (-1.340867,-1.830352)--(-1.275669,-1.833387)--(-1.035598,-1.890159)--(-1.088526,-1.890026)--cycle;
\path[fill=blue!36!white,draw=blue!36!white,line width=.12pt] (1.088526,-1.890026)--(1.035598,-1.890159)--(1.275669,-1.833387)--(1.340867,-1.830352)--cycle;
\path[fill=blue!37!white,draw=blue!37!white,line width=.12pt] (-0.785553,-1.934893)--(-0.741989,-1.902512)--(-0.498665,-1.932987)--(-0.527943,-1.967157)--cycle;
\path[fill=blue!37!white,draw=blue!37!white,line width=.12pt] (0.527943,-1.967157)--(0.498665,-1.932987)--(0.741989,-1.902512)--(0.785553,-1.934893)--cycle;
\path[fill=blue!14!white,draw=blue!14!white,line width=.12pt] (-1.855225,-0.068452)--(-1.975871,-0.131270)--(-1.807416,-0.224200)--(-1.697056,-0.155708)--cycle;
\path[fill=blue!14!white,draw=blue!14!white,line width=.12pt] (1.697056,-0.155708)--(1.807416,-0.224200)--(1.975871,-0.131270)--(1.855225,-0.068452)--cycle;
\path[fill=blue!29!white,draw=blue!29!white,line width=.12pt] (-1.944527,-1.569617)--(-1.881411,-1.636151)--(-1.647647,-1.722837)--(-1.702920,-1.659211)--cycle;
\path[fill=blue!33!white,draw=blue!33!white,line width=.12pt] (1.702920,-1.659211)--(1.647647,-1.722837)--(1.881411,-1.636151)--(1.944527,-1.569617)--cycle;
\path[fill=blue!32!white,draw=blue!32!white,line width=.12pt] (-1.647647,-1.722837)--(-1.580295,-1.758598)--(-1.340867,-1.830352)--(-1.398014,-1.797648)--cycle;
\path[fill=blue!35!white,draw=blue!35!white,line width=.12pt] (1.398014,-1.797648)--(1.340867,-1.830352)--(1.580295,-1.758598)--(1.647647,-1.722837)--cycle;
\path[fill=blue!22!white,draw=blue!22!white,line width=.12pt] (-0.701962,-0.357445)--(-0.748354,-0.327284)--(-0.567664,-0.359610)--(-0.532473,-0.387767)--cycle;
\path[fill=blue!20!white,draw=blue!20!white,line width=.12pt] (0.532473,-0.387767)--(0.567664,-0.359610)--(0.748354,-0.327284)--(0.701962,-0.357445)--cycle;
\path[fill=blue!20!white,draw=blue!20!white,line width=.12pt] (-1.086442,-0.236929)--(-1.163283,-0.230405)--(-0.987036,-0.283224)--(-0.921837,-0.286259)--cycle;
\path[fill=blue!18!white,draw=blue!18!white,line width=.12pt] (0.921837,-0.286259)--(0.987036,-0.283224)--(1.163283,-0.230405)--(1.086442,-0.236929)--cycle;
\path[fill=blue!37!white,draw=blue!37!white,line width=.12pt] (-0.527943,-1.967157)--(-0.498665,-1.932987)--(-0.250539,-1.951390)--(-0.265249,-1.986640)--cycle;
\path[fill=blue!37!white,draw=blue!37!white,line width=.12pt] (0.265249,-1.986640)--(0.250539,-1.951390)--(0.498665,-1.932987)--(0.527943,-1.967157)--cycle;
\path[fill=blue!15!white,draw=blue!15!white,line width=.12pt] (-1.586697,-0.113840)--(-1.697056,-0.155708)--(-1.522544,-0.234792)--(-1.423533,-0.187782)--cycle;
\path[fill=blue!15!white,draw=blue!15!white,line width=.12pt] (1.423533,-0.187782)--(1.522544,-0.234792)--(1.697056,-0.155708)--(1.586697,-0.113840)--cycle;
\path[fill=blue!17!white,draw=blue!17!white,line width=.12pt] (-1.328326,-0.169203)--(-1.423533,-0.187782)--(-1.246659,-0.253371)--(-1.163283,-0.230405)--cycle;
\path[fill=blue!16!white,draw=blue!16!white,line width=.12pt] (1.163283,-0.230405)--(1.246659,-0.253371)--(1.423533,-0.187782)--(1.328326,-0.169203)--cycle;
\path[fill=blue!37!white,draw=blue!37!white,line width=.12pt] (-0.265249,-1.986640)--(-0.250539,-1.951390)--(-0.000000,-1.957544)--(-0.000000,-1.993156)--cycle;
\path[fill=blue!37!white,draw=blue!37!white,line width=.12pt] (-0.000000,-1.993156)--(-0.000000,-1.957544)--(0.250539,-1.951390)--(0.265249,-1.986640)--cycle;
\path[fill=blue!18!white,draw=blue!18!white,line width=.12pt] (-2.461751,-0.874991)--(-2.473633,-1.015029)--(-2.262742,-1.131371)--(-2.251872,-0.990774)--cycle;
\path[fill=blue!25!white,draw=blue!25!white,line width=.12pt] (2.251872,-0.990774)--(2.262742,-1.131371)--(2.473633,-1.015029)--(2.461751,-0.874991)--cycle;
\path[fill=blue!22!white,draw=blue!22!white,line width=.12pt] (-0.532473,-0.387767)--(-0.567664,-0.359610)--(-0.381508,-0.382925)--(-0.357857,-0.409636)--cycle;
\path[fill=blue!21!white,draw=blue!21!white,line width=.12pt] (0.357857,-0.409636)--(0.381508,-0.382925)--(0.567664,-0.359610)--(0.532473,-0.387767)--cycle;
\path[fill=blue!34!white,draw=blue!34!white,line width=.12pt] (-1.088526,-1.890026)--(-1.035598,-1.890159)--(-0.785553,-1.934893)--(-0.825702,-1.937046)--cycle;
\path[fill=blue!36!white,draw=blue!36!white,line width=.12pt] (0.825702,-1.937046)--(0.785553,-1.934893)--(1.035598,-1.890159)--(1.088526,-1.890026)--cycle;
\path[fill=blue!13!white,draw=blue!13!white,line width=.12pt] (-1.975871,-0.131270)--(-2.091880,-0.218298)--(-1.913535,-0.316685)--(-1.807416,-0.224200)--cycle;
\path[fill=blue!15!white,draw=blue!15!white,line width=.12pt] (1.807416,-0.224200)--(1.913535,-0.316685)--(2.091880,-0.218298)--(1.975871,-0.131270)--cycle;
\path[fill=blue!23!white,draw=blue!23!white,line width=.12pt] (-2.251872,-1.261099)--(-2.219681,-1.374972)--(-1.991426,-1.478411)--(-2.020307,-1.366038)--cycle;
\path[fill=blue!29!white,draw=blue!29!white,line width=.12pt] (2.020307,-1.366038)--(1.991426,-1.478411)--(2.219681,-1.374972)--(2.251872,-1.261099)--cycle;
\path[fill=blue!22!white,draw=blue!22!white,line width=.12pt] (-0.357857,-0.409636)--(-0.381508,-0.382925)--(-0.191677,-0.397005)--(-0.179794,-0.422843)--cycle;
\path[fill=blue!21!white,draw=blue!21!white,line width=.12pt] (0.179794,-0.422843)--(0.191677,-0.397005)--(0.381508,-0.382925)--(0.357857,-0.409636)--cycle;
\path[fill=blue!20!white,draw=blue!20!white,line width=.12pt] (-0.921837,-0.286259)--(-0.987036,-0.283224)--(-0.801283,-0.327152)--(-0.748354,-0.327284)--cycle;
\path[fill=blue!18!white,draw=blue!18!white,line width=.12pt] (0.748354,-0.327284)--(0.801283,-0.327152)--(0.987036,-0.283224)--(0.921837,-0.286259)--cycle;
\path[fill=blue!22!white,draw=blue!22!white,line width=.12pt] (-0.179794,-0.422843)--(-0.191677,-0.397005)--(-0.000000,-0.401713)--(-0.000000,-0.427259)--cycle;
\path[fill=blue!21!white,draw=blue!21!white,line width=.12pt] (-0.000000,-0.427259)--(-0.000000,-0.401713)--(0.191677,-0.397005)--(0.179794,-0.422843)--cycle;
\path[fill=blue!16!white,draw=blue!16!white,line width=.12pt] (-2.426560,-0.730582)--(-2.461751,-0.874991)--(-2.251872,-0.990774)--(-2.219681,-0.844710)--cycle;
\path[fill=blue!22!white,draw=blue!22!white,line width=.12pt] (2.219681,-0.844710)--(2.251872,-0.990774)--(2.461751,-0.874991)--(2.426560,-0.730582)--cycle;
\path[fill=blue!12!white,draw=blue!12!white,line width=.12pt] (-2.091880,-0.218298)--(-2.198794,-0.326193)--(-2.011334,-0.429608)--(-1.913535,-0.316685)--cycle;
\path[fill=blue!15!white,draw=blue!15!white,line width=.12pt] (1.913535,-0.316685)--(2.011334,-0.429608)--(2.198794,-0.326193)--(2.091880,-0.218298)--cycle;
\path[fill=blue!32!white,draw=blue!32!white,line width=.12pt] (-1.398014,-1.797648)--(-1.340867,-1.830352)--(-1.088526,-1.890026)--(-1.134919,-1.859866)--cycle;
\path[fill=blue!34!white,draw=blue!34!white,line width=.12pt] (1.134919,-1.859866)--(1.088526,-1.890026)--(1.340867,-1.830352)--(1.398014,-1.797648)--cycle;
\path[fill=blue!26!white,draw=blue!26!white,line width=.12pt] (-1.991426,-1.478411)--(-1.944527,-1.569617)--(-1.702920,-1.659211)--(-1.743993,-1.570165)--cycle;
\path[fill=blue!31!white,draw=blue!31!white,line width=.12pt] (1.743993,-1.570165)--(1.702920,-1.659211)--(1.944527,-1.569617)--(1.991426,-1.478411)--cycle;
\path[fill=blue!35!white,draw=blue!35!white,line width=.12pt] (-0.825702,-1.937046)--(-0.785553,-1.934893)--(-0.527943,-1.967157)--(-0.554926,-1.970959)--cycle;
\path[fill=blue!36!white,draw=blue!36!white,line width=.12pt] (0.554926,-1.970959)--(0.527943,-1.967157)--(0.785553,-1.934893)--(0.825702,-1.937046)--cycle;
\path[fill=blue!14!white,draw=blue!14!white,line width=.12pt] (-1.697056,-0.155708)--(-1.807416,-0.224200)--(-1.621555,-0.308427)--(-1.522544,-0.234792)--cycle;
\path[fill=blue!14!white,draw=blue!14!white,line width=.12pt] (1.522544,-0.234792)--(1.621555,-0.308427)--(1.807416,-0.224200)--(1.697056,-0.155708)--cycle;
\path[fill=blue!29!white,draw=blue!29!white,line width=.12pt] (-1.702920,-1.659211)--(-1.647647,-1.722837)--(-1.398014,-1.797648)--(-1.444914,-1.736532)--cycle;
\path[fill=blue!32!white,draw=blue!32!white,line width=.12pt] (1.444914,-1.736532)--(1.398014,-1.797648)--(1.647647,-1.722837)--(1.702920,-1.659211)--cycle;
\path[fill=blue!14!white,draw=blue!14!white,line width=.12pt] (-2.369413,-0.587353)--(-2.426560,-0.730582)--(-2.219681,-0.844710)--(-2.167407,-0.698793)--cycle;
\path[fill=blue!20!white,draw=blue!20!white,line width=.12pt] (2.167407,-0.698793)--(2.219681,-0.844710)--(2.426560,-0.730582)--(2.369413,-0.587353)--cycle;
\path[fill=blue!17!white,draw=blue!17!white,line width=.12pt] (-1.163283,-0.230405)--(-1.246659,-0.253371)--(-1.057780,-0.309976)--(-0.987036,-0.283224)--cycle;
\path[fill=blue!16!white,draw=blue!16!white,line width=.12pt] (0.987036,-0.283224)--(1.057780,-0.309976)--(1.246659,-0.253371)--(1.163283,-0.230405)--cycle;
\path[fill=blue!12!white,draw=blue!12!white,line width=.12pt] (-2.198794,-0.326193)--(-2.292506,-0.450808)--(-2.097056,-0.558630)--(-2.011334,-0.429608)--cycle;
\path[fill=blue!16!white,draw=blue!16!white,line width=.12pt] (2.011334,-0.429608)--(2.097056,-0.558630)--(2.292506,-0.450808)--(2.198794,-0.326193)--cycle;
\path[fill=blue!18!white,draw=blue!18!white,line width=.12pt] (0.567664,-0.359610)--(0.607814,-0.361764)--(0.801283,-0.327152)--(0.748354,-0.327284)--cycle;
\path[fill=blue!20!white,draw=blue!20!white,line width=.12pt] (-0.748354,-0.327284)--(-0.801283,-0.327152)--(-0.607814,-0.361764)--(-0.567664,-0.359610)--cycle;
\path[fill=blue!13!white,draw=blue!13!white,line width=.12pt] (-2.292506,-0.450808)--(-2.369413,-0.587353)--(-2.167407,-0.698793)--(-2.097056,-0.558630)--cycle;
\path[fill=blue!18!white,draw=blue!18!white,line width=.12pt] (2.097056,-0.558630)--(2.167407,-0.698793)--(2.369413,-0.587353)--(2.292506,-0.450808)--cycle;
\path[fill=blue!20!white,draw=blue!20!white,line width=.12pt] (-2.262742,-1.131371)--(-2.251872,-1.261099)--(-2.020307,-1.366038)--(-2.030059,-1.236817)--cycle;
\path[fill=blue!26!white,draw=blue!26!white,line width=.12pt] (2.030059,-1.236817)--(2.020307,-1.366038)--(2.251872,-1.261099)--(2.262742,-1.131371)--cycle;
\path[fill=blue!15!white,draw=blue!15!white,line width=.12pt] (-1.423533,-0.187782)--(-1.522544,-0.234792)--(-1.333369,-0.304943)--(-1.246659,-0.253371)--cycle;
\path[fill=blue!15!white,draw=blue!15!white,line width=.12pt] (1.246659,-0.253371)--(1.333369,-0.304943)--(1.522544,-0.234792)--(1.423533,-0.187782)--cycle;
\path[fill=blue!35!white,draw=blue!35!white,line width=.12pt] (-0.554926,-1.970959)--(-0.527943,-1.967157)--(-0.265249,-1.986640)--(-0.278805,-1.991439)--cycle;
\path[fill=blue!35!white,draw=blue!35!white,line width=.12pt] (0.278805,-1.991439)--(0.265249,-1.986640)--(0.527943,-1.967157)--(0.554926,-1.970959)--cycle;
\path[fill=blue!35!white,draw=blue!35!white,line width=.12pt] (-0.278805,-1.991439)--(-0.265249,-1.986640)--(-0.000000,-1.993156)--(-0.000000,-1.998287)--cycle;
\path[fill=blue!35!white,draw=blue!35!white,line width=.12pt] (-0.000000,-1.998287)--(-0.000000,-1.993156)--(0.265249,-1.986640)--(0.278805,-1.991439)--cycle;
\path[fill=blue!19!white,draw=blue!19!white,line width=.12pt] (-0.567664,-0.359610)--(-0.607814,-0.361764)--(-0.408491,-0.386728)--(-0.381508,-0.382925)--cycle;
\path[fill=blue!18!white,draw=blue!18!white,line width=.12pt] (0.381508,-0.382925)--(0.408491,-0.386728)--(0.607814,-0.361764)--(0.567664,-0.359610)--cycle;
\path[fill=blue!34!white,draw=blue!34!white,line width=.12pt] (0.860893,-1.908890)--(0.825702,-1.937046)--(1.088526,-1.890026)--(1.134919,-1.859866)--cycle;
\path[fill=blue!32!white,draw=blue!32!white,line width=.12pt] (-1.134919,-1.859866)--(-1.088526,-1.890026)--(-0.825702,-1.937046)--(-0.860893,-1.908890)--cycle;
\path[fill=blue!13!white,draw=blue!13!white,line width=.12pt] (-1.807416,-0.224200)--(-1.913535,-0.316685)--(-1.716761,-0.405857)--(-1.621555,-0.308427)--cycle;
\path[fill=blue!14!white,draw=blue!14!white,line width=.12pt] (1.621555,-0.308427)--(1.716761,-0.405857)--(1.913535,-0.316685)--(1.807416,-0.224200)--cycle;
\path[fill=blue!17!white,draw=blue!17!white,line width=.12pt] (-0.987036,-0.283224)--(-1.057780,-0.309976)--(-0.858714,-0.357051)--(-0.801283,-0.327152)--cycle;
\path[fill=blue!16!white,draw=blue!16!white,line width=.12pt] (0.801283,-0.327152)--(0.858714,-0.357051)--(1.057780,-0.309976)--(0.987036,-0.283224)--cycle;
\path[fill=blue!23!white,draw=blue!23!white,line width=.12pt] (-2.020307,-1.366038)--(-1.991426,-1.478411)--(-1.743993,-1.570165)--(-1.769285,-1.459123)--cycle;
\path[fill=blue!28!white,draw=blue!28!white,line width=.12pt] (1.769285,-1.459123)--(1.743993,-1.570165)--(1.991426,-1.478411)--(2.020307,-1.366038)--cycle;
\path[fill=blue!19!white,draw=blue!19!white,line width=.12pt] (-0.381508,-0.382925)--(-0.408491,-0.386728)--(-0.205234,-0.401803)--(-0.191677,-0.397005)--cycle;
\path[fill=blue!19!white,draw=blue!19!white,line width=.12pt] (0.191677,-0.397005)--(0.205234,-0.401803)--(0.408491,-0.386728)--(0.381508,-0.382925)--cycle;
\path[fill=blue!18!white,draw=blue!18!white,line width=.12pt] (-2.251872,-0.990774)--(-2.262742,-1.131371)--(-2.030059,-1.236817)--(-2.020307,-1.095713)--cycle;
\path[fill=blue!24!white,draw=blue!24!white,line width=.12pt] (2.020307,-1.095713)--(2.030059,-1.236817)--(2.262742,-1.131371)--(2.251872,-0.990774)--cycle;
\path[fill=blue!19!white,draw=blue!19!white,line width=.12pt] (-0.191677,-0.397005)--(-0.205234,-0.401803)--(-0.000000,-0.406844)--(-0.000000,-0.401713)--cycle;
\path[fill=blue!19!white,draw=blue!19!white,line width=.12pt] (-0.000000,-0.401713)--(-0.000000,-0.406844)--(0.205234,-0.401803)--(0.191677,-0.397005)--cycle;
\path[fill=blue!29!white,draw=blue!29!white,line width=.12pt] (-1.444914,-1.736532)--(-1.398014,-1.797648)--(-1.134919,-1.859866)--(-1.172992,-1.800837)--cycle;
\path[fill=blue!32!white,draw=blue!32!white,line width=.12pt] (1.172992,-1.800837)--(1.134919,-1.859866)--(1.398014,-1.797648)--(1.444914,-1.736532)--cycle;
\path[fill=blue!26!white,draw=blue!26!white,line width=.12pt] (-1.743993,-1.570165)--(-1.702920,-1.659211)--(-1.444914,-1.736532)--(-1.479763,-1.649352)--cycle;
\path[fill=blue!30!white,draw=blue!30!white,line width=.12pt] (1.479763,-1.649352)--(1.444914,-1.736532)--(1.702920,-1.659211)--(1.743993,-1.570165)--cycle;
\path[fill=blue!14!white,draw=blue!14!white,line width=.12pt] (-1.522544,-0.234792)--(-1.621555,-0.308427)--(-1.420078,-0.383140)--(-1.333369,-0.304943)--cycle;
\path[fill=blue!14!white,draw=blue!14!white,line width=.12pt] (1.333369,-0.304943)--(1.420078,-0.383140)--(1.621555,-0.308427)--(1.522544,-0.234792)--cycle;
\path[fill=blue!15!white,draw=blue!15!white,line width=.12pt] (-1.246659,-0.253371)--(-1.333369,-0.304943)--(-1.131352,-0.365485)--(-1.057780,-0.309976)--cycle;
\path[fill=blue!15!white,draw=blue!15!white,line width=.12pt] (1.057780,-0.309976)--(1.131352,-0.365485)--(1.333369,-0.304943)--(1.246659,-0.253371)--cycle;
\path[fill=blue!12!white,draw=blue!12!white,line width=.12pt] (-1.913535,-0.316685)--(-2.011334,-0.429608)--(-1.804504,-0.523338)--(-1.716761,-0.405857)--cycle;
\path[fill=blue!15!white,draw=blue!15!white,line width=.12pt] (1.716761,-0.405857)--(1.804504,-0.523338)--(2.011334,-0.429608)--(1.913535,-0.316685)--cycle;
\path[fill=blue!32!white,draw=blue!32!white,line width=.12pt] (-0.860893,-1.908890)--(-0.825702,-1.937046)--(-0.554926,-1.970959)--(-0.578577,-1.944248)--cycle;
\path[fill=blue!33!white,draw=blue!33!white,line width=.12pt] (0.578577,-1.944248)--(0.554926,-1.970959)--(0.825702,-1.937046)--(0.860893,-1.908890)--cycle;
\path[fill=blue!16!white,draw=blue!16!white,line width=.12pt] (-2.219681,-0.844710)--(-2.251872,-0.990774)--(-2.020307,-1.095713)--(-1.991426,-0.948149)--cycle;
\path[fill=blue!21!white,draw=blue!21!white,line width=.12pt] (1.991426,-0.948149)--(2.020307,-1.095713)--(2.251872,-0.990774)--(2.219681,-0.844710)--cycle;
\path[fill=blue!17!white,draw=blue!17!white,line width=.12pt] (-0.801283,-0.327152)--(-0.858714,-0.357051)--(-0.651378,-0.394145)--(-0.607814,-0.361764)--cycle;
\path[fill=blue!16!white,draw=blue!16!white,line width=.12pt] (0.607814,-0.361764)--(0.651378,-0.394145)--(0.858714,-0.357051)--(0.801283,-0.327152)--cycle;
\path[fill=blue!12!white,draw=blue!12!white,line width=.12pt] (-2.011334,-0.429608)--(-2.097056,-0.558630)--(-1.881411,-0.656355)--(-1.804504,-0.523338)--cycle;
\path[fill=blue!16!white,draw=blue!16!white,line width=.12pt] (1.804504,-0.523338)--(1.881411,-0.656355)--(2.097056,-0.558630)--(2.011334,-0.429608)--cycle;
\path[fill=blue!14!white,draw=blue!14!white,line width=.12pt] (-2.167407,-0.698793)--(-2.219681,-0.844710)--(-1.991426,-0.948149)--(-1.944527,-0.799796)--cycle;
\path[fill=blue!19!white,draw=blue!19!white,line width=.12pt] (1.944527,-0.799796)--(1.991426,-0.948149)--(2.219681,-0.844710)--(2.167407,-0.698793)--cycle;
\path[fill=blue!32!white,draw=blue!32!white,line width=.12pt] (-0.578577,-1.944248)--(-0.554926,-1.970959)--(-0.278805,-1.991439)--(-0.290688,-1.965600)--cycle;
\path[fill=blue!33!white,draw=blue!33!white,line width=.12pt] (0.290688,-1.965600)--(0.278805,-1.991439)--(0.554926,-1.970959)--(0.578577,-1.944248)--cycle;
\path[fill=blue!20!white,draw=blue!20!white,line width=.12pt] (-2.030059,-1.236817)--(-2.020307,-1.366038)--(-1.769285,-1.459123)--(-1.777825,-1.330351)--cycle;
\path[fill=blue!26!white,draw=blue!26!white,line width=.12pt] (1.777825,-1.330351)--(1.769285,-1.459123)--(2.020307,-1.366038)--(2.030059,-1.236817)--cycle;
\path[fill=blue!13!white,draw=blue!13!white,line width=.12pt] (-2.097056,-0.558630)--(-2.167407,-0.698793)--(-1.944527,-0.799796)--(-1.881411,-0.656355)--cycle;
\path[fill=blue!17!white,draw=blue!17!white,line width=.12pt] (1.881411,-0.656355)--(1.944527,-0.799796)--(2.167407,-0.698793)--(2.097056,-0.558630)--cycle;
\path[fill=blue!17!white,draw=blue!17!white,line width=.12pt] (-0.607814,-0.361764)--(-0.651378,-0.394145)--(-0.437769,-0.420898)--(-0.408491,-0.386728)--cycle;
\path[fill=blue!17!white,draw=blue!17!white,line width=.12pt] (0.408491,-0.386728)--(0.437769,-0.420898)--(0.651378,-0.394145)--(0.607814,-0.361764)--cycle;
\path[fill=blue!32!white,draw=blue!32!white,line width=.12pt] (-0.290688,-1.965600)--(-0.278805,-1.991439)--(-0.000000,-1.998287)--(-0.000000,-1.972741)--cycle;
\path[fill=blue!33!white,draw=blue!33!white,line width=.12pt] (-0.000000,-1.972741)--(-0.000000,-1.998287)--(0.278805,-1.991439)--(0.290688,-1.965600)--cycle;
\path[fill=blue!29!white,draw=blue!29!white,line width=.12pt] (-1.172992,-1.800837)--(-1.134919,-1.859866)--(-0.860893,-1.908890)--(-0.889774,-1.851505)--cycle;
\path[fill=blue!31!white,draw=blue!31!white,line width=.12pt] (0.889774,-1.851505)--(0.860893,-1.908890)--(1.134919,-1.859866)--(1.172992,-1.800837)--cycle;
\path[fill=blue!13!white,draw=blue!13!white,line width=.12pt] (-1.621555,-0.308427)--(-1.716761,-0.405857)--(-1.503455,-0.484957)--(-1.420078,-0.383140)--cycle;
\path[fill=blue!14!white,draw=blue!14!white,line width=.12pt] (1.420078,-0.383140)--(1.503455,-0.484957)--(1.716761,-0.405857)--(1.621555,-0.308427)--cycle;
\path[fill=blue!15!white,draw=blue!15!white,line width=.12pt] (-1.057780,-0.309976)--(-1.131352,-0.365485)--(-0.918440,-0.415835)--(-0.858714,-0.357051)--cycle;
\path[fill=blue!15!white,draw=blue!15!white,line width=.12pt] (0.858714,-0.357051)--(0.918440,-0.415835)--(1.131352,-0.365485)--(1.057780,-0.309976)--cycle;
\path[fill=blue!23!white,draw=blue!23!white,line width=.12pt] (-1.769285,-1.459123)--(-1.743993,-1.570165)--(-1.479763,-1.649352)--(-1.501223,-1.539458)--cycle;
\path[fill=blue!27!white,draw=blue!27!white,line width=.12pt] (1.501223,-1.539458)--(1.479763,-1.649352)--(1.743993,-1.570165)--(1.769285,-1.459123)--cycle;
\path[fill=blue!14!white,draw=blue!14!white,line width=.12pt] (-1.333369,-0.304943)--(-1.420078,-0.383140)--(-1.204924,-0.447619)--(-1.131352,-0.365485)--cycle;
\path[fill=blue!14!white,draw=blue!14!white,line width=.12pt] (1.131352,-0.365485)--(1.204924,-0.447619)--(1.420078,-0.383140)--(1.333369,-0.304943)--cycle;
\path[fill=blue!17!white,draw=blue!17!white,line width=.12pt] (0.205234,-0.401803)--(0.219943,-0.437053)--(0.437769,-0.420898)--(0.408491,-0.386728)--cycle;
\path[fill=blue!17!white,draw=blue!17!white,line width=.12pt] (-0.408491,-0.386728)--(-0.437769,-0.420898)--(-0.219943,-0.437053)--(-0.205234,-0.401803)--cycle;
\path[fill=blue!26!white,draw=blue!26!white,line width=.12pt] (-1.479763,-1.649352)--(-1.444914,-1.736532)--(-1.172992,-1.800837)--(-1.201283,-1.715208)--cycle;
\path[fill=blue!29!white,draw=blue!29!white,line width=.12pt] (1.201283,-1.715208)--(1.172992,-1.800837)--(1.444914,-1.736532)--(1.479763,-1.649352)--cycle;
\path[fill=blue!17!white,draw=blue!17!white,line width=.12pt] (-0.205234,-0.401803)--(-0.219943,-0.437053)--(-0.000000,-0.442456)--(-0.000000,-0.406844)--cycle;
\path[fill=blue!17!white,draw=blue!17!white,line width=.12pt] (-0.000000,-0.406844)--(-0.000000,-0.442456)--(0.219943,-0.437053)--(0.205234,-0.401803)--cycle;
\path[fill=blue!18!white,draw=blue!18!white,line width=.12pt] (-2.020307,-1.095713)--(-2.030059,-1.236817)--(-1.777825,-1.330351)--(-1.769285,-1.188798)--cycle;
\path[fill=blue!23!white,draw=blue!23!white,line width=.12pt] (1.769285,-1.188798)--(1.777825,-1.330351)--(2.030059,-1.236817)--(2.020307,-1.095713)--cycle;
\path[fill=blue!29!white,draw=blue!29!white,line width=.12pt] (-0.889774,-1.851505)--(-0.860893,-1.908890)--(-0.578577,-1.944248)--(-0.597986,-1.888050)--cycle;
\path[fill=blue!31!white,draw=blue!31!white,line width=.12pt] (0.597986,-1.888050)--(0.578577,-1.944248)--(0.860893,-1.908890)--(0.889774,-1.851505)--cycle;
\path[fill=blue!12!white,draw=blue!12!white,line width=.12pt] (-1.716761,-0.405857)--(-1.804504,-0.523338)--(-1.580295,-0.606480)--(-1.503455,-0.484957)--cycle;
\path[fill=blue!15!white,draw=blue!15!white,line width=.12pt] (1.503455,-0.484957)--(1.580295,-0.606480)--(1.804504,-0.523338)--(1.716761,-0.405857)--cycle;
\path[fill=blue!15!white,draw=blue!15!white,line width=.12pt] (-0.858714,-0.357051)--(-0.918440,-0.415835)--(-0.696683,-0.455508)--(-0.651378,-0.394145)--cycle;
\path[fill=blue!15!white,draw=blue!15!white,line width=.12pt] (0.651378,-0.394145)--(0.696683,-0.455508)--(0.918440,-0.415835)--(0.858714,-0.357051)--cycle;
\path[fill=blue!16!white,draw=blue!16!white,line width=.12pt] (-1.991426,-0.948149)--(-2.020307,-1.095713)--(-1.769285,-1.188798)--(-1.743993,-1.039904)--cycle;
\path[fill=blue!21!white,draw=blue!21!white,line width=.12pt] (1.743993,-1.039904)--(1.769285,-1.188798)--(2.020307,-1.095713)--(1.991426,-0.948149)--cycle;
\path[fill=blue!12!white,draw=blue!12!white,line width=.12pt] (-1.804504,-0.523338)--(-1.881411,-0.656355)--(-1.647647,-0.743041)--(-1.580295,-0.606480)--cycle;
\path[fill=blue!16!white,draw=blue!16!white,line width=.12pt] (1.580295,-0.606480)--(1.647647,-0.743041)--(1.881411,-0.656355)--(1.804504,-0.523338)--cycle;
\path[fill=blue!29!white,draw=blue!29!white,line width=.12pt] (-0.597986,-1.888050)--(-0.578577,-1.944248)--(-0.290688,-1.965600)--(-0.300440,-1.910118)--cycle;
\path[fill=blue!30!white,draw=blue!30!white,line width=.12pt] (0.300440,-1.910118)--(0.290688,-1.965600)--(0.578577,-1.944248)--(0.597986,-1.888050)--cycle;
\path[fill=blue!20!white,draw=blue!20!white,line width=.12pt] (-1.777825,-1.330351)--(-1.769285,-1.459123)--(-1.501223,-1.539458)--(-1.508470,-1.411074)--cycle;
\path[fill=blue!25!white,draw=blue!25!white,line width=.12pt] (1.508470,-1.411074)--(1.501223,-1.539458)--(1.769285,-1.459123)--(1.777825,-1.330351)--cycle;
\path[fill=blue!13!white,draw=blue!13!white,line width=.12pt] (-1.420078,-0.383140)--(-1.503455,-0.484957)--(-1.275669,-0.553222)--(-1.204924,-0.447619)--cycle;
\path[fill=blue!14!white,draw=blue!14!white,line width=.12pt] (1.204924,-0.447619)--(1.275669,-0.553222)--(1.503455,-0.484957)--(1.420078,-0.383140)--cycle;
\path[fill=blue!14!white,draw=blue!14!white,line width=.12pt] (-1.131352,-0.365485)--(-1.204924,-0.447619)--(-0.978167,-0.501243)--(-0.918440,-0.415835)--cycle;
\path[fill=blue!14!white,draw=blue!14!white,line width=.12pt] (0.918440,-0.415835)--(0.978167,-0.501243)--(1.204924,-0.447619)--(1.131352,-0.365485)--cycle;
\path[fill=blue!26!white,draw=blue!26!white,line width=.12pt] (-1.201283,-1.715208)--(-1.172992,-1.800837)--(-0.889774,-1.851505)--(-0.911234,-1.767098)--cycle;
\path[fill=blue!29!white,draw=blue!29!white,line width=.12pt] (0.911234,-1.767098)--(0.889774,-1.851505)--(1.172992,-1.800837)--(1.201283,-1.715208)--cycle;
\path[fill=blue!14!white,draw=blue!14!white,line width=.12pt] (-1.944527,-0.799796)--(-1.991426,-0.948149)--(-1.743993,-1.039904)--(-1.702920,-0.889390)--cycle;
\path[fill=blue!19!white,draw=blue!19!white,line width=.12pt] (1.702920,-0.889390)--(1.743993,-1.039904)--(1.991426,-0.948149)--(1.944527,-0.799796)--cycle;
\path[fill=blue!15!white,draw=blue!15!white,line width=.12pt] (-0.651378,-0.394145)--(-0.696683,-0.455508)--(-0.468217,-0.484122)--(-0.437769,-0.420898)--cycle;
\path[fill=blue!15!white,draw=blue!15!white,line width=.12pt] (0.437769,-0.420898)--(0.468217,-0.484122)--(0.696683,-0.455508)--(0.651378,-0.394145)--cycle;
\path[fill=blue!13!white,draw=blue!13!white,line width=.12pt] (-1.881411,-0.656355)--(-1.944527,-0.799796)--(-1.702920,-0.889390)--(-1.647647,-0.743041)--cycle;
\path[fill=blue!17!white,draw=blue!17!white,line width=.12pt] (1.647647,-0.743041)--(1.702920,-0.889390)--(1.944527,-0.799796)--(1.881411,-0.656355)--cycle;
\path[fill=blue!30!white,draw=blue!30!white,line width=.12pt] (-0.300440,-1.910118)--(-0.290688,-1.965600)--(-0.000000,-1.972741)--(-0.000000,-1.917498)--cycle;
\path[fill=blue!30!white,draw=blue!30!white,line width=.12pt] (-0.000000,-1.917498)--(-0.000000,-1.972741)--(0.290688,-1.965600)--(0.300440,-1.910118)--cycle;
\path[fill=blue!23!white,draw=blue!23!white,line width=.12pt] (-1.501223,-1.539458)--(-1.479763,-1.649352)--(-1.201283,-1.715208)--(-1.218704,-1.606269)--cycle;
\path[fill=blue!27!white,draw=blue!27!white,line width=.12pt] (1.218704,-1.606269)--(1.201283,-1.715208)--(1.479763,-1.649352)--(1.501223,-1.539458)--cycle;
\path[fill=blue!15!white,draw=blue!15!white,line width=.12pt] (-0.437769,-0.420898)--(-0.468217,-0.484122)--(-0.235241,-0.501401)--(-0.219943,-0.437053)--cycle;
\path[fill=blue!15!white,draw=blue!15!white,line width=.12pt] (0.219943,-0.437053)--(0.235241,-0.501401)--(0.468217,-0.484122)--(0.437769,-0.420898)--cycle;
\path[fill=blue!15!white,draw=blue!15!white,line width=.12pt] (-0.219943,-0.437053)--(-0.235241,-0.501401)--(-0.000000,-0.507180)--(-0.000000,-0.442456)--cycle;
\path[fill=blue!15!white,draw=blue!15!white,line width=.12pt] (-0.000000,-0.442456)--(-0.000000,-0.507180)--(0.235241,-0.501401)--(0.219943,-0.437053)--cycle;
\path[fill=blue!18!white,draw=blue!18!white,line width=.12pt] (-1.769285,-1.188798)--(-1.777825,-1.330351)--(-1.508470,-1.411074)--(-1.501223,-1.269133)--cycle;
\path[fill=blue!22!white,draw=blue!22!white,line width=.12pt] (1.501223,-1.269133)--(1.508470,-1.411074)--(1.777825,-1.330351)--(1.769285,-1.188798)--cycle;
\path[fill=blue!12!white,draw=blue!12!white,line width=.12pt] (-1.503455,-0.484957)--(-1.580295,-0.606480)--(-1.340867,-0.678234)--(-1.275669,-0.553222)--cycle;
\path[fill=blue!14!white,draw=blue!14!white,line width=.12pt] (1.275669,-0.553222)--(1.340867,-0.678234)--(1.580295,-0.606480)--(1.503455,-0.484957)--cycle;
\path[fill=blue!14!white,draw=blue!14!white,line width=.12pt] (-0.918440,-0.415835)--(-0.978167,-0.501243)--(-0.741989,-0.543496)--(-0.696683,-0.455508)--cycle;
\path[fill=blue!14!white,draw=blue!14!white,line width=.12pt] (0.696683,-0.455508)--(0.741989,-0.543496)--(0.978167,-0.501243)--(0.918440,-0.415835)--cycle;
\path[fill=blue!26!white,draw=blue!26!white,line width=.12pt] (-0.911234,-1.767098)--(-0.889774,-1.851505)--(-0.597986,-1.888050)--(-0.612409,-1.804524)--cycle;
\path[fill=blue!28!white,draw=blue!28!white,line width=.12pt] (0.612409,-1.804524)--(0.597986,-1.888050)--(0.889774,-1.851505)--(0.911234,-1.767098)--cycle;
\path[fill=blue!13!white,draw=blue!13!white,line width=.12pt] (-1.204924,-0.447619)--(-1.275669,-0.553222)--(-1.035598,-0.609994)--(-0.978167,-0.501243)--cycle;
\path[fill=blue!14!white,draw=blue!14!white,line width=.12pt] (0.978167,-0.501243)--(1.035598,-0.609994)--(1.275669,-0.553222)--(1.204924,-0.447619)--cycle;
\path[fill=blue!20!white,draw=blue!20!white,line width=.12pt] (-1.508470,-1.411074)--(-1.501223,-1.539458)--(-1.218704,-1.606269)--(-1.224587,-1.478207)--cycle;
\path[fill=blue!24!white,draw=blue!24!white,line width=.12pt] (1.224587,-1.478207)--(1.218704,-1.606269)--(1.501223,-1.539458)--(1.508470,-1.411074)--cycle;
\path[fill=blue!23!white,draw=blue!23!white,line width=.12pt] (-1.218704,-1.606269)--(-1.201283,-1.715208)--(-0.911234,-1.767098)--(-0.924449,-1.658912)--cycle;
\path[fill=blue!26!white,draw=blue!26!white,line width=.12pt] (0.924449,-1.658912)--(0.911234,-1.767098)--(1.201283,-1.715208)--(1.218704,-1.606269)--cycle;
\path[fill=blue!16!white,draw=blue!16!white,line width=.12pt] (-1.743993,-1.039904)--(-1.769285,-1.188798)--(-1.501223,-1.269133)--(-1.479763,-1.119090)--cycle;
\path[fill=blue!20!white,draw=blue!20!white,line width=.12pt] (1.479763,-1.119090)--(1.501223,-1.269133)--(1.769285,-1.188798)--(1.743993,-1.039904)--cycle;
\path[fill=blue!12!white,draw=blue!12!white,line width=.12pt] (-1.580295,-0.606480)--(-1.647647,-0.743041)--(-1.398014,-0.817853)--(-1.340867,-0.678234)--cycle;
\path[fill=blue!15!white,draw=blue!15!white,line width=.12pt] (1.340867,-0.678234)--(1.398014,-0.817853)--(1.647647,-0.743041)--(1.580295,-0.606480)--cycle;
\path[fill=blue!26!white,draw=blue!26!white,line width=.12pt] (-0.612409,-1.804524)--(-0.597986,-1.888050)--(-0.300440,-1.910118)--(-0.307686,-1.827125)--cycle;
\path[fill=blue!28!white,draw=blue!28!white,line width=.12pt] (0.307686,-1.827125)--(0.300440,-1.910118)--(0.597986,-1.888050)--(0.612409,-1.804524)--cycle;
\path[fill=blue!14!white,draw=blue!14!white,line width=.12pt] (-0.696683,-0.455508)--(-0.741989,-0.543496)--(-0.498665,-0.573971)--(-0.468217,-0.484122)--cycle;
\path[fill=blue!14!white,draw=blue!14!white,line width=.12pt] (0.468217,-0.484122)--(0.498665,-0.573971)--(0.741989,-0.543496)--(0.696683,-0.455508)--cycle;
\path[fill=blue!14!white,draw=blue!14!white,line width=.12pt] (-1.702920,-0.889390)--(-1.743993,-1.039904)--(-1.479763,-1.119090)--(-1.444914,-0.966711)--cycle;
\path[fill=blue!18!white,draw=blue!18!white,line width=.12pt] (1.444914,-0.966711)--(1.479763,-1.119090)--(1.743993,-1.039904)--(1.702920,-0.889390)--cycle;
\path[fill=blue!13!white,draw=blue!13!white,line width=.12pt] (-1.647647,-0.743041)--(-1.702920,-0.889390)--(-1.444914,-0.966711)--(-1.398014,-0.817853)--cycle;
\path[fill=blue!16!white,draw=blue!16!white,line width=.12pt] (1.398014,-0.817853)--(1.444914,-0.966711)--(1.702920,-0.889390)--(1.647647,-0.743041)--cycle;
\path[fill=blue!27!white,draw=blue!27!white,line width=.12pt] (-0.307686,-1.827125)--(-0.300440,-1.910118)--(-0.000000,-1.917498)--(-0.000000,-1.834683)--cycle;
\path[fill=blue!27!white,draw=blue!27!white,line width=.12pt] (-0.000000,-1.834683)--(-0.000000,-1.917498)--(0.300440,-1.910118)--(0.307686,-1.827125)--cycle;
\path[fill=blue!14!white,draw=blue!14!white,line width=.12pt] (-0.468217,-0.484122)--(-0.498665,-0.573971)--(-0.250539,-0.592374)--(-0.235241,-0.501401)--cycle;
\path[fill=blue!14!white,draw=blue!14!white,line width=.12pt] (0.235241,-0.501401)--(0.250539,-0.592374)--(0.498665,-0.573971)--(0.468217,-0.484122)--cycle;
\path[fill=blue!12!white,draw=blue!12!white,line width=.12pt] (-1.275669,-0.553222)--(-1.340867,-0.678234)--(-1.088526,-0.737908)--(-1.035598,-0.609994)--cycle;
\path[fill=blue!14!white,draw=blue!14!white,line width=.12pt] (1.035598,-0.609994)--(1.088526,-0.737908)--(1.340867,-0.678234)--(1.275669,-0.553222)--cycle;
\path[fill=blue!13!white,draw=blue!13!white,line width=.12pt] (-0.978167,-0.501243)--(-1.035598,-0.609994)--(-0.785553,-0.654728)--(-0.741989,-0.543496)--cycle;
\path[fill=blue!14!white,draw=blue!14!white,line width=.12pt] (0.741989,-0.543496)--(0.785553,-0.654728)--(1.035598,-0.609994)--(0.978167,-0.501243)--cycle;
\path[fill=blue!18!white,draw=blue!18!white,line width=.12pt] (-1.501223,-1.269133)--(-1.508470,-1.411074)--(-1.224587,-1.478207)--(-1.218704,-1.335944)--cycle;
\path[fill=blue!22!white,draw=blue!22!white,line width=.12pt] (1.218704,-1.335944)--(1.224587,-1.478207)--(1.508470,-1.411074)--(1.501223,-1.269133)--cycle;
\path[fill=blue!14!white,draw=blue!14!white,line width=.12pt] (-0.235241,-0.501401)--(-0.250539,-0.592374)--(-0.000000,-0.598528)--(-0.000000,-0.507180)--cycle;
\path[fill=blue!14!white,draw=blue!14!white,line width=.12pt] (-0.000000,-0.507180)--(-0.000000,-0.598528)--(0.250539,-0.592374)--(0.235241,-0.501401)--cycle;
\path[fill=blue!23!white,draw=blue!23!white,line width=.12pt] (-0.924449,-1.658912)--(-0.911234,-1.767098)--(-0.612409,-1.804524)--(-0.621290,-1.696881)--cycle;
\path[fill=blue!25!white,draw=blue!25!white,line width=.12pt] (0.621290,-1.696881)--(0.612409,-1.804524)--(0.911234,-1.767098)--(0.924449,-1.658912)--cycle;
\path[fill=blue!20!white,draw=blue!20!white,line width=.12pt] (-1.224587,-1.478207)--(-1.218704,-1.606269)--(-0.924449,-1.658912)--(-0.928911,-1.531105)--cycle;
\path[fill=blue!23!white,draw=blue!23!white,line width=.12pt] (0.928911,-1.531105)--(0.924449,-1.658912)--(1.218704,-1.606269)--(1.224587,-1.478207)--cycle;
\path[fill=blue!12!white,draw=blue!12!white,line width=.12pt] (-1.340867,-0.678234)--(-1.398014,-0.817853)--(-1.134919,-0.880070)--(-1.088526,-0.737908)--cycle;
\path[fill=blue!15!white,draw=blue!15!white,line width=.12pt] (1.088526,-0.737908)--(1.134919,-0.880070)--(1.398014,-0.817853)--(1.340867,-0.678234)--cycle;
\path[fill=blue!16!white,draw=blue!16!white,line width=.12pt] (-1.479763,-1.119090)--(-1.501223,-1.269133)--(-1.218704,-1.335944)--(-1.201283,-1.184946)--cycle;
\path[fill=blue!19!white,draw=blue!19!white,line width=.12pt] (1.201283,-1.184946)--(1.218704,-1.335944)--(1.501223,-1.269133)--(1.479763,-1.119090)--cycle;
\path[fill=blue!13!white,draw=blue!13!white,line width=.12pt] (-0.741989,-0.543496)--(-0.785553,-0.654728)--(-0.527943,-0.686992)--(-0.498665,-0.573971)--cycle;
\path[fill=blue!13!white,draw=blue!13!white,line width=.12pt] (0.498665,-0.573971)--(0.527943,-0.686992)--(0.785553,-0.654728)--(0.741989,-0.543496)--cycle;
\path[fill=blue!24!white,draw=blue!24!white,line width=.12pt] (-0.621290,-1.696881)--(-0.612409,-1.804524)--(-0.307686,-1.827125)--(-0.312148,-1.719809)--cycle;
\path[fill=blue!25!white,draw=blue!25!white,line width=.12pt] (0.312148,-1.719809)--(0.307686,-1.827125)--(0.612409,-1.804524)--(0.621290,-1.696881)--cycle;
\path[fill=blue!13!white,draw=blue!13!white,line width=.12pt] (-1.398014,-0.817853)--(-1.444914,-0.966711)--(-1.172992,-1.031016)--(-1.134919,-0.880070)--cycle;
\path[fill=blue!16!white,draw=blue!16!white,line width=.12pt] (1.134919,-0.880070)--(1.172992,-1.031016)--(1.444914,-0.966711)--(1.398014,-0.817853)--cycle;
\path[fill=blue!12!white,draw=blue!12!white,line width=.12pt] (-1.035598,-0.609994)--(-1.088526,-0.737908)--(-0.825702,-0.784928)--(-0.785553,-0.654728)--cycle;
\path[fill=blue!14!white,draw=blue!14!white,line width=.12pt] (0.785553,-0.654728)--(0.825702,-0.784928)--(1.088526,-0.737908)--(1.035598,-0.609994)--cycle;
\path[fill=blue!14!white,draw=blue!14!white,line width=.12pt] (-1.444914,-0.966711)--(-1.479763,-1.119090)--(-1.201283,-1.184946)--(-1.172992,-1.031016)--cycle;
\path[fill=blue!17!white,draw=blue!17!white,line width=.12pt] (1.172992,-1.031016)--(1.201283,-1.184946)--(1.479763,-1.119090)--(1.444914,-0.966711)--cycle;
\path[fill=blue!24!white,draw=blue!24!white,line width=.12pt] (-0.312148,-1.719809)--(-0.307686,-1.827125)--(-0.000000,-1.834683)--(-0.000000,-1.727477)--cycle;
\path[fill=blue!24!white,draw=blue!24!white,line width=.12pt] (-0.000000,-1.727477)--(-0.000000,-1.834683)--(0.307686,-1.827125)--(0.312148,-1.719809)--cycle;
\path[fill=blue!13!white,draw=blue!13!white,line width=.12pt] (-0.498665,-0.573971)--(-0.527943,-0.686992)--(-0.265249,-0.706475)--(-0.250539,-0.592374)--cycle;
\path[fill=blue!13!white,draw=blue!13!white,line width=.12pt] (0.250539,-0.592374)--(0.265249,-0.706475)--(0.527943,-0.686992)--(0.498665,-0.573971)--cycle;
\path[fill=blue!18!white,draw=blue!18!white,line width=.12pt] (-1.218704,-1.335944)--(-1.224587,-1.478207)--(-0.928911,-1.531105)--(-0.924449,-1.388587)--cycle;
\path[fill=blue!21!white,draw=blue!21!white,line width=.12pt] (0.924449,-1.388587)--(0.928911,-1.531105)--(1.224587,-1.478207)--(1.218704,-1.335944)--cycle;
\path[fill=blue!21!white,draw=blue!21!white,line width=.12pt] (-0.928911,-1.531105)--(-0.924449,-1.658912)--(-0.621290,-1.696881)--(-0.624289,-1.569256)--cycle;
\path[fill=blue!23!white,draw=blue!23!white,line width=.12pt] (0.624289,-1.569256)--(0.621290,-1.696881)--(0.924449,-1.658912)--(0.928911,-1.531105)--cycle;
\path[fill=blue!13!white,draw=blue!13!white,line width=.12pt] (-0.250539,-0.592374)--(-0.265249,-0.706475)--(-0.000000,-0.712991)--(-0.000000,-0.598528)--cycle;
\path[fill=blue!13!white,draw=blue!13!white,line width=.12pt] (-0.000000,-0.598528)--(-0.000000,-0.712991)--(0.265249,-0.706475)--(0.250539,-0.592374)--cycle;
\path[fill=blue!14!white,draw=blue!14!white,line width=.12pt] (0.825702,-0.784928)--(0.860893,-0.929094)--(1.134919,-0.880070)--(1.088526,-0.737908)--cycle;
\path[fill=blue!12!white,draw=blue!12!white,line width=.12pt] (-1.088526,-0.737908)--(-1.134919,-0.880070)--(-0.860893,-0.929094)--(-0.825702,-0.784928)--cycle;
\path[fill=blue!12!white,draw=blue!12!white,line width=.12pt] (-0.785553,-0.654728)--(-0.825702,-0.784928)--(-0.554926,-0.818841)--(-0.527943,-0.686992)--cycle;
\path[fill=blue!13!white,draw=blue!13!white,line width=.12pt] (0.527943,-0.686992)--(0.554926,-0.818841)--(0.825702,-0.784928)--(0.785553,-0.654728)--cycle;
\path[fill=blue!16!white,draw=blue!16!white,line width=.12pt] (-1.201283,-1.184946)--(-1.218704,-1.335944)--(-0.924449,-1.388587)--(-0.911234,-1.236837)--cycle;
\path[fill=blue!19!white,draw=blue!19!white,line width=.12pt] (0.911234,-1.236837)--(0.924449,-1.388587)--(1.218704,-1.335944)--(1.201283,-1.184946)--cycle;
\path[fill=blue!21!white,draw=blue!21!white,line width=.12pt] (-0.624289,-1.569256)--(-0.621290,-1.696881)--(-0.312148,-1.719809)--(-0.313655,-1.592296)--cycle;
\path[fill=blue!22!white,draw=blue!22!white,line width=.12pt] (0.313655,-1.592296)--(0.312148,-1.719809)--(0.621290,-1.696881)--(0.624289,-1.569256)--cycle;
\path[fill=blue!13!white,draw=blue!13!white,line width=.12pt] (-1.134919,-0.880070)--(-1.172992,-1.031016)--(-0.889774,-1.081685)--(-0.860893,-0.929094)--cycle;
\path[fill=blue!15!white,draw=blue!15!white,line width=.12pt] (0.860893,-0.929094)--(0.889774,-1.081685)--(1.172992,-1.031016)--(1.134919,-0.880070)--cycle;
\path[fill=blue!14!white,draw=blue!14!white,line width=.12pt] (-1.172992,-1.031016)--(-1.201283,-1.184946)--(-0.911234,-1.236837)--(-0.889774,-1.081685)--cycle;
\path[fill=blue!17!white,draw=blue!17!white,line width=.12pt] (0.889774,-1.081685)--(0.911234,-1.236837)--(1.201283,-1.184946)--(1.172992,-1.031016)--cycle;
\path[fill=blue!18!white,draw=blue!18!white,line width=.12pt] (-0.924449,-1.388587)--(-0.928911,-1.531105)--(-0.624289,-1.569256)--(-0.621290,-1.426556)--cycle;
\path[fill=blue!20!white,draw=blue!20!white,line width=.12pt] (0.621290,-1.426556)--(0.624289,-1.569256)--(0.928911,-1.531105)--(0.924449,-1.388587)--cycle;
\path[fill=blue!21!white,draw=blue!21!white,line width=.12pt] (-0.313655,-1.592296)--(-0.312148,-1.719809)--(-0.000000,-1.727477)--(-0.000000,-1.600000)--cycle;
\path[fill=blue!22!white,draw=blue!22!white,line width=.12pt] (-0.000000,-1.600000)--(-0.000000,-1.727477)--(0.312148,-1.719809)--(0.313655,-1.592296)--cycle;
\path[fill=blue!12!white,draw=blue!12!white,line width=.12pt] (-0.527943,-0.686992)--(-0.554926,-0.818841)--(-0.278805,-0.839321)--(-0.265249,-0.706475)--cycle;
\path[fill=blue!13!white,draw=blue!13!white,line width=.12pt] (0.265249,-0.706475)--(0.278805,-0.839321)--(0.554926,-0.818841)--(0.527943,-0.686992)--cycle;
\path[fill=blue!12!white,draw=blue!12!white,line width=.12pt] (-0.825702,-0.784928)--(-0.860893,-0.929094)--(-0.578577,-0.964452)--(-0.554926,-0.818841)--cycle;
\path[fill=blue!14!white,draw=blue!14!white,line width=.12pt] (0.554926,-0.818841)--(0.578577,-0.964452)--(0.860893,-0.929094)--(0.825702,-0.784928)--cycle;
\path[fill=blue!13!white,draw=blue!13!white,line width=.12pt] (-0.000000,-0.712991)--(-0.000000,-0.846169)--(0.278805,-0.839321)--(0.265249,-0.706475)--cycle;
\path[fill=blue!13!white,draw=blue!13!white,line width=.12pt] (-0.265249,-0.706475)--(-0.278805,-0.839321)--(-0.000000,-0.846169)--(-0.000000,-0.712991)--cycle;
\path[fill=blue!16!white,draw=blue!16!white,line width=.12pt] (-0.911234,-1.236837)--(-0.924449,-1.388587)--(-0.621290,-1.426556)--(-0.612409,-1.274262)--cycle;
\path[fill=blue!18!white,draw=blue!18!white,line width=.12pt] (0.612409,-1.274262)--(0.621290,-1.426556)--(0.924449,-1.388587)--(0.911234,-1.236837)--cycle;
\path[fill=blue!18!white,draw=blue!18!white,line width=.12pt] (-0.621290,-1.426556)--(-0.624289,-1.569256)--(-0.313655,-1.592296)--(-0.312148,-1.449484)--cycle;
\path[fill=blue!20!white,draw=blue!20!white,line width=.12pt] (0.312148,-1.449484)--(0.313655,-1.592296)--(0.624289,-1.569256)--(0.621290,-1.426556)--cycle;
\path[fill=blue!13!white,draw=blue!13!white,line width=.12pt] (-0.860893,-0.929094)--(-0.889774,-1.081685)--(-0.597986,-1.118229)--(-0.578577,-0.964452)--cycle;
\path[fill=blue!15!white,draw=blue!15!white,line width=.12pt] (0.578577,-0.964452)--(0.597986,-1.118229)--(0.889774,-1.081685)--(0.860893,-0.929094)--cycle;
\path[fill=blue!14!white,draw=blue!14!white,line width=.12pt] (-0.889774,-1.081685)--(-0.911234,-1.236837)--(-0.612409,-1.274262)--(-0.597986,-1.118229)--cycle;
\path[fill=blue!16!white,draw=blue!16!white,line width=.12pt] (0.597986,-1.118229)--(0.612409,-1.274262)--(0.911234,-1.236837)--(0.889774,-1.081685)--cycle;
\path[fill=blue!13!white,draw=blue!13!white,line width=.12pt] (-0.554926,-0.818841)--(-0.578577,-0.964452)--(-0.290688,-0.985804)--(-0.278805,-0.839321)--cycle;
\path[fill=blue!13!white,draw=blue!13!white,line width=.12pt] (0.278805,-0.839321)--(0.290688,-0.985804)--(0.578577,-0.964452)--(0.554926,-0.818841)--cycle;
\path[fill=blue!19!white,draw=blue!19!white,line width=.12pt] (-0.312148,-1.449484)--(-0.313655,-1.592296)--(-0.000000,-1.600000)--(-0.000000,-1.457152)--cycle;
\path[fill=blue!19!white,draw=blue!19!white,line width=.12pt] (-0.000000,-1.457152)--(-0.000000,-1.600000)--(0.313655,-1.592296)--(0.312148,-1.449484)--cycle;
\path[fill=blue!13!white,draw=blue!13!white,line width=.12pt] (-0.278805,-0.839321)--(-0.290688,-0.985804)--(-0.000000,-0.992945)--(-0.000000,-0.846169)--cycle;
\path[fill=blue!13!white,draw=blue!13!white,line width=.12pt] (-0.000000,-0.846169)--(-0.000000,-0.992945)--(0.290688,-0.985804)--(0.278805,-0.839321)--cycle;
\path[fill=blue!16!white,draw=blue!16!white,line width=.12pt] (-0.612409,-1.274262)--(-0.621290,-1.426556)--(-0.312148,-1.449484)--(-0.307686,-1.296863)--cycle;
\path[fill=blue!18!white,draw=blue!18!white,line width=.12pt] (0.307686,-1.296863)--(0.312148,-1.449484)--(0.621290,-1.426556)--(0.612409,-1.274262)--cycle;
\path[fill=blue!13!white,draw=blue!13!white,line width=.12pt] (-0.578577,-0.964452)--(-0.597986,-1.118229)--(-0.300440,-1.140298)--(-0.290688,-0.985804)--cycle;
\path[fill=blue!14!white,draw=blue!14!white,line width=.12pt] (0.290688,-0.985804)--(0.300440,-1.140298)--(0.597986,-1.118229)--(0.578577,-0.964452)--cycle;
\path[fill=blue!15!white,draw=blue!15!white,line width=.12pt] (-0.597986,-1.118229)--(-0.612409,-1.274262)--(-0.307686,-1.296863)--(-0.300440,-1.140298)--cycle;
\path[fill=blue!16!white,draw=blue!16!white,line width=.12pt] (0.300440,-1.140298)--(0.307686,-1.296863)--(0.612409,-1.274262)--(0.597986,-1.118229)--cycle;
\path[fill=blue!17!white,draw=blue!17!white,line width=.12pt] (-0.000000,-1.304421)--(-0.000000,-1.457152)--(0.312148,-1.449484)--(0.307686,-1.296863)--cycle;
\path[fill=blue!17!white,draw=blue!17!white,line width=.12pt] (-0.307686,-1.296863)--(-0.312148,-1.449484)--(-0.000000,-1.457152)--(-0.000000,-1.304421)--cycle;
\path[fill=blue!14!white,draw=blue!14!white,line width=.12pt] (-0.290688,-0.985804)--(-0.300440,-1.140298)--(-0.000000,-1.147677)--(-0.000000,-0.992945)--cycle;
\path[fill=blue!14!white,draw=blue!14!white,line width=.12pt] (-0.000000,-0.992945)--(-0.000000,-1.147677)--(0.300440,-1.140298)--(0.290688,-0.985804)--cycle;
\path[fill=blue!15!white,draw=blue!15!white,line width=.12pt] (-0.300440,-1.140298)--(-0.307686,-1.296863)--(-0.000000,-1.304421)--(-0.000000,-1.147677)--cycle;
\path[fill=blue!15!white,draw=blue!15!white,line width=.12pt] (-0.000000,-1.147677)--(-0.000000,-1.304421)--(0.307686,-1.296863)--(0.300440,-1.140298)--cycle;

\draw[green!45!black,very thick,domain=10:85,samples=60]
plot({(2.4-.8*cos(\x))*cos(130)},
{.5*(2.4-.8*cos(\x))*sin(130)+.8660254038*.8*sin(\x)});
\draw[blue!75!black,very thick,domain=100:160,samples=60]
plot({(2.4-.8*cos(40))*cos(\x)},
{.5*(2.4-.8*cos(40))*sin(\x)+.8660254038*.8*sin(40)});
\tp{Ageom}{40}{130}
\fill(Ageom)circle(1.5pt)node[below right]{$A$};
\draw[green!45!black,->,thick](Ageom)--++
({.95*sin(40)*cos(130)},
{.95*(.5*sin(40)*sin(130)+.8660254038*cos(40))})
node[above left]{$\mathbf e_\theta,\ \sigma_2$};
\draw[blue!75!black,->,thick](Ageom)--++
({-.95*sin(130)},{.95*.5*cos(130)});
\draw[thin](-1.94,.80)--(-3.34,.58);
\node[anchor=east,blue!75!black]at(-3.38,.58){$\mathbf e_\varphi,\ \sigma_1$};
\draw[red!75!black,->,thick](Ageom)--++
({-.95*cos(40)*cos(130)},
{.95*(-.5*cos(40)*sin(130)+.8660254038*sin(40))})
node[above right]{$\mathbf n$};
\node at(0,-2.60){两个力均在膜的切平面内};
\end{scope}
\begin{scope}[xshift=6.4cm,yshift=.45cm]
\draw[gray!70](0,0)circle(1.12);
\draw[->](-1.55,0)--(1.65,0)node[right]{$\mathbf e_\rho$};
\draw[->](0,-1.4)--(0,1.72)node[above]{$z$};
\coordinate(Am)at({-1.12*cos(40)},{1.12*sin(40)});
\draw[thin](0,0)--(Am)node[midway,above right=3pt]{$r$};
\draw(-.37,0)arc(180:140:.37);
\node at(-.57,.21){$\theta$};
\fill(Am)circle(1.5pt)node[below left]{$A$};
\draw[green!45!black,->,thick](Am)--++({.85*sin(40)},{.85*cos(40)})
node[above]{$\mathbf e_\theta$};
\draw[red!75!black,->,thick](Am)--++({-.66*cos(40)},{.66*sin(40)})
node[left]{$\mathbf n$};
\node at(0,-1.95){$\kappa_\theta=\dfrac1r$};
\node at(0,-2.65){$\kappa_\varphi=-\dfrac{\cos\theta}{R-r\cos\theta}$};
\end{scope}
\end{tikzpicture}
% 6. The crack follows phi at theta=0; maximum tension acts along theta.
\begin{tikzpicture}[font=\small]
\begin{scope}
% Automatically generated vector mesh; parameters and camera in build/torus-mesh-q03.json.
\path[fill=blue!39!white,draw=blue!39!white,line width=.12pt] (0.000000,1.304421)--(0.000000,1.147677)--(-0.300440,1.140298)--(-0.307686,1.296863)--cycle;
\path[fill=blue!39!white,draw=blue!39!white,line width=.12pt] (0.307686,1.296863)--(0.300440,1.140298)--(0.000000,1.147677)--(0.000000,1.304421)--cycle;
\path[fill=blue!40!white,draw=blue!40!white,line width=.12pt] (0.300440,1.140298)--(0.290688,0.985804)--(0.000000,0.992945)--(0.000000,1.147677)--cycle;
\path[fill=blue!40!white,draw=blue!40!white,line width=.12pt] (0.000000,1.147677)--(0.000000,0.992945)--(-0.290688,0.985804)--(-0.300440,1.140298)--cycle;
\path[fill=blue!37!white,draw=blue!37!white,line width=.12pt] (0.312148,1.449484)--(0.307686,1.296863)--(0.000000,1.304421)--(0.000000,1.457152)--cycle;
\path[fill=blue!37!white,draw=blue!37!white,line width=.12pt] (0.000000,1.457152)--(0.000000,1.304421)--(-0.307686,1.296863)--(-0.312148,1.449484)--cycle;
\path[fill=blue!39!white,draw=blue!39!white,line width=.12pt] (0.612409,1.274262)--(0.597986,1.118229)--(0.300440,1.140298)--(0.307686,1.296863)--cycle;
\path[fill=blue!38!white,draw=blue!38!white,line width=.12pt] (-0.307686,1.296863)--(-0.300440,1.140298)--(-0.597986,1.118229)--(-0.612409,1.274262)--cycle;
\path[fill=blue!41!white,draw=blue!41!white,line width=.12pt] (0.597986,1.118229)--(0.578577,0.964452)--(0.290688,0.985804)--(0.300440,1.140298)--cycle;
\path[fill=blue!40!white,draw=blue!40!white,line width=.12pt] (-0.300440,1.140298)--(-0.290688,0.985804)--(-0.578577,0.964452)--(-0.597986,1.118229)--cycle;
\path[fill=blue!38!white,draw=blue!38!white,line width=.12pt] (0.621290,1.426556)--(0.612409,1.274262)--(0.307686,1.296863)--(0.312148,1.449484)--cycle;
\path[fill=blue!36!white,draw=blue!36!white,line width=.12pt] (-0.312148,1.449484)--(-0.307686,1.296863)--(-0.612409,1.274262)--(-0.621290,1.426556)--cycle;
\path[fill=blue!41!white,draw=blue!41!white,line width=.12pt] (0.290688,0.985804)--(0.278805,0.839321)--(0.000000,0.846169)--(0.000000,0.992945)--cycle;
\path[fill=blue!41!white,draw=blue!41!white,line width=.12pt] (0.000000,0.992945)--(0.000000,0.846169)--(-0.278805,0.839321)--(-0.290688,0.985804)--cycle;
\path[fill=blue!35!white,draw=blue!35!white,line width=.12pt] (0.313655,1.592296)--(0.312148,1.449484)--(0.000000,1.457152)--(0.000000,1.600000)--cycle;
\path[fill=blue!35!white,draw=blue!35!white,line width=.12pt] (0.000000,1.600000)--(0.000000,1.457152)--(-0.312148,1.449484)--(-0.313655,1.592296)--cycle;
\path[fill=blue!41!white,draw=blue!41!white,line width=.12pt] (-0.290688,0.985804)--(-0.278805,0.839321)--(-0.554926,0.818841)--(-0.578577,0.964452)--cycle;
\path[fill=blue!41!white,draw=blue!41!white,line width=.12pt] (0.578577,0.964452)--(0.554926,0.818841)--(0.278805,0.839321)--(0.290688,0.985804)--cycle;
\path[fill=blue!38!white,draw=blue!38!white,line width=.12pt] (-0.612409,1.274262)--(-0.597986,1.118229)--(-0.889774,1.081685)--(-0.911234,1.236837)--cycle;
\path[fill=blue!40!white,draw=blue!40!white,line width=.12pt] (0.911234,1.236837)--(0.889774,1.081685)--(0.597986,1.118229)--(0.612409,1.274262)--cycle;
\path[fill=blue!41!white,draw=blue!41!white,line width=.12pt] (0.889774,1.081685)--(0.860893,0.929094)--(0.578577,0.964452)--(0.597986,1.118229)--cycle;
\path[fill=blue!39!white,draw=blue!39!white,line width=.12pt] (-0.597986,1.118229)--(-0.578577,0.964452)--(-0.860893,0.929094)--(-0.889774,1.081685)--cycle;
\path[fill=blue!36!white,draw=blue!36!white,line width=.12pt] (0.624289,1.569256)--(0.621290,1.426556)--(0.312148,1.449484)--(0.313655,1.592296)--cycle;
\path[fill=blue!34!white,draw=blue!34!white,line width=.12pt] (-0.313655,1.592296)--(-0.312148,1.449484)--(-0.621290,1.426556)--(-0.624289,1.569256)--cycle;
\path[fill=blue!38!white,draw=blue!38!white,line width=.12pt] (0.924449,1.388587)--(0.911234,1.236837)--(0.612409,1.274262)--(0.621290,1.426556)--cycle;
\path[fill=blue!36!white,draw=blue!36!white,line width=.12pt] (-0.621290,1.426556)--(-0.612409,1.274262)--(-0.911234,1.236837)--(-0.924449,1.388587)--cycle;
\path[fill=blue!41!white,draw=blue!41!white,line width=.12pt] (0.278805,0.839321)--(0.265249,0.706475)--(0.000000,0.712991)--(0.000000,0.846169)--cycle;
\path[fill=blue!41!white,draw=blue!41!white,line width=.12pt] (0.000000,0.846169)--(0.000000,0.712991)--(-0.265249,0.706475)--(-0.278805,0.839321)--cycle;
\path[fill=blue!42!white,draw=blue!42!white,line width=.12pt] (0.860893,0.929094)--(0.825702,0.784928)--(0.554926,0.818841)--(0.578577,0.964452)--cycle;
\path[fill=blue!40!white,draw=blue!40!white,line width=.12pt] (-0.578577,0.964452)--(-0.554926,0.818841)--(-0.825702,0.784928)--(-0.860893,0.929094)--cycle;
\path[fill=blue!42!white,draw=blue!42!white,line width=.12pt] (0.554926,0.818841)--(0.527943,0.686992)--(0.265249,0.706475)--(0.278805,0.839321)--cycle;
\path[fill=blue!41!white,draw=blue!41!white,line width=.12pt] (-0.278805,0.839321)--(-0.265249,0.706475)--(-0.527943,0.686992)--(-0.554926,0.818841)--cycle;
\path[fill=blue!33!white,draw=blue!33!white,line width=.12pt] (0.312148,1.719809)--(0.313655,1.592296)--(0.000000,1.600000)--(0.000000,1.727477)--cycle;
\path[fill=blue!32!white,draw=blue!32!white,line width=.12pt] (0.000000,1.727477)--(0.000000,1.600000)--(-0.313655,1.592296)--(-0.312148,1.719809)--cycle;
\path[fill=blue!36!white,draw=blue!36!white,line width=.12pt] (0.928911,1.531105)--(0.924449,1.388587)--(0.621290,1.426556)--(0.624289,1.569256)--cycle;
\path[fill=blue!34!white,draw=blue!34!white,line width=.12pt] (-0.624289,1.569256)--(-0.621290,1.426556)--(-0.924449,1.388587)--(-0.928911,1.531105)--cycle;
\path[fill=blue!37!white,draw=blue!37!white,line width=.12pt] (-0.911234,1.236837)--(-0.889774,1.081685)--(-1.172992,1.031016)--(-1.201283,1.184946)--cycle;
\path[fill=blue!40!white,draw=blue!40!white,line width=.12pt] (1.201283,1.184946)--(1.172992,1.031016)--(0.889774,1.081685)--(0.911234,1.236837)--cycle;
\path[fill=blue!41!white,draw=blue!41!white,line width=.12pt] (1.172992,1.031016)--(1.134919,0.880070)--(0.860893,0.929094)--(0.889774,1.081685)--cycle;
\path[fill=blue!39!white,draw=blue!39!white,line width=.12pt] (-0.889774,1.081685)--(-0.860893,0.929094)--(-1.134919,0.880070)--(-1.172992,1.031016)--cycle;
\path[fill=blue!32!white,draw=blue!32!white,line width=.12pt] (-0.312148,1.719809)--(-0.313655,1.592296)--(-0.624289,1.569256)--(-0.621290,1.696881)--cycle;
\path[fill=blue!33!white,draw=blue!33!white,line width=.12pt] (0.621290,1.696881)--(0.624289,1.569256)--(0.313655,1.592296)--(0.312148,1.719809)--cycle;
\path[fill=blue!38!white,draw=blue!38!white,line width=.12pt] (1.218704,1.335944)--(1.201283,1.184946)--(0.911234,1.236837)--(0.924449,1.388587)--cycle;
\path[fill=blue!35!white,draw=blue!35!white,line width=.12pt] (-0.924449,1.388587)--(-0.911234,1.236837)--(-1.201283,1.184946)--(-1.218704,1.335944)--cycle;
\path[fill=blue!42!white,draw=blue!42!white,line width=.12pt] (0.825702,0.784928)--(0.785553,0.654728)--(0.527943,0.686992)--(0.554926,0.818841)--cycle;
\path[fill=blue!41!white,draw=blue!41!white,line width=.12pt] (-0.554926,0.818841)--(-0.527943,0.686992)--(-0.785553,0.654728)--(-0.825702,0.784928)--cycle;
\path[fill=blue!42!white,draw=blue!42!white,line width=.12pt] (1.134919,0.880070)--(1.088526,0.737908)--(0.825702,0.784928)--(0.860893,0.929094)--cycle;
\path[fill=blue!40!white,draw=blue!40!white,line width=.12pt] (-0.860893,0.929094)--(-0.825702,0.784928)--(-1.088526,0.737908)--(-1.134919,0.880070)--cycle;
\path[fill=blue!41!white,draw=blue!41!white,line width=.12pt] (0.265249,0.706475)--(0.250539,0.592374)--(0.000000,0.598528)--(0.000000,0.712991)--cycle;
\path[fill=blue!41!white,draw=blue!41!white,line width=.12pt] (0.000000,0.712991)--(0.000000,0.598528)--(-0.250539,0.592374)--(-0.265249,0.706475)--cycle;
\path[fill=blue!33!white,draw=blue!33!white,line width=.12pt] (0.924449,1.658912)--(0.928911,1.531105)--(0.624289,1.569256)--(0.621290,1.696881)--cycle;
\path[fill=blue!31!white,draw=blue!31!white,line width=.12pt] (-0.621290,1.696881)--(-0.624289,1.569256)--(-0.928911,1.531105)--(-0.924449,1.658912)--cycle;
\path[fill=blue!33!white,draw=blue!33!white,line width=.12pt] (-0.928911,1.531105)--(-0.924449,1.388587)--(-1.218704,1.335944)--(-1.224587,1.478207)--cycle;
\path[fill=blue!36!white,draw=blue!36!white,line width=.12pt] (1.224587,1.478207)--(1.218704,1.335944)--(0.924449,1.388587)--(0.928911,1.531105)--cycle;
\path[fill=blue!41!white,draw=blue!41!white,line width=.12pt] (-0.265249,0.706475)--(-0.250539,0.592374)--(-0.498665,0.573971)--(-0.527943,0.686992)--cycle;
\path[fill=blue!41!white,draw=blue!41!white,line width=.12pt] (0.527943,0.686992)--(0.498665,0.573971)--(0.250539,0.592374)--(0.265249,0.706475)--cycle;
\path[fill=blue!30!white,draw=blue!30!white,line width=.12pt] (0.000000,1.834683)--(0.000000,1.727477)--(-0.312148,1.719809)--(-0.307686,1.827125)--cycle;
\path[fill=blue!30!white,draw=blue!30!white,line width=.12pt] (0.307686,1.827125)--(0.312148,1.719809)--(0.000000,1.727477)--(0.000000,1.834683)--cycle;
\path[fill=blue!40!white,draw=blue!40!white,line width=.12pt] (1.479763,1.119090)--(1.444914,0.966711)--(1.172992,1.031016)--(1.201283,1.184946)--cycle;
\path[fill=blue!37!white,draw=blue!37!white,line width=.12pt] (-1.201283,1.184946)--(-1.172992,1.031016)--(-1.444914,0.966711)--(-1.479763,1.119090)--cycle;
\path[fill=blue!42!white,draw=blue!42!white,line width=.12pt] (1.088526,0.737908)--(1.035598,0.609994)--(0.785553,0.654728)--(0.825702,0.784928)--cycle;
\path[fill=blue!40!white,draw=blue!40!white,line width=.12pt] (-0.825702,0.784928)--(-0.785553,0.654728)--(-1.035598,0.609994)--(-1.088526,0.737908)--cycle;
\path[fill=blue!41!white,draw=blue!41!white,line width=.12pt] (1.444914,0.966711)--(1.398014,0.817853)--(1.134919,0.880070)--(1.172992,1.031016)--cycle;
\path[fill=blue!38!white,draw=blue!38!white,line width=.12pt] (-1.172992,1.031016)--(-1.134919,0.880070)--(-1.398014,0.817853)--(-1.444914,0.966711)--cycle;
\path[fill=blue!29!white,draw=blue!29!white,line width=.12pt] (-0.307686,1.827125)--(-0.312148,1.719809)--(-0.621290,1.696881)--(-0.612409,1.804524)--cycle;
\path[fill=blue!30!white,draw=blue!30!white,line width=.12pt] (0.612409,1.804524)--(0.621290,1.696881)--(0.312148,1.719809)--(0.307686,1.827125)--cycle;
\path[fill=blue!41!white,draw=blue!41!white,line width=.12pt] (0.785553,0.654728)--(0.741989,0.543496)--(0.498665,0.573971)--(0.527943,0.686992)--cycle;
\path[fill=blue!41!white,draw=blue!41!white,line width=.12pt] (-0.527943,0.686992)--(-0.498665,0.573971)--(-0.741989,0.543496)--(-0.785553,0.654728)--cycle;
\path[fill=blue!38!white,draw=blue!38!white,line width=.12pt] (1.501223,1.269133)--(1.479763,1.119090)--(1.201283,1.184946)--(1.218704,1.335944)--cycle;
\path[fill=blue!35!white,draw=blue!35!white,line width=.12pt] (-1.218704,1.335944)--(-1.201283,1.184946)--(-1.479763,1.119090)--(-1.501223,1.269133)--cycle;
\path[fill=blue!42!white,draw=blue!42!white,line width=.12pt] (1.398014,0.817853)--(1.340867,0.678234)--(1.088526,0.737908)--(1.134919,0.880070)--cycle;
\path[fill=blue!39!white,draw=blue!39!white,line width=.12pt] (-1.134919,0.880070)--(-1.088526,0.737908)--(-1.340867,0.678234)--(-1.398014,0.817853)--cycle;
\path[fill=blue!31!white,draw=blue!31!white,line width=.12pt] (-0.924449,1.658912)--(-0.928911,1.531105)--(-1.224587,1.478207)--(-1.218704,1.606269)--cycle;
\path[fill=blue!34!white,draw=blue!34!white,line width=.12pt] (1.218704,1.606269)--(1.224587,1.478207)--(0.928911,1.531105)--(0.924449,1.658912)--cycle;
\path[fill=blue!31!white,draw=blue!31!white,line width=.12pt] (0.911234,1.767098)--(0.924449,1.658912)--(0.621290,1.696881)--(0.612409,1.804524)--cycle;
\path[fill=blue!29!white,draw=blue!29!white,line width=.12pt] (-0.612409,1.804524)--(-0.621290,1.696881)--(-0.924449,1.658912)--(-0.911234,1.767098)--cycle;
\path[fill=blue!40!white,draw=blue!40!white,line width=.12pt] (0.250539,0.592374)--(0.235241,0.501401)--(0.000000,0.507180)--(0.000000,0.598528)--cycle;
\path[fill=blue!40!white,draw=blue!40!white,line width=.12pt] (0.000000,0.598528)--(0.000000,0.507180)--(-0.235241,0.501401)--(-0.250539,0.592374)--cycle;
\path[fill=blue!36!white,draw=blue!36!white,line width=.12pt] (1.508470,1.411074)--(1.501223,1.269133)--(1.218704,1.335944)--(1.224587,1.478207)--cycle;
\path[fill=blue!32!white,draw=blue!32!white,line width=.12pt] (-1.224587,1.478207)--(-1.218704,1.335944)--(-1.501223,1.269133)--(-1.508470,1.411074)--cycle;
\path[fill=blue!40!white,draw=blue!40!white,line width=.12pt] (-0.785553,0.654728)--(-0.741989,0.543496)--(-0.978167,0.501243)--(-1.035598,0.609994)--cycle;
\path[fill=blue!41!white,draw=blue!41!white,line width=.12pt] (1.035598,0.609994)--(0.978167,0.501243)--(0.741989,0.543496)--(0.785553,0.654728)--cycle;
\path[fill=blue!42!white,draw=blue!42!white,line width=.12pt] (1.340867,0.678234)--(1.275669,0.553222)--(1.035598,0.609994)--(1.088526,0.737908)--cycle;
\path[fill=blue!40!white,draw=blue!40!white,line width=.12pt] (-1.088526,0.737908)--(-1.035598,0.609994)--(-1.275669,0.553222)--(-1.340867,0.678234)--cycle;
\path[fill=blue!40!white,draw=blue!40!white,line width=.12pt] (-0.250539,0.592374)--(-0.235241,0.501401)--(-0.468217,0.484122)--(-0.498665,0.573971)--cycle;
\path[fill=blue!40!white,draw=blue!40!white,line width=.12pt] (0.498665,0.573971)--(0.468217,0.484122)--(0.235241,0.501401)--(0.250539,0.592374)--cycle;
\path[fill=blue!27!white,draw=blue!27!white,line width=.12pt] (0.000000,1.917498)--(0.000000,1.834683)--(-0.307686,1.827125)--(-0.300440,1.910118)--cycle;
\path[fill=blue!27!white,draw=blue!27!white,line width=.12pt] (0.300440,1.910118)--(0.307686,1.827125)--(0.000000,1.834683)--(0.000000,1.917498)--cycle;
\path[fill=blue!41!white,draw=blue!41!white,line width=.12pt] (1.702920,0.889390)--(1.647647,0.743041)--(1.398014,0.817853)--(1.444914,0.966711)--cycle;
\path[fill=blue!38!white,draw=blue!38!white,line width=.12pt] (-1.444914,0.966711)--(-1.398014,0.817853)--(-1.647647,0.743041)--(-1.702920,0.889390)--cycle;
\path[fill=blue!40!white,draw=blue!40!white,line width=.12pt] (1.743993,1.039904)--(1.702920,0.889390)--(1.444914,0.966711)--(1.479763,1.119090)--cycle;
\path[fill=blue!36!white,draw=blue!36!white,line width=.12pt] (-1.479763,1.119090)--(-1.444914,0.966711)--(-1.702920,0.889390)--(-1.743993,1.039904)--cycle;
\path[fill=blue!40!white,draw=blue!40!white,line width=.12pt] (0.741989,0.543496)--(0.696683,0.455508)--(0.468217,0.484122)--(0.498665,0.573971)--cycle;
\path[fill=blue!40!white,draw=blue!40!white,line width=.12pt] (-0.498665,0.573971)--(-0.468217,0.484122)--(-0.696683,0.455508)--(-0.741989,0.543496)--cycle;
\path[fill=blue!28!white,draw=blue!28!white,line width=.12pt] (0.597986,1.888050)--(0.612409,1.804524)--(0.307686,1.827125)--(0.300440,1.910118)--cycle;
\path[fill=blue!26!white,draw=blue!26!white,line width=.12pt] (-0.300440,1.910118)--(-0.307686,1.827125)--(-0.612409,1.804524)--(-0.597986,1.888050)--cycle;
\path[fill=blue!39!white,draw=blue!39!white,line width=.12pt] (-1.398014,0.817853)--(-1.340867,0.678234)--(-1.580295,0.606480)--(-1.647647,0.743041)--cycle;
\path[fill=blue!42!white,draw=blue!42!white,line width=.12pt] (1.647647,0.743041)--(1.580295,0.606480)--(1.340867,0.678234)--(1.398014,0.817853)--cycle;
\path[fill=blue!34!white,draw=blue!34!white,line width=.12pt] (-1.501223,1.269133)--(-1.479763,1.119090)--(-1.743993,1.039904)--(-1.769285,1.188798)--cycle;
\path[fill=blue!38!white,draw=blue!38!white,line width=.12pt] (1.769285,1.188798)--(1.743993,1.039904)--(1.479763,1.119090)--(1.501223,1.269133)--cycle;
\path[fill=blue!31!white,draw=blue!31!white,line width=.12pt] (1.201283,1.715208)--(1.218704,1.606269)--(0.924449,1.658912)--(0.911234,1.767098)--cycle;
\path[fill=blue!28!white,draw=blue!28!white,line width=.12pt] (-0.911234,1.767098)--(-0.924449,1.658912)--(-1.218704,1.606269)--(-1.201283,1.715208)--cycle;
\path[fill=blue!30!white,draw=blue!30!white,line width=.12pt] (-1.218704,1.606269)--(-1.224587,1.478207)--(-1.508470,1.411074)--(-1.501223,1.539458)--cycle;
\path[fill=blue!34!white,draw=blue!34!white,line width=.12pt] (1.501223,1.539458)--(1.508470,1.411074)--(1.224587,1.478207)--(1.218704,1.606269)--cycle;
\path[fill=blue!40!white,draw=blue!40!white,line width=.12pt] (-1.035598,0.609994)--(-0.978167,0.501243)--(-1.204924,0.447619)--(-1.275669,0.553222)--cycle;
\path[fill=blue!41!white,draw=blue!41!white,line width=.12pt] (1.275669,0.553222)--(1.204924,0.447619)--(0.978167,0.501243)--(1.035598,0.609994)--cycle;
\path[fill=blue!26!white,draw=blue!26!white,line width=.12pt] (-0.597986,1.888050)--(-0.612409,1.804524)--(-0.911234,1.767098)--(-0.889774,1.851505)--cycle;
\path[fill=blue!28!white,draw=blue!28!white,line width=.12pt] (0.889774,1.851505)--(0.911234,1.767098)--(0.612409,1.804524)--(0.597986,1.888050)--cycle;
\path[fill=blue!40!white,draw=blue!40!white,line width=.12pt] (0.978167,0.501243)--(0.918440,0.415835)--(0.696683,0.455508)--(0.741989,0.543496)--cycle;
\path[fill=blue!40!white,draw=blue!40!white,line width=.12pt] (-0.741989,0.543496)--(-0.696683,0.455508)--(-0.918440,0.415835)--(-0.978167,0.501243)--cycle;
\path[fill=blue!40!white,draw=blue!40!white,line width=.12pt] (-1.340867,0.678234)--(-1.275669,0.553222)--(-1.503455,0.484957)--(-1.580295,0.606480)--cycle;
\path[fill=blue!42!white,draw=blue!42!white,line width=.12pt] (1.580295,0.606480)--(1.503455,0.484957)--(1.275669,0.553222)--(1.340867,0.678234)--cycle;
\path[fill=blue!32!white,draw=blue!32!white,line width=.12pt] (-1.508470,1.411074)--(-1.501223,1.269133)--(-1.769285,1.188798)--(-1.777825,1.330351)--cycle;
\path[fill=blue!36!white,draw=blue!36!white,line width=.12pt] (1.777825,1.330351)--(1.769285,1.188798)--(1.501223,1.269133)--(1.508470,1.411074)--cycle;
\path[fill=blue!39!white,draw=blue!39!white,line width=.12pt] (0.000000,0.507180)--(0.000000,0.442456)--(-0.219943,0.437053)--(-0.235241,0.501401)--cycle;
\path[fill=blue!39!white,draw=blue!39!white,line width=.12pt] (0.235241,0.501401)--(0.219943,0.437053)--(0.000000,0.442456)--(0.000000,0.507180)--cycle;
\path[fill=blue!39!white,draw=blue!39!white,line width=.12pt] (-0.235241,0.501401)--(-0.219943,0.437053)--(-0.437769,0.420898)--(-0.468217,0.484122)--cycle;
\path[fill=blue!39!white,draw=blue!39!white,line width=.12pt] (0.468217,0.484122)--(0.437769,0.420898)--(0.219943,0.437053)--(0.235241,0.501401)--cycle;
\path[fill=blue!31!white,draw=blue!31!white,line width=.12pt] (1.479763,1.649352)--(1.501223,1.539458)--(1.218704,1.606269)--(1.201283,1.715208)--cycle;
\path[fill=blue!27!white,draw=blue!27!white,line width=.12pt] (-1.201283,1.715208)--(-1.218704,1.606269)--(-1.501223,1.539458)--(-1.479763,1.649352)--cycle;
\path[fill=blue!24!white,draw=blue!24!white,line width=.12pt] (0.290688,1.965600)--(0.300440,1.910118)--(0.000000,1.917498)--(0.000000,1.972741)--cycle;
\path[fill=blue!24!white,draw=blue!24!white,line width=.12pt] (0.000000,1.972741)--(0.000000,1.917498)--(-0.300440,1.910118)--(-0.290688,1.965600)--cycle;
\path[fill=blue!41!white,draw=blue!41!white,line width=.12pt] (1.944527,0.799796)--(1.881411,0.656355)--(1.647647,0.743041)--(1.702920,0.889390)--cycle;
\path[fill=blue!37!white,draw=blue!37!white,line width=.12pt] (-1.702920,0.889390)--(-1.647647,0.743041)--(-1.881411,0.656355)--(-1.944527,0.799796)--cycle;
\path[fill=blue!39!white,draw=blue!39!white,line width=.12pt] (0.696683,0.455508)--(0.651378,0.394145)--(0.437769,0.420898)--(0.468217,0.484122)--cycle;
\path[fill=blue!39!white,draw=blue!39!white,line width=.12pt] (-0.468217,0.484122)--(-0.437769,0.420898)--(-0.651378,0.394145)--(-0.696683,0.455508)--cycle;
\path[fill=blue!40!white,draw=blue!40!white,line width=.12pt] (1.991426,0.948149)--(1.944527,0.799796)--(1.702920,0.889390)--(1.743993,1.039904)--cycle;
\path[fill=blue!35!white,draw=blue!35!white,line width=.12pt] (-1.743993,1.039904)--(-1.702920,0.889390)--(-1.944527,0.799796)--(-1.991426,0.948149)--cycle;
\path[fill=blue!28!white,draw=blue!28!white,line width=.12pt] (1.172992,1.800837)--(1.201283,1.715208)--(0.911234,1.767098)--(0.889774,1.851505)--cycle;
\path[fill=blue!25!white,draw=blue!25!white,line width=.12pt] (-0.889774,1.851505)--(-0.911234,1.767098)--(-1.201283,1.715208)--(-1.172992,1.800837)--cycle;
\path[fill=blue!40!white,draw=blue!40!white,line width=.12pt] (1.204924,0.447619)--(1.131352,0.365485)--(0.918440,0.415835)--(0.978167,0.501243)--cycle;
\path[fill=blue!40!white,draw=blue!40!white,line width=.12pt] (-0.978167,0.501243)--(-0.918440,0.415835)--(-1.131352,0.365485)--(-1.204924,0.447619)--cycle;
\path[fill=blue!40!white,draw=blue!40!white,line width=.12pt] (-1.275669,0.553222)--(-1.204924,0.447619)--(-1.420078,0.383140)--(-1.503455,0.484957)--cycle;
\path[fill=blue!41!white,draw=blue!41!white,line width=.12pt] (1.503455,0.484957)--(1.420078,0.383140)--(1.204924,0.447619)--(1.275669,0.553222)--cycle;
\path[fill=blue!29!white,draw=blue!29!white,line width=.12pt] (-1.501223,1.539458)--(-1.508470,1.411074)--(-1.777825,1.330351)--(-1.769285,1.459123)--cycle;
\path[fill=blue!34!white,draw=blue!34!white,line width=.12pt] (1.769285,1.459123)--(1.777825,1.330351)--(1.508470,1.411074)--(1.501223,1.539458)--cycle;
\path[fill=blue!25!white,draw=blue!25!white,line width=.12pt] (0.578577,1.944248)--(0.597986,1.888050)--(0.300440,1.910118)--(0.290688,1.965600)--cycle;
\path[fill=blue!24!white,draw=blue!24!white,line width=.12pt] (-0.290688,1.965600)--(-0.300440,1.910118)--(-0.597986,1.888050)--(-0.578577,1.944248)--cycle;
\path[fill=blue!42!white,draw=blue!42!white,line width=.12pt] (1.881411,0.656355)--(1.804504,0.523338)--(1.580295,0.606480)--(1.647647,0.743041)--cycle;
\path[fill=blue!38!white,draw=blue!38!white,line width=.12pt] (-1.647647,0.743041)--(-1.580295,0.606480)--(-1.804504,0.523338)--(-1.881411,0.656355)--cycle;
\path[fill=blue!38!white,draw=blue!38!white,line width=.12pt] (2.020307,1.095713)--(1.991426,0.948149)--(1.743993,1.039904)--(1.769285,1.188798)--cycle;
\path[fill=blue!33!white,draw=blue!33!white,line width=.12pt] (-1.769285,1.188798)--(-1.743993,1.039904)--(-1.991426,0.948149)--(-2.020307,1.095713)--cycle;
\path[fill=blue!39!white,draw=blue!39!white,line width=.12pt] (0.918440,0.415835)--(0.858714,0.357051)--(0.651378,0.394145)--(0.696683,0.455508)--cycle;
\path[fill=blue!39!white,draw=blue!39!white,line width=.12pt] (-0.696683,0.455508)--(-0.651378,0.394145)--(-0.858714,0.357051)--(-0.918440,0.415835)--cycle;
\path[fill=blue!39!white,draw=blue!39!white,line width=.12pt] (-1.580295,0.606480)--(-1.503455,0.484957)--(-1.716761,0.405857)--(-1.804504,0.523338)--cycle;
\path[fill=blue!42!white,draw=blue!42!white,line width=.12pt] (1.804504,0.523338)--(1.716761,0.405857)--(1.503455,0.484957)--(1.580295,0.606480)--cycle;
\path[fill=blue!25!white,draw=blue!25!white,line width=.12pt] (0.860893,1.908890)--(0.889774,1.851505)--(0.597986,1.888050)--(0.578577,1.944248)--cycle;
\path[fill=blue!23!white,draw=blue!23!white,line width=.12pt] (-0.578577,1.944248)--(-0.597986,1.888050)--(-0.889774,1.851505)--(-0.860893,1.908890)--cycle;
\path[fill=blue!31!white,draw=blue!31!white,line width=.12pt] (-1.777825,1.330351)--(-1.769285,1.188798)--(-2.020307,1.095713)--(-2.030059,1.236817)--cycle;
\path[fill=blue!36!white,draw=blue!36!white,line width=.12pt] (2.030059,1.236817)--(2.020307,1.095713)--(1.769285,1.188798)--(1.777825,1.330351)--cycle;
\path[fill=blue!37!white,draw=blue!37!white,line width=.12pt] (0.000000,0.442456)--(0.000000,0.406844)--(-0.205234,0.401803)--(-0.219943,0.437053)--cycle;
\path[fill=blue!37!white,draw=blue!37!white,line width=.12pt] (0.219943,0.437053)--(0.205234,0.401803)--(0.000000,0.406844)--(0.000000,0.442456)--cycle;
\path[fill=blue!28!white,draw=blue!28!white,line width=.12pt] (1.444914,1.736532)--(1.479763,1.649352)--(1.201283,1.715208)--(1.172992,1.800837)--cycle;
\path[fill=blue!25!white,draw=blue!25!white,line width=.12pt] (-1.172992,1.800837)--(-1.201283,1.715208)--(-1.479763,1.649352)--(-1.444914,1.736532)--cycle;
\path[fill=blue!37!white,draw=blue!37!white,line width=.12pt] (-0.219943,0.437053)--(-0.205234,0.401803)--(-0.408491,0.386728)--(-0.437769,0.420898)--cycle;
\path[fill=blue!37!white,draw=blue!37!white,line width=.12pt] (0.437769,0.420898)--(0.408491,0.386728)--(0.205234,0.401803)--(0.219943,0.437053)--cycle;
\path[fill=blue!40!white,draw=blue!40!white,line width=.12pt] (-1.204924,0.447619)--(-1.131352,0.365485)--(-1.333369,0.304943)--(-1.420078,0.383140)--cycle;
\path[fill=blue!40!white,draw=blue!40!white,line width=.12pt] (1.420078,0.383140)--(1.333369,0.304943)--(1.131352,0.365485)--(1.204924,0.447619)--cycle;
\path[fill=blue!27!white,draw=blue!27!white,line width=.12pt] (-1.479763,1.649352)--(-1.501223,1.539458)--(-1.769285,1.459123)--(-1.743993,1.570165)--cycle;
\path[fill=blue!31!white,draw=blue!31!white,line width=.12pt] (1.743993,1.570165)--(1.769285,1.459123)--(1.501223,1.539458)--(1.479763,1.649352)--cycle;
\path[fill=blue!39!white,draw=blue!39!white,line width=.12pt] (1.131352,0.365485)--(1.057780,0.309976)--(0.858714,0.357051)--(0.918440,0.415835)--cycle;
\path[fill=blue!39!white,draw=blue!39!white,line width=.12pt] (-0.918440,0.415835)--(-0.858714,0.357051)--(-1.057780,0.309976)--(-1.131352,0.365485)--cycle;
\path[fill=blue!40!white,draw=blue!40!white,line width=.12pt] (-1.503455,0.484957)--(-1.420078,0.383140)--(-1.621555,0.308427)--(-1.716761,0.405857)--cycle;
\path[fill=blue!41!white,draw=blue!41!white,line width=.12pt] (1.716761,0.405857)--(1.621555,0.308427)--(1.420078,0.383140)--(1.503455,0.484957)--cycle;
\path[fill=blue!25!white,draw=blue!25!white,line width=.12pt] (1.134919,1.859866)--(1.172992,1.800837)--(0.889774,1.851505)--(0.860893,1.908890)--cycle;
\path[fill=blue!23!white,draw=blue!23!white,line width=.12pt] (-0.860893,1.908890)--(-0.889774,1.851505)--(-1.172992,1.800837)--(-1.134919,1.859866)--cycle;
\path[fill=blue!22!white,draw=blue!22!white,line width=.12pt] (0.278805,1.991439)--(0.290688,1.965600)--(0.000000,1.972741)--(0.000000,1.998287)--cycle;
\path[fill=blue!21!white,draw=blue!21!white,line width=.12pt] (0.000000,1.998287)--(0.000000,1.972741)--(-0.290688,1.965600)--(-0.278805,1.991439)--cycle;
\path[fill=blue!37!white,draw=blue!37!white,line width=.12pt] (0.651378,0.394145)--(0.607814,0.361764)--(0.408491,0.386728)--(0.437769,0.420898)--cycle;
\path[fill=blue!37!white,draw=blue!37!white,line width=.12pt] (-0.437769,0.420898)--(-0.408491,0.386728)--(-0.607814,0.361764)--(-0.651378,0.394145)--cycle;
\path[fill=blue!41!white,draw=blue!41!white,line width=.12pt] (2.167407,0.698793)--(2.097056,0.558630)--(1.881411,0.656355)--(1.944527,0.799796)--cycle;
\path[fill=blue!37!white,draw=blue!37!white,line width=.12pt] (-1.944527,0.799796)--(-1.881411,0.656355)--(-2.097056,0.558630)--(-2.167407,0.698793)--cycle;
\path[fill=blue!34!white,draw=blue!34!white,line width=.12pt] (2.020307,1.366038)--(2.030059,1.236817)--(1.777825,1.330351)--(1.769285,1.459123)--cycle;
\path[fill=blue!28!white,draw=blue!28!white,line width=.12pt] (-1.769285,1.459123)--(-1.777825,1.330351)--(-2.030059,1.236817)--(-2.020307,1.366038)--cycle;
\path[fill=blue!22!white,draw=blue!22!white,line width=.12pt] (0.554926,1.970959)--(0.578577,1.944248)--(0.290688,1.965600)--(0.278805,1.991439)--cycle;
\path[fill=blue!21!white,draw=blue!21!white,line width=.12pt] (-0.278805,1.991439)--(-0.290688,1.965600)--(-0.578577,1.944248)--(-0.554926,1.970959)--cycle;
\path[fill=blue!40!white,draw=blue!40!white,line width=.12pt] (2.219681,0.844710)--(2.167407,0.698793)--(1.944527,0.799796)--(1.991426,0.948149)--cycle;
\path[fill=blue!35!white,draw=blue!35!white,line width=.12pt] (-1.991426,0.948149)--(-1.944527,0.799796)--(-2.167407,0.698793)--(-2.219681,0.844710)--cycle;
\path[fill=blue!42!white,draw=blue!42!white,line width=.12pt] (2.097056,0.558630)--(2.011334,0.429608)--(1.804504,0.523338)--(1.881411,0.656355)--cycle;
\path[fill=blue!38!white,draw=blue!38!white,line width=.12pt] (-1.881411,0.656355)--(-1.804504,0.523338)--(-2.011334,0.429608)--(-2.097056,0.558630)--cycle;
\path[fill=blue!37!white,draw=blue!37!white,line width=.12pt] (0.858714,0.357051)--(0.801283,0.327152)--(0.607814,0.361764)--(0.651378,0.394145)--cycle;
\path[fill=blue!38!white,draw=blue!38!white,line width=.12pt] (-0.651378,0.394145)--(-0.607814,0.361764)--(-0.801283,0.327152)--(-0.858714,0.357051)--cycle;
\path[fill=blue!33!white,draw=blue!33!white,line width=.12pt] (-2.020307,1.095713)--(-1.991426,0.948149)--(-2.219681,0.844710)--(-2.251872,0.990774)--cycle;
\path[fill=blue!38!white,draw=blue!38!white,line width=.12pt] (2.251872,0.990774)--(2.219681,0.844710)--(1.991426,0.948149)--(2.020307,1.095713)--cycle;
\path[fill=blue!22!white,draw=blue!22!white,line width=.12pt] (0.825702,1.937046)--(0.860893,1.908890)--(0.578577,1.944248)--(0.554926,1.970959)--cycle;
\path[fill=blue!21!white,draw=blue!21!white,line width=.12pt] (-0.554926,1.970959)--(-0.578577,1.944248)--(-0.860893,1.908890)--(-0.825702,1.937046)--cycle;
\path[fill=blue!39!white,draw=blue!39!white,line width=.12pt] (-1.804504,0.523338)--(-1.716761,0.405857)--(-1.913535,0.316685)--(-2.011334,0.429608)--cycle;
\path[fill=blue!42!white,draw=blue!42!white,line width=.12pt] (2.011334,0.429608)--(1.913535,0.316685)--(1.716761,0.405857)--(1.804504,0.523338)--cycle;
\path[fill=blue!39!white,draw=blue!39!white,line width=.12pt] (-1.131352,0.365485)--(-1.057780,0.309976)--(-1.246659,0.253371)--(-1.333369,0.304943)--cycle;
\path[fill=blue!39!white,draw=blue!39!white,line width=.12pt] (1.333369,0.304943)--(1.246659,0.253371)--(1.057780,0.309976)--(1.131352,0.365485)--cycle;
\path[fill=blue!40!white,draw=blue!40!white,line width=.12pt] (-1.420078,0.383140)--(-1.333369,0.304943)--(-1.522544,0.234792)--(-1.621555,0.308427)--cycle;
\path[fill=blue!40!white,draw=blue!40!white,line width=.12pt] (1.621555,0.308427)--(1.522544,0.234792)--(1.333369,0.304943)--(1.420078,0.383140)--cycle;
\path[fill=blue!28!white,draw=blue!28!white,line width=.12pt] (1.702920,1.659211)--(1.743993,1.570165)--(1.479763,1.649352)--(1.444914,1.736532)--cycle;
\path[fill=blue!24!white,draw=blue!24!white,line width=.12pt] (-1.444914,1.736532)--(-1.479763,1.649352)--(-1.743993,1.570165)--(-1.702920,1.659211)--cycle;
\path[fill=blue!25!white,draw=blue!25!white,line width=.12pt] (1.398014,1.797648)--(1.444914,1.736532)--(1.172992,1.800837)--(1.134919,1.859866)--cycle;
\path[fill=blue!22!white,draw=blue!22!white,line width=.12pt] (-1.134919,1.859866)--(-1.172992,1.800837)--(-1.444914,1.736532)--(-1.398014,1.797648)--cycle;
\path[fill=blue!35!white,draw=blue!35!white,line width=.12pt] (0.205234,0.401803)--(0.191677,0.397005)--(0.000000,0.401713)--(0.000000,0.406844)--cycle;
\path[fill=blue!35!white,draw=blue!35!white,line width=.12pt] (0.000000,0.406844)--(0.000000,0.401713)--(-0.191677,0.397005)--(-0.205234,0.401803)--cycle;
\path[fill=blue!30!white,draw=blue!30!white,line width=.12pt] (-2.030059,1.236817)--(-2.020307,1.095713)--(-2.251872,0.990774)--(-2.262742,1.131371)--cycle;
\path[fill=blue!36!white,draw=blue!36!white,line width=.12pt] (2.262742,1.131371)--(2.251872,0.990774)--(2.020307,1.095713)--(2.030059,1.236817)--cycle;
\path[fill=blue!35!white,draw=blue!35!white,line width=.12pt] (0.408491,0.386728)--(0.381508,0.382925)--(0.191677,0.397005)--(0.205234,0.401803)--cycle;
\path[fill=blue!35!white,draw=blue!35!white,line width=.12pt] (-0.205234,0.401803)--(-0.191677,0.397005)--(-0.381508,0.382925)--(-0.408491,0.386728)--cycle;
\path[fill=blue!31!white,draw=blue!31!white,line width=.12pt] (1.991426,1.478411)--(2.020307,1.366038)--(1.769285,1.459123)--(1.743993,1.570165)--cycle;
\path[fill=blue!26!white,draw=blue!26!white,line width=.12pt] (-1.743993,1.570165)--(-1.769285,1.459123)--(-2.020307,1.366038)--(-1.991426,1.478411)--cycle;
\path[fill=blue!37!white,draw=blue!37!white,line width=.12pt] (1.057780,0.309976)--(0.987036,0.283224)--(0.801283,0.327152)--(0.858714,0.357051)--cycle;
\path[fill=blue!38!white,draw=blue!38!white,line width=.12pt] (-0.858714,0.357051)--(-0.801283,0.327152)--(-0.987036,0.283224)--(-1.057780,0.309976)--cycle;
\path[fill=blue!41!white,draw=blue!41!white,line width=.12pt] (1.913535,0.316685)--(1.807416,0.224200)--(1.621555,0.308427)--(1.716761,0.405857)--cycle;
\path[fill=blue!40!white,draw=blue!40!white,line width=.12pt] (-1.716761,0.405857)--(-1.621555,0.308427)--(-1.807416,0.224200)--(-1.913535,0.316685)--cycle;
\path[fill=blue!20!white,draw=blue!20!white,line width=.12pt] (-0.825702,1.937046)--(-0.860893,1.908890)--(-1.134919,1.859866)--(-1.088526,1.890026)--cycle;
\path[fill=blue!22!white,draw=blue!22!white,line width=.12pt] (1.088526,1.890026)--(1.134919,1.859866)--(0.860893,1.908890)--(0.825702,1.937046)--cycle;
\path[fill=blue!35!white,draw=blue!35!white,line width=.12pt] (0.607814,0.361764)--(0.567664,0.359610)--(0.381508,0.382925)--(0.408491,0.386728)--cycle;
\path[fill=blue!36!white,draw=blue!36!white,line width=.12pt] (-0.408491,0.386728)--(-0.381508,0.382925)--(-0.567664,0.359610)--(-0.607814,0.361764)--cycle;
\path[fill=blue!19!white,draw=blue!19!white,line width=.12pt] (0.265249,1.986640)--(0.278805,1.991439)--(0.000000,1.998287)--(0.000000,1.993156)--cycle;
\path[fill=blue!19!white,draw=blue!19!white,line width=.12pt] (0.000000,1.993156)--(0.000000,1.998287)--(-0.278805,1.991439)--(-0.265249,1.986640)--cycle;
\path[fill=blue!19!white,draw=blue!19!white,line width=.12pt] (0.527943,1.967157)--(0.554926,1.970959)--(0.278805,1.991439)--(0.265249,1.986640)--cycle;
\path[fill=blue!19!white,draw=blue!19!white,line width=.12pt] (-0.265249,1.986640)--(-0.278805,1.991439)--(-0.554926,1.970959)--(-0.527943,1.967157)--cycle;
\path[fill=blue!39!white,draw=blue!39!white,line width=.12pt] (-1.333369,0.304943)--(-1.246659,0.253371)--(-1.423533,0.187782)--(-1.522544,0.234792)--cycle;
\path[fill=blue!39!white,draw=blue!39!white,line width=.12pt] (1.522544,0.234792)--(1.423533,0.187782)--(1.246659,0.253371)--(1.333369,0.304943)--cycle;
\path[fill=blue!28!white,draw=blue!28!white,line width=.12pt] (-2.020307,1.366038)--(-2.030059,1.236817)--(-2.262742,1.131371)--(-2.251872,1.261099)--cycle;
\path[fill=blue!34!white,draw=blue!34!white,line width=.12pt] (2.251872,1.261099)--(2.262742,1.131371)--(2.030059,1.236817)--(2.020307,1.366038)--cycle;
\path[fill=blue!41!white,draw=blue!41!white,line width=.12pt] (2.369413,0.587353)--(2.292506,0.450808)--(2.097056,0.558630)--(2.167407,0.698793)--cycle;
\path[fill=blue!36!white,draw=blue!36!white,line width=.12pt] (-2.167407,0.698793)--(-2.097056,0.558630)--(-2.292506,0.450808)--(-2.369413,0.587353)--cycle;
\path[fill=blue!34!white,draw=blue!34!white,line width=.12pt] (0.801283,0.327152)--(0.748354,0.327284)--(0.567664,0.359610)--(0.607814,0.361764)--cycle;
\path[fill=blue!36!white,draw=blue!36!white,line width=.12pt] (-0.607814,0.361764)--(-0.567664,0.359610)--(-0.748354,0.327284)--(-0.801283,0.327152)--cycle;
\path[fill=blue!42!white,draw=blue!42!white,line width=.12pt] (2.292506,0.450808)--(2.198794,0.326193)--(2.011334,0.429608)--(2.097056,0.558630)--cycle;
\path[fill=blue!38!white,draw=blue!38!white,line width=.12pt] (-2.097056,0.558630)--(-2.011334,0.429608)--(-2.198794,0.326193)--(-2.292506,0.450808)--cycle;
\path[fill=blue!37!white,draw=blue!37!white,line width=.12pt] (1.246659,0.253371)--(1.163283,0.230405)--(0.987036,0.283224)--(1.057780,0.309976)--cycle;
\path[fill=blue!38!white,draw=blue!38!white,line width=.12pt] (-1.057780,0.309976)--(-0.987036,0.283224)--(-1.163283,0.230405)--(-1.246659,0.253371)--cycle;
\path[fill=blue!40!white,draw=blue!40!white,line width=.12pt] (2.426560,0.730582)--(2.369413,0.587353)--(2.167407,0.698793)--(2.219681,0.844710)--cycle;
\path[fill=blue!34!white,draw=blue!34!white,line width=.12pt] (-2.219681,0.844710)--(-2.167407,0.698793)--(-2.369413,0.587353)--(-2.426560,0.730582)--cycle;
\path[fill=blue!25!white,draw=blue!25!white,line width=.12pt] (1.647647,1.722837)--(1.702920,1.659211)--(1.444914,1.736532)--(1.398014,1.797648)--cycle;
\path[fill=blue!22!white,draw=blue!22!white,line width=.12pt] (-1.398014,1.797648)--(-1.444914,1.736532)--(-1.702920,1.659211)--(-1.647647,1.722837)--cycle;
\path[fill=blue!40!white,draw=blue!40!white,line width=.12pt] (-1.621555,0.308427)--(-1.522544,0.234792)--(-1.697056,0.155708)--(-1.807416,0.224200)--cycle;
\path[fill=blue!40!white,draw=blue!40!white,line width=.12pt] (1.807416,0.224200)--(1.697056,0.155708)--(1.522544,0.234792)--(1.621555,0.308427)--cycle;
\path[fill=blue!18!white,draw=blue!18!white,line width=.12pt] (-0.527943,1.967157)--(-0.554926,1.970959)--(-0.825702,1.937046)--(-0.785553,1.934893)--cycle;
\path[fill=blue!19!white,draw=blue!19!white,line width=.12pt] (0.785553,1.934893)--(0.825702,1.937046)--(0.554926,1.970959)--(0.527943,1.967157)--cycle;
\path[fill=blue!28!white,draw=blue!28!white,line width=.12pt] (1.944527,1.569617)--(1.991426,1.478411)--(1.743993,1.570165)--(1.702920,1.659211)--cycle;
\path[fill=blue!23!white,draw=blue!23!white,line width=.12pt] (-1.702920,1.659211)--(-1.743993,1.570165)--(-1.991426,1.478411)--(-1.944527,1.569617)--cycle;
\path[fill=blue!22!white,draw=blue!22!white,line width=.12pt] (1.340867,1.830352)--(1.398014,1.797648)--(1.134919,1.859866)--(1.088526,1.890026)--cycle;
\path[fill=blue!20!white,draw=blue!20!white,line width=.12pt] (-1.088526,1.890026)--(-1.134919,1.859866)--(-1.398014,1.797648)--(-1.340867,1.830352)--cycle;
\path[fill=blue!39!white,draw=blue!39!white,line width=.12pt] (-2.011334,0.429608)--(-1.913535,0.316685)--(-2.091880,0.218298)--(-2.198794,0.326193)--cycle;
\path[fill=blue!42!white,draw=blue!42!white,line width=.12pt] (2.198794,0.326193)--(2.091880,0.218298)--(1.913535,0.316685)--(2.011334,0.429608)--cycle;
\path[fill=blue!32!white,draw=blue!32!white,line width=.12pt] (-2.251872,0.990774)--(-2.219681,0.844710)--(-2.426560,0.730582)--(-2.461751,0.874991)--cycle;
\path[fill=blue!38!white,draw=blue!38!white,line width=.12pt] (2.461751,0.874991)--(2.426560,0.730582)--(2.219681,0.844710)--(2.251872,0.990774)--cycle;
\path[fill=blue!32!white,draw=blue!32!white,line width=.12pt] (0.191677,0.397005)--(0.179794,0.422843)--(0.000000,0.427259)--(0.000000,0.401713)--cycle;
\path[fill=blue!33!white,draw=blue!33!white,line width=.12pt] (0.000000,0.401713)--(0.000000,0.427259)--(-0.179794,0.422843)--(-0.191677,0.397005)--cycle;
\path[fill=blue!36!white,draw=blue!36!white,line width=.12pt] (-0.801283,0.327152)--(-0.748354,0.327284)--(-0.921837,0.286259)--(-0.987036,0.283224)--cycle;
\path[fill=blue!34!white,draw=blue!34!white,line width=.12pt] (0.987036,0.283224)--(0.921837,0.286259)--(0.748354,0.327284)--(0.801283,0.327152)--cycle;
\path[fill=blue!33!white,draw=blue!33!white,line width=.12pt] (-0.191677,0.397005)--(-0.179794,0.422843)--(-0.357857,0.409636)--(-0.381508,0.382925)--cycle;
\path[fill=blue!32!white,draw=blue!32!white,line width=.12pt] (0.381508,0.382925)--(0.357857,0.409636)--(0.179794,0.422843)--(0.191677,0.397005)--cycle;
\path[fill=blue!25!white,draw=blue!25!white,line width=.12pt] (-1.991426,1.478411)--(-2.020307,1.366038)--(-2.251872,1.261099)--(-2.219681,1.374972)--cycle;
\path[fill=blue!31!white,draw=blue!31!white,line width=.12pt] (2.219681,1.374972)--(2.251872,1.261099)--(2.020307,1.366038)--(1.991426,1.478411)--cycle;
\path[fill=blue!39!white,draw=blue!39!white,line width=.12pt] (-1.913535,0.316685)--(-1.807416,0.224200)--(-1.975871,0.131270)--(-2.091880,0.218298)--cycle;
\path[fill=blue!41!white,draw=blue!41!white,line width=.12pt] (2.091880,0.218298)--(1.975871,0.131270)--(1.807416,0.224200)--(1.913535,0.316685)--cycle;
\path[fill=blue!20!white,draw=blue!20!white,line width=.12pt] (1.035598,1.890159)--(1.088526,1.890026)--(0.825702,1.937046)--(0.785553,1.934893)--cycle;
\path[fill=blue!18!white,draw=blue!18!white,line width=.12pt] (-0.785553,1.934893)--(-0.825702,1.937046)--(-1.088526,1.890026)--(-1.035598,1.890159)--cycle;
\path[fill=blue!32!white,draw=blue!32!white,line width=.12pt] (0.567664,0.359610)--(0.532473,0.387767)--(0.357857,0.409636)--(0.381508,0.382925)--cycle;
\path[fill=blue!33!white,draw=blue!33!white,line width=.12pt] (-0.381508,0.382925)--(-0.357857,0.409636)--(-0.532473,0.387767)--(-0.567664,0.359610)--cycle;
\path[fill=blue!29!white,draw=blue!29!white,line width=.12pt] (-2.262742,1.131371)--(-2.251872,0.990774)--(-2.461751,0.874991)--(-2.473633,1.015029)--cycle;
\path[fill=blue!36!white,draw=blue!36!white,line width=.12pt] (2.473633,1.015029)--(2.461751,0.874991)--(2.251872,0.990774)--(2.262742,1.131371)--cycle;
\path[fill=blue!17!white,draw=blue!17!white,line width=.12pt] (0.000000,1.957544)--(0.000000,1.993156)--(-0.265249,1.986640)--(-0.250539,1.951390)--cycle;
\path[fill=blue!17!white,draw=blue!17!white,line width=.12pt] (0.250539,1.951390)--(0.265249,1.986640)--(0.000000,1.993156)--(0.000000,1.957544)--cycle;
\path[fill=blue!37!white,draw=blue!37!white,line width=.12pt] (1.423533,0.187782)--(1.328326,0.169203)--(1.163283,0.230405)--(1.246659,0.253371)--cycle;
\path[fill=blue!38!white,draw=blue!38!white,line width=.12pt] (-1.246659,0.253371)--(-1.163283,0.230405)--(-1.328326,0.169203)--(-1.423533,0.187782)--cycle;
\path[fill=blue!39!white,draw=blue!39!white,line width=.12pt] (-1.522544,0.234792)--(-1.423533,0.187782)--(-1.586697,0.113840)--(-1.697056,0.155708)--cycle;
\path[fill=blue!39!white,draw=blue!39!white,line width=.12pt] (1.697056,0.155708)--(1.586697,0.113840)--(1.423533,0.187782)--(1.522544,0.234792)--cycle;
\path[fill=blue!17!white,draw=blue!17!white,line width=.12pt] (0.498665,1.932987)--(0.527943,1.967157)--(0.265249,1.986640)--(0.250539,1.951390)--cycle;
\path[fill=blue!17!white,draw=blue!17!white,line width=.12pt] (-0.250539,1.951390)--(-0.265249,1.986640)--(-0.527943,1.967157)--(-0.498665,1.932987)--cycle;
\path[fill=blue!36!white,draw=blue!36!white,line width=.12pt] (-0.987036,0.283224)--(-0.921837,0.286259)--(-1.086442,0.236929)--(-1.163283,0.230405)--cycle;
\path[fill=blue!34!white,draw=blue!34!white,line width=.12pt] (1.163283,0.230405)--(1.086442,0.236929)--(0.921837,0.286259)--(0.987036,0.283224)--cycle;
\path[fill=blue!34!white,draw=blue!34!white,line width=.12pt] (-0.567664,0.359610)--(-0.532473,0.387767)--(-0.701962,0.357445)--(-0.748354,0.327284)--cycle;
\path[fill=blue!32!white,draw=blue!32!white,line width=.12pt] (0.748354,0.327284)--(0.701962,0.357445)--(0.532473,0.387767)--(0.567664,0.359610)--cycle;
\path[fill=blue!19!white,draw=blue!19!white,line width=.12pt] (-1.340867,1.830352)--(-1.398014,1.797648)--(-1.647647,1.722837)--(-1.580295,1.758598)--cycle;
\path[fill=blue!22!white,draw=blue!22!white,line width=.12pt] (1.580295,1.758598)--(1.647647,1.722837)--(1.398014,1.797648)--(1.340867,1.830352)--cycle;
\path[fill=blue!21!white,draw=blue!21!white,line width=.12pt] (-1.647647,1.722837)--(-1.702920,1.659211)--(-1.944527,1.569617)--(-1.881411,1.636151)--cycle;
\path[fill=blue!25!white,draw=blue!25!white,line width=.12pt] (1.881411,1.636151)--(1.944527,1.569617)--(1.702920,1.659211)--(1.647647,1.722837)--cycle;
\path[fill=blue!40!white,draw=blue!40!white,line width=.12pt] (1.975871,0.131270)--(1.855225,0.068452)--(1.697056,0.155708)--(1.807416,0.224200)--cycle;
\path[fill=blue!40!white,draw=blue!40!white,line width=.12pt] (-1.807416,0.224200)--(-1.697056,0.155708)--(-1.855225,0.068452)--(-1.975871,0.131270)--cycle;
\path[fill=blue!17!white,draw=blue!17!white,line width=.12pt] (0.741989,1.902512)--(0.785553,1.934893)--(0.527943,1.967157)--(0.498665,1.932987)--cycle;
\path[fill=blue!17!white,draw=blue!17!white,line width=.12pt] (-0.498665,1.932987)--(-0.527943,1.967157)--(-0.785553,1.934893)--(-0.741989,1.902512)--cycle;
\path[fill=blue!20!white,draw=blue!20!white,line width=.12pt] (1.275669,1.833387)--(1.340867,1.830352)--(1.088526,1.890026)--(1.035598,1.890159)--cycle;
\path[fill=blue!18!white,draw=blue!18!white,line width=.12pt] (-1.035598,1.890159)--(-1.088526,1.890026)--(-1.340867,1.830352)--(-1.275669,1.833387)--cycle;
\path[fill=blue!34!white,draw=blue!34!white,line width=.12pt] (2.461751,1.145316)--(2.473633,1.015029)--(2.262742,1.131371)--(2.251872,1.261099)--cycle;
\path[fill=blue!27!white,draw=blue!27!white,line width=.12pt] (-2.251872,1.261099)--(-2.262742,1.131371)--(-2.473633,1.015029)--(-2.461751,1.145316)--cycle;
\path[fill=blue!41!white,draw=blue!41!white,line width=.12pt] (2.548600,0.466550)--(2.465877,0.333925)--(2.292506,0.450808)--(2.369413,0.587353)--cycle;
\path[fill=blue!35!white,draw=blue!35!white,line width=.12pt] (-2.369413,0.587353)--(-2.292506,0.450808)--(-2.465877,0.333925)--(-2.548600,0.466550)--cycle;
\path[fill=blue!42!white,draw=blue!42!white,line width=.12pt] (2.465877,0.333925)--(2.365079,0.214089)--(2.198794,0.326193)--(2.292506,0.450808)--cycle;
\path[fill=blue!37!white,draw=blue!37!white,line width=.12pt] (-2.292506,0.450808)--(-2.198794,0.326193)--(-2.365079,0.214089)--(-2.465877,0.333925)--cycle;
\path[fill=blue!28!white,draw=blue!28!white,line width=.12pt] (2.167407,1.468614)--(2.219681,1.374972)--(1.991426,1.478411)--(1.944527,1.569617)--cycle;
\path[fill=blue!23!white,draw=blue!23!white,line width=.12pt] (-1.944527,1.569617)--(-1.991426,1.478411)--(-2.219681,1.374972)--(-2.167407,1.468614)--cycle;
\path[fill=blue!30!white,draw=blue!30!white,line width=.12pt] (0.179794,0.422843)--(0.170043,0.478325)--(0.000000,0.482502)--(0.000000,0.427259)--cycle;
\path[fill=blue!30!white,draw=blue!30!white,line width=.12pt] (0.000000,0.427259)--(0.000000,0.482502)--(-0.170043,0.478325)--(-0.179794,0.422843)--cycle;
\path[fill=blue!32!white,draw=blue!32!white,line width=.12pt] (0.921837,0.286259)--(0.864690,0.318963)--(0.701962,0.357445)--(0.748354,0.327284)--cycle;
\path[fill=blue!34!white,draw=blue!34!white,line width=.12pt] (-0.748354,0.327284)--(-0.701962,0.357445)--(-0.864690,0.318963)--(-0.921837,0.286259)--cycle;
\path[fill=blue!40!white,draw=blue!40!white,line width=.12pt] (2.610069,0.606865)--(2.548600,0.466550)--(2.369413,0.587353)--(2.426560,0.730582)--cycle;
\path[fill=blue!33!white,draw=blue!33!white,line width=.12pt] (-2.426560,0.730582)--(-2.369413,0.587353)--(-2.548600,0.466550)--(-2.610069,0.606865)--cycle;
\path[fill=blue!42!white,draw=blue!42!white,line width=.12pt] (2.365079,0.214089)--(2.250079,0.111645)--(2.091880,0.218298)--(2.198794,0.326193)--cycle;
\path[fill=blue!38!white,draw=blue!38!white,line width=.12pt] (-2.198794,0.326193)--(-2.091880,0.218298)--(-2.250079,0.111645)--(-2.365079,0.214089)--cycle;
\path[fill=blue!29!white,draw=blue!29!white,line width=.12pt] (0.357857,0.409636)--(0.338447,0.465835)--(0.170043,0.478325)--(0.179794,0.422843)--cycle;
\path[fill=blue!30!white,draw=blue!30!white,line width=.12pt] (-0.179794,0.422843)--(-0.170043,0.478325)--(-0.338447,0.465835)--(-0.357857,0.409636)--cycle;
\path[fill=blue!38!white,draw=blue!38!white,line width=.12pt] (-1.423533,0.187782)--(-1.328326,0.169203)--(-1.480578,0.100206)--(-1.586697,0.113840)--cycle;
\path[fill=blue!37!white,draw=blue!37!white,line width=.12pt] (1.586697,0.113840)--(1.480578,0.100206)--(1.328326,0.169203)--(1.423533,0.187782)--cycle;
\path[fill=blue!16!white,draw=blue!16!white,line width=.12pt] (-0.741989,1.902512)--(-0.785553,1.934893)--(-1.035598,1.890159)--(-0.978167,1.860259)--cycle;
\path[fill=blue!17!white,draw=blue!17!white,line width=.12pt] (0.978167,1.860259)--(1.035598,1.890159)--(0.785553,1.934893)--(0.741989,1.902512)--cycle;
\path[fill=blue!37!white,draw=blue!37!white,line width=.12pt] (-1.163283,0.230405)--(-1.086442,0.236929)--(-1.240584,0.179769)--(-1.328326,0.169203)--cycle;
\path[fill=blue!34!white,draw=blue!34!white,line width=.12pt] (1.328326,0.169203)--(1.240584,0.179769)--(1.086442,0.236929)--(1.163283,0.230405)--cycle;
\path[fill=blue!29!white,draw=blue!29!white,line width=.12pt] (0.532473,0.387767)--(0.503593,0.445151)--(0.338447,0.465835)--(0.357857,0.409636)--cycle;
\path[fill=blue!31!white,draw=blue!31!white,line width=.12pt] (-0.357857,0.409636)--(-0.338447,0.465835)--(-0.503593,0.445151)--(-0.532473,0.387767)--cycle;
\path[fill=blue!39!white,draw=blue!39!white,line width=.12pt] (1.855225,0.068452)--(1.734580,0.032258)--(1.586697,0.113840)--(1.697056,0.155708)--cycle;
\path[fill=blue!39!white,draw=blue!39!white,line width=.12pt] (-1.697056,0.155708)--(-1.586697,0.113840)--(-1.734580,0.032258)--(-1.855225,0.068452)--cycle;
\path[fill=blue!38!white,draw=blue!38!white,line width=.12pt] (2.647922,0.749480)--(2.610069,0.606865)--(2.426560,0.730582)--(2.461751,0.874991)--cycle;
\path[fill=blue!31!white,draw=blue!31!white,line width=.12pt] (-2.461751,0.874991)--(-2.426560,0.730582)--(-2.610069,0.606865)--(-2.647922,0.749480)--cycle;
\path[fill=blue!15!white,draw=blue!15!white,line width=.12pt] (0.000000,1.892820)--(0.000000,1.957544)--(-0.250539,1.951390)--(-0.235241,1.887042)--cycle;
\path[fill=blue!15!white,draw=blue!15!white,line width=.12pt] (0.235241,1.887042)--(0.250539,1.951390)--(0.000000,1.957544)--(0.000000,1.892820)--cycle;
\path[fill=blue!41!white,draw=blue!41!white,line width=.12pt] (2.250079,0.111645)--(2.125296,0.030531)--(1.975871,0.131270)--(2.091880,0.218298)--cycle;
\path[fill=blue!39!white,draw=blue!39!white,line width=.12pt] (-2.091880,0.218298)--(-1.975871,0.131270)--(-2.125296,0.030531)--(-2.250079,0.111645)--cycle;
\path[fill=blue!15!white,draw=blue!15!white,line width=.12pt] (-0.235241,1.887042)--(-0.250539,1.951390)--(-0.498665,1.932987)--(-0.468217,1.869763)--cycle;
\path[fill=blue!15!white,draw=blue!15!white,line width=.12pt] (0.468217,1.869763)--(0.498665,1.932987)--(0.250539,1.951390)--(0.235241,1.887042)--cycle;
\path[fill=blue!31!white,draw=blue!31!white,line width=.12pt] (2.426560,1.260844)--(2.461751,1.145316)--(2.251872,1.261099)--(2.219681,1.374972)--cycle;
\path[fill=blue!24!white,draw=blue!24!white,line width=.12pt] (-2.219681,1.374972)--(-2.251872,1.261099)--(-2.461751,1.145316)--(-2.426560,1.260844)--cycle;
\path[fill=blue!35!white,draw=blue!35!white,line width=.12pt] (-0.921837,0.286259)--(-0.864690,0.318963)--(-1.019091,0.272690)--(-1.086442,0.236929)--cycle;
\path[fill=blue!32!white,draw=blue!32!white,line width=.12pt] (1.086442,0.236929)--(1.019091,0.272690)--(0.864690,0.318963)--(0.921837,0.286259)--cycle;
\path[fill=blue!22!white,draw=blue!22!white,line width=.12pt] (1.804504,1.675456)--(1.881411,1.636151)--(1.647647,1.722837)--(1.580295,1.758598)--cycle;
\path[fill=blue!19!white,draw=blue!19!white,line width=.12pt] (-1.580295,1.758598)--(-1.647647,1.722837)--(-1.881411,1.636151)--(-1.804504,1.675456)--cycle;
\path[fill=blue!18!white,draw=blue!18!white,line width=.12pt] (-1.275669,1.833387)--(-1.340867,1.830352)--(-1.580295,1.758598)--(-1.503455,1.765122)--cycle;
\path[fill=blue!20!white,draw=blue!20!white,line width=.12pt] (1.503455,1.765122)--(1.580295,1.758598)--(1.340867,1.830352)--(1.275669,1.833387)--cycle;
\path[fill=blue!29!white,draw=blue!29!white,line width=.12pt] (0.701962,0.357445)--(0.663889,0.416474)--(0.503593,0.445151)--(0.532473,0.387767)--cycle;
\path[fill=blue!31!white,draw=blue!31!white,line width=.12pt] (-0.532473,0.387767)--(-0.503593,0.445151)--(-0.663889,0.416474)--(-0.701962,0.357445)--cycle;
\path[fill=blue!36!white,draw=blue!36!white,line width=.12pt] (2.660703,0.888912)--(2.647922,0.749480)--(2.461751,0.874991)--(2.473633,1.015029)--cycle;
\path[fill=blue!29!white,draw=blue!29!white,line width=.12pt] (-2.473633,1.015029)--(-2.461751,0.874991)--(-2.647922,0.749480)--(-2.660703,0.888912)--cycle;
\path[fill=blue!15!white,draw=blue!15!white,line width=.12pt] (0.696683,1.841149)--(0.741989,1.902512)--(0.498665,1.932987)--(0.468217,1.869763)--cycle;
\path[fill=blue!15!white,draw=blue!15!white,line width=.12pt] (-0.468217,1.869763)--(-0.498665,1.932987)--(-0.741989,1.902512)--(-0.696683,1.841149)--cycle;
\path[fill=blue!25!white,draw=blue!25!white,line width=.12pt] (2.097056,1.538426)--(2.167407,1.468614)--(1.944527,1.569617)--(1.881411,1.636151)--cycle;
\path[fill=blue!21!white,draw=blue!21!white,line width=.12pt] (-1.881411,1.636151)--(-1.944527,1.569617)--(-2.167407,1.468614)--(-2.097056,1.538426)--cycle;
\path[fill=blue!17!white,draw=blue!17!white,line width=.12pt] (1.204924,1.806635)--(1.275669,1.833387)--(1.035598,1.890159)--(0.978167,1.860259)--cycle;
\path[fill=blue!16!white,draw=blue!16!white,line width=.12pt] (-0.978167,1.860259)--(-1.035598,1.890159)--(-1.275669,1.833387)--(-1.204924,1.806635)--cycle;
\path[fill=blue!40!white,draw=blue!40!white,line width=.12pt] (2.125296,0.030531)--(1.995527,-0.026136)--(1.855225,0.068452)--(1.975871,0.131270)--cycle;
\path[fill=blue!40!white,draw=blue!40!white,line width=.12pt] (-1.975871,0.131270)--(-1.855225,0.068452)--(-1.995527,-0.026136)--(-2.125296,0.030531)--cycle;
\path[fill=blue!27!white,draw=blue!27!white,line width=.12pt] (0.000000,0.482502)--(0.000000,0.565317)--(-0.162796,0.561318)--(-0.170043,0.478325)--cycle;
\path[fill=blue!27!white,draw=blue!27!white,line width=.12pt] (0.170043,0.478325)--(0.162796,0.561318)--(0.000000,0.565317)--(0.000000,0.482502)--cycle;
\path[fill=blue!37!white,draw=blue!37!white,line width=.12pt] (-1.328326,0.169203)--(-1.240584,0.179769)--(-1.382778,0.115330)--(-1.480578,0.100206)--cycle;
\path[fill=blue!34!white,draw=blue!34!white,line width=.12pt] (1.480578,0.100206)--(1.382778,0.115330)--(1.240584,0.179769)--(1.328326,0.169203)--cycle;
\path[fill=blue!29!white,draw=blue!29!white,line width=.12pt] (0.864690,0.318963)--(0.817791,0.380079)--(0.663889,0.416474)--(0.701962,0.357445)--cycle;
\path[fill=blue!32!white,draw=blue!32!white,line width=.12pt] (-0.701962,0.357445)--(-0.663889,0.416474)--(-0.817791,0.380079)--(-0.864690,0.318963)--cycle;
\path[fill=blue!39!white,draw=blue!39!white,line width=.12pt] (-1.586697,0.113840)--(-1.480578,0.100206)--(-1.618570,0.024081)--(-1.734580,0.032258)--cycle;
\path[fill=blue!37!white,draw=blue!37!white,line width=.12pt] (1.734580,0.032258)--(1.618570,0.024081)--(1.480578,0.100206)--(1.586697,0.113840)--cycle;
\path[fill=blue!26!white,draw=blue!26!white,line width=.12pt] (0.338447,0.465835)--(0.324025,0.549360)--(0.162796,0.561318)--(0.170043,0.478325)--cycle;
\path[fill=blue!28!white,draw=blue!28!white,line width=.12pt] (-0.170043,0.478325)--(-0.162796,0.561318)--(-0.324025,0.549360)--(-0.338447,0.465835)--cycle;
\path[fill=blue!32!white,draw=blue!32!white,line width=.12pt] (1.240584,0.179769)--(1.163677,0.219074)--(1.019091,0.272690)--(1.086442,0.236929)--cycle;
\path[fill=blue!35!white,draw=blue!35!white,line width=.12pt] (-1.086442,0.236929)--(-1.019091,0.272690)--(-1.163677,0.219074)--(-1.240584,0.179769)--cycle;
\path[fill=blue!15!white,draw=blue!15!white,line width=.12pt] (0.918440,1.801476)--(0.978167,1.860259)--(0.741989,1.902512)--(0.696683,1.841149)--cycle;
\path[fill=blue!15!white,draw=blue!15!white,line width=.12pt] (-0.696683,1.841149)--(-0.741989,1.902512)--(-0.978167,1.860259)--(-0.918440,1.801476)--cycle;
\path[fill=blue!28!white,draw=blue!28!white,line width=.12pt] (2.369413,1.357174)--(2.426560,1.260844)--(2.219681,1.374972)--(2.167407,1.468614)--cycle;
\path[fill=blue!22!white,draw=blue!22!white,line width=.12pt] (-2.167407,1.468614)--(-2.219681,1.374972)--(-2.426560,1.260844)--(-2.369413,1.357174)--cycle;
\path[fill=blue!28!white,draw=blue!28!white,line width=.12pt] (-0.338447,0.465835)--(-0.324025,0.549360)--(-0.482133,0.529559)--(-0.503593,0.445151)--cycle;
\path[fill=blue!26!white,draw=blue!26!white,line width=.12pt] (0.503593,0.445151)--(0.482133,0.529559)--(0.324025,0.549360)--(0.338447,0.465835)--cycle;
\path[fill=blue!14!white,draw=blue!14!white,line width=.12pt] (0.219943,1.796069)--(0.235241,1.887042)--(0.000000,1.892820)--(0.000000,1.801472)--cycle;
\path[fill=blue!14!white,draw=blue!14!white,line width=.12pt] (0.000000,1.801472)--(0.000000,1.892820)--(-0.235241,1.887042)--(-0.219943,1.796069)--cycle;
\path[fill=blue!33!white,draw=blue!33!white,line width=.12pt] (2.647922,1.019805)--(2.660703,0.888912)--(2.473633,1.015029)--(2.461751,1.145316)--cycle;
\path[fill=blue!26!white,draw=blue!26!white,line width=.12pt] (-2.461751,1.145316)--(-2.473633,1.015029)--(-2.660703,0.888912)--(-2.647922,1.019805)--cycle;
\path[fill=blue!37!white,draw=blue!37!white,line width=.12pt] (-2.465877,0.333925)--(-2.365079,0.214089)--(-2.508586,0.094375)--(-2.615501,0.209109)--cycle;
\path[fill=blue!41!white,draw=blue!41!white,line width=.12pt] (2.615501,0.209109)--(2.508586,0.094375)--(2.365079,0.214089)--(2.465877,0.333925)--cycle;
\path[fill=blue!35!white,draw=blue!35!white,line width=.12pt] (-2.548600,0.466550)--(-2.465877,0.333925)--(-2.615501,0.209109)--(-2.703244,0.337547)--cycle;
\path[fill=blue!41!white,draw=blue!41!white,line width=.12pt] (2.703244,0.337547)--(2.615501,0.209109)--(2.465877,0.333925)--(2.548600,0.466550)--cycle;
\path[fill=blue!39!white,draw=blue!39!white,line width=.12pt] (1.995527,-0.026136)--(1.865758,-0.056178)--(1.734580,0.032258)--(1.855225,0.068452)--cycle;
\path[fill=blue!40!white,draw=blue!40!white,line width=.12pt] (-1.855225,0.068452)--(-1.734580,0.032258)--(-1.865758,-0.056178)--(-1.995527,-0.026136)--cycle;
\path[fill=blue!14!white,draw=blue!14!white,line width=.12pt] (0.437769,1.779914)--(0.468217,1.869763)--(0.235241,1.887042)--(0.219943,1.796069)--cycle;
\path[fill=blue!14!white,draw=blue!14!white,line width=.12pt] (-0.219943,1.796069)--(-0.235241,1.887042)--(-0.468217,1.869763)--(-0.437769,1.779914)--cycle;
\path[fill=blue!38!white,draw=blue!38!white,line width=.12pt] (-2.365079,0.214089)--(-2.250079,0.111645)--(-2.386608,-0.002248)--(-2.508586,0.094375)--cycle;
\path[fill=blue!42!white,draw=blue!42!white,line width=.12pt] (2.508586,0.094375)--(2.386608,-0.002248)--(2.250079,0.111645)--(2.365079,0.214089)--cycle;
\path[fill=blue!20!white,draw=blue!20!white,line width=.12pt] (1.716761,1.686022)--(1.804504,1.675456)--(1.580295,1.758598)--(1.503455,1.765122)--cycle;
\path[fill=blue!17!white,draw=blue!17!white,line width=.12pt] (-1.503455,1.765122)--(-1.580295,1.758598)--(-1.804504,1.675456)--(-1.716761,1.686022)--cycle;
\path[fill=blue!29!white,draw=blue!29!white,line width=.12pt] (1.019091,0.272690)--(0.963817,0.336316)--(0.817791,0.380079)--(0.864690,0.318963)--cycle;
\path[fill=blue!32!white,draw=blue!32!white,line width=.12pt] (-0.864690,0.318963)--(-0.817791,0.380079)--(-0.963817,0.336316)--(-1.019091,0.272690)--cycle;
\path[fill=blue!16!white,draw=blue!16!white,line width=.12pt] (-1.204924,1.806635)--(-1.275669,1.833387)--(-1.503455,1.765122)--(-1.420078,1.742156)--cycle;
\path[fill=blue!17!white,draw=blue!17!white,line width=.12pt] (1.420078,1.742156)--(1.503455,1.765122)--(1.275669,1.833387)--(1.204924,1.806635)--cycle;
\path[fill=blue!29!white,draw=blue!29!white,line width=.12pt] (-0.503593,0.445151)--(-0.482133,0.529559)--(-0.635598,0.502103)--(-0.663889,0.416474)--cycle;
\path[fill=blue!26!white,draw=blue!26!white,line width=.12pt] (0.663889,0.416474)--(0.635598,0.502103)--(0.482133,0.529559)--(0.503593,0.445151)--cycle;
\path[fill=blue!33!white,draw=blue!33!white,line width=.12pt] (-2.610069,0.606865)--(-2.548600,0.466550)--(-2.703244,0.337547)--(-2.768442,0.474751)--cycle;
\path[fill=blue!39!white,draw=blue!39!white,line width=.12pt] (2.768442,0.474751)--(2.703244,0.337547)--(2.548600,0.466550)--(2.610069,0.606865)--cycle;
\path[fill=blue!19!white,draw=blue!19!white,line width=.12pt] (-1.804504,1.675456)--(-1.881411,1.636151)--(-2.097056,1.538426)--(-2.011334,1.581726)--cycle;
\path[fill=blue!22!white,draw=blue!22!white,line width=.12pt] (2.011334,1.581726)--(2.097056,1.538426)--(1.881411,1.636151)--(1.804504,1.675456)--cycle;
\path[fill=blue!14!white,draw=blue!14!white,line width=.12pt] (0.651378,1.753160)--(0.696683,1.841149)--(0.468217,1.869763)--(0.437769,1.779914)--cycle;
\path[fill=blue!14!white,draw=blue!14!white,line width=.12pt] (-0.437769,1.779914)--(-0.468217,1.869763)--(-0.696683,1.841149)--(-0.651378,1.753160)--cycle;
\path[fill=blue!39!white,draw=blue!39!white,line width=.12pt] (-2.250079,0.111645)--(-2.125296,0.030531)--(-2.254254,-0.077046)--(-2.386608,-0.002248)--cycle;
\path[fill=blue!41!white,draw=blue!41!white,line width=.12pt] (2.386608,-0.002248)--(2.254254,-0.077046)--(2.125296,0.030531)--(2.250079,0.111645)--cycle;
\path[fill=blue!15!white,draw=blue!15!white,line width=.12pt] (1.131352,1.751126)--(1.204924,1.806635)--(0.978167,1.860259)--(0.918440,1.801476)--cycle;
\path[fill=blue!15!white,draw=blue!15!white,line width=.12pt] (-0.918440,1.801476)--(-0.978167,1.860259)--(-1.204924,1.806635)--(-1.131352,1.751126)--cycle;
\path[fill=blue!24!white,draw=blue!24!white,line width=.12pt] (0.162796,0.561318)--(0.158334,0.668634)--(0.000000,0.672523)--(0.000000,0.565317)--cycle;
\path[fill=blue!24!white,draw=blue!24!white,line width=.12pt] (0.000000,0.565317)--(0.000000,0.672523)--(-0.158334,0.668634)--(-0.162796,0.561318)--cycle;
\path[fill=blue!34!white,draw=blue!34!white,line width=.12pt] (1.618570,0.024081)--(1.511656,0.044233)--(1.382778,0.115330)--(1.480578,0.100206)--cycle;
\path[fill=blue!38!white,draw=blue!38!white,line width=.12pt] (-1.480578,0.100206)--(-1.382778,0.115330)--(-1.511656,0.044233)--(-1.618570,0.024081)--cycle;
\path[fill=blue!35!white,draw=blue!35!white,line width=.12pt] (-1.240584,0.179769)--(-1.163677,0.219074)--(-1.297056,0.158630)--(-1.382778,0.115330)--cycle;
\path[fill=blue!32!white,draw=blue!32!white,line width=.12pt] (1.382778,0.115330)--(1.297056,0.158630)--(1.163677,0.219074)--(1.240584,0.179769)--cycle;
\path[fill=blue!30!white,draw=blue!30!white,line width=.12pt] (-2.647922,0.749480)--(-2.610069,0.606865)--(-2.768442,0.474751)--(-2.808591,0.615449)--cycle;
\path[fill=blue!38!white,draw=blue!38!white,line width=.12pt] (2.808591,0.615449)--(2.768442,0.474751)--(2.610069,0.606865)--(2.647922,0.749480)--cycle;
\path[fill=blue!24!white,draw=blue!24!white,line width=.12pt] (0.324025,0.549360)--(0.315143,0.657004)--(0.158334,0.668634)--(0.162796,0.561318)--cycle;
\path[fill=blue!25!white,draw=blue!25!white,line width=.12pt] (-0.162796,0.561318)--(-0.158334,0.668634)--(-0.315143,0.657004)--(-0.324025,0.549360)--cycle;
\path[fill=blue!26!white,draw=blue!26!white,line width=.12pt] (0.817791,0.380079)--(0.782941,0.467259)--(0.635598,0.502103)--(0.663889,0.416474)--cycle;
\path[fill=blue!29!white,draw=blue!29!white,line width=.12pt] (-0.663889,0.416474)--(-0.635598,0.502103)--(-0.782941,0.467259)--(-0.817791,0.380079)--cycle;
\path[fill=blue!31!white,draw=blue!31!white,line width=.12pt] (2.610069,1.137127)--(2.647922,1.019805)--(2.461751,1.145316)--(2.426560,1.260844)--cycle;
\path[fill=blue!24!white,draw=blue!24!white,line width=.12pt] (-2.426560,1.260844)--(-2.461751,1.145316)--(-2.647922,1.019805)--(-2.610069,1.137127)--cycle;
\path[fill=blue!25!white,draw=blue!25!white,line width=.12pt] (2.292506,1.430603)--(2.369413,1.357174)--(2.167407,1.468614)--(2.097056,1.538426)--cycle;
\path[fill=blue!20!white,draw=blue!20!white,line width=.12pt] (-2.097056,1.538426)--(-2.167407,1.468614)--(-2.369413,1.357174)--(-2.292506,1.430603)--cycle;
\path[fill=blue!37!white,draw=blue!37!white,line width=.12pt] (1.865758,-0.056178)--(1.740975,-0.058441)--(1.618570,0.024081)--(1.734580,0.032258)--cycle;
\path[fill=blue!39!white,draw=blue!39!white,line width=.12pt] (-1.734580,0.032258)--(-1.618570,0.024081)--(-1.740975,-0.058441)--(-1.865758,-0.056178)--cycle;
\path[fill=blue!40!white,draw=blue!40!white,line width=.12pt] (2.254254,-0.077046)--(2.116611,-0.127144)--(1.995527,-0.026136)--(2.125296,0.030531)--cycle;
\path[fill=blue!40!white,draw=blue!40!white,line width=.12pt] (-2.125296,0.030531)--(-1.995527,-0.026136)--(-2.116611,-0.127144)--(-2.254254,-0.077046)--cycle;
\path[fill=blue!14!white,draw=blue!14!white,line width=.12pt] (0.858714,1.716067)--(0.918440,1.801476)--(0.696683,1.841149)--(0.651378,1.753160)--cycle;
\path[fill=blue!14!white,draw=blue!14!white,line width=.12pt] (-0.651378,1.753160)--(-0.696683,1.841149)--(-0.918440,1.801476)--(-0.858714,1.716067)--cycle;
\path[fill=blue!13!white,draw=blue!13!white,line width=.12pt] (0.000000,1.687009)--(0.000000,1.801472)--(-0.219943,1.796069)--(-0.205234,1.681968)--cycle;
\path[fill=blue!13!white,draw=blue!13!white,line width=.12pt] (0.205234,1.681968)--(0.219943,1.796069)--(0.000000,1.801472)--(0.000000,1.687009)--cycle;
\path[fill=blue!33!white,draw=blue!33!white,line width=.12pt] (-1.019091,0.272690)--(-0.963817,0.336316)--(-1.100561,0.285608)--(-1.163677,0.219074)--cycle;
\path[fill=blue!29!white,draw=blue!29!white,line width=.12pt] (1.163677,0.219074)--(1.100561,0.285608)--(0.963817,0.336316)--(1.019091,0.272690)--cycle;
\path[fill=blue!25!white,draw=blue!25!white,line width=.12pt] (-0.324025,0.549360)--(-0.315143,0.657004)--(-0.468918,0.637745)--(-0.482133,0.529559)--cycle;
\path[fill=blue!23!white,draw=blue!23!white,line width=.12pt] (0.482133,0.529559)--(0.468918,0.637745)--(0.315143,0.657004)--(0.324025,0.549360)--cycle;
\path[fill=blue!13!white,draw=blue!13!white,line width=.12pt] (0.408491,1.666893)--(0.437769,1.779914)--(0.219943,1.796069)--(0.205234,1.681968)--cycle;
\path[fill=blue!13!white,draw=blue!13!white,line width=.12pt] (-0.205234,1.681968)--(-0.219943,1.796069)--(-0.437769,1.779914)--(-0.408491,1.666893)--cycle;
\path[fill=blue!28!white,draw=blue!28!white,line width=.12pt] (-2.660703,0.888912)--(-2.647922,0.749480)--(-2.808591,0.615449)--(-2.822148,0.754235)--cycle;
\path[fill=blue!36!white,draw=blue!36!white,line width=.12pt] (2.822148,0.754235)--(2.808591,0.615449)--(2.647922,0.749480)--(2.660703,0.888912)--cycle;
\path[fill=blue!23!white,draw=blue!23!white,line width=.12pt] (0.635598,0.502103)--(0.618176,0.611042)--(0.468918,0.637745)--(0.482133,0.529559)--cycle;
\path[fill=blue!26!white,draw=blue!26!white,line width=.12pt] (-0.482133,0.529559)--(-0.468918,0.637745)--(-0.618176,0.611042)--(-0.635598,0.502103)--cycle;
\path[fill=blue!30!white,draw=blue!30!white,line width=.12pt] (-0.817791,0.380079)--(-0.782941,0.467259)--(-0.922745,0.425362)--(-0.963817,0.336316)--cycle;
\path[fill=blue!26!white,draw=blue!26!white,line width=.12pt] (0.963817,0.336316)--(0.922745,0.425362)--(0.782941,0.467259)--(0.817791,0.380079)--cycle;
\path[fill=blue!16!white,draw=blue!16!white,line width=.12pt] (-1.420078,1.742156)--(-1.503455,1.765122)--(-1.716761,1.686022)--(-1.621555,1.667443)--cycle;
\path[fill=blue!17!white,draw=blue!17!white,line width=.12pt] (1.621555,1.667443)--(1.716761,1.686022)--(1.503455,1.765122)--(1.420078,1.742156)--cycle;
\path[fill=blue!15!white,draw=blue!15!white,line width=.12pt] (-1.131352,1.751126)--(-1.204924,1.806635)--(-1.420078,1.742156)--(-1.333369,1.690584)--cycle;
\path[fill=blue!15!white,draw=blue!15!white,line width=.12pt] (1.333369,1.690584)--(1.420078,1.742156)--(1.204924,1.806635)--(1.131352,1.751126)--cycle;
\path[fill=blue!13!white,draw=blue!13!white,line width=.12pt] (0.607814,1.641929)--(0.651378,1.753160)--(0.437769,1.779914)--(0.408491,1.666893)--cycle;
\path[fill=blue!13!white,draw=blue!13!white,line width=.12pt] (-0.408491,1.666893)--(-0.437769,1.779914)--(-0.651378,1.753160)--(-0.607814,1.641929)--cycle;
\path[fill=blue!40!white,draw=blue!40!white,line width=.12pt] (-1.995527,-0.026136)--(-1.865758,-0.056178)--(-1.978968,-0.150618)--(-2.116611,-0.127144)--cycle;
\path[fill=blue!39!white,draw=blue!39!white,line width=.12pt] (2.116611,-0.127144)--(1.978968,-0.150618)--(1.865758,-0.056178)--(1.995527,-0.026136)--cycle;
\path[fill=blue!22!white,draw=blue!22!white,line width=.12pt] (0.000000,0.672523)--(0.000000,0.800000)--(-0.156827,0.796148)--(-0.158334,0.668634)--cycle;
\path[fill=blue!21!white,draw=blue!21!white,line width=.12pt] (0.158334,0.668634)--(0.156827,0.796148)--(0.000000,0.800000)--(0.000000,0.672523)--cycle;
\path[fill=blue!17!white,draw=blue!17!white,line width=.12pt] (-1.716761,1.686022)--(-1.804504,1.675456)--(-2.011334,1.581726)--(-1.913535,1.596850)--cycle;
\path[fill=blue!20!white,draw=blue!20!white,line width=.12pt] (1.913535,1.596850)--(2.011334,1.581726)--(1.804504,1.675456)--(1.716761,1.686022)--cycle;
\path[fill=blue!14!white,draw=blue!14!white,line width=.12pt] (1.057780,1.668992)--(1.131352,1.751126)--(0.918440,1.801476)--(0.858714,1.716067)--cycle;
\path[fill=blue!14!white,draw=blue!14!white,line width=.12pt] (-0.858714,1.716067)--(-0.918440,1.801476)--(-1.131352,1.751126)--(-1.057780,1.668992)--cycle;
\path[fill=blue!21!white,draw=blue!21!white,line width=.12pt] (0.315143,0.657004)--(0.312145,0.784628)--(0.156827,0.796148)--(0.158334,0.668634)--cycle;
\path[fill=blue!22!white,draw=blue!22!white,line width=.12pt] (-0.158334,0.668634)--(-0.156827,0.796148)--(-0.312145,0.784628)--(-0.315143,0.657004)--cycle;
\path[fill=blue!28!white,draw=blue!28!white,line width=.12pt] (2.548600,1.236371)--(2.610069,1.137127)--(2.426560,1.260844)--(2.369413,1.357174)--cycle;
\path[fill=blue!21!white,draw=blue!21!white,line width=.12pt] (-2.369413,1.357174)--(-2.426560,1.260844)--(-2.610069,1.137127)--(-2.548600,1.236371)--cycle;
\path[fill=blue!32!white,draw=blue!32!white,line width=.12pt] (1.511656,0.044233)--(1.417944,0.091940)--(1.297056,0.158630)--(1.382778,0.115330)--cycle;
\path[fill=blue!36!white,draw=blue!36!white,line width=.12pt] (-1.382778,0.115330)--(-1.297056,0.158630)--(-1.417944,0.091940)--(-1.511656,0.044233)--cycle;
\path[fill=blue!13!white,draw=blue!13!white,line width=.12pt] (0.191677,1.549123)--(0.205234,1.681968)--(0.000000,1.687009)--(0.000000,1.553831)--cycle;
\path[fill=blue!13!white,draw=blue!13!white,line width=.12pt] (0.000000,1.553831)--(0.000000,1.687009)--(-0.205234,1.681968)--(-0.191677,1.549123)--cycle;
\path[fill=blue!23!white,draw=blue!23!white,line width=.12pt] (0.782941,0.467259)--(0.761481,0.577153)--(0.618176,0.611042)--(0.635598,0.502103)--cycle;
\path[fill=blue!27!white,draw=blue!27!white,line width=.12pt] (-0.635598,0.502103)--(-0.618176,0.611042)--(-0.761481,0.577153)--(-0.782941,0.467259)--cycle;
\path[fill=blue!29!white,draw=blue!29!white,line width=.12pt] (1.297056,0.158630)--(1.226706,0.228443)--(1.100561,0.285608)--(1.163677,0.219074)--cycle;
\path[fill=blue!33!white,draw=blue!33!white,line width=.12pt] (-1.163677,0.219074)--(-1.100561,0.285608)--(-1.226706,0.228443)--(-1.297056,0.158630)--cycle;
\path[fill=blue!35!white,draw=blue!35!white,line width=.12pt] (1.740975,-0.058441)--(1.625975,-0.032838)--(1.511656,0.044233)--(1.618570,0.024081)--cycle;
\path[fill=blue!38!white,draw=blue!38!white,line width=.12pt] (-1.618570,0.024081)--(-1.511656,0.044233)--(-1.625975,-0.032838)--(-1.740975,-0.058441)--cycle;
\path[fill=blue!23!white,draw=blue!23!white,line width=.12pt] (-0.315143,0.657004)--(-0.312145,0.784628)--(-0.464455,0.765552)--(-0.468918,0.637745)--cycle;
\path[fill=blue!21!white,draw=blue!21!white,line width=.12pt] (0.468918,0.637745)--(0.464455,0.765552)--(0.312145,0.784628)--(0.315143,0.657004)--cycle;
\path[fill=blue!13!white,draw=blue!13!white,line width=.12pt] (0.801283,1.607317)--(0.858714,1.716067)--(0.651378,1.753160)--(0.607814,1.641929)--cycle;
\path[fill=blue!14!white,draw=blue!14!white,line width=.12pt] (-0.607814,1.641929)--(-0.651378,1.753160)--(-0.858714,1.716067)--(-0.801283,1.607317)--cycle;
\path[fill=blue!36!white,draw=blue!36!white,line width=.12pt] (-2.615501,0.209109)--(-2.508586,0.094375)--(-2.627935,-0.031796)--(-2.739936,0.077561)--cycle;
\path[fill=blue!41!white,draw=blue!41!white,line width=.12pt] (2.739936,0.077561)--(2.627935,-0.031796)--(2.508586,0.094375)--(2.615501,0.209109)--cycle;
\path[fill=blue!12!white,draw=blue!12!white,line width=.12pt] (0.381508,1.535043)--(0.408491,1.666893)--(0.205234,1.681968)--(0.191677,1.549123)--cycle;
\path[fill=blue!13!white,draw=blue!13!white,line width=.12pt] (-0.191677,1.549123)--(-0.205234,1.681968)--(-0.408491,1.666893)--(-0.381508,1.535043)--cycle;
\path[fill=blue!33!white,draw=blue!33!white,line width=.12pt] (2.808591,0.885774)--(2.822148,0.754235)--(2.660703,0.888912)--(2.647922,1.019805)--cycle;
\path[fill=blue!25!white,draw=blue!25!white,line width=.12pt] (-2.647922,1.019805)--(-2.660703,0.888912)--(-2.822148,0.754235)--(-2.808591,0.885774)--cycle;
\path[fill=blue!18!white,draw=blue!18!white,line width=.12pt] (-2.011334,1.581726)--(-2.097056,1.538426)--(-2.292506,1.430603)--(-2.198794,1.478311)--cycle;
\path[fill=blue!22!white,draw=blue!22!white,line width=.12pt] (2.198794,1.478311)--(2.292506,1.430603)--(2.097056,1.538426)--(2.011334,1.581726)--cycle;
\path[fill=blue!41!white,draw=blue!41!white,line width=.12pt] (2.627935,-0.031796)--(2.500154,-0.122284)--(2.386608,-0.002248)--(2.508586,0.094375)--cycle;
\path[fill=blue!38!white,draw=blue!38!white,line width=.12pt] (-2.508586,0.094375)--(-2.386608,-0.002248)--(-2.500154,-0.122284)--(-2.627935,-0.031796)--cycle;
\path[fill=blue!34!white,draw=blue!34!white,line width=.12pt] (-2.703244,0.337547)--(-2.615501,0.209109)--(-2.739936,0.077561)--(-2.831853,0.201586)--cycle;
\path[fill=blue!40!white,draw=blue!40!white,line width=.12pt] (2.831853,0.201586)--(2.739936,0.077561)--(2.615501,0.209109)--(2.703244,0.337547)--cycle;
\path[fill=blue!31!white,draw=blue!31!white,line width=.12pt] (-0.963817,0.336316)--(-0.922745,0.425362)--(-1.053662,0.376814)--(-1.100561,0.285608)--cycle;
\path[fill=blue!26!white,draw=blue!26!white,line width=.12pt] (1.100561,0.285608)--(1.053662,0.376814)--(0.922745,0.425362)--(0.963817,0.336316)--cycle;
\path[fill=blue!39!white,draw=blue!39!white,line width=.12pt] (-2.386608,-0.002248)--(-2.254254,-0.077046)--(-2.361503,-0.190425)--(-2.500154,-0.122284)--cycle;
\path[fill=blue!41!white,draw=blue!41!white,line width=.12pt] (2.500154,-0.122284)--(2.361503,-0.190425)--(2.254254,-0.077046)--(2.386608,-0.002248)--cycle;
\path[fill=blue!19!white,draw=blue!19!white,line width=.12pt] (0.156827,0.796148)--(0.158334,0.938959)--(0.000000,0.942848)--(0.000000,0.800000)--cycle;
\path[fill=blue!19!white,draw=blue!19!white,line width=.12pt] (0.000000,0.800000)--(0.000000,0.942848)--(-0.158334,0.938959)--(-0.156827,0.796148)--cycle;
\path[fill=blue!20!white,draw=blue!20!white,line width=.12pt] (0.618176,0.611042)--(0.612293,0.739104)--(0.464455,0.765552)--(0.468918,0.637745)--cycle;
\path[fill=blue!23!white,draw=blue!23!white,line width=.12pt] (-0.468918,0.637745)--(-0.464455,0.765552)--(-0.612293,0.739104)--(-0.618176,0.611042)--cycle;
\path[fill=blue!13!white,draw=blue!13!white,line width=.12pt] (-0.381508,1.535043)--(-0.408491,1.666893)--(-0.607814,1.641929)--(-0.567664,1.511728)--cycle;
\path[fill=blue!12!white,draw=blue!12!white,line width=.12pt] (0.567664,1.511728)--(0.607814,1.641929)--(0.408491,1.666893)--(0.381508,1.535043)--cycle;
\path[fill=blue!37!white,draw=blue!37!white,line width=.12pt] (1.978968,-0.150618)--(1.846614,-0.146565)--(1.740975,-0.058441)--(1.865758,-0.056178)--cycle;
\path[fill=blue!39!white,draw=blue!39!white,line width=.12pt] (-1.865758,-0.056178)--(-1.740975,-0.058441)--(-1.846614,-0.146565)--(-1.978968,-0.150618)--cycle;
\path[fill=blue!18!white,draw=blue!18!white,line width=.12pt] (0.312145,0.784628)--(0.315143,0.927329)--(0.158334,0.938959)--(0.156827,0.796148)--cycle;
\path[fill=blue!20!white,draw=blue!20!white,line width=.12pt] (-0.156827,0.796148)--(-0.158334,0.938959)--(-0.315143,0.927329)--(-0.312145,0.784628)--cycle;
\path[fill=blue!32!white,draw=blue!32!white,line width=.12pt] (-2.768442,0.474751)--(-2.703244,0.337547)--(-2.831853,0.201586)--(-2.900154,0.335511)--cycle;
\path[fill=blue!39!white,draw=blue!39!white,line width=.12pt] (2.900154,0.335511)--(2.831853,0.201586)--(2.703244,0.337547)--(2.768442,0.474751)--cycle;
\path[fill=blue!27!white,draw=blue!27!white,line width=.12pt] (-0.782941,0.467259)--(-0.761481,0.577153)--(-0.897452,0.536404)--(-0.922745,0.425362)--cycle;
\path[fill=blue!23!white,draw=blue!23!white,line width=.12pt] (0.922745,0.425362)--(0.897452,0.536404)--(0.761481,0.577153)--(0.782941,0.467259)--cycle;
\path[fill=blue!13!white,draw=blue!13!white,line width=.12pt] (0.179794,1.402639)--(0.191677,1.549123)--(0.000000,1.553831)--(0.000000,1.407055)--cycle;
\path[fill=blue!13!white,draw=blue!13!white,line width=.12pt] (0.000000,1.407055)--(0.000000,1.553831)--(-0.191677,1.549123)--(-0.179794,1.402639)--cycle;
\path[fill=blue!14!white,draw=blue!14!white,line width=.12pt] (-1.057780,1.668992)--(-1.131352,1.751126)--(-1.333369,1.690584)--(-1.246659,1.612387)--cycle;
\path[fill=blue!14!white,draw=blue!14!white,line width=.12pt] (1.246659,1.612387)--(1.333369,1.690584)--(1.131352,1.751126)--(1.057780,1.668992)--cycle;
\path[fill=blue!15!white,draw=blue!15!white,line width=.12pt] (1.522544,1.620433)--(1.621555,1.667443)--(1.420078,1.742156)--(1.333369,1.690584)--cycle;
\path[fill=blue!15!white,draw=blue!15!white,line width=.12pt] (-1.333369,1.690584)--(-1.420078,1.742156)--(-1.621555,1.667443)--(-1.522544,1.620433)--cycle;
\path[fill=blue!14!white,draw=blue!14!white,line width=.12pt] (-0.801283,1.607317)--(-0.858714,1.716067)--(-1.057780,1.668992)--(-0.987036,1.563389)--cycle;
\path[fill=blue!13!white,draw=blue!13!white,line width=.12pt] (0.987036,1.563389)--(1.057780,1.668992)--(0.858714,1.716067)--(0.801283,1.607317)--cycle;
\path[fill=blue!13!white,draw=blue!13!white,line width=.12pt] (0.357857,1.389432)--(0.381508,1.535043)--(0.191677,1.549123)--(0.179794,1.402639)--cycle;
\path[fill=blue!13!white,draw=blue!13!white,line width=.12pt] (-0.179794,1.402639)--(-0.191677,1.549123)--(-0.381508,1.535043)--(-0.357857,1.389432)--cycle;
\path[fill=blue!40!white,draw=blue!40!white,line width=.12pt] (2.361503,-0.190425)--(2.217311,-0.233600)--(2.116611,-0.127144)--(2.254254,-0.077046)--cycle;
\path[fill=blue!39!white,draw=blue!39!white,line width=.12pt] (-2.254254,-0.077046)--(-2.116611,-0.127144)--(-2.217311,-0.233600)--(-2.361503,-0.190425)--cycle;
\path[fill=blue!20!white,draw=blue!20!white,line width=.12pt] (-0.312145,0.784628)--(-0.315143,0.927329)--(-0.468918,0.908070)--(-0.464455,0.765552)--cycle;
\path[fill=blue!18!white,draw=blue!18!white,line width=.12pt] (0.464455,0.765552)--(0.468918,0.908070)--(0.315143,0.927329)--(0.312145,0.784628)--cycle;
\path[fill=blue!16!white,draw=blue!16!white,line width=.12pt] (-1.621555,1.667443)--(-1.716761,1.686022)--(-1.913535,1.596850)--(-1.807416,1.583216)--cycle;
\path[fill=blue!17!white,draw=blue!17!white,line width=.12pt] (1.807416,1.583216)--(1.913535,1.596850)--(1.716761,1.686022)--(1.621555,1.667443)--cycle;
\path[fill=blue!25!white,draw=blue!25!white,line width=.12pt] (2.465877,1.313721)--(2.548600,1.236371)--(2.369413,1.357174)--(2.292506,1.430603)--cycle;
\path[fill=blue!19!white,draw=blue!19!white,line width=.12pt] (-2.292506,1.430603)--(-2.369413,1.357174)--(-2.548600,1.236371)--(-2.465877,1.313721)--cycle;
\path[fill=blue!17!white,draw=blue!17!white,line width=.12pt] (0.158334,0.938959)--(0.162796,1.091580)--(0.000000,1.095579)--(0.000000,0.942848)--cycle;
\path[fill=blue!17!white,draw=blue!17!white,line width=.12pt] (0.000000,0.942848)--(0.000000,1.095579)--(-0.162796,1.091580)--(-0.158334,0.938959)--cycle;
\path[fill=blue!12!white,draw=blue!12!white,line width=.12pt] (0.748354,1.479402)--(0.801283,1.607317)--(0.607814,1.641929)--(0.567664,1.511728)--cycle;
\path[fill=blue!14!white,draw=blue!14!white,line width=.12pt] (-0.567664,1.511728)--(-0.607814,1.641929)--(-0.801283,1.607317)--(-0.748354,1.479402)--cycle;
\path[fill=blue!24!white,draw=blue!24!white,line width=.12pt] (-0.618176,0.611042)--(-0.612293,0.739104)--(-0.754235,0.705537)--(-0.761481,0.577153)--cycle;
\path[fill=blue!20!white,draw=blue!20!white,line width=.12pt] (0.761481,0.577153)--(0.754235,0.705537)--(0.612293,0.739104)--(0.618176,0.611042)--cycle;
\path[fill=blue!30!white,draw=blue!30!white,line width=.12pt] (2.768442,1.005012)--(2.808591,0.885774)--(2.647922,1.019805)--(2.610069,1.137127)--cycle;
\path[fill=blue!23!white,draw=blue!23!white,line width=.12pt] (-2.610069,1.137127)--(-2.647922,1.019805)--(-2.808591,0.885774)--(-2.768442,1.005012)--cycle;
\path[fill=blue!14!white,draw=blue!14!white,line width=.12pt] (0.000000,1.252323)--(0.000000,1.407055)--(-0.179794,1.402639)--(-0.170043,1.248146)--cycle;
\path[fill=blue!14!white,draw=blue!14!white,line width=.12pt] (0.170043,1.248146)--(0.179794,1.402639)--(0.000000,1.407055)--(0.000000,1.252323)--cycle;
\path[fill=blue!29!white,draw=blue!29!white,line width=.12pt] (1.417944,0.091940)--(1.341037,0.165370)--(1.226706,0.228443)--(1.297056,0.158630)--cycle;
\path[fill=blue!34!white,draw=blue!34!white,line width=.12pt] (-1.297056,0.158630)--(-1.226706,0.228443)--(-1.341037,0.165370)--(-1.417944,0.091940)--cycle;
\path[fill=blue!16!white,draw=blue!16!white,line width=.12pt] (0.315143,0.927329)--(0.324025,1.079622)--(0.162796,1.091580)--(0.158334,0.938959)--cycle;
\path[fill=blue!18!white,draw=blue!18!white,line width=.12pt] (-0.158334,0.938959)--(-0.162796,1.091580)--(-0.324025,1.079622)--(-0.315143,0.927329)--cycle;
\path[fill=blue!30!white,draw=blue!30!white,line width=.12pt] (-2.808591,0.615449)--(-2.768442,0.474751)--(-2.900154,0.335511)--(-2.942213,0.474190)--cycle;
\path[fill=blue!37!white,draw=blue!37!white,line width=.12pt] (2.942213,0.474190)--(2.900154,0.335511)--(2.768442,0.474751)--(2.808591,0.615449)--cycle;
\path[fill=blue!32!white,draw=blue!32!white,line width=.12pt] (1.625975,-0.032838)--(1.525177,0.019647)--(1.417944,0.091940)--(1.511656,0.044233)--cycle;
\path[fill=blue!36!white,draw=blue!36!white,line width=.12pt] (-1.511656,0.044233)--(-1.417944,0.091940)--(-1.525177,0.019647)--(-1.625975,-0.032838)--cycle;
\path[fill=blue!14!white,draw=blue!14!white,line width=.12pt] (-0.357857,1.389432)--(-0.381508,1.535043)--(-0.567664,1.511728)--(-0.532473,1.367563)--cycle;
\path[fill=blue!12!white,draw=blue!12!white,line width=.12pt] (0.532473,1.367563)--(0.567664,1.511728)--(0.381508,1.535043)--(0.357857,1.389432)--cycle;
\path[fill=blue!15!white,draw=blue!15!white,line width=.12pt] (0.162796,1.091580)--(0.170043,1.248146)--(0.000000,1.252323)--(0.000000,1.095579)--cycle;
\path[fill=blue!15!white,draw=blue!15!white,line width=.12pt] (0.000000,1.095579)--(0.000000,1.252323)--(-0.170043,1.248146)--(-0.162796,1.091580)--cycle;
\path[fill=blue!26!white,draw=blue!26!white,line width=.12pt] (1.226706,0.228443)--(1.174431,0.322085)--(1.053662,0.376814)--(1.100561,0.285608)--cycle;
\path[fill=blue!31!white,draw=blue!31!white,line width=.12pt] (-1.100561,0.285608)--(-1.053662,0.376814)--(-1.174431,0.322085)--(-1.226706,0.228443)--cycle;
\path[fill=blue!13!white,draw=blue!13!white,line width=.12pt] (0.338447,1.235655)--(0.357857,1.389432)--(0.179794,1.402639)--(0.170043,1.248146)--cycle;
\path[fill=blue!14!white,draw=blue!14!white,line width=.12pt] (-0.170043,1.248146)--(-0.179794,1.402639)--(-0.357857,1.389432)--(-0.338447,1.235655)--cycle;
\path[fill=blue!18!white,draw=blue!18!white,line width=.12pt] (0.612293,0.739104)--(0.618176,0.881367)--(0.468918,0.908070)--(0.464455,0.765552)--cycle;
\path[fill=blue!21!white,draw=blue!21!white,line width=.12pt] (-0.464455,0.765552)--(-0.468918,0.908070)--(-0.618176,0.881367)--(-0.612293,0.739104)--cycle;
\path[fill=blue!16!white,draw=blue!16!white,line width=.12pt] (-0.162796,1.091580)--(-0.170043,1.248146)--(-0.338447,1.235655)--(-0.324025,1.079622)--cycle;
\path[fill=blue!15!white,draw=blue!15!white,line width=.12pt] (0.324025,1.079622)--(0.338447,1.235655)--(0.170043,1.248146)--(0.162796,1.091580)--cycle;
\path[fill=blue!16!white,draw=blue!16!white,line width=.12pt] (-1.913535,1.596850)--(-2.011334,1.581726)--(-2.198794,1.478311)--(-2.091880,1.498463)--cycle;
\path[fill=blue!20!white,draw=blue!20!white,line width=.12pt] (2.091880,1.498463)--(2.198794,1.478311)--(2.011334,1.581726)--(1.913535,1.596850)--cycle;
\path[fill=blue!16!white,draw=blue!16!white,line width=.12pt] (0.468918,0.908070)--(0.482133,1.059820)--(0.324025,1.079622)--(0.315143,0.927329)--cycle;
\path[fill=blue!18!white,draw=blue!18!white,line width=.12pt] (-0.315143,0.927329)--(-0.324025,1.079622)--(-0.482133,1.059820)--(-0.468918,0.908070)--cycle;
\path[fill=blue!40!white,draw=blue!40!white,line width=.12pt] (-2.116611,-0.127144)--(-1.978968,-0.150618)--(-2.073119,-0.250151)--(-2.217311,-0.233600)--cycle;
\path[fill=blue!39!white,draw=blue!39!white,line width=.12pt] (2.217311,-0.233600)--(2.073119,-0.250151)--(1.978968,-0.150618)--(2.116611,-0.127144)--cycle;
\path[fill=blue!23!white,draw=blue!23!white,line width=.12pt] (1.053662,0.376814)--(1.024781,0.489187)--(0.897452,0.536404)--(0.922745,0.425362)--cycle;
\path[fill=blue!28!white,draw=blue!28!white,line width=.12pt] (-0.922745,0.425362)--(-0.897452,0.536404)--(-1.024781,0.489187)--(-1.053662,0.376814)--cycle;
\path[fill=blue!38!white,draw=blue!38!white,line width=.12pt] (-1.740975,-0.058441)--(-1.625975,-0.032838)--(-1.724636,-0.115141)--(-1.846614,-0.146565)--cycle;
\path[fill=blue!35!white,draw=blue!35!white,line width=.12pt] (1.846614,-0.146565)--(1.724636,-0.115141)--(1.625975,-0.032838)--(1.740975,-0.058441)--cycle;
\path[fill=blue!15!white,draw=blue!15!white,line width=.12pt] (-0.338447,1.235655)--(-0.357857,1.389432)--(-0.532473,1.367563)--(-0.503593,1.214972)--cycle;
\path[fill=blue!13!white,draw=blue!13!white,line width=.12pt] (0.503593,1.214972)--(0.532473,1.367563)--(0.357857,1.389432)--(0.338447,1.235655)--cycle;
\path[fill=blue!14!white,draw=blue!14!white,line width=.12pt] (0.482133,1.059820)--(0.503593,1.214972)--(0.338447,1.235655)--(0.324025,1.079622)--cycle;
\path[fill=blue!16!white,draw=blue!16!white,line width=.12pt] (-0.324025,1.079622)--(-0.338447,1.235655)--(-0.503593,1.214972)--(-0.482133,1.059820)--cycle;
\path[fill=blue!12!white,draw=blue!12!white,line width=.12pt] (0.701962,1.337241)--(0.748354,1.479402)--(0.567664,1.511728)--(0.532473,1.367563)--cycle;
\path[fill=blue!14!white,draw=blue!14!white,line width=.12pt] (-0.532473,1.367563)--(-0.567664,1.511728)--(-0.748354,1.479402)--(-0.701962,1.337241)--cycle;
\path[fill=blue!12!white,draw=blue!12!white,line width=.12pt] (0.921837,1.438377)--(0.987036,1.563389)--(0.801283,1.607317)--(0.748354,1.479402)--cycle;
\path[fill=blue!14!white,draw=blue!14!white,line width=.12pt] (-0.748354,1.479402)--(-0.801283,1.607317)--(-0.987036,1.563389)--(-0.921837,1.438377)--cycle;
\path[fill=blue!20!white,draw=blue!20!white,line width=.12pt] (0.897452,0.536404)--(0.888912,0.665176)--(0.754235,0.705537)--(0.761481,0.577153)--cycle;
\path[fill=blue!25!white,draw=blue!25!white,line width=.12pt] (-0.761481,0.577153)--(-0.754235,0.705537)--(-0.888912,0.665176)--(-0.897452,0.536404)--cycle;
\path[fill=blue!13!white,draw=blue!13!white,line width=.12pt] (1.163283,1.510570)--(1.246659,1.612387)--(1.057780,1.668992)--(0.987036,1.563389)--cycle;
\path[fill=blue!14!white,draw=blue!14!white,line width=.12pt] (-0.987036,1.563389)--(-1.057780,1.668992)--(-1.246659,1.612387)--(-1.163283,1.510570)--cycle;
\path[fill=blue!27!white,draw=blue!27!white,line width=.12pt] (-2.822148,0.754235)--(-2.808591,0.615449)--(-2.942213,0.474190)--(-2.956415,0.612293)--cycle;
\path[fill=blue!35!white,draw=blue!35!white,line width=.12pt] (2.956415,0.612293)--(2.942213,0.474190)--(2.808591,0.615449)--(2.822148,0.754235)--cycle;
\path[fill=blue!19!white,draw=blue!19!white,line width=.12pt] (-0.468918,0.908070)--(-0.482133,1.059820)--(-0.635598,1.032365)--(-0.618176,0.881367)--cycle;
\path[fill=blue!16!white,draw=blue!16!white,line width=.12pt] (0.618176,0.881367)--(0.635598,1.032365)--(0.482133,1.059820)--(0.468918,0.908070)--cycle;
\path[fill=blue!22!white,draw=blue!22!white,line width=.12pt] (-0.612293,0.739104)--(-0.618176,0.881367)--(-0.761481,0.847478)--(-0.754235,0.705537)--cycle;
\path[fill=blue!18!white,draw=blue!18!white,line width=.12pt] (0.754235,0.705537)--(0.761481,0.847478)--(0.618176,0.881367)--(0.612293,0.739104)--cycle;
\path[fill=blue!14!white,draw=blue!14!white,line width=.12pt] (1.423533,1.546798)--(1.522544,1.620433)--(1.333369,1.690584)--(1.246659,1.612387)--cycle;
\path[fill=blue!14!white,draw=blue!14!white,line width=.12pt] (-1.246659,1.612387)--(-1.333369,1.690584)--(-1.522544,1.620433)--(-1.423533,1.546798)--cycle;
\path[fill=blue!15!white,draw=blue!15!white,line width=.12pt] (-0.503593,1.214972)--(-0.532473,1.367563)--(-0.701962,1.337241)--(-0.663889,1.186295)--cycle;
\path[fill=blue!13!white,draw=blue!13!white,line width=.12pt] (0.663889,1.186295)--(0.701962,1.337241)--(0.532473,1.367563)--(0.503593,1.214972)--cycle;
\path[fill=blue!14!white,draw=blue!14!white,line width=.12pt] (0.635598,1.032365)--(0.663889,1.186295)--(0.503593,1.214972)--(0.482133,1.059820)--cycle;
\path[fill=blue!17!white,draw=blue!17!white,line width=.12pt] (-0.482133,1.059820)--(-0.503593,1.214972)--(-0.663889,1.186295)--(-0.635598,1.032365)--cycle;
\path[fill=blue!15!white,draw=blue!15!white,line width=.12pt] (-1.522544,1.620433)--(-1.621555,1.667443)--(-1.807416,1.583216)--(-1.697056,1.541348)--cycle;
\path[fill=blue!15!white,draw=blue!15!white,line width=.12pt] (1.697056,1.541348)--(1.807416,1.583216)--(1.621555,1.667443)--(1.522544,1.620433)--cycle;
\path[fill=blue!21!white,draw=blue!21!white,line width=.12pt] (-2.548600,1.236371)--(-2.610069,1.137127)--(-2.768442,1.005012)--(-2.703244,1.107367)--cycle;
\path[fill=blue!28!white,draw=blue!28!white,line width=.12pt] (2.703244,1.107367)--(2.768442,1.005012)--(2.610069,1.137127)--(2.548600,1.236371)--cycle;
\path[fill=blue!18!white,draw=blue!18!white,line width=.12pt] (-2.198794,1.478311)--(-2.292506,1.430603)--(-2.465877,1.313721)--(-2.365079,1.366207)--cycle;
\path[fill=blue!22!white,draw=blue!22!white,line width=.12pt] (2.365079,1.366207)--(2.465877,1.313721)--(2.292506,1.430603)--(2.198794,1.478311)--cycle;
\path[fill=blue!39!white,draw=blue!39!white,line width=.12pt] (-1.978968,-0.150618)--(-1.846614,-0.146565)--(-1.934468,-0.239441)--(-2.073119,-0.250151)--cycle;
\path[fill=blue!37!white,draw=blue!37!white,line width=.12pt] (2.073119,-0.250151)--(1.934468,-0.239441)--(1.846614,-0.146565)--(1.978968,-0.150618)--cycle;
\path[fill=blue!12!white,draw=blue!12!white,line width=.12pt] (0.864690,1.298758)--(0.921837,1.438377)--(0.748354,1.479402)--(0.701962,1.337241)--cycle;
\path[fill=blue!15!white,draw=blue!15!white,line width=.12pt] (-0.701962,1.337241)--(-0.748354,1.479402)--(-0.921837,1.438377)--(-0.864690,1.298758)--cycle;
\path[fill=blue!32!white,draw=blue!32!white,line width=.12pt] (-1.226706,0.228443)--(-1.174431,0.322085)--(-1.283890,0.261700)--(-1.341037,0.165370)--cycle;
\path[fill=blue!26!white,draw=blue!26!white,line width=.12pt] (1.341037,0.165370)--(1.283890,0.261700)--(1.174431,0.322085)--(1.226706,0.228443)--cycle;
\path[fill=blue!29!white,draw=blue!29!white,line width=.12pt] (1.525177,0.019647)--(1.442454,0.096998)--(1.341037,0.165370)--(1.417944,0.091940)--cycle;
\path[fill=blue!35!white,draw=blue!35!white,line width=.12pt] (-1.417944,0.091940)--(-1.341037,0.165370)--(-1.442454,0.096998)--(-1.525177,0.019647)--cycle;
\path[fill=blue!23!white,draw=blue!23!white,line width=.12pt] (1.174431,0.322085)--(1.142240,0.435958)--(1.024781,0.489187)--(1.053662,0.376814)--cycle;
\path[fill=blue!29!white,draw=blue!29!white,line width=.12pt] (-1.053662,0.376814)--(-1.024781,0.489187)--(-1.142240,0.435958)--(-1.174431,0.322085)--cycle;
\path[fill=blue!16!white,draw=blue!16!white,line width=.12pt] (0.761481,0.847478)--(0.782941,0.997521)--(0.635598,1.032365)--(0.618176,0.881367)--cycle;
\path[fill=blue!19!white,draw=blue!19!white,line width=.12pt] (-0.618176,0.881367)--(-0.635598,1.032365)--(-0.782941,0.997521)--(-0.761481,0.847478)--cycle;
\path[fill=blue!37!white,draw=blue!37!white,line width=.12pt] (-2.627935,-0.031796)--(-2.500154,-0.122284)--(-2.589621,-0.247306)--(-2.721975,-0.163208)--cycle;
\path[fill=blue!41!white,draw=blue!41!white,line width=.12pt] (2.721975,-0.163208)--(2.589621,-0.247306)--(2.500154,-0.122284)--(2.627935,-0.031796)--cycle;
\path[fill=blue!20!white,draw=blue!20!white,line width=.12pt] (1.024781,0.489187)--(1.015029,0.618408)--(0.888912,0.665176)--(0.897452,0.536404)--cycle;
\path[fill=blue!26!white,draw=blue!26!white,line width=.12pt] (-0.897452,0.536404)--(-0.888912,0.665176)--(-1.015029,0.618408)--(-1.024781,0.489187)--cycle;
\path[fill=blue!41!white,draw=blue!41!white,line width=.12pt] (2.837984,-0.059451)--(2.721975,-0.163208)--(2.627935,-0.031796)--(2.739936,0.077561)--cycle;
\path[fill=blue!36!white,draw=blue!36!white,line width=.12pt] (-2.739936,0.077561)--(-2.627935,-0.031796)--(-2.721975,-0.163208)--(-2.837984,-0.059451)--cycle;
\path[fill=blue!22!white,draw=blue!22!white,line width=.12pt] (-0.754235,0.705537)--(-0.761481,0.847478)--(-0.897452,0.806729)--(-0.888912,0.665176)--cycle;
\path[fill=blue!18!white,draw=blue!18!white,line width=.12pt] (0.888912,0.665176)--(0.897452,0.806729)--(0.761481,0.847478)--(0.754235,0.705537)--cycle;
\path[fill=blue!12!white,draw=blue!12!white,line width=.12pt] (1.086442,1.389047)--(1.163283,1.510570)--(0.987036,1.563389)--(0.921837,1.438377)--cycle;
\path[fill=blue!14!white,draw=blue!14!white,line width=.12pt] (-0.921837,1.438377)--(-0.987036,1.563389)--(-1.163283,1.510570)--(-1.086442,1.389047)--cycle;
\path[fill=blue!32!white,draw=blue!32!white,line width=.12pt] (1.724636,-0.115141)--(1.617721,-0.057553)--(1.525177,0.019647)--(1.625975,-0.032838)--cycle;
\path[fill=blue!37!white,draw=blue!37!white,line width=.12pt] (-1.625975,-0.032838)--(-1.525177,0.019647)--(-1.617721,-0.057553)--(-1.724636,-0.115141)--cycle;
\path[fill=blue!38!white,draw=blue!38!white,line width=.12pt] (-2.500154,-0.122284)--(-2.361503,-0.190425)--(-2.446009,-0.308514)--(-2.589621,-0.247306)--cycle;
\path[fill=blue!41!white,draw=blue!41!white,line width=.12pt] (2.589621,-0.247306)--(2.446009,-0.308514)--(2.361503,-0.190425)--(2.500154,-0.122284)--cycle;
\path[fill=blue!17!white,draw=blue!17!white,line width=.12pt] (1.975871,1.490286)--(2.091880,1.498463)--(1.913535,1.596850)--(1.807416,1.583216)--cycle;
\path[fill=blue!15!white,draw=blue!15!white,line width=.12pt] (-1.807416,1.583216)--(-1.913535,1.596850)--(-2.091880,1.498463)--(-1.975871,1.490286)--cycle;
\path[fill=blue!13!white,draw=blue!13!white,line width=.12pt] (0.817791,1.149900)--(0.864690,1.298758)--(0.701962,1.337241)--(0.663889,1.186295)--cycle;
\path[fill=blue!16!white,draw=blue!16!white,line width=.12pt] (-0.663889,1.186295)--(-0.701962,1.337241)--(-0.864690,1.298758)--(-0.817791,1.149900)--cycle;
\path[fill=blue!25!white,draw=blue!25!white,line width=.12pt] (-2.808591,0.885774)--(-2.822148,0.754235)--(-2.956415,0.612293)--(-2.942213,0.744515)--cycle;
\path[fill=blue!33!white,draw=blue!33!white,line width=.12pt] (2.942213,0.744515)--(2.956415,0.612293)--(2.822148,0.754235)--(2.808591,0.885774)--cycle;
\path[fill=blue!17!white,draw=blue!17!white,line width=.12pt] (-0.635598,1.032365)--(-0.663889,1.186295)--(-0.817791,1.149900)--(-0.782941,0.997521)--cycle;
\path[fill=blue!14!white,draw=blue!14!white,line width=.12pt] (0.782941,0.997521)--(0.817791,1.149900)--(0.663889,1.186295)--(0.635598,1.032365)--cycle;
\path[fill=blue!33!white,draw=blue!33!white,line width=.12pt] (-2.831853,0.201586)--(-2.739936,0.077561)--(-2.837984,-0.059451)--(-2.933190,0.059976)--cycle;
\path[fill=blue!40!white,draw=blue!40!white,line width=.12pt] (2.933190,0.059976)--(2.837984,-0.059451)--(2.739936,0.077561)--(2.831853,0.201586)--cycle;
\path[fill=blue!14!white,draw=blue!14!white,line width=.12pt] (-1.163283,1.510570)--(-1.246659,1.612387)--(-1.423533,1.546798)--(-1.328326,1.449368)--cycle;
\path[fill=blue!13!white,draw=blue!13!white,line width=.12pt] (1.328326,1.449368)--(1.423533,1.546798)--(1.246659,1.612387)--(1.163283,1.510570)--cycle;
\path[fill=blue!39!white,draw=blue!39!white,line width=.12pt] (-2.361503,-0.190425)--(-2.217311,-0.233600)--(-2.296657,-0.344479)--(-2.446009,-0.308514)--cycle;
\path[fill=blue!40!white,draw=blue!40!white,line width=.12pt] (2.446009,-0.308514)--(2.296657,-0.344479)--(2.217311,-0.233600)--(2.361503,-0.190425)--cycle;
\path[fill=blue!39!white,draw=blue!39!white,line width=.12pt] (-1.846614,-0.146565)--(-1.724636,-0.115141)--(-1.806687,-0.201882)--(-1.934468,-0.239441)--cycle;
\path[fill=blue!35!white,draw=blue!35!white,line width=.12pt] (1.934468,-0.239441)--(1.806687,-0.201882)--(1.724636,-0.115141)--(1.846614,-0.146565)--cycle;
\path[fill=blue!20!white,draw=blue!20!white,line width=.12pt] (-0.761481,0.847478)--(-0.782941,0.997521)--(-0.922745,0.955623)--(-0.897452,0.806729)--cycle;
\path[fill=blue!16!white,draw=blue!16!white,line width=.12pt] (0.897452,0.806729)--(0.922745,0.955623)--(0.782941,0.997521)--(0.761481,0.847478)--cycle;
\path[fill=blue!12!white,draw=blue!12!white,line width=.12pt] (1.019091,1.252486)--(1.086442,1.389047)--(0.921837,1.438377)--(0.864690,1.298758)--cycle;
\path[fill=blue!15!white,draw=blue!15!white,line width=.12pt] (-0.864690,1.298758)--(-0.921837,1.438377)--(-1.086442,1.389047)--(-1.019091,1.252486)--cycle;
\path[fill=blue!14!white,draw=blue!14!white,line width=.12pt] (1.586697,1.472856)--(1.697056,1.541348)--(1.522544,1.620433)--(1.423533,1.546798)--cycle;
\path[fill=blue!14!white,draw=blue!14!white,line width=.12pt] (-1.423533,1.546798)--(-1.522544,1.620433)--(-1.697056,1.541348)--(-1.586697,1.472856)--cycle;
\path[fill=blue!31!white,draw=blue!31!white,line width=.12pt] (-2.900154,0.335511)--(-2.831853,0.201586)--(-2.933190,0.059976)--(-3.003935,0.190486)--cycle;
\path[fill=blue!39!white,draw=blue!39!white,line width=.12pt] (3.003935,0.190486)--(2.933190,0.059976)--(2.831853,0.201586)--(2.900154,0.335511)--cycle;
\path[fill=blue!25!white,draw=blue!25!white,line width=.12pt] (2.615501,1.188905)--(2.703244,1.107367)--(2.548600,1.236371)--(2.465877,1.313721)--cycle;
\path[fill=blue!19!white,draw=blue!19!white,line width=.12pt] (-2.465877,1.313721)--(-2.548600,1.236371)--(-2.703244,1.107367)--(-2.615501,1.188905)--cycle;
\path[fill=blue!19!white,draw=blue!19!white,line width=.12pt] (2.250079,1.391810)--(2.365079,1.366207)--(2.198794,1.478311)--(2.091880,1.498463)--cycle;
\path[fill=blue!16!white,draw=blue!16!white,line width=.12pt] (-2.091880,1.498463)--(-2.198794,1.478311)--(-2.365079,1.366207)--(-2.250079,1.391810)--cycle;
\path[fill=blue!23!white,draw=blue!23!white,line width=.12pt] (-0.888912,0.665176)--(-0.897452,0.806729)--(-1.024781,0.759512)--(-1.015029,0.618408)--cycle;
\path[fill=blue!18!white,draw=blue!18!white,line width=.12pt] (1.015029,0.618408)--(1.024781,0.759512)--(0.897452,0.806729)--(0.888912,0.665176)--cycle;
\path[fill=blue!40!white,draw=blue!40!white,line width=.12pt] (-2.217311,-0.233600)--(-2.073119,-0.250151)--(-2.147305,-0.353819)--(-2.296657,-0.344479)--cycle;
\path[fill=blue!39!white,draw=blue!39!white,line width=.12pt] (2.296657,-0.344479)--(2.147305,-0.353819)--(2.073119,-0.250151)--(2.217311,-0.233600)--cycle;
\path[fill=blue!14!white,draw=blue!14!white,line width=.12pt] (0.922745,0.955623)--(0.963817,1.106137)--(0.817791,1.149900)--(0.782941,0.997521)--cycle;
\path[fill=blue!18!white,draw=blue!18!white,line width=.12pt] (-0.782941,0.997521)--(-0.817791,1.149900)--(-0.963817,1.106137)--(-0.922745,0.955623)--cycle;
\path[fill=blue!13!white,draw=blue!13!white,line width=.12pt] (0.963817,1.106137)--(1.019091,1.252486)--(0.864690,1.298758)--(0.817791,1.149900)--cycle;
\path[fill=blue!16!white,draw=blue!16!white,line width=.12pt] (-0.817791,1.149900)--(-0.864690,1.298758)--(-1.019091,1.252486)--(-0.963817,1.106137)--cycle;
\path[fill=blue!26!white,draw=blue!26!white,line width=.12pt] (-1.024781,0.489187)--(-1.015029,0.618408)--(-1.131371,0.565685)--(-1.142240,0.435958)--cycle;
\path[fill=blue!20!white,draw=blue!20!white,line width=.12pt] (1.142240,0.435958)--(1.131371,0.565685)--(1.015029,0.618408)--(1.024781,0.489187)--cycle;
\path[fill=blue!30!white,draw=blue!30!white,line width=.12pt] (-1.174431,0.322085)--(-1.142240,0.435958)--(-1.248699,0.377228)--(-1.283890,0.261700)--cycle;
\path[fill=blue!23!white,draw=blue!23!white,line width=.12pt] (1.283890,0.261700)--(1.248699,0.377228)--(1.142240,0.435958)--(1.174431,0.322085)--cycle;
\path[fill=blue!26!white,draw=blue!26!white,line width=.12pt] (1.442454,0.096998)--(1.380985,0.196241)--(1.283890,0.261700)--(1.341037,0.165370)--cycle;
\path[fill=blue!33!white,draw=blue!33!white,line width=.12pt] (-1.341037,0.165370)--(-1.283890,0.261700)--(-1.380985,0.196241)--(-1.442454,0.096998)--cycle;
\path[fill=blue!15!white,draw=blue!15!white,line width=.12pt] (-1.086442,1.389047)--(-1.163283,1.510570)--(-1.328326,1.449368)--(-1.240584,1.331887)--cycle;
\path[fill=blue!12!white,draw=blue!12!white,line width=.12pt] (1.240584,1.331887)--(1.328326,1.449368)--(1.163283,1.510570)--(1.086442,1.389047)--cycle;
\path[fill=blue!22!white,draw=blue!22!white,line width=.12pt] (-2.768442,1.005012)--(-2.808591,0.885774)--(-2.942213,0.744515)--(-2.900154,0.865772)--cycle;
\path[fill=blue!30!white,draw=blue!30!white,line width=.12pt] (2.900154,0.865772)--(2.942213,0.744515)--(2.808591,0.885774)--(2.768442,1.005012)--cycle;
\path[fill=blue!29!white,draw=blue!29!white,line width=.12pt] (1.617721,-0.057553)--(1.529978,0.023985)--(1.442454,0.096998)--(1.525177,0.019647)--cycle;
\path[fill=blue!35!white,draw=blue!35!white,line width=.12pt] (-1.525177,0.019647)--(-1.442454,0.096998)--(-1.529978,0.023985)--(-1.617721,-0.057553)--cycle;
\path[fill=blue!15!white,draw=blue!15!white,line width=.12pt] (1.855225,1.454092)--(1.975871,1.490286)--(1.807416,1.583216)--(1.697056,1.541348)--cycle;
\path[fill=blue!15!white,draw=blue!15!white,line width=.12pt] (-1.697056,1.541348)--(-1.807416,1.583216)--(-1.975871,1.490286)--(-1.855225,1.454092)--cycle;
\path[fill=blue!29!white,draw=blue!29!white,line width=.12pt] (-2.942213,0.474190)--(-2.900154,0.335511)--(-3.003935,0.190486)--(-3.047499,0.327062)--cycle;
\path[fill=blue!37!white,draw=blue!37!white,line width=.12pt] (3.047499,0.327062)--(3.003935,0.190486)--(2.900154,0.335511)--(2.942213,0.474190)--cycle;
\path[fill=blue!21!white,draw=blue!21!white,line width=.12pt] (-0.897452,0.806729)--(-0.922745,0.955623)--(-1.053662,0.907076)--(-1.024781,0.759512)--cycle;
\path[fill=blue!16!white,draw=blue!16!white,line width=.12pt] (1.024781,0.759512)--(1.053662,0.907076)--(0.922745,0.955623)--(0.897452,0.806729)--cycle;
\path[fill=blue!40!white,draw=blue!40!white,line width=.12pt] (-2.073119,-0.250151)--(-1.934468,-0.239441)--(-2.003693,-0.336176)--(-2.147305,-0.353819)--cycle;
\path[fill=blue!37!white,draw=blue!37!white,line width=.12pt] (2.147305,-0.353819)--(2.003693,-0.336176)--(1.934468,-0.239441)--(2.073119,-0.250151)--cycle;
\path[fill=blue!14!white,draw=blue!14!white,line width=.12pt] (-1.328326,1.449368)--(-1.423533,1.546798)--(-1.586697,1.472856)--(-1.480578,1.380371)--cycle;
\path[fill=blue!13!white,draw=blue!13!white,line width=.12pt] (1.480578,1.380371)--(1.586697,1.472856)--(1.423533,1.546798)--(1.328326,1.449368)--cycle;
\path[fill=blue!12!white,draw=blue!12!white,line width=.12pt] (1.163677,1.198870)--(1.240584,1.331887)--(1.086442,1.389047)--(1.019091,1.252486)--cycle;
\path[fill=blue!16!white,draw=blue!16!white,line width=.12pt] (-1.019091,1.252486)--(-1.086442,1.389047)--(-1.240584,1.331887)--(-1.163677,1.198870)--cycle;
\path[fill=blue!37!white,draw=blue!37!white,line width=.12pt] (-1.724636,-0.115141)--(-1.617721,-0.057553)--(-1.694686,-0.138917)--(-1.806687,-0.201882)--cycle;
\path[fill=blue!32!white,draw=blue!32!white,line width=.12pt] (1.806687,-0.201882)--(1.694686,-0.138917)--(1.617721,-0.057553)--(1.724636,-0.115141)--cycle;
\path[fill=blue!19!white,draw=blue!19!white,line width=.12pt] (-0.922745,0.955623)--(-0.963817,1.106137)--(-1.100561,1.055429)--(-1.053662,0.907076)--cycle;
\path[fill=blue!14!white,draw=blue!14!white,line width=.12pt] (1.053662,0.907076)--(1.100561,1.055429)--(0.963817,1.106137)--(0.922745,0.955623)--cycle;
\path[fill=blue!24!white,draw=blue!24!white,line width=.12pt] (-1.015029,0.618408)--(-1.024781,0.759512)--(-1.142240,0.706283)--(-1.131371,0.565685)--cycle;
\path[fill=blue!18!white,draw=blue!18!white,line width=.12pt] (1.131371,0.565685)--(1.142240,0.706283)--(1.024781,0.759512)--(1.015029,0.618408)--cycle;
\path[fill=blue!17!white,draw=blue!17!white,line width=.12pt] (-0.963817,1.106137)--(-1.019091,1.252486)--(-1.163677,1.198870)--(-1.100561,1.055429)--cycle;
\path[fill=blue!13!white,draw=blue!13!white,line width=.12pt] (1.100561,1.055429)--(1.163677,1.198870)--(1.019091,1.252486)--(0.963817,1.106137)--cycle;
\path[fill=blue!22!white,draw=blue!22!white,line width=.12pt] (2.508586,1.246493)--(2.615501,1.188905)--(2.465877,1.313721)--(2.365079,1.366207)--cycle;
\path[fill=blue!17!white,draw=blue!17!white,line width=.12pt] (-2.365079,1.366207)--(-2.465877,1.313721)--(-2.615501,1.188905)--(-2.508586,1.246493)--cycle;
\path[fill=blue!15!white,draw=blue!15!white,line width=.12pt] (-1.975871,1.490286)--(-2.091880,1.498463)--(-2.250079,1.391810)--(-2.125296,1.389547)--cycle;
\path[fill=blue!17!white,draw=blue!17!white,line width=.12pt] (2.125296,1.389547)--(2.250079,1.391810)--(2.091880,1.498463)--(1.975871,1.490286)--cycle;
\path[fill=blue!27!white,draw=blue!27!white,line width=.12pt] (-1.142240,0.435958)--(-1.131371,0.565685)--(-1.236817,0.507515)--(-1.248699,0.377228)--cycle;
\path[fill=blue!20!white,draw=blue!20!white,line width=.12pt] (1.248699,0.377228)--(1.236817,0.507515)--(1.131371,0.565685)--(1.142240,0.435958)--cycle;
\path[fill=blue!26!white,draw=blue!26!white,line width=.12pt] (-2.956415,0.612293)--(-2.942213,0.474190)--(-3.047499,0.327062)--(-3.062209,0.464455)--cycle;
\path[fill=blue!35!white,draw=blue!35!white,line width=.12pt] (3.062209,0.464455)--(3.047499,0.327062)--(2.942213,0.474190)--(2.956415,0.612293)--cycle;
\path[fill=blue!23!white,draw=blue!23!white,line width=.12pt] (1.380985,0.196241)--(1.343133,0.313564)--(1.248699,0.377228)--(1.283890,0.261700)--cycle;
\path[fill=blue!30!white,draw=blue!30!white,line width=.12pt] (-1.283890,0.261700)--(-1.248699,0.377228)--(-1.343133,0.313564)--(-1.380985,0.196241)--cycle;
\path[fill=blue!20!white,draw=blue!20!white,line width=.12pt] (-2.703244,1.107367)--(-2.768442,1.005012)--(-2.900154,0.865772)--(-2.831853,0.971406)--cycle;
\path[fill=blue!27!white,draw=blue!27!white,line width=.12pt] (2.831853,0.971406)--(2.900154,0.865772)--(2.768442,1.005012)--(2.703244,1.107367)--cycle;
\path[fill=blue!14!white,draw=blue!14!white,line width=.12pt] (1.734580,1.391274)--(1.855225,1.454092)--(1.697056,1.541348)--(1.586697,1.472856)--cycle;
\path[fill=blue!14!white,draw=blue!14!white,line width=.12pt] (-1.586697,1.472856)--(-1.697056,1.541348)--(-1.855225,1.454092)--(-1.734580,1.391274)--cycle;
\path[fill=blue!15!white,draw=blue!15!white,line width=.12pt] (-1.240584,1.331887)--(-1.328326,1.449368)--(-1.480578,1.380371)--(-1.382778,1.267448)--cycle;
\path[fill=blue!12!white,draw=blue!12!white,line width=.12pt] (1.382778,1.267448)--(1.480578,1.380371)--(1.328326,1.449368)--(1.240584,1.331887)--cycle;
\path[fill=blue!33!white,draw=blue!33!white,line width=.12pt] (-1.442454,0.096998)--(-1.380985,0.196241)--(-1.464780,0.126340)--(-1.529978,0.023985)--cycle;
\path[fill=blue!26!white,draw=blue!26!white,line width=.12pt] (1.529978,0.023985)--(1.464780,0.126340)--(1.380985,0.196241)--(1.442454,0.096998)--cycle;
\path[fill=blue!39!white,draw=blue!39!white,line width=.12pt] (-1.934468,-0.239441)--(-1.806687,-0.201882)--(-1.871339,-0.292227)--(-2.003693,-0.336176)--cycle;
\path[fill=blue!35!white,draw=blue!35!white,line width=.12pt] (2.003693,-0.336176)--(1.871339,-0.292227)--(1.806687,-0.201882)--(1.934468,-0.239441)--cycle;
\path[fill=blue!37!white,draw=blue!37!white,line width=.12pt] (-2.721975,-0.163208)--(-2.589621,-0.247306)--(-2.654149,-0.376111)--(-2.789801,-0.298596)--cycle;
\path[fill=blue!41!white,draw=blue!41!white,line width=.12pt] (2.789801,-0.298596)--(2.654149,-0.376111)--(2.589621,-0.247306)--(2.721975,-0.163208)--cycle;
\path[fill=blue!16!white,draw=blue!16!white,line width=.12pt] (1.142240,0.706283)--(1.174431,0.852346)--(1.053662,0.907076)--(1.024781,0.759512)--cycle;
\path[fill=blue!21!white,draw=blue!21!white,line width=.12pt] (-1.024781,0.759512)--(-1.053662,0.907076)--(-1.174431,0.852346)--(-1.142240,0.706283)--cycle;
\path[fill=blue!38!white,draw=blue!38!white,line width=.12pt] (-2.589621,-0.247306)--(-2.446009,-0.308514)--(-2.506958,-0.430175)--(-2.654149,-0.376111)--cycle;
\path[fill=blue!41!white,draw=blue!41!white,line width=.12pt] (2.654149,-0.376111)--(2.506958,-0.430175)--(2.446009,-0.308514)--(2.589621,-0.247306)--cycle;
\path[fill=blue!35!white,draw=blue!35!white,line width=.12pt] (-2.837984,-0.059451)--(-2.721975,-0.163208)--(-2.789801,-0.298596)--(-2.908701,-0.200610)--cycle;
\path[fill=blue!41!white,draw=blue!41!white,line width=.12pt] (2.908701,-0.200610)--(2.789801,-0.298596)--(2.721975,-0.163208)--(2.837984,-0.059451)--cycle;
\path[fill=blue!39!white,draw=blue!39!white,line width=.12pt] (-2.446009,-0.308514)--(-2.296657,-0.344479)--(-2.353885,-0.458712)--(-2.506958,-0.430175)--cycle;
\path[fill=blue!40!white,draw=blue!40!white,line width=.12pt] (2.506958,-0.430175)--(2.353885,-0.458712)--(2.296657,-0.344479)--(2.446009,-0.308514)--cycle;
\path[fill=blue!36!white,draw=blue!36!white,line width=.12pt] (-1.617721,-0.057553)--(-1.529978,0.023985)--(-1.602769,-0.052966)--(-1.694686,-0.138917)--cycle;
\path[fill=blue!30!white,draw=blue!30!white,line width=.12pt] (1.694686,-0.138917)--(1.602769,-0.052966)--(1.529978,0.023985)--(1.617721,-0.057553)--cycle;
\path[fill=blue!16!white,draw=blue!16!white,line width=.12pt] (-1.163677,1.198870)--(-1.240584,1.331887)--(-1.382778,1.267448)--(-1.297056,1.138426)--cycle;
\path[fill=blue!12!white,draw=blue!12!white,line width=.12pt] (1.297056,1.138426)--(1.382778,1.267448)--(1.240584,1.331887)--(1.163677,1.198870)--cycle;
\path[fill=blue!14!white,draw=blue!14!white,line width=.12pt] (1.174431,0.852346)--(1.226706,0.998263)--(1.100561,1.055429)--(1.053662,0.907076)--cycle;
\path[fill=blue!19!white,draw=blue!19!white,line width=.12pt] (-1.053662,0.907076)--(-1.100561,1.055429)--(-1.226706,0.998263)--(-1.174431,0.852346)--cycle;
\path[fill=blue!18!white,draw=blue!18!white,line width=.12pt] (1.236817,0.507515)--(1.248699,0.647553)--(1.142240,0.706283)--(1.131371,0.565685)--cycle;
\path[fill=blue!25!white,draw=blue!25!white,line width=.12pt] (-1.131371,0.565685)--(-1.142240,0.706283)--(-1.248699,0.647553)--(-1.236817,0.507515)--cycle;
\path[fill=blue!24!white,draw=blue!24!white,line width=.12pt] (-2.942213,0.744515)--(-2.956415,0.612293)--(-3.062209,0.464455)--(-3.047499,0.597387)--cycle;
\path[fill=blue!32!white,draw=blue!32!white,line width=.12pt] (3.047499,0.597387)--(3.062209,0.464455)--(2.956415,0.612293)--(2.942213,0.744515)--cycle;
\path[fill=blue!17!white,draw=blue!17!white,line width=.12pt] (-1.100561,1.055429)--(-1.163677,1.198870)--(-1.297056,1.138426)--(-1.226706,0.998263)--cycle;
\path[fill=blue!13!white,draw=blue!13!white,line width=.12pt] (1.226706,0.998263)--(1.297056,1.138426)--(1.163677,1.198870)--(1.100561,1.055429)--cycle;
\path[fill=blue!33!white,draw=blue!33!white,line width=.12pt] (-2.933190,0.059976)--(-2.837984,-0.059451)--(-2.908701,-0.200610)--(-3.006279,-0.085917)--cycle;
\path[fill=blue!40!white,draw=blue!40!white,line width=.12pt] (3.006279,-0.085917)--(2.908701,-0.200610)--(2.837984,-0.059451)--(2.933190,0.059976)--cycle;
\path[fill=blue!40!white,draw=blue!40!white,line width=.12pt] (-2.296657,-0.344479)--(-2.147305,-0.353819)--(-2.200811,-0.460624)--(-2.353885,-0.458712)--cycle;
\path[fill=blue!39!white,draw=blue!39!white,line width=.12pt] (2.353885,-0.458712)--(2.200811,-0.460624)--(2.147305,-0.353819)--(2.296657,-0.344479)--cycle;
\path[fill=blue!19!white,draw=blue!19!white,line width=.12pt] (2.386608,1.277917)--(2.508586,1.246493)--(2.365079,1.366207)--(2.250079,1.391810)--cycle;
\path[fill=blue!16!white,draw=blue!16!white,line width=.12pt] (-2.250079,1.391810)--(-2.365079,1.366207)--(-2.508586,1.246493)--(-2.386608,1.277917)--cycle;
\path[fill=blue!14!white,draw=blue!14!white,line width=.12pt] (-1.855225,1.454092)--(-1.975871,1.490286)--(-2.125296,1.389547)--(-1.995527,1.359505)--cycle;
\path[fill=blue!15!white,draw=blue!15!white,line width=.12pt] (1.995527,1.359505)--(2.125296,1.389547)--(1.975871,1.490286)--(1.855225,1.454092)--cycle;
\path[fill=blue!15!white,draw=blue!15!white,line width=.12pt] (-1.480578,1.380371)--(-1.586697,1.472856)--(-1.734580,1.391274)--(-1.618570,1.304246)--cycle;
\path[fill=blue!13!white,draw=blue!13!white,line width=.12pt] (1.618570,1.304246)--(1.734580,1.391274)--(1.586697,1.472856)--(1.480578,1.380371)--cycle;
\path[fill=blue!28!white,draw=blue!28!white,line width=.12pt] (-1.248699,0.377228)--(-1.236817,0.507515)--(-1.330351,0.444456)--(-1.343133,0.313564)--cycle;
\path[fill=blue!21!white,draw=blue!21!white,line width=.12pt] (1.343133,0.313564)--(1.330351,0.444456)--(1.236817,0.507515)--(1.248699,0.377228)--cycle;
\path[fill=blue!18!white,draw=blue!18!white,line width=.12pt] (-2.615501,1.188905)--(-2.703244,1.107367)--(-2.831853,0.971406)--(-2.739936,1.057357)--cycle;
\path[fill=blue!24!white,draw=blue!24!white,line width=.12pt] (2.739936,1.057357)--(2.831853,0.971406)--(2.703244,1.107367)--(2.615501,1.188905)--cycle;
\path[fill=blue!38!white,draw=blue!38!white,line width=.12pt] (-1.806687,-0.201882)--(-1.694686,-0.138917)--(-1.755330,-0.223661)--(-1.871339,-0.292227)--cycle;
\path[fill=blue!33!white,draw=blue!33!white,line width=.12pt] (1.871339,-0.292227)--(1.755330,-0.223661)--(1.694686,-0.138917)--(1.806687,-0.201882)--cycle;
\path[fill=blue!24!white,draw=blue!24!white,line width=.12pt] (1.464780,0.126340)--(1.424631,0.245578)--(1.343133,0.313564)--(1.380985,0.196241)--cycle;
\path[fill=blue!31!white,draw=blue!31!white,line width=.12pt] (-1.380985,0.196241)--(-1.343133,0.313564)--(-1.424631,0.245578)--(-1.464780,0.126340)--cycle;
\path[fill=blue!31!white,draw=blue!31!white,line width=.12pt] (-3.003935,0.190486)--(-2.933190,0.059976)--(-3.006279,-0.085917)--(-3.078787,0.041074)--cycle;
\path[fill=blue!38!white,draw=blue!38!white,line width=.12pt] (3.078787,0.041074)--(3.006279,-0.085917)--(2.933190,0.059976)--(3.003935,0.190486)--cycle;
\path[fill=blue!40!white,draw=blue!40!white,line width=.12pt] (-2.147305,-0.353819)--(-2.003693,-0.336176)--(-2.053620,-0.435837)--(-2.200811,-0.460624)--cycle;
\path[fill=blue!37!white,draw=blue!37!white,line width=.12pt] (2.200811,-0.460624)--(2.053620,-0.435837)--(2.003693,-0.336176)--(2.147305,-0.353819)--cycle;
\path[fill=blue!16!white,draw=blue!16!white,line width=.12pt] (1.248699,0.647553)--(1.283890,0.791962)--(1.174431,0.852346)--(1.142240,0.706283)--cycle;
\path[fill=blue!22!white,draw=blue!22!white,line width=.12pt] (-1.142240,0.706283)--(-1.174431,0.852346)--(-1.283890,0.791962)--(-1.248699,0.647553)--cycle;
\path[fill=blue!15!white,draw=blue!15!white,line width=.12pt] (-1.382778,1.267448)--(-1.480578,1.380371)--(-1.618570,1.304246)--(-1.511656,1.196351)--cycle;
\path[fill=blue!12!white,draw=blue!12!white,line width=.12pt] (1.511656,1.196351)--(1.618570,1.304246)--(1.480578,1.380371)--(1.382778,1.267448)--cycle;
\path[fill=blue!34!white,draw=blue!34!white,line width=.12pt] (-1.529978,0.023985)--(-1.464780,0.126340)--(-1.534468,0.052668)--(-1.602769,-0.052966)--cycle;
\path[fill=blue!27!white,draw=blue!27!white,line width=.12pt] (1.602769,-0.052966)--(1.534468,0.052668)--(1.464780,0.126340)--(1.529978,0.023985)--cycle;
\path[fill=blue!21!white,draw=blue!21!white,line width=.12pt] (-2.900154,0.865772)--(-2.942213,0.744515)--(-3.047499,0.597387)--(-3.003935,0.720748)--cycle;
\path[fill=blue!30!white,draw=blue!30!white,line width=.12pt] (3.003935,0.720748)--(3.047499,0.597387)--(2.942213,0.744515)--(2.900154,0.865772)--cycle;
\path[fill=blue!14!white,draw=blue!14!white,line width=.12pt] (-1.734580,1.391274)--(-1.855225,1.454092)--(-1.995527,1.359505)--(-1.865758,1.302838)--cycle;
\path[fill=blue!14!white,draw=blue!14!white,line width=.12pt] (1.865758,1.302838)--(1.995527,1.359505)--(1.855225,1.454092)--(1.734580,1.391274)--cycle;
\path[fill=blue!14!white,draw=blue!14!white,line width=.12pt] (1.283890,0.791962)--(1.341037,0.935191)--(1.226706,0.998263)--(1.174431,0.852346)--cycle;
\path[fill=blue!20!white,draw=blue!20!white,line width=.12pt] (-1.174431,0.852346)--(-1.226706,0.998263)--(-1.341037,0.935191)--(-1.283890,0.791962)--cycle;
\path[fill=blue!16!white,draw=blue!16!white,line width=.12pt] (-1.297056,1.138426)--(-1.382778,1.267448)--(-1.511656,1.196351)--(-1.417944,1.071736)--cycle;
\path[fill=blue!12!white,draw=blue!12!white,line width=.12pt] (1.417944,1.071736)--(1.511656,1.196351)--(1.382778,1.267448)--(1.297056,1.138426)--cycle;
\path[fill=blue!25!white,draw=blue!25!white,line width=.12pt] (-1.236817,0.507515)--(-1.248699,0.647553)--(-1.343133,0.583889)--(-1.330351,0.444456)--cycle;
\path[fill=blue!18!white,draw=blue!18!white,line width=.12pt] (1.330351,0.444456)--(1.343133,0.583889)--(1.248699,0.647553)--(1.236817,0.507515)--cycle;
\path[fill=blue!15!white,draw=blue!15!white,line width=.12pt] (-2.125296,1.389547)--(-2.250079,1.391810)--(-2.386608,1.277917)--(-2.254254,1.281970)--cycle;
\path[fill=blue!17!white,draw=blue!17!white,line width=.12pt] (2.254254,1.281970)--(2.386608,1.277917)--(2.250079,1.391810)--(2.125296,1.389547)--cycle;
\path[fill=blue!13!white,draw=blue!13!white,line width=.12pt] (1.341037,0.935191)--(1.417944,1.071736)--(1.297056,1.138426)--(1.226706,0.998263)--cycle;
\path[fill=blue!18!white,draw=blue!18!white,line width=.12pt] (-1.226706,0.998263)--(-1.297056,1.138426)--(-1.417944,1.071736)--(-1.341037,0.935191)--cycle;
\path[fill=blue!28!white,draw=blue!28!white,line width=.12pt] (-3.047499,0.327062)--(-3.003935,0.190486)--(-3.078787,0.041074)--(-3.123436,0.175483)--cycle;
\path[fill=blue!36!white,draw=blue!36!white,line width=.12pt] (3.123436,0.175483)--(3.078787,0.041074)--(3.003935,0.190486)--(3.047499,0.327062)--cycle;
\path[fill=blue!39!white,draw=blue!39!white,line width=.12pt] (-2.003693,-0.336176)--(-1.871339,-0.292227)--(-1.917969,-0.385305)--(-2.053620,-0.435837)--cycle;
\path[fill=blue!35!white,draw=blue!35!white,line width=.12pt] (2.053620,-0.435837)--(1.917969,-0.385305)--(1.871339,-0.292227)--(2.003693,-0.336176)--cycle;
\path[fill=blue!37!white,draw=blue!37!white,line width=.12pt] (-1.694686,-0.138917)--(-1.602769,-0.052966)--(-1.660123,-0.133114)--(-1.755330,-0.223661)--cycle;
\path[fill=blue!30!white,draw=blue!30!white,line width=.12pt] (1.755330,-0.223661)--(1.660123,-0.133114)--(1.602769,-0.052966)--(1.694686,-0.138917)--cycle;
\path[fill=blue!17!white,draw=blue!17!white,line width=.12pt] (-2.508586,1.246493)--(-2.615501,1.188905)--(-2.739936,1.057357)--(-2.627935,1.120322)--cycle;
\path[fill=blue!22!white,draw=blue!22!white,line width=.12pt] (2.627935,1.120322)--(2.739936,1.057357)--(2.615501,1.188905)--(2.508586,1.246493)--cycle;
\path[fill=blue!29!white,draw=blue!29!white,line width=.12pt] (-1.343133,0.313564)--(-1.330351,0.444456)--(-1.411074,0.377117)--(-1.424631,0.245578)--cycle;
\path[fill=blue!21!white,draw=blue!21!white,line width=.12pt] (1.424631,0.245578)--(1.411074,0.377117)--(1.330351,0.444456)--(1.343133,0.313564)--cycle;
\path[fill=blue!13!white,draw=blue!13!white,line width=.12pt] (1.740975,1.221724)--(1.865758,1.302838)--(1.734580,1.391274)--(1.618570,1.304246)--cycle;
\path[fill=blue!15!white,draw=blue!15!white,line width=.12pt] (-1.618570,1.304246)--(-1.734580,1.391274)--(-1.865758,1.302838)--(-1.740975,1.221724)--cycle;
\path[fill=blue!16!white,draw=blue!16!white,line width=.12pt] (1.343133,0.583889)--(1.380985,0.726503)--(1.283890,0.791962)--(1.248699,0.647553)--cycle;
\path[fill=blue!23!white,draw=blue!23!white,line width=.12pt] (-1.248699,0.647553)--(-1.283890,0.791962)--(-1.380985,0.726503)--(-1.343133,0.583889)--cycle;
\path[fill=blue!32!white,draw=blue!32!white,line width=.12pt] (-1.464780,0.126340)--(-1.424631,0.245578)--(-1.492409,0.173925)--(-1.534468,0.052668)--cycle;
\path[fill=blue!24!white,draw=blue!24!white,line width=.12pt] (1.534468,0.052668)--(1.492409,0.173925)--(1.424631,0.245578)--(1.464780,0.126340)--cycle;
\path[fill=blue!19!white,draw=blue!19!white,line width=.12pt] (-2.831853,0.971406)--(-2.900154,0.865772)--(-3.003935,0.720748)--(-2.933190,0.829797)--cycle;
\path[fill=blue!27!white,draw=blue!27!white,line width=.12pt] (2.933190,0.829797)--(3.003935,0.720748)--(2.900154,0.865772)--(2.831853,0.971406)--cycle;
\path[fill=blue!38!white,draw=blue!38!white,line width=.12pt] (-2.654149,-0.376111)--(-2.506958,-0.430175)--(-2.543764,-0.554239)--(-2.693116,-0.507458)--cycle;
\path[fill=blue!41!white,draw=blue!41!white,line width=.12pt] (2.693116,-0.507458)--(2.543764,-0.554239)--(2.506958,-0.430175)--(2.654149,-0.376111)--cycle;
\path[fill=blue!39!white,draw=blue!39!white,line width=.12pt] (-2.506958,-0.430175)--(-2.353885,-0.458712)--(-2.388443,-0.575200)--(-2.543764,-0.554239)--cycle;
\path[fill=blue!40!white,draw=blue!40!white,line width=.12pt] (2.543764,-0.554239)--(2.388443,-0.575200)--(2.353885,-0.458712)--(2.506958,-0.430175)--cycle;
\path[fill=blue!26!white,draw=blue!26!white,line width=.12pt] (-3.062209,0.464455)--(-3.047499,0.327062)--(-3.123436,0.175483)--(-3.138513,0.312145)--cycle;
\path[fill=blue!34!white,draw=blue!34!white,line width=.12pt] (3.138513,0.312145)--(3.123436,0.175483)--(3.047499,0.327062)--(3.062209,0.464455)--cycle;
\path[fill=blue!38!white,draw=blue!38!white,line width=.12pt] (-1.871339,-0.292227)--(-1.755330,-0.223661)--(-1.799069,-0.310969)--(-1.917969,-0.385305)--cycle;
\path[fill=blue!33!white,draw=blue!33!white,line width=.12pt] (1.917969,-0.385305)--(1.799069,-0.310969)--(1.755330,-0.223661)--(1.871339,-0.292227)--cycle;
\path[fill=blue!14!white,draw=blue!14!white,line width=.12pt] (-1.995527,1.359505)--(-2.125296,1.389547)--(-2.254254,1.281970)--(-2.116611,1.258496)--cycle;
\path[fill=blue!15!white,draw=blue!15!white,line width=.12pt] (2.116611,1.258496)--(2.254254,1.281970)--(2.125296,1.389547)--(1.995527,1.359505)--cycle;
\path[fill=blue!36!white,draw=blue!36!white,line width=.12pt] (-2.789801,-0.298596)--(-2.654149,-0.376111)--(-2.693116,-0.507458)--(-2.830759,-0.436656)--cycle;
\path[fill=blue!41!white,draw=blue!41!white,line width=.12pt] (2.830759,-0.436656)--(2.693116,-0.507458)--(2.654149,-0.376111)--(2.789801,-0.298596)--cycle;
\path[fill=blue!12!white,draw=blue!12!white,line width=.12pt] (1.625975,1.119280)--(1.740975,1.221724)--(1.618570,1.304246)--(1.511656,1.196351)--cycle;
\path[fill=blue!16!white,draw=blue!16!white,line width=.12pt] (-1.511656,1.196351)--(-1.618570,1.304246)--(-1.740975,1.221724)--(-1.625975,1.119280)--cycle;
\path[fill=blue!40!white,draw=blue!40!white,line width=.12pt] (-2.353885,-0.458712)--(-2.200811,-0.460624)--(-2.233123,-0.569536)--(-2.388443,-0.575200)--cycle;
\path[fill=blue!39!white,draw=blue!39!white,line width=.12pt] (2.388443,-0.575200)--(2.233123,-0.569536)--(2.200811,-0.460624)--(2.353885,-0.458712)--cycle;
\path[fill=blue!14!white,draw=blue!14!white,line width=.12pt] (1.380985,0.726503)--(1.442454,0.866819)--(1.341037,0.935191)--(1.283890,0.791962)--cycle;
\path[fill=blue!21!white,draw=blue!21!white,line width=.12pt] (-1.283890,0.791962)--(-1.341037,0.935191)--(-1.442454,0.866819)--(-1.380985,0.726503)--cycle;
\path[fill=blue!35!white,draw=blue!35!white,line width=.12pt] (-1.602769,-0.052966)--(-1.534468,0.052668)--(-1.589379,-0.024064)--(-1.660123,-0.133114)--cycle;
\path[fill=blue!27!white,draw=blue!27!white,line width=.12pt] (1.660123,-0.133114)--(1.589379,-0.024064)--(1.534468,0.052668)--(1.602769,-0.052966)--cycle;
\path[fill=blue!35!white,draw=blue!35!white,line width=.12pt] (-2.908701,-0.200610)--(-2.789801,-0.298596)--(-2.830759,-0.436656)--(-2.951405,-0.344554)--cycle;
\path[fill=blue!40!white,draw=blue!40!white,line width=.12pt] (2.951405,-0.344554)--(2.830759,-0.436656)--(2.789801,-0.298596)--(2.908701,-0.200610)--cycle;
\path[fill=blue!26!white,draw=blue!26!white,line width=.12pt] (-1.330351,0.444456)--(-1.343133,0.583889)--(-1.424631,0.515903)--(-1.411074,0.377117)--cycle;
\path[fill=blue!18!white,draw=blue!18!white,line width=.12pt] (1.411074,0.377117)--(1.424631,0.515903)--(1.343133,0.583889)--(1.330351,0.444456)--cycle;
\path[fill=blue!15!white,draw=blue!15!white,line width=.12pt] (-2.386608,1.277917)--(-2.508586,1.246493)--(-2.627935,1.120322)--(-2.500154,1.157881)--cycle;
\path[fill=blue!19!white,draw=blue!19!white,line width=.12pt] (2.500154,1.157881)--(2.627935,1.120322)--(2.508586,1.246493)--(2.386608,1.277917)--cycle;
\path[fill=blue!12!white,draw=blue!12!white,line width=.12pt] (1.525177,0.999443)--(1.625975,1.119280)--(1.511656,1.196351)--(1.417944,1.071736)--cycle;
\path[fill=blue!17!white,draw=blue!17!white,line width=.12pt] (-1.417944,1.071736)--(-1.511656,1.196351)--(-1.625975,1.119280)--(-1.525177,0.999443)--cycle;
\path[fill=blue!19!white,draw=blue!19!white,line width=.12pt] (-1.341037,0.935191)--(-1.417944,1.071736)--(-1.525177,0.999443)--(-1.442454,0.866819)--cycle;
\path[fill=blue!13!white,draw=blue!13!white,line width=.12pt] (1.442454,0.866819)--(1.525177,0.999443)--(1.417944,1.071736)--(1.341037,0.935191)--cycle;
\path[fill=blue!40!white,draw=blue!40!white,line width=.12pt] (-2.200811,-0.460624)--(-2.053620,-0.435837)--(-2.083771,-0.537466)--(-2.233123,-0.569536)--cycle;
\path[fill=blue!37!white,draw=blue!37!white,line width=.12pt] (2.233123,-0.569536)--(2.083771,-0.537466)--(2.053620,-0.435837)--(2.200811,-0.460624)--cycle;
\path[fill=blue!14!white,draw=blue!14!white,line width=.12pt] (-1.865758,1.302838)--(-1.995527,1.359505)--(-2.116611,1.258496)--(-1.978968,1.208398)--cycle;
\path[fill=blue!14!white,draw=blue!14!white,line width=.12pt] (1.978968,1.208398)--(2.116611,1.258496)--(1.995527,1.359505)--(1.865758,1.302838)--cycle;
\path[fill=blue!32!white,draw=blue!32!white,line width=.12pt] (-3.006279,-0.085917)--(-2.908701,-0.200610)--(-2.951405,-0.344554)--(-3.050416,-0.234690)--cycle;
\path[fill=blue!39!white,draw=blue!39!white,line width=.12pt] (3.050416,-0.234690)--(2.951405,-0.344554)--(2.908701,-0.200610)--(3.006279,-0.085917)--cycle;
\path[fill=blue!29!white,draw=blue!29!white,line width=.12pt] (-1.424631,0.245578)--(-1.411074,0.377117)--(-1.478207,0.306147)--(-1.492409,0.173925)--cycle;
\path[fill=blue!21!white,draw=blue!21!white,line width=.12pt] (1.492409,0.173925)--(1.478207,0.306147)--(1.411074,0.377117)--(1.424631,0.245578)--cycle;
\path[fill=blue!23!white,draw=blue!23!white,line width=.12pt] (-3.047499,0.597387)--(-3.062209,0.464455)--(-3.138513,0.312145)--(-3.123436,0.445808)--cycle;
\path[fill=blue!32!white,draw=blue!32!white,line width=.12pt] (3.123436,0.445808)--(3.138513,0.312145)--(3.062209,0.464455)--(3.047499,0.597387)--cycle;
\path[fill=blue!37!white,draw=blue!37!white,line width=.12pt] (-1.755330,-0.223661)--(-1.660123,-0.133114)--(-1.701490,-0.215687)--(-1.799069,-0.310969)--cycle;
\path[fill=blue!30!white,draw=blue!30!white,line width=.12pt] (1.799069,-0.310969)--(1.701490,-0.215687)--(1.660123,-0.133114)--(1.755330,-0.223661)--cycle;
\path[fill=blue!17!white,draw=blue!17!white,line width=.12pt] (-2.739936,1.057357)--(-2.831853,0.971406)--(-2.933190,0.829797)--(-2.837984,0.920344)--cycle;
\path[fill=blue!24!white,draw=blue!24!white,line width=.12pt] (2.837984,0.920344)--(2.933190,0.829797)--(2.831853,0.971406)--(2.739936,1.057357)--cycle;
\path[fill=blue!24!white,draw=blue!24!white,line width=.12pt] (-1.343133,0.583889)--(-1.380985,0.726503)--(-1.464780,0.656601)--(-1.424631,0.515903)--cycle;
\path[fill=blue!16!white,draw=blue!16!white,line width=.12pt] (1.424631,0.515903)--(1.464780,0.656601)--(1.380985,0.726503)--(1.343133,0.583889)--cycle;
\path[fill=blue!40!white,draw=blue!40!white,line width=.12pt] (-2.053620,-0.435837)--(-1.917969,-0.385305)--(-1.946127,-0.480221)--(-2.083771,-0.537466)--cycle;
\path[fill=blue!36!white,draw=blue!36!white,line width=.12pt] (2.083771,-0.537466)--(1.946127,-0.480221)--(1.917969,-0.385305)--(2.053620,-0.435837)--cycle;
\path[fill=blue!24!white,draw=blue!24!white,line width=.12pt] (1.589379,-0.024064)--(1.545814,0.099296)--(1.492409,0.173925)--(1.534468,0.052668)--cycle;
\path[fill=blue!33!white,draw=blue!33!white,line width=.12pt] (-1.534468,0.052668)--(-1.492409,0.173925)--(-1.545814,0.099296)--(-1.589379,-0.024064)--cycle;
\path[fill=blue!15!white,draw=blue!15!white,line width=.12pt] (-2.254254,1.281970)--(-2.386608,1.277917)--(-2.500154,1.157881)--(-2.361503,1.168591)--cycle;
\path[fill=blue!17!white,draw=blue!17!white,line width=.12pt] (2.361503,1.168591)--(2.500154,1.157881)--(2.386608,1.277917)--(2.254254,1.281970)--cycle;
\path[fill=blue!15!white,draw=blue!15!white,line width=.12pt] (-1.740975,1.221724)--(-1.865758,1.302838)--(-1.978968,1.208398)--(-1.846614,1.133600)--cycle;
\path[fill=blue!13!white,draw=blue!13!white,line width=.12pt] (1.846614,1.133600)--(1.978968,1.208398)--(1.865758,1.302838)--(1.740975,1.221724)--cycle;
\path[fill=blue!30!white,draw=blue!30!white,line width=.12pt] (-3.078787,0.041074)--(-3.006279,-0.085917)--(-3.050416,-0.234690)--(-3.123988,-0.111288)--cycle;
\path[fill=blue!38!white,draw=blue!38!white,line width=.12pt] (3.123988,-0.111288)--(3.050416,-0.234690)--(3.006279,-0.085917)--(3.078787,0.041074)--cycle;
\path[fill=blue!15!white,draw=blue!15!white,line width=.12pt] (1.464780,0.656601)--(1.529978,0.793806)--(1.442454,0.866819)--(1.380985,0.726503)--cycle;
\path[fill=blue!21!white,draw=blue!21!white,line width=.12pt] (-1.380985,0.726503)--(-1.442454,0.866819)--(-1.529978,0.793806)--(-1.464780,0.656601)--cycle;
\path[fill=blue!27!white,draw=blue!27!white,line width=.12pt] (-1.411074,0.377117)--(-1.424631,0.515903)--(-1.492409,0.444251)--(-1.478207,0.306147)--cycle;
\path[fill=blue!19!white,draw=blue!19!white,line width=.12pt] (1.478207,0.306147)--(1.492409,0.444251)--(1.424631,0.515903)--(1.411074,0.377117)--cycle;
\path[fill=blue!12!white,draw=blue!12!white,line width=.12pt] (1.724636,1.036978)--(1.846614,1.133600)--(1.740975,1.221724)--(1.625975,1.119280)--cycle;
\path[fill=blue!16!white,draw=blue!16!white,line width=.12pt] (-1.625975,1.119280)--(-1.740975,1.221724)--(-1.846614,1.133600)--(-1.724636,1.036978)--cycle;
\path[fill=blue!21!white,draw=blue!21!white,line width=.12pt] (-3.003935,0.720748)--(-3.047499,0.597387)--(-3.123436,0.445808)--(-3.078787,0.571335)--cycle;
\path[fill=blue!29!white,draw=blue!29!white,line width=.12pt] (3.078787,0.571335)--(3.123436,0.445808)--(3.047499,0.597387)--(3.003935,0.720748)--cycle;
\path[fill=blue!39!white,draw=blue!39!white,line width=.12pt] (-1.917969,-0.385305)--(-1.799069,-0.310969)--(-1.825482,-0.400001)--(-1.946127,-0.480221)--cycle;
\path[fill=blue!33!white,draw=blue!33!white,line width=.12pt] (1.946127,-0.480221)--(1.825482,-0.400001)--(1.799069,-0.310969)--(1.917969,-0.385305)--cycle;
\path[fill=blue!35!white,draw=blue!35!white,line width=.12pt] (-1.660123,-0.133114)--(-1.589379,-0.024064)--(-1.628983,-0.103118)--(-1.701490,-0.215687)--cycle;
\path[fill=blue!28!white,draw=blue!28!white,line width=.12pt] (1.701490,-0.215687)--(1.628983,-0.103118)--(1.589379,-0.024064)--(1.660123,-0.133114)--cycle;
\path[fill=blue!19!white,draw=blue!19!white,line width=.12pt] (-1.442454,0.866819)--(-1.525177,0.999443)--(-1.617721,0.922243)--(-1.529978,0.793806)--cycle;
\path[fill=blue!13!white,draw=blue!13!white,line width=.12pt] (1.529978,0.793806)--(1.617721,0.922243)--(1.525177,0.999443)--(1.442454,0.866819)--cycle;
\path[fill=blue!21!white,draw=blue!21!white,line width=.12pt] (2.721975,0.988910)--(2.837984,0.920344)--(2.739936,1.057357)--(2.627935,1.120322)--cycle;
\path[fill=blue!16!white,draw=blue!16!white,line width=.12pt] (-2.627935,1.120322)--(-2.739936,1.057357)--(-2.837984,0.920344)--(-2.721975,0.988910)--cycle;
\path[fill=blue!17!white,draw=blue!17!white,line width=.12pt] (-1.525177,0.999443)--(-1.625975,1.119280)--(-1.724636,1.036978)--(-1.617721,0.922243)--cycle;
\path[fill=blue!13!white,draw=blue!13!white,line width=.12pt] (1.617721,0.922243)--(1.724636,1.036978)--(1.625975,1.119280)--(1.525177,0.999443)--cycle;
\path[fill=blue!14!white,draw=blue!14!white,line width=.12pt] (-2.116611,1.258496)--(-2.254254,1.281970)--(-2.361503,1.168591)--(-2.217311,1.152040)--cycle;
\path[fill=blue!15!white,draw=blue!15!white,line width=.12pt] (2.217311,1.152040)--(2.361503,1.168591)--(2.254254,1.281970)--(2.116611,1.258496)--cycle;
\path[fill=blue!27!white,draw=blue!27!white,line width=.12pt] (-3.123436,0.175483)--(-3.078787,0.041074)--(-3.123988,-0.111288)--(-3.169293,0.020912)--cycle;
\path[fill=blue!36!white,draw=blue!36!white,line width=.12pt] (3.169293,0.020912)--(3.123988,-0.111288)--(3.078787,0.041074)--(3.123436,0.175483)--cycle;
\path[fill=blue!30!white,draw=blue!30!white,line width=.12pt] (-1.492409,0.173925)--(-1.478207,0.306147)--(-1.531105,0.232228)--(-1.545814,0.099296)--cycle;
\path[fill=blue!22!white,draw=blue!22!white,line width=.12pt] (1.545814,0.099296)--(1.531105,0.232228)--(1.478207,0.306147)--(1.492409,0.173925)--cycle;
\path[fill=blue!39!white,draw=blue!39!white,line width=.12pt] (-2.543764,-0.554239)--(-2.388443,-0.575200)--(-2.400000,-0.692820)--(-2.556072,-0.679508)--cycle;
\path[fill=blue!40!white,draw=blue!40!white,line width=.12pt] (2.556072,-0.679508)--(2.400000,-0.692820)--(2.388443,-0.575200)--(2.543764,-0.554239)--cycle;
\path[fill=blue!40!white,draw=blue!40!white,line width=.12pt] (-2.388443,-0.575200)--(-2.233123,-0.569536)--(-2.243928,-0.679508)--(-2.400000,-0.692820)--cycle;
\path[fill=blue!39!white,draw=blue!39!white,line width=.12pt] (2.400000,-0.692820)--(2.243928,-0.679508)--(2.233123,-0.569536)--(2.388443,-0.575200)--cycle;
\path[fill=blue!38!white,draw=blue!38!white,line width=.12pt] (-2.693116,-0.507458)--(-2.543764,-0.554239)--(-2.556072,-0.679508)--(-2.706147,-0.640083)--cycle;
\path[fill=blue!40!white,draw=blue!40!white,line width=.12pt] (2.706147,-0.640083)--(2.556072,-0.679508)--(2.543764,-0.554239)--(2.693116,-0.507458)--cycle;
\path[fill=blue!24!white,draw=blue!24!white,line width=.12pt] (-1.424631,0.515903)--(-1.464780,0.656601)--(-1.534468,0.582930)--(-1.492409,0.444251)--cycle;
\path[fill=blue!17!white,draw=blue!17!white,line width=.12pt] (1.492409,0.444251)--(1.534468,0.582930)--(1.464780,0.656601)--(1.424631,0.515903)--cycle;
\path[fill=blue!38!white,draw=blue!38!white,line width=.12pt] (-1.799069,-0.310969)--(-1.701490,-0.215687)--(-1.726471,-0.299889)--(-1.825482,-0.400001)--cycle;
\path[fill=blue!31!white,draw=blue!31!white,line width=.12pt] (1.825482,-0.400001)--(1.726471,-0.299889)--(1.701490,-0.215687)--(1.799069,-0.310969)--cycle;
\path[fill=blue!40!white,draw=blue!40!white,line width=.12pt] (-2.233123,-0.569536)--(-2.083771,-0.537466)--(-2.093853,-0.640083)--(-2.243928,-0.679508)--cycle;
\path[fill=blue!38!white,draw=blue!38!white,line width=.12pt] (2.243928,-0.679508)--(2.093853,-0.640083)--(2.083771,-0.537466)--(2.233123,-0.569536)--cycle;
\path[fill=blue!36!white,draw=blue!36!white,line width=.12pt] (-2.830759,-0.436656)--(-2.693116,-0.507458)--(-2.706147,-0.640083)--(-2.844456,-0.576059)--cycle;
\path[fill=blue!40!white,draw=blue!40!white,line width=.12pt] (2.844456,-0.576059)--(2.706147,-0.640083)--(2.693116,-0.507458)--(2.830759,-0.436656)--cycle;
\path[fill=blue!19!white,draw=blue!19!white,line width=.12pt] (-2.933190,0.829797)--(-3.003935,0.720748)--(-3.078787,0.571335)--(-3.006279,0.683903)--cycle;
\path[fill=blue!26!white,draw=blue!26!white,line width=.12pt] (3.006279,0.683903)--(3.078787,0.571335)--(3.003935,0.720748)--(2.933190,0.829797)--cycle;
\path[fill=blue!15!white,draw=blue!15!white,line width=.12pt] (-1.978968,1.208398)--(-2.116611,1.258496)--(-2.217311,1.152040)--(-2.073119,1.108865)--cycle;
\path[fill=blue!14!white,draw=blue!14!white,line width=.12pt] (2.073119,1.108865)--(2.217311,1.152040)--(2.116611,1.258496)--(1.978968,1.208398)--cycle;
\path[fill=blue!33!white,draw=blue!33!white,line width=.12pt] (-1.589379,-0.024064)--(-1.545814,0.099296)--(-1.584333,0.022409)--(-1.628983,-0.103118)--cycle;
\path[fill=blue!25!white,draw=blue!25!white,line width=.12pt] (1.628983,-0.103118)--(1.584333,0.022409)--(1.545814,0.099296)--(1.589379,-0.024064)--cycle;
\path[fill=blue!15!white,draw=blue!15!white,line width=.12pt] (-2.500154,1.157881)--(-2.627935,1.120322)--(-2.721975,0.988910)--(-2.589621,1.032859)--cycle;
\path[fill=blue!19!white,draw=blue!19!white,line width=.12pt] (2.589621,1.032859)--(2.721975,0.988910)--(2.627935,1.120322)--(2.500154,1.157881)--cycle;
\path[fill=blue!25!white,draw=blue!25!white,line width=.12pt] (-3.138513,0.312145)--(-3.123436,0.175483)--(-3.169293,0.020912)--(-3.184591,0.156827)--cycle;
\path[fill=blue!34!white,draw=blue!34!white,line width=.12pt] (3.184591,0.156827)--(3.169293,0.020912)--(3.123436,0.175483)--(3.138513,0.312145)--cycle;
\path[fill=blue!22!white,draw=blue!22!white,line width=.12pt] (-1.464780,0.656601)--(-1.529978,0.793806)--(-1.602769,0.716855)--(-1.534468,0.582930)--cycle;
\path[fill=blue!15!white,draw=blue!15!white,line width=.12pt] (1.534468,0.582930)--(1.602769,0.716855)--(1.529978,0.793806)--(1.464780,0.656601)--cycle;
\path[fill=blue!40!white,draw=blue!40!white,line width=.12pt] (-2.083771,-0.537466)--(-1.946127,-0.480221)--(-1.955544,-0.576059)--(-2.093853,-0.640083)--cycle;
\path[fill=blue!36!white,draw=blue!36!white,line width=.12pt] (2.093853,-0.640083)--(1.955544,-0.576059)--(1.946127,-0.480221)--(2.083771,-0.537466)--cycle;
\path[fill=blue!34!white,draw=blue!34!white,line width=.12pt] (-2.951405,-0.344554)--(-2.830759,-0.436656)--(-2.844456,-0.576059)--(-2.965685,-0.489898)--cycle;
\path[fill=blue!40!white,draw=blue!40!white,line width=.12pt] (2.965685,-0.489898)--(2.844456,-0.576059)--(2.830759,-0.436656)--(2.951405,-0.344554)--cycle;
\path[fill=blue!28!white,draw=blue!28!white,line width=.12pt] (-1.478207,0.306147)--(-1.492409,0.444251)--(-1.545814,0.369621)--(-1.531105,0.232228)--cycle;
\path[fill=blue!19!white,draw=blue!19!white,line width=.12pt] (1.531105,0.232228)--(1.545814,0.369621)--(1.492409,0.444251)--(1.478207,0.306147)--cycle;
\path[fill=blue!15!white,draw=blue!15!white,line width=.12pt] (-1.846614,1.133600)--(-1.978968,1.208398)--(-2.073119,1.108865)--(-1.934468,1.040724)--cycle;
\path[fill=blue!13!white,draw=blue!13!white,line width=.12pt] (1.934468,1.040724)--(2.073119,1.108865)--(1.978968,1.208398)--(1.846614,1.133600)--cycle;
\path[fill=blue!20!white,draw=blue!20!white,line width=.12pt] (-1.529978,0.793806)--(-1.617721,0.922243)--(-1.694686,0.840879)--(-1.602769,0.716855)--cycle;
\path[fill=blue!14!white,draw=blue!14!white,line width=.12pt] (1.602769,0.716855)--(1.694686,0.840879)--(1.617721,0.922243)--(1.529978,0.793806)--cycle;
\path[fill=blue!36!white,draw=blue!36!white,line width=.12pt] (-1.701490,-0.215687)--(-1.628983,-0.103118)--(-1.652899,-0.183733)--(-1.726471,-0.299889)--cycle;
\path[fill=blue!28!white,draw=blue!28!white,line width=.12pt] (1.726471,-0.299889)--(1.652899,-0.183733)--(1.628983,-0.103118)--(1.701490,-0.215687)--cycle;
\path[fill=blue!16!white,draw=blue!16!white,line width=.12pt] (-1.724636,1.036978)--(-1.846614,1.133600)--(-1.934468,1.040724)--(-1.806687,0.950236)--cycle;
\path[fill=blue!13!white,draw=blue!13!white,line width=.12pt] (1.806687,0.950236)--(1.934468,1.040724)--(1.846614,1.133600)--(1.724636,1.036978)--cycle;
\path[fill=blue!18!white,draw=blue!18!white,line width=.12pt] (-1.617721,0.922243)--(-1.724636,1.036978)--(-1.806687,0.950236)--(-1.694686,0.840879)--cycle;
\path[fill=blue!13!white,draw=blue!13!white,line width=.12pt] (1.694686,0.840879)--(1.806687,0.950236)--(1.724636,1.036978)--(1.617721,0.922243)--cycle;
\path[fill=blue!39!white,draw=blue!39!white,line width=.12pt] (-1.946127,-0.480221)--(-1.825482,-0.400001)--(-1.834315,-0.489898)--(-1.955544,-0.576059)--cycle;
\path[fill=blue!34!white,draw=blue!34!white,line width=.12pt] (1.955544,-0.576059)--(1.834315,-0.489898)--(1.825482,-0.400001)--(1.946127,-0.480221)--cycle;
\path[fill=blue!14!white,draw=blue!14!white,line width=.12pt] (-2.361503,1.168591)--(-2.500154,1.157881)--(-2.589621,1.032859)--(-2.446009,1.050502)--cycle;
\path[fill=blue!17!white,draw=blue!17!white,line width=.12pt] (2.446009,1.050502)--(2.589621,1.032859)--(2.500154,1.157881)--(2.361503,1.168591)--cycle;
\path[fill=blue!17!white,draw=blue!17!white,line width=.12pt] (-2.837984,0.920344)--(-2.933190,0.829797)--(-3.006279,0.683903)--(-2.908701,0.779186)--cycle;
\path[fill=blue!24!white,draw=blue!24!white,line width=.12pt] (2.908701,0.779186)--(3.006279,0.683903)--(2.933190,0.829797)--(2.837984,0.920344)--cycle;
\path[fill=blue!32!white,draw=blue!32!white,line width=.12pt] (-3.050416,-0.234690)--(-2.951405,-0.344554)--(-2.965685,-0.489898)--(-3.065176,-0.384910)--cycle;
\path[fill=blue!39!white,draw=blue!39!white,line width=.12pt] (3.065176,-0.384910)--(2.965685,-0.489898)--(2.951405,-0.344554)--(3.050416,-0.234690)--cycle;
\path[fill=blue!31!white,draw=blue!31!white,line width=.12pt] (-1.545814,0.099296)--(-1.531105,0.232228)--(-1.569256,0.156072)--(-1.584333,0.022409)--cycle;
\path[fill=blue!22!white,draw=blue!22!white,line width=.12pt] (1.584333,0.022409)--(1.569256,0.156072)--(1.531105,0.232228)--(1.545814,0.099296)--cycle;
\path[fill=blue!22!white,draw=blue!22!white,line width=.12pt] (-3.123436,0.445808)--(-3.138513,0.312145)--(-3.184591,0.156827)--(-3.169293,0.291237)--cycle;
\path[fill=blue!31!white,draw=blue!31!white,line width=.12pt] (3.169293,0.291237)--(3.184591,0.156827)--(3.138513,0.312145)--(3.123436,0.445808)--cycle;
\path[fill=blue!25!white,draw=blue!25!white,line width=.12pt] (-1.492409,0.444251)--(-1.534468,0.582930)--(-1.589379,0.506197)--(-1.545814,0.369621)--cycle;
\path[fill=blue!17!white,draw=blue!17!white,line width=.12pt] (1.545814,0.369621)--(1.589379,0.506197)--(1.534468,0.582930)--(1.492409,0.444251)--cycle;
\path[fill=blue!38!white,draw=blue!38!white,line width=.12pt] (-1.825482,-0.400001)--(-1.726471,-0.299889)--(-1.734824,-0.384910)--(-1.834315,-0.489898)--cycle;
\path[fill=blue!31!white,draw=blue!31!white,line width=.12pt] (1.834315,-0.489898)--(1.734824,-0.384910)--(1.726471,-0.299889)--(1.825482,-0.400001)--cycle;
\path[fill=blue!14!white,draw=blue!14!white,line width=.12pt] (-2.217311,1.152040)--(-2.361503,1.168591)--(-2.446009,1.050502)--(-2.296657,1.041162)--cycle;
\path[fill=blue!15!white,draw=blue!15!white,line width=.12pt] (2.296657,1.041162)--(2.446009,1.050502)--(2.361503,1.168591)--(2.217311,1.152040)--cycle;
\path[fill=blue!34!white,draw=blue!34!white,line width=.12pt] (-1.628983,-0.103118)--(-1.584333,0.022409)--(-1.607593,-0.055995)--(-1.652899,-0.183733)--cycle;
\path[fill=blue!25!white,draw=blue!25!white,line width=.12pt] (1.652899,-0.183733)--(1.607593,-0.055995)--(1.584333,0.022409)--(1.628983,-0.103118)--cycle;
\path[fill=blue!29!white,draw=blue!29!white,line width=.12pt] (-3.123988,-0.111288)--(-3.050416,-0.234690)--(-3.065176,-0.384910)--(-3.139104,-0.265131)--cycle;
\path[fill=blue!37!white,draw=blue!37!white,line width=.12pt] (3.139104,-0.265131)--(3.065176,-0.384910)--(3.050416,-0.234690)--(3.123988,-0.111288)--cycle;
\path[fill=blue!23!white,draw=blue!23!white,line width=.12pt] (-1.534468,0.582930)--(-1.602769,0.716855)--(-1.660123,0.636707)--(-1.589379,0.506197)--cycle;
\path[fill=blue!15!white,draw=blue!15!white,line width=.12pt] (1.589379,0.506197)--(1.660123,0.636707)--(1.602769,0.716855)--(1.534468,0.582930)--cycle;
\path[fill=blue!16!white,draw=blue!16!white,line width=.12pt] (-2.721975,0.988910)--(-2.837984,0.920344)--(-2.908701,0.779186)--(-2.789801,0.853522)--cycle;
\path[fill=blue!21!white,draw=blue!21!white,line width=.12pt] (2.789801,0.853522)--(2.908701,0.779186)--(2.837984,0.920344)--(2.721975,0.988910)--cycle;
\path[fill=blue!28!white,draw=blue!28!white,line width=.12pt] (-1.531105,0.232228)--(-1.545814,0.369621)--(-1.584333,0.292734)--(-1.569256,0.156072)--cycle;
\path[fill=blue!20!white,draw=blue!20!white,line width=.12pt] (1.569256,0.156072)--(1.584333,0.292734)--(1.545814,0.369621)--(1.531105,0.232228)--cycle;
\path[fill=blue!40!white,draw=blue!40!white,line width=.12pt] (-2.400000,-0.692820)--(-2.243928,-0.679508)--(-2.233123,-0.789480)--(-2.388443,-0.810441)--cycle;
\path[fill=blue!39!white,draw=blue!39!white,line width=.12pt] (2.388443,-0.810441)--(2.233123,-0.789480)--(2.243928,-0.679508)--(2.400000,-0.692820)--cycle;
\path[fill=blue!39!white,draw=blue!39!white,line width=.12pt] (-2.556072,-0.679508)--(-2.400000,-0.692820)--(-2.388443,-0.810441)--(-2.543764,-0.804777)--cycle;
\path[fill=blue!40!white,draw=blue!40!white,line width=.12pt] (2.543764,-0.804777)--(2.388443,-0.810441)--(2.400000,-0.692820)--(2.556072,-0.679508)--cycle;
\path[fill=blue!40!white,draw=blue!40!white,line width=.12pt] (-2.243928,-0.679508)--(-2.093853,-0.640083)--(-2.083771,-0.742699)--(-2.233123,-0.789480)--cycle;
\path[fill=blue!38!white,draw=blue!38!white,line width=.12pt] (2.233123,-0.789480)--(2.083771,-0.742699)--(2.093853,-0.640083)--(2.243928,-0.679508)--cycle;
\path[fill=blue!20!white,draw=blue!20!white,line width=.12pt] (-3.078787,0.571335)--(-3.123436,0.445808)--(-3.169293,0.291237)--(-3.123988,0.418974)--cycle;
\path[fill=blue!29!white,draw=blue!29!white,line width=.12pt] (3.123988,0.418974)--(3.169293,0.291237)--(3.123436,0.445808)--(3.078787,0.571335)--cycle;
\path[fill=blue!15!white,draw=blue!15!white,line width=.12pt] (-2.073119,1.108865)--(-2.217311,1.152040)--(-2.296657,1.041162)--(-2.147305,1.005197)--cycle;
\path[fill=blue!14!white,draw=blue!14!white,line width=.12pt] (2.147305,1.005197)--(2.296657,1.041162)--(2.217311,1.152040)--(2.073119,1.108865)--cycle;
\path[fill=blue!38!white,draw=blue!38!white,line width=.12pt] (-2.706147,-0.640083)--(-2.556072,-0.679508)--(-2.543764,-0.804777)--(-2.693116,-0.772707)--cycle;
\path[fill=blue!40!white,draw=blue!40!white,line width=.12pt] (2.693116,-0.772707)--(2.543764,-0.804777)--(2.556072,-0.679508)--(2.706147,-0.640083)--cycle;
\path[fill=blue!21!white,draw=blue!21!white,line width=.12pt] (-1.602769,0.716855)--(-1.694686,0.840879)--(-1.755330,0.756135)--(-1.660123,0.636707)--cycle;
\path[fill=blue!14!white,draw=blue!14!white,line width=.12pt] (1.660123,0.636707)--(1.755330,0.756135)--(1.694686,0.840879)--(1.602769,0.716855)--cycle;
\path[fill=blue!40!white,draw=blue!40!white,line width=.12pt] (-2.093853,-0.640083)--(-1.955544,-0.576059)--(-1.946127,-0.671897)--(-2.083771,-0.742699)--cycle;
\path[fill=blue!36!white,draw=blue!36!white,line width=.12pt] (2.083771,-0.742699)--(1.946127,-0.671897)--(1.955544,-0.576059)--(2.093853,-0.640083)--cycle;
\path[fill=blue!37!white,draw=blue!37!white,line width=.12pt] (-1.726471,-0.299889)--(-1.652899,-0.183733)--(-1.660896,-0.265131)--(-1.734824,-0.384910)--cycle;
\path[fill=blue!29!white,draw=blue!29!white,line width=.12pt] (1.734824,-0.384910)--(1.660896,-0.265131)--(1.652899,-0.183733)--(1.726471,-0.299889)--cycle;
\path[fill=blue!16!white,draw=blue!16!white,line width=.12pt] (-1.934468,1.040724)--(-2.073119,1.108865)--(-2.147305,1.005197)--(-2.003693,0.943989)--cycle;
\path[fill=blue!13!white,draw=blue!13!white,line width=.12pt] (2.003693,0.943989)--(2.147305,1.005197)--(2.073119,1.108865)--(1.934468,1.040724)--cycle;
\path[fill=blue!27!white,draw=blue!27!white,line width=.12pt] (-3.169293,0.020912)--(-3.123988,-0.111288)--(-3.139104,-0.265131)--(-3.184628,-0.135163)--cycle;
\path[fill=blue!35!white,draw=blue!35!white,line width=.12pt] (3.184628,-0.135163)--(3.139104,-0.265131)--(3.123988,-0.111288)--(3.169293,0.020912)--cycle;
\path[fill=blue!18!white,draw=blue!18!white,line width=.12pt] (-1.694686,0.840879)--(-1.806687,0.950236)--(-1.871339,0.859891)--(-1.755330,0.756135)--cycle;
\path[fill=blue!13!white,draw=blue!13!white,line width=.12pt] (1.755330,0.756135)--(1.871339,0.859891)--(1.806687,0.950236)--(1.694686,0.840879)--cycle;
\path[fill=blue!32!white,draw=blue!32!white,line width=.12pt] (-1.584333,0.022409)--(-1.569256,0.156072)--(-1.592296,0.078414)--(-1.607593,-0.055995)--cycle;
\path[fill=blue!23!white,draw=blue!23!white,line width=.12pt] (1.607593,-0.055995)--(1.592296,0.078414)--(1.569256,0.156072)--(1.584333,0.022409)--cycle;
\path[fill=blue!17!white,draw=blue!17!white,line width=.12pt] (-1.806687,0.950236)--(-1.934468,1.040724)--(-2.003693,0.943989)--(-1.871339,0.859891)--cycle;
\path[fill=blue!13!white,draw=blue!13!white,line width=.12pt] (1.871339,0.859891)--(2.003693,0.943989)--(1.934468,1.040724)--(1.806687,0.950236)--cycle;
\path[fill=blue!36!white,draw=blue!36!white,line width=.12pt] (-2.844456,-0.576059)--(-2.706147,-0.640083)--(-2.693116,-0.772707)--(-2.830759,-0.715462)--cycle;
\path[fill=blue!40!white,draw=blue!40!white,line width=.12pt] (2.830759,-0.715462)--(2.693116,-0.772707)--(2.706147,-0.640083)--(2.844456,-0.576059)--cycle;
\path[fill=blue!15!white,draw=blue!15!white,line width=.12pt] (-2.589621,1.032859)--(-2.721975,0.988910)--(-2.789801,0.853522)--(-2.654149,0.904054)--cycle;
\path[fill=blue!19!white,draw=blue!19!white,line width=.12pt] (2.654149,0.904054)--(2.789801,0.853522)--(2.721975,0.988910)--(2.589621,1.032859)--cycle;
\path[fill=blue!26!white,draw=blue!26!white,line width=.12pt] (-1.545814,0.369621)--(-1.589379,0.506197)--(-1.628983,0.427143)--(-1.584333,0.292734)--cycle;
\path[fill=blue!18!white,draw=blue!18!white,line width=.12pt] (1.584333,0.292734)--(1.628983,0.427143)--(1.589379,0.506197)--(1.545814,0.369621)--cycle;
\path[fill=blue!40!white,draw=blue!40!white,line width=.12pt] (-1.955544,-0.576059)--(-1.834315,-0.489898)--(-1.825482,-0.579795)--(-1.946127,-0.671897)--cycle;
\path[fill=blue!34!white,draw=blue!34!white,line width=.12pt] (1.946127,-0.671897)--(1.825482,-0.579795)--(1.834315,-0.489898)--(1.955544,-0.576059)--cycle;
\path[fill=blue!18!white,draw=blue!18!white,line width=.12pt] (-3.006279,0.683903)--(-3.078787,0.571335)--(-3.123988,0.418974)--(-3.050416,0.535130)--cycle;
\path[fill=blue!26!white,draw=blue!26!white,line width=.12pt] (3.050416,0.535130)--(3.123988,0.418974)--(3.078787,0.571335)--(3.006279,0.683903)--cycle;
\path[fill=blue!35!white,draw=blue!35!white,line width=.12pt] (-1.652899,-0.183733)--(-1.607593,-0.055995)--(-1.615372,-0.135163)--(-1.660896,-0.265131)--cycle;
\path[fill=blue!26!white,draw=blue!26!white,line width=.12pt] (1.660896,-0.265131)--(1.615372,-0.135163)--(1.607593,-0.055995)--(1.652899,-0.183733)--cycle;
\path[fill=blue!34!white,draw=blue!34!white,line width=.12pt] (-2.965685,-0.489898)--(-2.844456,-0.576059)--(-2.830759,-0.715462)--(-2.951405,-0.635242)--cycle;
\path[fill=blue!39!white,draw=blue!39!white,line width=.12pt] (2.951405,-0.635242)--(2.830759,-0.715462)--(2.844456,-0.576059)--(2.965685,-0.489898)--cycle;
\path[fill=blue!39!white,draw=blue!39!white,line width=.12pt] (-1.834315,-0.489898)--(-1.734824,-0.384910)--(-1.726471,-0.469932)--(-1.825482,-0.579795)--cycle;
\path[fill=blue!32!white,draw=blue!32!white,line width=.12pt] (1.825482,-0.579795)--(1.726471,-0.469932)--(1.734824,-0.384910)--(1.834315,-0.489898)--cycle;
\path[fill=blue!24!white,draw=blue!24!white,line width=.12pt] (-3.184591,0.156827)--(-3.169293,0.020912)--(-3.184628,-0.135163)--(-3.200000,0.000000)--cycle;
\path[fill=blue!33!white,draw=blue!33!white,line width=.12pt] (3.200000,0.000000)--(3.184628,-0.135163)--(3.169293,0.020912)--(3.184591,0.156827)--cycle;
\path[fill=blue!14!white,draw=blue!14!white,line width=.12pt] (-2.446009,1.050502)--(-2.589621,1.032859)--(-2.654149,0.904054)--(-2.506958,0.928840)--cycle;
\path[fill=blue!17!white,draw=blue!17!white,line width=.12pt] (2.506958,0.928840)--(2.654149,0.904054)--(2.589621,1.032859)--(2.446009,1.050502)--cycle;
\path[fill=blue!23!white,draw=blue!23!white,line width=.12pt] (-1.589379,0.506197)--(-1.660123,0.636707)--(-1.701490,0.554134)--(-1.628983,0.427143)--cycle;
\path[fill=blue!16!white,draw=blue!16!white,line width=.12pt] (1.628983,0.427143)--(1.701490,0.554134)--(1.660123,0.636707)--(1.589379,0.506197)--cycle;
\path[fill=blue!29!white,draw=blue!29!white,line width=.12pt] (-1.569256,0.156072)--(-1.584333,0.292734)--(-1.607593,0.214330)--(-1.592296,0.078414)--cycle;
\path[fill=blue!20!white,draw=blue!20!white,line width=.12pt] (1.592296,0.078414)--(1.607593,0.214330)--(1.584333,0.292734)--(1.569256,0.156072)--cycle;
\path[fill=blue!16!white,draw=blue!16!white,line width=.12pt] (-2.908701,0.779186)--(-3.006279,0.683903)--(-3.050416,0.535130)--(-2.951405,0.635242)--cycle;
\path[fill=blue!23!white,draw=blue!23!white,line width=.12pt] (2.951405,0.635242)--(3.050416,0.535130)--(3.006279,0.683903)--(2.908701,0.779186)--cycle;
\path[fill=blue!14!white,draw=blue!14!white,line width=.12pt] (-2.296657,1.041162)--(-2.446009,1.050502)--(-2.506958,0.928840)--(-2.353885,0.926929)--cycle;
\path[fill=blue!15!white,draw=blue!15!white,line width=.12pt] (2.353885,0.926929)--(2.506958,0.928840)--(2.446009,1.050502)--(2.296657,1.041162)--cycle;
\path[fill=blue!31!white,draw=blue!31!white,line width=.12pt] (-3.065176,-0.384910)--(-2.965685,-0.489898)--(-2.951405,-0.635242)--(-3.050416,-0.535130)--cycle;
\path[fill=blue!38!white,draw=blue!38!white,line width=.12pt] (3.050416,-0.535130)--(2.951405,-0.635242)--(2.965685,-0.489898)--(3.065176,-0.384910)--cycle;
\path[fill=blue!21!white,draw=blue!21!white,line width=.12pt] (-1.660123,0.636707)--(-1.755330,0.756135)--(-1.799069,0.668826)--(-1.701490,0.554134)--cycle;
\path[fill=blue!14!white,draw=blue!14!white,line width=.12pt] (1.701490,0.554134)--(1.799069,0.668826)--(1.755330,0.756135)--(1.660123,0.636707)--cycle;
\path[fill=blue!37!white,draw=blue!37!white,line width=.12pt] (-1.734824,-0.384910)--(-1.660896,-0.265131)--(-1.652899,-0.346529)--(-1.726471,-0.469932)--cycle;
\path[fill=blue!29!white,draw=blue!29!white,line width=.12pt] (1.726471,-0.469932)--(1.652899,-0.346529)--(1.660896,-0.265131)--(1.734824,-0.384910)--cycle;
\path[fill=blue!32!white,draw=blue!32!white,line width=.12pt] (-1.607593,-0.055995)--(-1.592296,0.078414)--(-1.600000,-0.000000)--(-1.615372,-0.135163)--cycle;
\path[fill=blue!23!white,draw=blue!23!white,line width=.12pt] (1.615372,-0.135163)--(1.600000,-0.000000)--(1.592296,0.078414)--(1.607593,-0.055995)--cycle;
\path[fill=blue!41!white,draw=blue!41!white,line width=.12pt] (-2.233123,-0.789480)--(-2.083771,-0.742699)--(-2.053620,-0.844328)--(-2.200811,-0.898392)--cycle;
\path[fill=blue!38!white,draw=blue!38!white,line width=.12pt] (2.200811,-0.898392)--(2.053620,-0.844328)--(2.083771,-0.742699)--(2.233123,-0.789480)--cycle;
\path[fill=blue!40!white,draw=blue!40!white,line width=.12pt] (-2.388443,-0.810441)--(-2.233123,-0.789480)--(-2.200811,-0.898392)--(-2.353885,-0.926929)--cycle;
\path[fill=blue!39!white,draw=blue!39!white,line width=.12pt] (2.353885,-0.926929)--(2.200811,-0.898392)--(2.233123,-0.789480)--(2.388443,-0.810441)--cycle;
\path[fill=blue!22!white,draw=blue!22!white,line width=.12pt] (-3.169293,0.291237)--(-3.184591,0.156827)--(-3.200000,0.000000)--(-3.184628,0.135163)--cycle;
\path[fill=blue!31!white,draw=blue!31!white,line width=.12pt] (3.184628,0.135163)--(3.200000,0.000000)--(3.184591,0.156827)--(3.169293,0.291237)--cycle;
\path[fill=blue!41!white,draw=blue!41!white,line width=.12pt] (-2.083771,-0.742699)--(-1.946127,-0.671897)--(-1.917969,-0.766813)--(-2.053620,-0.844328)--cycle;
\path[fill=blue!36!white,draw=blue!36!white,line width=.12pt] (2.053620,-0.844328)--(1.917969,-0.766813)--(1.946127,-0.671897)--(2.083771,-0.742699)--cycle;
\path[fill=blue!27!white,draw=blue!27!white,line width=.12pt] (-1.584333,0.292734)--(-1.628983,0.427143)--(-1.652899,0.346529)--(-1.607593,0.214330)--cycle;
\path[fill=blue!18!white,draw=blue!18!white,line width=.12pt] (1.607593,0.214330)--(1.652899,0.346529)--(1.628983,0.427143)--(1.584333,0.292734)--cycle;
\path[fill=blue!15!white,draw=blue!15!white,line width=.12pt] (-2.147305,1.005197)--(-2.296657,1.041162)--(-2.353885,0.926929)--(-2.200811,0.898392)--cycle;
\path[fill=blue!14!white,draw=blue!14!white,line width=.12pt] (2.200811,0.898392)--(2.353885,0.926929)--(2.296657,1.041162)--(2.147305,1.005197)--cycle;
\path[fill=blue!19!white,draw=blue!19!white,line width=.12pt] (-1.755330,0.756135)--(-1.871339,0.859891)--(-1.917969,0.766813)--(-1.799069,0.668826)--cycle;
\path[fill=blue!13!white,draw=blue!13!white,line width=.12pt] (1.799069,0.668826)--(1.917969,0.766813)--(1.871339,0.859891)--(1.755330,0.756135)--cycle;
\path[fill=blue!39!white,draw=blue!39!white,line width=.12pt] (-2.543764,-0.804777)--(-2.388443,-0.810441)--(-2.353885,-0.926929)--(-2.506958,-0.928840)--cycle;
\path[fill=blue!40!white,draw=blue!40!white,line width=.12pt] (2.506958,-0.928840)--(2.353885,-0.926929)--(2.388443,-0.810441)--(2.543764,-0.804777)--cycle;
\path[fill=blue!16!white,draw=blue!16!white,line width=.12pt] (-2.003693,0.943989)--(-2.147305,1.005197)--(-2.200811,0.898392)--(-2.053620,0.844328)--cycle;
\path[fill=blue!13!white,draw=blue!13!white,line width=.12pt] (2.053620,0.844328)--(2.200811,0.898392)--(2.147305,1.005197)--(2.003693,0.943989)--cycle;
\path[fill=blue!17!white,draw=blue!17!white,line width=.12pt] (-1.871339,0.859891)--(-2.003693,0.943989)--(-2.053620,0.844328)--(-1.917969,0.766813)--cycle;
\path[fill=blue!13!white,draw=blue!13!white,line width=.12pt] (1.917969,0.766813)--(2.053620,0.844328)--(2.003693,0.943989)--(1.871339,0.859891)--cycle;
\path[fill=blue!40!white,draw=blue!40!white,line width=.12pt] (-1.946127,-0.671897)--(-1.825482,-0.579795)--(-1.799069,-0.668826)--(-1.917969,-0.766813)--cycle;
\path[fill=blue!35!white,draw=blue!35!white,line width=.12pt] (1.917969,-0.766813)--(1.799069,-0.668826)--(1.825482,-0.579795)--(1.946127,-0.671897)--cycle;
\path[fill=blue!15!white,draw=blue!15!white,line width=.12pt] (-2.789801,0.853522)--(-2.908701,0.779186)--(-2.951405,0.635242)--(-2.830759,0.715462)--cycle;
\path[fill=blue!21!white,draw=blue!21!white,line width=.12pt] (2.830759,0.715462)--(2.951405,0.635242)--(2.908701,0.779186)--(2.789801,0.853522)--cycle;
\path[fill=blue!29!white,draw=blue!29!white,line width=.12pt] (-3.139104,-0.265131)--(-3.065176,-0.384910)--(-3.050416,-0.535130)--(-3.123988,-0.418974)--cycle;
\path[fill=blue!37!white,draw=blue!37!white,line width=.12pt] (3.123988,-0.418974)--(3.050416,-0.535130)--(3.065176,-0.384910)--(3.139104,-0.265131)--cycle;
\path[fill=blue!37!white,draw=blue!37!white,line width=.12pt] (-2.693116,-0.772707)--(-2.543764,-0.804777)--(-2.506958,-0.928840)--(-2.654149,-0.904054)--cycle;
\path[fill=blue!40!white,draw=blue!40!white,line width=.12pt] (2.654149,-0.904054)--(2.506958,-0.928840)--(2.543764,-0.804777)--(2.693116,-0.772707)--cycle;
\path[fill=blue!35!white,draw=blue!35!white,line width=.12pt] (-1.660896,-0.265131)--(-1.615372,-0.135163)--(-1.607593,-0.214330)--(-1.652899,-0.346529)--cycle;
\path[fill=blue!27!white,draw=blue!27!white,line width=.12pt] (1.652899,-0.346529)--(1.607593,-0.214330)--(1.615372,-0.135163)--(1.660896,-0.265131)--cycle;
\path[fill=blue!30!white,draw=blue!30!white,line width=.12pt] (-1.592296,0.078414)--(-1.607593,0.214330)--(-1.615372,0.135163)--(-1.600000,0.000000)--cycle;
\path[fill=blue!21!white,draw=blue!21!white,line width=.12pt] (1.600000,0.000000)--(1.615372,0.135163)--(1.607593,0.214330)--(1.592296,0.078414)--cycle;
\path[fill=blue!24!white,draw=blue!24!white,line width=.12pt] (-1.628983,0.427143)--(-1.701490,0.554134)--(-1.726471,0.469932)--(-1.652899,0.346529)--cycle;
\path[fill=blue!16!white,draw=blue!16!white,line width=.12pt] (1.652899,0.346529)--(1.726471,0.469932)--(1.701490,0.554134)--(1.628983,0.427143)--cycle;
\path[fill=blue!39!white,draw=blue!39!white,line width=.12pt] (-1.825482,-0.579795)--(-1.726471,-0.469932)--(-1.701490,-0.554134)--(-1.799069,-0.668826)--cycle;
\path[fill=blue!32!white,draw=blue!32!white,line width=.12pt] (1.799069,-0.668826)--(1.701490,-0.554134)--(1.726471,-0.469932)--(1.825482,-0.579795)--cycle;
\path[fill=blue!19!white,draw=blue!19!white,line width=.12pt] (-3.123988,0.418974)--(-3.169293,0.291237)--(-3.184628,0.135163)--(-3.139104,0.265131)--cycle;
\path[fill=blue!28!white,draw=blue!28!white,line width=.12pt] (3.139104,0.265131)--(3.184628,0.135163)--(3.169293,0.291237)--(3.123988,0.418974)--cycle;
\path[fill=blue!14!white,draw=blue!14!white,line width=.12pt] (-2.654149,0.904054)--(-2.789801,0.853522)--(-2.830759,0.715462)--(-2.693116,0.772707)--cycle;
\path[fill=blue!18!white,draw=blue!18!white,line width=.12pt] (2.693116,0.772707)--(2.830759,0.715462)--(2.789801,0.853522)--(2.654149,0.904054)--cycle;
\path[fill=blue!36!white,draw=blue!36!white,line width=.12pt] (-2.830759,-0.715462)--(-2.693116,-0.772707)--(-2.654149,-0.904054)--(-2.789801,-0.853522)--cycle;
\path[fill=blue!40!white,draw=blue!40!white,line width=.12pt] (2.789801,-0.853522)--(2.654149,-0.904054)--(2.693116,-0.772707)--(2.830759,-0.715462)--cycle;
\path[fill=blue!26!white,draw=blue!26!white,line width=.12pt] (-3.184628,-0.135163)--(-3.139104,-0.265131)--(-3.123988,-0.418974)--(-3.169293,-0.291237)--cycle;
\path[fill=blue!35!white,draw=blue!35!white,line width=.12pt] (3.169293,-0.291237)--(3.123988,-0.418974)--(3.139104,-0.265131)--(3.184628,-0.135163)--cycle;
\path[fill=blue!22!white,draw=blue!22!white,line width=.12pt] (-1.701490,0.554134)--(-1.799069,0.668826)--(-1.825482,0.579795)--(-1.726471,0.469932)--cycle;
\path[fill=blue!15!white,draw=blue!15!white,line width=.12pt] (1.726471,0.469932)--(1.825482,0.579795)--(1.799069,0.668826)--(1.701490,0.554134)--cycle;
\path[fill=blue!38!white,draw=blue!38!white,line width=.12pt] (-1.726471,-0.469932)--(-1.652899,-0.346529)--(-1.628983,-0.427143)--(-1.701490,-0.554134)--cycle;
\path[fill=blue!30!white,draw=blue!30!white,line width=.12pt] (1.701490,-0.554134)--(1.628983,-0.427143)--(1.652899,-0.346529)--(1.726471,-0.469932)--cycle;
\path[fill=blue!33!white,draw=blue!33!white,line width=.12pt] (-1.615372,-0.135163)--(-1.600000,-0.000000)--(-1.592296,-0.078414)--(-1.607593,-0.214330)--cycle;
\path[fill=blue!24!white,draw=blue!24!white,line width=.12pt] (1.607593,-0.214330)--(1.592296,-0.078414)--(1.600000,-0.000000)--(1.615372,-0.135163)--cycle;
\path[fill=blue!27!white,draw=blue!27!white,line width=.12pt] (-1.607593,0.214330)--(-1.652899,0.346529)--(-1.660896,0.265131)--(-1.615372,0.135163)--cycle;
\path[fill=blue!19!white,draw=blue!19!white,line width=.12pt] (1.615372,0.135163)--(1.660896,0.265131)--(1.652899,0.346529)--(1.607593,0.214330)--cycle;
\path[fill=blue!14!white,draw=blue!14!white,line width=.12pt] (-2.506958,0.928840)--(-2.654149,0.904054)--(-2.693116,0.772707)--(-2.543764,0.804777)--cycle;
\path[fill=blue!17!white,draw=blue!17!white,line width=.12pt] (2.543764,0.804777)--(2.693116,0.772707)--(2.654149,0.904054)--(2.506958,0.928840)--cycle;
\path[fill=blue!17!white,draw=blue!17!white,line width=.12pt] (-3.050416,0.535130)--(-3.123988,0.418974)--(-3.139104,0.265131)--(-3.065176,0.384910)--cycle;
\path[fill=blue!25!white,draw=blue!25!white,line width=.12pt] (3.065176,0.384910)--(3.139104,0.265131)--(3.123988,0.418974)--(3.050416,0.535130)--cycle;
\path[fill=blue!33!white,draw=blue!33!white,line width=.12pt] (-2.951405,-0.635242)--(-2.830759,-0.715462)--(-2.789801,-0.853522)--(-2.908701,-0.779186)--cycle;
\path[fill=blue!39!white,draw=blue!39!white,line width=.12pt] (2.908701,-0.779186)--(2.789801,-0.853522)--(2.830759,-0.715462)--(2.951405,-0.635242)--cycle;
\path[fill=blue!19!white,draw=blue!19!white,line width=.12pt] (-1.799069,0.668826)--(-1.917969,0.766813)--(-1.946127,0.671897)--(-1.825482,0.579795)--cycle;
\path[fill=blue!14!white,draw=blue!14!white,line width=.12pt] (1.825482,0.579795)--(1.946127,0.671897)--(1.917969,0.766813)--(1.799069,0.668826)--cycle;
\path[fill=blue!41!white,draw=blue!41!white,line width=.12pt] (-2.053620,-0.844328)--(-1.917969,-0.766813)--(-1.871339,-0.859891)--(-2.003693,-0.943989)--cycle;
\path[fill=blue!37!white,draw=blue!37!white,line width=.12pt] (2.003693,-0.943989)--(1.871339,-0.859891)--(1.917969,-0.766813)--(2.053620,-0.844328)--cycle;
\path[fill=blue!41!white,draw=blue!41!white,line width=.12pt] (-2.200811,-0.898392)--(-2.053620,-0.844328)--(-2.003693,-0.943989)--(-2.147305,-1.005197)--cycle;
\path[fill=blue!38!white,draw=blue!38!white,line width=.12pt] (2.147305,-1.005197)--(2.003693,-0.943989)--(2.053620,-0.844328)--(2.200811,-0.898392)--cycle;
\path[fill=blue!14!white,draw=blue!14!white,line width=.12pt] (-2.353885,0.926929)--(-2.506958,0.928840)--(-2.543764,0.804777)--(-2.388443,0.810441)--cycle;
\path[fill=blue!15!white,draw=blue!15!white,line width=.12pt] (2.388443,0.810441)--(2.543764,0.804777)--(2.506958,0.928840)--(2.353885,0.926929)--cycle;
\path[fill=blue!41!white,draw=blue!41!white,line width=.12pt] (-1.917969,-0.766813)--(-1.799069,-0.668826)--(-1.755330,-0.756135)--(-1.871339,-0.859891)--cycle;
\path[fill=blue!35!white,draw=blue!35!white,line width=.12pt] (1.871339,-0.859891)--(1.755330,-0.756135)--(1.799069,-0.668826)--(1.917969,-0.766813)--cycle;
\path[fill=blue!40!white,draw=blue!40!white,line width=.12pt] (-2.353885,-0.926929)--(-2.200811,-0.898392)--(-2.147305,-1.005197)--(-2.296657,-1.041162)--cycle;
\path[fill=blue!39!white,draw=blue!39!white,line width=.12pt] (2.296657,-1.041162)--(2.147305,-1.005197)--(2.200811,-0.898392)--(2.353885,-0.926929)--cycle;
\path[fill=blue!36!white,draw=blue!36!white,line width=.12pt] (-1.652899,-0.346529)--(-1.607593,-0.214330)--(-1.584333,-0.292734)--(-1.628983,-0.427143)--cycle;
\path[fill=blue!27!white,draw=blue!27!white,line width=.12pt] (1.628983,-0.427143)--(1.584333,-0.292734)--(1.607593,-0.214330)--(1.652899,-0.346529)--cycle;
\path[fill=blue!18!white,draw=blue!18!white,line width=.12pt] (-1.917969,0.766813)--(-2.053620,0.844328)--(-2.083771,0.742699)--(-1.946127,0.671897)--cycle;
\path[fill=blue!13!white,draw=blue!13!white,line width=.12pt] (1.946127,0.671897)--(2.083771,0.742699)--(2.053620,0.844328)--(1.917969,0.766813)--cycle;
\path[fill=blue!23!white,draw=blue!23!white,line width=.12pt] (-3.200000,0.000000)--(-3.184628,-0.135163)--(-3.169293,-0.291237)--(-3.184591,-0.156827)--cycle;
\path[fill=blue!32!white,draw=blue!32!white,line width=.12pt] (3.184591,-0.156827)--(3.169293,-0.291237)--(3.184628,-0.135163)--(3.200000,-0.000000)--cycle;
\path[fill=blue!15!white,draw=blue!15!white,line width=.12pt] (-2.200811,0.898392)--(-2.353885,0.926929)--(-2.388443,0.810441)--(-2.233123,0.789480)--cycle;
\path[fill=blue!14!white,draw=blue!14!white,line width=.12pt] (2.233123,0.789480)--(2.388443,0.810441)--(2.353885,0.926929)--(2.200811,0.898392)--cycle;
\path[fill=blue!16!white,draw=blue!16!white,line width=.12pt] (-2.053620,0.844328)--(-2.200811,0.898392)--(-2.233123,0.789480)--(-2.083771,0.742699)--cycle;
\path[fill=blue!13!white,draw=blue!13!white,line width=.12pt] (2.083771,0.742699)--(2.233123,0.789480)--(2.200811,0.898392)--(2.053620,0.844328)--cycle;
\path[fill=blue!31!white,draw=blue!31!white,line width=.12pt] (-1.600000,0.000000)--(-1.615372,0.135163)--(-1.607593,0.055995)--(-1.592296,-0.078414)--cycle;
\path[fill=blue!22!white,draw=blue!22!white,line width=.12pt] (1.592296,-0.078414)--(1.607593,0.055995)--(1.615372,0.135163)--(1.600000,-0.000000)--cycle;
\path[fill=blue!25!white,draw=blue!25!white,line width=.12pt] (-1.652899,0.346529)--(-1.726471,0.469932)--(-1.734824,0.384910)--(-1.660896,0.265131)--cycle;
\path[fill=blue!17!white,draw=blue!17!white,line width=.12pt] (1.660896,0.265131)--(1.734824,0.384910)--(1.726471,0.469932)--(1.652899,0.346529)--cycle;
\path[fill=blue!40!white,draw=blue!40!white,line width=.12pt] (-1.799069,-0.668826)--(-1.701490,-0.554134)--(-1.660123,-0.636707)--(-1.755330,-0.756135)--cycle;
\path[fill=blue!33!white,draw=blue!33!white,line width=.12pt] (1.755330,-0.756135)--(1.660123,-0.636707)--(1.701490,-0.554134)--(1.799069,-0.668826)--cycle;
\path[fill=blue!16!white,draw=blue!16!white,line width=.12pt] (-2.951405,0.635242)--(-3.050416,0.535130)--(-3.065176,0.384910)--(-2.965685,0.489898)--cycle;
\path[fill=blue!23!white,draw=blue!23!white,line width=.12pt] (2.965685,0.489898)--(3.065176,0.384910)--(3.050416,0.535130)--(2.951405,0.635242)--cycle;
\path[fill=blue!39!white,draw=blue!39!white,line width=.12pt] (-2.506958,-0.928840)--(-2.353885,-0.926929)--(-2.296657,-1.041162)--(-2.446009,-1.050502)--cycle;
\path[fill=blue!40!white,draw=blue!40!white,line width=.12pt] (2.446009,-1.050502)--(2.296657,-1.041162)--(2.353885,-0.926929)--(2.506958,-0.928840)--cycle;
\path[fill=blue!31!white,draw=blue!31!white,line width=.12pt] (-3.050416,-0.535130)--(-2.951405,-0.635242)--(-2.908701,-0.779186)--(-3.006279,-0.683903)--cycle;
\path[fill=blue!38!white,draw=blue!38!white,line width=.12pt] (3.006279,-0.683903)--(2.908701,-0.779186)--(2.951405,-0.635242)--(3.050416,-0.535130)--cycle;
\path[fill=blue!34!white,draw=blue!34!white,line width=.12pt] (-1.607593,-0.214330)--(-1.592296,-0.078414)--(-1.569256,-0.156072)--(-1.584333,-0.292734)--cycle;
\path[fill=blue!25!white,draw=blue!25!white,line width=.12pt] (1.584333,-0.292734)--(1.569256,-0.156072)--(1.592296,-0.078414)--(1.607593,-0.214330)--cycle;
\path[fill=blue!38!white,draw=blue!38!white,line width=.12pt] (-1.701490,-0.554134)--(-1.628983,-0.427143)--(-1.589379,-0.506197)--(-1.660123,-0.636707)--cycle;
\path[fill=blue!31!white,draw=blue!31!white,line width=.12pt] (1.660123,-0.636707)--(1.589379,-0.506197)--(1.628983,-0.427143)--(1.701490,-0.554134)--cycle;
\path[fill=blue!37!white,draw=blue!37!white,line width=.12pt] (-2.654149,-0.904054)--(-2.506958,-0.928840)--(-2.446009,-1.050502)--(-2.589621,-1.032859)--cycle;
\path[fill=blue!40!white,draw=blue!40!white,line width=.12pt] (2.589621,-1.032859)--(2.446009,-1.050502)--(2.506958,-0.928840)--(2.654149,-0.904054)--cycle;
\path[fill=blue!21!white,draw=blue!21!white,line width=.12pt] (-3.184628,0.135163)--(-3.200000,0.000000)--(-3.184591,-0.156827)--(-3.169293,-0.020912)--cycle;
\path[fill=blue!30!white,draw=blue!30!white,line width=.12pt] (3.169293,-0.020912)--(3.184591,-0.156827)--(3.200000,-0.000000)--(3.184628,0.135163)--cycle;
\path[fill=blue!22!white,draw=blue!22!white,line width=.12pt] (-1.726471,0.469932)--(-1.825482,0.579795)--(-1.834315,0.489898)--(-1.734824,0.384910)--cycle;
\path[fill=blue!15!white,draw=blue!15!white,line width=.12pt] (1.734824,0.384910)--(1.834315,0.489898)--(1.825482,0.579795)--(1.726471,0.469932)--cycle;
\path[fill=blue!15!white,draw=blue!15!white,line width=.12pt] (-2.830759,0.715462)--(-2.951405,0.635242)--(-2.965685,0.489898)--(-2.844456,0.576059)--cycle;
\path[fill=blue!20!white,draw=blue!20!white,line width=.12pt] (2.844456,0.576059)--(2.965685,0.489898)--(2.951405,0.635242)--(2.830759,0.715462)--cycle;
\path[fill=blue!28!white,draw=blue!28!white,line width=.12pt] (-1.615372,0.135163)--(-1.660896,0.265131)--(-1.652899,0.183733)--(-1.607593,0.055995)--cycle;
\path[fill=blue!19!white,draw=blue!19!white,line width=.12pt] (1.607593,0.055995)--(1.652899,0.183733)--(1.660896,0.265131)--(1.615372,0.135163)--cycle;
\path[fill=blue!28!white,draw=blue!28!white,line width=.12pt] (-3.123988,-0.418974)--(-3.050416,-0.535130)--(-3.006279,-0.683903)--(-3.078787,-0.571335)--cycle;
\path[fill=blue!36!white,draw=blue!36!white,line width=.12pt] (3.078787,-0.571335)--(3.006279,-0.683903)--(3.050416,-0.535130)--(3.123988,-0.418974)--cycle;
\path[fill=blue!20!white,draw=blue!20!white,line width=.12pt] (-1.825482,0.579795)--(-1.946127,0.671897)--(-1.955544,0.576059)--(-1.834315,0.489898)--cycle;
\path[fill=blue!14!white,draw=blue!14!white,line width=.12pt] (1.834315,0.489898)--(1.955544,0.576059)--(1.946127,0.671897)--(1.825482,0.579795)--cycle;
\path[fill=blue!36!white,draw=blue!36!white,line width=.12pt] (-1.628983,-0.427143)--(-1.584333,-0.292734)--(-1.545814,-0.369621)--(-1.589379,-0.506197)--cycle;
\path[fill=blue!28!white,draw=blue!28!white,line width=.12pt] (1.589379,-0.506197)--(1.545814,-0.369621)--(1.584333,-0.292734)--(1.628983,-0.427143)--cycle;
\path[fill=blue!35!white,draw=blue!35!white,line width=.12pt] (-2.789801,-0.853522)--(-2.654149,-0.904054)--(-2.589621,-1.032859)--(-2.721975,-0.988910)--cycle;
\path[fill=blue!39!white,draw=blue!39!white,line width=.12pt] (2.721975,-0.988910)--(2.589621,-1.032859)--(2.654149,-0.904054)--(2.789801,-0.853522)--cycle;
\path[fill=blue!14!white,draw=blue!14!white,line width=.12pt] (-2.693116,0.772707)--(-2.830759,0.715462)--(-2.844456,0.576059)--(-2.706147,0.640083)--cycle;
\path[fill=blue!18!white,draw=blue!18!white,line width=.12pt] (2.706147,0.640083)--(2.844456,0.576059)--(2.830759,0.715462)--(2.693116,0.772707)--cycle;
\path[fill=blue!41!white,draw=blue!41!white,line width=.12pt] (-2.003693,-0.943989)--(-1.871339,-0.859891)--(-1.806687,-0.950236)--(-1.934468,-1.040724)--cycle;
\path[fill=blue!37!white,draw=blue!37!white,line width=.12pt] (1.934468,-1.040724)--(1.806687,-0.950236)--(1.871339,-0.859891)--(2.003693,-0.943989)--cycle;
\path[fill=blue!31!white,draw=blue!31!white,line width=.12pt] (-1.592296,-0.078414)--(-1.607593,0.055995)--(-1.584333,-0.022409)--(-1.569256,-0.156072)--cycle;
\path[fill=blue!22!white,draw=blue!22!white,line width=.12pt] (1.569256,-0.156072)--(1.584333,-0.022409)--(1.607593,0.055995)--(1.592296,-0.078414)--cycle;
\path[fill=blue!41!white,draw=blue!41!white,line width=.12pt] (-1.871339,-0.859891)--(-1.755330,-0.756135)--(-1.694686,-0.840879)--(-1.806687,-0.950236)--cycle;
\path[fill=blue!36!white,draw=blue!36!white,line width=.12pt] (1.806687,-0.950236)--(1.694686,-0.840879)--(1.755330,-0.756135)--(1.871339,-0.859891)--cycle;
\path[fill=blue!19!white,draw=blue!19!white,line width=.12pt] (-3.139104,0.265131)--(-3.184628,0.135163)--(-3.169293,-0.020912)--(-3.123988,0.111288)--cycle;
\path[fill=blue!27!white,draw=blue!27!white,line width=.12pt] (3.123988,0.111288)--(3.169293,-0.020912)--(3.184628,0.135163)--(3.139104,0.265131)--cycle;
\path[fill=blue!41!white,draw=blue!41!white,line width=.12pt] (-2.147305,-1.005197)--(-2.003693,-0.943989)--(-1.934468,-1.040724)--(-2.073119,-1.108865)--cycle;
\path[fill=blue!38!white,draw=blue!38!white,line width=.12pt] (2.073119,-1.108865)--(1.934468,-1.040724)--(2.003693,-0.943989)--(2.147305,-1.005197)--cycle;
\path[fill=blue!25!white,draw=blue!25!white,line width=.12pt] (-1.660896,0.265131)--(-1.734824,0.384910)--(-1.726471,0.299889)--(-1.652899,0.183733)--cycle;
\path[fill=blue!17!white,draw=blue!17!white,line width=.12pt] (1.652899,0.183733)--(1.726471,0.299889)--(1.734824,0.384910)--(1.660896,0.265131)--cycle;
\path[fill=blue!18!white,draw=blue!18!white,line width=.12pt] (-1.946127,0.671897)--(-2.083771,0.742699)--(-2.093853,0.640083)--(-1.955544,0.576059)--cycle;
\path[fill=blue!14!white,draw=blue!14!white,line width=.12pt] (1.955544,0.576059)--(2.093853,0.640083)--(2.083771,0.742699)--(1.946127,0.671897)--cycle;
\path[fill=blue!40!white,draw=blue!40!white,line width=.12pt] (-1.755330,-0.756135)--(-1.660123,-0.636707)--(-1.602769,-0.716855)--(-1.694686,-0.840879)--cycle;
\path[fill=blue!33!white,draw=blue!33!white,line width=.12pt] (1.694686,-0.840879)--(1.602769,-0.716855)--(1.660123,-0.636707)--(1.755330,-0.756135)--cycle;
\path[fill=blue!14!white,draw=blue!14!white,line width=.12pt] (-2.543764,0.804777)--(-2.693116,0.772707)--(-2.706147,0.640083)--(-2.556072,0.679508)--cycle;
\path[fill=blue!16!white,draw=blue!16!white,line width=.12pt] (2.556072,0.679508)--(2.706147,0.640083)--(2.693116,0.772707)--(2.543764,0.804777)--cycle;
\path[fill=blue!40!white,draw=blue!40!white,line width=.12pt] (-2.296657,-1.041162)--(-2.147305,-1.005197)--(-2.073119,-1.108865)--(-2.217311,-1.152040)--cycle;
\path[fill=blue!39!white,draw=blue!39!white,line width=.12pt] (2.217311,-1.152040)--(2.073119,-1.108865)--(2.147305,-1.005197)--(2.296657,-1.041162)--cycle;
\path[fill=blue!25!white,draw=blue!25!white,line width=.12pt] (-3.169293,-0.291237)--(-3.123988,-0.418974)--(-3.078787,-0.571335)--(-3.123436,-0.445808)--cycle;
\path[fill=blue!34!white,draw=blue!34!white,line width=.12pt] (3.123436,-0.445808)--(3.078787,-0.571335)--(3.123988,-0.418974)--(3.169293,-0.291237)--cycle;
\path[fill=blue!16!white,draw=blue!16!white,line width=.12pt] (-2.083771,0.742699)--(-2.233123,0.789480)--(-2.243928,0.679508)--(-2.093853,0.640083)--cycle;
\path[fill=blue!14!white,draw=blue!14!white,line width=.12pt] (2.093853,0.640083)--(2.243928,0.679508)--(2.233123,0.789480)--(2.083771,0.742699)--cycle;
\path[fill=blue!14!white,draw=blue!14!white,line width=.12pt] (-2.388443,0.810441)--(-2.543764,0.804777)--(-2.556072,0.679508)--(-2.400000,0.692820)--cycle;
\path[fill=blue!15!white,draw=blue!15!white,line width=.12pt] (2.400000,0.692820)--(2.556072,0.679508)--(2.543764,0.804777)--(2.388443,0.810441)--cycle;
\path[fill=blue!15!white,draw=blue!15!white,line width=.12pt] (-2.233123,0.789480)--(-2.388443,0.810441)--(-2.400000,0.692820)--(-2.243928,0.679508)--cycle;
\path[fill=blue!14!white,draw=blue!14!white,line width=.12pt] (2.243928,0.679508)--(2.400000,0.692820)--(2.388443,0.810441)--(2.233123,0.789480)--cycle;
\path[fill=blue!34!white,draw=blue!34!white,line width=.12pt] (-1.584333,-0.292734)--(-1.569256,-0.156072)--(-1.531105,-0.232228)--(-1.545814,-0.369621)--cycle;
\path[fill=blue!26!white,draw=blue!26!white,line width=.12pt] (1.545814,-0.369621)--(1.531105,-0.232228)--(1.569256,-0.156072)--(1.584333,-0.292734)--cycle;
\path[fill=blue!33!white,draw=blue!33!white,line width=.12pt] (-2.908701,-0.779186)--(-2.789801,-0.853522)--(-2.721975,-0.988910)--(-2.837984,-0.920344)--cycle;
\path[fill=blue!38!white,draw=blue!38!white,line width=.12pt] (2.837984,-0.920344)--(2.721975,-0.988910)--(2.789801,-0.853522)--(2.908701,-0.779186)--cycle;
\path[fill=blue!39!white,draw=blue!39!white,line width=.12pt] (-1.660123,-0.636707)--(-1.589379,-0.506197)--(-1.534468,-0.582930)--(-1.602769,-0.716855)--cycle;
\path[fill=blue!31!white,draw=blue!31!white,line width=.12pt] (1.602769,-0.716855)--(1.534468,-0.582930)--(1.589379,-0.506197)--(1.660123,-0.636707)--cycle;
\path[fill=blue!17!white,draw=blue!17!white,line width=.12pt] (-3.065176,0.384910)--(-3.139104,0.265131)--(-3.123988,0.111288)--(-3.050416,0.234690)--cycle;
\path[fill=blue!25!white,draw=blue!25!white,line width=.12pt] (3.050416,0.234690)--(3.123988,0.111288)--(3.139104,0.265131)--(3.065176,0.384910)--cycle;
\path[fill=blue!29!white,draw=blue!29!white,line width=.12pt] (-1.607593,0.055995)--(-1.652899,0.183733)--(-1.628983,0.103118)--(-1.584333,-0.022409)--cycle;
\path[fill=blue!20!white,draw=blue!20!white,line width=.12pt] (1.584333,-0.022409)--(1.628983,0.103118)--(1.652899,0.183733)--(1.607593,0.055995)--cycle;
\path[fill=blue!39!white,draw=blue!39!white,line width=.12pt] (-2.446009,-1.050502)--(-2.296657,-1.041162)--(-2.217311,-1.152040)--(-2.361503,-1.168591)--cycle;
\path[fill=blue!40!white,draw=blue!40!white,line width=.12pt] (2.361503,-1.168591)--(2.217311,-1.152040)--(2.296657,-1.041162)--(2.446009,-1.050502)--cycle;
\path[fill=blue!23!white,draw=blue!23!white,line width=.12pt] (-1.734824,0.384910)--(-1.834315,0.489898)--(-1.825482,0.400001)--(-1.726471,0.299889)--cycle;
\path[fill=blue!16!white,draw=blue!16!white,line width=.12pt] (1.726471,0.299889)--(1.825482,0.400001)--(1.834315,0.489898)--(1.734824,0.384910)--cycle;
\path[fill=blue!37!white,draw=blue!37!white,line width=.12pt] (-1.589379,-0.506197)--(-1.545814,-0.369621)--(-1.492409,-0.444251)--(-1.534468,-0.582930)--cycle;
\path[fill=blue!29!white,draw=blue!29!white,line width=.12pt] (1.534468,-0.582930)--(1.492409,-0.444251)--(1.545814,-0.369621)--(1.589379,-0.506197)--cycle;
\path[fill=blue!23!white,draw=blue!23!white,line width=.12pt] (-3.184591,-0.156827)--(-3.169293,-0.291237)--(-3.123436,-0.445808)--(-3.138513,-0.312145)--cycle;
\path[fill=blue!32!white,draw=blue!32!white,line width=.12pt] (3.138513,-0.312145)--(3.123436,-0.445808)--(3.169293,-0.291237)--(3.184591,-0.156827)--cycle;
\path[fill=blue!32!white,draw=blue!32!white,line width=.12pt] (-1.569256,-0.156072)--(-1.584333,-0.022409)--(-1.545814,-0.099296)--(-1.531105,-0.232228)--cycle;
\path[fill=blue!23!white,draw=blue!23!white,line width=.12pt] (1.531105,-0.232228)--(1.545814,-0.099296)--(1.584333,-0.022409)--(1.569256,-0.156072)--cycle;
\path[fill=blue!15!white,draw=blue!15!white,line width=.12pt] (-2.965685,0.489898)--(-3.065176,0.384910)--(-3.050416,0.234690)--(-2.951405,0.344554)--cycle;
\path[fill=blue!22!white,draw=blue!22!white,line width=.12pt] (2.951405,0.344554)--(3.050416,0.234690)--(3.065176,0.384910)--(2.965685,0.489898)--cycle;
\path[fill=blue!30!white,draw=blue!30!white,line width=.12pt] (-3.006279,-0.683903)--(-2.908701,-0.779186)--(-2.837984,-0.920344)--(-2.933190,-0.829797)--cycle;
\path[fill=blue!37!white,draw=blue!37!white,line width=.12pt] (2.933190,-0.829797)--(2.837984,-0.920344)--(2.908701,-0.779186)--(3.006279,-0.683903)--cycle;
\path[fill=blue!37!white,draw=blue!37!white,line width=.12pt] (-2.589621,-1.032859)--(-2.446009,-1.050502)--(-2.361503,-1.168591)--(-2.500154,-1.157881)--cycle;
\path[fill=blue!40!white,draw=blue!40!white,line width=.12pt] (2.500154,-1.157881)--(2.361503,-1.168591)--(2.446009,-1.050502)--(2.589621,-1.032859)--cycle;
\path[fill=blue!20!white,draw=blue!20!white,line width=.12pt] (-1.834315,0.489898)--(-1.955544,0.576059)--(-1.946127,0.480221)--(-1.825482,0.400001)--cycle;
\path[fill=blue!15!white,draw=blue!15!white,line width=.12pt] (1.825482,0.400001)--(1.946127,0.480221)--(1.955544,0.576059)--(1.834315,0.489898)--cycle;
\path[fill=blue!41!white,draw=blue!41!white,line width=.12pt] (-1.806687,-0.950236)--(-1.694686,-0.840879)--(-1.617721,-0.922243)--(-1.724636,-1.036978)--cycle;
\path[fill=blue!36!white,draw=blue!36!white,line width=.12pt] (1.724636,-1.036978)--(1.617721,-0.922243)--(1.694686,-0.840879)--(1.806687,-0.950236)--cycle;
\path[fill=blue!41!white,draw=blue!41!white,line width=.12pt] (-1.934468,-1.040724)--(-1.806687,-0.950236)--(-1.724636,-1.036978)--(-1.846614,-1.133600)--cycle;
\path[fill=blue!38!white,draw=blue!38!white,line width=.12pt] (1.846614,-1.133600)--(1.724636,-1.036978)--(1.806687,-0.950236)--(1.934468,-1.040724)--cycle;
\path[fill=blue!26!white,draw=blue!26!white,line width=.12pt] (-1.652899,0.183733)--(-1.726471,0.299889)--(-1.701490,0.215687)--(-1.628983,0.103118)--cycle;
\path[fill=blue!18!white,draw=blue!18!white,line width=.12pt] (1.628983,0.103118)--(1.701490,0.215687)--(1.726471,0.299889)--(1.652899,0.183733)--cycle;
\path[fill=blue!40!white,draw=blue!40!white,line width=.12pt] (-1.694686,-0.840879)--(-1.602769,-0.716855)--(-1.529978,-0.793806)--(-1.617721,-0.922243)--cycle;
\path[fill=blue!34!white,draw=blue!34!white,line width=.12pt] (1.617721,-0.922243)--(1.529978,-0.793806)--(1.602769,-0.716855)--(1.694686,-0.840879)--cycle;
\path[fill=blue!41!white,draw=blue!41!white,line width=.12pt] (-2.073119,-1.108865)--(-1.934468,-1.040724)--(-1.846614,-1.133600)--(-1.978968,-1.208398)--cycle;
\path[fill=blue!39!white,draw=blue!39!white,line width=.12pt] (1.978968,-1.208398)--(1.846614,-1.133600)--(1.934468,-1.040724)--(2.073119,-1.108865)--cycle;
\path[fill=blue!35!white,draw=blue!35!white,line width=.12pt] (-1.545814,-0.369621)--(-1.531105,-0.232228)--(-1.478207,-0.306147)--(-1.492409,-0.444251)--cycle;
\path[fill=blue!26!white,draw=blue!26!white,line width=.12pt] (1.492409,-0.444251)--(1.478207,-0.306147)--(1.531105,-0.232228)--(1.545814,-0.369621)--cycle;
\path[fill=blue!14!white,draw=blue!14!white,line width=.12pt] (-2.844456,0.576059)--(-2.965685,0.489898)--(-2.951405,0.344554)--(-2.830759,0.436656)--cycle;
\path[fill=blue!20!white,draw=blue!20!white,line width=.12pt] (2.830759,0.436656)--(2.951405,0.344554)--(2.965685,0.489898)--(2.844456,0.576059)--cycle;
\path[fill=blue!18!white,draw=blue!18!white,line width=.12pt] (-1.955544,0.576059)--(-2.093853,0.640083)--(-2.083771,0.537466)--(-1.946127,0.480221)--cycle;
\path[fill=blue!14!white,draw=blue!14!white,line width=.12pt] (1.946127,0.480221)--(2.083771,0.537466)--(2.093853,0.640083)--(1.955544,0.576059)--cycle;
\path[fill=blue!39!white,draw=blue!39!white,line width=.12pt] (-1.602769,-0.716855)--(-1.534468,-0.582930)--(-1.464780,-0.656601)--(-1.529978,-0.793806)--cycle;
\path[fill=blue!32!white,draw=blue!32!white,line width=.12pt] (1.529978,-0.793806)--(1.464780,-0.656601)--(1.534468,-0.582930)--(1.602769,-0.716855)--cycle;
\path[fill=blue!20!white,draw=blue!20!white,line width=.12pt] (-3.169293,-0.020912)--(-3.184591,-0.156827)--(-3.138513,-0.312145)--(-3.123436,-0.175483)--cycle;
\path[fill=blue!29!white,draw=blue!29!white,line width=.12pt] (3.123436,-0.175483)--(3.138513,-0.312145)--(3.184591,-0.156827)--(3.169293,-0.020912)--cycle;
\path[fill=blue!35!white,draw=blue!35!white,line width=.12pt] (-2.721975,-0.988910)--(-2.589621,-1.032859)--(-2.500154,-1.157881)--(-2.627935,-1.120322)--cycle;
\path[fill=blue!39!white,draw=blue!39!white,line width=.12pt] (2.627935,-1.120322)--(2.500154,-1.157881)--(2.589621,-1.032859)--(2.721975,-0.988910)--cycle;
\path[fill=blue!29!white,draw=blue!29!white,line width=.12pt] (-1.584333,-0.022409)--(-1.628983,0.103118)--(-1.589379,0.024064)--(-1.545814,-0.099296)--cycle;
\path[fill=blue!21!white,draw=blue!21!white,line width=.12pt] (1.545814,-0.099296)--(1.589379,0.024064)--(1.628983,0.103118)--(1.584333,-0.022409)--cycle;
\path[fill=blue!40!white,draw=blue!40!white,line width=.12pt] (-2.217311,-1.152040)--(-2.073119,-1.108865)--(-1.978968,-1.208398)--(-2.116611,-1.258496)--cycle;
\path[fill=blue!39!white,draw=blue!39!white,line width=.12pt] (2.116611,-1.258496)--(1.978968,-1.208398)--(2.073119,-1.108865)--(2.217311,-1.152040)--cycle;
\path[fill=blue!28!white,draw=blue!28!white,line width=.12pt] (-3.078787,-0.571335)--(-3.006279,-0.683903)--(-2.933190,-0.829797)--(-3.003935,-0.720748)--cycle;
\path[fill=blue!35!white,draw=blue!35!white,line width=.12pt] (3.003935,-0.720748)--(2.933190,-0.829797)--(3.006279,-0.683903)--(3.078787,-0.571335)--cycle;
\path[fill=blue!14!white,draw=blue!14!white,line width=.12pt] (-2.706147,0.640083)--(-2.844456,0.576059)--(-2.830759,0.436656)--(-2.693116,0.507458)--cycle;
\path[fill=blue!18!white,draw=blue!18!white,line width=.12pt] (2.693116,0.507458)--(2.830759,0.436656)--(2.844456,0.576059)--(2.706147,0.640083)--cycle;
\path[fill=blue!16!white,draw=blue!16!white,line width=.12pt] (-2.093853,0.640083)--(-2.243928,0.679508)--(-2.233123,0.569536)--(-2.083771,0.537466)--cycle;
\path[fill=blue!14!white,draw=blue!14!white,line width=.12pt] (2.083771,0.537466)--(2.233123,0.569536)--(2.243928,0.679508)--(2.093853,0.640083)--cycle;
\path[fill=blue!23!white,draw=blue!23!white,line width=.12pt] (-1.726471,0.299889)--(-1.825482,0.400001)--(-1.799069,0.310969)--(-1.701490,0.215687)--cycle;
\path[fill=blue!16!white,draw=blue!16!white,line width=.12pt] (1.701490,0.215687)--(1.799069,0.310969)--(1.825482,0.400001)--(1.726471,0.299889)--cycle;
\path[fill=blue!37!white,draw=blue!37!white,line width=.12pt] (-1.534468,-0.582930)--(-1.492409,-0.444251)--(-1.424631,-0.515903)--(-1.464780,-0.656601)--cycle;
\path[fill=blue!30!white,draw=blue!30!white,line width=.12pt] (1.464780,-0.656601)--(1.424631,-0.515903)--(1.492409,-0.444251)--(1.534468,-0.582930)--cycle;
\path[fill=blue!14!white,draw=blue!14!white,line width=.12pt] (-2.556072,0.679508)--(-2.706147,0.640083)--(-2.693116,0.507458)--(-2.543764,0.554239)--cycle;
\path[fill=blue!16!white,draw=blue!16!white,line width=.12pt] (2.543764,0.554239)--(2.693116,0.507458)--(2.706147,0.640083)--(2.556072,0.679508)--cycle;
\path[fill=blue!15!white,draw=blue!15!white,line width=.12pt] (-2.243928,0.679508)--(-2.400000,0.692820)--(-2.388443,0.575200)--(-2.233123,0.569536)--cycle;
\path[fill=blue!14!white,draw=blue!14!white,line width=.12pt] (2.233123,0.569536)--(2.388443,0.575200)--(2.400000,0.692820)--(2.243928,0.679508)--cycle;
\path[fill=blue!14!white,draw=blue!14!white,line width=.12pt] (-2.400000,0.692820)--(-2.556072,0.679508)--(-2.543764,0.554239)--(-2.388443,0.575200)--cycle;
\path[fill=blue!15!white,draw=blue!15!white,line width=.12pt] (2.388443,0.575200)--(2.543764,0.554239)--(2.556072,0.679508)--(2.400000,0.692820)--cycle;
\path[fill=blue!32!white,draw=blue!32!white,line width=.12pt] (-1.531105,-0.232228)--(-1.545814,-0.099296)--(-1.492409,-0.173925)--(-1.478207,-0.306147)--cycle;
\path[fill=blue!24!white,draw=blue!24!white,line width=.12pt] (1.478207,-0.306147)--(1.492409,-0.173925)--(1.545814,-0.099296)--(1.531105,-0.232228)--cycle;
\path[fill=blue!18!white,draw=blue!18!white,line width=.12pt] (-3.123988,0.111288)--(-3.169293,-0.020912)--(-3.123436,-0.175483)--(-3.078787,-0.041074)--cycle;
\path[fill=blue!27!white,draw=blue!27!white,line width=.12pt] (3.078787,-0.041074)--(3.123436,-0.175483)--(3.169293,-0.020912)--(3.123988,0.111288)--cycle;
\path[fill=blue!39!white,draw=blue!39!white,line width=.12pt] (-2.361503,-1.168591)--(-2.217311,-1.152040)--(-2.116611,-1.258496)--(-2.254254,-1.281970)--cycle;
\path[fill=blue!40!white,draw=blue!40!white,line width=.12pt] (2.254254,-1.281970)--(2.116611,-1.258496)--(2.217311,-1.152040)--(2.361503,-1.168591)--cycle;
\path[fill=blue!41!white,draw=blue!41!white,line width=.12pt] (-1.724636,-1.036978)--(-1.617721,-0.922243)--(-1.525177,-0.999443)--(-1.625975,-1.119280)--cycle;
\path[fill=blue!37!white,draw=blue!37!white,line width=.12pt] (1.625975,-1.119280)--(1.525177,-0.999443)--(1.617721,-0.922243)--(1.724636,-1.036978)--cycle;
\path[fill=blue!33!white,draw=blue!33!white,line width=.12pt] (-2.837984,-0.920344)--(-2.721975,-0.988910)--(-2.627935,-1.120322)--(-2.739936,-1.057357)--cycle;
\path[fill=blue!38!white,draw=blue!38!white,line width=.12pt] (2.739936,-1.057357)--(2.627935,-1.120322)--(2.721975,-0.988910)--(2.837984,-0.920344)--cycle;
\path[fill=blue!41!white,draw=blue!41!white,line width=.12pt] (-1.617721,-0.922243)--(-1.529978,-0.793806)--(-1.442454,-0.866819)--(-1.525177,-0.999443)--cycle;
\path[fill=blue!35!white,draw=blue!35!white,line width=.12pt] (1.525177,-0.999443)--(1.442454,-0.866819)--(1.529978,-0.793806)--(1.617721,-0.922243)--cycle;
\path[fill=blue!26!white,draw=blue!26!white,line width=.12pt] (-1.628983,0.103118)--(-1.701490,0.215687)--(-1.660123,0.133114)--(-1.589379,0.024064)--cycle;
\path[fill=blue!19!white,draw=blue!19!white,line width=.12pt] (1.589379,0.024064)--(1.660123,0.133114)--(1.701490,0.215687)--(1.628983,0.103118)--cycle;
\path[fill=blue!21!white,draw=blue!21!white,line width=.12pt] (-1.825482,0.400001)--(-1.946127,0.480221)--(-1.917969,0.385305)--(-1.799069,0.310969)--cycle;
\path[fill=blue!15!white,draw=blue!15!white,line width=.12pt] (1.799069,0.310969)--(1.917969,0.385305)--(1.946127,0.480221)--(1.825482,0.400001)--cycle;
\path[fill=blue!25!white,draw=blue!25!white,line width=.12pt] (-3.123436,-0.445808)--(-3.078787,-0.571335)--(-3.003935,-0.720748)--(-3.047499,-0.597387)--cycle;
\path[fill=blue!33!white,draw=blue!33!white,line width=.12pt] (3.047499,-0.597387)--(3.003935,-0.720748)--(3.078787,-0.571335)--(3.123436,-0.445808)--cycle;
\path[fill=blue!42!white,draw=blue!42!white,line width=.12pt] (-1.846614,-1.133600)--(-1.724636,-1.036978)--(-1.625975,-1.119280)--(-1.740975,-1.221724)--cycle;
\path[fill=blue!38!white,draw=blue!38!white,line width=.12pt] (1.740975,-1.221724)--(1.625975,-1.119280)--(1.724636,-1.036978)--(1.846614,-1.133600)--cycle;
\path[fill=blue!35!white,draw=blue!35!white,line width=.12pt] (-1.492409,-0.444251)--(-1.478207,-0.306147)--(-1.411074,-0.377117)--(-1.424631,-0.515903)--cycle;
\path[fill=blue!27!white,draw=blue!27!white,line width=.12pt] (1.424631,-0.515903)--(1.411074,-0.377117)--(1.478207,-0.306147)--(1.492409,-0.444251)--cycle;
\path[fill=blue!39!white,draw=blue!39!white,line width=.12pt] (-1.529978,-0.793806)--(-1.464780,-0.656601)--(-1.380985,-0.726503)--(-1.442454,-0.866819)--cycle;
\path[fill=blue!33!white,draw=blue!33!white,line width=.12pt] (1.442454,-0.866819)--(1.380985,-0.726503)--(1.464780,-0.656601)--(1.529978,-0.793806)--cycle;
\path[fill=blue!16!white,draw=blue!16!white,line width=.12pt] (-3.050416,0.234690)--(-3.123988,0.111288)--(-3.078787,-0.041074)--(-3.006279,0.085917)--cycle;
\path[fill=blue!24!white,draw=blue!24!white,line width=.12pt] (3.006279,0.085917)--(3.078787,-0.041074)--(3.123988,0.111288)--(3.050416,0.234690)--cycle;
\path[fill=blue!41!white,draw=blue!41!white,line width=.12pt] (-1.978968,-1.208398)--(-1.846614,-1.133600)--(-1.740975,-1.221724)--(-1.865758,-1.302838)--cycle;
\path[fill=blue!39!white,draw=blue!39!white,line width=.12pt] (1.865758,-1.302838)--(1.740975,-1.221724)--(1.846614,-1.133600)--(1.978968,-1.208398)--cycle;
\path[fill=blue!37!white,draw=blue!37!white,line width=.12pt] (-2.500154,-1.157881)--(-2.361503,-1.168591)--(-2.254254,-1.281970)--(-2.386608,-1.277917)--cycle;
\path[fill=blue!39!white,draw=blue!39!white,line width=.12pt] (2.386608,-1.277917)--(2.254254,-1.281970)--(2.361503,-1.168591)--(2.500154,-1.157881)--cycle;
\path[fill=blue!30!white,draw=blue!30!white,line width=.12pt] (-1.545814,-0.099296)--(-1.589379,0.024064)--(-1.534468,-0.052668)--(-1.492409,-0.173925)--cycle;
\path[fill=blue!21!white,draw=blue!21!white,line width=.12pt] (1.492409,-0.173925)--(1.534468,-0.052668)--(1.589379,0.024064)--(1.545814,-0.099296)--cycle;
\path[fill=blue!18!white,draw=blue!18!white,line width=.12pt] (-1.946127,0.480221)--(-2.083771,0.537466)--(-2.053620,0.435837)--(-1.917969,0.385305)--cycle;
\path[fill=blue!14!white,draw=blue!14!white,line width=.12pt] (1.917969,0.385305)--(2.053620,0.435837)--(2.083771,0.537466)--(1.946127,0.480221)--cycle;
\path[fill=blue!38!white,draw=blue!38!white,line width=.12pt] (-1.464780,-0.656601)--(-1.424631,-0.515903)--(-1.343133,-0.583889)--(-1.380985,-0.726503)--cycle;
\path[fill=blue!30!white,draw=blue!30!white,line width=.12pt] (1.380985,-0.726503)--(1.343133,-0.583889)--(1.424631,-0.515903)--(1.464780,-0.656601)--cycle;
\path[fill=blue!30!white,draw=blue!30!white,line width=.12pt] (-2.933190,-0.829797)--(-2.837984,-0.920344)--(-2.739936,-1.057357)--(-2.831853,-0.971406)--cycle;
\path[fill=blue!37!white,draw=blue!37!white,line width=.12pt] (2.831853,-0.971406)--(2.739936,-1.057357)--(2.837984,-0.920344)--(2.933190,-0.829797)--cycle;
\path[fill=blue!24!white,draw=blue!24!white,line width=.12pt] (-1.701490,0.215687)--(-1.799069,0.310969)--(-1.755330,0.223661)--(-1.660123,0.133114)--cycle;
\path[fill=blue!17!white,draw=blue!17!white,line width=.12pt] (1.660123,0.133114)--(1.755330,0.223661)--(1.799069,0.310969)--(1.701490,0.215687)--cycle;
\path[fill=blue!22!white,draw=blue!22!white,line width=.12pt] (-3.138513,-0.312145)--(-3.123436,-0.445808)--(-3.047499,-0.597387)--(-3.062209,-0.464455)--cycle;
\path[fill=blue!31!white,draw=blue!31!white,line width=.12pt] (3.062209,-0.464455)--(3.047499,-0.597387)--(3.123436,-0.445808)--(3.138513,-0.312145)--cycle;
\path[fill=blue!33!white,draw=blue!33!white,line width=.12pt] (-1.478207,-0.306147)--(-1.492409,-0.173925)--(-1.424631,-0.245578)--(-1.411074,-0.377117)--cycle;
\path[fill=blue!25!white,draw=blue!25!white,line width=.12pt] (1.411074,-0.377117)--(1.424631,-0.245578)--(1.492409,-0.173925)--(1.478207,-0.306147)--cycle;
\path[fill=blue!15!white,draw=blue!15!white,line width=.12pt] (-2.951405,0.344554)--(-3.050416,0.234690)--(-3.006279,0.085917)--(-2.908701,0.200610)--cycle;
\path[fill=blue!22!white,draw=blue!22!white,line width=.12pt] (2.908701,0.200610)--(3.006279,0.085917)--(3.050416,0.234690)--(2.951405,0.344554)--cycle;
\path[fill=blue!40!white,draw=blue!40!white,line width=.12pt] (-2.116611,-1.258496)--(-1.978968,-1.208398)--(-1.865758,-1.302838)--(-1.995527,-1.359505)--cycle;
\path[fill=blue!40!white,draw=blue!40!white,line width=.12pt] (1.995527,-1.359505)--(1.865758,-1.302838)--(1.978968,-1.208398)--(2.116611,-1.258496)--cycle;
\path[fill=blue!17!white,draw=blue!17!white,line width=.12pt] (-2.083771,0.537466)--(-2.233123,0.569536)--(-2.200811,0.460624)--(-2.053620,0.435837)--cycle;
\path[fill=blue!14!white,draw=blue!14!white,line width=.12pt] (2.053620,0.435837)--(2.200811,0.460624)--(2.233123,0.569536)--(2.083771,0.537466)--cycle;
\path[fill=blue!41!white,draw=blue!41!white,line width=.12pt] (-1.525177,-0.999443)--(-1.442454,-0.866819)--(-1.341037,-0.935191)--(-1.417944,-1.071736)--cycle;
\path[fill=blue!35!white,draw=blue!35!white,line width=.12pt] (1.417944,-1.071736)--(1.341037,-0.935191)--(1.442454,-0.866819)--(1.525177,-0.999443)--cycle;
\path[fill=blue!42!white,draw=blue!42!white,line width=.12pt] (-1.625975,-1.119280)--(-1.525177,-0.999443)--(-1.417944,-1.071736)--(-1.511656,-1.196351)--cycle;
\path[fill=blue!37!white,draw=blue!37!white,line width=.12pt] (1.511656,-1.196351)--(1.417944,-1.071736)--(1.525177,-0.999443)--(1.625975,-1.119280)--cycle;
\path[fill=blue!35!white,draw=blue!35!white,line width=.12pt] (-2.627935,-1.120322)--(-2.500154,-1.157881)--(-2.386608,-1.277917)--(-2.508586,-1.246493)--cycle;
\path[fill=blue!39!white,draw=blue!39!white,line width=.12pt] (2.508586,-1.246493)--(2.386608,-1.277917)--(2.500154,-1.157881)--(2.627935,-1.120322)--cycle;
\path[fill=blue!36!white,draw=blue!36!white,line width=.12pt] (-1.424631,-0.515903)--(-1.411074,-0.377117)--(-1.330351,-0.444456)--(-1.343133,-0.583889)--cycle;
\path[fill=blue!28!white,draw=blue!28!white,line width=.12pt] (1.343133,-0.583889)--(1.330351,-0.444456)--(1.411074,-0.377117)--(1.424631,-0.515903)--cycle;
\path[fill=blue!14!white,draw=blue!14!white,line width=.12pt] (-2.830759,0.436656)--(-2.951405,0.344554)--(-2.908701,0.200610)--(-2.789801,0.298596)--cycle;
\path[fill=blue!19!white,draw=blue!19!white,line width=.12pt] (2.789801,0.298596)--(2.908701,0.200610)--(2.951405,0.344554)--(2.830759,0.436656)--cycle;
\path[fill=blue!27!white,draw=blue!27!white,line width=.12pt] (-1.589379,0.024064)--(-1.660123,0.133114)--(-1.602769,0.052966)--(-1.534468,-0.052668)--cycle;
\path[fill=blue!19!white,draw=blue!19!white,line width=.12pt] (1.534468,-0.052668)--(1.602769,0.052966)--(1.660123,0.133114)--(1.589379,0.024064)--cycle;
\path[fill=blue!40!white,draw=blue!40!white,line width=.12pt] (-1.442454,-0.866819)--(-1.380985,-0.726503)--(-1.283890,-0.791962)--(-1.341037,-0.935191)--cycle;
\path[fill=blue!33!white,draw=blue!33!white,line width=.12pt] (1.341037,-0.935191)--(1.283890,-0.791962)--(1.380985,-0.726503)--(1.442454,-0.866819)--cycle;
\path[fill=blue!15!white,draw=blue!15!white,line width=.12pt] (-2.233123,0.569536)--(-2.388443,0.575200)--(-2.353885,0.458712)--(-2.200811,0.460624)--cycle;
\path[fill=blue!14!white,draw=blue!14!white,line width=.12pt] (2.200811,0.460624)--(2.353885,0.458712)--(2.388443,0.575200)--(2.233123,0.569536)--cycle;
\path[fill=blue!42!white,draw=blue!42!white,line width=.12pt] (-1.740975,-1.221724)--(-1.625975,-1.119280)--(-1.511656,-1.196351)--(-1.618570,-1.304246)--cycle;
\path[fill=blue!38!white,draw=blue!38!white,line width=.12pt] (1.618570,-1.304246)--(1.511656,-1.196351)--(1.625975,-1.119280)--(1.740975,-1.221724)--cycle;
\path[fill=blue!13!white,draw=blue!13!white,line width=.12pt] (-2.693116,0.507458)--(-2.830759,0.436656)--(-2.789801,0.298596)--(-2.654149,0.376111)--cycle;
\path[fill=blue!18!white,draw=blue!18!white,line width=.12pt] (2.654149,0.376111)--(2.789801,0.298596)--(2.830759,0.436656)--(2.693116,0.507458)--cycle;
\path[fill=blue!39!white,draw=blue!39!white,line width=.12pt] (-2.254254,-1.281970)--(-2.116611,-1.258496)--(-1.995527,-1.359505)--(-2.125296,-1.389547)--cycle;
\path[fill=blue!40!white,draw=blue!40!white,line width=.12pt] (2.125296,-1.389547)--(1.995527,-1.359505)--(2.116611,-1.258496)--(2.254254,-1.281970)--cycle;
\path[fill=blue!21!white,draw=blue!21!white,line width=.12pt] (-1.799069,0.310969)--(-1.917969,0.385305)--(-1.871339,0.292227)--(-1.755330,0.223661)--cycle;
\path[fill=blue!16!white,draw=blue!16!white,line width=.12pt] (1.755330,0.223661)--(1.871339,0.292227)--(1.917969,0.385305)--(1.799069,0.310969)--cycle;
\path[fill=blue!20!white,draw=blue!20!white,line width=.12pt] (-3.123436,-0.175483)--(-3.138513,-0.312145)--(-3.062209,-0.464455)--(-3.047499,-0.327062)--cycle;
\path[fill=blue!28!white,draw=blue!28!white,line width=.12pt] (3.047499,-0.327062)--(3.062209,-0.464455)--(3.138513,-0.312145)--(3.123436,-0.175483)--cycle;
\path[fill=blue!14!white,draw=blue!14!white,line width=.12pt] (-2.388443,0.575200)--(-2.543764,0.554239)--(-2.506958,0.430175)--(-2.353885,0.458712)--cycle;
\path[fill=blue!15!white,draw=blue!15!white,line width=.12pt] (2.353885,0.458712)--(2.506958,0.430175)--(2.543764,0.554239)--(2.388443,0.575200)--cycle;
\path[fill=blue!13!white,draw=blue!13!white,line width=.12pt] (-2.543764,0.554239)--(-2.693116,0.507458)--(-2.654149,0.376111)--(-2.506958,0.430175)--cycle;
\path[fill=blue!16!white,draw=blue!16!white,line width=.12pt] (2.506958,0.430175)--(2.654149,0.376111)--(2.693116,0.507458)--(2.543764,0.554239)--cycle;
\path[fill=blue!27!white,draw=blue!27!white,line width=.12pt] (-3.003935,-0.720748)--(-2.933190,-0.829797)--(-2.831853,-0.971406)--(-2.900154,-0.865772)--cycle;
\path[fill=blue!35!white,draw=blue!35!white,line width=.12pt] (2.900154,-0.865772)--(2.831853,-0.971406)--(2.933190,-0.829797)--(3.003935,-0.720748)--cycle;
\path[fill=blue!30!white,draw=blue!30!white,line width=.12pt] (-1.492409,-0.173925)--(-1.534468,-0.052668)--(-1.464780,-0.126340)--(-1.424631,-0.245578)--cycle;
\path[fill=blue!22!white,draw=blue!22!white,line width=.12pt] (1.424631,-0.245578)--(1.464780,-0.126340)--(1.534468,-0.052668)--(1.492409,-0.173925)--cycle;
\path[fill=blue!38!white,draw=blue!38!white,line width=.12pt] (-1.380985,-0.726503)--(-1.343133,-0.583889)--(-1.248699,-0.647553)--(-1.283890,-0.791962)--cycle;
\path[fill=blue!31!white,draw=blue!31!white,line width=.12pt] (1.283890,-0.791962)--(1.248699,-0.647553)--(1.343133,-0.583889)--(1.380985,-0.726503)--cycle;
\path[fill=blue!41!white,draw=blue!41!white,line width=.12pt] (-1.865758,-1.302838)--(-1.740975,-1.221724)--(-1.618570,-1.304246)--(-1.734580,-1.391274)--cycle;
\path[fill=blue!39!white,draw=blue!39!white,line width=.12pt] (1.734580,-1.391274)--(1.618570,-1.304246)--(1.740975,-1.221724)--(1.865758,-1.302838)--cycle;
\path[fill=blue!33!white,draw=blue!33!white,line width=.12pt] (-1.411074,-0.377117)--(-1.424631,-0.245578)--(-1.343133,-0.313564)--(-1.330351,-0.444456)--cycle;
\path[fill=blue!25!white,draw=blue!25!white,line width=.12pt] (1.330351,-0.444456)--(1.343133,-0.313564)--(1.424631,-0.245578)--(1.411074,-0.377117)--cycle;
\path[fill=blue!32!white,draw=blue!32!white,line width=.12pt] (-2.739936,-1.057357)--(-2.627935,-1.120322)--(-2.508586,-1.246493)--(-2.615501,-1.188905)--cycle;
\path[fill=blue!37!white,draw=blue!37!white,line width=.12pt] (2.615501,-1.188905)--(2.508586,-1.246493)--(2.627935,-1.120322)--(2.739936,-1.057357)--cycle;
\path[fill=blue!24!white,draw=blue!24!white,line width=.12pt] (-1.660123,0.133114)--(-1.755330,0.223661)--(-1.694686,0.138917)--(-1.602769,0.052966)--cycle;
\path[fill=blue!17!white,draw=blue!17!white,line width=.12pt] (1.602769,0.052966)--(1.694686,0.138917)--(1.755330,0.223661)--(1.660123,0.133114)--cycle;
\path[fill=blue!19!white,draw=blue!19!white,line width=.12pt] (-1.917969,0.385305)--(-2.053620,0.435837)--(-2.003693,0.336176)--(-1.871339,0.292227)--cycle;
\path[fill=blue!15!white,draw=blue!15!white,line width=.12pt] (1.871339,0.292227)--(2.003693,0.336176)--(2.053620,0.435837)--(1.917969,0.385305)--cycle;
\path[fill=blue!18!white,draw=blue!18!white,line width=.12pt] (-3.078787,-0.041074)--(-3.123436,-0.175483)--(-3.047499,-0.327062)--(-3.003935,-0.190486)--cycle;
\path[fill=blue!26!white,draw=blue!26!white,line width=.12pt] (3.003935,-0.190486)--(3.047499,-0.327062)--(3.123436,-0.175483)--(3.078787,-0.041074)--cycle;
\path[fill=blue!41!white,draw=blue!41!white,line width=.12pt] (-1.417944,-1.071736)--(-1.341037,-0.935191)--(-1.226706,-0.998263)--(-1.297056,-1.138426)--cycle;
\path[fill=blue!36!white,draw=blue!36!white,line width=.12pt] (1.297056,-1.138426)--(1.226706,-0.998263)--(1.341037,-0.935191)--(1.417944,-1.071736)--cycle;
\path[fill=blue!37!white,draw=blue!37!white,line width=.12pt] (-2.386608,-1.277917)--(-2.254254,-1.281970)--(-2.125296,-1.389547)--(-2.250079,-1.391810)--cycle;
\path[fill=blue!39!white,draw=blue!39!white,line width=.12pt] (2.250079,-1.391810)--(2.125296,-1.389547)--(2.254254,-1.281970)--(2.386608,-1.277917)--cycle;
\path[fill=blue!36!white,draw=blue!36!white,line width=.12pt] (-1.343133,-0.583889)--(-1.330351,-0.444456)--(-1.236817,-0.507515)--(-1.248699,-0.647553)--cycle;
\path[fill=blue!29!white,draw=blue!29!white,line width=.12pt] (1.248699,-0.647553)--(1.236817,-0.507515)--(1.330351,-0.444456)--(1.343133,-0.583889)--cycle;
\path[fill=blue!42!white,draw=blue!42!white,line width=.12pt] (-1.511656,-1.196351)--(-1.417944,-1.071736)--(-1.297056,-1.138426)--(-1.382778,-1.267448)--cycle;
\path[fill=blue!38!white,draw=blue!38!white,line width=.12pt] (1.382778,-1.267448)--(1.297056,-1.138426)--(1.417944,-1.071736)--(1.511656,-1.196351)--cycle;
\path[fill=blue!40!white,draw=blue!40!white,line width=.12pt] (-1.341037,-0.935191)--(-1.283890,-0.791962)--(-1.174431,-0.852346)--(-1.226706,-0.998263)--cycle;
\path[fill=blue!34!white,draw=blue!34!white,line width=.12pt] (1.226706,-0.998263)--(1.174431,-0.852346)--(1.283890,-0.791962)--(1.341037,-0.935191)--cycle;
\path[fill=blue!40!white,draw=blue!40!white,line width=.12pt] (-1.995527,-1.359505)--(-1.865758,-1.302838)--(-1.734580,-1.391274)--(-1.855225,-1.454092)--cycle;
\path[fill=blue!40!white,draw=blue!40!white,line width=.12pt] (1.855225,-1.454092)--(1.734580,-1.391274)--(1.865758,-1.302838)--(1.995527,-1.359505)--cycle;
\path[fill=blue!24!white,draw=blue!24!white,line width=.12pt] (-3.047499,-0.597387)--(-3.003935,-0.720748)--(-2.900154,-0.865772)--(-2.942213,-0.744515)--cycle;
\path[fill=blue!33!white,draw=blue!33!white,line width=.12pt] (2.942213,-0.744515)--(2.900154,-0.865772)--(3.003935,-0.720748)--(3.047499,-0.597387)--cycle;
\path[fill=blue!27!white,draw=blue!27!white,line width=.12pt] (-1.534468,-0.052668)--(-1.602769,0.052966)--(-1.529978,-0.023985)--(-1.464780,-0.126340)--cycle;
\path[fill=blue!20!white,draw=blue!20!white,line width=.12pt] (1.464780,-0.126340)--(1.529978,-0.023985)--(1.602769,0.052966)--(1.534468,-0.052668)--cycle;
\path[fill=blue!42!white,draw=blue!42!white,line width=.12pt] (-1.618570,-1.304246)--(-1.511656,-1.196351)--(-1.382778,-1.267448)--(-1.480578,-1.380371)--cycle;
\path[fill=blue!39!white,draw=blue!39!white,line width=.12pt] (1.480578,-1.380371)--(1.382778,-1.267448)--(1.511656,-1.196351)--(1.618570,-1.304246)--cycle;
\path[fill=blue!38!white,draw=blue!38!white,line width=.12pt] (-1.283890,-0.791962)--(-1.248699,-0.647553)--(-1.142240,-0.706283)--(-1.174431,-0.852346)--cycle;
\path[fill=blue!32!white,draw=blue!32!white,line width=.12pt] (1.174431,-0.852346)--(1.142240,-0.706283)--(1.248699,-0.647553)--(1.283890,-0.791962)--cycle;
\path[fill=blue!17!white,draw=blue!17!white,line width=.12pt] (-2.053620,0.435837)--(-2.200811,0.460624)--(-2.147305,0.353819)--(-2.003693,0.336176)--cycle;
\path[fill=blue!14!white,draw=blue!14!white,line width=.12pt] (2.003693,0.336176)--(2.147305,0.353819)--(2.200811,0.460624)--(2.053620,0.435837)--cycle;
\path[fill=blue!16!white,draw=blue!16!white,line width=.12pt] (-3.006279,0.085917)--(-3.078787,-0.041074)--(-3.003935,-0.190486)--(-2.933190,-0.059976)--cycle;
\path[fill=blue!23!white,draw=blue!23!white,line width=.12pt] (2.933190,-0.059976)--(3.003935,-0.190486)--(3.078787,-0.041074)--(3.006279,0.085917)--cycle;
\path[fill=blue!30!white,draw=blue!30!white,line width=.12pt] (-1.424631,-0.245578)--(-1.464780,-0.126340)--(-1.380985,-0.196241)--(-1.343133,-0.313564)--cycle;
\path[fill=blue!23!white,draw=blue!23!white,line width=.12pt] (1.343133,-0.313564)--(1.380985,-0.196241)--(1.464780,-0.126340)--(1.424631,-0.245578)--cycle;
\path[fill=blue!21!white,draw=blue!21!white,line width=.12pt] (-1.755330,0.223661)--(-1.871339,0.292227)--(-1.806687,0.201882)--(-1.694686,0.138917)--cycle;
\path[fill=blue!16!white,draw=blue!16!white,line width=.12pt] (1.694686,0.138917)--(1.806687,0.201882)--(1.871339,0.292227)--(1.755330,0.223661)--cycle;
\path[fill=blue!30!white,draw=blue!30!white,line width=.12pt] (-2.831853,-0.971406)--(-2.739936,-1.057357)--(-2.615501,-1.188905)--(-2.703244,-1.107367)--cycle;
\path[fill=blue!36!white,draw=blue!36!white,line width=.12pt] (2.703244,-1.107367)--(2.615501,-1.188905)--(2.739936,-1.057357)--(2.831853,-0.971406)--cycle;
\path[fill=blue!33!white,draw=blue!33!white,line width=.12pt] (-1.330351,-0.444456)--(-1.343133,-0.313564)--(-1.248699,-0.377228)--(-1.236817,-0.507515)--cycle;
\path[fill=blue!26!white,draw=blue!26!white,line width=.12pt] (1.236817,-0.507515)--(1.248699,-0.377228)--(1.343133,-0.313564)--(1.330351,-0.444456)--cycle;
\path[fill=blue!41!white,draw=blue!41!white,line width=.12pt] (-1.734580,-1.391274)--(-1.618570,-1.304246)--(-1.480578,-1.380371)--(-1.586697,-1.472856)--cycle;
\path[fill=blue!39!white,draw=blue!39!white,line width=.12pt] (1.586697,-1.472856)--(1.480578,-1.380371)--(1.618570,-1.304246)--(1.734580,-1.391274)--cycle;
\path[fill=blue!39!white,draw=blue!39!white,line width=.12pt] (-2.125296,-1.389547)--(-1.995527,-1.359505)--(-1.855225,-1.454092)--(-1.975871,-1.490286)--cycle;
\path[fill=blue!40!white,draw=blue!40!white,line width=.12pt] (1.975871,-1.490286)--(1.855225,-1.454092)--(1.995527,-1.359505)--(2.125296,-1.389547)--cycle;
\path[fill=blue!35!white,draw=blue!35!white,line width=.12pt] (-2.508586,-1.246493)--(-2.386608,-1.277917)--(-2.250079,-1.391810)--(-2.365079,-1.366207)--cycle;
\path[fill=blue!38!white,draw=blue!38!white,line width=.12pt] (2.365079,-1.366207)--(2.250079,-1.391810)--(2.386608,-1.277917)--(2.508586,-1.246493)--cycle;
\path[fill=blue!15!white,draw=blue!15!white,line width=.12pt] (-2.200811,0.460624)--(-2.353885,0.458712)--(-2.296657,0.344479)--(-2.147305,0.353819)--cycle;
\path[fill=blue!14!white,draw=blue!14!white,line width=.12pt] (2.147305,0.353819)--(2.296657,0.344479)--(2.353885,0.458712)--(2.200811,0.460624)--cycle;
\path[fill=blue!14!white,draw=blue!14!white,line width=.12pt] (-2.908701,0.200610)--(-3.006279,0.085917)--(-2.933190,-0.059976)--(-2.837984,0.059451)--cycle;
\path[fill=blue!21!white,draw=blue!21!white,line width=.12pt] (2.837984,0.059451)--(2.933190,-0.059976)--(3.006279,0.085917)--(2.908701,0.200610)--cycle;
\path[fill=blue!41!white,draw=blue!41!white,line width=.12pt] (-1.297056,-1.138426)--(-1.226706,-0.998263)--(-1.100561,-1.055429)--(-1.163677,-1.198870)--cycle;
\path[fill=blue!37!white,draw=blue!37!white,line width=.12pt] (1.163677,-1.198870)--(1.100561,-1.055429)--(1.226706,-0.998263)--(1.297056,-1.138426)--cycle;
\path[fill=blue!22!white,draw=blue!22!white,line width=.12pt] (-3.062209,-0.464455)--(-3.047499,-0.597387)--(-2.942213,-0.744515)--(-2.956415,-0.612293)--cycle;
\path[fill=blue!30!white,draw=blue!30!white,line width=.12pt] (2.956415,-0.612293)--(2.942213,-0.744515)--(3.047499,-0.597387)--(3.062209,-0.464455)--cycle;
\path[fill=blue!36!white,draw=blue!36!white,line width=.12pt] (-1.248699,-0.647553)--(-1.236817,-0.507515)--(-1.131371,-0.565685)--(-1.142240,-0.706283)--cycle;
\path[fill=blue!29!white,draw=blue!29!white,line width=.12pt] (1.142240,-0.706283)--(1.131371,-0.565685)--(1.236817,-0.507515)--(1.248699,-0.647553)--cycle;
\path[fill=blue!40!white,draw=blue!40!white,line width=.12pt] (-1.226706,-0.998263)--(-1.174431,-0.852346)--(-1.053662,-0.907076)--(-1.100561,-1.055429)--cycle;
\path[fill=blue!35!white,draw=blue!35!white,line width=.12pt] (1.100561,-1.055429)--(1.053662,-0.907076)--(1.174431,-0.852346)--(1.226706,-0.998263)--cycle;
\path[fill=blue!42!white,draw=blue!42!white,line width=.12pt] (-1.382778,-1.267448)--(-1.297056,-1.138426)--(-1.163677,-1.198870)--(-1.240584,-1.331887)--cycle;
\path[fill=blue!38!white,draw=blue!38!white,line width=.12pt] (1.240584,-1.331887)--(1.163677,-1.198870)--(1.297056,-1.138426)--(1.382778,-1.267448)--cycle;
\path[fill=blue!24!white,draw=blue!24!white,line width=.12pt] (-1.602769,0.052966)--(-1.694686,0.138917)--(-1.617721,0.057553)--(-1.529978,-0.023985)--cycle;
\path[fill=blue!18!white,draw=blue!18!white,line width=.12pt] (1.529978,-0.023985)--(1.617721,0.057553)--(1.694686,0.138917)--(1.602769,0.052966)--cycle;
\path[fill=blue!14!white,draw=blue!14!white,line width=.12pt] (-2.353885,0.458712)--(-2.506958,0.430175)--(-2.446009,0.308514)--(-2.296657,0.344479)--cycle;
\path[fill=blue!15!white,draw=blue!15!white,line width=.12pt] (2.296657,0.344479)--(2.446009,0.308514)--(2.506958,0.430175)--(2.353885,0.458712)--cycle;
\path[fill=blue!13!white,draw=blue!13!white,line width=.12pt] (-2.789801,0.298596)--(-2.908701,0.200610)--(-2.837984,0.059451)--(-2.721975,0.163208)--cycle;
\path[fill=blue!19!white,draw=blue!19!white,line width=.12pt] (2.721975,0.163208)--(2.837984,0.059451)--(2.908701,0.200610)--(2.789801,0.298596)--cycle;
\path[fill=blue!13!white,draw=blue!13!white,line width=.12pt] (-2.506958,0.430175)--(-2.654149,0.376111)--(-2.589621,0.247306)--(-2.446009,0.308514)--cycle;
\path[fill=blue!16!white,draw=blue!16!white,line width=.12pt] (2.446009,0.308514)--(2.589621,0.247306)--(2.654149,0.376111)--(2.506958,0.430175)--cycle;
\path[fill=blue!38!white,draw=blue!38!white,line width=.12pt] (-1.174431,-0.852346)--(-1.142240,-0.706283)--(-1.024781,-0.759512)--(-1.053662,-0.907076)--cycle;
\path[fill=blue!33!white,draw=blue!33!white,line width=.12pt] (1.053662,-0.907076)--(1.024781,-0.759512)--(1.142240,-0.706283)--(1.174431,-0.852346)--cycle;
\path[fill=blue!13!white,draw=blue!13!white,line width=.12pt] (-2.654149,0.376111)--(-2.789801,0.298596)--(-2.721975,0.163208)--(-2.589621,0.247306)--cycle;
\path[fill=blue!17!white,draw=blue!17!white,line width=.12pt] (2.589621,0.247306)--(2.721975,0.163208)--(2.789801,0.298596)--(2.654149,0.376111)--cycle;
\path[fill=blue!19!white,draw=blue!19!white,line width=.12pt] (-1.871339,0.292227)--(-2.003693,0.336176)--(-1.934468,0.239441)--(-1.806687,0.201882)--cycle;
\path[fill=blue!15!white,draw=blue!15!white,line width=.12pt] (1.806687,0.201882)--(1.934468,0.239441)--(2.003693,0.336176)--(1.871339,0.292227)--cycle;
\path[fill=blue!28!white,draw=blue!28!white,line width=.12pt] (-1.464780,-0.126340)--(-1.529978,-0.023985)--(-1.442454,-0.096998)--(-1.380985,-0.196241)--cycle;
\path[fill=blue!21!white,draw=blue!21!white,line width=.12pt] (1.380985,-0.196241)--(1.442454,-0.096998)--(1.529978,-0.023985)--(1.464780,-0.126340)--cycle;
\path[fill=blue!42!white,draw=blue!42!white,line width=.12pt] (-1.480578,-1.380371)--(-1.382778,-1.267448)--(-1.240584,-1.331887)--(-1.328326,-1.449368)--cycle;
\path[fill=blue!39!white,draw=blue!39!white,line width=.12pt] (1.328326,-1.449368)--(1.240584,-1.331887)--(1.382778,-1.267448)--(1.480578,-1.380371)--cycle;
\path[fill=blue!40!white,draw=blue!40!white,line width=.12pt] (-1.855225,-1.454092)--(-1.734580,-1.391274)--(-1.586697,-1.472856)--(-1.697056,-1.541348)--cycle;
\path[fill=blue!40!white,draw=blue!40!white,line width=.12pt] (1.697056,-1.541348)--(1.586697,-1.472856)--(1.734580,-1.391274)--(1.855225,-1.454092)--cycle;
\path[fill=blue!27!white,draw=blue!27!white,line width=.12pt] (-2.900154,-0.865772)--(-2.831853,-0.971406)--(-2.703244,-1.107367)--(-2.768442,-1.005012)--cycle;
\path[fill=blue!34!white,draw=blue!34!white,line width=.12pt] (2.768442,-1.005012)--(2.703244,-1.107367)--(2.831853,-0.971406)--(2.900154,-0.865772)--cycle;
\path[fill=blue!31!white,draw=blue!31!white,line width=.12pt] (-1.343133,-0.313564)--(-1.380985,-0.196241)--(-1.283890,-0.261700)--(-1.248699,-0.377228)--cycle;
\path[fill=blue!24!white,draw=blue!24!white,line width=.12pt] (1.248699,-0.377228)--(1.283890,-0.261700)--(1.380985,-0.196241)--(1.343133,-0.313564)--cycle;
\path[fill=blue!19!white,draw=blue!19!white,line width=.12pt] (-3.047499,-0.327062)--(-3.062209,-0.464455)--(-2.956415,-0.612293)--(-2.942213,-0.474190)--cycle;
\path[fill=blue!28!white,draw=blue!28!white,line width=.12pt] (2.942213,-0.474190)--(2.956415,-0.612293)--(3.062209,-0.464455)--(3.047499,-0.327062)--cycle;
\path[fill=blue!34!white,draw=blue!34!white,line width=.12pt] (-1.236817,-0.507515)--(-1.248699,-0.377228)--(-1.142240,-0.435958)--(-1.131371,-0.565685)--cycle;
\path[fill=blue!27!white,draw=blue!27!white,line width=.12pt] (1.131371,-0.565685)--(1.142240,-0.435958)--(1.248699,-0.377228)--(1.236817,-0.507515)--cycle;
\path[fill=blue!37!white,draw=blue!37!white,line width=.12pt] (-2.250079,-1.391810)--(-2.125296,-1.389547)--(-1.975871,-1.490286)--(-2.091880,-1.498463)--cycle;
\path[fill=blue!39!white,draw=blue!39!white,line width=.12pt] (2.091880,-1.498463)--(1.975871,-1.490286)--(2.125296,-1.389547)--(2.250079,-1.391810)--cycle;
\path[fill=blue!32!white,draw=blue!32!white,line width=.12pt] (-2.615501,-1.188905)--(-2.508586,-1.246493)--(-2.365079,-1.366207)--(-2.465877,-1.313721)--cycle;
\path[fill=blue!37!white,draw=blue!37!white,line width=.12pt] (2.465877,-1.313721)--(2.365079,-1.366207)--(2.508586,-1.246493)--(2.615501,-1.188905)--cycle;
\path[fill=blue!41!white,draw=blue!41!white,line width=.12pt] (-1.163677,-1.198870)--(-1.100561,-1.055429)--(-0.963817,-1.106137)--(-1.019091,-1.252486)--cycle;
\path[fill=blue!37!white,draw=blue!37!white,line width=.12pt] (1.019091,-1.252486)--(0.963817,-1.106137)--(1.100561,-1.055429)--(1.163677,-1.198870)--cycle;
\path[fill=blue!36!white,draw=blue!36!white,line width=.12pt] (-1.142240,-0.706283)--(-1.131371,-0.565685)--(-1.015029,-0.618408)--(-1.024781,-0.759512)--cycle;
\path[fill=blue!30!white,draw=blue!30!white,line width=.12pt] (1.024781,-0.759512)--(1.015029,-0.618408)--(1.131371,-0.565685)--(1.142240,-0.706283)--cycle;
\path[fill=blue!40!white,draw=blue!40!white,line width=.12pt] (-1.100561,-1.055429)--(-1.053662,-0.907076)--(-0.922745,-0.955623)--(-0.963817,-1.106137)--cycle;
\path[fill=blue!35!white,draw=blue!35!white,line width=.12pt] (0.963817,-1.106137)--(0.922745,-0.955623)--(1.053662,-0.907076)--(1.100561,-1.055429)--cycle;
\path[fill=blue!22!white,draw=blue!22!white,line width=.12pt] (-1.694686,0.138917)--(-1.806687,0.201882)--(-1.724636,0.115141)--(-1.617721,0.057553)--cycle;
\path[fill=blue!17!white,draw=blue!17!white,line width=.12pt] (1.617721,0.057553)--(1.724636,0.115141)--(1.806687,0.201882)--(1.694686,0.138917)--cycle;
\path[fill=blue!42!white,draw=blue!42!white,line width=.12pt] (-1.240584,-1.331887)--(-1.163677,-1.198870)--(-1.019091,-1.252486)--(-1.086442,-1.389047)--cycle;
\path[fill=blue!38!white,draw=blue!38!white,line width=.12pt] (1.086442,-1.389047)--(1.019091,-1.252486)--(1.163677,-1.198870)--(1.240584,-1.331887)--cycle;
\path[fill=blue!41!white,draw=blue!41!white,line width=.12pt] (-1.586697,-1.472856)--(-1.480578,-1.380371)--(-1.328326,-1.449368)--(-1.423533,-1.546798)--cycle;
\path[fill=blue!40!white,draw=blue!40!white,line width=.12pt] (1.423533,-1.546798)--(1.328326,-1.449368)--(1.480578,-1.380371)--(1.586697,-1.472856)--cycle;
\path[fill=blue!17!white,draw=blue!17!white,line width=.12pt] (-2.003693,0.336176)--(-2.147305,0.353819)--(-2.073119,0.250151)--(-1.934468,0.239441)--cycle;
\path[fill=blue!14!white,draw=blue!14!white,line width=.12pt] (1.934468,0.239441)--(2.073119,0.250151)--(2.147305,0.353819)--(2.003693,0.336176)--cycle;
\path[fill=blue!38!white,draw=blue!38!white,line width=.12pt] (-1.053662,-0.907076)--(-1.024781,-0.759512)--(-0.897452,-0.806729)--(-0.922745,-0.955623)--cycle;
\path[fill=blue!33!white,draw=blue!33!white,line width=.12pt] (0.922745,-0.955623)--(0.897452,-0.806729)--(1.024781,-0.759512)--(1.053662,-0.907076)--cycle;
\path[fill=blue!17!white,draw=blue!17!white,line width=.12pt] (-3.003935,-0.190486)--(-3.047499,-0.327062)--(-2.942213,-0.474190)--(-2.900154,-0.335511)--cycle;
\path[fill=blue!25!white,draw=blue!25!white,line width=.12pt] (2.900154,-0.335511)--(2.942213,-0.474190)--(3.047499,-0.327062)--(3.003935,-0.190486)--cycle;
\path[fill=blue!39!white,draw=blue!39!white,line width=.12pt] (-1.975871,-1.490286)--(-1.855225,-1.454092)--(-1.697056,-1.541348)--(-1.807416,-1.583216)--cycle;
\path[fill=blue!39!white,draw=blue!39!white,line width=.12pt] (1.807416,-1.583216)--(1.697056,-1.541348)--(1.855225,-1.454092)--(1.975871,-1.490286)--cycle;
\path[fill=blue!25!white,draw=blue!25!white,line width=.12pt] (-1.529978,-0.023985)--(-1.617721,0.057553)--(-1.525177,-0.019647)--(-1.442454,-0.096998)--cycle;
\path[fill=blue!19!white,draw=blue!19!white,line width=.12pt] (1.442454,-0.096998)--(1.525177,-0.019647)--(1.617721,0.057553)--(1.529978,-0.023985)--cycle;
\path[fill=blue!24!white,draw=blue!24!white,line width=.12pt] (-2.942213,-0.744515)--(-2.900154,-0.865772)--(-2.768442,-1.005012)--(-2.808591,-0.885774)--cycle;
\path[fill=blue!32!white,draw=blue!32!white,line width=.12pt] (2.808591,-0.885774)--(2.768442,-1.005012)--(2.900154,-0.865772)--(2.942213,-0.744515)--cycle;
\path[fill=blue!42!white,draw=blue!42!white,line width=.12pt] (-1.328326,-1.449368)--(-1.240584,-1.331887)--(-1.086442,-1.389047)--(-1.163283,-1.510570)--cycle;
\path[fill=blue!39!white,draw=blue!39!white,line width=.12pt] (1.163283,-1.510570)--(1.086442,-1.389047)--(1.240584,-1.331887)--(1.328326,-1.449368)--cycle;
\path[fill=blue!28!white,draw=blue!28!white,line width=.12pt] (-1.380985,-0.196241)--(-1.442454,-0.096998)--(-1.341037,-0.165370)--(-1.283890,-0.261700)--cycle;
\path[fill=blue!21!white,draw=blue!21!white,line width=.12pt] (1.283890,-0.261700)--(1.341037,-0.165370)--(1.442454,-0.096998)--(1.380985,-0.196241)--cycle;
\path[fill=blue!31!white,draw=blue!31!white,line width=.12pt] (-1.248699,-0.377228)--(-1.283890,-0.261700)--(-1.174431,-0.322085)--(-1.142240,-0.435958)--cycle;
\path[fill=blue!24!white,draw=blue!24!white,line width=.12pt] (1.142240,-0.435958)--(1.174431,-0.322085)--(1.283890,-0.261700)--(1.248699,-0.377228)--cycle;
\path[fill=blue!34!white,draw=blue!34!white,line width=.12pt] (-1.131371,-0.565685)--(-1.142240,-0.435958)--(-1.024781,-0.489187)--(-1.015029,-0.618408)--cycle;
\path[fill=blue!28!white,draw=blue!28!white,line width=.12pt] (1.015029,-0.618408)--(1.024781,-0.489187)--(1.142240,-0.435958)--(1.131371,-0.565685)--cycle;
\path[fill=blue!41!white,draw=blue!41!white,line width=.12pt] (-1.019091,-1.252486)--(-0.963817,-1.106137)--(-0.817791,-1.149900)--(-0.864690,-1.298758)--cycle;
\path[fill=blue!38!white,draw=blue!38!white,line width=.12pt] (0.864690,-1.298758)--(0.817791,-1.149900)--(0.963817,-1.106137)--(1.019091,-1.252486)--cycle;
\path[fill=blue!40!white,draw=blue!40!white,line width=.12pt] (-0.963817,-1.106137)--(-0.922745,-0.955623)--(-0.782941,-0.997521)--(-0.817791,-1.149900)--cycle;
\path[fill=blue!36!white,draw=blue!36!white,line width=.12pt] (0.817791,-1.149900)--(0.782941,-0.997521)--(0.922745,-0.955623)--(0.963817,-1.106137)--cycle;
\path[fill=blue!15!white,draw=blue!15!white,line width=.12pt] (-2.147305,0.353819)--(-2.296657,0.344479)--(-2.217311,0.233600)--(-2.073119,0.250151)--cycle;
\path[fill=blue!14!white,draw=blue!14!white,line width=.12pt] (2.073119,0.250151)--(2.217311,0.233600)--(2.296657,0.344479)--(2.147305,0.353819)--cycle;
\path[fill=blue!36!white,draw=blue!36!white,line width=.12pt] (-1.024781,-0.759512)--(-1.015029,-0.618408)--(-0.888912,-0.665176)--(-0.897452,-0.806729)--cycle;
\path[fill=blue!31!white,draw=blue!31!white,line width=.12pt] (0.897452,-0.806729)--(0.888912,-0.665176)--(1.015029,-0.618408)--(1.024781,-0.759512)--cycle;
\path[fill=blue!35!white,draw=blue!35!white,line width=.12pt] (-2.365079,-1.366207)--(-2.250079,-1.391810)--(-2.091880,-1.498463)--(-2.198794,-1.478311)--cycle;
\path[fill=blue!38!white,draw=blue!38!white,line width=.12pt] (2.198794,-1.478311)--(2.091880,-1.498463)--(2.250079,-1.391810)--(2.365079,-1.366207)--cycle;
\path[fill=blue!29!white,draw=blue!29!white,line width=.12pt] (-2.703244,-1.107367)--(-2.615501,-1.188905)--(-2.465877,-1.313721)--(-2.548600,-1.236371)--cycle;
\path[fill=blue!35!white,draw=blue!35!white,line width=.12pt] (2.548600,-1.236371)--(2.465877,-1.313721)--(2.615501,-1.188905)--(2.703244,-1.107367)--cycle;
\path[fill=blue!15!white,draw=blue!15!white,line width=.12pt] (-2.933190,-0.059976)--(-3.003935,-0.190486)--(-2.900154,-0.335511)--(-2.831853,-0.201586)--cycle;
\path[fill=blue!23!white,draw=blue!23!white,line width=.12pt] (2.831853,-0.201586)--(2.900154,-0.335511)--(3.003935,-0.190486)--(2.933190,-0.059976)--cycle;
\path[fill=blue!40!white,draw=blue!40!white,line width=.12pt] (-1.697056,-1.541348)--(-1.586697,-1.472856)--(-1.423533,-1.546798)--(-1.522544,-1.620433)--cycle;
\path[fill=blue!40!white,draw=blue!40!white,line width=.12pt] (1.522544,-1.620433)--(1.423533,-1.546798)--(1.586697,-1.472856)--(1.697056,-1.541348)--cycle;
\path[fill=blue!42!white,draw=blue!42!white,line width=.12pt] (-1.086442,-1.389047)--(-1.019091,-1.252486)--(-0.864690,-1.298758)--(-0.921837,-1.438377)--cycle;
\path[fill=blue!39!white,draw=blue!39!white,line width=.12pt] (0.921837,-1.438377)--(0.864690,-1.298758)--(1.019091,-1.252486)--(1.086442,-1.389047)--cycle;
\path[fill=blue!38!white,draw=blue!38!white,line width=.12pt] (-0.922745,-0.955623)--(-0.897452,-0.806729)--(-0.761481,-0.847478)--(-0.782941,-0.997521)--cycle;
\path[fill=blue!34!white,draw=blue!34!white,line width=.12pt] (0.782941,-0.997521)--(0.761481,-0.847478)--(0.897452,-0.806729)--(0.922745,-0.955623)--cycle;
\path[fill=blue!19!white,draw=blue!19!white,line width=.12pt] (-1.806687,0.201882)--(-1.934468,0.239441)--(-1.846614,0.146565)--(-1.724636,0.115141)--cycle;
\path[fill=blue!15!white,draw=blue!15!white,line width=.12pt] (1.724636,0.115141)--(1.846614,0.146565)--(1.934468,0.239441)--(1.806687,0.201882)--cycle;
\path[fill=blue!14!white,draw=blue!14!white,line width=.12pt] (-2.296657,0.344479)--(-2.446009,0.308514)--(-2.361503,0.190425)--(-2.217311,0.233600)--cycle;
\path[fill=blue!15!white,draw=blue!15!white,line width=.12pt] (2.217311,0.233600)--(2.361503,0.190425)--(2.446009,0.308514)--(2.296657,0.344479)--cycle;
\path[fill=blue!41!white,draw=blue!41!white,line width=.12pt] (-1.423533,-1.546798)--(-1.328326,-1.449368)--(-1.163283,-1.510570)--(-1.246659,-1.612387)--cycle;
\path[fill=blue!40!white,draw=blue!40!white,line width=.12pt] (1.246659,-1.612387)--(1.163283,-1.510570)--(1.328326,-1.449368)--(1.423533,-1.546798)--cycle;
\path[fill=blue!14!white,draw=blue!14!white,line width=.12pt] (-2.837984,0.059451)--(-2.933190,-0.059976)--(-2.831853,-0.201586)--(-2.739936,-0.077561)--cycle;
\path[fill=blue!21!white,draw=blue!21!white,line width=.12pt] (2.739936,-0.077561)--(2.831853,-0.201586)--(2.933190,-0.059976)--(2.837984,0.059451)--cycle;
\path[fill=blue!40!white,draw=blue!40!white,line width=.12pt] (-0.817791,-1.149900)--(-0.782941,-0.997521)--(-0.635598,-1.032365)--(-0.663889,-1.186295)--cycle;
\path[fill=blue!37!white,draw=blue!37!white,line width=.12pt] (0.663889,-1.186295)--(0.635598,-1.032365)--(0.782941,-0.997521)--(0.817791,-1.149900)--cycle;
\path[fill=blue!21!white,draw=blue!21!white,line width=.12pt] (-2.956415,-0.612293)--(-2.942213,-0.744515)--(-2.808591,-0.885774)--(-2.822148,-0.754235)--cycle;
\path[fill=blue!29!white,draw=blue!29!white,line width=.12pt] (2.822148,-0.754235)--(2.808591,-0.885774)--(2.942213,-0.744515)--(2.956415,-0.612293)--cycle;
\path[fill=blue!41!white,draw=blue!41!white,line width=.12pt] (-0.864690,-1.298758)--(-0.817791,-1.149900)--(-0.663889,-1.186295)--(-0.701962,-1.337241)--cycle;
\path[fill=blue!38!white,draw=blue!38!white,line width=.12pt] (0.701962,-1.337241)--(0.663889,-1.186295)--(0.817791,-1.149900)--(0.864690,-1.298758)--cycle;
\path[fill=blue!37!white,draw=blue!37!white,line width=.12pt] (-2.091880,-1.498463)--(-1.975871,-1.490286)--(-1.807416,-1.583216)--(-1.913535,-1.596850)--cycle;
\path[fill=blue!39!white,draw=blue!39!white,line width=.12pt] (1.913535,-1.596850)--(1.807416,-1.583216)--(1.975871,-1.490286)--(2.091880,-1.498463)--cycle;
\path[fill=blue!13!white,draw=blue!13!white,line width=.12pt] (-2.446009,0.308514)--(-2.589621,0.247306)--(-2.500154,0.122284)--(-2.361503,0.190425)--cycle;
\path[fill=blue!16!white,draw=blue!16!white,line width=.12pt] (2.361503,0.190425)--(2.500154,0.122284)--(2.589621,0.247306)--(2.446009,0.308514)--cycle;
\path[fill=blue!22!white,draw=blue!22!white,line width=.12pt] (-1.617721,0.057553)--(-1.724636,0.115141)--(-1.625975,0.032838)--(-1.525177,-0.019647)--cycle;
\path[fill=blue!17!white,draw=blue!17!white,line width=.12pt] (1.525177,-0.019647)--(1.625975,0.032838)--(1.724636,0.115141)--(1.617721,0.057553)--cycle;
\path[fill=blue!42!white,draw=blue!42!white,line width=.12pt] (-1.163283,-1.510570)--(-1.086442,-1.389047)--(-0.921837,-1.438377)--(-0.987036,-1.563389)--cycle;
\path[fill=blue!40!white,draw=blue!40!white,line width=.12pt] (0.987036,-1.563389)--(0.921837,-1.438377)--(1.086442,-1.389047)--(1.163283,-1.510570)--cycle;
\path[fill=blue!36!white,draw=blue!36!white,line width=.12pt] (-0.897452,-0.806729)--(-0.888912,-0.665176)--(-0.754235,-0.705537)--(-0.761481,-0.847478)--cycle;
\path[fill=blue!32!white,draw=blue!32!white,line width=.12pt] (0.761481,-0.847478)--(0.754235,-0.705537)--(0.888912,-0.665176)--(0.897452,-0.806729)--cycle;
\path[fill=blue!13!white,draw=blue!13!white,line width=.12pt] (-2.721975,0.163208)--(-2.837984,0.059451)--(-2.739936,-0.077561)--(-2.627935,0.031796)--cycle;
\path[fill=blue!18!white,draw=blue!18!white,line width=.12pt] (2.627935,0.031796)--(2.739936,-0.077561)--(2.837984,0.059451)--(2.721975,0.163208)--cycle;
\path[fill=blue!34!white,draw=blue!34!white,line width=.12pt] (-1.015029,-0.618408)--(-1.024781,-0.489187)--(-0.897452,-0.536404)--(-0.888912,-0.665176)--cycle;
\path[fill=blue!28!white,draw=blue!28!white,line width=.12pt] (0.888912,-0.665176)--(0.897452,-0.536404)--(1.024781,-0.489187)--(1.015029,-0.618408)--cycle;
\path[fill=blue!13!white,draw=blue!13!white,line width=.12pt] (-2.589621,0.247306)--(-2.721975,0.163208)--(-2.627935,0.031796)--(-2.500154,0.122284)--cycle;
\path[fill=blue!17!white,draw=blue!17!white,line width=.12pt] (2.500154,0.122284)--(2.627935,0.031796)--(2.721975,0.163208)--(2.589621,0.247306)--cycle;
\path[fill=blue!38!white,draw=blue!38!white,line width=.12pt] (-0.782941,-0.997521)--(-0.761481,-0.847478)--(-0.618176,-0.881367)--(-0.635598,-1.032365)--cycle;
\path[fill=blue!35!white,draw=blue!35!white,line width=.12pt] (0.635598,-1.032365)--(0.618176,-0.881367)--(0.761481,-0.847478)--(0.782941,-0.997521)--cycle;
\path[fill=blue!31!white,draw=blue!31!white,line width=.12pt] (-1.142240,-0.435958)--(-1.174431,-0.322085)--(-1.053662,-0.376814)--(-1.024781,-0.489187)--cycle;
\path[fill=blue!25!white,draw=blue!25!white,line width=.12pt] (1.024781,-0.489187)--(1.053662,-0.376814)--(1.174431,-0.322085)--(1.142240,-0.435958)--cycle;
\path[fill=blue!25!white,draw=blue!25!white,line width=.12pt] (-1.442454,-0.096998)--(-1.525177,-0.019647)--(-1.417944,-0.091940)--(-1.341037,-0.165370)--cycle;
\path[fill=blue!19!white,draw=blue!19!white,line width=.12pt] (1.341037,-0.165370)--(1.417944,-0.091940)--(1.525177,-0.019647)--(1.442454,-0.096998)--cycle;
\path[fill=blue!28!white,draw=blue!28!white,line width=.12pt] (-1.283890,-0.261700)--(-1.341037,-0.165370)--(-1.226706,-0.228443)--(-1.174431,-0.322085)--cycle;
\path[fill=blue!22!white,draw=blue!22!white,line width=.12pt] (1.174431,-0.322085)--(1.226706,-0.228443)--(1.341037,-0.165370)--(1.283890,-0.261700)--cycle;
\path[fill=blue!42!white,draw=blue!42!white,line width=.12pt] (-0.921837,-1.438377)--(-0.864690,-1.298758)--(-0.701962,-1.337241)--(-0.748354,-1.479402)--cycle;
\path[fill=blue!39!white,draw=blue!39!white,line width=.12pt] (0.748354,-1.479402)--(0.701962,-1.337241)--(0.864690,-1.298758)--(0.921837,-1.438377)--cycle;
\path[fill=blue!17!white,draw=blue!17!white,line width=.12pt] (-1.934468,0.239441)--(-2.073119,0.250151)--(-1.978968,0.150618)--(-1.846614,0.146565)--cycle;
\path[fill=blue!15!white,draw=blue!15!white,line width=.12pt] (1.846614,0.146565)--(1.978968,0.150618)--(2.073119,0.250151)--(1.934468,0.239441)--cycle;
\path[fill=blue!32!white,draw=blue!32!white,line width=.12pt] (-2.465877,-1.313721)--(-2.365079,-1.366207)--(-2.198794,-1.478311)--(-2.292506,-1.430603)--cycle;
\path[fill=blue!36!white,draw=blue!36!white,line width=.12pt] (2.292506,-1.430603)--(2.198794,-1.478311)--(2.365079,-1.366207)--(2.465877,-1.313721)--cycle;
\path[fill=blue!26!white,draw=blue!26!white,line width=.12pt] (-2.768442,-1.005012)--(-2.703244,-1.107367)--(-2.548600,-1.236371)--(-2.610069,-1.137127)--cycle;
\path[fill=blue!33!white,draw=blue!33!white,line width=.12pt] (2.610069,-1.137127)--(2.548600,-1.236371)--(2.703244,-1.107367)--(2.768442,-1.005012)--cycle;
\path[fill=blue!39!white,draw=blue!39!white,line width=.12pt] (-1.807416,-1.583216)--(-1.697056,-1.541348)--(-1.522544,-1.620433)--(-1.621555,-1.667443)--cycle;
\path[fill=blue!39!white,draw=blue!39!white,line width=.12pt] (1.621555,-1.667443)--(1.522544,-1.620433)--(1.697056,-1.541348)--(1.807416,-1.583216)--cycle;
\path[fill=blue!40!white,draw=blue!40!white,line width=.12pt] (-0.663889,-1.186295)--(-0.635598,-1.032365)--(-0.482133,-1.059820)--(-0.503593,-1.214972)--cycle;
\path[fill=blue!37!white,draw=blue!37!white,line width=.12pt] (0.503593,-1.214972)--(0.482133,-1.059820)--(0.635598,-1.032365)--(0.663889,-1.186295)--cycle;
\path[fill=blue!41!white,draw=blue!41!white,line width=.12pt] (-0.701962,-1.337241)--(-0.663889,-1.186295)--(-0.503593,-1.214972)--(-0.532473,-1.367563)--cycle;
\path[fill=blue!39!white,draw=blue!39!white,line width=.12pt] (0.532473,-1.367563)--(0.503593,-1.214972)--(0.663889,-1.186295)--(0.701962,-1.337241)--cycle;
\path[fill=blue!40!white,draw=blue!40!white,line width=.12pt] (-1.522544,-1.620433)--(-1.423533,-1.546798)--(-1.246659,-1.612387)--(-1.333369,-1.690584)--cycle;
\path[fill=blue!40!white,draw=blue!40!white,line width=.12pt] (1.333369,-1.690584)--(1.246659,-1.612387)--(1.423533,-1.546798)--(1.522544,-1.620433)--cycle;
\path[fill=blue!36!white,draw=blue!36!white,line width=.12pt] (-0.761481,-0.847478)--(-0.754235,-0.705537)--(-0.612293,-0.739104)--(-0.618176,-0.881367)--cycle;
\path[fill=blue!32!white,draw=blue!32!white,line width=.12pt] (0.618176,-0.881367)--(0.612293,-0.739104)--(0.754235,-0.705537)--(0.761481,-0.847478)--cycle;
\path[fill=blue!38!white,draw=blue!38!white,line width=.12pt] (-0.635598,-1.032365)--(-0.618176,-0.881367)--(-0.468918,-0.908070)--(-0.482133,-1.059820)--cycle;
\path[fill=blue!35!white,draw=blue!35!white,line width=.12pt] (0.482133,-1.059820)--(0.468918,-0.908070)--(0.618176,-0.881367)--(0.635598,-1.032365)--cycle;
\path[fill=blue!19!white,draw=blue!19!white,line width=.12pt] (-2.942213,-0.474190)--(-2.956415,-0.612293)--(-2.822148,-0.754235)--(-2.808591,-0.615449)--cycle;
\path[fill=blue!27!white,draw=blue!27!white,line width=.12pt] (2.808591,-0.615449)--(2.822148,-0.754235)--(2.956415,-0.612293)--(2.942213,-0.474190)--cycle;
\path[fill=blue!41!white,draw=blue!41!white,line width=.12pt] (-1.246659,-1.612387)--(-1.163283,-1.510570)--(-0.987036,-1.563389)--(-1.057780,-1.668992)--cycle;
\path[fill=blue!40!white,draw=blue!40!white,line width=.12pt] (1.057780,-1.668992)--(0.987036,-1.563389)--(1.163283,-1.510570)--(1.246659,-1.612387)--cycle;
\path[fill=blue!34!white,draw=blue!34!white,line width=.12pt] (-0.888912,-0.665176)--(-0.897452,-0.536404)--(-0.761481,-0.577153)--(-0.754235,-0.705537)--cycle;
\path[fill=blue!29!white,draw=blue!29!white,line width=.12pt] (0.754235,-0.705537)--(0.761481,-0.577153)--(0.897452,-0.536404)--(0.888912,-0.665176)--cycle;
\path[fill=blue!42!white,draw=blue!42!white,line width=.12pt] (-0.987036,-1.563389)--(-0.921837,-1.438377)--(-0.748354,-1.479402)--(-0.801283,-1.607317)--cycle;
\path[fill=blue!40!white,draw=blue!40!white,line width=.12pt] (0.801283,-1.607317)--(0.748354,-1.479402)--(0.921837,-1.438377)--(0.987036,-1.563389)--cycle;
\path[fill=blue!42!white,draw=blue!42!white,line width=.12pt] (-0.748354,-1.479402)--(-0.701962,-1.337241)--(-0.532473,-1.367563)--(-0.567664,-1.511728)--cycle;
\path[fill=blue!40!white,draw=blue!40!white,line width=.12pt] (0.567664,-1.511728)--(0.532473,-1.367563)--(0.701962,-1.337241)--(0.748354,-1.479402)--cycle;
\path[fill=blue!40!white,draw=blue!40!white,line width=.12pt] (-0.503593,-1.214972)--(-0.482133,-1.059820)--(-0.324025,-1.079622)--(-0.338447,-1.235655)--cycle;
\path[fill=blue!38!white,draw=blue!38!white,line width=.12pt] (0.338447,-1.235655)--(0.324025,-1.079622)--(0.482133,-1.059820)--(0.503593,-1.214972)--cycle;
\path[fill=blue!41!white,draw=blue!41!white,line width=.12pt] (-0.532473,-1.367563)--(-0.503593,-1.214972)--(-0.338447,-1.235655)--(-0.357857,-1.389432)--cycle;
\path[fill=blue!39!white,draw=blue!39!white,line width=.12pt] (0.357857,-1.389432)--(0.338447,-1.235655)--(0.503593,-1.214972)--(0.532473,-1.367563)--cycle;
\path[fill=blue!19!white,draw=blue!19!white,line width=.12pt] (-1.724636,0.115141)--(-1.846614,0.146565)--(-1.740975,0.058441)--(-1.625975,0.032838)--cycle;
\path[fill=blue!16!white,draw=blue!16!white,line width=.12pt] (1.625975,0.032838)--(1.740975,0.058441)--(1.846614,0.146565)--(1.724636,0.115141)--cycle;
\path[fill=blue!31!white,draw=blue!31!white,line width=.12pt] (-1.024781,-0.489187)--(-1.053662,-0.376814)--(-0.922745,-0.425362)--(-0.897452,-0.536404)--cycle;
\path[fill=blue!26!white,draw=blue!26!white,line width=.12pt] (0.897452,-0.536404)--(0.922745,-0.425362)--(1.053662,-0.376814)--(1.024781,-0.489187)--cycle;
\path[fill=blue!15!white,draw=blue!15!white,line width=.12pt] (-2.073119,0.250151)--(-2.217311,0.233600)--(-2.116611,0.127144)--(-1.978968,0.150618)--cycle;
\path[fill=blue!14!white,draw=blue!14!white,line width=.12pt] (1.978968,0.150618)--(2.116611,0.127144)--(2.217311,0.233600)--(2.073119,0.250151)--cycle;
\path[fill=blue!38!white,draw=blue!38!white,line width=.12pt] (-0.482133,-1.059820)--(-0.468918,-0.908070)--(-0.315143,-0.927329)--(-0.324025,-1.079622)--cycle;
\path[fill=blue!36!white,draw=blue!36!white,line width=.12pt] (0.324025,-1.079622)--(0.315143,-0.927329)--(0.468918,-0.908070)--(0.482133,-1.059820)--cycle;
\path[fill=blue!34!white,draw=blue!34!white,line width=.12pt] (-2.198794,-1.478311)--(-2.091880,-1.498463)--(-1.913535,-1.596850)--(-2.011334,-1.581726)--cycle;
\path[fill=blue!38!white,draw=blue!38!white,line width=.12pt] (2.011334,-1.581726)--(1.913535,-1.596850)--(2.091880,-1.498463)--(2.198794,-1.478311)--cycle;
\path[fill=blue!39!white,draw=blue!39!white,line width=.12pt] (-0.338447,-1.235655)--(-0.324025,-1.079622)--(-0.162796,-1.091580)--(-0.170043,-1.248146)--cycle;
\path[fill=blue!38!white,draw=blue!38!white,line width=.12pt] (0.170043,-1.248146)--(0.162796,-1.091580)--(0.324025,-1.079622)--(0.338447,-1.235655)--cycle;
\path[fill=blue!36!white,draw=blue!36!white,line width=.12pt] (-0.618176,-0.881367)--(-0.612293,-0.739104)--(-0.464455,-0.765552)--(-0.468918,-0.908070)--cycle;
\path[fill=blue!33!white,draw=blue!33!white,line width=.12pt] (0.468918,-0.908070)--(0.464455,-0.765552)--(0.612293,-0.739104)--(0.618176,-0.881367)--cycle;
\path[fill=blue!41!white,draw=blue!41!white,line width=.12pt] (-0.357857,-1.389432)--(-0.338447,-1.235655)--(-0.170043,-1.248146)--(-0.179794,-1.402639)--cycle;
\path[fill=blue!40!white,draw=blue!40!white,line width=.12pt] (0.179794,-1.402639)--(0.170043,-1.248146)--(0.338447,-1.235655)--(0.357857,-1.389432)--cycle;
\path[fill=blue!28!white,draw=blue!28!white,line width=.12pt] (-1.174431,-0.322085)--(-1.226706,-0.228443)--(-1.100561,-0.285608)--(-1.053662,-0.376814)--cycle;
\path[fill=blue!23!white,draw=blue!23!white,line width=.12pt] (1.053662,-0.376814)--(1.100561,-0.285608)--(1.226706,-0.228443)--(1.174431,-0.322085)--cycle;
\path[fill=blue!39!white,draw=blue!39!white,line width=.12pt] (-0.170043,-1.248146)--(-0.162796,-1.091580)--(-0.000000,-1.095579)--(-0.000000,-1.252323)--cycle;
\path[fill=blue!39!white,draw=blue!39!white,line width=.12pt] (-0.000000,-1.252323)--(-0.000000,-1.095579)--(0.162796,-1.091580)--(0.170043,-1.248146)--cycle;
\path[fill=blue!42!white,draw=blue!42!white,line width=.12pt] (-0.567664,-1.511728)--(-0.532473,-1.367563)--(-0.357857,-1.389432)--(-0.381508,-1.535043)--cycle;
\path[fill=blue!40!white,draw=blue!40!white,line width=.12pt] (0.381508,-1.535043)--(0.357857,-1.389432)--(0.532473,-1.367563)--(0.567664,-1.511728)--cycle;
\path[fill=blue!22!white,draw=blue!22!white,line width=.12pt] (-1.525177,-0.019647)--(-1.625975,0.032838)--(-1.511656,-0.044233)--(-1.417944,-0.091940)--cycle;
\path[fill=blue!18!white,draw=blue!18!white,line width=.12pt] (1.417944,-0.091940)--(1.511656,-0.044233)--(1.625975,0.032838)--(1.525177,-0.019647)--cycle;
\path[fill=blue!17!white,draw=blue!17!white,line width=.12pt] (-2.900154,-0.335511)--(-2.942213,-0.474190)--(-2.808591,-0.615449)--(-2.768442,-0.474751)--cycle;
\path[fill=blue!24!white,draw=blue!24!white,line width=.12pt] (2.768442,-0.474751)--(2.808591,-0.615449)--(2.942213,-0.474190)--(2.900154,-0.335511)--cycle;
\path[fill=blue!38!white,draw=blue!38!white,line width=.12pt] (-0.324025,-1.079622)--(-0.315143,-0.927329)--(-0.158334,-0.938959)--(-0.162796,-1.091580)--cycle;
\path[fill=blue!36!white,draw=blue!36!white,line width=.12pt] (0.162796,-1.091580)--(0.158334,-0.938959)--(0.315143,-0.927329)--(0.324025,-1.079622)--cycle;
\path[fill=blue!25!white,draw=blue!25!white,line width=.12pt] (-1.341037,-0.165370)--(-1.417944,-0.091940)--(-1.297056,-0.158630)--(-1.226706,-0.228443)--cycle;
\path[fill=blue!20!white,draw=blue!20!white,line width=.12pt] (1.226706,-0.228443)--(1.297056,-0.158630)--(1.417944,-0.091940)--(1.341037,-0.165370)--cycle;
\path[fill=blue!40!white,draw=blue!40!white,line width=.12pt] (-0.179794,-1.402639)--(-0.170043,-1.248146)--(-0.000000,-1.252323)--(-0.000000,-1.407055)--cycle;
\path[fill=blue!40!white,draw=blue!40!white,line width=.12pt] (-0.000000,-1.407055)--(-0.000000,-1.252323)--(0.170043,-1.248146)--(0.179794,-1.402639)--cycle;
\path[fill=blue!24!white,draw=blue!24!white,line width=.12pt] (-2.808591,-0.885774)--(-2.768442,-1.005012)--(-2.610069,-1.137127)--(-2.647922,-1.019805)--cycle;
\path[fill=blue!31!white,draw=blue!31!white,line width=.12pt] (2.647922,-1.019805)--(2.610069,-1.137127)--(2.768442,-1.005012)--(2.808591,-0.885774)--cycle;
\path[fill=blue!34!white,draw=blue!34!white,line width=.12pt] (-0.754235,-0.705537)--(-0.761481,-0.577153)--(-0.618176,-0.611042)--(-0.612293,-0.739104)--cycle;
\path[fill=blue!30!white,draw=blue!30!white,line width=.12pt] (0.612293,-0.739104)--(0.618176,-0.611042)--(0.761481,-0.577153)--(0.754235,-0.705537)--cycle;
\path[fill=blue!42!white,draw=blue!42!white,line width=.12pt] (-0.801283,-1.607317)--(-0.748354,-1.479402)--(-0.567664,-1.511728)--(-0.607814,-1.641929)--cycle;
\path[fill=blue!40!white,draw=blue!40!white,line width=.12pt] (0.607814,-1.641929)--(0.567664,-1.511728)--(0.748354,-1.479402)--(0.801283,-1.607317)--cycle;
\path[fill=blue!37!white,draw=blue!37!white,line width=.12pt] (-0.162796,-1.091580)--(-0.158334,-0.938959)--(-0.000000,-0.942848)--(-0.000000,-1.095579)--cycle;
\path[fill=blue!37!white,draw=blue!37!white,line width=.12pt] (-0.000000,-1.095579)--(-0.000000,-0.942848)--(0.158334,-0.938959)--(0.162796,-1.091580)--cycle;
\path[fill=blue!29!white,draw=blue!29!white,line width=.12pt] (-2.548600,-1.236371)--(-2.465877,-1.313721)--(-2.292506,-1.430603)--(-2.369413,-1.357174)--cycle;
\path[fill=blue!35!white,draw=blue!35!white,line width=.12pt] (2.369413,-1.357174)--(2.292506,-1.430603)--(2.465877,-1.313721)--(2.548600,-1.236371)--cycle;
\path[fill=blue!37!white,draw=blue!37!white,line width=.12pt] (-1.913535,-1.596850)--(-1.807416,-1.583216)--(-1.621555,-1.667443)--(-1.716761,-1.686022)--cycle;
\path[fill=blue!38!white,draw=blue!38!white,line width=.12pt] (1.716761,-1.686022)--(1.621555,-1.667443)--(1.807416,-1.583216)--(1.913535,-1.596850)--cycle;
\path[fill=blue!36!white,draw=blue!36!white,line width=.12pt] (-0.468918,-0.908070)--(-0.464455,-0.765552)--(-0.312145,-0.784628)--(-0.315143,-0.927329)--cycle;
\path[fill=blue!34!white,draw=blue!34!white,line width=.12pt] (0.315143,-0.927329)--(0.312145,-0.784628)--(0.464455,-0.765552)--(0.468918,-0.908070)--cycle;
\path[fill=blue!14!white,draw=blue!14!white,line width=.12pt] (-2.217311,0.233600)--(-2.361503,0.190425)--(-2.254254,0.077046)--(-2.116611,0.127144)--cycle;
\path[fill=blue!15!white,draw=blue!15!white,line width=.12pt] (2.116611,0.127144)--(2.254254,0.077046)--(2.361503,0.190425)--(2.217311,0.233600)--cycle;
\path[fill=blue!41!white,draw=blue!41!white,line width=.12pt] (-0.381508,-1.535043)--(-0.357857,-1.389432)--(-0.179794,-1.402639)--(-0.191677,-1.549123)--cycle;
\path[fill=blue!41!white,draw=blue!41!white,line width=.12pt] (0.191677,-1.549123)--(0.179794,-1.402639)--(0.357857,-1.389432)--(0.381508,-1.535043)--cycle;
\path[fill=blue!41!white,draw=blue!41!white,line width=.12pt] (-1.057780,-1.668992)--(-0.987036,-1.563389)--(-0.801283,-1.607317)--(-0.858714,-1.716067)--cycle;
\path[fill=blue!40!white,draw=blue!40!white,line width=.12pt] (0.858714,-1.716067)--(0.801283,-1.607317)--(0.987036,-1.563389)--(1.057780,-1.668992)--cycle;
\path[fill=blue!39!white,draw=blue!39!white,line width=.12pt] (-1.621555,-1.667443)--(-1.522544,-1.620433)--(-1.333369,-1.690584)--(-1.420078,-1.742156)--cycle;
\path[fill=blue!39!white,draw=blue!39!white,line width=.12pt] (1.420078,-1.742156)--(1.333369,-1.690584)--(1.522544,-1.620433)--(1.621555,-1.667443)--cycle;
\path[fill=blue!40!white,draw=blue!40!white,line width=.12pt] (-1.333369,-1.690584)--(-1.246659,-1.612387)--(-1.057780,-1.668992)--(-1.131352,-1.751126)--cycle;
\path[fill=blue!40!white,draw=blue!40!white,line width=.12pt] (1.131352,-1.751126)--(1.057780,-1.668992)--(1.246659,-1.612387)--(1.333369,-1.690584)--cycle;
\path[fill=blue!41!white,draw=blue!41!white,line width=.12pt] (-0.000000,-1.553831)--(-0.000000,-1.407055)--(0.179794,-1.402639)--(0.191677,-1.549123)--cycle;
\path[fill=blue!41!white,draw=blue!41!white,line width=.12pt] (-0.191677,-1.549123)--(-0.179794,-1.402639)--(-0.000000,-1.407055)--(-0.000000,-1.553831)--cycle;
\path[fill=blue!31!white,draw=blue!31!white,line width=.12pt] (-0.897452,-0.536404)--(-0.922745,-0.425362)--(-0.782941,-0.467259)--(-0.761481,-0.577153)--cycle;
\path[fill=blue!27!white,draw=blue!27!white,line width=.12pt] (0.761481,-0.577153)--(0.782941,-0.467259)--(0.922745,-0.425362)--(0.897452,-0.536404)--cycle;
\path[fill=blue!15!white,draw=blue!15!white,line width=.12pt] (-2.831853,-0.201586)--(-2.900154,-0.335511)--(-2.768442,-0.474751)--(-2.703244,-0.337547)--cycle;
\path[fill=blue!22!white,draw=blue!22!white,line width=.12pt] (2.703244,-0.337547)--(2.768442,-0.474751)--(2.900154,-0.335511)--(2.831853,-0.201586)--cycle;
\path[fill=blue!36!white,draw=blue!36!white,line width=.12pt] (-0.315143,-0.927329)--(-0.312145,-0.784628)--(-0.156827,-0.796148)--(-0.158334,-0.938959)--cycle;
\path[fill=blue!34!white,draw=blue!34!white,line width=.12pt] (0.158334,-0.938959)--(0.156827,-0.796148)--(0.312145,-0.784628)--(0.315143,-0.927329)--cycle;
\path[fill=blue!17!white,draw=blue!17!white,line width=.12pt] (-1.846614,0.146565)--(-1.978968,0.150618)--(-1.865758,0.056178)--(-1.740975,0.058441)--cycle;
\path[fill=blue!15!white,draw=blue!15!white,line width=.12pt] (1.740975,0.058441)--(1.865758,0.056178)--(1.978968,0.150618)--(1.846614,0.146565)--cycle;
\path[fill=blue!42!white,draw=blue!42!white,line width=.12pt] (-0.607814,-1.641929)--(-0.567664,-1.511728)--(-0.381508,-1.535043)--(-0.408491,-1.666893)--cycle;
\path[fill=blue!41!white,draw=blue!41!white,line width=.12pt] (0.408491,-1.666893)--(0.381508,-1.535043)--(0.567664,-1.511728)--(0.607814,-1.641929)--cycle;
\path[fill=blue!34!white,draw=blue!34!white,line width=.12pt] (-0.612293,-0.739104)--(-0.618176,-0.611042)--(-0.468918,-0.637745)--(-0.464455,-0.765552)--cycle;
\path[fill=blue!31!white,draw=blue!31!white,line width=.12pt] (0.464455,-0.765552)--(0.468918,-0.637745)--(0.618176,-0.611042)--(0.612293,-0.739104)--cycle;
\path[fill=blue!35!white,draw=blue!35!white,line width=.12pt] (-0.158334,-0.938959)--(-0.156827,-0.796148)--(-0.000000,-0.800000)--(-0.000000,-0.942848)--cycle;
\path[fill=blue!35!white,draw=blue!35!white,line width=.12pt] (-0.000000,-0.942848)--(-0.000000,-0.800000)--(0.156827,-0.796148)--(0.158334,-0.938959)--cycle;
\path[fill=blue!13!white,draw=blue!13!white,line width=.12pt] (-2.361503,0.190425)--(-2.500154,0.122284)--(-2.386608,0.002248)--(-2.254254,0.077046)--cycle;
\path[fill=blue!15!white,draw=blue!15!white,line width=.12pt] (2.254254,0.077046)--(2.386608,0.002248)--(2.500154,0.122284)--(2.361503,0.190425)--cycle;
\path[fill=blue!28!white,draw=blue!28!white,line width=.12pt] (-1.053662,-0.376814)--(-1.100561,-0.285608)--(-0.963817,-0.336316)--(-0.922745,-0.425362)--cycle;
\path[fill=blue!23!white,draw=blue!23!white,line width=.12pt] (0.922745,-0.425362)--(0.963817,-0.336316)--(1.100561,-0.285608)--(1.053662,-0.376814)--cycle;
\path[fill=blue!14!white,draw=blue!14!white,line width=.12pt] (-2.739936,-0.077561)--(-2.831853,-0.201586)--(-2.703244,-0.337547)--(-2.615501,-0.209109)--cycle;
\path[fill=blue!20!white,draw=blue!20!white,line width=.12pt] (2.615501,-0.209109)--(2.703244,-0.337547)--(2.831853,-0.201586)--(2.739936,-0.077561)--cycle;
\path[fill=blue!13!white,draw=blue!13!white,line width=.12pt] (-2.500154,0.122284)--(-2.627935,0.031796)--(-2.508586,-0.094375)--(-2.386608,0.002248)--cycle;
\path[fill=blue!16!white,draw=blue!16!white,line width=.12pt] (2.386608,0.002248)--(2.508586,-0.094375)--(2.627935,0.031796)--(2.500154,0.122284)--cycle;
\path[fill=blue!32!white,draw=blue!32!white,line width=.12pt] (-2.292506,-1.430603)--(-2.198794,-1.478311)--(-2.011334,-1.581726)--(-2.097056,-1.538426)--cycle;
\path[fill=blue!36!white,draw=blue!36!white,line width=.12pt] (2.097056,-1.538426)--(2.011334,-1.581726)--(2.198794,-1.478311)--(2.292506,-1.430603)--cycle;
\path[fill=blue!21!white,draw=blue!21!white,line width=.12pt] (-2.822148,-0.754235)--(-2.808591,-0.885774)--(-2.647922,-1.019805)--(-2.660703,-0.888912)--cycle;
\path[fill=blue!29!white,draw=blue!29!white,line width=.12pt] (2.660703,-0.888912)--(2.647922,-1.019805)--(2.808591,-0.885774)--(2.822148,-0.754235)--cycle;
\path[fill=blue!42!white,draw=blue!42!white,line width=.12pt] (-0.408491,-1.666893)--(-0.381508,-1.535043)--(-0.191677,-1.549123)--(-0.205234,-1.681968)--cycle;
\path[fill=blue!41!white,draw=blue!41!white,line width=.12pt] (0.205234,-1.681968)--(0.191677,-1.549123)--(0.381508,-1.535043)--(0.408491,-1.666893)--cycle;
\path[fill=blue!13!white,draw=blue!13!white,line width=.12pt] (-2.627935,0.031796)--(-2.739936,-0.077561)--(-2.615501,-0.209109)--(-2.508586,-0.094375)--cycle;
\path[fill=blue!18!white,draw=blue!18!white,line width=.12pt] (2.508586,-0.094375)--(2.615501,-0.209109)--(2.739936,-0.077561)--(2.627935,0.031796)--cycle;
\path[fill=blue!41!white,draw=blue!41!white,line width=.12pt] (-0.858714,-1.716067)--(-0.801283,-1.607317)--(-0.607814,-1.641929)--(-0.651378,-1.753160)--cycle;
\path[fill=blue!40!white,draw=blue!40!white,line width=.12pt] (0.651378,-1.753160)--(0.607814,-1.641929)--(0.801283,-1.607317)--(0.858714,-1.716067)--cycle;
\path[fill=blue!33!white,draw=blue!33!white,line width=.12pt] (-0.464455,-0.765552)--(-0.468918,-0.637745)--(-0.315143,-0.657004)--(-0.312145,-0.784628)--cycle;
\path[fill=blue!31!white,draw=blue!31!white,line width=.12pt] (0.312145,-0.784628)--(0.315143,-0.657004)--(0.468918,-0.637745)--(0.464455,-0.765552)--cycle;
\path[fill=blue!19!white,draw=blue!19!white,line width=.12pt] (-1.625975,0.032838)--(-1.740975,0.058441)--(-1.618570,-0.024081)--(-1.511656,-0.044233)--cycle;
\path[fill=blue!16!white,draw=blue!16!white,line width=.12pt] (1.511656,-0.044233)--(1.618570,-0.024081)--(1.740975,0.058441)--(1.625975,0.032838)--cycle;
\path[fill=blue!25!white,draw=blue!25!white,line width=.12pt] (-1.226706,-0.228443)--(-1.297056,-0.158630)--(-1.163677,-0.219074)--(-1.100561,-0.285608)--cycle;
\path[fill=blue!21!white,draw=blue!21!white,line width=.12pt] (1.100561,-0.285608)--(1.163677,-0.219074)--(1.297056,-0.158630)--(1.226706,-0.228443)--cycle;
\path[fill=blue!31!white,draw=blue!31!white,line width=.12pt] (-0.761481,-0.577153)--(-0.782941,-0.467259)--(-0.635598,-0.502103)--(-0.618176,-0.611042)--cycle;
\path[fill=blue!27!white,draw=blue!27!white,line width=.12pt] (0.618176,-0.611042)--(0.635598,-0.502103)--(0.782941,-0.467259)--(0.761481,-0.577153)--cycle;
\path[fill=blue!41!white,draw=blue!41!white,line width=.12pt] (-0.205234,-1.681968)--(-0.191677,-1.549123)--(-0.000000,-1.553831)--(-0.000000,-1.687009)--cycle;
\path[fill=blue!41!white,draw=blue!41!white,line width=.12pt] (-0.000000,-1.687009)--(-0.000000,-1.553831)--(0.191677,-1.549123)--(0.205234,-1.681968)--cycle;
\path[fill=blue!22!white,draw=blue!22!white,line width=.12pt] (-1.417944,-0.091940)--(-1.511656,-0.044233)--(-1.382778,-0.115330)--(-1.297056,-0.158630)--cycle;
\path[fill=blue!18!white,draw=blue!18!white,line width=.12pt] (1.297056,-0.158630)--(1.382778,-0.115330)--(1.511656,-0.044233)--(1.417944,-0.091940)--cycle;
\path[fill=blue!26!white,draw=blue!26!white,line width=.12pt] (-2.610069,-1.137127)--(-2.548600,-1.236371)--(-2.369413,-1.357174)--(-2.426560,-1.260844)--cycle;
\path[fill=blue!33!white,draw=blue!33!white,line width=.12pt] (2.426560,-1.260844)--(2.369413,-1.357174)--(2.548600,-1.236371)--(2.610069,-1.137127)--cycle;
\path[fill=blue!33!white,draw=blue!33!white,line width=.12pt] (-0.312145,-0.784628)--(-0.315143,-0.657004)--(-0.158334,-0.668634)--(-0.156827,-0.796148)--cycle;
\path[fill=blue!32!white,draw=blue!32!white,line width=.12pt] (0.156827,-0.796148)--(0.158334,-0.668634)--(0.315143,-0.657004)--(0.312145,-0.784628)--cycle;
\path[fill=blue!40!white,draw=blue!40!white,line width=.12pt] (-1.131352,-1.751126)--(-1.057780,-1.668992)--(-0.858714,-1.716067)--(-0.918440,-1.801476)--cycle;
\path[fill=blue!40!white,draw=blue!40!white,line width=.12pt] (0.918440,-1.801476)--(0.858714,-1.716067)--(1.057780,-1.668992)--(1.131352,-1.751126)--cycle;
\path[fill=blue!34!white,draw=blue!34!white,line width=.12pt] (-2.011334,-1.581726)--(-1.913535,-1.596850)--(-1.716761,-1.686022)--(-1.804504,-1.675456)--cycle;
\path[fill=blue!37!white,draw=blue!37!white,line width=.12pt] (1.804504,-1.675456)--(1.716761,-1.686022)--(1.913535,-1.596850)--(2.011334,-1.581726)--cycle;
\path[fill=blue!33!white,draw=blue!33!white,line width=.12pt] (-0.156827,-0.796148)--(-0.158334,-0.668634)--(-0.000000,-0.672523)--(-0.000000,-0.800000)--cycle;
\path[fill=blue!32!white,draw=blue!32!white,line width=.12pt] (-0.000000,-0.800000)--(-0.000000,-0.672523)--(0.158334,-0.668634)--(0.156827,-0.796148)--cycle;
\path[fill=blue!15!white,draw=blue!15!white,line width=.12pt] (-1.978968,0.150618)--(-2.116611,0.127144)--(-1.995527,0.026136)--(-1.865758,0.056178)--cycle;
\path[fill=blue!14!white,draw=blue!14!white,line width=.12pt] (1.865758,0.056178)--(1.995527,0.026136)--(2.116611,0.127144)--(1.978968,0.150618)--cycle;
\path[fill=blue!41!white,draw=blue!41!white,line width=.12pt] (-0.651378,-1.753160)--(-0.607814,-1.641929)--(-0.408491,-1.666893)--(-0.437769,-1.779914)--cycle;
\path[fill=blue!41!white,draw=blue!41!white,line width=.12pt] (0.437769,-1.779914)--(0.408491,-1.666893)--(0.607814,-1.641929)--(0.651378,-1.753160)--cycle;
\path[fill=blue!39!white,draw=blue!39!white,line width=.12pt] (-1.420078,-1.742156)--(-1.333369,-1.690584)--(-1.131352,-1.751126)--(-1.204924,-1.806635)--cycle;
\path[fill=blue!39!white,draw=blue!39!white,line width=.12pt] (1.204924,-1.806635)--(1.131352,-1.751126)--(1.333369,-1.690584)--(1.420078,-1.742156)--cycle;
\path[fill=blue!37!white,draw=blue!37!white,line width=.12pt] (-1.716761,-1.686022)--(-1.621555,-1.667443)--(-1.420078,-1.742156)--(-1.503455,-1.765122)--cycle;
\path[fill=blue!38!white,draw=blue!38!white,line width=.12pt] (1.503455,-1.765122)--(1.420078,-1.742156)--(1.621555,-1.667443)--(1.716761,-1.686022)--cycle;
\path[fill=blue!28!white,draw=blue!28!white,line width=.12pt] (-0.922745,-0.425362)--(-0.963817,-0.336316)--(-0.817791,-0.380079)--(-0.782941,-0.467259)--cycle;
\path[fill=blue!24!white,draw=blue!24!white,line width=.12pt] (0.782941,-0.467259)--(0.817791,-0.380079)--(0.963817,-0.336316)--(0.922745,-0.425362)--cycle;
\path[fill=blue!31!white,draw=blue!31!white,line width=.12pt] (-0.618176,-0.611042)--(-0.635598,-0.502103)--(-0.482133,-0.529559)--(-0.468918,-0.637745)--cycle;
\path[fill=blue!28!white,draw=blue!28!white,line width=.12pt] (0.468918,-0.637745)--(0.482133,-0.529559)--(0.635598,-0.502103)--(0.618176,-0.611042)--cycle;
\path[fill=blue!18!white,draw=blue!18!white,line width=.12pt] (-2.808591,-0.615449)--(-2.822148,-0.754235)--(-2.660703,-0.888912)--(-2.647922,-0.749480)--cycle;
\path[fill=blue!26!white,draw=blue!26!white,line width=.12pt] (2.647922,-0.749480)--(2.660703,-0.888912)--(2.822148,-0.754235)--(2.808591,-0.615449)--cycle;
\path[fill=blue!41!white,draw=blue!41!white,line width=.12pt] (-0.437769,-1.779914)--(-0.408491,-1.666893)--(-0.205234,-1.681968)--(-0.219943,-1.796069)--cycle;
\path[fill=blue!41!white,draw=blue!41!white,line width=.12pt] (0.219943,-1.796069)--(0.205234,-1.681968)--(0.408491,-1.666893)--(0.437769,-1.779914)--cycle;
\path[fill=blue!31!white,draw=blue!31!white,line width=.12pt] (-0.468918,-0.637745)--(-0.482133,-0.529559)--(-0.324025,-0.549360)--(-0.315143,-0.657004)--cycle;
\path[fill=blue!29!white,draw=blue!29!white,line width=.12pt] (0.315143,-0.657004)--(0.324025,-0.549360)--(0.482133,-0.529559)--(0.468918,-0.637745)--cycle;
\path[fill=blue!25!white,draw=blue!25!white,line width=.12pt] (-1.100561,-0.285608)--(-1.163677,-0.219074)--(-1.019091,-0.272690)--(-0.963817,-0.336316)--cycle;
\path[fill=blue!21!white,draw=blue!21!white,line width=.12pt] (0.963817,-0.336316)--(1.019091,-0.272690)--(1.163677,-0.219074)--(1.100561,-0.285608)--cycle;
\path[fill=blue!41!white,draw=blue!41!white,line width=.12pt] (-0.219943,-1.796069)--(-0.205234,-1.681968)--(-0.000000,-1.687009)--(-0.000000,-1.801472)--cycle;
\path[fill=blue!41!white,draw=blue!41!white,line width=.12pt] (-0.000000,-1.801472)--(-0.000000,-1.687009)--(0.205234,-1.681968)--(0.219943,-1.796069)--cycle;
\path[fill=blue!40!white,draw=blue!40!white,line width=.12pt] (-0.918440,-1.801476)--(-0.858714,-1.716067)--(-0.651378,-1.753160)--(-0.696683,-1.841149)--cycle;
\path[fill=blue!40!white,draw=blue!40!white,line width=.12pt] (0.696683,-1.841149)--(0.651378,-1.753160)--(0.858714,-1.716067)--(0.918440,-1.801476)--cycle;
\path[fill=blue!14!white,draw=blue!14!white,line width=.12pt] (-2.116611,0.127144)--(-2.254254,0.077046)--(-2.125296,-0.030531)--(-1.995527,0.026136)--cycle;
\path[fill=blue!14!white,draw=blue!14!white,line width=.12pt] (1.995527,0.026136)--(2.125296,-0.030531)--(2.254254,0.077046)--(2.116611,0.127144)--cycle;
\path[fill=blue!17!white,draw=blue!17!white,line width=.12pt] (-1.740975,0.058441)--(-1.865758,0.056178)--(-1.734580,-0.032258)--(-1.618570,-0.024081)--cycle;
\path[fill=blue!15!white,draw=blue!15!white,line width=.12pt] (1.618570,-0.024081)--(1.734580,-0.032258)--(1.865758,0.056178)--(1.740975,0.058441)--cycle;
\path[fill=blue!29!white,draw=blue!29!white,line width=.12pt] (-2.369413,-1.357174)--(-2.292506,-1.430603)--(-2.097056,-1.538426)--(-2.167407,-1.468614)--cycle;
\path[fill=blue!34!white,draw=blue!34!white,line width=.12pt] (2.167407,-1.468614)--(2.097056,-1.538426)--(2.292506,-1.430603)--(2.369413,-1.357174)--cycle;
\path[fill=blue!23!white,draw=blue!23!white,line width=.12pt] (-2.647922,-1.019805)--(-2.610069,-1.137127)--(-2.426560,-1.260844)--(-2.461751,-1.145316)--cycle;
\path[fill=blue!30!white,draw=blue!30!white,line width=.12pt] (2.461751,-1.145316)--(2.426560,-1.260844)--(2.610069,-1.137127)--(2.647922,-1.019805)--cycle;
\path[fill=blue!28!white,draw=blue!28!white,line width=.12pt] (-0.782941,-0.467259)--(-0.817791,-0.380079)--(-0.663889,-0.416474)--(-0.635598,-0.502103)--cycle;
\path[fill=blue!25!white,draw=blue!25!white,line width=.12pt] (0.635598,-0.502103)--(0.663889,-0.416474)--(0.817791,-0.380079)--(0.782941,-0.467259)--cycle;
\path[fill=blue!30!white,draw=blue!30!white,line width=.12pt] (-0.315143,-0.657004)--(-0.324025,-0.549360)--(-0.162796,-0.561318)--(-0.158334,-0.668634)--cycle;
\path[fill=blue!29!white,draw=blue!29!white,line width=.12pt] (0.158334,-0.668634)--(0.162796,-0.561318)--(0.324025,-0.549360)--(0.315143,-0.657004)--cycle;
\path[fill=blue!16!white,draw=blue!16!white,line width=.12pt] (-2.768442,-0.474751)--(-2.808591,-0.615449)--(-2.647922,-0.749480)--(-2.610069,-0.606865)--cycle;
\path[fill=blue!24!white,draw=blue!24!white,line width=.12pt] (2.610069,-0.606865)--(2.647922,-0.749480)--(2.808591,-0.615449)--(2.768442,-0.474751)--cycle;
\path[fill=blue!22!white,draw=blue!22!white,line width=.12pt] (-1.297056,-0.158630)--(-1.382778,-0.115330)--(-1.240584,-0.179769)--(-1.163677,-0.219074)--cycle;
\path[fill=blue!19!white,draw=blue!19!white,line width=.12pt] (1.163677,-0.219074)--(1.240584,-0.179769)--(1.382778,-0.115330)--(1.297056,-0.158630)--cycle;
\path[fill=blue!20!white,draw=blue!20!white,line width=.12pt] (-1.511656,-0.044233)--(-1.618570,-0.024081)--(-1.480578,-0.100206)--(-1.382778,-0.115330)--cycle;
\path[fill=blue!16!white,draw=blue!16!white,line width=.12pt] (1.382778,-0.115330)--(1.480578,-0.100206)--(1.618570,-0.024081)--(1.511656,-0.044233)--cycle;
\path[fill=blue!30!white,draw=blue!30!white,line width=.12pt] (-0.158334,-0.668634)--(-0.162796,-0.561318)--(-0.000000,-0.565317)--(-0.000000,-0.672523)--cycle;
\path[fill=blue!30!white,draw=blue!30!white,line width=.12pt] (-0.000000,-0.672523)--(-0.000000,-0.565317)--(0.162796,-0.561318)--(0.158334,-0.668634)--cycle;
\path[fill=blue!39!white,draw=blue!39!white,line width=.12pt] (-1.204924,-1.806635)--(-1.131352,-1.751126)--(-0.918440,-1.801476)--(-0.978167,-1.860259)--cycle;
\path[fill=blue!39!white,draw=blue!39!white,line width=.12pt] (0.978167,-1.860259)--(0.918440,-1.801476)--(1.131352,-1.751126)--(1.204924,-1.806635)--cycle;
\path[fill=blue!13!white,draw=blue!13!white,line width=.12pt] (-2.254254,0.077046)--(-2.386608,0.002248)--(-2.250079,-0.111645)--(-2.125296,-0.030531)--cycle;
\path[fill=blue!15!white,draw=blue!15!white,line width=.12pt] (2.125296,-0.030531)--(2.250079,-0.111645)--(2.386608,0.002248)--(2.254254,0.077046)--cycle;
\path[fill=blue!40!white,draw=blue!40!white,line width=.12pt] (-0.696683,-1.841149)--(-0.651378,-1.753160)--(-0.437769,-1.779914)--(-0.468217,-1.869763)--cycle;
\path[fill=blue!40!white,draw=blue!40!white,line width=.12pt] (0.468217,-1.869763)--(0.437769,-1.779914)--(0.651378,-1.753160)--(0.696683,-1.841149)--cycle;
\path[fill=blue!32!white,draw=blue!32!white,line width=.12pt] (-2.097056,-1.538426)--(-2.011334,-1.581726)--(-1.804504,-1.675456)--(-1.881411,-1.636151)--cycle;
\path[fill=blue!35!white,draw=blue!35!white,line width=.12pt] (1.881411,-1.636151)--(1.804504,-1.675456)--(2.011334,-1.581726)--(2.097056,-1.538426)--cycle;
\path[fill=blue!15!white,draw=blue!15!white,line width=.12pt] (-2.703244,-0.337547)--(-2.768442,-0.474751)--(-2.610069,-0.606865)--(-2.548600,-0.466550)--cycle;
\path[fill=blue!21!white,draw=blue!21!white,line width=.12pt] (2.548600,-0.466550)--(2.610069,-0.606865)--(2.768442,-0.474751)--(2.703244,-0.337547)--cycle;
\path[fill=blue!28!white,draw=blue!28!white,line width=.12pt] (-0.635598,-0.502103)--(-0.663889,-0.416474)--(-0.503593,-0.445151)--(-0.482133,-0.529559)--cycle;
\path[fill=blue!25!white,draw=blue!25!white,line width=.12pt] (0.482133,-0.529559)--(0.503593,-0.445151)--(0.663889,-0.416474)--(0.635598,-0.502103)--cycle;
\path[fill=blue!37!white,draw=blue!37!white,line width=.12pt] (-1.503455,-1.765122)--(-1.420078,-1.742156)--(-1.204924,-1.806635)--(-1.275669,-1.833387)--cycle;
\path[fill=blue!38!white,draw=blue!38!white,line width=.12pt] (1.275669,-1.833387)--(1.204924,-1.806635)--(1.420078,-1.742156)--(1.503455,-1.765122)--cycle;
\path[fill=blue!25!white,draw=blue!25!white,line width=.12pt] (-0.963817,-0.336316)--(-1.019091,-0.272690)--(-0.864690,-0.318963)--(-0.817791,-0.380079)--cycle;
\path[fill=blue!22!white,draw=blue!22!white,line width=.12pt] (0.817791,-0.380079)--(0.864690,-0.318963)--(1.019091,-0.272690)--(0.963817,-0.336316)--cycle;
\path[fill=blue!34!white,draw=blue!34!white,line width=.12pt] (-1.804504,-1.675456)--(-1.716761,-1.686022)--(-1.503455,-1.765122)--(-1.580295,-1.758598)--cycle;
\path[fill=blue!37!white,draw=blue!37!white,line width=.12pt] (1.580295,-1.758598)--(1.503455,-1.765122)--(1.716761,-1.686022)--(1.804504,-1.675456)--cycle;
\path[fill=blue!12!white,draw=blue!12!white,line width=.12pt] (-2.386608,0.002248)--(-2.508586,-0.094375)--(-2.365079,-0.214089)--(-2.250079,-0.111645)--cycle;
\path[fill=blue!16!white,draw=blue!16!white,line width=.12pt] (2.250079,-0.111645)--(2.365079,-0.214089)--(2.508586,-0.094375)--(2.386608,0.002248)--cycle;
\path[fill=blue!40!white,draw=blue!40!white,line width=.12pt] (-0.468217,-1.869763)--(-0.437769,-1.779914)--(-0.219943,-1.796069)--(-0.235241,-1.887042)--cycle;
\path[fill=blue!40!white,draw=blue!40!white,line width=.12pt] (0.235241,-1.887042)--(0.219943,-1.796069)--(0.437769,-1.779914)--(0.468217,-1.869763)--cycle;
\path[fill=blue!15!white,draw=blue!15!white,line width=.12pt] (-1.865758,0.056178)--(-1.995527,0.026136)--(-1.855225,-0.068452)--(-1.734580,-0.032258)--cycle;
\path[fill=blue!14!white,draw=blue!14!white,line width=.12pt] (1.734580,-0.032258)--(1.855225,-0.068452)--(1.995527,0.026136)--(1.865758,0.056178)--cycle;
\path[fill=blue!13!white,draw=blue!13!white,line width=.12pt] (-2.615501,-0.209109)--(-2.703244,-0.337547)--(-2.548600,-0.466550)--(-2.465877,-0.333925)--cycle;
\path[fill=blue!19!white,draw=blue!19!white,line width=.12pt] (2.465877,-0.333925)--(2.548600,-0.466550)--(2.703244,-0.337547)--(2.615501,-0.209109)--cycle;
\path[fill=blue!13!white,draw=blue!13!white,line width=.12pt] (-2.508586,-0.094375)--(-2.615501,-0.209109)--(-2.465877,-0.333925)--(-2.365079,-0.214089)--cycle;
\path[fill=blue!17!white,draw=blue!17!white,line width=.12pt] (2.365079,-0.214089)--(2.465877,-0.333925)--(2.615501,-0.209109)--(2.508586,-0.094375)--cycle;
\path[fill=blue!21!white,draw=blue!21!white,line width=.12pt] (-2.660703,-0.888912)--(-2.647922,-1.019805)--(-2.461751,-1.145316)--(-2.473633,-1.015029)--cycle;
\path[fill=blue!28!white,draw=blue!28!white,line width=.12pt] (2.473633,-1.015029)--(2.461751,-1.145316)--(2.647922,-1.019805)--(2.660703,-0.888912)--cycle;
\path[fill=blue!40!white,draw=blue!40!white,line width=.12pt] (-0.235241,-1.887042)--(-0.219943,-1.796069)--(-0.000000,-1.801472)--(-0.000000,-1.892820)--cycle;
\path[fill=blue!40!white,draw=blue!40!white,line width=.12pt] (-0.000000,-1.892820)--(-0.000000,-1.801472)--(0.219943,-1.796069)--(0.235241,-1.887042)--cycle;
\path[fill=blue!28!white,draw=blue!28!white,line width=.12pt] (-0.482133,-0.529559)--(-0.503593,-0.445151)--(-0.338447,-0.465835)--(-0.324025,-0.549360)--cycle;
\path[fill=blue!26!white,draw=blue!26!white,line width=.12pt] (0.324025,-0.549360)--(0.338447,-0.465835)--(0.503593,-0.445151)--(0.482133,-0.529559)--cycle;
\path[fill=blue!26!white,draw=blue!26!white,line width=.12pt] (-2.426560,-1.260844)--(-2.369413,-1.357174)--(-2.167407,-1.468614)--(-2.219681,-1.374972)--cycle;
\path[fill=blue!32!white,draw=blue!32!white,line width=.12pt] (2.219681,-1.374972)--(2.167407,-1.468614)--(2.369413,-1.357174)--(2.426560,-1.260844)--cycle;
\path[fill=blue!39!white,draw=blue!39!white,line width=.12pt] (-0.978167,-1.860259)--(-0.918440,-1.801476)--(-0.696683,-1.841149)--(-0.741989,-1.902512)--cycle;
\path[fill=blue!39!white,draw=blue!39!white,line width=.12pt] (0.741989,-1.902512)--(0.696683,-1.841149)--(0.918440,-1.801476)--(0.978167,-1.860259)--cycle;
\path[fill=blue!22!white,draw=blue!22!white,line width=.12pt] (-1.163677,-0.219074)--(-1.240584,-0.179769)--(-1.086442,-0.236929)--(-1.019091,-0.272690)--cycle;
\path[fill=blue!19!white,draw=blue!19!white,line width=.12pt] (1.019091,-0.272690)--(1.086442,-0.236929)--(1.240584,-0.179769)--(1.163677,-0.219074)--cycle;
\path[fill=blue!28!white,draw=blue!28!white,line width=.12pt] (-0.324025,-0.549360)--(-0.338447,-0.465835)--(-0.170043,-0.478325)--(-0.162796,-0.561318)--cycle;
\path[fill=blue!26!white,draw=blue!26!white,line width=.12pt] (0.162796,-0.561318)--(0.170043,-0.478325)--(0.338447,-0.465835)--(0.324025,-0.549360)--cycle;
\path[fill=blue!17!white,draw=blue!17!white,line width=.12pt] (-1.618570,-0.024081)--(-1.734580,-0.032258)--(-1.586697,-0.113840)--(-1.480578,-0.100206)--cycle;
\path[fill=blue!15!white,draw=blue!15!white,line width=.12pt] (1.480578,-0.100206)--(1.586697,-0.113840)--(1.734580,-0.032258)--(1.618570,-0.024081)--cycle;
\path[fill=blue!25!white,draw=blue!25!white,line width=.12pt] (-0.817791,-0.380079)--(-0.864690,-0.318963)--(-0.701962,-0.357445)--(-0.663889,-0.416474)--cycle;
\path[fill=blue!22!white,draw=blue!22!white,line width=.12pt] (0.663889,-0.416474)--(0.701962,-0.357445)--(0.864690,-0.318963)--(0.817791,-0.380079)--cycle;
\path[fill=blue!20!white,draw=blue!20!white,line width=.12pt] (-1.382778,-0.115330)--(-1.480578,-0.100206)--(-1.328326,-0.169203)--(-1.240584,-0.179769)--cycle;
\path[fill=blue!17!white,draw=blue!17!white,line width=.12pt] (1.240584,-0.179769)--(1.328326,-0.169203)--(1.480578,-0.100206)--(1.382778,-0.115330)--cycle;
\path[fill=blue!27!white,draw=blue!27!white,line width=.12pt] (-0.162796,-0.561318)--(-0.170043,-0.478325)--(-0.000000,-0.482502)--(-0.000000,-0.565317)--cycle;
\path[fill=blue!27!white,draw=blue!27!white,line width=.12pt] (-0.000000,-0.565317)--(-0.000000,-0.482502)--(0.170043,-0.478325)--(0.162796,-0.561318)--cycle;
\path[fill=blue!14!white,draw=blue!14!white,line width=.12pt] (-1.995527,0.026136)--(-2.125296,-0.030531)--(-1.975871,-0.131270)--(-1.855225,-0.068452)--cycle;
\path[fill=blue!14!white,draw=blue!14!white,line width=.12pt] (1.855225,-0.068452)--(1.975871,-0.131270)--(2.125296,-0.030531)--(1.995527,0.026136)--cycle;
\path[fill=blue!37!white,draw=blue!37!white,line width=.12pt] (-1.275669,-1.833387)--(-1.204924,-1.806635)--(-0.978167,-1.860259)--(-1.035598,-1.890159)--cycle;
\path[fill=blue!38!white,draw=blue!38!white,line width=.12pt] (1.035598,-1.890159)--(0.978167,-1.860259)--(1.204924,-1.806635)--(1.275669,-1.833387)--cycle;
\path[fill=blue!29!white,draw=blue!29!white,line width=.12pt] (-2.167407,-1.468614)--(-2.097056,-1.538426)--(-1.881411,-1.636151)--(-1.944527,-1.569617)--cycle;
\path[fill=blue!33!white,draw=blue!33!white,line width=.12pt] (1.944527,-1.569617)--(1.881411,-1.636151)--(2.097056,-1.538426)--(2.167407,-1.468614)--cycle;
\path[fill=blue!39!white,draw=blue!39!white,line width=.12pt] (-0.741989,-1.902512)--(-0.696683,-1.841149)--(-0.468217,-1.869763)--(-0.498665,-1.932987)--cycle;
\path[fill=blue!39!white,draw=blue!39!white,line width=.12pt] (0.498665,-1.932987)--(0.468217,-1.869763)--(0.696683,-1.841149)--(0.741989,-1.902512)--cycle;
\path[fill=blue!18!white,draw=blue!18!white,line width=.12pt] (-2.647922,-0.749480)--(-2.660703,-0.888912)--(-2.473633,-1.015029)--(-2.461751,-0.874991)--cycle;
\path[fill=blue!25!white,draw=blue!25!white,line width=.12pt] (2.461751,-0.874991)--(2.473633,-1.015029)--(2.660703,-0.888912)--(2.647922,-0.749480)--cycle;
\path[fill=blue!25!white,draw=blue!25!white,line width=.12pt] (-0.663889,-0.416474)--(-0.701962,-0.357445)--(-0.532473,-0.387767)--(-0.503593,-0.445151)--cycle;
\path[fill=blue!23!white,draw=blue!23!white,line width=.12pt] (0.503593,-0.445151)--(0.532473,-0.387767)--(0.701962,-0.357445)--(0.663889,-0.416474)--cycle;
\path[fill=blue!34!white,draw=blue!34!white,line width=.12pt] (-1.580295,-1.758598)--(-1.503455,-1.765122)--(-1.275669,-1.833387)--(-1.340867,-1.830352)--cycle;
\path[fill=blue!36!white,draw=blue!36!white,line width=.12pt] (1.340867,-1.830352)--(1.275669,-1.833387)--(1.503455,-1.765122)--(1.580295,-1.758598)--cycle;
\path[fill=blue!32!white,draw=blue!32!white,line width=.12pt] (-1.881411,-1.636151)--(-1.804504,-1.675456)--(-1.580295,-1.758598)--(-1.647647,-1.722837)--cycle;
\path[fill=blue!35!white,draw=blue!35!white,line width=.12pt] (1.647647,-1.722837)--(1.580295,-1.758598)--(1.804504,-1.675456)--(1.881411,-1.636151)--cycle;
\path[fill=blue!22!white,draw=blue!22!white,line width=.12pt] (-1.019091,-0.272690)--(-1.086442,-0.236929)--(-0.921837,-0.286259)--(-0.864690,-0.318963)--cycle;
\path[fill=blue!19!white,draw=blue!19!white,line width=.12pt] (0.864690,-0.318963)--(0.921837,-0.286259)--(1.086442,-0.236929)--(1.019091,-0.272690)--cycle;
\path[fill=blue!23!white,draw=blue!23!white,line width=.12pt] (-2.461751,-1.145316)--(-2.426560,-1.260844)--(-2.219681,-1.374972)--(-2.251872,-1.261099)--cycle;
\path[fill=blue!30!white,draw=blue!30!white,line width=.12pt] (2.251872,-1.261099)--(2.219681,-1.374972)--(2.426560,-1.260844)--(2.461751,-1.145316)--cycle;
\path[fill=blue!39!white,draw=blue!39!white,line width=.12pt] (-0.498665,-1.932987)--(-0.468217,-1.869763)--(-0.235241,-1.887042)--(-0.250539,-1.951390)--cycle;
\path[fill=blue!39!white,draw=blue!39!white,line width=.12pt] (0.250539,-1.951390)--(0.235241,-1.887042)--(0.468217,-1.869763)--(0.498665,-1.932987)--cycle;
\path[fill=blue!13!white,draw=blue!13!white,line width=.12pt] (-2.125296,-0.030531)--(-2.250079,-0.111645)--(-2.091880,-0.218298)--(-1.975871,-0.131270)--cycle;
\path[fill=blue!15!white,draw=blue!15!white,line width=.12pt] (1.975871,-0.131270)--(2.091880,-0.218298)--(2.250079,-0.111645)--(2.125296,-0.030531)--cycle;
\path[fill=blue!39!white,draw=blue!39!white,line width=.12pt] (-0.250539,-1.951390)--(-0.235241,-1.887042)--(-0.000000,-1.892820)--(-0.000000,-1.957544)--cycle;
\path[fill=blue!39!white,draw=blue!39!white,line width=.12pt] (-0.000000,-1.957544)--(-0.000000,-1.892820)--(0.235241,-1.887042)--(0.250539,-1.951390)--cycle;
\path[fill=blue!16!white,draw=blue!16!white,line width=.12pt] (-2.610069,-0.606865)--(-2.647922,-0.749480)--(-2.461751,-0.874991)--(-2.426560,-0.730582)--cycle;
\path[fill=blue!23!white,draw=blue!23!white,line width=.12pt] (2.426560,-0.730582)--(2.461751,-0.874991)--(2.647922,-0.749480)--(2.610069,-0.606865)--cycle;
\path[fill=blue!15!white,draw=blue!15!white,line width=.12pt] (-1.734580,-0.032258)--(-1.855225,-0.068452)--(-1.697056,-0.155708)--(-1.586697,-0.113840)--cycle;
\path[fill=blue!15!white,draw=blue!15!white,line width=.12pt] (1.586697,-0.113840)--(1.697056,-0.155708)--(1.855225,-0.068452)--(1.734580,-0.032258)--cycle;
\path[fill=blue!25!white,draw=blue!25!white,line width=.12pt] (-0.503593,-0.445151)--(-0.532473,-0.387767)--(-0.357857,-0.409636)--(-0.338447,-0.465835)--cycle;
\path[fill=blue!23!white,draw=blue!23!white,line width=.12pt] (0.338447,-0.465835)--(0.357857,-0.409636)--(0.532473,-0.387767)--(0.503593,-0.445151)--cycle;
\path[fill=blue!20!white,draw=blue!20!white,line width=.12pt] (-1.240584,-0.179769)--(-1.328326,-0.169203)--(-1.163283,-0.230405)--(-1.086442,-0.236929)--cycle;
\path[fill=blue!17!white,draw=blue!17!white,line width=.12pt] (1.086442,-0.236929)--(1.163283,-0.230405)--(1.328326,-0.169203)--(1.240584,-0.179769)--cycle;
\path[fill=blue!37!white,draw=blue!37!white,line width=.12pt] (-1.035598,-1.890159)--(-0.978167,-1.860259)--(-0.741989,-1.902512)--(-0.785553,-1.934893)--cycle;
\path[fill=blue!38!white,draw=blue!38!white,line width=.12pt] (0.785553,-1.934893)--(0.741989,-1.902512)--(0.978167,-1.860259)--(1.035598,-1.890159)--cycle;
\path[fill=blue!17!white,draw=blue!17!white,line width=.12pt] (-1.480578,-0.100206)--(-1.586697,-0.113840)--(-1.423533,-0.187782)--(-1.328326,-0.169203)--cycle;
\path[fill=blue!16!white,draw=blue!16!white,line width=.12pt] (1.328326,-0.169203)--(1.423533,-0.187782)--(1.586697,-0.113840)--(1.480578,-0.100206)--cycle;
\path[fill=blue!25!white,draw=blue!25!white,line width=.12pt] (-0.338447,-0.465835)--(-0.357857,-0.409636)--(-0.179794,-0.422843)--(-0.170043,-0.478325)--cycle;
\path[fill=blue!24!white,draw=blue!24!white,line width=.12pt] (0.170043,-0.478325)--(0.179794,-0.422843)--(0.357857,-0.409636)--(0.338447,-0.465835)--cycle;
\path[fill=blue!12!white,draw=blue!12!white,line width=.12pt] (-2.250079,-0.111645)--(-2.365079,-0.214089)--(-2.198794,-0.326193)--(-2.091880,-0.218298)--cycle;
\path[fill=blue!16!white,draw=blue!16!white,line width=.12pt] (2.091880,-0.218298)--(2.198794,-0.326193)--(2.365079,-0.214089)--(2.250079,-0.111645)--cycle;
\path[fill=blue!14!white,draw=blue!14!white,line width=.12pt] (-2.548600,-0.466550)--(-2.610069,-0.606865)--(-2.426560,-0.730582)--(-2.369413,-0.587353)--cycle;
\path[fill=blue!21!white,draw=blue!21!white,line width=.12pt] (2.369413,-0.587353)--(2.426560,-0.730582)--(2.610069,-0.606865)--(2.548600,-0.466550)--cycle;
\path[fill=blue!22!white,draw=blue!22!white,line width=.12pt] (-0.864690,-0.318963)--(-0.921837,-0.286259)--(-0.748354,-0.327284)--(-0.701962,-0.357445)--cycle;
\path[fill=blue!20!white,draw=blue!20!white,line width=.12pt] (0.701962,-0.357445)--(0.748354,-0.327284)--(0.921837,-0.286259)--(0.864690,-0.318963)--cycle;
\path[fill=blue!24!white,draw=blue!24!white,line width=.12pt] (-0.170043,-0.478325)--(-0.179794,-0.422843)--(-0.000000,-0.427259)--(-0.000000,-0.482502)--cycle;
\path[fill=blue!24!white,draw=blue!24!white,line width=.12pt] (-0.000000,-0.482502)--(-0.000000,-0.427259)--(0.179794,-0.422843)--(0.170043,-0.478325)--cycle;
\path[fill=blue!26!white,draw=blue!26!white,line width=.12pt] (-2.219681,-1.374972)--(-2.167407,-1.468614)--(-1.944527,-1.569617)--(-1.991426,-1.478411)--cycle;
\path[fill=blue!31!white,draw=blue!31!white,line width=.12pt] (1.991426,-1.478411)--(1.944527,-1.569617)--(2.167407,-1.468614)--(2.219681,-1.374972)--cycle;
\path[fill=blue!12!white,draw=blue!12!white,line width=.12pt] (-2.365079,-0.214089)--(-2.465877,-0.333925)--(-2.292506,-0.450808)--(-2.198794,-0.326193)--cycle;
\path[fill=blue!17!white,draw=blue!17!white,line width=.12pt] (2.198794,-0.326193)--(2.292506,-0.450808)--(2.465877,-0.333925)--(2.365079,-0.214089)--cycle;
\path[fill=blue!13!white,draw=blue!13!white,line width=.12pt] (-2.465877,-0.333925)--(-2.548600,-0.466550)--(-2.369413,-0.587353)--(-2.292506,-0.450808)--cycle;
\path[fill=blue!19!white,draw=blue!19!white,line width=.12pt] (2.292506,-0.450808)--(2.369413,-0.587353)--(2.548600,-0.466550)--(2.465877,-0.333925)--cycle;
\path[fill=blue!20!white,draw=blue!20!white,line width=.12pt] (-2.473633,-1.015029)--(-2.461751,-1.145316)--(-2.251872,-1.261099)--(-2.262742,-1.131371)--cycle;
\path[fill=blue!27!white,draw=blue!27!white,line width=.12pt] (2.262742,-1.131371)--(2.251872,-1.261099)--(2.461751,-1.145316)--(2.473633,-1.015029)--cycle;
\path[fill=blue!34!white,draw=blue!34!white,line width=.12pt] (-1.340867,-1.830352)--(-1.275669,-1.833387)--(-1.035598,-1.890159)--(-1.088526,-1.890026)--cycle;
\path[fill=blue!36!white,draw=blue!36!white,line width=.12pt] (1.088526,-1.890026)--(1.035598,-1.890159)--(1.275669,-1.833387)--(1.340867,-1.830352)--cycle;
\path[fill=blue!37!white,draw=blue!37!white,line width=.12pt] (-0.785553,-1.934893)--(-0.741989,-1.902512)--(-0.498665,-1.932987)--(-0.527943,-1.967157)--cycle;
\path[fill=blue!37!white,draw=blue!37!white,line width=.12pt] (0.527943,-1.967157)--(0.498665,-1.932987)--(0.741989,-1.902512)--(0.785553,-1.934893)--cycle;
\path[fill=blue!14!white,draw=blue!14!white,line width=.12pt] (-1.855225,-0.068452)--(-1.975871,-0.131270)--(-1.807416,-0.224200)--(-1.697056,-0.155708)--cycle;
\path[fill=blue!14!white,draw=blue!14!white,line width=.12pt] (1.697056,-0.155708)--(1.807416,-0.224200)--(1.975871,-0.131270)--(1.855225,-0.068452)--cycle;
\path[fill=blue!29!white,draw=blue!29!white,line width=.12pt] (-1.944527,-1.569617)--(-1.881411,-1.636151)--(-1.647647,-1.722837)--(-1.702920,-1.659211)--cycle;
\path[fill=blue!33!white,draw=blue!33!white,line width=.12pt] (1.702920,-1.659211)--(1.647647,-1.722837)--(1.881411,-1.636151)--(1.944527,-1.569617)--cycle;
\path[fill=blue!32!white,draw=blue!32!white,line width=.12pt] (-1.647647,-1.722837)--(-1.580295,-1.758598)--(-1.340867,-1.830352)--(-1.398014,-1.797648)--cycle;
\path[fill=blue!35!white,draw=blue!35!white,line width=.12pt] (1.398014,-1.797648)--(1.340867,-1.830352)--(1.580295,-1.758598)--(1.647647,-1.722837)--cycle;
\path[fill=blue!22!white,draw=blue!22!white,line width=.12pt] (-0.701962,-0.357445)--(-0.748354,-0.327284)--(-0.567664,-0.359610)--(-0.532473,-0.387767)--cycle;
\path[fill=blue!20!white,draw=blue!20!white,line width=.12pt] (0.532473,-0.387767)--(0.567664,-0.359610)--(0.748354,-0.327284)--(0.701962,-0.357445)--cycle;
\path[fill=blue!20!white,draw=blue!20!white,line width=.12pt] (-1.086442,-0.236929)--(-1.163283,-0.230405)--(-0.987036,-0.283224)--(-0.921837,-0.286259)--cycle;
\path[fill=blue!18!white,draw=blue!18!white,line width=.12pt] (0.921837,-0.286259)--(0.987036,-0.283224)--(1.163283,-0.230405)--(1.086442,-0.236929)--cycle;
\path[fill=blue!37!white,draw=blue!37!white,line width=.12pt] (-0.527943,-1.967157)--(-0.498665,-1.932987)--(-0.250539,-1.951390)--(-0.265249,-1.986640)--cycle;
\path[fill=blue!37!white,draw=blue!37!white,line width=.12pt] (0.265249,-1.986640)--(0.250539,-1.951390)--(0.498665,-1.932987)--(0.527943,-1.967157)--cycle;
\path[fill=blue!15!white,draw=blue!15!white,line width=.12pt] (-1.586697,-0.113840)--(-1.697056,-0.155708)--(-1.522544,-0.234792)--(-1.423533,-0.187782)--cycle;
\path[fill=blue!15!white,draw=blue!15!white,line width=.12pt] (1.423533,-0.187782)--(1.522544,-0.234792)--(1.697056,-0.155708)--(1.586697,-0.113840)--cycle;
\path[fill=blue!17!white,draw=blue!17!white,line width=.12pt] (-1.328326,-0.169203)--(-1.423533,-0.187782)--(-1.246659,-0.253371)--(-1.163283,-0.230405)--cycle;
\path[fill=blue!16!white,draw=blue!16!white,line width=.12pt] (1.163283,-0.230405)--(1.246659,-0.253371)--(1.423533,-0.187782)--(1.328326,-0.169203)--cycle;
\path[fill=blue!37!white,draw=blue!37!white,line width=.12pt] (-0.265249,-1.986640)--(-0.250539,-1.951390)--(-0.000000,-1.957544)--(-0.000000,-1.993156)--cycle;
\path[fill=blue!37!white,draw=blue!37!white,line width=.12pt] (-0.000000,-1.993156)--(-0.000000,-1.957544)--(0.250539,-1.951390)--(0.265249,-1.986640)--cycle;
\path[fill=blue!18!white,draw=blue!18!white,line width=.12pt] (-2.461751,-0.874991)--(-2.473633,-1.015029)--(-2.262742,-1.131371)--(-2.251872,-0.990774)--cycle;
\path[fill=blue!25!white,draw=blue!25!white,line width=.12pt] (2.251872,-0.990774)--(2.262742,-1.131371)--(2.473633,-1.015029)--(2.461751,-0.874991)--cycle;
\path[fill=blue!22!white,draw=blue!22!white,line width=.12pt] (-0.532473,-0.387767)--(-0.567664,-0.359610)--(-0.381508,-0.382925)--(-0.357857,-0.409636)--cycle;
\path[fill=blue!21!white,draw=blue!21!white,line width=.12pt] (0.357857,-0.409636)--(0.381508,-0.382925)--(0.567664,-0.359610)--(0.532473,-0.387767)--cycle;
\path[fill=blue!34!white,draw=blue!34!white,line width=.12pt] (-1.088526,-1.890026)--(-1.035598,-1.890159)--(-0.785553,-1.934893)--(-0.825702,-1.937046)--cycle;
\path[fill=blue!36!white,draw=blue!36!white,line width=.12pt] (0.825702,-1.937046)--(0.785553,-1.934893)--(1.035598,-1.890159)--(1.088526,-1.890026)--cycle;
\path[fill=blue!13!white,draw=blue!13!white,line width=.12pt] (-1.975871,-0.131270)--(-2.091880,-0.218298)--(-1.913535,-0.316685)--(-1.807416,-0.224200)--cycle;
\path[fill=blue!15!white,draw=blue!15!white,line width=.12pt] (1.807416,-0.224200)--(1.913535,-0.316685)--(2.091880,-0.218298)--(1.975871,-0.131270)--cycle;
\path[fill=blue!23!white,draw=blue!23!white,line width=.12pt] (-2.251872,-1.261099)--(-2.219681,-1.374972)--(-1.991426,-1.478411)--(-2.020307,-1.366038)--cycle;
\path[fill=blue!29!white,draw=blue!29!white,line width=.12pt] (2.020307,-1.366038)--(1.991426,-1.478411)--(2.219681,-1.374972)--(2.251872,-1.261099)--cycle;
\path[fill=blue!22!white,draw=blue!22!white,line width=.12pt] (-0.357857,-0.409636)--(-0.381508,-0.382925)--(-0.191677,-0.397005)--(-0.179794,-0.422843)--cycle;
\path[fill=blue!21!white,draw=blue!21!white,line width=.12pt] (0.179794,-0.422843)--(0.191677,-0.397005)--(0.381508,-0.382925)--(0.357857,-0.409636)--cycle;
\path[fill=blue!20!white,draw=blue!20!white,line width=.12pt] (-0.921837,-0.286259)--(-0.987036,-0.283224)--(-0.801283,-0.327152)--(-0.748354,-0.327284)--cycle;
\path[fill=blue!18!white,draw=blue!18!white,line width=.12pt] (0.748354,-0.327284)--(0.801283,-0.327152)--(0.987036,-0.283224)--(0.921837,-0.286259)--cycle;
\path[fill=blue!22!white,draw=blue!22!white,line width=.12pt] (-0.179794,-0.422843)--(-0.191677,-0.397005)--(-0.000000,-0.401713)--(-0.000000,-0.427259)--cycle;
\path[fill=blue!21!white,draw=blue!21!white,line width=.12pt] (-0.000000,-0.427259)--(-0.000000,-0.401713)--(0.191677,-0.397005)--(0.179794,-0.422843)--cycle;
\path[fill=blue!16!white,draw=blue!16!white,line width=.12pt] (-2.426560,-0.730582)--(-2.461751,-0.874991)--(-2.251872,-0.990774)--(-2.219681,-0.844710)--cycle;
\path[fill=blue!22!white,draw=blue!22!white,line width=.12pt] (2.219681,-0.844710)--(2.251872,-0.990774)--(2.461751,-0.874991)--(2.426560,-0.730582)--cycle;
\path[fill=blue!12!white,draw=blue!12!white,line width=.12pt] (-2.091880,-0.218298)--(-2.198794,-0.326193)--(-2.011334,-0.429608)--(-1.913535,-0.316685)--cycle;
\path[fill=blue!15!white,draw=blue!15!white,line width=.12pt] (1.913535,-0.316685)--(2.011334,-0.429608)--(2.198794,-0.326193)--(2.091880,-0.218298)--cycle;
\path[fill=blue!32!white,draw=blue!32!white,line width=.12pt] (-1.398014,-1.797648)--(-1.340867,-1.830352)--(-1.088526,-1.890026)--(-1.134919,-1.859866)--cycle;
\path[fill=blue!34!white,draw=blue!34!white,line width=.12pt] (1.134919,-1.859866)--(1.088526,-1.890026)--(1.340867,-1.830352)--(1.398014,-1.797648)--cycle;
\path[fill=blue!26!white,draw=blue!26!white,line width=.12pt] (-1.991426,-1.478411)--(-1.944527,-1.569617)--(-1.702920,-1.659211)--(-1.743993,-1.570165)--cycle;
\path[fill=blue!31!white,draw=blue!31!white,line width=.12pt] (1.743993,-1.570165)--(1.702920,-1.659211)--(1.944527,-1.569617)--(1.991426,-1.478411)--cycle;
\path[fill=blue!35!white,draw=blue!35!white,line width=.12pt] (-0.825702,-1.937046)--(-0.785553,-1.934893)--(-0.527943,-1.967157)--(-0.554926,-1.970959)--cycle;
\path[fill=blue!36!white,draw=blue!36!white,line width=.12pt] (0.554926,-1.970959)--(0.527943,-1.967157)--(0.785553,-1.934893)--(0.825702,-1.937046)--cycle;
\path[fill=blue!14!white,draw=blue!14!white,line width=.12pt] (-1.697056,-0.155708)--(-1.807416,-0.224200)--(-1.621555,-0.308427)--(-1.522544,-0.234792)--cycle;
\path[fill=blue!14!white,draw=blue!14!white,line width=.12pt] (1.522544,-0.234792)--(1.621555,-0.308427)--(1.807416,-0.224200)--(1.697056,-0.155708)--cycle;
\path[fill=blue!29!white,draw=blue!29!white,line width=.12pt] (-1.702920,-1.659211)--(-1.647647,-1.722837)--(-1.398014,-1.797648)--(-1.444914,-1.736532)--cycle;
\path[fill=blue!32!white,draw=blue!32!white,line width=.12pt] (1.444914,-1.736532)--(1.398014,-1.797648)--(1.647647,-1.722837)--(1.702920,-1.659211)--cycle;
\path[fill=blue!14!white,draw=blue!14!white,line width=.12pt] (-2.369413,-0.587353)--(-2.426560,-0.730582)--(-2.219681,-0.844710)--(-2.167407,-0.698793)--cycle;
\path[fill=blue!20!white,draw=blue!20!white,line width=.12pt] (2.167407,-0.698793)--(2.219681,-0.844710)--(2.426560,-0.730582)--(2.369413,-0.587353)--cycle;
\path[fill=blue!17!white,draw=blue!17!white,line width=.12pt] (-1.163283,-0.230405)--(-1.246659,-0.253371)--(-1.057780,-0.309976)--(-0.987036,-0.283224)--cycle;
\path[fill=blue!16!white,draw=blue!16!white,line width=.12pt] (0.987036,-0.283224)--(1.057780,-0.309976)--(1.246659,-0.253371)--(1.163283,-0.230405)--cycle;
\path[fill=blue!12!white,draw=blue!12!white,line width=.12pt] (-2.198794,-0.326193)--(-2.292506,-0.450808)--(-2.097056,-0.558630)--(-2.011334,-0.429608)--cycle;
\path[fill=blue!16!white,draw=blue!16!white,line width=.12pt] (2.011334,-0.429608)--(2.097056,-0.558630)--(2.292506,-0.450808)--(2.198794,-0.326193)--cycle;
\path[fill=blue!18!white,draw=blue!18!white,line width=.12pt] (0.567664,-0.359610)--(0.607814,-0.361764)--(0.801283,-0.327152)--(0.748354,-0.327284)--cycle;
\path[fill=blue!20!white,draw=blue!20!white,line width=.12pt] (-0.748354,-0.327284)--(-0.801283,-0.327152)--(-0.607814,-0.361764)--(-0.567664,-0.359610)--cycle;
\path[fill=blue!13!white,draw=blue!13!white,line width=.12pt] (-2.292506,-0.450808)--(-2.369413,-0.587353)--(-2.167407,-0.698793)--(-2.097056,-0.558630)--cycle;
\path[fill=blue!18!white,draw=blue!18!white,line width=.12pt] (2.097056,-0.558630)--(2.167407,-0.698793)--(2.369413,-0.587353)--(2.292506,-0.450808)--cycle;
\path[fill=blue!20!white,draw=blue!20!white,line width=.12pt] (-2.262742,-1.131371)--(-2.251872,-1.261099)--(-2.020307,-1.366038)--(-2.030059,-1.236817)--cycle;
\path[fill=blue!26!white,draw=blue!26!white,line width=.12pt] (2.030059,-1.236817)--(2.020307,-1.366038)--(2.251872,-1.261099)--(2.262742,-1.131371)--cycle;
\path[fill=blue!15!white,draw=blue!15!white,line width=.12pt] (-1.423533,-0.187782)--(-1.522544,-0.234792)--(-1.333369,-0.304943)--(-1.246659,-0.253371)--cycle;
\path[fill=blue!15!white,draw=blue!15!white,line width=.12pt] (1.246659,-0.253371)--(1.333369,-0.304943)--(1.522544,-0.234792)--(1.423533,-0.187782)--cycle;
\path[fill=blue!35!white,draw=blue!35!white,line width=.12pt] (-0.554926,-1.970959)--(-0.527943,-1.967157)--(-0.265249,-1.986640)--(-0.278805,-1.991439)--cycle;
\path[fill=blue!35!white,draw=blue!35!white,line width=.12pt] (0.278805,-1.991439)--(0.265249,-1.986640)--(0.527943,-1.967157)--(0.554926,-1.970959)--cycle;
\path[fill=blue!35!white,draw=blue!35!white,line width=.12pt] (-0.278805,-1.991439)--(-0.265249,-1.986640)--(-0.000000,-1.993156)--(-0.000000,-1.998287)--cycle;
\path[fill=blue!35!white,draw=blue!35!white,line width=.12pt] (-0.000000,-1.998287)--(-0.000000,-1.993156)--(0.265249,-1.986640)--(0.278805,-1.991439)--cycle;
\path[fill=blue!19!white,draw=blue!19!white,line width=.12pt] (-0.567664,-0.359610)--(-0.607814,-0.361764)--(-0.408491,-0.386728)--(-0.381508,-0.382925)--cycle;
\path[fill=blue!18!white,draw=blue!18!white,line width=.12pt] (0.381508,-0.382925)--(0.408491,-0.386728)--(0.607814,-0.361764)--(0.567664,-0.359610)--cycle;
\path[fill=blue!34!white,draw=blue!34!white,line width=.12pt] (0.860893,-1.908890)--(0.825702,-1.937046)--(1.088526,-1.890026)--(1.134919,-1.859866)--cycle;
\path[fill=blue!32!white,draw=blue!32!white,line width=.12pt] (-1.134919,-1.859866)--(-1.088526,-1.890026)--(-0.825702,-1.937046)--(-0.860893,-1.908890)--cycle;
\path[fill=blue!13!white,draw=blue!13!white,line width=.12pt] (-1.807416,-0.224200)--(-1.913535,-0.316685)--(-1.716761,-0.405857)--(-1.621555,-0.308427)--cycle;
\path[fill=blue!14!white,draw=blue!14!white,line width=.12pt] (1.621555,-0.308427)--(1.716761,-0.405857)--(1.913535,-0.316685)--(1.807416,-0.224200)--cycle;
\path[fill=blue!17!white,draw=blue!17!white,line width=.12pt] (-0.987036,-0.283224)--(-1.057780,-0.309976)--(-0.858714,-0.357051)--(-0.801283,-0.327152)--cycle;
\path[fill=blue!16!white,draw=blue!16!white,line width=.12pt] (0.801283,-0.327152)--(0.858714,-0.357051)--(1.057780,-0.309976)--(0.987036,-0.283224)--cycle;
\path[fill=blue!23!white,draw=blue!23!white,line width=.12pt] (-2.020307,-1.366038)--(-1.991426,-1.478411)--(-1.743993,-1.570165)--(-1.769285,-1.459123)--cycle;
\path[fill=blue!28!white,draw=blue!28!white,line width=.12pt] (1.769285,-1.459123)--(1.743993,-1.570165)--(1.991426,-1.478411)--(2.020307,-1.366038)--cycle;
\path[fill=blue!19!white,draw=blue!19!white,line width=.12pt] (-0.381508,-0.382925)--(-0.408491,-0.386728)--(-0.205234,-0.401803)--(-0.191677,-0.397005)--cycle;
\path[fill=blue!19!white,draw=blue!19!white,line width=.12pt] (0.191677,-0.397005)--(0.205234,-0.401803)--(0.408491,-0.386728)--(0.381508,-0.382925)--cycle;
\path[fill=blue!18!white,draw=blue!18!white,line width=.12pt] (-2.251872,-0.990774)--(-2.262742,-1.131371)--(-2.030059,-1.236817)--(-2.020307,-1.095713)--cycle;
\path[fill=blue!24!white,draw=blue!24!white,line width=.12pt] (2.020307,-1.095713)--(2.030059,-1.236817)--(2.262742,-1.131371)--(2.251872,-0.990774)--cycle;
\path[fill=blue!19!white,draw=blue!19!white,line width=.12pt] (-0.191677,-0.397005)--(-0.205234,-0.401803)--(-0.000000,-0.406844)--(-0.000000,-0.401713)--cycle;
\path[fill=blue!19!white,draw=blue!19!white,line width=.12pt] (-0.000000,-0.401713)--(-0.000000,-0.406844)--(0.205234,-0.401803)--(0.191677,-0.397005)--cycle;
\path[fill=blue!29!white,draw=blue!29!white,line width=.12pt] (-1.444914,-1.736532)--(-1.398014,-1.797648)--(-1.134919,-1.859866)--(-1.172992,-1.800837)--cycle;
\path[fill=blue!32!white,draw=blue!32!white,line width=.12pt] (1.172992,-1.800837)--(1.134919,-1.859866)--(1.398014,-1.797648)--(1.444914,-1.736532)--cycle;
\path[fill=blue!26!white,draw=blue!26!white,line width=.12pt] (-1.743993,-1.570165)--(-1.702920,-1.659211)--(-1.444914,-1.736532)--(-1.479763,-1.649352)--cycle;
\path[fill=blue!30!white,draw=blue!30!white,line width=.12pt] (1.479763,-1.649352)--(1.444914,-1.736532)--(1.702920,-1.659211)--(1.743993,-1.570165)--cycle;
\path[fill=blue!14!white,draw=blue!14!white,line width=.12pt] (-1.522544,-0.234792)--(-1.621555,-0.308427)--(-1.420078,-0.383140)--(-1.333369,-0.304943)--cycle;
\path[fill=blue!14!white,draw=blue!14!white,line width=.12pt] (1.333369,-0.304943)--(1.420078,-0.383140)--(1.621555,-0.308427)--(1.522544,-0.234792)--cycle;
\path[fill=blue!15!white,draw=blue!15!white,line width=.12pt] (-1.246659,-0.253371)--(-1.333369,-0.304943)--(-1.131352,-0.365485)--(-1.057780,-0.309976)--cycle;
\path[fill=blue!15!white,draw=blue!15!white,line width=.12pt] (1.057780,-0.309976)--(1.131352,-0.365485)--(1.333369,-0.304943)--(1.246659,-0.253371)--cycle;
\path[fill=blue!12!white,draw=blue!12!white,line width=.12pt] (-1.913535,-0.316685)--(-2.011334,-0.429608)--(-1.804504,-0.523338)--(-1.716761,-0.405857)--cycle;
\path[fill=blue!15!white,draw=blue!15!white,line width=.12pt] (1.716761,-0.405857)--(1.804504,-0.523338)--(2.011334,-0.429608)--(1.913535,-0.316685)--cycle;
\path[fill=blue!32!white,draw=blue!32!white,line width=.12pt] (-0.860893,-1.908890)--(-0.825702,-1.937046)--(-0.554926,-1.970959)--(-0.578577,-1.944248)--cycle;
\path[fill=blue!33!white,draw=blue!33!white,line width=.12pt] (0.578577,-1.944248)--(0.554926,-1.970959)--(0.825702,-1.937046)--(0.860893,-1.908890)--cycle;
\path[fill=blue!16!white,draw=blue!16!white,line width=.12pt] (-2.219681,-0.844710)--(-2.251872,-0.990774)--(-2.020307,-1.095713)--(-1.991426,-0.948149)--cycle;
\path[fill=blue!21!white,draw=blue!21!white,line width=.12pt] (1.991426,-0.948149)--(2.020307,-1.095713)--(2.251872,-0.990774)--(2.219681,-0.844710)--cycle;
\path[fill=blue!17!white,draw=blue!17!white,line width=.12pt] (-0.801283,-0.327152)--(-0.858714,-0.357051)--(-0.651378,-0.394145)--(-0.607814,-0.361764)--cycle;
\path[fill=blue!16!white,draw=blue!16!white,line width=.12pt] (0.607814,-0.361764)--(0.651378,-0.394145)--(0.858714,-0.357051)--(0.801283,-0.327152)--cycle;
\path[fill=blue!12!white,draw=blue!12!white,line width=.12pt] (-2.011334,-0.429608)--(-2.097056,-0.558630)--(-1.881411,-0.656355)--(-1.804504,-0.523338)--cycle;
\path[fill=blue!16!white,draw=blue!16!white,line width=.12pt] (1.804504,-0.523338)--(1.881411,-0.656355)--(2.097056,-0.558630)--(2.011334,-0.429608)--cycle;
\path[fill=blue!14!white,draw=blue!14!white,line width=.12pt] (-2.167407,-0.698793)--(-2.219681,-0.844710)--(-1.991426,-0.948149)--(-1.944527,-0.799796)--cycle;
\path[fill=blue!19!white,draw=blue!19!white,line width=.12pt] (1.944527,-0.799796)--(1.991426,-0.948149)--(2.219681,-0.844710)--(2.167407,-0.698793)--cycle;
\path[fill=blue!32!white,draw=blue!32!white,line width=.12pt] (-0.578577,-1.944248)--(-0.554926,-1.970959)--(-0.278805,-1.991439)--(-0.290688,-1.965600)--cycle;
\path[fill=blue!33!white,draw=blue!33!white,line width=.12pt] (0.290688,-1.965600)--(0.278805,-1.991439)--(0.554926,-1.970959)--(0.578577,-1.944248)--cycle;
\path[fill=blue!20!white,draw=blue!20!white,line width=.12pt] (-2.030059,-1.236817)--(-2.020307,-1.366038)--(-1.769285,-1.459123)--(-1.777825,-1.330351)--cycle;
\path[fill=blue!26!white,draw=blue!26!white,line width=.12pt] (1.777825,-1.330351)--(1.769285,-1.459123)--(2.020307,-1.366038)--(2.030059,-1.236817)--cycle;
\path[fill=blue!13!white,draw=blue!13!white,line width=.12pt] (-2.097056,-0.558630)--(-2.167407,-0.698793)--(-1.944527,-0.799796)--(-1.881411,-0.656355)--cycle;
\path[fill=blue!17!white,draw=blue!17!white,line width=.12pt] (1.881411,-0.656355)--(1.944527,-0.799796)--(2.167407,-0.698793)--(2.097056,-0.558630)--cycle;
\path[fill=blue!17!white,draw=blue!17!white,line width=.12pt] (-0.607814,-0.361764)--(-0.651378,-0.394145)--(-0.437769,-0.420898)--(-0.408491,-0.386728)--cycle;
\path[fill=blue!17!white,draw=blue!17!white,line width=.12pt] (0.408491,-0.386728)--(0.437769,-0.420898)--(0.651378,-0.394145)--(0.607814,-0.361764)--cycle;
\path[fill=blue!32!white,draw=blue!32!white,line width=.12pt] (-0.290688,-1.965600)--(-0.278805,-1.991439)--(-0.000000,-1.998287)--(-0.000000,-1.972741)--cycle;
\path[fill=blue!33!white,draw=blue!33!white,line width=.12pt] (-0.000000,-1.972741)--(-0.000000,-1.998287)--(0.278805,-1.991439)--(0.290688,-1.965600)--cycle;
\path[fill=blue!29!white,draw=blue!29!white,line width=.12pt] (-1.172992,-1.800837)--(-1.134919,-1.859866)--(-0.860893,-1.908890)--(-0.889774,-1.851505)--cycle;
\path[fill=blue!31!white,draw=blue!31!white,line width=.12pt] (0.889774,-1.851505)--(0.860893,-1.908890)--(1.134919,-1.859866)--(1.172992,-1.800837)--cycle;
\path[fill=blue!13!white,draw=blue!13!white,line width=.12pt] (-1.621555,-0.308427)--(-1.716761,-0.405857)--(-1.503455,-0.484957)--(-1.420078,-0.383140)--cycle;
\path[fill=blue!14!white,draw=blue!14!white,line width=.12pt] (1.420078,-0.383140)--(1.503455,-0.484957)--(1.716761,-0.405857)--(1.621555,-0.308427)--cycle;
\path[fill=blue!15!white,draw=blue!15!white,line width=.12pt] (-1.057780,-0.309976)--(-1.131352,-0.365485)--(-0.918440,-0.415835)--(-0.858714,-0.357051)--cycle;
\path[fill=blue!15!white,draw=blue!15!white,line width=.12pt] (0.858714,-0.357051)--(0.918440,-0.415835)--(1.131352,-0.365485)--(1.057780,-0.309976)--cycle;
\path[fill=blue!23!white,draw=blue!23!white,line width=.12pt] (-1.769285,-1.459123)--(-1.743993,-1.570165)--(-1.479763,-1.649352)--(-1.501223,-1.539458)--cycle;
\path[fill=blue!27!white,draw=blue!27!white,line width=.12pt] (1.501223,-1.539458)--(1.479763,-1.649352)--(1.743993,-1.570165)--(1.769285,-1.459123)--cycle;
\path[fill=blue!14!white,draw=blue!14!white,line width=.12pt] (-1.333369,-0.304943)--(-1.420078,-0.383140)--(-1.204924,-0.447619)--(-1.131352,-0.365485)--cycle;
\path[fill=blue!14!white,draw=blue!14!white,line width=.12pt] (1.131352,-0.365485)--(1.204924,-0.447619)--(1.420078,-0.383140)--(1.333369,-0.304943)--cycle;
\path[fill=blue!17!white,draw=blue!17!white,line width=.12pt] (0.205234,-0.401803)--(0.219943,-0.437053)--(0.437769,-0.420898)--(0.408491,-0.386728)--cycle;
\path[fill=blue!17!white,draw=blue!17!white,line width=.12pt] (-0.408491,-0.386728)--(-0.437769,-0.420898)--(-0.219943,-0.437053)--(-0.205234,-0.401803)--cycle;
\path[fill=blue!26!white,draw=blue!26!white,line width=.12pt] (-1.479763,-1.649352)--(-1.444914,-1.736532)--(-1.172992,-1.800837)--(-1.201283,-1.715208)--cycle;
\path[fill=blue!29!white,draw=blue!29!white,line width=.12pt] (1.201283,-1.715208)--(1.172992,-1.800837)--(1.444914,-1.736532)--(1.479763,-1.649352)--cycle;
\path[fill=blue!17!white,draw=blue!17!white,line width=.12pt] (-0.205234,-0.401803)--(-0.219943,-0.437053)--(-0.000000,-0.442456)--(-0.000000,-0.406844)--cycle;
\path[fill=blue!17!white,draw=blue!17!white,line width=.12pt] (-0.000000,-0.406844)--(-0.000000,-0.442456)--(0.219943,-0.437053)--(0.205234,-0.401803)--cycle;
\path[fill=blue!18!white,draw=blue!18!white,line width=.12pt] (-2.020307,-1.095713)--(-2.030059,-1.236817)--(-1.777825,-1.330351)--(-1.769285,-1.188798)--cycle;
\path[fill=blue!23!white,draw=blue!23!white,line width=.12pt] (1.769285,-1.188798)--(1.777825,-1.330351)--(2.030059,-1.236817)--(2.020307,-1.095713)--cycle;
\path[fill=blue!29!white,draw=blue!29!white,line width=.12pt] (-0.889774,-1.851505)--(-0.860893,-1.908890)--(-0.578577,-1.944248)--(-0.597986,-1.888050)--cycle;
\path[fill=blue!31!white,draw=blue!31!white,line width=.12pt] (0.597986,-1.888050)--(0.578577,-1.944248)--(0.860893,-1.908890)--(0.889774,-1.851505)--cycle;
\path[fill=blue!12!white,draw=blue!12!white,line width=.12pt] (-1.716761,-0.405857)--(-1.804504,-0.523338)--(-1.580295,-0.606480)--(-1.503455,-0.484957)--cycle;
\path[fill=blue!15!white,draw=blue!15!white,line width=.12pt] (1.503455,-0.484957)--(1.580295,-0.606480)--(1.804504,-0.523338)--(1.716761,-0.405857)--cycle;
\path[fill=blue!15!white,draw=blue!15!white,line width=.12pt] (-0.858714,-0.357051)--(-0.918440,-0.415835)--(-0.696683,-0.455508)--(-0.651378,-0.394145)--cycle;
\path[fill=blue!15!white,draw=blue!15!white,line width=.12pt] (0.651378,-0.394145)--(0.696683,-0.455508)--(0.918440,-0.415835)--(0.858714,-0.357051)--cycle;
\path[fill=blue!16!white,draw=blue!16!white,line width=.12pt] (-1.991426,-0.948149)--(-2.020307,-1.095713)--(-1.769285,-1.188798)--(-1.743993,-1.039904)--cycle;
\path[fill=blue!21!white,draw=blue!21!white,line width=.12pt] (1.743993,-1.039904)--(1.769285,-1.188798)--(2.020307,-1.095713)--(1.991426,-0.948149)--cycle;
\path[fill=blue!12!white,draw=blue!12!white,line width=.12pt] (-1.804504,-0.523338)--(-1.881411,-0.656355)--(-1.647647,-0.743041)--(-1.580295,-0.606480)--cycle;
\path[fill=blue!16!white,draw=blue!16!white,line width=.12pt] (1.580295,-0.606480)--(1.647647,-0.743041)--(1.881411,-0.656355)--(1.804504,-0.523338)--cycle;
\path[fill=blue!29!white,draw=blue!29!white,line width=.12pt] (-0.597986,-1.888050)--(-0.578577,-1.944248)--(-0.290688,-1.965600)--(-0.300440,-1.910118)--cycle;
\path[fill=blue!30!white,draw=blue!30!white,line width=.12pt] (0.300440,-1.910118)--(0.290688,-1.965600)--(0.578577,-1.944248)--(0.597986,-1.888050)--cycle;
\path[fill=blue!20!white,draw=blue!20!white,line width=.12pt] (-1.777825,-1.330351)--(-1.769285,-1.459123)--(-1.501223,-1.539458)--(-1.508470,-1.411074)--cycle;
\path[fill=blue!25!white,draw=blue!25!white,line width=.12pt] (1.508470,-1.411074)--(1.501223,-1.539458)--(1.769285,-1.459123)--(1.777825,-1.330351)--cycle;
\path[fill=blue!13!white,draw=blue!13!white,line width=.12pt] (-1.420078,-0.383140)--(-1.503455,-0.484957)--(-1.275669,-0.553222)--(-1.204924,-0.447619)--cycle;
\path[fill=blue!14!white,draw=blue!14!white,line width=.12pt] (1.204924,-0.447619)--(1.275669,-0.553222)--(1.503455,-0.484957)--(1.420078,-0.383140)--cycle;
\path[fill=blue!14!white,draw=blue!14!white,line width=.12pt] (-1.131352,-0.365485)--(-1.204924,-0.447619)--(-0.978167,-0.501243)--(-0.918440,-0.415835)--cycle;
\path[fill=blue!14!white,draw=blue!14!white,line width=.12pt] (0.918440,-0.415835)--(0.978167,-0.501243)--(1.204924,-0.447619)--(1.131352,-0.365485)--cycle;
\path[fill=blue!26!white,draw=blue!26!white,line width=.12pt] (-1.201283,-1.715208)--(-1.172992,-1.800837)--(-0.889774,-1.851505)--(-0.911234,-1.767098)--cycle;
\path[fill=blue!29!white,draw=blue!29!white,line width=.12pt] (0.911234,-1.767098)--(0.889774,-1.851505)--(1.172992,-1.800837)--(1.201283,-1.715208)--cycle;
\path[fill=blue!14!white,draw=blue!14!white,line width=.12pt] (-1.944527,-0.799796)--(-1.991426,-0.948149)--(-1.743993,-1.039904)--(-1.702920,-0.889390)--cycle;
\path[fill=blue!19!white,draw=blue!19!white,line width=.12pt] (1.702920,-0.889390)--(1.743993,-1.039904)--(1.991426,-0.948149)--(1.944527,-0.799796)--cycle;
\path[fill=blue!15!white,draw=blue!15!white,line width=.12pt] (-0.651378,-0.394145)--(-0.696683,-0.455508)--(-0.468217,-0.484122)--(-0.437769,-0.420898)--cycle;
\path[fill=blue!15!white,draw=blue!15!white,line width=.12pt] (0.437769,-0.420898)--(0.468217,-0.484122)--(0.696683,-0.455508)--(0.651378,-0.394145)--cycle;
\path[fill=blue!13!white,draw=blue!13!white,line width=.12pt] (-1.881411,-0.656355)--(-1.944527,-0.799796)--(-1.702920,-0.889390)--(-1.647647,-0.743041)--cycle;
\path[fill=blue!17!white,draw=blue!17!white,line width=.12pt] (1.647647,-0.743041)--(1.702920,-0.889390)--(1.944527,-0.799796)--(1.881411,-0.656355)--cycle;
\path[fill=blue!30!white,draw=blue!30!white,line width=.12pt] (-0.300440,-1.910118)--(-0.290688,-1.965600)--(-0.000000,-1.972741)--(-0.000000,-1.917498)--cycle;
\path[fill=blue!30!white,draw=blue!30!white,line width=.12pt] (-0.000000,-1.917498)--(-0.000000,-1.972741)--(0.290688,-1.965600)--(0.300440,-1.910118)--cycle;
\path[fill=blue!23!white,draw=blue!23!white,line width=.12pt] (-1.501223,-1.539458)--(-1.479763,-1.649352)--(-1.201283,-1.715208)--(-1.218704,-1.606269)--cycle;
\path[fill=blue!27!white,draw=blue!27!white,line width=.12pt] (1.218704,-1.606269)--(1.201283,-1.715208)--(1.479763,-1.649352)--(1.501223,-1.539458)--cycle;
\path[fill=blue!15!white,draw=blue!15!white,line width=.12pt] (-0.437769,-0.420898)--(-0.468217,-0.484122)--(-0.235241,-0.501401)--(-0.219943,-0.437053)--cycle;
\path[fill=blue!15!white,draw=blue!15!white,line width=.12pt] (0.219943,-0.437053)--(0.235241,-0.501401)--(0.468217,-0.484122)--(0.437769,-0.420898)--cycle;
\path[fill=blue!15!white,draw=blue!15!white,line width=.12pt] (-0.219943,-0.437053)--(-0.235241,-0.501401)--(-0.000000,-0.507180)--(-0.000000,-0.442456)--cycle;
\path[fill=blue!15!white,draw=blue!15!white,line width=.12pt] (-0.000000,-0.442456)--(-0.000000,-0.507180)--(0.235241,-0.501401)--(0.219943,-0.437053)--cycle;
\path[fill=blue!18!white,draw=blue!18!white,line width=.12pt] (-1.769285,-1.188798)--(-1.777825,-1.330351)--(-1.508470,-1.411074)--(-1.501223,-1.269133)--cycle;
\path[fill=blue!22!white,draw=blue!22!white,line width=.12pt] (1.501223,-1.269133)--(1.508470,-1.411074)--(1.777825,-1.330351)--(1.769285,-1.188798)--cycle;
\path[fill=blue!12!white,draw=blue!12!white,line width=.12pt] (-1.503455,-0.484957)--(-1.580295,-0.606480)--(-1.340867,-0.678234)--(-1.275669,-0.553222)--cycle;
\path[fill=blue!14!white,draw=blue!14!white,line width=.12pt] (1.275669,-0.553222)--(1.340867,-0.678234)--(1.580295,-0.606480)--(1.503455,-0.484957)--cycle;
\path[fill=blue!14!white,draw=blue!14!white,line width=.12pt] (-0.918440,-0.415835)--(-0.978167,-0.501243)--(-0.741989,-0.543496)--(-0.696683,-0.455508)--cycle;
\path[fill=blue!14!white,draw=blue!14!white,line width=.12pt] (0.696683,-0.455508)--(0.741989,-0.543496)--(0.978167,-0.501243)--(0.918440,-0.415835)--cycle;
\path[fill=blue!26!white,draw=blue!26!white,line width=.12pt] (-0.911234,-1.767098)--(-0.889774,-1.851505)--(-0.597986,-1.888050)--(-0.612409,-1.804524)--cycle;
\path[fill=blue!28!white,draw=blue!28!white,line width=.12pt] (0.612409,-1.804524)--(0.597986,-1.888050)--(0.889774,-1.851505)--(0.911234,-1.767098)--cycle;
\path[fill=blue!13!white,draw=blue!13!white,line width=.12pt] (-1.204924,-0.447619)--(-1.275669,-0.553222)--(-1.035598,-0.609994)--(-0.978167,-0.501243)--cycle;
\path[fill=blue!14!white,draw=blue!14!white,line width=.12pt] (0.978167,-0.501243)--(1.035598,-0.609994)--(1.275669,-0.553222)--(1.204924,-0.447619)--cycle;
\path[fill=blue!20!white,draw=blue!20!white,line width=.12pt] (-1.508470,-1.411074)--(-1.501223,-1.539458)--(-1.218704,-1.606269)--(-1.224587,-1.478207)--cycle;
\path[fill=blue!24!white,draw=blue!24!white,line width=.12pt] (1.224587,-1.478207)--(1.218704,-1.606269)--(1.501223,-1.539458)--(1.508470,-1.411074)--cycle;
\path[fill=blue!23!white,draw=blue!23!white,line width=.12pt] (-1.218704,-1.606269)--(-1.201283,-1.715208)--(-0.911234,-1.767098)--(-0.924449,-1.658912)--cycle;
\path[fill=blue!26!white,draw=blue!26!white,line width=.12pt] (0.924449,-1.658912)--(0.911234,-1.767098)--(1.201283,-1.715208)--(1.218704,-1.606269)--cycle;
\path[fill=blue!16!white,draw=blue!16!white,line width=.12pt] (-1.743993,-1.039904)--(-1.769285,-1.188798)--(-1.501223,-1.269133)--(-1.479763,-1.119090)--cycle;
\path[fill=blue!20!white,draw=blue!20!white,line width=.12pt] (1.479763,-1.119090)--(1.501223,-1.269133)--(1.769285,-1.188798)--(1.743993,-1.039904)--cycle;
\path[fill=blue!12!white,draw=blue!12!white,line width=.12pt] (-1.580295,-0.606480)--(-1.647647,-0.743041)--(-1.398014,-0.817853)--(-1.340867,-0.678234)--cycle;
\path[fill=blue!15!white,draw=blue!15!white,line width=.12pt] (1.340867,-0.678234)--(1.398014,-0.817853)--(1.647647,-0.743041)--(1.580295,-0.606480)--cycle;
\path[fill=blue!26!white,draw=blue!26!white,line width=.12pt] (-0.612409,-1.804524)--(-0.597986,-1.888050)--(-0.300440,-1.910118)--(-0.307686,-1.827125)--cycle;
\path[fill=blue!28!white,draw=blue!28!white,line width=.12pt] (0.307686,-1.827125)--(0.300440,-1.910118)--(0.597986,-1.888050)--(0.612409,-1.804524)--cycle;
\path[fill=blue!14!white,draw=blue!14!white,line width=.12pt] (-0.696683,-0.455508)--(-0.741989,-0.543496)--(-0.498665,-0.573971)--(-0.468217,-0.484122)--cycle;
\path[fill=blue!14!white,draw=blue!14!white,line width=.12pt] (0.468217,-0.484122)--(0.498665,-0.573971)--(0.741989,-0.543496)--(0.696683,-0.455508)--cycle;
\path[fill=blue!14!white,draw=blue!14!white,line width=.12pt] (-1.702920,-0.889390)--(-1.743993,-1.039904)--(-1.479763,-1.119090)--(-1.444914,-0.966711)--cycle;
\path[fill=blue!18!white,draw=blue!18!white,line width=.12pt] (1.444914,-0.966711)--(1.479763,-1.119090)--(1.743993,-1.039904)--(1.702920,-0.889390)--cycle;
\path[fill=blue!13!white,draw=blue!13!white,line width=.12pt] (-1.647647,-0.743041)--(-1.702920,-0.889390)--(-1.444914,-0.966711)--(-1.398014,-0.817853)--cycle;
\path[fill=blue!16!white,draw=blue!16!white,line width=.12pt] (1.398014,-0.817853)--(1.444914,-0.966711)--(1.702920,-0.889390)--(1.647647,-0.743041)--cycle;
\path[fill=blue!27!white,draw=blue!27!white,line width=.12pt] (-0.307686,-1.827125)--(-0.300440,-1.910118)--(-0.000000,-1.917498)--(-0.000000,-1.834683)--cycle;
\path[fill=blue!27!white,draw=blue!27!white,line width=.12pt] (-0.000000,-1.834683)--(-0.000000,-1.917498)--(0.300440,-1.910118)--(0.307686,-1.827125)--cycle;
\path[fill=blue!14!white,draw=blue!14!white,line width=.12pt] (-0.468217,-0.484122)--(-0.498665,-0.573971)--(-0.250539,-0.592374)--(-0.235241,-0.501401)--cycle;
\path[fill=blue!14!white,draw=blue!14!white,line width=.12pt] (0.235241,-0.501401)--(0.250539,-0.592374)--(0.498665,-0.573971)--(0.468217,-0.484122)--cycle;
\path[fill=blue!12!white,draw=blue!12!white,line width=.12pt] (-1.275669,-0.553222)--(-1.340867,-0.678234)--(-1.088526,-0.737908)--(-1.035598,-0.609994)--cycle;
\path[fill=blue!14!white,draw=blue!14!white,line width=.12pt] (1.035598,-0.609994)--(1.088526,-0.737908)--(1.340867,-0.678234)--(1.275669,-0.553222)--cycle;
\path[fill=blue!13!white,draw=blue!13!white,line width=.12pt] (-0.978167,-0.501243)--(-1.035598,-0.609994)--(-0.785553,-0.654728)--(-0.741989,-0.543496)--cycle;
\path[fill=blue!14!white,draw=blue!14!white,line width=.12pt] (0.741989,-0.543496)--(0.785553,-0.654728)--(1.035598,-0.609994)--(0.978167,-0.501243)--cycle;
\path[fill=blue!18!white,draw=blue!18!white,line width=.12pt] (-1.501223,-1.269133)--(-1.508470,-1.411074)--(-1.224587,-1.478207)--(-1.218704,-1.335944)--cycle;
\path[fill=blue!22!white,draw=blue!22!white,line width=.12pt] (1.218704,-1.335944)--(1.224587,-1.478207)--(1.508470,-1.411074)--(1.501223,-1.269133)--cycle;
\path[fill=blue!14!white,draw=blue!14!white,line width=.12pt] (-0.235241,-0.501401)--(-0.250539,-0.592374)--(-0.000000,-0.598528)--(-0.000000,-0.507180)--cycle;
\path[fill=blue!14!white,draw=blue!14!white,line width=.12pt] (-0.000000,-0.507180)--(-0.000000,-0.598528)--(0.250539,-0.592374)--(0.235241,-0.501401)--cycle;
\path[fill=blue!23!white,draw=blue!23!white,line width=.12pt] (-0.924449,-1.658912)--(-0.911234,-1.767098)--(-0.612409,-1.804524)--(-0.621290,-1.696881)--cycle;
\path[fill=blue!25!white,draw=blue!25!white,line width=.12pt] (0.621290,-1.696881)--(0.612409,-1.804524)--(0.911234,-1.767098)--(0.924449,-1.658912)--cycle;
\path[fill=blue!20!white,draw=blue!20!white,line width=.12pt] (-1.224587,-1.478207)--(-1.218704,-1.606269)--(-0.924449,-1.658912)--(-0.928911,-1.531105)--cycle;
\path[fill=blue!23!white,draw=blue!23!white,line width=.12pt] (0.928911,-1.531105)--(0.924449,-1.658912)--(1.218704,-1.606269)--(1.224587,-1.478207)--cycle;
\path[fill=blue!12!white,draw=blue!12!white,line width=.12pt] (-1.340867,-0.678234)--(-1.398014,-0.817853)--(-1.134919,-0.880070)--(-1.088526,-0.737908)--cycle;
\path[fill=blue!15!white,draw=blue!15!white,line width=.12pt] (1.088526,-0.737908)--(1.134919,-0.880070)--(1.398014,-0.817853)--(1.340867,-0.678234)--cycle;
\path[fill=blue!16!white,draw=blue!16!white,line width=.12pt] (-1.479763,-1.119090)--(-1.501223,-1.269133)--(-1.218704,-1.335944)--(-1.201283,-1.184946)--cycle;
\path[fill=blue!19!white,draw=blue!19!white,line width=.12pt] (1.201283,-1.184946)--(1.218704,-1.335944)--(1.501223,-1.269133)--(1.479763,-1.119090)--cycle;
\path[fill=blue!13!white,draw=blue!13!white,line width=.12pt] (-0.741989,-0.543496)--(-0.785553,-0.654728)--(-0.527943,-0.686992)--(-0.498665,-0.573971)--cycle;
\path[fill=blue!13!white,draw=blue!13!white,line width=.12pt] (0.498665,-0.573971)--(0.527943,-0.686992)--(0.785553,-0.654728)--(0.741989,-0.543496)--cycle;
\path[fill=blue!24!white,draw=blue!24!white,line width=.12pt] (-0.621290,-1.696881)--(-0.612409,-1.804524)--(-0.307686,-1.827125)--(-0.312148,-1.719809)--cycle;
\path[fill=blue!25!white,draw=blue!25!white,line width=.12pt] (0.312148,-1.719809)--(0.307686,-1.827125)--(0.612409,-1.804524)--(0.621290,-1.696881)--cycle;
\path[fill=blue!13!white,draw=blue!13!white,line width=.12pt] (-1.398014,-0.817853)--(-1.444914,-0.966711)--(-1.172992,-1.031016)--(-1.134919,-0.880070)--cycle;
\path[fill=blue!16!white,draw=blue!16!white,line width=.12pt] (1.134919,-0.880070)--(1.172992,-1.031016)--(1.444914,-0.966711)--(1.398014,-0.817853)--cycle;
\path[fill=blue!12!white,draw=blue!12!white,line width=.12pt] (-1.035598,-0.609994)--(-1.088526,-0.737908)--(-0.825702,-0.784928)--(-0.785553,-0.654728)--cycle;
\path[fill=blue!14!white,draw=blue!14!white,line width=.12pt] (0.785553,-0.654728)--(0.825702,-0.784928)--(1.088526,-0.737908)--(1.035598,-0.609994)--cycle;
\path[fill=blue!14!white,draw=blue!14!white,line width=.12pt] (-1.444914,-0.966711)--(-1.479763,-1.119090)--(-1.201283,-1.184946)--(-1.172992,-1.031016)--cycle;
\path[fill=blue!17!white,draw=blue!17!white,line width=.12pt] (1.172992,-1.031016)--(1.201283,-1.184946)--(1.479763,-1.119090)--(1.444914,-0.966711)--cycle;
\path[fill=blue!24!white,draw=blue!24!white,line width=.12pt] (-0.312148,-1.719809)--(-0.307686,-1.827125)--(-0.000000,-1.834683)--(-0.000000,-1.727477)--cycle;
\path[fill=blue!24!white,draw=blue!24!white,line width=.12pt] (-0.000000,-1.727477)--(-0.000000,-1.834683)--(0.307686,-1.827125)--(0.312148,-1.719809)--cycle;
\path[fill=blue!13!white,draw=blue!13!white,line width=.12pt] (-0.498665,-0.573971)--(-0.527943,-0.686992)--(-0.265249,-0.706475)--(-0.250539,-0.592374)--cycle;
\path[fill=blue!13!white,draw=blue!13!white,line width=.12pt] (0.250539,-0.592374)--(0.265249,-0.706475)--(0.527943,-0.686992)--(0.498665,-0.573971)--cycle;
\path[fill=blue!18!white,draw=blue!18!white,line width=.12pt] (-1.218704,-1.335944)--(-1.224587,-1.478207)--(-0.928911,-1.531105)--(-0.924449,-1.388587)--cycle;
\path[fill=blue!21!white,draw=blue!21!white,line width=.12pt] (0.924449,-1.388587)--(0.928911,-1.531105)--(1.224587,-1.478207)--(1.218704,-1.335944)--cycle;
\path[fill=blue!21!white,draw=blue!21!white,line width=.12pt] (-0.928911,-1.531105)--(-0.924449,-1.658912)--(-0.621290,-1.696881)--(-0.624289,-1.569256)--cycle;
\path[fill=blue!23!white,draw=blue!23!white,line width=.12pt] (0.624289,-1.569256)--(0.621290,-1.696881)--(0.924449,-1.658912)--(0.928911,-1.531105)--cycle;
\path[fill=blue!13!white,draw=blue!13!white,line width=.12pt] (-0.250539,-0.592374)--(-0.265249,-0.706475)--(-0.000000,-0.712991)--(-0.000000,-0.598528)--cycle;
\path[fill=blue!13!white,draw=blue!13!white,line width=.12pt] (-0.000000,-0.598528)--(-0.000000,-0.712991)--(0.265249,-0.706475)--(0.250539,-0.592374)--cycle;
\path[fill=blue!14!white,draw=blue!14!white,line width=.12pt] (0.825702,-0.784928)--(0.860893,-0.929094)--(1.134919,-0.880070)--(1.088526,-0.737908)--cycle;
\path[fill=blue!12!white,draw=blue!12!white,line width=.12pt] (-1.088526,-0.737908)--(-1.134919,-0.880070)--(-0.860893,-0.929094)--(-0.825702,-0.784928)--cycle;
\path[fill=blue!12!white,draw=blue!12!white,line width=.12pt] (-0.785553,-0.654728)--(-0.825702,-0.784928)--(-0.554926,-0.818841)--(-0.527943,-0.686992)--cycle;
\path[fill=blue!13!white,draw=blue!13!white,line width=.12pt] (0.527943,-0.686992)--(0.554926,-0.818841)--(0.825702,-0.784928)--(0.785553,-0.654728)--cycle;
\path[fill=blue!16!white,draw=blue!16!white,line width=.12pt] (-1.201283,-1.184946)--(-1.218704,-1.335944)--(-0.924449,-1.388587)--(-0.911234,-1.236837)--cycle;
\path[fill=blue!19!white,draw=blue!19!white,line width=.12pt] (0.911234,-1.236837)--(0.924449,-1.388587)--(1.218704,-1.335944)--(1.201283,-1.184946)--cycle;
\path[fill=blue!21!white,draw=blue!21!white,line width=.12pt] (-0.624289,-1.569256)--(-0.621290,-1.696881)--(-0.312148,-1.719809)--(-0.313655,-1.592296)--cycle;
\path[fill=blue!22!white,draw=blue!22!white,line width=.12pt] (0.313655,-1.592296)--(0.312148,-1.719809)--(0.621290,-1.696881)--(0.624289,-1.569256)--cycle;
\path[fill=blue!13!white,draw=blue!13!white,line width=.12pt] (-1.134919,-0.880070)--(-1.172992,-1.031016)--(-0.889774,-1.081685)--(-0.860893,-0.929094)--cycle;
\path[fill=blue!15!white,draw=blue!15!white,line width=.12pt] (0.860893,-0.929094)--(0.889774,-1.081685)--(1.172992,-1.031016)--(1.134919,-0.880070)--cycle;
\path[fill=blue!14!white,draw=blue!14!white,line width=.12pt] (-1.172992,-1.031016)--(-1.201283,-1.184946)--(-0.911234,-1.236837)--(-0.889774,-1.081685)--cycle;
\path[fill=blue!17!white,draw=blue!17!white,line width=.12pt] (0.889774,-1.081685)--(0.911234,-1.236837)--(1.201283,-1.184946)--(1.172992,-1.031016)--cycle;
\path[fill=blue!18!white,draw=blue!18!white,line width=.12pt] (-0.924449,-1.388587)--(-0.928911,-1.531105)--(-0.624289,-1.569256)--(-0.621290,-1.426556)--cycle;
\path[fill=blue!20!white,draw=blue!20!white,line width=.12pt] (0.621290,-1.426556)--(0.624289,-1.569256)--(0.928911,-1.531105)--(0.924449,-1.388587)--cycle;
\path[fill=blue!21!white,draw=blue!21!white,line width=.12pt] (-0.313655,-1.592296)--(-0.312148,-1.719809)--(-0.000000,-1.727477)--(-0.000000,-1.600000)--cycle;
\path[fill=blue!22!white,draw=blue!22!white,line width=.12pt] (-0.000000,-1.600000)--(-0.000000,-1.727477)--(0.312148,-1.719809)--(0.313655,-1.592296)--cycle;
\path[fill=blue!12!white,draw=blue!12!white,line width=.12pt] (-0.527943,-0.686992)--(-0.554926,-0.818841)--(-0.278805,-0.839321)--(-0.265249,-0.706475)--cycle;
\path[fill=blue!13!white,draw=blue!13!white,line width=.12pt] (0.265249,-0.706475)--(0.278805,-0.839321)--(0.554926,-0.818841)--(0.527943,-0.686992)--cycle;
\path[fill=blue!12!white,draw=blue!12!white,line width=.12pt] (-0.825702,-0.784928)--(-0.860893,-0.929094)--(-0.578577,-0.964452)--(-0.554926,-0.818841)--cycle;
\path[fill=blue!14!white,draw=blue!14!white,line width=.12pt] (0.554926,-0.818841)--(0.578577,-0.964452)--(0.860893,-0.929094)--(0.825702,-0.784928)--cycle;
\path[fill=blue!13!white,draw=blue!13!white,line width=.12pt] (-0.000000,-0.712991)--(-0.000000,-0.846169)--(0.278805,-0.839321)--(0.265249,-0.706475)--cycle;
\path[fill=blue!13!white,draw=blue!13!white,line width=.12pt] (-0.265249,-0.706475)--(-0.278805,-0.839321)--(-0.000000,-0.846169)--(-0.000000,-0.712991)--cycle;
\path[fill=blue!16!white,draw=blue!16!white,line width=.12pt] (-0.911234,-1.236837)--(-0.924449,-1.388587)--(-0.621290,-1.426556)--(-0.612409,-1.274262)--cycle;
\path[fill=blue!18!white,draw=blue!18!white,line width=.12pt] (0.612409,-1.274262)--(0.621290,-1.426556)--(0.924449,-1.388587)--(0.911234,-1.236837)--cycle;
\path[fill=blue!18!white,draw=blue!18!white,line width=.12pt] (-0.621290,-1.426556)--(-0.624289,-1.569256)--(-0.313655,-1.592296)--(-0.312148,-1.449484)--cycle;
\path[fill=blue!20!white,draw=blue!20!white,line width=.12pt] (0.312148,-1.449484)--(0.313655,-1.592296)--(0.624289,-1.569256)--(0.621290,-1.426556)--cycle;
\path[fill=blue!13!white,draw=blue!13!white,line width=.12pt] (-0.860893,-0.929094)--(-0.889774,-1.081685)--(-0.597986,-1.118229)--(-0.578577,-0.964452)--cycle;
\path[fill=blue!15!white,draw=blue!15!white,line width=.12pt] (0.578577,-0.964452)--(0.597986,-1.118229)--(0.889774,-1.081685)--(0.860893,-0.929094)--cycle;
\path[fill=blue!14!white,draw=blue!14!white,line width=.12pt] (-0.889774,-1.081685)--(-0.911234,-1.236837)--(-0.612409,-1.274262)--(-0.597986,-1.118229)--cycle;
\path[fill=blue!16!white,draw=blue!16!white,line width=.12pt] (0.597986,-1.118229)--(0.612409,-1.274262)--(0.911234,-1.236837)--(0.889774,-1.081685)--cycle;
\path[fill=blue!13!white,draw=blue!13!white,line width=.12pt] (-0.554926,-0.818841)--(-0.578577,-0.964452)--(-0.290688,-0.985804)--(-0.278805,-0.839321)--cycle;
\path[fill=blue!13!white,draw=blue!13!white,line width=.12pt] (0.278805,-0.839321)--(0.290688,-0.985804)--(0.578577,-0.964452)--(0.554926,-0.818841)--cycle;
\path[fill=blue!19!white,draw=blue!19!white,line width=.12pt] (-0.312148,-1.449484)--(-0.313655,-1.592296)--(-0.000000,-1.600000)--(-0.000000,-1.457152)--cycle;
\path[fill=blue!19!white,draw=blue!19!white,line width=.12pt] (-0.000000,-1.457152)--(-0.000000,-1.600000)--(0.313655,-1.592296)--(0.312148,-1.449484)--cycle;
\path[fill=blue!13!white,draw=blue!13!white,line width=.12pt] (-0.278805,-0.839321)--(-0.290688,-0.985804)--(-0.000000,-0.992945)--(-0.000000,-0.846169)--cycle;
\path[fill=blue!13!white,draw=blue!13!white,line width=.12pt] (-0.000000,-0.846169)--(-0.000000,-0.992945)--(0.290688,-0.985804)--(0.278805,-0.839321)--cycle;
\path[fill=blue!16!white,draw=blue!16!white,line width=.12pt] (-0.612409,-1.274262)--(-0.621290,-1.426556)--(-0.312148,-1.449484)--(-0.307686,-1.296863)--cycle;
\path[fill=blue!18!white,draw=blue!18!white,line width=.12pt] (0.307686,-1.296863)--(0.312148,-1.449484)--(0.621290,-1.426556)--(0.612409,-1.274262)--cycle;
\path[fill=blue!13!white,draw=blue!13!white,line width=.12pt] (-0.578577,-0.964452)--(-0.597986,-1.118229)--(-0.300440,-1.140298)--(-0.290688,-0.985804)--cycle;
\path[fill=blue!14!white,draw=blue!14!white,line width=.12pt] (0.290688,-0.985804)--(0.300440,-1.140298)--(0.597986,-1.118229)--(0.578577,-0.964452)--cycle;
\path[fill=blue!15!white,draw=blue!15!white,line width=.12pt] (-0.597986,-1.118229)--(-0.612409,-1.274262)--(-0.307686,-1.296863)--(-0.300440,-1.140298)--cycle;
\path[fill=blue!16!white,draw=blue!16!white,line width=.12pt] (0.300440,-1.140298)--(0.307686,-1.296863)--(0.612409,-1.274262)--(0.597986,-1.118229)--cycle;
\path[fill=blue!17!white,draw=blue!17!white,line width=.12pt] (-0.000000,-1.304421)--(-0.000000,-1.457152)--(0.312148,-1.449484)--(0.307686,-1.296863)--cycle;
\path[fill=blue!17!white,draw=blue!17!white,line width=.12pt] (-0.307686,-1.296863)--(-0.312148,-1.449484)--(-0.000000,-1.457152)--(-0.000000,-1.304421)--cycle;
\path[fill=blue!14!white,draw=blue!14!white,line width=.12pt] (-0.290688,-0.985804)--(-0.300440,-1.140298)--(-0.000000,-1.147677)--(-0.000000,-0.992945)--cycle;
\path[fill=blue!14!white,draw=blue!14!white,line width=.12pt] (-0.000000,-0.992945)--(-0.000000,-1.147677)--(0.300440,-1.140298)--(0.290688,-0.985804)--cycle;
\path[fill=blue!15!white,draw=blue!15!white,line width=.12pt] (-0.300440,-1.140298)--(-0.307686,-1.296863)--(-0.000000,-1.304421)--(-0.000000,-1.147677)--cycle;
\path[fill=blue!15!white,draw=blue!15!white,line width=.12pt] (-0.000000,-1.147677)--(-0.000000,-1.304421)--(0.307686,-1.296863)--(0.300440,-1.140298)--cycle;

\draw[red!75!black,thick,domain=10:170,samples=90]
plot({1.6*cos(\x)},{.8*sin(\x)});
\tp{Acrack}{0}{90}
\draw[red!80!black,very thick]
({(2.4-.8*cos(0))*cos(70)},{.5*(2.4-.8*cos(0))*sin(70)})
--({(2.4-.8*cos(2))*cos(80)},{.5*(2.4-.8*cos(2))*sin(80)+.8660254038*.8*sin(2)})
--({(2.4-.8*cos(-2))*cos(90)},{.5*(2.4-.8*cos(-2))*sin(90)+.8660254038*.8*sin(-2)})
--({(2.4-.8*cos(2))*cos(100)},{.5*(2.4-.8*cos(2))*sin(100)+.8660254038*.8*sin(2)})
--({1.6*cos(110)},{.8*sin(110)});
\draw[force,->](Acrack)--++(0,.70)node[above]{$\sigma_{2,\max}$};
\draw[force,->](Acrack)--++(0,-.65);
\draw[thin](.48,.75)--(1.55,2.05);
\node[anchor=west]at(1.62,2.08){裂缝沿内圈环向};
\node at(0,-2.6){最大拉力沿子午方向};
\end{scope}
\begin{scope}[xshift=5.8cm,yshift=-.6cm]
\draw[->](0,0)--(4.95,0)node[right]{$\theta$};
\draw[->](0,0)--(0,2.95)node[above]{$\sigma_i/(pr)$};
\draw[blue!75!black,thick](0,1.0)--(4.4,1.0)node[right]{$\sigma_1$};
\draw[green!45!black,thick,domain=0:360,samples=160]
plot({4.4*\x/360},{2*(6-cos(\x))/(2*(3-cos(\x)))});
\node[green!45!black,anchor=west]at(4.43,2.5){$\sigma_2$};
\foreach \xx/\lab in {0/0,1.1/\frac\pi2,2.2/\pi,3.3/\frac{3\pi}2,4.4/2\pi}{
\draw(\xx,0)--(\xx,-.06)node[below=3pt]{$\lab$};
}
\draw(0,1)--(-.06,1)node[left]{$\frac12$};
\draw(0,2)--(-.06,2)node[left]{$1$};
\draw[aux](0,2)--(4.4,2);
\node at(2.2,3.50){示意取 $R=3r$};
\end{scope}
\end{tikzpicture}
\end{document}
